\documentclass[11pt,oneside]{book}
\usepackage[margin=1in]{geometry}
\usepackage{amsmath,amssymb,amsthm,mathtools}
\usepackage{booktabs,longtable,array}
\usepackage{enumitem}
\usepackage[colorlinks=true,linkcolor=blue,urlcolor=blue,citecolor=blue]{hyperref}
\usepackage[nameinlink,noabbrev]{cleveref}

\newtheorem{theorem}{Theorem}[chapter]
\newtheorem{lemma}[theorem]{Lemma}
\newtheorem{proposition}[theorem]{Proposition}
\newtheorem{corollary}[theorem]{Corollary}
\newtheorem{assessment}[theorem]{Assessment}
\newtheorem{firewall}[theorem]{Claim firewall}
\newtheorem{openproblem}[theorem]{Open problem}
\newtheorem{record}[theorem]{Record}
\newtheorem{programme}[theorem]{Programme}
\theoremstyle{definition}
\newtheorem{definition}[theorem]{Definition}
\theoremstyle{remark}
\newtheorem{remark}[theorem]{Remark}

\newcommand{\MGclosed}{\textsc{closed}}
\newcommand{\MGopen}{\textsc{open}}
\newcommand{\MGpartial}{\textsc{partial}}
\newcommand{\opnorm}[1]{\lVert#1\rVert_{\mathrm{op}}}
\newcommand{\trnorm}[1]{\lVert#1\rVert_{1}}
\newcommand{\Tr}{\operatorname{Tr}}

\title{Renewal Yang--Mills Mass-Gap Research Notes\\
Fourth Series, Version V1}
\author{Proof-indexed continuation after archives f1, f2, and f3}
\date{September 2026}

\begin{document}
\maketitle
\tableofcontents

\chapter{Context, archive map, and present proof boundary}
\label{ch:V1-context}

\section{What this notebook is doing}
\label{sec:V1-purpose}

This is a new compact active notebook for the continuing Yang--Mills
mass-gap programme.  The preceding active file reached V100 and had
again become too large for reliable navigation, so it has been frozen
as archive f3.  The mathematics is not being restarted.  This note
records the major proved interfaces, identifies every withdrawn route,
and continues at the exact remaining global-quintic bottleneck.

The ultimate target is a four-dimensional continuum gauge theory
satisfying the Osterwalder--Schrader axioms whose reconstructed
Hamiltonian has a strictly positive spectral gap above the vacuum.
The current work is much earlier in that implication chain: it seeks a
support- and representation-uniform marked-tree estimate for connected
Wilson carriers.  No mass-gap conclusion is assumed.

\begin{record}[Exact frozen archives]
\label{rec:V1-archives}
The three frozen source sequences are:
\begin{center}
\begin{tabular}{p{0.58\textwidth}rr}
\toprule
archive & source lines & bytes\\
\midrule
\texttt{Renewal\_Yang\_Mills\_Mass\_Gap\_Research\_Notes\_V135\_f1.tex}
& \(88\,517\) & \(3\,083\,467\)\\
\texttt{Renewal\_Yang\_Mills\_Mass\_Gap\_Research\_Notes\_V100\_f2.tex}
& \(49\,581\) & \(1\,706\,928\)\\
\texttt{Renewal\_Yang\_Mills\_Mass\_Gap\_Research\_Notes\_V100\_f3.tex}
& \(48\,381\) & \(1\,720\,945\)\\
\bottomrule
\end{tabular}
\end{center}
Their SHA--256 digests are respectively
\[
\begin{split}
 &\texttt{82f6bbffeecdad3c5c63551660b19a0bdc7f3d6dffe4c68b4e9bba3a783a2b71},\\
 &\texttt{6444a1044e74dfcb1a72ed9c7248bc0f234efc4a0e4248c736cd7db212dfca9a},\\
 &\texttt{a5e05a144c753426d8a461c8290e9538a4cd19d8672bab8b08b509d82faf892f}.
\end{split}
\]
A theorem name written in typewriter font below is a proof index into
one of these exact archives, not a local cross-reference.
\end{record}

\begin{record}[Version protocol]
\label{rec:V1-protocol}
Each successor is made by copying the complete active note, appending
one new proof chapter, validating the append-only result, and retaining
only the newest active version.  Frozen archives remain.  At versions
divisible by five the newest active Standard-Model--Einstein and
Navier--Stokes notes are inspected read-only.  At V100 this fourth
series will be frozen as f4 and a new contextual V1 will begin.
\end{record}

\section{Logical route}
\label{sec:V1-route}

The intended implication chain is
\begin{equation}
\begin{split}
\text{exact Wilson cumulants}
&\Longrightarrow \text{marked-tree atlas (MTA)}
\Longrightarrow \text{weighted uniform fusion control (WUFC)}\\
&\Longrightarrow \text{connected source/Perron bounds (CPM)}
\Longrightarrow \text{uniform finite-edge estimates}\\
&\Longrightarrow \text{cutoff removal and reflection positivity}
\Longrightarrow \text{Osterwalder--Schrader reconstruction}\\
&\Longrightarrow \text{a positive physical spectral gap}.
\end{split}
\label{eq:V1-route}
\end{equation}
The programme is still proving the first arrow beyond row order four.

\section{Major mathematical results obtained so far}
\label{sec:V1-summary}

\subsection{Exact harmonic normalization and one-payer structure}

Archive f1 proves complete Peter--Weyl output pinching: output
dimensions cancel, multiplicity spaces remain orthogonal, and trace
norms sum before physical output compression.  The normalized
nonnegative row-pair stocks are
\begin{equation}
 \Xi_i:=\sum_e a_{i,e},
 \qquad
 H_{ij}:=\sum_e a_{i,e}a_{j,e},
 \qquad
 \theta_{ij}:={H_{ij}\over\Xi_i\Xi_j},
 \qquad 0\le\theta_{ij}\le1.
\label{eq:V1-stocks}
\end{equation}
Binary and ternary connected Wilson carriers obey their complete
marked-tree bounds.

Archive f2 replaces topology-priority propagation by exact first
connectivity.  Every connected coordinate word has one first coordinate
\(r\) where the cumulative row partition becomes connected; the proper
prefix factors over its blocks, the joining letter is the complete
cumulant of the block products, and the unrestricted suffix is a
complete moment contraction.  The final resolvent payer is allocated
once, never once per connected block.

\subsection{Local cumulants and sharp representation tests}

The archives establish representation-uniform complete covariance and
third-cumulant response cards, the selected complete fourth-cumulant
card \texttt{thm:V38-GSEL4}, and the one-coordinate fifth-cumulant
tree-valence decomposition
\texttt{thm:V56-uniform-one-link-quintic}.  At order five the local
decomposition has \(5^{5-2}=125\) labelled tree fibres.

Several initially plausible stronger estimates are false.  A bare
one-Casimir-per-input quintic bound fails on a coherent \(SU(3)\)
family.  An individual V58 graph merge-difference cannot be assigned
one normalized mark per displayed edge: V59 exhibits endpoint
covariance leakage along a path.  For \(SU(2)\), a complete equal-spin
heat square retains its singlet with product-state capacity
\[
 {1\over2j+1}\asymp (1+C_2(j))^{-1/2},
\]
so one complete heat slot cannot be promoted to a stronger inverse
Casimir power.

\subsection{Restored complete quartic MTA}

Archive f3 V89 corrects the normalized \(3+1\) atlas.  V92 and V94
construct physical suffix-rescue counterpackets to two literal local
weighted cards.  V93 and V95 replace those cards by finite globally
rescued but coordinate-summable square atlases and close respectively
all \(3+1\) and \(2+1+1\) sources.  Together with the established
\(2+2\) sources and the selected discrete fourth-cumulant source, f3
V96 proves
\begin{equation}
 \kappa_\otimes^{\rm abs}(\mathcal C_4^{\rm pay})
 \le C_{G,s}\sum_{\tau\in\mathfrak T_4}
             \prod_{ij\in E(\tau)}\theta_{ij}.
\label{eq:V1-quartic-MTA}
\end{equation}
The proof is an exact census of all fourteen proper prefix partitions:
four of type \(3+1\), three of type \(2+2\), six of type
\(2+1+1\), and the discrete partition.

\subsection{Global quintic reduction}

F3 V97 proves the coefficient-one first-connectivity decomposition of
the five-row carrier into all \(B_5-1=51\) proper prefix partitions.
Their types and counts are
\[
\begin{array}{c|rrrrrr}
\text{type}&4+1&3+2&3+1+1&2+2+1&2+1+1+1&1^5\\
\hline
\text{count}&5&10&10&15&10&1.
\end{array}
\]
Internal block trees and a quotient tree lift to one ordinary
five-row tree.  The finite rescue atlas retains at least one actual
joining edge, and its coordinate sum is bounded by a fixed multiple of
the \(125\)-tree global atlas.

F3 V98 enlarges this to an exchange-rescue atlas, permitting a strong
global pair label outside the literal source tree while retaining a
local joining anchor.  Two maximal-scale strong edges close the raw
five-row normalization whenever they are anchor-compatible.  Failed
exchange is classified exactly: the joining set is covered by the two
strong labels and possibly their unique closing chord.

F3 V99--V100 proves that this obstruction is short by \(M^{-1}\) at
word level and therefore needs \(M^{-2}\) in the oriented square.
Four source-occurring diffuse heat slots or four individually cofinal
complete heat slots supply that response.  The four-slot threshold is
sharp on tensor products of the \(SU(2)\) singlet test.  Three distinct
comparable partners of a maximal row always contain a compatible strong
pair.  The residual one-/two-partner core has at most three usable
cofinal slots and the exact remaining square exponents
\begin{equation}
 \delta_q={4-q\over2},
 \qquad q=0,1,2,3.
\label{eq:V1-residual-exponents}
\end{equation}

\section{Current status}
\label{sec:V1-status}

\begin{longtable}{>{\raggedright\arraybackslash}p{0.31\textwidth}
                  >{\raggedright\arraybackslash}p{0.13\textwidth}
                  >{\raggedright\arraybackslash}p{0.48\textwidth}}
\toprule
gate & status & exact boundary\\
\midrule
\endfirsthead
\toprule
gate & status & exact boundary\\
\midrule
\endhead
Binary and ternary normalized MTA
& \MGclosed
& Complete Fourier normalization, probability marks, and one-payer
  response are proved.\\ \addlinespace[0.35em]
Global quartic MTA
& \MGclosed
& All fourteen first-connectivity sources and every quartic coordinate
  topology are included in \eqref{eq:V1-quartic-MTA}.\\ \addlinespace[0.35em]
Local one-coordinate quintic card
& \MGclosed
& The complete fifth cumulant has a representation-uniform
  \(125\)-tree decomposition.\\ \addlinespace[0.35em]
Global quintic MTA
& \MGpartial
& Exact \(51\)-source census, scalar rescue summation, diffuse branch,
  compatible strong pairs, three-partner branch, and four-slot
  supplements are closed.  The finite \(q<4\) one-/two-partner core is
  open.\\ \addlinespace[0.35em]
All-order MTA and WUFC
& \MGopen
& No arity-uniform induction or constant-growth theorem has been
  proved.\\ \addlinespace[0.35em]
CPM, continuum construction, and gap
& \MGopen
& These downstream stages have not been reached.\\
\bottomrule
\end{longtable}

\begin{firewall}[Claim boundary at the fourth-series restart]
\label{fire:V1-initial-boundary}
The restored quartic theorem and the quintic reductions are finite
cutoff, finite-row constructive estimates.  They do not construct the
continuum Yang--Mills measure, prove reflection positivity after cutoff
removal, build the Osterwalder--Schrader Hilbert space or Hamiltonian,
or prove a positive physical spectral gap.  The full Yang--Mills mass
gap has not been proved.
\end{firewall}

\chapter{The cofinal-pair position inside the quintic prefix}
\label{ch:V1-pair-position}

\section{Exact \(37/14\) source census}
\label{sec:V1-pair-census}

Fix two row labels \(a\ne c\) and a proper partition
\(\rho<\widehat1\) of \([5]\).  Say that \(\rho\) is
\emph{pair-internal} if \(a,c\) belong to the same block and
\emph{pair-crossing} otherwise.

\begin{lemma}[Cofinal-pair census by prefix type]
\label{lem:V1-pair-census}
Among the \(51\) proper quintic prefix partitions, \(14\) are
pair-internal and \(37\) are pair-crossing.  The typewise counts are:
\begin{center}
\begin{tabular}{c|rrrrrr|r}
type&\(4+1\)&\(3+2\)&\(3+1+1\)&\(2+2+1\)
&\(2+1+1+1\)&\(1^5\)&total\\
\hline
pair-internal&3&4&3&3&1&0&14\\
pair-crossing&2&6&7&12&9&1&37.
\end{tabular}
\end{center}
\end{lemma}

\begin{proof}
For type \(4+1\), the pair is separated precisely when the singleton
is \(a\) or \(c\), giving crossing/internal counts \(2/3\).
For type \(3+2\), the pair is internal either in the two-block
(one choice) or in the three-block (choose its third member in three
ways), giving \(4/6\).
For type \(3+1+1\), choose the third member of the internal triple in
three ways, giving \(3/7\).
For type \(2+2+1\), an internal pair is one two-block and the remaining
three labels split as pair and singleton in three ways, giving \(3/12\).
For type \(2+1+1+1\), only \(\{a,c\}\) is internal, giving \(1/9\).
The discrete partition is crossing.  Summing gives \(14/37\).
\end{proof}

\begin{remark}[Bell-number check]
\label{rem:V1-Bell-check}
Merging \(a,c\) into one super-label bijects all partitions in which
they are together with \(\Pi_4\), of cardinality \(B_4=15\).
Removing the one-block partition leaves \(14\), agreeing with
\Cref{lem:V1-pair-census}.
\end{remark}

\section{The two different operator roles}
\label{sec:V1-pair-role}

For a fixed first-connectivity coordinate \(r\), let \(P\) be the
union of prefix block-tree edges and \(Q\) a lifted quotient tree for
the complete joining cumulant.  Thus \(|P|=5-|\rho|\),
\(|Q|=|\rho|-1\), and \(T_0=P\sqcup Q\) is a five-row spanning tree.

\begin{proposition}[Internal versus crossing placement]
\label{prop:V1-pair-role}
Let \(ac\) be the sole cofinal pair label.
\begin{enumerate}[label=\textup{(\roman*)},leftmargin=4em]
\item If \(\rho\) is pair-internal, \(ac\) belongs to the connected
      prefix carrier of that block and cannot be a local quotient edge.
\item If \(\rho\) is pair-crossing, the quotient vertices containing
      \(a,c\) are distinct.  Their edge extends to a quotient spanning
      tree and may be lifted by the row edge \(ac\).
\end{enumerate}
The second statement identifies an available local tree fibre; it does
not assert that the global cofinal occurrence lies at \(r\).
\end{proposition}

\begin{proof}
In the internal case, \(ac\) becomes a loop after quotient contraction.
In the crossing case, the edge between the two distinct blocks is an
acyclic quotient forest.  Any forest in a complete graph extends to a
spanning tree; lift this specified quotient edge by \(ac\).
\end{proof}

\section{A locally cofinal anchor}
\label{sec:V1-local-anchor}

Retain the f3 exchange-rescue interface.  Its nonnegative weight
\(\mathfrak x_{T_0,r}\) contains every five-row tree which keeps at
least one edge of \(Q\) local at \(r\), permits other edges of
\(P\cup Q\) in actual or global roles, and permits edges outside
\(T_0\) only globally.  It satisfies
\begin{equation}
 \sum_r\mathfrak x_{T_0,r}
 \le15\sum_{\tau\in\mathfrak T_5}
          \prod_{e\in E(\tau)}\theta_e.
\label{eq:V1-exchange-sum}
\end{equation}
Every fixed-payer word has
\begin{equation}
 \opnorm{\widetilde W_{\rho,T_0,r,\pi}}
 \le {C_{G,s}\over\prod_{i=1}^5\Xi_i}.
\label{eq:V1-raw}
\end{equation}

\begin{definition}[Locally cofinal joining edge]
\label{def:V1-local-cofinal}
An edge \(q=uv\in Q\) is locally
\((\eta,m_0,M)\)-cofinal at \(r\) if
\[
 \theta_{q,r}\ge\eta,\qquad
 \Xi_u,\Xi_v\ge m_0M,\qquad M=\max_i\Xi_i.
\]
A distinct global edge \(s=xy\) is maximal-scale strong if
\(\theta_s\ge\eta\) and \(\Xi_x,\Xi_y\ge m_0M\).
\end{definition}

\begin{lemma}[Local--global strong basis exchange]
\label{lem:V1-local-global-exchange}
Suppose \(q\in Q\) is locally cofinal and \(s\ne q\) is a global
maximal-scale strong edge.  There are actual source edges
\(f,g\in E(T_0)\) such that
\(\tau=\{q,s,f,g\}\) is a five-row spanning tree, and
\[
 \theta_{q,r}\theta_s\theta_f^{\rm act}\theta_g^{\rm act}
\]
is a term of \(\mathfrak x_{T_0,r}\).
\end{lemma}

\begin{proof}
Two distinct simple edges form a forest.  Starting from its components,
add an edge of the connected tree \(T_0\) crossing two current
components, and repeat.  Exactly two additions give a connected
five-vertex graph with four edges, hence a tree.  In the exchange
weight keep \(q\) local, \(s\) global, and the two additions in their
actual prefix or joining roles.
\end{proof}

\begin{theorem}[Local cofinality removes chord and coverage
obstructions]
\label{thm:V1-local-cofinal-close}
Under the hypotheses of
\Cref{lem:V1-local-global-exchange}, the fixed resolved word satisfies
\begin{align}
 \kappa_\otimes(\widetilde W^*\widetilde W)
 &\le C\mathfrak x_{T_0,r}^{\,2},&
 \kappa_\otimes(\widetilde W\widetilde W^*)
 &\le C\mathfrak x_{T_0,r}^{\,2}.
\label{eq:V1-local-cofinal-card}
\end{align}
Thus the f3 V98 chord-or-coverage obstruction closes whenever a
retained local joining edge is itself cofinal and the other strong pair
label is distinct.
\end{theorem}

\begin{proof}
Let \(F=\{f,g\}\).  Removing \(q,s\) from the spanning tree leaves a
two-edge forest.  If it is a matching, its endpoint-mass product divided
by \(\prod_i\Xi_i\) is at most one.  If it is a two-edge path, its two
isolated vertices must each be incident to \(q\) or \(s\), so their
masses are at least \(m_0M\); the path-centre mass is at most \(M\).
Hence
\[
 {\prod_{uv\in F}\Xi_u\Xi_v\over\prod_i\Xi_i}
 \le\max\{1,m_0^{-2}\}.
\]
The two actual edges obey
\(\theta_{uv}^{\rm act}\ge a_*^2/(\Xi_u\Xi_v)\).
Together with
\(\theta_{q,r}\theta_s\ge\eta^2\), this gives
\[
 \mathfrak x_{T_0,r}
 \ge {c_{\eta,m_0,a_*}\over\prod_i\Xi_i}.
\]
Compare with \eqref{eq:V1-raw}; ordinary norm squared bounds either
oriented positive product-state capacity.
\end{proof}

\begin{firewall}[The same pair label is not two strong edges]
\label{fire:V1-no-duplicate-pair}
If the sole global strong label and the local joining edge are both
\(ac\), \Cref{thm:V1-local-cofinal-close} does not apply.  A global
sum and one local summand of the same normalized pair stock are not two
edges of a simple spanning tree.  The one-partner core remains governed
by the \(q\)-slot exponents \eqref{eq:V1-residual-exponents}.
\end{firewall}

\begin{corollary}[Refined post-rollover residual]
\label{cor:V1-refined-residual}
In the supplement-free two-partner chord-or-coverage class, every word
with a locally cofinal retained joining edge closes.  Any persistent
failure has every retained local joining anchor noncofinal.  In the
one-partner class, the \(37\) pair-crossing prefixes permit alignment
with \(ac\) but obtain no second simple strong edge from that label;
the \(14\) pair-internal prefixes permit no local alignment.
\end{corollary}

\begin{proof}
Use \Cref{thm:V1-local-cofinal-close} and its contrapositive inside the
fixed finite fibre selection, followed by
\Cref{lem:V1-pair-census,prop:V1-pair-role,fire:V1-no-duplicate-pair}.
\end{proof}

\section{V1 ledger and V2 acquisition order}
\label{sec:V1-ledger}

\begin{theorem}[Fourth-series V1 result ledger]
\label{thm:V1-results}
This restarted note:
\begin{enumerate}[label=\textup{(V1.\arabic*)},leftmargin=5.8em]
\item records exact identities for archives f1--f3 and the present
      proof boundary;
\item summarizes the restored global quartic MTA and finite quintic
      \(q\)-slot core without reviving withdrawn claims;
\item proves the exact \(37/14\) crossing/internal census;
\item separates prefix-internal from quotient-crossing pair roles;
\item proves local--global strong basis exchange; and
\item closes the locally cofinal two-partner chord/coverage sector,
      reducing any survivor to local-anchor depletion.
\end{enumerate}
\end{theorem}

\begin{proof}
These are
\Cref{rec:V1-archives,
lem:V1-pair-census,
prop:V1-pair-role,
lem:V1-local-global-exchange,
thm:V1-local-cofinal-close,
cor:V1-refined-residual}.
\end{proof}

\begin{programme}[V2: locally depleted quotient response]
\label{prog:V1-V2}
\begin{enumerate}[leftmargin=3em]
\item Fix one of the \(37\) pair-crossing prefixes and one locally
      depleted joining-tree fibre.  Keep the quotient cumulant assembled.
\item Separate depletion of the selected coordinate from depletion due
      only to its endpoint lift; do not norm moment partitions separately.
\item For \(q=3\), seek exactly the missing \(M^{-1/2}\) response first.
      This is the smallest unresolved exponent and is \(SU(2)\)-sharp.
\item If that response fails, construct a physical counterpacket and
      enlarge the analytic carrier while retaining a local summation
      anchor.
\item Only afterward move to the \(14\) pair-internal prefixes, whose
      strong pair lies in a lower-order prefix block.
\item The next compulsory companion audit is V5.
\end{enumerate}
\end{programme}

\begin{firewall}[Claim boundary after fourth-series V1]
\label{fire:V1-claim-boundary}
V1 closes the locally cofinal two-partner obstruction and splits the
remaining pair geometry exactly.  It does not prove the depleted
\(q\)-slot response cards, global quintic or all-order MTA, WUFC, CPM,
cutoff removal, reflection positivity, Osterwalder--Schrader
reconstruction, construction of the continuum Yang--Mills measure, or
a uniform positive continuum spectral gap.  The full Yang--Mills mass
gap has not been proved.
\end{firewall}

\paragraph{Rollover record.}
Fourth-series V1 was started from the validated f3 archive
\texttt{Renewal\_Yang\_Mills\_Mass\_Gap\_Research\_Notes\_V100\_f3.tex},
whose exact identity is given in \Cref{rec:V1-archives}.  The next
compulsory companion audit is V5.

% =============================================================================
% BEGIN APPENDED FOURTH-SERIES V2 MATERIAL
% =============================================================================

\clearpage
\chapter{V2: Fresh strong-coordinate heat response and closure of the
\texorpdfstring{\(q=3\)}{q=3} core}
\label{ch:V2}

\section{Why the strong coordinate must be audited as an operator}
\label{sec:V2-purpose}

The fourth-series V1 programme asked first for the missing
\(M^{-1/2}\) response in the \(q=3\) core.  A balanced cofinal atom
already has exactly that heat scale, but f3 contains two binding
warnings:
\begin{enumerate}[label=\textup{(\roman*)},leftmargin=4em]
\item two occurrences of one protected projection do not square its
      product-state capacity; and
\item a bounded noncommuting separator can map a product vector into
      the protected singlet and amplify capacity from \(1/d\) to one.
\end{enumerate}
Thus the cofinal coordinate may be used only after its exact position
in the source word has been resolved.  This chapter proves the legal
branch: the coordinate is a complete unmarked central slot, disjoint
from the joining and payer roles.  It then supplies the fourth
half-power needed by the three already usable slots.

\begin{assessment}[The distinction from a repeated estimate]
\label{ass:V2-no-double-use}
At the strong coordinate below, the complete local operator receives
one analytic estimate: its central heat-square capacity.  The inequality
\(\theta_{ac}\ge\eta\) is used only to lower-bound one scalar monomial
on the right side of the desired tree card.  The local operator is not
replaced by a product of a heat estimate and an edge estimate, and the
protected projection is not counted twice.
\end{assessment}

\section{Fresh strong atoms}
\label{sec:V2-fresh}

Fix one resolved quintic first-connectivity word, its literal tree
\(T_0=P\sqcup Q\), payer placement, and square orientation.  Let
\(a\) be a maximal row, \(\Xi_a=M\), and let \(c\ne a\).

\begin{definition}[Fresh balanced strong coordinate]
\label{def:V2-fresh-strong}
A coordinate \(e_0\) is a fresh
\((\sigma,\lambda,M)\)-strong coordinate for \(a,c\) if:
\begin{enumerate}[label=\textup{(\roman*)},leftmargin=4em]
\item
\begin{equation}
 a_{a,e_0}\ge\sigma M,
 \qquad
 a_{c,e_0}\ge\lambda a_{a,e_0};
\label{eq:V2-fresh-masses}
\end{equation}
\item \(e_0\) is disjoint from the joining coordinate, payer
      coordinate, and the fixed finite marked/role-owner set;
\item the common exact source-slot resolution exposes at \(e_0\) one
      complete moment-partition heat sector
      \(0\preccurlyeq\mathcal M_{e_0}\preccurlyeq I\); and
\item in the required orientation the resolved normalized word has
\begin{equation}
 \widetilde W
 ={1\over\prod_i\Xi_i}\,
  L\,\mathcal M_{e_0}\,Q_0,
 \qquad
 [\mathcal M_{e_0},Q_0]=0,
 \qquad
 \|L\|+\|Q_0\|\le C_{G,s}.
\label{eq:V2-fresh-factorization}
\end{equation}
For the opposite square orientation, the analogous factorization is
required after taking adjoints.
\end{enumerate}
\end{definition}

\begin{lemma}[An unmarked complete source slot is central]
\label{lem:V2-fresh-central}
Suppose \(e_0\) satisfies
\eqref{eq:V2-fresh-masses} and lies outside the joining, payer, and
finite marked coordinates.  In every fixed moment-partition sector of
the common exact source-slot resolution, conditions
\textup{(iii)}--\textup{(iv)} of
\Cref{def:V2-fresh-strong} hold.
\end{lemma}

\begin{proof}
The exact source-slot resolution separates the complete coordinate
\(e_0\) from its complement before taking a norm.  On each local
isotypic block, its normalized heat multiplier is a scalar spectral
function with values in \([0,1]\).  Hence it is a positive contraction.

Every complete moment or cumulant intertwiner, canonical equivariant
recoupling, and local physical isotypic compression commutes with that
spectral function.  Operators at other coordinates commute tensorially.
This is precisely the central isotypic functional calculus proved in f3
\(\texttt{lem:V59-central-isotypic}\) and
\(\texttt{lem:V59-bank-commutation}\).

Move only this central factor to the displayed position.  All unpaid
complete moments and physical compressions are contractions; the
fixed-order joining cumulant and separators have representation-uniform
raw norms; and the sole payer, which is outside \(e_0\), has its
one-resolvent bound.  Complete physical pinching precedes the finite
triangle inequality, so no multiplicity sum occurs.  These are the
factors \(L,Q_0\) in \eqref{eq:V2-fresh-factorization}.  The adjoint
argument gives the other orientation.
\end{proof}

\begin{firewall}[Why joining and payer coordinates are excluded]
\label{fire:V2-role-exclusion}
If \(e_0\) is the joining coordinate, carries the payer, or belongs to
the marked role-owner packet, its complete local factor contains those
operations in their original order.  The proof of
\Cref{lem:V2-fresh-central} may not split that factor into independently
estimated pieces.  Such coincident-role atoms form the V3 residual.
\end{firewall}

\section{One coordinate, one half-power}
\label{sec:V2-half-power}

\begin{lemma}[Fresh-coordinate heat response]
\label{lem:V2-fresh-half}
For a fresh strong coordinate,
\begin{equation}
 \boxed{\;
 \kappa_{a|a^c}(\mathcal M_{e_0}^{\,2})
 \le C_{G,s,\sigma}M^{-1/2}.\;}
\label{eq:V2-fresh-half}
\end{equation}
The exponent \(1/2\) cannot be increased uniformly for one protected
\(SU(2)\) equal-spin slot.
\end{lemma}

\begin{proof}
The complete reducible-complement heat card in f3
\(\texttt{cor:V87-local-slot}\) gives
\[
 \|\mathcal M_{e_0}^{\Gamma_{a^c}}\|
 \le C_{G,s}a_{a,e_0}^{-1/2}
 \le C_{G,s,\sigma}M^{-1/2}.
\]
Product-state capacity is bounded by the appropriate
partial-transpose norm.  Since
\(0\preccurlyeq\mathcal M_{e_0}\preccurlyeq I\),
\(\mathcal M_{e_0}^2\preccurlyeq\mathcal M_{e_0}\), which proves
\eqref{eq:V2-fresh-half}.  The f3 V91 equal-spin singlet has capacity
\((2j+1)^{-1}\asymp a_j^{-1/2}\), proving sharpness.
\end{proof}

\begin{lemma}[The same atom supplies the scalar strong edge]
\label{lem:V2-fresh-edge}
Under \eqref{eq:V2-fresh-masses},
\begin{equation}
 \Xi_c\ge\lambda\sigma M,
 \qquad
 \theta_{ac}\ge\lambda\sigma^2.
\label{eq:V2-fresh-edge}
\end{equation}
\end{lemma}

\begin{proof}
The first inequality follows from
\(\Xi_c\ge a_{c,e_0}\).  Moreover,
\[
 H_{ac}\ge a_{a,e_0}a_{c,e_0}
 \ge\lambda\sigma^2M^2.
\]
Since \(\Xi_a=M\) and \(\Xi_c\le M\),
division by \(\Xi_a\Xi_c\le M^2\) proves the second inequality.
\end{proof}

\begin{proposition}[Simultaneous use is not capacity squaring]
\label{prop:V2-simultaneous-legal}
The conclusions
\eqref{eq:V2-fresh-half} and
\eqref{eq:V2-fresh-edge} may be used in the same tree-card comparison
without violating the one-coordinate/one-estimate rule.
\end{proposition}

\begin{proof}
The left side is estimated from the single factorization
\eqref{eq:V2-fresh-factorization} by the single positive operator
\(\mathcal M_{e_0}^2\).  No second copy of its capacity is introduced.
Independently, the numerical data
\eqref{eq:V2-fresh-masses} imply the scalar inequality
\eqref{eq:V2-fresh-edge}, which lower-bounds a right-side global
tree stock.  A proof would double-use the coordinate only if it
replaced the one local operator by a product of two separately normed
operators or squared its protected capacity.  Neither operation occurs.
\end{proof}

\section{Closure of the fresh \texorpdfstring{\(q=3\)}{q=3} core}
\label{sec:V2-q3}

Let \(e_1,e_2,e_3\) be the three usable cofinal heat slots retained by
the f3 V100 \(q=3\) resolution.  They are distinct from \(e_0\) and
from one another, and their complete product is denoted
\(\mathcal M_{D_3}\).

\begin{lemma}[Four distinct half-powers]
\label{lem:V2-four-half}
For a fresh \(e_0\) and the three usable slots,
\begin{equation}
 \boxed{\;
 \kappa_{a|a^c}
  \bigl((\mathcal M_{e_0}\otimes
         \mathcal M_{D_3})^2\bigr)
 \le C_{G,s,\sigma}M^{-2}.\;}
\label{eq:V2-four-half}
\end{equation}
\end{lemma}

\begin{proof}
Each usable slot obeys the same complete-slot
\(CM^{-1/2}\) bound as in \Cref{lem:V2-fresh-half}.  The four
coordinates are distinct tensor factors.  Partial transpose across the
fixed row cut and operator norm tensorize, even though the test state
may be entangled among coordinates within either side of the cut.
Multiplication gives \(M^{-4/2}=M^{-2}\).
\end{proof}

\begin{lemma}[One-strong exchange floor]
\label{lem:V2-one-strong-floor}
Let \(s=ac\) be the strong edge from
\Cref{lem:V2-fresh-edge}.  For any literal source tree
\(T_0=P\sqcup Q\) for which \(Q\setminus\{s\}\ne\varnothing\),
its exchange atlas contains a spanning-tree
monomial satisfying
\begin{equation}
 \mathfrak x_{T_0,r}
 \ge {c_{G,s,\sigma,\lambda}\over
 M\prod_{i=1}^5\Xi_i}.
\label{eq:V2-one-strong-floor}
\end{equation}
\end{lemma}

\begin{proof}
Choose a local joining edge \(q\in Q\setminus\{s\}\).
The two-edge forest \(\{s,q\}\) extends to a spanning tree by two
actual edges \(f,g\in T_0\).  Keep \(s\) global, \(q\) local, and
\(f,g\) in their actual roles.  Deleting \(s\) leaves a three-edge
forest with two components.  Its endpoint-mass product divided by
\(\prod_i\Xi_i\) is at most \(C_{\lambda,\sigma}M\): for component
sizes \(2+3\) the sole excess is one path-centre mass; for sizes
\(1+4\) the isolated vertex is a maximal-scale endpoint of \(s\), and
it cancels one of the at most two excess powers of the four-vertex
tree.  Multiplying the three actual-edge floors and
\(\theta_s\ge\lambda\sigma^2\) gives
\eqref{eq:V2-one-strong-floor}.
\end{proof}

\begin{firewall}[The coverage subcase \(Q=\{ac\}\)]
\label{fire:V2-Q-single-strong}
If the sole local joining edge is the same label \(ac\), a simple
tree cannot use a separate global copy of \(ac\).  The claimed
one-strong floor in \Cref{lem:V2-one-strong-floor} then requires either
a distinct strong label or a different quotient-tree fibre.  Without
one, only the literal \(M^{-8}\) floor is available.  This subcase is
excluded from the fresh \(q=3\) theorem below and remains a
coincident-label residual.
\end{firewall}

\begin{theorem}[Fresh-coordinate \(q=3\) closure]
\label{thm:V2-fresh-q3-close}
Consider a \(q=3\) residual word with a fresh balanced strong coordinate
\(e_0\).  Assume its fixed quotient-tree fibre contains a local joining
edge \(q\ne ac\).  Then both orientations of the exchange-rescue
quintic square card hold.
\end{theorem}

\begin{proof}
Apply the exact central insertion algebra to the product of the four
complete coordinate factors in
\Cref{lem:V2-four-half}.  All remaining source factors have the
uniform raw bound used in \eqref{eq:V1-raw}.  Hence either oriented
square has capacity at most
\[
 {C_{G,s,\sigma}M^{-2}\over(\prod_i\Xi_i)^2}.
\]
By \Cref{lem:V2-one-strong-floor}, this is at most
\(C\mathfrak x_{T_0,r}^{\,2}\).
\end{proof}

\begin{corollary}[Global aggregation of the fresh \(q=3\) branch]
\label{cor:V2-fresh-q3-global}
The sum of all fresh-coordinate \(q=3\) fibres satisfying
\(Q\setminus\{ac\}\ne\varnothing\) obeys its contribution to the
global quintic marked-tree atlas, uniformly in support,
representations, physical multiplicities, and payer placement.
\end{corollary}

\begin{proof}
Apply \Cref{thm:V2-fresh-q3-close} to each fixed member of the
finite source, tree-fibre, endpoint-lift, payer, and orientation list.
Product-state Cauchy--Schwarz gives the linear exchange weight, whose
coordinate sum is \eqref{eq:V1-exchange-sum}.  Fixed finite
recombination changes only the constant.
\end{proof}

\section{The noncentral no-go and exact residual}
\label{sec:V2-residual}

\begin{proposition}[No heat-capacity insertion through an arbitrary
separator]
\label{prop:V2-entangler-no-go}
There is no representation-uniform constant \(C\) such that
\begin{equation}
 \kappa_\otimes(U^*P_{\rm sing}U)
 \le C\,\kappa_\otimes(P_{\rm sing})
\label{eq:V2-false-insertion}
\end{equation}
for every unitary \(U\) on \(V\otimes V^*\), where
\(P_{\rm sing}\) is the invariant rank-one projection.
\end{proposition}

\begin{proof}
If \(d=\dim V\), then
\(\kappa_\otimes(P_{\rm sing})=1/d\).  Choose \(U\) mapping one
product unit vector to the normalized invariant vector.  Testing
\(U^*P_{\rm sing}U\) on that product vector gives capacity one.
Then \eqref{eq:V2-false-insertion} would imply \(1\le C/d\), which
fails as \(d\to\infty\).
\end{proof}

\begin{theorem}[Exact \(q=3\) residual after V2]
\label{thm:V2-q3-residual}
After the fixed source-slot resolution, every persistent \(q=3\)
counterpacket is confined to at least one of:
\begin{enumerate}[label=\textup{(\roman*)},leftmargin=4em]
\item every balanced cofinal strong coordinate coincides with the
      joining, payer, or finite marked role-owner set;
\item the fixed quotient-tree joining set is \(Q=\{ac\}\), the sole
      cofinal pair label, so no distinct local anchor is available;
\item the source restriction fails to expose the strong coordinate as
      a complete central moment-partition sector; or
\item the exact local factor contains a noncentral cut-entangling
      separator, for which
      \Cref{prop:V2-entangler-no-go} forbids a capacity insertion based
      only on operator norm.
\end{enumerate}
Every \(q=3\) fibre outside this list is closed by
\Cref{thm:V2-fresh-q3-close}.
\end{theorem}

\begin{proof}
If a cofinal balanced coordinate is disjoint from all fixed roles, the
common source-slot resolution and
\Cref{lem:V2-fresh-central} make it fresh unless
\textup{(iii)} or \textup{(iv)} occurs.  If a distinct local anchor
exists, \Cref{thm:V2-fresh-q3-close} applies.  Its failure therefore
forces \textup{(i)} or \textup{(ii)}.  These alternatives exhaust the
hypotheses used by that theorem.  The final no-go explains why
\textup{(iv)} cannot be discarded by a norm-one estimate.
\end{proof}

\begin{assessment}[Progress relative to the f3 core]
\label{ass:V2-progress}
The \(q=3\) exponent is acquired whenever the strong coordinate is
fresh and a distinct local joining anchor exists.  The remaining
difficulty is finite and operator-bearing: coincident joining/payer
roles, the single-label coverage fibre, or a genuinely noncentral
local separator.  The next step is therefore not a stronger scalar
tree inequality.  It is a complete one-coordinate analysis of those
role owners.
\end{assessment}

\section{V2 result ledger and V3 acquisition order}
\label{sec:V2-ledger}

\begin{theorem}[Fourth-series V2 result ledger]
\label{thm:V2-results}
Relative to V1, V2 proves:
\begin{enumerate}[label=\textup{(V2.\arabic*)},leftmargin=5.8em]
\item exact central factorization of a fresh strong coordinate;
\item its sharp \(M^{-1/2}\) complete-heat response;
\item simultaneous but nonduplicative use of that response and its
      scalar strong-edge lower bound;
\item the four-distinct-half-power \(M^{-2}\) supplement;
\item closure and global aggregation of every fresh \(q=3\) fibre with
      a distinct local joining anchor;
\item a finite-dimensional entangler counterexample to arbitrary
      noncentral capacity insertion; and
\item the exact four-class residual
      \Cref{thm:V2-q3-residual}.
\end{enumerate}
\end{theorem}

\begin{proof}
These are
\Cref{lem:V2-fresh-central,
lem:V2-fresh-half,
lem:V2-fresh-edge,
prop:V2-simultaneous-legal,
lem:V2-four-half,
thm:V2-fresh-q3-close,
cor:V2-fresh-q3-global,
prop:V2-entangler-no-go,
thm:V2-q3-residual}.
\end{proof}

\begin{programme}[V3: coincident-role \(q=3\) atoms]
\label{prog:V2-V3}
\begin{enumerate}[leftmargin=3em]
\item Enumerate the finite role owners at a cofinal coordinate:
      selected joining, local payer, prefix-tree witness, and physical
      separator.
\item For each role, form the complete local moment--cumulant operator
      before applying a norm.  Do not multiply estimates attached to
      coincident names.
\item Treat the single-label fibre \(Q=\{ac\}\) at quotient order two
      by its assembled block covariance, retaining its local coordinate
      sum.
\item Prove a central or projective bound for the actual physical
      separator; a generic norm-one bound is ruled out by
      \Cref{prop:V2-entangler-no-go}.
\item Test every claimed half-power on the equal-spin \(SU(2)\)
      protected channel and include a payer at the same coordinate.
\item If a coincident role lacks the response, construct the smallest
      physical counterpacket and move upstream to an assembled
      covariance--payer carrier.
\end{enumerate}
\end{programme}

\begin{firewall}[Claim boundary after fourth-series V2]
\label{fire:V2-claim-boundary}
V2 closes only the fresh-coordinate \(q=3\) core with a distinct local
joining anchor.  It does not close coincident-role or single-label
coverage fibres, the \(q=0,1,2\) response cards, global quintic or
all-order MTA, WUFC, CPM, cutoff removal, reflection positivity,
Osterwalder--Schrader reconstruction, construction of the continuum
Yang--Mills measure, or a uniform positive continuum spectral gap.
The full Yang--Mills mass gap has not been proved.
\end{firewall}

\paragraph{Parent-version record.}
V2 was formed from
\texttt{Renewal\_Yang\_Mills\_Mass\_Gap\_Research\_Notes\_V1.tex}.
The V1 parent had \(536\) logical source lines, \(20\,867\) bytes, and
SHA--256
\[
 \texttt{9ef1f1c3b4e1481d43d0fd9722194847d98bcc5fda14e3f0f8abbfb00ddd35e6}.
\]
The next compulsory companion audit is V5.

% =============================================================================
% END APPENDED FOURTH-SERIES V2 MATERIAL
% =============================================================================

% =============================================================================
% BEGIN APPENDED FOURTH-SERIES V3 MATERIAL
% =============================================================================

\clearpage
\chapter{V3: Selected block covariance and the single-label
\texorpdfstring{\(q=3\)}{q=3} coverage fibre}
\label{ch:V3}

\section{Purpose and exact scope}
\label{sec:V3-purpose}

V2 closes the fresh-coordinate \(q=3\) core when the literal quotient
tree has a local joining edge distinct from the global strong edge.
The first residual is \(Q=\{ac\}\): quotient order is two, the only
local joining edge has the same pair label as the cofinal strong stock,
and a simple tree cannot contain separate local and global copies of
that label.

At quotient order two the joining carrier has a special exact form.  It
is one complete block covariance, hence the difference of two positive
complete heat moments.  This chapter resolves both moments in one
nested-isotypic decomposition and estimates the square of their
difference once.  With an outside-local payer, the selected covariance
then supplies the missing \(M^{-1/2}\) response itself.

\begin{firewall}[No claim for a local payer]
\label{fire:V3-local-payer}
The payer is assumed not to be at the selected coordinate \(r\).
When the local covariance pays, its numerator and resolvent change the
operator before the square is formed.  V3 does not insert the unpaid
covariance estimate twice around that resolvent.  The local-payer
carrier remains for V4.
\end{firewall}

\section{The exact two-block joining operator}
\label{sec:V3-covariance}

Let \(\rho=A|B\) be a proper two-block partition of \([5]\), of type
\(4+1\) or \(3+2\).  At the first joining coordinate \(r\), let
\(\mathbb M_{A\cup B,r}\) be the normalized complete moment heat
operator with all five rows in one local block, and let
\(\mathbb M_{A|B,r}:=\mathbb M_{A,r}\otimes\mathbb M_{B,r}\) be the
normalized product of the two block moments.

\begin{lemma}[Exact selected block covariance]
\label{lem:V3-exact-covariance}
The complete quotient-order-two joining carrier is
\begin{equation}
 \boxed{\;
 \mathbb K_{A|B,r}
 =
 \mathbb M_{A\cup B,r}-\mathbb M_{A|B,r}.\;}
\label{eq:V3-exact-covariance}
\end{equation}
Both moment operators are positive contractions.  No local
moment-partition term may be removed before their difference is formed.
\end{lemma}

\begin{proof}
For the two quotient variables
\[
 Y_{A,r}:=\prod_{i\in A}F_{i,r},
 \qquad
 Y_{B,r}:=\prod_{i\in B}F_{i,r},
\]
the joining carrier is their second cumulant:
\[
 \kappa_2(Y_{A,r},Y_{B,r})
 =
 \mathbb E(Y_{A,r}Y_{B,r})
 -\mathbb EY_{A,r}\,\mathbb EY_{B,r}.
\]
The two terms are precisely the displayed complete moments.  A
normalized compact-group heat moment is a positive spectral multiplier
with spectrum in \([0,1]\); tensor products preserve positivity and
the contraction bound.
\end{proof}

\begin{lemma}[Common nested-isotypic resolution]
\label{lem:V3-common-resolution}
There is an orthogonal resolution at coordinate \(r\) which
simultaneously diagonalizes the Casimirs of \(A\), \(B\), and
\(A\cup B\).  On every resolved block,
\(\mathbb M_{A\cup B,r}\),
\(\mathbb M_{A|B,r}\), and
\(\mathbb K_{A|B,r}\) are scalar spectral multipliers on each
irreducible factor and commute with all physical intertwiners preserving
the resolved labels.  Summing the physical multiplicity blocks costs a
maximum, not their number.
\end{lemma}

\begin{proof}
Let \(J_A,J_B\) be the infinitesimal diagonal actions on the two
disjoint row blocks.  Their components commute with one another.  The
block Casimirs \(C_A=J_A^2\) and \(C_B=J_B^2\) commute with all
components of their own actions, hence with
\[
 C_{A\cup B}=(J_A+J_B)^2.
\]
Thus the three self-adjoint Casimirs admit a common finite-dimensional
spectral resolution.  Each heat moment is a bounded scalar function of
the corresponding Casimir data, so the three displayed operators are
simultaneously block diagonal.

Within a fixed joint spectral block, Schur's lemma makes each heat
multiplier the identity on the irreducible representation factor
tensored with a scalar identity on its resolved multiplicity label.
It therefore commutes with every label-preserving physical intertwiner.
Complete Peter--Weyl pinching is orthogonal; its trace-norm inequality
takes the maximum over multiplicities before the finite source
recombination.
\end{proof}

\begin{remark}[Resolution does not expand the covariance prematurely]
\label{rem:V3-resolution-order}
The common spectral resolution is inserted into both terms of
\eqref{eq:V3-exact-covariance} simultaneously.  The signed difference
is then formed on each block before its square or capacity is taken.
This differs from assigning separate edge marks to the two moment
terms.
\end{remark}

\section{A sharp selected-covariance half-power}
\label{sec:V3-half}

Fix \(a\in A\) without loss of generality and let
\(a_{a,r}\) be its local representation Casimir weight.

\begin{lemma}[Partial-transpose bounds for both covariance moments]
\label{lem:V3-two-moment-PT}
Across the row cut \(a|a^c\),
\begin{align}
 \|\mathbb M_{A\cup B,r}^{\Gamma_{a^c}}\|
 &\le C_{G,s}a_{a,r}^{-1/2},
\label{eq:V3-full-moment-PT}\\
 \|\mathbb M_{A|B,r}^{\Gamma_{a^c}}\|
 &\le C_{G,s}a_{a,r}^{-1/2}.
\label{eq:V3-split-moment-PT}
\end{align}
All complementary modules may be reducible and retain their complete
physical multiplicities.
\end{lemma}

\begin{proof}
For the full moment, group the local module as the irreducible row-\(a\)
module tensored with the complete reducible complementary module.  The
f3 reducible-complement heat card
\(\texttt{thm:V87-reducible-OIPT}\), followed by the Weyl dimension
bound, gives \eqref{eq:V3-full-moment-PT}.

For the split moment, the factor \(\mathbb M_{A,r}\) obeys the same
estimate across \(a|A\setminus\{a\}\).  The whole \(B\)-factor lies on
the complementary side; full transpose preserves its operator norm,
which is at most one.  Tensorization gives
\eqref{eq:V3-split-moment-PT}.
\end{proof}

\begin{theorem}[Selected block-covariance square card]
\label{thm:V3-selected-covariance-half}
The complete operator \eqref{eq:V3-exact-covariance} satisfies
\begin{equation}
 \boxed{\;
 \kappa_{a|a^c}
  (\mathbb K_{A|B,r}^*\mathbb K_{A|B,r})
 \le C_{G,s}a_{a,r}^{-1/2}.\;}
\label{eq:V3-selected-covariance-half}
\end{equation}
If \(a_{a,r}\ge\sigma M\), the right side is at most
\(C_{G,s,\sigma}M^{-1/2}\).
\end{theorem}

\begin{proof}
Put \(U=\mathbb M_{A\cup B,r}\) and
\(V=\mathbb M_{A|B,r}\).  Both are self-adjoint contractions.  The
elementary operator inequality
\[
 (U-V)^*(U-V)\preccurlyeq2U^2+2V^2
\]
follows by positivity of \((U+V)^*(U+V)\).  Since
\(0\preccurlyeq U,V\preccurlyeq I\),
\(U^2\preccurlyeq U\) and \(V^2\preccurlyeq V\).  Monotonicity and
subadditivity of product-state capacity give
\[
 \kappa_{a|a^c}(\mathbb K_{A|B,r}^*\mathbb K_{A|B,r})
 \le2\kappa_{a|a^c}(U)+2\kappa_{a|a^c}(V).
\]
Each capacity is bounded by the appropriate partial-transpose norm.
Apply \Cref{lem:V3-two-moment-PT}.  The cofinal conclusion follows by
substitution.
\end{proof}

\begin{assessment}[Sharpness and cancellation]
\label{ass:V3-half-scope}
The estimate keeps the covariance difference assembled until its
square is formed, but its upper bound does not claim an additional
cancellation between the two moment terms.  Its half-power is exactly
the protected one-slot currency and is sharp on the equal-spin singlet
sector.  This is precisely the one remaining half-power needed when
three distinct complete heat slots have already been extracted.
\end{assessment}

\section{Physical compression and an outside-local payer}
\label{sec:V3-physical}

\begin{lemma}[Selected-covariance central square insertion]
\label{lem:V3-central-square-insertion}
Resolve the selected coordinate as in
\Cref{lem:V3-common-resolution}, retain the complete covariance, and
place the unique payer outside \(r\).  For every fixed prefix-tree
fibre, output block, and square orientation, the normalized word can
be written so that
\begin{equation}
 \kappa_\otimes(\widetilde W^*\widetilde W)
 \le {C_{G,s}\over(\prod_i\Xi_i)^2}\,
 \kappa_{a|a^c}
 \left[
  (\mathbb K_{A|B,r}\otimes\mathcal M_D)^*
  (\mathbb K_{A|B,r}\otimes\mathcal M_D)
 \right],
\label{eq:V3-central-square-insertion}
\end{equation}
or the corresponding left-square inequality.  Here \(D\) is any
disjoint complete source-occurring heat bank already exposed in the
same exact word.
\end{lemma}

\begin{proof}
Perform complete physical pinching after the common nested-isotypic
resolution.  On each resolved block,
\(\mathbb K_{A|B,r}\) is a scalar spectral multiplier by
\Cref{lem:V3-common-resolution}; the complete bank
\(\mathcal M_D\) is central by the f3 bank-commutation theorem.
Both commute with every remaining label-preserving source intertwiner,
while factors on disjoint coordinates commute tensorially.

Move only these commuting factors to the centre of the word.  The
complete prefix carrier, unpaid moments, suffix, and physical
compressions have the fixed-order raw contraction bounds.  The payer
is outside \(r\) and occurs exactly once with its established
one-resolvent norm.  If the word is \(LZQ\), with
\(Z=\mathbb K_{A|B,r}\otimes\mathcal M_D\), centrality gives
\[
 (LZQ)^*(LZQ)
 \preccurlyeq \|L\|^2\|Q\|^2 Z^*Z.
\]
Divide once by the five row masses.  Bipartite capacity across
\(a|a^c\) enlarges the fully row-product test class, proving the
display.  The adjoint orientation is identical.
\end{proof}

\begin{firewall}[Why the common resolution is essential]
\label{fire:V3-common-essential}
A physical compression or recoupling which mixes unresolved nested
Casimir labels need not commute with the covariance square.  Applying
only its ordinary norm would be the forbidden entangler insertion of
\Cref{prop:V2-entangler-no-go}.  The block resolution in
\Cref{lem:V3-common-resolution} must therefore precede the capacity
estimate.
\end{firewall}

\section{The \texorpdfstring{\(Q=\{ac\}\)}{Q=ac} tree floor}
\label{sec:V3-tree-floor}

Assume \(Q=\{ac\}\), so the quotient order is two, and suppose
\begin{equation}
 a_{a,r}\ge\sigma M,
 \qquad
 a_{c,r}\ge\lambda a_{a,r}.
\label{eq:V3-local-cofinal-pair}
\end{equation}
Then \(\theta_{ac,r}\ge\lambda\sigma^2\) and both endpoint row masses
are at least fixed fractions of \(M\).

\begin{lemma}[Locally strong singleton-quotient floor]
\label{lem:V3-single-Q-floor}
Let \(P\) be the three-edge internal prefix forest for \(A|B\).
The literal tree \(T_0=P\sqcup\{ac\}\) satisfies
\begin{equation}
 \boxed{\;
 \theta_{ac,r}\prod_{e\in P}\theta_e^{<r}
 \ge {c_{\sigma,\lambda,a_*}\over
 M\prod_{i=1}^5\Xi_i}.\;}
\label{eq:V3-single-Q-floor}
\end{equation}
This monomial occurs in the exchange-rescue weight
\(\mathfrak x_{T_0,r}\).
\end{lemma}

\begin{proof}
The forest \(P\) consists of a tree on each of \(A,B\).
For type \(3+2\), its endpoint-mass product divided by
\(\prod_i\Xi_i\) is the mass of the centre of the three-vertex path,
at most \(M\).

For type \(4+1\), the singleton is isolated in \(P\) and is one
endpoint of the joining edge \(ac\).  Its mass is at least
\(\min\{\sigma,\lambda\sigma\}M\).  A four-vertex tree has total
positive degree excess at most two, so the numerator excess is at most
\(M^2\); division by the isolated maximal-scale mass leaves at most
\(C_{\sigma,\lambda}M\).

Each prefix edge has an actual occurrence below \(r\), hence its stock
is at least \(a_*^2\) divided by its endpoint row masses.  Multiply
the three floors and
\(\theta_{ac,r}\ge\lambda\sigma^2\), then apply the two displayed
mass-ratio bounds.  The term is the literal role choice in the exchange
atlas.
\end{proof}

\section{Closure of the single-label \texorpdfstring{\(q=3\)}{q=3}
fibre}
\label{sec:V3-closure}

\begin{theorem}[Outside-local \(Q=\{ac\}\), \(q=3\) closure]
\label{thm:V3-single-Q-q3-close}
Fix a quotient-order-two first-connectivity word with
\(Q=\{ac\}\), local cofinality
\eqref{eq:V3-local-cofinal-pair}, three distinct usable cofinal heat
slots \(D_3\), and a payer outside \(r\).  Then both orientations of
the exchange-rescue square card hold:
\begin{equation}
 \kappa_\otimes(\widetilde W^*\widetilde W)
 \le C_{G,s,\sigma,\lambda}\mathfrak x_{T_0,r}^{\,2},
\label{eq:V3-single-Q-q3-close}
\end{equation}
and analogously for \(\widetilde W\widetilde W^*\).
\end{theorem}

\begin{proof}
The three usable slots contribute \(CM^{-3/2}\) across the maximal-row
cut.  The selected covariance contributes the missing \(CM^{-1/2}\)
by \Cref{thm:V3-selected-covariance-half}.  They occupy distinct
tensor coordinates, so their partially transposed norms multiply:
\[
 \kappa_{a|a^c}
 \left[
  (\mathbb K_{A|B,r}\otimes\mathcal M_{D_3})^*
  (\mathbb K_{A|B,r}\otimes\mathcal M_{D_3})
 \right]
 \le C M^{-2}.
\]
Insert this in \eqref{eq:V3-central-square-insertion}.  The result is
at most
\[
 {CM^{-2}\over(\prod_i\Xi_i)^2}.
\]
By \Cref{lem:V3-single-Q-floor}, the last display is bounded by
\(C\mathfrak x_{T_0,r}^2\).
\end{proof}

\begin{corollary}[Completion of the outside-local \(q=3\) cofinal
joining class]
\label{cor:V3-outside-q3-complete}
Every \(q=3\) fibre with an outside-local payer and a locally cofinal
joining endpoint is closed:
\begin{enumerate}[label=\textup{(\roman*)},leftmargin=4em]
\item if a distinct local joining edge exists, use
      \Cref{thm:V2-fresh-q3-close} or
      \Cref{thm:V1-local-cofinal-close};
\item if the sole joining edge is the cofinal label \(ac\), use
      \Cref{thm:V3-single-Q-q3-close}.
\end{enumerate}
\end{corollary}

\begin{proof}
The two alternatives are exhaustive according as
\(Q\setminus\{ac\}\) is nonempty or empty.  The cited theorems cover
them at the fixed payer and orientation; finite tree/output
recombination and \eqref{eq:V1-exchange-sum} give aggregation.
\end{proof}

\begin{assessment}[Remaining \(q=3\) mechanism]
\label{ass:V3-residual}
After V3, a persistent \(q=3\) word must have at least one of:
\begin{enumerate}[label=\textup{(\roman*)},leftmargin=4em]
\item a local payer at the cofinal selected covariance;
\item every cofinal joining endpoint depleted at \(r\);
\item all cofinal mass trapped in another marked role-owner coordinate;
      or
\item failure of the common nested-isotypic central resolution for the
      actual physical separator.
\end{enumerate}
The first is a paid covariance algebra problem; the middle two require
role-owner routing; the last demands an explicit projective bound or a
counterexample.  No additional spanning-tree enumeration is needed.
\end{assessment}

\section{V3 result ledger and V4 acquisition order}
\label{sec:V3-ledger}

\begin{theorem}[Fourth-series V3 result ledger]
\label{thm:V3-results}
Relative to V1--V2, V3 proves:
\begin{enumerate}[label=\textup{(V3.\arabic*)},leftmargin=5.8em]
\item the exact selected two-block covariance identity;
\item its common nested-isotypic resolution without a multiplicity
      count;
\item half-power partial-transpose bounds for both moment terms;
\item the sharp selected-covariance square card;
\item central square insertion with every outside-local payer;
\item the exact one-power literal tree floor for \(Q=\{ac\}\);
\item closure of the single-label \(q=3\) coverage fibre; and
\item completion of every outside-local \(q=3\) fibre having a locally
      cofinal joining endpoint.
\end{enumerate}
\end{theorem}

\begin{proof}
These are
\Cref{lem:V3-exact-covariance,
lem:V3-common-resolution,
lem:V3-two-moment-PT,
thm:V3-selected-covariance-half,
lem:V3-central-square-insertion,
lem:V3-single-Q-floor,
thm:V3-single-Q-q3-close,
cor:V3-outside-q3-complete}.
\end{proof}

\begin{programme}[V4: paid selected covariance]
\label{prog:V3-V4}
\begin{enumerate}[leftmargin=3em]
\item Write the paid quotient-order-two carrier before squaring:
      include the selected output numerator and the sole resolvent on
      the same common nested-isotypic block.
\item Determine whether the paid multiplier is a scalar function of the
      same commuting Casimirs.  If it is, prove the analogue of
      \Cref{thm:V3-selected-covariance-half} directly.
\item If the numerator changes the moment subtraction, derive the exact
      paid covariance identity rather than inserting a resolvent into
      the unpaid formula.
\item Retain all physical multiplicities through orthogonal pinching and
      test both output-dual square orientations.
\item Test sharpness on the protected \(SU(2)\) singlet and on a gapped
      high-output channel.
\item Isolate locally depleted marked-role owners only after the paid
      selected case is settled.  The next companion audit remains V5.
\end{enumerate}
\end{programme}

\begin{firewall}[Claim boundary after fourth-series V3]
\label{fire:V3-claim-boundary}
V3 closes the outside-local \(q=3\) fibres with a locally cofinal
joining endpoint.  It does not close a local payer, locally depleted
role-owner fibres, the \(q=0,1,2\) response cards, global quintic or
all-order MTA, WUFC, CPM, cutoff removal, reflection positivity,
Osterwalder--Schrader reconstruction, construction of the continuum
Yang--Mills measure, or a uniform positive continuum spectral gap.
The full Yang--Mills mass gap has not been proved.
\end{firewall}

\paragraph{Parent-version record.}
V3 was formed from
\texttt{Renewal\_Yang\_Mills\_Mass\_Gap\_Research\_Notes\_V2.tex}.
The V2 parent had \(987\) logical source lines, \(38\,454\) bytes, and
SHA--256
\[
 \texttt{c5e3419ded26a65f18e73e973935e6d0abdcc1d6817aca45f531683f8fe00801}.
\]
The next compulsory companion audit is V5.

% =============================================================================
% END APPENDED FOURTH-SERIES V3 MATERIAL
% =============================================================================

% The V3 document terminator is intentionally disabled here:
% \end{document}

% =============================================================================
% BEGIN APPENDED FOURTH-SERIES V4 MATERIAL
% =============================================================================

\clearpage
\chapter{V4: The paid selected covariance and closure of the local-payer
\texorpdfstring{\(q=3\)}{q=3} fibre}
\label{ch:V4}

\section{Why the payer must be reconstructed before the square}
\label{sec:V4-purpose}

V3 proves the sharp half-power estimate for the unpaid
quotient-order-two covariance
\[
 \mathbb K_{A|B,r}
 =
 \mathbb M_{A\cup B,r}-\mathbb M_{A|B,r},
\]
but deliberately leaves open the case in which the selected coordinate
\(r\) carries the sole output-mass/resolvent debit.  The missing step is
not obtained by multiplying the V3 estimate by an abstract bounded
operator.  The debit is first assigned by the exact scalar
one-resolvent inequality and is then a function of the same local
output Casimir used in the common nested-isotypic resolution.

This chapter writes that operator explicitly.  The nontrivial-output
debit has a polynomial Casimir factor, but the two heat moments in the
complete covariance absorb its square.  On the trivial local output
the allocation numerator vanishes, so that block cannot carry the
local payer.  These two facts close the local-payer \(q=3\) branch
without inserting a denominator into an unpaid square.

\section{The exact locally allocated payer}
\label{sec:V4-local-payer}

Retain the notation of
\Cref{sec:V3-covariance}.  On the common nested-isotypic block
\(\omega=(J_A,J_B,U,\mu)\), write
\[
 c_A=c(J_A),\qquad c_B=c(J_B),\qquad c_U=c(U),
\]
where \(c(\boldsymbol1)=0\), and put \(s=s_\alpha>0\).
The two normalized moment multipliers are
\begin{equation}
 u_\omega=e^{-sc_U},
 \qquad
 v_\omega=e^{-s(c_A+c_B)}.
\label{eq:V4-two-heat-scalars}
\end{equation}
Thus the selected covariance multiplier is
\[
 k_\omega=u_\omega-v_\omega.
\]
Harmless group-dependent normalizations of the Casimir can be absorbed
into \(s\); all constants below are uniform in the representations and
the physical multiplicity label \(\mu\).

\begin{lemma}[Fusion comparison on a nested block]
\label{lem:V4-fusion-comparison}
There is a group-dependent constant \(C_G\) such that every block in
the common resolution obeys
\begin{equation}
 0\le c_U\le C_G(c_A+c_B).
\label{eq:V4-fusion-comparison}
\end{equation}
\end{lemma}

\begin{proof}
Let \(J_A^m,J_B^m\) denote the infinitesimal generators on the two
disjoint tensor factors in an orthonormal basis of the compact
semisimple Lie algebra.  They commute across the factors, and the
diagonal action on \(A\cup B\) is \(J_A^m+J_B^m\).  For every vector
\(\xi\),
\[
 \sum_m\|(J_A^m+J_B^m)\xi\|^2
 \le
 2\sum_m\bigl(\|J_A^m\xi\|^2+\|J_B^m\xi\|^2\bigr).
\]
Restricting this quadratic-form inequality to the joint
\((J_A,J_B,U)\)-isotypic block gives
\eqref{eq:V4-fusion-comparison}, with \(C_G=2\) in the common
normalization.
\end{proof}

\begin{definition}[Local Casimir-resolvent multiplier]
\label{def:V4-local-multiplier}
Let \(P_{0,r}\) be the projection onto the trivial local output at
\(r\).  Define
\begin{equation}
 \mathbb R_{s,r}
 :=
 \sum_{\omega:\,c_U>0}
 {c_U\over1-e^{-sc_U}}P_\omega,
 \qquad
 \mathbb R_{s,r}P_{0,r}=0.
\label{eq:V4-local-multiplier}
\end{equation}
The signed once-paid selected covariance is
\begin{equation}
 \boxed{\;
 \mathbb K^{\rm pay}_{A|B,r}
 :=
 \mathbb R_{s,r}\mathbb K_{A|B,r}.\;}
\label{eq:V4-paid-selected}
\end{equation}
\end{definition}

\begin{lemma}[This is the exact local term in the one-resolvent
allocation]
\label{lem:V4-exact-allocation}
Fix a completely resolved global output
\(\boldsymbol U=(U_e)_{e\in X}\), with
\[
 Q_{\boldsymbol U}=\prod_{e\in X}e^{-sc(U_e)},
 \qquad
 \Xi(\boldsymbol U)=\sum_{e\in X}c(U_e).
\]
The scalar allocation
\begin{equation}
 {\Xi(\boldsymbol U)\over1-Q_{\boldsymbol U}}
 \le
 \sum_{e\in X}{c(U_e)\over1-e^{-sc(U_e)}}
\label{eq:V4-global-allocation}
\end{equation}
assigns to \(r\) precisely the multiplier
\eqref{eq:V4-local-multiplier}.  If \(U_r=\boldsymbol1\), its
\(r\)-summand is zero and the payer is assigned to another coordinate.
\end{lemma}

\begin{proof}
For \(c(U_e)>0\), the product inequality
\[
 Q_{\boldsymbol U}\le e^{-sc(U_e)}
\]
implies
\[
 {c(U_e)\over1-Q_{\boldsymbol U}}
 \le
 {c(U_e)\over1-e^{-sc(U_e)}}.
\]
Sum over \(e\).  If \(c(U_e)=0\), the numerator of the corresponding
summand in \(\Xi(\boldsymbol U)\) is zero; it is omitted before any
division.  On the common local spectral block this is exactly
\eqref{eq:V4-local-multiplier}.  Thus the protected local block is not
given the indeterminate expression \(0/(1-1)\).
\end{proof}

\begin{firewall}[The protected block is not regularized locally]
\label{fire:V4-no-local-regularization}
The value of \(c/(1-e^{-sc})\) at \(c=0\) is not filled in by
continuity.  The exact global numerator has a zero \(r\)-coordinate
summand there, and the one-resolvent allocation has no local-payer
term.  Replacing that zero summand by \(s^{-1}P_{0,r}\) would create a
new payer which is absent from the physical carrier.
\end{firewall}

\section{Heat absorption of the payer square}
\label{sec:V4-absorption}

Let
\[
 c_*:=\inf\{c(U):U\ne\boldsymbol1\}>0.
\]
The positivity of \(c_*\) uses compact semisimplicity and the fixed
Casimir normalization.

\begin{lemma}[Two scalar square absorptions]
\label{lem:V4-two-absorptions}
For fixed \(s>0\), uniformly on every common block with \(c_U>0\),
\begin{align}
 \left({c_U\over1-e^{-sc_U}}\right)^2e^{-2sc_U}
 &\le C_{G,s}e^{-sc_U},
\label{eq:V4-full-absorption}\\
 \left({c_U\over1-e^{-sc_U}}\right)^2
 e^{-2s(c_A+c_B)}
 &\le C_{G,s}e^{-s(c_A+c_B)}.
\label{eq:V4-split-absorption}
\end{align}
\end{lemma}

\begin{proof}
Since \(c_U\ge c_*\),
\[
 1-e^{-sc_U}\ge1-e^{-sc_*}=:\delta_s>0.
\]
Consequently the left side of
\eqref{eq:V4-full-absorption}, divided by its right-side heat factor,
is at most
\[
 \delta_s^{-2}c_U^2e^{-sc_U},
\]
whose supremum on \([c_*,\infty)\) is finite.

For the second inequality put \(h=c_A+c_B\).  By
\Cref{lem:V4-fusion-comparison}, \(c_U\le C_Gh\).  After division by
the right-side heat factor, the left side is at most
\[
 C_G^2\delta_s^{-2}h^2e^{-sh},
\]
which is uniformly bounded on \([0,\infty)\).  This proves both
claims.
\end{proof}

\begin{proposition}[Operator form of square absorption]
\label{prop:V4-operator-absorption}
On the complete common nested-isotypic resolution,
\begin{align}
 \mathbb R_{s,r}^{\,2}\mathbb M_{A\cup B,r}^{\,2}
 &\preccurlyeq
 C_{G,s}\mathbb M_{A\cup B,r},
\label{eq:V4-full-operator-absorption}\\
 \mathbb R_{s,r}^{\,2}\mathbb M_{A|B,r}^{\,2}
 &\preccurlyeq
 C_{G,s}\mathbb M_{A|B,r}.
\label{eq:V4-split-operator-absorption}
\end{align}
Both inequalities remain valid after complete physical multiplicity
pinching.
\end{proposition}

\begin{proof}
All three operators are scalar on every block by
\Cref{lem:V3-common-resolution}.  Their blockwise coefficients are
exactly the two sides of
\eqref{eq:V4-full-absorption} and
\eqref{eq:V4-split-absorption}; on \(P_{0,r}\) the left sides vanish
by definition.  Scalar comparison on every orthogonal block gives the
Loewner inequalities.  Orthogonal physical pinching preserves positive
order and does not sum the block constants.
\end{proof}

\section{The paid covariance half-power}
\label{sec:V4-paid-half}

\begin{theorem}[Paid selected block-covariance square card]
\label{thm:V4-paid-selected-half}
For every \(a\in A\),
\begin{equation}
 \boxed{\;
 \kappa_{a|a^c}\!
 \left[
  (\mathbb K^{\rm pay}_{A|B,r})^*
  \mathbb K^{\rm pay}_{A|B,r}
 \right]
 \le C_{G,s}a_{a,r}^{-1/2}.\;}
\label{eq:V4-paid-selected-half}
\end{equation}
The same estimate holds for the left square.  In particular, if
\(a_{a,r}\ge\sigma M\), the right side is
\(C_{G,s,\sigma}M^{-1/2}\).
\end{theorem}

\begin{proof}
Put
\[
 U=\mathbb M_{A\cup B,r},
 \qquad
 V=\mathbb M_{A|B,r},
 \qquad
 R=\mathbb R_{s,r}.
\]
They commute on the common resolution.  Since
\((U-V)^*(U-V)\preccurlyeq2U^2+2V^2\),
\Cref{prop:V4-operator-absorption} gives
\begin{align*}
 (\mathbb K^{\rm pay}_{A|B,r})^*
 \mathbb K^{\rm pay}_{A|B,r}
 &=
 R^2(U-V)^*(U-V)\\
 &\preccurlyeq
 2R^2U^2+2R^2V^2\\
 &\preccurlyeq C_{G,s}(U+V).
\end{align*}
Take product-state capacity across \(a|a^c\) and use monotonicity and
subadditivity.  The two moment capacities are bounded by their
partial-transpose norms, which satisfy
\eqref{eq:V3-full-moment-PT}--%
\eqref{eq:V3-split-moment-PT}.  This proves
\eqref{eq:V4-paid-selected-half}.

Every operator is self-adjoint and mutually commuting on the resolved
blocks, so the left square has the same block coefficients.  The final
assertion follows by \(a_{a,r}\ge\sigma M\).
\end{proof}

\begin{remark}[Where the half-power comes from]
\label{rem:V4-half-origin}
The payer does not itself produce row decay.  Its squared polynomial
Casimir factor is absorbed by one half of the heat already present in
each complete endpoint moment.  The surviving heat moment then carries
the V3 partial-transpose half-power.  This is why the proof must keep
the exact paid covariance on its common spectral blocks.
\end{remark}

\section{Two mandatory regressions}
\label{sec:V4-regressions}

\begin{proposition}[Equal-spin protected-output regression]
\label{prop:V4-singlet-regression}
For an \(SU(2)\) equal-spin pair \(j\otimes j\), the local singlet
block \(U=0\) makes no contribution to
\(\mathbb K^{\rm pay}_{A|B,r}\).  Every nontrivial output has
\(c_U\ge2\), so the denominator in
\eqref{eq:V4-local-multiplier} is bounded below by
\(1-e^{-2s}\).  Hence no hidden singular singlet debit occurs in
\Cref{thm:V4-paid-selected-half}.
\end{proposition}

\begin{proof}
The first statement is
\(\mathbb R_{s,r}P_{0,r}=0\).  In the usual \(SU(2)\) normalization,
\(c_U=U(U+1)\), and \(U\ge1\) on the remaining integer-spin output
channels.  The asserted gap and denominator bound follow.
\end{proof}

\begin{proposition}[Gapped high-output regression]
\label{prop:V4-high-output-regression}
On blocks with \(sc_U\ge1\),
\begin{equation}
 \left({c_U\over1-e^{-sc_U}}\right)^2e^{-2sc_U}
 \le C s^{-2}e^{-sc_U}.
\label{eq:V4-high-output-regression}
\end{equation}
Thus an arbitrarily large output Casimir improves rather than degrades
the paid square after heat absorption.
\end{proposition}

\begin{proof}
For \(sc_U\ge1\), the denominator is at least \(1-e^{-1}\).  With
\(x=sc_U\),
\[
 c_U^2e^{-2sc_U}
 =
 s^{-2}x^2e^{-x}e^{-sc_U}
 \le Cs^{-2}e^{-sc_U}.
\]
\end{proof}

\section{Central insertion and closure of the local-payer branch}
\label{sec:V4-closure}

\begin{lemma}[Paid selected-covariance central square insertion]
\label{lem:V4-paid-central-insertion}
In the setting of
\Cref{lem:V3-central-square-insertion}, assign the unique payer to
\(r\) by \eqref{eq:V4-global-allocation} and retain all other factors
unpaid.  For every fixed prefix-tree fibre, physical output block, and
square orientation, the normalized word can be written with
\[
 Z^{\rm pay}
 =
 \mathbb K^{\rm pay}_{A|B,r}\otimes\mathcal M_D
\]
at its centre, and
\begin{equation}
 \kappa_\otimes(\widetilde W^*\widetilde W)
 \le
 {C_{G,s}\over(\prod_i\Xi_i)^2}
 \kappa_{a|a^c}\!
 \left[(Z^{\rm pay})^*Z^{\rm pay}\right].
\label{eq:V4-paid-central-insertion}
\end{equation}
The corresponding left-square estimate also holds.
\end{lemma}

\begin{proof}
The local payer \(\mathbb R_{s,r}\) is a scalar function of \(C_U\)
on the same nested-isotypic blocks as the covariance.  It therefore
commutes with \(\mathbb K_{A|B,r}\), the physical block projections,
and every label-preserving intertwiner.  The disjoint complete bank
\(\mathcal M_D\) is central by the earlier bank-commutation theorem.
Consequently \(Z^{\rm pay}\) may be moved as one assembled operator to
the centre of the exact word.

All remaining prefix, suffix, moment, and compression factors are
unpaid contractions in their established order.  The proof of
\Cref{lem:V3-central-square-insertion} now applies verbatim with
\(Z\) replaced by \(Z^{\rm pay}\).  The debit appears inside
\(Z^{\rm pay}\) exactly once; no external resolvent norm is used.
\end{proof}

\begin{theorem}[Local-payer \(Q=\{ac\}\), \(q=3\) closure]
\label{thm:V4-local-payer-q3-close}
Fix the V3 quotient-order-two first-connectivity word with
\[
 Q=\{ac\},\qquad
 a_{a,r}\ge\sigma M,\qquad
 a_{c,r}\ge\lambda a_{a,r},
\]
and three distinct usable cofinal heat slots \(D_3\).  If the sole
payer is allocated to the selected coordinate \(r\), then both
orientations of the exchange-rescue square card hold:
\begin{equation}
 \boxed{\;
 \kappa_\otimes(\widetilde W^*\widetilde W)
 \le C_{G,s,\sigma,\lambda}\mathfrak x_{T_0,r}^{\,2},\;}
\label{eq:V4-local-payer-q3-close}
\end{equation}
and likewise for \(\widetilde W\widetilde W^*\).
\end{theorem}

\begin{proof}
Tensorize
\Cref{thm:V4-paid-selected-half} with the three disjoint slot
half-powers from V2.  Across the maximal-row cut,
\[
 \kappa_{a|a^c}
 \left[(Z^{\rm pay})^*Z^{\rm pay}\right]
 \le
 C M^{-1/2}M^{-3/2}
 =
 CM^{-2}.
\]
Insert this estimate in
\eqref{eq:V4-paid-central-insertion}.  The normalized square is at
most
\[
 {CM^{-2}\over(\prod_i\Xi_i)^2}.
\]
The literal tree floor
\eqref{eq:V3-single-Q-floor} bounds this by
\(C\mathfrak x_{T_0,r}^2\).  The left-square orientation follows
from the left-square part of
\Cref{thm:V4-paid-selected-half,lem:V4-paid-central-insertion}.
\end{proof}

\begin{corollary}[All locally cofinal \(q=3\) selected covariances are
closed]
\label{cor:V4-all-local-cofinal-q3}
Every \(q=3\) fibre having a locally cofinal joining endpoint is
closed, independently of the payer location.
\end{corollary}

\begin{proof}
If the payer allocation chooses \(r\) and \(U_r\ne\boldsymbol1\), use
\Cref{thm:V4-local-payer-q3-close}.  If \(U_r=\boldsymbol1\), the
local numerator vanishes by
\Cref{lem:V4-exact-allocation}, so the corresponding nonzero payer
summand is outside \(r\); use
\Cref{cor:V3-outside-q3-complete}.  Every genuinely outside-local
allocation was already covered by that corollary.  These alternatives
are exhaustive in the scalar one-resolvent sum.
\end{proof}

\begin{firewall}[No multiplication of paid and unpaid covariance
estimates]
\label{fire:V4-no-double-covariance}
The proof uses the paid operator
\(\mathbb R_{s,r}(\mathbb M_{A\cup B,r}-\mathbb M_{A|B,r})\)
once.  It neither multiplies the V3 unpaid capacity by a resolvent norm
nor assigns separate denominators to the two endpoint moments.
The square absorption occurs only after the common signed covariance
and the unique payer have been assembled.
\end{firewall}

\section{V4 result ledger and next role-owner problem}
\label{sec:V4-ledger}

\begin{theorem}[Fourth-series V4 result ledger]
\label{thm:V4-results}
Relative to V1--V3, this chapter proves:
\begin{enumerate}[label=\textup{(V4.\arabic*)},leftmargin=5.8em]
\item the exact coordinatewise form of the locally allocated
      Casimir-resolvent payer;
\item the forced reassignment of a protected trivial-output block;
\item two scalar and operator square-absorption inequalities;
\item the paid selected-covariance half-power in both square
      orientations;
\item protected-singlet and high-output regressions;
\item central insertion of the assembled paid covariance;
\item closure of the local-payer \(Q=\{ac\}\), \(q=3\) fibre; and
\item completion of all locally cofinal \(q=3\) selected-covariance
      fibres, for every payer location.
\end{enumerate}
\end{theorem}

\begin{proof}
These are
\Cref{lem:V4-exact-allocation,
fire:V4-no-local-regularization,
lem:V4-two-absorptions,
prop:V4-operator-absorption,
thm:V4-paid-selected-half,
prop:V4-singlet-regression,
prop:V4-high-output-regression,
lem:V4-paid-central-insertion,
thm:V4-local-payer-q3-close,
cor:V4-all-local-cofinal-q3}.
\end{proof}

\begin{programme}[V5: compulsory audit and depleted cofinal owners]
\label{prog:V4-V5}
\begin{enumerate}[leftmargin=3em]
\item Perform the compulsory read-only audit of the newest active SM
      and NS notes and record exact filenames, sizes, and SHA--256
      digests.
\item Transfer only proof order which is type-correct for the present
      finite-cutoff capacity problem.
\item Return to the V3 residual in which every cofinal joining endpoint
      is depleted at \(r\), or its cofinal mass belongs to another
      marked role-owner coordinate.
\item Write the exact finite role-owner partition before applying a
      norm: selected coordinate, tree witness, payer, separator, and
      the three usable slot owners.
\item Use conservation of the maximal-row mass to obtain either a
      fresh cofinal slot, a locally cofinal joining endpoint, or a
      quantified residual with at most three occupied owner classes.
\item If the last alternative cannot supply four half-powers, move
      upstream to the first coordinate at which two owner classes
      merge and keep that complete covariance assembled.
\end{enumerate}
\end{programme}

\begin{firewall}[Claim boundary after fourth-series V4]
\label{fire:V4-claim-boundary}
V4 closes the local-payer obstruction and hence every \(q=3\) fibre
with a locally cofinal joining endpoint.  It does not close locally
depleted or displaced role-owner fibres, the \(q=0,1,2\) response
cards, global quintic or all-order MTA, WUFC, CPM, cutoff removal,
reflection positivity, Osterwalder--Schrader reconstruction,
construction of the continuum Yang--Mills measure, or a uniform
positive continuum spectral gap.  The full Yang--Mills mass gap has
not been proved.
\end{firewall}

\paragraph{Parent-version record.}
V4 was formed from
\texttt{Renewal\_Yang\_Mills\_Mass\_Gap\_Research\_Notes\_V3.tex}.
The V3 parent had \(1453\) logical source lines, \(56\,226\) bytes,
and SHA--256
\[
 \texttt{e3b1c1f79af360e95aaf621ef0ddc04a2f20f71364147cf5b493df667aac3432}.
\]
The next compulsory companion audit is V5.

% =============================================================================
% END APPENDED FOURTH-SERIES V4 MATERIAL
% =============================================================================

% The V4 document terminator is intentionally disabled here:
% \end{document}

% =============================================================================
% BEGIN APPENDED FOURTH-SERIES V5 MATERIAL
% =============================================================================

\clearpage
\chapter{V5: Companion audit, central role owners, and the exact
three-slot residual}
\label{ch:V5}

\section{Compulsory five-version companion audit}
\label{sec:V5-audit}

The required read-only audit was performed after V4 had validated and
before the V5 proof direction was chosen.

\begin{record}[Companion snapshots used by V5]
\label{rec:V5-companion-snapshots}
The newest active companion notes were:
\begin{enumerate}[leftmargin=3em]
\item
\(\texttt{Renewal\_SM\_Einstein\_Continuum\_Research\_Notes\_V82.tex}\),
with \(31\,974\) logical source lines, \(1\,176\,823\) bytes, and
SHA--256
\[
 \texttt{a12db87fb6ff9910cb4f54b592ba74cee1fde8f1838ce160e5957ad223a7c8b5};
\]
\item
\(\texttt{Renewal\_NS\_Stability\_Research\_Notes\_V31.tex}\),
with \(23\,052\) logical source lines, \(887\,537\) bytes, and
SHA--256
\[
 \texttt{f1121b62238ad5491bb7cf03635ea9934942edf573ed8faaa9ee79e1d38a6696}.
\]
\end{enumerate}
Neither companion file was changed.
\end{record}

\begin{assessment}[Type-correct transfer from SM V82]
\label{ass:V5-SM-transfer}
SM V82 first assigns every small chart to one stopping-cover vertex,
pays the finite nerve by a quantitative energy lower bound, retains
the exact overlap cocycle, and calls its finite descent complex the
physical index only under an explicit comparison theorem.  Counting
charts or writing an abstract complex is not itself the comparison.

The applicable Yang--Mills proof order is analogous but purely
algebraic: assign every coincident role to one physical coordinate,
retain that coordinate's exact complete operator, and use it for a
half-power only after a local centrality or covariance certificate has
been proved.  A finite list of role names is not itself decay.  No
SM determinant, index, probability, or energy estimate is imported.
\end{assessment}

\begin{assessment}[Type-correct transfer from NS V31]
\label{ass:V5-NS-transfer}
NS V31 is unchanged from the preceding audits.  It proves that a flat
probability mark removes stopping-time multiplicity, but exact
unmarking restores the missing first radial moment.  It also records
that a one-use parent sum can remain analytically circular.

For the present proof this forbids replacing the missing fourth
half-power by the number of excluded role owners.  Each physical
coordinate may be used once, and the exact \(M^{-1/2}\) response must
come from the complete operator at that coordinate or from an
assembled upstream covariance.  No Navier--Stokes estimate is imported.
\end{assessment}

\begin{record}[Audit-adjusted direction]
\label{rec:V5-audit-direction}
The audit confirms a two-stage attack.  First enlarge the usable-slot
notion to excluded coordinates whose actual complete local factor is
still central after all coincident roles are assembled.  Second write a
normalized owner-mass ledger for every three-slot survivor.  The latter
separates genuine three-slot saturation from a cofinal noncentral owner
and from mass trapped in a non-source-compatible remainder.
\end{record}

\section{A central role-owner heat card}
\label{sec:V5-central-owner}

V4 treats an excluded coordinate whose exact local factor is a paid
two-block covariance.  We now treat the other central possibility:
the local factor is one complete heat moment, possibly multiplied by
the locally allocated payer and by bounded central operations belonging
to other coincident roles.

\begin{definition}[Central heat owner]
\label{def:V5-central-heat-owner}
Fix a physical coordinate \(d\) and the maximal-row cut \(a|a^c\).
On one complete common physical resolution
\((P_\omega)\), a central heat owner is an exact local factor
\begin{equation}
 \mathbb L_d
 =
 \mathbb S_d\mathbb H_d,
 \qquad
 \mathbb H_d
 =
 \sum_\omega e^{-sh_\omega}P_\omega,
 \qquad
 \mathbb S_d
 =
 \sum_\omega s_\omega P_\omega,
\label{eq:V5-central-owner}
\end{equation}
where
\[
 h_\omega\ge0,\qquad |s_\omega|\le1,
\]
all physical multiplicities are retained, and
\begin{equation}
 \|\mathbb H_d^{\Gamma_{a^c}}\|
 \le C_{G,s}a_{a,d}^{-1/2}.
\label{eq:V5-owner-heat-half}
\end{equation}
If \(c_\omega\) is the output Casimir assigned to \(d\), require
\begin{equation}
 0\le c_\omega\le C_Gh_\omega.
\label{eq:V5-owner-fusion}
\end{equation}
The once-paid factor is
\begin{equation}
 \mathbb L_d^{\rm pay}
 :=
 \mathbb R_{s,d}\mathbb L_d,
 \qquad
 \mathbb R_{s,d}
 =
 \sum_{\omega:\,c_\omega>0}
 {c_\omega\over1-e^{-sc_\omega}}P_\omega,
\label{eq:V5-owner-paid}
\end{equation}
with zero multiplier on \(c_\omega=0\).
\end{definition}

\begin{remark}[What the definition certifies]
\label{rem:V5-owner-definition-scope}
The operator \(\mathbb S_d\) is the product, in its original local
order, of coincident operations already proved to be scalar
contractions on the chosen resolution.  A tree mark, recoupler, or
physical separator which mixes those blocks does not satisfy the
definition merely because its ordinary norm is at most one.
\end{remark}

\begin{theorem}[Unpaid and paid central-owner half-powers]
\label{thm:V5-central-owner-half}
Every central heat owner satisfies both square orientations
\begin{align}
 \kappa_{a|a^c}(\mathbb L_d^*\mathbb L_d)
 &\le C_{G,s}a_{a,d}^{-1/2},
\label{eq:V5-owner-unpaid-half}\\
 \kappa_{a|a^c}
 \bigl((\mathbb L_d^{\rm pay})^*
       \mathbb L_d^{\rm pay}\bigr)
 &\le C_{G,s}a_{a,d}^{-1/2}.
\label{eq:V5-owner-paid-half}
\end{align}
The identical bounds hold for the left squares.
\end{theorem}

\begin{proof}
All factors commute on the complete resolution.  For the unpaid
factor,
\[
 \mathbb L_d^*\mathbb L_d
 =
 \sum_\omega |s_\omega|^2e^{-2sh_\omega}P_\omega
 \preccurlyeq
 \sum_\omega e^{-sh_\omega}P_\omega
 =
 \mathbb H_d.
\]
Capacity monotonicity and
\eqref{eq:V5-owner-heat-half} prove
\eqref{eq:V5-owner-unpaid-half}.

On \(c_\omega>0\), compact semisimplicity gives
\(c_\omega\ge c_*>0\).  By
\eqref{eq:V5-owner-fusion},
\begin{align*}
 &\left({c_\omega\over1-e^{-sc_\omega}}\right)^2
 |s_\omega|^2e^{-2sh_\omega}\\
 &\qquad\le
 {C_G^2h_\omega^2e^{-sh_\omega}
  \over(1-e^{-sc_*})^2}\,
 e^{-sh_\omega}
 \le C_{G,s}e^{-sh_\omega}.
\end{align*}
The coefficient is zero on \(c_\omega=0\).  Hence
\[
 (\mathbb L_d^{\rm pay})^*
 \mathbb L_d^{\rm pay}
 \preccurlyeq C_{G,s}\mathbb H_d,
\]
which proves \eqref{eq:V5-owner-paid-half}.  Scalarity makes the two
square orientations identical.
\end{proof}

\begin{corollary}[Unified certified-owner list]
\label{cor:V5-certified-owner-list}
An excluded cofinal coordinate supplies one legitimate
\(M^{-1/2}\) square response if its assembled local operator is one of:
\begin{enumerate}[label=\textup{(\roman*)},leftmargin=4em]
\item an unpaid central heat owner;
\item a once-paid central heat owner;
\item the unpaid selected two-block covariance of V3; or
\item the once-paid selected two-block covariance of V4.
\end{enumerate}
No role attached to the same coordinate is estimated a second time.
\end{corollary}

\begin{proof}
The first two cases are
\Cref{thm:V5-central-owner-half}; the last two are
\Cref{thm:V3-selected-covariance-half,
thm:V4-paid-selected-half}.  In each case the complete assembled local
operator is the sole object whose square is estimated.
\end{proof}

\section{Closure from a certified fourth coordinate}
\label{sec:V5-fourth-owner}

\begin{theorem}[Certified excluded owner closes a three-slot core]
\label{thm:V5-certified-fourth-close}
Fix a resolved \(q=3\) quintic obstruction and its three distinct
usable cofinal slots \(E_3=\{e_1,e_2,e_3\}\).  Suppose a coordinate
\(d\notin E_3\) obeys
\[
 a_{a,d}\ge\tau M
\]
and its exact local factor belongs to the certified list in
\Cref{cor:V5-certified-owner-list}.  Assume also that \(d\) is not the
coordinate occurrence used to certify the fixed global strong-edge
stock in the comparison tree.  Then the oriented square has the
required \(M^{-2}\) response and satisfies \(XQTRSC_5\).  This
includes coordinates excluded from the original usable-slot definition
solely because they own the selected or payer role.
\end{theorem}

\begin{proof}
The three extracted slots give
\[
 \prod_{\ell=1}^3
 \|\mathbb M_{e_\ell}^{\Gamma_{a^c}}\|
 \le C_{G,s,\sigma}M^{-3/2}.
\]
The certified owner at \(d\) gives
\[
 \kappa_{a|a^c}(\mathbb L_d^*\mathbb L_d)
 \le C_{G,s,\tau}M^{-1/2},
\]
with \(\mathbb L_d\) interpreted as the appropriate paid or unpaid
heat/covariance owner.  These are distinct tensor coordinates.
Tensorization and central insertion therefore give \(CM^{-2}\) for
the complete square factor.  Insert that assembled factor directly in
the normalized raw-word estimate, giving
\[
 {C M^{-2}\over(\prod_i\Xi_i)^2}.
\]
The archived one-strong exchange floor is
\(c/(M\prod_i\Xi_i)\); its square dominates the last display.
This proves \(XQTRSC_5\).

The last assertion concerns bookkeeping only.  Exclusion from the old
usable-slot set prevents role duplication; the present proof instead
assembles every role at \(d\) and applies one certified-owner estimate.
The fixed-strong occurrence is excluded in the hypothesis precisely
because using its scalar strong stock and its heat response together
requires a separate same-coordinate proof.
\end{proof}

\begin{firewall}[An excluded coordinate is not automatically certified]
\label{fire:V5-excluded-not-certified}
Selected, payer, and fixed-strong are role names, not operator bounds.
\Cref{thm:V5-certified-fourth-close} applies only after the complete
local operator at the physical coordinate has been shown to have one
of the four forms in
\Cref{cor:V5-certified-owner-list}.  A noncentral separator cannot be
discarded by an ordinary contraction estimate inside the product-state
square.
\end{firewall}

\begin{firewall}[Fixed-strong coincidence remains open]
\label{fire:V5-fixed-strong-open}
\Cref{thm:V5-certified-fourth-close} does not count the heat
half-power at the same coordinate whose occurrence certifies the fixed
global strong-edge stock.  Such a same-coordinate comparison may be
valid, but it must be proved from the complete local strong packet.
Treating the scalar stock and the heat factor as two independent roles
would violate the one-coordinate/one-estimate rule.
\end{firewall}

\section{The exact owner-mass ledger for a surviving three-slot core}
\label{sec:V5-mass-ledger}

The archived f3 V85 owner theorem assigns all coincident role names to
one physical coordinate without changing the word or the payer.  Use
that assignment here, but do not repeat its proof.  Starting from a
fixed \(q=3\) core, partition the row-\(a\) support into four disjoint
sets:
\begin{align}
 E&:=E_3,
\notag\\
 C&:=\{\hbox{unextracted coordinates certified by
             \Cref{cor:V5-certified-owner-list}}\},
\notag\\
 N&:=\{\hbox{marked owner coordinates not certified there}\},
\notag\\
 R&:=\operatorname{supp}(a)\setminus(E\sqcup C\sqcup N).
\label{eq:V5-owner-partition}
\end{align}
The set \(N\) has cardinality at most a fixed quintic constant \(L_5\):
there are finitely many joining, payer, tree, strong, and separator
role names in one fixed first-connectivity fibre.

Put
\begin{equation}
 \mu_X:={1\over M}\sum_{e\in X}a_{a,e},
 \qquad X\in\{E,C,N,R\}.
\label{eq:V5-owner-masses}
\end{equation}

\begin{lemma}[Exact normalized owner identity]
\label{lem:V5-owner-identity}
The four nonnegative masses satisfy
\begin{equation}
 \boxed{\quad
 \mu_E+\mu_C+\mu_N+\mu_R=1.
 \quad}
\label{eq:V5-owner-identity}
\end{equation}
Moreover, if \(\mu_N\ge\eta\), some \(d\in N\) obeys
\begin{equation}
 a_{a,d}\ge{\eta\over L_5}M.
\label{eq:V5-noncentral-cofinal}
\end{equation}
\end{lemma}

\begin{proof}
The coordinate sets in \eqref{eq:V5-owner-partition} are disjoint and
exhaust the row support, while
\(\sum_ea_{a,e}=\Xi_a=M\).  This proves
\eqref{eq:V5-owner-identity}.  If \(\mu_N\ge\eta\), the sum of at most
\(L_5\) nonnegative masses is at least \(\eta M\), so one summand is
at least its average.
\end{proof}

\begin{theorem}[Five-way alternative for every \(q=3\) core]
\label{thm:V5-five-way}
Fix \(0<\eta<1/3\).  Every resolved \(q=3\) core has at least one of:
\begin{enumerate}[label=\textup{(\alph*)},leftmargin=4em]
\item a certified central atom \(d\in C\) with
      \(a_{a,d}\ge\tau M\) for some fixed \(\tau>0\);
\item a central diffuse counterpacket:
      \(\mu_C\ge\eta\) and
      \(\max_{d\in C}a_{a,d}/M\to0\);
\item a cofinal noncentral owner \(d\in N\) satisfying
      \eqref{eq:V5-noncentral-cofinal};
\item a non-owner remainder with \(\mu_R\ge\eta\); or
\item three-slot saturation,
\begin{equation}
 \mu_E>1-3\eta.
\label{eq:V5-three-slot-saturation}
\end{equation}
\end{enumerate}
In a persistent counterpacket, alternative \textup{(a)} closes by
\Cref{thm:V5-certified-fourth-close}.  If the complete \(C\)-bank in
\textup{(b)} occurs as the central source-compatible bank of the f3
diffuse theorem, then \textup{(b)} closes there.
\end{theorem}

\begin{proof}
If a fixed-fraction atom occurs in \(C\), this is
\textup{(a)}.  Otherwise, on a counterpacket for which
\(\mu_C\ge\eta\), pass through thresholds
\(\tau_m\downarrow0\).  Failure of \textup{(a)} at every fixed
threshold gives
\[
 \max_{d\in C}{a_{a,d}\over M}\longrightarrow0
\]
along a diagonal subsequence, which is \textup{(b)}.

Now suppose \(\mu_C<\eta\).  If
\(\mu_N\ge\eta\), use
\Cref{lem:V5-owner-identity} to obtain \textup{(c)}.  If
\(\mu_R\ge\eta\), this is \textup{(d)}.  If both fail, then
\eqref{eq:V5-owner-identity} gives
\[
 \mu_E
 =
 1-\mu_C-\mu_N-\mu_R
 >
 1-3\eta,
\]
which is \textup{(e)}.

The closure statement for \textup{(a)} is the cited theorem.
Alternative \textup{(b)} has positive row-mass fraction, vanishing
largest coordinate, and the required complete source occurrence under
the additional stated hypothesis; these are exactly the hypotheses of
the archived diffuse terminal.
\end{proof}

\begin{corollary}[Exact form of a surviving \(q=3\) obstruction]
\label{cor:V5-surviving-q3}
After the V3--V5 central certificates and the established
source-compatible diffuse terminal are removed, every \(q=3\) survivor
has at least one of:
\begin{enumerate}[label=\textup{(\roman*)},leftmargin=4em]
\item a cofinal physical coordinate carrying a genuinely noncentral
      assembled role-owner operator;
\item a fixed fraction of maximal-row mass in the non-owner remainder
      \(R\), whose source-compatible complete heat occurrence has not
      been proved; or
\item three-slot saturation \eqref{eq:V5-three-slot-saturation}.
\end{enumerate}
Thus a depleted selected coordinate by itself is no longer the
bottleneck.
\end{corollary}

\begin{proof}
Alternatives \textup{(a)} and the source-compatible part of
\textup{(b)} in
\Cref{thm:V5-five-way} are closed.  The remaining alternatives are
exactly \textup{(c)}--\textup{(e)}.  The selected coordinate, paid or
unpaid, is already certified by V3--V4 whenever it is cofinal.
\end{proof}

\begin{assessment}[Why the saturation branch is genuinely different]
\label{ass:V5-saturation-sharp}
If three extracted coordinates each carry a fixed fraction of \(M\)
and no fourth certified factor is present, their heat squares supply
only \(M^{-3/2}\).  The archived equal-spin \(SU(2)\) test proves this
slot count sharp.  Therefore
\eqref{eq:V5-three-slot-saturation} cannot be closed by subdividing the
same three coordinates into more role names.  Its missing half-power
must come from the assembled quotient/strong-coordinate packet.
\end{assessment}

\section{V5 result ledger and V6 acquisition order}
\label{sec:V5-ledger}

\begin{theorem}[Fourth-series V5 result ledger]
\label{thm:V5-results}
Relative to V1--V4, V5 proves:
\begin{enumerate}[label=\textup{(V5.\arabic*)},leftmargin=5.8em]
\item the compulsory hashed audit of active SM V82 and NS V31;
\item a precise central heat-owner definition retaining all coincident
      roles;
\item unpaid and once-paid central-owner half-power estimates;
\item a unified four-type list of certified cofinal owners;
\item closure of every three-slot core with a distinct certified
      cofinal owner not used as its fixed-strong occurrence;
\item the exact normalized owner-mass identity;
\item an exhaustive five-way three-slot alternative; and
\item reduction of every surviving \(q=3\) core to a cofinal
      noncentral owner, a non-source-compatible remainder, or sharp
      three-slot saturation.
\end{enumerate}
\end{theorem}

\begin{proof}
These are
\Cref{rec:V5-companion-snapshots,
ass:V5-SM-transfer,
ass:V5-NS-transfer,
def:V5-central-heat-owner,
thm:V5-central-owner-half,
cor:V5-certified-owner-list,
thm:V5-certified-fourth-close,
lem:V5-owner-identity,
thm:V5-five-way,
cor:V5-surviving-q3}.
\end{proof}

\begin{programme}[V6: three-slot saturation and the fixed strong
coordinate]
\label{prog:V5-V6}
\begin{enumerate}[leftmargin=3em]
\item In the saturated branch, retain the three extracted heat slots
      as one complete product and do not count their coincident role
      names again.
\item Resolve the fixed strong coordinate together with the
      first-connectivity quotient cumulant.  Determine whether its
      complete local operator is a central owner, a selected
      covariance, or a genuinely noncentral intertwiner packet.
\item If noncentral, write the exact two-sided square with that packet
      before applying an ordinary norm.  Seek a projective
      partial-transpose estimate, not a second scalar strong-edge mark.
\item On the cofinal noncentral-owner branch, classify the finite local
      operator by quotient order and square orientation.
\item On the \(R\)-mass branch, prove or refute source occurrence of a
      complete central subbank before invoking diffuse heat powers.
\item Test all proposed fourth half-powers on the equal-spin singlet
      channel and on the V2 entangler regression.
\end{enumerate}
\end{programme}

\begin{firewall}[Claim boundary after fourth-series V5]
\label{fire:V5-claim-boundary}
V5 closes only \(q=3\) cores with a fourth certified cofinal owner and
reduces the remaining \(q=3\) geometry to three explicit alternatives.
It does not close those alternatives, the \(q=0,1,2\) response cards,
global quintic or all-order MTA, WUFC, CPM, cutoff removal, reflection
positivity, Osterwalder--Schrader reconstruction, construction of the
continuum Yang--Mills measure, or a uniform positive continuum spectral
gap.  The full Yang--Mills mass gap has not been proved.
\end{firewall}

\paragraph{Parent-version record.}
V5 was formed from
\texttt{Renewal\_Yang\_Mills\_Mass\_Gap\_Research\_Notes\_V4.tex}.
The V4 parent had \(1996\) logical source lines, \(74\,337\) bytes,
and SHA--256
\[
 \texttt{3949ac5aa90563ada7d8a4866018e1ae2eb809a040721e6811733d630c57bba9}.
\]
The next compulsory companion audit is V10.

% =============================================================================
% END APPENDED FOURTH-SERIES V5 MATERIAL
% =============================================================================

% The V5 document terminator is intentionally disabled here:
% \end{document}

% =============================================================================
% BEGIN APPENDED FOURTH-SERIES V6 MATERIAL
% =============================================================================

\clearpage
\chapter{V6: Temporal localization of the fixed strong stock}
\label{ch:V6}

\section{Purpose}
\label{sec:V6-purpose}

V5 isolates three possible \(q=3\) survivors.  The sharp saturation
branch already contains three usable cofinal heat slots, so it needs
exactly one more \(M^{-1/2}\) square response.  A fixed strong pair
\(ac\) is also present, but the old usable-slot definition deletes the
coordinate occurrence used to witness that pair.

The deletion is necessary until the actual local operator is known.
It is not necessary when the witness lies in the complete heat suffix:
there the exact local operator is a central heat moment, with the sole
payer included if necessary.  Its half-power may be used once while
the strong overlap remains only a scalar lower bound for the target
tree weight.  This chapter proves that same-coordinate comparison and
then localizes every surviving strong stock to the three extracted
slots, the prefix, or a higher-order selected quotient cumulant.

\section{Temporal decomposition of a strong overlap}
\label{sec:V6-temporal}

Fix a maximal row \(a\), its comparable partner \(c\), and the
first-connectivity coordinate \(r\).  Write
\[
 M=\Xi_a,\qquad \Xi_c\ge m_0M,
\]
and define
\begin{equation}
 \theta_{ac,e}
 :=
 {a_{a,e}a_{c,e}\over M\Xi_c},
 \qquad
 \theta_{ac}=\sum_e\theta_{ac,e}.
\label{eq:V6-local-overlap}
\end{equation}
Its temporal pieces are
\begin{equation}
 \theta_{ac}^{-}(r)
 :=
 \sum_{e<r}\theta_{ac,e},
 \qquad
 \theta_{ac}^{0}(r)
 :=
 \theta_{ac,r},
 \qquad
 \theta_{ac}^{+}(r)
 :=
 \sum_{e>r}\theta_{ac,e}.
\label{eq:V6-temporal-pieces}
\end{equation}

For a fixed \(q=3\) core, let \(E_3\) be its three extracted usable
slots and set
\begin{align}
 \theta_{ac}^{E}
 &:=
 \sum_{e\in E_3}\theta_{ac,e},
\notag\\
 \overline\theta_{ac}^{-}(r)
 &:=
 \sum_{\substack{e<r\\e\notin E_3}}\theta_{ac,e},
\qquad
 \overline\theta_{ac}^{+}(r)
 :=
 \sum_{\substack{e>r\\e\notin E_3}}\theta_{ac,e}.
\label{eq:V6-deleted-temporal-pieces}
\end{align}
The selected coordinate is not in \(E_3\) by definition.

\begin{lemma}[Exact temporal trichotomy]
\label{lem:V6-temporal-trichotomy}
If \(\theta_{ac}\ge\eta\), then
\begin{equation}
 \max\{\theta_{ac}^{-}(r),
       \theta_{ac}^{0}(r),
       \theta_{ac}^{+}(r)\}
 \ge{\eta\over3}.
\label{eq:V6-temporal-trichotomy}
\end{equation}
\end{lemma}

\begin{proof}
The three terms in
\eqref{eq:V6-temporal-pieces} are nonnegative and sum exactly to
\(\theta_{ac}\).  If all were smaller than \(\eta/3\), their sum
would be smaller than \(\eta\).
\end{proof}

\begin{lemma}[Extracted-slot-deleted temporal trichotomy]
\label{lem:V6-deleted-temporal}
If \(\theta_{ac}\ge\eta\), then
\begin{equation}
 \max\{\theta_{ac}^{E},
       \overline\theta_{ac}^{-}(r),
       \theta_{ac}^{0}(r),
       \overline\theta_{ac}^{+}(r)\}
 \ge{\eta\over4}.
\label{eq:V6-deleted-temporal}
\end{equation}
\end{lemma}

\begin{proof}
The four displayed nonnegative quantities partition the complete
coordinate sum \(\theta_{ac}\): every coordinate is either one of the
three extracted slots, lies strictly before \(r\) outside those slots,
equals \(r\), or lies strictly after \(r\) outside those slots.
Their sum is therefore at least \(\eta\), so one is at least
\(\eta/4\).
\end{proof}

\begin{lemma}[Overlap forces row mass]
\label{lem:V6-overlap-forces-mass}
For every coordinate set \(D\),
\begin{align}
 \sum_{e\in D}a_{a,e}
 &\ge M\sum_{e\in D}\theta_{ac,e},
\label{eq:V6-a-mass-from-overlap}\\
 \sum_{e\in D}a_{c,e}
 &\ge \Xi_c\sum_{e\in D}\theta_{ac,e}.
\label{eq:V6-c-mass-from-overlap}
\end{align}
At the selected coordinate, if
\(\theta_{ac,r}\ge\beta\), then
\begin{equation}
 a_{a,r}\ge\beta M,
 \qquad
 a_{c,r}\ge\beta\Xi_c\ge\beta m_0M.
\label{eq:V6-local-two-row-cofinal}
\end{equation}
\end{lemma}

\begin{proof}
Since \(a_{c,e}\le\Xi_c\),
\[
 \theta_{ac,e}\le{a_{a,e}\over M}.
\]
Sum over \(D\) to obtain
\eqref{eq:V6-a-mass-from-overlap}; the other inequality is symmetric.
For \(D=\{r\}\), write
\[
 \theta_{ac,r}
 =
 {a_{a,r}\over M}{a_{c,r}\over\Xi_c}.
\]
Both factors lie in \([0,1]\).  If their product is at least
\(\beta\), each factor is at least \(\beta\), proving
\eqref{eq:V6-local-two-row-cofinal}.
\end{proof}

\begin{lemma}[A strong overlap contains a cofinal physical occurrence]
\label{lem:V6-strong-occurrence}
Fix \(0<\sigma<\eta\) and suppose \(\theta_{ac}\ge\eta\).  Among the
coordinates satisfying \(a_{a,e}\ge\sigma M\), one has
\begin{equation}
 {a_{c,e}\over\Xi_c}\ge\eta-\sigma.
\label{eq:V6-strong-occurrence}
\end{equation}
There are at most \(\lfloor\sigma^{-1}\rfloor\) candidate
coordinates.
\end{lemma}

\begin{proof}
Let
\[
 E_\sigma=\{e:a_{a,e}\ge\sigma M\}.
\]
On its complement,
\[
 \sum_{e\notin E_\sigma}\theta_{ac,e}
 =
 \sum_{e\notin E_\sigma}
 {a_{a,e}\over M}{a_{c,e}\over\Xi_c}
 \le
 \sigma\sum_e{a_{c,e}\over\Xi_c}
 =
 \sigma.
\]
Hence
\[
 \sum_{e\in E_\sigma}
 {a_{a,e}\over M}{a_{c,e}\over\Xi_c}
 \ge\eta-\sigma.
\]
Because \(\sum_{e\in E_\sigma}a_{a,e}/M\le1\), the largest
\(a_{c,e}/\Xi_c\) on this set is at least \(\eta-\sigma\).
Finally
\(|E_\sigma|\sigma M\le\sum_ea_{a,e}=M\).
\end{proof}

\section{Same-coordinate comparison is not duplicated decay}
\label{sec:V6-same-coordinate}

\begin{proposition}[Central strong occurrence can also carry its heat
half-power]
\label{prop:V6-same-coordinate}
Fix a \(q=3\) obstruction, its three usable slots \(E_3\), and a
strong pair with
\[
 \theta_{ac}\ge\eta,\qquad
 \Xi_a,\Xi_c\ge m_0M.
\]
Let \(d\notin E_3\) obey \(a_{a,d}\ge\tau M\).  Suppose the exact
local factor at \(d\), after all coincident roles are assembled, is
either:
\begin{enumerate}[label=\textup{(\roman*)},leftmargin=4em]
\item an unpaid or once-paid central heat owner of
      \Cref{def:V5-central-heat-owner}; or
\item an unpaid or once-paid selected two-block covariance of V3--V4.
\end{enumerate}
Then \(XQTRSC_5\) holds even when the occurrence \(d\) is also used to
certify the scalar strong stock \(\theta_{ac}\).
\end{proposition}

\begin{proof}
Estimate the complete assembled operator at \(d\) once.  By
\Cref{thm:V5-central-owner-half} in case \textup{(i)}, and by
\Cref{thm:V3-selected-covariance-half,
thm:V4-paid-selected-half} in case \textup{(ii)}, its oriented square
capacity is at most \(C_{\tau}M^{-1/2}\).  The three disjoint usable
slots contribute \(CM^{-3/2}\).  Tensorization and the applicable
central insertion theorem give the normalized word-square bound
\begin{equation}
 {C M^{-2}\over(\prod_i\Xi_i)^2}.
\label{eq:V6-same-coordinate-upper}
\end{equation}

Independently, the exchange atlas contains a spanning-tree monomial
with the global scalar edge \(\theta_{ac}\) and three actual source
edges.  The archived one-strong forest calculation gives
\[
 \mathfrak x_{T_0,r}
 \ge {c_{\eta,m_0}\over M\prod_i\Xi_i}.
\]
Its square dominates
\eqref{eq:V6-same-coordinate-upper}.  The occurrence \(d\) is not
subjected to two operator estimates: its complete local factor is
estimated once.  The inequality \(\theta_{ac}\ge\eta\) is only a
lower bound for the scalar target on the right side.  No product of
two local capacities at \(d\) is formed.
\end{proof}

\begin{firewall}[Exact local form is still mandatory]
\label{fire:V6-same-coordinate-scope}
\Cref{prop:V6-same-coordinate} removes the bookkeeping exclusion only
for the two displayed complete local forms.  It does not turn a
noncentral prefix recoupler, marked separator, or higher quotient
cumulant into a heat owner.  For those factors the square cannot be
reordered by citing the scalar overlap.
\end{firewall}

\section{The complete suffix strong branch closes}
\label{sec:V6-suffix}

\begin{theorem}[Suffix localization of the strong stock closes every
\(q=3\) counterpacket]
\label{thm:V6-suffix-close}
Consider a persistent resolved \(q=3\) counterpacket with extracted
slot set \(E_3\) and
\(\theta_{ac}\ge\eta\).  If, along a subsequence,
\begin{equation}
 \overline\theta_{ac}^{+}(r)\ge\beta>0,
\label{eq:V6-suffix-strong}
\end{equation}
then the packet satisfies \(XQTRSC_5\), a contradiction.  The
conclusion includes a payer inside the suffix.
\end{theorem}

\begin{proof}
Let
\[
 D=\{e:e>r,\ e\notin E_3,\ a_{a,e}>0\}.
\]
By
\Cref{lem:V6-overlap-forces-mass},
\begin{equation}
 S_a(D):=\sum_{e\in D}a_{a,e}\ge\beta M.
\label{eq:V6-suffix-row-mass}
\end{equation}
Let \(p\) be the payer coordinate when it belongs to \(D\).

First suppose that, along a subsequence,
\[
 a_{a,p}\ge{\beta\over2}M.
\]
The exact suffix factor at \(p\) is a once-paid central heat owner.
Apply
\Cref{prop:V6-same-coordinate}.

Otherwise delete \(p\); along a further subsequence,
\[
 S_a(D\setminus\{p\})\ge{\beta\over2}M.
\]
If
\[
 \limsup\max_{e\in D\setminus\{p\}}
 {a_{a,e}\over M}>0,
\]
choose a cofinal coordinate \(d\) realizing a fixed positive
subsequential lower bound.  Its exact suffix factor is an unpaid
central heat owner and
\Cref{prop:V6-same-coordinate} applies.

In the remaining case,
\[
 \max_{e\in D\setminus\{p\}}
 {a_{a,e}\over M}\longrightarrow0.
\]
The complete product over \(D\setminus\{p\}\) occurs centrally in the
same suffix word, is disjoint from the selected and payer coordinates,
and carries at least \(\beta M/2\) row mass.  It is therefore exactly
the source-compatible diffuse bank of the archived f3 V97 terminal.
That theorem supplies an arbitrarily high square power and closes the
word.  The three alternatives exhaust
\eqref{eq:V6-suffix-row-mass}.
\end{proof}

\begin{remark}[Why suffix occurrence is decisive]
\label{rem:V6-suffix-decisive}
For \(e>r\), the exact first-connectivity formula contains the full
moment suffix \(\mathcal M_{>r}\).  Thus source occurrence, tensor
separation from the selected coordinate, and physical multiplicity
completion are identities, not inferred from the scalar overlap.
They are unavailable automatically for \(e<r\), where the coordinate
lies inside a connected prefix carrier.
\end{remark}

\section{The selected covariance strong branch}
\label{sec:V6-selected}

\begin{theorem}[Selected strong phase closes at quotient order two]
\label{thm:V6-selected-k2-close}
Let \(q=3\), let \(\rho=A|B\), and assume
\begin{equation}
 \theta_{ac,r}\ge\beta>0.
\label{eq:V6-selected-strong}
\end{equation}
Then the quotient-order-two word satisfies \(XQTRSC_5\) for every
payer location, even when \(r\) is the fixed strong occurrence.
\end{theorem}

\begin{proof}
By \eqref{eq:V6-local-two-row-cofinal},
\[
 a_{a,r}\ge\beta M.
\]
The exact selected local factor is the covariance
\(\mathbb K_{A|B,r}\).  It has the unpaid V3 half-power when the payer
is outside \(r\), and the paid V4 half-power when the payer is at
\(r\).  Apply
\Cref{prop:V6-same-coordinate} with \(d=r\).
The argument uses only the maximal row \(a\); it does not require the
chosen lift of the quotient edge to have label \(ac\).
\end{proof}

\begin{firewall}[Higher quotient order is not a covariance]
\label{fire:V6-higher-selected-open}
When \(|\rho|=3,4,5\), the selected factor is respectively a complete
third, fourth, or fifth quotient cumulant.  V3--V4 prove no
partial-transpose half-power for those assembled operators.
\Cref{thm:V6-selected-k2-close} may not be extended by expanding such
a cumulant into moment-partition terms and estimating one term.
\end{firewall}

\section{Upstream localization theorem}
\label{sec:V6-upstream}

\begin{theorem}[Every surviving strong stock is upstream or owned by
the extracted slots]
\label{thm:V6-upstream-localization}
Fix a persistent \(q=3\) obstruction with
\(\theta_{ac}\ge\eta\).  Then, after passage to a subsequence:
\begin{enumerate}[label=\textup{(\roman*)},leftmargin=4em]
\item if \(|\rho|=2\), every survivor satisfies
\begin{equation}
 \boxed{\quad
 \max\{\theta_{ac}^{E},
       \overline\theta_{ac}^{-}(r)\}
 \ge{\eta\over4};
 \quad}
\label{eq:V6-k2-prefix-strong}
\end{equation}
\item if \(|\rho|\in\{3,4,5\}\), every survivor satisfies at least one
      of
\begin{equation}
 \theta_{ac}^{E}\ge{\eta\over4},
 \qquad
 \overline\theta_{ac}^{-}(r)\ge{\eta\over4},
 \qquad
 \theta_{ac}^{0}(r)\ge{\eta\over4},
\label{eq:V6-higher-upstream}
\end{equation}
where the second alternative is a cofinal higher quotient cumulant.
\end{enumerate}
\end{theorem}

\begin{proof}
By
\Cref{lem:V6-deleted-temporal}, one of the four owner-disjoint pieces
is at least \(\eta/4\).  The deleted-suffix alternative is impossible
for a survivor by
\Cref{thm:V6-suffix-close}.  At quotient order two, the selected
alternative is impossible by
\Cref{thm:V6-selected-k2-close}; hence only the extracted-slot-owned
or deleted-prefix alternative remains.  At higher quotient order the
selected alternative remains open by
\Cref{fire:V6-higher-selected-open}, yielding
\eqref{eq:V6-higher-upstream}.
\end{proof}

\begin{corollary}[The 37/14 census for the deleted-prefix alternative]
\label{cor:V6-census-meaning}
For the fixed strong pair \(ac\), the deleted-prefix branch splits
into the \(14\) proper partitions in which \(a,c\) lie in one prefix
block and the \(37\) partitions in which they lie in different prefix
blocks, as counted in V1.  In the first class the strong prefix stock
belongs to one lower connected block of arity at most four.  In the
second it spans two prefix blocks and must remain coupled to the
selected quotient cumulant; it cannot be assigned to either prefix
carrier separately.
\end{corollary}

\begin{proof}
The numerical split is the exact V1 fixed-pair partition census.  If
\(a,c\) lie in one block \(B\in\rho\), both rows belong to its complete
connected prefix carrier.  If they lie in distinct blocks, the exact
first-connectivity source contains a tensor product of separate
connected prefix carriers until \(r\); no cross-block connected prefix
factor exists before the selected quotient cumulant joins the blocks.
\end{proof}

\begin{assessment}[The actual V7 target]
\label{ass:V6-V7-target}
The suffix outside the three extracted slots is no longer part of the
\(q=3\) obstruction.  At quotient order two, neither is the selected
coordinate.  The remaining work is:
\begin{enumerate}[label=\textup{(\alph*)},leftmargin=4em]
\item strong overlap carried by the same three extracted slots;
\item internal strong stock inside one already-constructed lower-order
      connected prefix block; or
\item crossing strong stock distributed between distinct prefix blocks
      before the first joining covariance.
\end{enumerate}
At higher quotient order there is one additional finite local target:
the cofinal selected third-, fourth-, or fifth-cumulant square card.
\end{assessment}

\section{V6 result ledger and V7 acquisition order}
\label{sec:V6-ledger}

\begin{theorem}[Fourth-series V6 result ledger]
\label{thm:V6-results}
Relative to V1--V5, V6 proves:
\begin{enumerate}[label=\textup{(V6.\arabic*)},leftmargin=5.8em]
\item the exact prefix/selected/suffix decomposition of every fixed
      strong overlap and its extracted-slot-deleted refinement;
\item quantitative conversion from overlap stock to row mass;
\item localization of a strong overlap to finitely many cofinal
      physical occurrences;
\item a same-coordinate theorem permitting one certified local
      operator response together with its scalar strong-tree target;
\item closure of every \(q=3\) branch with fixed strong stock in the
      complete suffix;
\item closure of the selected strong phase at quotient order two for
      all payer locations;
\item localization of every surviving strong stock to the extracted
      slots, the prefix, or a higher selected cumulant; and
\item the analytic interpretation of the V1 \(37/14\) partition
      census.
\end{enumerate}
\end{theorem}

\begin{proof}
These are
\Cref{lem:V6-temporal-trichotomy,
lem:V6-deleted-temporal,
lem:V6-overlap-forces-mass,
lem:V6-strong-occurrence,
prop:V6-same-coordinate,
thm:V6-suffix-close,
thm:V6-selected-k2-close,
thm:V6-upstream-localization,
cor:V6-census-meaning}.
\end{proof}

\begin{programme}[V7: the prefix strong-stock interface]
\label{prog:V6-V7}
\begin{enumerate}[leftmargin=3em]
\item First analyze the extracted-slot-owned alternative
      \(\theta_{ac}^{E}\ge\eta/4\).  Determine whether the two-row
      structure inside one of the same three heat slots improves its
      one-row partial-transpose card; do not count it as a fourth
      coordinate.
\item In the \(14\) internal-pair partitions, isolate the unique prefix
      block containing \(a,c\), retain its complete connected carrier,
      and test whether the proved order-\(\le4\) square atlas supplies
      one maximal-row half-power without spending its tree stock twice.
\item In the \(37\) crossing-pair partitions, keep the two relevant
      prefix carriers and the selected quotient cumulant as one packet.
      Write the exact common-output resolution before any product-state
      norm.
\item Separate \(k=2\) from \(k=3,4,5\).  At \(k=2\), use the exact
      covariance difference; at higher order retain the complete
      Möbius sum.
\item Test a proposed prefix half-power on an equal-spin protected
      channel placed entirely below \(r\).
\item If the lower connected carrier has no standalone half-power,
      move to its own first-connectivity coordinate and use the proved
      arity-\(\le4\) atlas in the original square orientation.
\end{enumerate}
\end{programme}

\begin{firewall}[Claim boundary after fourth-series V6]
\label{fire:V6-claim-boundary}
V6 closes the strong-stock branch in the suffix after the three
extracted slots are deleted, and the selected quotient-order-two
strong branch of the \(q=3\) core.  It does not close strong stock
owned by those same three slots, prefix strong stock, cofinal higher
quotient cumulants,
noncentral owner or non-source-compatible remainder branches, the
\(q=0,1,2\) response cards, global quintic or all-order MTA, WUFC,
CPM, cutoff removal, reflection positivity,
Osterwalder--Schrader reconstruction, construction of the continuum
Yang--Mills measure, or a uniform positive continuum spectral gap.
The full Yang--Mills mass gap has not been proved.
\end{firewall}

\paragraph{Parent-version record.}
V6 was formed from
\texttt{Renewal\_Yang\_Mills\_Mass\_Gap\_Research\_Notes\_V5.tex}.
The V5 parent had \(2522\) logical source lines, \(93\,484\) bytes,
and SHA--256
\[
 \texttt{ae1727e36347e4e7f1c1e922c45cc1900cdca2d6a9ca912f4710128f5edc4a51}.
\]
The next compulsory companion audit is V10.

% =============================================================================
% END APPENDED FOURTH-SERIES V6 MATERIAL
% =============================================================================

% The V6 document terminator is intentionally disabled here:
% \end{document}

% =============================================================================
% BEGIN APPENDED FOURTH-SERIES V7 MATERIAL
% =============================================================================

\clearpage
\chapter{V7: Sharp failure of same-slot reuse and finite prefix
descent}
\label{ch:V7}

\section{Purpose}
\label{sec:V7-purpose}

The corrected V6 localization leaves three possible homes for a fixed
strong stock: the same three already extracted heat slots, the prefix,
or a higher-order selected quotient cumulant.  This chapter first
settles the same-slot question negatively.  An explicit \(SU(2)\)
singlet family has fixed strong overlap entirely on the three slots but
only the sharp \(M^{-3/2}\) three-slot square response.  Strong
two-row input geometry inside an already counted heat slot is therefore
not a fourth response.

For the prefix-internal \(14\)-partition class, we then go upstream to
the first-connectivity coordinate of the unique prefix block containing
the strong pair.  The strong stock has an exact three-way decomposition
relative to that nested coordinate.  A complete nested suffix closes;
a selected covariance closes; otherwise the block size decreases or
the pair becomes crossing.  Since the starting block has at most four
rows, this descent terminates after at most two strict block reductions.

\section{The three-slot strong-overlap countertest}
\label{sec:V7-three-slot-test}

Work with \(G=SU(2)\).  Let \(j\) be a positive integer spin,
\[
 d_j=2j+1,\qquad A_j=1+j(j+1),
\]
and take three distinct coordinates \(E_3=\{e_1,e_2,e_3\}\).  At
each coordinate rows \(a,c\) carry spin \(j\), and all other rows are
trivial.  Put
\begin{equation}
 \Xi_a=\Xi_c=3A_j=:M_j.
\label{eq:V7-three-slot-mass}
\end{equation}
Let \(P_{0,j}\) be the normalized singlet projection in
\(V_j\otimes V_j\).

\begin{lemma}[Fixed strong stock is owned entirely by the three slots]
\label{lem:V7-three-slot-strong}
For the family above,
\begin{equation}
 \theta_{ac}^{E}
 =
 \theta_{ac}
 =
 {1\over3}.
\label{eq:V7-three-slot-strong}
\end{equation}
Every coordinate is \((1/3,M_j)\)-cofinal on both rows.
\end{lemma}

\begin{proof}
At each \(e_\ell\),
\[
 {a_{a,e_\ell}\over\Xi_a}
 =
 {a_{c,e_\ell}\over\Xi_c}
 =
 {A_j\over3A_j}
 =
 {1\over3}.
\]
Therefore
\[
 \theta_{ac}
 =
 \sum_{\ell=1}^{3}
 {A_j^2\over(3A_j)^2}
 =
 {1\over3}.
\]
There are no other nontrivial coordinates.
\end{proof}

\begin{lemma}[Exact product-state capacity of the protected sector]
\label{lem:V7-singlet-capacity}
Across the row cut \(a|c\),
\begin{equation}
 \kappa_{a|c}(P_{0,j})={1\over d_j},
 \qquad
 \kappa_{a|c}
 \left(\bigotimes_{\ell=1}^{3}P_{0,j}\right)
 ={1\over d_j^3}.
\label{eq:V7-singlet-capacity}
\end{equation}
\end{lemma}

\begin{proof}
Choose an orthonormal basis and write the normalized invariant vector
as
\[
 \Omega_j=d_j^{-1/2}\sum_{m=1}^{d_j}e_m\otimes e_m^*
\]
after the standard unitary identification of the dual module.  For
unit vectors \(x,y\),
\[
 |\langle x\otimes y,\Omega_j\rangle|^2
 \le {1\over d_j}
\]
by Cauchy--Schwarz, and equality is attained by a matched basis pair.
This is the first capacity identity.  Product-state capacity
tensorizes for positive tensor products, giving the second.
\end{proof}

\begin{proposition}[Strong overlap does not improve three counted heat
slots]
\label{prop:V7-no-same-slot-fourth}
Let \(\mathbb H_{e_\ell}(s)\) be the complete two-row heat moment at
\(e_\ell\).  Its protected sector contains \(P_{0,j}\) with multiplier
one.  Consequently
\begin{equation}
 \kappa_{a|c}
 \left[
  \left(\bigotimes_{\ell=1}^{3}\mathbb H_{e_\ell}(s)\right)^*
  \left(\bigotimes_{\ell=1}^{3}\mathbb H_{e_\ell}(s)\right)
 \right]
 \ge {1\over d_j^3}
 \asymp M_j^{-3/2}.
\label{eq:V7-three-slot-lower}
\end{equation}
There is no representation-uniform heat-only estimate by
\(CM_j^{-2}\), despite
\(\theta_{ac}^{E}=1/3\).
\end{proposition}

\begin{proof}
On the total-spin-zero block the Casimir heat multiplier is
\(e^{-s0}=1\), so
\[
 \mathbb H_{e_\ell}(s)\succeq P_{0,j}.
\]
The square has the same action on that projection.  Positive
tensorization and
\Cref{lem:V7-singlet-capacity} give the lower bound.  Since
\(M_j=3(1+j(j+1))\asymp j^2\) and \(d_j\asymp j\),
\[
 M_j^2d_j^{-3}\asymp j\longrightarrow\infty.
\]
Thus no fixed constant can replace the right side of
\eqref{eq:V7-three-slot-lower} by \(CM_j^{-2}\).
\end{proof}

\begin{firewall}[Scope of the countertest]
\label{fire:V7-countertest-scope}
\Cref{prop:V7-no-same-slot-fourth} does not disprove
\(XQTRSC_5\).  It proves that the fixed strong scalar
\(\theta_{ac}^{E}\) cannot upgrade three already counted complete heat
slots to four heat half-powers.  Any closure of this branch must use an
additional source-occurring operator, normally the assembled quotient
cumulant or an upstream connected-prefix response.
\end{firewall}

\section{Why global lower-order MTA is not the missing half-power}
\label{sec:V7-MTA-not-half}

\begin{proposition}[Tree control does not imply a row half-power]
\label{prop:V7-MTA-nonimplication}
The proved arity-\(\le4\) marked-tree atlas
\[
 \kappa_\otimes^{\rm abs}(\mathcal C_B)
 \le
 C\sum_{\tau\in\mathfrak T_B}
 \prod_{ij\in E(\tau)}\theta_{ij}
\]
does not, by itself, imply
\[
 \kappa_{a|a^c}(\mathcal C_B^*\mathcal C_B)
 \le C M^{-1/2}
\]
when one internal edge stock is strong.
\end{proposition}

\begin{proof}
The tree-atlas right side is dimensionless.  Its hypotheses permit a
fixed tree monomial with all of its normalized edge stocks bounded
below independently of \(M\); in particular a strong edge
\(\theta_{ac}\ge\eta\) makes that right side no smaller.  The displayed
MTA estimate then supplies only an \(O(1)\) upper bound.  There is no
algebraic implication from an \(O(1)\) upper bound to
\(O(M^{-1/2})\).

This is a statement about the information content of the proved
inequality, not a counterexample to a stronger prefix theorem.  To
obtain the half-power one must retain a heat or covariance operator
inside the prefix before applying the tree atlas.
\end{proof}

\section{Nested first connectivity inside an internal prefix block}
\label{sec:V7-nested}

Fix one of the \(14\) global prefix partitions in which \(a,c\) belong
to the same block \(B\in\rho\).  Then \(2\le |B|\le4\).  In one fixed
coordinate word of the complete connected carrier
\(\mathcal K_B^{<r,\mathrm{conn}}\), let \(s<r\) be the first
coordinate at which \(B\) becomes connected and let
\(\sigma<\widehat1_B\) be its cumulative partition immediately before
\(s\).  The exact nested first-connectivity factor is
\begin{equation}
 \left(\bigotimes_{D\in\sigma}
       \mathcal K_D^{<s,\mathrm{conn}}\right)
 \otimes\mathcal J_s^\sigma
 \otimes\mathcal M_B^{(s,r)}.
\label{eq:V7-nested-source}
\end{equation}
Singleton \(D\)'s carry complete singleton moments, and
\(\mathcal M_B^{(s,r)}\) is the complete \(B\)-moment suffix on
\(s<e<r\).

Delete the three already extracted slots \(E_3\).  Define
\begin{align}
 \vartheta_{<s}
 &:=
 \sum_{\substack{e<s\\e\notin E_3}}\theta_{ac,e},
\notag\\
 \vartheta_s
 &:=
 \begin{cases}
 \theta_{ac,s},&s\notin E_3,\\
 0,&s\in E_3,
 \end{cases}
\notag\\
 \vartheta_{(s,r)}
 &:=
 \sum_{\substack{s<e<r\\e\notin E_3}}\theta_{ac,e}.
\label{eq:V7-nested-overlap}
\end{align}

\begin{lemma}[Exact nested temporal split]
\label{lem:V7-nested-split}
One has
\begin{equation}
 \overline\theta_{ac}^{-}(r)
 =
 \vartheta_{<s}+\vartheta_s+\vartheta_{(s,r)}.
\label{eq:V7-nested-split}
\end{equation}
If the left side is at least \(\beta\), one term on the right is at
least \(\beta/3\).
\end{lemma}

\begin{proof}
The coordinate set
\(\{e<r:e\notin E_3\}\) is the disjoint union of
\(\{e<s\}\setminus E_3\), \(\{s\}\setminus E_3\), and
\(\{s<e<r\}\setminus E_3\).  Sum
\eqref{eq:V6-local-overlap} over this partition.  The final assertion
is the pigeonhole principle.
\end{proof}

\section{Closure of the nested complete-suffix alternative}
\label{sec:V7-nested-suffix}

\begin{lemma}[A cofinal complete nested-suffix owner supplies the fourth
half-power]
\label{lem:V7-nested-owner}
Let \(d\in(s,r)\setminus E_3\) and suppose
\(a_{a,d}\ge\tau M\).  If \(d\) is not the payer, its exact factor in
\eqref{eq:V7-nested-source} is an unpaid central heat owner; if it is
the payer, its exact allocated factor is a once-paid central heat
owner.  In either case the \(q=3\) word closes.
\end{lemma}

\begin{proof}
The interval \((s,r)\) is a complete moment suffix for the block
\(B\).  At a fixed coordinate it is a normalized complete heat moment
on the row-\(a\) module and its full reducible complementary
\(B\setminus\{a\}\) module.  The other global prefix blocks and all
later coordinates are tensor factors.  Complete physical resolution
therefore gives
\eqref{eq:V5-owner-heat-half}; any coincident payer is the central
multiplier
\eqref{eq:V5-owner-paid}.  Apply
\Cref{thm:V5-central-owner-half} and tensorize with the three
extracted slot half-powers.  Direct central insertion and the
one-strong exchange floor give the closure exactly as in
\Cref{prop:V6-same-coordinate}.
\end{proof}

\begin{theorem}[Nested complete-suffix strong stock closes]
\label{thm:V7-nested-suffix-close}
If, along a persistent \(q=3\) counterpacket,
\begin{equation}
 \vartheta_{(s,r)}\ge\gamma>0,
\label{eq:V7-nested-suffix-strong}
\end{equation}
then the packet satisfies \(XQTRSC_5\), a contradiction.
\end{theorem}

\begin{proof}
By
\Cref{lem:V6-overlap-forces-mass}, the bank
\[
 D=\{e:s<e<r,\ e\notin E_3,\ a_{a,e}>0\}
\]
carries row-\(a\) mass at least \(\gamma M\).  If its payer coordinate
carries a fixed fraction, use the paid case of
\Cref{lem:V7-nested-owner}.  Otherwise delete the payer and retain at
least \(\gamma M/2\).

If the largest remaining coordinate has a positive normalized
limsup, pass to a cofinal subsequence and use the unpaid case of
\Cref{lem:V7-nested-owner}.  Otherwise the complete product bank has
positive row-mass fraction and vanishing largest atom.  It occurs
centrally in the exact nested suffix, is disjoint from \(E_3\), the
selected coordinates, and the payer, and hence satisfies the archived
source-compatible diffuse terminal.  These alternatives are
exhaustive.
\end{proof}

\section{Selected nested covariance and strict descent}
\label{sec:V7-descent}

\begin{lemma}[Selected nested covariance alternative]
\label{lem:V7-nested-covariance}
Suppose \(\vartheta_s\ge\gamma>0\) and
\(|\sigma|=2\).  Then the selected nested factor
\(\mathcal J_s^\sigma\) supplies the fourth half-power, for every
payer location, and the \(q=3\) word closes.
\end{lemma}

\begin{proof}
The overlap bound implies \(a_{a,s}\ge\gamma M\).
When \(|\sigma|=2\), the complete selected factor is the two-block
covariance
\[
 \mathbb M_{D_1\cup D_2,s}
 -
 \mathbb M_{D_1|D_2,s}.
\]
Use the V3 unpaid square card if the payer is outside \(s\), and the
V4 paid card if it is at \(s\).  All other factors in
\eqref{eq:V7-nested-source} lie at distinct coordinates or in the
same common nested-isotypic resolution.  The central square insertion,
the three \(E_3\) half-powers, and the scalar one-strong exchange floor
then give the conclusion by
\Cref{prop:V6-same-coordinate}.
\end{proof}

\begin{definition}[Internal descent step]
\label{def:V7-descent-step}
Assume \(\vartheta_{<s}\ge\gamma>0\).
If \(a,c\) lie in one block \(B'\in\sigma\), replace
\[
 (B,r,\overline\theta_{ac}^{-}(r))
 \quad\hbox{by}\quad
 (B',s,\vartheta_{<s}).
\label{eq:V7-descent-step}
\]
If they lie in distinct blocks of \(\sigma\), stop and call the
result a \emph{crossing-prefix packet}.  If
\(\vartheta_s\ge\gamma\) but \(|\sigma|\ge3\), stop and call it a
\emph{cofinal higher selected cumulant}.
\end{definition}

\begin{lemma}[Strict decrease of the internal block]
\label{lem:V7-strict-decrease}
Every internal descent step satisfies
\begin{equation}
 2\le|B'|\le|B|-1.
\label{eq:V7-strict-decrease}
\end{equation}
Starting from a quintic proper-prefix block, there are at most two
successive internal descent steps.
\end{lemma}

\begin{proof}
The block \(B'\) contains the two distinct rows \(a,c\), so
\(|B'|\ge2\).  The partition \(\sigma\) is proper and has at least one
other nonempty block, so \(B'\subsetneq B\) and
\(|B'|\le|B|-1\).  Initially \(|B|\le4\); after two strict decreases a
block containing two rows has size two and cannot decrease again.
\end{proof}

\begin{theorem}[Finite internal-prefix descent]
\label{thm:V7-finite-descent}
Fix a persistent \(q=3\) packet in the V1 \(14\)-partition
prefix-internal class with
\[
 \overline\theta_{ac}^{-}(r)\ge\beta>0.
\]
After at most two strict internal descent steps, one of the following
occurs:
\begin{enumerate}[label=\textup{(\roman*)},leftmargin=4em]
\item the word closes through a complete nested suffix;
\item the word closes through a selected nested covariance;
\item a crossing-prefix packet remains, carrying strong stock at least
      \(\beta/3^3\); or
\item a cofinal selected quotient cumulant of order \(3\) or \(4\)
      remains, with local overlap at least \(\beta/3^3\).
\end{enumerate}
Thus prefix-internal geometry is not an additional infinite-depth
obstruction.
\end{theorem}

\begin{proof}
At each current node apply
\Cref{lem:V7-nested-split}.  A nested-suffix share closes by
\Cref{thm:V7-nested-suffix-close}.  A selected share with quotient
order two closes by
\Cref{lem:V7-nested-covariance}.  A selected share of higher order is
\textup{(iv)}.  For an earlier-prefix share, either the pair is
separated by the predecessor partition, giving \textup{(iii)}, or it
lies in one smaller block and
\Cref{def:V7-descent-step} applies.

There are at most two strict decreases by
\Cref{lem:V7-strict-decrease}.  Including the initial three-way split
and two possible later splits loses at most a factor \(3^3\) in the
fixed positive overlap.  At a size-two block the predecessor partition
separates \(a,c\), so the next earlier-prefix alternative is crossing,
while its selected alternative has quotient order two and closes.
This exhausts the process.
\end{proof}

\begin{corollary}[Unified residual after V7]
\label{cor:V7-unified-residual}
After V7, every \(q=3\) strong-stock survivor belongs to one of:
\begin{enumerate}[label=\textup{(\alph*)},leftmargin=4em]
\item the extracted-slot-owned branch
      \(\theta_{ac}^{E}\ge\eta/4\), where heat-only reuse is ruled out
      by \Cref{prop:V7-no-same-slot-fourth};
\item a crossing-prefix packet, whether arising directly from the V1
      \(37\)-partition class or after the finite internal descent; or
\item a cofinal selected third-, fourth-, or fifth-cumulant packet.
\end{enumerate}
\end{corollary}

\begin{proof}
V6 removes the deleted suffix and global selected covariance branches.
Its deleted-prefix alternative is either crossing initially or
internal.  The internal case reduces by
\Cref{thm:V7-finite-descent}.  The only other V6 alternatives are
strong stock on \(E_3\) and higher selected cumulants.
\end{proof}

\begin{firewall}[No recursive invocation of quintic MTA]
\label{fire:V7-no-circularity}
The descent uses only coefficient-one first connectivity for blocks of
size at most four, complete heat/covariance cards, and the already
proved source-compatible diffuse theorem.  It does not apply the
desired quintic MTA to a quintic prefix.  The block cardinality strictly
decreases whenever the argument recurses.
\end{firewall}

\section{V7 result ledger and V8 acquisition order}
\label{sec:V7-ledger}

\begin{theorem}[Fourth-series V7 result ledger]
\label{thm:V7-results}
Relative to V1--V6, V7 proves:
\begin{enumerate}[label=\textup{(V7.\arabic*)},leftmargin=5.8em]
\item an exact \(SU(2)\) three-slot family with
      \(\theta_{ac}^{E}=1/3\);
\item the exact \(d_j^{-3}\asymp M^{-3/2}\) protected-sector capacity;
\item failure of every heat-only attempt to obtain a fourth response
      from strong stock on the same three slots;
\item nonimplication of a row half-power from global lower-order MTA
      alone;
\item the nested first-connectivity and temporal-overlap identities;
\item closure of a positive deleted strong share in a complete nested
      suffix;
\item closure of a cofinal selected nested covariance;
\item strict termination of internal block descent; and
\item reduction of all \(q=3\) strong-stock survivors to three finite
      analytic classes.
\end{enumerate}
\end{theorem}

\begin{proof}
These are
\Cref{lem:V7-three-slot-strong,
lem:V7-singlet-capacity,
prop:V7-no-same-slot-fourth,
prop:V7-MTA-nonimplication,
lem:V7-nested-split,
thm:V7-nested-suffix-close,
lem:V7-nested-covariance,
lem:V7-strict-decrease,
thm:V7-finite-descent,
cor:V7-unified-residual}.
\end{proof}

\begin{programme}[V8: crossing-prefix covariance ancestry]
\label{prog:V7-V8}
\begin{enumerate}[leftmargin=3em]
\item Fix the two predecessor blocks containing \(a\) and \(c\) in a
      crossing-prefix packet and keep their complete prefix carriers
      together with the selected quotient cumulant.
\item At quotient order two, use
      \[
      \mathbb M_{A\cup B,r}-\mathbb M_{A|B,r}
      \]
      without separately norming the two prefix histories.  Determine
      whether the strong prefix overlap appears as a common heat
      ancestry or as a difference response.
\item At quotient orders \(3,4,5\), write the full Möbius coefficients
      and group terms by whether the \(a\)- and \(c\)-blocks have
      merged.
\item Seek one projective \(M^{-1/2}\) bound for that grouped
      difference.  Test it on the V7 three-slot singlet family.
\item If the crossing prefix remains norm-one, identify the smallest
      exact packet and move to the first coordinate at which the two
      ancestry blocks share a physical output label.
\end{enumerate}
\end{programme}

\begin{firewall}[Claim boundary after fourth-series V7]
\label{fire:V7-claim-boundary}
V7 proves that strong stock on the same three slots supplies no fourth
heat response and reduces the internal-prefix branch by a finite
noncircular descent.  It does not close the resulting crossing-prefix
or higher selected-cumulant packets, the noncentral owner or
non-source-compatible remainder branches, the \(q=0,1,2\) response
cards, global quintic or all-order MTA, WUFC, CPM, cutoff removal,
reflection positivity, Osterwalder--Schrader reconstruction,
construction of the continuum Yang--Mills measure, or a uniform
positive continuum spectral gap.  The full Yang--Mills mass gap has
not been proved.
\end{firewall}

\paragraph{Parent-version record.}
V7 was formed from
\texttt{Renewal\_Yang\_Mills\_Mass\_Gap\_Research\_Notes\_V6.tex}.
The V6 parent had \(3073\) logical source lines, \(112\,369\) bytes,
and SHA--256
\[
 \texttt{1c5292e474fc1bd1ca80a4fa8c33825a935a2c9160c05fe5c35779af5f9e311d}.
\]
The next compulsory companion audit is V10.

% =============================================================================
% END APPENDED FOURTH-SERIES V7 MATERIAL
% =============================================================================

% The V7 document terminator is intentionally disabled here:
% \end{document}

% =============================================================================
% BEGIN APPENDED FOURTH-SERIES V8 MATERIAL
% =============================================================================

\clearpage
\chapter{V8: Fixed-order selected-cumulant half-powers and one-row
prefix descent}
\label{ch:V8}

\section{Purpose and proof-order change}
\label{sec:V8-purpose}

V7 reduces every deleted-prefix strong branch to crossing ancestry or
a cofinal selected cumulant.  The order-two selected factor was
handled in V3--V4 by a common nested-isotypic resolution.  For
quotient order \(k\ge3\), the Casimirs belonging to different
moment partitions need not commute, so there is no single resolution
which diagonalizes every term of the cumulant.

There is nevertheless a fixed-order square estimate.  One first uses
Hilbert-space Cauchy--Schwarz on the complete finite Möbius sum.  Each
resulting square belongs to one positive moment-partition term.  The
output Casimir commutes with every such term, and the heat in that term
absorbs the locally allocated payer.  The row-\(a\) block of every
partition then supplies the same partial-transpose half-power.

The other change is physical: retain the entire multiplicity space of
one total output isotype.  Refining it by an incompatible coupling-path
projection before the square would reintroduce the V2 entangler
obstruction.  No such refinement is needed for the operator-norm
carrier.

\section{Moment-partition form of a local cumulant}
\label{sec:V8-partition}

Let \(\mathcal Q=\{1,\ldots,k\}\), \(2\le k\le5\), denote the
quotient blocks present at one selected coordinate.  For
\(\pi\in\Pi(\mathcal Q)\), let
\begin{equation}
 \mathbb M_{\pi}
 :=
 \bigotimes_{D\in\pi}\mathbb M_D(s),
\label{eq:V8-partition-moment}
\end{equation}
where \(\mathbb M_D(s)\) is the normalized complete heat moment of
the quotient rows in \(D\), with its full reducible complement and
physical multiplicity space retained.  Then
\[
 0\preccurlyeq\mathbb M_\pi\preccurlyeq I.
\]

\begin{lemma}[Exact Möbius sum]
\label{lem:V8-Mobius-sum}
The complete selected quotient cumulant is
\begin{equation}
 \boxed{\quad
 \mathbb K_k
 =
 \sum_{\pi\in\Pi(\mathcal Q)}
 \mu_\pi\mathbb M_\pi,
 \qquad
 \mu_\pi=(-1)^{|\pi|-1}(|\pi|-1)!.
 \quad}
\label{eq:V8-Mobius-sum}
\end{equation}
Define the fixed arity constant
\begin{equation}
 L_k:=\sum_{\pi\in\Pi(\mathcal Q)}|\mu_\pi|.
\label{eq:V8-Lk}
\end{equation}
For every vector \(\xi\),
\begin{equation}
 \|\mathbb K_k\xi\|^2
 \le
 L_k\sum_{\pi}|\mu_\pi|
 \|\mathbb M_\pi\xi\|^2.
\label{eq:V8-vector-CS}
\end{equation}
\end{lemma}

\begin{proof}
The first identity is Möbius inversion on the set-partition lattice,
with the standard cumulant coefficient
\(\mu(\pi,\widehat1)=(-1)^{|\pi|-1}(|\pi|-1)!\).
For the second, write
\[
 \mathbb K_k\xi
 =
 \sum_\pi
 |\mu_\pi|^{1/2}
 \left(\operatorname{sgn}\mu_\pi\,
       |\mu_\pi|^{1/2}\mathbb M_\pi\xi\right)
\]
and apply Cauchy--Schwarz to the finite direct sum.  Its scalar factor
is \(\sum_\pi|\mu_\pi|=L_k\).
\end{proof}

\begin{corollary}[Positive square envelope]
\label{cor:V8-square-envelope}
One has the operator inequality
\begin{equation}
 \boxed{\quad
 \mathbb K_k^*\mathbb K_k
 \preccurlyeq
 L_k\sum_{\pi}|\mu_\pi|\mathbb M_\pi^2.
 \quad}
\label{eq:V8-square-envelope}
\end{equation}
The cumulant is self-adjoint, so the same formula controls its left
square.
\end{corollary}

\begin{proof}
Equation \eqref{eq:V8-vector-CS} holds for every vector and is exactly
the quadratic-form statement
\eqref{eq:V8-square-envelope}.  Each moment is self-adjoint, hence so
is their real linear combination \(\mathbb K_k\).
\end{proof}

\begin{firewall}[The cumulant is assembled before the inequality]
\label{fire:V8-assembled-first}
Equation \eqref{eq:V8-square-envelope} is applied after the exact
Möbius sum has been formed.  It does not assign independent quotient
tree marks to its partition terms.  The finite positive envelope is
used only to prove a heat half-power for the selected operator.
\end{firewall}

\section{Output Casimir and payer absorption}
\label{sec:V8-payer}

Let \(C_{\rm out}\) be the Casimir of the total diagonal output of all
\(k\) quotient rows.  For \(D\in\pi\), let \(C_D\) be its block
Casimir and put
\[
 H_\pi:=\sum_{D\in\pi}C_D.
\]

\begin{lemma}[Every partition moment commutes with the total output]
\label{lem:V8-total-commutation}
For every \(\pi\),
\begin{equation}
 [C_{\rm out},C_D]=0\quad(D\in\pi),
 \qquad
 [C_{\rm out},\mathbb M_\pi]=0.
\label{eq:V8-total-commutation}
\end{equation}
Moreover, in quadratic-form order,
\begin{equation}
 0\preccurlyeq C_{\rm out}\preccurlyeq kH_\pi.
\label{eq:V8-total-fusion}
\end{equation}
\end{lemma}

\begin{proof}
Write the total infinitesimal action as
\[
 J_{\rm out}^m=\sum_{D\in\pi}J_D^m.
\]
The actions belonging to distinct blocks commute, and the block
Casimir \(C_D=\sum_m(J_D^m)^2\) commutes with every component of
\(J_D\) and with every other block action.  Hence it commutes with
\(C_{\rm out}\), proving
\eqref{eq:V8-total-commutation}.

For every vector \(\xi\),
\[
 \sum_m\left\|\sum_{D\in\pi}J_D^m\xi\right\|^2
 \le
 |\pi|\sum_{D\in\pi}\sum_m\|J_D^m\xi\|^2
 \le
 k\langle\xi,H_\pi\xi\rangle.
\]
This is \eqref{eq:V8-total-fusion}.
\end{proof}

Let \(P_0\) project onto the trivial total output and define the
locally allocated payer
\begin{equation}
 \mathbb R_s
 :=
 \boldsymbol1_{\{C_{\rm out}>0\}}
 {C_{\rm out}\over1-e^{-sC_{\rm out}}},
 \qquad
 \mathbb R_sP_0=0.
\label{eq:V8-general-payer}
\end{equation}

\begin{lemma}[Partitionwise paid square absorption]
\label{lem:V8-partition-absorption}
For every \(\pi\in\Pi(\mathcal Q)\),
\begin{equation}
 \boxed{\quad
 \mathbb R_s^{\,2}\mathbb M_\pi^2
 \preccurlyeq C_{G,s,k}\mathbb M_\pi.
 \quad}
\label{eq:V8-partition-absorption}
\end{equation}
\end{lemma}

\begin{proof}
The block Casimirs \(C_D\), \(D\in\pi\), commute because their row
sets are disjoint.  By
\Cref{lem:V8-total-commutation}, they also commute with
\(C_{\rm out}\).  Resolve them simultaneously.  On a nontrivial
output block let \(c>0\) and \(h\ge0\) be the eigenvalues of
\(C_{\rm out}\) and \(H_\pi\).  Then
\[
 c\le kh,\qquad
 \mathbb M_\pi=e^{-sh}.
\]
The nonzero Casimir gap gives \(c\ge c_*>0\), and hence
\begin{align*}
 \left({c\over1-e^{-sc}}\right)^2e^{-2sh}
 &\le
 {k^2h^2e^{-sh}\over(1-e^{-sc_*})^2}\,e^{-sh}\\
 &\le C_{G,s,k}e^{-sh}.
\end{align*}
On the trivial output the left coefficient is zero.  Spectral
comparison proves the operator inequality.
\end{proof}

\begin{definition}[Paid complete selected cumulant]
\label{def:V8-paid-cumulant}
The exact once-paid local cumulant is
\begin{equation}
 \mathbb K_k^{\rm pay}:=\mathbb R_s\mathbb K_k.
\label{eq:V8-paid-cumulant}
\end{equation}
This is the local term produced by scalar one-resolvent allocation;
the trivial output contributes no local numerator.
\end{definition}

\begin{proposition}[Paid cumulant square envelope]
\label{prop:V8-paid-square-envelope}
One has
\begin{equation}
 \boxed{\quad
 (\mathbb K_k^{\rm pay})^*
 \mathbb K_k^{\rm pay}
 \preccurlyeq
 C_{G,s,k}
 \sum_{\pi\in\Pi(\mathcal Q)}
 |\mu_\pi|\mathbb M_\pi.
 \quad}
\label{eq:V8-paid-square-envelope}
\end{equation}
The same bound holds for the left square.
\end{proposition}

\begin{proof}
By
\Cref{lem:V8-total-commutation}, \(\mathbb R_s\) commutes with every
\(\mathbb M_\pi\), hence with \(\mathbb K_k\).  Multiply
\eqref{eq:V8-square-envelope} by \(\mathbb R_s^2\) and use
\Cref{lem:V8-partition-absorption} term by term.  Self-adjointness and
commutation make the two square orientations identical.
\end{proof}

\section{The fixed-order selected-cumulant half-power}
\label{sec:V8-half}

Let the original maximal row \(a\) belong to the unique quotient block
\(Q_a\in\mathcal Q\), and let \(a_{a,r}\) be its local Casimir mass.

\begin{lemma}[Every partition moment has the same row half-power]
\label{lem:V8-every-moment-half}
For every \(\pi\in\Pi(\mathcal Q)\),
\begin{equation}
 \|\mathbb M_\pi^{\Gamma_{a^c}}\|
 \le C_{G,s,k}a_{a,r}^{-1/2}.
\label{eq:V8-every-moment-half}
\end{equation}
\end{lemma}

\begin{proof}
There is a unique block \(D_a\in\pi\) containing \(Q_a\).  In
\(\mathbb M_{D_a}(s)\), group the irreducible row-\(a\) module against
the complete reducible module formed by every other row in \(D_a\).
The reducible-complement heat card used in V3 gives the required
partial-transpose norm.  Every block moment
\(\mathbb M_D(s)\), \(D\ne D_a\), acts entirely on the complementary
side of the cut; full transpose preserves its operator norm, which is
at most one.  Tensorization proves the claim.
\end{proof}

\begin{theorem}[Unpaid and paid selected-cumulant square cards]
\label{thm:V8-selected-cumulant-half}
For every fixed \(2\le k\le5\),
\begin{align}
 \kappa_{a|a^c}(\mathbb K_k^*\mathbb K_k)
 &\le C_{G,s,k}a_{a,r}^{-1/2},
\label{eq:V8-unpaid-cumulant-half}\\
 \kappa_{a|a^c}
 \bigl((\mathbb K_k^{\rm pay})^*
       \mathbb K_k^{\rm pay}\bigr)
 &\le C_{G,s,k}a_{a,r}^{-1/2}.
\label{eq:V8-paid-cumulant-half}
\end{align}
Both left-square analogues hold.
\end{theorem}

\begin{proof}
For the unpaid card, use
\Cref{cor:V8-square-envelope}, capacity monotonicity and
subadditivity, \(0\preccurlyeq\mathbb M_\pi^2\preccurlyeq
\mathbb M_\pi\), and
\Cref{lem:V8-every-moment-half}.  The coefficient sum depends only on
\(k\).

For the paid card, use
\Cref{prop:V8-paid-square-envelope} and the same moment half-power.
Self-adjointness gives the left-square statements.
\end{proof}

\section{Physical multiplicities are retained, not projected apart}
\label{sec:V8-multiplicity}

\begin{lemma}[Multiplicity-coarse physical pinching]
\label{lem:V8-coarse-pinching}
Decompose the selected local Hilbert space under the total diagonal
group as
\begin{equation}
 \mathcal H
 =
 \bigoplus_U(V_U\otimes\mathcal M_U),
\label{eq:V8-isotypic-decomposition}
\end{equation}
where \(\mathcal M_U\) is the full multiplicity space.  Every
\(\mathbb M_\pi\), \(\mathbb K_k\), and \(\mathbb R_s\) preserves each
summand.  If \(P_U\) is the full isotypic projection, then the square
envelopes of
\Cref{cor:V8-square-envelope,
prop:V8-paid-square-envelope} remain valid after compression by
\(P_U\), with the same constant and no multiplicity count.
\end{lemma}

\begin{proof}
Each block heat moment is equivariant for the total diagonal action;
so are their tensor products, finite linear combination, and every
function of \(C_{\rm out}\).  Hence they commute with \(P_U\).
Compressing a positive operator inequality by \(P_U\) preserves it.
Because of commutation,
\[
 (P_U\mathbb K_kP_U)^*(P_U\mathbb K_kP_U)
 =
 P_U\mathbb K_k^*\mathbb K_kP_U,
\]
and similarly in the paid case.

No basis or rank-one projection is inserted in \(\mathcal M_U\).
The operator norm on \(V_U\otimes\mathcal M_U\) already takes the
maximum over all multiplicity directions, rather than a sum of their
norms.
\end{proof}

\begin{firewall}[No incompatible coupling-path projection]
\label{fire:V8-no-fine-projection}
Different moment partitions can use noncommuting intermediate
Casimirs.  V8 does not project onto their separate coupling paths.
Such a projection need not commute with
\(\mathbb K_k\), and the capacity of \(PXP\) is not controlled by the
capacity of \(X\) for an arbitrary entangling \(P\).  The valid
physical block is the full total-output isotype with its complete
multiplicity space retained.
\end{firewall}

\begin{corollary}[Cofinal selected cumulants close the \(q=3\) core]
\label{cor:V8-cofinal-selected-close}
Fix a \(q=3\) obstruction.  If its selected coordinate \(r\) obeys
\[
 a_{a,r}\ge\tau M,
\]
then the word satisfies \(XQTRSC_5\) for every quotient order
\(2\le k\le5\) and every payer location, provided physical pinching is
performed at the multiplicity-coarse level of
\Cref{lem:V8-coarse-pinching}.
\end{corollary}

\begin{proof}
Use the unpaid or paid estimate in
\Cref{thm:V8-selected-cumulant-half} according as the payer is outside
or at \(r\).  It supplies \(CM^{-1/2}\); the three distinct extracted
slots supply \(CM^{-3/2}\).  All slot coordinates are disjoint from
\(r\).  Tensorization, coarse physical pinching, and the direct
same-coordinate scalar comparison of
\Cref{prop:V6-same-coordinate} give
\[
 {CM^{-2}\over(\prod_i\Xi_i)^2}
 \le C\mathfrak x_{T_0,r}^{\,2}.
\]
\end{proof}

\section{One-row descent in a crossing-prefix packet}
\label{sec:V8-one-row-descent}

Consider a deleted-prefix alternative
\begin{equation}
 \overline\theta_{ac}^{-}(r)\ge\beta>0
\label{eq:V8-prefix-overlap}
\end{equation}
in which \(a,c\) lie in distinct global prefix blocks.  Let
\(A\in\rho\) be the block containing \(a\).  By
\Cref{lem:V6-overlap-forces-mass},
\begin{equation}
 \sum_{\substack{e<r\\e\notin E_3}}a_{a,e}
 \ge\beta M.
\label{eq:V8-prefix-a-mass}
\end{equation}

\begin{lemma}[Singleton-block terminal]
\label{lem:V8-singleton-terminal}
If \(A=\{a\}\), the crossing-prefix packet closes.
\end{lemma}

\begin{proof}
The exact global first-connectivity source gives row \(a\) a complete
singleton moment prefix on all \(e<r\).  Delete \(E_3\) and, when
present, the payer coordinate.  If the payer itself carries a fixed
fraction of the mass in
\eqref{eq:V8-prefix-a-mass}, it is a paid central heat owner and
closes by
\Cref{thm:V5-central-owner-half}.  Otherwise the remaining complete
prefix bank retains a fixed mass fraction.  A cofinal atom is an unpaid
central heat owner; if no such atom persists, the bank is diffuse and
closes by the archived source-compatible diffuse terminal.  Every
owner used here is outside the already extracted set \(E_3\).
\end{proof}

Suppose now \(|A|\ge2\).  In a fixed word of
\(\mathcal K_A^{<r,\mathrm{conn}}\), let \(s<r\) be its first
connectivity coordinate and let \(\sigma<\widehat1_A\) be the
predecessor partition.  Decompose the deleted row mass as
\begin{align}
 m_{<s}
 &:=
 {1\over M}
 \sum_{\substack{e<s\\e\notin E_3}}a_{a,e},
\notag\\
 m_s
 &:=
 \begin{cases}
 a_{a,s}/M,&s\notin E_3,\\
 0,&s\in E_3,
 \end{cases}
\notag\\
 m_{(s,r)}
 &:=
 {1\over M}
 \sum_{\substack{s<e<r\\e\notin E_3}}a_{a,e}.
\label{eq:V8-row-mass-split}
\end{align}

\begin{lemma}[Exact one-row ancestry split]
\label{lem:V8-row-mass-split}
At the current descent node,
\begin{equation}
 m_{<s}+m_s+m_{(s,r)}
 =
 {1\over M}
 \sum_{\substack{e<r\\e\notin E_3}}a_{a,e}.
\label{eq:V8-row-mass-identity}
\end{equation}
Hence a current mass share \(\gamma\) yields a subshare at least
\(\gamma/3\).
\end{lemma}

\begin{proof}
This is the disjoint partition of the deleted coordinate interval into
\(e<s\), \(e=s\), and \(s<e<r\).
\end{proof}

\begin{lemma}[Later and selected one-row alternatives close]
\label{lem:V8-later-selected-close}
If \(m_{(s,r)}\ge\gamma\), the complete \(A\)-moment suffix closes by
the atom/diffuse argument of
\Cref{thm:V7-nested-suffix-close}.  If
\(m_s\ge\gamma\), the complete selected quotient cumulant at \(s\)
closes by
\Cref{cor:V8-cofinal-selected-close}, with \(k=|\sigma|\le|A|\).
\end{lemma}

\begin{proof}
For the later alternative, the exact nested first-connectivity word
contains a complete \(A\)-moment suffix on \((s,r)\).  The proof of
\Cref{thm:V7-nested-suffix-close} uses only positive row mass, removal
of \(E_3\) and the payer, and complete source occurrence, so it applies
without using row \(c\).

For the selected alternative,
\(a_{a,s}\ge\gamma M\).  The selected local factor is the complete
\(|\sigma|\)-th quotient cumulant, paid locally or unpaid according to
the exact payer allocation.  Apply
\Cref{thm:V8-selected-cumulant-half} and then the three-slot comparison
from
\Cref{cor:V8-cofinal-selected-close}.
\end{proof}

\begin{definition}[One-row strict descent]
\label{def:V8-one-row-step}
If \(m_{<s}\ge\gamma\), let \(A'\in\sigma\) be the unique predecessor
block containing \(a\).  If \(A'=\{a\}\), apply
\Cref{lem:V8-singleton-terminal} on the interval \(e<s\).  Otherwise
replace the current node \((A,r)\) by \((A',s)\).
\end{definition}

\begin{lemma}[One-row descent terminates]
\label{lem:V8-one-row-terminates}
Every nonterminal step has
\[
 2\le|A'|\le|A|-1.
\]
Starting with a proper quintic prefix block, there are at most three
nonterminal steps.
\end{lemma}

\begin{proof}
The predecessor partition \(\sigma\) is proper, so its block \(A'\)
is a strict subset of \(A\).  A nonterminal \(A'\) contains \(a\) and
at least one other row.  Initially \(|A|\le4\), proving the stated
bound.
\end{proof}

\begin{theorem}[Every deleted crossing-prefix strong branch closes]
\label{thm:V8-crossing-prefix-close}
Every persistent \(q=3\) crossing-prefix packet satisfying
\eqref{eq:V8-prefix-overlap} obeys \(XQTRSC_5\), for all payer
locations and both square orientations.
\end{theorem}

\begin{proof}
Equation \eqref{eq:V8-prefix-a-mass} supplies the initial positive
row-\(a\) share.  If \(A\) is a singleton, use
\Cref{lem:V8-singleton-terminal}.  Otherwise apply
\Cref{lem:V8-row-mass-split}.  A later or selected share closes by
\Cref{lem:V8-later-selected-close}; an earlier share invokes
\Cref{def:V8-one-row-step}.

The block size strictly decreases by
\Cref{lem:V8-one-row-terminates}; after at most three steps the
singleton terminal is reached unless an earlier later/selected
alternative has already closed.  Every pigeonhole loss is a fixed
power of three and therefore preserves a positive cofinal fraction.
No quintic conclusion is used recursively.
\end{proof}

\begin{corollary}[All deleted-prefix \(q=3\) strong branches close]
\label{cor:V8-all-prefix-close}
The deleted-prefix alternatives in both the V1 \(14\)-internal and
\(37\)-crossing partition classes satisfy \(XQTRSC_5\).
\end{corollary}

\begin{proof}
The direct crossing class is
\Cref{thm:V8-crossing-prefix-close}.  By
\Cref{thm:V7-finite-descent}, an internal class either already closes
or reduces to a crossing-prefix packet or a cofinal selected cumulant.
Use
\Cref{thm:V8-crossing-prefix-close} or
\Cref{cor:V8-cofinal-selected-close} in the two remaining cases.
\end{proof}

\section{Residual after the selected-cumulant theorem}
\label{sec:V8-residual}

\begin{theorem}[Strong-stock \(q=3\) reduction after V8]
\label{thm:V8-q3-reduction}
After V1--V8, every \(q=3\) strong-stock survivor must have
\begin{equation}
 \boxed{\quad
 \theta_{ac}^{E}\ge{\eta\over4},
 \quad}
\label{eq:V8-only-E3}
\end{equation}
so the fixed strong overlap is carried by the same three extracted
slots.  The heat-only fourth response is impossible by
\Cref{prop:V7-no-same-slot-fourth}.  Any remaining proof must use the
assembled quotient/prefix operator despite the selected row being
locally depleted, or a separate noncentral/remainder owner from V5.
\end{theorem}

\begin{proof}
The extracted-slot-deleted V6 trichotomy gives \(E_3\), deleted prefix,
deleted suffix, or selected cofinal stock.  The deleted suffix closes
by V6, every selected quotient order closes by
\Cref{cor:V8-cofinal-selected-close}, and every deleted prefix closes
by
\Cref{cor:V8-all-prefix-close}.  Thus only
\eqref{eq:V8-only-E3} remains.  The sharpness statement is V7.
\end{proof}

\begin{firewall}[Multiplicity scope of the V8 closure]
\label{fire:V8-closure-scope}
The selected-cumulant closure retains each full total-output
multiplicity space as required by
\Cref{lem:V8-coarse-pinching}.  It does not prove the same estimate
after arbitrary rank-one coupling-path projections.  If a later
source formula genuinely requires such a finer projection rather than
the operator norm on the full multiplicity block, that projection is a
separate noncentral-owner obstruction and must be returned to V5's
residual list.
\end{firewall}

\section{V8 result ledger and V9 acquisition order}
\label{sec:V8-ledger}

\begin{theorem}[Fourth-series V8 result ledger]
\label{thm:V8-results}
Relative to V1--V7, V8 proves:
\begin{enumerate}[label=\textup{(V8.\arabic*)},leftmargin=5.8em]
\item the exact finite Möbius form and positive square envelope for
      every selected quotient cumulant of order at most five;
\item commutation and fusion comparison of every partition moment with
      the total output Casimir;
\item paid square absorption for every partition term;
\item unpaid and once-paid selected-cumulant half-powers;
\item multiplicity-coarse physical pinching without an output count;
\item closure of every cofinal selected cumulant at orders \(2\)--\(5\);
\item a terminating one-row descent for crossing-prefix ancestry;
\item closure of all deleted-prefix \(q=3\) strong branches; and
\item reduction of the strong-stock \(q=3\) residual to overlap owned
      by the same three extracted slots.
\end{enumerate}
\end{theorem}

\begin{proof}
These are
\Cref{lem:V8-Mobius-sum,
cor:V8-square-envelope,
lem:V8-total-commutation,
lem:V8-partition-absorption,
prop:V8-paid-square-envelope,
thm:V8-selected-cumulant-half,
lem:V8-coarse-pinching,
cor:V8-cofinal-selected-close,
thm:V8-crossing-prefix-close,
cor:V8-all-prefix-close,
thm:V8-q3-reduction}.
\end{proof}

\begin{programme}[V9: depleted selected response against an
\(E_3\)-owned strong pair]
\label{prog:V8-V9}
\begin{enumerate}[leftmargin=3em]
\item Fix the V8 residual
      \(\theta_{ac}^{E}\ge\eta/4\) and retain the three-slot product as
      one sharp protected bank.
\item Write the complete selected quotient cumulant on the complementary
      locally depleted row-\(a\) sector.  A mere heat half-power there is
      unavailable.
\item Group the Möbius sum by whether the quotient blocks containing
      \(a\) and \(c\) have merged.  Seek a two-row covariance response
      proportional to the already present cross-slot overlap without
      treating that overlap as a fourth heat coordinate.
\item Test the grouped operator on the exact V7 triple-singlet sector.
      If it remains norm one, record a physical counterpacket to the
      proposed response and move upstream to the joint three-slot plus
      selected-coordinate carrier.
\item Reconcile the result with the V5 noncentral-owner and
      non-source-compatible remainder alternatives.
\end{enumerate}
\end{programme}

\begin{firewall}[Claim boundary after fourth-series V8]
\label{fire:V8-claim-boundary}
V8 closes all cofinal selected-cumulant and deleted-prefix
\(q=3\) branches at fixed order five, subject to multiplicity-coarse
physical pinching.  It does not close the sharp \(E_3\)-owned strong
branch, arbitrary fine multiplicity projections, the V5 noncentral or
remainder alternatives, the \(q=0,1,2\) response cards, global quintic
or all-order MTA, WUFC, CPM, cutoff removal, reflection positivity,
Osterwalder--Schrader reconstruction, construction of the continuum
Yang--Mills measure, or a uniform positive continuum spectral gap.
The full Yang--Mills mass gap has not been proved.
\end{firewall}

\paragraph{Parent-version record.}
V8 was formed from
\texttt{Renewal\_Yang\_Mills\_Mass\_Gap\_Research\_Notes\_V7.tex}.
The V7 parent had \(3619\) logical source lines, \(131\,421\) bytes,
and SHA--256
\[
 \texttt{40c82dd4f390b34a11be1274762148f29578edb05d65fa4cb57384e4ab3099dd}.
\]
The next compulsory companion audit is V10.

% =============================================================================
% END APPENDED FOURTH-SERIES V8 MATERIAL
% =============================================================================

% The V8 document terminator is intentionally disabled here:
% \end{document}

% =============================================================================
% BEGIN APPENDED FOURTH-SERIES V9 MATERIAL
% =============================================================================

\clearpage
\chapter{V9: The depleted-selected no-go and the two-sided
heat-sandwich gate}
\label{ch:V9}

\section{Purpose}
\label{sec:V9-purpose}

V8 leaves a \(q=3\) branch in which the fixed strong overlap is carried
by the same three already counted heat slots.  The selected coordinate
is locally depleted, so its fixed-dimensional local cumulant cannot
produce a power of the large row mass by itself.  We first prove this
as an exact tensor countertest.

The core still has a physical coordinate outside the three slots which
owns the fixed strong role.  Its local operator can be noncentral
because a marked separator or recoupling lies at the same coordinate.
An ordinary norm bound is insufficient, but full commutation is
stronger than necessary.  The correct intermediate condition is a
two-sided heat sandwich
\[
 H^{1/2}SH^{1/2},
\]
with \(S\) a bounded equivariant operator.  Such a packet preserves the
row half-power in both square orientations and also absorbs the local
payer.  A rank-one entangling example proves that a one-sided heat
factor does not suffice.

\section{A fixed depleted selected factor cannot repair three slots}
\label{sec:V9-depleted-no-go}

Retain the \(SU(2)\) three-slot family of
\Cref{sec:V7-three-slot-test}.  Thus
\[
 M_j=3(1+j(j+1)),\qquad
 \theta_{ac}^{E}={1\over3},
\]
and the protected three-slot projection is
\[
 \mathbb P_j:=P_{0,j}^{\otimes3}.
\]

\begin{lemma}[A nonzero positive operator is detected by a product
vector]
\label{lem:V9-product-detection}
Let \(X\ge0\) be a nonzero operator on a finite tensor product
\(\mathcal H_1\otimes\cdots\otimes\mathcal H_m\).  Then there are unit
vectors \(v_i\in\mathcal H_i\) such that
\begin{equation}
 \left\langle\bigotimes_i v_i,\,
 X\bigotimes_i v_i\right\rangle>0.
\label{eq:V9-product-detection}
\end{equation}
\end{lemma}

\begin{proof}
Suppose the displayed quadratic form vanished on every product vector.
Fix product vectors in all but one factor and polarize in the remaining
factor; then polarize successively in every factor.  All matrix
elements of \(X\) between elementary tensors vanish.  Elementary
tensors span the full finite tensor product, so \(X=0\), a
contradiction.
\end{proof}

Let \(\mathbb K_0\) be any fixed, nonzero selected local cumulant
acting on representations independent of \(j\), at a coordinate
disjoint from \(E_3\).  All other fixed prefix and suffix factors may
be absorbed into \(\mathbb K_0\) provided the resulting operator
remains nonzero.  Define
\begin{equation}
 \mathbb Z_j
 :=
 \left(\bigotimes_{\ell=1}^{3}\mathbb H_{e_\ell}(s)\right)
 \otimes\mathbb K_0.
\label{eq:V9-separated-packet}
\end{equation}

\begin{proposition}[Depleted selected-factor tensor no-go]
\label{prop:V9-depleted-selected-no-go}
There is a constant \(c_0>0\), independent of \(j\), such that
\begin{equation}
 \kappa_\otimes(\mathbb Z_j^*\mathbb Z_j)
 \ge {c_0\over d_j^3}
 \asymp M_j^{-3/2}.
\label{eq:V9-depleted-selected-lower}
\end{equation}
Consequently no estimate of the form
\[
 \kappa_\otimes(\mathbb Z_j^*\mathbb Z_j)
 \le C M_j^{-2}
\]
can follow from three-slot heat, fixed strong overlap on those slots,
and a disjoint selected factor whose representations remain bounded.
\end{proposition}

\begin{proof}
By
\Cref{lem:V9-product-detection}, choose a fixed product vector \(v\)
with
\[
 \langle v,\mathbb K_0^*\mathbb K_0v\rangle=c_0>0.
\]
At the three heat coordinates use the product vectors attaining the
singlet capacity in
\Cref{lem:V7-singlet-capacity}.  Their tensor product with \(v\) is
an admissible row-product test for the complete packet.  Since each
heat moment is the identity on its singlet projection, the expectation
is \(c_0d_j^{-3}\).  The asymptotic comparison and failure of
\(CM_j^{-2}\) are exactly those in
\Cref{prop:V7-no-same-slot-fourth}.
\end{proof}

\begin{remark}[A concrete fixed selected covariance]
\label{rem:V9-concrete-K0}
One may take \(\mathbb K_0\) to be a two-block \(SU(2)\) covariance on
fixed nontrivial spins:
\[
 e^{-sC_{A\cup B}}-e^{-s(C_A+C_B)}.
\]
Its two scalar multipliers differ on at least one fusion channel, so it
is nonzero.  The proposition therefore does not rely on an abstract
choice of operator.
\end{remark}

\begin{firewall}[What the no-go proves]
\label{fire:V9-no-go-scope}
\Cref{prop:V9-depleted-selected-no-go} disproves only a separated
supplement estimate.  It does not by itself disprove the complete
quintic MTA: the exact word can contain an additional cofinal
fixed-strong role owner, noncentral separator, or source response not
present in \(\mathbb Z_j\).  Those operations must be retained before
a conclusion about \(XQTRSC_5\) is drawn.
\end{firewall}

\section{Abstract two-sided heat sandwiches}
\label{sec:V9-sandwich}

\begin{definition}[Finite equivariant heat-sandwich packet]
\label{def:V9-heat-sandwich}
On one physical coordinate \(d\), a finite equivariant heat-sandwich
packet across \(a|a^c\) is
\begin{equation}
 \mathbb L_d
 =
 \sum_{\nu=1}^{N}
 \mathbb H_{\nu}^{1/2}
 \mathbb S_\nu
 \mathbb H_{\nu}^{1/2},
\label{eq:V9-heat-sandwich}
\end{equation}
where:
\begin{enumerate}[label=\textup{(\roman*)},leftmargin=4em]
\item \(N\le N_5\) is a fixed arity bound;
\item \(0\preccurlyeq\mathbb H_\nu\preccurlyeq I\) is a complete
      local heat moment and
\begin{equation}
 \|\mathbb H_\nu^{\Gamma_{a^c}}\|
 \le C_{G,s}a_{a,d}^{-1/2};
\label{eq:V9-sandwich-heat-half}
\end{equation}
\item \(\|\mathbb S_\nu\|\le C_G\);
\item \(\mathbb S_\nu\) is equivariant for the total diagonal output,
      so it commutes with the total output Casimir \(C_{\rm out}\);
      and
\item if \(h_\nu\) is the heat exponent on a common block and \(c\) the
      nontrivial total-output Casimir, then
\begin{equation}
 c\le C_G h_\nu.
\label{eq:V9-sandwich-fusion}
\end{equation}
\end{enumerate}
The complete multiplicity space is retained inside \(\mathbb S_\nu\).
\end{definition}

\begin{lemma}[Each sandwich has two positive square envelopes]
\label{lem:V9-two-square-envelopes}
Put
\[
 T_\nu
 =
 \mathbb H_\nu^{1/2}\mathbb S_\nu
 \mathbb H_\nu^{1/2}.
\]
Then
\begin{equation}
 T_\nu^*T_\nu
 \preccurlyeq C_G^2\mathbb H_\nu,
 \qquad
 T_\nu T_\nu^*
 \preccurlyeq C_G^2\mathbb H_\nu.
\label{eq:V9-two-square-envelopes}
\end{equation}
\end{lemma}

\begin{proof}
Since \(0\preccurlyeq\mathbb H_\nu\preccurlyeq I\),
\[
 \mathbb S_\nu^*\mathbb H_\nu\mathbb S_\nu
 \preccurlyeq
 \|\mathbb S_\nu\|^2I.
\]
Multiplying on both sides by \(\mathbb H_\nu^{1/2}\) gives the first
inequality.  Apply the same argument to
\(\mathbb S_\nu\mathbb H_\nu\mathbb S_\nu^*\) for the second.
\end{proof}

\begin{theorem}[Unpaid two-sided sandwich half-power]
\label{thm:V9-unpaid-sandwich-half}
Every packet in
\Cref{def:V9-heat-sandwich} satisfies
\begin{align}
 \kappa_{a|a^c}(\mathbb L_d^*\mathbb L_d)
 &\le C_{G,s,N_5}a_{a,d}^{-1/2},
\label{eq:V9-unpaid-sandwich-right}\\
 \kappa_{a|a^c}(\mathbb L_d\mathbb L_d^*)
 &\le C_{G,s,N_5}a_{a,d}^{-1/2}.
\label{eq:V9-unpaid-sandwich-left}
\end{align}
\end{theorem}

\begin{proof}
Hilbert-space Cauchy--Schwarz for the finite sum gives
\[
 \mathbb L_d^*\mathbb L_d
 \preccurlyeq
 N\sum_{\nu=1}^N T_\nu^*T_\nu,
\qquad
 \mathbb L_d\mathbb L_d^*
 \preccurlyeq
 N\sum_{\nu=1}^N T_\nu T_\nu^*.
\]
Use
\Cref{lem:V9-two-square-envelopes}, capacity monotonicity and
subadditivity, and
\eqref{eq:V9-sandwich-heat-half}.  Since \(N\le N_5\), all finite
factors enter the constant.
\end{proof}

\section{The locally paid sandwich}
\label{sec:V9-paid-sandwich}

Define
\begin{equation}
 \mathbb R_s
 =
 \boldsymbol1_{\{C_{\rm out}>0\}}
 {C_{\rm out}\over1-e^{-sC_{\rm out}}},
 \qquad
 \mathbb L_d^{\rm pay}:=\mathbb R_s\mathbb L_d.
\label{eq:V9-paid-sandwich}
\end{equation}

\begin{lemma}[Payer times one heat is absorbed at half time]
\label{lem:V9-half-time-absorption}
Under
\eqref{eq:V9-sandwich-fusion},
\begin{equation}
 \mathbb R_s^{\,2}\mathbb H_\nu(s)
 \preccurlyeq
 C_{G,s}\mathbb H_\nu(s/2).
\label{eq:V9-half-time-absorption}
\end{equation}
\end{lemma}

\begin{proof}
Resolve the commuting total output Casimir and the heat exponent.
On \(c>0\), the nonzero Casimir gap and \(c\le C_Gh\) give
\begin{align*}
 \left({c\over1-e^{-sc}}\right)^2e^{-sh}
 &\le
 {C_G^2h^2e^{-sh/2}\over(1-e^{-sc_*})^2}
 e^{-sh/2}\\
 &\le C_{G,s}e^{-(s/2)h}.
\end{align*}
The left side vanishes on the protected total output.
\end{proof}

\begin{theorem}[Once-paid two-sided sandwich half-power]
\label{thm:V9-paid-sandwich-half}
Every packet in
\Cref{def:V9-heat-sandwich} satisfies
\begin{align}
 \kappa_{a|a^c}
 \bigl((\mathbb L_d^{\rm pay})^*
       \mathbb L_d^{\rm pay}\bigr)
 &\le C_{G,s,N_5}a_{a,d}^{-1/2},
\label{eq:V9-paid-sandwich-right}\\
 \kappa_{a|a^c}
 \bigl(\mathbb L_d^{\rm pay}
       (\mathbb L_d^{\rm pay})^*\bigr)
 &\le C_{G,s,N_5}a_{a,d}^{-1/2}.
\label{eq:V9-paid-sandwich-left}
\end{align}
\end{theorem}

\begin{proof}
Equivariance gives
\[
 [\mathbb R_s,\mathbb S_\nu]=0.
\]
The total output Casimir also commutes with every complete heat moment,
so \(\mathbb R_s\) commutes with \(T_\nu\) and with their finite sum.
Multiply the two Cauchy--Schwarz square envelopes in the proof of
\Cref{thm:V9-unpaid-sandwich-half} by \(\mathbb R_s^2\).  By
\Cref{lem:V9-two-square-envelopes,
lem:V9-half-time-absorption}, both are bounded by
\[
 C\sum_\nu\mathbb H_\nu(s/2).
\]
The reducible-complement heat card at the fixed positive time \(s/2\)
has the same \(a_{a,d}^{-1/2}\) power.  Capacity monotonicity and
subadditivity prove both assertions.
\end{proof}

\begin{remark}[Full centrality is unnecessary]
\label{rem:V9-not-central}
The operators \(\mathbb S_\nu\) need not commute with
\(\mathbb H_\nu\).  The heat appears on both sides of
\(\mathbb S_\nu\), and that is exactly what gives both inequalities in
\eqref{eq:V9-two-square-envelopes}.  Only total-output equivariance is
used to pass the payer through the complete packet.
\end{remark}

\section{Sharp failure of a one-sided heat owner}
\label{sec:V9-one-sided}

\begin{proposition}[One-sided entangling heat loses one orientation]
\label{prop:V9-one-sided-fails}
Let
\[
 H=P_{\Omega}
\]
be the normalized maximally entangled rank-one projection on
\(\mathbb C^d\otimes\mathbb C^d\), and let \(U\) be a unitary sending
\(\Omega\) to a product vector \(x\otimes y\).  For
\[
 L=UH
\]
one has
\begin{align}
 L^*L&=P_\Omega,
 &
 \kappa_\otimes(L^*L)&={1\over d},
\label{eq:V9-one-sided-good}\\
 LL^*&=|x\otimes y\rangle\langle x\otimes y|,
 &
 \kappa_\otimes(LL^*)&=1.
\label{eq:V9-one-sided-bad}
\end{align}
Thus a one-sided heat factor and a norm-one separator do not imply
two oriented half-power cards.
\end{proposition}

\begin{proof}
The two operator identities follow from \(H^2=H\) and
\(U\Omega=x\otimes y\).  The first capacity is
\Cref{lem:V7-singlet-capacity}; the second projection is itself a
product state and therefore has capacity one.
\end{proof}

\begin{firewall}[The sandwich condition is sharp at the algebraic
level]
\label{fire:V9-sandwich-sharp}
Replacing
\(H^{1/2}SH^{1/2}\) by \(SH\) proves at most the right-square card;
replacing it by \(HS\) proves at most the left-square card.  Ordinary
norm control of \(S\) cannot repair the missing orientation.  Any
actual Yang--Mills owner which is only one-sided must be combined with
an adjacent heat factor before a two-orientation claim is made.
\end{firewall}

\section{Application to the fixed strong-role owner}
\label{sec:V9-owner-application}

\begin{definition}[Sandwich-certified strong owner]
\label{def:V9-sandwich-owner}
In a \(q=3\) core, let \(d_*\notin E_3\) be a cofinal physical
coordinate assigned to the fixed strong role by the owner partition:
\begin{equation}
 a_{a,d_*}\ge\sigma M.
\label{eq:V9-fixed-owner-cofinal}
\end{equation}
It is sandwich-certified if its complete local operator, with every
joining, payer, tree-mark, separator, and physical compression role at
\(d_*\) retained in original order, has the form
\eqref{eq:V9-heat-sandwich}; if it pays, the debit is
\eqref{eq:V9-paid-sandwich}.
\end{definition}

\begin{theorem}[A sandwich-certified fixed strong owner closes the
\(E_3\)-owned branch]
\label{thm:V9-sandwich-owner-close}
Assume
\[
 \theta_{ac}^{E}\ge\eta/4
\]
and that the fixed strong-role coordinate \(d_*\notin E_3\) satisfies
\eqref{eq:V9-fixed-owner-cofinal}.  If \(d_*\) is
sandwich-certified, the \(q=3\) word satisfies both orientations of
\(XQTRSC_5\), for every payer location.
\end{theorem}

\begin{proof}
Use
\Cref{thm:V9-unpaid-sandwich-half} or
\Cref{thm:V9-paid-sandwich-half} at \(d_*\), obtaining
\(C_{\sigma}M^{-1/2}\).  The three distinct \(E_3\) heat slots give
\(CM^{-3/2}\).  The coordinates are disjoint, and the complete
sandwich already contains every noncentral operation at \(d_*\), so
tensorization gives \(CM^{-2}\) in both square orientations.

Insert this complete four-factor packet in the normalized word.  The
scalar exchange-atlas floor uses the fixed strong pair but no second
operator estimate at \(d_*\).  The comparison is exactly
\eqref{eq:V6-same-coordinate-upper}, proving the square card.
\end{proof}

\begin{corollary}[Exact residual owner after V9]
\label{cor:V9-residual-owner}
Within the standing \(q=3\) construction, every survivor of the V8
\(E_3\)-owned branch must have a cofinal fixed strong-role coordinate
\(d_*\notin E_3\) whose complete local packet fails the finite
equivariant two-sided heat-sandwich condition in at least one square
orientation.  Equivalently, it contains an unsandwiched one-sided
separator or a physical operation not commuting with total output.
\end{corollary}

\begin{proof}
The cofinal coordinate \(d_*\) and its disjointness from \(E_3\) are
part of the fixed-strong-coordinate deletion performed before the
number \(q\) is defined.  If its complete packet were
sandwich-certified, the word would close by
\Cref{thm:V9-sandwich-owner-close}.  The negation of
\Cref{def:V9-heat-sandwich} is precisely the stated factorization or
equivariance failure.
\end{proof}

\begin{assessment}[The bottleneck is now an exact local factorization]
\label{ass:V9-factorization-bottleneck}
The remaining \(q=3\) issue is no longer a scalar mass or graph
question.  The exact question is whether every fixed-strong local
packet produced by the quintic source can be written as a bounded
finite sum of equivariant two-sided heat sandwiches after coincident
roles are consolidated.  If yes, V9 closes the sharp three-slot
branch.  If no, the smallest failing packet contains a one-sided
entangler of the type detected by
\Cref{prop:V9-one-sided-fails} and must be joined to an adjacent heat
coordinate before taking either square.
\end{assessment}

\section{V9 result ledger and V10 acquisition order}
\label{sec:V9-ledger}

\begin{theorem}[Fourth-series V9 result ledger]
\label{thm:V9-results}
Relative to V1--V8, V9 proves:
\begin{enumerate}[label=\textup{(V9.\arabic*)},leftmargin=5.8em]
\item a product-vector detection lemma for nonzero positive local
      factors;
\item an exact \(M^{-3/2}\) lower bound for three singlet slots tensored
      with a fixed depleted selected cumulant;
\item failure of the separated depleted-selected repair;
\item two positive square envelopes for a bounded two-sided heat
      sandwich;
\item unpaid and once-paid half-powers for finite equivariant
      heat-sandwich packets;
\item a sharp one-sided entangler counterexample;
\item closure of the \(E_3\)-owned branch whenever its cofinal fixed
      strong owner is sandwich-certified; and
\item reduction of the surviving \(q=3\) problem to one exact local
      heat-sandwich factorization.
\end{enumerate}
\end{theorem}

\begin{proof}
These are
\Cref{lem:V9-product-detection,
prop:V9-depleted-selected-no-go,
lem:V9-two-square-envelopes,
thm:V9-unpaid-sandwich-half,
lem:V9-half-time-absorption,
thm:V9-paid-sandwich-half,
prop:V9-one-sided-fails,
thm:V9-sandwich-owner-close,
cor:V9-residual-owner}.
\end{proof}

\begin{programme}[V10: compulsory audit and marked-owner
factorization]
\label{prog:V9-V10}
\begin{enumerate}[leftmargin=3em]
\item Perform the compulsory read-only audit of the newest active SM
      and NS notes and record filenames, sizes, and SHA--256 digests.
\item Return to the exact finite list of local role operations:
      normalized tree generators, joining cumulant, payer, recoupler,
      separator, and physical compression.
\item Use centrality of the quadratic Casimir and the heat semigroup
      law to place half of an actual heat factor on each side of every
      Casimir-normalized generator which commutes with its block
      action.
\item For a recoupler changing block partitions, decide whether an
      adjacent complete heat factor supplies the missing second side.
\item If a packet is only one-sided, retain it together with the
      nearest source-occurring heat coordinate and test the resulting
      two-coordinate sandwich.
\item After the \(q=3\) owner is settled, return to the exact
      \(q=2,1,0\) exponent table.
\end{enumerate}
\end{programme}

\begin{firewall}[Claim boundary after fourth-series V9]
\label{fire:V9-claim-boundary}
V9 does not prove that every actual marked fixed-strong packet is a
two-sided heat sandwich.  It closes the sharp \(q=3\) branch only under
that exact factorization certificate and proves that a separated
depleted selected factor cannot substitute for it.  The unsandwiched
owner, fine-multiplicity exceptions, \(q=0,1,2\) cards, global quintic
or all-order MTA, WUFC, CPM, cutoff removal, reflection positivity,
Osterwalder--Schrader reconstruction, construction of the continuum
Yang--Mills measure, and a uniform positive continuum spectral gap
remain open.  The full Yang--Mills mass gap has not been proved.
\end{firewall}

\paragraph{Parent-version record.}
V9 was formed from
\texttt{Renewal\_Yang\_Mills\_Mass\_Gap\_Research\_Notes\_V8.tex}.
The V8 parent had \(4311\) logical source lines, \(154\,779\) bytes,
and SHA--256
\[
 \texttt{2d4f77d9b36eba26f3fbed78e741ea2f424a4c4d7ec80c9177e3b54d53da6fae}.
\]
The next compulsory companion audit is V10.

% =============================================================================
% END APPENDED FOURTH-SERIES V9 MATERIAL
% =============================================================================

% =============================================================================
% BEGIN APPENDED FOURTH-SERIES V10 MATERIAL
% =============================================================================

\clearpage
\chapter{V10: Weighted projection sandwiches and the exact marked-owner
gate}
\label{ch:V10}

\section{Purpose and correction of the operator inventory}
\label{sec:V10-purpose}

V9 reduced the sharp \(q=3\) branch to a two-sided heat-sandwich
certificate for the complete owner outside the three already counted
heat slots.  Its programme described tree marks and the fixed strong
role as though each role name necessarily supplied an additional
operator.  That wording is too coarse.

In the representation-uniform tree card, the tree \(\tau\) indexes an
already formed summand \(\mathbb K_{5,\tau}\), and its valence monomial
appears on the scalar right-hand side of the norm estimate.  Likewise,
the fixed strong pair selects a scalar overlap stock in the exchange
atlas.  The labels \(\tau\) and strong do not, merely by being assigned
to a coordinate, compose the local word with a new Hilbert-space map.
The differentiated heat coefficients inside \(\mathbb K_{5,\tau}\) are
genuine operators, but V8 controls their assembled cumulant rather than
estimating the labels separately.

Consequently the remaining V9 obstruction must be sought among actual
maps in the physical word: a one-sided separator, output compression,
partial trace, or change of coupling resolution.  This chapter gives an
exact necessary-and-sufficient test for a one-sided projection to be a
uniform heat sandwich.  It closes every projection compatible with the
heat spectral resolution and isolates a weighted commutator for every
incompatible one.

\section{Compulsory V10 companion audit}
\label{sec:V10-audit}

The newest active companion files at the time of this audit were
\begin{center}
\begin{tabular}{@{}p{0.48\linewidth}rr@{}}
\toprule
file & logical lines & bytes\\
\midrule
\texttt{Renewal\_SM\_Einstein\_Continuum\_Research\_Notes\_V87.tex}
 & \(33\,433\) & \(1\,230\,732\)\\
\texttt{Renewal\_NS\_Stability\_Research\_Notes\_V31.tex}
 & \(23\,052\) & \(887\,537\)\\
\bottomrule
\end{tabular}
\end{center}
Their SHA--256 digests are, respectively,
\[
\begin{split}
&\texttt{b5582e8a6e62b2a6d3198dee3e6426593d313528b87f064abca12b671ebe1b4e},\\
&\texttt{f1121b62238ad5491bb7cf03635ea9934942edf573ed8faaa9ee79e1d38a6696}.
\end{split}
\]

SM V87 avoids differentiating a parameter-dependent Uhlenbeck gauge.
It freezes the reference connection at each chart centre, differentiates
only intrinsic tensor marks, and sews different centres through the
complete determinant line.  The transferable lesson is exact: a role
which merely chooses a local frame should be removed by an invariant
identity before operator norms are taken.  It must not be represented by
a fictitious marked operator.

NS V31 retains the complementary warning.  Exact mark removal restores a
first moment, but a one-use parent sum can remain analytically circular
when its absolute estimate is used to pay the very mark whose cancellation
created the sum.  In the present YM problem this forbids replacing the
actual recoupling packet by a scalar tree count and then spending the same
tree decay once more.

\begin{audit}[V10 transfer rule]
\label{audit:V10-transfer}
Invariant or scalar ownership labels are erased before the local operator
is factorized.  Genuine projections and recouplers remain in their exact
order.  An unmarked norm bound does not imply a two-sided marked heat
bound; the latter is certified only by the weighted similarity developed
below.
\end{audit}

\section{Role labels do not create operators}
\label{sec:V10-role-erasure}

\begin{definition}[Operator word and estimate labels]
\label{def:V10-word-label}
Let \(\mathbb W\) be a local operator obtained from the exact quintic
source by differentiation, integration, physical projection, and
contraction.  An \emph{estimate label} is auxiliary finite data used to
write an exact decomposition
\begin{equation}
 \mathbb W=\sum_{\lambda\in\Lambda}\mathbb W_\lambda
\label{eq:V10-labelled-decomp}
\end{equation}
or a scalar majorant
\begin{equation}
 \|\mathbb W_\lambda\|
 \le b_\lambda(a,\theta).
\label{eq:V10-labelled-majorant}
\end{equation}
It is not an operator factor unless the source formula for
\(\mathbb W_\lambda\) contains a corresponding linear map.
\end{definition}

\begin{lemma}[Scalar role erasure]
\label{lem:V10-role-erasure}
Suppose \(\lambda=(\tau,e_*,o)\) records a tree, a distinguished strong
edge, and an owner choice.  Assume that \(\tau,e_*,o\) occur only in the
indexing of \eqref{eq:V10-labelled-decomp} and in the scalar majorant
\eqref{eq:V10-labelled-majorant}.  Then changing or deleting these labels
does not change the ordered operator product defining any fixed
\(\mathbb W_\lambda\).  In particular, no bounded map
\(\mathbb S_\lambda\) may be inserted into a V9 heat sandwich solely from
the existence of one of these labels.
\end{lemma}

\begin{proof}
An ordered operator product is determined by the composable linear maps
appearing in the source formula.  By hypothesis the three labels occur
only in the domain of the finite sum and in a scalar function on the
right-hand side of an inequality.  Neither occurrence is composition by
a linear map.  Thus deletion of the labels leaves the product unchanged.
Conversely, inserting \(\mathbb S_\lambda\) would change the product unless
the source formula already contained that map.  Hence such an insertion
cannot be justified by bookkeeping alone.
\end{proof}

\begin{corollary}[Correct use of the V8 cumulant card]
\label{cor:V10-cumulant-use}
For the selected and prefix cumulants treated in V8, the Möbius sum is
assembled before its square envelope is taken.  Tree labels used to prove
the representation-uniform summand bounds do not create further
noncentral factors after that assembly.  Any remaining V9 sandwich
obstruction must therefore be an actual map not already contained in the
V8 cumulant.
\end{corollary}

\begin{proof}
Apply \Cref{lem:V10-role-erasure} to the tree decomposition and use the
assembled-first firewall \Cref{fire:V8-assembled-first}.  The Möbius
cumulant itself remains an operator and is controlled by the two
orientations proved in V8.  Only source-occurring maps outside that
assembled packet remain to be classified.
\end{proof}

\section{The exact weighted-similarity criterion}
\label{sec:V10-weighted-similarity}

Let \(\mathscr H\) be a finite-dimensional Hilbert space and let
\begin{equation}
 0<\mathbb H\preccurlyeq I,
 \qquad
 \mathbb A:=\mathbb H^{1/2}.
\label{eq:V10-positive-heat}
\end{equation}
Every finite Peter--Weyl block satisfies strict positivity; uniformity as
the representation labels grow is the issue.

\begin{theorem}[Unique sandwich coefficient for a one-sided map]
\label{thm:V10-unique-sandwich}
For any \(\mathbb T\in\mathcal B(\mathscr H)\), the factorizations
\begin{align}
 \mathbb T\mathbb H
 &=\mathbb H^{1/2}\mathbb S_R\mathbb H^{1/2},
\label{eq:V10-right-sandwich}\\
 \mathbb H\mathbb T
 &=\mathbb H^{1/2}\mathbb S_L\mathbb H^{1/2}
\label{eq:V10-left-sandwich}
\end{align}
have the unique coefficients
\begin{align}
 \mathbb S_R
 &=\mathbb H^{-1/2}\mathbb T\mathbb H^{1/2},
\label{eq:V10-SR}\\
 \mathbb S_L
 &=\mathbb H^{1/2}\mathbb T\mathbb H^{-1/2}.
\label{eq:V10-SL}
\end{align}
If \(\mathbb T=\mathbb T^*\), then
\(\mathbb S_L=\mathbb S_R^*\), so the two coefficient norms are equal.
\end{theorem}

\begin{proof}
Multiply \eqref{eq:V10-right-sandwich} on the left and right by
\(\mathbb H^{-1/2}\).  This forces \eqref{eq:V10-SR}; direct
multiplication proves existence.  The same argument gives
\eqref{eq:V10-SL}.  For self-adjoint \(\mathbb T\),
\[
 \mathbb S_R^*
 =\mathbb H^{1/2}\mathbb T\mathbb H^{-1/2}
 =\mathbb S_L.
\]
\end{proof}

\begin{corollary}[Uniform family criterion]
\label{cor:V10-uniform-criterion}
Let \((\mathbb H_\alpha,\mathbb T_\alpha)\) be any family of finite
blocks with \(0<\mathbb H_\alpha\preccurlyeq I\).  The right one-sided
packets \(\mathbb T_\alpha\mathbb H_\alpha\) are heat sandwiches with a
uniformly bounded middle coefficient if and only if
\begin{equation}
 \sup_\alpha
 \bigl\|
 \mathbb H_\alpha^{-1/2}
 \mathbb T_\alpha
 \mathbb H_\alpha^{1/2}
 \bigr\|<\infty.
\label{eq:V10-weighted-similarity}
\end{equation}
For self-adjoint \(\mathbb T_\alpha\), this one condition certifies both
one-sided orientations.
\end{corollary}

\begin{proof}
The coefficient in each block is unique by
\Cref{thm:V10-unique-sandwich}.  Therefore a uniformly bounded choice
exists precisely when the displayed supremum is finite.  The final
statement follows from equality of the two coefficient norms.
\end{proof}

\begin{lemma}[Weighted commutator form]
\label{lem:V10-commutator-form}
For every \(\mathbb T\),
\begin{align}
 \mathbb H^{-1/2}\mathbb T\mathbb H^{1/2}
 &=
 \mathbb T+
 \mathbb H^{-1/2}
 [\mathbb T,\mathbb H^{1/2}],
\label{eq:V10-commutator-right}\\
 \mathbb H^{1/2}\mathbb T\mathbb H^{-1/2}
 &=
 \mathbb T-
 [\mathbb T,\mathbb H^{1/2}]
 \mathbb H^{-1/2}.
\label{eq:V10-commutator-left}
\end{align}
Consequently
\[
 \sup_\alpha\|\mathbb T_\alpha\|<\infty,
 \qquad
 \sup_\alpha
 \bigl\|
 \mathbb H_\alpha^{-1/2}
 [\mathbb T_\alpha,\mathbb H_\alpha^{1/2}]
 \bigr\|<\infty
\]
is a sufficient right-sandwich certificate.  For self-adjoint
\(\mathbb T_\alpha\) it certifies both orientations.
\end{lemma}

\begin{proof}
Use
\(\mathbb T\mathbb H^{1/2}
=\mathbb H^{1/2}\mathbb T+
[\mathbb T,\mathbb H^{1/2}]\)
for the first identity and its rearrangement for the second.  The norm
claim follows from the triangle inequality and
\Cref{thm:V10-unique-sandwich}.
\end{proof}

\begin{corollary}[Compatible projections]
\label{cor:V10-compatible-projection}
If \(\mathbb P=\mathbb P^*=\mathbb P^2\) and
\([\mathbb P,\mathbb H]=0\), then
\[
 \mathbb P\mathbb H
 =\mathbb H\mathbb P
 =\mathbb H^{1/2}\mathbb P\mathbb H^{1/2}.
\]
Thus the middle coefficient has norm at most one in both orientations.
\end{corollary}

\begin{proof}
Functional calculus gives
\([\mathbb P,\mathbb H^{1/2}]=0\).  Now apply
\Cref{thm:V10-unique-sandwich}, for which both coefficients equal
\(\mathbb P\).
\end{proof}

\section{Sharp obstruction for an incompatible projection}
\label{sec:V10-projection-countertest}

\begin{proposition}[Norm-one projection with unbounded heat similarity]
\label{prop:V10-projection-counterexample}
For \(0<\varepsilon\le1\), set
\[
 \mathbb H_\varepsilon=
 \begin{pmatrix}1&0\\0&\varepsilon\end{pmatrix},
 \qquad
 \mathbb P={1\over2}
 \begin{pmatrix}1&1\\1&1\end{pmatrix}.
\]
Then \(\mathbb P\) is an orthogonal projection of norm one, but the
unique coefficient in
\[
 \mathbb P\mathbb H_\varepsilon
 =
 \mathbb H_\varepsilon^{1/2}
 \mathbb S_\varepsilon
 \mathbb H_\varepsilon^{1/2}
\]
is
\begin{equation}
 \mathbb S_\varepsilon
 ={1\over2}
 \begin{pmatrix}
 1&\sqrt\varepsilon\\
 \varepsilon^{-1/2}&1
 \end{pmatrix},
\label{eq:V10-bad-similarity}
\end{equation}
and
\begin{equation}
 \|\mathbb S_\varepsilon\|
 \ge {1\over2}\sqrt{1+\varepsilon^{-1}}.
\label{eq:V10-bad-lower}
\end{equation}
Hence ordinary contraction of a physical projection does not imply a
uniform heat sandwich.
\end{proposition}

\begin{proof}
Direct multiplication gives \(\mathbb P^2=\mathbb P=\mathbb P^*\).
Formula \eqref{eq:V10-SR} gives
\eqref{eq:V10-bad-similarity}.  Applying this matrix to
\((1,0)^{\mathsf T}\) yields
\[
 \|\mathbb S_\varepsilon(1,0)^{\mathsf T}\|
 ={1\over2}\sqrt{1+\varepsilon^{-1}},
\]
which proves \eqref{eq:V10-bad-lower}.  Uniqueness in
\Cref{thm:V10-unique-sandwich} rules out a different bounded middle
coefficient.
\end{proof}

\begin{firewall}[Why a basis change is not harmless]
\label{fire:V10-basis-change}
A recoupling matrix can be unitary in the unweighted Hilbert norm and
still mix eigenspaces on which the heat has exponentially different
weights.  \Cref{prop:V10-projection-counterexample} shows that unitarity
or projection norm one alone gives no uniform bound after conjugation by
\(\mathbb H^{1/2}\).  The exact weighted similarity, not the ordinary
operator norm, is the required datum.
\end{firewall}

\section{Casimir-compatible physical compression}
\label{sec:V10-compatible-compression}

\begin{definition}[Heat-compatible resolution]
\label{def:V10-compatible-resolution}
On a finite Peter--Weyl block, a physical compression
\(\mathbb P_\lambda\) is heat-compatible if it is a joint spectral
projection of self-adjoint block Casimirs which strongly commute with
the heat exponent defining \(\mathbb H_\lambda\).
\end{definition}

\begin{theorem}[Compatible compression preserves the V9 half-power]
\label{thm:V10-compatible-closes}
Let
\[
 \mathbb L_\lambda
 =\mathbb P_\lambda\mathbb H_\lambda
 \quad\hbox{or}\quad
 \mathbb L_\lambda
 =\mathbb H_\lambda\mathbb P_\lambda,
\]
where \(\mathbb P_\lambda\) is heat-compatible.  Assume the complete heat
moment \(\mathbb H_\lambda\) has the row-\(a\) partial-transpose
half-power \eqref{eq:V9-sandwich-heat-half}.  Then
\begin{align}
 \kappa_{a|a^c}(\mathbb L_\lambda^*\mathbb L_\lambda)
 &\le C a_{a,d}^{-1/2},\\
 \kappa_{a|a^c}(\mathbb L_\lambda\mathbb L_\lambda^*)
 &\le C a_{a,d}^{-1/2}.
\end{align}
If the compression also commutes with the total output Casimir, the same
conclusion holds with the V9 payer applied once.
\end{theorem}

\begin{proof}
By \Cref{cor:V10-compatible-projection},
\[
 \mathbb L_\lambda
 =\mathbb H_\lambda^{1/2}
  \mathbb P_\lambda
  \mathbb H_\lambda^{1/2}.
\]
This is a one-term packet in \Cref{def:V9-heat-sandwich}, with middle
norm at most one.  Apply \Cref{thm:V9-unpaid-sandwich-half}.  If the
compression commutes with the total output Casimir, it is equivariant
for the payer calculation, so
\Cref{thm:V9-paid-sandwich-half} applies.
\end{proof}

\begin{corollary}[Coarse total-isotypic pinching is safe]
\label{cor:V10-total-pinching}
The total-output isotypic projections retained in V8 commute with the
diagonal total Casimir and with every invariant heat moment.  They are
therefore heat-compatible, retain the full multiplicity space, and do
not obstruct either square orientation or the once-paid estimate.
\end{corollary}

\begin{proof}
Each total-isotypic projection is a spectral projection of the total
Casimir.  Invariance of the heat moment gives commutation.  Apply
\Cref{thm:V10-compatible-closes}; retention of the full multiplicity
space is exactly the multiplicity-coarse convention imposed in V8.
\end{proof}

\begin{corollary}[Same-tree coupling projections are safe]
\label{cor:V10-same-tree}
Suppose a coupling-path projection is a joint spectral projection of the
intermediate Casimirs used to build the heat exponent on that same
coupling tree.  Then it is heat-compatible and satisfies the conclusions
of \Cref{thm:V10-compatible-closes}.
\end{corollary}

\begin{proof}
The relevant intermediate Casimirs commute on the fixed coupling tree,
and the heat exponent is a Borel function of them.  Joint spectral
projections commute with every such function.  Apply
\Cref{thm:V10-compatible-closes}.
\end{proof}

\section{The crossed-resolution gate}
\label{sec:V10-crossed-gate}

\begin{definition}[Weighted recoupling constant]
\label{def:V10-weighted-recoupling}
For a physical projection or self-adjoint separator
\(\mathbb P_\lambda\) and a complete local heat moment
\(\mathbb H_\lambda\), define
\begin{equation}
 \mathfrak R(\mathbb P,\mathbb H)
 :=
 \sup_\lambda
 \left\|
 \mathbb H_\lambda^{-1/2}
 \mathbb P_\lambda
 \mathbb H_\lambda^{1/2}
 \right\|.
\label{eq:V10-recoupling-constant}
\end{equation}
The supremum ranges over the representation labels allowed by the
fixed-arity source packet.
\end{definition}

\begin{theorem}[Exact crossed-resolution alternative]
\label{thm:V10-crossed-alternative}
For a self-adjoint physical compression in a one-sided heat packet,
exactly one of the following holds:
\begin{enumerate}[label=\textup{(\alph*)},leftmargin=3em]
\item \(\mathfrak R(\mathbb P,\mathbb H)<\infty\), in which case the
      unpaid packet has both V9 square half-powers;
\item \(\mathfrak R(\mathbb P,\mathbb H)=\infty\), in which case no
      representation-uniform one-term heat-sandwich coefficient exists.
\end{enumerate}
In case \textup{(a)}, the once-paid conclusion additionally requires
equivariance for the total output Casimir.
\end{theorem}

\begin{proof}
\Cref{cor:V10-uniform-criterion} proves the exhaustive alternative and,
for self-adjoint compression, certifies both coefficient orientations.
The V9 two-square envelope then gives the unpaid half-powers.  The paid
theorem commutes the payer through the middle coefficient; this is
justified by total-output equivariance and is not a consequence of
weighted boundedness alone.
\end{proof}

\begin{corollary}[Exact residual \(q=3\) certificate after V10]
\label{cor:V10-q3-residual}
After scalar role erasure and the V8 assembled-cumulant estimate, every
V9 fixed-owner packet built only from
\begin{enumerate}[label=\textup{(\roman*)},leftmargin=3.6em]
\item complete heat moments,
\item total-isotypic pinching,
\item same-tree coupling projections, and
\item a once-used equivariant payer
\end{enumerate}
is sandwich-certified and closes by
\Cref{thm:V9-sandwich-owner-close}.  A surviving \(q=3\) packet must
contain an actual crossed-resolution map for which either
\begin{equation}
 \mathfrak R(\mathbb P,\mathbb H)=\infty
\label{eq:V10-residual-similarity}
\end{equation}
or total-output equivariance needed by the payer has not been proved.
\end{corollary}

\begin{proof}
The cumulant contributes its already proved V8 square card by
\Cref{cor:V10-cumulant-use}.  Items \textup{(i)--(iii)} are certified by
\Cref{thm:V10-compatible-closes,cor:V10-total-pinching,
cor:V10-same-tree}; item \textup{(iv)} is certified by the paid part of
\Cref{thm:V10-compatible-closes}.  V9 then supplies the missing fourth
half-power.  Negating this finite list leaves precisely a crossed
resolution or a non-equivariant payer interaction.  The former is
classified by \Cref{thm:V10-crossed-alternative}.
\end{proof}

\begin{assessment}[What has and has not been localized]
\label{ass:V10-localized}
V10 replaces the phrase noncentral owner by a verifiable matrix
quantity.  It does not prove that the actual quintic source needs a
crossed-resolution projection, nor that such a projection has finite or
infinite \(\mathfrak R\).  The next task is source inspection: list the
actual output compressions in the fixed-owner word and determine whether
multiplicity-coarse total pinching already suffices.  Only if a genuinely
crossed fine projection survives is a Racah-coefficient estimate needed.
\end{assessment}

\section{V10 result ledger and V11 acquisition order}
\label{sec:V10-ledger}

\begin{theorem}[Fourth-series V10 result ledger]
\label{thm:V10-results}
Relative to V1--V9, V10 proves:
\begin{enumerate}[label=\textup{(V10.\arabic*)},leftmargin=6.3em]
\item the compulsory SM/NS companion audit;
\item scalar erasure of tree, strong-edge, and owner labels which do not
      occur as maps in the source word;
\item a unique weighted-similarity formula for each one-sided heat
      sandwich;
\item an equivalent weighted-commutator certificate;
\item a sharp \(2\times2\) contraction counterexample to any
      ordinary-norm shortcut;
\item two-sided unpaid and paid half-powers for Casimir-compatible
      physical projections;
\item safety of total-isotypic and same-tree spectral pinching; and
\item reduction of the remaining \(q=3\) marked-owner issue to an actual
      crossed-resolution weighted similarity or a payer-equivariance
      failure.
\end{enumerate}
\end{theorem}

\begin{proof}
These are
\Cref{audit:V10-transfer,
lem:V10-role-erasure,
cor:V10-cumulant-use,
thm:V10-unique-sandwich,
cor:V10-uniform-criterion,
lem:V10-commutator-form,
prop:V10-projection-counterexample,
thm:V10-compatible-closes,
cor:V10-total-pinching,
cor:V10-same-tree,
thm:V10-crossed-alternative,
cor:V10-q3-residual}.
\end{proof}

\begin{programme}[V11: exact source-word projection census]
\label{prog:V10-V11}
\begin{enumerate}[leftmargin=3em]
\item Return to the quintic source formula before inequalities and list
      every actual separator, partial trace, output projection, and
      recoupling in the fixed-owner packet.
\item Delete every label-only role by
      \Cref{lem:V10-role-erasure}; retain differentiated heat factors as
      part of their assembled cumulant.
\item Test whether multiplicity-coarse total-isotypic pinching gives the
      same physical contraction as the fine resolution used in the old
      card.
\item For every unavoidable crossed projection, compute
      \(\mathfrak R(\mathbb P,\mathbb H)\) in the smallest \(SU(2)\)
      recoupling block before attempting a general compact-group bound.
\item If the weighted constant diverges, join the projection to both
      adjacent source heat factors and repeat the exact two-sided test.
\item Once the \(q=3\) word is closed or refuted, return to the
      \(q=2,1,0\) exponent table without reusing any owner.
\end{enumerate}
\end{programme}

\begin{firewall}[Claim boundary after fourth-series V10]
\label{fire:V10-claim-boundary}
V10 does not prove a representation-uniform bound for a crossed
recoupling matrix and does not prove the sharp \(q=3\) card in any packet
where such a matrix is indispensable.  It proves an exact criterion,
closes all compatible projection packets, and rules out ordinary norm
control as a substitute.  The residual \(q=3\) source census,
\(q=0,1,2\) cards, global quintic or all-order MTA, WUFC, CPM, cutoff
removal, reflection positivity, Osterwalder--Schrader reconstruction,
construction of the continuum Yang--Mills measure, and a uniform
positive continuum spectral gap remain open.  The full Yang--Mills mass
gap has not been proved.
\end{firewall}

\paragraph{Parent-version record.}
V10 was formed from
\texttt{Renewal\_Yang\_Mills\_Mass\_Gap\_Research\_Notes\_V9.tex}.
The V9 parent had \(4855\) logical source lines, \(173\,287\) bytes, and
SHA--256
\[
 \texttt{c767f12967f362ebe22db298d6c4743a6d2e0d2fb1d8dbfde63a9d99d0554be0}.
\]
The next compulsory companion audit is V15.

% =============================================================================
% END APPENDED FOURTH-SERIES V10 MATERIAL
% =============================================================================

% =============================================================================
% BEGIN APPENDED FOURTH-SERIES V11 MATERIAL
% =============================================================================

\clearpage
\chapter{V11: Exact source-word census, complete output duality, and
correction of the three-slot residual}
\label{ch:V11}

\section{Why the V10 projection gate is not the first live gate}
\label{sec:V11-purpose}

V10 gives the exact weighted-similarity criterion for a one-sided fine
projection.  The criterion is correct, but the compulsory upstream
source census changes its role.  The exact quintic first-connectivity
identity in archive f3 V97 is written before physical output resolution.
Its letters are complete prefix connected carriers, one complete
quotient cumulant, and a complete moment suffix.  Fine multiplicity
projections are not additional source maps.

Archive f2 V94 already explains how to estimate the complete output
Fourier carrier without pinching into multiplicity copies: dual
contractions are absorbed into the bounded equivariant separator on the
full output block.  Thus a crossed fine projection is an artefact of a
discarded pinched-absolute route unless the literal source formula is
shown to contain such a map.

The same audit finds a separate logical gap.  V9 conditionally closes a
packet having a cofinal fixed-strong owner \(d_*\notin E_3\), but its
residual corollary treats existence of \(d_*\) as automatic.  Three-slot
saturation does not imply that existence.  This chapter records the
correction and restores the exact V5 residual.

\section{The literal quintic source has no fine output projector}
\label{sec:V11-source-census}

\begin{record}[Imported exact first-connectivity identity]
\label{rec:V11-V97-source}
The proof indexed as \texttt{thm:V97-exact-quintic-source} in frozen
archive f3 states, with coefficient one,
\begin{equation}
 \mathcal C_5
 =
 \sum_r\ \sum_{\rho<\widehat1}
 \left(
   \bigotimes_{B\in\rho}
   \mathcal K_B^{<r,\mathrm{conn}}
 \right)
 \otimes\mathcal J_r^\rho
 \otimes\mathcal M_{>r}.
\label{eq:V11-exact-source}
\end{equation}
Here \(\mathcal K_B^{<r,\mathrm{conn}}\) is the complete connected
prefix carrier on \(B\), \(\mathcal J_r^\rho\) is the complete
quotient cumulant at the first-connectivity coordinate, and
\(\mathcal M_{>r}\) is the complete moment suffix.  Every connected
coordinate word occurs exactly once.
\end{record}

\begin{proposition}[Exact source-map census]
\label{prop:V11-source-map-census}
Before the optional output resolution used to estimate
\eqref{eq:V11-exact-source}, the operator factors in a fixed
\((r,\rho)\) summand are exactly:
\begin{enumerate}[label=\textup{(\roman*)},leftmargin=3.5em]
\item the complete prefix connected carriers, including singleton
      moments;
\item the complete joining cumulant; and
\item the complete suffix moment.
\end{enumerate}
The prefix-tree label, quotient-tree label, endpoint lift, strong-edge
designation, physical output label, and payer-location label are finite
decomposition or allocation data.  None is a further projection or
recoupling map in \eqref{eq:V11-exact-source}.
\end{proposition}

\begin{proof}
The right-hand side of \eqref{eq:V11-exact-source} explicitly displays
all compositional factors of each summand.  The tree and endpoint data
are introduced only after the complete prefix and joining cumulants have
been expanded into finitely many estimate fibres.  A physical output
label is introduced by inserting the spectral resolution of the
identity in a complete equivariant carrier; summing all its blocks
restores that carrier.  Payer location chooses where the single central
output multiplier is allocated and does not create a second multiplier.
Finally apply the role-erasure principle
\Cref{lem:V10-role-erasure}.  No other source map appears.
\end{proof}

\begin{firewall}[Resolution is not source incidence]
\label{fire:V11-resolution-not-source}
A spectral identity may be inserted to prove a bound, but its individual
summands cannot then be treated as compulsory maps in the original
word.  In particular, failure of a weighted bound for one fine pinched
summand does not obstruct the complete word when an exact complete-block
duality estimate is available.
\end{firewall}

\section{Complete output duality without fine pinching}
\label{sec:V11-output-duality}

Let an equivariant finite block decompose as
\begin{equation}
 \mathscr H
 =
 \bigoplus_{U\in\mathfrak U}
 H_U\otimes M_U,
\label{eq:V11-Schur-blocks}
\end{equation}
where \(H_U\) is irreducible and \(M_U\) is the full multiplicity space.
An equivariant operator has the form
\[
 \mathbb A
 =\bigoplus_U I_{H_U}\otimes a_U.
\]

\begin{lemma}[Dual contractions remain in the complete separator]
\label{lem:V11-dual-contraction}
Let
\begin{align}
 \mathbb A&=\bigoplus_U I_{H_U}\otimes a_U,\\
 \mathbb B&=\bigoplus_U I_{H_U}\otimes b_U,\\
 \mathbb C&=\bigoplus_U I_{H_U}\otimes c_U
\end{align}
be equivariant, with \(\sup_U\|c_U\|\le L\).  For any family of
contractions \(D_U\in\mathcal B(H_U)\), define
\begin{equation}
 \mathbb C_{\boldsymbol D}
 :=
 \bigoplus_U D_U^*\otimes c_U.
\label{eq:V11-dual-separator}
\end{equation}
Then
\begin{equation}
 \|\mathbb C_{\boldsymbol D}\|\le L
\label{eq:V11-dual-separator-norm}
\end{equation}
and the exact multiplicity multiplier in
\(\mathbb A\mathbb C_{\boldsymbol D}\mathbb B\) is
\(a_Uc_Ub_U\), with the dual contraction acting only on \(H_U\).
No multiplicity basis projection and no sum of block norms is created.
\end{lemma}

\begin{proof}
On the \(U\)-block,
\[
 \mathbb A\mathbb C_{\boldsymbol D}\mathbb B
 =
 D_U^*\otimes a_Uc_Ub_U.
\]
The norm of the direct sum is the supremum of its block norms, and
\[
 \|D_U^*\otimes c_U\|
 =\|D_U\|\,\|c_U\|
 \le L.
\]
This proves \eqref{eq:V11-dual-separator-norm}.  Since the entire
\(M_U\) factor remains inside \(a_Uc_Ub_U\), no choice of a
multiplicity vector has been made.
\end{proof}

\begin{theorem}[Nonsymmetric complete dual-sandwich inequality]
\label{thm:V11-nonsymmetric-dual}
Let \(\mathbb A,\mathbb B\) be arbitrary operators and let
\(\mathbb C_{\boldsymbol D}\) satisfy
\eqref{eq:V11-dual-separator-norm}.  For every density matrix
\(\varrho\),
\begin{equation}
 \left|
 \Tr\bigl(
  \varrho\mathbb A\mathbb C_{\boldsymbol D}\mathbb B
 \bigr)
 \right|
 \le
 L
 \sqrt{
  \Tr(\varrho\mathbb A\mathbb A^*)
  \Tr(\varrho\mathbb B^*\mathbb B)}.
\label{eq:V11-nonsymmetric-dual}
\end{equation}
Consequently, for product-state capacity across any fixed row cut,
\begin{equation}
 \sup_{\boldsymbol D}
 \kappa_\otimes^{\mathrm{abs}}
   (\mathbb A\mathbb C_{\boldsymbol D}\mathbb B)
 \le
 L
 \sqrt{
  \kappa_\otimes(\mathbb A\mathbb A^*)
  \kappa_\otimes(\mathbb B^*\mathbb B)}.
\label{eq:V11-capacity-dual}
\end{equation}
\end{theorem}

\begin{proof}
Schatten Hölder gives
\begin{align*}
 \left|
 \Tr(
  \varrho\mathbb A\mathbb C_{\boldsymbol D}\mathbb B)
 \right|
 &\le
 \|
  \varrho^{1/2}
  \mathbb A\mathbb C_{\boldsymbol D}
  \mathbb B\varrho^{1/2}
 \|_1\\
 &\le
 \|\varrho^{1/2}\mathbb A\|_2
 \|\mathbb C_{\boldsymbol D}\|
 \|\mathbb B\varrho^{1/2}\|_2.
\end{align*}
The squared Hilbert--Schmidt factors are
\[
 \|\varrho^{1/2}\mathbb A\|_2^2
 =\Tr(\varrho\mathbb A\mathbb A^*),
 \qquad
 \|\mathbb B\varrho^{1/2}\|_2^2
 =\Tr(\varrho\mathbb B^*\mathbb B).
\]
Insert \eqref{eq:V11-dual-separator-norm}.  Taking the supremum over
product densities and then over the contraction family proves
\eqref{eq:V11-capacity-dual}.
\end{proof}

\begin{corollary}[Two square cards replace a weighted recoupling
estimate]
\label{cor:V11-two-square-compiler}
Suppose a complete output-dual carrier has the exact form
\(\mathbb A\mathbb C_{\boldsymbol D}\mathbb B\), with the separator
bound in \Cref{lem:V11-dual-contraction}, and
\begin{align}
 \kappa_\otimes(\mathbb A\mathbb A^*)
 &\le C_A M^{-1/2},\\
 \kappa_\otimes(\mathbb B^*\mathbb B)
 &\le C_B M^{-1/2}.
\end{align}
Then
\begin{equation}
 \sup_{\boldsymbol D}
 \kappa_\otimes^{\mathrm{abs}}
 (\mathbb A\mathbb C_{\boldsymbol D}\mathbb B)
 \le L\sqrt{C_AC_B}\,M^{-1/2}.
\label{eq:V11-dual-half-power}
\end{equation}
No weighted similarity of an individual output projector is required.
\end{corollary}

\begin{proof}
Insert the two hypotheses into
\eqref{eq:V11-capacity-dual}.
\end{proof}

\begin{remark}[Relation to the frozen V94 theorem]
\label{rem:V11-V94}
\Cref{lem:V11-dual-contraction,thm:V11-nonsymmetric-dual} are the
nonsymmetric square-card form of the complete-block method indexed by
\texttt{lem:V94-separator-equivariant} and
\texttt{thm:V94-dual-sandwich} in archive f2.  The new formulation
matches the two oriented square estimates used in V8--V10 and avoids
an unnecessary self-adjointness hypothesis on the outer factors.
\end{remark}

\begin{corollary}[Complete quintic output resolution is not a crossed
projection obstruction]
\label{cor:V11-no-crossed-output}
For a fixed source summand in \eqref{eq:V11-exact-source}, retain every
complete equivariant prefix, joining, and suffix carrier and perform
output trace-norm duality only after assembly.  Canonical replica
regroupings and physical contractions then belong to
\(\mathbb C_{\boldsymbol D}\), with a fixed-arity norm bound.  Whenever
the two relevant oriented square cards have the required half-power,
the complete output carrier has the same half-power by
\Cref{cor:V11-two-square-compiler}.  Fine output pinching and its
weighted recoupling constant are absent.
\end{corollary}

\begin{proof}
Complete moments and cumulants intertwine the diagonal action, and a
canonical regrouping merely permutes tensor factors carrying that
action.  Their fixed-arity products are therefore equivariant and have
the Schur form in \eqref{eq:V11-Schur-blocks}.  The archived separator
size theorem \texttt{lem:V92-separator-size} supplies the uniform norm.
Apply \Cref{lem:V11-dual-contraction} and then
\Cref{cor:V11-two-square-compiler}.
\end{proof}

\begin{firewall}[The compiler does not manufacture its outer cards]
\label{fire:V11-outer-cards}
\Cref{cor:V11-no-crossed-output} removes a spurious output-resolution
loss.  It does not prove the two \(M^{-1/2}\) square cards for an outer
factor which lacks them, and it does not create a new source-occurring
coordinate.  Those are analytic and incidence questions, respectively.
\end{firewall}

\section{A countermodel to the distinct-owner inference}
\label{sec:V11-distinct-owner}

\begin{proposition}[Three-slot saturation does not force an outside
cofinal coordinate]
\label{prop:V11-no-outside-owner}
The hypotheses
\[
 q=3,
 \qquad
 \sum_{e\in E_3}a_{a,e}\ge(1-3\eta)M,
 \qquad
 \theta_{ac}^{E_3}\ge\eta/4
\]
do not imply that there exists \(d_*\notin E_3\) with
\(a_{a,d_*}\ge\sigma M\) for any fixed \(\sigma>0\).
\end{proposition}

\begin{proof}
Let the row-\(a\) support be exactly
\(E_3=\{e_1,e_2,e_3\}\), put
\[
 a_{a,e_i}=M/3
 \quad(1\le i\le3),
 \qquad
 a_{a,d}=0
 \quad(d\notin E_3),
\]
and choose nonnegative overlap stocks on the same three coordinates
with
\(\theta_{ac,e_i}=1/3\).  Then the displayed saturation and overlap
hypotheses hold for every \(0<\eta<1/3\), while every coordinate
outside \(E_3\) has zero row-\(a\) mass.  Hence the asserted \(d_*\)
does not exist.

These scalar data are also realized at the capacity level by the
equal-spin three-singlet family of
\Cref{prop:V7-no-same-slot-fourth}; that family has exactly the sharp
three-slot \(M^{-3/2}\) heat-square response.  Thus the countermodel is
consistent with the local representation obstruction already proved.
\end{proof}

\begin{correction}[Scope correction to V9 and V10]
\label{corr:V11-V9-scope}
\Cref{thm:V9-sandwich-owner-close} remains valid as a conditional
theorem when a distinct cofinal \(d_*\notin E_3\) is part of the
hypotheses.  However, the automatic-existence assertion in
\Cref{cor:V9-residual-owner} and the corresponding global reading of
\Cref{cor:V10-q3-residual} are withdrawn.  They may be used only on
packets for which the exact source incidence independently supplies
such a \(d_*\).
\end{correction}

\begin{proof}
The conditional V9 proof tensorizes the three \(E_3\) factors with the
distinct \(d_*\) factor, so it is valid under its stated hypothesis.
\Cref{prop:V11-no-outside-owner} disproves derivation of that hypothesis
from saturation.  V10 classifies the operator at \(d_*\) but likewise
does not prove its existence.  Therefore precisely the asserted scope
restriction is necessary.
\end{proof}

\section{Restored exact \(q=3\) boundary}
\label{sec:V11-restored-boundary}

\begin{theorem}[Correct residual after output-duality cleanup]
\label{thm:V11-correct-residual}
After all established central, covariance, deleted-prefix, selected
cumulant, compatible-compression, and complete output-duality closures
are applied, a persistent \(q=3\) counterpacket must retain at least one
of:
\begin{enumerate}[label=\textup{(\alph*)},leftmargin=3.5em]
\item a cofinal source-occurring assembled owner whose required oriented
      square card has not been proved;
\item a fixed fraction of maximal-row mass in the V5 non-owner
      remainder \(R\), without a proved complete source occurrence; or
\item three-slot saturation on \(E_3\), with no distinct cofinal
      source-occurring owner outside \(E_3\).
\end{enumerate}
Fine output multiplicity pinching is not an additional alternative.
\end{theorem}

\begin{proof}
\Cref{cor:V5-surviving-q3} gives the exhaustive alternatives before
V6--V10.  V6--V8 close the suffix, quotient-order-two, selected, and
deleted-prefix subbranches.  V9--V10 close any distinct owner carrying
their explicit sandwich certificate.  On the complete output carrier,
\Cref{cor:V11-no-crossed-output} replaces fine pinching by the exact
dual separator and hence deletes that estimating artefact.  The
remaining possibilities are exactly \textup{(a)--(c)}.  Their
exhaustiveness does not use the invalid distinct-owner inference,
which has been withdrawn in \Cref{corr:V11-V9-scope}.
\end{proof}

\begin{corollary}[Necessary nature of the next gain]
\label{cor:V11-necessary-gain}
In alternative \textup{(c)} of
\Cref{thm:V11-correct-residual}, no argument using only the three
already counted positive heat slots and scalar strong overlap can
supply the missing fourth \(M^{-1/2}\) factor.  Any closure must instead
prove one of:
\begin{enumerate}[label=\textup{(\roman*)},leftmargin=3.5em]
\item an additional cancellation gain in the assembled prefix or
      quotient cumulant on the same three coordinates; or
\item an exact source-incidence theorem exposing a further complete
      coordinate or bank.
\end{enumerate}
\end{corollary}

\begin{proof}
The equal-spin three-singlet lower bound
\Cref{prop:V7-no-same-slot-fourth} shows that the three positive heat
slots have sharp response \(M^{-3/2}\), even with fixed strong overlap.
The target square scale is \(M^{-2}\).  A scalar label cannot alter an
operator by \Cref{lem:V10-role-erasure}, and a fine output resolution
cannot create the missing gain by
\Cref{cor:V11-no-crossed-output}.  Thus the additional half-power must
come from cancellation in an existing assembled connected carrier or
from a genuinely additional source occurrence.
\end{proof}

\begin{assessment}[Upstream direction after V11]
\label{ass:V11-upstream}
The weighted \(6j\) problem is now secondary.  The first live question
is incidence and cancellation inside the coefficient-one
first-connectivity summand.  For each of the six prefix types one must
locate the three saturated coordinates relative to the prefix carriers,
the joining cumulant, and the suffix.  If all three are positive suffix
moments, the joining/prefix factor must provide a centered cancellation
gain.  If one lies inside a connected carrier, its complete square card
must be used there rather than counted again as heat.
\end{assessment}

\section{V11 result ledger and V12 acquisition order}
\label{sec:V11-ledger}

\begin{theorem}[Fourth-series V11 result ledger]
\label{thm:V11-results}
Relative to V1--V10, V11 proves:
\begin{enumerate}[label=\textup{(V11.\arabic*)},leftmargin=6.3em]
\item the exact source-map census of the coefficient-one quintic
      first-connectivity formula;
\item removal of tree, endpoint, strong, output, and payer labels from
      the list of literal source maps;
\item a complete-block dual-contraction lemma retaining all
      multiplicity spaces;
\item a nonsymmetric two-square dual-sandwich inequality;
\item elimination of fine output pinching as a quintic
      crossed-projection obstruction;
\item a countermodel to automatic existence of a distinct cofinal
      strong owner outside the saturated three slots;
\item the necessary scope correction to the V9--V10 residual; and
\item the restored exhaustive \(q=3\) boundary and its two possible
      sources of the missing gain.
\end{enumerate}
\end{theorem}

\begin{proof}
These are
\Cref{rec:V11-V97-source,
prop:V11-source-map-census,
lem:V11-dual-contraction,
thm:V11-nonsymmetric-dual,
cor:V11-no-crossed-output,
prop:V11-no-outside-owner,
corr:V11-V9-scope,
thm:V11-correct-residual,
cor:V11-necessary-gain}.
\end{proof}

\begin{programme}[V12: six-type incidence and centered-gain census]
\label{prog:V11-V12}
\begin{enumerate}[leftmargin=3em]
\item For each prefix type
      \(4+1,3+2,3+1+1,2+2+1,2+1+1+1,1^5\), place the three
      saturated coordinates among prefix, joining, and suffix letters
      of \eqref{eq:V11-exact-source}.
\item Consolidate repeated coordinate visits before assigning a
      positive heat or cumulant response.
\item Whenever a saturated coordinate lies in a connected prefix or
      joining carrier, use the complete oriented square card of that
      carrier and do not also count its heat marginal.
\item Isolate the cases in which all three saturated responses belong
      only to the positive suffix.  In those cases write the exact
      centered prefix--joining operator before any absolute value.
\item Test the smallest \(SU(2)\) block for a cancellation half-power.
      A lower-bound counterexample ends that route; otherwise formulate
      the representation-uniform centered-gain lemma actually needed.
\item Keep the non-owner remainder as a separate source-incidence
      problem; do not fold its mass into the marked-owner set.
\end{enumerate}
\end{programme}

\begin{firewall}[Claim boundary after fourth-series V11]
\label{fire:V11-claim-boundary}
V11 removes a spurious fine-projection obstruction and corrects an
unsupported distinct-owner inference.  It does not prove the centered
same-three-coordinate gain, a new source-incidence theorem, or the full
\(q=3\) card.  The remaining \(q=3\) alternatives, \(q=0,1,2\) cards,
global quintic or all-order MTA, WUFC, CPM, cutoff removal, reflection
positivity, Osterwalder--Schrader reconstruction, construction of the
continuum Yang--Mills measure, and a uniform positive continuum
spectral gap remain open.  The full Yang--Mills mass gap has not been
proved.
\end{firewall}

\paragraph{Parent-version record.}
V11 was formed from
\texttt{Renewal\_Yang\_Mills\_Mass\_Gap\_Research\_Notes\_V10.tex}.
The V10 parent had \(5461\) logical source lines, \(195\,691\) bytes,
and SHA--256
\[
 \texttt{66b4756faee43f344b0c5e07759852a9abed221a902988a4efcb88fa5c1942cf}.
\]
The next compulsory companion audit is V15.

% =============================================================================
% END APPENDED FOURTH-SERIES V11 MATERIAL
% =============================================================================

% =============================================================================
% BEGIN APPENDED FOURTH-SERIES V12 MATERIAL
% =============================================================================

\clearpage
\chapter{V12: Temporal incidence and the protected-diagonal test for a
same-support centered gain}
\label{ch:V12}

\section{Purpose}
\label{sec:V12-purpose}

V11 restores the true saturated \(q=3\) problem: the maximal row may
be supported entirely on the three already counted heat coordinates.
There need not be a fourth cofinal owner.  Since the three-slot
\(SU(2)\) example is sharp, the missing half-power must come either
from a source occurrence not visible in the scalar mass ledger or from
genuine cancellation in a complete connected carrier acting on the
same support.

This chapter first enumerates the temporal positions of the three slots
relative to first connectivity.  It then proves a necessary
protected-sector condition for any same-support cancellation theorem.
Finally it shows by a Bernoulli--singlet construction that Möbius
centering alone does not provide the missing decay.  The live
calculation is reduced to a finite protected matrix coefficient of the
actual prefix or quotient carrier.

\section{Seven temporal distributions}
\label{sec:V12-temporal}

Fix the unique first-connectivity coordinate \(r\) in
\eqref{eq:V11-exact-source}.  For the three distinct saturated
coordinates \(E_3\), put
\begin{equation}
 p:=|E_3\cap\{e:e<r\}|,
 \qquad
 j:=\boldsymbol1_{\{r\in E_3\}},
 \qquad
 s:=|E_3\cap\{e:e>r\}|.
\label{eq:V12-pjs}
\end{equation}

\begin{lemma}[Exact seven-pattern temporal census]
\label{lem:V12-seven-patterns}
The triple \((p,j,s)\) belongs to the following seven-element set:
\begin{equation}
 \boxed{
 \begin{split}
 \mathfrak I_3
 =\{&
 (0,0,3),(1,0,2),(2,0,1),(3,0,0),\\
 & (0,1,2),(1,1,1),(2,1,0)\}.
 \end{split}}
\label{eq:V12-seven-patterns}
\end{equation}
Every possible temporal placement of three distinct coordinates occurs
in exactly one displayed class.
\end{lemma}

\begin{proof}
Exactly one of the relations \(e<r\), \(e=r\), \(e>r\) holds for
each \(e\in E_3\).  Since the coordinates are distinct, at most one
equals \(r\), so
\[
 p+j+s=3,
 \qquad
 j\in\{0,1\}.
\]
For \(j=0\), the four nonnegative solutions are
\((p,s)=(0,3),(1,2),(2,1),(3,0)\).  For \(j=1\), the three solutions
are \((p,s)=(0,2),(1,1),(2,0)\).  The defining strict inequalities make
the classes disjoint, and the displayed equation makes them exhaustive.
\end{proof}

\begin{corollary}[Finite type-incidence table]
\label{cor:V12-type-incidence}
Combining \(\mathfrak I_3\) with the six prefix block-size types in
\Cref{rec:V11-V97-source} produces \(6\cdot7=42\) raw
type-incidence classes.  Combining with all \(51\) proper prefix
partitions produces \(51\cdot7=357\) raw
partition-incidence classes before tree, endpoint-lift, payer, and row
symmetries.
\end{corollary}

\begin{proof}
Temporal order of physical coordinates and block-size type of the row
partition are independent pieces of finite indexing data.  Apply the
product rule to the two finite sets.
\end{proof}

\begin{remark}[What the census does not assert]
\label{rem:V12-incidence-scope}
\Cref{cor:V12-type-incidence} counts possible index pairs.  It does not
assert that every pair supports a nonzero source word.  Vanishing is
decided only after the exact prefix and quotient cumulants are evaluated.
\end{remark}

\section{The three-singlet protected line}
\label{sec:V12-protected-line}

For \(G=SU(2)\), let \(V_j\) be the spin-\(j\) irreducible module,
\[
 d_j:=2j+1,
 \qquad
 M_j:=1+j(j+1).
\]
Let \(\Omega_j\in V_j\otimes V_j^*\) be the normalized invariant
vector and
\[
 P_j:=|\Omega_j\rangle\langle\Omega_j|.
\]
On three distinct coordinates define
\begin{equation}
 \Omega_j^{(3)}:=\Omega_j^{\otimes3},
 \qquad
 P_j^{(3)}:=P_j^{\otimes3}.
\label{eq:V12-three-singlets}
\end{equation}

\begin{lemma}[Exact product overlap]
\label{lem:V12-product-overlap}
Across the cut consisting of the three first tensor legs versus the
three dual tensor legs,
\begin{equation}
 \kappa_\otimes(P_j^{(3)})=d_j^{-3}.
\label{eq:V12-product-overlap}
\end{equation}
\end{lemma}

\begin{proof}
For one coordinate, the largest Schmidt coefficient of
\(\Omega_j\) is \(d_j^{-1/2}\), so the largest squared overlap with a
product unit vector is \(d_j^{-1}\).  The Schmidt coefficients of the
threefold tensor product are all \(d_j^{-3/2}\).  Hence its largest
squared product overlap is \(d_j^{-3}\), which is the product-state
capacity of its rank-one projector.
\end{proof}

\begin{theorem}[Protected-diagonal lower bound]
\label{thm:V12-protected-diagonal}
Let \(\mathbb K_j\) be any operator on the three-coordinate Hilbert
space and put
\begin{equation}
 \gamma_j
 :=
 \left|
 \langle
  \Omega_j^{(3)},
  \mathbb K_j\Omega_j^{(3)}
 \rangle
 \right|.
\label{eq:V12-gamma}
\end{equation}
For the two one-sided packets
\[
 \mathbb Z_{R,j}:=\mathbb K_jP_j^{(3)},
 \qquad
 \mathbb Z_{L,j}:=P_j^{(3)}\mathbb K_j,
\]
one has
\begin{align}
 \kappa_\otimes(\mathbb Z_{R,j}^*\mathbb Z_{R,j})
 &\ge \gamma_j^2d_j^{-3},
\label{eq:V12-right-lower}\\
 \kappa_\otimes(\mathbb Z_{L,j}\mathbb Z_{L,j}^*)
 &\ge \gamma_j^2d_j^{-3}.
\label{eq:V12-left-lower}
\end{align}
\end{theorem}

\begin{proof}
Choose a product unit vector \(\phi_j\) attaining the largest Schmidt
overlap in \Cref{lem:V12-product-overlap}.  Then
\[
 P_j^{(3)}\phi_j
 =
 \alpha_j\Omega_j^{(3)},
 \qquad
 |\alpha_j|^2=d_j^{-3}.
\]
Therefore
\begin{align*}
 \langle\phi_j,
  \mathbb Z_{R,j}^*\mathbb Z_{R,j}\phi_j\rangle
 &=
 |\alpha_j|^2
 \|\mathbb K_j\Omega_j^{(3)}\|^2\\
 &\ge
 d_j^{-3}
 \left|
 \langle\Omega_j^{(3)},
  \mathbb K_j\Omega_j^{(3)}
 \rangle
 \right|^2.
\end{align*}
This proves \eqref{eq:V12-right-lower}.  Apply the same calculation to
\(\mathbb K_j^*\) to obtain
\eqref{eq:V12-left-lower}; its protected diagonal is the complex
conjugate and has the same modulus.
\end{proof}

\begin{corollary}[Necessary decay of the protected coefficient]
\label{cor:V12-necessary-gamma}
If either square in
\Cref{thm:V12-protected-diagonal} obeys a uniform target estimate
\begin{equation}
 \kappa_\otimes(\mathbb Z^*\mathbb Z)
 \le C M_j^{-2}
\label{eq:V12-target-M2}
\end{equation}
in its corresponding orientation, then
\begin{equation}
 \gamma_j\le C' M_j^{-1/4}.
\label{eq:V12-gamma-decay}
\end{equation}
In particular, a representation-independent lower bound
\(\gamma_j\ge\gamma_0>0\) contradicts the \(M_j^{-2}\) card.
\end{corollary}

\begin{proof}
Since \(d_j\asymp M_j^{1/2}\),
\Cref{thm:V12-protected-diagonal} and
\eqref{eq:V12-target-M2} give
\[
 \gamma_j^2
 \le C M_j^{-2}d_j^3
 \le C' M_j^{-1/2}.
\]
Taking square roots proves \eqref{eq:V12-gamma-decay}.
\end{proof}

\begin{firewall}[The test is necessary, not sufficient]
\label{fire:V12-gamma-not-sufficient}
Decay or vanishing of the single coefficient \(\gamma_j\) does not
prove the desired upper bound.  Off-diagonal protected-to-gapped matrix
elements can still be large.  A failed diagonal test disproves the
route, while a passed test only permits the full square analysis to
continue.
\end{firewall}

\section{Exact protected coefficient of a Möbius cumulant}
\label{sec:V12-Mobius-coefficient}

Let \(m\ge2\), let \(\Pi_m\) be the partition lattice, and let
\[
 \mu(\pi,\widehat1)
 =(-1)^{|\pi|-1}(|\pi|-1)!.
\]
Consider an assembled local cumulant
\begin{equation}
 \mathbb K_m
 =
 \sum_{\pi\in\Pi_m}
 \mu(\pi,\widehat1)\mathbb M_\pi.
\label{eq:V12-Mobius-cumulant}
\end{equation}

\begin{lemma}[Common-eigenline reduction]
\label{lem:V12-common-eigenline}
Suppose a unit vector \(\Omega\) is a common eigenvector of all moment
partitions,
\[
 \mathbb M_\pi\Omega=m_\pi\Omega.
\]
Then
\begin{equation}
 \langle\Omega,\mathbb K_m\Omega\rangle
 =
 \sum_{\pi\in\Pi_m}
 \mu(\pi,\widehat1)m_\pi.
\label{eq:V12-protected-Mobius}
\end{equation}
If \(m_\pi=1\) for every \(\pi\), this coefficient is zero.
\end{lemma}

\begin{proof}
Apply \eqref{eq:V12-Mobius-cumulant} to \(\Omega\) and take its inner
product with \(\Omega\), obtaining
\eqref{eq:V12-protected-Mobius}.  If all eigenvalues are one, the
coefficient is
\[
 \sum_{\pi\in\Pi_m}\mu(\pi,\widehat1)=0,
\]
the Möbius cancellation identity on every nontrivial interval
\([\widehat0,\widehat1]\) with \(m\ge2\).
\end{proof}

\begin{assessment}[The finite calculation now required]
\label{ass:V12-finite-calculation}
For each nonzero class in the \(42\)-type temporal table, the first
protected test is to identify a common invariant line or multiplicity
block for the actual prefix--joining carrier and compute the finite
Möbius coefficient \eqref{eq:V12-protected-Mobius}.  It is not enough
to say that the carrier is centered: intermediate nontrivial block
outputs can fuse back to total output zero, making the numbers
\(m_\pi\) partition-dependent.
\end{assessment}

\section{Centering alone has no representation gain}
\label{sec:V12-centering-no-go}

\begin{proposition}[Bernoulli--singlet cumulant countertest]
\label{prop:V12-Bernoulli-no-go}
For each fixed \(m\in\{2,3,4,5\}\), there are bounded commuting random
operators \(Y_{1,j},\ldots,Y_{m,j}\) whose assembled order-\(m\)
cumulant is a nonzero fixed scalar multiple of \(P_j\).  Consequently
its two square capacities are comparable to
\[
 d_j^{-1}\asymp M_j^{-1/2}.
\]
Thus Möbius centering by itself does not improve the sharp one-slot
half-power.
\end{proposition}

\begin{proof}
Let \(B\) be a Bernoulli random variable with parameter
\(p\in(0,1)\), and set
\[
 Y_{1,j}=\cdots=Y_{m,j}:=B P_j.
\]
The operators commute and have norm at most one.  For every nonempty
block \(D\subseteq[m]\),
\[
 \mathbb E\prod_{i\in D}Y_{i,j}
 =pP_j,
\]
because \(B^{|D|}=B\) and \(P_j^{|D|}=P_j\).  Hence the
moment--cumulant formula gives
\begin{equation}
 \kappa_m(Y_{1,j},\ldots,Y_{m,j})
 =
 \left[
  \sum_{\pi\in\Pi_m}
  \mu(\pi,\widehat1)p^{|\pi|}
 \right]P_j
 =
 \kappa_m(B)P_j.
\label{eq:V12-Bernoulli-cumulant}
\end{equation}
For fixed \(m\), the polynomial \(\kappa_m(B)\) is not identically
zero: its coefficient of \(p\) is one.  Choose \(p\in(0,1)\) away
from its finitely many zeros and put
\(c_m:=\kappa_m(B)\ne0\).

The square of \eqref{eq:V12-Bernoulli-cumulant} is
\(|c_m|^2P_j\).  Its product-state capacity is
\(|c_m|^2d_j^{-1}\), by the one-coordinate version of
\Cref{lem:V12-product-overlap}.  Since
\(d_j\asymp M_j^{1/2}\), this has exactly the stated scale in both
orientations.
\end{proof}

\begin{corollary}[Three heat slots plus abstract centering remain
short]
\label{cor:V12-three-plus-centering}
Tensor the countertest in
\Cref{prop:V12-Bernoulli-no-go} with protected singlet heat factors
on the same three-coordinate support so that the cumulant acts as
\(c_m\) on their protected line.  Then the complete square has
product-state capacity
\[
 |c_m|^2d_j^{-3}
 \asymp M_j^{-3/2},
\]
not \(O(M_j^{-2})\).
\end{corollary}

\begin{proof}
On the range of \(P_j^{(3)}\), arrange the cumulant factor as
\(c_mP_j\) on one slot and the identity on the other two.  Its product
with \(P_j^{(3)}\) is \(c_mP_j^{(3)}\).  Apply
\Cref{lem:V12-product-overlap}.
\end{proof}

\begin{firewall}[Abstract countertest versus the YM source]
\label{fire:V12-abstract-scope}
\Cref{prop:V12-Bernoulli-no-go} is not asserted to be an actual
Yang--Mills source packet.  It proves that no theorem whose only
hypotheses are boundedness, commutativity, and Möbius centering can
supply the missing gain.  A positive YM result must use additional
source-specific heat eigenvalues, representation selection rules, or
incidence.
\end{firewall}

\section{A precise same-support acquisition card}
\label{sec:V12-acquisition-card}

\begin{definition}[Protected-decay test \(PDT_{\rho,\iota}\)]
\label{def:V12-PDT}
Fix a prefix type \(\rho\) and temporal pattern
\(\iota=(p,j,s)\in\mathfrak I_3\).  Let
\(\mathbb K_{\rho,\iota}(j)\) be the complete prefix--joining operator
remaining after the three protected heat slots are retained in their
literal source order.  The first protected-decay test is
\begin{equation}
 PDT_{\rho,\iota}:\qquad
 \left|
 \langle
 \Omega_j^{(3)},
 \mathbb K_{\rho,\iota}(j)
 \Omega_j^{(3)}
 \rangle
 \right|
 \le C_{\rho,\iota}M_j^{-1/4}.
\label{eq:V12-PDT}
\end{equation}
\end{definition}

\begin{proposition}[Failure of \(PDT\) excludes a same-support
\(M^{-2}\) square card]
\label{prop:V12-PDT-necessary}
If, along a representation sequence, the left side of
\eqref{eq:V12-PDT} is bounded below by a positive constant, then the
corresponding one-sided same-support square cannot satisfy a uniform
\(CM_j^{-2}\) estimate.
\end{proposition}

\begin{proof}
Use the complete operator as \(\mathbb K_j\) in
\Cref{thm:V12-protected-diagonal} and apply
\Cref{cor:V12-necessary-gamma}.
\end{proof}

\begin{programme}[How the card is to be used]
\label{prog:V12-card-use}
The finite source inspection must proceed in this order:
\begin{enumerate}[leftmargin=3em]
\item delete zero source classes and quotient by row symmetries;
\item compute \eqref{eq:V12-protected-Mobius} on every surviving
      protected multiplicity block;
\item stop a same-support route immediately if
      \Cref{prop:V12-PDT-necessary} applies;
\item if every protected diagonal vanishes or decays sufficiently,
      estimate the protected-to-gapped off-diagonal blocks;
\item only after the complete two-orientation square is controlled,
      aggregate the finite source fibres.
\end{enumerate}
\end{programme}

\section{V12 result ledger and V13 acquisition order}
\label{sec:V12-ledger}

\begin{theorem}[Fourth-series V12 result ledger]
\label{thm:V12-results}
Relative to V1--V11, V12 proves:
\begin{enumerate}[label=\textup{(V12.\arabic*)},leftmargin=6.3em]
\item the exhaustive seven-pattern temporal incidence census for three
      saturated coordinates;
\item the resulting \(42\) type-incidence and \(357\)
      partition-incidence raw counts;
\item the exact three-singlet product overlap \(d_j^{-3}\);
\item a protected-diagonal lower bound for both one-sided square
      orientations;
\item the necessary coefficient decay \(O(M_j^{-1/4})\) for an
      \(M_j^{-2}\) target;
\item reduction of a common-eigenline cumulant coefficient to a finite
      Möbius sum;
\item a Bernoulli--singlet countertest proving that abstract centering
      gives no representation gain; and
\item the source-specific protected-decay test for every surviving
      type-incidence class.
\end{enumerate}
\end{theorem}

\begin{proof}
These are
\Cref{lem:V12-seven-patterns,
cor:V12-type-incidence,
lem:V12-product-overlap,
thm:V12-protected-diagonal,
cor:V12-necessary-gamma,
lem:V12-common-eigenline,
prop:V12-Bernoulli-no-go,
def:V12-PDT,
prop:V12-PDT-necessary}.
\end{proof}

\begin{programme}[V13: protected block of the actual quotient
cumulant]
\label{prog:V12-V13}
\begin{enumerate}[leftmargin=3em]
\item Start with quotient order \(k=2\), where
      \(\mathcal J_r^\rho\) is a complete covariance, and compute its
      action on every total-trivial \(SU(2)\) multiplicity channel.
\item Separate the one-dimensional two-row singlet from higher
      invariant multiplicity created by intermediate nontrivial
      outputs.
\item Evaluate the two moment-partition eigenvalues before subtraction;
      do not infer cancellation merely from the word covariance.
\item Repeat for \(k=3\) only after the covariance block is settled.
\item Use the seven temporal classes to decide whether the tested
      protected line lies in the prefix, joining, or suffix carrier.
\item If a nondecaying protected coefficient survives, abandon
      same-support centering and return to exact source-incidence
      acquisition for the non-owner remainder.
\end{enumerate}
\end{programme}

\begin{firewall}[Claim boundary after fourth-series V12]
\label{fire:V12-claim-boundary}
V12 proves necessary protected-sector tests and rules out abstract
centering as the missing half-power.  It does not evaluate the actual
YM quotient-cumulant protected blocks, control their off-diagonal
parts, prove a new source occurrence, or close \(q=3\).  The remaining
\(q=3\) alternatives, \(q=0,1,2\) cards, global quintic or all-order
MTA, WUFC, CPM, cutoff removal, reflection positivity,
Osterwalder--Schrader reconstruction, construction of the continuum
Yang--Mills measure, and a uniform positive continuum spectral gap
remain open.  The full Yang--Mills mass gap has not been proved.
\end{firewall}

\paragraph{Parent-version record.}
V12 was formed from
\texttt{Renewal\_Yang\_Mills\_Mass\_Gap\_Research\_Notes\_V11.tex}.
The V11 parent had \(5976\) logical source lines, \(215\,221\) bytes,
and SHA--256
\[
 \texttt{2c0b75efdc2fac5f06952a58cea7ffc14c04a93439ac5bb6d1134be3865c3b2a}.
\]
The next compulsory companion audit is V15.

% =============================================================================
% END APPENDED FOURTH-SERIES V12 MATERIAL
% =============================================================================

% =============================================================================
% BEGIN APPENDED FOURTH-SERIES V13 MATERIAL
% =============================================================================

\clearpage
\chapter{V13: A payer-safe low-output covariance obstruction at
quotient order two}
\label{ch:V13}

\section{Purpose}
\label{sec:V13-purpose}

The first two rows of the V97 prefix table, \(4+1\) and \(3+2\),
have quotient order two.  Their joining carrier is a complete
covariance.  V12 asks whether this centered carrier, when it shares
the three saturated coordinates, can provide a fourth inverse
half-power.

The total singlet is unsuitable as a universal countertest when the
sole payer acts on the joining output, because the payer vanishes on
output zero.  The fixed spin-one output avoids that issue.  Its heat
weight and payer are fixed nonzero constants, while its overlap with
an explicit product vector is of order \(d_j^{-1}\).  This chapter
computes the complete covariance on that channel and proves that it is
sharp at the ordinary one-slot half-power, paid or unpaid.

\section{Exact two-block covariance spectrum}
\label{sec:V13-covariance-spectrum}

Fix \(G=SU(2)\) and heat time \(s>0\).  Write
\begin{equation}
 q_L:=e^{-sL(L+1)}
\label{eq:V13-qL}
\end{equation}
for the normalized heat eigenvalue on spin \(L\).  Let two quotient
blocks have output modules \(V_j\) and \(V_j\).  Their tensor product
decomposes multiplicity freely as
\[
 V_j\otimes V_j
 =
 \bigoplus_{K=0}^{2j}V_K,
\]
and \(P_{j,K}\) denotes the complete spin-\(K\) projection.

\begin{lemma}[Complete covariance multiplier]
\label{lem:V13-covariance-multiplier}
Let \(\mathbb H_{12}(s)\) be the common-variable heat moment of the
two blocks and let
\(\mathbb H_1(s)\otimes\mathbb H_2(s)\) be their cut-replicated
endpoint moment.  The complete covariance
\begin{equation}
 \mathbb D_j(s)
 :=
 \mathbb H_{12}(s)
 -
 \mathbb H_1(s)\otimes\mathbb H_2(s)
\label{eq:V13-covariance}
\end{equation}
has the spectral resolution
\begin{equation}
 \boxed{
 \mathbb D_j(s)
 =
 \sum_{K=0}^{2j}
 (q_K-q_j^2)P_{j,K}.}
\label{eq:V13-covariance-spectrum}
\end{equation}
\end{lemma}

\begin{proof}
The common-variable heat moment is the heat functional calculus of
the total diagonal Casimir.  It therefore acts by \(q_K\) on the
spin-\(K\) summand.  In the cut-replicated moment the two quotient
blocks are heated independently, so their eigenvalue on the fixed
block outputs \(j,j\) is \(q_jq_j=q_j^2\), independently of their
subsequent total coupling.  Subtract the two diagonal multipliers.
\end{proof}

\begin{corollary}[Nondecaying fixed spin-one coefficient]
\label{cor:V13-spin-one-coefficient}
On \(P_{j,1}\), for \(j\ge1\),
\begin{equation}
 \mathbb D_j(s)P_{j,1}
 =
 \Delta_j(s)P_{j,1},
 \qquad
 \Delta_j(s):=e^{-2s}-e^{-2sj(j+1)}.
\label{eq:V13-Delta}
\end{equation}
There is \(j_s<\infty\) such that
\begin{equation}
 \Delta_j(s)\ge {1\over2}e^{-2s}
 \qquad(j\ge j_s).
\label{eq:V13-Delta-lower}
\end{equation}
\end{corollary}

\begin{proof}
Set \(K=1\) in \eqref{eq:V13-covariance-spectrum}.  Since
\(e^{-2sj(j+1)}\to0\), choose \(j_s\) so that this term is at most
\(\frac12e^{-2s}\).
\end{proof}

\section{An explicit product vector in the spin-one output}
\label{sec:V13-product-vector}

Let
\[
 \phi_j:=|j,j\rangle\otimes|j,-j\rangle.
\]

\begin{lemma}[Exact spin-\(K\) product weight]
\label{lem:V13-CG-weight}
For every integer \(K\) with \(0\le K\le2j\),
\begin{equation}
 \langle\phi_j,P_{j,K}\phi_j\rangle
 =
 (2K+1)
 \frac{[(2j)!]^2}
 {(2j-K)!(2j+K+1)!}.
\label{eq:V13-CG-weight}
\end{equation}
In particular,
\begin{equation}
 p_{j,1}
 :=
 \langle\phi_j,P_{j,1}\phi_j\rangle
 =
 \frac{6j}{(2j+1)(2j+2)}.
\label{eq:V13-pj1}
\end{equation}
\end{lemma}

\begin{proof}
The vector \(\phi_j\) has total magnetic number zero.  Hence its
projection on the spin-\(K\) isotype is its squared
Clebsch--Gordan coefficient
\[
 |\langle j,j;j,-j\mid K,0\rangle|^2.
\]
Substitution of
\(j_1=j_2=j\), \(m_1=j\), \(m_2=-j\), \(M=0\)
in the Racah factorial formula leaves one summation term and gives
\eqref{eq:V13-CG-weight}.  For \(K=1\), cancel
\[
 \frac{[(2j)!]^2}{(2j-1)!(2j+2)!}
 =
 \frac{2j}{(2j+1)(2j+2)}
\]
and multiply by three, proving \eqref{eq:V13-pj1}.
\end{proof}

\begin{corollary}[Sharp inverse-dimension scale]
\label{cor:V13-pj1-scale}
For all \(j\ge1\),
\begin{equation}
 {1\over 2j+2}
 \le p_{j,1}
 \le {3\over2j+1}.
\label{eq:V13-pj1-scale}
\end{equation}
Thus \(p_{j,1}\asymp d_j^{-1}\asymp M_j^{-1/2}\).
\end{corollary}

\begin{proof}
From \eqref{eq:V13-pj1},
\[
 p_{j,1}
 =
 {3\over2j+1}{2j\over2j+2}.
\]
For \(j\ge1\), the second factor lies between \(1/2\) and one.
This gives a lower bound \(3/[2(2j+1)]\), which is stronger than the
displayed lower bound, and the stated upper bound.  Finally
\(d_j=2j+1\) and \(M_j=1+j(j+1)\) are comparable at the claimed
scales.
\end{proof}

\section{Paid and unpaid square lower bounds}
\label{sec:V13-square-lower}

\begin{theorem}[Sharp spin-one covariance half-power]
\label{thm:V13-covariance-sharp}
For \(j\ge j_s\),
\begin{equation}
 \kappa_\otimes\!\left(
  (\mathbb D_j(s)P_{j,1})^*
  (\mathbb D_j(s)P_{j,1})
 \right)
 \ge c_s d_j^{-1}
 \asymp c_s M_j^{-1/2}.
\label{eq:V13-unpaid-lower}
\end{equation}
The same lower bound holds in the opposite square orientation.
\end{theorem}

\begin{proof}
The operator in either square is
\(\Delta_j(s)^2P_{j,1}\).  Test it on the product vector
\(\phi_j\).  By
\Cref{cor:V13-spin-one-coefficient,cor:V13-pj1-scale},
\[
 \langle\phi_j,
  \Delta_j(s)^2P_{j,1}\phi_j\rangle
 =
 \Delta_j(s)^2p_{j,1}
 \ge
 {e^{-4s}\over4(2j+2)}.
\]
This proves the stated capacity lower bound.  The compressed covariance
is self-adjoint, so the two square orientations coincide.
\end{proof}

The selected output payer of V4 is
\[
 \mathbb R_s
 =
 \boldsymbol1_{\{C_{\rm out}>0\}}
 {C_{\rm out}\over1-e^{-sC_{\rm out}}}.
\]
On spin one it has the fixed value
\begin{equation}
 r_1(s):={2\over1-e^{-2s}}.
\label{eq:V13-r1}
\end{equation}

\begin{theorem}[The output payer does not remove the spin-one
obstruction]
\label{thm:V13-paid-sharp}
For \(j\ge j_s\),
\begin{equation}
 \kappa_\otimes\!\left(
  (\mathbb R_s\mathbb D_j(s)P_{j,1})^*
  (\mathbb R_s\mathbb D_j(s)P_{j,1})
 \right)
 \ge c_s^{\rm pay}d_j^{-1}.
\label{eq:V13-paid-lower}
\end{equation}
Again the opposite orientation has the same lower bound.
\end{theorem}

\begin{proof}
On the selected block the paid operator is
\(r_1(s)\Delta_j(s)P_{j,1}\).  The number \(r_1(s)\) is finite and
strictly positive for fixed \(s>0\).  Repeat the product-vector test in
\Cref{thm:V13-covariance-sharp}, multiplying its lower bound by
\(r_1(s)^2\).
\end{proof}

\begin{remark}[Why output zero was not used]
\label{rem:V13-not-singlet}
On the total singlet, the unpaid covariance coefficient is
\(1-q_j^2\), also nondecaying.  But the selected payer vanishes there
by definition.  Spin one proves a single lower bound valid on both the
unpaid and paid joining carriers.
\end{remark}

\section{Three same-support coordinates}
\label{sec:V13-three-coordinates}

Define the resolved fixed-output heat factor
\[
 \mathbb H_{j,1}(s):=q_1P_{j,1}.
\]
On three distinct coordinates put
\begin{equation}
 \mathbb Z_j
 :=
 \bigl(\mathbb D_j(s)P_{j,1}\bigr)
 \otimes
 \mathbb H_{j,1}(s)
 \otimes
 \mathbb H_{j,1}(s).
\label{eq:V13-three-packet}
\end{equation}

\begin{proposition}[A same-support covariance is not a fourth
half-power]
\label{prop:V13-no-fourth}
There is \(c_s>0\) such that
\begin{equation}
 \kappa_\otimes(\mathbb Z_j^*\mathbb Z_j)
 \ge c_s d_j^{-3}
 \asymp c_sM_j^{-3/2}
\label{eq:V13-three-lower}
\end{equation}
for all \(j\ge j_s\).  The assertion remains true if the covariance
factor is paid by \(\mathbb R_s\).  Consequently these three
coordinates do not satisfy a uniform \(CM_j^{-2}\) square estimate.
\end{proposition}

\begin{proof}
Test the tensor product on
\(\phi_j^{\otimes3}\).  The first-coordinate expectation is
\(\Delta_j(s)^2p_{j,1}\), while each remaining expectation is
\(q_1^2p_{j,1}\).  Thus
\[
 \langle\phi_j^{\otimes3},
  \mathbb Z_j^*\mathbb Z_j
  \phi_j^{\otimes3}\rangle
 =
 \Delta_j(s)^2q_1^4p_{j,1}^3
 \ge c_s d_j^{-3}.
\]
For the paid factor, multiply by the fixed positive number
\(r_1(s)^2\).  Since
\(d_j^{-3}\asymp M_j^{-3/2}\), this cannot be bounded by
\(CM_j^{-2}\) uniformly in \(j\).
\end{proof}

\begin{corollary}[Quotient-order-two same-support centering is not a
universal closure]
\label{cor:V13-k2-no-universal}
In the \(4+1\) and \(3+2\) prefix types, the fact that the joining
carrier is a complete covariance cannot by itself provide the missing
same-support half-power in the saturated \(q=3\) branch.  This remains
true for every payer placement at or away from the joining carrier.
\end{corollary}

\begin{proof}
Both prefix types have quotient order two by the V97 interface table,
so their joining carrier contains the covariance spectrum
\eqref{eq:V13-covariance-spectrum}.  The unpaid and joining-paid
possibilities are covered by
\Cref{prop:V13-no-fourth}; a payer elsewhere only multiplies a different
complete factor and cannot turn the covariance coefficient into an
additional inverse power.  Hence a theorem using only the three heat
responses plus covariance centering is contradicted by the displayed
channel.
\end{proof}

\begin{firewall}[This is a channel obstruction, not yet a full source
counterexample]
\label{fire:V13-channel-scope}
The full \(4+1\) or \(3+2\) source summand also contains its complete
prefix connected carriers and suffix.  Those factors may annihilate
the spin-one channel chosen above.  V13 therefore rules out a universal
local covariance-gain lemma; it does not prove that a nonzero complete
quintic source word survives on the channel.
\end{firewall}

\section{The exact surviving-source gate}
\label{sec:V13-survival-gate}

\begin{definition}[Spin-one channel survival]
\label{def:V13-channel-survival}
For a fixed quotient-order-two source fibre, let
\(\mathbb A_j\) and \(\mathbb B_j\) denote the complete prefix and
suffix factors on the two sides of the joining covariance after exact
output dualization.  The spin-one channel survives if there are
row-product unit vectors \(\xi_j,\zeta_j\) and \(c_0>0\) such that,
for all large \(j\),
\begin{equation}
 \left|
 \langle\xi_j,
  \mathbb A_jP_{j,1}\mathbb B_j\zeta_j
 \rangle
 \right|
 \ge c_0
 \left|
 \langle\xi_j,P_{j,1}\zeta_j\rangle
 \right|.
\label{eq:V13-channel-survival}
\end{equation}
\end{definition}

\begin{proposition}[A surviving channel transfers the obstruction]
\label{prop:V13-survival-transfers}
If \eqref{eq:V13-channel-survival} holds with product vectors whose
spin-one overlap is bounded below by \(c\,d_j^{-1/2}\), then insertion
of the complete covariance, paid or unpaid, cannot give a square upper
bound \(o(d_j^{-1})\) on that coordinate.
\end{proposition}

\begin{proof}
On \(P_{j,1}\), the covariance and payer are scalar by
\eqref{eq:V13-Delta} and \eqref{eq:V13-r1}.  The survival inequality
therefore lower-bounds the relevant matrix coefficient by
\[
 c_0|\Delta_j(s)|\,c\,d_j^{-1/2}
\]
in the unpaid case, and by the same expression times \(r_1(s)\) in
the paid case.  The square expectation is at least the squared modulus
of this coefficient.  Use \eqref{eq:V13-Delta-lower}.
\end{proof}

\begin{assessment}[Upstream consequence]
\label{ass:V13-upstream}
For quotient order two the next question is no longer the covariance
itself.  It is whether the quartic prefix in type \(4+1\), or the
ternary--binary prefix product in type \(3+2\), annihilates every fixed
low total-output channel compatible with three-slot saturation.  If
one channel survives, V13 ends the same-support centered route for that
source fibre and forces an additional source-incidence response.
\end{assessment}

\section{V13 result ledger and V14 acquisition order}
\label{sec:V13-ledger}

\begin{theorem}[Fourth-series V13 result ledger]
\label{thm:V13-results}
Relative to V1--V12, V13 proves:
\begin{enumerate}[label=\textup{(V13.\arabic*)},leftmargin=6.3em]
\item the exact complete two-block covariance spectrum;
\item a nondecaying spin-one covariance coefficient;
\item the exact highest/lowest-weight product overlap with each
      spin-\(K\) output;
\item inverse-dimension spin-one product weight;
\item sharp paid and unpaid covariance-square lower bounds;
\item a three-coordinate \(M^{-3/2}\) lower bound showing that the
      same-support covariance is not a fourth half-power;
\item exclusion of a universal covariance-centering closure for both
      quotient-order-two prefix types; and
\item reduction to an exact low-output channel-survival test for the
      surrounding prefix and suffix.
\end{enumerate}
\end{theorem}

\begin{proof}
These are
\Cref{lem:V13-covariance-multiplier,
cor:V13-spin-one-coefficient,
lem:V13-CG-weight,
cor:V13-pj1-scale,
thm:V13-covariance-sharp,
thm:V13-paid-sharp,
prop:V13-no-fourth,
cor:V13-k2-no-universal,
def:V13-channel-survival,
prop:V13-survival-transfers}.
\end{proof}

\begin{programme}[V14: source survival in the \(4+1\) and \(3+2\)
prefixes]
\label{prog:V13-V14}
\begin{enumerate}[leftmargin=3em]
\item Insert the complete restored quartic MTA carrier in the \(4+1\)
      prefix and determine its spin-\(j\) output support before taking
      norms.
\item Do the same for the ternary and binary prefix carriers in the
      \(3+2\) type.
\item Use only complete output isotypes and the V11 dual separator; do
      not split multiplicity copies.
\item If a fixed spin-one joining channel survives, combine
      \Cref{prop:V13-survival-transfers} with the three-slot packet and
      move to source-incidence acquisition.
\item If every fixed low output is annihilated, prove the exact
      selection rule and use the forced growing output to seek
      exponential heat decay.
\item Keep all payer positions, noting that the spin-one payer is a
      fixed nonzero constant.
\end{enumerate}
\end{programme}

\begin{firewall}[Claim boundary after fourth-series V13]
\label{fire:V13-claim-boundary}
V13 proves that a complete quotient covariance has no universal
same-support cancellation gain, even when it pays.  It does not prove
survival of the obstructing channel through a complete quintic prefix,
nor does it settle quotient orders three through five, source
occurrence, or \(q=3\).  The remaining \(q=3\) alternatives,
\(q=0,1,2\) cards, global quintic or all-order MTA, WUFC, CPM, cutoff
removal, reflection positivity, Osterwalder--Schrader reconstruction,
construction of the continuum Yang--Mills measure, and a uniform
positive continuum spectral gap remain open.  The full Yang--Mills
mass gap has not been proved.
\end{firewall}

\paragraph{Parent-version record.}
V13 was formed from
\texttt{Renewal\_Yang\_Mills\_Mass\_Gap\_Research\_Notes\_V12.tex}.
The V12 parent had \(6502\) logical source lines, \(232\,192\) bytes,
and SHA--256
\[
 \texttt{8b8451873304e33e647f8fcf8623f431c834689ab076621edfdaf6c127926c02}.
\]
The next compulsory companion audit is V15.

% =============================================================================
% END APPENDED FOURTH-SERIES V13 MATERIAL
% =============================================================================

% =============================================================================
% BEGIN APPENDED FOURTH-SERIES V14 MATERIAL
% =============================================================================

\clearpage
\chapter{V14: The singleton-prefix exponential barrier in the
\texorpdfstring{\(4+1\)}{4+1} source}
\label{ch:V14}

\section{Purpose and exact scope}
\label{sec:V14-purpose}

V13 shows that the quotient-order-two covariance has a fixed low-output
channel with only the sharp inverse-dimension half-power.  To decide
whether that channel survives the complete source, one must inspect the
prefix.  The \(4+1\) source has a special factor: its singleton prefix
block is the complete one-row prefix moment.  On a nontrivial
representation this factor is scalar heat, so a single cofinal prefix
coordinate gives exponential suppression.

This does not treat the four source partitions in which the maximal row
belongs to the four-block, and it does not apply to \(3+2\), which has
no singleton.  It does close the maximal-singleton \(4+1\) sector
whenever one of the three extracted cofinal coordinates lies before
first connectivity.  In the seven-pattern table of V12 this removes
five temporal classes.

\section{One-row heat is scalar}
\label{sec:V14-one-row}

At physical coordinate \(e\), let the maximal row \(a\) carry a
nontrivial irreducible representation \(\lambda_{a,e}\), and write
\[
 c_{a,e}:=C_2(\lambda_{a,e}),
 \qquad
 a_{a,e}:=1+c_{a,e}.
\]
The normalized one-row heat moment at time \(s>0\) is denoted
\(\mathbb h_{a,e}(s)\).

\begin{lemma}[Exact one-row multiplier]
\label{lem:V14-one-row-multiplier}
On \(V_{\lambda_{a,e}}\),
\begin{equation}
 \mathbb h_{a,e}(s)
 =
 e^{-sc_{a,e}}I.
\label{eq:V14-one-row-multiplier}
\end{equation}
If \(a_{a,e}\ge\sigma M\), then, for all sufficiently large \(M\),
\begin{equation}
 \|\mathbb h_{a,e}(s)\|
 \le
 e^{-(s\sigma/2)M}.
\label{eq:V14-one-row-exponential}
\end{equation}
\end{lemma}

\begin{proof}
The Fourier transform of the normalized compact-group heat kernel is
the scalar \(e^{-sC_2(\lambda)}\) on the irreducible module
\(V_\lambda\), proving \eqref{eq:V14-one-row-multiplier}.  The cofinal
hypothesis gives
\[
 c_{a,e}=a_{a,e}-1\ge\sigma M-1.
\]
For \(M\ge2/\sigma\), this is at least \((\sigma/2)M\), proving
\eqref{eq:V14-one-row-exponential}.
\end{proof}

\begin{corollary}[A singleton prefix retains the scalar factor]
\label{cor:V14-singleton-prefix}
Let \(\rho=\{B,\{a\}\}\) be a \(4+1\) prefix partition whose singleton
is the maximal row.  In the exact source summand
\eqref{eq:V11-exact-source}, the factor
\(\mathcal K_{\{a\}}^{<r,\mathrm{conn}}\) is the complete one-row
prefix moment.  Hence for every \(e<r\) it contains
\(\mathbb h_{a,e}(s)\) as a tensor factor.
\end{corollary}

\begin{proof}
The imported V97 source identity specifies that singleton blocks carry
their complete one-row prefix moments.  Coordinate independence makes
that moment the tensor product of the one-row heat multipliers at all
coordinates below \(r\).  The factor at a fixed \(e<r\) is therefore
\eqref{eq:V14-one-row-multiplier}.
\end{proof}

\section{Exponential absorption of every fixed-order debit}
\label{sec:V14-absorption}

\begin{lemma}[Polynomial times cofinal one-row heat]
\label{lem:V14-polynomial-absorption}
For \(c,\sigma>0\) and nonnegative integers \(m,N\), there is
\(C_{c,\sigma,m,N}<\infty\) such that
\begin{equation}
 M^m e^{-c\sigma M}
 \le C_{c,\sigma,m,N}M^{-N}
\label{eq:V14-polynomial-absorption}
\end{equation}
for every \(M\ge1\).
\end{lemma}

\begin{proof}
The continuous function
\[
 x\longmapsto x^{m+N}e^{-c\sigma x}
\]
is bounded on \([1,\infty)\): it tends to zero at infinity and is
continuous on every compact interval.  Its supremum is the desired
constant.
\end{proof}

\begin{lemma}[The sole payer is only polynomial on a cofinal block]
\label{lem:V14-payer-polynomial}
Suppose the payer acts on an output whose Casimir \(c_{\rm out}\) is
bounded by the fixed-arity fusion estimate \(c_{\rm out}\le C_GM\).
Then
\begin{equation}
 \left\|
 \boldsymbol1_{\{c_{\rm out}>0\}}
 {c_{\rm out}\over1-e^{-sc_{\rm out}}}
 \right\|
 \le C_{G,s}(1+M).
\label{eq:V14-payer-polynomial}
\end{equation}
If it acts at the cofinal singleton heat coordinate, its product with
\(\mathbb h_{a,e}(s)\) is \(O(M^{-N})\) for every \(N\).
\end{lemma}

\begin{proof}
On nonzero compact-group outputs the Casimir has a positive spectral
gap \(c_*>0\), so
\[
 1-e^{-sc_{\rm out}}
 \ge1-e^{-sc_*}>0.
\]
Together with \(c_{\rm out}\le C_GM\), this proves
\eqref{eq:V14-payer-polynomial}.  The second assertion follows from
\Cref{lem:V14-one-row-multiplier,
lem:V14-polynomial-absorption}.
\end{proof}

\begin{theorem}[Complete-source exponential barrier]
\label{thm:V14-complete-barrier}
Fix a \(4+1\) source fibre with singleton \(\{a\}\), where
\(\Xi_a=M\).  Suppose some coordinate \(e_*<r\) satisfies
\begin{equation}
 a_{a,e_*}\ge\sigma M.
\label{eq:V14-prefix-cofinal}
\end{equation}
Retain the complete quartic prefix carrier on the other block, the
complete quotient covariance, the complete suffix, all canonical
separators, physical output duality, and any one payer placement.
After the established row normalization, the contribution of this
source fibre has, in either square orientation, an additional factor
\begin{equation}
 O_{G,s,\sigma,N}(M^{-N})
\label{eq:V14-arbitrary-power}
\end{equation}
for every fixed \(N\).
\end{theorem}

\begin{proof}
By \Cref{cor:V14-singleton-prefix}, the literal source word contains
the tensor factor \(\mathbb h_{a,e_*}(s)\).  It is scalar on its entire
row factor, so it commutes through every operation at other coordinates
and through the complete output dual separator.  By
\eqref{eq:V14-one-row-exponential}, its norm is at most
\(e^{-(s\sigma/2)M}\).

All complete fixed-order moment and cumulant factors have their
established representation-uniform raw norms.  The fixed number of
canonical separators has a uniform norm.  There is only one payer.  If
its output grows with \(M\), use
\eqref{eq:V14-payer-polynomial}; if its output remains bounded, its norm
is constant.  Thus the square of the entire remaining raw packet costs
at most a fixed polynomial \(C_{G,s}M^m\), with \(m\) independent of
the support and representation labels.  The complete square is
therefore bounded by
\[
 C_{G,s}M^m e^{-s\sigma M}.
\]
Apply \Cref{lem:V14-polynomial-absorption}.  The argument is unchanged
after taking adjoints, proving both orientations.
\end{proof}

\begin{corollary}[The fixed spin-one covariance channel is killed]
\label{cor:V14-kills-spin-one}
Under the hypotheses of
\Cref{thm:V14-complete-barrier}, the channel in
\Cref{prop:V13-no-fourth} does not survive the complete \(4+1\)
source.  In particular, it cannot obstruct the required \(M^{-2}\)
square card.
\end{corollary}

\begin{proof}
Choose \(N=2\) in \eqref{eq:V14-arbitrary-power}.  The estimate is for
the complete source word and hence includes the spin-one output
component.  It is already at or below the target square scale.
\end{proof}

\section{Five temporal classes close in the maximal-singleton sector}
\label{sec:V14-temporal-closure}

The three extracted coordinates are individually cofinal: after the
fixed finite slot selection there is \(\sigma>0\) such that
\begin{equation}
 a_{a,e}\ge\sigma M
 \qquad(e\in E_3).
\label{eq:V14-E3-cofinal}
\end{equation}

\begin{theorem}[Maximal-singleton \(4+1\) temporal closure]
\label{thm:V14-five-pattern-close}
Let the \(4+1\) prefix singleton be the maximal row \(a\).  If the
temporal pattern \((p,j,s)\) from V12 has \(p\ge1\), the complete
source fibre satisfies the missing \(q=3\) square response for every
payer placement and both orientations.  The five closed patterns are
\begin{equation}
 (1,0,2),(2,0,1),(3,0,0),(1,1,1),(2,1,0).
\label{eq:V14-five-patterns}
\end{equation}
\end{theorem}

\begin{proof}
If \(p\ge1\), choose
\[
 e_*\in E_3\cap\{e:e<r\}.
\]
Then \eqref{eq:V14-E3-cofinal} gives
\eqref{eq:V14-prefix-cofinal}.  Apply
\Cref{thm:V14-complete-barrier} with \(N=2\).  The list
\eqref{eq:V14-five-patterns} is exactly the sublist with \(p\ge1\) in
\eqref{eq:V12-seven-patterns}.
\end{proof}

\begin{corollary}[Remaining maximal-singleton temporal classes]
\label{cor:V14-two-patterns}
In the maximal-singleton \(4+1\) sector, the singleton-prefix argument
leaves only
\begin{equation}
 (p,j,s)=(0,0,3)
 \quad\hbox{or}\quad
 (p,j,s)=(0,1,2)
\label{eq:V14-two-patterns}
\end{equation}
unless a different cofinal row-\(a\) coordinate below \(r\) is exposed.
\end{corollary}

\begin{proof}
These are the two elements of
\(\mathfrak I_3\) with \(p=0\).  If another cofinal coordinate
\(e_*<r\) exists outside \(E_3\), then
\Cref{thm:V14-complete-barrier} still applies; absent such a source
occurrence, the present argument has no cofinal one-row prefix factor.
\end{proof}

\section{Why the result does not extend formally to the other
quotient-order-two sources}
\label{sec:V14-limit}

\begin{firewall}[The maximal row must be the singleton]
\label{fire:V14-maximal-singleton}
If \(a\) belongs to the four-block, the singleton prefix moment acts
on another row.  No lower bound on that row's prefix Casimir follows
from \(\Xi_a=M\).  Moving the scalar heat factor from row \(a\) to the
singleton would therefore be an invalid mass transfer.
\end{firewall}

\begin{firewall}[The \(3+2\) prefix has no one-row factor]
\label{fire:V14-no-singleton}
Both prefix blocks in type \(3+2\) are connected nonsingleton
carriers.  Their Möbius sums contain moment partitions with different
intermediate outputs, so neither block may be replaced by the scalar
one-row multiplier \eqref{eq:V14-one-row-multiplier}.  A separate
protected-channel or incidence argument is required.
\end{firewall}

\begin{assessment}[Exact residual inside quotient order two]
\label{ass:V14-k2-residual}
After V14, the quotient-order-two same-support problem remains only in:
\begin{enumerate}[label=\textup{(\roman*)},leftmargin=3.5em]
\item \(4+1\) sources with the maximal row in the four-block;
\item maximal-singleton \(4+1\) sources with no cofinal prefix
      occurrence, including the two patterns
      \eqref{eq:V14-two-patterns}; and
\item all \(3+2\) sources.
\end{enumerate}
For these sectors V13 says that covariance centering alone is
insufficient.  The surrounding connected prefix must either annihilate
the fixed low-output channel or expose another complete response.
\end{assessment}

\section{V14 result ledger and V15 acquisition order}
\label{sec:V14-ledger}

\begin{theorem}[Fourth-series V14 result ledger]
\label{thm:V14-results}
Relative to V1--V13, V14 proves:
\begin{enumerate}[label=\textup{(V14.\arabic*)},leftmargin=6.3em]
\item the exact scalar one-row heat multiplier;
\item exponential decay from one cofinal one-row coordinate;
\item absorption of every fixed-order polynomial and the sole payer;
\item a complete-source exponential barrier retaining the quartic
      prefix, joining covariance, suffix, separators, outputs, and
      payer;
\item removal of the V13 spin-one obstruction in that sector;
\item closure of five of seven temporal patterns for \(4+1\) sources
      whose singleton is the maximal row; and
\item the exact three-part quotient-order-two residual.
\end{enumerate}
\end{theorem}

\begin{proof}
These are
\Cref{lem:V14-one-row-multiplier,
cor:V14-singleton-prefix,
lem:V14-polynomial-absorption,
lem:V14-payer-polynomial,
thm:V14-complete-barrier,
cor:V14-kills-spin-one,
thm:V14-five-pattern-close,
cor:V14-two-patterns,
ass:V14-k2-residual}.
\end{proof}

\begin{programme}[V15: compulsory audit and the \(3+2\) prefix]
\label{prog:V14-V15}
\begin{enumerate}[leftmargin=3em]
\item Audit the newest active SM and NS companion notes and record
      filenames, sizes, and SHA--256 digests.
\item Inspect the exact ternary and binary connected prefix factors in
      the \(3+2\) source before any positive envelope.
\item On the V13 spin-one joining channel, resolve the two prefix block
      outputs but retain their full multiplicity spaces.
\item Decide whether connectedness forces a cofinal one-row heat factor
      inside either prefix block.  Do not assume this from the block
      output alone.
\item If no scalar factor is forced, compute the protected diagonal of
      the complete ternary--binary prefix product and apply the V12
      lower-bound test.
\item Keep the maximal-row-in-four-block \(4+1\) residual separate.
\end{enumerate}
\end{programme}

\begin{firewall}[Claim boundary after fourth-series V14]
\label{fire:V14-claim-boundary}
V14 closes five temporal classes only in the \(4+1\) sector whose
singleton is the maximal row.  It does not close the other four
\(4+1\) placements, the two empty-cofinal-prefix patterns, any
\(3+2\) source, quotient orders three through five, the non-owner
remainder, or the full \(q=3\) card.  The remaining \(q=3\)
alternatives, \(q=0,1,2\) cards, global quintic or all-order MTA,
WUFC, CPM, cutoff removal, reflection positivity,
Osterwalder--Schrader reconstruction, construction of the continuum
Yang--Mills measure, and a uniform positive continuum spectral gap
remain open.  The full Yang--Mills mass gap has not been proved.
\end{firewall}

\paragraph{Parent-version record.}
V14 was formed from
\texttt{Renewal\_Yang\_Mills\_Mass\_Gap\_Research\_Notes\_V13.tex}.
The V13 parent had \(6984\) logical source lines, \(247\,446\) bytes,
and SHA--256
\[
 \texttt{f097121d0c2eaf7b86b66559dcf11df668cdfb69d0d17cc551a7d996d8c41db1}.
\]
The next compulsory companion audit is V15.

% =============================================================================
% END APPENDED FOURTH-SERIES V14 MATERIAL
% =============================================================================

% =============================================================================
% BEGIN APPENDED FOURTH-SERIES V15 MATERIAL
% =============================================================================

\clearpage
\chapter{V15: Companion audit and recursive minimal connected owners}
\label{ch:V15}

\section{Compulsory companion audit}
\label{sec:V15-audit}

The newest active companion files at the time of the V15 audit were
\begin{center}
\begin{tabular}{@{}p{0.49\linewidth}rr@{}}
\toprule
file & logical lines & bytes\\
\midrule
\texttt{Renewal\_SM\_Einstein\_Continuum\_Research\_Notes\_V93.tex}
 & \(35\,055\) & \(1\,291\,910\)\\
\texttt{Renewal\_NS\_Stability\_Research\_Notes\_V31.tex}
 & \(23\,052\) & \(887\,537\)\\
\bottomrule
\end{tabular}
\end{center}
Their SHA--256 digests are
\[
\begin{split}
&\texttt{98477aa07b4c32cacccf629b9d28570751a36185862af83e297819a774d3d302},\\
&\texttt{f1121b62238ad5491bb7cf03635ea9934942edf573ed8faaa9ee79e1d38a6696}.
\end{split}
\]

SM V93 proves that a critical fixed-scale \(L^4\) estimate fails by a
physical logarithm and replaces it with scale-covariant evaluation of
one running theory.  For connected words meeting several stopping
cores, it assigns the local residue to the smallest node containing the
whole ordered word; charging its ancestors would manufacture
multiplicity.

The type-correct YM transfer is combinatorial, not renormalization
group theory.  A saturated coordinate inside a connected prefix carrier
must be assigned to the smallest recursively connected row block
containing its actual local letter.  The same occurrence cannot then be
charged again at the parent prefix and at the five-row root.  No SM
beta function, determinant, or continuum estimate is imported.

NS V31 remains unchanged.  Its exact mark-removal identity again warns
that a restored parent sum is not automatically subcritical.  Here the
analogous forbidden step would be to use the binary or ternary
minimal-owner card and then count the same local heat occurrence again
as one of the three independent slots.

\begin{audit}[V15 direction after both companion reads]
\label{audit:V15-direction}
Recursively decompose the \(3+2\) prefix before testing a protected
channel.  Assign every physical occurrence to its unique smallest
connected owner.  Only a response exceeding the ordinary heat response
of the coordinates already owned may count as the missing fourth
half-power.
\end{audit}

\section{The recursive first-connectivity tree}
\label{sec:V15-recursive-tree}

Start with a fixed nonzero term of the coefficient-one source
\eqref{eq:V11-exact-source}.  For every nonsingleton prefix block
\(B\), apply the same exact first-connectivity identity to the complete
connected carrier \(\mathcal K_B^{<r,\mathrm{conn}}\), using its own
ordered coordinate interval.  Repeat on every nonsingleton child block.

\begin{definition}[Recursive row-block family]
\label{def:V15-row-tree}
Let \(\mathfrak L\) contain the root \(I=[5]\), all blocks of the
root prefix partition \(\rho\), and recursively all blocks of the
proper predecessor partition used in each nonsingleton connected
prefix carrier.  Add every singleton leaf.  Order
\(\mathfrak L\) by inclusion.
\end{definition}

\begin{lemma}[Laminarity and finite depth]
\label{lem:V15-laminar}
Any two members of \(\mathfrak L\) are disjoint or one contains the
other.  Every strict descending chain has length at most four.
\end{lemma}

\begin{proof}
Children introduced at one recursion step form a partition of their
parent.  Thus two children of the same parent are disjoint.  Nodes
introduced in different branches remain in disjoint parent blocks,
while nodes on the same branch are nested.  This proves laminarity.
Along a strict descending chain, the positive integer cardinality
strictly decreases from at most five to one, so there are at most four
strict inclusions.
\end{proof}

\begin{definition}[Minimal connected owner]
\label{def:V15-minimal-owner}
For a nonempty set \(\mathcal S\) of row labels contained in at least
one node of \(\mathfrak L\), define
\[
 \operatorname{own}(\mathcal S)
 :=
 \text{the smallest }B\in\mathfrak L
 \text{ such that }\mathcal S\subseteq B.
\]
For an exact operator subword, \(\mathcal S\) is the union of the row
supports of all its literal local maps, including the endpoints of
every joining letter.
\end{definition}

\begin{lemma}[Existence and uniqueness of the minimal owner]
\label{lem:V15-owner-unique}
The owner in \Cref{def:V15-minimal-owner} exists and is unique whenever
the root is admitted.
\end{lemma}

\begin{proof}
The root contains every row support, so the containing-node set is
nonempty and finite.  Choose a member of minimal cardinality.  If two
different minimal members existed, laminarity in
\Cref{lem:V15-laminar} would make them disjoint or nested.  They cannot
be disjoint because both contain the nonempty set \(\mathcal S\), and
strict nesting contradicts minimality.  Hence the owner is unique.
\end{proof}

\begin{proposition}[Exact owner partition of operator occurrences]
\label{prop:V15-owner-partition}
Assign each literal local operator occurrence in the recursively
resolved source word to the minimal node containing its row support.
Then:
\begin{enumerate}[label=\textup{(\roman*)},leftmargin=3.5em]
\item every occurrence is assigned exactly once;
\item an occurrence joining two children is owned by their least common
      ancestor and by neither child;
\item summing owner-restricted occurrence sets over
      \(\mathfrak L\) restores the exact occurrence list without a
      coefficient or ancestor multiplicity.
\end{enumerate}
\end{proposition}

\begin{proof}
Existence and uniqueness are
\Cref{lem:V15-owner-unique}.  A joining occurrence has row support
meeting two disjoint children.  No descendant inside either child can
contain that support; the first ancestor containing both is its unique
owner.  Finally, unique assignment partitions a finite set, so the
disjoint union of all owner fibres is the original occurrence set.
\end{proof}

\begin{firewall}[An owner label still is not an operator]
\label{fire:V15-owner-label}
\Cref{prop:V15-owner-partition} assigns literal maps; it does not
insert a projection, counterterm, heat factor, or scalar response at
the owner node.  The operator content remains the exact source content,
as required by \Cref{lem:V10-role-erasure}.
\end{firewall}

\section{Application to the \texorpdfstring{\(3+2\)}{3+2} prefix}
\label{sec:V15-32}

Let the root predecessor partition be
\[
 \rho=\{B_3,B_2\},
 \qquad
 |B_3|=3,\quad |B_2|=2,
\]
and let \(B_a\) be the unique block containing the maximal row \(a\).
The exact prefix is
\begin{equation}
 \mathcal K_{B_3}^{<r,\mathrm{conn}}
 \otimes
 \mathcal K_{B_2}^{<r,\mathrm{conn}}.
\label{eq:V15-32-prefix}
\end{equation}

\begin{theorem}[Lower-order localization of every prefix occurrence]
\label{thm:V15-lower-order-localization}
Every literal operator occurrence below \(r\) which contains row \(a\)
has a unique minimal owner contained in \(B_a\).  That owner has arity
one, two, or three.  More precisely:
\begin{enumerate}[label=\textup{(\alph*)},leftmargin=3.5em]
\item if \(|B_a|=2\), its nontrivial connected owner is binary;
\item if \(|B_a|=3\), its nontrivial connected owner is ternary or a
      binary descendant; and
\item a map meeting both \(B_3\) and \(B_2\) is not a prefix map but is
      owned at the root joining carrier.
\end{enumerate}
\end{theorem}

\begin{proof}
The two root children are disjoint.  A prefix occurrence belongs to
one tensor factor in \eqref{eq:V15-32-prefix}; if its row support
contains \(a\), that factor is the one indexed by \(B_a\).  Recursive
children of \(B_a\) have strictly smaller cardinality.  Thus their
possible nontrivial sizes are only two when \(|B_a|=3\), and none below
two when \(|B_a|=2\).  This proves \textup{(a)} and \textup{(b)}.
A map meeting both root children has support not contained in either,
so \Cref{lem:V15-owner-unique} assigns it to their least common
ancestor, the root; in the first-connectivity formula this is part of
the quotient joining carrier.
\end{proof}

\begin{corollary}[No new high-arity estimate is hidden in the
\(3+2\) prefix]
\label{cor:V15-no-high-arity}
The complete analytic factors at every minimal prefix owner in type
\(3+2\) are one-row moments, complete covariances, or complete third
cumulants.  Their representation-uniform oriented square cards are
already established.  The remaining issue is whether one of those
cards supplies a surplus beyond the heat responses already assigned to
the same physical coordinates.
\end{corollary}

\begin{proof}
\Cref{thm:V15-lower-order-localization} bounds the owner arity by
three.  The archive summary in V1 records the complete binary and
ternary marked-tree cards.  Unique ownership
\Cref{prop:V15-owner-partition} prevents any additional factor from
being inferred merely from the recursive label.
\end{proof}

\section{The no-double-response rule}
\label{sec:V15-no-double-response}

Let \(D\subseteq E_3\) be the subset of saturated coordinates assigned
to one minimal connected owner \(\mathcal O\), and put \(k:=|D|\).
The ordinary positive heat response of these coordinates in an
oriented square is \(M^{-k/2}\).

\begin{definition}[Owner-surplus card \(OSC(\mathcal O,D)\)]
\label{def:V15-OSC}
The owner-surplus card is the pair of bounds
\begin{align}
 \kappa_\otimes(
  \mathbb L_{\mathcal O,D}^*
  \mathbb L_{\mathcal O,D})
 &\le C M^{-(k+1)/2},
\label{eq:V15-OSC-right}\\
 \kappa_\otimes(
  \mathbb L_{\mathcal O,D}
  \mathbb L_{\mathcal O,D}^*)
 &\le C M^{-(k+1)/2},
\label{eq:V15-OSC-left}
\end{align}
where \(\mathbb L_{\mathcal O,D}\) is the complete literal owner
packet on \(D\), with every coincident joining and payer role retained.
\end{definition}

\begin{lemma}[A surplus card supplies exactly the missing response]
\label{lem:V15-OSC-compiler}
Assume the other \(3-k\) saturated slots are distinct tensor
coordinates and have their ordinary \(M^{-(3-k)/2}\) heat-square
response.  If \(OSC(\mathcal O,D)\) holds and the complete separator
between owner and complementary factors satisfies the V11 dual bound,
then the assembled square has response \(CM^{-2}\).
\end{lemma}

\begin{proof}
The owner card gives \(M^{-(k+1)/2}\); the complementary slots give
\(M^{-(3-k)/2}\).  Their exponents add to
\[
 {k+1\over2}+{3-k\over2}=2.
\]
Tensorize genuinely distinct coordinates and use
\Cref{thm:V11-nonsymmetric-dual} for the complete bounded separator.
No multiplicity projection or second owner charge is introduced.
\end{proof}

\begin{proposition}[Ordinary owner cards cannot be multiplied by the
same heat response]
\label{prop:V15-no-double-response}
Suppose the proved square card for
\(\mathbb L_{\mathcal O,D}\) has only the ordinary scale
\(CM^{-k/2}\).  If its proof already retains every heat factor at
coordinates in \(D\), then those same coordinates cannot also
contribute a separate factor \(M^{-k/2}\) in the complete-word square.
\end{proposition}

\begin{proof}
The complete owner packet and its square already contain the literal
maps at every coordinate in \(D\).  By the exact owner partition
\Cref{prop:V15-owner-partition}, no second occurrence of those maps
exists at an ancestor.  Multiplying by a second heat response would
therefore replace one source occurrence by two analytically independent
copies, changing the source word.  This is precisely the forbidden
operation in \Cref{fire:V5-fixed-strong-open} and the scalar-role
erasure of \Cref{lem:V10-role-erasure}.
\end{proof}

\begin{corollary}[V13 is failure of the one-coordinate binary surplus]
\label{cor:V15-V13-as-OSC}
For the spin-one covariance channel of V13 with \(k=1\),
\[
 \kappa_\otimes(
  \mathbb L^*\mathbb L)
 \ge cM^{-1/2},
\]
whereas \(OSC\) would require \(CM^{-1}\).  Thus the binary
owner-surplus card fails on that channel, paid and unpaid.
\end{corollary}

\begin{proof}
This is
\Cref{thm:V13-covariance-sharp,thm:V13-paid-sharp}, compared with
\eqref{eq:V15-OSC-right} at \(k=1\).
\end{proof}

\section{The audited residual for the \texorpdfstring{\(3+2\)}{3+2}
source}
\label{sec:V15-32-residual}

\begin{theorem}[Exact minimal-owner alternative]
\label{thm:V15-32-alternative}
In a saturated \(q=3\), type-\(3+2\) source fibre, assign the prefix
occurrences meeting \(E_3\) to minimal owners.  Then at least one of:
\begin{enumerate}[label=\textup{(\roman*)},leftmargin=3.5em]
\item a one-row minimal owner contains a cofinal coordinate and V14
      supplies exponential decay;
\item a binary or ternary minimal owner satisfies its owner-surplus
      card and the fibre closes by
      \Cref{lem:V15-OSC-compiler};
\item every nontrivial minimal owner has only its ordinary already
      counted response; or
\item the saturated coordinate is not contained in a source-supported
      prefix owner and belongs to the root joining/suffix or the V5
      non-owner remainder.
\end{enumerate}
\end{theorem}

\begin{proof}
Unique minimal-owner assignment is
\Cref{prop:V15-owner-partition}.  The possible prefix owner arities are
one, two, and three by
\Cref{thm:V15-lower-order-localization}.  A cofinal one-row occurrence
is handled by \Cref{thm:V14-complete-barrier}.  At arity two or three,
either the strictly stronger bounds
\eqref{eq:V15-OSC-right}--\eqref{eq:V15-OSC-left} hold or they do not.
If they hold, apply \Cref{lem:V15-OSC-compiler}; if not, only the
ordinary card is available and
\Cref{prop:V15-no-double-response} forbids counting it twice.  Any
occurrence not assigned to a prefix node is, by the root source
decomposition, in the joining carrier, suffix, or outside the proved
source-supported owner set.  These alternatives are exhaustive.
\end{proof}

\begin{assessment}[Audit-adjusted analytic target]
\label{ass:V15-target}
The \(3+2\) problem is not a new quartic estimate and not a fine
recoupling estimate.  Its only possible same-support gain is a binary
or ternary \(OSC\).  V13 rules out the universal binary version.
V16 must test the ternary card on fixed low-output \(SU(2)\) channels
and determine whether the exact ternary cumulant cancels its protected
diagonal.  If it does not, the same-support route is exhausted for
\(3+2\) and the programme must return to source incidence.
\end{assessment}

\section{V15 result ledger and V16 acquisition order}
\label{sec:V15-ledger}

\begin{theorem}[Fourth-series V15 result ledger]
\label{thm:V15-results}
Relative to V1--V14, V15 proves:
\begin{enumerate}[label=\textup{(V15.\arabic*)},leftmargin=6.3em]
\item the compulsory hashed audit of SM V93 and NS V31;
\item construction of the finite laminar recursive
      first-connectivity row tree;
\item existence and uniqueness of the smallest connected owner;
\item exact one-use assignment of every literal operator occurrence;
\item localization of every \(3+2\) prefix occurrence meeting the
      maximal row to arity at most three;
\item exclusion of any hidden high-arity prefix estimate;
\item the owner-surplus card and its exact \(M^{-2}\) compiler;
\item a proof that the same heat occurrence cannot be counted once in
      its owner and again at an ancestor; and
\item the exhaustive minimal-owner alternative for the \(3+2\)
      saturated source.
\end{enumerate}
\end{theorem}

\begin{proof}
These are
\Cref{audit:V15-direction,
lem:V15-laminar,
lem:V15-owner-unique,
prop:V15-owner-partition,
thm:V15-lower-order-localization,
cor:V15-no-high-arity,
def:V15-OSC,
lem:V15-OSC-compiler,
prop:V15-no-double-response,
cor:V15-V13-as-OSC,
thm:V15-32-alternative}.
\end{proof}

\begin{programme}[V16: the ternary owner-surplus test]
\label{prog:V15-V16}
\begin{enumerate}[leftmargin=3em]
\item Write the complete third cumulant
      \(H_{123}-H_{12}H_3-H_{13}H_2-H_{23}H_1
       +2H_1H_2H_3\)
      on a fixed simultaneous \(SU(2)\) coupling path.
\item Choose fixed low total output and cofinal input spins compatible
      with an explicit row-product vector.
\item Compute the protected diagonal before taking absolute values.
\item Include the payer on the same nonzero low-output block so that
      the test covers every allocation.
\item If the diagonal is nondecaying, use V12 to disprove the ternary
      \(OSC\).  If it cancels, compute the protected-to-gapped blocks
      before claiming an upper bound.
\item Keep the root joining covariance separate from the ternary
      prefix owner; minimal-owner uniqueness forbids merging their
      scalar stocks.
\end{enumerate}
\end{programme}

\begin{firewall}[Claim boundary after fourth-series V15]
\label{fire:V15-claim-boundary}
V15 supplies the required companion audit and a nonduplicating
recursive owner calculus.  It does not prove the binary or ternary
owner-surplus card, source occurrence for the V5 remainder, or the full
\(q=3\) estimate.  The remaining \(q=3\) alternatives,
\(q=0,1,2\) cards, global quintic or all-order MTA, WUFC, CPM, cutoff
removal, reflection positivity, Osterwalder--Schrader reconstruction,
construction of the continuum Yang--Mills measure, and a uniform
positive continuum spectral gap remain open.  The full Yang--Mills
mass gap has not been proved.
\end{firewall}

\paragraph{Parent-version record.}
V15 was formed from
\texttt{Renewal\_Yang\_Mills\_Mass\_Gap\_Research\_Notes\_V14.tex}.
The V14 parent had \(7358\) logical source lines, \(261\,208\) bytes,
and SHA--256
\[
 \texttt{2985c1e0b699ec6ac91bd16a861e166afeaa49436674da77b98d36e5fd3987a1}.
\]
The next compulsory companion audit is V20.

% =============================================================================
% END APPENDED FOURTH-SERIES V15 MATERIAL
% =============================================================================

% =============================================================================
% BEGIN APPENDED FOURTH-SERIES V16 MATERIAL
% =============================================================================

\clearpage
\chapter{V16: A complete ternary-cumulant obstruction to the
one-coordinate owner surplus}
\label{ch:V16}

\section{Purpose}
\label{sec:V16-purpose}

V15 reduces every \(3+2\) prefix occurrence meeting the maximal row to
an owner of arity at most three.  V13 disproves the universal binary
one-coordinate owner-surplus card.  This chapter performs the ternary
test without replacing the complete cumulant by one of its output
summands.

Take single-row spins \(j,j,1\).  On total output spin one, the terms
containing an independently heated spin-\(j\) row are exponentially
small.  The complete third cumulant therefore approaches
\[
 q_1(I-H_{12}),
\]
which has a fixed gap on the intermediate \(J_{12}=1\) sector.  An
explicit product vector has probability comparable to \(d_j^{-1}\) in
that sector.  Orthogonal projection is used only inside the proof to
lower-bound the norm of the complete cumulant applied to that product
vector; it is not inserted into the source word.

\section{The complete third cumulant on
\texorpdfstring{\(V_j\otimes V_j\otimes V_1\)}{Vj tensor Vj tensor V1}}
\label{sec:V16-ternary}

Continue to write
\[
 q_L=e^{-sL(L+1)}.
\]
For a nonempty row set \(D\subseteq\{1,2,3\}\), let
\(\mathbb H_D(s)\) be the complete common-variable heat moment on those
rows, tensor-extended by the identity on the complement.  Put
\[
 \mathbb H_i(s)=q_{\lambda_i}I
\]
on a single row of spin \(\lambda_i\).

\begin{definition}[Complete local third cumulant]
\label{def:V16-K3}
On row spins \((j,j,1)\), set
\begin{equation}
 \begin{split}
 \mathbb K_{3,j}(s)
 :=\;&
 \mathbb H_{123}
 -\mathbb H_{12}\mathbb H_3
 -\mathbb H_{13}\mathbb H_2
 -\mathbb H_{23}\mathbb H_1\\
 &+2\mathbb H_1\mathbb H_2\mathbb H_3.
 \end{split}
\label{eq:V16-K3}
\end{equation}
\end{definition}

Every term is self-adjoint: factors in each product act on disjoint row
blocks and commute.  Every term also intertwines the total diagonal
\(SU(2)\)-action.

Let \(P^{\mathrm{tot}}_1\) be the complete total-spin-one projection.
Inside its range, couple the first two rows first.  The possible
intermediate spins are \(L=0,1,2\).  Let \(Q_L\) be the corresponding
orthogonal projection, including the magnetic factor, so that
\[
 P^{\mathrm{tot}}_1=Q_0+Q_1+Q_2.
\]

\begin{lemma}[Exact total-spin-one block]
\label{lem:V16-K3-block}
On the total-spin-one block,
\begin{equation}
 \boxed{
 P^{\mathrm{tot}}_1\mathbb K_{3,j}P^{\mathrm{tot}}_1
 =
 q_1(I-\mathbb H_{12})
 -
 q_j(\mathbb H_{13}+\mathbb H_{23})
 +
 2q_j^2q_1 I.}
\label{eq:V16-K3-block}
\end{equation}
Here \(\mathbb H_{12}\) is diagonal in the \(Q_L\) resolution:
\begin{equation}
 \mathbb H_{12}Q_L=q_LQ_L,
 \qquad L=0,1,2.
\label{eq:V16-H12-diagonal}
\end{equation}
\end{lemma}

\begin{proof}
On total output spin one,
\(\mathbb H_{123}=q_1I\).  The third row has spin one, hence
\(\mathbb H_3=q_1I\); the first and second rows have spin \(j\), hence
\(\mathbb H_1=\mathbb H_2=q_jI\).  Substitute these scalar identities
into \eqref{eq:V16-K3}.  The heat moment of rows \(1,2\) is the
functional calculus of their intermediate Casimir, which proves
\eqref{eq:V16-H12-diagonal}.
\end{proof}

\begin{lemma}[Uniform gap away from the \(L=0\) channel]
\label{lem:V16-gap}
Put
\[
 Q:=Q_1+Q_2,
 \qquad
 \delta_s:=q_1(1-q_1)>0.
\]
Then
\begin{equation}
 q_1(I-\mathbb H_{12})
 \succeq
 \delta_s Q
\label{eq:V16-leading-gap}
\end{equation}
on the total-spin-one block.  Moreover, the remainder
\begin{equation}
 \mathbb E_j
 :=
 -q_j(\mathbb H_{13}+\mathbb H_{23})
 +2q_j^2q_1I
\label{eq:V16-remainder}
\end{equation}
satisfies
\begin{equation}
 \|\mathbb E_j\|\le2q_j+2q_j^2.
\label{eq:V16-remainder-bound}
\end{equation}
\end{lemma}

\begin{proof}
On \(Q_0\), the leading operator has eigenvalue
\(q_1(1-q_0)=0\).  On \(Q_1\), its eigenvalue is
\(\delta_s\).  On \(Q_2\), the eigenvalue is
\(q_1(1-q_2)\ge q_1(1-q_1)=\delta_s\), because
\(0<q_2<q_1<1\).  This proves
\eqref{eq:V16-leading-gap}.  Every complete heat moment is a positive
contraction, so the triangle inequality gives
\eqref{eq:V16-remainder-bound}.
\end{proof}

\section{A product vector sees the gapped intermediate channel}
\label{sec:V16-product}

Let
\begin{equation}
 \Phi_j
 :=
 |j,j\rangle\otimes|j,-j\rangle\otimes|1,1\rangle.
\label{eq:V16-product-vector}
\end{equation}
This is a product vector across the three physical rows.  Put
\[
 v_j:=P^{\mathrm{tot}}_1\Phi_j.
\]

\begin{lemma}[Quantitative \(Q_1\) component]
\label{lem:V16-Q1-component}
For every \(j\ge1\),
\begin{equation}
 \|Q_1v_j\|^2
 =
 {1\over2}p_{j,1}
 =
 {3j\over(2j+1)(2j+2)}
 \ge {c\over d_j}.
\label{eq:V16-Q1-weight}
\end{equation}
\end{lemma}

\begin{proof}
The squared projection of the first two tensor factors of
\(\Phi_j\) onto intermediate spin one and magnetic number zero is
\(p_{j,1}\), by \eqref{eq:V13-pj1}.  Coupling that vector with
\(|1,1\rangle\) to total spin one and total magnetic number one
contributes
\[
 |\langle1,0;1,1\mid1,1\rangle|^2={1\over2}.
\]
Multiplication proves the equality.  The lower bound follows from
\Cref{cor:V13-pj1-scale}.
\end{proof}

\begin{theorem}[Complete ternary square has only the ordinary
half-power]
\label{thm:V16-ternary-lower}
There are \(j_s<\infty\) and \(c_s>0\) such that
\begin{equation}
 \boxed{
 \kappa_\otimes(\mathbb K_{3,j}(s)^2)
 \ge {c_s\over d_j}
 \asymp c_sM_j^{-1/2}}
\label{eq:V16-ternary-lower}
\end{equation}
for every \(j\ge j_s\).
\end{theorem}

\begin{proof}
Because the complete cumulant commutes with total spin,
\[
 P^{\mathrm{tot}}_1\mathbb K_{3,j}\Phi_j
 =
 \bigl[
  q_1(I-\mathbb H_{12})+\mathbb E_j
 \bigr]v_j.
\]
By \Cref{lem:V16-gap} and orthogonality of \(Q_0,Q_1,Q_2\),
\[
 \|q_1(I-\mathbb H_{12})v_j\|
 \ge\delta_s\|Qv_j\|
 \ge\delta_s\|Q_1v_j\|.
\]
Therefore
\begin{align*}
 \|\mathbb K_{3,j}\Phi_j\|
 &\ge
 \|P^{\mathrm{tot}}_1\mathbb K_{3,j}\Phi_j\|\\
 &\ge
 \delta_s\|Q_1v_j\|
 -(2q_j+2q_j^2)\|v_j\|.
\end{align*}
Now \(\|v_j\|\le1\), while
\(\|Q_1v_j\|\ge c d_j^{-1/2}\) by
\Cref{lem:V16-Q1-component}.  The number
\[
 q_j=e^{-sj(j+1)}
\]
decays faster than \(d_j^{-1/2}\).  Hence, after increasing \(j_s\),
\[
 2q_j+2q_j^2
 \le{\delta_s\over2}\|Q_1v_j\|.
\]
It follows that
\[
 \|\mathbb K_{3,j}\Phi_j\|^2
 \ge{\delta_s^2\over4}\|Q_1v_j\|^2
 \ge {c_s\over d_j}.
\]
The left side is
\(\langle\Phi_j,\mathbb K_{3,j}^2\Phi_j\rangle\), an admissible
product-state expectation.  This proves
\eqref{eq:V16-ternary-lower}.
\end{proof}

\begin{firewall}[No fine projection was inserted]
\label{fire:V16-no-pinching}
The tested operator in
\Cref{thm:V16-ternary-lower} is the complete
\(\mathbb K_{3,j}\), not
\(Q_1\mathbb K_{3,j}Q_1\).  The orthogonal projections are used only
to discard components of the norm
\(\|\mathbb K_{3,j}\Phi_j\|\) when proving a lower bound.  Thus the
argument is compatible with the complete-output rule of V11.
\end{firewall}

\section{The payer-safe form}
\label{sec:V16-paid}

\begin{theorem}[The total-output payer preserves the ternary
obstruction]
\label{thm:V16-paid-lower}
Let \(\mathbb R_s\) be the unique selected payer.  There are
\(j_s<\infty\) and \(c_s^{\rm pay}>0\) such that
\begin{equation}
 \kappa_\otimes\!\left(
  (\mathbb R_s\mathbb K_{3,j})^*
  (\mathbb R_s\mathbb K_{3,j})
 \right)
 \ge {c_s^{\rm pay}\over d_j}
\label{eq:V16-paid-lower}
\end{equation}
for all \(j\ge j_s\).  The opposite square orientation obeys the same
lower bound.
\end{theorem}

\begin{proof}
The payer and the cumulant commute because both are equivariant
functions or intertwiners for total output.  On the total-spin-one
block,
\[
 \mathbb R_s=r_1(s)I,
 \qquad
 r_1(s)={2\over1-e^{-2s}}>0.
\]
Repeat the proof of \Cref{thm:V16-ternary-lower} after applying
\(P^{\mathrm{tot}}_1\).  The retained norm is multiplied by
\(r_1(s)\), giving the displayed bound.  Both operators are
self-adjoint and commute, so the two square orientations coincide.
\end{proof}

\begin{corollary}[The one-coordinate ternary \(OSC\) fails]
\label{cor:V16-ternary-OSC-fails}
For \(k=1\), the owner-surplus card
\eqref{eq:V15-OSC-right}--\eqref{eq:V15-OSC-left} would require
\(CM_j^{-1}\).  The complete third cumulant on the family above has
square capacity at least \(c_sM_j^{-1/2}\), paid and unpaid.  Hence no
representation-uniform one-coordinate ternary \(OSC\) follows from the
complete local heat-cumulant structure alone.
\end{corollary}

\begin{proof}
Compare
\Cref{thm:V16-ternary-lower,thm:V16-paid-lower} with
\Cref{def:V15-OSC} at \(k=1\).
\end{proof}

\section{Embedding in a three-slot analytic counterpacket}
\label{sec:V16-three-slot}

At two further distinct coordinates take the fixed spin-one resolved
heat operators
\(\mathbb H_{j,1}=q_1P_{j,1}\) of V13.  Define
\[
 \mathbb Z_{3,j}
 :=
 \mathbb K_{3,j}
 \otimes\mathbb H_{j,1}
 \otimes\mathbb H_{j,1}.
\]

\begin{proposition}[Ternary centering plus three occupied coordinates
remains short]
\label{prop:V16-three-slot-short}
For all sufficiently large \(j\),
\begin{equation}
 \kappa_\otimes(\mathbb Z_{3,j}^*\mathbb Z_{3,j})
 \ge c_s d_j^{-3}
 \asymp c_sM_j^{-3/2}.
\label{eq:V16-three-slot-short}
\end{equation}
The same is true when the ternary factor pays.
\end{proposition}

\begin{proof}
Use the product vector \(\Phi_j\) at the ternary coordinate and
\(\phi_j\) from V13 at each of the other two coordinates.  The first
expectation is at least \(c_sd_j^{-1}\) by
\Cref{thm:V16-ternary-lower}.  Each fixed spin-one heat square has
expectation
\[
 q_1^2p_{j,1}\ge c_sd_j^{-1}.
\]
Tensorization gives \eqref{eq:V16-three-slot-short}.  The paid claim
uses \Cref{thm:V16-paid-lower}.
\end{proof}

\begin{firewall}[Analytic family versus a complete prefix word]
\label{fire:V16-prefix-scope}
\Cref{prop:V16-three-slot-short} is an exact local analytic
counterpacket.  A complete binary--ternary \(3+2\) prefix contains
additional earlier and later coordinate factors.  Those factors may
annihilate the displayed product component or provide exponential
decay.  V16 therefore disproves a universal ternary surplus lemma; it
does not yet construct a full nonzero quintic source counterexample.
\end{firewall}

\section{Consequence for the minimal-owner alternative}
\label{sec:V16-consequence}

\begin{theorem}[Same-support local owner cards are exhausted at arity
at most three]
\label{thm:V16-local-exhaustion}
For a \(3+2\) prefix occurrence assigned to a one-coordinate minimal
owner:
\begin{enumerate}[label=\textup{(\roman*)},leftmargin=3.5em]
\item a one-row owner is exponentially small only under the cofinal
      scalar-heat hypothesis of V14;
\item the universal binary owner-surplus card fails by V13; and
\item the universal ternary owner-surplus card fails by V16.
\end{enumerate}
Thus a remaining \(3+2\) closure cannot be deduced from owner arity and
one-coordinate cofinality alone.  It must use surrounding source
incidence, a multi-coordinate owner-surplus theorem, or annihilation of
the explicit low-output channels.
\end{theorem}

\begin{proof}
Item \textup{(i)} is
\Cref{thm:V14-complete-barrier} together with its stated hypotheses.
Item \textup{(ii)} is \Cref{cor:V15-V13-as-OSC}; item \textup{(iii)}
is \Cref{cor:V16-ternary-OSC-fails}.  The list of possible minimal
prefix-owner arities is exhaustive by
\Cref{thm:V15-lower-order-localization}.  Therefore any further gain
must use data not present in the one-coordinate local cards.
\end{proof}

\begin{assessment}[Upstream move forced by the countertest]
\label{ass:V16-upstream}
The \(3+2\) programme should no longer seek a stronger local covariance
or third-cumulant bound.  The next object is the complete
multi-coordinate binary or ternary prefix carrier surrounding the
minimal owner.  One must determine whether first-connectivity forces a
second cofinal source occurrence in that carrier.  If not, the V13/V16
low-output components must be propagated through the exact prefix to
test whether a full source counterpacket exists.
\end{assessment}

\section{V16 result ledger and V17 acquisition order}
\label{sec:V16-ledger}

\begin{theorem}[Fourth-series V16 result ledger]
\label{thm:V16-results}
Relative to V1--V15, V16 proves:
\begin{enumerate}[label=\textup{(V16.\arabic*)},leftmargin=6.3em]
\item the exact total-spin-one block formula for the complete third
      cumulant on spins \(j,j,1\);
\item a uniform leading gap on intermediate spins one and two;
\item an exponentially small remainder estimate;
\item an explicit row-product vector with
      \(c\,d_j^{-1}\) weight in the gapped channel;
\item a complete, unpinched ternary square lower bound of order
      \(M_j^{-1/2}\);
\item the same lower bound with the total-output payer;
\item failure of the universal one-coordinate ternary owner-surplus
      card;
\item a three-occupied-coordinate \(M_j^{-3/2}\) analytic
      counterpacket; and
\item exhaustion of all one-coordinate minimal-owner surplus routes
      at arity at most three.
\end{enumerate}
\end{theorem}

\begin{proof}
These are
\Cref{lem:V16-K3-block,
lem:V16-gap,
lem:V16-Q1-component,
thm:V16-ternary-lower,
fire:V16-no-pinching,
thm:V16-paid-lower,
cor:V16-ternary-OSC-fails,
prop:V16-three-slot-short,
thm:V16-local-exhaustion}.
\end{proof}

\begin{programme}[V17: multi-coordinate prefix incidence]
\label{prog:V16-V17}
\begin{enumerate}[leftmargin=3em]
\item Apply exact first connectivity inside the binary and ternary
      \(3+2\) prefix blocks and retain their complete coordinate words.
\item For the block containing the maximal row, partition its row mass
      among the internal prefix, internal joining, and internal suffix.
\item If two distinct cofinal source coordinates occur inside one
      minimal owner, formulate the \(k=2\) owner-surplus card and test it
      before using any scalar overlap.
\item If only one cofinal coordinate occurs, propagate the V13 or V16
      low-output vector through the remaining bounded prefix factors.
\item Treat exact annihilation as a selection rule to be proved, not as
      a consequence of connectedness.
\item Keep all ancestor owner sets uncharged, as required by V15.
\end{enumerate}
\end{programme}

\begin{firewall}[Claim boundary after fourth-series V16]
\label{fire:V16-claim-boundary}
V16 rigorously disproves a universal one-coordinate ternary surplus,
paid and unpaid.  It does not determine survival through the full
multi-coordinate prefix, prove a two-coordinate surplus, produce a
full quintic counterexample, or close \(q=3\).  The remaining \(q=3\)
alternatives, \(q=0,1,2\) cards, global quintic or all-order MTA,
WUFC, CPM, cutoff removal, reflection positivity,
Osterwalder--Schrader reconstruction, construction of the continuum
Yang--Mills measure, and a uniform positive continuum spectral gap
remain open.  The full Yang--Mills mass gap has not been proved.
\end{firewall}

\paragraph{Parent-version record.}
V16 was formed from
\texttt{Renewal\_Yang\_Mills\_Mass\_Gap\_Research\_Notes\_V15.tex}.
The V15 parent had \(7800\) logical source lines, \(278\,923\) bytes,
and SHA--256
\[
 \texttt{5974aa6844f0cb4eef7f13c4e1aac8047ad9ea31daca5ebaa09e2b963c8e2e44}.
\]
The next compulsory companion audit is V20.

% =============================================================================
% END APPENDED FOURTH-SERIES V16 MATERIAL
% =============================================================================

% =============================================================================
% BEGIN APPENDED FOURTH-SERIES V17 MATERIAL
% =============================================================================

\clearpage
\chapter{V17: Tensor-product obstruction and the necessity of assembled
prefix cancellation}
\label{ch:V17}

\section{Purpose}
\label{sec:V17-purpose}

V13 and V16 give complete paid and unpaid low-output packets whose
square capacities are bounded below by \(c\,d_j^{-1}\), or
equivalently \(cM_j^{-1/2}\).  A \(3+2\) prefix can contain two or
three saturated physical coordinates inside the same recursively
connected owner.  Merely grouping those coordinates into one owner does
not create a surplus: a fixed source word is a tensor product over
physical coordinates.

This chapter proves the resulting \(k\)-coordinate lower bound for
\(k=1,2,3\).  It follows that no termwise positive-envelope proof can
establish the owner-surplus card.  If the complete prefix has extra
decay, it must arise from cancellation among different exact coordinate
words after their signed sum is assembled.

\section{A general tensor-product lower bound}
\label{sec:V17-tensor-lower}

\begin{lemma}[Product tests tensorize lower bounds]
\label{lem:V17-product-lower}
For \(\ell=1,\ldots,k\), let \(\mathbb L_{\ell}\) act at distinct
physical coordinates and suppose a row-product unit vector
\(\phi_\ell\) satisfies
\begin{equation}
 \langle\phi_\ell,
  \mathbb L_\ell^*\mathbb L_\ell
  \phi_\ell\rangle
 \ge b_\ell.
\label{eq:V17-local-lower}
\end{equation}
Then
\begin{equation}
 \kappa_\otimes\!\left[
  \left(\bigotimes_{\ell=1}^k\mathbb L_\ell\right)^*
  \left(\bigotimes_{\ell=1}^k\mathbb L_\ell\right)
 \right]
 \ge\prod_{\ell=1}^kb_\ell.
\label{eq:V17-tensor-lower}
\end{equation}
The analogous statement holds for the opposite square orientation.
\end{lemma}

\begin{proof}
The tensor product \(\phi:=\bigotimes_\ell\phi_\ell\) remains a
product vector across the global row cut: regrouping its factors by row
does not change that fact.  The expectation of the tensor-product square
factorizes into the product of
\eqref{eq:V17-local-lower}.  Product-state capacity is the supremum over
all admissible product tests, proving
\eqref{eq:V17-tensor-lower}.  Apply the same proof to
\(\mathbb L_\ell^*\) for the other orientation.
\end{proof}

\begin{theorem}[No tensor-product owner surplus]
\label{thm:V17-no-tensor-surplus}
Let \(1\le k\le3\).  At each of \(k\) distinct coordinates choose
either the complete binary packet of V13 or the complete ternary packet
of V16, using the same large spin \(j\), and allow the sole payer on any
one chosen nonzero spin-one output.  Then there are row-product vectors
for which the complete tensor-product square has expectation at least
\begin{equation}
 c_{s,k}d_j^{-k}
 \asymp c_{s,k}M_j^{-k/2}.
\label{eq:V17-k-lower}
\end{equation}
Consequently a representation-uniform owner-surplus estimate
\begin{equation}
 C M_j^{-(k+1)/2}
\label{eq:V17-OSC-target}
\end{equation}
is false for this family.
\end{theorem}

\begin{proof}
For a binary coordinate use
\Cref{thm:V13-covariance-sharp,thm:V13-paid-sharp}; for a ternary
coordinate use
\Cref{thm:V16-ternary-lower,thm:V16-paid-lower}.  Each supplies a
product vector and a square expectation at least \(c_sd_j^{-1}\).
The payer multiplier on spin one is a fixed positive number depending
only on \(s\), so it changes only the constant.  Apply
\Cref{lem:V17-product-lower}.

Since \(M_j^{1/2}\asymp d_j\), the ratio of
\eqref{eq:V17-k-lower} to
\eqref{eq:V17-OSC-target} is comparable to
\(M_j^{1/2}\), which diverges.  Thus no uniform surplus bound can hold.
\end{proof}

\begin{corollary}[More occupied coordinates do not solve the missing
half-power]
\label{cor:V17-more-not-surplus}
Moving from one to two or three saturated coordinates inside the same
minimal owner supplies exactly the corresponding two or three ordinary
heat half-powers.  It never supplies the additional fourth half-power
required by \(OSC\) solely through tensorization.
\end{corollary}

\begin{proof}
At \(k\) coordinates the lower bound is the ordinary scale
\(M^{-k/2}\), whereas \(OSC\) asks for
\(M^{-(k+1)/2}\).  Apply
\Cref{thm:V17-no-tensor-surplus}.
\end{proof}

\begin{firewall}[The theorem concerns fixed coordinate words]
\label{fire:V17-fixed-word}
A complete connected prefix carrier is a signed sum of fixed coordinate
words.  \Cref{thm:V17-no-tensor-surplus} applies to each tensor-product
word carrying the displayed low-output packets.  It does not preclude
cancellation after all words in the connected carrier are summed.
\end{firewall}

\section{Why a positive envelope cannot see the needed cancellation}
\label{sec:V17-positive-envelope}

Let a fixed recursively resolved connected prefix have the exact finite
or coordinate-summable expansion
\begin{equation}
 \mathbb K_B
 =
 \sum_{\omega\in\Omega_B}c_\omega\mathbb W_\omega,
\label{eq:V17-prefix-sum}
\end{equation}
where \(c_\omega\) includes the local Möbius signs and every
\(\mathbb W_\omega\) is a literal coordinate tensor product.

\begin{lemma}[Exact square and its cross terms]
\label{lem:V17-exact-square}
\begin{equation}
 \mathbb K_B^*\mathbb K_B
 =
 \sum_{\omega,\nu\in\Omega_B}
 \overline{c_\omega}c_\nu
 \mathbb W_\omega^*\mathbb W_\nu.
\label{eq:V17-exact-square}
\end{equation}
The diagonal positive terms
\(|c_\omega|^2\mathbb W_\omega^*\mathbb W_\omega\)
cannot cancel one another.  Any reduction below all termwise square
scales must use the off-diagonal signed terms before a positive
majorant is imposed.
\end{lemma}

\begin{proof}
Expand the product of the two sums in
\eqref{eq:V17-prefix-sum}.  The diagonal coefficients are nonnegative
moduli squared.  Therefore cancellation of their contribution to a
quadratic form can occur only through the real parts of off-diagonal
terms.
\end{proof}

\begin{proposition}[Cauchy square envelopes discard the only possible
surplus mechanism]
\label{prop:V17-envelope-loss}
For finite \(\Omega_B\),
\begin{equation}
 \mathbb K_B^*\mathbb K_B
 \preccurlyeq
 \left(\sum_{\omega}|c_\omega|\right)
 \left(
  \sum_{\omega}|c_\omega|
  \mathbb W_\omega^*\mathbb W_\omega
 \right).
\label{eq:V17-Cauchy-envelope}
\end{equation}
If the only available termwise estimates on the right have scale
\(M^{-k/2}\), then
\eqref{eq:V17-Cauchy-envelope} proves at most that same scale.  It
cannot prove \(M^{-(k+1)/2}\) with fixed coefficients.
\end{proposition}

\begin{proof}
For every vector \(v\),
\begin{align*}
 \|\mathbb K_Bv\|^2
 &=
 \left\|
  \sum_\omega c_\omega\mathbb W_\omega v
 \right\|^2\\
 &\le
 \left(\sum_\omega|c_\omega|\right)
 \left(\sum_\omega|c_\omega|
       \|\mathbb W_\omega v\|^2\right)
\end{align*}
by weighted Cauchy--Schwarz.  This is
\eqref{eq:V17-Cauchy-envelope}.  A fixed sum of
\(CM^{-k/2}\) bounds retains the exponent \(k/2\), proving the last
assertion.
\end{proof}

\begin{example}[Exact cancellation invisible termwise]
\label{ex:V17-cancellation}
Let \(P\ne0\) be any projection and set
\(\mathbb W_1=P\), \(\mathbb W_2=-P\).  Then
\[
 \mathbb W_1+\mathbb W_2=0,
\]
but each individual square equals \(P\), and the positive envelope is
\(4P\).  This elementary model is not a YM source formula; it shows why
an upper bound formed after absolute values cannot certify a
cancellation gain.
\end{example}

\section{The assembled owner-surplus card}
\label{sec:V17-assembled-card}

\begin{definition}[Assembled surplus \(AOSC(B,D)\)]
\label{def:V17-AOSC}
Let \(D\subseteq E_3\) be the \(k\)-coordinate set assigned to a
minimal connected prefix owner \(B\).  Retain the complete signed sum
\(\mathbb K_B\) in \eqref{eq:V17-prefix-sum}, all its local moment
partitions, and the single payer if allocated there.  The assembled
owner-surplus card is
\begin{align}
 \kappa_\otimes(\mathbb K_B^*\mathbb K_B)
 &\le C M^{-(k+1)/2},
\label{eq:V17-AOSC-right}\\
 \kappa_\otimes(\mathbb K_B\mathbb K_B^*)
 &\le C M^{-(k+1)/2}.
\label{eq:V17-AOSC-left}
\end{align}
\end{definition}

\begin{theorem}[Only an assembled card can close the same-owner
branch]
\label{thm:V17-only-assembled}
Suppose no new cofinal source coordinate outside \(D\) occurs in the
minimal owner.  A proof using only:
\begin{enumerate}[label=\textup{(\roman*)},leftmargin=3.5em]
\item separate square bounds for the fixed coordinate words;
\item ordinary norms of the separators; and
\item positive Cauchy envelopes
\end{enumerate}
cannot establish \(AOSC(B,D)\) uniformly for all complete binary and
ternary low-output packets.  A successful same-owner proof must retain
and estimate off-diagonal cross terms in
\eqref{eq:V17-exact-square}, or use an additional source-specific
identity implying their cancellation.
\end{theorem}

\begin{proof}
\Cref{thm:V17-no-tensor-surplus} gives fixed coordinate words whose
squares have the ordinary scale \(M^{-k/2}\) and violate the desired
surplus scale.  Ordinary separator norms do not change the exponent,
and \Cref{prop:V17-envelope-loss} retains the same termwise exponent.
The only part of the exact square discarded by this method is the
signed off-diagonal sum.  Therefore a proof of the stronger exponent
must use that sum or an equivalent exact source identity.
\end{proof}

\begin{corollary}[Correct upstream object for \(3+2\)]
\label{cor:V17-upstream-object}
In the \(3+2\) prefix, the next analytic object is the complete
binary or ternary connected prefix carrier after its internal
first-connectivity sum, not a stronger local covariance or
third-cumulant card and not a product of several such local cards.
\end{corollary}

\begin{proof}
The possible minimal-owner arities are at most three by
\Cref{thm:V15-lower-order-localization}.  Their termwise local surplus
cards fail by V13, V16, and
\Cref{thm:V17-no-tensor-surplus}.  Apply
\Cref{thm:V17-only-assembled}.
\end{proof}

\section{A finite protected-vector reduction for the assembled sum}
\label{sec:V17-protected-reduction}

\begin{definition}[Bad-channel scalar word]
\label{def:V17-bad-scalar}
For each coordinate word \(\mathbb W_\omega\), choose the V13/V16
row-product test vector \(\Phi\) on its occupied coordinates and fixed
unit product vectors on its spectators.  Define
\begin{equation}
 z_\omega(j)
 :=
 \langle\Phi,\mathbb W_\omega\Phi\rangle.
\label{eq:V17-zomega}
\end{equation}
\end{definition}

\begin{lemma}[Necessary scalar cancellation on a common product test]
\label{lem:V17-scalar-cancellation}
For the assembled carrier,
\begin{equation}
 \langle\Phi,\mathbb K_B\Phi\rangle
 =
 \sum_{\omega\in\Omega_B}c_\omega z_\omega(j).
\label{eq:V17-scalar-sum}
\end{equation}
If along a family the right side has magnitude
\(\ge cM^{-k/4}\), then
\begin{equation}
 \kappa_\otimes(\mathbb K_B^*\mathbb K_B)
 \ge c^2M^{-k/2},
\label{eq:V17-scalar-lower}
\end{equation}
and the assembled surplus card fails.
\end{lemma}

\begin{proof}
The identity follows from linearity and
\eqref{eq:V17-prefix-sum}.  Cauchy--Schwarz gives
\[
 |\langle\Phi,\mathbb K_B\Phi\rangle|^2
 \le
 \langle\Phi,\mathbb K_B^*\mathbb K_B\Phi\rangle
\]
for a unit vector \(\Phi\).  Since \(\Phi\) is a row-product test,
the latter is bounded above by product-state capacity, proving
\eqref{eq:V17-scalar-lower}.  This scale is larger than
\(M^{-(k+1)/2}\) by \(M^{1/2}\).
\end{proof}

\begin{assessment}[Concrete V18 calculation]
\label{ass:V17-V18}
The first V18 task is now a scalar calculation with exact ancestry.
For a binary or ternary prefix block, evaluate
\eqref{eq:V17-zomega} on the fixed low-output product family for every
internal first-connectivity word, preserving its Möbius sign.  If the
sum in \eqref{eq:V17-scalar-sum} stays at ordinary scale, the assembled
same-owner route fails.  If its leading term cancels, the next task is
the full vector norm, because scalar cancellation alone remains only a
necessary condition.
\end{assessment}

\section{V17 result ledger and V18 acquisition order}
\label{sec:V17-ledger}

\begin{theorem}[Fourth-series V17 result ledger]
\label{thm:V17-results}
Relative to V1--V16, V17 proves:
\begin{enumerate}[label=\textup{(V17.\arabic*)},leftmargin=6.3em]
\item tensorization of row-product lower bounds across physical
      coordinates;
\item failure of the owner-surplus exponent for every tensor product
      of one to three V13/V16 low-output packets;
\item proof that grouping more occupied coordinates does not create
      the missing half-power;
\item the exact cross-term expansion of a connected prefix square;
\item identification of the exponent loss in every positive Cauchy
      envelope;
\item the assembled owner-surplus card;
\item proof that only assembled signed cancellation can establish that
      card without a new source coordinate; and
\item a scalar protected-vector test whose failure disproves the
      assembled card.
\end{enumerate}
\end{theorem}

\begin{proof}
These are
\Cref{lem:V17-product-lower,
thm:V17-no-tensor-surplus,
cor:V17-more-not-surplus,
lem:V17-exact-square,
prop:V17-envelope-loss,
def:V17-AOSC,
thm:V17-only-assembled,
cor:V17-upstream-object,
lem:V17-scalar-cancellation}.
\end{proof}

\begin{programme}[V18: exact scalar ancestry on the bad channel]
\label{prog:V17-V18}
\begin{enumerate}[leftmargin=3em]
\item Begin with a binary connected prefix carrier and write its exact
      internal first-connectivity sum without absolute values.
\item Evaluate every complete moment and covariance letter on the V13
      spin-one product vector, retaining the physical coordinate order.
\item Sum the scalar word coefficients before estimating their
      magnitude.
\item Repeat for the ternary carrier using the V16 product vector.
\item If leading scalar cancellation occurs, compute
      \(\|\mathbb K_B\Phi\|\), not merely its diagonal expectation.
\item If it does not, record a full assembled-prefix obstruction and
      move upstream to source incidence for the non-owner remainder.
\end{enumerate}
\end{programme}

\begin{firewall}[Claim boundary after fourth-series V17]
\label{fire:V17-claim-boundary}
V17 proves that no termwise or tensor-product lower-order estimate can
supply the owner surplus.  It does not evaluate the exact assembled
prefix cancellation, prove survival through the full quintic word, or
close \(q=3\).  The remaining \(q=3\) alternatives,
\(q=0,1,2\) cards, global quintic or all-order MTA, WUFC, CPM, cutoff
removal, reflection positivity, Osterwalder--Schrader reconstruction,
construction of the continuum Yang--Mills measure, and a uniform
positive continuum spectral gap remain open.  The full Yang--Mills
mass gap has not been proved.
\end{firewall}

\paragraph{Parent-version record.}
V17 was formed from
\texttt{Renewal\_Yang\_Mills\_Mass\_Gap\_Research\_Notes\_V16.tex}.
The V16 parent had \(8287\) logical source lines, \(294\,414\) bytes,
and SHA--256
\[
 \texttt{749fafd13c045fd556a8388c001c5b9b1cf476cde77afff65b745b5e34e5d417}.
\]
The next compulsory companion audit is V20.

% =============================================================================
% END APPENDED FOURTH-SERIES V17 MATERIAL
% =============================================================================

% =============================================================================
% BEGIN APPENDED FOURTH-SERIES V18 MATERIAL
% =============================================================================

\clearpage
\chapter{V18: Exact binary-bank ancestry and failure of the assembled
owner surplus}
\label{ch:V18}

\section{Purpose}
\label{sec:V18-purpose}

V17 leaves open the possibility that the signed internal
first-connectivity sum cancels the low-output tensor-product
obstruction.  For a binary connected prefix, that sum can be evaluated
exactly.  The complete connected carrier over a coordinate bank is the
difference between the common-variable product moment and the
cut-replicated endpoint product.  Its first-connectivity expansion is
the ordinary telescoping identity.

On the equal-spin family, the cut-replicated endpoint is exponentially
small, whereas the common heat square has \(c\,d_j^{-1}\) product
expectation at every occupied coordinate.  Thus the complete assembled
binary carrier over \(k=1,2,3\) coordinates has square capacity at least
\(c\,d_j^{-k}\).  This disproves the assembled owner-surplus card, not
merely its termwise proof.  The same conclusion holds when a nonzero
output at one coordinate pays.

\section{Complete binary covariance over a bank}
\label{sec:V18-bank}

Let \(D=\{1,\ldots,k\}\) be an ordered bank of distinct physical
coordinates, with \(1\le k\le3\).  At every coordinate the two rows
carry \(V_j,V_j\).  Put
\begin{align}
 \mathbb A_j
 &:=
 \mathbb H_{12}(s)
 =
 \sum_{K=0}^{2j}q_KP_{j,K},
\label{eq:V18-Aj}\\
 \mathbb B_j
 &:=
 \mathbb H_1(s)\otimes\mathbb H_2(s)
 =
 q_j^2I.
\label{eq:V18-Bj}
\end{align}
Define
\begin{equation}
 \mathbb A_D:=\mathbb A_j^{\otimes k},
 \qquad
 \mathbb B_D:=\mathbb B_j^{\otimes k}
 =q_j^{2k}I.
\label{eq:V18-ADB}
\end{equation}

\begin{definition}[Complete binary bank carrier]
\label{def:V18-binary-bank}
The complete connected two-row carrier on \(D\) is
\begin{equation}
 \mathbb C_{D,j}
 :=
 \mathbb A_D-\mathbb B_D.
\label{eq:V18-bank-covariance}
\end{equation}
\end{definition}

\begin{lemma}[Exact first-connectivity telescope]
\label{lem:V18-telescope}
\begin{equation}
 \boxed{
 \mathbb C_{D,j}
 =
 \sum_{r=1}^{k}
 \mathbb A_j^{\otimes(r-1)}
 \otimes(\mathbb A_j-\mathbb B_j)
 \otimes\mathbb B_j^{\otimes(k-r)}.}
\label{eq:V18-telescope}
\end{equation}
Each binary connected coordinate word occurs at its unique first
coordinate \(r\) where the two rows use a common-variable moment.
\end{lemma}

\begin{proof}
For any two operators \(A,B\) acting on one coordinate,
\[
 A^{\otimes k}-B^{\otimes k}
 =
 \sum_{r=1}^{k}
 A^{\otimes(r-1)}
 \otimes(A-B)
 \otimes B^{\otimes(k-r)}.
\]
Indeed, expanding the right side makes adjacent terms telescope:
the negative term with \(r\) cancels the positive endpoint preceding
the term with \(r+1\), leaving only \(A^{\otimes k}-B^{\otimes k}\).
Apply this with \(A=\mathbb A_j\), \(B=\mathbb B_j\).
Before \(r\) the common moment has already connected the two rows; at
\(r\) the covariance is the joining letter; after \(r\) the displayed
endpoint convention completes the exact two-row first-connectivity
decomposition.
\end{proof}

\begin{firewall}[The complete bank, not its summands, is tested]
\label{fire:V18-complete-bank}
All cancellations among the \(k\) first-connectivity summands in
\eqref{eq:V18-telescope} have already been performed in
\eqref{eq:V18-bank-covariance}.  The lower bounds below apply to
\(\mathbb C_{D,j}\) itself and therefore cannot be evaded by asking to
assemble the source before taking its square.
\end{firewall}

\section{Unpaid assembled lower bound}
\label{sec:V18-unpaid}

Use the one-coordinate product vector
\[
 \phi_j=|j,j\rangle\otimes|j,-j\rangle
\]
and put \(\Phi_{D,j}:=\phi_j^{\otimes k}\).  Define
\begin{align}
 \alpha_j
 &:=
 \langle\phi_j,\mathbb A_j\phi_j\rangle,
\label{eq:V18-alpha}\\
 u_j
 &:=
 \langle\phi_j,\mathbb A_j^2\phi_j\rangle
 =
 \sum_{K=0}^{2j}p_{j,K}q_K^2.
\label{eq:V18-u}
\end{align}

\begin{lemma}[One-coordinate square mass]
\label{lem:V18-u-lower}
For all \(j\ge1\),
\begin{equation}
 u_j\ge p_{j,1}q_1^2\ge c_s d_j^{-1}.
\label{eq:V18-u-lower}
\end{equation}
Moreover \(0\le\alpha_j\le1\).
\end{lemma}

\begin{proof}
Every term in \eqref{eq:V18-u} is nonnegative.  Retain \(K=1\) and
use \Cref{cor:V13-pj1-scale}.  The operator \(\mathbb A_j\) is a
positive contraction, giving the last assertion.
\end{proof}

\begin{theorem}[Exact assembled binary-bank lower bound]
\label{thm:V18-bank-lower}
For each fixed \(k\in\{1,2,3\}\), there are \(j_{s,k}\) and
\(c_{s,k}>0\) such that
\begin{equation}
 \boxed{
 \kappa_\otimes(\mathbb C_{D,j}^2)
 \ge c_{s,k}d_j^{-k}
 \asymp c_{s,k}M_j^{-k/2}}
\label{eq:V18-bank-lower}
\end{equation}
for every \(j\ge j_{s,k}\).
\end{theorem}

\begin{proof}
Since \(\mathbb A_D,\mathbb B_D\) commute and are self-adjoint,
\begin{align}
 \langle\Phi_{D,j},\mathbb C_{D,j}^2\Phi_{D,j}\rangle
 &=
 \langle\Phi_{D,j},\mathbb A_D^2\Phi_{D,j}\rangle
 -2q_j^{2k}
  \langle\Phi_{D,j},\mathbb A_D\Phi_{D,j}\rangle
 +q_j^{4k}\notag\\
 &=
 u_j^k-2q_j^{2k}\alpha_j^k+q_j^{4k}.
\label{eq:V18-exact-expectation}
\end{align}
By \Cref{lem:V18-u-lower},
\[
 u_j^k\ge c_{s,k}d_j^{-k}.
\]
The middle term has absolute value at most \(2q_j^{2k}\).  Since
\(q_j=e^{-sj(j+1)}\), it is eventually at most
\(\frac12c_{s,k}d_j^{-k}\).  The last term is nonnegative.  Thus
\eqref{eq:V18-exact-expectation} is at least
\(\frac12c_{s,k}d_j^{-k}\).  The vector
\(\Phi_{D,j}\) is a row-product vector, so its expectation lower-bounds
product-state capacity.  Finally \(d_j\asymp M_j^{1/2}\).
\end{proof}

\begin{corollary}[The assembled binary \(AOSC\) fails]
\label{cor:V18-binary-AOSC-fails}
For every \(k\in\{1,2,3\}\), the complete binary bank carrier
\(\mathbb C_{D,j}\) fails the assembled owner-surplus target
\[
 CM_j^{-(k+1)/2}.
\]
\end{corollary}

\begin{proof}
The lower bound in \eqref{eq:V18-bank-lower} is larger than the
displayed target by a factor comparable to \(M_j^{1/2}\).
\end{proof}

\section{A payer on one nonzero output}
\label{sec:V18-paid}

Let \(\mathbb R_s^{(1)}\) act as the V4 payer on the total output of
the first coordinate and as the identity elsewhere.  It commutes with
\(\mathbb A_D\) and \(\mathbb B_D\).  Put
\begin{equation}
 \mathbb C_{D,j}^{\mathrm{pay}}
 :=
 \mathbb R_s^{(1)}\mathbb C_{D,j}.
\label{eq:V18-paid-bank}
\end{equation}

\begin{lemma}[Paid common moment retains the spin-one mass]
\label{lem:V18-paid-common}
There is \(c_{s,k}>0\) such that
\begin{equation}
 \|
  \mathbb R_s^{(1)}\mathbb A_D\Phi_{D,j}
 \|^2
 \ge c_{s,k}d_j^{-k}.
\label{eq:V18-paid-common}
\end{equation}
\end{lemma}

\begin{proof}
Project the first coordinate output onto spin one.  Orthogonal
projection can only decrease the norm.  On that block
\[
 \mathbb R_s\mathbb A_j
 =
 r_1(s)q_1P_{j,1},
\]
so the squared first-coordinate norm is
\[
 r_1(s)^2q_1^2p_{j,1}\ge c_sd_j^{-1}.
\]
At each of the other \(k-1\) coordinates the squared norm is \(u_j\),
which is at least \(c_sd_j^{-1}\) by
\Cref{lem:V18-u-lower}.  Tensorization proves the claim.
\end{proof}

\begin{lemma}[The paid replicated endpoint is exponentially small]
\label{lem:V18-paid-endpoint}
\begin{equation}
 \|
  \mathbb R_s^{(1)}\mathbb B_D
 \|
 \le C_s(1+j^2)q_j^{2k}.
\label{eq:V18-paid-endpoint}
\end{equation}
\end{lemma}

\begin{proof}
On \(V_j\otimes V_j\), every total output has spin
\(0\le K\le2j\).  The payer vanishes at \(K=0\), while for \(K>0\)
\[
 {K(K+1)\over1-e^{-sK(K+1)}}
 \le C_s(1+j^2)
\]
by the positive Casimir gap.  Since
\(\mathbb B_D=q_j^{2k}I\), the result follows.
\end{proof}

\begin{theorem}[Paid assembled binary-bank lower bound]
\label{thm:V18-paid-bank-lower}
For each \(k\in\{1,2,3\}\), there are \(j_{s,k}\) and \(c_{s,k}>0\)
such that
\begin{equation}
 \kappa_\otimes\!\left[
  (\mathbb C_{D,j}^{\mathrm{pay}})^*
  \mathbb C_{D,j}^{\mathrm{pay}}
 \right]
 \ge c_{s,k}d_j^{-k}
\label{eq:V18-paid-bank-lower}
\end{equation}
for all \(j\ge j_{s,k}\).  The opposite square orientation has the
same lower bound.
\end{theorem}

\begin{proof}
By the reverse triangle inequality,
\begin{align*}
 \|\mathbb C_{D,j}^{\mathrm{pay}}\Phi_{D,j}\|
 &\ge
 \|\mathbb R_s^{(1)}\mathbb A_D\Phi_{D,j}\|
 -
 \|\mathbb R_s^{(1)}\mathbb B_D\Phi_{D,j}\|\\
 &\ge
 c_{s,k}d_j^{-k/2}
 -
 C_s(1+j^2)q_j^{2k}.
\end{align*}
The second term is exponentially smaller than the first.  For all
large \(j\), the difference is at least
\(\frac12c_{s,k}d_j^{-k/2}\).  Square and use the row-product test
\(\Phi_{D,j}\).

The payer commutes with the self-adjoint bank covariance, so the paid
operator is self-adjoint.  Hence the two square orientations coincide.
\end{proof}

\begin{corollary}[Payer allocation cannot restore the binary surplus]
\label{cor:V18-paid-AOSC-fails}
The complete assembled binary \(AOSC\) fails whether the unique payer
is outside the binary bank or is allocated to any fixed coordinate
inside it.
\end{corollary}

\begin{proof}
An external payer does not alter the binary bank lower bound.
Coordinate symmetry reduces an internal payer to the first coordinate,
which is \Cref{thm:V18-paid-bank-lower}.
\end{proof}

\section{Compatibility with the saturated mass and strong-stock
ledger}
\label{sec:V18-ledger-compatibility}

\begin{proposition}[The equal-bank family is cofinal and strongly
overlapping]
\label{prop:V18-stock-family}
Let each of the two rows have the same weight
\[
 a_j:=1+j(j+1)
\]
at every coordinate in \(D\), and no weight elsewhere.  Then
\begin{equation}
 \Xi_1=\Xi_2=ka_j,
 \qquad
 \theta_{12}
 =
 {\sum_{e\in D}a_j^2\over(ka_j)^2}
 ={1\over k}.
\label{eq:V18-stock-family}
\end{equation}
Every coordinate carries the fixed fraction \(1/k\) of each row mass.
For \(k=3\), this realizes exact three-slot saturation with a fixed
strong pair.
\end{proposition}

\begin{proof}
All identities follow by direct substitution in
\eqref{eq:V1-stocks}.  For \(k=3\), the entire row mass lies in the
three distinct cofinal coordinates and the normalized overlap is
\(1/3\).
\end{proof}

\begin{corollary}[Failure is at the correct normalized mass scale]
\label{cor:V18-normalized-scale}
Since the total maximal-row mass is
\(M=ka_j\asymp_kM_j\), the lower bounds
\eqref{eq:V18-bank-lower} and
\eqref{eq:V18-paid-bank-lower} have scale \(M^{-k/2}\), whereas
the assembled owner-surplus target has scale \(M^{-(k+1)/2}\).
The failure is therefore not caused by using a local mass different
from the row mass.
\end{corollary}

\begin{proof}
For fixed \(k\), \(ka_j\) and \(M_j\) differ by a fixed factor.  Insert
this comparison into the two exponents.
\end{proof}

\section{Consequences for the \texorpdfstring{\(3+2\)}{3+2} prefix}
\label{sec:V18-consequence}

\begin{theorem}[Binary assembled-cancellation route is exhausted]
\label{thm:V18-binary-route-exhausted}
Suppose the minimal connected owner in a \(3+2\) prefix is binary and
contains \(k\in\{1,2,3\}\) saturated coordinates.  No theorem based
only on:
\begin{enumerate}[label=\textup{(\roman*)},leftmargin=3.5em]
\item exact binary first connectivity,
\item complete heat moments,
\item the unique output payer, and
\item strong overlap of the two rows
\end{enumerate}
can prove its assembled owner-surplus card uniformly in
representations.  A positive quintic proof must use an additional
surrounding source factor which annihilates or suppresses the family
of \Cref{prop:V18-stock-family}, or expose another source occurrence.
\end{theorem}

\begin{proof}
\Cref{lem:V18-telescope} includes the exact binary
first-connectivity assembly.  The unpaid and paid lower bounds are
\Cref{thm:V18-bank-lower,thm:V18-paid-bank-lower}.  The cofinality and
strong-overlap hypotheses are met by
\Cref{prop:V18-stock-family}, yet the required exponent fails by
\Cref{cor:V18-normalized-scale}.  Therefore the listed hypotheses are
insufficient.  Only a factor not present in this complete binary bank
can remove the counterfamily.
\end{proof}

\begin{firewall}[Not yet a quintic MTA counterexample]
\label{fire:V18-not-quintic-counterexample}
The binary bank is a literal lower-order prefix carrier, but a complete
five-row source also contains the complementary prefix block, the root
joining covariance, the suffix, and normalization.  One of those
factors may annihilate the product vector used here.  V18 disproves the
binary assembled-surplus route; it does not disprove global quintic
MTA.
\end{firewall}

\begin{assessment}[Remaining same-owner calculation]
\label{ass:V18-remaining}
The only unresolved assembled lower-order surplus inside a \(3+2\)
prefix is now ternary across multiple physical coordinates.  Its exact
carrier is a third cumulant of three coordinate products, not a tensor
product of local third cumulants.  V19 should write its seven-term
moment formula, evaluate it on the V16 low-output family, and determine
whether the independently heated endpoints are again exponentially
small.
\end{assessment}

\section{V18 result ledger and V19 acquisition order}
\label{sec:V18-ledger}

\begin{theorem}[Fourth-series V18 result ledger]
\label{thm:V18-results}
Relative to V1--V17, V18 proves:
\begin{enumerate}[label=\textup{(V18.\arabic*)},leftmargin=6.3em]
\item the exact binary-bank covariance and its internal
      first-connectivity telescope;
\item a complete assembled unpaid square lower bound
      \(c\,d_j^{-k}\) for \(k=1,2,3\);
\item failure of the binary assembled owner-surplus card;
\item retention of the same lower bound when one nonzero output pays;
\item a cofinal equal-bank family with strong overlap
      \(\theta_{12}=1/k\);
\item agreement of the local and total-row mass scales;
\item exhaustion of binary assembled cancellation as a universal
      \(3+2\) prefix repair; and
\item reduction of the remaining lower-order calculation to the
      multi-coordinate ternary carrier.
\end{enumerate}
\end{theorem}

\begin{proof}
These are
\Cref{lem:V18-telescope,
lem:V18-u-lower,
thm:V18-bank-lower,
cor:V18-binary-AOSC-fails,
lem:V18-paid-common,
lem:V18-paid-endpoint,
thm:V18-paid-bank-lower,
prop:V18-stock-family,
cor:V18-normalized-scale,
thm:V18-binary-route-exhausted}.
\end{proof}

\begin{programme}[V19: complete ternary bank cumulant]
\label{prog:V18-V19}
\begin{enumerate}[leftmargin=3em]
\item For three rows and a coordinate bank \(D\), write the complete
      seven-term third cumulant of the coordinate-product moments.
\item Use spins \(j,j,1\) at each occupied coordinate and the V16
      product vector.
\item On total spin one, separate the leading
      \(H_{123,D}-H_{12,D}H_{3,D}\) difference from terms containing
      an independently heated spin-\(j\) endpoint.
\item Prove or disprove a bank-level norm lower bound
      \(c\,d_j^{-|D|}\) after the complete seven-term assembly.
\item Insert the payer on one spin-one total output and repeat.
\item If the lower bound persists, abandon all lower-order
      owner-surplus routes and inspect the complementary prefix/root
      incidence of the full quintic word.
\end{enumerate}
\end{programme}

\begin{firewall}[Claim boundary after fourth-series V18]
\label{fire:V18-claim-boundary}
V18 settles the complete multi-coordinate binary bank negatively for
the owner-surplus programme.  It does not settle the complete ternary
bank, survival through the full quintic source, source incidence for
the V5 remainder, or \(q=3\).  The remaining \(q=3\) alternatives,
\(q=0,1,2\) cards, global quintic or all-order MTA, WUFC, CPM, cutoff
removal, reflection positivity, Osterwalder--Schrader reconstruction,
construction of the continuum Yang--Mills measure, and a uniform
positive continuum spectral gap remain open.  The full Yang--Mills
mass gap has not been proved.
\end{firewall}

\paragraph{Parent-version record.}
V18 was formed from
\texttt{Renewal\_Yang\_Mills\_Mass\_Gap\_Research\_Notes\_V17.tex}.
The V17 parent had \(8707\) logical source lines, \(309\,547\) bytes,
and SHA--256
\[
 \texttt{2cc58fd13e88446963d0c0fa04e114304f238c67a27fa21ec26e1d6178e3445d}.
\]
The next compulsory companion audit is V20.

% =============================================================================
% END APPENDED FOURTH-SERIES V18 MATERIAL
% =============================================================================

% =============================================================================
% BEGIN APPENDED FOURTH-SERIES V19 MATERIAL
% =============================================================================

\clearpage
\chapter{V19: Exact ternary-bank cumulant and exhaustion of all
lower-order assembled surplus routes}
\label{ch:V19}

\section{Purpose and a counting correction}
\label{sec:V19-purpose}

The V18 programme described the complete third cumulant as
seven-term.  That count is incorrect and is corrected here.  The
partition lattice \(\Pi_3\) has Bell number \(B_3=5\): one joint
moment, three two-block products, and one three-singleton product with
coefficient two.  The exact bank cumulant therefore has five
moment-partition terms.

This chapter evaluates that complete five-term operator over a bank of
\(k=1,2,3\) physical coordinates.  On the product of total-spin-one
blocks, its leading part is
\[
 q_1^k(I-\mathbb H_{12}^{\otimes k}).
\]
The other three moment-partition contributions contain an independently
heated spin-\(j\) row at every coordinate and are exponentially small.
The V16 product vector has \(c\,d_j^{-k}\) weight in the
intermediate-spin-one product block.  Hence the complete assembled
ternary square retains the ordinary \(M_j^{-k/2}\) scale, with or
without the payer.

\begin{correction}[Third-cumulant term count]
\label{corr:V19-five-terms}
The phrase seven-term in
\Cref{prog:V18-V19} is withdrawn.  The complete third cumulant has the
five terms displayed in \eqref{eq:V19-bank-cumulant} below.
\end{correction}

\section{Complete ternary cumulant over a coordinate bank}
\label{sec:V19-bank}

Let \(D=\{1,\ldots,k\}\), \(1\le k\le3\).  At every coordinate the
three rows carry spins \((j,j,1)\).  Define the one-coordinate moment
operators
\begin{align}
 A&:=\mathbb H_{123},&
 B_{12}&:=\mathbb H_{12}\mathbb H_3,\notag\\
 B_{13}&:=\mathbb H_{13}\mathbb H_2,&
 B_{23}&:=\mathbb H_{23}\mathbb H_1,\notag\\
 C&:=\mathbb H_1\mathbb H_2\mathbb H_3.
\label{eq:V19-local-moments}
\end{align}
Their complete bank products are
\[
 A_D:=A^{\otimes k},
 \quad
 B_{uv,D}:=B_{uv}^{\otimes k},
 \quad
 C_D:=C^{\otimes k}.
\]

\begin{definition}[Complete ternary bank cumulant]
\label{def:V19-bank-cumulant}
\begin{equation}
 \boxed{
 \mathbb K_{3,D}
 :=
 A_D-B_{12,D}-B_{13,D}-B_{23,D}+2C_D.}
\label{eq:V19-bank-cumulant}
\end{equation}
\end{definition}

\begin{lemma}[Exact moment--cumulant identity]
\label{lem:V19-exact-cumulant}
The operator in \eqref{eq:V19-bank-cumulant} is the complete connected
three-row carrier over \(D\).  It includes the entire internal
first-connectivity sum and is not the tensor product of the local
third cumulants \(\mathbb K_{3,j}\).
\end{lemma}

\begin{proof}
Let the three bank variables be the coordinate products
\[
 Y_i:=\prod_{e\in D}F_{i,e}.
\]
Independence across physical coordinates makes the moment of a row
block \(S\subseteq\{1,2,3\}\) equal to
\[
 \mathbb E\prod_{i\in S}Y_i
 =
 \bigotimes_{e\in D}\mathbb H_{S,e}.
\]
The order-three moment--cumulant identity is
\begin{align*}
 \kappa_3(Y_1,Y_2,Y_3)
 ={}&
 \mathbb E(Y_1Y_2Y_3)
 -\mathbb E(Y_1Y_2)\mathbb E(Y_3)\\
 &-\mathbb E(Y_1Y_3)\mathbb E(Y_2)
 -\mathbb E(Y_2Y_3)\mathbb E(Y_1)\\
 &+2\mathbb E(Y_1)\mathbb E(Y_2)\mathbb E(Y_3).
\end{align*}
Substitution gives \eqref{eq:V19-bank-cumulant}.  Exact internal first
connectivity is a coefficient-one partition of the connected
coordinate words in this cumulant, so summing it restores the same
operator.  In contrast, tensoring local cumulants would expand to
\(5^k\) terms with different coefficients and is not the displayed
five-term bank cumulant.
\end{proof}

\begin{firewall}[All ancestry is already assembled]
\label{fire:V19-assembled}
The lower bound below is applied to the five-term complete carrier
\(\mathbb K_{3,D}\).  It therefore includes every cancellation among
internal first-connectivity coordinates and every cross term in the
square.  No termwise positive envelope is used.
\end{firewall}

\section{Leading total-spin-one product block}
\label{sec:V19-leading}

At each coordinate let \(P^{\mathrm{tot}}_1\) and \(Q_1\) be the
projections from V16: \(P^{\mathrm{tot}}_1\) selects total spin one,
and \(Q_1\) further selects intermediate spin one for rows \(1,2\).
Put
\[
 P_D:=(P^{\mathrm{tot}}_1)^{\otimes k},
 \qquad
 Q_D:=Q_1^{\otimes k}.
\]

\begin{lemma}[Exact leading-remainder split]
\label{lem:V19-leading-split}
On \(P_D\mathscr H\),
\begin{equation}
 P_D\mathbb K_{3,D}P_D
 =
 q_1^k\bigl(I-\mathbb H_{12}^{\otimes k}\bigr)
 +
 \mathbb E_{D,j},
\label{eq:V19-leading-split}
\end{equation}
where
\begin{equation}
 \|\mathbb E_{D,j}\|
 \le2q_j^k+2q_j^{2k}.
\label{eq:V19-bank-remainder}
\end{equation}
\end{lemma}

\begin{proof}
On total spin one at every coordinate,
\[
 A_D=q_1^kI.
\]
Since row three has spin one,
\[
 B_{12,D}
 =
 (q_1\mathbb H_{12})^{\otimes k}
 =
 q_1^k\mathbb H_{12}^{\otimes k}.
\]
At each coordinate \(B_{13}\) contains the independent one-row factor
\(\mathbb H_2=q_jI\), while \(B_{23}\) contains
\(\mathbb H_1=q_jI\).  Their bank-product norms are therefore at most
\(q_j^k\).  Finally
\[
 C_D=(q_j^2q_1)^kI,
\]
whose contribution with coefficient two has norm at most
\(2q_j^{2k}\).  This proves both displays.
\end{proof}

\begin{lemma}[Uniform leading gap on \(Q_D\)]
\label{lem:V19-leading-gap}
On the range of \(Q_D\),
\begin{equation}
 q_1^k\bigl(I-\mathbb H_{12}^{\otimes k}\bigr)
 =
 \delta_{s,k}I,
 \qquad
 \delta_{s,k}:=q_1^k(1-q_1^k)>0.
\label{eq:V19-leading-gap}
\end{equation}
\end{lemma}

\begin{proof}
On \(Q_1\), the intermediate heat moment satisfies
\(\mathbb H_{12}=q_1I\).  Its \(k\)-fold tensor product therefore has
eigenvalue \(q_1^k\) on \(Q_D\).  Substitute this into the leading
operator.
\end{proof}

\section{Complete unpinched bank lower bound}
\label{sec:V19-lower}

Let \(\Phi_j\) be the row-product vector from
\eqref{eq:V16-product-vector} and set
\[
 \Phi_{D,j}:=\Phi_j^{\otimes k},
 \qquad
 v_{D,j}:=P_D\Phi_{D,j}.
\]

\begin{lemma}[Product weight in the leading block]
\label{lem:V19-QD-weight}
\begin{equation}
 \|Q_Dv_{D,j}\|^2
 =
 \left({p_{j,1}\over2}\right)^k
 \ge c_kd_j^{-k}.
\label{eq:V19-QD-weight}
\end{equation}
\end{lemma}

\begin{proof}
At one coordinate,
\(\|Q_1P^{\mathrm{tot}}_1\Phi_j\|^2=p_{j,1}/2\) by
\Cref{lem:V16-Q1-component}.  Tensorization over the \(k\) distinct
coordinates gives the equality.  Use
\Cref{cor:V13-pj1-scale} for the lower bound.
\end{proof}

\begin{theorem}[Complete ternary-bank square lower bound]
\label{thm:V19-bank-lower}
For each \(k\in\{1,2,3\}\), there are \(j_{s,k}<\infty\) and
\(c_{s,k}>0\) such that
\begin{equation}
 \boxed{
 \kappa_\otimes(\mathbb K_{3,D}^2)
 \ge c_{s,k}d_j^{-k}
 \asymp c_{s,k}M_j^{-k/2}}
\label{eq:V19-bank-lower}
\end{equation}
for every \(j\ge j_{s,k}\).
\end{theorem}

\begin{proof}
The complete bank cumulant commutes with total spin separately at every
coordinate.  Hence
\[
 P_D\mathbb K_{3,D}\Phi_{D,j}
 =
 \left[
  q_1^k(I-\mathbb H_{12}^{\otimes k})
  +\mathbb E_{D,j}
 \right]v_{D,j}.
\]
The leading operator is diagonal in the product intermediate-spin
resolution.  Therefore
\[
 \|
  q_1^k(I-\mathbb H_{12}^{\otimes k})v_{D,j}
 \|
 \ge
 \delta_{s,k}\|Q_Dv_{D,j}\|.
\]
By the reverse triangle inequality and
\eqref{eq:V19-bank-remainder},
\begin{align*}
 \|\mathbb K_{3,D}\Phi_{D,j}\|
 &\ge
 \|P_D\mathbb K_{3,D}\Phi_{D,j}\|\\
 &\ge
 \delta_{s,k}\|Q_Dv_{D,j}\|
 -(2q_j^k+2q_j^{2k}).
\end{align*}
The first term is at least
\(c_{s,k}d_j^{-k/2}\) by
\Cref{lem:V19-QD-weight}; the second is exponentially smaller.  For
all sufficiently large \(j\), the difference is at least half the
first term.  Squaring gives
\[
 \langle\Phi_{D,j},
  \mathbb K_{3,D}^2\Phi_{D,j}\rangle
 \ge c_{s,k}d_j^{-k}.
\]
The vector \(\Phi_{D,j}\) is a row-product vector, and
\(d_j\asymp M_j^{1/2}\), proving the result.
\end{proof}

\begin{firewall}[Analytic projections only]
\label{fire:V19-no-source-projection}
The tested operator is the complete five-term
\(\mathbb K_{3,D}\).  The projections \(P_D,Q_D\) occur only on the
right side of the inequality
\(\|\mathbb K_{3,D}\Phi\|\ge
\|P_D\mathbb K_{3,D}\Phi\|\), and then to retain one orthogonal
component of a norm.  They are not inserted into the source carrier.
\end{firewall}

\section{The paid bank}
\label{sec:V19-paid}

Let \(\mathbb R_s^{(1)}\) be the selected total-output payer at the
first coordinate and define
\[
 \mathbb K_{3,D}^{\mathrm{pay}}
 :=
 \mathbb R_s^{(1)}\mathbb K_{3,D}.
\]

\begin{theorem}[Paid complete ternary-bank lower bound]
\label{thm:V19-paid-lower}
For each \(k\in\{1,2,3\}\), after increasing \(j_{s,k}\),
\begin{equation}
 \kappa_\otimes\!\left[
  (\mathbb K_{3,D}^{\mathrm{pay}})^*
  \mathbb K_{3,D}^{\mathrm{pay}}
 \right]
 \ge c_{s,k}^{\mathrm{pay}}d_j^{-k}.
\label{eq:V19-paid-lower}
\end{equation}
The opposite orientation obeys the same lower bound.
\end{theorem}

\begin{proof}
On \(P_D\), the first-coordinate total output is spin one, so
\(\mathbb R_s^{(1)}=r_1(s)I\), with
\(r_1(s)>0\) given by \eqref{eq:V13-r1}.  Orthogonality of total-output
blocks gives
\[
 \|\mathbb R_s^{(1)}\mathbb K_{3,D}\Phi_{D,j}\|
 \ge
 r_1(s)\|P_D\mathbb K_{3,D}\Phi_{D,j}\|.
\]
The proof of \Cref{thm:V19-bank-lower} lower-bounds the last norm by
\(c_{s,k}d_j^{-k/2}\).  Square.  The payer commutes with the complete
cumulant; both are self-adjoint, so the paid packet is self-adjoint and
the two square orientations coincide.
\end{proof}

\begin{corollary}[The assembled ternary \(AOSC\) fails]
\label{cor:V19-ternary-AOSC-fails}
For every \(k\in\{1,2,3\}\), the complete ternary bank cumulant fails
the assembled owner-surplus target
\[
 CM_j^{-(k+1)/2},
\]
both paid and unpaid.
\end{corollary}

\begin{proof}
\Cref{thm:V19-bank-lower,thm:V19-paid-lower} give the larger scale
\(cM_j^{-k/2}\).
\end{proof}

\section{Mass and overlap compatibility}
\label{sec:V19-stocks}

\begin{proposition}[The ternary family has a saturated strong pair]
\label{prop:V19-stock-family}
At every coordinate in \(D\), assign weight
\(a_j=1+j(j+1)\) to rows one and two and the fixed weight
\(a_1=3\) to row three.  Put no weight elsewhere.  Then
\[
 \Xi_1=\Xi_2=ka_j,
 \qquad
 \Xi_3=3k,
 \qquad
 \theta_{12}={1\over k}.
\]
Every coordinate in \(D\) carries \(1/k\) of each maximal-row mass.
For \(k=3\), rows one and two realize a fixed strong pair entirely on
three saturated coordinates.
\end{proposition}

\begin{proof}
Substitute the displayed weights into
\eqref{eq:V1-stocks}, exactly as in
\Cref{prop:V18-stock-family}.  The fixed third-row weight does not
enter \(\theta_{12}\).
\end{proof}

\section{Exhaustion of lower-order assembled owner surplus}
\label{sec:V19-exhaustion}

\begin{theorem}[No universal lower-order owner-surplus route remains]
\label{thm:V19-lower-order-exhausted}
For a minimal prefix owner of arity at most three and a bank of
\(k\in\{1,2,3\}\) saturated coordinates:
\begin{enumerate}[label=\textup{(\roman*)},leftmargin=3.5em]
\item a one-row cofinal prefix is exponentially small under the V14
      scalar-heat hypothesis;
\item a complete assembled binary bank has no owner surplus by V18;
      and
\item a complete assembled ternary bank has no owner surplus by V19.
\end{enumerate}
Therefore the remaining \(3+2\) saturated source cannot be closed from
the intrinsic lower-order prefix carrier alone.  Any valid closure must
use its complementary prefix block together with the root joining
carrier, prove an exact annihilation of the displayed low-output
family, or expose an additional source-occurring coordinate.
\end{theorem}

\begin{proof}
The possible owner arities are exhaustive by
\Cref{thm:V15-lower-order-localization}.  The one-row result is
\Cref{thm:V14-complete-barrier}.  The complete binary and ternary
negative results are
\Cref{thm:V18-binary-route-exhausted,
cor:V19-ternary-AOSC-fails}.  These leave only data belonging to the
rest of the root source word or a genuine selection rule.
\end{proof}

\begin{assessment}[The next upstream object]
\label{ass:V19-next}
V20 should no longer analyze a prefix block in isolation.  After its
compulsory companion audit, it must compose both \(3+2\) prefix blocks
with the root covariance on the V19 product family.  The key question
is whether the second prefix block and the root joining covariance
annihilate the fixed total-spin-one product channel.  If the channel
survives, the same-support programme yields a full source-level
obstruction and the only remaining route is new source incidence.
\end{assessment}

\section{V19 result ledger and V20 acquisition order}
\label{sec:V19-ledger}

\begin{theorem}[Fourth-series V19 result ledger]
\label{thm:V19-results}
Relative to V1--V18, V19 proves:
\begin{enumerate}[label=\textup{(V19.\arabic*)},leftmargin=6.3em]
\item correction of the third-cumulant bank formula to five
      moment-partition terms;
\item the exact complete ternary bank cumulant and its distinction from
      a tensor product of local cumulants;
\item the leading total-spin-one product block and exponentially small
      remainder;
\item a \(c\,d_j^{-k}\) product weight in the gapped intermediate
      channel;
\item a complete unpinched assembled square lower bound
      \(c\,d_j^{-k}\) for \(k=1,2,3\);
\item the same lower bound with one nonzero total-output payer;
\item a saturated strong-pair realization of the analytic family; and
\item exhaustion of every universal assembled lower-order
      owner-surplus route at arity at most three.
\end{enumerate}
\end{theorem}

\begin{proof}
These are
\Cref{corr:V19-five-terms,
lem:V19-exact-cumulant,
lem:V19-leading-split,
lem:V19-leading-gap,
lem:V19-QD-weight,
thm:V19-bank-lower,
thm:V19-paid-lower,
cor:V19-ternary-AOSC-fails,
prop:V19-stock-family,
thm:V19-lower-order-exhausted}.
\end{proof}

\begin{programme}[V20: compulsory audit and full \(3+2\) source
composition]
\label{prog:V19-V20}
\begin{enumerate}[leftmargin=3em]
\item Audit the newest active SM and NS notes with exact sizes and
      SHA--256 digests.
\item Choose a fixed \(3+2\) root partition and write both complete
      prefix block carriers, the root covariance, and the suffix in
      their literal order.
\item Propagate the V19 product family through the complementary prefix
      block without resolving multiplicity copies.
\item Evaluate the root covariance on the resulting fixed low-output
      block, including every payer placement.
\item If a nonzero channel survives, compare its full source-word
      square with the exchange-rescue target and record the exact
      obstruction.
\item If it vanishes, identify and prove the representation selection
      rule responsible before attempting an upper bound.
\end{enumerate}
\end{programme}

\begin{firewall}[Claim boundary after fourth-series V19]
\label{fire:V19-claim-boundary}
V19 settles the complete multi-coordinate ternary bank negatively for
the owner-surplus programme and exhausts all isolated lower-order
prefix-owner routes.  It does not determine survival after composition
with the complementary prefix and root joining carrier, prove a full
quintic counterexample, or close \(q=3\).  The remaining \(q=3\)
alternatives, \(q=0,1,2\) cards, global quintic or all-order MTA,
WUFC, CPM, cutoff removal, reflection positivity,
Osterwalder--Schrader reconstruction, construction of the continuum
Yang--Mills measure, and a uniform positive continuum spectral gap
remain open.  The full Yang--Mills mass gap has not been proved.
\end{firewall}

\paragraph{Parent-version record.}
V19 was formed from
\texttt{Renewal\_Yang\_Mills\_Mass\_Gap\_Research\_Notes\_V18.tex}.
The V18 parent had \(9210\) logical source lines, \(325\,793\) bytes,
and SHA--256
\[
 \texttt{0eba0b4f66f9c0fd31ef2aa786a7f21bb0603c0320dfa5da2f27e38e46f241a1}.
\]
The next compulsory companion audit is V20.

% =============================================================================
% END APPENDED FOURTH-SERIES V19 MATERIAL
% =============================================================================

% =============================================================================
% BEGIN APPENDED FOURTH-SERIES V20 MATERIAL
% =============================================================================

\clearpage
\chapter{V20: Companion audit and transverse-block closure in the
balanced pair-crossing \texorpdfstring{\(3+2\)}{3+2} source}
\label{ch:V20}

\section{Compulsory companion audit}
\label{sec:V20-audit}

The newest active companion files at the time of this audit were
\begin{center}
\begin{tabular}{@{}p{0.49\linewidth}rr@{}}
\toprule
file & logical lines & bytes\\
\midrule
\texttt{Renewal\_SM\_Einstein\_Continuum\_Research\_Notes\_V97.tex}
 & \(36\,040\) & \(1\,329\,523\)\\
\texttt{Renewal\_NS\_Stability\_Research\_Notes\_V31.tex}
 & \(23\,052\) & \(887\,537\)\\
\bottomrule
\end{tabular}
\end{center}
Their SHA--256 digests are
\[
\begin{split}
&\texttt{5fc11be539e67c7224f68c736bfdb828b59e24d369560a58dd5436473c62b20c},\\
&\texttt{f1121b62238ad5491bb7cf03635ea9934942edf573ed8faaa9ee79e1d38a6696}.
\end{split}
\]

SM V97 constructs graphwise shell smoothing only after local Taylor
owners and subdivergences have been extracted.  A connected external
kernel gains smoothing from a tree path joining its external supports;
an internal loop which does not lie on such a path is not counted as
external smoothing.  Every forest term has one owner.

The type-correct YM transfer is an incidence rule.  A complementary
prefix block may supply the missing response only because it is a
literal tensor factor in the coefficient-one \(3+2\) source.  A row
label, a strong-edge label, or a factor from a different source fibre
cannot be substituted.  No SM shell, RG, determinant, or compactness
estimate is imported.

NS V31 is unchanged.  Its parent-sum warning again requires the
transverse response to remain in the same exact source word; a global
overlap sum cannot be multiplied into every first-connectivity fibre.

\begin{audit}[V20 source-level direction]
\label{audit:V20-direction}
Compose the two literal \(3+2\) prefix factors before seeking stronger
local cumulant estimates.  When the strong partner lies in the
complementary block, a certified partner response is transverse to the
maximal-row prefix packet and may be multiplied exactly once.  The
response must be localized to an actual prefix coordinate of that
same fibre.
\end{audit}

\section{Transverse row-block tensorization}
\label{sec:V20-transverse}

\begin{lemma}[Exact capacity product on disjoint row blocks]
\label{lem:V20-capacity-product}
Let \(B,C\) be disjoint row sets.  Let \(X\ge0\) act on all coordinate
factors of rows in \(B\), and let \(Y\ge0\) act on all coordinate
factors of rows in \(C\).  Then
\begin{equation}
 \boxed{
 \kappa_{B\sqcup C}(X\otimes Y)
 =
 \kappa_B(X)\kappa_C(Y).}
\label{eq:V20-capacity-product}
\end{equation}
The equality remains valid when \(X,Y\) occupy some of the same
physical coordinates, because their row tensor factors are disjoint.
\end{lemma}

\begin{proof}
A product density across all rows in \(B\sqcup C\) has the form
\[
 \left(\bigotimes_{i\in B}\rho_i\right)
 \otimes
 \left(\bigotimes_{i\in C}\rho_i\right).
\]
Its expectation on \(X\otimes Y\) is the product of the two
expectations.  Taking the supremum is at most the product of the two
separate suprema.  Conversely, choose product densities in the two
blocks which approach their respective suprema and tensor them.  This
proves equality.  Physical-coordinate labels do not enter the
factorization; only disjointness of row Hilbert factors is used.
\end{proof}

\begin{corollary}[A same-coordinate transverse response is not a
duplicate]
\label{cor:V20-same-coordinate-legal}
Suppose the \(3+2\) prefix is
\[
 \mathcal K_{B_3}^{<r,\mathrm{conn}}
 \otimes
 \mathcal K_{B_2}^{<r,\mathrm{conn}}.
\]
An operator response in the \(B_3\) factor and one in the \(B_2\)
factor may be tensorized even when both occur at the same physical
coordinate.  This uses two literal maps on disjoint row factors, not
two estimates of one map.
\end{corollary}

\begin{proof}
The tensor product is literal in
\eqref{eq:V11-exact-source}.  Apply
\Cref{lem:V20-capacity-product}.  Unique minimal ownership from V15
assigns the two maps to disjoint child blocks, so neither is an
ancestor copy of the other.
\end{proof}

\section{Localization of a balanced crossing partner}
\label{sec:V20-partner}

Let \(a\) be the maximal row, \(\Xi_a=M\), and let \(c\) be a
comparable strong partner:
\begin{equation}
 \Xi_c\ge\lambda M
\label{eq:V20-balanced-partner}
\end{equation}
for fixed \(\lambda>0\).  For a set \(D\) of physical coordinates
define
\[
 \theta_{ac}(D)
 :=
 {1\over\Xi_a\Xi_c}
 \sum_{e\in D}a_{a,e}a_{c,e}.
\]

\begin{lemma}[Prefix overlap forces a cofinal partner coordinate]
\label{lem:V20-partner-localization}
Let
\[
 D:=E_3\cap\{e:e<r\}.
\]
If
\begin{equation}
 \theta_{ac}(D)\ge\beta>0,
\label{eq:V20-prefix-overlap}
\end{equation}
then some \(e_*\in D\) satisfies
\begin{equation}
 a_{c,e_*}\ge{\beta\lambda\over3}M.
\label{eq:V20-partner-cofinal}
\end{equation}
\end{lemma}

\begin{proof}
The overlap hypothesis and \(\Xi_a=M\) give
\[
 \sum_{e\in D}a_{a,e}a_{c,e}
 \ge\beta M\Xi_c.
\]
Since \(a_{a,e}\le\Xi_a=M\),
\[
 \sum_{e\in D}a_{c,e}\ge\beta\Xi_c\ge\beta\lambda M.
\]
The set \(D\) contains at most the three coordinates of \(E_3\).
One summand is therefore at least one third of the last quantity.
\end{proof}

\begin{remark}[Why comparability is explicit]
\label{rem:V20-comparability}
A fixed normalized overlap does not by itself imply
\(\Xi_c\asymp M\).  A row of bounded total mass can have normalized
overlap one with a cofinal row if both are supported at one coordinate.
Thus \eqref{eq:V20-balanced-partner} cannot be omitted from the
localization lemma.
\end{remark}

\section{The complementary-block certificate}
\label{sec:V20-certificate}

Let the root prefix partition be
\[
 \rho=\{B_a,B_c\},
\]
where \(a\in B_a\), \(c\in B_c\), and
\(\{|B_a|,|B_c|\}=\{3,2\}\).  This is the pair-crossing
\(3+2\) geometry.

\begin{definition}[Certified transverse partner]
\label{def:V20-certified-partner}
The coordinate \(e_*<r\) is a certified transverse partner if the
complete literal operator at \(e_*\) inside
\(\mathcal K_{B_c}^{<r,\mathrm{conn}}\), with every coincident
joining and payer role retained, obeys both square estimates
\begin{align}
 \kappa_{B_c}(\mathbb T_{e_*}^*\mathbb T_{e_*})
 &\le C a_{c,e_*}^{-1/2},
\label{eq:V20-T-right}\\
 \kappa_{B_c}(\mathbb T_{e_*}\mathbb T_{e_*}^*)
 &\le C a_{c,e_*}^{-1/2}.
\label{eq:V20-T-left}
\end{align}
The paid version is used if the sole payer is allocated at \(e_*\).
\end{definition}

\begin{lemma}[Transverse fourth half-power]
\label{lem:V20-fourth-half}
Assume \eqref{eq:V20-prefix-overlap} and that the coordinate from
\Cref{lem:V20-partner-localization} is certified.  Then its complete
transverse factor has both square capacities bounded by
\begin{equation}
 C_{\beta,\lambda}M^{-1/2}.
\label{eq:V20-fourth-half}
\end{equation}
\end{lemma}

\begin{proof}
Insert \eqref{eq:V20-partner-cofinal} into
\eqref{eq:V20-T-right} and
\eqref{eq:V20-T-left}.
\end{proof}

\begin{theorem}[Balanced pair-crossing \(3+2\) prefix closure]
\label{thm:V20-32-crossing-close}
Consider a saturated \(q=3\) source fibre of type \(3+2\).  Assume:
\begin{enumerate}[label=\textup{(\roman*)},leftmargin=3.5em]
\item the fixed strong pair \(ac\) is pair-crossing, with
      \(a\in B_a\), \(c\in B_c\);
\item the partner is balanced as in
      \eqref{eq:V20-balanced-partner};
\item its prefix overlap satisfies
      \eqref{eq:V20-prefix-overlap}; and
\item the localized \(e_*\) is a certified transverse partner.
\end{enumerate}
Then the complete source fibre satisfies the missing \(M^{-2}\)
oriented square response for every payer placement.
\end{theorem}

\begin{proof}
The three already extracted maximal-row heat slots give the complete
\(B_a\)-side square response
\[
 C M^{-3/2}.
\]
They are counted exactly once inside their assembled maximal-row
packet.  By \Cref{lem:V20-fourth-half}, the literal complementary
\(B_c\) factor gives \(CM^{-1/2}\).  The exact root prefix is the
tensor product of the \(B_a\) and \(B_c\) connected carriers.
Therefore \Cref{lem:V20-capacity-product} gives
\[
 CM^{-3/2}M^{-1/2}=CM^{-2}.
\]

The complete root joining covariance, suffix, and output dual
contractions are retained in their literal order.  Their fixed-arity
separators are handled by
\Cref{thm:V11-nonsymmetric-dual}.  If the sole payer lies in the
maximal-row packet or the transverse packet, use the corresponding
paid square card; if it lies in the root separator or suffix, use its
established one-payer bound.  Only one paid estimate is used.  The
argument is symmetric under adjoint, so both orientations follow.
\end{proof}

\begin{corollary}[No physical-coordinate disjointness is required
between the fourth response and \(E_3\)]
\label{cor:V20-no-coordinate-disjointness}
The theorem remains valid when \(e_*\in E_3\), as it necessarily is
under \Cref{lem:V20-partner-localization}.  Disjointness of the two root
row blocks replaces physical-coordinate disjointness.
\end{corollary}

\begin{proof}
At \(e_*\), the maximal-row local map lies in the \(B_a\) tensor
factor and the partner map lies in the \(B_c\) tensor factor.
\Cref{cor:V20-same-coordinate-legal} applies.
\end{proof}

\begin{firewall}[Certification is still an operator condition]
\label{fire:V20-certification}
The scalar overlap proof produces a cofinal partner coordinate but does
not prove
\eqref{eq:V20-T-right}--\eqref{eq:V20-T-left}.  Those estimates follow
only when the complete local map is one of the already certified heat,
covariance, selected-cumulant, or compatible complete-block packets.
A noncentral unclassified map cannot be replaced by its ordinary norm.
\end{firewall}

\section{Exact residual after transverse closure}
\label{sec:V20-residual}

\begin{theorem}[Pair-crossing \(3+2\) residual]
\label{thm:V20-32-residual}
After \Cref{thm:V20-32-crossing-close}, a pair-crossing \(3+2\)
counterpacket must retain at least one of:
\begin{enumerate}[label=\textup{(\alph*)},leftmargin=3.5em]
\item the strong partner is not comparable to the maximal row;
\item the prefix share
      \(\theta_{ac}(E_3\cap\{e<r\})\) tends to zero;
\item the localized complementary-block coordinate has no certified
      complete operator response; or
\item the exact source factorization needed to retain the transverse
      packet fails.
\end{enumerate}
Pair-internal \(3+2\) sources remain governed by the earlier
prefix-internal descent and are not included in this theorem.
\end{theorem}

\begin{proof}
If none of \textup{(a)--(d)} occurs, pass to a subsequence on which
fixed \(\lambda,\beta>0\) satisfy
\eqref{eq:V20-balanced-partner} and
\eqref{eq:V20-prefix-overlap}, and on which the localized coordinate
has the certified source occurrence.  Then all hypotheses of
\Cref{thm:V20-32-crossing-close} hold, contradicting persistence of a
counterpacket.  The final sentence records the mutually exclusive
pair placement.
\end{proof}

\begin{assessment}[Audit-adjusted next direction]
\label{ass:V20-next}
The isolated lower-order owner-surplus programme is exhausted, but the
full \(3+2\) source supplies a legal transverse response in the
balanced pair-crossing prefix-overlap sector.  V21 should localize the
remaining strong stock among prefix, joining coordinate, and suffix.
If the prefix share vanishes, the fixed overlap must concentrate at the
joining coordinate or after it.  Those are literal root/suffix factors
and should be tested before any new local estimate is invented.
\end{assessment}

\section{V20 result ledger and V21 acquisition order}
\label{sec:V20-ledger}

\begin{theorem}[Fourth-series V20 result ledger]
\label{thm:V20-results}
Relative to V1--V19, V20 proves:
\begin{enumerate}[label=\textup{(V20.\arabic*)},leftmargin=6.3em]
\item the compulsory hashed audit of SM V97 and NS V31;
\item exact product-state capacity factorization on disjoint row
      blocks, including at the same physical coordinate;
\item localization of fixed prefix overlap to a cofinal balanced
      partner coordinate;
\item the certified transverse partner card;
\item closure of the balanced pair-crossing \(3+2\)
      prefix-overlap sector for every payer placement;
\item legality of using a complementary-block response on an
      \(E_3\) coordinate without duplicating an operator; and
\item the exhaustive four-part residual for pair-crossing \(3+2\)
      counterpackets.
\end{enumerate}
\end{theorem}

\begin{proof}
These are
\Cref{audit:V20-direction,
lem:V20-capacity-product,
cor:V20-same-coordinate-legal,
lem:V20-partner-localization,
def:V20-certified-partner,
lem:V20-fourth-half,
thm:V20-32-crossing-close,
cor:V20-no-coordinate-disjointness,
thm:V20-32-residual}.
\end{proof}

\begin{programme}[V21: joining/suffix localization of the residual
strong stock]
\label{prog:V20-V21}
\begin{enumerate}[leftmargin=3em]
\item Split \(\theta_{ac}(E_3)\) exactly into prefix,
      first-connectivity-coordinate, and suffix shares.
\item In the balanced pair-crossing branch, use V20 immediately when
      the prefix share is bounded below.
\item If the joining-coordinate share is bounded below, retain the
      complete root covariance and the complementary block response in
      one paid/unpaid packet.
\item If the suffix share is bounded below, use the complete suffix
      tensor factor before deleting any marked coordinates.
\item Treat a noncomparable partner by returning to the normalized row
      denominator rather than pretending its local Casimir is
      maximal-scale.
\item Keep certification failure as an exact operator list, following
      the V10--V11 complete-block rules.
\end{enumerate}
\end{programme}

\begin{firewall}[Claim boundary after fourth-series V20]
\label{fire:V20-claim-boundary}
V20 closes only the balanced, pair-crossing \(3+2\) sector with a
fixed prefix share of strong overlap and a certified complementary
operator.  It does not close joining/suffix overlap, noncomparable
partners, certification failures, other prefix types, the complete
\(q=3\) residual, or \(q=0,1,2\).  Global quintic or all-order MTA,
WUFC, CPM, cutoff removal, reflection positivity,
Osterwalder--Schrader reconstruction, construction of the continuum
Yang--Mills measure, and a uniform positive continuum spectral gap
remain open.  The full Yang--Mills mass gap has not been proved.
\end{firewall}

\paragraph{Parent-version record.}
V20 was formed from
\texttt{Renewal\_Yang\_Mills\_Mass\_Gap\_Research\_Notes\_V19.tex}.
The V19 parent had \(9715\) logical source lines, \(341\,822\) bytes,
and SHA--256
\[
 \texttt{579dab6aad5a1530f54d8ae6baf2327ecc9a73ae25d1ef5ee07f0d78bc4564d1}.
\]
The next compulsory companion audit is V25.

% =============================================================================
% END APPENDED FOURTH-SERIES V20 MATERIAL
% =============================================================================

% =============================================================================
% BEGIN APPENDED FOURTH-SERIES V21 MATERIAL
% =============================================================================

\clearpage
\chapter{V21: Source-order correction and exact prefix--joining--suffix
localization}
\label{ch:V21}

\section{Correction of the V20 closure scope}
\label{sec:V21-correction}

The transverse tensor identity in V20 is exact inside the \(3+2\)
prefix.  However, the proof of
\Cref{thm:V20-32-crossing-close} used all three \(E_3\) heat responses
as though they belonged to that prefix tensor factor.  In a general
temporal class, one slot may be the root joining coordinate and other
slots may lie in the suffix.  Complete output duality controls a bounded
separator between two already estimated outer factors; it does not move
a product-state-capacity gain from inside that separator to the prefix.
The theorem therefore needs an explicit source-order certificate.

\begin{definition}[Transverse outer-factorization certificate]
\label{def:V21-TOF}
A pair-crossing \(3+2\) source fibre has \(TOF_3\) if, in each square
orientation and after complete output dualization, its normalized word
can be written with two occurrences of an outer packet
\[
 \mathbb P
 =
 \mathbb P_{B_a,E_3}\otimes\mathbb T_{B_c,e_*},
\]
where:
\begin{enumerate}[label=\textup{(\roman*)},leftmargin=3.5em]
\item \(\mathbb P_{B_a,E_3}\) contains all three certified
      maximal-row slot responses and has both square capacities
      \(O(M^{-3/2})\);
\item \(\mathbb T_{B_c,e_*}\) is the certified transverse partner of
      V20 and has both square capacities \(O(M^{-1/2})\);
\item the two factors act on disjoint root row blocks; and
\item everything between the two outer occurrences is one complete
      bounded equivariant separator to which
      \Cref{thm:V11-nonsymmetric-dual} applies.
\end{enumerate}
\end{definition}

\begin{theorem}[Corrected transverse closure]
\label{thm:V21-corrected-transverse}
The conclusion of
\Cref{thm:V20-32-crossing-close} holds under its four displayed
hypotheses together with \(TOF_3\).
\end{theorem}

\begin{proof}
By disjoint row-block tensorization
\Cref{lem:V20-capacity-product}, the two square capacities of
\(\mathbb P\) are
\[
 O(M^{-3/2})O(M^{-1/2})=O(M^{-2}).
\]
Apply \Cref{thm:V11-nonsymmetric-dual} to the two outer occurrences
and their complete separator.  Its geometric mean retains the
\(M^{-2}\) exponent.  Paid outer packets use their paid square card,
and a payer in the separator is retained in its established complete
paid norm.  This proves both orientations with one payer.
\end{proof}

\begin{correction}[Withdrawal of the unconditional V20 reading]
\label{corr:V21-V20}
The unconditional wording of
\Cref{thm:V20-32-crossing-close,
cor:V20-no-coordinate-disjointness,
thm:V20-32-residual}
is withdrawn.  Those statements are valid only when \(TOF_3\), or an
equivalent exact two-orientation source factorization, is verified.
The prefix tensorization and partner-localization lemmas themselves
remain valid.
\end{correction}

\begin{proof}
\Cref{thm:V21-corrected-transverse} proves the claims with the missing
hypothesis.  If a heat response lies at the root or in the suffix, it
can be separated from the prefix by a noncentral operator.  The
one-sided entangler counterexample
\Cref{prop:V9-one-sided-fails} shows that ordinary norm control cannot
justify moving its capacity through that operator.  Hence deletion of
\(TOF_3\) is not proved and the unconditional reading must be
withdrawn.
\end{proof}

\section{Exact temporal split of the strong overlap}
\label{sec:V21-overlap-split}

Define
\begin{align}
 \theta_{ac}^{<}
 &:=
 \theta_{ac}(E_3\cap\{e:e<r\}),\notag\\
 \theta_{ac}^{=}
 &:=
 \theta_{ac}(E_3\cap\{r\}),\notag\\
 \theta_{ac}^{>}
 &:=
 \theta_{ac}(E_3\cap\{e:e>r\}).
\label{eq:V21-three-shares}
\end{align}

\begin{lemma}[Exact three-share identity]
\label{lem:V21-three-share}
\begin{equation}
 \boxed{
 \theta_{ac}(E_3)
 =
 \theta_{ac}^{<}
 +\theta_{ac}^{=}
 +\theta_{ac}^{>}.}
\label{eq:V21-three-share}
\end{equation}
All three terms are nonnegative.  If
\(\theta_{ac}(E_3)\ge\eta/4\), then
\begin{equation}
 \max\{\theta_{ac}^{<},\theta_{ac}^{=},\theta_{ac}^{>}\}
 \ge{\eta\over12}.
\label{eq:V21-one-share}
\end{equation}
\end{lemma}

\begin{proof}
The three coordinate sets in \eqref{eq:V21-three-shares} are disjoint
and partition \(E_3\).  Additivity of the nonnegative overlap sum gives
\eqref{eq:V21-three-share}.  If all three terms were less than
\(\eta/12\), their sum would be less than \(\eta/4\).
\end{proof}

\section{Balanced-partner localization in all three sectors}
\label{sec:V21-partner-localization}

Retain the balanced hypothesis
\(\Xi_c\ge\lambda M\).

\begin{lemma}[Cofinal partner in the dominant temporal sector]
\label{lem:V21-three-sector-localization}
If one of the three shares in
\eqref{eq:V21-three-shares} is at least \(\beta>0\), then:
\begin{enumerate}[label=\textup{(\roman*)},leftmargin=3.5em]
\item for the prefix share, some
      \(e_<\in E_3\), \(e_<<r\), satisfies
      \(a_{c,e_<}\ge(\beta\lambda/3)M\);
\item for the joining share,
      \(r\in E_3\) and
      \(a_{c,r}\ge\beta\lambda M\); or
\item for the suffix share, some
      \(e_>\in E_3\), \(e_>>r\), satisfies
      \(a_{c,e_>}\ge(\beta\lambda/3)M\).
\end{enumerate}
\end{lemma}

\begin{proof}
For any of the three coordinate sets \(D\), the overlap lower bound
gives
\[
 \sum_{e\in D}a_{a,e}a_{c,e}
 \ge\beta M\Xi_c.
\]
Since \(a_{a,e}\le M\),
\[
 \sum_{e\in D}a_{c,e}
 \ge\beta\Xi_c
 \ge\beta\lambda M.
\]
The prefix and suffix sets have cardinality at most three, giving
\textup{(i)} and \textup{(iii)} by averaging.  The joining set is
either empty or the singleton \(\{r\}\), giving \textup{(ii)}.
\end{proof}

\begin{corollary}[Exact cofinal temporal alternative]
\label{cor:V21-cofinal-alternative}
If
\(\theta_{ac}(E_3)\ge\eta/4\) and
\(\Xi_c\ge\lambda M\), then a partner coordinate of size at least
\((\eta\lambda/36)M\) occurs in the prefix, at the joining coordinate,
or in the suffix.  At the joining coordinate the stronger lower bound
\((\eta\lambda/12)M\) holds.
\end{corollary}

\begin{proof}
Use \Cref{lem:V21-three-share} with
\(\beta=\eta/12\), followed by
\Cref{lem:V21-three-sector-localization}.
\end{proof}

\section{The all-prefix subbranch}
\label{sec:V21-all-prefix}

\begin{proposition}[Automatic temporal part of \(TOF_3\)]
\label{prop:V21-all-prefix}
If the temporal class is \((p,j,s)=(3,0,0)\), then all three extracted
slot coordinates occur below \(r\).  In a pair-crossing \(3+2\)
prefix, their literal prefix maps factor between \(B_a\) and \(B_c\).
Thus the temporal placement and disjoint-row requirements in
\Cref{def:V21-TOF} hold.  Only the complete two-orientation outer
packet and partner certification remain to be checked.
\end{proposition}

\begin{proof}
The temporal assertion is the definition of the class
\((3,0,0)\).  Below \(r\), the root prefix is exactly
\[
 \mathcal K_{B_a}^{<r,\mathrm{conn}}
 \otimes
 \mathcal K_{B_c}^{<r,\mathrm{conn}},
\]
so literal maps are separated by root row block.  This gives the
claimed tensor placement.  It does not by itself prove the square cards
inside either connected factor, which are the remaining certificate
requirements.
\end{proof}

\begin{firewall}[Mixed temporal classes remain noncentral]
\label{fire:V21-mixed-classes}
For \(p=1\) or \(p=2\), some extracted heat responses occur at or after
the root joining coordinate.  The prefix tensor product alone does not
place those responses in the outer packet of
\Cref{def:V21-TOF}.  They require a new root/suffix two-sided square
factorization.  V20 does not close these mixed classes.
\end{firewall}

\section{Correct residual after V21}
\label{sec:V21-residual}

\begin{theorem}[Balanced pair-crossing temporal residual]
\label{thm:V21-residual}
In a balanced pair-crossing \(3+2\), saturated \(q=3\)
counterpacket, at least one of:
\begin{enumerate}[label=\textup{(\alph*)},leftmargin=3.5em]
\item the prefix share is bounded below but the transverse partner or
      \(TOF_3\) certificate fails;
\item the joining-coordinate share is bounded below and the complete
      root packet has no proved fourth response;
\item the suffix share is bounded below and the complete suffix packet
      has no proved fourth response; or
\item the assumed strong overlap or partner comparability degenerates.
\end{enumerate}
If the temporal class is all-prefix and the operator certificates hold,
the fibre closes by \Cref{thm:V21-corrected-transverse}.
\end{theorem}

\begin{proof}
When the overlap and comparability remain fixed, one of the three
shares is bounded below by
\Cref{lem:V21-three-share}.  Its cofinal partner is supplied by
\Cref{lem:V21-three-sector-localization}.  In the prefix case, closure
follows if both missing certificates hold.  Failure therefore gives
\textup{(a)}.  The other two temporal sectors are exactly
\textup{(b)} and \textup{(c)}.  Loss of either scalar premise gives
\textup{(d)}.  The all-prefix conclusion uses
\Cref{prop:V21-all-prefix,thm:V21-corrected-transverse}.
\end{proof}

\begin{assessment}[Upstream target after the correction]
\label{ass:V21-next}
The next proof should attack the joining-coordinate sector first.  It
has one physical coordinate, both strong endpoints are cofinal there,
and the literal root operator is a complete quotient covariance
between the \(3\)-block and \(2\)-block products.  V13 shows that a
bare covariance gives only one half-power, so the required question is
whether its two block inputs provide a transverse two-row factorization
inside both square orientations.  If not, the spin-one lower-output
family should be propagated through the complete root packet as a
source-level countertest.
\end{assessment}

\section{V21 result ledger and V22 acquisition order}
\label{sec:V21-ledger}

\begin{theorem}[Fourth-series V21 result ledger]
\label{thm:V21-results}
Relative to V1--V20, V21 proves:
\begin{enumerate}[label=\textup{(V21.\arabic*)},leftmargin=6.3em]
\item the exact transverse outer-factorization certificate missing
      from the V20 closure;
\item the corrected conditional transverse theorem;
\item withdrawal of the unjustified unconditional V20 reading;
\item the exact prefix--joining--suffix overlap identity;
\item localization of a balanced strong partner in all three temporal
      sectors with explicit constants;
\item the automatic temporal factorization in the all-prefix class;
      and
\item the corrected exhaustive residual for balanced pair-crossing
      \(3+2\) counterpackets.
\end{enumerate}
\end{theorem}

\begin{proof}
These are
\Cref{def:V21-TOF,
thm:V21-corrected-transverse,
corr:V21-V20,
lem:V21-three-share,
lem:V21-three-sector-localization,
cor:V21-cofinal-alternative,
prop:V21-all-prefix,
thm:V21-residual}.
\end{proof}

\begin{programme}[V22: complete root joining packet]
\label{prog:V21-V22}
\begin{enumerate}[leftmargin=3em]
\item Fix the joining-dominant branch
      \(\theta_{ac}^{=}\ge\eta/12\).
\item Write the complete root covariance between the \(3\)-block and
      \(2\)-block products before output resolution.
\item Retain the two block-input factors and test whether they are a
      genuine disjoint-row sandwich in both square orientations.
\item Include the payer on a fixed nonzero root output.
\item If the factorization fails, evaluate the V13 spin-one channel
      after the two complete prefix inputs have been inserted.
\item Treat suffix dominance only after the root packet is classified.
\end{enumerate}
\end{programme}

\begin{firewall}[Claim boundary after fourth-series V21]
\label{fire:V21-claim-boundary}
V21 corrects the V20 source-order gap and proves exact temporal
localization.  It does not prove \(TOF_3\) in mixed classes, classify
the joining or suffix packets, close the full \(3+2\) source, or close
\(q=3\).  The remaining \(q=3\) alternatives, \(q=0,1,2\) cards,
global quintic or all-order MTA, WUFC, CPM, cutoff removal, reflection
positivity, Osterwalder--Schrader reconstruction, construction of the
continuum Yang--Mills measure, and a uniform positive continuum
spectral gap remain open.  The full Yang--Mills mass gap has not been
proved.
\end{firewall}

\paragraph{Parent-version record.}
V21 was formed from
\texttt{Renewal\_Yang\_Mills\_Mass\_Gap\_Research\_Notes\_V20.tex}.
The V20 parent had \(10\,126\) logical source lines,
\(356\,849\) bytes, and SHA--256
\[
 \texttt{862f59b0ff981f3d6ddd128d12d01fd808651cf931675d64e8957dd3d9eb56dd}.
\]
The next compulsory companion audit is V25.

% =============================================================================
% END APPENDED FOURTH-SERIES V21 MATERIAL
% =============================================================================

\clearpage
% =============================================================================
% BEGIN APPENDED FOURTH-SERIES V22 MATERIAL
% =============================================================================

\chapter{V22: The complete joining covariance has only one root
response}
\label{ch:V22}

\section{Purpose and point of entry}
\label{sec:V22-purpose}

V21 reduced the balanced pair-crossing \(3+2\) fibre to three temporal
sectors and identified the joining-dominant branch
\[
 \theta_{ac}^{=}\ge {\eta\over12}.
\]
At its joining coordinate \(r\), both strong rows are cofinal.  The
remaining question is operator-theoretic: can the two quotient-block
inputs of the complete root covariance be charged as two independent
row-block responses?  This chapter answers that question negatively
for the natural positive-envelope route.  The fixed spin-one channel
of V13 survives in the \emph{complete}, unresolved root covariance and
has only the ordinary \(d_j^{-1}\asymp M_j^{-1/2}\) square scale.
Consequently, the fact that two heavy rows meet at \(r\) does not
create two inverse half-powers there.

\section{The unresolved quotient-block covariance}
\label{sec:V22-root-spectrum}

Let the two root blocks be \(B\) and \(C\).  After their complete
prefix carriers have been evaluated, fix irreducible block-output
modules \(V_\lambda\) and \(V_\mu\).  Write
\(\Pi_{\lambda,\mu}^{\nu}\) for the complete isotypic projection onto
all copies of \(V_\nu\) in \(V_\lambda\otimes V_\mu\), and put
\[
 q_\alpha(s):=e^{-sC_\alpha},
\]
where \(C_\alpha\) is the Casimir eigenvalue of \(V_\alpha\).

\begin{lemma}[Complete root covariance before output resolution]
\label{lem:V22-root-spectrum}
On \(V_\lambda\otimes V_\mu\), the quotient-order-two joining
covariance is
\begin{equation}
 \boxed{
 \mathbb D_{\lambda,\mu}(s)
 =
 \sum_{\nu\subset\lambda\otimes\mu}
 \bigl(q_\nu(s)-q_\lambda(s)q_\mu(s)\bigr)
 \Pi_{\lambda,\mu}^{\nu}.}
\label{eq:V22-root-spectrum}
\end{equation}
The sum includes the full multiplicity space of every output type.
In particular, no output projection has been inserted into the
source word in \eqref{eq:V22-root-spectrum}.
\end{lemma}

\begin{proof}
The common-variable term is the heat functional calculus of the
diagonal Casimir on \(V_\lambda\otimes V_\mu\); hence it has
eigenvalue \(q_\nu(s)\) on the complete \(\nu\)-isotype.  In the
cut-replicated term the two block variables are independent.  Schur
centrality gives the scalar
\(q_\lambda(s)q_\mu(s)\) on the entire tensor product, including every
multiplicity copy.  Subtraction gives
\eqref{eq:V22-root-spectrum}.  Because the complete isotypic
projections sum to the identity, the formula is the unresolved
operator itself rather than a selected-output replacement.
\end{proof}

\begin{corollary}[The \(SU(2)\) spin-one block is present in the
complete root]
\label{cor:V22-spin-one-present}
For \(G=SU(2)\), \(\lambda=\mu=j\), and \(j\ge1\),
\[
 \mathbb D_{j,j}(s)
 =
 \sum_{K=0}^{2j}
 \bigl(e^{-sK(K+1)}-e^{-2sj(j+1)}\bigr)P_{j,K}.
\]
Its complete square therefore satisfies, on the product vector
\(\phi_j=|j,j\rangle\otimes|j,-j\rangle\),
\begin{equation}
 \langle\phi_j,\mathbb D_{j,j}(s)^2\phi_j\rangle
 \ge
 \Delta_j(s)^2p_{j,1}
 \ge c_s d_j^{-1}
\label{eq:V22-complete-root-lower}
\end{equation}
for all sufficiently large \(j\).
\end{corollary}

\begin{proof}
The spectral formula is
\Cref{lem:V22-root-spectrum} with \(C_K=K(K+1)\).  All summands in
the spectral resolution of the square are nonnegative, so retaining
only \(K=1\) gives the first inequality.  The second is
\Cref{cor:V13-spin-one-coefficient,cor:V13-pj1-scale}.  Notice that
the proof tests the complete operator; it does not replace it by
\(\mathbb D_{j,j}P_{j,1}\).
\end{proof}

\section{Failure of a two-block positive response envelope}
\label{sec:V22-envelope}

The following formulation isolates exactly the inference that would
be needed to turn the two heavy root inputs into two separate
half-powers.

\begin{definition}[Root two-block response envelope]
\label{def:V22-RBRE}
A family \(T_j\) on \(V_j\otimes V_j\) has a root two-block response
envelope if there are positive operators \(X_j\) on the first block
and \(Y_j\) on the second block, and constants \(A,B_1,B_2\)
independent of \(j\), such that
\begin{align}
 T_j^*T_j
 &\preccurlyeq A\,X_j\otimes Y_j,
\label{eq:V22-RBRE-order}\\
 \kappa_B(X_j)&\le B_1d_j^{-1},
 \qquad
 \kappa_C(Y_j)\le B_2d_j^{-1}.
\label{eq:V22-RBRE-cards}
\end{align}
Here the capacities use the inherited row-product cuts inside the two
disjoint root blocks.  The definition is also applied after local
compression to fixed block-output modules.
\end{definition}

\begin{theorem}[No root two-block response envelope for the complete
covariance]
\label{thm:V22-no-RBRE}
Neither \(\mathbb D_{j,j}(s)\) nor its once-paid version has a root
two-block response envelope with constants uniform in \(j\).  The
same conclusion holds in the opposite square orientation.
\end{theorem}

\begin{proof}
Assume first that
\eqref{eq:V22-RBRE-order}--\eqref{eq:V22-RBRE-cards} hold for
\(T_j=\mathbb D_{j,j}(s)\).  Monotonicity of product-state capacity
and the exact disjoint-block product rule
\Cref{lem:V20-capacity-product} give
\[
 \kappa_{B\sqcup C}\bigl(\mathbb D_{j,j}(s)^2\bigr)
 \le
 AB_1B_2d_j^{-2}.
\]
On the other hand, \(\phi_j\) is a block-product vector, and
\eqref{eq:V22-complete-root-lower} gives
\[
 \kappa_{B\sqcup C}\bigl(\mathbb D_{j,j}(s)^2\bigr)
 \ge c_sd_j^{-1}.
\]
These inequalities are incompatible as \(j\to\infty\).

For the joining-paid operator
\(\mathbb R_s\mathbb D_{j,j}(s)\), the spin-one term in the same
complete square is multiplied by the fixed strictly positive number
\(r_1(s)^2\); use \Cref{thm:V13-paid-sharp}.  Both operators are
self-adjoint and commute with the payer on their common spectral
resolution, so the opposite square orientation has the same
contradiction.
\end{proof}

\begin{corollary}[Two cofinal endpoints do not give two root
half-powers]
\label{cor:V22-two-endpoints-one-response}
There is no uniform argument which, solely from the facts that the
two strong endpoints \(a\in B\) and \(c\in C\) are cofinal at the
joining coordinate and that the root carrier is a complete
covariance, assigns one \(M^{-1/2}\) square response to each endpoint.
On the equal-spin family the complete root square has the larger
scale
\[
 d_j^{-1}\asymp M_j^{-1/2},
\]
not \(d_j^{-2}\asymp M_j^{-1}\), paid or unpaid.
\end{corollary}

\begin{proof}
Put the varying representation \(V_j\) on one row in each root block
and use fixed, in particular trivial where allowed, representations
on the other root-coordinate rows.  Each block then has output
\(V_j\).  Charging one inverse half-power to each disjoint block
would give a root two-block response envelope and hence a
\(Cd_j^{-2}\) square bound.  This is excluded by
\Cref{thm:V22-no-RBRE}.  Since
\(M_j=1+j(j+1)\asymp d_j^2\), the displayed mass scales follow.
\end{proof}

\begin{firewall}[What the envelope no-go does and does not exclude]
\label{fire:V22-envelope-scope}
\Cref{thm:V22-no-RBRE} excludes a positive product majorant carrying
one response card on each root block.  It also excludes any proposed
two-sided sandwich theorem whose conclusion would imply the same
\(Cd_j^{-2}\) capacity bound.  It does not exclude an estimate using
an additional, genuinely \(j\)-dependent prefix or suffix factor.
Such a factor must be displayed in the exact source order; it cannot
be inferred from the two endpoints of the covariance alone.
\end{firewall}

\section{A complete joining--suffix counterpacket}
\label{sec:V22-complete-packet}

To avoid the fine-output problem completely, define the full
common-variable heat moment on \(V_j\otimes V_j\) by
\begin{equation}
 \mathbb A_j(s)
 :=
 \sum_{K=0}^{2j}q_K(s)P_{j,K}.
\label{eq:V22-complete-heat}
\end{equation}
On three distinct physical coordinates, the first designated as the
joining coordinate and the other two as suffix coordinates, put
\begin{equation}
 \mathbb Z_j^{\rm comp}
 :=
 \mathbb D_{j,j}(s)
 \otimes \mathbb A_j(s)
 \otimes \mathbb A_j(s).
\label{eq:V22-complete-packet}
\end{equation}

\begin{proposition}[Complete three-coordinate lower packet]
\label{prop:V22-complete-packet}
For all sufficiently large \(j\),
\begin{equation}
 \kappa_\otimes\!\left(
  (\mathbb Z_j^{\rm comp})^*\mathbb Z_j^{\rm comp}
 \right)
 \ge c_sd_j^{-3}
 \asymp c_sM_j^{-3/2}.
\label{eq:V22-complete-packet-lower}
\end{equation}
The conclusion remains true when the sole output payer is placed on
any one of the three displayed coordinates.
\end{proposition}

\begin{proof}
Test on the row-product vector \(\phi_j^{\otimes3}\), with tensor
factors regrouped by row.  By
\eqref{eq:V22-complete-root-lower}, the joining-coordinate
expectation is at least \(\Delta_j(s)^2p_{j,1}\).  The complete heat
square has nonnegative spectral resolution
\[
 \mathbb A_j(s)^2
 =
 \sum_{K=0}^{2j}q_K(s)^2P_{j,K};
\]
hence each suffix expectation is at least \(q_1(s)^2p_{j,1}\).
Multiplication gives
\[
 \langle\phi_j^{\otimes3},
  (\mathbb Z_j^{\rm comp})^*\mathbb Z_j^{\rm comp}
  \phi_j^{\otimes3}\rangle
 \ge
 \Delta_j(s)^2q_1(s)^4p_{j,1}^3
 \ge c_sd_j^{-3}.
\]
If the payer is on the joining coordinate, retain the spin-one term
and multiply by \(r_1(s)^2\).  If it is on a suffix coordinate, that
coordinate's retained spin-one term is likewise multiplied by
\(r_1(s)^2\).  Thus every displayed payer placement changes only the
positive constant.
\end{proof}

\begin{proposition}[Insertion of a fixed nonzero complete prefix]
\label{prop:V22-fixed-prefix}
Let \(\mathbb F\) be a fixed nonzero complete prefix carrier, acting at
physical coordinates disjoint from the three coordinates in
\eqref{eq:V22-complete-packet}, with all its representation labels
independent of \(j\).  Then
\[
 \mathbb W_j
 :=
 \mathbb F\otimes\mathbb Z_j^{\rm comp}
\]
satisfies
\begin{equation}
 \kappa_\otimes(\mathbb W_j^*\mathbb W_j)
 \ge c_{\mathbb F,s}d_j^{-3}.
\label{eq:V22-prefix-packet-lower}
\end{equation}
The same statement holds for a fixed nonzero paid prefix carrier.
\end{proposition}

\begin{proof}
By \Cref{lem:V9-product-detection}, there is a row-product unit vector
\(v\) in the fixed prefix factors such that
\[
 c_{\mathbb F}
 :=
 \langle v,\mathbb F^*\mathbb Fv\rangle>0.
\]
Tensor \(v\) with the product vector used in
\Cref{prop:V22-complete-packet}.  Operators at distinct physical
coordinates tensor exactly before the harmless canonical
rebracketing by rows.  The square expectation is therefore
\(c_{\mathbb F}\) times the expectation in
\eqref{eq:V22-complete-packet-lower}.  This proves the claim.  The
paid case is identical once the fixed paid prefix carrier is assumed
nonzero.
\end{proof}

\begin{corollary}[The \((0,1,2)\) temporal pattern has no automatic
fourth response]
\label{cor:V22-012-no-automatic}
Consider a pair-crossing \(3+2\) fibre whose fixed prefix carrier is
nonzero.  Place one varying \(V_j\) row in each root block at the
joining coordinate and at two suffix coordinates, keeping the other
labels fixed.  The two varying rows then have three strongly
overlapping cofinal coordinates, but the complete
prefix--joining--suffix packet has square scale at least
\(c d_j^{-3}\asymp cM_j^{-3/2}\).  Thus the temporal pattern
\((p,j,s)=(0,1,2)\) cannot be closed at scale \(CM_j^{-2}\) by the
three heat responses plus an alleged second response inside the root
covariance.
\end{corollary}

\begin{proof}
The mass of each varying row is comparable to \(3M_j\), and their
overlap at the three varying coordinates has a fixed positive
fraction.  The complete operator lower bound is
\Cref{prop:V22-fixed-prefix}.  A fourth inverse half-power would imply
a \(CM_j^{-2}\) bound for this packet, whereas
\(M_j^{-3/2}/M_j^{-2}=M_j^{1/2}\to\infty\).
\end{proof}

\begin{firewall}[Not yet a counterexample to the normalized source
theorem]
\label{fire:V22-source-scope}
The preceding packet is a complete operator countertest to the
proposed \emph{analytic response mechanism}.  It is not yet a
counterexample to XQTRSC\(_5\) or to the full quintic estimate.  Such a
counterexample requires an exact admissible five-row source fibre,
its normalization, the chosen prefix-tree rescue weight, all
spectator factors, and proof that no source sum cancels the packet.
V22 has deliberately not replaced that scalar and combinatorial audit
by an inference from the local operator lower bound.
\end{firewall}

\section{Corrected residual after the root test}
\label{sec:V22-residual}

\begin{theorem}[Joining-sector classification]
\label{thm:V22-joining-classification}
In the joining-dominant alternative of
\Cref{thm:V21-residual}, the root covariance supplies at most the one
already-counted heat-scale response unless one proves an additional
\(j\)-dependent prefix or suffix estimate in exact source order.
In particular:
\begin{enumerate}[label=\textup{(\alph*)},leftmargin=3.5em]
\item the two cofinal endpoints at the root cannot be separated into
      two positive block response cards;
\item the once-paid output does not repair that failure;
\item fixed nonzero prefix insertion does not repair the
      \((0,1,2)\) three-coordinate packet; and
\item any valid closure must come from source incidence, cancellation,
      or a revised normalized target, rather than from bare root
      covariance centering.
\end{enumerate}
\end{theorem}

\begin{proof}
Items \textup{(a)} and \textup{(b)} are
\Cref{thm:V22-no-RBRE,cor:V22-two-endpoints-one-response}.
Item \textup{(c)} is
\Cref{prop:V22-fixed-prefix,cor:V22-012-no-automatic}.  These exhaust
the local root and fixed-prefix response mechanisms considered in the
V21 programme.  Therefore any remaining inverse half-power must
depend on further \(j\)-dependent source structure or on a change in
the target comparison, which is \textup{(d)}.
\end{proof}

\begin{assessment}[Upstream move forced by V22]
\label{ass:V22-upstream}
The root operator problem is now sharp: the complete covariance
retains the V13 spin-one obstruction, even after an arbitrary fixed
nonzero prefix.  Repeating local covariance estimates would be
circular.  The next note should therefore go upstream to the exact
normalized \(3+2\) source atlas.  It should instantiate the
\((0,1,2)\) packet on five rows, compute every row normalization and
prefix-tree rescue denominator, and decide whether the desired
XQTRSC\(_5\) right-hand side is genuinely \(M^{-2}\) on that family.
If it is, the source word and cancellation signs must be evaluated
exactly.  If it is not, the response target rather than the root
analysis requires correction.
\end{assessment}

\section{V22 result ledger and V23 acquisition order}
\label{sec:V22-ledger}

\begin{theorem}[Fourth-series V22 result ledger]
\label{thm:V22-results}
Relative to V1--V21, V22 proves:
\begin{enumerate}[label=\textup{(V22.\arabic*)},leftmargin=6.3em]
\item the complete root covariance spectrum with all multiplicities
      retained;
\item persistence of the spin-one lower bound in the unresolved
      operator;
\item failure of every uniform two-block positive response envelope,
      paid and unpaid;
\item the complete, unresolved joining--suffix
      \(d_j^{-3}\asymp M_j^{-3/2}\) counterpacket;
\item stability of that packet under any fixed nonzero complete
      prefix insertion; and
\item the corrected joining-sector residual and the forced upstream
      normalization audit.
\end{enumerate}
\end{theorem}

\begin{proof}
These are
\Cref{lem:V22-root-spectrum,
cor:V22-spin-one-present,
thm:V22-no-RBRE,
prop:V22-complete-packet,
prop:V22-fixed-prefix,
cor:V22-012-no-automatic,
thm:V22-joining-classification,
ass:V22-upstream}.
\end{proof}

\begin{programme}[V23: exact normalized \(3+2\) packet]
\label{prog:V22-V23}
\begin{enumerate}[leftmargin=3em]
\item Extract the exact XQTRSC\(_5\) normalization and rescue weight
      for one pair-crossing \(3+2\) prefix tree.
\item Build a legal five-row \(SU(2)\) family with fixed nonzero
      \(3+2\) prefix, one joining coordinate, and two suffix
      coordinates carrying the equal-spin packet.
\item Compute all five row masses, the strong-overlap stock, and every
      tree-edge denominator without asymptotic placeholders.
\item Evaluate the complete source signs on the spin-one test vector;
      do not insert a fine output projection.
\item Compare the resulting lower scale with the exact normalized
      right-hand side.
\item If cancellation kills the packet, identify its first exact
      source ancestor and descend there; if it survives, record the
      necessary repair to the target atlas.
\end{enumerate}
\end{programme}

\begin{firewall}[Claim boundary after fourth-series V22]
\label{fire:V22-claim-boundary}
V22 closes the proposed root two-endpoint factorization route and
constructs a complete operator counterpacket for the
\((0,1,2)\) response mechanism.  It does not yet compare that packet
with the exact normalized five-row source theorem, prove a nonzero
uncancelled quintic source counterexample, close the all-prefix or
other mixed temporal classes, or close \(q=3\).  The remaining
\(q=3\) alternatives, \(q=0,1,2\) cards, global quintic or all-order
MTA, WUFC, CPM, cutoff removal, reflection positivity,
Osterwalder--Schrader reconstruction, construction of the continuum
Yang--Mills measure, and a uniform positive continuum spectral gap
remain open.  The full Yang--Mills mass gap has not been proved.
\end{firewall}

\paragraph{Parent-version record.}
V22 was formed from
\texttt{Renewal\_Yang\_Mills\_Mass\_Gap\_Research\_Notes\_V21.tex}.
The V21 parent had \(10\,474\) logical source lines,
\(369\,658\) bytes, and SHA--256
\[
 \texttt{469246d0b670bf87f03da957051f4ac03e5ff79451bd937923f036561bd51871}.
\]
The next compulsory companion audit is V25.

% =============================================================================
% END APPENDED FOURTH-SERIES V22 MATERIAL
% =============================================================================

\clearpage
% =============================================================================
% BEGIN APPENDED FOURTH-SERIES V23 MATERIAL
% =============================================================================

\chapter{V23: Exact normalization of a five-row
\texorpdfstring{\((0,1,2)\)}{(0,1,2)} packet}
\label{ch:V23}

\section{Purpose}
\label{sec:V23-purpose}

V22 proves that the joining covariance in the mixed temporal pattern
\((p,j,s)=(0,1,2)\) supplies only one square half-power.  Its final
firewall left open whether row normalization and the exchange-rescue
atlas nevertheless make the complete five-row source estimate easier.
This chapter performs that scalar audit exactly.

There is a legal pair-support geometry in which the exchange atlas has
only one nonzero labelled spanning tree.  If the maximal row is the
centre of the three-block prefix tree, its degree excess creates
exactly the missing \(M^{-2}\) oriented-square demand.  The normalized
complete operator packet from V22 is then larger than the exchange
target by a factor comparable to \(d_j\asymp M^{1/2}\).  Thus
normalization does not remove the obstruction.  The only gate not
settled here is quantitative survival through the chosen physical
output-dual block.

\section{A legal six-coordinate five-row skeleton}
\label{sec:V23-skeleton}

Label the rows
\[
 \{a,b,c,d,e\}
\]
and order six physical coordinates by
\[
 u<v<w<r<s_1<s_2.
\]
Fix a nontrivial \(SU(2)\) spin \(\ell\), independent of \(j\), and
write
\[
 B:=1+\ell(\ell+1),
 \qquad
 A_j:=1+j(j+1),
 \qquad
 d_j:=2j+1.
\]
Use the following active representation supports:
\begin{center}
\begin{tabular}{@{}ccl@{}}
\toprule
coordinate & active rows & spin\\
\midrule
\(u\) & \(a,b\) & \(\ell,\ell\)\\
\(v\) & \(a,d\) & \(\ell,\ell\)\\
\(w\) & \(c,e\) & \(\ell,\ell\)\\
\(r\) & \(a,c\) & \(j,j\)\\
\(s_1\) & \(a,c\) & \(j,j\)\\
\(s_2\) & \(a,c\) & \(j,j\)\\
\bottomrule
\end{tabular}
\end{center}
All undisplayed row-coordinate representations are trivial and carry
zero active stock.  The root prefix partition is
\begin{equation}
 \rho=\{\{a,b,d\},\{c,e\}\}.
\label{eq:V23-rho}
\end{equation}
Choose the literal lifted tree
\begin{equation}
 T_*=\{ab,ad,ce,ac\},
 \qquad
 P=\{ab,ad,ce\},
 \qquad
 Q=\{ac\}.
\label{eq:V23-tree}
\end{equation}
The edge \(ac\) is joined locally at \(r\).

\begin{lemma}[First-connectivity legality]
\label{lem:V23-legality}
The support in \Cref{sec:V23-skeleton} has cumulative partition
\(\rho\) immediately below \(r\), and \(r\) is its unique first
coordinate of full five-row connectivity.  The prefix forest and
quotient lift are exactly those in \eqref{eq:V23-tree}.
\end{lemma}

\begin{proof}
At \(u\), rows \(a,b\) become connected.  At \(v\), row \(d\) joins
that component, producing \(\{a,b,d\}\).  At \(w\), rows \(c,e\)
form the other component.  Thus the cumulative partition below \(r\)
is \eqref{eq:V23-rho}.  The coordinate \(r\) contains \(a\) and \(c\),
one in each block, so it connects the two blocks.  No earlier
coordinate crosses them.  The three prefix connections are
\(ab,ad,ce\), and the quotient edge \(ac\) lifts the two-block
quotient tree.  This proves every assertion.
\end{proof}

\section{Exact masses and stocks}
\label{sec:V23-stocks}

\begin{lemma}[Exact row masses]
\label{lem:V23-masses}
The five row masses are
\begin{equation}
 \boxed{
 \Xi_a=3A_j+2B,\qquad
 \Xi_c=3A_j+B,\qquad
 \Xi_b=\Xi_d=\Xi_e=B.}
\label{eq:V23-masses}
\end{equation}
In particular,
\[
 M:=\max_i\Xi_i=\Xi_a,
 \qquad
 \Xi_c/M\longrightarrow1.
\]
\end{lemma}

\begin{proof}
Row \(a\) occurs with spin \(j\) at \(r,s_1,s_2\) and with spin
\(\ell\) at \(u,v\).  Row \(c\) has the same three spin-\(j\)
occurrences and one spin-\(\ell\) occurrence at \(w\).  Each of
\(b,d,e\) has one spin-\(\ell\) occurrence.  Summing the active
weights gives \eqref{eq:V23-masses}.  The remaining statements are
immediate.
\end{proof}

\begin{lemma}[Exact nonzero pair stocks]
\label{lem:V23-pair-stocks}
The only nonzero pair stocks are the four labels in \(T_*\).  Their
relevant values are
\begin{align}
 \theta_{ab}^{<r}=\theta_{ab}
 &= {B\over\Xi_a},
 &
 \theta_{ad}^{<r}=\theta_{ad}
 &= {B\over\Xi_a},
\label{eq:V23-prefix-a-stocks}\\
 \theta_{ce}^{<r}=\theta_{ce}
 &= {B\over\Xi_c},
 &
 \theta_{ac,r}
 &= {A_j^2\over\Xi_a\Xi_c},
\label{eq:V23-prefix-c-root-stock}\\
 \theta_{ac}
 &= {3A_j^2\over\Xi_a\Xi_c}.
\label{eq:V23-global-strong}
\end{align}
Consequently \(ac\) is a balanced strong pair, its entire stock is
carried by the three coordinates \(r,s_1,s_2\), and each of those
coordinates carries asymptotic share \(1/3\) of both heavy rows.
\end{lemma}

\begin{proof}
Two distinct rows overlap only where they are displayed together in
the support table.  At \(u,v,w\), the overlap numerator is \(B^2\);
division by the corresponding row masses and cancellation of the
fixed spectator mass \(B\) gives
\eqref{eq:V23-prefix-a-stocks}--the first part of
\eqref{eq:V23-prefix-c-root-stock}.  At each of
\(r,s_1,s_2\), the \(ac\) overlap numerator is \(A_j^2\).  This gives
the local and global formulas.  Using \eqref{eq:V23-masses},
\[
 \theta_{ac}\longrightarrow {1\over3},
 \qquad
 {\Xi_c\over\Xi_a}\longrightarrow1,
 \qquad
 {A_j\over\Xi_a},{A_j\over\Xi_c}\longrightarrow {1\over3}.
\]
These are the claimed strong-overlap and cofinality statements.
\end{proof}

\section{The exchange-rescue weight is exact}
\label{sec:V23-exchange}

\begin{theorem}[Exact exchange weight for the skeleton]
\label{thm:V23-exact-exchange}
For the literal tree \eqref{eq:V23-tree} and joining coordinate \(r\),
the exchange-rescue weight of archive f3 V98 is
\begin{equation}
 \boxed{
 \mathfrak x_{T_*,r}
 =
 {8B^3A_j^2\over\Xi_a^3\Xi_c^2}.}
\label{eq:V23-exact-exchange}
\end{equation}
Hence there are constants \(0<c_{\ell}\le C_{\ell}<\infty\) such
that
\begin{equation}
 c_{\ell}A_j^{-3}
 \le\mathfrak x_{T_*,r}
 \le C_{\ell}A_j^{-3}.
\label{eq:V23-exchange-scale}
\end{equation}
\end{theorem}

\begin{proof}
By \Cref{lem:V23-pair-stocks}, the support graph of the pair stocks is
the tree \(T_*\).  Therefore every labelled spanning tree
\(\tau\ne T_*\) in the exchange-rescue sum has a zero edge factor.
For \(\tau=T_*\), the nonempty local subset of
\(E(\tau)\cap Q=\{ac\}\) is forced to be \(L=\{ac\}\).  There are
exactly \(2^{|P|}=8\) choices of
\(A\subseteq P\).  A prefix edge has no occurrence at or above \(r\),
so its global and prefix stocks agree.  Thus all eight terms are
identical and
\[
 \mathfrak x_{T_*,r}
 =
 8\theta_{ab}\theta_{ad}\theta_{ce}\theta_{ac,r}.
\]
Insert
\eqref{eq:V23-prefix-a-stocks}--%
\eqref{eq:V23-prefix-c-root-stock} to obtain
\eqref{eq:V23-exact-exchange}.  Finally
\(\Xi_a,\Xi_c\asymp_\ell A_j\), which proves
\eqref{eq:V23-exchange-scale}.
\end{proof}

\begin{remark}[The degree excess is on the maximal row]
\label{rem:V23-degree-excess}
Inside the three-row prefix tree, \(a\) is the centre of the path
\(b-a-d\).  Its prefix degree is two, while every other row has prefix
degree one.  After the strong joining factor is treated at constant
scale, this single excess degree is exactly the factor \(A_j^{-1}\)
between the raw normalization \(A_j^{-2}\) and the rescue weight
\(A_j^{-3}\).
\end{remark}

\section{A fixed nonzero prefix carrier exists}
\label{sec:V23-prefix}

\begin{lemma}[Nonzero fixed prefix realization]
\label{lem:V23-nonzero-prefix}
The fixed-spin coordinates \(u,v,w\) may be resolved so that the
complete prefix factor
\[
 \mathbb F_\ell
 :=
 \mathcal K_{\{a,b,d\}}^{<r,\mathrm{conn}}
 \otimes
 \mathcal K_{\{c,e\}}^{<r,\mathrm{conn}}
\]
has a nonzero tree fibre with prefix edges \(ab,ad,ce\).  Its operator
and every nonzero singular value are independent of \(j\).
\end{lemma}

\begin{proof}
At \(u\), retain a nonzero spectral block of the equal-spin binary
covariance on \(a,b\).  It is nonzero because on its singlet block the
covariance multiplier is
\[
 1-e^{-2s\ell(\ell+1)}\ne0.
\]
At \(v\), the cumulative block \(\{a,b\}\) joins \(d\); because only
the \(a,d\) coordinate legs are nontrivial there, the same binary
covariance calculation provides a nonzero block with endpoint lift
\(ad\).  At \(w\), retain a nonzero binary covariance block on
\(c,e\).  The coefficient-one recursive first-connectivity formula
places these three complete blocks on distinct physical coordinates,
so their tensor product is nonzero.  All involved representations
are the fixed spin \(\ell\), proving independence of \(j\).
\end{proof}

\begin{firewall}[Fixed tree fibre, not a moment-term split]
\label{fire:V23-prefix-fibre}
\Cref{lem:V23-nonzero-prefix} resolves complete covariance output
blocks and an endpoint lift already allowed in the fixed finite
tree-fibre census.  It does not norm or retain one term of a
moment--cumulant difference.  Each binary covariance remains assembled.
\end{firewall}

\section{The normalized pre-dual deficit}
\label{sec:V23-normalized-deficit}

Put the sole payer on the joining coordinate.  Let
\(\widehat{\mathbb W}_j\) denote the resulting resolved word after the
five row normalizations and before the last physical output-dual
contraction.  On the fibre of
\Cref{lem:V23-nonzero-prefix}, its varying part is the complete packet
of V22:
\begin{equation}
 \widehat{\mathbb W}_j
 =
 {1\over\Xi_a\Xi_b\Xi_c\Xi_d\Xi_e}\,
 \mathbb F_\ell
 \otimes
 \bigl(\mathbb R_s\mathbb D_{j,j}(s)\bigr)
 \otimes
 \mathbb A_j(s)
 \otimes
 \mathbb A_j(s),
\label{eq:V23-predual-word}
\end{equation}
up to the canonical tensor rebracketing by rows.

\begin{theorem}[Exact normalized pre-dual lower bound]
\label{thm:V23-predual-lower}
There is \(c_{\ell,s}>0\) such that, for all sufficiently large \(j\),
\begin{equation}
 \kappa_\otimes\!\left(
  \widehat{\mathbb W}_j^*\widehat{\mathbb W}_j
 \right)
 \ge
 {c_{\ell,s}d_j^{-3}\over
  \Xi_a^2\Xi_c^2B^6}.
\label{eq:V23-predual-lower}
\end{equation}
Moreover,
\begin{equation}
 { \kappa_\otimes(
    \widehat{\mathbb W}_j^*\widehat{\mathbb W}_j)}
  { \mathfrak x_{T_*,r}^{\,2}}
 \ge
 c_{\ell,s}\,
 {d_j^{-3}\Xi_a^4\Xi_c^2\over A_j^4}
 \ge c'_{\ell,s}d_j.
\label{eq:V23-deficit-ratio}
\end{equation}
Thus the pre-dual word violates a uniform exchange-rescue square card
by precisely one inverse half-power.
\end{theorem}

\begin{proof}
The fixed prefix is nonzero, so
\Cref{prop:V22-fixed-prefix} applies to
\eqref{eq:V23-predual-word}, including the joining payer.  By
\eqref{eq:V23-masses},
\[
 (\Xi_a\Xi_b\Xi_c\Xi_d\Xi_e)^2
 =
 \Xi_a^2\Xi_c^2B^6.
\]
This proves \eqref{eq:V23-predual-lower}.

Square \eqref{eq:V23-exact-exchange} and divide the first lower bound
by it.  The fixed powers of \(B\) are absorbed into
\(c_{\ell,s}\), giving the first inequality in
\eqref{eq:V23-deficit-ratio}.  Since
\(\Xi_a,\Xi_c\asymp_\ell A_j\), its varying factor is
\[
 {A_j^2\over d_j^3}.
\]
The exact identities
\[
 A_j=1+j(j+1),
 \qquad d_j=2j+1
\]
give \(A_j\asymp d_j^2\), and hence
\(A_j^2d_j^{-3}\asymp d_j\).  This proves the second inequality.
\end{proof}

\begin{corollary}[Normalization does not close the V22 packet]
\label{cor:V23-normalization-no-close}
The fixed spectator masses and all five row denominators are already
included in \eqref{eq:V23-deficit-ratio}.  Therefore neither unequal
row masses nor the finite exchange enlargement supplies the missing
factor.  Any valid \(XQTRSC_5\) proof on this family must gain at least
\(Cd_j^{-1}\asymp CM^{-1/2}\) in the \emph{square} after
\eqref{eq:V23-predual-word}.
\end{corollary}

\begin{proof}
The exchange weight was evaluated as the entire nonnegative V98 sum,
not merely one of its terms, in
\Cref{thm:V23-exact-exchange}.  The ratio in
\eqref{eq:V23-deficit-ratio} tends to infinity like \(d_j\).
Consequently a subsequent operation must decrease the square capacity
of this packet by at least a constant multiple of \(d_j^{-1}\) if the
card is to hold.
\end{proof}

\section{The physical-output survival quotient}
\label{sec:V23-output-gate}

Let
\(\widetilde{\mathbb W}_{j,\boldsymbol D}\) be the fixed physical
output-dual block obtained from
\(\widehat{\mathbb W}_j\), with contraction family
\(\boldsymbol D\) in the notation of
\Cref{lem:V11-dual-contraction}.  Whenever the denominator is nonzero,
define the right-oriented survival quotient
\begin{equation}
 \sigma_j^{R}(\boldsymbol D)
 :=
 {\kappa_\otimes(
   \widetilde{\mathbb W}_{j,\boldsymbol D}^{*}
   \widetilde{\mathbb W}_{j,\boldsymbol D})\over
  \kappa_\otimes(
   \widehat{\mathbb W}_j^{*}\widehat{\mathbb W}_j)}.
\label{eq:V23-survival-R}
\end{equation}
Define \(\sigma_j^L(\boldsymbol D)\) analogously with the opposite
square.

\begin{theorem}[Exact output-survival alternative]
\label{thm:V23-output-alternative}
For the legal source skeleton above:
\begin{enumerate}[label=\textup{(\roman*)},leftmargin=3.5em]
\item if some admissible fixed output-dual family satisfies
      \(\liminf_j\sigma_j^R(\boldsymbol D)>0\) or
      \(\liminf_j\sigma_j^L(\boldsymbol D)>0\), then the corresponding
      \(XQTRSC_5\) orientation is false;
\item if \(XQTRSC_5\) holds uniformly, every admissible output-dual
      family which retains the V22 test packet must satisfy
\begin{equation}
 \sigma_j^R(\boldsymbol D)\le C_{\ell,s}d_j^{-1},
 \qquad
 \sigma_j^L(\boldsymbol D)\le C_{\ell,s}d_j^{-1}
\label{eq:V23-necessary-output-loss}
\end{equation}
      in the respective orientation.
\end{enumerate}
\end{theorem}

\begin{proof}
The \(XQTRSC_5\) right side is
\(C_{G,s}\mathfrak x_{T_*,r}^{\,2}\).  Multiply
\eqref{eq:V23-deficit-ratio} by
\(\sigma_j^R(\boldsymbol D)\).  If the latter has a fixed positive
lower bound, the ratio of the physical square to the target diverges
like \(d_j\), proving \textup{(i)}.  Conversely, boundedness of that
ratio forces
\(\sigma_j^R(\boldsymbol D)\le C d_j^{-1}\).  The left-oriented
argument is identical.
\end{proof}

\begin{proposition}[Coefficient-one source separation removes
inter-fibre cancellation]
\label{prop:V23-no-interfibre-cancellation}
The loss required by
\eqref{eq:V23-necessary-output-loss} cannot be attributed to
cancellation between different first-connectivity coordinates,
prefix partitions, lifted trees, endpoint lifts, or payer placements.
It must occur within the selected physical output-dual block of the
single resolved fibre.
\end{proposition}

\begin{proof}
The exact identity
\eqref{eq:V11-exact-source} assigns every connected coordinate word
to one unique pair \((r,\rho)\) with coefficient one.  The definition
of \(XQTRSC_5\) then fixes a prefix-tree fibre, quotient-tree fibre,
endpoint lifts, physical output block, and payer placement before the
square estimate is asserted.  The skeleton above fixes each of those
labels except the output-dual contraction now under examination.
Thus no other fibre is present on the left side of
\eqref{eq:V23-survival-R}; cancellation with such a fibre is
unavailable.
\end{proof}

\begin{firewall}[The remaining gap is quantitative]
\label{fire:V23-output-gap}
\Cref{lem:V11-dual-contraction,thm:V11-nonsymmetric-dual} give
dimension-free \emph{upper} bounds for a physical dual separator.
They do not give a lower bound for a chosen contraction family.
Therefore V23 does not assert that an admissible physical output block
has survival quotient bounded below.  It proves the sharper necessary
statement \eqref{eq:V23-necessary-output-loss}: saving the current
exchange card requires a full additional \(d_j^{-1}\) square loss
from the last output-dual step.
\end{firewall}

\begin{assessment}[The next calculation is finite and explicit]
\label{ass:V23-next}
No further scalar rescue or root covariance estimate can decide the
packet.  V24 should write the actual \(SU(2)\) output-dual contraction
on the retained spin-one root and suffix blocks and the fixed
spin-\(\ell\) prefix block.  It should compute the relevant
Clebsch--Gordan coefficient in both square orientations.  A lower
bound larger than \(Cd_j^{-1}\) disproves the present \(XQTRSC_5\)
interface; an exact \(O(d_j^{-1})\) factor identifies a genuine
physical-output fourth response.
\end{assessment}

\section{V23 result ledger and V24 acquisition order}
\label{sec:V23-ledger}

\begin{theorem}[Fourth-series V23 result ledger]
\label{thm:V23-results}
Relative to V1--V22, V23 proves:
\begin{enumerate}[label=\textup{(V23.\arabic*)},leftmargin=6.3em]
\item a legal five-row, six-coordinate pair-crossing \(3+2\) skeleton
      with temporal pattern \((0,1,2)\);
\item its exact row masses and complete pair-stock support;
\item the exact exchange-rescue weight, including all eight nonzero
      role assignments and no omitted trees;
\item existence of a fixed nonzero complete prefix-tree fibre;
\item the normalized pre-dual deficit
      \(\kappa/\mathfrak x^2\gtrsim d_j\);
\item the necessary \(d_j^{-1}\) physical-output square loss; and
\item removal of every inter-fibre cancellation explanation.
\end{enumerate}
\end{theorem}

\begin{proof}
These are
\Cref{lem:V23-legality,
lem:V23-masses,
lem:V23-pair-stocks,
thm:V23-exact-exchange,
lem:V23-nonzero-prefix,
thm:V23-predual-lower,
cor:V23-normalization-no-close,
thm:V23-output-alternative,
prop:V23-no-interfibre-cancellation}.
\end{proof}

\begin{programme}[V24: physical-output survival coefficient]
\label{prog:V23-V24}
\begin{enumerate}[leftmargin=3em]
\item Extract the actual rowwise output-dual map for the resolved
      \(3+2\) word; distinguish it from an arbitrary contraction.
\item Retain the root and two suffix total-spin-one blocks and the
      fixed nonzero prefix block.
\item Express the physical dual coefficient as an \(SU(2)\)
      recoupling coefficient, including all normalization factors.
\item Compute its exact \(j\)-dependence in both square orientations.
\item Compare its square with the necessary threshold \(d_j^{-1}\)
      from \eqref{eq:V23-necessary-output-loss}.
\item Depending on the result, either record a counterexample to the
      present exchange card and enlarge the target, or certify the
      physical-output fourth response and return to the remaining
      temporal classes.
\end{enumerate}
\end{programme}

\begin{firewall}[Claim boundary after fourth-series V23]
\label{fire:V23-claim-boundary}
V23 proves that the V22 packet is legal at the mass, stock,
first-connectivity, tree-fibre, payer, and normalization levels, and
that the full exchange weight remains short by one square half-power.
It does not decide the physical-output survival quotient, and hence
does not yet disprove or prove \(XQTRSC_5\), close \(q=3\), or prove
global quintic MTA.  The remaining \(q=3\) classes, \(q=0,1,2\) cards,
global quintic or all-order MTA, WUFC, CPM, cutoff removal, reflection
positivity, Osterwalder--Schrader reconstruction, construction of the
continuum Yang--Mills measure, and a uniform positive continuum
spectral gap remain open.  The full Yang--Mills mass gap has not been
proved.
\end{firewall}

\paragraph{Parent-version record.}
V23 was formed from
\texttt{Renewal\_Yang\_Mills\_Mass\_Gap\_Research\_Notes\_V22.tex}.
The V22 parent had \(10\,948\) logical source lines,
\(387\,672\) bytes, and SHA--256
\[
 \texttt{743e14ac4369e14eb1b630665ce707681de025268c43e879e775ee8805552bdd}.
\]
The next compulsory companion audit is V25.

% =============================================================================
% END APPENDED FOURTH-SERIES V23 MATERIAL
% =============================================================================

\clearpage
% =============================================================================
% BEGIN APPENDED FOURTH-SERIES V24 MATERIAL
% =============================================================================

\chapter{V24: Identity output duality disproves the pointwise
exchange-rescue square card}
\label{ch:V24}

\section{Purpose and correction of the V23 acquisition plan}
\label{sec:V24-purpose}

V23 reduced the explicit pair-crossing \(3+2\) fibre to one question:
does the last physical output-dual contraction necessarily remove a
factor \(d_j^{-1}\) from the oriented square?  The proposed V24 plan
asked for a new Clebsch--Gordan computation.  Going upstream to the
exact Fourier duality of frozen archive f2 V94 shows that this
calculation is unnecessary.

The dual formula takes a supremum over \emph{all} contractions on each
irreducible output factor.  The identity contraction is admissible and
leaves the retained complete output block unchanged.  Hence the
survival quotient of V23 is exactly one for an allowed dual family.
Combined with the exact deficit ratio
\eqref{eq:V23-deficit-ratio}, this gives a counterexample to the
pointwise sufficient interface \(XQTRSC_5\).

\section{The exact dual family contains the identity}
\label{sec:V24-dual}

We first record only the part of the f2 V94 theorem needed here.  For a
fixed input tensor product, let
\[
 \mathcal H_{\rm in}
 =
 \bigoplus_{U\in\mathcal S}
 W_U(H_U\otimes M_U),
\]
and let \(K_U\) be the exact complete multiplier on \(M_U\).  For
contractions \(D_U\in\mathcal B(H_U)\), the unsplit dual operator is
\begin{equation}
 \mathcal T_K(\boldsymbol D)
 :=
 \sum_{U\in\mathcal S}
 W_U(D_U^*\otimes K_U)W_U^*.
\label{eq:V24-dual-operator}
\end{equation}
The exact block-duality identity takes the supremum over the full ball
\(\|D_U\|\le1\), independently for each \(U\).

\begin{lemma}[Identity-dual retention on a fixed output block]
\label{lem:V24-identity-retention}
Fix \(U_0\in\mathcal S\) and define
\begin{equation}
 D_{U_0}^{\rm id}=I_{H_{U_0}},
 \qquad
 D_U^{\rm id}=0\quad(U\ne U_0).
\label{eq:V24-identity-family}
\end{equation}
Then \(\boldsymbol D^{\rm id}\) is an admissible dual family and
\begin{equation}
 \mathcal T_K(\boldsymbol D^{\rm id})
 =
 W_{U_0}(I_{H_{U_0}}\otimes K_{U_0})W_{U_0}^*.
\label{eq:V24-identity-block}
\end{equation}
Thus on a resolved word already restricted to the complete
\(U_0\)-block, the output-dual operation is the identity.  It changes
neither oriented square nor its product-state capacity.
\end{lemma}

\begin{proof}
Both \(I_{H_{U_0}}\) and the zero operators have norm at most one, so
\eqref{eq:V24-identity-family} belongs to the contraction family in
the exact dual supremum.  Substitute it into
\eqref{eq:V24-dual-operator}; every term except \(U_0\) vanishes, and
the remaining term is \eqref{eq:V24-identity-block}.  The
\(U_0\)-block inclusion \(W_{U_0}\) is already part of the resolved
pre-dual operator.  Therefore applying this dual family reproduces
that operator exactly, not merely up to unitary equivalence.  Equality
of the operators implies equality of both squares and their
capacities.
\end{proof}

\begin{firewall}[Identity duality is not a favourable extra choice]
\label{fire:V24-identity-legality}
The identity family is not an output map added to the source.  It is
one member of the contraction ball over which the exact trace-norm
duality takes its supremum.  Any uniform square estimate used before
that supremum must therefore control it.  Excluding
\(\boldsymbol D^{\rm id}\) would no longer prove the archived exact
Fourier coefficient bound.
\end{firewall}

\section{The V23 survival quotient equals one}
\label{sec:V24-survival}

\begin{proposition}[No physical-output loss on the test fibre]
\label{prop:V24-survival-one}
Choose the complete output block \(U_0\) used in the V23 resolved
word and the identity family
\eqref{eq:V24-identity-family}.  Then
\begin{equation}
 \boxed{
 \sigma_j^R(\boldsymbol D^{\rm id})
 =
 \sigma_j^L(\boldsymbol D^{\rm id})
 =1}
\label{eq:V24-survival-one}
\end{equation}
for every \(j\) for which the fixed prefix block is nonzero.
\end{proposition}

\begin{proof}
By construction,
\(\widehat{\mathbb W}_j\) in
\eqref{eq:V23-predual-word} is the complete \(U_0\)-block before its
dual contraction.  By
\Cref{lem:V24-identity-retention},
\[
 \widetilde{\mathbb W}_{j,\boldsymbol D^{\rm id}}
 =
 \widehat{\mathbb W}_j.
\]
Substitution into the two definitions of the survival quotient proves
\eqref{eq:V24-survival-one}.  Nonvanishing of the denominator follows
from \Cref{thm:V23-predual-lower}.
\end{proof}

\begin{corollary}[The necessary output loss cannot hold]
\label{cor:V24-no-output-loss}
The necessary condition
\eqref{eq:V23-necessary-output-loss} fails for the admissible identity
family, since \(1\not\le Cd_j^{-1}\) uniformly as \(j\to\infty\).
\end{corollary}

\begin{proof}
Use \eqref{eq:V24-survival-one} and \(d_j\to\infty\).
\end{proof}

\section{Counterexample to the pointwise card}
\label{sec:V24-counterexample}

\begin{theorem}[The archived pointwise \(XQTRSC_5\) interface is
false]
\label{thm:V24-XQTRSC-false}
For \(G=SU(2)\) and fixed heat time \(s>0\), there is no constant
\(C_{G,s}\) for which both pointwise exchange-rescue square estimates
hold for every fixed resolved quintic fibre, payer, output-dual family,
and representation height.  More precisely, the right-oriented card
fails on the legal family of V23 with the joining payer and identity
dual family:
\begin{equation}
 {\kappa_\otimes(
  \widetilde{\mathbb W}_{j,\boldsymbol D^{\rm id}}^*
  \widetilde{\mathbb W}_{j,\boldsymbol D^{\rm id}})}
 {\mathfrak x_{T_*,r}^{\,2}}
 \ge c_{\ell,s}d_j
 \longrightarrow\infty.
\label{eq:V24-XQTRSC-counterexample}
\end{equation}
The same construction gives the opposite-orientation failure.
\end{theorem}

\begin{proof}
The source support, first-connectivity coordinate, prefix partition,
lifted tree, payer placement, and exact exchange weight were proved
legal in
\Cref{lem:V23-legality,
lem:V23-masses,
lem:V23-pair-stocks,
thm:V23-exact-exchange,
lem:V23-nonzero-prefix}.
By \Cref{prop:V24-survival-one}, the numerator in
\eqref{eq:V24-XQTRSC-counterexample} equals the pre-dual numerator in
\eqref{eq:V23-deficit-ratio}.  That equation supplies the displayed
lower bound.  It diverges, contradicting any representation-uniform
constant.  All varying complete heat and covariance factors are
self-adjoint and the fixed prefix fibre is tested in the corresponding
opposite square, so the V22 lower test and identity-dual argument apply
to the left orientation as well.
\end{proof}

\begin{corollary}[Withdraw the pointwise closure programme]
\label{cor:V24-withdraw-pointwise}
Any earlier conditional theorem whose conclusion was obtained by
verifying \(XQTRSC_5\) for the V23 fibre remains a valid implication
from its stated hypotheses, but those hypotheses cannot be universal.
In particular, no completion of the global quintic programme may cite
the pointwise card as holding for every resolved fibre.
\end{corollary}

\begin{proof}
A valid conditional implication is not altered by a counterexample to
the universality of its hypothesis.  The universal premise is excluded
by \Cref{thm:V24-XQTRSC-false}.
\end{proof}

\begin{firewall}[The global quintic MTA is not disproved]
\label{fire:V24-global-not-disproved}
Archive f3 V98 proves that universal \(XQTRSC_5\) would be
\emph{sufficient} for global quintic MTA.  It does not prove the
converse.  The counterexample above fixes a tree/cumulant resolution
before squaring.  The complete connected quintic source sums the
finite signed resolution before its absolute Fourier coefficient is
taken.  Cancellation among those terms can therefore make the
complete carrier smaller even though one resolved square is too
large.  V24 disproves the pointwise proof interface, not global
quintic MTA and not Yang--Mills theory.
\end{firewall}

\section{Where the missing cancellation can live}
\label{sec:V24-aggregate}

For fixed \((r,\rho)\), fixed physical output block, payer, and dual
family, let \(\mathfrak F(r,\rho)\) be the finite set of prefix-tree,
quotient-tree, and endpoint-lift resolution labels.  Define the
assembled resolved source
\begin{equation}
 \mathbb S_{r,\rho,\pi,\boldsymbol D}
 :=
 \sum_{\alpha\in\mathfrak F(r,\rho)}
 \widetilde{\mathbb W}_{r,\rho,\alpha,\pi,\boldsymbol D}.
\label{eq:V24-assembled-source}
\end{equation}
The signs and coefficients in this sum are those of the exact complete
cumulants; absolute values are not inserted termwise.

\begin{definition}[Aggregated exchange-rescue square target]
\label{def:V24-AXQTRSC}
The minimally repaired finite-fibre target \(AXQTRSC_5\) is
\begin{align}
 \kappa_\otimes(
  \mathbb S_{r,\rho,\pi,\boldsymbol D}^*
  \mathbb S_{r,\rho,\pi,\boldsymbol D})
 &\le
 C_{G,s}
 \left(
  \sum_{\alpha\in\mathfrak F(r,\rho)}
  \mathfrak x_{T_\alpha,r}
 \right)^2,
\label{eq:V24-AXQTRSC-right}\\
 \kappa_\otimes(
  \mathbb S_{r,\rho,\pi,\boldsymbol D}
  \mathbb S_{r,\rho,\pi,\boldsymbol D}^*)
 &\le
 C_{G,s}
 \left(
  \sum_{\alpha\in\mathfrak F(r,\rho)}
  \mathfrak x_{T_\alpha,r}
 \right)^2.
\label{eq:V24-AXQTRSC-left}
\end{align}
Only a fixed arity-five family is assembled, so this repair creates no
support-size factor.
\end{definition}

\begin{proposition}[Aggregated target remains sufficient]
\label{prop:V24-AXQTRSC-sufficient}
If \(AXQTRSC_5\) holds for every
\((r,\rho,\pi,\boldsymbol D)\), then the same exchange-summation
argument as archive f3 V98 yields global quintic MTA.
\end{proposition}

\begin{proof}
Product-state Cauchy--Schwarz applied to either orientation gives
\[
 \kappa_\otimes^{\rm abs}
  (\mathbb S_{r,\rho,\pi,\boldsymbol D})
 \le
 C_{G,s}
 \sum_{\alpha\in\mathfrak F(r,\rho)}
 \mathfrak x_{T_\alpha,r}.
\]
The resolution family has cardinality bounded solely by order five.
Sum in \(r\) and apply the archived exchange-rescue summation to each
fixed tree label.  The remaining prefix partitions, payer placements,
physical blocks, and positive-parts choices are also fixed finite
families.  Their multiplicities are absorbed into \(C_{G,s}\), leaving
the complete \(125\)-tree global atlas.
\end{proof}

\begin{assessment}[Exact new bottleneck]
\label{ass:V24-bottleneck}
The V23 packet shows what \(AXQTRSC_5\) must accomplish: within the
assembled \(3+2\) source at the same \((r,\rho)\), the signed
tree/cumulant sum must cancel the identity-dual spin-one coefficient
to relative square order \(d_j^{-1}\).  The next calculation is
therefore not another heat estimate.  It is the finite Möbius and
endpoint-lift coefficient of the complete \(3+2\) prefix--joining
source on the explicit support skeleton.
\end{assessment}

\section{V24 result ledger and V25 acquisition order}
\label{sec:V24-ledger}

\begin{theorem}[Fourth-series V24 result ledger]
\label{thm:V24-results}
Relative to V1--V23, V24 proves:
\begin{enumerate}[label=\textup{(V24.\arabic*)},leftmargin=6.3em]
\item the identity-dual retention lemma for every complete output
      block;
\item exact survival quotients
      \(\sigma_j^R=\sigma_j^L=1\) on the V23 packet;
\item an explicit legal counterexample to both universal pointwise
      \(XQTRSC_5\) orientations;
\item the precise withdrawal of the pointwise closure route without
      withdrawing global quintic MTA;
\item the aggregated finite-fibre target \(AXQTRSC_5\); and
\item sufficiency of that repaired target for the global quintic
      marked-tree atlas.
\end{enumerate}
\end{theorem}

\begin{proof}
These are
\Cref{lem:V24-identity-retention,
prop:V24-survival-one,
thm:V24-XQTRSC-false,
cor:V24-withdraw-pointwise,
def:V24-AXQTRSC,
prop:V24-AXQTRSC-sufficient}.
\end{proof}

\begin{programme}[V25: audit and assembled \(3+2\) coefficient]
\label{prog:V24-V25}
\begin{enumerate}[leftmargin=3em]
\item Perform the compulsory audit of the newest active SM and NS
      notes and import only type-correct structural lessons.
\item Return to the exact coefficient-one first-connectivity source,
      before termwise tree absolute values.
\item On the V23 six-coordinate skeleton, list every nonzero
      prefix-tree, quotient-tree, and endpoint-lift term at the fixed
      \((r,\rho)\).
\item Evaluate their signed complete spin-one coefficient under the
      identity output dual.
\item Decide whether the sum gains the necessary amplitude factor
      \(d_j^{-1/2}\), equivalently square factor \(d_j^{-1}\).
\item If not, disprove \(AXQTRSC_5\) and move farther upstream to the
      complete quintic carrier before the \((r,\rho)\) split.
\end{enumerate}
\end{programme}

\begin{firewall}[Claim boundary after fourth-series V24]
\label{fire:V24-claim-boundary}
V24 disproves the universal pointwise \(XQTRSC_5\) sufficient
interface and replaces it by a still-unproved aggregated target.  It
does not disprove global quintic MTA, prove \(AXQTRSC_5\), close
\(q=3\), or prove any later analytic or constructive gate.  The
remaining quintic aggregation, \(q=0,1,2\) sectors, all-order MTA,
WUFC, CPM, cutoff removal, reflection positivity,
Osterwalder--Schrader reconstruction, construction of the continuum
Yang--Mills measure, and a uniform positive continuum spectral gap
remain open.  The full Yang--Mills mass gap has not been proved.
\end{firewall}

\paragraph{Parent-version record.}
V24 was formed from
\texttt{Renewal\_Yang\_Mills\_Mass\_Gap\_Research\_Notes\_V23.tex}.
The V23 parent had \(11\,506\) logical source lines,
\(407\,051\) bytes, and SHA--256
\[
 \texttt{1e8528a10785e0c8e0b0197854ed03b90a1e602992bb002154b66d16b9dee9a9}.
\]
The next compulsory companion audit is V25.

% =============================================================================
% END APPENDED FOURTH-SERIES V24 MATERIAL
% =============================================================================

\clearpage
% =============================================================================
% BEGIN APPENDED FOURTH-SERIES V25 MATERIAL
% =============================================================================

\chapter{V25: Companion audit and direct absolute closure of the
spin-one packet}
\label{ch:V25}

\section{Compulsory companion audit}
\label{sec:V25-audit}

At this checkpoint the newest active companion notes were:
\begin{center}
\begin{tabular}{@{}p{0.60\linewidth}rr@{}}
\toprule
file & logical lines & bytes\\
\midrule
\texttt{Renewal\_SM\_Einstein\_Continuum\_Research\_Notes\_V4.tex}
 & \(1\,209\) & \(42\,226\)\\
\texttt{Renewal\_NS\_Stability\_Research\_Notes\_V31.tex}
 & \(23\,052\) & \(887\,537\)\\
\bottomrule
\end{tabular}
\end{center}
Their SHA--256 digests were, respectively,
\[
\begin{split}
&\texttt{543c807098c7568dd83debb9bae46db762d6ffc1b10c11d1ba0ff0b9b7e80d4c},\\
&\texttt{f1121b62238ad5491bb7cf03635ea9934942edf573ed8faaa9ee79e1d38a6696}.
\end{split}
\]

SM V4 constructs an exact finite-dimensional shooting map from a
local ultraviolet stable graph to infrared constraints.  Its decisive
firewall is that equality of dimensions only permits an isolated
transverse match; it does not prove existence.  Existence requires a
quantitative residual and inverse-derivative certificate.

The type-correct YM transfer is immediate.  The fact that the
tree/cumulant resolution family in
\eqref{eq:V24-assembled-source} is finite does not imply that its
signed sum cancels the V23 packet.  A cancellation claim needs an
explicit coefficient or a quantitative inverse/transversality
statement.  Counting resolution labels is not evidence of decay.

NS V31 proves an exact dyadic first-moment formation ledger and then
exhibits a balanced-helicity nullspace for its sector-difference hinge.
Its V32 programme requires complete signed triad work to be assembled
before positive parts are taken.  It also insists that a mechanism
which merely returns the continuation-class bound is not progress.

The type-correct YM transfer is the order of operations, not an NS
estimate: retain the complete signed Wilson coefficient before taking
a square or an absolute value, and test any proposed gain on the sharp
V23 packet.  Neither the NS heat ladder nor its helicity identities are
imported.

\begin{audit}[V25 direction after both companion checks]
\label{audit:V25-direction}
Do not seek cancellation from finite aggregation alone.  First test
the linear absolute product coefficient actually required by global
MTA.  The false pointwise square route may lose a square root that the
exact signed coefficient never loses.
\end{audit}

\section{The aggregated square repair also fails}
\label{sec:V25-aggregate-square}

\begin{lemma}[Only one nonzero tree/lift fibre on the V23 support]
\label{lem:V25-one-fibre}
Fix the V23 support, first-connectivity pair \((r,\rho)\), joining
payer, retained output block, and identity dual family.  Among the
prefix-tree, quotient-tree, and endpoint-lift labels in
\(\mathfrak F(r,\rho)\), the only nonzero edge-labelled fibre has
\[
 P=\{ab,ad,ce\},
 \qquad Q=\{ac\},
\]
with the endpoint lifts displayed in \(T_*\).
\end{lemma}

\begin{proof}
The only nonzero pair stocks are
\(ab,ad,ce,ac\), by
\Cref{lem:V23-pair-stocks}.  The three-block \(\{a,b,d\}\) has only
the two available internal edges \(ab,ad\); hence its unique nonzero
spanning tree is the path \(b-a-d\).  The two-block \(\{c,e\}\) has
the forced edge \(ce\).  The quotient has two vertices and the forced
edge \(ac\).  At each involved coordinate, only the two rows naming
that edge are nontrivial, so every endpoint lift is forced.  Every
other resolved edge fibre contains a zero pair occurrence and
vanishes.
\end{proof}

\begin{corollary}[The proposed \(AXQTRSC_5\) square target is false]
\label{cor:V25-AXQTRSC-false}
The aggregated source
\(\mathbb S_{r,\rho,\pi,\boldsymbol D^{\rm id}}\) on the V23 support
equals its unique nonzero resolved fibre.  Consequently the ratio in
\eqref{eq:V24-XQTRSC-counterexample} also disproves both orientations
of \eqref{eq:V24-AXQTRSC-right}--%
\eqref{eq:V24-AXQTRSC-left}.
\end{corollary}

\begin{proof}
Use \Cref{lem:V25-one-fibre} in
\eqref{eq:V24-assembled-source}.  There is no second nonzero summand
with which the V23 word could cancel.  Theorem
\ref{thm:V24-XQTRSC-false} then applies verbatim.
\end{proof}

\begin{firewall}[Why this does not revive termwise absolute values]
\label{fire:V25-no-termwise}
The one-fibre conclusion is a special property of the explicit sparse
support.  It is used only to test a universal theorem.  V25 does not
replace a general assembled source by the sum of absolute values of
its tree fibres.
\end{firewall}

\section{A partial-transpose bound for the spin-one projection}
\label{sec:V25-spin-one-PT}

Let \(|1,M\rangle\), \(M=-1,0,1\), be the standard coupled basis of the
total-spin-one copy in \(V_j\otimes V_j\), and
\[
 P_{j,1}:=\sum_{M=-1}^{1}|1,M\rangle\langle1,M|.
\]
Partial transpose below is taken in the second spin-\(j\) factor.

\begin{lemma}[Uniform low-spin Clebsch--Gordan coefficient]
\label{lem:V25-CG-upper}
For \(j\ge1\), every Schmidt coefficient of every
\(|1,M\rangle\) is at most \(\sqrt{3/d_j}\).  Consequently
\begin{equation}
 \left\|
  (|1,M\rangle\langle1,M|)^\Gamma
 \right\|
 \le {3\over d_j}.
\label{eq:V25-rank-one-PT}
\end{equation}
\end{lemma}

\begin{proof}
In the uncoupled basis, the nonzero coefficients for \(M=0\) are
\[
 \langle j,m;j,-m\mid1,0\rangle
 =
 (-1)^{j-m}
 \sqrt{\frac{3}{j(j+1)(2j+1)}}\,m.
\]
Their squares are at most \(3/d_j\).  For \(M=1\),
\[
 \left|
 \langle j,m;j,1-m\mid1,1\rangle
 \right|^2
 =
 \frac{3(j+m)(j-m+1)}
 {2j(j+1)(2j+1)}
 \le {3\over d_j}
\]
for \(j\ge1\); the \(M=-1\) formula is obtained by replacing
\(m\) by \(-m\).  These uncoupled expansions are Schmidt
decompositions because distinct \(m\) give distinct orthonormal basis
vectors in each tensor factor.

If a unit vector has Schmidt coefficients \(s_k\), the eigenvalues of
the partial transpose of its rank-one projector are
\(s_k^2\) and \(\pm s_ks_l\).  Its operator norm is therefore
\(\max_k s_k^2\).  The coefficient bounds give
\eqref{eq:V25-rank-one-PT}.
\end{proof}

\begin{corollary}[Complete spin-one projection has inverse-dimension
partial-transpose norm]
\label{cor:V25-P1-PT}
For \(j\ge1\),
\begin{equation}
 \boxed{
 \|P_{j,1}^{\Gamma}\|
 \le {9\over d_j}.}
\label{eq:V25-P1-PT}
\end{equation}
On three distinct coordinates,
\begin{equation}
 \left\|
  (P_{j,1}^{\otimes3})^\Gamma
 \right\|
 \le {9^3\over d_j^3}.
\label{eq:V25-three-P1-PT}
\end{equation}
\end{corollary}

\begin{proof}
Apply the triangle inequality to the three rank-one summands of
\(P_{j,1}\) and use
\Cref{lem:V25-CG-upper}.  Partial transpose and operator norm
tensorize on distinct physical coordinates, giving the second
display.
\end{proof}

\begin{remark}[Sharp scale]
\label{rem:V25-P1-sharp}
\Cref{cor:V13-pj1-scale} gives a product vector with
\(\langle\phi_j,P_{j,1}\phi_j\rangle\gtrsim d_j^{-1}\).
Thus the inverse-dimension scale in
\eqref{eq:V25-P1-PT} is sharp up to a fixed constant.
\end{remark}

\section{Direct absolute closure of the counterpacket}
\label{sec:V25-direct}

Resolve the root covariance and the two suffix heats of the V23 word
on their complete total-spin-one isotypes.  On this fixed output
block their varying multiplier is
\begin{equation}
 \mathbb G_j
 =
 r_1(s)\Delta_j(s)q_1(s)^2
 P_{j,1}^{\otimes3}.
\label{eq:V25-high-block}
\end{equation}
This is the assembled paid covariance multiplier at \(r\), followed
by the two complete suffix heat multipliers; it is not one term of the
covariance.

\begin{theorem}[Direct absolute estimate for the V23 fibre]
\label{thm:V25-direct-close}
For the identity output dual and the fixed nonzero prefix fibre of
\Cref{lem:V23-nonzero-prefix},
\begin{equation}
 \kappa_\otimes^{\rm abs}
  (\widetilde{\mathbb W}_{j,\boldsymbol D^{\rm id}})
 \le
 {C_{\ell,s}d_j^{-3}\over
  \Xi_a\Xi_cB^3}.
\label{eq:V25-direct-upper}
\end{equation}
Consequently,
\begin{equation}
 \kappa_\otimes^{\rm abs}
  (\widetilde{\mathbb W}_{j,\boldsymbol D^{\rm id}})
 \le
 C_{\ell,s}\mathfrak x_{T_*,r}
\label{eq:V25-direct-tree-card}
\end{equation}
uniformly in \(j\).
\end{theorem}

\begin{proof}
For a product density matrix across the five rows and any operator
\(Z\), the usual partial-transpose estimate gives
\[
 |\Tr(\varrho Z)|
 =
 |\Tr(\varrho^\Gamma Z^\Gamma)|
 \le\|Z^\Gamma\|,
\]
because \(\varrho^\Gamma\) is again a density matrix.  Apply this
across the cut \(a|\{b,c,d,e\}\).  The fixed prefix operator
\(\mathbb F_\ell\) has a finite partial-transpose norm depending only
on \(\ell,s\).  The scalar coefficient in
\eqref{eq:V25-high-block} is uniformly bounded: \(q_1(s)\) and
\(r_1(s)\) are fixed, while
\(|\Delta_j(s)|\le1\).  Hence
\Cref{cor:V25-P1-PT} and tensorization give
\[
 \left\|
  (\mathbb F_\ell\otimes\mathbb G_j)^\Gamma
 \right\|
 \le C_{\ell,s}d_j^{-3}.
\]
Divide once by the row-mass product from
\eqref{eq:V23-predual-word}; this proves
\eqref{eq:V25-direct-upper}.

By \eqref{eq:V23-masses},
\(\Xi_a,\Xi_c\asymp_\ell A_j\), and
\(d_j^{-3}\lesssim A_j^{-3/2}\).  Therefore the left side of
\eqref{eq:V25-direct-upper} is
\[
 O_{\ell,s}(A_j^{-7/2}).
\]
On the other hand,
\Cref{thm:V23-exact-exchange} gives
\(\mathfrak x_{T_*,r}\asymp_\ell A_j^{-3}\).
Since \(A_j\ge1\), the former is bounded by a fixed multiple of the
latter, proving \eqref{eq:V25-direct-tree-card}.
\end{proof}

\begin{corollary}[The square-root loss was artificial on the test
fibre]
\label{cor:V25-square-root-artificial}
On the same legal fibre, both \(XQTRSC_5\) and \(AXQTRSC_5\) fail,
but the direct absolute tree card
\eqref{eq:V25-direct-tree-card} holds with a spare factor
\(A_j^{-1/2}\).  Therefore failure of those square interfaces is not
evidence against global quintic MTA.
\end{corollary}

\begin{proof}
The two failures are
\Cref{thm:V24-XQTRSC-false,cor:V25-AXQTRSC-false}; the direct bound is
\Cref{thm:V25-direct-close}.  Comparing
\(A_j^{-7/2}\) with \(A_j^{-3}\) gives the spare factor.
\end{proof}

\section{The corrected analytic interface}
\label{sec:V25-interface}

\begin{definition}[Direct exchange-rescue coefficient card]
\label{def:V25-DXQTRC}
For the exact assembled source
\(\mathbb S_{r,\rho,\pi,\boldsymbol D}\) of
\eqref{eq:V24-assembled-source}, define \(DXQTRC_5\) by
\begin{equation}
 \boxed{
 \kappa_\otimes^{\rm abs}
  (\mathbb S_{r,\rho,\pi,\boldsymbol D})
 \le
 C_{G,s}
 \sum_{\alpha\in\mathfrak F(r,\rho)}
 \mathfrak x_{T_\alpha,r}.}
\label{eq:V25-DXQTRC}
\end{equation}
The signed finite source is assembled before the absolute product
coefficient is taken.  No oriented square is introduced.
\end{definition}

\begin{proposition}[\(DXQTRC_5\) is exactly sufficient]
\label{prop:V25-DXQTRC-sufficient}
If \eqref{eq:V25-DXQTRC} holds for every fixed
\((r,\rho,\pi,\boldsymbol D)\), then the global quintic marked-tree
atlas follows by the archived exchange-rescue summation.
\end{proposition}

\begin{proof}
Sum \eqref{eq:V25-DXQTRC} over the fixed finite prefix partitions,
payer placements, physical positive-parts choices, and output dual
families in the exact Fourier formula.  Then sum in \(r\).  For each
tree fibre, the archived estimate
\[
 \sum_r\mathfrak x_{T_\alpha,r}
 \le
 15\sum_{\tau\in\mathfrak T_5}
 \prod_{e\in E(\tau)}\theta_e
\]
applies.  All remaining multiplicities depend only on arity five and
are absorbed into \(C_{G,s}\).  This is precisely the desired global
tree atlas.
\end{proof}

\begin{proposition}[A sufficient direct partial-transpose compiler]
\label{prop:V25-direct-compiler}
Fix a source fibre with maximal row \(a\), \(M=\Xi_a\), and normalized
raw prefactor \((\prod_i\Xi_i)^{-1}\).  Suppose its assembled
unnormalized source admits, in one output-dual block, an exact
factorization
\[
 \mathbb L\,
 \bigotimes_{e\in E_3}\mathbb H_e\,
 \mathbb Q,
 \qquad |E_3|=3,
\]
such that, after partial transpose across \(a|a^c\),
\[
 \|\mathbb L^\Gamma\|\,
 \|\mathbb Q^\Gamma\|\le C,
 \qquad
 \|\mathbb H_e^\Gamma\|\le C M^{-1/2}.
\]
Then
\begin{equation}
 \kappa_\otimes^{\rm abs}(\mathbb S)
 \le
 {C M^{-3/2}\over\prod_i\Xi_i}.
\label{eq:V25-direct-compiler}
\end{equation}
In every one-strong-edge tree geometry whose exchange floor is
\[
 \mathfrak x\ge {c\over M\prod_i\Xi_i},
\]
this proves the direct card with a spare \(M^{-1/2}\).
\end{proposition}

\begin{proof}
Partial transpose is multiplicative on the displayed tensor factors.
Its operator norm is therefore at most \(CM^{-3/2}\).  The
partial-transpose trace estimate used in
\Cref{thm:V25-direct-close}, followed by the raw normalization, gives
\eqref{eq:V25-direct-compiler}.  Since \(M\ge1\),
\[
 {M^{-3/2}\over\prod_i\Xi_i}
 \le
 {M^{-1}\over\prod_i\Xi_i}
 \le c^{-1}\mathfrak x.
\]
\end{proof}

\begin{firewall}[The remaining separator problem]
\label{fire:V25-direct-separator}
The compiler requires a bound for the partial transpose of the
\emph{assembled} outer factors.  Ordinary operator norm does not
control partial-transpose norm uniformly in dimension, and a
noncentral recoupler can invalidate a factorwise argument.  Thus V25
does not apply \Cref{prop:V25-direct-compiler} automatically to all
\(q=3\) sources.  It replaces the false square target by the correct
linear target and isolates the precise whole-word separator
certificate still needed.
\end{firewall}

\begin{assessment}[Direction after the V25 audit]
\label{ass:V25-next}
The sparse counterpacket has served its purpose: it disproves the
square interfaces while passing the direct coefficient card.  The
next step is to classify which exact \(q=3\) first-connectivity
sources have a dimension-uniform partial-transpose separator in the
maximal-row cut.  For failures, one must use complete dual-sandwich
factorization or signed source assembly before partial transpose,
following the companion audit's common order-of-operations lesson.
\end{assessment}

\section{V25 result ledger and V26 acquisition order}
\label{sec:V25-ledger}

\begin{theorem}[Fourth-series V25 result ledger]
\label{thm:V25-results}
Relative to V1--V24, V25 proves:
\begin{enumerate}[label=\textup{(V25.\arabic*)},leftmargin=6.3em]
\item the compulsory SM V4 and NS V31 companion audit;
\item failure of the aggregated square repair on the sparse V23
      support;
\item an explicit inverse-dimension partial-transpose bound for the
      complete spin-one projection;
\item direct absolute closure of the very packet which disproves both
      square interfaces;
\item the corrected \(DXQTRC_5\) analytic target and its exact
      sufficiency for global quintic MTA; and
\item a three-slot direct partial-transpose compiler with the precise
      remaining separator hypothesis.
\end{enumerate}
\end{theorem}

\begin{proof}
These are
\Cref{audit:V25-direction,
lem:V25-one-fibre,
cor:V25-AXQTRSC-false,
lem:V25-CG-upper,
cor:V25-P1-PT,
thm:V25-direct-close,
def:V25-DXQTRC,
prop:V25-DXQTRC-sufficient,
prop:V25-direct-compiler}.
\end{proof}

\begin{programme}[V26: direct-separator census for \(q=3\)]
\label{prog:V25-V26}
\begin{enumerate}[leftmargin=3em]
\item Return to the seven temporal classes
      \((p,j,s)\) with \(p+j+s=3\).
\item For each first-connectivity source type, write the exact factors
      outside the three retained heat slots before any square.
\item Mark the classes in which those outer factors are local
      unitaries, complete central moments, or fixed-dimensional maps;
      prove their partial-transpose norm is uniform.
\item For a dimension-growing prefix or joining recoupler, apply the
      complete output-dual sandwich directly to the linear coefficient
      rather than taking an oriented square.
\item Test every proposed separator bound on the V9 entangler and V23
      identity-dual packets.
\item Close every certified class by
      \Cref{prop:V25-direct-compiler} and state the exact residual.
\end{enumerate}
\end{programme}

\begin{firewall}[Claim boundary after fourth-series V25]
\label{fire:V25-claim-boundary}
V25 disproves both pointwise and finitely aggregated square
interfaces, but proves the correct direct absolute card on their sharp
counterpacket and formulates an exactly sufficient replacement.  It
does not yet prove \(DXQTRC_5\) for all source fibres, close \(q=3\),
or prove global quintic MTA.  The remaining direct-separator classes,
\(q=0,1,2\) sectors, all-order MTA, WUFC, CPM, cutoff removal,
reflection positivity, Osterwalder--Schrader reconstruction,
construction of the continuum Yang--Mills measure, and a uniform
positive continuum spectral gap remain open.  The full Yang--Mills
mass gap has not been proved.
\end{firewall}

\paragraph{Parent-version record.}
V25 was formed from
\texttt{Renewal\_Yang\_Mills\_Mass\_Gap\_Research\_Notes\_V24.tex}.
The V24 parent had \(11\,876\) logical source lines,
\(420\,877\) bytes, and SHA--256
\[
 \texttt{aa3b66d6c2df04422e6ed4de0b66dcb446c36ea3ab8b57eb290cc5e87f1734bc}.
\]
The next compulsory companion audit is V30.

% =============================================================================
% END APPENDED FOURTH-SERIES V25 MATERIAL
% =============================================================================

\clearpage
% =============================================================================
% BEGIN APPENDED FOURTH-SERIES V26 MATERIAL
% =============================================================================

\chapter{V26: A checkable direct output-dual separator certificate}
\label{ch:V26}

\section{Purpose}
\label{sec:V26-purpose}

V25 replaces the false oriented-square interface by the direct
absolute coefficient card \(DXQTRC_5\).  Its explicit proof used the
identity output dual.  The exact Fourier formula, however, takes a
supremum over arbitrary contractions \(D_U\) on irreducible output
factors.  A general completion therefore needs a uniform way to pass
the three heat responses through the assembled output dualizer.

This chapter proves two facts.  First, ordinary norm bounds on the
outer separator are insufficient: a norm-one entangler can erase all
three inverse-dimension gains.  Second, a finite row-cut decomposable
sandwich gives an exact sufficient certificate.  The certificate is
stable under sums and source assembly, applies in either orientation,
and closes every \(q=3\) fibre for which it is verified.

\section{A threefold entangler obstruction}
\label{sec:V26-entangler}

Let
\[
 \Omega_d=d^{-1/2}\sum_{m=1}^d e_m\otimes e_m,
 \qquad
 P_d=|\Omega_d\rangle\langle\Omega_d|
\]
on \(\mathbb C^d\otimes\mathbb C^d\).

\begin{lemma}[Threefold ordinary-norm amplification]
\label{lem:V26-three-entangler}
There is a unitary \(U_d\) and a product unit vector \(x_d\) such that
\[
 U_dx_d=\Omega_d^{\otimes3}
\]
on
\((\mathbb C^d)^{\otimes3}\otimes
 (\mathbb C^d)^{\otimes3}\).  Consequently
\begin{align}
 \kappa_\otimes^{\rm abs}
 \left[
  U_d^*P_d^{\otimes3}U_d
 \right]&=1,
\label{eq:V26-entangled-capacity}\\
 \left\|
  (P_d^{\otimes3})^\Gamma
 \right\|&=d^{-3}.
\label{eq:V26-protected-PT}
\end{align}
\end{lemma}

\begin{proof}
The full Hilbert space is finite-dimensional, so a unitary can map
any prescribed unit vector to any other; choose a product unit vector
\(x_d\) and impose \(U_dx_d=\Omega_d^{\otimes3}\).  Then
\[
 U_d^*P_d^{\otimes3}U_d
 =
 |x_d\rangle\langle x_d|,
\]
whose product-state absolute capacity is one.  For one maximally
entangled rank-one projection, partial transpose has operator norm
\(d^{-1}\).  Partial transpose and operator norm tensorize over the
three copies, proving \eqref{eq:V26-protected-PT}.
\end{proof}

\begin{corollary}[Ordinary separator norm cannot prove the direct
card]
\label{cor:V26-norm-no-go}
There is no dimension-independent inequality of the form
\[
 \kappa_\otimes^{\rm abs}(LZR)
 \le C\|L\|\,\|R\|\,\|Z^\Gamma\|
\]
for arbitrary \(L,R,Z\), even when \(Z\) is a positive threefold
protected heat block and \(L=R^*\) is unitary.
\end{corollary}

\begin{proof}
Take
\[
 Z=P_d^{\otimes3},\qquad L=U_d^*,\qquad R=U_d.
\]
The proposed left side is one by
\eqref{eq:V26-entangled-capacity}, while its right side is
\(Cd^{-3}\) by \eqref{eq:V26-protected-PT}.
\end{proof}

\begin{firewall}[Abstract entangler versus the exact YM dualizer]
\label{fire:V26-entangler-scope}
\Cref{cor:V26-norm-no-go} disproves an inference from ordinary norm
data; it does not assert that the unitary \(U_d\) occurs in a Yang--Mills
source fibre.  The exact output dualizer has representation-theoretic
structure.  That structure, rather than the bound \(\|D_U\|\le1\)
alone, must be used.
\end{firewall}

\section{Row-cut decomposable sandwich norm}
\label{sec:V26-decomposable}

Fix the bipartition
\[
 \mathcal H=\mathcal H_a\otimes\mathcal H_{\bar a}
\]
given by a maximal row and its complement.

\begin{definition}[Decomposable sandwich cost]
\label{def:V26-DSC}
Let \(Z\) be an operator on \(\mathcal H\).  A representation
\begin{equation}
 \mathcal W
 =
 \sum_{\nu=1}^{N}
 (L_{a,\nu}\otimes L_{\bar a,\nu})\,
 Z\,
 (R_{a,\nu}\otimes R_{\bar a,\nu})
\label{eq:V26-decomp-sandwich}
\end{equation}
has row-cut cost
\begin{equation}
 \mathfrak d_a(\mathcal W;Z)
 :=
 \sum_{\nu=1}^{N}
 \|L_{a,\nu}\|\|L_{\bar a,\nu}\|
 \|R_{a,\nu}\|\|R_{\bar a,\nu}\|.
\label{eq:V26-DSC}
\end{equation}
The optimal cost is the infimum over all finite representations
\eqref{eq:V26-decomp-sandwich}.  Tensor factors on which an operator
acts trivially are suppressed.
\end{definition}

\begin{lemma}[Partial transpose passes a decomposable sandwich]
\label{lem:V26-DSC-PT}
Every representation \eqref{eq:V26-decomp-sandwich} satisfies
\begin{equation}
 \|\mathcal W^\Gamma\|
 \le
 \mathfrak d_a(\mathcal W;Z)\,
 \|Z^\Gamma\|.
\label{eq:V26-DSC-PT}
\end{equation}
The same cost controls the adjoint orientation.
\end{lemma}

\begin{proof}
Write \(Z=\sum_iA_i\otimes B_i\).  For one summand indexed by \(\nu\),
partial transpose on \(\mathcal H_{\bar a}\) gives
\begin{align*}
 &\left[
 (L_{a,\nu}\otimes L_{\bar a,\nu})
 Z
 (R_{a,\nu}\otimes R_{\bar a,\nu})
 \right]^\Gamma\\
 &\quad=
 (L_{a,\nu}\otimes R_{\bar a,\nu}^{\mathsf T})
 Z^\Gamma
 (R_{a,\nu}\otimes L_{\bar a,\nu}^{\mathsf T}).
\end{align*}
The operator norm of a transpose equals the original norm.
Submultiplicativity and then the triangle inequality yield
\eqref{eq:V26-DSC-PT}.  Taking adjoints exchanges left and right
factors without changing the product of their norms.
\end{proof}

\begin{corollary}[Stability under finite source assembly]
\label{cor:V26-DSC-sum}
If
\(\mathcal W=\sum_{\alpha}\mathcal W_\alpha\) and each
\(\mathcal W_\alpha\) has a sandwich representation around the same
\(Z\), then
\[
 \mathfrak d_a(\mathcal W;Z)
 \le
 \sum_\alpha\mathfrak d_a(\mathcal W_\alpha;Z).
\]
In particular, a fixed arity-five family of uniformly bounded costs
has uniformly bounded assembled cost.
\end{corollary}

\begin{proof}
Concatenate the finite sandwich representations and take the
infimum.
\end{proof}

\begin{example}[Safe separators]
\label{ex:V26-safe-separators}
The cost is at most one for a single row-local sandwich by
contractions.  It remains uniformly bounded for:
\begin{enumerate}[label=\textup{(\roman*)},leftmargin=3.5em]
\item rowwise unitaries and canonical regroupings which do not cross
      the cut;
\item maps acting only on the maximal row or only on its complement;
\item any fixed-dimensional prefix or spectator operator; and
\item fixed finite sums and products of these maps.
\end{enumerate}
\end{example}

\begin{proof}
Items \textup{(i)} and \textup{(ii)} use one term in
\eqref{eq:V26-decomp-sandwich}.  A fixed-dimensional operator has a
finite elementary-tensor expansion whose projective coefficient norm
is a fixed constant, proving \textup{(iii)}.  Expanding finite sums
and products proves \textup{(iv)}.
\end{proof}

\section{The direct dual-separator certificate}
\label{sec:V26-certificate}

Let \(E_3=\{e_1,e_2,e_3\}\) be the three distinct cofinal coordinates
of a \(q=3\) fibre, and let
\[
 Z_{E_3}:=\mathbb H_{e_1}\otimes
            \mathbb H_{e_2}\otimes
            \mathbb H_{e_3}
\]
denote their complete assembled local operators, including a selected
or payer role when it occurs at one of these coordinates.

\begin{definition}[Direct dual-separator certificate]
\label{def:V26-DDSC}
A resolved \(q=3\) source has \(DDSC_3(C,\tau)\) across the maximal
row cut if:
\begin{enumerate}[label=\textup{(\roman*)},leftmargin=3.5em]
\item \(a_{a,e_k}\ge\tau M\) for \(k=1,2,3\);
\item in one fixed output-dual block,
\begin{equation}
 \|Z_{E_3}^{\Gamma}\|
 \le C M^{-3/2};
\label{eq:V26-three-PT}
\end{equation}
\item the complete unnormalized source, with its actual dualizer and
      all signs assembled, has a representation around \(Z_{E_3}\)
      of cost
\begin{equation}
 \mathfrak d_a(\mathbb W^{\rm raw};Z_{E_3})\le C;
\label{eq:V26-source-cost}
\end{equation}
\item the adjoint word has the same property.
\end{enumerate}
\end{definition}

\begin{remark}[How the local heat hypothesis is verified]
\label{rem:V26-local-heat}
For an unpaid complete two-input heat moment,
\eqref{eq:V26-three-PT} follows coordinatewise from the archived
V83 complete output-isotypic partial-transpose theorem and the Weyl
dimension lower bound.  The V25 fixed-spin-one calculation verifies
the selected paid covariance packet.  Other selected or paid local
forms must use their own complete partial-transpose theorem; a square
capacity estimate alone does not imply
\eqref{eq:V26-three-PT}.
\end{remark}

\begin{theorem}[\(DDSC_3\) closes a one-strong \(q=3\) fibre]
\label{thm:V26-DDSC-close}
Suppose a normalized resolved source satisfies \(DDSC_3(C,\tau)\)
and its exchange-rescue weight has the one-strong floor
\begin{equation}
 \mathfrak x_{T_0,r}
 \ge {c\over M\prod_{i=1}^5\Xi_i}.
\label{eq:V26-one-strong-floor}
\end{equation}
Then both direct coefficient orientations obey
\begin{equation}
 \kappa_\otimes^{\rm abs}(\widetilde{\mathbb W})
 \le C_{G,s,\tau,c}\mathfrak x_{T_0,r}.
\label{eq:V26-direct-close}
\end{equation}
\end{theorem}

\begin{proof}
For a row-product density \(\varrho\), partial transpose on the
complement of \(a\) remains a density and
\[
 |\Tr(\varrho\mathbb W^{\rm raw})|
 \le\|(\mathbb W^{\rm raw})^\Gamma\|.
\]
Apply \Cref{lem:V26-DSC-PT} and
\eqref{eq:V26-three-PT}--%
\eqref{eq:V26-source-cost} to obtain
\[
 \kappa_\otimes^{\rm abs}(\mathbb W^{\rm raw})
 \le C^2M^{-3/2}.
\]
After the five row normalizations,
\[
 \kappa_\otimes^{\rm abs}(\widetilde{\mathbb W})
 \le {C^2M^{-3/2}\over\prod_i\Xi_i}
 \le
 {C^2M^{-1}\over\prod_i\Xi_i}.
\]
Compare with \eqref{eq:V26-one-strong-floor}.  The adjoint certificate
gives the opposite orientation with the same constants.
\end{proof}

\begin{corollary}[Certified direct source assembly]
\label{cor:V26-assembled-close}
If every nonzero tree/lift term at fixed \((r,\rho,\pi)\) has a
\(DDSC_3\) representation around the same three-coordinate packet,
with uniformly bounded total cost, then the assembled source satisfies
\(DXQTRC_5\) in that fibre.
\end{corollary}

\begin{proof}
Use \Cref{cor:V26-DSC-sum} before applying
\Cref{thm:V26-DDSC-close}, and sum the corresponding finite
exchange-rescue weights on the right.
\end{proof}

\section{Temporal placement does not certify the dualizer}
\label{sec:V26-temporal}

The seven \(q=3\) temporal classes are
\[
\begin{split}
 &(3,0,0),(2,0,1),(1,0,2),(0,0,3),\\
 &(2,1,0),(1,1,1),(0,1,2).
\end{split}
\]

\begin{proposition}[Temporal class is insufficient for \(DDSC_3\)]
\label{prop:V26-temporal-insufficient}
The triple \((p,j,s)\) determines the physical-coordinate locations
of the three responses but does not determine the row-cut sandwich
cost of the physical output dualizer.  No one of the seven temporal
labels alone implies \eqref{eq:V26-source-cost}.
\end{proposition}

\begin{proof}
The temporal label records only whether each coordinate is below, at,
or above the first-connectivity coordinate.  It contains no matrix
data for the contraction \(D_U\) in
\eqref{eq:V24-dual-operator}.  The identity dualizer and the abstract
entangling conjugation in
\Cref{lem:V26-three-entangler} can be attached to the same
three-coordinate protected packet without changing any of its
temporal locations, yet their row-cut effects differ by a factor
\(d^3\).  This proves logical nonimplication.  The abstract entangler
is used only to separate the data sets, as stated in
\Cref{fire:V26-entangler-scope}.
\end{proof}

\begin{corollary}[The V23 identity-dual fibre is \(DDSC_3\)-safe]
\label{cor:V26-V23-safe}
The V23 fibre with identity dualizer satisfies \(DDSC_3\) and hence
its direct closure in \Cref{thm:V25-direct-close} is an instance of
\Cref{thm:V26-DDSC-close}.
\end{corollary}

\begin{proof}
The threefold partial-transpose bound is
\eqref{eq:V25-three-P1-PT}.  The identity dualizer is row local.  The
remaining prefix operator is fixed-dimensional and therefore has
bounded cost by \Cref{ex:V26-safe-separators}.  The same statements
hold after taking adjoints.
\end{proof}

\begin{assessment}[Exact residual after V26]
\label{ass:V26-residual}
The \(q=3\) problem is no longer a question about the number of heat
slots.  Three slots are more than sufficient for the direct linear
target once their partial-transpose gain survives the output dualizer.
The residual consists precisely of dimension-growing output blocks
whose embedded contraction \(W_U(D_U^*\otimes I)W_U^*\) has no proved
uniform row-cut decomposable sandwich cost, or local selected/payer
packets for which the complete partial-transpose estimate itself is
missing.
\end{assessment}

\section{V26 result ledger and V27 acquisition order}
\label{sec:V26-ledger}

\begin{theorem}[Fourth-series V26 result ledger]
\label{thm:V26-results}
Relative to V1--V25, V26 proves:
\begin{enumerate}[label=\textup{(V26.\arabic*)},leftmargin=6.3em]
\item a threefold entangler counterexample to every ordinary-norm
      direct-separator shortcut;
\item the row-cut decomposable sandwich cost and its exact
      partial-transpose inequality;
\item stability of that certificate under finite signed source
      assembly;
\item the \(DDSC_3\) certificate;
\item direct closure of every one-strong \(q=3\) fibre satisfying that
      certificate;
\item the insufficiency of temporal labels alone; and
\item certification of the V23 identity-dual packet as the first
      nontrivial safe example.
\end{enumerate}
\end{theorem}

\begin{proof}
These are
\Cref{lem:V26-three-entangler,
cor:V26-norm-no-go,
lem:V26-DSC-PT,
cor:V26-DSC-sum,
def:V26-DDSC,
thm:V26-DDSC-close,
cor:V26-assembled-close,
prop:V26-temporal-insufficient,
cor:V26-V23-safe}.
\end{proof}

\begin{programme}[V27: actual output-block decomposable norm]
\label{prog:V26-V27}
\begin{enumerate}[leftmargin=3em]
\item Write
      \(W_U(D_U^*\otimes I_{M_U})W_U^*\) across the maximal-row cut
      using Clebsch--Gordan intertwiners, without replacing \(D_U\) by
      its ordinary norm.
\item Compute or bound its optimal decomposable sandwich cost uniformly
      over \(\|D_U\|\le1\).
\item Start with \(SU(2)\), quotient order two, and fixed low total
      output \(U=1\), where multiplicity is absent.
\item Test the bound on the V26 entangler and the V23 identity family.
\item If the cost grows with \(d_j\), combine the dualizer with the
      adjacent heat block before decomposing; do not estimate them
      separately.
\item Propagate every proved block estimate through
      \Cref{thm:V26-DDSC-close} and update the exact \(q=3\) residual.
\end{enumerate}
\end{programme}

\begin{firewall}[Claim boundary after fourth-series V26]
\label{fire:V26-claim-boundary}
V26 gives an exact sufficient separator certificate and proves why an
ordinary contraction bound cannot replace it.  It does not prove the
certificate uniformly for all physical output dualizers or all local
selected/payer packets, and therefore does not close \(q=3\) or global
quintic MTA.  The unresolved output-block costs, \(q=0,1,2\) sectors,
all-order MTA, WUFC, CPM, cutoff removal, reflection positivity,
Osterwalder--Schrader reconstruction, construction of the continuum
Yang--Mills measure, and a uniform positive continuum spectral gap
remain open.  The full Yang--Mills mass gap has not been proved.
\end{firewall}

\paragraph{Parent-version record.}
V26 was formed from
\texttt{Renewal\_Yang\_Mills\_Mass\_Gap\_Research\_Notes\_V25.tex}.
The V25 parent had \(12\,380\) logical source lines,
\(438\,116\) bytes, and SHA--256
\[
 \texttt{9de6fc8ebde4c3ab6b5aa7659053b13e1f740797d9580484be1517a87b02ed36}.
\]
The next compulsory companion audit is V30.

% =============================================================================
% END APPENDED FOURTH-SERIES V26 MATERIAL
% =============================================================================

\clearpage
% =============================================================================
% BEGIN APPENDED FOURTH-SERIES V27 MATERIAL
% =============================================================================

\chapter{V27: Absorption of every output dualizer on the three
spin-one block}
\label{ch:V27}

\section{Purpose and change of tactic}
\label{sec:V27-purpose}

V26 asks for a uniform decomposable cost of the embedded output
dualizer.  That cost need not be the right quantity: a low-output
Clebsch--Gordan isometry is highly entangled when viewed by itself.
The correct object is the \emph{dualizer already compressed to the
low-output heat block}.

For the spin-one channel, every Clebsch--Gordan basis vector has
largest Schmidt coefficient \(O(d_j^{-1/2})\).  A matrix unit between
two such vectors has partial-transpose norm \(O(d_j^{-1})\).
Therefore an arbitrary contraction on the fixed three-dimensional
spin-one output, after embedding into \(V_j\otimes V_j\), retains the
full inverse-dimension gain.  Tensoring three coordinates gives
\(O(d_j^{-3})\).  This proves the direct V23 tree card uniformly over
\emph{every} physical output dualizer, not only the identity.

\section{Partial transpose of a bipartite rank-one operator}
\label{sec:V27-rank-one}

Choose bases and identify a vector
\[
 \psi=\sum_{i,\alpha}A_{i\alpha}e_i\otimes f_\alpha
\]
with its coefficient matrix \(A:\mathcal H_B\to\mathcal H_A\).
Its Schmidt coefficients are the singular values of \(A\).

\begin{lemma}[Exact cross-rank-one partial-transpose norm]
\label{lem:V27-cross-rank-one}
Let \(\psi,\varphi\in\mathcal H_A\otimes\mathcal H_B\) have
coefficient matrices \(A,B\).  Then
\begin{equation}
 \boxed{
 \left\|
  (|\psi\rangle\langle\varphi|)^{\Gamma_B}
 \right\|
 =
 \|A\|\,\|B\|.}
\label{eq:V27-cross-rank-one}
\end{equation}
Equivalently, the right side is the product of the largest Schmidt
coefficients of \(\psi\) and \(\varphi\).
\end{lemma}

\begin{proof}
On basis vectors, direct expansion gives
\[
 (|\psi\rangle\langle\varphi|)^{\Gamma_B}
 =
 \sum_{i,j,\alpha,\beta}
 A_{i\alpha}\overline{B_{j\beta}}\,
 |e_i\rangle\langle e_j|
 \otimes
 |f_\beta\rangle\langle f_\alpha|.
\]
Let \(S(e_j\otimes f_\alpha)=f_\alpha\otimes e_j\) be the tensor
swap, with the evident source and target identifications.  The
displayed operator is
\[
 (A\otimes B^*)S.
\]
The swap is unitary, so its norm is
\(\|A\otimes B^*\|=\|A\|\|B\|\), proving
\eqref{eq:V27-cross-rank-one}.
\end{proof}

\begin{corollary}[A fixed output matrix inherits two Schmidt
coefficients]
\label{cor:V27-fixed-output-matrix}
Let \(T:E\to\mathcal H_A\otimes\mathcal H_B\) be an isometry, choose
an orthonormal basis \((u_\alpha)_{\alpha=1}^m\) of \(E\), and put
\(\psi_\alpha=Tu_\alpha\).  If every \(\psi_\alpha\) has largest
Schmidt coefficient at most \(\varepsilon\), then for every
\(D\in\mathcal B(E)\),
\begin{equation}
 \|(TDT^*)^\Gamma\|
 \le m^2\varepsilon^2\|D\|.
\label{eq:V27-fixed-output-matrix}
\end{equation}
\end{corollary}

\begin{proof}
Write
\[
 D=\sum_{\alpha,\beta}D_{\alpha\beta}
 |u_\alpha\rangle\langle u_\beta|.
\]
Then
\[
 TDT^*
 =
 \sum_{\alpha,\beta}D_{\alpha\beta}
 |\psi_\alpha\rangle\langle\psi_\beta|.
\]
Every matrix entry obeys
\(|D_{\alpha\beta}|\le\|D\|\).  Apply
\Cref{lem:V27-cross-rank-one} to each of the \(m^2\) terms and use the
triangle inequality.
\end{proof}

\section{One spin-one output with an arbitrary contraction}
\label{sec:V27-one-spin}

Let
\[
 T_j:V_1\longrightarrow V_j\otimes V_j,
 \qquad
 T_j|1,M\rangle=|1,M\rangle_j
\]
be the multiplicity-one total-spin-one Clebsch--Gordan isometry.

\begin{theorem}[Arbitrary spin-one dualizer retains the half-power]
\label{thm:V27-one-dualizer}
For \(j\ge1\) and every \(D\in\mathcal B(V_1)\),
\begin{equation}
 \boxed{
 \|(T_jDT_j^*)^\Gamma\|
 \le {27\over d_j}\|D\|.}
\label{eq:V27-one-dualizer}
\end{equation}
The transpose may be taken on either spin-\(j\) input.
\end{theorem}

\begin{proof}
By \Cref{lem:V25-CG-upper}, each of the three vectors
\(|1,M\rangle_j\) has largest Schmidt coefficient at most
\(\sqrt{3/d_j}\).  Apply
\Cref{cor:V27-fixed-output-matrix} with
\(m=3\) and \(\varepsilon^2=3/d_j\), giving
\eqref{eq:V27-one-dualizer}.  Partial transpose on the other input is
the full transpose of the displayed partial transpose, up to the
chosen bases; full transpose preserves operator norm.
\end{proof}

\begin{remark}[Why this is stronger than a dualizer cost]
\label{rem:V27-combined}
The row-cut decomposable norm of
\(T_jDT_j^*\) may be large when its entangled isometry is expanded
first.  Theorem \ref{thm:V27-one-dualizer} never estimates \(T_j\) or
\(D\) separately.  The two Schmidt coefficients in each matrix unit
are retained together, producing \(d_j^{-1}\).
\end{remark}

\section{Three spin-one outputs}
\label{sec:V27-three-spin}

Put
\[
 \mathbf T_j:=T_j^{\otimes3}:
 V_1^{\otimes3}\longrightarrow
 (V_j\otimes V_j)^{\otimes3}.
\]
The bipartition groups the three first spin-\(j\) inputs against the
three second inputs.

\begin{theorem}[Every three-coordinate dualizer retains
\(d_j^{-3}\)]
\label{thm:V27-three-dualizer}
There is an absolute constant \(C_1\) such that for every
\(D\in\mathcal B(V_1^{\otimes3})\),
\begin{equation}
 \boxed{
 \|(\mathbf T_jD\mathbf T_j^*)^\Gamma\|
 \le C_1d_j^{-3}\|D\|.}
\label{eq:V27-three-dualizer}
\end{equation}
\end{theorem}

\begin{proof}
Use the tensor basis
\(|1,M_1\rangle\otimes|1,M_2\rangle\otimes|1,M_3\rangle\).
There are \(m=3^3=27\) basis vectors.  Each embedded vector is the
tensor product of three Clebsch--Gordan vectors from
\Cref{thm:V27-one-dualizer}; its largest Schmidt coefficient is at
most
\[
 \left({3\over d_j}\right)^{3/2}.
\]
Apply \Cref{cor:V27-fixed-output-matrix} with
\(m=27\).  One may take
\[
 C_1=27^2\,3^3.
\]
The numerical constant is immaterial and independent of \(j,D\).
\end{proof}

\begin{corollary}[Fixed additional output factors are harmless]
\label{cor:V27-fixed-outputs}
Let \(E_0\) be any fixed finite-dimensional prefix/spectator output
space and let
\[
 T_0:E_0\to\mathcal H_{0,a}\otimes\mathcal H_{0,\bar a}
\]
be a fixed isometry.  There is \(C_{E_0,T_0}<\infty\) such that, for
every contraction \(D\) on \(E_0\otimes V_1^{\otimes3}\),
\begin{equation}
 \left\|
  \left[
  (T_0\otimes\mathbf T_j)
  D
  (T_0\otimes\mathbf T_j)^*
  \right]^\Gamma
 \right\|
 \le C_{E_0,T_0}d_j^{-3}.
\label{eq:V27-fixed-prefix-dual}
\end{equation}
\end{corollary}

\begin{proof}
Choose a fixed basis of \(E_0\).  The embedded basis vectors of
\(T_0\) have largest Schmidt coefficients at most one, while the
three-spin factor contributes
\((3/d_j)^{3/2}\).  Apply
\Cref{cor:V27-fixed-output-matrix} on the tensor output space.  Its
dimension is fixed independently of \(j\), so all resulting factors
are absorbed into \(C_{E_0,T_0}\).
\end{proof}

\section{Application to the exact physical output duality}
\label{sec:V27-physical}

On the V23 source fibre, resolve the joining coordinate and the two
suffix coordinates on their complete spin-one isotypes.  All prefix
representations have fixed spin \(\ell\), so their retained output
space is a fixed \(E_0\).  Any final irreducible output block
\(H_U\otimes M_U\) is embedded in
\[
 E_0\otimes V_1^{\otimes3}
\]
through a fixed-dimensional isometry \(J_U\).  For a dual contraction
\(D_U\) on \(H_U\), extend
\[
 J_U(D_U^*\otimes I_{M_U})J_U^*
\]
by zero to the full fixed output space.  Its norm is at most one.

\begin{theorem}[Uniform physical-dual closure of the V23 packet]
\label{thm:V27-V23-all-duals}
For every admissible physical output block and every dual contraction
\(\|D_U\|\le1\), the normalized V23 word obeys
\begin{equation}
 \kappa_\otimes^{\rm abs}
  (\widetilde{\mathbb W}_{j,\boldsymbol D})
 \le
 {C_{\ell,s}d_j^{-3}\over\Xi_a\Xi_cB^3}
 \le
 C_{\ell,s}\mathfrak x_{T_*,r}.
\label{eq:V27-V23-all-duals}
\end{equation}
The estimate is uniform over the complete dual family appearing in
the exact Fourier trace-norm identity.
\end{theorem}

\begin{proof}
On the three retained spin-one blocks, the paid root covariance and
the two suffix heats have scalar multiplier
\[
 r_1(s)\Delta_j(s)q_1(s)^2,
\]
whose absolute value is bounded uniformly in \(j\).  Absorb this
scalar, the fixed prefix multiplier, and the embedded contraction
\(J_U(D_U^*\otimes I)J_U^*\) into the operator \(D\) in
\Cref{cor:V27-fixed-outputs}.  Its norm is bounded by a constant
depending only on \(\ell,s\).  Equation
\eqref{eq:V27-fixed-prefix-dual} therefore bounds the full
unnormalized partial-transpose norm by
\(C_{\ell,s}d_j^{-3}\).

The partial-transpose trace estimate and division by
\[
 \Xi_a\Xi_b\Xi_c\Xi_d\Xi_e
 =
 \Xi_a\Xi_cB^3
\]
give the first inequality.  The comparison with the exact exchange
weight is the final step of
\Cref{thm:V25-direct-close}; it is uniform in the dualizer because
the preceding constant is.
\end{proof}

\begin{corollary}[The exact Fourier supremum costs no power on this
packet]
\label{cor:V27-Fourier-sup}
Taking the supremum over all dual contractions and summing the fixed
finite output blocks preserves
\eqref{eq:V27-V23-all-duals}.  Thus the complete physical Fourier
coefficient of the V23 fibre satisfies the direct tree card.
\end{corollary}

\begin{proof}
The estimate in
\Cref{thm:V27-V23-all-duals} is uniform over \(\|D_U\|\le1\), exactly
the family in the archived block-duality identity.  The retained
prefix and three spin-one output tensor product has fixed dimension,
so the number of its final irreducible blocks and multiplicities is
independent of \(j\).  Absorb that finite sum into \(C_{\ell,s}\).
\end{proof}

\begin{firewall}[Fixed low output is essential]
\label{fire:V27-low-output}
The proof uses that the three local outputs have fixed spin one.
If an output spin grows with \(j\), the output dimension and the
largest Clebsch--Gordan coefficient need not obey the constants used
above.  V27 does not sum all local output spins before establishing a
heat-weighted version of
\eqref{eq:V27-three-dualizer}.
\end{firewall}

\section{Upgrade of the V26 certificate}
\label{sec:V27-upgrade}

\begin{definition}[Low-output dual absorption]
\label{def:V27-LODA}
A \(q=3\) fibre has \(LODA_3\) if its three cofinal local outputs range
over a fixed finite set, its remaining output factors are
fixed-dimensional, and the complete physical dualizer compressed to
that output space obeys
\[
 \|(\mathbf T_{\rm out}D\mathbf T_{\rm out}^*)^\Gamma\|
 \le C M^{-3/2}
\]
uniformly for all admissible dual contractions.
\end{definition}

\begin{proposition}[\(LODA_3\) implies the direct card]
\label{prop:V27-LODA-close}
If a one-strong \(q=3\) source has \(LODA_3\), uniformly bounded
fixed-dimensional prefix factors, and the exchange floor
\eqref{eq:V26-one-strong-floor}, then it satisfies \(DXQTRC_5\).
\end{proposition}

\begin{proof}
The \(LODA_3\) inequality is already the full dualizer-plus-response
partial-transpose bound, so no separate decomposable cost is needed.
Apply the normalization and tree-floor comparison in the proof of
\Cref{thm:V26-DDSC-close}.
\end{proof}

\begin{assessment}[Exact residual after V27]
\label{ass:V27-residual}
All fixed-low-output versions of the sharp V23 packet are now closed
uniformly over physical dualizers.  A surviving \(q=3\) obstruction
must therefore have at least one growing local output, a
dimension-growing prefix output, or a source separator not absorbed
into the same low-output embedding.  The next task is the
heat-weighted sum over growing output spins; estimating its dualizer
before the heat weight would repeat the V26 error.
\end{assessment}

\section{V27 result ledger and V28 acquisition order}
\label{sec:V27-ledger}

\begin{theorem}[Fourth-series V27 result ledger]
\label{thm:V27-results}
Relative to V1--V26, V27 proves:
\begin{enumerate}[label=\textup{(V27.\arabic*)},leftmargin=6.3em]
\item the exact partial-transpose norm of a cross rank-one bipartite
      operator;
\item a fixed-output matrix lemma retaining two Schmidt coefficients;
\item the uniform \(d_j^{-1}\) bound for every contraction on a
      spin-one output;
\item the uniform \(d_j^{-3}\) bound for an arbitrary contraction
      coupling three spin-one outputs;
\item stability under every fixed prefix/spectator output;
\item direct closure of the V23 source for the full physical Fourier
      dual supremum; and
\item the \(LODA_3\) certificate and its direct-card compiler.
\end{enumerate}
\end{theorem}

\begin{proof}
These are
\Cref{lem:V27-cross-rank-one,
cor:V27-fixed-output-matrix,
thm:V27-one-dualizer,
thm:V27-three-dualizer,
cor:V27-fixed-outputs,
thm:V27-V23-all-duals,
cor:V27-Fourier-sup,
def:V27-LODA,
prop:V27-LODA-close}.
\end{proof}

\begin{programme}[V28: heat-summed growing outputs]
\label{prog:V27-V28}
\begin{enumerate}[leftmargin=3em]
\item Replace the fixed spin-one isometry \(T_j\) by the complete
      family \(T_{j,K}:V_K\to V_j\otimes V_j\).
\item Prove a bound for
      \((T_{j,K}D_KT_{j,K}^*)^\Gamma\) retaining both the input
      dimension \(d_j^{-1}\) and explicit polynomial dependence on
      \(2K+1\).
\item Insert the heat or paid covariance multiplier before summing in
      \(K\).
\item Use Gaussian Casimir summability to show that the complete
      output-dual block retains \(Cd_j^{-1}\) at one coordinate.
\item Tensor three coordinates only after the complete one-coordinate
      output sum is controlled.
\item Propagate the result through \(LODA_3\) and update the exact
      growing-output \(q=3\) residual.
\end{enumerate}
\end{programme}

\begin{firewall}[Claim boundary after fourth-series V27]
\label{fire:V27-claim-boundary}
V27 closes the direct coefficient card for the complete physical
dual family on every fixed-low-output V23-type packet.  It does not
control growing output spins, arbitrary dimension-growing prefix
outputs, all \(q=3\) fibres, or global quintic MTA.  The unresolved
heat-summed output blocks, \(q=0,1,2\) sectors, all-order MTA, WUFC,
CPM, cutoff removal, reflection positivity, Osterwalder--Schrader
reconstruction, construction of the continuum Yang--Mills measure,
and a uniform positive continuum spectral gap remain open.  The full
Yang--Mills mass gap has not been proved.
\end{firewall}

\paragraph{Parent-version record.}
V27 was formed from
\texttt{Renewal\_Yang\_Mills\_Mass\_Gap\_Research\_Notes\_V26.tex}.
The V26 parent had \(12\,839\) logical source lines,
\(453\,693\) bytes, and SHA--256
\[
 \texttt{bdad4e52745aa0e44d846750f2925aed713c1bb04a6a43358c79db003d57911e}.
\]
The next compulsory companion audit is V30.

% =============================================================================
% END APPENDED FOURTH-SERIES V27 MATERIAL
% =============================================================================

\clearpage
% =============================================================================
% BEGIN APPENDED FOURTH-SERIES V28 MATERIAL
% =============================================================================

\chapter{V28: Heat-summed output dualizers and the local-Casimir
preserving class}
\label{ch:V28}

\section{Purpose}
\label{sec:V28-purpose}

V27 absorbs every dual contraction on three fixed spin-one outputs.
The next obstruction is a local output spin \(K\) which grows with the
input spin \(j\).  A direct Racah-formula estimate is unnecessary.
The complete spin-\(K\) projection has a scalar partial trace, and
that identity bounds the largest Schmidt coefficient of \emph{every}
vector in its range.  The resulting cost is polynomial in
\(d_K=2K+1\), hence summable against the heat multiplier and also
against the paid covariance multiplier.

This proves a complete one-coordinate output-dual card and its
three-coordinate tensor version whenever the assembled separator
preserves the three local Casimir labels.  The exact residual is
therefore cross-label mixing in multiplicity space, not growth of an
individual output spin.

\section{A basis-free Schmidt bound for a complete isotype}
\label{sec:V28-Schmidt}

Let
\[
 T_{j,K}:V_K\longrightarrow V_j\otimes V_j
\]
be the multiplicity-one \(SU(2)\) Clebsch--Gordan isometry, and let
\[
 P_{j,K}:=T_{j,K}T_{j,K}^*,
 \qquad d_K:=2K+1.
\]

\begin{lemma}[Scalar marginal of a total-spin projection]
\label{lem:V28-scalar-marginal}
For \(0\le K\le2j\),
\begin{equation}
 \operatorname{Tr}_2P_{j,K}
 =
 {d_K\over d_j}I_{V_j},
 \qquad
 \operatorname{Tr}_1P_{j,K}
 =
 {d_K\over d_j}I_{V_j}.
\label{eq:V28-scalar-marginal}
\end{equation}
\end{lemma}

\begin{proof}
The projection \(P_{j,K}\) commutes with the diagonal \(SU(2)\)
action.  Its partial trace over either input therefore commutes with
the irreducible spin-\(j\) action on the other input, and is scalar by
Schur's lemma.  Its trace is
\(\operatorname{rank}P_{j,K}=d_K\).  Since
\(\operatorname{Tr}I_{V_j}=d_j\), the scalar is \(d_K/d_j\).
\end{proof}

\begin{corollary}[Uniform Schmidt coefficient on the whole isotype]
\label{cor:V28-Schmidt}
Every unit vector \(\psi\in\operatorname{Ran}P_{j,K}\) has largest
Schmidt coefficient at most
\begin{equation}
 \sqrt{d_K/d_j}.
\label{eq:V28-Schmidt}
\end{equation}
\end{corollary}

\begin{proof}
The positive rank-one operator
\(|\psi\rangle\langle\psi|\) is bounded above by \(P_{j,K}\).
Positive partial trace and
\Cref{lem:V28-scalar-marginal} give
\[
 \operatorname{Tr}_2|\psi\rangle\langle\psi|
 \preccurlyeq {d_K\over d_j}I.
\]
The eigenvalues of the reduced density matrix are the squared Schmidt
coefficients of \(\psi\).  Its largest eigenvalue is therefore at most
\(d_K/d_j\).
\end{proof}

\begin{theorem}[One fixed output with arbitrary dual contraction]
\label{thm:V28-fixed-K}
For every \(D_K\in\mathcal B(V_K)\),
\begin{equation}
 \boxed{
 \left\|
  (T_{j,K}D_KT_{j,K}^*)^\Gamma
 \right\|
 \le {d_K^3\over d_j}\|D_K\|.}
\label{eq:V28-fixed-K}
\end{equation}
\end{theorem}

\begin{proof}
Choose any orthonormal basis of \(V_K\).  Its \(d_K\) embedded basis
vectors all satisfy
\eqref{eq:V28-Schmidt}.  Apply
\Cref{cor:V27-fixed-output-matrix} with
\[
 m=d_K,\qquad \varepsilon^2=d_K/d_j.
\]
\end{proof}

\begin{remark}[Consistency with the sharp spin-one result]
\label{rem:V28-spin-one}
At \(K=1\), \eqref{eq:V28-fixed-K} gives
\(27d_j^{-1}\|D_1\|\), exactly
\Cref{thm:V27-one-dualizer}.  At \(K=2j\), the polynomial factor is
large enough to allow the coherent product vector in the top-spin
isotype; the heat weight will suppress this sector.
\end{remark}

\section{The complete one-coordinate heat sum}
\label{sec:V28-one-sum}

For a contraction family
\(\boldsymbol D=(D_K)_{K=0}^{2j}\), define
\begin{align}
 \mathcal H_j(\boldsymbol D)
 &:=
 \sum_{K=0}^{2j}q_K(s)
 T_{j,K}D_KT_{j,K}^*,
\label{eq:V28-dual-heat}\\
 \mathcal D_j(\boldsymbol D)
 &:=
 \sum_{K=0}^{2j}
 \bigl(q_K(s)-q_j(s)^2\bigr)
 T_{j,K}D_KT_{j,K}^*.
\label{eq:V28-dual-covariance}
\end{align}
For the paid covariance put
\begin{equation}
 \mathcal D_j^{\rm pay}(\boldsymbol D)
 :=
 \sum_{K=1}^{2j}
 r_K(s)\bigl(q_K(s)-q_j(s)^2\bigr)
 T_{j,K}D_KT_{j,K}^*,
\quad
 r_K(s):={K(K+1)\over1-e^{-sK(K+1)}}.
\label{eq:V28-dual-paid}
\end{equation}

\begin{theorem}[Heat-summed arbitrary-output dualizer]
\label{thm:V28-one-summed}
For fixed \(s>0\), there is \(C_s<\infty\) such that, uniformly in
\(j\) and in all families \(\|D_K\|\le1\),
\begin{align}
 \|\mathcal H_j(\boldsymbol D)^\Gamma\|
 &\le {C_s\over d_j},
\label{eq:V28-heat-summed}\\
 \|\mathcal D_j(\boldsymbol D)^\Gamma\|
 &\le {C_s\over d_j},
\label{eq:V28-cov-summed}\\
 \|(\mathcal D_j^{\rm pay}(\boldsymbol D))^\Gamma\|
 &\le {C_s\over d_j}.
\label{eq:V28-paid-summed}
\end{align}
The same conclusion holds for the paid heat family with multiplier
\(r_K(s)q_K(s)\), \(K\ge1\).
\end{theorem}

\begin{proof}
By \Cref{thm:V28-fixed-K} and the triangle inequality,
\[
 \|\mathcal H_j(\boldsymbol D)^\Gamma\|
 \le
 {1\over d_j}
 \sum_{K=0}^{2j}d_K^3e^{-sK(K+1)}
 \le {C_s\over d_j},
\]
because the infinite Gaussian-polynomial series converges.

For the unpaid covariance,
\[
 \sum_{K=0}^{2j}
 d_K^3|q_K-q_j^2|
 \le
 \sum_{K=0}^{\infty}d_K^3q_K
 +
 q_j^2\sum_{K=0}^{2j}d_K^3.
\]
The first term is finite.  The second is uniformly bounded because
\[
 \sum_{K=0}^{2j}d_K^3\le C(1+j)^4,
 \qquad
 q_j^2=e^{-2sj(j+1)}.
\]
This proves \eqref{eq:V28-cov-summed}.

For \(K\ge1\),
\[
 r_K(s)
 \le {K(K+1)\over1-e^{-2s}}
 \le C_s(1+K)^2.
\]
Therefore
\[
 \sum_{K=1}^{\infty}d_K^3r_Kq_K<\infty,
\]
while
\[
 q_j^2\sum_{K=1}^{2j}d_K^3r_K
 \le
 C_se^{-2sj(j+1)}(1+j)^6
\]
is uniformly bounded.  This proves
\eqref{eq:V28-paid-summed}.  Omitting the replicated endpoint term
gives the paid heat assertion.
\end{proof}

\begin{firewall}[The output sum is taken after dualizer absorption]
\label{fire:V28-order}
Theorem \ref{thm:V28-one-summed} first combines \(D_K\) with the
spin-\(K\) isometry and only then sums the heat-weighted
partial-transpose norms.  Replacing \(D_K\) by its norm before the
embedding would discard the two Schmidt coefficients responsible for
\(d_j^{-1}\).
\end{firewall}

\section{Three coordinates with preserved local Casimirs}
\label{sec:V28-three}

For a tuple
\(\boldsymbol K=(K_1,K_2,K_3)\), put
\[
 E_{\boldsymbol K}:=
 V_{K_1}\otimes V_{K_2}\otimes V_{K_3},
 \qquad
 m_{\boldsymbol K}:=\prod_{\nu=1}^3d_{K_\nu},
\]
and
\[
 T_{j,\boldsymbol K}
 :=
 T_{j,K_1}\otimes T_{j,K_2}\otimes T_{j,K_3}.
\]

\begin{lemma}[Three-output tuple bound]
\label{lem:V28-tuple}
For every
\(D_{\boldsymbol K}\in\mathcal B(E_{\boldsymbol K})\),
\begin{equation}
 \left\|
  (T_{j,\boldsymbol K}D_{\boldsymbol K}
   T_{j,\boldsymbol K}^*)^\Gamma
 \right\|
 \le
 {m_{\boldsymbol K}^3\over d_j^3}
 \|D_{\boldsymbol K}\|.
\label{eq:V28-tuple}
\end{equation}
\end{lemma}

\begin{proof}
In a tensor product of bases used in
\Cref{cor:V28-Schmidt}, every embedded basis vector has largest
Schmidt coefficient at most
\[
 \prod_{\nu=1}^3\sqrt{d_{K_\nu}/d_j}
 =
 \sqrt{m_{\boldsymbol K}/d_j^3}.
\]
The output space has dimension \(m_{\boldsymbol K}\).  Apply
\Cref{cor:V27-fixed-output-matrix}.
\end{proof}

\begin{definition}[Local-Casimir preserving dual separator]
\label{def:V28-LCP}
An assembled three-coordinate output dual separator is
local-Casimir preserving if, after the common resolution of the three
local total Casimirs, it is block diagonal:
\begin{equation}
 D=\bigoplus_{\boldsymbol K}D_{\boldsymbol K},
 \qquad
 \sup_{\boldsymbol K}\|D_{\boldsymbol K}\|\le L.
\label{eq:V28-LCP}
\end{equation}
Fixed-dimensional prefix output factors may be included in each block
and absorbed into \(L\).
\end{definition}

\begin{theorem}[Three complete output sums with an \(LCP\) separator]
\label{thm:V28-three-summed}
At each of three coordinates choose one of the complete heat,
covariance, paid heat, or paid covariance coefficient families in
\Cref{thm:V28-one-summed}, with only one payer overall.  If the
remaining assembled dual separator is \(LCP\) with bound \(L\), then
the complete three-coordinate dual operator satisfies
\begin{equation}
 \boxed{
 \|\mathcal Z_j^\Gamma\|
 \le {C_{s,L}\over d_j^3}.}
\label{eq:V28-three-summed}
\end{equation}
\end{theorem}

\begin{proof}
Let \(b_\nu(K)\) be the scalar multiplier at coordinate \(\nu\).
The proof of \Cref{thm:V28-one-summed} shows
\begin{equation}
 \sum_{K=0}^{2j}d_K^3|b_\nu(K)|
 \le C_s
\label{eq:V28-one-polynomial-sum}
\end{equation}
for each allowed family, uniformly in \(j\).  By the \(LCP\)
decomposition, the full operator is the orthogonal sum over tuples of
\[
 \left(\prod_{\nu=1}^3b_\nu(K_\nu)\right)
 T_{j,\boldsymbol K}D_{\boldsymbol K}
 T_{j,\boldsymbol K}^*.
\]
Apply \Cref{lem:V28-tuple}, then enlarge the norm of the direct sum by
the sum of its block norms:
\begin{align*}
 \|\mathcal Z_j^\Gamma\|
 &\le
 {L\over d_j^3}
 \sum_{\boldsymbol K}
 \prod_{\nu=1}^3
 d_{K_\nu}^3|b_\nu(K_\nu)|\\
 &=
 {L\over d_j^3}
 \prod_{\nu=1}^3
 \left(
  \sum_{K_\nu=0}^{2j}
  d_{K_\nu}^3|b_\nu(K_\nu)|
 \right)
 \le {C_{s,L}\over d_j^3}.
\end{align*}
\end{proof}

\begin{corollary}[Growing outputs close when local Casimirs are
preserved]
\label{cor:V28-LCP-close}
A one-strong \(q=3\) source whose three cofinal factors have the forms
listed in \Cref{thm:V28-three-summed}, whose other output factors are
fixed-dimensional, and whose assembled separator is \(LCP\), satisfies
\(LODA_3\) and hence \(DXQTRC_5\).
\end{corollary}

\begin{proof}
Because \(d_j\asymp M^{1/2}\) on a cofinal equal-spin family,
\eqref{eq:V28-three-summed} is the \(CM^{-3/2}\) bound in
\Cref{def:V27-LODA}.  Fixed-dimensional factors alter only the
constant.  Apply \Cref{prop:V27-LODA-close}.
\end{proof}

\section{The precise cross-label residual}
\label{sec:V28-residual}

\begin{proposition}[Individual growing outputs are no longer an
obstruction]
\label{prop:V28-no-individual-output}
Failure of a direct three-slot proof cannot be attributed solely to
large local output spins or to the polynomial growth of their
Clebsch--Gordan spaces.  Heat summation absorbs those costs uniformly.
Any remaining output obstruction must involve off-diagonal mixing
between distinct local Casimir tuples, a non-heat-weighted
dimension-growing prefix factor, or a local operator outside the four
families in \Cref{thm:V28-three-summed}.
\end{proposition}

\begin{proof}
Within every fixed tuple the bound is
\Cref{lem:V28-tuple}; summation of all tuples under \(LCP\) is
\Cref{thm:V28-three-summed}.  Therefore a failure must violate one of
the hypotheses used there, and the listed three mechanisms are
exactly those hypotheses.
\end{proof}

\begin{firewall}[Cross-label mixing is not bounded by block norms]
\label{fire:V28-cross-label}
If a separator has blocks
\(D_{\boldsymbol K,\boldsymbol L}\) with
\(\boldsymbol K\ne\boldsymbol L\), the heat multiplier may occur on
only one side of an off-diagonal matrix unit.  Summing
\(\|D_{\boldsymbol K,\boldsymbol L}\|\le1\) over the unheated label
can create a polynomial output count.  The diagonal proof above may
not be applied after deleting those off-diagonal blocks by ordinary
norm.
\end{firewall}

\begin{assessment}[Next upstream question]
\label{ass:V28-next}
The exact source census should now decide whether local total Casimirs
are genuine conserved labels of the physical output separator.  If
they are, V28 closes the entire heat-summed output issue.  If not, the
off-diagonal block must be paired with heat on both sides or controlled
by a weighted Schur test.  No additional fixed-\(K\) recoupling
calculation is needed.
\end{assessment}

\section{V28 result ledger and V29 acquisition order}
\label{sec:V28-ledger}

\begin{theorem}[Fourth-series V28 result ledger]
\label{thm:V28-results}
Relative to V1--V27, V28 proves:
\begin{enumerate}[label=\textup{(V28.\arabic*)},leftmargin=6.3em]
\item the exact scalar marginal of every total-spin-\(K\) projection;
\item a basis-free Schmidt bound for its entire range;
\item the \(d_K^3/d_j\) arbitrary-dualizer estimate;
\item complete heat, covariance, paid-heat, and paid-covariance output
      summation at scale \(d_j^{-1}\);
\item a three-output tuple theorem at scale \(d_j^{-3}\);
\item closure after every local-Casimir preserving separator; and
\item reduction of the growing-output problem to exact cross-label
      mixing.
\end{enumerate}
\end{theorem}

\begin{proof}
These are
\Cref{lem:V28-scalar-marginal,
cor:V28-Schmidt,
thm:V28-fixed-K,
thm:V28-one-summed,
lem:V28-tuple,
thm:V28-three-summed,
cor:V28-LCP-close,
prop:V28-no-individual-output}.
\end{proof}

\begin{programme}[V29: exact local-Casimir commutator census]
\label{prog:V28-V29}
\begin{enumerate}[leftmargin=3em]
\item List every literal map in the first-connectivity source between
      the three cofinal coordinate factors and the final output dual.
\item For each map, compute its commutator with the three local total
      Casimirs.
\item Prove \(LCP\) for canonical row regroupings, complete physical
      pinching, and the identity/multiplicity dual maps wherever the
      commutators vanish.
\item Isolate the first actual source map with a nonzero off-diagonal
      \(\boldsymbol K,\boldsymbol L\) block.
\item If such a map exists, write its weighted block matrix and apply a
      two-sided heat Schur test before taking the output supremum.
\item Use the result to close or sharply restate the complete \(q=3\)
      direct-output residual before the V30 companion audit.
\end{enumerate}
\end{programme}

\begin{firewall}[Claim boundary after fourth-series V28]
\label{fire:V28-claim-boundary}
V28 controls every individual and heat-summed growing output under an
exact local-Casimir preserving separator.  It does not prove that all
physical source separators have this property, control their
off-diagonal local-label blocks, close all \(q=3\) fibres, or prove
global quintic MTA.  The cross-label census, \(q=0,1,2\) sectors,
all-order MTA, WUFC, CPM, cutoff removal, reflection positivity,
Osterwalder--Schrader reconstruction, construction of the continuum
Yang--Mills measure, and a uniform positive continuum spectral gap
remain open.  The full Yang--Mills mass gap has not been proved.
\end{firewall}

\paragraph{Parent-version record.}
V28 was formed from
\texttt{Renewal\_Yang\_Mills\_Mass\_Gap\_Research\_Notes\_V27.tex}.
The V27 parent had \(13\,279\) logical source lines,
\(468\,446\) bytes, and SHA--256
\[
 \texttt{7200d52756fa6c11c4c854d67e4cfa660f449db47241f246f403aa93d7df3a8e}.
\]
The next compulsory companion audit is V30.

% =============================================================================
% END APPENDED FOURTH-SERIES V28 MATERIAL
% =============================================================================

\clearpage
% =============================================================================
% BEGIN APPENDED FOURTH-SERIES V29 MATERIAL
% =============================================================================

\chapter{V29: Exact local-Casimir conservation and unequal-spin
two-row closure}
\label{ch:V29}

\section{Purpose}
\label{sec:V29-purpose}

V28 leaves one apparent output obstruction: an exact source separator
might mix different triples of local total-Casimir labels.  The
coefficient-one source census shows that this does not occur.  Every
local moment partition is a function of block Casimirs, all of which
commute with the full local Casimir.  Cumulant assembly, prefix
products, the payer, and coordinate tensoring preserve that
commutation.  In the final Schur decomposition the local Casimirs act
on multiplicity space, while the Fourier dual contractions act on the
irreducible output factor.  They therefore commute.

After proving this \(LCP\) property, the heat-summed argument is
extended from equal spins \(j,j\) to arbitrary \(SU(2)\) input spins
\(p,q\).  This closes the direct coefficient card for every \(q=3\)
fibre whose three cofinal response coordinates have incidence two and
whose remaining source factors are fixed-dimensional.  All seven
temporal placements are included.

\section{Local block Casimirs commute with the full local Casimir}
\label{sec:V29-local-commutators}

At a physical coordinate \(e\), let \(S_e\) be the active row set.
For \(B\subseteq S_e\), write
\[
 J_{B,e}^m:=\sum_{i\in B}J_{i,e}^m,
 \qquad
 C_{B,e}:=\sum_m(J_{B,e}^m)^2,
\]
and put \(C_e:=C_{S_e,e}\).  Dual orientations only replace some
generators by their contragredient actions and do not change the
commutator argument.

\begin{lemma}[Nested or crossing block Casimirs preserve total local
spin]
\label{lem:V29-block-total}
For every \(B\subseteq S_e\),
\begin{equation}
 [C_{B,e},C_e]=0.
\label{eq:V29-block-total}
\end{equation}
Consequently all block Casimirs belonging to any one moment partition
commute with \(C_e\), even when block systems from two different
moment-partition terms are not mutually commuting.
\end{lemma}

\begin{proof}
The Casimir \(C_{B,e}\) commutes with every component
\(J_{B,e}^m\) of its own diagonal action.  It also commutes with every
\(J_{S_e\setminus B,e}^m\), because those generators act on disjoint
row tensor factors.  Hence it commutes with each component
\[
 J_{S_e,e}^m
 =
 J_{B,e}^m+J_{S_e\setminus B,e}^m.
\]
It therefore commutes with the sum of their squares, namely \(C_e\).
The last statement applies this calculation separately to every
block in every partition; mutual commutation between two different
partition systems is not needed.
\end{proof}

\begin{corollary}[Every complete local moment and cumulant preserves
\(C_e\)]
\label{cor:V29-local-cumulant}
Every complete local moment-partition operator
\(\mathbb M_{\pi,e}\), every complete local cumulant
\[
 \mathbb K_e
 =
 \sum_{\pi}\mu(\pi,\widehat1)\mathbb M_{\pi,e},
\]
and every selected or once-paid version assembled from them commutes
with \(C_e\).
\end{corollary}

\begin{proof}
A moment-partition operator is a product of heat functions
\(e^{-sC_{B,e}}\) over the disjoint blocks \(B\in\pi\).  Each factor
commutes with \(C_e\) by functional calculus and
\Cref{lem:V29-block-total}; so does their product.  Linear combination
with the Möbius coefficients proves the cumulant statement.  A
selected output projection and the local payer are spectral functions
of \(C_e\), hence preserve the same commutator.
\end{proof}

\section{The complete first-connectivity source is \(LCP\)}
\label{sec:V29-global-commutators}

Let \(X\) be the finite physical coordinate set in one resolved word.
The coordinate actions commute, and the full diagonal action is
\[
 J_{\rm tot}^m:=\sum_{e\in X}J_e^m.
\]

\begin{lemma}[Each local Casimir lies in the global commutant]
\label{lem:V29-local-global}
For every \(e\in X\),
\begin{equation}
 [C_e,J_{\rm tot}^m]=0
 \quad\hbox{for all }m.
\label{eq:V29-local-global}
\end{equation}
Therefore, in the complete Schur decomposition
\[
 \mathcal H_{\rm in}
 =
 \bigoplus_U H_U\otimes M_U,
\]
there is a self-adjoint \(c_{e,U}\) on \(M_U\) such that
\begin{equation}
 C_e
 =
 \bigoplus_U I_{H_U}\otimes c_{e,U}.
\label{eq:V29-local-multiplicity}
\end{equation}
\end{lemma}

\begin{proof}
The local Casimir commutes with every component of \(J_e\), and with
every \(J_f\), \(f\ne e\), by disjoint tensor support.  This proves
\eqref{eq:V29-local-global}.  An operator in the commutant of a
completely reducible compact-group action has the block form
\eqref{eq:V29-local-multiplicity} by Schur's lemma.
\end{proof}

\begin{theorem}[Exact local-Casimir preserving source census]
\label{thm:V29-LCP-source}
Every literal map in the exact first-connectivity source
\eqref{eq:V11-exact-source}, including:
\begin{enumerate}[label=\textup{(\roman*)},leftmargin=3.5em]
\item complete prefix connected carriers;
\item the complete joining cumulant;
\item the complete moment suffix;
\item the selected local output and the sole payer;
\item canonical coordinate/row regroupings; and
\item complete physical isotypic pinching,
\end{enumerate}
preserves every local total-Casimir label \(C_e\), after labels are
carried through the canonical regrouping.  The exact Fourier dual
operator
\[
 \bigoplus_U D_U^*\otimes I_{M_U}
\]
also commutes with all \(C_e\).  Thus the fully assembled physical
output separator is \(LCP\) in the sense of
\Cref{def:V28-LCP}; it has no off-diagonal
\(\boldsymbol K,\boldsymbol L\) block between distinct local
total-Casimir tuples.
\end{theorem}

\begin{proof}
Items \textup{(i)--(iii)} are sums and products, over physical
coordinates, of the complete local operators in
\Cref{cor:V29-local-cumulant}; hence they commute with every \(C_e\).
Item \textup{(iv)} is functional calculus of one of those same
Casimirs.  A canonical regrouping is the natural tensor
identification; conjugating by it carries \(C_e\) to the identically
labelled Casimir in the regrouped tensor order and creates no linear
combination of different eigenvalues.  Complete isotypic pinching is
spectral pinching by the global action and commutes with the
commutant elements \eqref{eq:V29-local-multiplicity}.

Finally, on each global output block,
\[
 (D_U^*\otimes I_{M_U})(I_{H_U}\otimes c_{e,U})
 =
 (I_{H_U}\otimes c_{e,U})(D_U^*\otimes I_{M_U}).
\]
All literal maps and the dualizer therefore reduce the joint spectral
resolution of the commuting family \((C_e)_{e\in X}\).  Restriction to
the three cofinal coordinates is exactly the block-diagonal condition
\eqref{eq:V28-LCP}.
\end{proof}

\begin{corollary}[The V28 cross-label residual is absent from the
literal source]
\label{cor:V29-no-cross-label}
Off-diagonal mixing of distinct local total-Casimir tuples can arise
only from an auxiliary fine projection or recoupling not present in
\eqref{eq:V11-exact-source}.  It is not an obstruction for the
complete physical Fourier dual carrier.
\end{corollary}

\begin{proof}
This is the last assertion of
\Cref{thm:V29-LCP-source}, combined with the source-map census
\Cref{prop:V11-source-map-census}.
\end{proof}

\begin{firewall}[Intermediate Casimirs may still fail to commute]
\label{fire:V29-intermediate}
Two block Casimirs coming from crossing, nonnested row partitions at
the same coordinate need not commute.  V29 uses only the full local
Casimir \(C_e\), which commutes with every one of them.  No simultaneous
resolution of all intermediate coupling labels is asserted.
\end{firewall}

\section{Unequal-spin output dualizers}
\label{sec:V29-unequal}

Let \(p,q\in\frac12\mathbb N_0\).  The multiplicity-free decomposition
is
\[
 V_p\otimes V_q
 =
 \bigoplus_{K=|p-q|}^{p+q}V_K.
\]
Write \(T_{p,q;K}:V_K\to V_p\otimes V_q\) for the corresponding
isometry and \(P_{p,q;K}=T_{p,q;K}T_{p,q;K}^*\).

\begin{lemma}[Unequal-spin scalar marginals]
\label{lem:V29-unequal-marginal}
For every allowed \(K\),
\begin{equation}
 \operatorname{Tr}_{V_q}P_{p,q;K}
 =
 {d_K\over d_p}I_{V_p},
 \qquad
 \operatorname{Tr}_{V_p}P_{p,q;K}
 =
 {d_K\over d_q}I_{V_q}.
\label{eq:V29-unequal-marginal}
\end{equation}
Every unit vector in this output isotype has largest Schmidt
coefficient at most
\begin{equation}
 \sqrt{\frac{d_K}{\max\{d_p,d_q\}}}.
\label{eq:V29-unequal-Schmidt}
\end{equation}
\end{lemma}

\begin{proof}
The Schur and trace argument in
\Cref{lem:V28-scalar-marginal} applies separately to the irreducible
input modules \(V_p,V_q\), giving
\eqref{eq:V29-unequal-marginal}.  For a unit vector
\(\psi\) in the isotype,
\(|\psi\rangle\langle\psi|\preccurlyeq P_{p,q;K}\).
Taking either partial trace bounds the largest eigenvalue of either
reduced density matrix by \(d_K/d_p\) and by \(d_K/d_q\).  The
nonzero spectra of those reduced matrices coincide, so the better
bound is \eqref{eq:V29-unequal-Schmidt}.
\end{proof}

\begin{theorem}[Unequal-spin fixed-output dualizer]
\label{thm:V29-unequal-fixed}
For every \(D_K\in\mathcal B(V_K)\),
\begin{equation}
 \left\|
  (T_{p,q;K}D_KT_{p,q;K}^*)^\Gamma
 \right\|
 \le
 {d_K^3\over\max\{d_p,d_q\}}\|D_K\|.
\label{eq:V29-unequal-fixed}
\end{equation}
\end{theorem}

\begin{proof}
Apply \Cref{cor:V27-fixed-output-matrix} with output dimension
\(m=d_K\) and squared Schmidt bound
\(d_K/\max\{d_p,d_q\}\) from
\Cref{lem:V29-unequal-marginal}.
\end{proof}

Put
\[
 q_L=e^{-sL(L+1)},\qquad
 q_pq:=e^{-s[p(p+1)+q(q+1)]}.
\]

\begin{theorem}[Unequal-spin complete heat and covariance sums]
\label{thm:V29-unequal-summed}
For fixed \(s>0\), uniformly in \(p,q\) and in contraction families
\(\|D_K\|\le1\),
\begin{align}
 \left\|
  \sum_{K=|p-q|}^{p+q}
  q_K(T_{p,q;K}D_KT_{p,q;K}^*)^\Gamma
 \right\|
 &\le
 {C_s\over\max\{d_p,d_q\}},
\label{eq:V29-unequal-heat}\\
 \left\|
  \sum_{K=|p-q|}^{p+q}
  (q_K-q_{pq})
  (T_{p,q;K}D_KT_{p,q;K}^*)^\Gamma
 \right\|
 &\le
 {C_s\over\max\{d_p,d_q\}}.
\label{eq:V29-unequal-cov}
\end{align}
The paid heat and paid covariance versions obey the same bound.
\end{theorem}

\begin{proof}
The heat part is bounded by
\[
 {1\over\max\{d_p,d_q\}}
 \sum_{K=0}^{\infty}d_K^3e^{-sK(K+1)}.
\]
For the independently heated endpoint term, assume without loss that
\(p\ge q\).  Then
\[
 \sum_{K=|p-q|}^{p+q}d_K^3
 \le C(1+p+q)^4\le C(1+p)^4,
\]
while
\[
 q_{pq}\le e^{-sp(p+1)}.
\]
Their product is uniformly bounded.  This proves
\eqref{eq:V29-unequal-cov}.  Multiplication by the payer
\(r_K(s)\le C_s(1+K)^2\) changes the polynomial powers to six; the
Gaussian heat term and the endpoint factor still dominate them
uniformly, exactly as in
\Cref{thm:V28-one-summed}.
\end{proof}

\section{The incidence-two \(q=3\) closure}
\label{sec:V29-q3}

\begin{theorem}[Complete \(SU(2)\) incidence-two direct closure]
\label{thm:V29-incidence-two-close}
Fix a one-strong \(q=3\) first-connectivity source and a maximal row
\(a\), \(\Xi_a=M\).  Suppose:
\begin{enumerate}[label=\textup{(\roman*)},leftmargin=3.5em]
\item at each of the three cofinal response coordinates \(e_k\), the
      varying local factor has exactly two nontrivial input rows, of
      spins \(p_k,q_k\), with row \(a\) carrying \(p_k\);
\item
      \(1+p_k(p_k+1)\ge\tau M\) for \(k=1,2,3\);
\item each assembled local factor is a complete heat, covariance,
      paid heat, or paid covariance family from
      \Cref{thm:V29-unequal-summed}, with only one payer overall;
\item all remaining nontrivial source/output factors are
      fixed-dimensional; and
\item the exchange floor
      \eqref{eq:V26-one-strong-floor} holds.
\end{enumerate}
Then the complete physical Fourier coefficient, uniformly over every
dual contraction and in all seven temporal placements, satisfies
\(DXQTRC_5\).
\end{theorem}

\begin{proof}
By \Cref{thm:V29-LCP-source}, the complete source and every physical
dual contraction preserve the three local output labels.  Apply
\Cref{thm:V29-unequal-summed} at each coordinate and the tuple
argument of \Cref{thm:V28-three-summed}.  Since
\[
 d_{p_k}=2p_k+1
 \ge c_\tau M^{1/2},
\]
the resulting full partial-transpose bound is
\[
 \prod_{k=1}^3{C_s\over\max\{d_{p_k},d_{q_k}\}}
 \le C_{s,\tau}M^{-3/2}.
\]
Fixed-dimensional factors alter only the constant.  Normalize by the
five row masses and compare with the exchange floor as in
\Cref{thm:V26-DDSC-close}.  The commutator proof does not depend on
whether a coordinate lies in the prefix, at the joining coordinate,
or in the suffix, so all seven temporal triples are included.
\end{proof}

\begin{corollary}[The entire V23-type family is closed, not just its
spin-one block]
\label{cor:V29-V23-complete}
On the sparse support geometry of V23, sum all root and suffix output
spins and take the complete physical Fourier dual supremum.  The
resulting direct coefficient obeys the exact exchange-tree target
uniformly in \(j\).
\end{corollary}

\begin{proof}
At the three varying coordinates only rows \(a,c\) are nontrivial,
with \(p_k=q_k=j\).  The fixed prefix satisfies the remaining
hypotheses.  Apply
\Cref{thm:V29-incidence-two-close} and the exact exchange weight
\Cref{thm:V23-exact-exchange}.
\end{proof}

\begin{firewall}[Remaining \(q=3\) incidence]
\label{fire:V29-higher-incidence}
V29 does not treat a cofinal coordinate at which three or more
dimension-growing row inputs enter the same complete local moment or
cumulant, nor a dimension-growing prefix factor outside the three
certified coordinates.  The two-input scalar marginal has no automatic
replacement by choosing a pair inside a higher-incidence local
operator.  Such a replacement would split an assembled cumulant.
\end{firewall}

\begin{assessment}[Exact residual after V29]
\label{ass:V29-residual}
Cross-local-Casimir output mixing and arbitrary two-row output height
are closed.  A persistent \(q=3\) obstruction must now contain either
a higher-incidence cofinal local operator, a dimension-growing
uncertified prefix/spectator output, or a local selected/payer form
outside the complete heat/covariance families.  V30 should audit the
companions and then attack the minimal higher-incidence case by a
scalar marginal on one irreducible row versus the complete reducible
complement.
\end{assessment}

\section{V29 result ledger and V30 acquisition order}
\label{sec:V29-ledger}

\begin{theorem}[Fourth-series V29 result ledger]
\label{thm:V29-results}
Relative to V1--V28, V29 proves:
\begin{enumerate}[label=\textup{(V29.\arabic*)},leftmargin=6.3em]
\item exact commutation of every block Casimir with the full local
      Casimir;
\item preservation of each local total-Casimir label by the complete
      first-connectivity source, payer, physical pinching, and Fourier
      dualizer;
\item removal of the V28 cross-label residual from the literal source;
\item unequal-spin scalar marginals and Schmidt bounds;
\item complete unequal-spin heat, covariance, and paid output sums;
\item \(DXQTRC_5\) for every incidence-two \(q=3\) fibre with only
      fixed-dimensional remaining factors, in all temporal classes;
      and
\item full output-summed closure of the V23 sparse family.
\end{enumerate}
\end{theorem}

\begin{proof}
These are
\Cref{lem:V29-block-total,
cor:V29-local-cumulant,
lem:V29-local-global,
thm:V29-LCP-source,
cor:V29-no-cross-label,
lem:V29-unequal-marginal,
thm:V29-unequal-fixed,
thm:V29-unequal-summed,
thm:V29-incidence-two-close,
cor:V29-V23-complete}.
\end{proof}

\begin{programme}[V30: audit and minimal higher-incidence output]
\label{prog:V29-V30}
\begin{enumerate}[leftmargin=3em]
\item Perform the compulsory audit of the newest active SM and NS
      notes.
\item Let one cofinal row module \(V_p\) face the complete reducible
      tensor product \(W\) of all other active rows at a
      higher-incidence coordinate.
\item Compute the partial traces of each total-output isotype,
      retaining its full multiplicity in \(V_p\otimes W\).
\item Bound the arbitrary output dualizer after embedding, using the
      scalar marginal on \(V_p\) and a heat-weighted multiplicity sum.
\item Include the selected cumulant and payer before summing outputs;
      do not reduce the operator to a chosen pair.
\item Propagate any complete one-coordinate half-power through the
      three-slot direct compiler and update the \(q=3\) residual.
\end{enumerate}
\end{programme}

\begin{firewall}[Claim boundary after fourth-series V29]
\label{fire:V29-claim-boundary}
V29 proves the complete direct card for the incidence-two \(q=3\)
class and removes local-Casimir output mixing from the exact source.
It does not control higher-incidence local cumulants or
dimension-growing uncertified factors, close all \(q=3\), or prove
global quintic MTA.  The higher-incidence output problem,
\(q=0,1,2\) sectors, all-order MTA, WUFC, CPM, cutoff removal,
reflection positivity, Osterwalder--Schrader reconstruction,
construction of the continuum Yang--Mills measure, and a uniform
positive continuum spectral gap remain open.  The full Yang--Mills
mass gap has not been proved.
\end{firewall}

\paragraph{Parent-version record.}
V29 was formed from
\texttt{Renewal\_Yang\_Mills\_Mass\_Gap\_Research\_Notes\_V28.tex}.
The V28 parent had \(13\,756\) logical source lines,
\(483\,246\) bytes, and SHA--256
\[
 \texttt{63dbfa9673192971071f7dd5ffc653be77f85c69968acaca96b956280b602ae9}.
\]
The next compulsory companion audit is V30.

% =============================================================================
% END APPENDED FOURTH-SERIES V29 MATERIAL
% =============================================================================

\clearpage
% =============================================================================
% BEGIN APPENDED FOURTH-SERIES V30 MATERIAL
% =============================================================================

\chapter{V30: Companion audit and higher-incidence local output
closure}
\label{ch:V30}

\section{Compulsory companion audit}
\label{sec:V30-audit}

The newest active companion notes at this checkpoint were:
\begin{center}
\begin{tabular}{@{}p{0.60\linewidth}rr@{}}
\toprule
file & logical lines & bytes\\
\midrule
\texttt{Renewal\_SM\_Einstein\_Continuum\_Research\_Notes\_V9.tex}
 & \(3\,257\) & \(111\,955\)\\
\texttt{Renewal\_NS\_Stability\_Research\_Notes\_V31.tex}
 & \(23\,052\) & \(887\,537\)\\
\bottomrule
\end{tabular}
\end{center}
Their SHA--256 digests were
\[
\begin{split}
&\texttt{261459ea200c0aa7e908cc79f0a1e5e7b066f6924c3e76ee3abd7b5196ea337e},\\
&\texttt{f1121b62238ad5491bb7cf03635ea9934942edf573ed8faaa9ee79e1d38a6696}.
\end{split}
\]

SM V9 proves that the gauge-invariant Hessian descends exactly to an
admissible Faddeev--Popov slice.  The determinant cancels the
gauge-orbit Jacobian but does not change the sign of the physical
quadratic form.  A conformal negative mode which is not gauge remains
negative, and a positive auxiliary block cannot repair it by Schur
complement.

The type-correct YM lesson is a quotient firewall.  Physical output
duality and multiplicity bookkeeping cannot repair a failed local
operator estimate merely by being part of a quotient construction.
The higher-incidence cumulant must retain its inverse-dimension
response after the actual dualizer is included.

NS V31 is unchanged.  Its balanced-helicity example again shows that
a cancellation mechanism can have an exact nullspace outside the
signed sector for which it was designed.  Its acquisition order
requires signed triad work to be assembled before positive parts.

The type-correct YM transfer is to keep the full moment--cumulant sum
assembled.  Selecting a favourable pair inside a higher-incidence
coordinate would leave the signed cumulant sector and is forbidden.
The estimate below instead treats one irreducible maximal row against
the complete reducible complement and retains every moment-partition
term.

\begin{audit}[V30 direction]
\label{audit:V30-direction}
Prove an output-dual partial-transpose estimate for the complete
higher-incidence local cumulant.  Do not infer it from gauge/output
quotienting, a chosen input pair, or a determinant cancellation.
\end{audit}

\section{An ancilla-stable output embedding lemma}
\label{sec:V30-ancilla}

Let \(T:E\to\mathcal H_A\otimes\mathcal H_B\) be an isometry.  Fix an
orthonormal basis \((u_\alpha)_{\alpha=1}^m\) of \(E\) and put
\(\psi_\alpha=Tu_\alpha\).

\begin{lemma}[Ancilla-stable dualizer absorption]
\label{lem:V30-ancilla}
Assume every \(\psi_\alpha\) has largest Schmidt coefficient at most
\(\varepsilon\).  Let \(\mathcal F\) be any finite-dimensional
ancillary space placed entirely on the \(B\)-side of the cut.  Then
for every \(D\in\mathcal B(E\otimes\mathcal F)\),
\begin{equation}
 \boxed{
 \left\|
  \left[
  (T\otimes I_{\mathcal F})D
  (T^*\otimes I_{\mathcal F})
  \right]^{\Gamma_{B\otimes\mathcal F}}
 \right\|
 \le m^2\varepsilon^2\|D\|.}
\label{eq:V30-ancilla}
\end{equation}
The constant is independent of \(\dim\mathcal F\).
\end{lemma}

\begin{proof}
Write \(D=[D_{\alpha\beta}]_{\alpha,\beta=1}^m\) as an
operator-valued matrix on \(\mathcal F\).  Then
\[
 (T\otimes I)D(T^*\otimes I)
 =
 \sum_{\alpha,\beta}
 |\psi_\alpha\rangle\langle\psi_\beta|
 \otimes D_{\alpha\beta}.
\]
Each block satisfies
\(\|D_{\alpha\beta}\|\le\|D\|\).  After partial transpose, the norm of
one summand is
\[
 \left\|
 (|\psi_\alpha\rangle\langle\psi_\beta|)^{\Gamma_B}
 \right\|
 \|D_{\alpha\beta}^{\mathsf T}\|
 \le\varepsilon^2\|D\|
\]
by \Cref{lem:V27-cross-rank-one}.  Sum the \(m^2\) blocks.  No basis
sum over \(\mathcal F\) was taken.
\end{proof}

\begin{corollary}[Unequal-spin output with arbitrary complementary
ancilla]
\label{cor:V30-unequal-ancilla}
For the \(SU(2)\) isometry
\(T_{p,q;K}:V_K\to V_p\otimes V_q\), any ancillary
\(\mathcal F\) on the \(q\)-side, and every
\(D\in\mathcal B(V_K\otimes\mathcal F)\),
\begin{equation}
 \left\|
  \left[
  (T_{p,q;K}\otimes I)D
  (T_{p,q;K}^*\otimes I)
  \right]^\Gamma
 \right\|
 \le
 {d_K^3\over\max\{d_p,d_q\}}\|D\|.
\label{eq:V30-unequal-ancilla}
\end{equation}
\end{corollary}

\begin{proof}
Use the squared Schmidt bound
\eqref{eq:V29-unequal-Schmidt} and \(m=d_K\) in
\Cref{lem:V30-ancilla}.
\end{proof}

\section{One irreducible row versus a reducible complement}
\label{sec:V30-reducible}

Let \(V_p\) be the representation on the maximal row at one physical
coordinate and let \(W\) be the tensor product of all other active row
modules in the same moment block.  Decompose
\[
 W
 =
 \bigoplus_q\mathcal N_q\otimes V_q.
\]
The multiplicity spaces \(\mathcal N_q\) are placed on the
complementary side of the row cut.

\begin{theorem}[Reducible-complement heat with every output dualizer]
\label{thm:V30-reducible-dual-heat}
Let the complete output heat on \(V_p\otimes W\) be followed and
preceded by arbitrary contractions which preserve the complement
type \(q\) and the total output \(K\), but may act arbitrarily on all
remaining complementary ancillary factors.  Then the fully embedded
dual operator \(\mathcal H_{p,W}^{\rm dual}(s)\) satisfies
\begin{equation}
 \boxed{
 \|(\mathcal H_{p,W}^{\rm dual}(s))^\Gamma\|
 \le {C_s\over d_p}.}
\label{eq:V30-reducible-dual-heat}
\end{equation}
The constant is independent of \(W\), its irreducible support, and all
multiplicities \(\dim\mathcal N_q\).
\end{theorem}

\begin{proof}
The complement decomposition is an orthogonal direct sum preserved by
the complete heat and by the local-Casimir census
\Cref{thm:V29-LCP-source}.  Partial transpose preserves this direct
sum, so its operator norm is the maximum over \(q\), not a sum.

Fix \(q\).  Decompose \(V_p\otimes V_q\) into the
multiplicity-free outputs \(V_K\).  All factors
\(\mathcal N_q\) and all other complement-only outputs are ancillary
spaces in \Cref{cor:V30-unequal-ancilla}.  After multiplying by the
heat scalar \(q_K(s)\) and summing \(K\), that corollary gives
\[
 {1\over\max\{d_p,d_q\}}
 \sum_{K=|p-q|}^{p+q}d_K^3e^{-sK(K+1)}
 \le {C_s\over d_p}.
\]
This is uniform in \(q\).  Taking the direct-sum maximum proves
\eqref{eq:V30-reducible-dual-heat}.
\end{proof}

\begin{remark}[Strengthening of the archived V87 card]
\label{rem:V30-V87}
Archive f3 V87 proves the same \(d_p^{-1}\) scale for the complete
positive heat operator.  Theorem \ref{thm:V30-reducible-dual-heat}
retains an arbitrary output contraction and arbitrary complementary
ancilla inside the embedding.  It is therefore the form needed by the
direct Fourier coefficient, not merely by a positive heat bank.
\end{remark}

\section{Complete higher-incidence moment partitions}
\label{sec:V30-partitions}

Fix one local moment partition \(\pi\) of the active rows at coordinate
\(e\).  Let \(B_a\in\pi\) be the unique block containing the maximal
row \(a\).  Its module is
\[
 V_p\otimes W_{B_a},
\]
while every other block lies wholly on the complementary side.

\begin{theorem}[Moment-partition dual response at arbitrary incidence]
\label{thm:V30-partition-response}
Let \(\mathbb M_{\pi,e}^{\rm dual}\) be the complete local
moment-partition operator with an arbitrary admissible physical
output dualizer retained.  Then
\begin{equation}
 \|(\mathbb M_{\pi,e}^{\rm dual})^{\Gamma_{a^c}}\|
 \le {C_{s,m}\over d_p},
\label{eq:V30-partition-response}
\end{equation}
where \(m\le5\) is the local arity.  The bound is uniform in every
other input representation and multiplicity.
\end{theorem}

\begin{proof}
All partition blocks other than \(B_a\) act on the complementary side
and are heat contractions.  Their output spaces, multiplicities, and
all later complement-only maps form the ancillary space in
\Cref{thm:V30-reducible-dual-heat}.  By
\Cref{thm:V29-LCP-source}, the dualizer preserves the output label of
the \(B_a\)-block, so the theorem applies without deleting any
off-diagonal source map.  Its \(d_p^{-1}\) estimate is unaffected by
the ordinary norms of the other heat contractions.
\end{proof}

\begin{corollary}[Complete local cumulant dual response]
\label{cor:V30-cumulant-response}
Let
\[
 \mathbb K_e
 =
 \sum_{\pi\in\Pi_m}
 \mu(\pi,\widehat1)\mathbb M_{\pi,e}
\]
be any complete local cumulant of arity \(2\le m\le5\), with every
physical output block and dualizer retained.  Then
\begin{equation}
 \boxed{
 \|(\mathbb K_e^{\rm dual})^{\Gamma_{a^c}}\|
 \le {C_{s,m}\over d_p}.}
\label{eq:V30-cumulant-response}
\end{equation}
Selected complete output projections preserve the bound.
\end{corollary}

\begin{proof}
Apply \Cref{thm:V30-partition-response} to every partition term and
sum the fixed Möbius coefficients:
\[
 \sum_{\pi\in\Pi_m}|\mu(\pi,\widehat1)|<\infty.
\]
This constant depends only on \(m\le5\).  A selected complete output
projection is part of the contraction retained in the ancilla-stable
estimate and has norm at most one.
\end{proof}

\section{The payer at arbitrary incidence}
\label{sec:V30-payer}

\begin{lemma}[One linear payer absorption across a moment partition]
\label{lem:V30-payer-absorption}
Let \(h_\pi=\sum_{B\in\pi}C_{B,e}\) be the sum of the partition-block
Casimirs and let \(C_e\) be the full local output Casimir.  There is a
fixed arity constant \(c_m\) such that
\[
 0\le C_e\le c_mh_\pi.
\]
For fixed \(s>0\),
\begin{equation}
 {C_e\over1-e^{-sC_e}}\,
 \boldsymbol1_{\{C_e>0\}}\,
 e^{-sh_\pi}
 \le C_{s,m}e^{-(s/2)h_\pi}
\label{eq:V30-linear-payer-absorption}
\end{equation}
on every common total-output block.
\end{lemma}

\begin{proof}
The generator of the full action is the sum of the block generators.
Cauchy--Schwarz gives
\[
 \left|\sum_{B\in\pi}J_B\right|^2
 \le |\pi|\sum_{B\in\pi}|J_B|^2,
\]
so one may take \(c_m=m\).  On \(C_e>0\), the compact \(SU(2)\)
Casimir has a fixed positive spectral floor, hence
\((1-e^{-sC_e})^{-1}\le C_s\).  Therefore the left side of
\eqref{eq:V30-linear-payer-absorption} is bounded by
\[
 C_s c_m h_\pi e^{-sh_\pi}
 \le C_{s,m}e^{-(s/2)h_\pi}.
\]
It vanishes on \(C_e=0\).
\end{proof}

\begin{theorem}[Paid complete local cumulant response]
\label{thm:V30-paid-cumulant}
Under the hypotheses of
\Cref{cor:V30-cumulant-response}, the once-paid complete local
cumulant satisfies
\begin{equation}
 \boxed{
 \|((\mathbb K_e^{\rm pay})^{\rm dual})^{\Gamma_{a^c}}\|
 \le {C_{s,m}\over d_p}.}
\label{eq:V30-paid-cumulant}
\end{equation}
\end{theorem}

\begin{proof}
Apply \Cref{lem:V30-payer-absorption} separately to each assembled
moment-partition term before summing the Möbius coefficients.  It
replaces the paid heat product at time \(s\) by an unpaid heat product
at time \(s/2\), with a uniform constant.  Then use
\Cref{thm:V30-partition-response} at time \(s/2\) and sum the fixed
partition coefficients.  No input pair or moment term is retained
after the sum; the triangle inequality is only over the fixed complete
Möbius polynomial.
\end{proof}

\begin{firewall}[One payer, one linear absorption]
\label{fire:V30-one-payer}
Theorem \ref{thm:V30-paid-cumulant} absorbs the unique payer once in
the complete local output.  It does not assign a payer separately to
each partition block or multiply a paid estimate by an unpaid one.
\end{firewall}

\section{Three higher-incidence coordinates}
\label{sec:V30-three}

\begin{theorem}[Three complete local cumulants with a joint dualizer]
\label{thm:V30-three-cumulants}
Let \(e_1,e_2,e_3\) be distinct coordinates.  At \(e_k\), let the
maximal row \(a\) carry spin \(p_k\), and let the assembled local
factor be a complete heat moment or a complete cumulant of arity at
most five, paid at at most one coordinate.  Retain the full reducible
complements, all local output sums, and an arbitrary joint physical
dualizer preserving the three local total-Casimir labels.  Then
\begin{equation}
 \boxed{
 \|\mathcal Z^\Gamma\|
 \le
 C_{s}\prod_{k=1}^3d_{p_k}^{-1}.}
\label{eq:V30-three-cumulants}
\end{equation}
\end{theorem}

\begin{proof}
Resolve each complement into irreducibles and each local total output.
The joint label preservation is
\Cref{thm:V29-LCP-source}.  On one fixed label tuple, expand the joint
dualizer as an operator-valued matrix in the tensor product of the
three local output bases.  Iterating
\Cref{lem:V30-ancilla} gives the product of the three squared Schmidt
bounds; all complement multiplicity spaces remain operator-valued
ancillas and create no basis count.

For each moment-partition term, the polynomial output-dimension costs
factor by coordinate and are summed against its heat multipliers as in
\Cref{thm:V30-reducible-dual-heat}.  The three resulting convergent
series multiply.  Sum last over the fixed moment partitions of the
three local cumulants.  If one coordinate is paid, first use
\Cref{lem:V30-payer-absorption} there.  This proves
\eqref{eq:V30-three-cumulants}.
\end{proof}

\begin{corollary}[Higher-incidence local factors are certified for
\(q=3\)]
\label{cor:V30-higher-incidence-certified}
If
\[
 1+p_k(p_k+1)\ge\tau M
 \qquad(k=1,2,3),
\]
then the complete three-coordinate higher-incidence packet in
\Cref{thm:V30-three-cumulants} has partial-transpose norm at most
\[
 C_{s,\tau}M^{-3/2}.
\]
Thus higher local incidence and the selected/payer roles are not,
by themselves, residual obstructions to \(DXQTRC_5\).
\end{corollary}

\begin{proof}
The cofinality inequalities and
\(d_{p_k}=2p_k+1\) give
\(d_{p_k}\ge c_\tau M^{1/2}\).  Insert these three estimates into
\eqref{eq:V30-three-cumulants}.
\end{proof}

\begin{firewall}[Dimension-growing row-\(a\) factors outside
\(E_3\)]
\label{fire:V30-outside-row}
The ancilla spaces in V30 lie on the complement of the maximal-row
cut.  A dimension-growing physical coordinate outside \(E_3\) at
which row \(a\) is also active is not such an ancilla.  If its exact
map is placed around the three-coordinate packet using only an
ordinary norm, the V26 entangler obstruction applies.  V30 therefore
closes higher incidence \emph{at the three certified coordinates}; it
does not yet control an unbounded residual row-\(a\) prefix or suffix.
\end{firewall}

\begin{assessment}[Exact residual after V30]
\label{ass:V30-residual}
The local-operator part of \(q=3\) is now complete for \(SU(2)\):
arbitrary incidence, reducible complements, all output heights and
multiplicities, selected cumulants, the sole payer, and every output
dual contraction retain the product \(M^{-3/2}\) at the three
cofinal coordinates.  The remaining question is global in physical
coordinates: can the residual row-\(a\) source outside \(E_3\) be
shown fixed-dimensional, diffuse and damped, or jointly
partial-transpose bounded without destroying that packet?
\end{assessment}

\section{V30 result ledger and V31 acquisition order}
\label{sec:V30-ledger}

\begin{theorem}[Fourth-series V30 result ledger]
\label{thm:V30-results}
Relative to V1--V29, V30 proves:
\begin{enumerate}[label=\textup{(V30.\arabic*)},leftmargin=6.3em]
\item the compulsory SM V9 and NS V31 companion audit;
\item an ancilla-stable arbitrary-output embedding theorem with no
      ancillary dimension loss;
\item the reducible-complement heat-dual card \(C_sd_p^{-1}\);
\item the complete moment-partition and cumulant response at arbitrary
      local incidence;
\item linear absorption of the sole payer before output summation;
\item the paid complete-cumulant response;
\item the three-coordinate joint-dual bound
      \(C\prod_kd_{p_k}^{-1}\); and
\item removal of all higher-incidence local operators from the
      \(q=3\) residual.
\end{enumerate}
\end{theorem}

\begin{proof}
These are
\Cref{audit:V30-direction,
lem:V30-ancilla,
thm:V30-reducible-dual-heat,
thm:V30-partition-response,
cor:V30-cumulant-response,
lem:V30-payer-absorption,
thm:V30-paid-cumulant,
thm:V30-three-cumulants,
cor:V30-higher-incidence-certified}.
\end{proof}

\begin{programme}[V31: residual row-mass exterior to \(E_3\)]
\label{prog:V30-V31}
\begin{enumerate}[leftmargin=3em]
\item Return to the exact V5 owner-mass partition
      \(E\sqcup C\sqcup N\sqcup R\) under three-slot saturation.
\item Keep the complete \(E_3\) higher-incidence packet from V30
      assembled.
\item Split the residual row-\(a\) mass outside \(E_3\) into fixed
      atoms, diffuse banks, and finitely many marked noncentral owners.
\item Absorb complement-only factors as V30 ancillas and apply the
      archived diffuse-bank theorem to source-occurring central
      exterior banks.
\item For a cofinal exterior owner, use its already classified complete
      local form or descend to its first connected ancestor.
\item Prove a whole-word partial-transpose bound or state the exact
      surviving noncentral exterior packet; do not discard it by
      ordinary norm.
\end{enumerate}
\end{programme}

\begin{firewall}[Claim boundary after fourth-series V30]
\label{fire:V30-claim-boundary}
V30 closes the complete local higher-incidence and output-duality
problem at the three \(q=3\) coordinates.  It does not yet control
dimension-growing maximal-row factors outside those coordinates,
close all \(q=3\), or prove global quintic MTA.  The exterior
row-mass problem, \(q=0,1,2\) sectors, all-order MTA, WUFC, CPM,
cutoff removal, reflection positivity, Osterwalder--Schrader
reconstruction, construction of the continuum Yang--Mills measure,
and a uniform positive continuum spectral gap remain open.  The full
Yang--Mills mass gap has not been proved.
\end{firewall}

\paragraph{Parent-version record.}
V30 was formed from
\texttt{Renewal\_Yang\_Mills\_Mass\_Gap\_Research\_Notes\_V29.tex}.
The V29 parent had \(14\,254\) logical source lines,
\(500\,772\) bytes, and SHA--256
\[
 \texttt{532dab0101ede46235c774190ad4d397a512cf80e613a975da7ec61cf7ab5b8c}.
\]
The next compulsory companion audit is V35.

% =============================================================================
% END APPENDED FOURTH-SERIES V30 MATERIAL
% =============================================================================

% The V30 document terminator is intentionally disabled here:
% \end{document}

% =============================================================================
% BEGIN APPENDED FOURTH-SERIES V31 MATERIAL
% =============================================================================

\clearpage
\chapter{V31: Whole-word output absorption and large-heat closure of
the \texorpdfstring{\(q=3\)}{q=3} exterior}
\label{ch:V31}

\section{Purpose and upstream correction}
\label{sec:V31-purpose}

V30 closed the complete higher-incidence operator at each of the three
cofinal response coordinates but left dimension-growing row-\(a\)
factors outside that set.  The V30 acquisition order proposed returning
to the V5 partition \(E\sqcup C\sqcup N\sqcup R\).  That partition is
useful for positive central insertions, but it is not the sharp
description of the literal Fourier operator.

The exact source-map census of V11 and the local-Casimir census of V29
give a more upstream description.  After recursive first-connectivity,
there are only finitely many joining cumulants.  Expanding those
cumulants at the \emph{row-partition} level leaves coordinatewise
moment-partition heat maps, followed by one joint contraction which
preserves all local output labels.  Tree, lift, strong-edge, and owner
names are estimation data, not extra source maps.

This chapter estimates that joint contraction all at once.  The useful
local quantity is not the ordinary norm of an exterior factor and not
the decomposable norm of a separated dualizer.  It is the weighted sum
of the largest Schmidt coefficients of all heat-output basis vectors.
That sum factorizes over physical coordinates before the joint
dualizer is estimated.

For \(SU(2)\), the weighted sum gives two simultaneous facts.  At a
cofinal coordinate it is \(O(d_p^{-1/2})\), hence its square is the
required \(O(d_p^{-1})\).  At every nontrivial row-\(a\) coordinate it
is a strict contraction once the heat time is larger than a universal
threshold.  Consequently an arbitrarily long exterior row-\(a\)
support cannot amplify the three cofinal gains in that heat-time
regime.

\section{The weighted Schmidt sum for one output block}
\label{sec:V31-local}

Write
\[
 d_j:=2j+1,\qquad c_j:=j(j+1),
\]
and let
\[
 T_{p,q;K}:V_K\longrightarrow V_p\otimes V_q
\]
be the multiplicity-one Clebsch--Gordan isometry.  The allowed spins
\(K\) run from \(|p-q|\) to \(p+q\) in unit steps.  Define
\begin{equation}
 \ell_s(p,q)
 :=
 \sum_{K=|p-q|}^{p+q}
 d_K e^{-sc_K/2}
 \left({d_K\over\max\{d_p,d_q\}}\right)^{1/2}.
\label{eq:V31-ell}
\end{equation}
The factor \(d_K\) counts an output basis and the square root is the
Schmidt bound from V29.

\begin{lemma}[Operator-valued output matrix estimate]
\label{lem:V31-operator-matrix}
Let \(\mathcal F\) be any finite-dimensional Hilbert space.  For every
family
\[
 D_K\in\mathcal B(V_K\otimes\mathcal F),
 \qquad
 \sup_K\|D_K\|\le1,
\]
one has
\begin{equation}
 \left\|
  \sum_K e^{-sc_K}
  \left[
   (T_{p,q;K}\otimes I_{\mathcal F})
   D_K
   (T_{p,q;K}^*\otimes I_{\mathcal F})
  \right]^{\Gamma_{V_q}}
 \right\|
 \le \ell_s(p,q)^2.
\label{eq:V31-operator-matrix}
\end{equation}
No factor depending on \(\dim\mathcal F\) occurs.
\end{lemma}

\begin{proof}
Choose an orthonormal basis \((u_{K,m})_{m=1}^{d_K}\) of \(V_K\) and
put
\[
 \psi_{K,m}:=T_{p,q;K}u_{K,m}.
\]
By the scalar-marginal theorem of V29, every \(\psi_{K,m}\) has largest
Schmidt coefficient at most
\[
 \varepsilon_K
 :=
 \left({d_K\over\max\{d_p,d_q\}}\right)^{1/2}.
\]
Write \(D_K=[D_{K,mn}]\) as an operator-valued matrix on
\(\mathcal F\).  Each entry satisfies
\(\|D_{K,mn}\|\le1\).  The exact cross-rank-one identity of V27 gives
\[
 \left\|
  (|\psi_{K,m}\rangle\langle\psi_{K,n}|)^{\Gamma_{V_q}}
 \right\|
 \le\varepsilon_K^2.
\]
Triangle inequality therefore gives the sharper blockwise sum
\[
 \sum_K d_K^2e^{-sc_K}\varepsilon_K^2.
\]
This is at most
\[
 \left(\sum_Kd_Ke^{-sc_K/2}\varepsilon_K\right)^2
 =\ell_s(p,q)^2,
\]
because all summands are nonnegative.  The expansion uses only the
\(V_K\) basis; the entries remain operators on \(\mathcal F\), so its
dimension is never counted.
\end{proof}

\begin{remark}[Why the larger square is retained]
\label{rem:V31-larger-square}
For one coordinate the blockwise sum in the proof is better than
\(\ell_s(p,q)^2\).  The latter is the stable quantity for a joint
dualizer acting on many coordinates: after expanding one joint output
basis, the sum of weighted Schmidt coefficients factorizes, and its
square includes cross-output pairs.  Those pairs cannot be discarded
for an arbitrary contraction.
\end{remark}

\begin{lemma}[Cofinal gain and uniform strict contraction]
\label{lem:V31-two-bounds}
For every \(s>0\) there is \(A_s<\infty\) such that
\begin{equation}
 \ell_s(p,q)^2\le {A_s\over d_p}
\label{eq:V31-cofinal-local}
\end{equation}
for all \(p,q\).  There also exists a universal
\(s_\sharp<\infty\) such that, for every \(s\ge s_\sharp\), there is
\(\rho_s<1\) satisfying
\begin{equation}
 p>0\quad\Longrightarrow\quad
 \ell_s(p,q)^2\le\rho_s
\label{eq:V31-strict-local}
\end{equation}
uniformly in \(q\).
\end{lemma}

\begin{proof}
Since \(\max\{d_p,d_q\}\ge d_p\),
\[
 \ell_s(p,q)
 \le {1\over\sqrt{d_p}}
 \sum_{K\in\widehat{SU(2)}}d_K^{3/2}e^{-sc_K/2}.
\]
The full representation sum is finite for every \(s>0\).  Squaring
proves \eqref{eq:V31-cofinal-local}.

For the strict estimate, isolate \(K=0\).  It occurs only when \(p=q\).
If \(p>0\), then \(d_p\ge2\), and its contribution to
\(\ell_s(p,p)\) is exactly
\[
 {1\over\sqrt{d_p}}\le {1\over\sqrt2}.
\]
For \(K>0\), the triangle inequality \(K\le p+q\) implies
\[
 d_K\le d_p+d_q-1<2\max\{d_p,d_q\}.
\]
Hence
\[
 d_K
 \left({d_K\over\max\{d_p,d_q\}}\right)^{1/2}
 \le\sqrt2\,d_K.
\]
Put
\[
 \delta_s
 :=
 \sqrt2\sum_{\substack{K\in\widehat{SU(2)}\\K>0}}
 d_Ke^{-sc_K/2}.
\]
Then
\[
 \ell_s(p,q)\le {1\over\sqrt2}+\delta_s.
\]
The series is dominated at every positive time and decreases to zero
as \(s\to\infty\).  Choose \(s_\sharp\) so that
\(\delta_{s_\sharp}<1-2^{-1/2}\), and set
\[
 \rho_s:=\left(2^{-1/2}+\delta_s\right)^2<1.
\]
This proves \eqref{eq:V31-strict-local}.
\end{proof}

\begin{corollary}[Reducible complement and multiplicities]
\label{cor:V31-reducible}
Let
\[
 W=\bigoplus_q\mathcal N_q\otimes V_q
\]
be an arbitrary finite-dimensional \(SU(2)\)-module on the complement
of row \(a\).  Suppose the output contraction preserves \(q\) and the
total output \(K\), but acts arbitrarily on \(\mathcal N_q\) and on all
other complement-side output factors.  The one-coordinate heat-output
map obeys \eqref{eq:V31-cofinal-local}; if \(p>0\) and
\(s\ge s_\sharp\), it also obeys
\eqref{eq:V31-strict-local}.  Both constants are independent of \(W\)
and of every multiplicity.
\end{corollary}

\begin{proof}
Fix \(q\).  Include \(\mathcal N_q\) and all other complement-side
outputs in the operator space \(\mathcal F\) of
\Cref{lem:V31-operator-matrix}.  The output-label census of V29 makes
different \(q\)-blocks orthogonal and invariant.  Thus the norm of
their direct sum is the maximum, not a sum.  Apply
\Cref{lem:V31-two-bounds} and take that maximum.
\end{proof}

\section{All-at-once absorption of a joint dualizer}
\label{sec:V31-joint}

Let \(A\) be a finite set of physical coordinates at which row \(a\)
is active.  On a fixed moment-partition word, decompose the complement
of \(V_{p_e}\) in the partition block containing \(a\) as
\[
 W_e=\bigoplus_{q_e}\mathcal N_{e,q_e}\otimes V_{q_e}.
\]
Fix a tuple \(\boldsymbol q=(q_e)_{e\in A}\).  Every output basis vector
is indexed by
\[
 \boldsymbol\alpha=((K_e,m_e))_{e\in A}.
\]

\begin{theorem}[Joint heat-output dualizer estimate]
\label{thm:V31-joint}
Let \(D\) be any contraction which preserves the tuple of complement
types \(\boldsymbol q\) and the tuple of local total outputs
\(\boldsymbol K\), and which may otherwise couple all physical output
spaces and all complement multiplicities.  Let \(\mathbb H_A(D)\) be
the complete tensor product of the local moment-partition heat maps,
with \(D\) retained between their output isometries.  Then
\begin{equation}
 \boxed{
 \|\mathbb H_A(D)^{\Gamma_{a^c}}\|
 \le
 \prod_{e\in A}\ell_s(p_e,q_e)^2.}
\label{eq:V31-joint}
\end{equation}
The same estimate holds after adjoining arbitrary heat contractions
on partition blocks not containing row \(a\), and after taking the
adjoint orientation.
\end{theorem}

\begin{proof}
On the fixed \(\boldsymbol q\)-block, expand \(D\) as an
operator-valued matrix in the joint output basis
\((u_{\boldsymbol\alpha})\).  Its entries act on the tensor product of
all complement multiplicities and all output spaces belonging to
partition blocks not containing \(a\).  This entire residual space lies
on the \(a^c\)-side of the cut.  Every matrix entry has norm at most
one, and full transpose on that residual space preserves its norm.

The embedded vector for \(\boldsymbol\alpha\) is
\[
 \Psi_{\boldsymbol\alpha}
 =
 \bigotimes_{e\in A}
 T_{p_e,q_e;K_e}u_{K_e,m_e}.
\]
Its largest Schmidt coefficient across \(a|a^c\) is at most
\[
 \prod_{e\in A}
 \left({d_{K_e}\over
 \max\{d_{p_e},d_{q_e}\}}\right)^{1/2}.
\]
Insert the square-root heat weight
\(\prod_e e^{-sc_{K_e}/2}\) on both sides of \(D\).  The
cross-rank-one partial-transpose identity of V27 and triangle
inequality bound the resulting matrix by the square of
\[
 \sum_{\boldsymbol\alpha}
 \prod_{e\in A}
 e^{-sc_{K_e}/2}
 \left({d_{K_e}\over
 \max\{d_{p_e},d_{q_e}\}}\right)^{1/2}.
\]
The joint basis sum factorizes exactly over \(e\), giving the right
side of \eqref{eq:V31-joint}.  Orthogonality in
\(\boldsymbol q\) again replaces the sum of complement types by a
maximum.

Blocks not containing \(a\) are entirely on the transposed side.
Their output isometries become conjugate isometries after full
transpose, while their heat multipliers have norm at most one.  They
therefore do not change the bound.  Replacing \(D\) by \(D^*\) proves
the other orientation.
\end{proof}

\begin{corollary}[Three gains survive arbitrary exterior support]
\label{cor:V31-three-survive}
Assume \(s\ge s_\sharp\), and let
\[
 E_3=\{e_1,e_2,e_3\}\subseteq A.
\]
Then
\begin{equation}
 \|\mathbb H_A(D)^{\Gamma_{a^c}}\|
 \le
 A_s^3
 \prod_{k=1}^3d_{p_{e_k}}^{-1}.
\label{eq:V31-three-survive}
\end{equation}
The constant is independent of \(|A|\), the spins and multiplicities
outside \(E_3\), and the joint dualizer.
\end{corollary}

\begin{proof}
Use \eqref{eq:V31-cofinal-local} at the three displayed coordinates
and \eqref{eq:V31-strict-local} at every coordinate in
\(A\setminus E_3\).  Then
\[
 \prod_{e\in A\setminus E_3}\ell_s(p_e,q_e)^2
 \le \rho_s^{|A\setminus E_3|}\le1.
\]
Insert these estimates into \eqref{eq:V31-joint}.
\end{proof}

\section{Recursive first-connectivity has only finitely many
non-moment letters}
\label{sec:V31-recursion}

\begin{lemma}[At most four recursive joining nodes]
\label{lem:V31-four-joins}
In a fully recursively resolved connected five-row prefix, the
laminar row-block tree has at most four nonsingleton nodes.  Hence at
most four physical-coordinate occurrences carry joining cumulants;
all remaining occurrences in the fixed recursive fibre are complete
moment letters.
\end{lemma}

\begin{proof}
The leaves are nonempty disjoint row blocks and there are at most five
of them.  Every nonsingleton node has at least two children.  For a
finite rooted tree with \(L\) leaves and \(I\) internal nodes, counting
edges gives
\[
 L+I-1\ge2I,
\]
so \(I\le L-1\le4\).  The coefficient-one first-connectivity identity
assigns exactly one joining coordinate to each nonsingleton node.  Its
prefix children are resolved recursively and its suffix is a complete
moment.  Thus, after all nodes have been resolved, every coordinate not
used by one of these joining occurrences belongs to a moment letter.
\end{proof}

\begin{lemma}[Fixed row-partition expansion]
\label{lem:V31-fixed-expansion}
Fix every recursive first-connectivity coordinate and predecessor
partition.  Expanding the joining cumulant at each nonsingleton node
produces a signed sum
\begin{equation}
 \mathbb W^{\rm raw}
 =
 \sum_{\nu\in\mathcal R_5}c_\nu\mathbb H_\nu
\label{eq:V31-fixed-expansion}
\end{equation}
such that:
\begin{enumerate}[label=\textup{(\roman*)},leftmargin=4em]
\item every \(\mathbb H_\nu\) is a coordinatewise
      moment-partition heat word;
\item
\[
 |\mathcal R_5|+\sum_{\nu\in\mathcal R_5}|c_\nu|
 \le C_5,
\]
where \(C_5\) is independent of physical support and representation
labels; and
\item the unique payer, if present, occurs once in every term and can
      be absorbed at its physical coordinate by the V30 linear-payer
      lemma.
\end{enumerate}
\end{lemma}

\begin{proof}
At a node with \(m\le5\) children, the complete joining cumulant is
the moment--cumulant sum over \(\Pi_m\).  Its number of terms and the
sum of the absolute M\"obius coefficients depend only on \(m\).
By \Cref{lem:V31-four-joins}, at most four such expansions are made.
Their product is therefore bounded by a constant depending only on
row order five.  Every summand is a moment partition at each joining
coordinate, while the exact first-connectivity formula already gives
moments at all other coordinates.  This proves
\textup{(i)}--\textup{(ii)}.

The exact source has one payer allocation.  Expanding a cumulant does
not duplicate that multiplier.  V30 proves, term by term, that the
linear payer times the local heat product at time \(s\) is bounded by
a constant times the corresponding heat product at time \(s/2\).
This is \textup{(iii)}.
\end{proof}

\begin{firewall}[Finite expansion is not coordinatewise absolution]
\label{fire:V31-no-coordinate-expansion}
\Cref{lem:V31-fixed-expansion} expands at most four row-partition
cumulants after their recursive first-connectivity coordinates have
been fixed.  It does not expand independently at every physical
coordinate and does not turn a coordinate-summable connected carrier
into exponentially many absolute words.  The coefficient cost is
therefore support independent.
\end{firewall}

\section{Large-heat closure of the complete \texorpdfstring{\(q=3\)}
{q=3} exterior}
\label{sec:V31-q3}

\begin{theorem}[Whole-word three-slot partial-transpose bound]
\label{thm:V31-whole-word}
Let \(G=SU(2)\), \(s\ge s_\sharp\), and fix a fully recursively
resolved quintic first-connectivity fibre.  Let \(a\) be a maximal row
and let \(E_3=\{e_1,e_2,e_3\}\) be three distinct response coordinates.
Suppose row \(a\) carries spin \(p_k\) at \(e_k\) and
\begin{equation}
 1+p_k(p_k+1)\ge\tau M,
 \qquad
 M=\Xi_a,
 \qquad k=1,2,3.
\label{eq:V31-three-cofinal}
\end{equation}
Retain all recursive prefix carriers, the complete joining carrier,
the complete suffix, every local output and multiplicity, every
canonical equivariant separator, an arbitrary admissible physical
Fourier dualizer, and the unique payer.  Then the unnormalized raw
operator satisfies both orientations of
\begin{equation}
 \boxed{
 \|(\mathbb W^{\rm raw})^{\Gamma_{a^c}}\|
 \le C_{s,\tau}M^{-3/2}.}
\label{eq:V31-whole-word}
\end{equation}
The constant is independent of the physical support and of every
representation label outside \(E_3\).
\end{theorem}

\begin{proof}
Use \Cref{lem:V31-fixed-expansion}.  On an unpaid term, the literal
source-map census of V11 places all output recouplings and the physical
dual contraction in one joint contraction \(D_\nu\) of uniformly
bounded norm.  The V29 commutator census says that \(D_\nu\) preserves
every complement type and local total-output label required by
\Cref{thm:V31-joint}.  Apply
\Cref{cor:V31-three-survive}.

If the payer is present, first apply the V30 linear-payer absorption at
its unique coordinate.  When that coordinate belongs to \(E_3\), the
same weighted-Schmidt proof at time \(s/2\) gives
\(C_s/d_{p_k}\).  When it lies outside \(E_3\), it contributes only a
constant \(C_s\); the strict factors at all other exterior coordinates
remain unchanged.  There is only one such loss.

Now sum the fixed absolute coefficient cost \(C_5\) from
\eqref{eq:V31-fixed-expansion}.  By
\eqref{eq:V31-three-cofinal},
\[
 d_{p_k}=2p_k+1\ge c_\tau M^{1/2},
\]
so \eqref{eq:V31-three-survive} is at most
\(C_{s,\tau}M^{-3/2}\).  The argument with adjoints replaces every
dualizer and separator by its adjoint without changing any norm or
label-preservation property, proving the opposite orientation.
\end{proof}

\begin{corollary}[Large-heat \(q=3\) direct card]
\label{cor:V31-q3-close}
For \(SU(2)\) and \(s\ge s_\sharp\), every one-strong \(q=3\) fibre
with exchange-rescue floor
\[
 \mathfrak x_{T_0,r}
 \ge {c\over M\prod_{i=1}^5\Xi_i}
\]
satisfies \(DXQTRC_5\), including arbitrary dimension-growing
row-\(a\) factors outside \(E_3\).
\end{corollary}

\begin{proof}
After the five row normalizations,
\Cref{thm:V31-whole-word} gives
\[
 \kappa_\otimes^{\rm abs}(\widetilde{\mathbb W})
 \le
 {C_{s,\tau}M^{-3/2}\over\prod_i\Xi_i}
 \le
 {C_{s,\tau}M^{-1}\over\prod_i\Xi_i}.
\]
Compare with the displayed exchange-rescue floor.  The finite
tree/lift resolution changes only the row-order-five constant, and the
adjoint estimate supplies the other coefficient orientation.
\end{proof}

\begin{assessment}[What happened to the V5 exterior partition]
\label{ass:V31-owner-partition}
In the literal source, a cofinal exterior coordinate does not carry an
additional noncentral owner map merely because it was marked during a
previous estimate.  Recursive joining letters are among the at most
four cumulants in \Cref{lem:V31-four-joins}; after their fixed
row-partition expansion they are covered by the joint output estimate.
Tree, lift, selected, strong, and owner labels create no further
operator by V11.  Thus, in the large-heat regime, the old
\(C\), \(N\), and \(R\) exterior alternatives are simultaneously
absorbed without a diffuse-bank extraction.  This statement is about
the exact recursively resolved source and does not assert that an
arbitrary externally inserted noncentral operator is harmless.
\end{assessment}

\section{Why the heat-time restriction is genuine for this route}
\label{sec:V31-small-time}

\begin{proposition}[A small-time block-diagonal dualizer is not a
strict partial-transpose contraction]
\label{prop:V31-small-time}
Take \(p=q=\tfrac12\), let \(P_0,P_1\) be the total-spin zero and one
projections in \(V_{1/2}\otimes V_{1/2}\), and put
\[
 t=e^{-2s},
 \qquad
 D=P_0-P_1.
\]
Then \(D\) is a total-Casimir-preserving contraction, while the
heat-weighted embedded dual operator is
\[
 X_s=P_0-tP_1=(1+t)P_0-tI
\]
and
\begin{equation}
 \|X_s^\Gamma\|={1+3t\over2}.
\label{eq:V31-small-time-norm}
\end{equation}
Consequently
\[
 0<s<{\log3\over2}
 \quad\Longrightarrow\quad
 \|X_s^\Gamma\|>1.
\]
Tensor powers amplify this norm exponentially.
\end{proposition}

\begin{proof}
The singlet is maximally entangled in dimension two, so
\(P_0^\Gamma\) has eigenvalues \(1/2\) on a three-dimensional subspace
and \(-1/2\) on a one-dimensional subspace.  Therefore
\[
 X_s^\Gamma=(1+t)P_0^\Gamma-tI
\]
has eigenvalues
\[
 {1-t\over2}
 \quad\hbox{and}\quad
 -{1+3t\over2}.
\]
This proves \eqref{eq:V31-small-time-norm}.  The last assertion follows
from
\[
 \|(X_s^{\otimes N})^\Gamma\|
 =\|X_s^\Gamma\|^N.
\]
\end{proof}

\begin{firewall}[The countertest is to the compiler, not to
\(DXQTRC_5\)]
\label{fire:V31-small-time-scope}
\Cref{prop:V31-small-time} proves that local total-Casimir preservation
and ordinary dualizer contraction do not yield a support-uniform
strict partial-transpose factor for every \(s>0\).  It does not prove
that the corresponding tensor dualizers occur with the required signs
in a connected quintic coefficient, and it does not compute the exact
row-product capacity after exterior coordinates are coupled.  Thus it
blocks an all-\(s\) extension of the V31 \emph{compiler}; it is not a
counterexample to global quintic MTA.
\end{firewall}

\begin{assessment}[Exact residual after V31]
\label{ass:V31-residual}
The dimension-growing row-\(a\) exterior is closed for
\(SU(2)\) at heat times \(s\ge s_\sharp\).  At arbitrary positive heat
time, V31 reduces the obstruction to a precise question: does the
physical Fourier dual algebra and the signed connected
first-connectivity sum permit repeated realization of the local
block-sign amplification in
\Cref{prop:V31-small-time}, or does source compatibility force an
injective-norm contraction which ordinary partial transpose misses?
This is upstream of the V5 owner-mass ledger and cannot be answered by
calling exterior factors contractions.
\end{assessment}

\section{V31 result ledger and V32 acquisition order}
\label{sec:V31-ledger}

\begin{theorem}[Fourth-series V31 result ledger]
\label{thm:V31-results}
Relative to V1--V30, V31 proves:
\begin{enumerate}[label=\textup{(V31.\arabic*)},leftmargin=6.3em]
\item the one-coordinate weighted Schmidt-sum estimate with arbitrary
      operator-valued output blocks;
\item its simultaneous \(C_s/d_p\) cofinal gain and uniform strict
      contraction above a finite heat-time threshold;
\item a multiplicity-uniform reducible-complement version;
\item an all-at-once joint-dualizer estimate factorizing over every
      row-\(a\) coordinate;
\item survival of the three \(M^{-1/2}\) coordinate gains through
      arbitrary exterior row-\(a\) support;
\item the four-node bound and fixed-cost row-partition expansion for a
      fully recursive quintic source;
\item the whole-word \(CM^{-3/2}\) partial-transpose estimate;
\item \(DXQTRC_5\) for all one-strong \(q=3\) fibres when
      \(s\ge s_\sharp\); and
\item an exact spin-\(\tfrac12\) small-time obstruction to extending
      the same strict partial-transpose compiler using only
      local-Casimir preservation.
\end{enumerate}
\end{theorem}

\begin{proof}
These are
\Cref{lem:V31-operator-matrix,
lem:V31-two-bounds,
cor:V31-reducible,
thm:V31-joint,
cor:V31-three-survive,
lem:V31-four-joins,
lem:V31-fixed-expansion,
thm:V31-whole-word,
cor:V31-q3-close,
prop:V31-small-time}.
\end{proof}

\begin{programme}[V32: the exact small-time exterior algebra]
\label{prog:V31-V32}
\begin{enumerate}[leftmargin=3em]
\item Return to the exact f2 V94 complete-output trace-norm duality and
      characterize the allowed joint dual algebra before any
      partial-transpose majorant.
\item Decide whether tensor products of the block sign
      \(P_0-P_1\) from \Cref{prop:V31-small-time} are admissible in one
      physical connected quintic fibre, with all source signs retained.
\item If they are not admissible, state and prove the missing
      source-compatibility restriction and insert it into
      \Cref{thm:V31-joint}.
\item If they are admissible, compute the row-product injective norm of
      the heat-weighted Werner block and its tensor powers; do not infer
      it from the larger partial-transpose norm.
\item Use the exact signed recursive first-connectivity sum before
      taking positive parts.  A cancellation which appears only after
      summing different fibres must be kept at that level.
\item Extend \Cref{cor:V31-q3-close} to every \(s>0\), or isolate an
      explicit physical small-time exterior packet which remains.
\end{enumerate}
\end{programme}

\begin{firewall}[Claim boundary after fourth-series V31]
\label{fire:V31-claim-boundary}
V31 closes the complete dimension-growing exterior and hence the
one-strong \(q=3\) direct card only for \(SU(2)\) heat times
\(s\ge s_\sharp\).  It does not close the arbitrary-small-time
\(q=3\) problem, the \(q=0,1,2\) sectors, global quintic MTA,
all-order MTA, WUFC, CPM, cutoff removal, reflection positivity,
Osterwalder--Schrader reconstruction, construction of the continuum
Yang--Mills measure, or a uniform positive continuum spectral gap.
The full Yang--Mills mass gap has not been proved.
\end{firewall}

\paragraph{Parent-version record.}
V31 was formed from
\texttt{Renewal\_Yang\_Mills\_Mass\_Gap\_Research\_Notes\_V30.tex}.
The V30 parent had \(14\,763\) logical source lines,
\(518\,925\) bytes, and SHA--256
\[
 \texttt{69940730ffc793d014225869a6532feea7e2ef30660b94762150305be4ede522}.
\]
The next compulsory companion audit is V35.

% =============================================================================
% END APPENDED FOURTH-SERIES V31 MATERIAL
% =============================================================================

% The V31 document terminator is intentionally disabled here:
% \end{document}

% =============================================================================
% BEGIN APPENDED FOURTH-SERIES V32 MATERIAL
% =============================================================================

\clearpage
\chapter{V32: Exact final-output dual algebra and all-heat-time
\texorpdfstring{\(q=3\)}{q=3} closure}
\label{ch:V32}

\section{The V31 enlargement is not the physical dual algebra}
\label{sec:V32-purpose}

V31 deliberately enlarged the physical output contraction to an
arbitrary contraction on every tuple of local output spaces.  That
enlargement is safe for an upper bound and closes the exterior at
large heat time, but the spin-\(\tfrac12\) sign test shows that it is
too large for a support-uniform strict partial-transpose argument at
small heat time.

The exact f2 V94 duality is substantially more rigid.  For the final
global Schur decomposition
\begin{equation}
 \mathcal H_{\rm in}
 =
 \bigoplus_U W_U(H_U\otimes M_U),
\label{eq:V32-global-Schur}
\end{equation}
the polar dualizer has the form
\begin{equation}
 \mathbb D
 =
 \bigoplus_U W_U(D_U^*\otimes I_{M_U})W_U^*,
 \qquad \|D_U\|\le1.
\label{eq:V32-physical-dualizer}
\end{equation}
Every complete equivariant source multiplier acts on the opposite
Schur factor:
\begin{equation}
 \mathbb H
 =
 \bigoplus_U W_U(I_{H_U}\otimes h_U)W_U^*.
\label{eq:V32-multiplicity-operator}
\end{equation}
Thus \(\mathbb D\) and \(\mathbb H\) commute.  V29 used only the weaker
consequence that \(\mathbb D\) preserves local total-Casimir labels.
Here the full tensor-factor identity in
\eqref{eq:V32-physical-dualizer} is decisive.

\section{The positive exact-dual sandwich}
\label{sec:V32-positive-sandwich}

\begin{lemma}[Exact-dual positivity domination]
\label{lem:V32-positive-dual}
Let \(\mathbb H\ge0\) have the equivariant block form
\eqref{eq:V32-multiplicity-operator}, and let \(\mathbb D\) be any
exact physical dualizer \eqref{eq:V32-physical-dualizer}.  Then
\begin{equation}
 \mathbb H\mathbb D
 =
 \mathbb H^{1/2}\mathbb D\mathbb H^{1/2}
\label{eq:V32-positive-factorization}
\end{equation}
and, for every density matrix \(\varrho\),
\begin{equation}
 \boxed{
 \left|\Tr(\varrho\mathbb H\mathbb D)\right|
 \le \Tr(\varrho\mathbb H).}
\label{eq:V32-positive-domination}
\end{equation}
The assertion holds with \(\mathbb D\mathbb H\), adjoints, and
arbitrary finite direct sums of global output blocks.
\end{lemma}

\begin{proof}
On the \(U\)-block,
\[
 (I\otimes h_U)(D_U^*\otimes I)
 =
 D_U^*\otimes h_U
 =
 (I\otimes h_U^{1/2})(D_U^*\otimes I)
   (I\otimes h_U^{1/2}).
\]
This proves \eqref{eq:V32-positive-factorization}.  Put
\[
 Y:=\mathbb H^{1/2}\varrho\mathbb H^{1/2}\ge0.
\]
Since \(\|\mathbb D\|\le1\), trace-norm duality gives
\[
 |\Tr(Y\mathbb D)|
 \le\|Y\|_1\|\mathbb D\|
 \le\Tr Y
 =\Tr(\varrho\mathbb H).
\]
The same calculation is invariant under adjoints and block direct
sums.
\end{proof}

\begin{corollary}[Exact dualizers do not amplify a positive moment
word]
\label{cor:V32-no-dual-amplification}
For every positive coordinatewise moment-partition word
\(\mathbb H_\nu\),
\begin{equation}
 \kappa_\otimes^{\rm abs}(\mathbb H_\nu\mathbb D)
 \le \kappa_\otimes(\mathbb H_\nu).
\label{eq:V32-no-dual-amplification}
\end{equation}
No heat-time threshold and no partial transpose of \(\mathbb D\) is
required.
\end{corollary}

\begin{proof}
Every complete moment-partition heat factor is an intertwiner, so its
coordinate tensor product has
\eqref{eq:V32-multiplicity-operator}.  Apply
\Cref{lem:V32-positive-dual} to every row-product density and take the
supremum.
\end{proof}

\begin{remark}[Why this does not revive the false square card]
\label{rem:V32-not-square-card}
The domination is linear in the exact dualized coefficient.  It does
not assert
\[
 (\mathbb H_\nu\mathbb D)^*(\mathbb H_\nu\mathbb D)
 \preccurlyeq \mathbb H_\nu^2
\]
for a noncommuting row cut, nor does it imply the false pointwise
\(XQTRSC_5\) square interface excluded in V24.  It is precisely the
direct route required by \(DXQTRC_5\).
\end{remark}

\section{The V31 tensor sign is excluded after one coordinate}
\label{sec:V32-sign}

For spin-\(\tfrac12\) input pairs, write
\[
 E:=V_0\oplus V_1,
 \qquad
 S:=P_0-P_1.
\]
For \(N\) physical coordinates, resolve
\[
 E^{\otimes N}
 =
 \bigoplus_U V_U\otimes M_U^{(N)}.
\]

\begin{proposition}[Multiplicity consistency of the physical dual
algebra]
\label{prop:V32-dual-algebra}
The exact physical dual algebra on \(E^{\otimes N}\) is
\begin{equation}
 \mathfrak A_N
 :=
 \bigoplus_U\mathcal B(V_U)\otimes I_{M_U^{(N)}}.
\label{eq:V32-dual-algebra}
\end{equation}
An operator which is scalar on every local-output tuple belongs to
\(\mathfrak A_N\) only if its scalar is the same on all multiplicity
copies of each fixed global \(U\).
\end{proposition}

\begin{proof}
Equation \eqref{eq:V32-dual-algebra} is exactly the family
\eqref{eq:V32-physical-dualizer} furnished by f2 V94 block trace-norm
duality.  On a fixed global \(U\), its restriction is one operator on
\(V_U\) tensored with the identity on the whole multiplicity space.
If it is also scalar on each local-output tuple, comparison on two
tuple subspaces contributing to the same \(U\) forces those scalars to
be equal.
\end{proof}

\begin{corollary}[Repeated local block signs are not exact Fourier
dualizers]
\label{cor:V32-sign-excluded}
The one-coordinate sign \(S\) belongs to \(\mathfrak A_1\), but
\begin{equation}
 S^{\otimes N}\notin\mathfrak A_N
 \qquad(N\ge2).
\label{eq:V32-sign-excluded}
\end{equation}
\end{corollary}

\begin{proof}
For \(N=1\), local total spin is the global output, so the choices
\(D_0=1\), \(D_1=-I_{V_1}\) realize \(S\).

For \(N=2\), the tuples \((K_1,K_2)=(0,1)\) and \((1,1)\) both contain
global output \(U=1\).  The operator \(S\otimes S\) has scalar
\(-1\) on the first tuple and \(+1\) on the second.  This violates the
multiplicity consistency in
\Cref{prop:V32-dual-algebra}.  For \(N>2\), tensor both comparison
subspaces with \(N-2\) local singlets.  Their global output remains
\(U=1\), and the same contradiction persists.
\end{proof}

\begin{assessment}[Resolution of the V31 small-time test]
\label{ass:V32-sign-resolution}
The one-coordinate computation
\(\|X_s^\Gamma\|=(1+3e^{-2s})/2\) remains correct and proves that the
V31 enlarged local-dual compiler is not strictly contractive for all
\(s>0\).  Its tensor-power amplification is not an exact V94 dual
family once two coordinates are present.  More importantly,
\Cref{lem:V32-positive-dual} controls every exact family, including
non-scalar matrices \(D_U\), without classifying their partial
transpose.
\end{assessment}

\section{Exterior erasure for one positive moment word}
\label{sec:V32-exterior}

Fix a maximal row \(a\), three distinct coordinates
\[
 E_3=\{e_1,e_2,e_3\},
\]
and one term \(\mathbb H_\nu\) in the fixed row-partition expansion
\eqref{eq:V31-fixed-expansion}.  At each coordinate, its local factor
is a positive complete moment-partition heat operator.  Factors on
different coordinates tensorize.

\begin{lemma}[Unpaid exterior factors erase by positive order]
\label{lem:V32-exterior-order}
Let \(\mathbb H_{\nu,E_3}\) be the product of the three local factors
on \(E_3\).  If the term is unpaid, then
\begin{equation}
 0\preccurlyeq\mathbb H_\nu
 \preccurlyeq
 \mathbb H_{\nu,E_3}\otimes I_{E_3^c}.
\label{eq:V32-exterior-order}
\end{equation}
Consequently
\begin{equation}
 \kappa_\otimes(\mathbb H_\nu)
 \le
 \left\|
  \mathbb H_{\nu,E_3}^{\Gamma_{a^c}}
 \right\|.
\label{eq:V32-exterior-capacity}
\end{equation}
\end{lemma}

\begin{proof}
Every normalized heat multiplier lies in \([0,1]\); hence every local
moment-partition product is a positive contraction.  Tensoring the
coordinatewise inequalities \(0\preccurlyeq H_e\preccurlyeq I\)
outside \(E_3\) proves \eqref{eq:V32-exterior-order}.

Expectation against a positive density preserves this order.  A
row-product density is, in particular, a product density across
\(a|a^c\), and the partial transpose of such a density is again a
density.  Therefore its expectation against
\(\mathbb H_{\nu,E_3}\) is at most the operator norm of the displayed
partial transpose.  The identity exterior does not change that norm.
\end{proof}

\begin{lemma}[The unique payer costs only one constant]
\label{lem:V32-paid-exterior}
For fixed \(s>0\), the conclusion of
\Cref{lem:V32-exterior-order} remains valid up to a constant \(C_s\)
when the unique payer lies outside \(E_3\).  If it lies in \(E_3\), the
local factor there satisfies the paid V30 response
\begin{equation}
 \|H_{e_k}^{\rm pay,\Gamma_{a^c}}\|
 \le {C_s\over d_{p_k}}.
\label{eq:V32-paid-local}
\end{equation}
\end{lemma}

\begin{proof}
At an exterior payer coordinate, V30 linear absorption gives
\[
 0\preccurlyeq
 R(C_e)\boldsymbol1_{\{C_e>0\}}\mathbb M_{\pi,e}(s)
 \preccurlyeq
 C_s\mathbb M_{\pi,e}(s/2)
 \preccurlyeq C_sI.
\]
All other exterior factors are positive contractions.  This proves the
first assertion.  At a coordinate in \(E_3\), use the paid complete
moment-partition case of V30.  A selected complete output projection
commutes with the local heat factors and is already included in that
estimate.
\end{proof}

\begin{theorem}[Three-coordinate positive-word response at every heat
time]
\label{thm:V32-positive-three}
Let \(s>0\).  Suppose row \(a\) carries spin \(p_k\) at \(e_k\) and
\begin{equation}
 1+p_k(p_k+1)\ge\tau M,
 \qquad
 M=\Xi_a,
 \qquad k=1,2,3.
\label{eq:V32-three-cofinal}
\end{equation}
For every term of
\eqref{eq:V31-fixed-expansion}, paid or unpaid,
\begin{equation}
 \boxed{
 \kappa_\otimes^{\rm abs}(\mathbb H_\nu\mathbb D)
 \le C_{s,\tau}M^{-3/2}.}
\label{eq:V32-positive-three}
\end{equation}
The estimate is uniform in all exterior representations,
multiplicities, and physical support.
\end{theorem}

\begin{proof}
First apply
\Cref{cor:V32-no-dual-amplification}.  Remove the exterior by
\Cref{lem:V32-exterior-order,lem:V32-paid-exterior}.  At each of the
three retained coordinates, the V30 moment-partition response with
identity output contraction gives
\[
 \|H_{e_k}^{\Gamma_{a^c}}\|
 \le {C_s\over d_{p_k}},
\]
with \eqref{eq:V32-paid-local} at the payer coordinate if it is among
them.  The three factors lie at distinct physical coordinates, so
their partial-transpose norms multiply.  Finally
\eqref{eq:V32-three-cofinal} implies
\[
 d_{p_k}\ge c_\tau M^{1/2}.
\]
This proves \eqref{eq:V32-positive-three}.
\end{proof}

\section{All-\texorpdfstring{\(s\)}{s} closure of the direct
\texorpdfstring{\(q=3\)}{q=3} card}
\label{sec:V32-q3}

\begin{theorem}[Complete whole-word \(q=3\) direct bound]
\label{thm:V32-whole-word}
Let \(G=SU(2)\) and \(s>0\).  Under the hypotheses of
\Cref{thm:V32-positive-three}, retain the complete recursively resolved
quintic source, all prefix and joining cumulants, the suffix, every
physical output and multiplicity, the exact V94 dual family, and the
unique payer.  Then
\begin{equation}
 \boxed{
 \kappa_\otimes^{\rm abs}(\mathbb W^{\rm raw})
 \le C_{s,\tau}M^{-3/2}.}
\label{eq:V32-whole-word}
\end{equation}
The same estimate holds for the adjoint coefficient orientation.
\end{theorem}

\begin{proof}
Use the exact fixed-cost expansion
\eqref{eq:V31-fixed-expansion}.  Each
\(\mathbb H_\nu\) is positive and equivariant.  Apply
\Cref{thm:V32-positive-three} term by term and sum
\[
 \sum_{\nu\in\mathcal R_5}|c_\nu|\le C_5.
\]
This is a fixed row-partition expansion at at most four recursive
joining nodes, not a support-sized positive envelope.  The exact V94
dualizer is kept through
\Cref{lem:V32-positive-dual}; no local-output dual family is introduced.
Replacing the source and dualizer by their adjoints leaves the positive
moment word and all norms unchanged.
\end{proof}

\begin{corollary}[The one-strong \(q=3\) sector is closed]
\label{cor:V32-q3-closed}
For \(SU(2)\) and every \(s>0\), every one-strong \(q=3\) fibre with
\begin{equation}
 \mathfrak x_{T_0,r}
 \ge {c\over M\prod_{i=1}^5\Xi_i}
\label{eq:V32-exchange-floor}
\end{equation}
satisfies \(DXQTRC_5\).  Dimension-growing maximal-row factors outside
the three response coordinates create no residual branch.
\end{corollary}

\begin{proof}
After row normalization,
\eqref{eq:V32-whole-word} gives
\[
 \kappa_\otimes^{\rm abs}(\widetilde{\mathbb W})
 \le
 {C_{s,\tau}M^{-3/2}\over\prod_i\Xi_i}
 \le
 {C_{s,\tau}M^{-1}\over\prod_i\Xi_i}.
\]
Compare with \eqref{eq:V32-exchange-floor}.  The fixed tree/lift
resolution and both orientations are already included in
\Cref{thm:V32-whole-word}.
\end{proof}

\begin{firewall}[Scope of the \(q=3\) closure]
\label{fire:V32-q3-scope}
\Cref{cor:V32-q3-closed} closes the residual one-strong
\(q=3\) direct fibres identified in the fourth-series reduction.  It
does not restore the false pointwise square card of V24 and does not
address the larger missing responses in the \(q=0,1,2\) sectors.
Consequently global quintic MTA is not yet proved.
\end{firewall}

\begin{assessment}[Exact residual after V32]
\label{ass:V32-residual}
The fourth-series \(q=3\) bottleneck is removed for \(SU(2)\) at every
positive heat time.  The next finite-order obstruction is \(q=2\):
two cofinal complete heat coordinates supply \(M^{-1}\), whereas the
one-strong raw normalization needs a direct \(M^{-1}\) response and
the old oriented-square formulation asked for one further
half-power.  V24 requires working at the direct coefficient level.
The first question is therefore whether the exact-dual positivity
sandwich already makes the two ordinary coordinate responses
sufficient after the exchange floor is normalized, before seeking any
new owner surplus.
\end{assessment}

\section{V32 result ledger and V33 acquisition order}
\label{sec:V32-ledger}

\begin{theorem}[Fourth-series V32 result ledger]
\label{thm:V32-results}
Relative to V1--V31, V32 proves:
\begin{enumerate}[label=\textup{(V32.\arabic*)},leftmargin=6.3em]
\item the exact final-output/multiplicity tensor-factor separation of
      the V94 physical dual algebra;
\item positive domination of every exact dualized equivariant moment
      word;
\item the multiplicity-consistency criterion for local-output tuple
      scalars;
\item exclusion of repeated V31 local block signs from the exact
      Fourier dual algebra;
\item positive-order erasure of arbitrary exterior support;
\item a one-constant treatment of every payer placement;
\item the three-coordinate \(CM^{-3/2}\) direct response for every
      \(s>0\);
\item the complete whole-word bound after the fixed recursive
      row-partition expansion; and
\item closure of all remaining one-strong \(q=3\) direct fibres for
      \(SU(2)\).
\end{enumerate}
\end{theorem}

\begin{proof}
These are
\Cref{lem:V32-positive-dual,
cor:V32-no-dual-amplification,
prop:V32-dual-algebra,
cor:V32-sign-excluded,
lem:V32-exterior-order,
lem:V32-paid-exterior,
thm:V32-positive-three,
thm:V32-whole-word,
cor:V32-q3-closed}.
\end{proof}

\begin{programme}[V33: direct normalization of the \(q=2\) sector]
\label{prog:V32-V33}
\begin{enumerate}[leftmargin=3em]
\item Return to the exact f3 V99--V100 \(q\)-slot normalization and
      distinguish its old square deficit from the direct
      \(DXQTRC_5\) target.
\item For two cofinal coordinates, apply the V32 positive exact-dual
      sandwich before any square or Cauchy--Schwarz step.
\item Compare the resulting
      \(M^{-1}/\prod_i\Xi_i\) direct bound with the exact one-strong
      exchange-rescue floor, including every endpoint-mass factor.
\item If the exponents match, close \(q=2\) without inventing a third
      owner.  If they do not, state the exact remaining scalar debit.
\item Check every payer placement using one V30 absorption and retain
      the fixed recursive row-partition expansion.
\item Only after the exact comparison move to \(q=1\), whose direct
      deficit may still require an additional response.
\end{enumerate}
\end{programme}

\begin{firewall}[Claim boundary after fourth-series V32]
\label{fire:V32-claim-boundary}
V32 closes the complete \(SU(2)\) one-strong \(q=3\) direct sector for
all \(s>0\).  It does not close \(q=0,1,2\), global quintic MTA,
all-order MTA, WUFC, CPM, cutoff removal, reflection positivity,
Osterwalder--Schrader reconstruction, construction of the continuum
Yang--Mills measure, or a uniform positive continuum spectral gap.
The full Yang--Mills mass gap has not been proved.
\end{firewall}

\paragraph{Parent-version record.}
V32 was formed from
\texttt{Renewal\_Yang\_Mills\_Mass\_Gap\_Research\_Notes\_V31.tex}.
The V31 parent had \(15\,457\) logical source lines,
\(543\,284\) bytes, and SHA--256
\[
 \texttt{3a74b4f20574de28d3a1fe7b36bc7c110d891065421d849788536f155cdb768d}.
\]
The next compulsory companion audit is V35.

% =============================================================================
% END APPENDED FOURTH-SERIES V32 MATERIAL
% =============================================================================

% The V32 document terminator is intentionally disabled here:
% \end{document}

% =============================================================================
% BEGIN APPENDED FOURTH-SERIES V33 MATERIAL
% =============================================================================

\clearpage
\chapter{V33: Direct normalization closes the
\texorpdfstring{\(q=2\)}{q=2} sector}
\label{ch:V33}

\section{The old four-slot budget belongs to a square interface}
\label{sec:V33-purpose}

Archive f3 V99--V100 correctly computes the one-strong exchange floor
\begin{equation}
 \mathfrak x_{T_0,r}
 \ge {c\over M\prod_{i=1}^5\Xi_i}.
\label{eq:V33-exchange-floor}
\end{equation}
It then seeks an oriented-square comparison.  Since the normalized raw
word is only known there at scale
\((\prod_i\Xi_i)^{-1}\), squaring and comparing with
\eqref{eq:V33-exchange-floor} asks for \(M^{-2}\) in the square.  This
is the origin of the four-slot table
\[
 \delta_q={4-q\over2}.
\]

V24 proves that the universal pointwise square interface is false.
V25 replaces it by the direct coefficient target \(DXQTRC_5\).
For that target, \eqref{eq:V33-exchange-floor} asks for only
\begin{equation}
 {M^{-1}\over\prod_i\Xi_i}
\label{eq:V33-direct-target}
\end{equation}
at normalized word level.  Therefore two ordinary cofinal coordinate
responses, each of size \(M^{-1/2}\), already have the exact direct
power.  No third owner and no square surplus are required.

\section{The corrected direct slot budget}
\label{sec:V33-budget}

\begin{definition}[Direct \(q\)-slot deficit]
\label{def:V33-direct-deficit}
For \(q\in\{0,1,2,3\}\), define
\begin{equation}
 \varepsilon_q
 :=
 \max\left\{0,1-{q\over2}\right\}.
\label{eq:V33-epsilonq}
\end{equation}
It is the additional raw inverse-\(M\) power needed after
\(q\) direct coordinate responses of total size \(M^{-q/2}\) have
been retained.
\end{definition}

\begin{lemma}[Exact direct response table]
\label{lem:V33-direct-table}
The direct one-strong response budget is
\begin{center}
\begin{tabular}{c|cccc}
\(q\)&0&1&2&3\\
\hline
coordinate response&\(1\)&\(M^{-1/2}\)&\(M^{-1}\)&\(M^{-3/2}\)\\
additional direct power&\(M^{-1}\)&\(M^{-1/2}\)&\(1\)&\(1\).
\end{tabular}
\end{center}
Equivalently, the remaining source must provide
\(M^{-\varepsilon_q}\).  In particular, \(q=2\) and \(q=3\) have no
remaining direct exponent.
\end{lemma}

\begin{proof}
The scalar target is \(M^{-1}\), not the square target \(M^{-2}\).
Division by the available factor \(M^{-q/2}\) leaves
\[
 M^{-(1-q/2)}
\]
when \(q<2\), and no negative power when \(q\ge2\).  This is exactly
\eqref{eq:V33-epsilonq} and the displayed table.
\end{proof}

\begin{proposition}[Relation to the archived square budget]
\label{prop:V33-budget-relation}
For \(q=0,1,2,3\), the archived square deficit and the direct deficit
satisfy
\begin{equation}
 \delta_q=1-{q\over2}+1,
 \qquad
 \varepsilon_q=\max\{0,\delta_q-1\}.
\label{eq:V33-budget-relation}
\end{equation}
The lost unit is exactly the passage from a linear coefficient target
to its oriented square.
\end{proposition}

\begin{proof}
Substitute
\(\delta_q=(4-q)/2=2-q/2\).  The direct target has exponent one,
whereas the square target has exponent two.  Their difference is one.
Truncation at zero records that extra decay is harmless but not owed.
\end{proof}

\begin{firewall}[No theorem from the false square card is reused]
\label{fire:V33-no-square}
The scalar exchange floor
\eqref{eq:V33-exchange-floor}, the exact source decomposition, and
the local heat response theorems remain valid.  V33 does not cite the
false implication from a resolved oriented square to a universal
pointwise card.  It compares the direct absolute coefficient with the
nonnegative exchange weight exactly as required by \(DXQTRC_5\).
\end{firewall}

\section{Two positive response coordinates}
\label{sec:V33-two}

\begin{theorem}[Two-coordinate positive-word response]
\label{thm:V33-positive-two}
Let \(G=SU(2)\), \(s>0\), and let
\[
 E_2=\{e_1,e_2\}
\]
be two distinct coordinates in one moment-partition word
\(\mathbb H_\nu\).  Suppose row \(a\) carries spin \(p_k\) at \(e_k\)
and
\begin{equation}
 1+p_k(p_k+1)\ge\tau M,
 \qquad
 M=\Xi_a,
 \qquad k=1,2.
\label{eq:V33-two-cofinal}
\end{equation}
Retain an arbitrary exact V94 final-output dual family and the unique
payer.  Then
\begin{equation}
 \boxed{
 \kappa_\otimes^{\rm abs}(\mathbb H_\nu\mathbb D)
 \le C_{s,\tau}M^{-1}.}
\label{eq:V33-positive-two}
\end{equation}
The bound is uniform in all exterior representations, multiplicities,
and physical support.
\end{theorem}

\begin{proof}
By the exact-dual positivity sandwich
\Cref{lem:V32-positive-dual},
\[
 \kappa_\otimes^{\rm abs}(\mathbb H_\nu\mathbb D)
 \le \kappa_\otimes(\mathbb H_\nu).
\]
All unpaid exterior local factors are positive contractions and erase
by the positive-order proof of
\Cref{lem:V32-exterior-order}, now with \(E_2\) in place of \(E_3\).
If the payer lies outside \(E_2\), V30 absorption bounds its complete
local factor by \(C_sI\).  If it lies at \(e_k\), retain the paid local
response there.

At each retained coordinate, the complete V30 moment-partition theorem
gives
\[
 \|H_{e_k}^{\Gamma_{a^c}}\|
 \le {C_s\over d_{p_k}},
\]
including every local incidence, output multiplicity, selected
projection, and the sole payer when present.  The two coordinates
tensorize, so
\[
 \kappa_\otimes(\mathbb H_\nu)
 \le {C_s\over d_{p_1}d_{p_2}}.
\]
From \eqref{eq:V33-two-cofinal},
\(d_{p_k}\ge c_\tau M^{1/2}\), proving
\eqref{eq:V33-positive-two}.
\end{proof}

\begin{theorem}[Complete whole-word \(q=2\) response]
\label{thm:V33-whole-word}
Under the hypotheses of
\Cref{thm:V33-positive-two}, let
\(\mathbb W^{\rm raw}\) be the complete recursively resolved quintic
source fibre, with all prefix and joining cumulants, suffix moments,
outputs, multiplicities, exact duality, and payer retained.  Then
\begin{equation}
 \boxed{
 \kappa_\otimes^{\rm abs}(\mathbb W^{\rm raw})
 \le C_{s,\tau}M^{-1}.}
\label{eq:V33-whole-word}
\end{equation}
The adjoint orientation obeys the same bound.
\end{theorem}

\begin{proof}
Use the support-independent signed expansion
\[
 \mathbb W^{\rm raw}
 =
 \sum_{\nu\in\mathcal R_5}c_\nu\mathbb H_\nu
\]
from \Cref{lem:V31-fixed-expansion}.  Apply
\Cref{thm:V33-positive-two} to every positive equivariant moment word
before taking absolute values, and then use
\[
 \sum_\nu|c_\nu|\le C_5.
\]
The exact dual family is controlled inside each positive sandwich.
Taking adjoints preserves positivity of the moment words and the norm
of every \(D_U\).
\end{proof}

\begin{corollary}[The one-strong \(q=2\) sector is closed]
\label{cor:V33-q2-closed}
Every one-strong \(q=2\) \(SU(2)\) fibre satisfying the exchange floor
\eqref{eq:V33-exchange-floor} obeys \(DXQTRC_5\) for every \(s>0\).
\end{corollary}

\begin{proof}
Normalize \eqref{eq:V33-whole-word} by the five row masses:
\[
 \kappa_\otimes^{\rm abs}(\widetilde{\mathbb W})
 \le
 {C_{s,\tau}M^{-1}\over\prod_i\Xi_i}.
\]
This is a constant multiple of the right side of
\eqref{eq:V33-exchange-floor}.  The fixed tree/lift family sums into
the exchange-rescue weight, exactly as in the direct compiler of V25.
\end{proof}

\begin{corollary}[All \(q\ge2\) one-strong fibres are closed]
\label{cor:V33-qge2}
For \(SU(2)\) and \(s>0\), no one-strong direct obstruction remains in
the \(q=2\) or \(q=3\) sectors.
\end{corollary}

\begin{proof}
Use \Cref{cor:V33-q2-closed} for \(q=2\) and
\Cref{cor:V32-q3-closed} for \(q=3\).
\end{proof}

\section{The first genuinely unpaid direct exponent}
\label{sec:V33-residual}

\begin{assessment}[Exact residual after V33]
\label{ass:V33-residual}
Only \(q=0\) and \(q=1\) remain in the one-strong direct core.  Their
exact additional raw powers are respectively
\[
 M^{-1}
 \qquad\hbox{and}\qquad
 M^{-1/2}.
\]
For \(q=1\), one more ordinary complete coordinate response would
close the fibre, even if that coordinate owns a joining or payer role;
V30--V32 show that such roles must be assembled but do not require a
surplus exponent.  If no second cofinal coordinate exists, the
owner-mass ledger must force either a source-occurring diffuse bank or
mass concentration inside the single already-used coordinate.  That
is now the correct use of the V5 partition.
\end{assessment}

\begin{firewall}[The global quintic theorem is still open]
\label{fire:V33-global-open}
The result removes two rows of the direct slot table.  It does not
prove that every one-strong residual has \(q\ge2\), and it does not
control the \(M^{-1/2}\) or \(M^{-1}\) direct deficits in the \(q=1\)
or \(q=0\) sectors.  Hence global quintic MTA is not yet proved.
\end{firewall}

\section{V33 result ledger and V34 acquisition order}
\label{sec:V33-ledger}

\begin{theorem}[Fourth-series V33 result ledger]
\label{thm:V33-results}
Relative to V1--V32, V33 proves:
\begin{enumerate}[label=\textup{(V33.\arabic*)},leftmargin=6.3em]
\item the corrected direct slot deficit
      \(\varepsilon_q=\max\{0,1-q/2\}\);
\item the exact direct \(q=0,1,2,3\) response table;
\item identification of the extra unit in the archived table as the
      cost of the withdrawn oriented-square interface;
\item a two-coordinate positive-word response \(CM^{-1}\) with exact
      final-output duality and every payer placement;
\item the complete recursively assembled whole-word \(CM^{-1}\)
      response; and
\item closure of every one-strong \(q=2\) fibre, hence of all
      \(q\ge2\) one-strong direct fibres, for \(SU(2)\).
\end{enumerate}
\end{theorem}

\begin{proof}
These are
\Cref{def:V33-direct-deficit,
lem:V33-direct-table,
prop:V33-budget-relation,
thm:V33-positive-two,
thm:V33-whole-word,
cor:V33-q2-closed,
cor:V33-qge2}.
\end{proof}

\begin{programme}[V34: the \(q=1\) owner-mass dichotomy]
\label{prog:V33-V34}
\begin{enumerate}[leftmargin=3em]
\item Fix a \(q=1\) one-strong fibre and one cofinal response coordinate
      \(e_1\).  Return to the exact row-\(a\) mass outside \(e_1\).
\item Remove the finite recursive joining/payer set only after its
      complete local factor has been classified by V30.
\item If another exterior coordinate carries a fixed fraction of
      \(M\), apply the V33 two-coordinate theorem even when it is a
      marked role owner.
\item If exterior mass is diffuse, expose a complete source-occurring
      moment bank and use V87 with only the one additional direct
      half-power required, not the old square exponent.
\item If almost all mass lies at \(e_1\), descend to the exact local
      one-coordinate quotient cumulant and test whether its direct
      response exceeds the ordinary \(d_{p_1}^{-1}\) scale.  Do not
      count the same coordinate twice.
\item State the surviving concentrated packet explicitly before the
      compulsory V35 companion audit.
\end{enumerate}
\end{programme}

\begin{firewall}[Claim boundary after fourth-series V33]
\label{fire:V33-claim-boundary}
V33 closes all one-strong \(q\ge2\) direct fibres for \(SU(2)\).  It
does not close \(q=0,1\), global quintic MTA, all-order MTA, WUFC,
CPM, cutoff removal, reflection positivity, Osterwalder--Schrader
reconstruction, construction of the continuum Yang--Mills measure, or
a uniform positive continuum spectral gap.  The full Yang--Mills mass
gap has not been proved.
\end{firewall}

\paragraph{Parent-version record.}
V33 was formed from
\texttt{Renewal\_Yang\_Mills\_Mass\_Gap\_Research\_Notes\_V32.tex}.
The V32 parent had \(15\,958\) logical source lines,
\(560\,272\) bytes, and SHA--256
\[
 \texttt{f79931f91f3375c949c9a137fed92aefe3bcae10281f8e0d1b7dd86c452747a1}.
\]
The next compulsory companion audit is V35.

% =============================================================================
% END APPENDED FOURTH-SERIES V33 MATERIAL
% =============================================================================

% The V33 document terminator is intentionally disabled here:
% \end{document}

% =============================================================================
% BEGIN APPENDED FOURTH-SERIES V34 MATERIAL
% =============================================================================

\clearpage
\chapter{V34: The \texorpdfstring{\(q=1\)}{q=1} exterior-mass
dichotomy and the single-coordinate terminal}
\label{ch:V34}

\section{Set-up}
\label{sec:V34-setup}

Fix a sequence of one-strong \(q=1\) fibres with maximal row \(a\),
\[
 \Xi_a=M\longrightarrow\infty,
\]
and one usable cofinal response coordinate \(e_1\), so that
\begin{equation}
 a_{a,e_1}\ge\tau M.
\label{eq:V34-first-slot}
\end{equation}
Resolve every connected prefix recursively as in V31.  Let \(J\) be
the set of physical coordinates outside \(e_1\) carrying at least one
recursive joining occurrence or the unique payer.  By
\Cref{lem:V31-four-joins},
\begin{equation}
 |J|\le5.
\label{eq:V34-marked-count}
\end{equation}
Tree, lift, selected, and strong labels do not enlarge \(J\), because
V11 proves that they create no additional source map.

Let
\begin{equation}
 D:=
 \operatorname{supp}(a)\setminus(\{e_1\}\sqcup J)
\label{eq:V34-D}
\end{equation}
be the unmarked exterior.  In every fixed term of the V31
row-partition expansion, all coordinates in \(D\) occur as a complete
source moment bank disjoint from the payer.

\section{An exhaustive three-way alternative}
\label{sec:V34-trichotomy}

\begin{theorem}[\(q=1\) exterior-mass trichotomy]
\label{thm:V34-trichotomy}
After passage to a subsequence, at least one of the following holds:
\begin{enumerate}[label=\textup{(\roman*)},leftmargin=4em]
\item there is a coordinate \(f\ne e_1\) and a fixed
      \(\sigma>0\) such that
\begin{equation}
 a_{a,f}\ge\sigma M;
\label{eq:V34-second-atom}
\end{equation}
\item there are fixed \(\sigma>0\) and unmarked exterior banks
      \(D_M\subseteq D\) such that
\begin{equation}
 \sum_{e\in D_M}a_{a,e}\ge\sigma M,
 \qquad
 \max_{e\in D_M}{a_{a,e}\over M}\longrightarrow0;
\label{eq:V34-diffuse}
\end{equation}
\item the first coordinate carries asymptotically all maximal-row
      mass:
\begin{equation}
 {a_{a,e_1}\over M}\longrightarrow1.
\label{eq:V34-saturation}
\end{equation}
\end{enumerate}
\end{theorem}

\begin{proof}
If
\[
 \limsup_M\max_{f\ne e_1}{a_{a,f}\over M}>0,
\]
select a coordinate attaining the maximum and pass to a subsequence.
This gives \textup{(i)}.

Suppose \textup{(i)} fails.  Since \(J\) has at most five elements,
\[
 {1\over M}\sum_{f\in J}a_{a,f}\longrightarrow0.
\]
If
\[
 \limsup_M {1\over M}\sum_{e\in D}a_{a,e}>0,
\]
pass to a subsequence on which this sum is at least \(\sigma M\) and
take \(D_M=D\).  Failure of \textup{(i)} gives the atomwise limit in
\eqref{eq:V34-diffuse}, so \textup{(ii)} holds.

If this second limsup is zero, the disjoint mass identity
\[
 M
 =
 a_{a,e_1}
 +\sum_{f\in J}a_{a,f}
 +\sum_{e\in D}a_{a,e}
\]
gives \eqref{eq:V34-saturation}.
\end{proof}

\begin{remark}[Why the finite marked set is harmless in the
dichotomy]
\label{rem:V34-marked}
A positive mass fraction cannot hide among joining and payer
coordinates without creating a second cofinal atom, because their
number is bounded by \eqref{eq:V34-marked-count}.  No centrality is
used for this counting statement.  Analytic control is applied only
after the complete marked local factor has been expanded and retained.
\end{remark}

\section{The second-atom branch}
\label{sec:V34-atom}

\begin{theorem}[Every second cofinal atom closes \(q=1\)]
\label{thm:V34-second-close}
Alternative \textup{(i)} of
\Cref{thm:V34-trichotomy} satisfies \(DXQTRC_5\) for \(SU(2)\) and
every \(s>0\), whether \(f\) is unmarked, a recursive joining
coordinate, or the payer coordinate.
\end{theorem}

\begin{proof}
Expand the at most four joining cumulants by
\Cref{lem:V31-fixed-expansion}.  In each positive moment word, the
complete local factor at \(f\) is a moment-partition heat factor,
possibly with the sole payer or selected output retained.  V30 gives
its uniform response
\[
 {C_s\over d_{p_f}}
 \le C_{s,\sigma}M^{-1/2}.
\]
The factor at \(e_1\) gives the same estimate by
\eqref{eq:V34-first-slot}.  Thus
\Cref{thm:V33-positive-two} applies termwise.  The exact V94 dualizer
is removed by the V32 positive sandwich, and the fixed
row-partition coefficient sum is harmless.  Apply
\Cref{cor:V33-q2-closed}.
\end{proof}

\section{The diffuse exterior branch}
\label{sec:V34-diffuse}

\begin{lemma}[The \(k=1\) V87 proof supplies a PT majorant]
\label{lem:V34-V87-majorant}
Under \eqref{eq:V34-diffuse}, for all sufficiently large \(M\) and
every positive moment word there is a positive coordinate tensor
\(\mathbb B_{D_M}\) such that
\begin{equation}
 0\preccurlyeq\mathbb H_{D_M}
 \preccurlyeq\mathbb B_{D_M},
 \qquad
 \|\mathbb B_{D_M}^{\Gamma_{a^c}}\|
 \le C_{s,\sigma}M^{-1/2}.
\label{eq:V34-bank-majorant}
\end{equation}
\end{lemma}

\begin{proof}
The atomwise limit in \eqref{eq:V34-diffuse} eventually gives
\[
 \max_{e\in D_M}a_{a,e}\le{\sigma\over2}M.
\]
Apply the construction in the proof of the archived f3 theorem
\texttt{thm:V87-diffuse-bank-power} with \(k=1\).  At the largest bank
coordinate it chooses the smaller of the complete heat
partial-transpose bound and the strict protected/gapped majorant; at
all remaining coordinates it uses the latter.  Their tensor product
is a positive majorant of the complete moment bank.  The
\(k\)-largest-coordinate optimization gives exactly
\[
 C_{s,\sigma}M^{-1/2}
\]
for its full-row partial-transpose norm.  This is
\eqref{eq:V34-bank-majorant}.
\end{proof}

\begin{theorem}[A diffuse exterior half-power closes \(q=1\)]
\label{thm:V34-diffuse-close}
Alternative \textup{(ii)} of
\Cref{thm:V34-trichotomy} satisfies \(DXQTRC_5\) for \(SU(2)\) and
every \(s>0\).
\end{theorem}

\begin{proof}
Fix a positive word in
\eqref{eq:V31-fixed-expansion}.  Use
\Cref{lem:V34-V87-majorant} on \(D_M\).  At \(e_1\), V30 gives
\[
 \|H_{e_1}^{\Gamma_{a^c}}\|
 \le C_{s,\tau}M^{-1/2}.
\]
The bank is disjoint from \(e_1\) and the payer.  All remaining
unpaid exterior factors are positive contractions.  The payer factor
is bounded by \(C_sI\) if it lies outside \(e_1\), and is included in
the V30 local response otherwise.  Positive order therefore gives
\[
 \mathbb H_\nu
 \preccurlyeq
 C_s\,
 H_{e_1}\otimes\mathbb B_{D_M}\otimes I.
\]
Partial transpose and tensorization yield
\[
 \kappa_\otimes(\mathbb H_\nu)
 \le C_{s,\tau,\sigma}M^{-1}.
\]
Apply the exact-dual positivity sandwich, sum the fixed
row-partition coefficients, normalize by \(\prod_i\Xi_i\), and compare
with the one-strong floor
\eqref{eq:V33-exchange-floor}.
\end{proof}

\begin{corollary}[Only single-coordinate saturation survives in
\(q=1\)]
\label{cor:V34-q1-residual}
Every persistent one-strong \(q=1\) failure has a subsequence obeying
\eqref{eq:V34-saturation}.
\end{corollary}

\begin{proof}
The trichotomy is exhaustive.  Its first two alternatives close by
\Cref{thm:V34-second-close,thm:V34-diffuse-close}.
\end{proof}

\section{Saturation forces the strong partner onto the same
coordinate}
\label{sec:V34-strong-collapse}

\begin{lemma}[Overlap concentration]
\label{lem:V34-overlap-concentration}
Suppose \eqref{eq:V34-saturation} holds and \(c\ne a\) is a
maximal-scale strong partner:
\begin{equation}
 \theta_{ac}\ge\eta,
 \qquad
 \Xi_c\ge m_0M.
\label{eq:V34-strong-partner}
\end{equation}
Then, for all sufficiently large \(M\),
\begin{equation}
 {a_{c,e_1}\over\Xi_c}\ge{\eta\over2},
 \qquad
 a_{c,e_1}\ge{\eta m_0\over2}M.
\label{eq:V34-partner-at-e1}
\end{equation}
\end{lemma}

\begin{proof}
Use the probability weights
\[
 w_i(e):={a_{i,e}\over\Xi_i}.
\]
The exact overlap identity gives
\[
 \theta_{ac}
 =
 w_a(e_1)w_c(e_1)
 +\sum_{e\ne e_1}w_a(e)w_c(e).
\]
Since \(0\le w_c(e)\le1\),
\[
 \sum_{e\ne e_1}w_a(e)w_c(e)
 \le1-w_a(e_1)\longrightarrow0.
\]
For large \(M\), this exterior sum is at most \(\eta/2\).  The strong
lower bound then gives
\[
 w_a(e_1)w_c(e_1)\ge\eta/2.
\]
Since \(w_a(e_1)\le1\), the first inequality in
\eqref{eq:V34-partner-at-e1} follows.  Multiply by
\(\Xi_c\ge m_0M\) for the second.
\end{proof}

\begin{assessment}[The \(q=1\) terminal is really a coincident
strong-coordinate core]
\label{ass:V34-terminal}
In the only surviving \(q=1\) branch, row \(a\) and its strong partner
are both cofinal at \(e_1\), and every other row-\(a\) coordinate
carries \(o(M)\) mass.  Thus the global strong stock, the sole ordinary
heat response, and any selected or payer role occurring at \(e_1\)
must be treated as one complete local packet.  Reassigning the strong
coordinate to \(e_1\) moves the bookkeeping to the \(q=0\)
role-owned core; it does not manufacture a second response.
\end{assessment}

\begin{firewall}[One coordinate is still only one response]
\label{fire:V34-one-coordinate}
The complete local heat response at \(e_1\) is
\(O(M^{-1/2})\) and is sharp on the equal-spin singlet family.  The
direct one-strong target still requires \(O(M^{-1})\).  The scalar fact
that the same coordinate realizes the strong overlap is used on the
right side of the exchange comparison and cannot be counted as a
second analytic factor on the left.  V34 therefore reduces the
problem; it does not close the saturated packet.
\end{firewall}

\begin{assessment}[Exact residual after V34]
\label{ass:V34-residual}
All \(q=1\) fibres with exterior mass close.  The remaining
\(q=0,1\) obstruction is one physical coordinate \(e_1\) carrying
asymptotically all maximal-row mass and a cofinal strong partner on
that same coordinate.  Its exact complete local moment/cumulant,
selected output, payer, and final V94 dual action must be compared with
the local contribution of the strong edge.  This finite local packet,
not a diffuse bank or an unclassified exterior operator, is the next
object.
\end{assessment}

\section{V34 result ledger and V35 acquisition order}
\label{sec:V34-ledger}

\begin{theorem}[Fourth-series V34 result ledger]
\label{thm:V34-results}
Relative to V1--V33, V34 proves:
\begin{enumerate}[label=\textup{(V34.\arabic*)},leftmargin=6.3em]
\item the exact five-coordinate bound for recursive joining and payer
      roles outside the first response coordinate;
\item an exhaustive second-atom/diffuse-bank/single-atom trichotomy;
\item closure of every second cofinal atom, including marked role
      coordinates;
\item extraction of the \(k=1\) partial-transpose majorant from the
      archived V87 diffuse-bank proof;
\item closure of every \(q=1\) diffuse exterior branch;
\item reduction of every persistent \(q=1\) failure to
      single-coordinate saturation; and
\item concentration of every maximal-scale strong partner at that same
      physical coordinate.
\end{enumerate}
\end{theorem}

\begin{proof}
These are
\Cref{eq:V34-marked-count,
thm:V34-trichotomy,
thm:V34-second-close,
lem:V34-V87-majorant,
thm:V34-diffuse-close,
cor:V34-q1-residual,
lem:V34-overlap-concentration}.
\end{proof}

\begin{programme}[V35: companion audit and the coincident local
terminal]
\label{prog:V34-V35}
\begin{enumerate}[leftmargin=3em]
\item Perform the compulsory read-only audit of the newest active SM
      and NS notes before selecting the proof direction.
\item Fix the saturated coordinate \(e_1\), its strong partner \(c\),
      the exact local row partition, selected output, and payer status.
\item Separate the cases in which row \(a\) is a singleton moment
      block, belongs to a two-row block, or belongs to a
      higher-incidence block.
\item Use the exponential singleton response immediately.  In the
      nonsingleton cases retain the complete local cumulant and exact
      V94 positive sandwich.
\item Compare the local strong-edge contribution before replacing it
      by the generic one-strong forest floor; the concentrated support
      may improve that scalar floor.
\item Either prove the missing direct \(M^{-1/2}\) response or exhibit
      the exact low-output packet which survives all complete local
      sums.
\end{enumerate}
\end{programme}

\begin{firewall}[Claim boundary after fourth-series V34]
\label{fire:V34-claim-boundary}
V34 closes all \(q=1\) branches with nonnegligible exterior row mass
and reduces every survivor to a coincident single-coordinate strong
packet.  It does not close that packet, all \(q=0\), global quintic
MTA, all-order MTA, WUFC, CPM, cutoff removal, reflection positivity,
Osterwalder--Schrader reconstruction, construction of the continuum
Yang--Mills measure, or a uniform positive continuum spectral gap.
The full Yang--Mills mass gap has not been proved.
\end{firewall}

\paragraph{Parent-version record.}
V34 was formed from
\texttt{Renewal\_Yang\_Mills\_Mass\_Gap\_Research\_Notes\_V33.tex}.
The V33 parent had \(16\,306\) logical source lines,
\(572\,103\) bytes, and SHA--256
\[
 \texttt{efad39097a75e6186ada84f5ae52ce7e8207429ac7c460328199434af3852446}.
\]
The next compulsory companion audit is V35.

% =============================================================================
% END APPENDED FOURTH-SERIES V34 MATERIAL
% =============================================================================

% The V34 document terminator is intentionally disabled here:
% \end{document}

% =============================================================================
% BEGIN APPENDED FOURTH-SERIES V35 MATERIAL
% =============================================================================

\clearpage
\chapter{V35: Companion audit and the balanced low-output
single-coordinate terminal}
\label{ch:V35}

\section{Compulsory companion audit}
\label{sec:V35-audit}

The V35 read-only checkpoint was taken while the companion programmes
were running in parallel.  The newest complete active snapshots read
were:
\begin{center}
\begin{tabular}{@{}p{0.60\linewidth}rr@{}}
\toprule
file & logical lines & bytes\\
\midrule
\texttt{Renewal\_SM\_Einstein\_Continuum\_Research\_Notes\_V20.tex}
 & \(7\,166\) & \(246\,890\)\\
\texttt{Renewal\_NS\_Stability\_Research\_Notes\_V31.tex}
 & \(23\,052\) & \(887\,537\)\\
\bottomrule
\end{tabular}
\end{center}
Their SHA--256 digests were
\[
\begin{split}
&\texttt{485685878eac8aea4497179101075f7db3ab5e40bd51db923ce181ee47307f0b},\\
&\texttt{f1121b62238ad5491bb7cf03635ea9934942edf573ed8faaa9ee79e1d38a6696}.
\end{split}
\]
Neither companion file was modified.

SM V20 constructs negative Hamiltonian-energy spaces for time-zero
sources which are forms but not Hilbert vectors.  It computes exact
loss thresholds and proves a conditional constraint contraction only
when the nonlinear maps close at that same loss.  A weaker bounded
form norm cannot be cited as a stronger response, and an uncompensated
loss prevents a fixed-space iteration.

The type-correct Yang--Mills lesson concerns the response ledger, not
Hamiltonian forms.  The sharp \(M^{-1/2}\) response of one protected
coordinate cannot be relabelled as the missing \(M^{-1}\) direct
response.  Any improvement must be an actual extra decay estimate for
the complete coincident packet.

NS V31 proves that flat annular marks reconstruct critical coarea with
an exact first radial moment.  Heat formation pays that moment only
from the scale-matched critical source; a weaker rung restores the
missing scale in its coefficient.  A balanced-helicity solution also
shows that signed hinge cancellation controls only its own
polarization sector.

The type-correct Yang--Mills lesson is again exact matching.  Heat
damping may close a high-output spectral tail, but the low-output
protected sector must retain the complete local moment--cumulant sum.
A cancellation from a different sector or an ordinary norm of the
same local factor cannot pay the missing half-power.

\begin{record}[Audit-adjusted V35 direction]
\label{rec:V35-direction}
At the saturated coordinate, first remove singleton and high-output
terms by their actual heat scale.  Then isolate the balanced
low-output spectral window of the complete local packet.  Do not claim
closure from a norm living at the wrong response exponent.
\end{record}

\section{The saturated one-partner geometry}
\label{sec:V35-geometry}

Retain the terminal sequence of V34.  Thus
\begin{equation}
 {a_{a,e_*}\over M}\longrightarrow1,
 \qquad
 \Xi_a=M,
\label{eq:V35-a-saturation}
\end{equation}
and a unique maximal-scale partner \(c\) satisfies
\begin{equation}
 a_{c,e_*}\ge\lambda M
\label{eq:V35-c-cofinal}
\end{equation}
for a fixed \(\lambda>0\).  Write \(p\) for the spin of row \(a\) at
\(e_*\).  Then
\begin{equation}
 d_p\asymp M^{1/2}.
\label{eq:V35-dp}
\end{equation}

\begin{lemma}[Every other local row is subcofinal]
\label{lem:V35-other-small}
On a persistent one-partner subsequence,
\begin{equation}
 {a_{b,e_*}\over M}\longrightarrow0
 \qquad(b\notin\{a,c\}).
\label{eq:V35-other-small}
\end{equation}
Equivalently, if row \(b\) carries spin \(r_b\) at \(e_*\), then
\[
 r_b=o(p).
\]
\end{lemma}

\begin{proof}
If \eqref{eq:V35-other-small} failed, some \(b\ne a,c\) would have
\(a_{b,e_*}\ge\epsilon M\) on a subsequence.  Together with
\eqref{eq:V35-a-saturation}, this gives
\[
 \theta_{ab}
 \ge
 {a_{a,e_*}a_{b,e_*}\over\Xi_a\Xi_b}
 \ge \epsilon(1-o(1)){M\over\Xi_b}.
\]
Since \(\Xi_b\le M\), row \(b\) would be a second comparable strong
partner.  That leaves the one-partner terminal and enters the already
closed multi-partner branch of f3 V99.  Hence every persistent
one-partner subsequence satisfies \eqref{eq:V35-other-small}.
For \(SU(2)\), \(a_{b,e_*}=1+r_b(r_b+1)\) and
\(a_{a,e_*}=1+p(p+1)\asymp M\), proving \(r_b=o(p)\).
\end{proof}

\section{Singleton and partner-excluding moment blocks}
\label{sec:V35-blocks}

Fix one positive row-partition word in the expansion
\eqref{eq:V31-fixed-expansion}, and let \(B_a\) be the local partition
block containing row \(a\) at \(e_*\).

\begin{lemma}[A saturated singleton closes with arbitrary power]
\label{lem:V35-singleton}
If \(B_a=\{a\}\), then for every \(N\ge1\) the complete dualized word
has an additional factor \(O_{s,N}(M^{-N})\).  The same holds if the
unique payer is allocated at \(e_*\).
\end{lemma}

\begin{proof}
The row-\(a\) factor is the scalar
\[
 e^{-sp(p+1)}I_{V_p}.
\]
By \eqref{eq:V35-a-saturation},
\(p(p+1)\ge M/2\) for all sufficiently large \(M\).  The payer grows
at most polynomially in the local Casimir, while every other
fixed-order factor has its established representation-uniform norm.
Polynomial times \(e^{-sM/2}\) is \(O(M^{-N})\) for every \(N\).
The scalar commutes through the exact V94 dualizer and every other
coordinate.
\end{proof}

\begin{lemma}[A block which omits the unique partner is high-output]
\label{lem:V35-omit-c}
Suppose \(B_a\ne\{a\}\) but \(c\notin B_a\).  Decompose
\[
 \bigotimes_{b\in B_a\setminus\{a\}}V_{r_b}
 =
 \bigoplus_q\mathcal N_q\otimes V_q.
\]
Then, uniformly over every occurring \(q\),
\begin{equation}
 q=o(p),
 \qquad
 K\ge|p-q|\ge {p\over2}
\label{eq:V35-high-from-omit}
\end{equation}
for every total output \(K\) and all sufficiently large \(M\).
Consequently these terms satisfy the direct \(O(M^{-1})\) target.
\end{lemma}

\begin{proof}
There are at most three rows in
\(B_a\setminus\{a\}\), and
\Cref{lem:V35-other-small} gives \(r_b=o(p)\) for each.  Every
irreducible output of their tensor product satisfies
\[
 q\le\sum_{b\in B_a\setminus\{a\}}r_b=o(p).
\]
Thus \(q\le p/2\) eventually, and the Clebsch--Gordan triangle
inequality gives \eqref{eq:V35-high-from-omit}.

The local heat multiplier is at most
\[
 e^{-sK(K+1)}\le e^{-sp^2/4}.
\]
The sole payer and every fixed-arity output sum cost only a polynomial
in \(p\).  Hence the local factor is \(O(p^{-N})\) for every \(N\).
Take \(N=4\), use \(p^2\asymp M\), and apply the V32 exact-dual
positive sandwich and exterior erasure.
\end{proof}

\begin{corollary}[Every surviving moment word couples \(a\) to \(c\)
locally]
\label{cor:V35-ac-block}
In the saturated one-partner terminal, every positive
row-partition term not already satisfying the direct target has
\[
 a,c\in B_a.
\]
\end{corollary}

\begin{proof}
The singleton case is
\Cref{lem:V35-singleton}; the nonsingleton partner-excluding case is
\Cref{lem:V35-omit-c}.
\end{proof}

\section{High total output has the missing direct power}
\label{sec:V35-high-output}

\begin{lemma}[Gaussian tail with a polynomial payer]
\label{lem:V35-Gaussian-tail}
For every \(s>0\) and integer \(m\ge0\), there is \(C_{s,m}\) such
that, for \(R\ge1\),
\begin{equation}
 \sum_{\substack{K\in\widehat{SU(2)}\\c_K\ge R}}
 d_K^m e^{-(s/2)c_K}
 \le C_{s,m}e^{-(s/4)R}.
\label{eq:V35-Gaussian-tail}
\end{equation}
\end{lemma}

\begin{proof}
The number of spins in each interval
\(n\le K<n+1\) is bounded, \(d_K^m\) is polynomial in \(n\), and
\(c_K\ge n^2\).  Absorb the polynomial into
\(e^{(s/4)c_K}\) outside a fixed compact set and sum the remaining
Gaussian series.  Enlarging the constant handles the compact part.
\end{proof}

\begin{definition}[Low-output window]
\label{def:V35-low-window}
Put
\begin{equation}
 R_p(s):={8\over s}\log(2+d_p)
\label{eq:V35-Rp}
\end{equation}
and let \(P_{\rm lo}\) be the sum of local total-output projections
with \(c_K<R_p(s)\).  Put \(P_{\rm hi}=I-P_{\rm lo}\).
\end{definition}

\begin{theorem}[The high-output packet closes]
\label{thm:V35-high-close}
For every positive moment-partition term at \(e_*\), including its
selected projection and the unique payer when allocated there,
\begin{equation}
 \left\|
  (P_{\rm hi}H_{e_*}^{\rm pay})^{\Gamma_{a^c}}
 \right\|
 \le {C_s\over d_p^2}.
\label{eq:V35-high-close}
\end{equation}
The same estimate holds without the payer.  Hence the complete
high-output part satisfies the missing direct \(O(M^{-1})\) response.
\end{theorem}

\begin{proof}
Decompose the complement of \(V_p\) in \(B_a\) into
\(\bigoplus_q\mathcal N_q\otimes V_q\).  On a fixed \(q,K\)-block,
the V29 scalar-marginal argument gives
\[
 \|(T_{p,q;K}D_KT_{p,q;K}^*)^\Gamma\|
 \le {d_K^3\over d_p}\|D_K\|.
\]
Multiplicity spaces lie wholly on the complementary side and cause
no loss.  In the unpaid heat word, sum this estimate with
\(e^{-sc_K}\).  In the paid word, first use the V30 payer absorption,
which replaces \(s\) by \(s/2\) and introduces a fixed constant.
Thus both cases are bounded by
\[
 {C_s\over d_p}
 \sum_{c_K\ge R_p(s)}d_K^3e^{-(s/2)c_K}.
\]
By \Cref{lem:V35-Gaussian-tail} and
\eqref{eq:V35-Rp}, the sum is at most
\[
 C_s e^{-(s/4)R_p(s)}
 =
 C_s(2+d_p)^{-2}.
\]
This proves a bound \(C_sd_p^{-3}\), which is stronger than
\eqref{eq:V35-high-close}.  Selected output projection only restricts
the sum.  Since \(d_p^2\asymp M\), the final assertion follows from
the V32 positive sandwich and exterior erasure.
\end{proof}

\section{Exact form of the surviving low-output packet}
\label{sec:V35-low-output}

\begin{theorem}[Balanced logarithmic output reduction]
\label{thm:V35-low-reduction}
Every saturated one-partner term not closed by
\Cref{lem:V35-singleton,lem:V35-omit-c,thm:V35-high-close} has:
\begin{enumerate}[label=\textup{(\roman*)},leftmargin=4em]
\item \(a,c\in B_a\);
\item a complementary resultant \(q\) and total output \(K\) satisfying
\begin{equation}
 |p-q|\le K,
 \qquad
 c_K<{8\over s}\log(2+d_p);
\label{eq:V35-low-constraints}
\end{equation}
\item consequently,
\begin{equation}
 K=O_s(\sqrt{\log d_p}),
 \qquad
 |p-q|=O_s(\sqrt{\log d_p});
\label{eq:V35-near-match}
\end{equation}
and
\item the general paid partial-transpose bound
\begin{equation}
 \left\|
  (P_{\rm lo}H_{e_*}^{\rm pay})^\Gamma
 \right\|
 \le
 {C_s(1+\log(2+d_p))^3\over d_p},
\label{eq:V35-low-upper}
\end{equation}
with exponent \(2\) in place of \(3\) in the unpaid case.
\end{enumerate}
\end{theorem}

\begin{proof}
Item \textup{(i)} is
\Cref{cor:V35-ac-block}.  The low projection gives the second
inequality in \eqref{eq:V35-low-constraints}; the first is the
Clebsch--Gordan triangle inequality.  Since
\(c_K=K(K+1)\), these imply
\eqref{eq:V35-near-match}.

For the unpaid term, the V29 fixed-output estimate and triangle
inequality give
\[
 {1\over d_p}
 \sum_{c_K<R_p(s)}d_K^3
 \le
 {C_s(1+R_p(s))^2\over d_p}.
\]
Indeed the largest allowed \(K\) is \(O(\sqrt{R_p})\), and the sum of
\((2K+1)^3\) up to that value is \(O(R_p^2)\).  A payer contributes
at most \(C_s(1+c_K)\) on this window, increasing the sum to
\(O(R_p^3)\).  Insert \eqref{eq:V35-Rp}.
\end{proof}

\begin{proposition}[The protected endpoint is still sharp]
\label{prop:V35-protected-sharp}
The low window contains \(K=0\).  For \(p=q\), the corresponding
equal-spin invariant projection has product-state capacity
\begin{equation}
 {1\over d_p}\asymp M^{-1/2}.
\label{eq:V35-protected-sharp}
\end{equation}
Therefore no estimate using only the positive low-output moment word
can supply the required \(d_p^{-2}\asymp M^{-1}\) response uniformly.
\end{proposition}

\begin{proof}
The invariant vector in \(V_p\otimes V_p\) is maximally entangled, so
the largest squared overlap with a product vector is \(d_p^{-1}\).
Its heat weight at \(K=0\) is one.  This is the sharp singlet test
already established in f3 V91 and used in V13.  Since
\(d_p^{-1}/d_p^{-2}=d_p\to\infty\), the stronger positive-word bound
is impossible.
\end{proof}

\begin{firewall}[Where cancellation is still allowed]
\label{fire:V35-cancellation}
\Cref{prop:V35-protected-sharp} concerns one positive
moment-partition term.  The complete local joining or prefix cumulant
is a signed sum of such terms.  Cancellation of its common
low-output protected component is not excluded, but it must be
computed before the fixed row-partition triangle inequality in
\eqref{eq:V31-fixed-expansion}.  Neither the SM negative-energy form
nor the NS signed hinge supplies that cancellation.
\end{firewall}

\begin{assessment}[Exact residual after V35]
\label{ass:V35-residual}
The single-coordinate terminal is now reduced to a balanced
low-output packet.  Row \(a\) and its unique strong partner \(c\) lie
in the same local moment block; their effective spins match to within
\(O_s(\sqrt{\log d_p})\); and the total output has the same logarithmic
height.  Singleton, partner-excluding, and high-output terms already
have the full direct \(M^{-1}\) response.  The remaining question is
whether the \emph{complete signed local moment--cumulant symbol}
annihilates or suppresses the \(K=O(\sqrt{\log d_p})\) protected
window by one additional factor \(d_p^{-1}\).
\end{assessment}

\section{V35 result ledger and V36 acquisition order}
\label{sec:V35-ledger}

\begin{theorem}[Fourth-series V35 result ledger]
\label{thm:V35-results}
Relative to V1--V34, V35 proves:
\begin{enumerate}[label=\textup{(V35.\arabic*)},leftmargin=6.3em]
\item the compulsory SM V20 and NS V31 companion audit;
\item subcofinality of every row other than the unique strong partner
      at the saturated coordinate;
\item arbitrary-power closure of saturated singleton moment terms;
\item arbitrary-power closure of every local block which omits the
      unique partner;
\item a Gaussian output-tail estimate retaining a polynomial payer;
\item the full \(d_p^{-2}\) direct response on the high-output packet;
\item reduction of every survivor to
      \(K,|p-q|=O_s(\sqrt{\log d_p})\);
\item an explicit polylogarithmic-over-\(d_p\) upper bound on that
      low window; and
\item the sharp \(d_p^{-1}\) protected lower test which locates the
      missing factor inside the signed local cumulant.
\end{enumerate}
\end{theorem}

\begin{proof}
These are
\Cref{rec:V35-direction,
lem:V35-other-small,
lem:V35-singleton,
lem:V35-omit-c,
lem:V35-Gaussian-tail,
thm:V35-high-close,
thm:V35-low-reduction,
prop:V35-protected-sharp}.
\end{proof}

\begin{programme}[V36: exact protected symbol of the coincident
cumulant]
\label{prog:V35-V36}
\begin{enumerate}[leftmargin=3em]
\item Fix each possible recursive node containing \(a,c\) and write
      the complete local moment--cumulant polynomial at \(e_*\) before
      taking absolute values.
\item Compress that polynomial to a fixed
      \((q,K)\) block in the V35 logarithmic window, retaining all
      other row outputs and multiplicities.
\item Compute the coefficient of the equal-spin \(K=0\) projection.
      Determine whether M\"obius cancellation removes it for a joining
      cumulant, a prefix cumulant, or neither.
\item If the \(K=0\) coefficient vanishes, bound the first surviving
      finite difference in \(K\) or in the independently heated
      endpoint and test for one factor \(d_p^{-1}\).
\item If it does not vanish, construct the complete physical
      single-coordinate counterpacket including the prefix, suffix,
      exact V94 duality, and exchange weight; do not infer failure from
      the isolated positive projection.
\item Propagate every closed local node through the \(q=0,1\) census.
\end{enumerate}
\end{programme}

\begin{firewall}[Claim boundary after fourth-series V35]
\label{fire:V35-claim-boundary}
V35 closes every part of the saturated single-coordinate terminal
except the balanced logarithmic-output component of the complete
signed local cumulant.  It does not prove that component has the
missing \(d_p^{-1}\) factor, close all \(q=0,1\), prove global quintic
MTA, all-order MTA, WUFC, CPM, cutoff removal, reflection positivity,
Osterwalder--Schrader reconstruction, construction of the continuum
Yang--Mills measure, or a uniform positive continuum spectral gap.
The full Yang--Mills mass gap has not been proved.
\end{firewall}

\paragraph{Parent-version record.}
V35 was formed from
\texttt{Renewal\_Yang\_Mills\_Mass\_Gap\_Research\_Notes\_V34.tex}.
The V34 parent had \(16\,702\) logical source lines,
\(585\,267\) bytes, and SHA--256
\[
 \texttt{fa0fa4c753f3293d24fd84a02f8853310fb7b6279c73c94a380c2712e114afd6}.
\]
The next compulsory companion audit is V40.

% =============================================================================
% END APPENDED FOURTH-SERIES V35 MATERIAL
% =============================================================================

% \end{document} % disabled by fourth-series V36 append

% =============================================================================
% BEGIN APPENDED FOURTH-SERIES V36 MATERIAL
% =============================================================================

\chapter{V36: High-pair M\"obius contraction and the absence of a
generic protected surplus}
\label{chap:V36-high-pair-contraction}

\section{Purpose and correction of the V35 acquisition order}
\label{sec:V36-purpose}

V35 isolated the only packet not controlled by its direct estimates:
at one saturated local coordinate there is a unique cofinal pair
\((a,c)\), its spin labels \(p,q\) satisfy
\[
 |p-q|=O_s\!\left(\sqrt{\log d_p}\right),
\]
and the coupled output \(L\) lies in the same logarithmic window.
The positive protected projection has only the ordinary
\(d_p^{-1}\) response.  The V35 programme therefore proposed looking
for another factor \(d_p^{-1}\) in the local M\"obius polynomial.

That proposal requires a location correction.  A saturated coordinate
need not carry a cumulant.  If it is an unmarked prefix or suffix
coordinate, its source factor is one moment and there is no local
M\"obius sum at all.  If it is a joining coordinate, the part of the
M\"obius sum which keeps the unique high pair together does not
generically cancel: it is exactly a cumulant with one fewer input.
The complementary part, which separates the high pair, is
arbitrary-power small.  Thus local centering gives an
\emph{arity descent}, not a new representation-dimension power.

This chapter proves that assertion before any absolute value is taken.
It supersedes items 1--4 of
\Cref{prog:V35-V36}; it does not alter any estimate proved in V35.

\section{Partition contraction}
\label{sec:V36-partition-contraction}

\begin{definition}[Moment and cumulant polynomials]
\label{def:V36-cumulant}
Let \(I_m=\{1,\ldots,m\}\), let \(\Pi(I_m)\) be its partition lattice,
and let \(\widehat 1_m\) be the one-block partition.  For every
nonempty \(B\subset I_m\), let \(H_B\) denote the physical heat-moment
operator on the tensor factors indexed by \(B\).  For
\(\pi\in\Pi(I_m)\), put
\[
 M_\pi:=\prod_{B\in\pi}H_B.
\]
The factors in this product have disjoint supports.  The complete
local cumulant is
\[
 {\mathcal K}_m
 :=
 \sum_{\pi\in\Pi(I_m)}
 \mu_m(\pi,\widehat 1_m)M_\pi,
 \qquad
 \mu_m(\pi,\widehat 1_m)
 =
 (-1)^{|\pi|-1}(|\pi|-1)! .
\]
\end{definition}

The formula is the same one used in V12.  It is recalled here because
the dependence of the coefficient only on the number of blocks is the
entire combinatorial mechanism.

\begin{definition}[Together and separated partition sectors]
\label{def:V36-sectors}
For two distinguished labels \(1,2\), define
\[
 \Pi_m^{\mathrm{tog}}
 :=
 \{\pi\in\Pi(I_m):1\sim_\pi2\},
 \qquad
 \Pi_m^{\mathrm{sep}}
 :=
 \{\pi\in\Pi(I_m):1\not\sim_\pi2\}.
\]
Let
\[
 \overline I_{m-1}:=\{\star,3,\ldots,m\}.
\]
For \(\pi\in\Pi_m^{\mathrm{tog}}\), let \(c_{12}\pi\) be the partition
of \(\overline I_{m-1}\) obtained by replacing \(1,2\) in their common
block by the single label \(\star\).
\end{definition}

\begin{lemma}[M\"obius-preserving pair contraction]
\label{lem:V36-Mobius-contraction}
The map
\[
 c_{12}:\Pi_m^{\mathrm{tog}}\longrightarrow
 \Pi(\overline I_{m-1})
\]
is a bijection, preserves the number of blocks, and satisfies
\[
 \mu_m(\pi,\widehat 1_m)
 =
 \mu_{m-1}(c_{12}\pi,\widehat 1_{m-1}).
\]
\end{lemma}

\begin{proof}
Expansion of the label \(\star\) into the two labels \(1,2\) in the
same block is an inverse to \(c_{12}\).  Contraction changes neither
that block nor any other block except for its list of labels, so
\[
 |c_{12}\pi|=|\pi|.
\]
The displayed equality follows from
\[
 \mu_r(\sigma,\widehat 1_r)
 =(-1)^{|\sigma|-1}(|\sigma|-1)!,
\]
which depends only on \(|\sigma|\), not on \(r\).
\end{proof}

\section{The physical high-pair compression}
\label{sec:V36-physical-compression}

Fix one local joining coordinate and one quotient fibre.  Relabel the
unique cofinal rows \(a,c\) as \(1,2\), with spins \(p,q\).
All remaining local row spins will be denoted \(r_3,\ldots,r_m\).
The V35 unique-partner lemma gives
\begin{equation}
\label{eq:V36-subcofinal}
 R_p:=\sum_{\nu=3}^m r_\nu=o(p),
 \qquad
 q/p\longrightarrow1
\end{equation}
along every persistent saturated sequence.  The number \(m\) is
bounded by the fixed perturbative row order.

Resolve \(V_p\otimes V_q\) into irreducibles and let \(Q_{p,q;L}\)
denote the orthogonal projection onto an intermediate spin \(L\).
The resolution is auxiliary recoupling data.  It is not a new source
map, and it is summed when the fibre is reconstructed.  Let \(P_p\)
be any further physical output compression contained in the V35 low
window
\begin{equation}
\label{eq:V36-low-window}
 L,\ |p-q|\le A_s\sqrt{\log d_p}.
\end{equation}
We take \(P_p\le Q_{p,q;L}\), including all multiplicity identities
required by exact V94 duality.

\begin{definition}[Contracted moment system]
\label{def:V36-contracted-system}
On \(Q_{p,q;L}(V_p\otimes V_q)\), replace the pair of inputs \(1,2\)
by one composite input \(\star\) carrying \(V_L\).  If a block
\(\overline B\subset\overline I_{m-1}\) contains \(\star\), define
\(\overline H_{\overline B}\) by coupling \(V_L\) to the factors
indexed by \(\overline B\setminus\{\star\}\) and applying the same
physical heat multiplier as in the original moment.  If
\(\star\notin\overline B\), put
\(\overline H_{\overline B}=H_{\overline B}\).
Set
\[
 \overline M_\sigma
 :=\prod_{\overline B\in\sigma}\overline H_{\overline B},
 \qquad
 \overline{\mathcal K}_{m-1}^{(L)}
 :=
 \sum_{\sigma\in\Pi(\overline I_{m-1})}
 \mu_{m-1}(\sigma,\widehat 1_{m-1})
 \overline M_\sigma .
\]
\end{definition}

\begin{lemma}[Exact equality in the together sector]
\label{lem:V36-together-exact}
After the fixed-\(L\) recoupling identification,
\[
 P_p M_\pi P_p
 =
 P_p\overline M_{c_{12}\pi}P_p
 \qquad
 (\pi\in\Pi_m^{\mathrm{tog}}).
\]
Consequently,
\begin{equation}
\label{eq:V36-together-cumulant}
 P_p
 \sum_{\pi\in\Pi_m^{\mathrm{tog}}}
 \mu_m(\pi,\widehat 1_m)M_\pi
 P_p
 =
 P_p\overline{\mathcal K}_{m-1}^{(L)}P_p.
\end{equation}
\end{lemma}

\begin{proof}
Let \(B_\pi\) be the unique block of \(\pi\) containing \(1,2\).
Associativity of the tensor product and the unitary Racah recoupling
map allow the first two factors to be coupled first.  On
\(Q_{p,q;L}\), their tensor product is then the composite irreducible
\(V_L\), with its multiplicity identity retained.  The diagonal heat
functional calculus of the total Casimir on \(B_\pi\) is therefore
exactly the functional calculus defining
\(\overline H_{c_{12}B_\pi}\).  Blocks other than \(B_\pi\) are
unchanged and have disjoint supports.  This proves the first equality.

Sum it over \(\Pi_m^{\mathrm{tog}}\).  By
\Cref{lem:V36-Mobius-contraction}, contraction is a bijection and
preserves every M\"obius coefficient.  The resulting sum is precisely
\(\overline{\mathcal K}_{m-1}^{(L)}\).
\end{proof}

\begin{lemma}[Separated sectors are arbitrary-power small]
\label{lem:V36-separated-small}
Assume \eqref{eq:V36-subcofinal}.  For every \(N\ge0\), there are
\(p_N\) and \(C_{s,m,N}\) such that, for \(p\ge p_N\),
\begin{equation}
\label{eq:V36-separated-bound}
 \left\|
 P_p\sum_{\pi\in\Pi_m^{\mathrm{sep}}}
 \mu_m(\pi,\widehat 1_m)M_\pi P_p
 \right\|
 \le C_{s,m,N}d_p^{-N}.
\end{equation}
The same assertion holds after the exact local duality and with any
one polynomial output payer allowed in the terminal estimate.
\end{lemma}

\begin{proof}
Fix \(\pi\in\Pi_m^{\mathrm{sep}}\) and let \(B_1\) be the block
containing label \(1\).  It does not contain label \(2\).  Hence the
sum of all spins in \(B_1\setminus\{1\}\) is at most \(R_p=o(p)\).
Every irreducible output \(J\) of that block obeys the reverse triangle
inequality
\[
 J\ge p-R_p.
\]
For large \(p\), this is at least \(p/2\).  The heat multiplier of this
block is consequently bounded by
\[
 \exp\{-sJ(J+1)\}
 \le \exp\{-s p^2/4\}.
\]
Every other heat factor is a contraction.  Since the number of
partitions and all M\"obius coefficients depend only on the fixed
integer \(m\), summation gives a constant multiple of this Gaussian.

An output payer and exact local partial transpose can contribute only
a fixed polynomial in the representation dimensions at fixed row
order.  Each \(SU(2)\) representation dimension is linear in its spin,
and all local spins are \(O(p)\) on the selected subsequence.
Thus the paid or dualized norm is bounded by
\(C_{s,m,D}p^D e^{-sp^2/4}\) for a fixed \(D\).  A Gaussian dominates
every power of \(d_p=2p+1\), proving \eqref{eq:V36-separated-bound}
for arbitrary \(N\).
\end{proof}

\begin{theorem}[High-pair cumulant contraction and arity descent]
\label{thm:V36-arity-descent}
Under the V35 saturated unique-partner hypotheses, on every fixed
low-window intermediate channel \(L\),
\begin{equation}
\label{eq:V36-arity-descent}
 P_p{\mathcal K}_mP_p
 =
 P_p\overline{\mathcal K}_{m-1}^{(L)}P_p+E_{p,L},
\end{equation}
where, for every \(N\ge0\),
\[
 \|E_{p,L}\|\le C_{s,m,N}d_p^{-N}.
\]
The remainder has the same arbitrary-power estimate in the paid
exact-duality response norm.
\end{theorem}

\begin{proof}
Split the defining partition sum for \({\mathcal K}_m\) into
\(\Pi_m^{\mathrm{tog}}\) and \(\Pi_m^{\mathrm{sep}}\).
\Cref{lem:V36-together-exact} identifies the first sum exactly with
the contracted cumulant.  \Cref{lem:V36-separated-small} supplies the
remainder estimate for the second sum.
\end{proof}

\begin{remark}[Uniformity over the logarithmic window]
\label{rem:V36-window-uniformity}
The proof uses only the subcofinal bound for the rows other than
\(1,2\); it does not spend any negative power of \(L\).  Its constants
are therefore uniform for all \(L\) satisfying
\eqref{eq:V36-low-window}.  Summing the
\(O_s((\log d_p)^{1/2})\) allowed values of \(L\), and their polynomial
multiplicities, preserves arbitrary-power smallness of \(E_{p,L}\).
\end{remark}

\section{The binary and ternary witnesses}
\label{sec:V36-witnesses}

\begin{proposition}[Binary descent ends at a first moment]
\label{prop:V36-binary}
For two inputs,
\[
 {\mathcal K}_2=H_{\{1,2\}}-H_{\{1\}}H_{\{2\}}.
\]
On the spin-\(L\) channel of \(V_p\otimes V_q\),
\begin{equation}
\label{eq:V36-binary-symbol}
 Q_{p,q;L}{\mathcal K}_2Q_{p,q;L}
 =
 \left(
 e^{-sL(L+1)}
 -
 e^{-sp(p+1)}e^{-sq(q+1)}
 \right)Q_{p,q;L}.
\end{equation}
If \(p,q\to\infty\), \(|p-q|=O_s(\sqrt{\log d_p})\), and \(L\) is
fixed, the coefficient tends to
\(e^{-sL(L+1)}\).  In particular it tends to \(1\) at \(L=0\).
\end{proposition}

\begin{proof}
The partition with \(1,2\) together contracts to the sole partition
of \(\{\star\}\), so its contracted cumulant is the first moment
\(e^{-sL(L+1)}Q_{p,q;L}\).  The separated partition is the product of
the two endpoint heat operators, giving the second term in
\eqref{eq:V36-binary-symbol}.  That term tends to zero faster than any
power of \(p\).
\end{proof}

\begin{corollary}[No binary protected cancellation]
\label{cor:V36-no-binary-cancellation}
On the equal-spin protected singlet,
\[
 Q_{p,p;0}{\mathcal K}_2Q_{p,p;0}
 =
 \left(1-e^{-2sp(p+1)}\right)Q_{p,p;0}.
\]
Thus the complete signed binary cumulant has a nonzero limiting
coefficient.  The exact protected partial-transpose response proved
in V13 remains of order \(d_p^{-1}\), not \(d_p^{-2}\).
\end{corollary}

\begin{proof}
Set \(q=p\) and \(L=0\) in
\eqref{eq:V36-binary-symbol}.  The response statement is then the
sharp V13 protected-block calculation applied to a coefficient
tending to one.
\end{proof}

\begin{proposition}[Ternary descent agrees with the V16 obstruction]
\label{prop:V36-ternary}
For three inputs,
\[
\begin{split}
 {\mathcal K}_3={}&H_{123}
 -H_{12}H_3-H_{13}H_2-H_{23}H_1
 +2H_1H_2H_3.
\end{split}
\]
If inputs \(1,2\) form the unique high pair, then on their fixed
low-spin channel \(L\),
\[
 P_p{\mathcal K}_3P_p
 =
 P_p\left(
 \overline H_{\star3}
 -
 \overline H_{\star}\overline H_3
 \right)P_p+O_N(d_p^{-N}).
\]
For the physical \((p,p,1)\) family and the total-spin-one packet used
in V16, this contracted covariance has the nonzero leading symbol
\[
 e^{-2s}\bigl(I-H_{12}\bigr),
\]
up to the exponentially small endpoint terms already bounded there.
\end{proposition}

\begin{proof}
The together partitions are
\(\{123\}\) and \(\{12\}\{3\}\), with coefficients \(1\) and \(-1\).
They contract to the two partitions of \(\{\star,3\}\), hence give the
displayed covariance.  The remaining three partitions separate
labels \(1,2\) and are arbitrary-power small by
\Cref{lem:V36-separated-small}.  The final formula is the exact V16
block calculation, now identified as the \(m=3\) instance of
\Cref{thm:V36-arity-descent}.
\end{proof}

\section{Location census at the saturated coordinate}
\label{sec:V36-location-census}

\begin{theorem}[Moment-or-descent dichotomy]
\label{thm:V36-location-dichotomy}
Let \(e_*\) be the single saturated coordinate left by the V34--V35
\(q=1\) reduction.
\begin{enumerate}[label=\textup{(\alph*)},leftmargin=3em]
\item If \(e_*\) is an unmarked prefix or suffix coordinate, the local
      source factor is one moment operator.  There is no M\"obius sum
      at \(e_*\), so the protected coefficient cannot be removed by
      local centering there.
\item If \(e_*\) is a joining coordinate carrying an \(m\)-input
      cumulant, its low-window symbol is, modulo arbitrary-power small
      error, the \((m-1)\)-input cumulant obtained by replacing the
      unique high pair by its coupled output.
\item Tree, lift, payer, duality, and strong-coordinate labels do not
      create additional moment partitions.  They can change routing,
      norms, or weights, but they are not new source maps.
\end{enumerate}
\end{theorem}

\begin{proof}
Part (a) is the prefix/suffix source grammar: an unmarked coordinate
appears through one heat moment, whereas the partition sum is attached
only to a joining cumulant.  Part (b) is
\Cref{thm:V36-arity-descent}.  For (c), the construction assigns these
labels after the physical moments and cumulants have been formed.
Exact duality acts on the resulting operator, and a payer is a norm
weight.  Neither operation alters the partition lattice in
\Cref{def:V36-cumulant}.  This is also why the auxiliary recoupling
projection in \Cref{def:V36-contracted-system} must not be counted as
a source.
\end{proof}

\begin{corollary}[Finite arity termination]
\label{cor:V36-finite-arity}
Repeated high-pair contraction at one joining coordinate terminates
after at most \(m-1\) steps at either:
\begin{enumerate}[label=\textup{(\roman*)},leftmargin=3em]
\item a lower-arity cumulant with a nonzero protected symbol;
\item an identically vanishing contracted cumulant; or
\item a first moment of the final composite input.
\end{enumerate}
Only case \textup{(ii)} produces cancellation merely from the local
M\"obius algebra.  In cases \textup{(i)} and \textup{(iii)}, a missing
representation-dimension power must come from outside this local
centering operation.
\end{corollary}

\begin{proof}
Each application of \Cref{thm:V36-arity-descent} lowers the number of
inputs by one.  Since \(m\) is finite, the process terminates.  At
arity one the cumulant is the first moment.  If termination occurs
earlier, the contracted cumulant is either zero or nonzero on the
selected channel.  The arbitrary-power errors cannot turn a fixed
nonzero limiting symbol into an additional algebraic
\(d_p^{-1}\) gain.
\end{proof}

\section{A no-go theorem for a universal local surplus}
\label{sec:V36-no-go}

\begin{theorem}[No universal protected M\"obius surplus]
\label{thm:V36-no-universal-surplus}
There is no estimate, valid for every physical saturated
unique-partner packet solely by virtue of local M\"obius centering, of
the form
\[
 \|\mathfrak R_p({\mathcal K}_m)\|
 \le C_s d_p^{-2},
\]
where \(\mathfrak R_p\) is the exact V94 protected local response.
More precisely, suppose a contracted lower-arity symbol admits
protected test functionals \(\Lambda_p\), of at most polynomial norm,
such that
\[
 \left|
 \Lambda_p\!\left(
 P_p\overline{\mathcal K}_{m-1}^{(L)}P_p
 \right)
 \right|
 \ge c_s d_p^{-1}.
\]
Then, for all sufficiently large \(p\),
\[
 \left|\Lambda_p(P_p{\mathcal K}_mP_p)\right|
 \ge \frac{c_s}{2}d_p^{-1}.
\]
\end{theorem}

\begin{proof}
By \Cref{thm:V36-arity-descent},
\[
\begin{split}
 \left|\Lambda_p(P_p{\mathcal K}_mP_p)\right|
 &\ge
 \left|
 \Lambda_p(P_p\overline{\mathcal K}_{m-1}^{(L)}P_p)
 \right|
 -
 \|\Lambda_p\|\,\|E_{p,L}\|.
\end{split}
\]
If \(\|\Lambda_p\|\le C d_p^D\), choose the arbitrary-power estimate
for \(E_{p,L}\) with \(N>D+2\).  The error is then
\(o(d_p^{-1})\), which gives the asserted lower bound.  The binary
family of \Cref{cor:V36-no-binary-cancellation} supplies an
unconditional physical witness; the ternary V16 family supplies a
second one.
\end{proof}

\begin{assessment}[Meaning of the no-go theorem]
\label{ass:V36-no-go-meaning}
\Cref{thm:V36-no-universal-surplus} is a local interface theorem, not
a counterexample to global quintic MTA.  It excludes only the proposed
generic inference that centering at the saturated coordinate gives
one additional factor \(d_p^{-1}\).
The full physical packet still contains exact exchange coefficients,
unequal row-mass factors, the normalization by the admissible label
volume, and aggregation over first-connectivity fibres.  Any of those
may supply the missing summability, but none may be replaced by a
fictitious local cancellation.
\end{assessment}

\section{V36 result ledger and V37 upstream programme}
\label{sec:V36-ledger}

\begin{theorem}[Fourth-series V36 result ledger]
\label{thm:V36-results}
Relative to V1--V35, V36 proves:
\begin{enumerate}[label=\textup{(V36.\arabic*)},leftmargin=6.3em]
\item a M\"obius-preserving bijection between partitions which keep
      the unique high pair together and partitions of one lower arity;
\item exact equality of the together sector with the contracted
      lower-arity physical cumulant on every fixed intermediate
      channel;
\item arbitrary-power smallness, including one polynomial payer and
      exact local duality, of every partition which separates the
      unique high pair;
\item the high-pair arity-descent identity
      \eqref{eq:V36-arity-descent};
\item its binary endpoint at a nonvanishing first moment and its
      agreement with the ternary V16 obstruction;
\item the moment-or-descent location dichotomy for the saturated
      coordinate; and
\item a no-go theorem excluding a universal additional
      \(d_p^{-1}\) merely from local M\"obius centering.
\end{enumerate}
\end{theorem}

\begin{proof}
These are respectively
\Cref{lem:V36-Mobius-contraction,
lem:V36-together-exact,
lem:V36-separated-small,
thm:V36-arity-descent,
prop:V36-binary,
prop:V36-ternary,
thm:V36-location-dichotomy,
thm:V36-no-universal-surplus}.
\end{proof}

\begin{programme}[V37: return to the exact exchange interface]
\label{prog:V36-V37}
The remaining balanced packet must now be tested before the crude
terminal majorants discard its global weights.
\begin{enumerate}[leftmargin=3em]
\item Return to the exact pre-majorization formula for one
      first-connectivity fibre with one saturated coordinate.
\item Retain separately the protected response
      \(d_p^{-1}\), the normalization by the admissible label volume,
      the exact exchange monomial, and every unequal row-mass factor.
\item Determine whether the same-coordinate constraint forcing
      \(q=p+O_s(\sqrt{\log d_p})\) removes an independent spin sum or
      supplies a compensating inverse label volume.
\item Compare the result with the generic
      \(1/(M\prod\Xi)\) substitute only after the exact powers have
      been counted.
\item If the exact fibre still has order \(d_p^{-1}\), assemble the
      complete physical packet, including prefix, suffix, exchange,
      and all contracted cumulant sectors, before applying absolute
      values.
\item Propagate a proved gain through the \(q=0,1\) census; otherwise
      isolate the precise global aggregation inequality that remains.
\end{enumerate}
\end{programme}

\begin{firewall}[Claim boundary after fourth-series V36]
\label{fire:V36-claim-boundary}
V36 proves that the hoped-for generic local centering gain does not
exist and replaces it by an exact arity-descent theorem.  It does not
show that every contracted lower-order symbol is nonzero, nor that the
full single-coordinate packet fails: exact exchange and aggregation
weights have not yet been restored.  It does not close all \(q=0,1\),
prove global quintic MTA, all-order MTA, WUFC, CPM, cutoff removal,
reflection positivity, Osterwalder--Schrader reconstruction,
construction of the continuum Yang--Mills measure, or a uniform
positive continuum spectral gap.  The full Yang--Mills mass gap has
not been proved.
\end{firewall}

\paragraph{Parent-version record.}
V36 was formed from
\texttt{Renewal\_Yang\_Mills\_Mass\_Gap\_Research\_Notes\_V35.tex}.
The V35 parent had \(17\,182\) logical source lines,
\(602\,033\) bytes, and SHA--256
\[
 \texttt{37472d4ef58415b9f6cd310614691f53f76b6336175cdc4411b2d8d081d5716c}.
\]
The next compulsory companion audit is V40.

% =============================================================================
% END APPENDED FOURTH-SERIES V36 MATERIAL
% =============================================================================

% \end{document} % disabled by fourth-series V37 append

% =============================================================================
% BEGIN APPENDED FOURTH-SERIES V37 MATERIAL
% =============================================================================

\clearpage
\chapter{V37: A saturated sparse fibre defeats the classical
quintic tree atlas}
\label{chap:V37-sparse-counterpacket}

\section{Purpose: test the exact target, not another local estimate}
\label{sec:V37-purpose}

V36 proves that the protected high-pair block cannot acquire a
generic extra inverse dimension from local M\"obius centering.  The
V36 programme therefore returned to the exact exchange-rescue weight.
This chapter performs that upstream test on a source which is the
single-coordinate contraction of the V23 skeleton: the two high rows
meet only at the first-connectivity coordinate, and every other
representation is fixed.

The result is decisive for the present proof architecture.  The
same-coordinate constraint removes all independent high-spin sums,
but the marked-tree target is pointwise in the representation labels
and therefore receives no volume gain.  The exact classical tree
weight is \(A_j^{-3}\), while an admissible paid connected coefficient
is \(A_j^{-5/2}\).  Their ratio diverges like \(d_j\).
Consequently not only the \(DXQTRC_5\) compiler but the proposed
uniform classical quintic marked-tree inequality itself is false in
this normalization.

This is a no-go theorem for that intermediate estimate.  It is not a
counterexample to Yang--Mills theory or to the existence of a mass
gap.  It says that the continuum programme must retain a quantum
bridge response or a representation-summed norm before demanding a
classical four-edge tree product.

\section{The four-coordinate saturated skeleton}
\label{sec:V37-skeleton}

Let the five rows be \(a,b,c,d,e\), and order four physical
coordinates by
\[
 u<v<w<r.
\]
Fix a nontrivial spin \(\ell\), independent of \(j\), and use the
following nontrivial row representations:
\begin{center}
\begin{tabular}{c|ccccc}
coordinate & \(a\)&\(b\)&\(c\)&\(d\)&\(e\)\\
\hline
\(u\)&\(\ell\)&\(\ell\)&0&0&0\\
\(v\)&\(\ell\)&0&0&\(\ell\)&0\\
\(w\)&0&0&\(\ell\)&0&\(\ell\)\\
\(r\)&\(j\)&0&\(j\)&0&0 .
\end{tabular}
\end{center}
All omitted representations are trivial.  Put
\[
 A_j:=1+j(j+1),
 \qquad
 B:=1+\ell(\ell+1).
\]

\begin{lemma}[Unique first connectivity and literal tree]
\label{lem:V37-legality}
Immediately below \(r\), the cumulative row partition is
\[
 \rho=\{\{a,b,d\},\{c,e\}\}.
\]
The coordinate \(r\) is the unique first coordinate at which all five
rows are connected.  The recursive prefix edges and quotient edge
are forced to be
\[
 P=\{ab,ad,ce\},
 \qquad
 Q=\{ac\},
 \qquad
 T_*=\{ab,ad,ce,ac\}.
\]
Every connected local-cumulant word on this support uses all four
displayed binary connections, so the coefficient-one
first-connectivity formula has one nonzero \((r,\rho)\) term.
\end{lemma}

\begin{proof}
At \(u\), the component \(\{a,b\}\) is formed.  At \(v\), row \(d\)
joins it through \(a\), giving \(\{a,b,d\}\).  At \(w\), the component
\(\{c,e\}\) is formed.  No coordinate below \(r\) crosses these two
components.  At \(r\), the only nontrivial cross-block pair is \(ac\),
so it connects them for the first time.  This gives the displayed
partition and literal tree.

The graph of all available nontrivial row pairs is itself the tree
\(T_*\).  Removing any one of its four edges disconnects the five row
labels.  A local-cumulant word missing that edge therefore has
non-full cumulative join and is annihilated by the outer connected
cumulant.  Hence the connected source has the single coefficient-one
word asserted.
\end{proof}

\begin{lemma}[Exact masses and saturation]
\label{lem:V37-masses}
The row masses are
\begin{equation}
\label{eq:V37-masses}
 \Xi_a=A_j+2B,\qquad
 \Xi_c=A_j+B,\qquad
 \Xi_b=\Xi_d=\Xi_e=B.
\end{equation}
Thus \(M=\Xi_a\), \(\Xi_c/M\to1\), and the coordinate \(r\) carries
asymptotically all the mass of both \(a\) and \(c\).
\end{lemma}

\begin{proof}
Sum the active weights in each row of the table.  Since \(B\) is
fixed and \(A_j\to\infty\),
\[
 {A_j\over\Xi_a}\longrightarrow1,
 \qquad
 {A_j\over\Xi_c}\longrightarrow1.
\]
This is precisely the single-coordinate saturation alternative of
V34.
\end{proof}

\section{The entire classical tree atlas on the skeleton}
\label{sec:V37-tree-atlas}

\begin{lemma}[Exact nonzero pair stocks]
\label{lem:V37-stocks}
The only nonzero global pair stocks are the four edges of \(T_*\), and
\begin{align}
 \theta_{ab}^{<r}=\theta_{ab}
 &= {B\over\Xi_a},
 &
 \theta_{ad}^{<r}=\theta_{ad}
 &= {B\over\Xi_a},
\label{eq:V37-a-stocks}\\
 \theta_{ce}^{<r}=\theta_{ce}
 &= {B\over\Xi_c},
 &
 \theta_{ac,r}=\theta_{ac}
 &= {A_j^2\over\Xi_a\Xi_c}.
\label{eq:V37-c-root-stocks}
\end{align}
In particular, \(ac\) is the unique strong pair and its full stock is
concentrated at \(r\).
\end{lemma}

\begin{proof}
Each pair in \(T_*\) meets at exactly its displayed coordinate and no
other pair ever meets.  At a fixed-spin edge the overlap numerator is
\(B^2\); one endpoint has mass \(B\), giving the first three
fractions.  At \(r\), the overlap numerator is \(A_j^2\), giving the
last fraction.  It tends to one by
\Cref{lem:V37-masses}.
\end{proof}

\begin{theorem}[Exact tree sum and exact exchange-rescue weight]
\label{thm:V37-exact-tree-weight}
On the V37 support,
\begin{equation}
\label{eq:V37-tree-sum}
 \sum_{\tau\in\mathfrak T_5}
 \prod_{xy\in E(\tau)}\theta_{xy}
 =
 {B^3A_j^2\over\Xi_a^3\Xi_c^2}.
\end{equation}
For the literal first-connectivity fibre,
\begin{equation}
\label{eq:V37-exchange-weight}
 \mathfrak x_{T_*,r}
 =
 {8B^3A_j^2\over\Xi_a^3\Xi_c^2}.
\end{equation}
Both quantities are comparable, with constants depending only on
\(\ell\), to \(A_j^{-3}\).
\end{theorem}

\begin{proof}
By \Cref{lem:V37-stocks}, the support graph of all nonzero stocks is
the four-edge tree \(T_*\).  Thus every labelled spanning-tree
monomial other than \(T_*\) vanishes.  Multiplying the four values in
\eqref{eq:V37-a-stocks}--\eqref{eq:V37-c-root-stocks} proves
\eqref{eq:V37-tree-sum}.

In the V98 exchange-rescue definition, \(\tau=T_*\) is forced and the
nonempty local subset of \(Q=\{ac\}\) is \(\{ac\}\).  Every one of the
\(2^{|P|}=8\) choices assigning a prefix edge its prefix or global
stock has the same value, because each prefix stock has no occurrence
at or above \(r\).  This proves \eqref{eq:V37-exchange-weight}.
Finally, \(\Xi_a,\Xi_c\asymp_\ell A_j\).
\end{proof}

\begin{assessment}[Why saturation supplies no exchange bonus]
\label{ass:V37-no-volume-bonus}
The equality \(\theta_{ac,r}=\theta_{ac}\asymp1\) uses the concentrated
coordinate exactly once.  A simple tree cannot use a second global
copy of the same edge label.  The remaining three actual edges have
stocks of order \(A_j^{-1}\), so their product is \(A_j^{-3}\).
The constraint that the two root spins are equal removes a label
variable but does not enlarge the pointwise tree monomial.
\end{assessment}

\section{An exact paid connected coefficient}
\label{sec:V37-coefficient}

Let \(\mathbb F_\ell\) be the complete recursive prefix carrier on
the three fixed coordinates \(u,v,w\) in the forced tree fibre.
It is the tensor product, after row rebracketing, of three fixed-spin
binary covariance carriers.  As proved for the V23 prefix,
\(\mathbb F_\ell\ne0\), is self-adjoint, and is independent of \(j\).
Let
\[
 \mathbb D_j(s)
 =
 \sum_{K=0}^{2j}
 \left(e^{-sK(K+1)}-e^{-2sj(j+1)}\right)P_{j,K}
\]
be the complete root covariance, and put the sole payer at \(r\):
\[
 \mathbb R_s
 =
 \boldsymbol1_{\{C_{\rm out}>0\}}
 {C_{\rm out}\over1-e^{-sC_{\rm out}}}.
\]

\begin{lemma}[Fixed prefix product detection]
\label{lem:V37-prefix-detection}
There is a row-product unit vector \(v_\ell\) and a constant
\(\beta_{\ell,s}>0\), both independent of \(j\), such that
\[
 \left|
 \langle v_\ell,\mathbb F_\ell v_\ell\rangle
 \right|
 =\beta_{\ell,s}.
\]
\end{lemma}

\begin{proof}
Each of the three fixed binary covariance carriers has a nonzero
singlet coefficient
\[
 1-e^{-2s\ell(\ell+1)}.
\]
Thus their tensor product is a nonzero self-adjoint operator.
If its expectation vanished on every row-product vector, polarization
and the fact that product rank-one projections span the
finite-dimensional operator space would make the operator zero.
Therefore a detecting row-product vector exists.  Its expectation is
a fixed positive absolute constant, denoted
\(\beta_{\ell,s}\).
\end{proof}

\begin{lemma}[Complete paid root has a direct protected lower bound]
\label{lem:V37-paid-root-direct}
For
\[
 \phi_j:=|j,j\rangle\otimes|j,-j\rangle
\]
and all sufficiently large \(j\),
\begin{equation}
\label{eq:V37-paid-root-direct}
 \langle\phi_j,\mathbb R_s\mathbb D_j(s)\phi_j\rangle
 \ge {c_s\over d_j}.
\end{equation}
\end{lemma}

\begin{proof}
Write
\[
 p_{j,K}:=\langle\phi_j,P_{j,K}\phi_j\rangle
\]
and
\[
 r_K(s):=
 \begin{cases}
 0,&K=0,\\
 \displaystyle {K(K+1)\over1-e^{-sK(K+1)}},&K\ge1.
 \end{cases}
\]
The commuting spectral resolutions give
\[
 \langle\phi_j,\mathbb R_s\mathbb D_j(s)\phi_j\rangle
 =
 \sum_{K=1}^{2j}
 r_K(s)\left(e^{-sK(K+1)}-e^{-2sj(j+1)}\right)p_{j,K}.
\]
Retain the positive \(K=1\) heat contribution and bound the absolute
value of the entire replicated-endpoint contribution:
\[
\begin{split}
 \langle\phi_j,\mathbb R_s\mathbb D_j(s)\phi_j\rangle
 &\ge
 r_1(s)e^{-2s}p_{j,1}
 -
 e^{-2sj(j+1)}
 \sum_{K=1}^{2j}r_K(s)p_{j,K}.
\end{split}
\]
Since \(K(K+1)\le2j(2j+1)\) and
\(1-e^{-sK(K+1)}\ge1-e^{-2s}\) for \(K\ge1\),
\[
 \sum_{K=1}^{2j}r_K(s)p_{j,K}
 \le C_sj^2.
\]
The second term is therefore \(O_s(j^2e^{-2sj(j+1)})\).
The exact Clebsch--Gordan formula from V13 gives
\[
 p_{j,1}
 =
 {6j\over(2j+1)(2j+2)}
 \ge {1\over2j+2}.
\]
The Gaussian endpoint error is \(o(d_j^{-1})\), proving
\eqref{eq:V37-paid-root-direct}.
\end{proof}

\begin{theorem}[Normalized connected coefficient lower bound]
\label{thm:V37-coefficient-lower}
Let \(\widetilde{\mathbb W}_j\) be the exact V37 connected word, with
the payer at \(r\), after the five row normalizations and with the
identity admissible output dual.  Then, for all sufficiently large
\(j\),
\begin{equation}
\label{eq:V37-coefficient-lower}
 \kappa_\otimes^{\rm abs}(\widetilde{\mathbb W}_j)
 \ge
 {c_{\ell,s}\over
   d_j\,\Xi_a\Xi_cB^3}.
\end{equation}
In particular,
\[
 \kappa_\otimes^{\rm abs}(\widetilde{\mathbb W}_j)
 \gtrsim_{\ell,s} A_j^{-5/2}.
\]
\end{theorem}

\begin{proof}
By \Cref{lem:V37-legality} and the exact coefficient-one
first-connectivity formula, the normalized word is
\[
 \widetilde{\mathbb W}_j
 =
 {1\over\Xi_a\Xi_cB^3}\,
 \mathbb F_\ell\otimes
 \bigl(\mathbb R_s\mathbb D_j(s)\bigr),
\]
up to canonical rebracketing by rows.  The identity output dual is
admissible by the exact V94 dual algebra and changes no factor.

Rebracket \(v_\ell\otimes\phi_j\) by rows.  It remains a row-product
unit vector because the physical coordinates are tensor factors
inside each row.  Its expectation factors exactly.  Apply
\Cref{lem:V37-prefix-detection,lem:V37-paid-root-direct} and take the
absolute value to obtain \eqref{eq:V37-coefficient-lower}.
Since \(d_j\asymp A_j^{1/2}\) and
\(\Xi_a,\Xi_c\asymp_\ell A_j\), the last estimate follows.
\end{proof}

\section{Failure of the pointwise and global classical tree targets}
\label{sec:V37-failure}

\begin{theorem}[The direct exchange-rescue card is false in the
saturated \(q=1\) sector]
\label{thm:V37-DXQTRC-false}
No constant \(C_{\ell,s}\), independent of \(j\), can satisfy
\[
 \kappa_\otimes^{\rm abs}(\widetilde{\mathbb W}_j)
 \le C_{\ell,s}\mathfrak x_{T_*,r}
\]
on the V37 family.  Quantitatively,
\begin{equation}
\label{eq:V37-exchange-ratio}
 {\kappa_\otimes^{\rm abs}(\widetilde{\mathbb W}_j)
  \over\mathfrak x_{T_*,r}}
 \ge
 c_{\ell,s}
 {\Xi_a^2\Xi_c\over B^6A_j^2d_j}
 \asymp_{\ell,s}{A_j\over d_j}
 \asymp d_j.
\end{equation}
\end{theorem}

\begin{proof}
Divide \eqref{eq:V37-coefficient-lower} by
\eqref{eq:V37-exchange-weight}.  Fixed numerical factors and powers
of \(B\) are absorbed into \(c_{\ell,s}\).  The mass asymptotics in
\Cref{lem:V37-masses} give the first comparison.  Finally
\[
 A_j=1+j(j+1)\asymp d_j^2,
\]
so \(A_j/d_j\asymp d_j\to\infty\).
\end{proof}

\begin{theorem}[Failure of the proposed uniform classical quintic
marked-tree inequality]
\label{thm:V37-MTA-false}
In the normalization of
\(\textup{\texttt{eq:V97-conditional-quintic-MTA}}\), there is no
constant \(C_{G,s}\) for which every admissible paid connected
five-row coefficient obeys
\[
 \kappa_\otimes^{\rm abs}(\mathcal C_5^{\rm pay})
 \le
 C_{G,s}
 \sum_{\tau\in\mathfrak T_5}
 \prod_{xy\in E(\tau)}\theta_{xy}.
\]
\end{theorem}

\begin{proof}
Use the V37 support and put the payer at \(r\).
\Cref{lem:V37-legality} shows that the exact coefficient-one
first-connectivity decomposition has only the displayed connected
word; there is no second first-connectivity fibre or connected
local-cumulant word with which it can cancel.  The identity output
dual is admissible.  Thus
\Cref{thm:V37-coefficient-lower} is a lower bound for the exact paid
connected coefficient on this support.

The complete tree sum, not a selected monomial or a lower floor, is
the exact quantity in \eqref{eq:V37-tree-sum}.  Dividing the lower
bound by that complete sum gives the same divergent ratio as
\eqref{eq:V37-exchange-ratio}, up to the fixed factor eight.  Hence no
uniform \(C_{G,s}\) exists.
\end{proof}

\begin{corollary}[No hidden label-volume or fibre-aggregation repair]
\label{cor:V37-no-aggregation}
The V37 failure cannot be repaired within the stated pointwise MTA
target by summing over first-connectivity coordinates, predecessor
partitions, tree fibres, endpoint lifts, or an independent second
high-spin label.
\end{corollary}

\begin{proof}
There is one first-connectivity coordinate \(r\), one predecessor
partition \(\rho\), one nonzero tree \(T_*\), and forced endpoint
lifts.  The two root spins are exactly equal to \(j\), so there is no
second high-spin label.  All eight exchange role assignments are
already included in \eqref{eq:V37-exchange-weight}; they change only a
fixed factor.  The complete tree atlas is already evaluated in
\eqref{eq:V37-tree-sum}.  Therefore none of the named finite or absent
aggregations can change the divergent power.
\end{proof}

\begin{firewall}[What has and has not been disproved]
\label{fire:V37-scope}
\Cref{thm:V37-MTA-false} disproves the proposed \emph{classical
coefficientwise quintic marked-tree atlas} for the heat-cumulant
carrier and normalization used in this programme.  It does not
disprove a representation-summed estimate, a Schatten or Hilbert-space
bound retaining Plancherel weights, an atlas augmented by a quantum
bridge or vertex response, Wilson ultraviolet finiteness, the
existence of continuum Yang--Mills theory, or a positive mass gap.
No conclusion about the truth or falsity of the Yang--Mills mass-gap
conjecture follows from this interface counterexample.
\end{firewall}

\section{V37 result ledger and V38 upstream replacement programme}
\label{sec:V37-ledger}

\begin{theorem}[Fourth-series V37 result ledger]
\label{thm:V37-results}
Relative to V1--V36, V37 proves:
\begin{enumerate}[label=\textup{(V37.\arabic*)},leftmargin=6.3em]
\item a legal four-coordinate \(3+2\) first-connectivity skeleton in
      the single-coordinate saturated \(q=1\) sector;
\item uniqueness of its connected word and literal five-row tree;
\item its exact row masses and all nonzero pair stocks;
\item exact evaluation of both its complete classical tree sum and
      its exchange-rescue weight at scale \(A_j^{-3}\);
\item a direct \(c_sd_j^{-1}\) lower bound for the complete paid root
      covariance, without fine output replacement;
\item an exact normalized connected-coefficient lower bound of scale
      \(A_j^{-5/2}\);
\item divergence by \(d_j\) of the direct coefficient/tree-weight
      ratio; and
\item failure of both \(DXQTRC_5\) in this sector and the proposed
      uniform classical coefficientwise quintic MTA.
\end{enumerate}
\end{theorem}

\begin{proof}
These statements are
\Cref{lem:V37-legality,
lem:V37-masses,
lem:V37-stocks,
thm:V37-exact-tree-weight,
lem:V37-paid-root-direct,
thm:V37-coefficient-lower,
thm:V37-DXQTRC-false,
thm:V37-MTA-false}.
\end{proof}

\begin{programme}[V38: identify the weakest sufficient quantum atlas]
\label{prog:V37-V38}
The false classical pointwise target must not remain a premise farther
downstream.
\begin{enumerate}[leftmargin=3em]
\item Trace the use of quintic MTA in WUFC and CPM and determine
      whether those arguments truly require a coefficientwise
      four-edge tree product, or only a representation-summed
      contraction estimate.
\item Restore the exact Peter--Weyl/Plancherel weights before taking
      the high-spin supremum.  Test whether their label sum absorbs
      the V37 factor \(d_j\).
\item If a pointwise statement is indispensable, define a quantum
      bridge atlas which retains the protected root response and prove
      a support-uniform coordinate-summation theorem for it.
\item Prove necessity on the V37 family: every replacement right side
      must be at least \(c_{\ell,s}A_j^{-5/2}\).
\item Recheck the already closed \(q\ge2\), diffuse \(q=1\), and
      second-atom sectors against the replacement target without
      weakening their valid estimates.
\item Only after a sufficient replacement is proved should the
      \(q=0\) census or higher perturbative order be resumed.
\end{enumerate}
\end{programme}

\begin{firewall}[Claim boundary after fourth-series V37]
\label{fire:V37-claim-boundary}
V37 identifies a false intermediate theorem and prevents it from being
used in a claimed mass-gap proof.  It has not yet supplied the
replacement quantum or representation-summed atlas, closed \(q=0,1\)
under such an atlas, proved a valid global quintic estimate, all-order
control, WUFC, CPM, cutoff removal, reflection positivity,
Osterwalder--Schrader reconstruction, construction of the continuum
Yang--Mills measure, or a uniform positive continuum spectral gap.
The full Yang--Mills mass gap has not been proved.
\end{firewall}

\paragraph{Parent-version record.}
V37 was formed from
\texttt{Renewal\_Yang\_Mills\_Mass\_Gap\_Research\_Notes\_V36.tex}.
The V36 parent had \(17\,784\) logical source lines,
\(623\,427\) bytes, and SHA--256
\[
 \texttt{adbd5ae1e5bfba03aa29c6c94b520f72e35bcb9591dcaa35aa6c1af2be46855d}.
\]
The next compulsory companion audit is V40.

% =============================================================================
% END APPENDED FOURTH-SERIES V37 MATERIAL
% =============================================================================

% \end{document} % disabled by fourth-series V38 append

% =============================================================================
% BEGIN APPENDED FOURTH-SERIES V38 MATERIAL
% =============================================================================

\clearpage
\chapter{V38: Restoring the balanced Fourier height quotient closes
the saturated protected terminal}
\label{chap:V38-height-quotient}

\section{Correction of scope after V37}
\label{sec:V38-scope-correction}

V37 proves a genuine no-go theorem for the scalar coefficientwise
atlas stated in f3 V97:
\[
 \kappa_\otimes^{\rm abs}(\mathcal C_5^{\rm pay})
 \not\lesssim
 \sum_{\tau\in\mathfrak T_5}
 \prod_{ij\in E(\tau)}\theta_{ij}.
\]
That scalar inequality was used as a sufficient surrogate after the
Fourier weight quotient had been bounded above by one.  The original
MTA gate which compiles into WUFC is more precise.  In the notation of
archive f1 \texttt{def:V31-MTA}, each input and output coefficient is
measured with the balanced Fourier weight
\[
 w_\rho(\boldsymbol R)
 =
 e^{\rho\operatorname{ht}(\boldsymbol R)}
 \bigl(1+C_E(\boldsymbol R)\bigr),
\]
together with its Peter--Weyl dimension and trace norm.  Therefore a
resolved fusion coefficient contains the exact quotient
\begin{equation}
\label{eq:V38-height-quotient}
 \exp\!\left\{
 \rho\left(
 \operatorname{ht}(\boldsymbol T)
 -
 \sum_i\operatorname{ht}(\boldsymbol R_i)
 \right)\right\}.
\end{equation}
Highest-weight subadditivity makes this quotient at most one, but it
can be exponentially smaller.

The protected packet is exactly such a case: two high inputs fuse to
a logarithmic-height output.  Hence the strongest wording of the V37
title must not be read as a disproof of the balanced-Fourier MTA.
What V37 disproves is the \emph{height-erased scalar V97 surrogate}.
This chapter proves that the retained height quotient closes the V35
balanced low-output packet pointwise, before any representation sum.

\begin{record}[Correct implication boundary]
\label{rec:V38-correct-boundary}
The valid logical statements are
\[
 \text{scalar V97 atlas}
 \Longrightarrow
 \text{balanced-Fourier MTA}
 \Longrightarrow
 \text{WUFC},
\]
where the first implication is a sufficient proof route, not an
equivalence.  V37 refutes its premise.  It does not refute the middle
statement, and no use of
\Cref{thm:V37-MTA-false} beyond the scalar V97 normalization is
permitted.
\end{record}

\section{The exact high-pair height deficit}
\label{sec:V38-height-deficit}

For \(SU(2)\), normalize representation height by
\[
 h(j):=j.
\]
Changing to the integer highest weight \(2j\) merely replaces
\(\rho\) by \(2\rho\).  At one physical coordinate, let the input
spins in a moment block be
\[
 p,r_2,\ldots,r_t.
\]
Couple the complementary factors first, and let \(q\) be one of their
resultant spins.  Then couple \(V_p\otimes V_q\) to total output
\(V_K\).

\begin{lemma}[Height deficit from a compressed high pair]
\label{lem:V38-height-deficit}
Let
\[
 H_{\rm in}:=p+\sum_{\nu=2}^t r_\nu.
\]
For every occurring complementary resultant \(q\) and final output
\(K\),
\begin{equation}
\label{eq:V38-height-deficit}
 H_{\rm in}-K
 \ge p+q-K.
\end{equation}
Consequently the local exponential part of the balanced Fourier
quotient obeys
\begin{equation}
\label{eq:V38-local-height-ratio}
 e^{\rho(K-H_{\rm in})}
 \le e^{-\rho(p+q-K)}
 \le1.
\end{equation}
\end{lemma}

\begin{proof}
Every irreducible resultant of
\(\bigotimes_{\nu=2}^tV_{r_\nu}\) satisfies
\[
 q\le\sum_{\nu=2}^t r_\nu.
\]
Subtract \(K\) and add \(p\) to obtain
\eqref{eq:V38-height-deficit}.  The first inequality in
\eqref{eq:V38-local-height-ratio} follows because \(\rho>0\).
The Clebsch--Gordan triangle inequality \(K\le p+q\) gives the second.
\end{proof}

\begin{lemma}[Exponential gain on the V35 low window]
\label{lem:V38-low-height-gain}
Assume the V35 saturated low-window constraints
\[
 |p-q|+K\le A_s\sqrt{\log(2+d_p)}.
\]
For every fixed \(\rho>0\), there is \(p_{\rho,s}\) such that
\begin{equation}
\label{eq:V38-low-height-gain}
 e^{\rho(K-H_{\rm in})}
 \le e^{-\rho p}
 \qquad(p\ge p_{\rho,s}).
\end{equation}
In particular, for every \(D,N\ge0\),
\begin{equation}
\label{eq:V38-exp-beats-power}
 p^D e^{\rho(K-H_{\rm in})}
 \le C_{\rho,s,D,N}d_p^{-N}.
\end{equation}
\end{lemma}

\begin{proof}
The window gives
\[
 q\ge p-A_s\sqrt{\log(2+d_p)}
\]
and
\[
 K\le A_s\sqrt{\log(2+d_p)}.
\]
Therefore
\[
 p+q-K
 \ge
 2p-2A_s\sqrt{\log(2+d_p)}
 \ge p
\]
for all sufficiently large \(p\).  Apply
\Cref{lem:V38-height-deficit}.  Exponential decay in \(p\) dominates
every polynomial in \(d_p=2p+1\), proving
\eqref{eq:V38-exp-beats-power}.
\end{proof}

\section{Balanced-Fourier closure of the protected local packet}
\label{sec:V38-local-closure}

\begin{lemma}[All nonexponential local costs are polynomial]
\label{lem:V38-polynomial-cost}
Fix row arity \(t\le5\), one moment-partition term, an output payer
placement, exact equivariant recoupling, and an admissible V94
dualizer.  On the V35 low window, the ratio of all balanced-Fourier
factors other than
\eqref{eq:V38-local-height-ratio}, including output dimensions,
Casimir weights, multiplicities, the payer, and the finite output
sum, is bounded by
\[
 C_{s,t}(1+p)^D
\]
for a fixed exponent \(D\) independent of the representation labels.
\end{lemma}

\begin{proof}
At fixed arity, an \(SU(2)\) dimension is linear in its spin, every
Clebsch--Gordan multiplicity is one, and the number of intermediate
recoupling labels is bounded by a fixed product of linear spin
ranges.  The balanced Casimir weight is quadratic in the spins.  An
unpaid heat multiplier is at most one.  On the low window, the payer
is \(O_s(1+K(K+1))=O_s(1+\log d_p)\).  Exact recoupling maps and V94
dualizers are contractions on their full multiplicity spaces.
Multiplying these finitely many costs gives one fixed polynomial in
\(p\); a logarithm can be absorbed into that polynomial.
\end{proof}

\begin{theorem}[Arbitrary-power balanced-Fourier closure of the low
window]
\label{thm:V38-low-BFA-close}
For every \(\rho>0\) and every \(N\ge0\), the complete V35 saturated
low-output packet, including the payer in any allowed position,
all output multiplicities, exact V94 duality, and the signed local
cumulant, has normalized balanced-Fourier response
\begin{equation}
\label{eq:V38-low-BFA-close}
 \mathfrak B_{\rho,s}^{\rm lo}(p)
 \le C_{\rho,s,N}d_p^{-N}.
\end{equation}
The estimate holds termwise for the fixed finite row-partition
expansion and hence also for its assembled signed sum.
\end{theorem}

\begin{proof}
For one fixed row-partition term and one fixed recoupling channel, the
exact normalized BFA coefficient contains
\eqref{eq:V38-local-height-ratio}.  By
\Cref{lem:V38-polynomial-cost}, every remaining norm and counting
factor is at most \(C(1+p)^D\).  Apply
\eqref{eq:V38-exp-beats-power} with \(N\) enlarged by \(D+2\).

The number of row partitions, recursive joining nodes, payer
placements, and source orientations depends only on arity five.
The allowed logarithmic output and mismatch channels contribute only
another polynomial or polylogarithmic factor, already absorbed by the
arbitrary choice of \(N\).  Summing these fixed and polynomial
families proves \eqref{eq:V38-low-BFA-close}.  Since the estimate is
valid before the triangle inequality, it also bounds the complete
signed cumulant; V36 arity descent is compatible with, but not needed
to manufacture, the height gain.
\end{proof}

\begin{theorem}[The complete saturated terminal closes in the actual
weighted interface]
\label{thm:V38-saturated-close}
In the \(SU(2)\) balanced-Fourier MTA interface with every fixed
\(\rho>0\), the single-coordinate saturated one-partner terminal of
V35 satisfies the full direct \(M^{-1}\) response required by the
one-strong quintic normalization.
\end{theorem}

\begin{proof}
Split into the V35 cases.  A singleton block and a block omitting the
unique partner are arbitrary-power small by
\Cref{lem:V35-singleton,lem:V35-omit-c}.  The high-output packet obeys
the \(C_sd_p^{-2}\) direct response by
\Cref{thm:V35-high-close}; the exponential Fourier quotient is at
most one by highest-weight subadditivity, so it cannot worsen that
bound.  The remaining low-output packet is arbitrary-power small by
\Cref{thm:V38-low-BFA-close}.  Since
\[
 d_p^2\asymp M,
\]
all branches have the required \(CM^{-1}\) response.
\end{proof}

\begin{remark}[Joining, prefix, and role-owned coordinates]
\label{rem:V38-location}
The proof of \Cref{thm:V38-low-BFA-close} uses only the actual input
and output representations.  If the saturated coordinate is an
unmarked moment, the same height quotient is present.  If it is a
joining cumulant, each moment-partition term has the quotient and the
finite sum is harmless.  Tree, lift, strong, selected, and payer roles
create no new source representation and therefore do not remove the
height deficit.  Thus the V36 moment-or-descent dichotomy changes the
local symbol but not this weighted closure.
\end{remark}

\section{Re-evaluation of the V37 sparse packet}
\label{sec:V38-V37-reevaluation}

\begin{proposition}[The V37 packet is superpolynomially small in BFA]
\label{prop:V38-V37-BFA}
Let the V37 root inputs be \(j,j\).  After division by the two input
balanced exponential weights, the complete paid root covariance
satisfies, for every \(N\ge0\),
\begin{equation}
\label{eq:V38-V37-BFA}
 \sum_{K=0}^{2j}
 e^{\rho(K-2j)}
 P_s(j,K)
 \left|
 e^{-sK(K+1)}-e^{-2sj(j+1)}
 \right|
 \le C_{\rho,s,N}d_j^{-N},
\end{equation}
where \(P_s(j,K)\) may be any fixed polynomial majorant for the
dimension, balanced Casimir, payer, and recoupling costs.
\end{proposition}

\begin{proof}
Split the sum at \(K=j\).  If \(K\le j\), then
\[
 e^{\rho(K-2j)}\le e^{-\rho j},
\]
while the heat difference is at most two and the sum of the fixed
polynomial majorant over \(O(j)\) values is polynomial in \(j\).
This part is therefore smaller than every inverse power of \(d_j\).

If \(K>j\), then
\[
 e^{-sK(K+1)}
 \le e^{-sj(j+1)}
\]
and the replicated endpoint term is
\(e^{-2sj(j+1)}\).  The height quotient is at most one, so the second
part is bounded by a polynomial in \(j\) times
\(e^{-sj(j+1)}\), again smaller than every inverse power.  Add the two
bounds.
\end{proof}

\begin{corollary}[The V37 obstruction is absent from the actual MTA
norm]
\label{cor:V38-V37-not-MTA}
The exact V37 connected coefficient violates the height-erased scalar
tree card by the factor \(d_j\), but satisfies its balanced-Fourier
marked-tree comparison with arbitrary polynomial room.  No
representation-label summation is required for this repair.
\end{corollary}

\begin{proof}
The three prefix coordinates and their balanced weights are fixed in
\(j\).  Apply \Cref{prop:V38-V37-BFA} to the root coordinate and then
divide by the same row-mass normalization used in V37.  The complete
nonzero tree sum is only polynomially small, of order \(A_j^{-3}\), by
\Cref{thm:V37-exact-tree-weight}.  Choose \(N\) large enough to
dominate that polynomial.  This proves the weighted comparison.
\end{proof}

\begin{assessment}[The indispensable order of estimates]
\label{ass:V38-order}
The height quotient may be replaced by one only after a
representation-uniform scalar estimate has already been proved.
V37 shows that making this replacement before treating the protected
low-output sector destroys a decisive exponential gain and creates a
false scalar target.  Conversely, height alone is not counted as an
extra local heat response: it is an exact quotient of the output and
input norms.
\end{assessment}

\section{Consequences for WUFC and the remaining frontier}
\label{sec:V38-consequences}

\begin{proposition}[The MTA-to-WUFC compiler requires the weighted
atlas, not the scalar V97 surrogate]
\label{prop:V38-compiler-scope}
The leaf-peeling proof indexed by
\texttt{thm:V31-MTA-compiler} in archive f1 uses:
\begin{enumerate}[label=\textup{(\roman*)},leftmargin=3em]
\item the balanced marked seminorms, including their exact
      probability normalization;
\item the BFA output norm, including the height and Casimir weights;
\item the tree marking of shared physical links; and
\item factorial-exponential control of the number of labelled trees.
\end{enumerate}
It does not require the separate height-erased scalar inequality
refuted in V37.
\end{proposition}

\begin{proof}
The compiler starts from f1
\texttt{eq:V31-MTA}, whose left side is a BFA norm and whose right side
is a product of marked BFA seminorms.  Its proof peels tree leaves by
\texttt{eq:V31-mark-normalization}; it never replaces the exact
Fourier weight quotient by one.  The scalar row-stock atlas is one
sufficient way to prove a coefficient bound inside that norm, but is
not an assumption of the spatial peeling argument.  Therefore its
failure does not invalidate the compiler.
\end{proof}

\begin{firewall}[Group and scope boundary]
\label{fire:V38-group-scope}
V38 proves the protected-terminal repair for the \(SU(2)\) programme
treated in V31--V37.  The analogous statement for \(SU(3)\) requires
the corresponding highest-weight deficit and multiplicity-polynomial
ledger and is not claimed here.  V38 closes the saturated local
terminal in the actual weighted interface; it does not by itself
establish every quintic BFA source family, the entire \(q=0\) sector,
an arity-uniform MTA constant, or WUFC.
\end{firewall}

\section{V38 result ledger and V39 acquisition order}
\label{sec:V38-ledger}

\begin{theorem}[Fourth-series V38 result ledger]
\label{thm:V38-results}
Relative to V1--V37, V38 proves:
\begin{enumerate}[label=\textup{(V38.\arabic*)},leftmargin=6.3em]
\item the exact distinction between the false height-erased V97
      scalar atlas and the balanced-Fourier MTA which compiles WUFC;
\item a high-pair height-deficit inequality for every recoupling
      channel;
\item exponential suppression of the complete V35 logarithmic-output
      window;
\item arbitrary-power balanced-Fourier closure including payer,
      multiplicities, exact duality, and the signed cumulant;
\item closure of the full single-coordinate saturated one-partner
      terminal in the actual weighted interface;
\item superpolynomial BFA closure of the exact V37 sparse packet; and
\item preservation of the archived MTA-to-WUFC compiler despite the
      V37 scalar no-go theorem.
\end{enumerate}
\end{theorem}

\begin{proof}
These are
\Cref{rec:V38-correct-boundary,
lem:V38-height-deficit,
lem:V38-low-height-gain,
thm:V38-low-BFA-close,
thm:V38-saturated-close,
prop:V38-V37-BFA,
prop:V38-compiler-scope}.
\end{proof}

\begin{programme}[V39: weighted census of the \(q=0\) role-owned core]
\label{prog:V38-V39}
\begin{enumerate}[leftmargin=3em]
\item Return to the exact \(q=0\) source-role census before the scalar
      height quotient was erased.
\item At every maximal-row coordinate, split outputs into a
      height-deficit sector and a height-near-saturated sector.
\item Close the height-deficit sector by the V38 exponential quotient.
      In the height-near-saturated sector, use the actual heat Casimir
      of the large output before payer absorption.
\item Prove a support-uniform dichotomy showing that any persistent
      word has either a source-occurring diffuse bank, a second
      cofinal coordinate, or one of those two local closures.
\item Retain the balanced multi-mark seminorms through the tree
      comparison; do not revert to the false scalar V97 target.
\item State exactly which \(q=0\) fibres remain after this weighted
      census and prepare the compulsory V40 companion audit.
\end{enumerate}
\end{programme}

\begin{firewall}[Claim boundary after fourth-series V38]
\label{fire:V38-claim-boundary}
V38 repairs the V37 overreach and closes the saturated protected
terminal in the balanced-Fourier \(SU(2)\) interface.  It does not
close the complete \(q=0\) sector, prove the full quintic marked-tree
atlas in BFA norm, an all-order atlas, WUFC, CPM, cutoff removal,
reflection positivity, Osterwalder--Schrader reconstruction,
construction of the continuum Yang--Mills measure, or a uniform
positive continuum spectral gap.  The full Yang--Mills mass gap has
not been proved.
\end{firewall}

\paragraph{Parent-version record.}
V38 was formed from
\texttt{Renewal\_Yang\_Mills\_Mass\_Gap\_Research\_Notes\_V37.tex}.
The V37 parent had \(18\,333\) logical source lines,
\(642\,174\) bytes, and SHA--256
\[
 \texttt{74c5155726682e8baa66494fd43530ff8829b62877fd906753600f3991855527}.
\]
The next compulsory companion audit is V40.

% =============================================================================
% END APPENDED FOURTH-SERIES V38 MATERIAL
% =============================================================================

% \end{document} % disabled by fourth-series V39 append

% =============================================================================
% BEGIN APPENDED FOURTH-SERIES V39 MATERIAL
% =============================================================================

\clearpage
\chapter{V39: Fixed-time closure of the role-owned core and the
uniform small-noise transition band}
\label{chap:V39-transition-band}

\section{Purpose and uniformity correction}
\label{sec:V39-purpose}

V38 restores the exponential representation-height quotient and
thereby closes the balanced protected output at every fixed heat time
\(s>0\).  The actual WUFC gate, however, is cofinal in the Wilson
parameter.  In the heat notation of the local model, this sends the
effective time \(s=s_\alpha\) to zero.  A constant denoted \(C_s\) in
a fixed-time theorem may then diverge and cannot be inserted into the
uniform WUFC constant.

This chapter separates the two statements.  First, it proves an exact
height--heat dichotomy and uses it to close the complete \(q=0\)
role-owned core for every fixed \(s>0\).  Second, it gives a sharp
countertest to uniformity and reduces the cofinal problem to a
height-near-saturated parabolic band.  Thus V39 gains a finite-time
theorem while preventing that theorem from being overclaimed as
WUFC.

\section{The height--heat envelope}
\label{sec:V39-envelope}

Consider one positive moment block at a physical coordinate.  Let
\[
 H:=\sum_{\nu=1}^t r_\nu
\]
be the sum of its \(SU(2)\) input heights and let \(K\) be one of its
total output spins.  The exponential part of the balanced Fourier
quotient times the common-variable heat multiplier is
\begin{equation}
\label{eq:V39-envelope-symbol}
 \mathfrak e_{\rho,s}(H,K)
 :=
 e^{-\rho(H-K)}e^{-sK(K+1)}.
\end{equation}
The Clebsch--Gordan inequalities give \(0\le K\le H\).

\begin{lemma}[Exact optimized height--heat envelope]
\label{lem:V39-optimized-envelope}
For all \(H,K\ge0\), \(\rho,s>0\),
\begin{equation}
\label{eq:V39-optimized-envelope}
 \mathfrak e_{\rho,s}(H,K)
 \le
 \exp\!\left\{
 -\rho H+{\rho^2\over4s}
 \right\}.
\end{equation}
More usefully, if one input spin is \(p\), then
\begin{equation}
\label{eq:V39-two-sector-envelope}
 \mathfrak e_{\rho,s}(H,K)
 \le
 \begin{cases}
 e^{-\rho p/2},&K\le p/2,\\
 e^{-sp^2/4},&K>p/2.
 \end{cases}
\end{equation}
\end{lemma}

\begin{proof}
Dropping the favorable term \(-sK\) and completing the square gives
\[
\begin{split}
 -\rho(H-K)-sK(K+1)
 &\le
 -\rho H+\rho K-sK^2\\
 &=
 -s\left(K-{\rho\over2s}\right)^2
 -\rho H+{\rho^2\over4s}.
\end{split}
\]
This proves \eqref{eq:V39-optimized-envelope}.

Since \(H\ge p\), the first sector has
\[
 H-K\ge p/2.
\]
In the second sector,
\[
 K(K+1)\ge K^2>p^2/4.
\]
Insert these two inequalities in
\eqref{eq:V39-envelope-symbol}.
\end{proof}

\begin{corollary}[Fixed-time arbitrary-power local response]
\label{cor:V39-fixed-time-local}
Fix \(s,\rho>0\) and one row arity \(t\le5\).  If a moment block
contains an input spin \(p\to\infty\), then, for every \(D,N\ge0\),
\begin{equation}
\label{eq:V39-fixed-time-power}
 (1+p)^D
 \sum_{K}
 P_t(p,K)\mathfrak e_{\rho,s}(H,K)
 \le C_{\rho,s,t,D,N}d_p^{-N},
\end{equation}
where \(P_t\) is any fixed polynomial accounting for dimensions,
recoupling labels, multiplicities, and one payer.
\end{corollary}

\begin{proof}
Split the output sum at \(K=p/2\) and apply
\eqref{eq:V39-two-sector-envelope}.  At fixed arity, the number of
allowed \(K\)'s and the sum of \(P_t(p,K)\) are polynomial in \(p\).
Both \(e^{-\rho p/2}\) and \(e^{-sp^2/4}\) dominate every inverse
power of \(d_p\) for fixed \(\rho,s>0\).
\end{proof}

\begin{lemma}[Endpoint-replicated terms obey the same conclusion]
\label{lem:V39-replicated}
If the moment partition separates the spin-\(p\) input into its own
heat block or into a block whose other input spins have total
\(o(p)\), then its endpoint heat factor is at most
\[
 e^{-sp^2/4}
\]
for all sufficiently large \(p\).  Hence the fixed-time
arbitrary-power conclusion of
\Cref{cor:V39-fixed-time-local} also holds for every separated term
of a local cumulant.
\end{lemma}

\begin{proof}
For a singleton, the heat factor is \(e^{-sp(p+1)}\).  More generally,
the reverse triangle inequality makes every output of the separated
block at least \(p-o(p)\ge p/2\).  Apply the heat functional calculus.
All remaining factors are contractions before the fixed polynomial
payer and dimension costs.
\end{proof}

\section{A complete fixed-time \(q=0\) census}
\label{sec:V39-q0}

Fix a one-strong \(q=0\) sequence with maximal row \(a\) and
\[
 \Xi_a=M\longrightarrow\infty.
\]
Let \(J\) contain every physical coordinate excluded from the usable
slot set because it carries a recursive joining occurrence, the
unique payer, the selected local output role, or the fixed strong
coordinate role.  Recursive first connectivity and the one-payer
ledger give
\begin{equation}
\label{eq:V39-J-bound}
 |J|\le C_5
\end{equation}
for an absolute row-order-five constant.  Let
\[
 D:=\operatorname{supp}(a)\setminus J.
\]
In each term of the fixed row-partition expansion, \(D\) is an
unmarked complete moment bank.

\begin{lemma}[\(q=0\) mass alternative]
\label{lem:V39-q0-alternative}
After passage to a subsequence, one of the following holds:
\begin{enumerate}[label=\textup{(\roman*)},leftmargin=3em]
\item there are unmarked banks \(D_M\subseteq D\) and
      \(\sigma>0\) such that
\[
 \sum_{e\in D_M}a_{a,e}\ge\sigma M,
 \qquad
 \max_{e\in D_M}{a_{a,e}\over M}\longrightarrow0;
\]
\item some role-owned coordinate \(e_*\in J\) satisfies
\[
 a_{a,e_*}\ge\sigma M
\]
for a fixed \(\sigma>0\).
\end{enumerate}
\end{lemma}

\begin{proof}
Because \(q=0\), no coordinate in the unmarked set \(D\) carries a
fixed positive fraction of \(M\); otherwise it would be a usable
cofinal slot.  Hence
\[
 \max_{e\in D}{a_{a,e}\over M}\longrightarrow0
\]
along a persistent subsequence.

If \(\sum_{e\in D}a_{a,e}\) has a fixed positive fraction of \(M\),
take \(D_M=D\) and obtain \textup{(i)}.  Otherwise the finite set
\(J\) carries a fixed positive fraction of \(M\).  By
\eqref{eq:V39-J-bound} and pigeonhole, one of its coordinates carries
at least that fraction divided by \(C_5\).  Passing to a fixed role
coordinate subsequence gives \textup{(ii)}.
\end{proof}

\begin{theorem}[Fixed-time balanced-Fourier closure of \(q=0\)]
\label{thm:V39-fixed-q0-close}
For \(SU(2)\), fixed \(s,\rho>0\), every one-strong \(q=0\) quintic
core satisfies the direct \(M^{-1}\) response in the balanced-Fourier
interface.
\end{theorem}

\begin{proof}
Apply \Cref{lem:V39-q0-alternative}.  In alternative \textup{(i)}, the
diffuse-bank theorem indexed by
\texttt{thm:V87-diffuse-bank-power} in archive f3, with \(k=2\),
gives the complete bank response
\[
 C_{\rho,s,\sigma}M^{-1}.
\]
The bank is unmarked and disjoint from the payer, so the V32 positive
exact-dual sandwich and exterior erasure apply exactly as in the
\(k=1\) proof of \Cref{thm:V34-diffuse-close}.

In alternative \textup{(ii)}, expand the finitely many joining
cumulants at the row-partition level, retaining the complete local
moment block at \(e_*\).  If the block containing row \(a\) has a
high output, use the heat half of
\eqref{eq:V39-two-sector-envelope}; if it has a low output, use the
height half.  Endpoint-replicated terms are covered by
\Cref{lem:V39-replicated}.  The payer, selected output, recoupling
labels, exact duality, and all fixed-arity sums cost only a polynomial
in the spin \(p_*\), while
\[
 p_*^2\asymp a_{a,e_*}\gtrsim M.
\]
\Cref{cor:V39-fixed-time-local} therefore supplies arbitrary inverse
powers of \(M\), in particular \(M^{-1}\).  The fixed signed
row-partition sum changes only the constant.
\end{proof}

\begin{corollary}[All fixed-time one-strong \(q=0,1\) cores close in
the weighted interface]
\label{cor:V39-fixed-q01}
For \(SU(2)\) and fixed \(s,\rho>0\), no \(q=0\) or \(q=1\)
one-strong obstruction remains in the balanced-Fourier quintic
interface.
\end{corollary}

\begin{proof}
Use \Cref{thm:V39-fixed-q0-close} for \(q=0\).  For \(q=1\), the
second-atom and diffuse alternatives are
\Cref{thm:V34-second-close,thm:V34-diffuse-close}; the saturated
alternative is \Cref{thm:V38-saturated-close}.
\end{proof}

\begin{firewall}[This is a fixed-time theorem]
\label{fire:V39-fixed-time}
The constants in
\Cref{cor:V39-fixed-time-local,thm:V39-fixed-q0-close} depend on
\(s\).  Neither theorem asserts a bound uniform as \(s\downarrow0\).
Accordingly \Cref{cor:V39-fixed-q01} does not yet prove the cofinal
Wilson MTA or WUFC.
\end{firewall}

\section{Sharp failure of uniformity}
\label{sec:V39-nonuniform}

\begin{proposition}[No uniform arbitrary-power heat response]
\label{prop:V39-no-uniform-heat}
For every \(N>0\), there is no constant \(C_N\) such that
\begin{equation}
\label{eq:V39-false-uniform}
 e^{-sp(p+1)}\le C_Nd_p^{-N}
\end{equation}
for all \(0<s\le1\) and all spins \(p\ge1\).
The balanced height quotient does not repair this example when the
output spin equals the input spin.
\end{proposition}

\begin{proof}
Take integer spins \(p_n=n\) and
\[
 s_n={1\over n(n+1)}.
\]
Then the left side of \eqref{eq:V39-false-uniform} is \(e^{-1}\),
while \(d_{p_n}^{-N}=(2n+1)^{-N}\to0\).  For a singleton moment the
output spin is \(K=p_n\), so its exact height quotient is
\[
 e^{\rho(K-p_n)}=1.
\]
Thus no fixed constant can make \eqref{eq:V39-false-uniform} hold.
\end{proof}

\begin{remark}[A payer only strengthens the countertest]
\label{rem:V39-payer}
On the sequence in \Cref{prop:V39-no-uniform-heat}, the payer
\[
 {p(p+1)\over1-e^{-sp(p+1)}}
\]
has order \(p^2\).  Payer absorption remains valid at each fixed
\(s\), but its constant necessarily records the small-time loss.
\end{remark}

\section{The uniform algebraic transition band}
\label{sec:V39-transition}

Let \(D_*\) be a fixed polynomial exponent large enough to dominate
all local dimension, Casimir, recoupling, multiplicity, and payer
costs at row order five.  To obtain the required \(d_p^{-2}\) direct
response uniformly in \(s\), the easy sectors are:
\begin{align}
 H-K
 &\ge {D_*+4\over\rho}\log d_p,
\label{eq:V39-height-easy}\\
 sK(K+1)
 &\ge (D_*+4)\log d_p.
\label{eq:V39-heat-easy}
\end{align}

\begin{lemma}[Uniform closure outside the transition band]
\label{lem:V39-outside-transition}
Every row-order-five local moment term satisfying either
\eqref{eq:V39-height-easy} or
\eqref{eq:V39-heat-easy} has balanced-Fourier response
\[
 O(d_p^{-3})
\]
after all fixed polynomial costs and one payer.  The same holds for
the corresponding finite cumulant sum.
\end{lemma}

\begin{proof}
Under \eqref{eq:V39-height-easy}, the height quotient is at most
\(d_p^{-(D_*+4)}\).  Under \eqref{eq:V39-heat-easy}, the heat
multiplier has the same bound.  Multiply by at most
\(Cd_p^{D_*}\), sum the polynomially many fixed-arity labels, and
retain one spare inverse power.  A local cumulant has a fixed
Bell-number coefficient sum.
\end{proof}

\begin{definition}[Height-near-saturated parabolic band]
\label{def:V39-transition-band}
The uniform transition band consists of channels satisfying both
\begin{equation}
\label{eq:V39-transition-band}
 H-K<
 {D_*+4\over\rho}\log d_p,
 \qquad
 sK(K+1)<(D_*+4)\log d_p.
\end{equation}
\end{definition}

\begin{proposition}[Geometry of a transition channel]
\label{prop:V39-transition-geometry}
If the input block contains spin \(p\) and
\eqref{eq:V39-transition-band} holds, then
\begin{align}
 K&\ge p-C_{\rho,D_*}\log d_p,
\label{eq:V39-K-near-p}\\
 sp^2&\le C_{\rho,D_*}\log d_p
\label{eq:V39-parabolic}
\end{align}
for all sufficiently large \(p\).
\end{proposition}

\begin{proof}
Since \(H\ge p\), the first inequality in
\eqref{eq:V39-transition-band} gives
\[
 K>p-{D_*+4\over\rho}\log d_p,
\]
which is \eqref{eq:V39-K-near-p}.  Hence \(K\ge p/2\) for large \(p\).
The second transition inequality then implies
\[
 {s p^2\over4}
 \le sK^2
 \le sK(K+1)
 <(D_*+4)\log d_p,
\]
proving \eqref{eq:V39-parabolic}.
\end{proof}

\begin{assessment}[Exact cofinal bottleneck after V39]
\label{ass:V39-bottleneck}
The balanced protected sector \(K\ll p\) is not the uniform problem:
its height deficit closes independently of \(s\).  The only local
channels not closed by height or heat are nearly height preserving,
\[
 K=p+O(\log d_p),
\]
at the small-noise parabolic scale
\[
 sp^2=O(\log d_p).
\]
This band includes the sharp singleton sequence of
\Cref{prop:V39-no-uniform-heat}.  It must be treated from the exact
Wilson fusion multiplier and the cumulant/resolvent cancellation,
not by a fixed-\(s\) Gaussian tail.
\end{assessment}

\section{V39 result ledger and V40 audit programme}
\label{sec:V39-ledger}

\begin{theorem}[Fourth-series V39 result ledger]
\label{thm:V39-results}
Relative to V1--V38, V39 proves:
\begin{enumerate}[label=\textup{(V39.\arabic*)},leftmargin=6.3em]
\item the optimized height--heat envelope and its two-sector form;
\item fixed-time arbitrary-power response for every cofinal local
      moment block, including separated cumulant terms;
\item an exhaustive diffuse-bank/role-owned-atom alternative for the
      \(q=0\) core;
\item fixed-time balanced-Fourier closure of all one-strong
      \(q=0,1\) quintic cores;
\item a sharp parabolic counterexample to heat-time-uniform inverse
      powers;
\item uniform algebraic closure outside the
      height-near-saturated parabolic band; and
\item exact localization of the cofinal residual to
      \eqref{eq:V39-transition-band}.
\end{enumerate}
\end{theorem}

\begin{proof}
These are
\Cref{lem:V39-optimized-envelope,
cor:V39-fixed-time-local,
lem:V39-replicated,
lem:V39-q0-alternative,
thm:V39-fixed-q0-close,
cor:V39-fixed-q01,
prop:V39-no-uniform-heat,
lem:V39-outside-transition,
prop:V39-transition-geometry}.
\end{proof}

\begin{programme}[V40: companion audit and exact transition
multiplier]
\label{prog:V39-V40}
\begin{enumerate}[leftmargin=3em]
\item Perform the compulsory read-only audit of the newest active SM
      and NS notes before selecting the transition-band route.
\item Restore the exact Wilson character multiplier
      \(\eta_j(\alpha)\) and reduced-resolvent denominator; do not
      replace them by a fixed heat kernel.
\item On the scaling \(s_\alpha p^2=O(\log d_p)\), compute the complete
      binary and higher joining finite differences before absolute
      values.
\item Separate an unmarked height-preserving moment, a binary
      covariance, and joining arity at least three.  The fixed-height
      small-noise theorem may not be cited uniformly when
      \(p\to\infty\).
\item Determine whether the exact resolvent-normalized cumulant gives
      the missing \(d_p^{-2}\) response; otherwise construct a
      transition-band counterpacket in the full BFA norm.
\item Propagate every closed transition type through the \(q=0\)
      source-role census.
\end{enumerate}
\end{programme}

\begin{firewall}[Claim boundary after fourth-series V39]
\label{fire:V39-claim-boundary}
V39 closes \(q=0,1\) only at fixed heat time and isolates the
small-noise transition band.  It does not prove uniform cofinal
quintic MTA, an all-order atlas, WUFC, CPM, cutoff removal, reflection
positivity, Osterwalder--Schrader reconstruction, construction of the
continuum Yang--Mills measure, or a uniform positive continuum
spectral gap.  The full Yang--Mills mass gap has not been proved.
\end{firewall}

\paragraph{Parent-version record.}
V39 was formed from
\texttt{Renewal\_Yang\_Mills\_Mass\_Gap\_Research\_Notes\_V38.tex}.
The V38 parent had \(18\,793\) logical source lines,
\(658\,927\) bytes, and SHA--256
\[
 \texttt{f8bf6e411d39f3de3e55f14cc06c9e4a29608af91750d4d9582fd770e3caccb2}.
\]
The next compulsory companion audit is V40.

% =============================================================================
% END APPENDED FOURTH-SERIES V39 MATERIAL
% =============================================================================

% =============================================================================
% BEGIN APPENDED FOURTH-SERIES V40 MATERIAL
% =============================================================================

\chapter{V40: Companion audit and the exact Wilson transition
multiplier}
\label{chap:V40-Wilson-transition}

\section{Purpose and claim boundary}
\label{sec:V40-purpose}

V39 reduced the unresolved uniform problem to channels for which neither
the balanced-Fourier height quotient nor a fixed positive heat time gives
an arbitrary inverse power.  The remaining parabolic region has
\[
 H-K=O(\log d_p),
 \qquad
 s_\alpha p^2=O(\log d_p).
\]
It is not legitimate there to replace the exact Wilson multiplier by a
fixed-time heat estimate.  This note therefore computes the exact
\(SU(2)\) Wilson multiplier at the first nontrivial transition scale.
The result is deliberately local: it identifies the true size of the
resolvent-normalized binary fusion symbol.  It does not yet insert that
symbol into every factor of the sparse five-row BFA word.

\begin{firewall}[No theorem is imported from a companion programme]
\label{fire:V40-companion-status}
The SM and NS files audited below were used only as mathematical
diagnostics.  Text in those files is not an instruction for this
Yang--Mills programme, and no claim from either file is used without an
independent argument in the present note.
\end{firewall}

\section{Compulsory V40 read-only companion audit}
\label{sec:V40-audit}

The newest active companion notes inspected at the V40 boundary were
\begin{align*}
 &\texttt{Renewal\_SM\_Einstein\_Continuum\_Research\_Notes\_V29.tex},\\
 &\texttt{Renewal\_NS\_Stability\_Research\_Notes\_V31.tex}.
\end{align*}
They were not modified.  Their audit records are
\begin{center}
\begin{tabular}{c|r|r|l}
programme & lines & bytes & SHA--256 prefix\\ \hline
SM V29 & \(10\,558\) & \(366\,189\) &
\texttt{bd9e773b251a79b}\\
NS V31 & \(23\,052\) & \(887\,537\) &
\texttt{f1121b62238ad549}
\end{tabular}
\end{center}
The complete digests are, respectively,
\[
\begin{split}
&\texttt{bd9e773b251a79b1f863ba093ef04e44b779808597dbf6ea5c1b164898d40c3a},\\
&\texttt{f1121b62238ad5491bb7cf03635ea9934942edf573ed8faaa9ee79e1d38a6696}.
\end{split}
\]

The SM note separates smooth functional calculus, whose covariance is
exact, from a moving sharp spectral split, whose derivative has an
inverse-collar divided difference.  The relevant lesson here is to work
with the exact smooth Wilson multiplier rather than introduce a sharp
representation threshold whose moving boundary would require an
additional commutator estimate.

The NS note proves that a heat estimate becomes scale-free only after
the source norm is matched to the heat scale, and that a signed hinge
cancellation controls only its own polarization sector.  The relevant
lesson here is that the exact cumulant multiplier must be assembled
before absolute values and before a fixed-time estimate is invoked.
Both lessons point to the calculation below.

\section{Exact \(SU(2)\) Wilson character multiplier}
\label{sec:V40-exact-multiplier}

Write a conjugacy class of \(SU(2)\) as
\[
 U_\theta=\operatorname{diag}(e^{i\theta},e^{-i\theta}),
 \qquad 0\leq\theta\leq\pi.
\]
The spin-\(j\) irreducible representation has
\[
 d_j=2j+1,
 \qquad
 \chi_j(\theta)={\sin(d_j\theta)\over\sin\theta},
 \qquad
 c_j=j(j+1).
\]
For \(\alpha>0\), let
\begin{equation}
\label{eq:V40-Wilson-law}
 d\nu_\alpha(U)
 =
 Z_\alpha^{-1}
 \exp\!\left\{\alpha {1\over2}\operatorname{ReTr}U\right\}\,dU.
\end{equation}
Thus the class density is proportional to
\(e^{\alpha\cos\theta}\sin^2\theta\,d\theta\).  Its normalized central
Fourier multiplier is
\begin{equation}
\label{eq:V40-eta-def}
 \eta_j(\alpha)
 =
 {1\over d_j}\int_{SU(2)}\chi_j(U)\,d\nu_\alpha(U).
\end{equation}

\begin{lemma}[Bessel formula for the exact Wilson multiplier]
\label{lem:V40-Bessel-formula}
For every \(j\in\frac12\mathbb N_0\),
\begin{equation}
\label{eq:V40-Bessel-ratio}
 \eta_j(\alpha)={I_{2j+1}(\alpha)\over I_1(\alpha)},
\end{equation}
where \(I_n\) is the modified Bessel function.
\end{lemma}

\begin{proof}
For integral \(n\geq0\),
\[
 I_n(\alpha)
 ={1\over\pi}\int_0^\pi e^{\alpha\cos\theta}
 \cos(n\theta)\,d\theta.
\]
Haar's class formula and
\(\sin(n\theta)\sin\theta
=\frac12\{\cos((n-1)\theta)-\cos((n+1)\theta)\}\) give
\[
 {2\over\pi}\int_0^\pi e^{\alpha\cos\theta}\sin^2\theta\,d\theta
 =I_0(\alpha)-I_2(\alpha)
 ={2\over\alpha}I_1(\alpha),
\]
and, with \(n=2j+1\),
\[
 {2\over\pi n}\int_0^\pi
 e^{\alpha\cos\theta}\sin(n\theta)\sin\theta\,d\theta
 ={I_{n-1}(\alpha)-I_{n+1}(\alpha)\over n}
 ={2\over\alpha}I_n(\alpha).
\]
The last equalities use
\(I_{n-1}-I_{n+1}=2nI_n/\alpha\).
Dividing the second display by the first proves
\eqref{eq:V40-Bessel-ratio}.
\end{proof}

\begin{corollary}[Wilson-to-heat calibration]
\label{cor:V40-calibration}
If
\[
 s_\alpha
 ={-\log\eta_{1/2}(\alpha)\over c_{1/2}},
\]
then
\begin{equation}
\label{eq:V40-calibration}
 s_\alpha={2\over\alpha}+O(\alpha^{-2}).
\end{equation}
\end{corollary}

\begin{proof}
The fixed-index estimate proved in
\Cref{lem:V40-transition-Bessel} below gives
\[
 1-\eta_{1/2}(\alpha)
 ={2c_{1/2}\over\alpha}+o(\alpha^{-1}).
\]
Since \(-\log(1-t)=t+O(t^2)\), division by \(c_{1/2}=3/4\)
gives \eqref{eq:V40-calibration}.  The usual fixed-index Laplace
expansion improves the displayed remainder from \(o(\alpha^{-1})\)
to \(O(\alpha^{-2})\); for completeness, this follows by retaining
the \(y^4/\alpha\) term in the proof below, with an
\(e^{-y^2/4}(1+y^8)\) majorant.
\end{proof}

\section{A uniform Laplace lemma at \(j\asymp\sqrt\alpha\)}
\label{sec:V40-Laplace}

\begin{lemma}[Transition Bessel ratios and fixed shifts]
\label{lem:V40-transition-Bessel}
Let \(n_\alpha\) be positive integers such that
\[
 {n_\alpha\over\sqrt\alpha}\longrightarrow x\in(0,\infty).
\]
For every fixed integer \(h\),
\begin{align}
 {I_{n_\alpha}(\alpha)\over I_1(\alpha)}
 &\longrightarrow e^{-x^2/2},
\label{eq:V40-Bessel-profile}\\
 \sqrt\alpha\,
 {I_{n_\alpha+h}(\alpha)-I_{n_\alpha}(\alpha)\over I_1(\alpha)}
 &\longrightarrow
 -hxe^{-x^2/2}.
\label{eq:V40-Bessel-shift}
\end{align}
For every fixed positive integer \(m\),
\begin{equation}
\label{eq:V40-fixed-index}
 \alpha\left(1-{I_m(\alpha)\over I_1(\alpha)}\right)
 \longrightarrow {m^2-1\over2}.
\end{equation}
\end{lemma}

\begin{proof}
Put \(\theta=y/\sqrt\alpha\) in the integral formula for \(I_n\).
After multiplication by \(e^{-\alpha}\sqrt\alpha\), its integrand is
\[
 {\bf1}_{[0,\pi\sqrt\alpha]}(y)
 \exp\{\alpha(\cos(y/\sqrt\alpha)-1)\}
 \cos(ny/\sqrt\alpha).
\]
On every fixed \(y\)-interval this converges to
\(e^{-y^2/2}\cos(xy)\) when
\(n/\sqrt\alpha\to x\).  Choose a fixed
\(\varepsilon\in(0,\pi/2)\).  On
\(0\leq y\leq\varepsilon\sqrt\alpha\), the inequality
\(1-\cos(y/\sqrt\alpha)\geq c_\varepsilon y^2/\alpha\)
provides a Gaussian majorant.  The complementary part of the
\(\theta\)-integral is \(O(e^{\alpha\cos\varepsilon})\), exponentially
smaller than \(e^\alpha/\sqrt\alpha\).  Dominated convergence and the
Gaussian Fourier transform therefore give
\[
 e^{-\alpha}\sqrt\alpha\,I_{n_\alpha}(\alpha)
 \longrightarrow
 {1\over\pi}\int_0^\infty e^{-y^2/2}\cos(xy)\,dy
 ={e^{-x^2/2}\over\sqrt{2\pi}}.
\]
For \(I_1\), the same limit is \(1/\sqrt{2\pi}\).  This proves
\eqref{eq:V40-Bessel-profile}.

For the fixed shift, the mean-value theorem gives the uniform
majorant
\[
 \sqrt\alpha\,
 \left|
 \cos((n_\alpha+h)y/\sqrt\alpha)
 -\cos(n_\alpha y/\sqrt\alpha)
 \right|
 \leq |h|y.
\]
The expression inside the absolute value, after multiplication by
\(\sqrt\alpha\), converges pointwise to
\(-hy\sin(xy)\).  The same localization and dominated-convergence
argument yields
\[
 {-\frac h\pi\int_0^\infty
 ye^{-y^2/2}\sin(xy)\,dy
 \over
 \frac1\pi\int_0^\infty e^{-y^2/2}\,dy}
 =-hxe^{-x^2/2},
\]
which is \eqref{eq:V40-Bessel-shift}.

Finally,
\[
 \alpha\{\cos(y/\sqrt\alpha)-\cos(my/\sqrt\alpha)\}
 \longrightarrow {m^2-1\over2}y^2.
\]
Taylor's theorem on the localized interval gives an integrable
majorant \(C_m(1+y^4)e^{-c_\varepsilon y^2}\); the exterior remains
exponentially small.  Division by the \(I_1\) asymptotic gives
\[
 {m^2-1\over2}
 {\int_0^\infty y^2e^{-y^2/2}\,dy
  \over
  \int_0^\infty e^{-y^2/2}\,dy}
 ={m^2-1\over2},
\]
proving \eqref{eq:V40-fixed-index}.
\end{proof}

\begin{corollary}[Exact Wilson transition profile]
\label{cor:V40-Wilson-profile}
If \(j_\alpha\to\infty\) and
\[
 {d_{j_\alpha}\over\sqrt\alpha}\longrightarrow x\in(0,\infty),
\]
then
\[
 \eta_{j_\alpha}(\alpha)\longrightarrow e^{-x^2/2},
 \qquad
 s_\alpha c_{j_\alpha}\longrightarrow {x^2\over2}.
\]
\end{corollary}

\begin{proof}
The first assertion is \eqref{eq:V40-Bessel-profile}.  Since
\[
 c_j={d_j^2-1\over4},
\]
the second follows from \eqref{eq:V40-calibration}.
\end{proof}

\section{The resolvent-normalized binary fusion symbol}
\label{sec:V40-binary-symbol}

Whenever \(K\) occurs in \(p\otimes r\), define the exact scalar
binary fusion quotient
\begin{equation}
\label{eq:V40-binary-symbol}
 {\mathfrak B}_\alpha(p,r;K)
 =
 {\eta_K(\alpha)-\eta_p(\alpha)\eta_r(\alpha)
  \over
  1-\eta_K(\alpha)}.
\end{equation}
This is the coefficient produced by first forming the two-input
connected difference and then applying the reduced Wilson resolvent
on the nontrivial output \(K\).  Thus the cancellation in the
numerator has already occurred; estimating the two moment terms
separately would lose information.

\begin{theorem}[One cofinal input and a fixed fusion partner]
\label{thm:V40-one-cofinal}
Fix \(r\in\frac12\mathbb N\).  Let \(p_\alpha\to\infty\) with
\[
 {d_{p_\alpha}\over\sqrt\alpha}\longrightarrow x\in(0,\infty).
\]
Fix a fusion displacement
\(\delta\in\{-r,-r+1,\ldots,r\}\) such that
\(K_\alpha=p_\alpha+\delta\) is admissible.  Put
\(a=e^{-x^2/2}\).  If \(\delta\ne0\), then
\begin{align}
 \sqrt\alpha\,
 {\mathfrak B}_\alpha(p_\alpha,r;K_\alpha)
 &\longrightarrow
 {-2\delta xa\over1-a},
\label{eq:V40-one-cofinal-sqrt}\\
 d_{p_\alpha}\,
 {\mathfrak B}_\alpha(p_\alpha,r;K_\alpha)
 &\longrightarrow
 {-2\delta x^2a\over1-a}.
\label{eq:V40-one-cofinal-dim}
\end{align}
If \(\delta=0\), then
\begin{align}
 \alpha\,
 {\mathfrak B}_\alpha(p_\alpha,r;p_\alpha)
 &\longrightarrow
 {2c_r a\over1-a},
\label{eq:V40-central-alpha}\\
 d_{p_\alpha}^{\,2}
 {\mathfrak B}_\alpha(p_\alpha,r;p_\alpha)
 &\longrightarrow
 {2c_r x^2a\over1-a}.
\label{eq:V40-central-dim}
\end{align}
\end{theorem}

\begin{proof}
Let \(n_\alpha=d_{p_\alpha}\).  Since
\(d_{K_\alpha}=n_\alpha+2\delta\),
\Cref{lem:V40-transition-Bessel} with \(h=2\delta\) gives
\[
 \sqrt\alpha\{\eta_{K_\alpha}-\eta_{p_\alpha}\}
 \longrightarrow -2\delta xa.
\]
For the fixed representation \(r\), take \(m=d_r\) in
\eqref{eq:V40-fixed-index}.  Since \(m^2-1=4c_r\),
\[
 \alpha(1-\eta_r)\longrightarrow 2c_r.
\]
Consequently
\[
 \sqrt\alpha\,
 \eta_{p_\alpha}(1-\eta_r)\longrightarrow0.
\]
When \(\delta\ne0\), decompose the numerator as
\[
 \eta_{K_\alpha}-\eta_{p_\alpha}\eta_r
 =
 \{\eta_{K_\alpha}-\eta_{p_\alpha}\}
 +\eta_{p_\alpha}(1-\eta_r).
\]
The first term has the stated nonzero
\(\alpha^{-1/2}\) limit and the second is lower order.  Moreover
\(1-\eta_{K_\alpha}\to1-a\).  This proves
\eqref{eq:V40-one-cofinal-sqrt}; multiplying by
\(d_{p_\alpha}/\sqrt\alpha\to x\) proves
\eqref{eq:V40-one-cofinal-dim}.

For \(\delta=0\), the numerator is exactly
\(\eta_{p_\alpha}(1-\eta_r)\).  Its product with \(\alpha\) tends to
\(2c_ra\), while the denominator tends to \(1-a\).  This proves
\eqref{eq:V40-central-alpha}, and multiplication by
\(d_{p_\alpha}^2/\alpha\to x^2\) gives
\eqref{eq:V40-central-dim}.
\end{proof}

\begin{corollary}[Sharp half-power loss in shifted channels]
\label{cor:V40-half-power}
For every fixed admissible \(\delta\ne0\), the exact Wilson quotient
in \eqref{eq:V40-binary-symbol} is of exact order
\[
 d_{p_\alpha}^{-1}
\]
on the parabolic scale.  It is not
\(O(d_{p_\alpha}^{-2})\).  The undisplaced channel \(\delta=0\) is of
order \(d_{p_\alpha}^{-2}\).
\end{corollary}

\begin{proof}
The limits in
\eqref{eq:V40-one-cofinal-dim} and
\eqref{eq:V40-central-dim} are finite, and the former is nonzero
when \(\delta\ne0\).
\end{proof}

\section{Two cofinal inputs: the height-saturated Cartan endpoint}
\label{sec:V40-Cartan}

\begin{theorem}[Nondecaying Cartan fusion quotient]
\label{thm:V40-Cartan-profile}
Let \(p_\alpha\to\infty\) satisfy
\[
 {d_{p_\alpha}\over\sqrt\alpha}\longrightarrow x\in(0,\infty).
\]
The highest-spin representation \(K_\alpha=2p_\alpha\) occurs in
\(p_\alpha\otimes p_\alpha\), and
\begin{equation}
\label{eq:V40-Cartan-limit}
 {\mathfrak B}_\alpha
 (p_\alpha,p_\alpha;2p_\alpha)
 \longrightarrow
 -{e^{-x^2}\over1+e^{-x^2}}.
\end{equation}
The limit is nonzero.
\end{theorem}

\begin{proof}
The output dimension satisfies
\[
 {d_{2p_\alpha}\over\sqrt\alpha}
 ={4p_\alpha+1\over\sqrt\alpha}\longrightarrow2x.
\]
Hence \Cref{cor:V40-Wilson-profile} gives
\[
 \eta_{2p_\alpha}\longrightarrow e^{-2x^2},
 \qquad
 \eta_{p_\alpha}^2\longrightarrow e^{-x^2}.
\]
Substitution in \eqref{eq:V40-binary-symbol} gives
\[
 {e^{-2x^2}-e^{-x^2}\over1-e^{-2x^2}}
 =-{e^{-x^2}\over1+e^{-x^2}},
\]
as claimed.
\end{proof}

\begin{proposition}[Height weights do not remove the Cartan endpoint]
\label{prop:V40-height-saturation}
For the Cartan channel \(p\otimes p\to2p\), the balanced exponential
quotient is exactly
\[
 \exp\{\rho(\operatorname{ht}(2p)
 -\operatorname{ht}(p)-\operatorname{ht}(p))\}=1.
\]
Thus neither the exact Wilson quotient nor the exponential height
weight supplies a vanishing factor in
\Cref{thm:V40-Cartan-profile}.
\end{proposition}

\begin{proof}
For \(SU(2)\), \(\operatorname{ht}(j)=j\) in the normalization used
throughout the fourth series.  The exponent therefore vanishes.
\end{proof}

\section{What the transition calculation proves, and what it does not}
\label{sec:V40-consequences}

\begin{theorem}[Failure of a representation-uniform local
\(d_p^{-2}\) binary card]
\label{thm:V40-local-card-no-go}
There is no constant \(C\), independent of \(\alpha,p,r,K\), for
which every admissible transition channel satisfies
\[
 |{\mathfrak B}_\alpha(p,r;K)|
 \leq C d_p^{-2}.
\]
There is also no representation-uniform \(o(1)\) estimate for the
two-cofinal Cartan channel.
\end{theorem}

\begin{proof}
The first assertion is contradicted by any fixed displaced channel
in \Cref{cor:V40-half-power}.  The second is contradicted by the
nonzero limit in \eqref{eq:V40-Cartan-limit}.
\end{proof}

\begin{firewall}[This is not yet a counterexample to quintic MTA]
\label{fire:V40-not-MTA-counterexample}
\Cref{thm:V40-local-card-no-go} concerns the scalar one-link fusion
quotient after its cumulant and resolvent cancellation.  The full BFA
word also contains output dimension and Casimir weights, input
normalizations, Clebsch--Gordan trace-norm factors, the payer, tree
marks at the remaining coordinates, and the final output dual.
Those factors can change the power ledger.  Therefore the theorem
invalidates a proposed local \(d_p^{-2}\) shortcut but does not, by
itself, disprove or prove the full quintic MTA.
\end{firewall}

\begin{assessment}[Direction correction after the V40 audit]
\label{ass:V40-direction}
The V39 fixed-time closure cannot be promoted uniformly by continuity
in \(s\).  The exact Wilson law has three distinct transition
behaviours:
\begin{enumerate}[leftmargin=3em]
\item an undisplaced one-cofinal channel has the desired
      \(d_p^{-2}\) binary response;
\item a fixed displaced channel has only \(d_p^{-1}\) response; and
\item a two-cofinal height-saturated Cartan channel has no scalar
      decay at all.
\end{enumerate}
Accordingly, the next step must be a full BFA power ledger for the
sparse five-row word, performed before any factor is discarded.  A
generic appeal to fixed-order Gaussian excess would be circular,
because its constants were proved only on fixed Peter--Weyl banks.
\end{assessment}

\section{V40 result ledger and V41 programme}
\label{sec:V40-ledger}

\begin{theorem}[Fourth-series V40 result ledger]
\label{thm:V40-results}
Relative to V1--V39, V40 proves:
\begin{enumerate}[label=\textup{(V40.\arabic*)},leftmargin=6.3em]
\item the exact Bessel-ratio formula for every \(SU(2)\) Wilson
      character multiplier;
\item a self-contained uniform Laplace theorem at
      \(d_j\asymp\sqrt\alpha\), including fixed representation shifts;
\item the exact calibration \(s_\alpha=2/\alpha+O(\alpha^{-2})\);
\item an exact \(d_p^{-1}\) asymptotic for every nonzero fixed fusion
      displacement and an exact \(d_p^{-2}\) asymptotic for the
      undisplaced channel;
\item a nonzero transition limit for the two-cofinal Cartan fusion
      quotient;
\item failure of a representation-uniform local
      \(d_p^{-2}\) binary shortcut; and
\item localization of the next decision to the complete BFA power
      ledger, rather than another heat envelope.
\end{enumerate}
\end{theorem}

\begin{proof}
These are
\Cref{lem:V40-Bessel-formula,
lem:V40-transition-Bessel,
cor:V40-calibration,
thm:V40-one-cofinal,
cor:V40-half-power,
thm:V40-Cartan-profile,
thm:V40-local-card-no-go,
fire:V40-not-MTA-counterexample}.
\end{proof}

\begin{programme}[V41: insert the transition symbol into the exact
five-row BFA word]
\label{prog:V40-V41}
\begin{enumerate}[leftmargin=3em]
\item Return to the four-coordinate sparse tree
      \(\{ab,ad,ce,ac\}\) of V37, but put the root pair in the
      transition scaling \(d_p\asymp\sqrt\alpha\).
\item Write one exact Fourier coefficient from inputs through the
      connected Wilson difference, reduced resolvent, payer, and
      output BFA weight.  Keep all dimension and Casimir ratios.
\item Evaluate the Cartan Clebsch--Gordan coefficient using
      highest-weight vectors, and evaluate its trace-norm contribution
      rather than bounding it by one prematurely.
\item Compare the resulting normalized coefficient with the exact
      four tree marks.  If it exceeds the marked right-hand side,
      promote the packet to a full MTA counterexample; if not, record
      the compensating factor as the missing transition card.
\item Repeat the ledger for one-cofinal displaced channels, where
      \Cref{cor:V40-half-power} leaves exactly one inverse dimension
      to be found.
\item Only after these exact ledgers, propagate the result through
      the \(q=0\) role census.
\end{enumerate}
\end{programme}

\begin{firewall}[Claim boundary after fourth-series V40]
\label{fire:V40-claim-boundary}
V40 rigorously determines the exact \(SU(2)\) Wilson transition
profile for the binary resolvent-normalized fusion scalar.  It does
not close the complete transition-band five-row BFA estimate, prove a
uniform quintic MTA, an all-order atlas, WUFC, CPM, cutoff removal,
reflection positivity, Osterwalder--Schrader reconstruction,
construction of the continuum Yang--Mills measure, or a uniform
positive continuum spectral gap.  The full Yang--Mills mass gap has
not been proved.
\end{firewall}

\paragraph{Parent-version record.}
V40 was formed from
\texttt{Renewal\_Yang\_Mills\_Mass\_Gap\_Research\_Notes\_V39.tex}.
The V39 parent had \(19\,274\) logical source lines,
\(674\,837\) bytes, and SHA--256
\[
 \texttt{b2ed2ca8ee0550c9f828b1629447688e8cea11001011600ee375414d33daa5bd}.
\]
The next compulsory companion audit is V45.

% =============================================================================
% END APPENDED FOURTH-SERIES V40 MATERIAL
% =============================================================================

% =============================================================================
% BEGIN APPENDED FOURTH-SERIES V41 MATERIAL
% =============================================================================

\chapter{V41: Exact BFA normalization closes the sparse Wilson
transition packet}
\label{chap:V41-sparse-transition}

\section{Purpose and nonduplication}
\label{sec:V41-purpose}

V40 found that the scalar Wilson fusion quotient does not decay in
the two-cofinal Cartan channel.  Taken alone, that fact suggests a
transition counterpacket based on the V37 four-coordinate tree.  Such
a conclusion would be premature for two separate reasons.

First, archive f1 V120--V125 already studied square-root Wilson
displacements, and f1 V126 showed that its particular off-cut family
is paid more efficiently by retaining an old-pair maximum-square
capacity.  We do not repeat that off-cut analysis.  The V37 support
has a different geometry: it has only one cofinal pair, at the final
joining coordinate.

Second, the three fixed-spin prefix connections of the V37 word are
themselves connected Wilson differences.  At fixed heat time they
were harmless constants; on the cofinal Wilson clock
\(\alpha\asymp d_j^2\), each is \(O(\alpha^{-1})\).  The present note
keeps those three factors through the exact Peter--Weyl
normalization.  Their product more than pays the nondecaying Cartan
root.

\begin{assessment}[Upstream correction of the V40 programme]
\label{ass:V41-upstream-correction}
The complete BFA ledger is required before selecting a single hot
coordinate as the obstruction.  A transition counterpacket must keep
every other connected coordinate hot enough to avoid its own
small-noise debit.  The original V37 skeleton does not do so.
\end{assessment}

\section{The parabolic corridor and fixed-prefix differences}
\label{sec:V41-prefix}

Retain the V37 notation
\[
 A_j=1+j(j+1),
 \qquad
 B=1+\ell(\ell+1),
\]
with fixed nontrivial
\(\ell\in\frac12\mathbb N\).  Work in a compact parabolic corridor:
\begin{equation}
\label{eq:V41-corridor}
 0<x_-\le {d_j\over\sqrt\alpha}\le x_+<\infty.
\end{equation}
Then
\begin{equation}
\label{eq:V41-alpha-A}
 c_{x_\pm}A_j\le\alpha\le C_{x_\pm}A_j.
\end{equation}

For an admissible output \(Q\subset\ell\otimes\ell\), put
\begin{equation}
\label{eq:V41-prefix-difference}
 \Delta_{\alpha,\ell}(Q)
 :=
 \eta_Q(\alpha)-\eta_\ell(\alpha)^2.
\end{equation}

\begin{lemma}[Each fixed prefix covariance pays one Wilson clock]
\label{lem:V41-fixed-prefix}
There are \(C_\ell<\infty\) and \(\alpha_\ell\) such that, for every
\(\alpha\ge\alpha_\ell\) and every
\(Q\in\{0,1,\ldots,2\ell\}\) occurring in
\(\ell\otimes\ell\),
\begin{equation}
\label{eq:V41-fixed-prefix-bound}
 |\Delta_{\alpha,\ell}(Q)|
 \le {C_\ell\over\alpha}.
\end{equation}
\end{lemma}

\begin{proof}
For each fixed spin \(r\), the exact Bessel formula and
\eqref{eq:V40-fixed-index} give
\[
 \eta_r(\alpha)
 =
 1-{2r(r+1)\over\alpha}+o_r(\alpha^{-1}).
\]
Hence
\[
 \eta_Q-\eta_\ell^2
 =
 {4\ell(\ell+1)-2Q(Q+1)\over\alpha}
 +o_{\ell,Q}(\alpha^{-1}).
\]
Only finitely many \(Q\)'s occur for fixed \(\ell\).  Enlarge the
constant to make the estimate simultaneous over that finite set and
over the remaining compact interval
\([\alpha_\ell,2\alpha_\ell]\).
\end{proof}

\begin{remark}[Connectedness is the source of the factor]
\label{rem:V41-connected-prefix}
The estimate is for the assembled difference
\(\eta_Q-\eta_\ell^2\).  Either moment term separately tends to one
and supplies no \(\alpha^{-1}\).  Thus the three prefix gains below
would disappear if the moment--cumulant subtraction were normed
before assembly.
\end{remark}

\section{Exact resolved word on the V37 support}
\label{sec:V41-word}

Use the four coordinates and row supports
\[
 u:\{a,b\},\qquad
 v:\{a,d\},\qquad
 w:\{c,e\},\qquad
 r:\{a,c\}
\]
from \Cref{lem:V37-legality}.  The input spins at \(u,v,w\) are
\(\ell,\ell\), and those at \(r\) are \(j,j\).
Let
\[
 \boldsymbol Q=(Q_u,Q_v,Q_w,K)
\]
be a complete output label, with
\[
 Q_u,Q_v,Q_w\subset\ell\otimes\ell,
 \qquad K\subset j\otimes j.
\]
Trivial \(Q\)'s are permitted; trivial coordinates are removed from
the exact output support.  Let
\begin{align}
 q_{\alpha,\boldsymbol Q}
 &:=
 \eta_{Q_u}\eta_{Q_v}\eta_{Q_w}\eta_K,
\label{eq:V41-output-multiplier}\\
 C_E(\boldsymbol Q)
 &:=
 c_{Q_u}+c_{Q_v}+c_{Q_w}+c_K,
\label{eq:V41-output-Casimir}\\
 \Xi(\boldsymbol Q)
 &:=
 |\operatorname{supp}\boldsymbol Q|+C_E(\boldsymbol Q).
\label{eq:V41-output-mass}
\end{align}
All multipliers in this chapter are evaluated at \(\alpha\).

\begin{lemma}[Unique resolved scalar multiplier]
\label{lem:V41-resolved-multiplier}
For the coefficient-one first-connectivity word of V37, the complete
resolved scalar before Peter--Weyl compression is
\begin{equation}
\label{eq:V41-resolved-scalar}
 {\bf1}_{\{\boldsymbol Q\ne\boldsymbol0\}}
 {1\over1-q_{\alpha,\boldsymbol Q}}
 \prod_{z\in\{u,v,w\}}
 \Delta_{\alpha,\ell}(Q_z)
 \bigl(\eta_K-\eta_j^2\bigr).
\end{equation}
\end{lemma}

\begin{proof}
Immediately below \(r\), the connected components are
\(\{a,b,d\}\) and \(\{c,e\}\).  At \(u,v,w\), the recursive
first-connectivity expansion therefore contains the three binary
connected differences associated with the forced prefix edges
\(ab,ad,ce\).  Centrality diagonalizes each one on its local output
spin and gives \eqref{eq:V41-prefix-difference}.  At \(r\), the two
components join through the binary \(ac\) difference
\(\eta_K-\eta_j^2\).  Independence of the four Wilson links
multiplies these four local scalars.

The final reduced resolvent acts on the product output by
\((1-q_{\alpha,\boldsymbol Q})^{-1}\), and \(Q_0\) deletes precisely
the completely trivial output.  The support graph is already the
tree \(T_*\), so
\Cref{lem:V37-legality} excludes a second connected word.
\end{proof}

\section{Uniform payment of the final BFA mass and resolvent}
\label{sec:V41-payer}

We use the representation-uniform product gap proved in archive f1
\texttt{thm:V30-two-sided-gap}.  Its proof uses only compact-group
Dirichlet geometry and therefore applies to the present \(SU(2)\)
Wilson law:
\begin{equation}
\label{eq:V41-product-gap}
 1-q_{\alpha,\boldsymbol Q}
 \ge
 c\min\{s_\alpha C_E(\boldsymbol Q),1\}
\end{equation}
for every nontrivial product output.  The clock
\eqref{eq:V40-calibration} gives
\(s_\alpha\asymp\alpha^{-1}\).

\begin{lemma}[One final payer costs at most one Wilson clock]
\label{lem:V41-one-clock-payer}
Under \eqref{eq:V41-corridor}, uniformly over every nontrivial output
\(\boldsymbol Q\),
\begin{equation}
\label{eq:V41-payer-bound}
 {\Xi(\boldsymbol Q)\over1-q_{\alpha,\boldsymbol Q}}
 |\eta_K-\eta_j^2|
 \le C_{\ell,x_\pm}\alpha.
\end{equation}
\end{lemma}

\begin{proof}
Because \(K\le2j\) and the three \(Q_z\)'s belong to a fixed finite
set,
\[
 0<C_E(\boldsymbol Q)\le C_\ell(1+j^2)\le C_{\ell,x_\pm}\alpha.
\]
The nontrivial \(SU(2)\) Casimir has a positive group-dependent
floor, so
\[
 \Xi(\boldsymbol Q)
 \le C\{1+C_E(\boldsymbol Q)\}
 \le C C_E(\boldsymbol Q).
\]
From \eqref{eq:V41-product-gap} and
\(C_E(\boldsymbol Q)\le C\alpha\),
\[
 {\Xi(\boldsymbol Q)\over1-q_{\alpha,\boldsymbol Q}}
 \le
 C\max\!\left\{
 {\alpha\Xi(\boldsymbol Q)\over C_E(\boldsymbol Q)},
 \Xi(\boldsymbol Q)
 \right\}
 \le C_{\ell,x_\pm}\alpha.
\]
Finally \(|\eta_K-\eta_j^2|\le2\).
\end{proof}

\begin{remark}[The Cartan scalar is included]
\label{rem:V41-Cartan-included}
At \(K=2j\), the quotient
\((\eta_K-\eta_j^2)/(1-\eta_K)\) has the nonzero V40 limit and
\(\Xi(\boldsymbol Q)\asymp A_j\).  Its full payer is therefore of
order \(A_j\asymp\alpha\), exactly the worst scale allowed by
\eqref{eq:V41-payer-bound}.  No Cartan decay has been inserted
implicitly.
\end{remark}

\section{Peter--Weyl compression and the complete output sum}
\label{sec:V41-pinching}

Choose one pure Peter--Weyl block in each input row and normalize it
to have unmarked BFA norm one.  For row \(i\), write its trace-class
coefficient as \(A_i\).  Thus
\begin{equation}
\label{eq:V41-input-normalization}
 e^{\rho h_i}\Xi_i d_i\|A_i\|_1=1,
\end{equation}
where \(h_i,\Xi_i,d_i\) are its total input height, balanced mass, and
representation dimension.

For a full product output \(\boldsymbol Q\), let
\(\omega_{\boldsymbol Q,\gamma}\) be the raw isotypic pinching mass,
including a multiplicity label \(\gamma\), after division by
\(\prod_i\|A_i\|_1\).  The exact product normalization of archive f1
\texttt{thm:V46-exact-product-normalization} gives
\begin{equation}
\label{eq:V41-pinching-probability}
 \omega_{\boldsymbol Q,\gamma}\ge0,
 \qquad
 \sum_{\boldsymbol Q,\gamma}
 \omega_{\boldsymbol Q,\gamma}\le1.
\end{equation}
The factor \(d_{\boldsymbol Q}\) in the output BFA norm cancels the
exact Fourier coefficient factor
\(\prod_i d_i/d_{\boldsymbol Q}\).  Consequently no output dimension
sum remains outside \eqref{eq:V41-pinching-probability}.

\begin{theorem}[Complete transition estimate on the sparse fibre]
\label{thm:V41-sparse-transition}
Let \(\mathcal W_{\alpha,j}^{\rm sp}\) be the complete connected V37
word for the exact \(SU(2)\) Wilson law, including all physical
outputs, multiplicities, the single final reduced resolvent, \(Q_0\),
and its output BFA norm.  Normalize each of its five pure input
blocks to have BFA norm one.  Under
\eqref{eq:V41-corridor},
\begin{equation}
\label{eq:V41-sparse-bound}
 \|\mathcal W_{\alpha,j}^{\rm sp}\|_{\rho,2}^{\rm BFA}
 \le
 {C_{\rho,\ell,x_\pm}\alpha^{-2}
  \over
  \Xi_a\Xi_cB^3}
 \le C_{\rho,\ell,x_\pm}A_j^{-4}.
\end{equation}
\end{theorem}

\begin{proof}
For each output, exact height subadditivity gives
\[
 e^{\rho(h_{\rm out}-\sum_i h_i)}\le1.
\]
Insert \eqref{eq:V41-input-normalization} into the exact
Peter--Weyl product formula.  After the dimension cancellation, the
normalized contribution of
\((\boldsymbol Q,\gamma)\) is at most
\[
 {1\over\Xi_a\Xi_cB^3}\,
 \omega_{\boldsymbol Q,\gamma}
 {\Xi(\boldsymbol Q)\over1-q_{\alpha,\boldsymbol Q}}
 |\eta_K-\eta_j^2|
 \prod_{z\in\{u,v,w\}}
 |\Delta_{\alpha,\ell}(Q_z)|.
\]
The three fixed prefix differences are bounded by
\((C_\ell/\alpha)^3\) by
\Cref{lem:V41-fixed-prefix}.  The remaining output mass, resolvent,
and root difference are bounded by \(C\alpha\) by
\Cref{lem:V41-one-clock-payer}.  Sum the raw masses using
\eqref{eq:V41-pinching-probability}.  This proves the first
inequality in \eqref{eq:V41-sparse-bound}.

On the corridor, \(\alpha\asymp A_j\), while
\(\Xi_a=A_j+2B\), \(\Xi_c=A_j+B\), and \(B\) is fixed.  Hence the
right side is \(O(A_j^{-4})\).
\end{proof}

\begin{corollary}[The complete V37 transition packet satisfies its
marked-tree target]
\label{cor:V41-tree-payment}
On the same corridor,
\begin{equation}
\label{eq:V41-tree-payment}
 \|\mathcal W_{\alpha,j}^{\rm sp}\|_{\rho,2}^{\rm BFA}
 \le
 {C_{\rho,\ell,x_\pm}\over A_j}
 \sum_{\tau\in\mathfrak T_5}
 \prod_{xy\in E(\tau)}\theta_{xy}.
\end{equation}
In particular the sparse packet has one full inverse-\(A_j\) surplus
over its required quintic marked-tree payment.
\end{corollary}

\begin{proof}
The exact tree sum is
\[
 {B^3A_j^2\over\Xi_a^3\Xi_c^2}
 \asymp_\ell A_j^{-3}
\]
by \eqref{eq:V37-tree-sum}.  The left side is
\(O(A_j^{-4})\) by
\Cref{thm:V41-sparse-transition}.  Comparing the two estimates gives
\eqref{eq:V41-tree-payment}.
\end{proof}

\begin{corollary}[No transition counterexample comes from the V37
Cartan endpoint]
\label{cor:V41-no-V37-counterexample}
The nonzero Cartan limit in
\eqref{eq:V40-Cartan-limit} cannot be promoted to a counterexample to
quintic BFA MTA on the V37 support.
\end{corollary}

\begin{proof}
The Cartan endpoint is one term of the complete output sum in
\Cref{thm:V41-sparse-transition}.  Even at its worst payer size
\(O(\alpha)\), it is multiplied by the three fixed connected prefix
differences, of total size \(O(\alpha^{-3})\).  The resulting complete
word is smaller than the tree target by
\Cref{cor:V41-tree-payment}.
\end{proof}

\section{A reusable sparse-transition principle}
\label{sec:V41-principle}

\begin{proposition}[Fixed connected prefixes neutralize one hot
payer]
\label{prop:V41-fixed-prefix-principle}
Consider a fixed-arity first-connectivity word with \(k\) distinct
fixed-representation binary prefix covariances and one remaining
joining coordinate in a compact parabolic corridor.  Suppose:
\begin{enumerate}[label=\textup{(\roman*)},leftmargin=3em]
\item its exact output dimensions cancel into a raw pinching
      subprobability;
\item its balanced height quotient is at most one;
\item its final output mass, reduced resolvent, and joining
      difference have combined norm at most \(C\alpha\); and
\item all other local factors are contractions.
\end{enumerate}
Then its normalized BFA response is bounded by
\begin{equation}
\label{eq:V41-k-prefix}
 {C\alpha^{1-k}\over\prod_i\Xi_i}.
\end{equation}
\end{proposition}

\begin{proof}
Each fixed connected prefix is \(O(\alpha^{-1})\) by the finite-bank
argument of \Cref{lem:V41-fixed-prefix}.  Multiply the \(k\) factors,
the one \(O(\alpha)\) payer, and the input mass denominators.  The
height quotient is at most one, and the raw output pinching masses
sum to at most one.  This gives
\eqref{eq:V41-k-prefix}.
\end{proof}

\begin{firewall}[Limits of the reusable principle]
\label{fire:V41-principle-scope}
The proposition does not apply when a prefix carrier is a complete
moment rather than a connected difference, when its representations
also grow with \(\alpha\), when two local roles share a noncommuting
coordinate before the cumulant is assembled, or when its coefficient
is normed termwise.  Those are precisely the ways a transition
residual can evade the V41 payment.
\end{firewall}

\section{Residual transition geometry after V41}
\label{sec:V41-residual}

\begin{theorem}[Necessary features of a persistent sparse residual]
\label{thm:V41-residual-necessary}
Within the one-strong \(q=0\) census, a persistent transition
counterpacket cannot be the V37 tree with three fixed connected
prefixes.  Any residual that retains a nondecaying two-cofinal
joining symbol must satisfy at least one of:
\begin{enumerate}[label=\textup{(\roman*)},leftmargin=3em]
\item fewer than three weak tree edges are realized as distinct
      fixed-representation connected prefix covariances;
\item some weak edge is carried by an uncentered moment block;
\item a prefix representation grows on the Wilson scale, so its
      connected difference is not \(O(\alpha^{-1})\); or
\item prefix and joining roles coincide in a single noncommuting
      local cumulant and cannot be multiplied as independent scalar
      differences.
\end{enumerate}
\end{theorem}

\begin{proof}
If none of \textup{(i)}--\textup{(iv)} holds, the word contains the
three independent fixed connected prefix factors and satisfies the
hypotheses of
\Cref{prop:V41-fixed-prefix-principle} with \(k=3\).
The exact V37 mass and tree ledger then gives
\Cref{cor:V41-tree-payment}, contradicting persistence.
\end{proof}

\begin{assessment}[New bottleneck]
\label{ass:V41-bottleneck}
The remaining transition problem is not the value of the hot Cartan
symbol.  It is the source-role question of whether a one-strong
\(q=0\) word can realize three small tree marks while fewer than
three independent connected prefix differences pay their
small-noise clocks.  This is a finite first-connectivity and
moment-partition classification, followed by a local
coincident-role estimate.
\end{assessment}

\section{V41 result ledger and V42 programme}
\label{sec:V41-ledger}

\begin{theorem}[Fourth-series V41 result ledger]
\label{thm:V41-results}
Relative to V1--V40, V41 proves:
\begin{enumerate}[label=\textup{(V41.\arabic*)},leftmargin=6.3em]
\item every fixed-spin prefix covariance contributes one exact
      \(O(\alpha^{-1})\) Wilson factor;
\item the complete resolved multiplier of the four-coordinate V37
      connected word;
\item a uniform \(O(\alpha)\) bound for its one final BFA
      mass--resolvent payer, including the nondecaying Cartan channel;
\item exact use of Peter--Weyl dimension cancellation and raw
      pinching subprobability before the output sum;
\item the complete \(O(A_j^{-4})\) BFA estimate for all sparse
      transition outputs and multiplicities;
\item an \(A_j^{-1}\) surplus over the exact quintic marked-tree
      target; and
\item four necessary mechanisms by which a genuine transition
      residual must evade this payment.
\end{enumerate}
\end{theorem}

\begin{proof}
These are
\Cref{lem:V41-fixed-prefix,
lem:V41-resolved-multiplier,
lem:V41-one-clock-payer,
thm:V41-sparse-transition,
cor:V41-tree-payment,
thm:V41-residual-necessary}.
\end{proof}

\begin{programme}[V42: classify uncentered and coincident-prefix
transition words]
\label{prog:V41-V42}
\begin{enumerate}[leftmargin=3em]
\item Return to the exact five-row first-connectivity formula and
      classify which of the three prefix tree edges can appear as
      complete moments rather than connected differences in a
      \(q=0\) source.
\item Preserve the outer fifth cumulant while contracting each
      predecessor block.  Use M\"obius identities to determine the
      minimum number of independent small-noise differences, not a
      termwise moment count.
\item For a role coincidence, form the complete local
      moment--cumulant--payer operator before taking its norm.
\item If at least two fixed connected prefixes are forced, compare
      their \(\alpha^{-2}\) gain with the exact row-mass and tree
      deficit.  Do not assume in advance that three are necessary.
\item If an uncentered weak prefix survives, construct its smallest
      legal source word and test it with the exact BFA
      normalization.
\item Keep the archived V126 old-pair maximum-square route available
      for geometries with a previously selected bi-heavy pair, but
      do not apply it to the single-strong-pair V37 geometry.
\end{enumerate}
\end{programme}

\begin{firewall}[Claim boundary after fourth-series V41]
\label{fire:V41-claim-boundary}
V41 closes the complete V37 sparse packet uniformly on compact
parabolic corridors.  It does not classify every uncentered or
coincident-role \(q=0\) word, close the full transition band, prove a
uniform quintic MTA, an all-order atlas, WUFC, CPM, cutoff removal,
reflection positivity, Osterwalder--Schrader reconstruction,
construction of the continuum Yang--Mills measure, or a uniform
positive continuum spectral gap.  The full Yang--Mills mass gap has
not been proved.
\end{firewall}

\paragraph{Parent-version record.}
V41 was formed from
\texttt{Renewal\_Yang\_Mills\_Mass\_Gap\_Research\_Notes\_V40.tex}.
The V40 parent had \(19\,880\) logical source lines,
\(694\,386\) bytes, and SHA--256
\[
 \texttt{45d22667aac756abfd3d33e8ceb0fdbb1f008841386190fc673e13cdb9f0542e}.
\]
The next compulsory companion audit is V45.

% =============================================================================
% END APPENDED FOURTH-SERIES V41 MATERIAL
% =============================================================================

% =============================================================================
% BEGIN APPENDED FOURTH-SERIES V42 MATERIAL
% =============================================================================

\chapter{V42: Local-cumulant connectivity forces the full fixed-prefix
Wilson rank}
\label{chap:V42-connectivity-rank}

\section{Purpose}
\label{sec:V42-purpose}

V41 left open whether an outer fifth-cumulant word could use an
uncentered moment in place of one of the three fixed prefix
covariances.  At the level of a selected moment-partition term it can:
singleton local blocks are ordinary moments.  At the level of the
assembled outer cumulant, however, a singleton block carries no
connectivity.  The exact product-cumulant formula forces enough
non-singleton local cumulant blocks to generate the predecessor
partition.

This note proves that statement as a partition-lattice identity.
The resulting rank count is insensitive to whether several prefix
connections occur at one coordinate.  On a fixed Peter--Weyl bank,
a local cumulant block of size \(t\) costs
\(\alpha^{-(t-1)}\); its exponent is exactly its connectivity rank.
Thus every fixed-representation prefix realizing the V37
\(3+2\) predecessor partition carries at least three Wilson clocks,
even with uncentered spectators or coincident prefix roles.

\section{An exact product-cumulant formula}
\label{sec:V42-product-cumulant}

Let \(I=\{1,\ldots,m\}\) index the input rows and let \(E\) be a finite
set of physical coordinates.  At coordinate \(e\), let
\[
 I_e:=\{i\in I:e\hbox{ belongs to the support of row }i\}.
\]
Write \(X_{i,e}\) for the local matrix coefficient carried by row
\(i\) at \(e\), and write
\[
 G_i:=\bigotimes_{e\in E}X_{i,e},
\]
with an identity factor when \(i\notin I_e\).  Wilson variables at
distinct coordinates are independent.

For \(\pi_e\in\Pi(I_e)\), define
\begin{equation}
\label{eq:V42-local-partition-carrier}
 {\mathcal K}_{e,\pi_e}
 :=
 \bigotimes_{B\in\pi_e}
 \kappa_{|B|}^{\alpha,e}\bigl((X_{i,e})_{i\in B}\bigr).
\end{equation}
Here a singleton cumulant is the first moment.  Let
\(\widetilde\pi_e\in\Pi(I)\) be obtained by adjoining every
\(i\notin I_e\) as a singleton.

\begin{theorem}[Connected local-partition formula]
\label{thm:V42-Leonov-Shiryaev}
The complete row cumulant has the exact expansion
\begin{equation}
\label{eq:V42-connected-local-formula}
 \kappa_m^\alpha(G_1,\ldots,G_m)
 =
 \sum_{\substack{\pi_e\in\Pi(I_e),\,e\in E\\
       \bigvee_{e\in E}\widetilde\pi_e=\widehat1_I}}
 \bigotimes_{e\in E}{\mathcal K}_{e,\pi_e}.
\end{equation}
Thus a tuple of local partitions occurs with coefficient one exactly
when the join of its local connectivity partitions is the one-block
row partition; every disconnected tuple cancels.
\end{theorem}

\begin{proof}
Start from the outer moment--cumulant formula
\[
 \kappa_m^\alpha(G_1,\ldots,G_m)
 =
 \sum_{\sigma\in\Pi(I)}
 \mu(\sigma,\widehat1_I)
 \prod_{D\in\sigma}
 \mathbb E_\alpha\!\left[\prod_{i\in D}G_i\right].
\]
Independence across physical coordinates gives
\[
 \prod_{D\in\sigma}
 \mathbb E_\alpha\!\left[\prod_{i\in D}G_i\right]
 =
 \bigotimes_{e\in E}
 \prod_{D\in\sigma}
 \mathbb E_\alpha\!\left[
 \prod_{i\in D\cap I_e}X_{i,e}\right].
\]
Expand each local moment into local cumulants.  A tuple
\((\pi_e)_{e\in E}\) occurs in the term indexed by \(\sigma\) if and
only if every block of every \(\widetilde\pi_e\) is contained in a
block of \(\sigma\).  Equivalently, with
\[
 \tau:=\bigvee_{e\in E}\widetilde\pi_e,
\]
one has \(\tau\le\sigma\).  The total coefficient of this tuple is
therefore
\[
 \sum_{\sigma:\,\tau\le\sigma\le\widehat1_I}
 \mu(\sigma,\widehat1_I).
\]
M\"obius inversion on the partition lattice makes this coefficient
equal to zero unless \(\tau=\widehat1_I\), in which case it equals
one.  This proves \eqref{eq:V42-connected-local-formula}.

The calculation is valid entrywise for matrix coefficients.
Operators belonging to different physical coordinates act on tensor
factors and therefore commute.  Reassembling the entries gives the
operator identity.
\end{proof}

\begin{corollary}[A two-row coordinate is necessarily centered]
\label{cor:V42-two-row-centered}
Suppose \(I_e=\{i,k\}\) and \(e\) is the only coordinate crossing a
cut which separates \(i\) from \(k\).  In every summand of
\eqref{eq:V42-connected-local-formula}, the local partition at \(e\)
is \(\{\{i,k\}\}\), and its factor is the binary covariance
\[
 \kappa_2^{\alpha,e}(X_{i,e},X_{k,e}).
\]
The uncentered partition
\(\{\{i\},\{k\}\}\) cancels from the assembled row cumulant.
\end{corollary}

\begin{proof}
If \(i\) and \(k\) are singleton-separated at \(e\), no local block
at that coordinate crosses the stated cut.  By hypothesis no other
coordinate crosses it, so the join of all lifted local partitions
cannot be \(\widehat1_I\).  The term is absent from
\eqref{eq:V42-connected-local-formula}.
\end{proof}

\section{Partition rank and the minimum number of clocks}
\label{sec:V42-rank}

For a partition \(\pi\in\Pi(J)\), define its connectivity rank by
\begin{equation}
\label{eq:V42-rank}
 r_J(\pi):=|J|-|\pi|
 =\sum_{B\in\pi}(|B|-1).
\end{equation}
Singleton blocks contribute zero.

\begin{lemma}[Subadditivity of partition rank under joins]
\label{lem:V42-rank-subadditive}
For partitions \(\pi_1,\ldots,\pi_N\in\Pi(I)\),
\begin{equation}
\label{eq:V42-rank-subadditive}
 r_I\!\left(\bigvee_{\nu=1}^N\pi_\nu\right)
 \le\sum_{\nu=1}^N r_I(\pi_\nu).
\end{equation}
\end{lemma}

\begin{proof}
Begin with the discrete partition, which has \(m\) blocks.  Adjoining
one block \(B\) to the current equivalence relation can merge at most
\(|B|\) existing components and hence can reduce their number by at
most \(|B|-1\).  Add successively every block of every
\(\pi_\nu\).  The total possible reduction in block number is at most
\[
 \sum_{\nu=1}^N\sum_{B\in\pi_\nu}(|B|-1)
 =\sum_{\nu=1}^N r_I(\pi_\nu).
\]
The actual reduction is the rank of the final join, proving
\eqref{eq:V42-rank-subadditive}.
\end{proof}

\begin{corollary}[Every connected local tuple has rank at least
\(m-1\)]
\label{cor:V42-total-rank}
Every tuple retained in
\eqref{eq:V42-connected-local-formula} satisfies
\begin{equation}
\label{eq:V42-total-rank}
 \sum_{e\in E}r_{I_e}(\pi_e)\ge m-1.
\end{equation}
\end{corollary}

\begin{proof}
Adjoining inactive singleton rows does not change rank, so
\[
 r_I(\widetilde\pi_e)=r_{I_e}(\pi_e).
\]
Apply \Cref{lem:V42-rank-subadditive} and use
\[
 r_I(\widehat1_I)=m-1.
\]
\end{proof}

\begin{lemma}[Rank below the first-connectivity coordinate]
\label{lem:V42-prefix-rank}
Fix an ordering of the physical coordinates and let \(r\) be a
first-connectivity coordinate.  Suppose the join of the local
partitions strictly below \(r\) is the predecessor partition
\[
 \rho=\{C_1,\ldots,C_b\}\in\Pi(I).
\]
Then
\begin{equation}
\label{eq:V42-prefix-rank}
 \sum_{e<r}r_{I_e}(\pi_e)
 \ge r_I(\rho)
 =m-b
 =\sum_{\nu=1}^b(|C_\nu|-1).
\end{equation}
\end{lemma}

\begin{proof}
Apply \Cref{lem:V42-rank-subadditive} to the lifted partitions with
\(e<r\).  Their join is \(\rho\) by hypothesis, and adjoining
inactive singleton rows again leaves every local rank unchanged.
\end{proof}

\begin{corollary}[The \(3+2\) predecessor forces three prefix ranks]
\label{cor:V42-three-prefix-ranks}
For five rows and
\[
 \rho=\{\{a,b,d\},\{c,e\}\},
\]
every local-partition tuple below the joining coordinate satisfies
\[
 \sum_{e<r}r_{I_e}(\pi_e)\ge3.
\]
This remains true if several connections occur at one coordinate and
if arbitrary singleton moment blocks are also present.
\end{corollary}

\begin{proof}
Here \(m=5\) and \(b=2\), so
\eqref{eq:V42-prefix-rank} gives \(5-2=3\).
Coincidence of coordinates changes the distribution of the summands
in the left side but not their sum.  Singleton blocks have rank zero
by definition.
\end{proof}

\section{Fixed-bank cumulant degree}
\label{sec:V42-fixed-bank}

The small-noise normal form in archive f1 V29 proves the following
finite-bank fact.  We restate its exact scope because it is the
analytic input paired with the new rank lemma.

\begin{record}[Fixed Peter--Weyl cumulant card]
\label{rec:V42-fixed-bank-card}
Fix \(t\ge2\) and a finite set \(\mathcal F\) of \(SU(2)\)
representations.  There are constants
\(C_{\mathcal F,t}\) and \(\alpha_{\mathcal F,t}\) such that, for all
\(\alpha\ge\alpha_{\mathcal F,t}\), all
\(R_1,\ldots,R_t\in\mathcal F\), and trace-class coefficient inputs,
\begin{equation}
\label{eq:V42-fixed-cumulant}
 \left\|
 \kappa_t^\alpha(X_1,\ldots,X_t)
 \right\|
 \le
 C_{\mathcal F,t}\alpha^{-(t-1)}
 \prod_{\nu=1}^t\|X_\nu\|.
\end{equation}
The same bound holds after complete isotypic compression.  Its
constant is not asserted uniformly when the representation bank
grows with \(\alpha\).
\end{record}

The proof recorded in f1 uses
\(U=\exp(\alpha^{-1/2}Z)\), the Gaussian Laplace normal form, and the
connected-graph degree of the cumulant.  A connected \(t\)-source
Gaussian diagram has at least \(t-1\) contractions, each of size
\(\alpha^{-1}\).  Taylor remainders and the complement of the
identity chart obey the same order on a fixed representation bank.
This is precisely the hypothesis needed below and avoids redoing the
archived normal-form construction.

\begin{lemma}[A local partition pays its rank]
\label{lem:V42-local-rank-payment}
Under the hypotheses of
\Cref{rec:V42-fixed-bank-card}, for every
\(\pi_e\in\Pi(I_e)\),
\begin{equation}
\label{eq:V42-local-rank-payment}
 \|{\mathcal K}_{e,\pi_e}\|
 \le
 C_{\mathcal F,m}
 \alpha^{-r_{I_e}(\pi_e)}
 \prod_{i\in I_e}\|X_{i,e}\|.
\end{equation}
\end{lemma}

\begin{proof}
Apply \eqref{eq:V42-fixed-cumulant} to each block
\(B\in\pi_e\) with \(|B|\ge2\).  Singleton first moments are
contractions and cost no inverse power.  Multiplication over the
blocks gives exponent
\[
 \sum_{B\in\pi_e}(|B|-1)=r_{I_e}(\pi_e).
\]
Only finitely many partitions occur at fixed \(m\), so their
constants are absorbed into \(C_{\mathcal F,m}\).
\end{proof}

\begin{theorem}[Fixed-prefix Wilson rank theorem]
\label{thm:V42-fixed-prefix-rank}
Fix a first-connectivity fibre with predecessor partition
\(\rho=\{C_1,\ldots,C_b\}\).  Suppose every representation at every
coordinate strictly below the first-connectivity coordinate belongs
to one fixed finite bank.  Then each assembled local-partition term
which realizes \(\rho\) has prefix operator norm at most
\begin{equation}
\label{eq:V42-fixed-prefix-main}
 C_{\mathcal F,m}\alpha^{-(m-b)}
 \prod_{i,e<r}\|X_{i,e}\|.
\end{equation}
The same estimate holds after summing all local partitions and all
fixed multiplicity data.
\end{theorem}

\begin{proof}
For one retained local tuple, multiply
\eqref{eq:V42-local-rank-payment} over \(e<r\).
\Cref{lem:V42-prefix-rank} makes its total exponent at least
\(m-b\).  The number of local partition tuples and multiplicity
labels is finite because the row order, physical prefix support, and
representation bank in this local statement are fixed.  Summing
changes only the constant.
\end{proof}

\begin{corollary}[Coincident prefix roles do not lose a clock]
\label{cor:V42-coincident-prefix}
Suppose one fixed-representation coordinate contains a local
cumulant block on \(t\) rows and that this block is used to connect
\(t-1\) predecessor relations.  Its norm is
\[
 O(\alpha^{-(t-1)}).
\]
It therefore pays exactly as many Wilson clocks as its maximum
connectivity rank, even though the corresponding tree-edge roles
coincide at one coordinate.
\end{corollary}

\begin{proof}
The one-block partition on those \(t\) rows has rank \(t-1\).
Apply \Cref{lem:V42-local-rank-payment}.
\end{proof}

\section{Completion of the fixed-prefix transition sector}
\label{sec:V42-completion}

\begin{theorem}[All fixed-prefix \(3+2\) transition words have the
V41 three-clock payment]
\label{thm:V42-fixed-32}
Consider a five-row first-connectivity fibre whose predecessor
partition is
\[
 \rho=\{\{a,b,d\},\{c,e\}\}.
\]
Assume:
\begin{enumerate}[label=\textup{(\roman*)},leftmargin=3em]
\item all representations below the joining coordinate lie in one
      fixed finite bank;
\item the joining coordinate is in a compact Wilson transition
      corridor and all of its input and output spins are
      \(O(\sqrt\alpha)\);
\item exact Peter--Weyl dimensions are compressed into raw pinching
      mass before the output trace norm; and
\item the final BFA mass and reduced resolvent are retained together.
\end{enumerate}
Then the complete prefix contributes \(O(\alpha^{-3})\), independently
of uncentered singleton moments and coincident prefix roles.  The
complete final mass--resolvent--joining carrier costs at most
\(O(\alpha)\).  Consequently the resolved word is bounded by
\begin{equation}
\label{eq:V42-32-word}
 {C\alpha^{-2}\over\prod_{i=1}^5\Xi_i}
\end{equation}
before its geometric tree marks are inserted.
\end{theorem}

\begin{proof}
The prefix exponent is \(m-b=5-2=3\) by
\Cref{thm:V42-fixed-prefix-rank}.  For the final carrier, all output
Casimirs are \(O(\alpha)\).  The representation-uniform product gap
\eqref{eq:V41-product-gap}, the Casimir floor on nontrivial output,
and the universal fixed-order cumulant norm give, exactly as in
\Cref{lem:V41-one-clock-payer},
\[
 {\Xi(\boldsymbol Q)\over1-q_{\alpha,\boldsymbol Q}}
 \|\hbox{joining cumulant on }\boldsymbol Q\|
 \le C\alpha.
\]
The balanced height quotient is at most one.  Exact Fourier
normalization makes the raw output masses a subprobability, so the
output sum introduces no label count.  Multiplication by the five
input mass denominators gives
\eqref{eq:V42-32-word}.
\end{proof}

\begin{corollary}[The uncentered-prefix loophole in V41 is closed]
\label{cor:V42-uncentered-closed}
An uncentered singleton moment can occur in an individual local
partition term, but it cannot replace the rank needed to form a weak
prefix tree edge.  Hence alternatives \textup{(i)}, \textup{(ii)},
and the fixed-representation part of \textup{(iv)} in
\Cref{thm:V41-residual-necessary} do not create a residual in a
fixed-prefix \(3+2\) fibre.
\end{corollary}

\begin{proof}
Singletons have rank zero.  The predecessor join nevertheless has
rank three, so the remaining nonsingleton local blocks have total
rank at least three by
\Cref{cor:V42-three-prefix-ranks}.  Their fixed-bank cumulants pay
the corresponding three powers by
\Cref{thm:V42-fixed-prefix-rank}.  Coincidence is covered by
\Cref{cor:V42-coincident-prefix}.
\end{proof}

\begin{firewall}[No support-uniform conclusion from a fixed bank]
\label{fire:V42-fixed-bank-scope}
\Cref{thm:V42-fixed-prefix-rank} is uniform over a fixed
representation bank and a fixed resolved prefix support.  It does not
control a number of prefix coordinates growing with the support, nor
a prefix representation of height comparable to \(\sqrt\alpha\).
Those cases require probability-mark summation and transition
estimates, respectively.  The theorem is used here only to remove the
claimed uncentered and coincident-role loopholes.
\end{firewall}

\section{The reduced transition frontier}
\label{sec:V42-frontier}

\begin{theorem}[Every persistent \(3+2\) transition prefix has a
growing local block]
\label{thm:V42-growing-prefix-necessary}
After the diffuse-bank and second-atom closures already recorded in
V34--V39, any persistent \(3+2\) transition fibre not paid by
\Cref{thm:V42-fixed-32} must contain, below its first-connectivity
coordinate, a nonsingleton local cumulant block with at least one
representation label escaping every fixed Peter--Weyl bank.
\end{theorem}

\begin{proof}
If all prefix labels remained in some fixed finite bank along a
subsequence, \Cref{thm:V42-fixed-32} would give the three-clock
payment on that subsequence.  Singleton moments cannot carry the
missing connectivity by
\Cref{cor:V42-uncentered-closed}.  Therefore a persistent unclosed
subsequence must contain a nonsingleton prefix block outside every
fixed bank.
\end{proof}

\begin{assessment}[Exact bottleneck after V42]
\label{ass:V42-bottleneck}
The remaining local obstruction has at least two hot regions: the
height-near-saturated joining coordinate and a growing
representation in a predecessor cumulant block.  The next estimate
must compare the growing prefix block's exact Wilson connected
difference with its incident probability marks.  A fixed-bank
Gaussian theorem is no longer available, and a termwise moment bound
would erase the only useful cancellation.
\end{assessment}

\section{V42 result ledger and V43 programme}
\label{sec:V42-ledger}

\begin{theorem}[Fourth-series V42 result ledger]
\label{thm:V42-results}
Relative to V1--V41, V42 proves:
\begin{enumerate}[label=\textup{(V42.\arabic*)},leftmargin=6.3em]
\item the exact connected local-partition expansion
      \eqref{eq:V42-connected-local-formula};
\item cancellation of every local-partition tuple whose lifted join
      is disconnected;
\item the partition-rank lower bound \(m-1\) for a full cumulant and
      \(m-b\) below a \(b\)-block predecessor partition;
\item the exact three-rank floor for a five-row \(3+2\) prefix;
\item conversion of local partition rank into Wilson powers on a
      fixed Peter--Weyl bank;
\item stability of that payment under singleton spectators and
      coincident higher-arity prefix roles;
\item the complete \(O(\alpha^{-2}/\prod_i\Xi_i)\) normalized
      transition word bound in the fixed-prefix \(3+2\) sector; and
\item necessity of a growing nonsingleton prefix block in every
      remaining persistent \(3+2\) transition fibre.
\end{enumerate}
\end{theorem}

\begin{proof}
These are
\Cref{thm:V42-Leonov-Shiryaev,
cor:V42-total-rank,
lem:V42-prefix-rank,
cor:V42-three-prefix-ranks,
lem:V42-local-rank-payment,
cor:V42-coincident-prefix,
thm:V42-fixed-32,
thm:V42-growing-prefix-necessary}.
\end{proof}

\begin{programme}[V43: growing-prefix transition debit]
\label{prog:V42-V43}
\begin{enumerate}[leftmargin=3em]
\item Select a minimal nonsingleton growing prefix block and denote
      by \(H\) its total incident Casimir mass.
\item Keep its complete signed Wilson cumulant and derive a bound in
      terms of \(\min\{s_\alpha H,1\}\), with the correct power for
      its partition rank.
\item Compare that bound with the product of the incident normalized
      link marks; use raw pinching rather than a representation-count
      measure.
\item Split into \(s_\alpha H\ll1\), transition, and hot regimes.
      The fixed-bank result is the first regime only after a
      quantitative remainder is proved.
\item If a rank-\(t-1\) bound by
      \(\min\{s_\alpha H,1\}^{t-1}\) is false, identify the sharp
      mixed-mass replacement before returning to the full word.
\item Propagate only a proved prefix debit through the \(3+2\)
      source-role census.
\end{enumerate}
\end{programme}

\begin{firewall}[Claim boundary after fourth-series V42]
\label{fire:V42-claim-boundary}
V42 closes every compact-transition \(3+2\) word whose predecessor
representations remain in a fixed finite bank and proves that
uncentered singleton moments or fixed coincident roles do not evade
the payment.  It does not control growing prefix cumulants, close the
full transition band, prove uniform quintic MTA, an all-order atlas,
WUFC, CPM, cutoff removal, reflection positivity,
Osterwalder--Schrader reconstruction, construction of the continuum
Yang--Mills measure, or a uniform positive continuum spectral gap.
The full Yang--Mills mass gap has not been proved.
\end{firewall}

\paragraph{Parent-version record.}
V42 was formed from
\texttt{Renewal\_Yang\_Mills\_Mass\_Gap\_Research\_Notes\_V41.tex}.
The V41 parent had \(20\,422\) logical source lines,
\(713\,427\) bytes, and SHA--256
\[
 \texttt{f2fb159d2f6725e054c41a69e9883c27176ff537f2d0a4577916a4a0a7d5908a}.
\]
The next compulsory companion audit is V45.

% =============================================================================
% END APPENDED FOURTH-SERIES V42 MATERIAL
% =============================================================================

% =============================================================================
% BEGIN APPENDED FOURTH-SERIES V43 MATERIAL
% =============================================================================

\chapter{V43: Dimension-free covariance debits close every binary-leaf
transition tree}
\label{chap:V43-binary-leaf}

\section{Purpose and relation to the archive}
\label{sec:V43-purpose}

V42 proves that a persistent transition prefix must leave every fixed
Peter--Weyl bank.  Archive f1 V125 already supplies a one-sided
support-free square-root Wilson displacement bound.  Rather than
repeat that theorem, this note uses its two-sided centered form and
applies it to the precise quintic normalization.

The resulting estimate is stronger than a fixed-bank expansion for
the four-coordinate binary tree: it is uniform for arbitrary growing
prefix representations.  Each binary prefix covariance contributes
the geometric mean of the two endpoint character gaps.  After the
leaf-row BFA masses are divided out, the three geometric means and
the one final payer leave an additional inverse root scale.  Hence
the entire binary-leaf transition class closes.

\section{A dimension-free centered covariance inequality}
\label{sec:V43-covariance}

Let \(\pi_R,\pi_S\) be unitary irreducible representations of a
compact group and let \(\nu\) be a central inversion-invariant
probability law.  Write
\[
 \mathbb E_\nu\pi_R=q_RI_R,
 \qquad
 \mathbb E_\nu\pi_S=q_SI_S,
\]
where \(q_R,q_S\in[-1,1]\) are real.  Define
\begin{equation}
\label{eq:V43-centered-covariance}
 {\mathcal C}_{R,S}
 :=
 \mathbb E_\nu(\pi_R\otimes\pi_S)
 -
 (\mathbb E_\nu\pi_R)\otimes(\mathbb E_\nu\pi_S).
\end{equation}

\begin{lemma}[Operator-valued Cauchy--Schwarz]
\label{lem:V43-operator-CS}
Let \(X(\omega)\in{\mathcal B}(H)\) and
\(Y(\omega)\in{\mathcal B}(K)\) be square-integrable operator fields.
Then
\begin{equation}
\label{eq:V43-operator-CS}
 \left\|\mathbb E[X\otimes Y]\right\|
 \le
 \left\|\mathbb E[XX^*]\right\|^{1/2}
 \left\|\mathbb E[Y^*Y]\right\|^{1/2}.
\end{equation}
\end{lemma}

\begin{proof}
Define bounded maps from \(H\otimes K\) to
\(L^2(\Omega;H\otimes K)\) by
\[
 (A\xi)(\omega)=(X(\omega)^*\otimes I)\xi,
 \qquad
 (B\xi)(\omega)=(I\otimes Y(\omega))\xi.
\]
Then
\[
 A^*B=\mathbb E[X\otimes Y],
\]
while
\[
 A^*A=\mathbb E[XX^*]\otimes I,
 \qquad
 B^*B=I\otimes\mathbb E[Y^*Y].
\]
The inequality \(\|A^*B\|\le\|A\|\|B\|\) proves
\eqref{eq:V43-operator-CS}.
\end{proof}

\begin{theorem}[Two-sided character-gap covariance debit]
\label{thm:V43-two-sided-debit}
The complete covariance in
\eqref{eq:V43-centered-covariance} obeys
\begin{equation}
\label{eq:V43-character-gap-debit}
 \boxed{
 \|{\mathcal C}_{R,S}\|
 \le
 \sqrt{1-q_R^2}\sqrt{1-q_S^2}
 \le
 2\sqrt{(1-q_R)(1-q_S)}.}
\end{equation}
Every multiplicity-resolved output compression satisfies the same
bound.  The constant is independent of representation dimensions and
output multiplicities.
\end{theorem}

\begin{proof}
Put
\[
 X=\pi_R-q_RI_R,
 \qquad
 Y=\pi_S-q_SI_S.
\]
Expansion and centrality give
\[
 {\mathcal C}_{R,S}=\mathbb E[X\otimes Y].
\]
Unitarity, reality of the character scalars, and
\(\mathbb E\pi_R=q_RI_R\) imply
\[
 \mathbb E[XX^*]=(1-q_R^2)I_R.
\]
Likewise
\(\mathbb E[Y^*Y]=(1-q_S^2)I_S\).
Apply \Cref{lem:V43-operator-CS}.  Since
\(1-q^2=(1-q)(1+q)\le2(1-q)\), the second inequality follows.
Orthogonal compression cannot increase operator norm.
\end{proof}

\begin{corollary}[Casimir form for the Wilson law]
\label{cor:V43-Casimir-debit}
For the exact Wilson law and all sufficiently large \(\alpha\),
\begin{equation}
\label{eq:V43-Casimir-debit}
 \|{\mathcal C}_{R,S}\|
 \le
 C
 \sqrt{\min\{s_\alpha c_R,1\}}
 \sqrt{\min\{s_\alpha c_S,1\}}
 \le
 Cs_\alpha\sqrt{c_Rc_S}.
\end{equation}
\end{corollary}

\begin{proof}
Use the representation-uniform upper character gap
\[
 1-q_R\le C\min\{s_\alpha c_R,1\}
\]
and its \(S\)-analogue in
\eqref{eq:V43-character-gap-debit}.  The final inequality uses
\(\min\{x,1\}\le x\).
\end{proof}

\begin{remark}[Consistency with V40]
\label{rem:V43-V40-consistency}
If one endpoint is fixed and the other lies at
\(c_R\asymp\alpha\), \eqref{eq:V43-Casimir-debit} gives
\(O(\alpha^{-1/2})\), exactly the sharp displaced scale of
\Cref{thm:V40-one-cofinal}.  If both endpoints are fixed it gives
\(O(\alpha^{-1})\), recovering the order used in V41.  If both are
at transition scale it allows order one, consistently with the
Cartan limit \eqref{eq:V40-Cartan-limit}.
\end{remark}

\section{The arbitrary binary-leaf prefix family}
\label{sec:V43-family}

Retain the V37 incidence graph but allow every prefix representation
to depend on \(\alpha\).  At the four coordinates write
\begin{center}
\begin{tabular}{c|ccccc}
coordinate & \(a\)&\(b\)&\(c\)&\(d\)&\(e\)\\
\hline
\(u\)&\(R_u\)&\(S_u\)&\(0\)&\(0\)&\(0\)\\
\(v\)&\(R_v\)&\(0\)&\(0\)&\(S_v\)&\(0\)\\
\(w\)&\(0\)&\(0\)&\(R_w\)&\(0\)&\(S_w\)\\
\(r\)&\(j\)&\(0\)&\(j\)&\(0\)&\(0\).
\end{tabular}
\end{center}
Put
\begin{align}
 A&:=1+c_j,
\label{eq:V43-root-mass}\\
 x_z&:=1+c_{R_z},
 &
 y_z&:=1+c_{S_z}
 \qquad(z=u,v,w).
\label{eq:V43-prefix-masses}
\end{align}
The row masses are
\begin{align}
 \Xi_a&=A+x_u+x_v,
 &
 \Xi_c&=A+x_w,
\label{eq:V43-center-masses}\\
 \Xi_b&=y_u,
 &
 \Xi_d&=y_v,
 &
 \Xi_e&=y_w.
\label{eq:V43-leaf-masses}
\end{align}
Assume root saturation and a compact parabolic clock:
\begin{equation}
\label{eq:V43-root-corridor}
 A\ge\sigma\max\{\Xi_a,\Xi_c\},
 \qquad
 0<x_-\le {d_j\over\sqrt\alpha}\le x_+<\infty.
\end{equation}
Thus
\begin{equation}
\label{eq:V43-comparable-scales}
 \Xi_a\asymp_{\sigma}A,
 \qquad
 \Xi_c\asymp_{\sigma}A,
 \qquad
 A\asymp_{x_\pm}\alpha,
 \qquad
 s_\alpha\asymp A^{-1}.
\end{equation}

\begin{lemma}[Exact tree monomial for arbitrary prefix masses]
\label{lem:V43-tree-monomial}
The support intersection graph is still the unique tree
\[
 T_*=\{ab,ad,ce,ac\}.
\]
Its marked monomial is
\begin{equation}
\label{eq:V43-tree-monomial}
 {\mathfrak t}_*
 =
 {x_u\over\Xi_a}
 {x_v\over\Xi_a}
 {x_w\over\Xi_c}
 {A^2\over\Xi_a\Xi_c}
 =
 {A^2x_ux_vx_w\over\Xi_a^3\Xi_c^2}.
\end{equation}
It is independent of the leaf masses \(y_u,y_v,y_w\).
\end{lemma}

\begin{proof}
At \(u\), the normalized pair stock is
\[
 {x_uy_u\over\Xi_a\Xi_b}
 ={x_u\over\Xi_a}
\]
because the exact support of leaf row \(b\) is \(\{u\}\).
The same calculation gives \(x_v/\Xi_a\) at \(v\) and
\(x_w/\Xi_c\) at \(w\).  The root stock is
\(A^2/(\Xi_a\Xi_c)\).  No other row pair shares a coordinate, so the
displayed tree is the only nonzero spanning tree.
\end{proof}

\section{The complete BFA estimate}
\label{sec:V43-BFA}

Let \(\mathcal W_{\alpha,j}^{\rm bl}\) denote the complete
first-connectivity word on this binary-leaf support, with all local
outputs, exact multiplicity resolution, \(Q_0\), one final reduced
resolvent, and the output BFA norm.  Normalize each pure input row
block to have BFA norm one.

\begin{lemma}[Uniform final carrier on the saturated corridor]
\label{lem:V43-final-carrier}
After the three prefix covariances have been left outside the
estimate, the combined output BFA mass, final reduced resolvent, and
root covariance have norm at most
\begin{equation}
\label{eq:V43-final-carrier}
 C_{\sigma,x_\pm}\alpha.
\end{equation}
\end{lemma}

\begin{proof}
Every input Casimir is at most its row mass.  By
\eqref{eq:V43-root-corridor}, the sum of all five-row input Casimirs
is \(O_{\sigma}(A)=O(\alpha)\).  Fusion subadditivity gives the same
bound for every complete output Casimir.

On a nontrivial output, the product character gap
\eqref{eq:V41-product-gap} and the Casimir floor imply
\[
 {\Xi(\boldsymbol Q)\over1-q_{\alpha,\boldsymbol Q}}
 \le C_{\sigma,x_\pm}\alpha.
\]
The root covariance has norm at most two.  The completely trivial
output is deleted by \(Q_0\).  This proves the claim.
\end{proof}

\begin{theorem}[Uniform closure of every saturated binary-leaf
transition tree]
\label{thm:V43-binary-leaf-close}
Under \eqref{eq:V43-root-corridor},
\begin{equation}
\label{eq:V43-binary-leaf-bound}
 \|\mathcal W_{\alpha,j}^{\rm bl}\|_{\rho,2}^{\rm BFA}
 \le
 {C_{\rho,\sigma,x_\pm}\over A}\,
 {\mathfrak t}_*.
\end{equation}
The constant is independent of all six growing prefix
representations, every output label and multiplicity, and \(\alpha\).
\end{theorem}

\begin{proof}
At \(u\), \Cref{cor:V43-Casimir-debit} gives
\[
 \|{\mathcal C}_{R_u,S_u}\|
 \le Cs_\alpha\sqrt{x_uy_u};
\]
we enlarged Casimirs by one, which only changes the constant.
The analogous estimates hold at \(v,w\).  Exact height
subadditivity makes the BFA exponential quotient at most one.

Normalize the five inputs.  The exact Fourier product factor cancels
the output representation dimension, and the raw isotypic pinching
masses sum to at most one.  Apply
\Cref{lem:V43-final-carrier} after the three covariance estimates.
Using \eqref{eq:V43-leaf-masses}, the result is
\begin{align}
 \|\mathcal W_{\alpha,j}^{\rm bl}\|_{\rho,2}^{\rm BFA}
 &\le
 {C\alpha s_\alpha^3
  \sqrt{x_uy_ux_vy_vx_wy_w}
  \over
  \Xi_a\Xi_c y_uy_vy_w}
\notag\\
 &=
 {C\alpha s_\alpha^3\over\Xi_a\Xi_c}
 \sqrt{{x_ux_vx_w\over y_uy_vy_w}}.
\label{eq:V43-word-intermediate}
\end{align}
Divide this by \eqref{eq:V43-tree-monomial}.  One obtains
\begin{align}
 {\|\mathcal W_{\alpha,j}^{\rm bl}\|_{\rho,2}^{\rm BFA}
  \over{\mathfrak t}_*}
 &\le
 C\alpha s_\alpha^3
 {\Xi_a^2\Xi_c\over A^2}
 {1\over\sqrt{x_uy_ux_vy_vx_wy_w}}.
\label{eq:V43-ratio}
\end{align}
Every \(x_z,y_z\ge1\).  By
\eqref{eq:V43-comparable-scales},
\[
 \alpha s_\alpha^3
 {\Xi_a^2\Xi_c\over A^2}
 \le {C\over A}.
\]
Insert this in \eqref{eq:V43-ratio} to prove
\eqref{eq:V43-binary-leaf-bound}.
\end{proof}

\begin{corollary}[Growing prefix labels do not revive the V37
packet]
\label{cor:V43-growing-V37}
The conclusion of
\Cref{cor:V41-no-V37-counterexample} remains true when all six prefix
representations grow arbitrarily with \(\alpha\), provided the root
remains saturated as in \eqref{eq:V43-root-corridor}.
\end{corollary}

\begin{proof}
\Cref{thm:V43-binary-leaf-close} is uniform in those labels and
controls the complete output sum, including the Cartan endpoint.
\end{proof}

\begin{remark}[Why no prefix heat-scale split was needed]
\label{rem:V43-no-prefix-split}
If a prefix mass \(x_z\) is small compared with \(A\), the covariance
debit contains the corresponding small factor
\(s_\alpha\sqrt{x_zy_z}\).  If it is comparable with \(A\), its tree
mark \(x_z/\Xi_a\) or \(x_z/\Xi_c\) is not small.  The ratio
\eqref{eq:V43-ratio} treats both cases simultaneously; splitting into
cold and hot prefix sectors would only weaken the ledger.
\end{remark}

\section{Consequences for the transition census}
\label{sec:V43-consequences}

\begin{theorem}[Binary-leaf exclusion from the persistent residual]
\label{thm:V43-binary-leaf-exclusion}
No one-strong compact-transition residual can have all three of its
weak predecessor connections realized on distinct two-row
coordinates whose opposite endpoints are exact-support leaf rows.
This remains true whether the prefix representations are fixed,
mesoscopic, or of transition size.
\end{theorem}

\begin{proof}
Such a residual is exactly the family of
\Cref{thm:V43-binary-leaf-close}.  Its complete BFA word is bounded
by the required tree monomial with an additional factor \(A^{-1}\).
\end{proof}

\begin{assessment}[Residual geometry after V43]
\label{ass:V43-residual}
Together, V42 and V43 remove fixed higher-arity prefixes and arbitrary
binary-leaf prefixes.  A persistent transition word must therefore
use at least one genuinely shared prefix geometry:
\begin{enumerate}[label=\textup{(\roman*)},leftmargin=3em]
\item a weak endpoint row has substantial mass outside the selected
      prefix coordinate, so its leaf denominator does not cancel;
\item one coordinate is incident to three or more rows with growing
      representations and requires a higher-cumulant displacement
      estimate; or
\item a prefix coordinate shares the joining or payer operation
      before the complete physical carrier is assembled.
\end{enumerate}
The first case is a distributed-mass problem and should be returned
to the existing diffuse/second-atom machinery before a new local
inequality is attempted.
\end{assessment}

\section{V43 result ledger and V44 programme}
\label{sec:V43-ledger}

\begin{theorem}[Fourth-series V43 result ledger]
\label{thm:V43-results}
Relative to V1--V42, V43 proves:
\begin{enumerate}[label=\textup{(V43.\arabic*)},leftmargin=6.3em]
\item a dimension-free operator Cauchy--Schwarz estimate for tensor
      covariances;
\item the two-sided character-gap covariance debit
      \eqref{eq:V43-character-gap-debit};
\item its representation-uniform Wilson--Casimir form
      \eqref{eq:V43-Casimir-debit};
\item the exact marked-tree monomial for arbitrary growing labels on
      the saturated binary-leaf skeleton;
\item a complete BFA/output/resolvent estimate retaining raw
      pinching mass;
\item uniform marked-tree closure with an additional inverse root
      scale; and
\item exclusion of every fixed, mesoscopic, and transition-size
      binary-leaf prefix from the persistent residual.
\end{enumerate}
\end{theorem}

\begin{proof}
These are
\Cref{lem:V43-operator-CS,
thm:V43-two-sided-debit,
cor:V43-Casimir-debit,
lem:V43-tree-monomial,
lem:V43-final-carrier,
thm:V43-binary-leaf-close,
thm:V43-binary-leaf-exclusion}.
\end{proof}

\begin{programme}[V44: distributed weak endpoints and shared-prefix
reduction]
\label{prog:V43-V44}
\begin{enumerate}[leftmargin=3em]
\item For each weak predecessor edge, compare the mass at its selected
      coordinate with the full mass of its opposite endpoint row.
\item If a fixed fraction lies at that coordinate, apply the V43
      binary-leaf calculation with a bounded saturation loss.
\item If not, use the off-coordinate mass to select either a second
      atom or a diffuse unmarked bank, respecting the existing
      q-census and all role exclusions.
\item Prove that three simultaneously distributed endpoints either
      yield the required three probability marks or force a
      higher-incidence coordinate.
\item Isolate the exact higher-incidence local carrier only after the
      distributed cases have been removed.
\item Prepare the compulsory V45 SM/NS audit without modifying either
      companion file.
\end{enumerate}
\end{programme}

\begin{firewall}[Claim boundary after fourth-series V43]
\label{fire:V43-claim-boundary}
V43 closes the complete saturated binary-leaf transition family
uniformly in all representation labels.  It does not close
distributed weak endpoints, growing higher-incidence prefix
cumulants, or every coincident joining/payer carrier; hence it does
not prove the full quintic MTA, an all-order atlas, WUFC, CPM, cutoff
removal, reflection positivity, Osterwalder--Schrader reconstruction,
construction of the continuum Yang--Mills measure, or a uniform
positive continuum spectral gap.  The full Yang--Mills mass gap has
not been proved.
\end{firewall}

\paragraph{Parent-version record.}
V43 was formed from
\texttt{Renewal\_Yang\_Mills\_Mass\_Gap\_Research\_Notes\_V42.tex}.
The V42 parent had \(20\,973\) logical source lines,
\(733\,415\) bytes, and SHA--256
\[
 \texttt{5dc0a36fd09c2fca68cf2cde93137af5bbf442f537e7388f69bfb0d70b972e14}.
\]
The next compulsory companion audit is V45.

% =============================================================================
% END APPENDED FOURTH-SERIES V43 MATERIAL
% =============================================================================

% =============================================================================
% BEGIN APPENDED FOURTH-SERIES V44 MATERIAL
% =============================================================================

\chapter{V44: Distributed endpoint masses cancel in the exact
binary-spine ledger}
\label{chap:V44-distributed-endpoints}

\section{Purpose and correction of the V43 frontier}
\label{sec:V44-purpose}

V43 assumed that the three rows opposite the two cofinal center rows
were exact-support leaves.  Its final assessment consequently listed
off-coordinate endpoint mass as a possible residual.  The exact BFA
ledger shows that this restriction is unnecessary.

If a selected binary prefix coordinate carries local endpoint masses
\(x\) and \(y\), its covariance contributes
\(s_\alpha\sqrt{xy}\).  The normalized input contributes the full
endpoint denominator \(\Xi_{\rm leaf}^{-1}\), while the selected tree
mark contains the same denominator.  In the ratio of the complete
word to that tree monomial, \(\Xi_{\rm leaf}\) cancels identically.
Arbitrary spectator mass on that row therefore does not cause a
loss.  This note performs that cancellation for all three prefix
edges simultaneously and reduces the residual to genuinely
higher-incidence local carriers.

\section{A binary spine with arbitrary spectator banks}
\label{sec:V44-spine}

Let the maximal row be \(a\), put \(M=\Xi_a\), and assume
\begin{equation}
\label{eq:V44-maximal-row}
 \Xi_i\le M
 \qquad(i=a,b,c,d,e).
\end{equation}
Select four distinct physical coordinates \(u,v,w,r\) with binary
spine
\[
 u:ab,\qquad v:ad,\qquad w:ce,\qquad r:ac.
\]
At \(u,v,w\), write the local balanced masses as
\begin{align}
 x_u&:=1+c_{R_{a,u}},
 &y_u&:=1+c_{R_{b,u}},
\notag\\
 x_v&:=1+c_{R_{a,v}},
 &y_v&:=1+c_{R_{d,v}},
\label{eq:V44-spine-masses}\\
 x_w&:=1+c_{R_{c,w}},
 &y_w&:=1+c_{R_{e,w}}.
\notag
\end{align}
At \(r\), take equal spin \(j\) on \(a,c\) and put
\[
 A:=1+c_j.
\]
Assume the root is saturated and parabolic:
\begin{equation}
\label{eq:V44-root-hypotheses}
 A\ge\sigma\max\{\Xi_a,\Xi_c\},
 \qquad
 0<x_-\le{d_j\over\sqrt\alpha}\le x_+<\infty.
\end{equation}
In particular
\begin{equation}
\label{eq:V44-root-scales}
 \Xi_a,\Xi_c,M,\alpha\asymp_{\sigma,x_\pm}A,
 \qquad
 s_\alpha\asymp A^{-1}.
\end{equation}

All other coordinates are allowed to carry arbitrary representation
labels and arbitrary row mass, subject only to
\eqref{eq:V44-maximal-row}.  In the resolved word considered below,
their complete local carriers are assumed to be unmarked moment
contractions.  They may change the final output representation and
its multiplicity, but they do not contain the selected joining
cumulants or a second payer.

\begin{definition}[Resolved binary-spine word]
\label{def:V44-spine-word}
A resolved binary-spine word is a first-connectivity summand in which:
\begin{enumerate}[label=\textup{(\roman*)},leftmargin=3em]
\item the assembled local factors at \(u,v,w\) are the complete
      binary covariances on \(ab,ad,ce\), respectively;
\item the coordinate \(r\) joins the two predecessor components
      through the complete \(ac\) carrier;
\item every coordinate outside \(\{u,v,w,r\}\) is a complete
      contraction disjoint from the joining and payer roles; and
\item the final \(Q_0\), reduced resolvent, output BFA weight, exact
      physical duality, and all output multiplicities are retained.
\end{enumerate}
\end{definition}

\section{The selected tree lower bound}
\label{sec:V44-tree}

Let
\[
 \theta_{ik,e}:=
 {a_{i,e}a_{k,e}\over\Xi_i\Xi_k}
\]
be the local contribution of coordinate \(e\) to the global pair
stock \(\theta_{ik}\).

\begin{lemma}[A selected binary spine gives a legal tree monomial]
\label{lem:V44-selected-tree}
The complete marked-tree sum is bounded below by
\begin{equation}
\label{eq:V44-selected-tree}
 {\mathfrak t}_{\rm sel}
 :=
 {A^2x_uy_ux_vy_vx_wy_w
  \over
  \Xi_a^3\Xi_c^2\Xi_b\Xi_d\Xi_e}.
\end{equation}
\end{lemma}

\begin{proof}
All pair stocks are nonnegative and dominate their selected
coordinate contributions.  Hence the tree
\(T_*=\{ab,ad,ce,ac\}\) has monomial at least
\[
 \theta_{ab,u}\theta_{ad,v}
 \theta_{ce,w}\theta_{ac,r}.
\]
Substitution of the four selected local masses gives exactly
\eqref{eq:V44-selected-tree}.  Extra coordinate intersections can
only add nonnegative terms to the complete tree sum.
\end{proof}

\begin{remark}[No endpoint saturation is used]
\label{rem:V44-no-leaf-saturation}
Unlike \eqref{eq:V43-tree-monomial}, the selected monomial retains
\(\Xi_b,\Xi_d,\Xi_e\).  There is no assumption that
\(y_u=\Xi_b\), \(y_v=\Xi_d\), or \(y_w=\Xi_e\).
\end{remark}

\section{Spectator-stable output compression}
\label{sec:V44-output}

\begin{lemma}[The final carrier remains one-clock bounded]
\label{lem:V44-spectator-payer}
For every resolved binary-spine word, after the three selected prefix
covariances have been left outside the estimate, the combined norm of
all spectator moments, the complete root carrier, the final BFA mass,
the reduced resolvent, \(Q_0\), and exact physical duality is at most
\begin{equation}
\label{eq:V44-spectator-payer}
 C_{\sigma,x_\pm}\alpha.
\end{equation}
The bound is uniform in the size and representation content of the
spectator banks.
\end{lemma}

\begin{proof}
Complete Wilson moments are contractions.  Their tensor product,
canonical row regrouping, exact isotypic pinching, and the physical
duality maps are contractions before the final scalar BFA weight.

By \eqref{eq:V44-maximal-row}, the total input Casimir over all five
rows is at most \(5M\).  Fusion subadditivity therefore gives
\[
 C_E(\boldsymbol Q)\le C M\le C_{\sigma,x_\pm}\alpha
\]
for every final output.  On a nontrivial output, use the
representation-uniform product gap
\eqref{eq:V41-product-gap}:
\[
 {\Xi(\boldsymbol Q)\over1-q_{\alpha,\boldsymbol Q}}
 \le C_{\sigma,x_\pm}\alpha.
\]
The root connected carrier has fixed-order norm bounded by its Bell
coefficient sum, and \(Q_0\) deletes the trivial output.  Exact
Fourier normalization cancels output dimensions into raw pinching
mass, whose total is at most one.  No spectator coordinate count is
created.  This proves \eqref{eq:V44-spectator-payer}.
\end{proof}

\section{Exact cancellation of distributed endpoint denominators}
\label{sec:V44-cancellation}

\begin{theorem}[Distributed-endpoint binary-spine estimate]
\label{thm:V44-distributed-close}
Normalize the five pure input row blocks of a resolved binary-spine
word to have BFA norm one.  Under
\eqref{eq:V44-maximal-row}--\eqref{eq:V44-root-hypotheses},
\begin{equation}
\label{eq:V44-distributed-bound}
 \|\mathcal W_{\alpha}^{\rm spine}\|_{\rho,2}^{\rm BFA}
 \le
 {C_{\rho,\sigma,x_\pm}\over A}\,
 {\mathfrak t}_{\rm sel}
 \le
 {C_{\rho,\sigma,x_\pm}\over A}
 \sum_{\tau\in\mathfrak T_5}
 \prod_{ik\in E(\tau)}\theta_{ik}.
\end{equation}
The constant is independent of every spectator support, endpoint row
mass, representation label, output multiplicity, and \(\alpha\).
\end{theorem}

\begin{proof}
At the three selected prefix coordinates,
\Cref{cor:V43-Casimir-debit} gives
\[
 \|{\mathcal C}_{a b,u}\|
 \le Cs_\alpha\sqrt{x_uy_u},
 \quad
 \|{\mathcal C}_{a d,v}\|
 \le Cs_\alpha\sqrt{x_vy_v},
 \quad
 \|{\mathcal C}_{c e,w}\|
 \le Cs_\alpha\sqrt{x_wy_w}.
\]
The harmless additions of one to the local Casimirs have been
absorbed in the constant.

After input normalization, the five row-mass denominators are
\(\Xi_a\Xi_b\Xi_c\Xi_d\Xi_e\).  Exact height subadditivity bounds the
exponential quotient by one.  Apply
\Cref{lem:V44-spectator-payer}, cancel all output dimensions by the
exact Fourier product formula, and sum raw pinching masses.  This
gives
\begin{equation}
\label{eq:V44-word-bound}
 \|\mathcal W_{\alpha}^{\rm spine}\|_{\rho,2}^{\rm BFA}
 \le
 {C\alpha s_\alpha^3
  \sqrt{x_uy_ux_vy_vx_wy_w}
  \over
  \Xi_a\Xi_b\Xi_c\Xi_d\Xi_e}.
\end{equation}
Divide by \eqref{eq:V44-selected-tree}.  Every endpoint denominator
\(\Xi_b,\Xi_d,\Xi_e\) cancels exactly, leaving
\begin{equation}
\label{eq:V44-exact-ratio}
 {\|\mathcal W_{\alpha}^{\rm spine}\|_{\rho,2}^{\rm BFA}
  \over{\mathfrak t}_{\rm sel}}
 \le
 C\alpha s_\alpha^3
 {\Xi_a^2\Xi_c\over A^2}
 {1\over\sqrt{x_uy_ux_vy_vx_wy_w}}.
\end{equation}
All six local masses are at least one.  The scale relations
\eqref{eq:V44-root-scales} make the remaining polynomial factor at
most \(C/A\).  This proves the first inequality in
\eqref{eq:V44-distributed-bound}.  The second follows from
\Cref{lem:V44-selected-tree}.
\end{proof}

\begin{corollary}[Spectator banks cannot create the transition
counterpacket]
\label{cor:V44-spectators-closed}
Arbitrary off-coordinate mass on \(b,d,e\), including diffuse or
cofinal spectator banks, does not reopen the binary-spine transition
family so long as those banks enter the resolved word through
complete moment contractions disjoint from the joining and payer.
\end{corollary}

\begin{proof}
The spectator masses enter both sides of
\eqref{eq:V44-exact-ratio} only through the row denominators, which
cancel.  Their local operators and raw output compression were
already included in \Cref{lem:V44-spectator-payer}.
\end{proof}

\begin{assessment}[Correction of the V43 residual list]
\label{ass:V44-correction}
Item \textup{(i)} in \Cref{ass:V43-residual} is not by itself a
residual: substantial mass outside a selected weak-edge coordinate
does not cause a normalized loss.  It becomes relevant only if that
mass carries another joining or payer role, so that it is no longer a
spectator contraction.
\end{assessment}

\section{A cut formulation}
\label{sec:V44-cut}

\begin{definition}[Binary-separable predecessor edge]
\label{def:V44-binary-separable}
A predecessor edge \(ik\) is binary-separable in a resolved word if
there is a coordinate \(e\) such that:
\begin{enumerate}[label=\textup{(\roman*)},leftmargin=3em]
\item the complete assembled factor selected at \(e\) is the binary
      covariance of rows \(i,k\);
\item no third row participates in that local cumulant block;
\item \(e\) is not the final payer coordinate; and
\item every other local factor at \(e\) tensor-separates from this
      covariance before the output pinching.
\end{enumerate}
\end{definition}

\begin{proposition}[Three binary-separable predecessor edges close]
\label{prop:V44-three-separable}
Suppose a saturated parabolic \(3+2\) first-connectivity fibre admits
a predecessor spanning forest with three binary-separable edges on
distinct coordinates, and its final joining coordinate supplies the
fourth tree edge.  Then every resolved word in that fibre satisfies
its marked-tree payment with an additional \(O(A^{-1})\) factor.
\end{proposition}

\begin{proof}
Relabel the two predecessor components and their internal tree edges
as \(ab,ad\) and \(ce\).  The proof of
\Cref{thm:V44-distributed-close} uses only binary separability,
nonnegativity of the selected pair-stock contributions, maximal-row
mass, root saturation, and exact output compression.  It does not use
the absence of additional spectator coordinates.  Apply that proof
after relabelling.
\end{proof}

\begin{theorem}[Higher-incidence necessity]
\label{thm:V44-higher-incidence-necessary}
After V42--V44, every persistent saturated parabolic \(3+2\)
transition fibre must fail to possess a three-edge
binary-separable predecessor forest.  Equivalently, in every
predecessor spanning forest at least one required connection is
carried by a local block which:
\begin{enumerate}[label=\textup{(\alph*)},leftmargin=3em]
\item has incidence at least three;
\item shares a coordinate with the final payer or joining carrier; or
\item remains entangled with another local block through the physical
output compression.
\end{enumerate}
\end{theorem}

\begin{proof}
If a three-edge binary-separable predecessor forest existed,
\Cref{prop:V44-three-separable} would pay the complete word.  Hence a
persistent residual must violate binary separability on at least one
edge.  The negation of
\Cref{def:V44-binary-separable}, after excluding a mere spectator
factor by \Cref{cor:V44-spectators-closed}, is exactly the three cases
listed.
\end{proof}

\begin{assessment}[Exact V45 local target]
\label{ass:V44-V45-target}
The next local object is no longer a distributed pair covariance.
It is the smallest growing higher-incidence block, beginning with a
three-row Wilson cumulant whose two required predecessor edges share
one coordinate.  The needed estimate must preserve both incident
marks, exact multiplicity compression, and at most one final payer.
\end{assessment}

\section{V44 result ledger and V45 audit programme}
\label{sec:V44-ledger}

\begin{theorem}[Fourth-series V44 result ledger]
\label{thm:V44-results}
Relative to V1--V43, V44 proves:
\begin{enumerate}[label=\textup{(V44.\arabic*)},leftmargin=6.3em]
\item a binary-spine formulation allowing arbitrary spectator
      supports and endpoint row masses;
\item a nonnegative selected-tree lower bound retaining every full
      endpoint denominator;
\item support-uniform stability of the one-clock final
      mass--resolvent carrier under spectator moments;
\item exact cancellation of all three distributed endpoint
      denominators in the word/tree ratio;
\item uniform binary-spine marked-tree closure with an additional
      inverse root scale;
\item a coordinate-free binary-separability criterion; and
\item necessity of a genuinely higher-incidence, payer-coincident, or
      compression-entangled predecessor connection in every
      persistent \(3+2\) transition fibre.
\end{enumerate}
\end{theorem}

\begin{proof}
These are
\Cref{lem:V44-selected-tree,
lem:V44-spectator-payer,
thm:V44-distributed-close,
cor:V44-spectators-closed,
prop:V44-three-separable,
thm:V44-higher-incidence-necessary}.
\end{proof}

\begin{programme}[V45: companion audit and growing ternary-prefix
carrier]
\label{prog:V44-V45}
\begin{enumerate}[leftmargin=3em]
\item Perform the compulsory read-only audit of the newest active SM
      and NS notes and record exact file identities.
\item Form the complete three-row centered Wilson cumulant on one
      coordinate before any norm or output-channel sum.
\item Derive a dimension-free two-mark estimate from its covariance
      decompositions, or exhibit a sharp obstruction if two
      square-root debits cannot be multiplied.
\item Retain the exact Peter--Weyl partial trace so that output
      multiplicities aggregate as raw pinching mass.
\item Insert the ternary carrier into a \(3+2\) predecessor word and
      compare it with the two incident selected pair stocks.
\item Treat payer coincidence only after the unpaid ternary estimate
      is settled.
\end{enumerate}
\end{programme}

\begin{firewall}[Claim boundary after fourth-series V44]
\label{fire:V44-claim-boundary}
V44 closes every saturated binary-separable \(3+2\) transition word,
including arbitrary distributed spectator mass.  It does not close
growing higher-incidence local cumulants, payer-coincident ternary
carriers, or every compression-entangled prefix; hence it does not
prove the full quintic MTA, an all-order atlas, WUFC, CPM, cutoff
removal, reflection positivity, Osterwalder--Schrader reconstruction,
construction of the continuum Yang--Mills measure, or a uniform
positive continuum spectral gap.  The full Yang--Mills mass gap has
not been proved.
\end{firewall}

\paragraph{Parent-version record.}
V44 was formed from
\texttt{Renewal\_Yang\_Mills\_Mass\_Gap\_Research\_Notes\_V43.tex}.
The V43 parent had \(21\,467\) logical source lines,
\(748\,970\) bytes, and SHA--256
\[
 \texttt{a0594b143d0c3d8b42e2e059418cba559615a0b20d9836b62e8121c54cfe9ef6}.
\]
The compulsory companion audit is due in V45.

% =============================================================================
% END APPENDED FOURTH-SERIES V44 MATERIAL
% =============================================================================

% =============================================================================
% BEGIN APPENDED FOURTH-SERIES V45 MATERIAL
% =============================================================================

\chapter{V45: Companion audit and a two-mark ternary Wilson carrier}
\label{chap:V45-ternary-carrier}

\section{Compulsory read-only companion audit}
\label{sec:V45-audit}

The newest active companion files at the V45 boundary are
\begin{align*}
 &\texttt{Renewal\_SM\_Einstein\_Continuum\_Research\_Notes\_V37.tex},\\
 &\texttt{Renewal\_NS\_Stability\_Research\_Notes\_V31.tex}.
\end{align*}
They were inspected read-only and were not modified.  Their exact
identities are
\begin{center}
\begin{tabular}{c|r|r|l}
programme & logical lines & bytes & SHA--256 prefix\\ \hline
SM V37 & \(13\,907\) & \(477\,719\) &
\texttt{0c9b1296eb10903f}\\
NS V31 & \(23\,052\) & \(887\,537\) &
\texttt{f1121b62238ad549}
\end{tabular}
\end{center}
The complete digests are
\[
\begin{split}
&\texttt{0c9b1296eb10903f065e3c3a83e27bc0fe2a8a8fd40a3fc60b0838c3985e07db},\\
&\texttt{f1121b62238ad5491bb7cf03635ea9934942edf573ed8faaa9ee79e1d38a6696}.
\end{split}
\]

SM V37 derives a connected interpolation identity and, at its
Gaussian endpoint, assigns one covariance to every external massive
leg before any norm is taken.  It also distinguishes rooted local
decoupling from a global trace-multiplicity statement.  The relevant
YM lesson is to expose selected centered legs inside the complete
local cumulant and to aggregate output multiplicities only after
exact physical compression.

NS V31 proves that heat formation alone is scale-sharp and that a
signed hinge controls only its polarization sector.  The relevant YM
lesson is that smoothing cannot replace the signed ternary
moment--cumulant identity.  Accordingly, the estimate below begins
with the complete centered ternary carrier and only then takes an
operator norm.

\begin{firewall}[Companion status]
\label{fire:V45-companion-status}
The companion files supply direction checks, not premises.  The
ternary identity and all norm estimates used below are proved
independently in this chapter.
\end{firewall}

\section{The complete ternary cumulant is a centered triple moment}
\label{sec:V45-centered-triple}

Let \(\pi_i\) be unitary irreducible representations on \(H_i\) for
\(i=1,2,3\), evaluated at one common Wilson variable \(U\).  Put
\[
 A_i(U):=\pi_i(U),
 \qquad
 \mathbb E_\alpha A_i=q_iI_i,
 \qquad
 X_i:=A_i-q_iI_i.
\]
Central inversion invariance makes \(q_i\in[-1,1]\).

\begin{lemma}[Exact centered form of the third cumulant]
\label{lem:V45-centered-triple}
After the canonical tensor rebracketing,
\begin{equation}
\label{eq:V45-centered-triple}
 \kappa_3^\alpha(A_1,A_2,A_3)
 =
 \mathbb E_\alpha[X_1\otimes X_2\otimes X_3].
\end{equation}
\end{lemma}

\begin{proof}
The third moment--cumulant formula is
\begin{align*}
 \kappa_3^\alpha(A_1,A_2,A_3)
 ={}&
 \mathbb E[A_1\otimes A_2\otimes A_3]\\
 &-\mathbb E[A_1\otimes A_2]\otimes\mathbb E[A_3]
 -\mathbb E[A_1\otimes A_3]\otimes\mathbb E[A_2]\\
 &-\mathbb E[A_2\otimes A_3]\otimes\mathbb E[A_1]
 +2\mathbb E[A_1]\otimes\mathbb E[A_2]\otimes\mathbb E[A_3].
\end{align*}
Expand the right side of
\eqref{eq:V45-centered-triple}.  Its one-constant terms give the
three negative pair-moment products.  Its two-constant terms occur
with positive sign three times, while its three-constant term occurs
with negative sign; their net coefficient is \(3-1=2\).
This is exactly the displayed cumulant formula.
\end{proof}

\section{A dimension-free two-selected-leg bound}
\label{sec:V45-two-leg}

\begin{theorem}[Two-mark ternary character-gap debit]
\label{thm:V45-two-mark-debit}
For any two distinct selected indices \(i,k\in\{1,2,3\}\),
\begin{equation}
\label{eq:V45-two-mark-debit}
 \boxed{
 \|\kappa_3^\alpha(A_1,A_2,A_3)\|
 \le
 2\sqrt{1-q_i^2}\sqrt{1-q_k^2}
 \le
 4\sqrt{(1-q_i)(1-q_k)}.}
\end{equation}
The same bound holds after every multiplicity-resolved isotypic
compression.  Its constant is independent of all three
representation dimensions and of the unselected representation.
\end{theorem}

\begin{proof}
By symmetry take \(i=1,k=2\).  The centered factors satisfy
\[
 \|X_3(U)\|\le\|\pi_3(U)\|+|q_3|\le2.
\]
Apply the operator-valued Cauchy--Schwarz construction from
\Cref{lem:V43-operator-CS} to
\[
 X=X_1,
 \qquad
 Y=X_2\otimes X_3.
\]
The first quadratic average is
\[
 \mathbb E[X_1X_1^*]=(1-q_1^2)I_1.
\]
Moreover,
\[
 (X_2\otimes X_3)^*(X_2\otimes X_3)
 =
 X_2^*X_2\otimes X_3^*X_3
 \preccurlyeq4X_2^*X_2\otimes I_3.
\]
Taking expectations gives operator norm at most
\(4(1-q_2^2)\).  Thus
\[
 \|\mathbb E[X_1\otimes X_2\otimes X_3]\|
 \le
 2\sqrt{1-q_1^2}\sqrt{1-q_2^2}.
\]
Use \Cref{lem:V45-centered-triple} and
\(1-q^2\le2(1-q)\).  Orthogonal compression cannot increase the
norm.
\end{proof}

\begin{corollary}[Wilson--Casimir two-mark form]
\label{cor:V45-two-mark-Casimir}
For every sufficiently large \(\alpha\), every three
representations, and any two selected legs \(i,k\),
\begin{equation}
\label{eq:V45-two-mark-Casimir}
 \|\kappa_3^\alpha(A_1,A_2,A_3)\|
 \le
 C
 \sqrt{\min\{s_\alpha c_i,1\}}
 \sqrt{\min\{s_\alpha c_k,1\}}
 \le
 Cs_\alpha\sqrt{c_ic_k}.
\end{equation}
\end{corollary}

\begin{proof}
Insert the representation-uniform upper character-gap estimate in
\eqref{eq:V45-two-mark-debit}, then use
\(\min\{x,1\}\le x\).
\end{proof}

\begin{remark}[The unselected center may be fully hot]
\label{rem:V45-hot-center}
No character gap of the third leg appears in
\eqref{eq:V45-two-mark-Casimir}.  Thus the estimate remains valid
when the unselected center representation has Casimir comparable to
or larger than \(\alpha\).  This is the precise advantage of spending
the two covariance debits on the two incident leaf legs.
\end{remark}

\section{A ternary--binary predecessor spine}
\label{sec:V45-spine}

Let the five rows be \(a,b,c,d,e\).  Choose distinct physical
coordinates \(t,w,r\) with:
\begin{enumerate}[label=\textup{(\roman*)},leftmargin=3em]
\item a complete ternary cumulant on rows \(a,b,d\) at \(t\);
\item a complete binary covariance on rows \(c,e\) at \(w\); and
\item the final \(ac\) joining carrier at \(r\).
\end{enumerate}
These local blocks form the predecessor partition
\[
 \rho=\{\{a,b,d\},\{c,e\}\}.
\]
Allow arbitrary spectator moment banks on all rows, disjoint from the
three displayed local carriers and the sole final payer.

Write the local balanced masses at \(t\) as
\[
 x:=1+c_{R_{a,t}},
 \qquad
 y:=1+c_{R_{b,t}},
 \qquad
 z:=1+c_{R_{d,t}},
\]
and those at \(w\) as
\[
 u:=1+c_{R_{c,w}},
 \qquad
 v:=1+c_{R_{e,w}}.
\]
At \(r\), use the equal-spin mass \(A=1+c_j\) on \(a,c\).  Assume
\begin{equation}
\label{eq:V45-spine-hypotheses}
 \Xi_i\le\Xi_a=:M,
 \qquad
 A\ge\sigma\max\{\Xi_a,\Xi_c\},
 \qquad
 0<x_-\le {d_j\over\sqrt\alpha}\le x_+<\infty.
\end{equation}
Consequently
\[
 \Xi_a,\Xi_c,M,\alpha\asymp A,
 \qquad
 s_\alpha\asymp A^{-1},
\]
with constants depending only on
\(\sigma,x_-,x_+\).

\begin{lemma}[Selected tree for the ternary--binary spine]
\label{lem:V45-selected-tree}
The complete tree sum is bounded below by
\begin{equation}
\label{eq:V45-selected-tree}
 {\mathfrak t}_{3,2}
 :=
 {A^2x^2yzuv
  \over
  \Xi_a^3\Xi_c^2\Xi_b\Xi_d\Xi_e}.
\end{equation}
\end{lemma}

\begin{proof}
At coordinate \(t\), select the two tree edges \(ab,ad\); repeated
use of the same physical coordinate by distinct incident tree edges
is allowed in the marked seminorm.  Their local pair-stock
contributions are
\[
 {xy\over\Xi_a\Xi_b},
 \qquad
 {xz\over\Xi_a\Xi_d}.
\]
At \(w\), select \(ce\), with contribution
\(uv/(\Xi_c\Xi_e)\).  At \(r\), select \(ac\), with contribution
\(A^2/(\Xi_a\Xi_c)\).  Multiply the four contributions.  Global pair
stocks dominate their selected local contributions, and all other
tree monomials are nonnegative.
\end{proof}

\section{Complete BFA payment}
\label{sec:V45-BFA}

\begin{theorem}[Uniform ternary--binary transition payment]
\label{thm:V45-ternary-binary-close}
Let \(\mathcal W_\alpha^{3,2}\) be a resolved word with the
ternary--binary predecessor spine above, including arbitrary
spectator moment banks, every physical output and multiplicity,
exact Fourier normalization and duality, \(Q_0\), one final reduced
resolvent, and the output BFA norm.  Normalize each pure input row
block to have BFA norm one.  Then
\begin{equation}
\label{eq:V45-ternary-binary-close}
 \|\mathcal W_\alpha^{3,2}\|_{\rho,2}^{\rm BFA}
 \le
 C_{\rho,\sigma,x_\pm}
 {\mathfrak t}_{3,2}.
\end{equation}
The constant is uniform in all representation labels, spectator
supports, output labels, multiplicities, and \(\alpha\).
\end{theorem}

\begin{proof}
Apply \Cref{cor:V45-two-mark-Casimir} to the ternary block, selecting
the \(b,d\) legs:
\[
 \|\kappa_3^\alpha(A_a,A_b,A_d)\|
 \le Cs_\alpha\sqrt{yz}.
\]
The additions of one to the Casimirs only enlarge the right side.
Apply the binary estimate
\Cref{cor:V43-Casimir-debit} at \(w\):
\[
 \|{\mathcal C}_{c e,w}\|
 \le Cs_\alpha\sqrt{uv}.
\]

All five row masses are at most \(M\asymp A\), so the total output
Casimir is \(O(A)=O(\alpha)\).  Exactly as in
\Cref{lem:V44-spectator-payer}, the complete final
BFA-mass--resolvent--root carrier is bounded by \(C\alpha\).
Spectator moments, physical duality, and height quotient are
contractions or bounded by one.  Exact Fourier dimension
cancellation leaves raw output pinching mass of total at most one.
After the five input BFA normalizations,
\begin{equation}
\label{eq:V45-word-bound}
 \|\mathcal W_\alpha^{3,2}\|_{\rho,2}^{\rm BFA}
 \le
 {C\alpha s_\alpha^2\sqrt{yzuv}
  \over
  \Xi_a\Xi_b\Xi_c\Xi_d\Xi_e}.
\end{equation}

Divide by \eqref{eq:V45-selected-tree}.  Every endpoint denominator
cancels, and
\begin{equation}
\label{eq:V45-word-tree-ratio}
 {\|\mathcal W_\alpha^{3,2}\|_{\rho,2}^{\rm BFA}
  \over{\mathfrak t}_{3,2}}
 \le
 C\alpha s_\alpha^2
 {\Xi_a^2\Xi_c\over A^2}
 {1\over x^2\sqrt{yzuv}}.
\end{equation}
All five local masses are at least one.  The scale relations in
\eqref{eq:V45-spine-hypotheses} give
\[
 \alpha s_\alpha^2
 {\Xi_a^2\Xi_c\over A^2}
 \le C.
\]
This proves \eqref{eq:V45-ternary-binary-close}.
\end{proof}

\begin{corollary}[A growing ternary prefix can pay two incident marks]
\label{cor:V45-ternary-two-marks}
In a saturated parabolic \(3+2\) fibre, a distinct-coordinate ternary
prefix block connecting the three-row predecessor component supplies
the two incident marked-tree edges, even when all three of its
representations grow with \(\alpha\).  No fixed-bank expansion or
output multiplicity count is required.
\end{corollary}

\begin{proof}
The two selected legs in
\Cref{thm:V45-two-mark-debit} supply the two character-gap square
roots.  Their exact normalized power ledger is
\Cref{thm:V45-ternary-binary-close}.
\end{proof}

\begin{remark}[The estimate is tight at the ledger level]
\label{rem:V45-no-surplus}
Unlike the binary-spine estimate of V44, the bound
\eqref{eq:V45-word-tree-ratio} generally leaves no inverse power of
\(A\).  A ternary block pays two tree edges with one factor
\(s_\alpha\), while the remaining binary prefix pays the third edge
with a second factor.  The final hot payer consumes one clock.  The
row-mass quotient consumes the remaining scale.  No unproved surplus
is claimed.
\end{remark}

\section{Remaining higher-incidence geometries}
\label{sec:V45-residual}

\begin{theorem}[Distinct-coordinate \(3+2\) predecessor carriers are
closed]
\label{thm:V45-distinct-32}
Suppose a saturated parabolic \(3+2\) predecessor fibre admits either:
\begin{enumerate}[label=\textup{(\alph*)},leftmargin=3em]
\item three binary-separable predecessor edges; or
\item one ternary block spanning the three-row component and one
      binary-separable block spanning the two-row component,
\end{enumerate}
all disjoint from the final joining/payer coordinate.  Then its
resolved BFA word satisfies the required marked-tree payment.
\end{theorem}

\begin{proof}
Case \textup{(a)} is
\Cref{prop:V44-three-separable}.  Case \textup{(b)} is
\Cref{thm:V45-ternary-binary-close}, after relabelling the rows.
\end{proof}

\begin{assessment}[Exact residual after V45]
\label{ass:V45-residual}
For a \(3+2\) predecessor partition, the only connected component
types are a three-row tree and a two-row edge.  V44 and V45 close both
ways of realizing the three-row tree on coordinates disjoint from the
final payer: two binary carriers or one ternary carrier.  A
persistent residual must therefore involve a true role coincidence:
the higher-incidence prefix shares the selected payer or final
joining coordinate, or its physical compression cannot be separated
from that operation before the norm.
\end{assessment}

\begin{firewall}[No automatic payer-coincident extension]
\label{fire:V45-payer-coincidence}
The proof leaves the ternary carrier outside the final
mass--resolvent estimate.  If the payer and ternary carrier act in the
same unresolved multiplicity space, multiplying their separate
operator norms may lose the two selected-leg structure.  Such a
coincident carrier must be assembled and estimated as one operator;
\Cref{thm:V45-ternary-binary-close} is not cited for it.
\end{firewall}

\section{V45 result ledger and V46 programme}
\label{sec:V45-ledger}

\begin{theorem}[Fourth-series V45 result ledger]
\label{thm:V45-results}
Relative to V1--V44, V45 proves:
\begin{enumerate}[label=\textup{(V45.\arabic*)},leftmargin=6.3em]
\item the required read-only audit of active SM V37 and NS V31;
\item the exact centered-triple identity for the complete third
      Wilson cumulant;
\item a dimension-free two-selected-leg character-gap debit;
\item its representation-uniform two-mark Casimir form;
\item the exact selected tree monomial for a ternary--binary
      \(3+2\) spine with arbitrary spectator mass;
\item complete BFA payment of that spine, including raw output
      pinching and the single final resolvent; and
\item closure of every distinct-coordinate \(3+2\) predecessor
      carrier, leaving only true joining/payer/compression
      coincidences.
\end{enumerate}
\end{theorem}

\begin{proof}
These are
\Cref{sec:V45-audit,
lem:V45-centered-triple,
thm:V45-two-mark-debit,
cor:V45-two-mark-Casimir,
lem:V45-selected-tree,
thm:V45-ternary-binary-close,
thm:V45-distinct-32}.
\end{proof}

\begin{programme}[V46: payer-coincident ternary carrier]
\label{prog:V45-V46}
\begin{enumerate}[leftmargin=3em]
\item Write the complete multiplicity-resolved operator obtained when
      the ternary prefix output also carries the selected Casimir
      payer or participates in the final joining projection.
\item Commute no factor through a Racah projection unless exact
      equivariance proves the commutation.
\item Seek a weighted operator Cauchy--Schwarz inequality in which the
      output Casimir divided by the product character gap remains
      inside the quadratic form.
\item Test the estimate on the \(SU(2)\) Cartan and protected endpoint
      channels from V40.
\item If the weighted two-leg estimate fails, return upstream to the
      global final-output Casimir sum and use raw pinching before the
      local norm.
\item Propagate the result through every remaining \(3+2\) role
      placement before considering other predecessor partitions.
\end{enumerate}
\end{programme}

\begin{firewall}[Claim boundary after fourth-series V45]
\label{fire:V45-claim-boundary}
V45 closes every distinct-coordinate binary or ternary realization of
the saturated parabolic \(3+2\) predecessor partition.  It does not
close payer-coincident or joining-entangled ternary carriers, every
other predecessor partition, or the full quintic MTA; hence it does
not prove an all-order atlas, WUFC, CPM, cutoff removal, reflection
positivity, Osterwalder--Schrader reconstruction, construction of the
continuum Yang--Mills measure, or a uniform positive continuum
spectral gap.  The full Yang--Mills mass gap has not been proved.
\end{firewall}

\paragraph{Parent-version record.}
V45 was formed from
\texttt{Renewal\_Yang\_Mills\_Mass\_Gap\_Research\_Notes\_V44.tex}.
The V44 parent had \(21\,895\) logical source lines,
\(764\,842\) bytes, and SHA--256
\[
 \texttt{5b411916ba69c74a3b44294546df8bab292506012126461d8e476a2238edc01d}.
\]
The next compulsory companion audit is V50.

% =============================================================================
% END APPENDED FOURTH-SERIES V45 MATERIAL
% =============================================================================

% =============================================================================
% BEGIN APPENDED FOURTH-SERIES V46 MATERIAL
% =============================================================================

\chapter{V46: Selected-leg cumulant bounds survive a coincident
Casimir payer}
\label{chap:V46-payer-coincidence}

\section{Purpose}
\label{sec:V46-purpose}

V45 proved the required two-mark estimate with the ternary prefix
carrier left outside the final mass--resolvent bound.  This did not
yet justify a placement in which the selected Casimir payer acts in
the same unresolved output multiplicity space.

The needed repair has two parts.  First, the two-selected-leg
character-gap estimate extends from the third cumulant to every fixed
cumulant order.  Second, an equivariant carrier is block diagonal on
complete isotypic spaces, and a nonnegative Casimir/resolvent weight
can be retained inside the multiplicity partial trace without a
multiplicity factor.  Together these facts show that payer
coincidence does not erase the V45 debits.

\section{Selected centered moments}
\label{sec:V46-centered-moments}

Let \(A_i(U)=\pi_i(U)\) be unitary representations of one common
central inversion-invariant random group element, and put
\[
 \mathbb E A_i=q_iI_i,
 \qquad
 X_i=A_i-q_iI_i.
\]
Then
\begin{equation}
\label{eq:V46-centered-variance}
 \|X_i(U)\|\le2,
 \qquad
 \mathbb E[X_iX_i^*]=(1-q_i^2)I_i,
 \qquad
 \mathbb E[X_i^*X_i]=(1-q_i^2)I_i.
\end{equation}

\begin{lemma}[One- and two-selected-leg moment bounds]
\label{lem:V46-selected-moments}
Let \(B\) be a nonempty ordered set of centered factors.
For \(i\in B\),
\begin{equation}
\label{eq:V46-one-selected-moment}
 \left\|\mathbb E\bigotimes_{\nu\in B}X_\nu\right\|
 \le
 2^{|B|-1}\sqrt{1-q_i^2}.
\end{equation}
For distinct \(i,k\in B\),
\begin{equation}
\label{eq:V46-two-selected-moment}
 \left\|\mathbb E\bigotimes_{\nu\in B}X_\nu\right\|
 \le
 2^{|B|-2}
 \sqrt{1-q_i^2}\sqrt{1-q_k^2}.
\end{equation}
\end{lemma}

\begin{proof}
For \eqref{eq:V46-one-selected-moment}, apply the operator
Cauchy--Schwarz construction of
\Cref{lem:V43-operator-CS} with \(X=X_i\) and with the tensor product
of all other centered factors as the second field.  The second field
has pointwise norm at most \(2^{|B|-1}\), while the first quadratic
average is given by \eqref{eq:V46-centered-variance}.

For \eqref{eq:V46-two-selected-moment}, take \(X=X_i\) and put
\(X_k\) into the second field, together with the remaining
\(|B|-2\) factors.  Pointwise,
\[
 \left(\bigotimes_{\nu\in B\setminus\{i\}}X_\nu\right)^*
 \left(\bigotimes_{\nu\in B\setminus\{i\}}X_\nu\right)
 \preccurlyeq
 2^{2(|B|-2)}
 X_k^*X_k\otimes I.
\]
The second quadratic average therefore has norm at most
\[
 2^{2(|B|-2)}(1-q_k^2).
\]
Combine this with the \(i\)-variance and take square roots.
\end{proof}

\section{Every fixed-order cumulant has two selectable debits}
\label{sec:V46-general-cumulant}

\begin{lemma}[Shift invariance]
\label{lem:V46-shift-invariance}
For \(m\ge2\),
\begin{equation}
\label{eq:V46-shift-invariance}
 \kappa_m^\alpha(A_1,\ldots,A_m)
 =
 \kappa_m^\alpha(X_1,\ldots,X_m).
\end{equation}
\end{lemma}

\begin{proof}
Cumulants are multilinear, and a cumulant of order at least two with
one deterministic scalar-identity argument is zero.  Expand
\(A_i=X_i+q_iI_i\) in every argument.
\end{proof}

\begin{theorem}[Dimension-free two-selected-leg cumulant theorem]
\label{thm:V46-general-two-leg}
For every fixed \(m\ge2\), there is a finite numerical constant
\(C_m\) such that, for any two distinct selected indices \(i,k\),
\begin{equation}
\label{eq:V46-general-two-leg}
 \boxed{
 \|\kappa_m^\alpha(A_1,\ldots,A_m)\|
 \le
 C_m
 \sqrt{1-q_i^2}\sqrt{1-q_k^2}.}
\end{equation}
The estimate is uniform in all representation dimensions and in
every multiplicity-resolved output compression.
\end{theorem}

\begin{proof}
Use \Cref{lem:V46-shift-invariance} and expand the centered cumulant:
\[
 \kappa_m^\alpha(X_1,\ldots,X_m)
 =
 \sum_{\pi\in\Pi(I_m)}
 \mu(\pi,\widehat1_m)
 \bigotimes_{B\in\pi}
 \mathbb E\bigotimes_{\nu\in B}X_\nu.
\]
Every partition with a singleton block vanishes because
\(\mathbb E X_\nu=0\).  Fix a remaining partition.

If \(i,k\) lie in the same block, apply
\eqref{eq:V46-two-selected-moment} to that block and bound every other
centered moment by the product of the pointwise norms, at most a
fixed power of two.

If \(i,k\) lie in distinct blocks, apply
\eqref{eq:V46-one-selected-moment} to each of those two blocks and
the pointwise norm bound to all other blocks.  In both cases the
partition term is at most a constant depending only on \(m\), times
the right side of \eqref{eq:V46-general-two-leg}.
The finite sum of absolute M\"obius coefficients defines \(C_m\).
Orthogonal compression cannot increase the operator norm.
\end{proof}

\begin{corollary}[Wilson--Casimir form at every fixed order]
\label{cor:V46-general-Casimir}
For every fixed \(m\ge2\), every sufficiently large \(\alpha\), and
any two selected legs \(i,k\),
\begin{equation}
\label{eq:V46-general-Casimir}
 \|\kappa_m^\alpha(A_1,\ldots,A_m)\|
 \le
 C_m s_\alpha\sqrt{c_ic_k}.
\end{equation}
\end{corollary}

\begin{proof}
Combine \eqref{eq:V46-general-two-leg} with the uniform upper
character-gap estimate, as in
\Cref{cor:V45-two-mark-Casimir}.
\end{proof}

\begin{firewall}[Only two debits are asserted]
\label{fire:V46-only-two}
For \(m\ge4\), \Cref{thm:V46-general-two-leg} does not supply
\(m-1\) independent Wilson clocks.  The proof uses quadratic
Cauchy--Schwarz on two selected legs and bounds every other centered
leg in operator norm.  Additional marked-tree edges require either
other local carriers, a fixed-bank connected degree, or a stronger
multi-leg estimate.
\end{firewall}

\section{Equivariant multiplicity blocks and weighted pinching}
\label{sec:V46-weighted-pinching}

Let
\[
 \mathscr H
 =
 \bigoplus_T W_T(H_T\otimes M_T)
\]
be a complete multiplicity-resolved decomposition for the diagonal
group action.

\begin{lemma}[Central cumulants are equivariant]
\label{lem:V46-equivariant}
Every complete local Wilson moment and cumulant \(K\) commutes with
the diagonal group action.  Consequently
\begin{equation}
\label{eq:V46-equivariant-block}
 W_T^*KW_T=I_{H_T}\otimes K_T
\end{equation}
on each isotypic block.  Total Casimir functional calculus and every
central Wilson multiplier commute with \(K\).
\end{lemma}

\begin{proof}
Conjugating every representation factor by the diagonal action
replaces the integration variable \(U\) by \(gUg^{-1}\).  The Wilson
law is central, so each complete moment is invariant under this
change of variable.  Finite linear combinations and products
preserve equivariance, hence so do cumulants.

Schur's lemma on a complete isotypic block gives
\eqref{eq:V46-equivariant-block}.  The total Casimir and central
Wilson multipliers are scalar on \(H_T\) and the identity on \(M_T\),
so they commute with every operator of the displayed form.
\end{proof}

\begin{lemma}[No-count weighted multiplicity partial trace]
\label{lem:V46-weighted-pinching}
Let \(X_T\) be trace class on \(H_T\otimes M_T\), let \(K_T\) be
bounded on \(M_T\), and let \(0\le w_T\le W\).  Then
\begin{equation}
\label{eq:V46-one-block-trace}
 \left\|
 \operatorname{Tr}_{M_T}
 \left[X_T(I_{H_T}\otimes K_T)\right]
 \right\|_1
 \le
 \|X_T\|_1\|K_T\|.
\end{equation}
Consequently,
\begin{equation}
\label{eq:V46-weighted-pinching}
 \sum_T w_T
 \left\|
 \operatorname{Tr}_{M_T}
 \left[X_T(I_{H_T}\otimes K_T)\right]
 \right\|_1
 \le
 W\left(\sup_T\|K_T\|\right)
 \sum_T\|X_T\|_1.
\end{equation}
No factor \(\dim M_T\) occurs.
\end{lemma}

\begin{proof}
By trace-norm duality, the left side of
\eqref{eq:V46-one-block-trace} is
\[
 \sup_{\|B_T\|\le1}
 \left|
 \operatorname{Tr}_{H_T\otimes M_T}
 \left[
 X_T(B_T\otimes K_T)
 \right]\right|.
\]
This is at most
\(\|X_T\|_1\|K_T\|\), proving the one-block estimate.  Multiply by
\(w_T\), use the two suprema, and sum to obtain
\eqref{eq:V46-weighted-pinching}.
\end{proof}

\begin{corollary}[Racah changes do not affect the weighted estimate]
\label{cor:V46-Racah-stability}
If the isotypic decomposition is changed by an equivariant unitary
recoupling, \eqref{eq:V46-weighted-pinching} is unchanged.
\end{corollary}

\begin{proof}
Equivariant recoupling conjugates \(X_T,K_T\) unitarily on the
multiplicity space.  Trace and operator norms are unitarily
invariant, and the Casimir/resolvent scalar \(w_T\) is unchanged.
\end{proof}

\section{The coincident payer}
\label{sec:V46-coincident-payer}

For a complete product output \(\boldsymbol T\ne\boldsymbol1\), set
\begin{equation}
\label{eq:V46-payer-weight}
 w_{\alpha,\boldsymbol T}
 :=
 {\Xi(\boldsymbol T)\over
  1-q_{\alpha,\boldsymbol T}}.
\end{equation}
In the saturated corridor of V45, the product character gap gives
\begin{equation}
\label{eq:V46-payer-sup}
 0\le w_{\alpha,\boldsymbol T}\le C\alpha.
\end{equation}

\begin{theorem}[Two selected legs survive payer coincidence]
\label{thm:V46-payer-two-leg}
Let a complete fixed-order local cumulant
\(\kappa_m^\alpha(A_1,\ldots,A_m)\) act in the same unresolved
isotypic multiplicity space as the selected Casimir payer, the final
reduced resolvent, and an equivariant Racah recoupling.  On a
saturated corridor where every complete output Casimir is
\(O(\alpha)\), its full weighted, multiplicity-summed response obeys
\begin{equation}
\label{eq:V46-payer-two-leg}
 \mathfrak P_\alpha
 \le
 C_m\alpha
 \sqrt{1-q_i^2}\sqrt{1-q_k^2}
 \sum_{\boldsymbol T,\gamma}
 \|X_{\boldsymbol T,\gamma}\|_1
\end{equation}
for any two selected legs \(i,k\).  Equivalently,
\begin{equation}
\label{eq:V46-payer-Casimir}
 \mathfrak P_\alpha
 \le
 C_m\alpha s_\alpha\sqrt{c_ic_k}
 \sum_{\boldsymbol T,\gamma}
 \|X_{\boldsymbol T,\gamma}\|_1.
\end{equation}
There is no output multiplicity factor.
\end{theorem}

\begin{proof}
By \Cref{lem:V46-equivariant}, the complete cumulant has block form
\(I\otimes K_{\boldsymbol T,\gamma}\) and commutes with the payer and
resolvent scalar.  Its block operator norm is bounded by
\eqref{eq:V46-general-two-leg}.  Apply
\Cref{lem:V46-weighted-pinching} with
\(W=C\alpha\) from \eqref{eq:V46-payer-sup}.  Equivariant recoupling
causes no loss by \Cref{cor:V46-Racah-stability}.  This proves
\eqref{eq:V46-payer-two-leg}; the Wilson--Casimir form follows from
\Cref{cor:V46-general-Casimir}.
\end{proof}

\begin{corollary}[Payer-coincident ternary--binary spines close]
\label{cor:V46-ternary-payer-close}
The conclusion of
\Cref{thm:V45-ternary-binary-close} remains valid if the selected
Casimir payer is assigned to the ternary prefix coordinate, or if the
ternary output is recoupled nontrivially before the final product
pinching.
\end{corollary}

\begin{proof}
In the V45 ledger, replace the separate use of the ternary bound and
the \(C\alpha\) final carrier by
\Cref{thm:V46-payer-two-leg} with \(m=3\) and selected legs \(b,d\).
Its right side is exactly
\[
 C\alpha s_\alpha\sqrt{yz}
\]
times raw pinching mass.  Multiply by the remaining binary covariance
\(Cs_\alpha\sqrt{uv}\), divide by the same five input masses, and
repeat \eqref{eq:V45-word-tree-ratio}.  The result is unchanged.
\end{proof}

\begin{corollary}[Compression entanglement caused only by recoupling
is removed]
\label{cor:V46-compression-removed}
A V45 ternary output is not a residual merely because its
multiplicity space is expressed in a different equivariant Racah
basis at final output.
\end{corollary}

\begin{proof}
This is \Cref{cor:V46-Racah-stability} inserted into
\Cref{cor:V46-ternary-payer-close}.
\end{proof}

\section{Residual after the payer repair}
\label{sec:V46-residual}

\begin{theorem}[Only genuine same-coordinate joining remains in the
\(3+2\) census]
\label{thm:V46-32-residual}
Within the saturated parabolic \(3+2\) predecessor census, all
binary-separable, distinct ternary--binary, payer-coincident
ternary--binary, and equivariantly recoupled versions are paid by
V44--V46.  Any remaining local residual must have the
first-connectivity joining operation itself at the same physical
coordinate as one or more predecessor connections.
\end{theorem}

\begin{proof}
Binary-separable spines are paid by
\Cref{prop:V44-three-separable}.  Distinct ternary--binary spines are
paid by \Cref{thm:V45-ternary-binary-close}.  Payer coincidence and
equivariant compression are paid by
\Cref{cor:V46-ternary-payer-close,
cor:V46-compression-removed}.  These exhaust the three alternatives
of \Cref{thm:V44-higher-incidence-necessary} except coincidence with
the joining operation itself.
\end{proof}

\begin{assessment}[The next local carrier]
\label{ass:V46-next-carrier}
At a genuine same-coordinate joining, the local partition block
simultaneously connects predecessor components and may also realize
internal predecessor edges.  It cannot be factored into a prefix
cumulant times a joining cumulant.  The correct V47 object is the
complete local cumulant block, with a selected spanning tree inside
its row incidence and with the final resolvent retained once.
\end{assessment}

\section{V46 result ledger and V47 programme}
\label{sec:V46-ledger}

\begin{theorem}[Fourth-series V46 result ledger]
\label{thm:V46-results}
Relative to V1--V45, V46 proves:
\begin{enumerate}[label=\textup{(V46.\arabic*)},leftmargin=6.3em]
\item one- and two-selected-leg bounds for arbitrary centered local
      moments;
\item a dimension-free two-selected-leg theorem for every fixed
      cumulant order;
\item equivariant block structure of complete Wilson cumulants;
\item a no-count weighted multiplicity partial-trace inequality;
\item stability of that inequality under exact Racah recoupling;
\item preservation of two selected Wilson debits when the Casimir
      payer and reduced resolvent act in the same multiplicity space;
\item closure of every payer-coincident ternary--binary spine; and
\item reduction of the \(3+2\) residual to genuine same-coordinate
      first connectivity.
\end{enumerate}
\end{theorem}

\begin{proof}
These are
\Cref{lem:V46-selected-moments,
thm:V46-general-two-leg,
lem:V46-equivariant,
lem:V46-weighted-pinching,
cor:V46-Racah-stability,
thm:V46-payer-two-leg,
cor:V46-ternary-payer-close,
thm:V46-32-residual}.
\end{proof}

\begin{programme}[V47: same-coordinate first-connectivity cumulant]
\label{prog:V46-V47}
\begin{enumerate}[leftmargin=3em]
\item Enumerate the local partition blocks at the first-connectivity
      coordinate which both connect the two predecessor components
      and realize one or more internal tree edges.
\item For each block, choose the maximum number of leaf or small-mass
      legs permitted by its incidence tree and apply
      \Cref{thm:V46-general-two-leg}.
\item Compare the two available character-gap square roots with all
      local tree marks realized by the same block.
\item Treat arity four and five separately; do not infer four marked
      debits from the two-selected-leg theorem.
\item Use the archived all-order one-link MTA when every row support
      is concentrated at the joining coordinate.
\item If distributed row mass leaves an unmatched mark, return to
      probability-mark aggregation rather than adding an unproved
      local clock.
\end{enumerate}
\end{programme}

\begin{firewall}[Claim boundary after fourth-series V46]
\label{fire:V46-claim-boundary}
V46 closes payer coincidence and equivariant compression for the
ternary--binary \(3+2\) carrier.  It does not close a local block
which simultaneously realizes predecessor and joining edges, every
other predecessor partition, or the full quintic MTA; hence it does
not prove an all-order atlas, WUFC, CPM, cutoff removal, reflection
positivity, Osterwalder--Schrader reconstruction, construction of the
continuum Yang--Mills measure, or a uniform positive continuum
spectral gap.  The full Yang--Mills mass gap has not been proved.
\end{firewall}

\paragraph{Parent-version record.}
V46 was formed from
\texttt{Renewal\_Yang\_Mills\_Mass\_Gap\_Research\_Notes\_V45.tex}.
The V45 parent had \(22\,386\) logical source lines,
\(781\,322\) bytes, and SHA--256
\[
 \texttt{17188b235e52f6199045644679680c8f00f07f315c2365cf88260d014c9159b5}.
\]
The next compulsory companion audit is V50.

% =============================================================================
% END APPENDED FOURTH-SERIES V46 MATERIAL
% =============================================================================

% =============================================================================
% BEGIN APPENDED FOURTH-SERIES V47 MATERIAL
% =============================================================================

\chapter{V47: The saturated five-row same-coordinate cumulant is
overpaid}
\label{chap:V47-five-row-local}

\section{Purpose}
\label{sec:V47-purpose}

V46 reduced the \(3+2\) transition census to a first-connectivity
coordinate which simultaneously realizes predecessor and joining
relations.  The maximal coincidence is a single local cumulant block
containing all five rows.  Although such a block realizes four tree
edges at once, the two-selected-leg theorem is sufficient: saturation
puts three powers of the cofinal root mass into a selected tree
monomial, while the BFA input and payer ledger requires only one
Wilson clock from the local cumulant.

This note proves the complete output-summed statement, including
arbitrary spectator mass on all rows and a payer in the same
multiplicity space.

\section{The saturated one-coordinate block}
\label{sec:V47-setup}

Let \(r\) be a physical coordinate on which all five rows
\(a,b,c,d,e\) are nontrivial.  On \(a,c\), take equal spin \(j\) and
put
\[
 A:=1+c_j.
\]
Write the other three local masses as
\[
 y:=1+c_{R_{b,r}},
 \qquad
 z:=1+c_{R_{d,r}},
 \qquad
 v:=1+c_{R_{e,r}}.
\]
Allow arbitrary additional spectator moment banks on every row.  Let
\(\Xi_i\) be the full row masses and assume
\begin{equation}
\label{eq:V47-hypotheses}
 \Xi_i\le\Xi_a=:M,
 \qquad
 A\ge\sigma\max\{\Xi_a,\Xi_c\},
 \qquad
 0<x_-\le {d_j\over\sqrt\alpha}\le x_+<\infty.
\end{equation}
Thus
\begin{equation}
\label{eq:V47-scales}
 \Xi_a,\Xi_c,M,\alpha\asymp_{\sigma,x_\pm}A,
 \qquad
 s_\alpha\asymp A^{-1}.
\end{equation}

Let \({\mathcal K}_{5,r}\) be the complete fifth local Wilson
cumulant on the five rows at \(r\), assembled before any output norm.
The final \(Q_0\), reduced resolvent, and selected Casimir payer may
act in the same multiplicity space.

\section{A fully local selected tree}
\label{sec:V47-tree}

\begin{lemma}[Selected four-edge monomial at one coordinate]
\label{lem:V47-tree}
The complete marked-tree sum is bounded below by the contribution of
\[
 T_*=\{ab,ad,ce,ac\}
\]
with every edge marked by \(r\):
\begin{equation}
\label{eq:V47-tree}
 {\mathfrak t}_{5,r}
 :=
 {A^5yzv\over
  \Xi_a^3\Xi_c^2\Xi_b\Xi_d\Xi_e}.
\end{equation}
\end{lemma}

\begin{proof}
The four local pair-stock contributions are
\[
 {Ay\over\Xi_a\Xi_b},
 \qquad
 {Az\over\Xi_a\Xi_d},
 \qquad
 {Av\over\Xi_c\Xi_e},
 \qquad
 {A^2\over\Xi_a\Xi_c}.
\]
Their product is \eqref{eq:V47-tree}.  A link marking may assign the
same physical coordinate to several distinct tree edges, and the
marked BFA seminorm permits repeated marks.  Global pair stocks
dominate the displayed local contributions, while every other tree
monomial is nonnegative.
\end{proof}

\begin{remark}[Why the cofinal power is \(A^5\)]
\label{rem:V47-A5}
The strong edge \(ac\) contributes \(A^2\).  Each of the other three
selected edges has one endpoint on a cofinal row and contributes one
more \(A\).  This local numerator is the compensation unavailable in
the separated weak-coordinate skeleton before its covariance debits
were restored.
\end{remark}

\section{Complete weighted response}
\label{sec:V47-response}

\begin{lemma}[Two leaf legs pay the local cumulant]
\label{lem:V47-two-leaves}
The complete five-row local carrier satisfies
\begin{equation}
\label{eq:V47-two-leaves}
 \|{\mathcal K}_{5,r}\|
 \le C s_\alpha\sqrt{yz}.
\end{equation}
The estimate remains valid after any equivariant output compression.
\end{lemma}

\begin{proof}
Apply \Cref{cor:V46-general-Casimir} with \(m=5\), selecting the
\(b,d\) legs.  Replacing \(c_{R_{b,r}},c_{R_{d,r}}\) by the larger
quantities \(y,z\) gives \eqref{eq:V47-two-leaves}.
\end{proof}

\begin{theorem}[Five-row same-coordinate BFA payment]
\label{thm:V47-five-row-close}
Let \(\mathcal W_{\alpha}^{5,r}\) be the complete resolved word whose
first-connectivity carrier at \(r\) is
\({\mathcal K}_{5,r}\), including arbitrary spectator moments,
all physical outputs and multiplicities, exact Fourier normalization
and duality, \(Q_0\), and the coincident final
mass--resolvent payer.  Normalize each pure input row block to have
BFA norm one.  Then
\begin{equation}
\label{eq:V47-five-row-close}
 \|\mathcal W_{\alpha}^{5,r}\|_{\rho,2}^{\rm BFA}
 \le
 {C_{\rho,\sigma,x_\pm}\over A^2}
 {\mathfrak t}_{5,r}.
\end{equation}
\end{theorem}

\begin{proof}
The total input Casimir is at most \(5M=O(A)=O(\alpha)\).
Every output Casimir has the same upper scale by fusion
subadditivity.  Thus the product character gap gives the payer bound
\[
 {\Xi(\boldsymbol T)\over1-q_{\alpha,\boldsymbol T}}
 \le C\alpha
\]
on every nontrivial output.

Use \Cref{thm:V46-payer-two-leg} directly with \(m=5\), selected legs
\(b,d\), and the raw isotypic input blocks.  Spectator moments and
height quotient are bounded by one.  Exact output dimensions cancel,
and the raw pinching masses sum to at most one.  After the five input
normalizations,
\begin{equation}
\label{eq:V47-word-bound}
 \|\mathcal W_{\alpha}^{5,r}\|_{\rho,2}^{\rm BFA}
 \le
 {C\alpha s_\alpha\sqrt{yz}\over
  \Xi_a\Xi_b\Xi_c\Xi_d\Xi_e}.
\end{equation}
Divide by \eqref{eq:V47-tree}.  Every noncofinal row denominator
cancels:
\begin{equation}
\label{eq:V47-ratio}
 {\|\mathcal W_{\alpha}^{5,r}\|_{\rho,2}^{\rm BFA}
  \over{\mathfrak t}_{5,r}}
 \le
 C\alpha s_\alpha
 {\Xi_a^2\Xi_c\over A^5}
 {1\over v\sqrt{yz}}.
\end{equation}
All local masses are at least one.  By
\eqref{eq:V47-scales},
\[
 \alpha s_\alpha
 {\Xi_a^2\Xi_c\over A^5}
 \le {C\over A^2}.
\]
This proves \eqref{eq:V47-five-row-close}.
\end{proof}

\begin{corollary}[The maximal same-coordinate coincidence is not a
residual]
\label{cor:V47-maximal-closed}
A local fifth cumulant containing the two cofinal rows and all three
remaining rows satisfies the quintic marked-tree target uniformly on
the parabolic corridor, even when the payer, joining, predecessor
connections, and output recoupling all occur at that coordinate.
\end{corollary}

\begin{proof}
Every named operation is included in
\Cref{thm:V47-five-row-close}.  Its right side is smaller than the
selected tree monomial by \(O(A^{-2})\).
\end{proof}

\section{A general local-star count}
\label{sec:V47-star-count}

\begin{proposition}[Two cofinal centers make a high-incidence local
tree cheap]
\label{prop:V47-local-star}
Suppose a single local cumulant contains the two saturated cofinal
rows \(a,c\) and \(k\ge2\) additional rows.  Choose a tree on those
\(k+2\) vertices which contains \(ac\) and attaches every additional
row directly to either \(a\) or \(c\).  Its local pair-stock numerator
contains
\[
 A^{k+2}\prod_{\nu=1}^k h_\nu,
\]
where \(h_\nu\) are the additional local row masses.  The complete
cumulant has, for any two selected additional rows,
\[
 O\!\left(s_\alpha\sqrt{h_i h_j}\right)
\]
operator norm.  Thus increasing local incidence strengthens the
tree numerator without increasing the one-payer scale.
\end{proposition}

\begin{proof}
The edge \(ac\) contributes \(A^2\).  Each of the \(k\) attachments
contributes \(Ah_\nu\), proving the numerator formula.  The cumulant
estimate is \Cref{cor:V46-general-Casimir}.
\end{proof}

\begin{firewall}[The proposition is not yet the arity-four ledger]
\label{fire:V47-not-arity-four}
When \(k=2\), a fifth row lies outside the local cumulant and must be
connected elsewhere.  Its additional carrier, row denominator, and
tree edge must be retained together.  V48 performs that ledger; it is
not inferred merely from \Cref{prop:V47-local-star}.
\end{firewall}

\section{Residual after V47}
\label{sec:V47-residual}

\begin{theorem}[Proper same-coordinate blocks remain]
\label{thm:V47-residual}
Within the \(3+2\) saturated parabolic census, V47 removes every
first-connectivity term whose cross-component local block contains
all five rows.  A persistent same-coordinate residual must have a
proper cross-component block of arity two, three, or four, possibly
accompanied by other disjoint local blocks at the same coordinate.
\end{theorem}

\begin{proof}
The arity-five case is
\Cref{cor:V47-maximal-closed}.  A cross-component block must contain
at least two rows and at most five.  Excluding arity five leaves the
listed possibilities.  Disjoint blocks are permitted by the local
partition formula \eqref{eq:V42-connected-local-formula}.
\end{proof}

\begin{assessment}[Next case]
\label{ass:V47-next}
The next strongest coincidence is a four-row block containing both
cofinal rows and two other rows.  The fifth row must be connected by
one prefix carrier.  The two selected-leg debit of the four-row block
and the covariance debit of that prefix supply two Wilson clocks,
while the local tree numerator supplies four cofinal powers.  This is
the V48 ledger.
\end{assessment}

\section{V47 result ledger and V48 programme}
\label{sec:V47-ledger}

\begin{theorem}[Fourth-series V47 result ledger]
\label{thm:V47-results}
Relative to V1--V46, V47 proves:
\begin{enumerate}[label=\textup{(V47.\arabic*)},leftmargin=6.3em]
\item the exact selected four-edge tree monomial for a five-row
      saturated local block;
\item a two-leaf character-gap estimate for the complete fifth local
      cumulant;
\item insertion of that estimate into the coincident
      mass--resolvent payer through weighted raw pinching;
\item uniform BFA marked-tree payment with an \(A^{-2}\) surplus;
\item closure of the maximal same-coordinate first-connectivity
      coincidence; and
\item reduction of the remaining local census to proper
      cross-component blocks of arity two, three, or four.
\end{enumerate}
\end{theorem}

\begin{proof}
These are
\Cref{lem:V47-tree,
lem:V47-two-leaves,
thm:V47-five-row-close,
cor:V47-maximal-closed,
thm:V47-residual}.
\end{proof}

\begin{programme}[V48: four-row cross block and one exterior
connection]
\label{prog:V47-V48}
\begin{enumerate}[leftmargin=3em]
\item Enumerate the four-row cross blocks which contain both cofinal
      rows and two of \(b,d,e\).
\item Retain the exterior fifth-row connection as one complete binary
      covariance or a fixed higher-incidence carrier.
\item Select the two noncofinal legs inside the four-row cumulant and
      apply the payer-coincident theorem.
\item Compare the resulting two Wilson clocks with one selected
      four-edge tree monomial, retaining every distributed row
      denominator.
\item Treat blocks omitting one cofinal row separately; their
      high-pair separation cannot be silently bounded by the
      high-pair-together estimate.
\item Proceed to arity three only after the arity-four ledger closes.
\end{enumerate}
\end{programme}

\begin{firewall}[Claim boundary after fourth-series V47]
\label{fire:V47-claim-boundary}
V47 closes the complete five-row same-coordinate cumulant in the
saturated transition corridor.  It does not close every proper
cross-component local block, every predecessor partition, or the full
quintic MTA; hence it does not prove an all-order atlas, WUFC, CPM,
cutoff removal, reflection positivity, Osterwalder--Schrader
reconstruction, construction of the continuum Yang--Mills measure,
or a uniform positive continuum spectral gap.  The full Yang--Mills
mass gap has not been proved.
\end{firewall}

\paragraph{Parent-version record.}
V47 was formed from
\texttt{Renewal\_Yang\_Mills\_Mass\_Gap\_Research\_Notes\_V46.tex}.
The V46 parent had \(22\,869\) logical source lines,
\(797\,734\) bytes, and SHA--256
\[
 \texttt{1cce93309c4709f290299009b8e51ae647a4e26634a8c6cbed07d2d99988b2f3}.
\]
The next compulsory companion audit is V50.

% =============================================================================
% END APPENDED FOURTH-SERIES V47 MATERIAL
% =============================================================================

% =============================================================================
% BEGIN APPENDED FOURTH-SERIES V48 MATERIAL
% =============================================================================

\chapter{V48: Four-row high-pair joining blocks and the exterior
covariance}
\label{chap:V48-four-row}

\section{Purpose}
\label{sec:V48-purpose}

V47 closes a five-row local block at the first-connectivity
coordinate.  This note treats the next case: a four-row block contains
both saturated cofinal rows and two further rows, while the fifth row
is attached by one binary-separable predecessor covariance.  The
four-row cumulant supplies one two-selected-leg Wilson clock, the
exterior covariance supplies a second, and the local tree numerator
contains four powers of the root mass.  The resulting word has an
\(A^{-2}\) surplus over its selected tree monomial.

\section{Normal form of a high-pair four-row block}
\label{sec:V48-normal-form}

Let \(r\) be the first-connectivity coordinate.  Suppose one local
cumulant block at \(r\) contains the two cofinal rows \(a,c\) and two
rows \(i,k\) from \(\{b,d,e\}\).  Let \(l\) be the remaining row.
Choose a row \(h\in\{a,c,i,k\}\) and a distinct coordinate
\(w\) at which a complete binary covariance joins \(h\) to \(l\).
The choices are restricted only by the requirement that the selected
four edges form a labelled spanning tree.

At \(r\), put equal spin \(j\) on \(a,c\), set
\[
 A:=1+c_j,
\]
and denote the other two local masses by
\[
 y:=1+c_{R_{i,r}},
 \qquad
 z:=1+c_{R_{k,r}}.
\]
At \(w\), denote the local masses on \(h,l\) by
\[
 u:=1+c_{R_{h,w}},
 \qquad
 v:=1+c_{R_{l,w}}.
\]
Let \(\Xi_s\) be the complete row masses, including arbitrary
spectator moment banks, and assume
\begin{equation}
\label{eq:V48-hypotheses}
 \Xi_s\le\Xi_a=:M,
 \qquad
 A\ge\sigma\max\{\Xi_a,\Xi_c\},
 \qquad
 0<x_-\le {d_j\over\sqrt\alpha}\le x_+<\infty.
\end{equation}
Then
\begin{equation}
\label{eq:V48-scales}
 \Xi_a,\Xi_c,M,\alpha\asymp A,
 \qquad
 s_\alpha\asymp A^{-1}.
\end{equation}

\begin{definition}[Admissible local attachment tree]
\label{def:V48-attachment-tree}
An admissible tree \(T_{4+1}\) consists of:
\begin{enumerate}[label=\textup{(\roman*)},leftmargin=3em]
\item the strong edge \(ac\) at \(r\);
\item two edges at \(r\), each attaching one of \(i,k\) to either
      \(a\) or \(c\), with the four vertices
      \(\{a,c,i,k\}\) connected; and
\item the exterior edge \(hl\) at \(w\).
\end{enumerate}
\end{definition}

\begin{lemma}[Every four-row high-pair block admits the normal form]
\label{lem:V48-normal-form}
If the four-row block contains \(a,c,i,k\) and the fifth row is
connected to the predecessor fibre by a binary-separable carrier,
then, after relabelling \(i,k,l\) and possibly exchanging \(a,c\),
there is an admissible tree \(T_{4+1}\).
\end{lemma}

\begin{proof}
The complete graph on the four rows in one local cumulant block
contains the star formed by \(ac\) and by attaching each of \(i,k\)
to either cofinal row.  Choose the attachments so those four vertices
are connected.  By hypothesis the remaining row \(l\) has one
binary-separable connection to a row already in that component.
Adding it produces a five-vertex tree.
\end{proof}

\section{Selected tree and local debits}
\label{sec:V48-tree}

\begin{lemma}[Selected \(4+1\) tree monomial]
\label{lem:V48-tree}
Let \(d_s\) be the degree of row \(s\) in the selected admissible
tree \(T_{4+1}\).  The complete marked-tree sum is bounded below by
\begin{equation}
\label{eq:V48-tree}
 {\mathfrak t}_{4+1}
 :=
 {A^4yzuv\over
  \prod_{s\in\{a,b,c,d,e\}}\Xi_s^{d_s}}.
\end{equation}
\end{lemma}

\begin{proof}
At \(r\), the strong edge contributes
\[
 {A^2\over\Xi_a\Xi_c}.
\]
The two other local edges contribute
\[
 {Ay\over\Xi_{g_i}\Xi_i},
 \qquad
 {Az\over\Xi_{g_k}\Xi_k},
\]
where \(g_i,g_k\in\{a,c\}\).
The exterior edge contributes
\[
 {uv\over\Xi_h\Xi_l}.
\]
Multiplication gives four factors of \(A\), the four remaining local
masses \(y,z,u,v\), and exactly one row denominator for every
incident selected edge.  This is \eqref{eq:V48-tree}.  Global pair
stocks dominate the selected local contributions.
\end{proof}

\begin{lemma}[Two clocks for the \(4+1\) carrier]
\label{lem:V48-two-clocks}
The complete four-row local cumulant and exterior covariance obey
\begin{equation}
\label{eq:V48-two-clocks}
 \|{\mathcal K}_{4,r}\|
 \,\|{\mathcal C}_{h l,w}\|
 \le
 C s_\alpha^2\sqrt{yzuv}.
\end{equation}
\end{lemma}

\begin{proof}
Select the two noncofinal legs \(i,k\) in
\Cref{cor:V46-general-Casimir}:
\[
 \|{\mathcal K}_{4,r}\|
 \le Cs_\alpha\sqrt{yz}.
\]
Apply \Cref{cor:V43-Casimir-debit} to the exterior binary covariance:
\[
 \|{\mathcal C}_{h l,w}\|
 \le Cs_\alpha\sqrt{uv}.
\]
The coordinates are distinct and the two complete carriers
tensor-separate before final output pinching, so their norms multiply.
\end{proof}

\section{Complete output-summed payment}
\label{sec:V48-payment}

\begin{theorem}[Uniform \(4+1\) same-coordinate transition closure]
\label{thm:V48-four-plus-one-close}
Let \(\mathcal W_\alpha^{4+1}\) be the complete resolved BFA word with
the four-row high-pair block at \(r\), the exterior covariance at
\(w\), arbitrary spectator moment banks, exact physical output and
multiplicity sums, \(Q_0\), one final reduced resolvent, and a payer
in any equivariant output position.  Normalize each pure input row
block to BFA norm one.  Under \eqref{eq:V48-hypotheses},
\begin{equation}
\label{eq:V48-four-plus-one-close}
 \|\mathcal W_\alpha^{4+1}\|_{\rho,2}^{\rm BFA}
 \le
 {C_{\rho,\sigma,x_\pm}\over A^2}
 {\mathfrak t}_{4+1}.
\end{equation}
\end{theorem}

\begin{proof}
The total input and output Casimir are \(O(M)=O(\alpha)\).  The
product character gap and
\Cref{lem:V46-weighted-pinching} therefore bound the complete
mass--resolvent payer, including output multiplicities, by
\(C\alpha\) times raw pinching mass.

Insert \eqref{eq:V48-two-clocks}, the five input mass denominators,
and the height quotient, which is at most one.  The raw pinching
masses sum to at most one, giving
\begin{equation}
\label{eq:V48-word}
 \|\mathcal W_\alpha^{4+1}\|_{\rho,2}^{\rm BFA}
 \le
 {C\alpha s_\alpha^2\sqrt{yzuv}\over
  \Xi_a\Xi_b\Xi_c\Xi_d\Xi_e}.
\end{equation}
Divide by \eqref{eq:V48-tree}.  One copy of every row denominator
cancels, and
\begin{equation}
\label{eq:V48-ratio}
 {\|\mathcal W_\alpha^{4+1}\|_{\rho,2}^{\rm BFA}
  \over{\mathfrak t}_{4+1}}
 \le
 C\alpha s_\alpha^2
 {\prod_s\Xi_s^{d_s-1}\over A^4}
 {1\over\sqrt{yzuv}}.
\end{equation}
Every local mass is at least one.  By
\eqref{eq:V48-scales}, every row mass is at most \(M\asymp A\).
Since a five-vertex tree has
\[
 \sum_s(d_s-1)=2|E(T_{4+1})|-5=3,
\]
the remaining scale is at most
\[
 C\alpha s_\alpha^2{M^3\over A^4}
 \le {C\over A^2}.
\]
\end{proof}

\begin{corollary}[All high-pair arity-four coincidences with a
binary exterior close]
\label{cor:V48-all-high-pair-four}
Every same-coordinate four-row cross block which contains both
cofinal rows and whose omitted row is attached by a distinct
binary-separable covariance satisfies its marked-tree payment,
uniformly in all representation and multiplicity labels.
\end{corollary}

\begin{proof}
Use \Cref{lem:V48-normal-form} and
\Cref{thm:V48-four-plus-one-close}.
\end{proof}

\section{Why high-pair separation is a different case}
\label{sec:V48-separated}

\begin{firewall}[A four-row cross block may omit one cofinal row]
\label{fire:V48-separated-high-pair}
At the first-connectivity coordinate, a four-row block can cross the
\(3+2\) predecessor cut while containing only one of the cofinal rows.
The other cofinal row may lie in a disjoint local block or singleton
of the same local partition.  In that case
\Cref{cor:V46-general-Casimir} cannot select two noncofinal legs while
simultaneously treating the high pair as one saturated root.  The
tree numerator and the Wilson multiplier have a different ledger.
No result of this chapter is applied to that separated-high-pair
case.
\end{firewall}

\begin{theorem}[Updated same-coordinate residual]
\label{thm:V48-residual}
After V47--V48, a persistent \(3+2\) same-coordinate residual must
belong to at least one of:
\begin{enumerate}[label=\textup{(\roman*)},leftmargin=3em]
\item a cross-component local block separates the two cofinal rows;
\item the cross block contains the high pair but has arity three,
      with two predecessor connections outside it;
\item the cross block is binary and all predecessor connections occur
      elsewhere; or
\item the omitted row is attached through a non-binary-separable
      carrier at the same coordinate.
\end{enumerate}
\end{theorem}

\begin{proof}
V47 removes arity five.  V48 removes every arity-four block containing
the high pair with a binary-separable exterior connection.  The
remaining block arities and the explicit failure of the two
hypotheses give exactly the four alternatives listed.
\end{proof}

\begin{assessment}[V49 target]
\label{ass:V48-V49}
The next controlled case is a ternary cross block containing the high
pair and one noncofinal row.  It realizes the strong joining edge and
one additional tree edge.  Two binary predecessor covariances remain
outside it.  One square-root debit from the extra ternary leg, paired
with a cofinal leg, and the two full covariance debits are expected to
overpay the selected tree.  That fractional clock must be counted
exactly.
\end{assessment}

\section{V48 result ledger and V49 programme}
\label{sec:V48-ledger}

\begin{theorem}[Fourth-series V48 result ledger]
\label{thm:V48-results}
Relative to V1--V47, V48 proves:
\begin{enumerate}[label=\textup{(V48.\arabic*)},leftmargin=6.3em]
\item the normal form of a four-row same-coordinate block containing
      both cofinal rows;
\item its selected \(A^4\) four-edge tree monomial;
\item two exact Wilson clocks from the four-row selected-leg cumulant
      and exterior covariance;
\item complete output-summed, payer-coincident BFA payment with an
      \(A^{-2}\) surplus;
\item closure of every high-pair arity-four coincidence with a
      binary-separable exterior; and
\item isolation of separated-high-pair and lower-arity
      same-coordinate residuals.
\end{enumerate}
\end{theorem}

\begin{proof}
These are
\Cref{lem:V48-normal-form,
lem:V48-tree,
lem:V48-two-clocks,
thm:V48-four-plus-one-close,
cor:V48-all-high-pair-four,
thm:V48-residual}.
\end{proof}

\begin{programme}[V49: high-pair ternary cross block]
\label{prog:V48-V49}
\begin{enumerate}[leftmargin=3em]
\item Put \(a,c\) and one noncofinal row in the local ternary joining
      block.
\item Select the noncofinal leg and one cofinal leg in the V45
      character-gap theorem; retain the resulting half-clock.
\item Keep the other two predecessor connections as complete binary
      covariances with arbitrary distributed endpoint mass.
\item Compute the exact four-edge tree monomial and compare all
      fractional powers without replacing local masses by one too
      early.
\item Include a coincident payer using the V46 weighted-pinching
      theorem.
\item Leave separated-high-pair blocks for the compulsory V50 audit
      unless the V49 ledger closes them automatically.
\end{enumerate}
\end{programme}

\begin{firewall}[Claim boundary after fourth-series V48]
\label{fire:V48-claim-boundary}
V48 closes high-pair arity-four same-coordinate blocks with a
binary-separable exterior connection.  It does not close ternary,
binary, separated-high-pair, or nonseparable exterior cases, every
predecessor partition, or the full quintic MTA; hence it does not
prove an all-order atlas, WUFC, CPM, cutoff removal, reflection
positivity, Osterwalder--Schrader reconstruction, construction of the
continuum Yang--Mills measure, or a uniform positive continuum
spectral gap.  The full Yang--Mills mass gap has not been proved.
\end{firewall}

\paragraph{Parent-version record.}
V48 was formed from
\texttt{Renewal\_Yang\_Mills\_Mass\_Gap\_Research\_Notes\_V47.tex}.
The V47 parent had \(23\,218\) logical source lines,
\(809\,613\) bytes, and SHA--256
\[
 \texttt{2dcf2fe375b57370774fd8b8ab6408fbadf8cb87b722479f85994dc243e006d9}.
\]
The next compulsory companion audit is V50.

% =============================================================================
% END APPENDED FOURTH-SERIES V48 MATERIAL
% =============================================================================

% =============================================================================
% BEGIN APPENDED FOURTH-SERIES V49 MATERIAL
% =============================================================================

\chapter{V49: A high-pair ternary joining block gains a half-clock}
\label{chap:V49-ternary-cross}

\section{Purpose}
\label{sec:V49-purpose}

V47--V48 close same-coordinate joining blocks of arity five and four
when they contain both cofinal rows.  This note treats arity three.
The local block contains the high pair and one additional row.  Its
two-selected-leg estimate uses that additional row and either
cofinal row, producing
\[
 s_\alpha\sqrt{A h}
 \asymp\sqrt{s_\alpha h}.
\]
Two exterior binary covariances supply two more Wilson clocks.  The
exact tree-degree ledger then leaves an \(A^{-3/2}\) surplus.

\section{Ternary cross block and two exterior connections}
\label{sec:V49-setup}

Let \(r\) be the first-connectivity coordinate.  Suppose a complete
local ternary cumulant block contains the two cofinal rows \(a,c\)
and one noncofinal row \(i\in\{b,d,e\}\).  Let \(k,l\) be the two
remaining rows.  At \(r\), put equal spin \(j\) on \(a,c\), let
\[
 A:=1+c_j,
 \qquad
 h:=1+c_{R_{i,r}},
\]
and choose at \(r\) the tree edges \(ac\) and \(gi\), where
\(g\in\{a,c\}\).

Let \(w_1,w_2\ne r\) be distinct coordinates carrying complete
binary-separable covariances with local endpoint masses
\[
 x_1,y_1
 \qquad\hbox{and}\qquad
 x_2,y_2,
\]
such that the two exterior edges attach \(k,l\) and the three
vertices already present at \(r\) into a five-row spanning tree
\(T_{3+1+1}\).  The attachment vertices may be any rows already
connected at that stage.

Allow arbitrary spectator moment banks.  With full row masses
\(\Xi_s\), assume
\begin{equation}
\label{eq:V49-hypotheses}
 \Xi_s\le\Xi_a=:M,
 \qquad
 A\ge\sigma\max\{\Xi_a,\Xi_c\},
 \qquad
 0<x_-\le {d_j\over\sqrt\alpha}\le x_+<\infty.
\end{equation}
Then
\begin{equation}
\label{eq:V49-scales}
 \Xi_a,\Xi_c,M,\alpha\asymp A,
 \qquad
 s_\alpha\asymp A^{-1}.
\end{equation}

\section{Selected tree monomial}
\label{sec:V49-tree}

\begin{lemma}[Exact degree-form tree lower bound]
\label{lem:V49-tree}
Let \(d_s\) be the degree of row \(s\) in
\(T_{3+1+1}\).  The complete marked-tree sum is bounded below by
\begin{equation}
\label{eq:V49-tree}
 {\mathfrak t}_{3+1+1}
 :=
 {A^3h\,x_1y_1x_2y_2
  \over
  \prod_{s\in\{a,b,c,d,e\}}\Xi_s^{d_s}}.
\end{equation}
\end{lemma}

\begin{proof}
The local strong edge \(ac\) contributes
\[
 {A^2\over\Xi_a\Xi_c},
\]
and the attachment \(gi\) contributes
\[
 {Ah\over\Xi_g\Xi_i}.
\]
The two exterior edges contribute their endpoint products
\[
 {x_1y_1\over\Xi_{p_1}\Xi_{q_1}},
 \qquad
 {x_2y_2\over\Xi_{p_2}\Xi_{q_2}}.
\]
Their product has exactly the denominator
\(\prod_s\Xi_s^{d_s}\) and the numerator in
\eqref{eq:V49-tree}.  Global pair stocks dominate these selected
coordinate contributions.
\end{proof}

\section{The fractional-clock carrier}
\label{sec:V49-half-clock}

\begin{lemma}[Ternary high-pair half-clock]
\label{lem:V49-half-clock}
The complete local ternary cumulant at \(r\) obeys
\begin{equation}
\label{eq:V49-half-clock}
 \|{\mathcal K}_{3,r}\|
 \le
 C_{\sigma,x_\pm}\sqrt{s_\alpha h}.
\end{equation}
\end{lemma}

\begin{proof}
Select row \(i\) and the cofinal row \(g\) in
\Cref{cor:V46-general-Casimir}.  This gives
\[
 \|{\mathcal K}_{3,r}\|
 \le
 Cs_\alpha\sqrt{Ah}
 =
 C\sqrt{s_\alpha A}\sqrt{s_\alpha h}.
\]
The corridor gives \(s_\alpha A\asymp1\), proving
\eqref{eq:V49-half-clock}.
\end{proof}

\begin{lemma}[Complete local product]
\label{lem:V49-local-product}
The ternary block and the two exterior binary covariances satisfy
\begin{equation}
\label{eq:V49-local-product}
 \|{\mathcal K}_{3,r}\|
 \prod_{\nu=1}^2\|{\mathcal C}_{\nu,w_\nu}\|
 \le
 C
 s_\alpha^{5/2}
 \sqrt{h\,x_1y_1x_2y_2}.
\end{equation}
\end{lemma}

\begin{proof}
Use \Cref{lem:V49-half-clock} for the ternary factor and
\Cref{cor:V43-Casimir-debit} for each exterior covariance:
\[
 \|{\mathcal C}_{\nu,w_\nu}\|
 \le Cs_\alpha\sqrt{x_\nu y_\nu}.
\]
The three distinct-coordinate carriers tensor-separate before the
final output pinching.
\end{proof}

\section{Complete BFA payment}
\label{sec:V49-payment}

\begin{theorem}[Uniform high-pair ternary cross-block closure]
\label{thm:V49-ternary-cross-close}
Let \(\mathcal W_\alpha^{3+1+1}\) be the complete resolved word with
the ternary high-pair block at \(r\), the two exterior
binary-separable covariances, arbitrary spectator moment banks,
exact physical output and multiplicity sums, \(Q_0\), one final
reduced resolvent, and a payer in any equivariant output position.
Normalize each pure input row block to BFA norm one.  Under
\eqref{eq:V49-hypotheses},
\begin{equation}
\label{eq:V49-close}
 \|\mathcal W_\alpha^{3+1+1}\|_{\rho,2}^{\rm BFA}
 \le
 {C_{\rho,\sigma,x_\pm}\over A^{3/2}}
 {\mathfrak t}_{3+1+1}.
\end{equation}
\end{theorem}

\begin{proof}
All row masses are at most \(M\asymp A\), so the total output
Casimir is \(O(A)=O(\alpha)\).  The product character gap and the
weighted pinching theorem give a \(C\alpha\) bound for the complete
final mass--resolvent payer, with no output multiplicity count.

Insert \eqref{eq:V49-local-product}, divide by the five normalized
input masses, and bound the height quotient by one:
\begin{equation}
\label{eq:V49-word}
 \|\mathcal W_\alpha^{3+1+1}\|_{\rho,2}^{\rm BFA}
 \le
 {C\alpha s_\alpha^{5/2}
  \sqrt{h\,x_1y_1x_2y_2}
  \over
  \Xi_a\Xi_b\Xi_c\Xi_d\Xi_e}.
\end{equation}
Divide by \eqref{eq:V49-tree}.  One copy of each row mass cancels:
\begin{equation}
\label{eq:V49-ratio}
 {\|\mathcal W_\alpha^{3+1+1}\|_{\rho,2}^{\rm BFA}
  \over{\mathfrak t}_{3+1+1}}
 \le
 C\alpha s_\alpha^{5/2}
 {\prod_s\Xi_s^{d_s-1}\over A^3}
 {1\over
  \sqrt{h\,x_1y_1x_2y_2}}.
\end{equation}
Every local mass in the last denominator is at least one.  Since a
five-vertex tree has
\[
 \sum_s(d_s-1)=3
\]
and every \(\Xi_s\le M\asymp A\), the remaining scale is bounded by
\[
 C\alpha s_\alpha^{5/2}{M^3\over A^3}
 \le {C\over A^{3/2}}.
\]
This proves \eqref{eq:V49-close}.
\end{proof}

\begin{corollary}[All high-pair ternary coincidences with two binary
exteriors close]
\label{cor:V49-all-ternary}
Every same-coordinate ternary cross block which contains both
cofinal rows and whose other two rows are attached through two
distinct binary-separable covariances satisfies its marked-tree
payment uniformly in all representation and multiplicity labels.
\end{corollary}

\begin{proof}
Relabel the noncofinal rows and the attachment tree, then apply
\Cref{thm:V49-ternary-cross-close}.
\end{proof}

\section{Disposition of high-pair-together cross blocks}
\label{sec:V49-disposition}

\begin{theorem}[High-pair-together same-coordinate disposition]
\label{thm:V49-high-pair-disposition}
Suppose a saturated \(3+2\) first-connectivity term has a
cross-component local cumulant block at the joining coordinate which
contains both cofinal rows \(a,c\).  If every connection not realized
inside that block is binary-separable at a distinct coordinate, then
the resolved word satisfies its marked-tree payment for every
possible block arity \(2,3,4,5\).
\end{theorem}

\begin{proof}
If the cross block has arity five, use
\Cref{thm:V47-five-row-close}.  At arity four, use
\Cref{thm:V48-four-plus-one-close}.  At arity three, use
\Cref{thm:V49-ternary-cross-close}.  At arity two, the joining block
is the binary high-pair carrier and all three predecessor connections
are outside it.  They are paid by
\Cref{prop:V44-three-separable}, or by
\Cref{thm:V45-ternary-binary-close} if two of the internal
connections are assembled into a distinct ternary prefix.  Payer
coincidence is harmless by
\Cref{thm:V46-payer-two-leg}.
\end{proof}

\begin{firewall}[Non-binary exterior coincidences]
\label{fire:V49-exterior-scope}
The disposition theorem requires each connection outside the joining
block either to be binary-separable or to belong to the distinct
ternary prefix already handled in V45.  A second higher-incidence
block sharing rows or multiplicity data with the joining block is not
automatically a tensor product and remains outside the theorem.
\end{firewall}

\section{The separated-high-pair frontier}
\label{sec:V49-separated-frontier}

\begin{assessment}[Exact residual before the V50 audit]
\label{ass:V49-residual}
The remaining same-coordinate obstruction is no longer a block
containing the high pair.  It is a local partition in which the two
cofinal rows lie in different blocks, while some other block crosses
the predecessor cut and creates first connectivity.  The strong
\(ac\) pair is then present in the marked-tree geometry but absent
from the same local cumulant block.  Its two endpoint moments and the
crossing block must be retained as one signed local-partition carrier.
V36's fixed-time separated-sector Gaussian is not uniform on the
parabolic clock, so it cannot close this case.
\end{assessment}

\section{V49 result ledger and V50 audit programme}
\label{sec:V49-ledger}

\begin{theorem}[Fourth-series V49 result ledger]
\label{thm:V49-results}
Relative to V1--V48, V49 proves:
\begin{enumerate}[label=\textup{(V49.\arabic*)},leftmargin=6.3em]
\item the exact degree-form selected tree for a ternary high-pair
      joining block and two exterior connections;
\item a sharp ledger-level half-clock from one cofinal and one
      noncofinal selected leg;
\item multiplication with two binary covariance clocks;
\item complete output-summed, payer-coincident BFA payment with an
      \(A^{-3/2}\) surplus;
\item closure of every high-pair-together joining block of arity
      two through five under the stated exterior separability; and
\item isolation of the parabolic separated-high-pair local partition
      as the next genuine obstruction.
\end{enumerate}
\end{theorem}

\begin{proof}
These are
\Cref{lem:V49-tree,
lem:V49-half-clock,
lem:V49-local-product,
thm:V49-ternary-cross-close,
thm:V49-high-pair-disposition,
ass:V49-residual}.
\end{proof}

\begin{programme}[V50: companion audit and separated-high-pair
carrier]
\label{prog:V49-V50}
\begin{enumerate}[leftmargin=3em]
\item Perform the compulsory read-only audit of the newest active SM
      and NS notes.
\item Enumerate the local partitions whose crossing block connects
      the \(3+2\) predecessor components while separating \(a,c\).
\item Assemble the sum of all such partitions before taking an
      operator norm; do not estimate the high endpoint moments
      separately.
\item Use exact \(SU(2)\) Wilson transition multipliers to determine
      whether the separated family has a signed cancellation at
      order one or only a square-root debit.
\item Compare any surviving coefficient with a tree containing the
      strong geometric edge \(ac\), noting that this edge need not
      belong to the crossing cumulant block.
\item Go upstream to the outer local-partition identity if the
      separated subfamily is not closed under its own sign.
\end{enumerate}
\end{programme}

\begin{firewall}[Claim boundary after fourth-series V49]
\label{fire:V49-claim-boundary}
V49 closes high-pair-together same-coordinate joining blocks with the
stated separable exterior carriers.  It does not close local
partitions separating the high pair, multiple entangled
higher-incidence blocks, every predecessor partition, or the full
quintic MTA; hence it does not prove an all-order atlas, WUFC, CPM,
cutoff removal, reflection positivity, Osterwalder--Schrader
reconstruction, construction of the continuum Yang--Mills measure,
or a uniform positive continuum spectral gap.  The full Yang--Mills
mass gap has not been proved.
\end{firewall}

\paragraph{Parent-version record.}
V49 was formed from
\texttt{Renewal\_Yang\_Mills\_Mass\_Gap\_Research\_Notes\_V48.tex}.
The V48 parent had \(23\,578\) logical source lines,
\(822\,002\) bytes, and SHA--256
\[
 \texttt{def53c96a29fdfa9e74058bd339ccd9dbb3d3efc9e139bc85c0753532ecc8509}.
\]
The compulsory companion audit is due in V50.

% =============================================================================
% END APPENDED FOURTH-SERIES V49 MATERIAL
% =============================================================================

% =============================================================================
% BEGIN APPENDED FOURTH-SERIES V50 MATERIAL
% =============================================================================

\chapter{V50: Companion audit and exact assembly of the relative
joining covariance}
\label{chap:V50-relative-joining}

\section{Compulsory read-only companion audit}
\label{sec:V50-audit}

The newest active companion notes at the V50 boundary are
\begin{align*}
 &\texttt{Renewal\_SM\_Einstein\_Continuum\_Research\_Notes\_V42.tex},\\
 &\texttt{Renewal\_NS\_Stability\_Research\_Notes\_V31.tex}.
\end{align*}
They were inspected read-only and were not modified.  Their exact
identities are
\begin{center}
\begin{tabular}{c|r|r|l}
programme & logical lines & bytes & SHA--256 prefix\\ \hline
SM V42 & \(16\,029\) & \(550\,018\) &
\texttt{d6e56c0bc7188096}\\
NS V31 & \(23\,052\) & \(887\,537\) &
\texttt{f1121b62238ad549}
\end{tabular}
\end{center}
The complete digests are
\[
\begin{split}
&\texttt{d6e56c0bc7188096e7da2bf17e793268d60cfe8180e74cb03db4aace4659197d},\\
&\texttt{f1121b62238ad5491bb7cf03635ea9934942edf573ed8faaa9ee79e1d38a6696}.
\end{split}
\]

SM V42 proves that a zero of a finite projected system is not even an
approximate zero of the full tower until the omitted residual,
infinite-tail inverse, and nonlinear remainder are controlled.  The
direct analogue here is that a selected high-pair-together or
high-pair-separated partition subfamily is not the complete local
joining carrier.

NS V31 remains unchanged.  Its sharp heat and packet tests again show
that a signed nonlinear source must be assembled before a smoothing
or absolute-value estimate is applied.  These two audits require an
upstream correction: sum the entire family of local partitions
indecomposable relative to the predecessor partition before applying
the transition bounds.

\begin{firewall}[Companion status]
\label{fire:V50-companion-status}
The companion notes determine the audit direction only.  The relative
partition identity and covariance estimate below are proved
self-containedly.
\end{firewall}

\section{Cumulants of block-products}
\label{sec:V50-product-cumulant}

Let \(I=\{1,\ldots,m\}\), let
\[
 \rho=\{C_1,\ldots,C_b\}\in\Pi(I),
\]
and let \(A_i\) be random matrix coefficients of one common group
variable.  Define the block-products
\begin{equation}
\label{eq:V50-block-products}
 Y_\nu:=\bigotimes_{i\in C_\nu}A_i
 \qquad(1\le\nu\le b).
\end{equation}
For \(\pi\in\Pi(I)\), put
\begin{equation}
\label{eq:V50-local-cumulant-product}
 {\mathcal K}_\pi
 :=
 \bigotimes_{B\in\pi}
 \kappa_{|B|}\bigl((A_i)_{i\in B}\bigr).
\end{equation}

\begin{theorem}[Relative connected-partition identity]
\label{thm:V50-relative-cumulant}
The cumulant of the block-products is
\begin{equation}
\label{eq:V50-relative-cumulant}
 \boxed{
 \kappa_b(Y_1,\ldots,Y_b)
 =
 \sum_{\substack{\pi\in\Pi(I)\\
                  \pi\vee\rho=\widehat1_I}}
 {\mathcal K}_\pi.}
\end{equation}
Every local partition indecomposable relative to \(\rho\) occurs once,
and every relatively disconnected partition cancels.
\end{theorem}

\begin{proof}
Expand the outer cumulant over partitions
\(\sigma\in\Pi(\{1,\ldots,b\})\):
\[
 \kappa_b(Y_1,\ldots,Y_b)
 =
 \sum_{\sigma}
 \mu(\sigma,\widehat1_b)
 \prod_{D\in\sigma}
 \mathbb E\!\left[\bigotimes_{\nu\in D}Y_\nu\right].
\]
Expand each displayed moment into cumulants of the original
\(A_i\)'s.  Fix \(\pi\in\Pi(I)\).  It occurs in the term \(\sigma\)
exactly when every block of \(\pi\) is contained in the union of the
\(\rho\)-blocks indexed by one block of \(\sigma\).

Let \(\overline\pi\) be the partition of
\(\{1,\ldots,b\}\) generated by declaring two \(\rho\)-blocks
equivalent when some block of \(\pi\) meets both, followed by
transitive closure.  The compatibility condition is
\(\overline\pi\le\sigma\).  Hence the total coefficient of
\({\mathcal K}_\pi\) is
\[
 \sum_{\sigma:\,\overline\pi\le\sigma\le\widehat1_b}
 \mu(\sigma,\widehat1_b).
\]
M\"obius inversion makes it one if
\(\overline\pi=\widehat1_b\) and zero otherwise.  The first condition
is equivalent to
\(\pi\vee\rho=\widehat1_I\), proving
\eqref{eq:V50-relative-cumulant}.
\end{proof}

\begin{corollary}[The two-block joining carrier is one covariance]
\label{cor:V50-two-block-covariance}
If \(\rho=\{C_-,C_+\}\), then
\begin{equation}
\label{eq:V50-two-block-covariance}
 \sum_{\pi:\,\pi\vee\rho=\widehat1_I}
 {\mathcal K}_\pi
 =
 \mathbb E[Y_-\otimes Y_+]
 -\mathbb E[Y_-]\otimes\mathbb E[Y_+].
\end{equation}
\end{corollary}

\begin{proof}
The right side is
\(\kappa_2(Y_-,Y_+)\).  Apply
\Cref{thm:V50-relative-cumulant}.
\end{proof}

\section{Dimension-free covariance of reducible block-products}
\label{sec:V50-block-covariance}

Each \(Y_\pm(U)\) in \eqref{eq:V50-two-block-covariance} is unitary,
although it need not be irreducible.  Put
\[
 Q_\pm:=\mathbb E_\alpha Y_\pm.
\]
Central inversion invariance makes \(Q_\pm\) self-adjoint
contractions which commute with the diagonal group action.

\begin{lemma}[Reducible block covariance bound]
\label{lem:V50-reducible-covariance}
The complete covariance
\[
 {\mathcal J}_\rho
 :=
 \mathbb E_\alpha[Y_-\otimes Y_+]-Q_-\otimes Q_+
\]
obeys
\begin{equation}
\label{eq:V50-reducible-covariance}
 \|{\mathcal J}_\rho\|
 \le
 \|I-Q_-^2\|^{1/2}\|I-Q_+^2\|^{1/2}
 \le1.
\end{equation}
It is equivariant, and every isotypic compression obeys the same
bound.
\end{lemma}

\begin{proof}
Let \(X_\pm=Y_\pm-Q_\pm\).  Then
\[
 {\mathcal J}_\rho=\mathbb E[X_-\otimes X_+].
\]
Unitarity and self-adjointness of the means give
\[
 \mathbb E[X_\pm X_\pm^*]=I-Q_\pm^2.
\]
Apply \Cref{lem:V43-operator-CS}.  Since \(Q_\pm\) are self-adjoint
contractions,
\[
 0\preccurlyeq I-Q_\pm^2\preccurlyeq I.
\]
Centrality of the law proves equivariance exactly as in
\Cref{lem:V46-equivariant}; compression does not increase the norm.
\end{proof}

\begin{corollary}[One-clock weighted joining bound]
\label{cor:V50-weighted-joining}
On a saturated parabolic corridor where the total output Casimir is
\(O(\alpha)\), the complete relative joining covariance, final BFA
mass, reduced resolvent, coincident payer, and multiplicity sum are
bounded by
\begin{equation}
\label{eq:V50-weighted-joining}
 C\alpha
 \sum_{\boldsymbol T,\gamma}
 \|X_{\boldsymbol T,\gamma}\|_1.
\end{equation}
\end{corollary}

\begin{proof}
The product character gap gives
\[
 {\Xi(\boldsymbol T)\over1-q_{\alpha,\boldsymbol T}}
 \le C\alpha
\]
on every nontrivial output.  Apply
\Cref{lem:V46-weighted-pinching} to the equivariant operator
\({\mathcal J}_\rho\), using
\eqref{eq:V50-reducible-covariance}.  The trivial output is removed
by \(Q_0\).
\end{proof}

\section{Application to the \(3+2\) first-connectivity fibre}
\label{sec:V50-32}

Let
\[
 \rho=\{\{a,b,d\},\{c,e\}\}
\]
be the partition immediately below the first-connectivity coordinate
\(r\).  At \(r\), define
\[
 Y_-:=A_{a,r}\otimes A_{b,r}\otimes A_{d,r},
 \qquad
 Y_+:=A_{c,r}\otimes A_{e,r},
\]
with identity factors for rows inactive at \(r\).

\begin{theorem}[Exact first-connectivity joining assembly]
\label{thm:V50-exact-joining}
The complete local joining sum at \(r\), including every
moment--cumulant partition which connects the two predecessor
components, is exactly
\begin{equation}
\label{eq:V50-exact-joining}
 {\mathcal J}_{3+2,r}
 =
 \operatorname{Cov}_\alpha(Y_-,Y_+).
\end{equation}
In particular, the subfamilies in which the cofinal rows \(a,c\) lie
in the same local cumulant block or in different blocks are not
separate source terms after exact assembly.
\end{theorem}

\begin{proof}
The coefficient-one first-connectivity formula retains precisely the
local partitions \(\pi\) for which
\[
 \pi\vee\rho=\widehat1_{\{a,b,c,d,e\}}.
\]
Their local carrier is the sum of
\({\mathcal K}_\pi\).  Apply
\Cref{cor:V50-two-block-covariance}.
Whether \(a,c\) share a block is merely a subdivision of that sum and
has no invariant status in \eqref{eq:V50-exact-joining}.
\end{proof}

\begin{firewall}[Why the separated subfamily must not be normed]
\label{fire:V50-no-separated-norm}
The M\"obius cancellation in
\eqref{eq:V50-exact-joining} occurs between high-pair-together and
high-pair-separated local partitions.  Bounding either subfamily in
absolute value before their sum replaces the covariance by a
collection of raw moments and can lose its centered structure.
No V50 proof uses such a termwise norm.
\end{firewall}

\begin{theorem}[Separated-high-pair partitions create no new
transition residual]
\label{thm:V50-separated-closed}
Consider a saturated parabolic \(3+2\) first-connectivity fibre whose
three-edge predecessor forest below \(r\) is of one of the following
types:
\begin{enumerate}[label=\textup{(\alph*)},leftmargin=3em]
\item three binary-separable carriers, as in V44; or
\item one ternary carrier on the three-row component and one binary
      carrier on the two-row component, as in V45,
\end{enumerate}
including the payer-coincident and equivariantly recoupled variants
of V46.  Then the complete joining sum
\eqref{eq:V50-exact-joining}, including all partitions separating
the cofinal pair, satisfies the required marked-tree payment.
\end{theorem}

\begin{proof}
The proofs of
\Cref{thm:V44-distributed-close,
thm:V45-ternary-binary-close}
use from the final joining carrier only:
\begin{enumerate}[label=\textup{(\roman*)},leftmargin=3em]
\item equivariance;
\item a representation-uniform operator bound;
\item a \(C\alpha\) bound after the one final BFA
      mass--resolvent payer; and
\item raw pinching subprobability for the output sum.
\end{enumerate}
The assembled covariance has \textup{(i)--(ii)} by
\Cref{lem:V50-reducible-covariance} and
\textup{(iii)--(iv)} by
\Cref{cor:V50-weighted-joining}.  Repeat the unchanged prefix and
tree ledgers of V44 or V45.  The V46 weighted-pinching theorem covers
payer coincidence and Racah recoupling.  Since
\eqref{eq:V50-exact-joining} already contains both high-pair
subfamilies, the conclusion includes the separated partitions.
\end{proof}

\begin{corollary}[The V49 separated-high-pair frontier is
superseded]
\label{cor:V50-V49-superseded}
The separated-high-pair family isolated in
\Cref{ass:V49-residual} is not an independent local bottleneck in
the prefix geometries covered by
\Cref{thm:V50-separated-closed}.
\end{corollary}

\begin{proof}
It is a proper subfamily of the exact covariance
\eqref{eq:V50-exact-joining}; the complete covariance is paid by
\Cref{thm:V50-separated-closed}.
\end{proof}

\section{What V47--V49 retain}
\label{sec:V50-prior-status}

\begin{assessment}[Termwise results are auxiliary checks]
\label{ass:V50-termwise-auxiliary}
The V47--V49 estimates remain valid for the explicitly assembled
local cumulant blocks in their stated scope.  They are useful
regression tests on maximal, arity-four, and arity-three channels.
They are no longer the preferred proof route for the whole joining
coordinate, because \Cref{thm:V50-exact-joining} controls its complete
relative covariance without splitting the partition family.
\end{assessment}

\section{The actual residual after V50}
\label{sec:V50-residual}

\begin{theorem}[Residual moves entirely into the predecessor forest]
\label{thm:V50-prefix-residual}
For a saturated parabolic \(3+2\) fibre, the local joining coordinate
creates no additional high-pair-together/separated obstruction once
its relative partition family is assembled.  Any residual not paid
by \Cref{thm:V50-separated-closed} must arise because the connected
carrier below \(r\) for at least one predecessor component is neither:
\begin{enumerate}[label=\textup{(\roman*)},leftmargin=3em]
\item a binary-separable tree; nor
\item a distinct ternary carrier of the type paid in V45--V46.
\end{enumerate}
\end{theorem}

\begin{proof}
The complete joining carrier, payer, and output sum are paid by
\Cref{cor:V50-weighted-joining}.  When the prefix has type
\textup{(i)} or \textup{(ii)}, the whole word is paid by
\Cref{thm:V50-separated-closed}.  Therefore a persistent unclosed
word must fail both prefix alternatives.
\end{proof}

\begin{assessment}[Audit-driven direction change]
\label{ass:V50-direction}
The next task is not another Wilson asymptotic at the joining
coordinate.  It is an exact relative-cumulant decomposition of the
three-row predecessor component across its whole prefix coordinate
bank, retaining probability marks and one output compression.  This
is the full residual demanded by the SM audit and the signed-source
assembly demanded by the NS audit.
\end{assessment}

\section{V50 result ledger and V51 programme}
\label{sec:V50-ledger}

\begin{theorem}[Fourth-series V50 result ledger]
\label{thm:V50-results}
Relative to V1--V49, V50 proves:
\begin{enumerate}[label=\textup{(V50.\arabic*)},leftmargin=6.3em]
\item the compulsory read-only audit of active SM V42 and NS V31;
\item the exact relative connected-partition identity
      \eqref{eq:V50-relative-cumulant};
\item identification of every two-block joining sum with one
      covariance of block-products;
\item a dimension-free bound for reducible block-product covariance;
\item its no-count insertion into the final BFA
      mass--resolvent payer;
\item exact assembly of the \(3+2\) first-connectivity local carrier;
\item closure of both high-pair-together and separated local
      partitions for the prefix geometries paid in V44--V46; and
\item relocation of the genuine residual from the joining coordinate
      to the complete predecessor forest.
\end{enumerate}
\end{theorem}

\begin{proof}
These are
\Cref{sec:V50-audit,
thm:V50-relative-cumulant,
cor:V50-two-block-covariance,
lem:V50-reducible-covariance,
cor:V50-weighted-joining,
thm:V50-exact-joining,
thm:V50-separated-closed,
thm:V50-prefix-residual}.
\end{proof}

\begin{programme}[V51: complete three-row predecessor bank]
\label{prog:V50-V51}
\begin{enumerate}[leftmargin=3em]
\item Regard the entire prefix of the component \(\{a,b,d\}\) as one
      three-row cumulant of row-products, not as a selected sequence
      of local blocks.
\item Apply the connected local-partition identity of V42 to its
      coordinate bank and split only by incidence topology after
      exact assembly.
\item Prove a two-mark BFA estimate for the complete prefix output,
      retaining raw pinching across all prefix representations and
      multiplicities.
\item Combine it with the binary covariance of the component
      \(\{c,e\}\) and the relative joining covariance of V50.
\item Check that no intermediate reduced resolvent is introduced;
      only the final output carries a payer.
\item If a support-uniform prefix estimate fails, identify whether
      the obstruction is coordinate multiplicity, representation
      mass, or noncommuting output compression.
\end{enumerate}
\end{programme}

\begin{firewall}[Claim boundary after fourth-series V50]
\label{fire:V50-claim-boundary}
V50 removes the separated-high-pair joining subfamily by exact
relative-cumulant assembly and closes every \(3+2\) fibre whose
predecessor forest has the V44 or V45 form.  It does not yet prove a
support-uniform complete three-row predecessor-bank estimate, every
predecessor partition, or the full quintic MTA; hence it does not
prove an all-order atlas, WUFC, CPM, cutoff removal, reflection
positivity, Osterwalder--Schrader reconstruction, construction of the
continuum Yang--Mills measure, or a uniform positive continuum
spectral gap.  The full Yang--Mills mass gap has not been proved.
\end{firewall}

\paragraph{Parent-version record.}
V50 was formed from
\texttt{Renewal\_Yang\_Mills\_Mass\_Gap\_Research\_Notes\_V49.tex}.
The V49 parent had \(23\,944\) logical source lines,
\(834\,194\) bytes, and SHA--256
\[
 \texttt{acd337cbc6241d6da6b8953020f67c6c4808353c92863bceba0d1de9ec653f40}.
\]
The next compulsory companion audit is V55.

% =============================================================================
% END APPENDED FOURTH-SERIES V50 MATERIAL
% =============================================================================

% =============================================================================
% BEGIN APPENDED FOURTH-SERIES V51 MATERIAL
% =============================================================================

\chapter{V51: A projection-complete ternary prefix numerator}
\label{chap:V51-ternary-prefix}

\section{Why a numerator theorem is required}
\label{sec:V51-purpose}

The three-row marked-tree atlas proved in archive f1 ends with
\(Q_0(I-\mathcal C_\alpha)^{-1}Q_0\).  It is therefore not a legal
black box inside the \(3+2\) predecessor component: the five-row word
has only one reduced resolvent, at its final output.  Inserting the
resolved ternary atlas in the prefix would spend that inverse twice.
Moreover, its \(Q_0\) suppresses the scalar ternary output, while that
output is a legitimate input to the later joining covariance.

Archive f1 also proves the complete common-bank operator numerator
\texttt{thm:V50-operator-numerator}.  The present chapter extracts the
precise projection-complete statement needed here.  It allows
arbitrary supports, includes the scalar output, and contains no
resolvent.

\section{Arbitrary-support operator cumulant}
\label{sec:V51-operator}

Let \(E\) be a finite coordinate set and let
\(F_i=\bigotimes_{e\in E}F_{i,e}\), \(i=1,2,3\), where
\(F_{i,e}\) is a unitary representation matrix on an active
coordinate and is the identity otherwise.  Put
\[
 q_{i,e}I:=\mathbb E_\alpha F_{i,e}.
\]
For an inactive leg set \(a_{i,e}=0\), and for an active leg let
\(a_{i,e}\ge1\) be its balanced representation mass.  Define
\begin{align}
 H_{ij}&:=\sum_{e\in E}a_{i,e}a_{j,e},
 \label{eq:V51-Hij}\\
 J_{123}&:=\sum_{e\in E}a_{1,e}a_{2,e}a_{3,e}.
 \label{eq:V51-J123}
\end{align}
Thus the summands vanish unless all displayed legs are active.

For independent replacement of coordinate \(e\), define the
two-sided influence
\begin{align}
 \gamma_{i,e}:=\max\biggl\{&
 {1\over\sqrt2}
 \left\|\mathbb E_\alpha[
 (F_i-F_i^{(e)})(F_i-F_i^{(e)})^*]\right\|_{\rm op}^{1/2},
 \notag\\
 &
 {1\over\sqrt2}
 \left\|\mathbb E_\alpha[
 (F_i-F_i^{(e)})^*(F_i-F_i^{(e)})]\right\|_{\rm op}^{1/2}
 \biggr\},
 \label{eq:V51-influence}
\end{align}
and put
\[
 W_{ij}:=\sum_e\gamma_{i,e}\gamma_{j,e},
 \qquad
 \overline W_{ij}:=\min\{W_{ij},1\}.
\]

\begin{lemma}[Inactive-leg completion of the local connected formula]
\label{lem:V51-inactive-completion}
The exact connected-coordinate formula
\(\textup{\Cref{thm:V42-Leonov-Shiryaev}}\), applied to
\(\kappa_3^\alpha(F_1,F_2,F_3)\), is unchanged if every row is
completed by identity matrices on inactive coordinates.  A local
nonsingleton cumulant containing an inactive identity leg vanishes.
Consequently every surviving local hyperedge is supported on rows
which are simultaneously active at that coordinate.
\end{lemma}

\begin{proof}
For \(t\ge2\), a joint cumulant with one argument equal to the
identity is zero.  This follows directly from the moment--cumulant
formula: partitions with the identity index as a singleton cancel
partitions obtained by adjoining that index to each other block.
Equivalently, the logarithm of the moment generating function is
affine in a scalar identity source and has no mixed derivative
containing it.  Hence identity completion adds only local singleton
blocks.  Such blocks have rank zero and do not change the lifted join,
which proves the assertion.
\end{proof}

\begin{lemma}[Wilson influence and junction cards]
\label{lem:V51-local-cards}
For \(s_\alpha=2/\alpha+O(\alpha^{-2})\), uniformly in the
representation and support,
\begin{align}
 \gamma_{i,e}
 &\le C\min\{s_\alpha^{1/2}a_{i,e}^{1/2},1\},
 \label{eq:V51-influence-card}\\
 \left\|
 \kappa_3^{\alpha,e}(F_{1,e},F_{2,e},F_{3,e})
 \right\|_{\rm op}
 &\le Cs_\alpha^2
 a_{1,e}a_{2,e}a_{3,e}.
 \label{eq:V51-junction-card}
\end{align}
The second line is read as zero unless all three legs are active.
\end{lemma}

\begin{proof}
These are respectively the geodesic displacement estimate used in
archive f1
\texttt{eq:V32-Wilson-coordinate-influence} and the
representation-uniform one-link third-cumulant estimate
\texttt{eq:V37-operator-one-link}.  Both estimates are operator
inequalities before isotypic compression.  Identity completion gives
zero influence and zero mixed cumulant on inactive legs by
\Cref{lem:V51-inactive-completion}.
\end{proof}

\begin{theorem}[Projection-complete ternary prefix numerator]
\label{thm:V51-prefix-numerator}
For arbitrary finite supports and arbitrary \(SU(2)\)
representations,
\begin{align}
 \left\|\kappa_3^\alpha(F_1,F_2,F_3)\right\|_{\rm op}
 &\le C\biggl[
  s_\alpha^2J_{123}
  +\sum_{\{i,j,k\}=\{1,2,3\}}
   \overline W_{ij}\,\overline W_{ik}
 \biggr]
 \label{eq:V51-capped-numerator}\\
 &\le Cs_\alpha^2
 \left[
 J_{123}
 H_{12}H_{13}+H_{12}H_{23}+H_{13}H_{23}
 \right].
 \label{eq:V51-polynomial-numerator}
\end{align}
Both inequalities hold before any output projection.  In particular
they hold for the trivial output and, separately, after every
nontrivial isotypic compression, with no output-multiplicity factor.
\end{theorem}

\begin{proof}
At one coordinate the exact moment decomposition is
\[
 Q_{123,e}
 =q_{1,e}q_{2,e}q_{3,e}I
  +\sum_{\{i,j,k\}=\{1,2,3\}}
   q_{k,e}\operatorname{Cov}_e(F_{i,e},F_{j,e})
  +\kappa_3^{\alpha,e}(F_{1,e},F_{2,e},F_{3,e}).
\]
By \Cref{lem:V51-inactive-completion} this identity also covers
coordinates with one or two inactive rows.  Tensor it over \(E\) and
apply the connected local-partition formula of V42.  The surviving
terms are precisely coordinate assignments whose pair bonds and
triple junctions connect the three row labels.

Let
\[
 V:=\sum_e
 \left\|\kappa_3^{\alpha,e}
 (F_{1,e},F_{2,e},F_{3,e})\right\|_{\rm op}.
\]
The operator Efron--Stein covariance inequality proved in archive f1
\texttt{lem:V50-operator-covariance} gives
\[
 \left\|\operatorname{Cov}_e(F_{i,e},F_{j,e})\right\|_{\rm op}
 \le\gamma_{i,e}\gamma_{j,e}.
\]
If \(V\le1\) and all \(W_{ij}\le1/4\), take operator norms in the
connected expansion and mark either its first junction or the first
bonds of two distinct types.  The unmarked remainder is at most
\[
 \exp\!\left(V+W_{12}+W_{13}+W_{23}\right)\le e^{7/4}.
\]
This proves
\begin{equation}
 \left\|\kappa_3^\alpha(F_1,F_2,F_3)\right\|_{\rm op}
 \le C\left[
 V+\overline W_{12}\overline W_{13}
  +\overline W_{12}\overline W_{23}
  +\overline W_{13}\overline W_{23}\right]
 \label{eq:V51-replicated-tree}
\end{equation}
in the small branch.

The five-term third-cumulant formula gives a universal operator cap.
It absorbs the case \(V>1\) and every case with two pair stocks at
least \(1/4\).  In the remaining large branch only one pair stock,
say \(W_{12}\), can exceed \(1/4\).  The exact identity
\[
 \kappa_3^\alpha(F_1,F_2,F_3)
 =\operatorname{Cov}_\alpha(F_1F_2,F_3)
  -\mathbb E_\alpha F_1\operatorname{Cov}_\alpha(F_2,F_3)
  -\mathbb E_\alpha F_2\operatorname{Cov}_\alpha(F_1,F_3)
\]
and the same operator covariance inequality give
\[
 \left\|\kappa_3^\alpha(F_1,F_2,F_3)\right\|_{\rm op}
 \le C(W_{13}+W_{23}).
\]
Since \(\overline W_{12}\ge1/4\), this is bounded by the two tree
terms incident to the strong pair.  Permuting the rows proves
\eqref{eq:V51-replicated-tree} in every branch.

\Cref{lem:V51-local-cards} gives
\[
 V\le Cs_\alpha^2J_{123},
 \qquad
 \overline W_{ij}\le Cs_\alpha H_{ij}.
\]
Substitution proves \eqref{eq:V51-capped-numerator} and
\eqref{eq:V51-polynomial-numerator}.  Orthogonal compression is
contractive, and the trivial projection is just one such compression.
No sum over output copies has been taken, so no multiplicity enters.
\end{proof}

\begin{corollary}[Junction terms already lie under ordinary trees]
\label{cor:V51-junction-tree}
For each choice of center \(i\), with
\(\{i,j,k\}=\{1,2,3\}\),
\[
 J_{123}\le H_{ij}H_{ik}.
\]
Thus the last display in
\eqref{eq:V51-polynomial-numerator} may be replaced by a fixed
multiple of the sum of its three ordinary labelled tree monomials.
\end{corollary}

\begin{proof}
For every coordinate active on all three rows, the product
\(H_{ij}H_{ik}\) contains the diagonal term
\(a_{i,e}^2a_{j,e}a_{k,e}\), which dominates
\(a_{i,e}a_{j,e}a_{k,e}\) because \(a_{i,e}\ge1\).
Sum these diagonal terms.
\end{proof}

\section{Raw coefficient lifting, including the scalar channel}
\label{sec:V51-coefficients}

For this section fix pure input Fourier blocks with trace-class
coefficients \(A_i\), dimensions \(d_i\), heights \(h_i\), and masses
\(\Xi_i>0\).  Normalize them by
\begin{equation}
 e^{\rho h_i}\Xi_i d_i\|A_i\|_1=1.
 \label{eq:V51-input-normalization}
\end{equation}
Let \(B_{\boldsymbol T,\mu}\) be the raw output coefficients obtained
by fusing the complete ternary numerator, including
\(\boldsymbol T=\boldsymbol1\).

\begin{theorem}[No-count unresolvented BFA numerator]
\label{thm:V51-raw-BFA}
With the normalization
\eqref{eq:V51-input-normalization},
\begin{align}
 \sum_{\boldsymbol T,\mu}
 e^{\rho h_{\boldsymbol T}}d_{\boldsymbol T}
 \|B_{\boldsymbol T,\mu}\|_1
 &\le
 {C\over\Xi_1\Xi_2\Xi_3}
 \biggl[
 s_\alpha^2J_{123}
 +\sum_{\{i,j,k\}=\{1,2,3\}}
  \overline W_{ij}\overline W_{ik}
 \biggr]
 \label{eq:V51-raw-BFA-capped}\\
 &\le
 Cs_\alpha^2
 \sum_{\{i,j,k\}=\{1,2,3\}}
 \Xi_i\theta_{ij}\theta_{ik}.
 \label{eq:V51-raw-BFA-tree}
\end{align}
Here
\(\theta_{ij}=H_{ij}/(\Xi_i\Xi_j)\).  The output sum is over
isotypic copies, not over a representation-counting measure.
\end{theorem}

\begin{proof}
In the exact Fourier product formula the coefficient preceding an
output copy contains \(d_1d_2d_3/d_{\boldsymbol T}\).
Multiplication by the raw output weight
\(d_{\boldsymbol T}\) cancels the output dimension.  Height
subadditivity gives
\(e^{\rho h_{\boldsymbol T}}\le
e^{\rho(h_1+h_2+h_3)}\).
Complete isotypic pinching and the multiplicity partial trace are
trace-norm contractions whose raw masses sum to at most one.  Hence
\[
 \sum_{\boldsymbol T,\mu}
 e^{\rho h_{\boldsymbol T}}d_{\boldsymbol T}
 \|B_{\boldsymbol T,\mu}\|_1
 \le
 e^{\rho(h_1+h_2+h_3)}
 d_1d_2d_3\prod_i\|A_i\|_1\,
 \left\|\kappa_3^\alpha(F_1,F_2,F_3)\right\|_{\rm op}.
\]
Use \eqref{eq:V51-input-normalization} and
\eqref{eq:V51-capped-numerator} to obtain
\eqref{eq:V51-raw-BFA-capped}.  Then use
\Cref{cor:V51-junction-tree} and
\[
 {H_{ij}H_{ik}\over\Xi_i\Xi_j\Xi_k}
 =\Xi_i\theta_{ij}\theta_{ik}
\]
to prove \eqref{eq:V51-raw-BFA-tree}.  The argument never removes
\(\boldsymbol T=\boldsymbol1\).
\end{proof}

\begin{firewall}[What has and has not spent a payer]
\label{fire:V51-no-prefix-resolvent}
\Cref{thm:V51-raw-BFA} is a numerator estimate.  Its left side has
neither the output mass \(\Xi(\boldsymbol T)\) nor
\((1-q_{\alpha,\boldsymbol T})^{-1}\).  It is therefore legal below
the final five-row joining coordinate.  It does not follow by
applying the archived resolved ternary MTA twice, and it includes the
scalar output which that MTA's \(Q_0\) removes.
\end{firewall}

\section{Insertion into the \(3+2\) predecessor}
\label{sec:V51-32}

Let the three-row predecessor component be
\(\{a,b,d\}\) and the two-row component be \(\{c,e\}\).
Write \(\theta_{ce}=H_{ce}/(\Xi_c\Xi_e)\).

\begin{lemma}[Complete binary predecessor numerator]
\label{lem:V51-binary-prefix}
After normalizing the \(c,e\) inputs as in
\eqref{eq:V51-input-normalization}, the complete binary covariance
prefix, summed by raw pinching over all outputs including the scalar
one, is bounded by
\[
 Cs_\alpha\theta_{ce}.
\]
\end{lemma}

\begin{proof}
The two-sided operator covariance inequality gives
\[
 \|\operatorname{Cov}_\alpha(F_c,F_e)\|_{\rm op}
 \le W_{ce}\le Cs_\alpha H_{ce}.
\]
The dimension cancellation, height subadditivity, and raw pinching
argument in \Cref{thm:V51-raw-BFA} with two inputs divides this by
\(\Xi_c\Xi_e\).  No output projection is excluded.
\end{proof}

\begin{theorem}[Projection-complete \(3+2\) prefix bound]
\label{thm:V51-32-prefix}
Before the first-connectivity joining coordinate and before the one
final payer, the complete \(3+2\) predecessor forest satisfies
\begin{equation}
 \boxed{
 \mathcal N_{3+2}^{\rm prefix}
 \le
 Cs_\alpha^3
 \sum_{\{i,j,k\}=\{a,b,d\}}
 \Xi_i\theta_{ij}\theta_{ik}\theta_{ce}.}
 \label{eq:V51-32-prefix}
\end{equation}
This includes every prefix output representation and multiplicity,
including a scalar output of either predecessor component.
\end{theorem}

\begin{proof}
Apply \Cref{thm:V51-raw-BFA} to the complete three-row block and
\Cref{lem:V51-binary-prefix} to the two-row block.  The predecessor
partition makes the two coordinate families independent before the
joining coordinate, so their raw coefficient masses multiply.
Multiplication of the two bounds gives
\eqref{eq:V51-32-prefix}.  Raw pinching may be performed separately
on the two tensor factors and then on their product; its total mass
remains at most one.
\end{proof}

\begin{corollary}[Thermal and transition \(3+2\) prefixes close]
\label{cor:V51-thermal-close}
Suppose
\[
 s_\alpha\max\{\Xi_a,\Xi_b,\Xi_d\}\le L
\]
for a fixed \(L\).  After appending the complete relative joining
covariance of V50 and the one final BFA mass--resolvent payer, the
resolved \(3+2\) word is bounded by
\[
 C_Ls_\alpha
 \sum_{\{i,j,k\}=\{a,b,d\}}
 \theta_{ij}\theta_{ik}\theta_{ce}.
\]
It therefore satisfies its five-row tree payment with one additional
Wilson clock.
\end{corollary}

\begin{proof}
The V50 weighted joining theorem costs at most
\(C/s_\alpha\) and sums its outputs by raw pinching.  Multiply that
cost by \eqref{eq:V51-32-prefix}.  The hypothesis gives
\(s_\alpha^2\Xi_i\le Ls_\alpha\) for each center \(i\).
Every monomial displayed is the tree on
\(\{a,b,d\}\) centered at \(i\), together with the edge \(ce\);
the joining covariance supplies the remaining quotient-tree edge
exactly as in \Cref{thm:V50-exact-joining}.  Hence the required
five-row tree monomial is present and the extra factor
\(s_\alpha\) is surplus.
\end{proof}

\section{The exact residual exposed by the capped estimate}
\label{sec:V51-residual}

\begin{assessment}[Hot-center information was discarded too early]
\label{ass:V51-hot-residual}
The polynomial relaxation
\(\overline W_{ij}\le Cs_\alpha H_{ij}\) loses information when a
three-row center has \(s_\alpha\Xi_i\gg1\).  The exact connected
coordinate expansion still contains singleton Wilson means on every
coordinate not selected as a bond or junction.  Those means include
the private and exclusive spectator banks of a hot center and have
quadratic product decay.  They were bounded by one in
\Cref{thm:V51-prefix-numerator}.  Consequently the remaining
hot-center factor \(s_\alpha^2\Xi_i\) is not evidence of a divergent
word; it identifies the precise factor that must be repaid by
retaining the spectator product before taking norms.
\end{assessment}

\begin{theorem}[Exhaustive next obstruction]
\label{thm:V51-next-obstruction}
After V51, a \(3+2\) first-connectivity fibre can remain outside the
proved marked-tree bound only if all of the following occur:
\begin{enumerate}[label=\textup{(\roman*)},leftmargin=3em]
\item the three-row prefix has a center \(i\) with
      \(s_\alpha\Xi_i\) unbounded;
\item the complete projection-free three-row cumulant is required,
      so neither its scalar nor nontrivial output may be discarded;
\item the prefix numerator is estimated after at least one hot
      spectator mean has been replaced by its unit bound; and
\item no already-proved binary-separable or single-ternary local
      realization from V44--V46 applies.
\end{enumerate}
Thus the next task is a spectator-retaining version of
\eqref{eq:V51-capped-numerator}, not another output resolvent
estimate.
\end{theorem}

\begin{proof}
If \textup{(i)} fails, \Cref{cor:V51-thermal-close} closes the word.
V50 proves that the joining coordinate, including its output sum and
single payer, is already closed.  V44--V46 close the geometries
excluded by \textup{(iv)}.  The only relaxation in the remaining
prefix proof which depends on a hot total mass is the deletion of
unmarked singleton means in the replicated-tree estimate.  The scalar
and nontrivial prefix projections were both retained by
\Cref{thm:V51-raw-BFA}, proving \textup{(ii)} and locating the loss in
\textup{(iii)}.
\end{proof}

\section{V51 result ledger and V52 programme}
\label{sec:V51-ledger}

\begin{theorem}[Fourth-series V51 result ledger]
\label{thm:V51-results}
Relative to V1--V50, V51 proves:
\begin{enumerate}[label=\textup{(V51.\arabic*)},leftmargin=6.3em]
\item identity completion of the exact local connected-partition
      formula on arbitrary three-row supports;
\item a projection-complete, representation-uniform ternary
      numerator estimate with capped pair influences;
\item the polynomial two-mark consequence in global pair stocks;
\item inclusion of both the scalar and every nontrivial output
      compression without a multiplicity count;
\item raw BFA coefficient lifting with exact output-dimension
      cancellation and no intermediate resolvent;
\item the complete projection-free \(3+2\) predecessor bound
      \eqref{eq:V51-32-prefix};
\item closure with a surplus Wilson clock of the thermal and
      transition three-row-center regimes; and
\item localization of the remaining hot-center loss to discarded
      singleton spectator means.
\end{enumerate}
\end{theorem}

\begin{proof}
These are
\Cref{lem:V51-inactive-completion,
thm:V51-prefix-numerator,
cor:V51-junction-tree,
thm:V51-raw-BFA,
fire:V51-no-prefix-resolvent,
thm:V51-32-prefix,
cor:V51-thermal-close,
thm:V51-next-obstruction}.
\end{proof}

\begin{programme}[V52: spectator-retaining hot prefix]
\label{prog:V51-V52}
\begin{enumerate}[leftmargin=3em]
\item In the exact connected-coordinate expansion, mark the first
      junction or the first two distinct pair bonds but retain every
      unmarked singleton multiplier.
\item Decompose the unmarked coordinates into private,
      exclusive-pair, and triple-overlap spectator banks.
\item Prove first- and second-moment bounds for their complete
      Wilson multiplier product, without comparing an input
      multiplier with a fused output multiplier.
\item Use the maximal hot center to recover the factor lost in
      \(s_\alpha^2\Xi_i\), keeping the scalar prefix output.
\item Insert the resulting prefix estimate into
      \Cref{thm:V50-exact-joining} with only the final payer.
\item If the hot mass lies entirely on coordinates already selected
      by the two marks, split into diffuse and concentrated mark
      banks before applying the multiplier debit.
\end{enumerate}
\end{programme}

\begin{firewall}[Claim boundary after fourth-series V51]
\label{fire:V51-claim-boundary}
V51 proves a support- and representation-uniform, projection-complete
three-row prefix numerator and closes the \(3+2\) source whenever the
three-row center is thermal or in the Wilson transition range.  It
does not yet close the hot-center spectator branch, every \(3+2\)
source, the full quintic MTA, an all-order atlas, WUFC, CPM, cutoff
removal, reflection positivity, Osterwalder--Schrader reconstruction,
construction of the continuum Yang--Mills measure, or a uniform
positive continuum spectral gap.  The full Yang--Mills mass gap has
not been proved.
\end{firewall}

\paragraph{Parent-version record.}
V51 was formed from
\texttt{Renewal\_Yang\_Mills\_Mass\_Gap\_Research\_Notes\_V50.tex}.
The V50 parent had \(24\,407\) logical source lines,
\(850\,226\) bytes, and SHA--256
\[
 \texttt{85d3c170864c7a7a6f62b4f588bddee697e1f0c1acb6b020ec6e794173bde7df}.
\]
The next compulsory companion audit is V55.

% =============================================================================
% END APPENDED FOURTH-SERIES V51 MATERIAL
% =============================================================================

% =============================================================================
% BEGIN APPENDED FOURTH-SERIES V52 MATERIAL
% =============================================================================

\chapter{V52: Spectator retention closes the pair-bond \(3+2\) source}
\label{chap:V52-hot-prefix}

\section{Private and overlap mass}
\label{sec:V52-split}

Keep the arbitrary-support three-row notation of V51.  For
\(\{i,j,k\}=\{1,2,3\}\), let
\begin{align}
 P_i&:=\sum_{\substack{e:\ i\ {\rm active}\\
                 j,k\ {\rm inactive}}}a_{i,e},
 \label{eq:V52-private-mass}\\
 Z_i&:=\Xi_i-P_i,
 \label{eq:V52-overlap-mass}\\
 b_i&:=\prod_{\substack{e:\ i\ {\rm active}\\
                 j,k\ {\rm inactive}}}q_{\alpha,R_{i,e}}.
 \label{eq:V52-private-multiplier}
\end{align}
Empty products equal one.  Pair-exclusive and triple-overlap
coordinates belong to \(Z_i\), not to \(P_i\).

\begin{lemma}[Exact private-bank factorization]
\label{lem:V52-private-factor}
Let \(G_i\) be obtained from \(F_i\) by deleting its private
coordinates.  Then
\begin{equation}
 \boxed{
 \kappa_3^\alpha(F_1,F_2,F_3)
 =b_1b_2b_3\,
 \kappa_3^\alpha(G_1,G_2,G_3).}
 \label{eq:V52-private-factor}
\end{equation}
The identity holds before every Fourier output projection, including
the trivial projection.
\end{lemma}

\begin{proof}
Fix a private coordinate of row \(i\).  Every term in the five-term
formula for \(\kappa_3(F_1,F_2,F_3)\) contains the first variable
exactly once, either in a joint moment or in a product of moments.
The private coordinate is independent of every other coordinate and
of the other two rows.  Its expectation therefore factors from every
one of the five terms as the same central scalar
\(q_{\alpha,R_{i,e}}\).  Remove the private coordinates one at a
time and multiply the resulting scalars.  Since this is an operator
identity before compression, all output projections inherit it.
\end{proof}

\begin{lemma}[Overlap mass forces a large incident stock]
\label{lem:V52-overlap-stock}
For every center \(i\),
\begin{equation}
 Z_i\le H_{ij}+H_{ik}.
 \label{eq:V52-overlap-stock}
\end{equation}
Consequently, if \(Z_i\ge\Xi_i/2\), then
\[
 \max\{H_{ij},H_{ik}\}\ge\Xi_i/4.
\]
\end{lemma}

\begin{proof}
At a coordinate counted by \(Z_i\), at least one of \(j,k\) is
active and its balanced mass is at least one.  Hence
\[
 a_{i,e}\le
 a_{i,e}a_{j,e}+a_{i,e}a_{k,e}.
\]
Sum over the active coordinates of row \(i\).  The last assertion
follows by the pigeonhole principle.
\end{proof}

\begin{record}[Quadratic private multiplier debit]
\label{rec:V52-private-debit}
The product second-moment theorem in archive f1
\texttt{eq:V32-balanced-second-moment} gives
\begin{equation}
 0\le b_i\le
 C\min\left\{1,{1\over(s_\alpha P_i)^2}\right\}.
 \label{eq:V52-private-debit}
\end{equation}
The constant is independent of the number of private coordinates and
their representations.
\end{record}

This is an exact Wilson multiplier statement, not a comparison
between an input multiplier and a fused output multiplier.  It is
therefore compatible with the no-monotonicity firewall of the
preceding archives.

\section{The fully capped ternary numerator}
\label{sec:V52-capped}

\begin{lemma}[Capping the junction stock]
\label{lem:V52-capped-junction}
The V51 numerator estimate improves to
\begin{equation}
 \left\|\kappa_3^\alpha(F_1,F_2,F_3)\right\|_{\rm op}
 \le C\biggl[
 \min\{s_\alpha^2J_{123},1\}
 +\sum_{\{i,j,k\}=\{1,2,3\}}
  \overline W_{ij}\overline W_{ik}
 \biggr].
 \label{eq:V52-fully-capped}
\end{equation}
After deleting private coordinates, the right side has the same
\(J_{123}\) and \(W_{ij}\).
\end{lemma}

\begin{proof}
The proof of
\eqref{eq:V51-replicated-tree} already uses the universal five-term
operator cap when the junction stock
\[
 V=\sum_e\|\kappa_3^{\alpha,e}\|_{\rm op}
\]
is at least one.  Thus \(V\) may be replaced there by
\(\min\{V,1\}\).  By \Cref{lem:V51-local-cards},
\[
 \min\{V,1\}\le C\min\{s_\alpha^2J_{123},1\}.
\]
Private coordinates contain only one active row, so they contribute
neither a pair influence product nor a ternary junction.
\end{proof}

\begin{corollary}[Centerwise spectator-retaining bound]
\label{cor:V52-centerwise}
For each center \(i\), every tree term incident to \(ij,ik\), and the
junction term when assigned to that center, may be multiplied by
\(b_i\) in \eqref{eq:V52-fully-capped}.
\end{corollary}

\begin{proof}
The exact factor in \eqref{eq:V52-private-factor} is
\(b_1b_2b_3\le b_i\), since the Wilson multipliers are nonnegative
and at most one.  The reduced cumulant has the same connected stocks
by \Cref{lem:V52-capped-junction}.
\end{proof}

\section{The hot-center payment lemma}
\label{sec:V52-payment}

For a center \(i\), abbreviate
\[
 x:=\Xi_i,\qquad h_1:=H_{ij},\qquad h_2:=H_{ik},
 \qquad
 u_\nu:=\min\{s_\alpha h_\nu,1\}.
\]
When \(h_1h_2=0\), the corresponding connected tree and junction
terms vanish, so the ratios below are read only for \(h_1h_2>0\).

\begin{lemma}[Private/overlap dichotomy pays a centered pair tree]
\label{lem:V52-pair-payment}
Uniformly in \(x,h_1,h_2\),
\begin{equation}
 b_i\,{x\,u_1u_2\over h_1h_2}\le C.
 \label{eq:V52-pair-payment}
\end{equation}
More precisely, the left side is \(O(s_\alpha)\) if
\(s_\alpha x\le L\) or \(Z_i\ge x/2\), and it is
\(O(x^{-1})\) if \(P_i>x/2\) and \(s_\alpha x>L\).
\end{lemma}

\begin{proof}
If \(s_\alpha x\le L\), use \(u_\nu\le s_\alpha h_\nu\):
\[
 b_i\,{x\,u_1u_2\over h_1h_2}
 \le s_\alpha^2x\le Ls_\alpha.
\]
If \(Z_i\ge x/2\), then
\Cref{lem:V52-overlap-stock} gives, after relabelling,
\(h_1\ge x/4\).  Use \(u_1\le1\) and
\(u_2\le s_\alpha h_2\) to obtain
\[
 b_i\,{x\,u_1u_2\over h_1h_2}
 \le {s_\alpha x\over h_1}\le4s_\alpha.
\]
The remaining case has \(P_i>x/2\).  By
\Cref{rec:V52-private-debit},
\[
 b_i\le {C\over s_\alpha^2P_i^2}
 \le {C\over s_\alpha^2x^2}.
\]
Using \(u_\nu\le s_\alpha h_\nu\) now gives
\[
 b_i\,{x\,u_1u_2\over h_1h_2}
 \le Cx^{-1}.
\]
The three cases are exhaustive.
\end{proof}

\begin{lemma}[Thermal and private-hot junction payment]
\label{lem:V52-junction-payment}
If either \(s_\alpha x\le L\) or \(P_i>x/2\), then
\begin{equation}
 b_i\,{x\,\min\{s_\alpha^2J_{123},1\}
              \over h_1h_2}
 \le C_Ls_\alpha.
 \label{eq:V52-junction-payment}
\end{equation}
\end{lemma}

\begin{proof}
Use
\[
 J_{123}\le h_1h_2
\]
from \Cref{cor:V51-junction-tree}.  In the thermal case the ratio is
at most \(s_\alpha^2x\le Ls_\alpha\).  In the private-hot case,
\Cref{rec:V52-private-debit} gives
\[
 b_i\,{x\,\min\{s_\alpha^2J_{123},1\}\over h_1h_2}
 \le b_i s_\alpha^2x\le Cx^{-1}\le C_Ls_\alpha.
\]
\end{proof}

\begin{remark}[Why pair-exclusive hot mass is harmless]
\label{rem:V52-exclusive}
Mass on a coordinate shared by \(i\) with only one other row is not a
private spectator.  It enlarges one of \(h_1,h_2\) and is paid by the
overlap branch of \Cref{lem:V52-pair-payment}.  No decay is demanded
from a multiplier which has already participated in a selected
covariance bond.
\end{remark}

\section{Rigorous closure and the surviving junction branch}
\label{sec:V52-32}

\begin{theorem}[All centered pair-bond \(3+2\) sources close]
\label{thm:V52-pair-source-close}
Consider the complete \(3+2\) first-connectivity source, with
three-row component \(\{a,b,d\}\) and binary component \(\{c,e\}\).
Every term in which the connected three-row numerator is carried by
two pair bonds satisfies its five-row marked-tree payment uniformly
in support, representations, outputs, and multiplicities, including
scalar predecessor outputs.
\end{theorem}

\begin{proof}
For the three-row tree centered at \(i\), the fully capped raw prefix
contribution after input normalization is
\[
 {C b_i u_1u_2\over\Xi_i\Xi_j\Xi_k}.
\]
The binary predecessor contributes
\(Cs_\alpha\theta_{ce}\), and the final joining covariance, output
mass, and single reduced resolvent cost at most
\(C/s_\alpha\) by V50.  The last two factors cancel, so division by
the desired internal forest monomial
\[
 \theta_{ij}\theta_{ik}\theta_{ce}
 ={h_1h_2H_{ce}\over
   \Xi_i^2\Xi_j\Xi_k\Xi_c\Xi_e}
\]
leaves exactly
\[
 Cb_i\,{\Xi_i u_1u_2\over h_1h_2}.
\]
This is uniformly bounded by
\Cref{lem:V52-pair-payment}.  The exact relative joining covariance
supplies the remaining quotient-tree edge as in
\Cref{thm:V50-exact-joining}.  All coefficient sums use the raw
pinching theorem of V51, so neither a representation count nor a
scalar-output exclusion occurs.
\end{proof}

\begin{theorem}[Thermal and private-hot junction sources close]
\label{thm:V52-junction-partial-close}
The same \(3+2\) source is closed when its three-row connected
numerator is carried by a local junction and, for at least one
assigned tree center \(i\), either
\[
 s_\alpha\Xi_i\le L
 \qquad\hbox{or}\qquad
 P_i>\Xi_i/2.
\]
\end{theorem}

\begin{proof}
Repeat the proof of
\Cref{thm:V52-pair-source-close}, replacing \(u_1u_2\) by
\(\min\{s_\alpha^2J_{123},1\}\).  Before the binary-prefix clock
cancels the final payer, the quotient by the selected ternary tree is
the left side of \eqref{eq:V52-junction-payment}.  That lemma supplies
the necessary \(O(s_\alpha)\), so the full quotient is uniformly
bounded.
\end{proof}

\begin{theorem}[Rigorous \(3+2\) reduction after V52]
\label{thm:V52-rigorous-reduction}
Every complete \(3+2\) first-connectivity source is closed except
possibly the following sharply specified branch:
\begin{enumerate}[label=\textup{(\roman*)},leftmargin=3em]
\item the three-row prefix contribution is carried by a local
      junction rather than by two pair bonds;
\item every candidate tree center has
      \(s_\alpha\Xi_i\gg1\) and overlap mass at least \(\Xi_i/2\);
\item no private multiplier debit is available at the assigned
      center.
\end{enumerate}
\end{theorem}

\begin{proof}
\Cref{thm:V52-pair-source-close} removes all pair-bond realizations.
\Cref{thm:V52-junction-partial-close} removes thermal and private-hot
junction assignments.  In the remaining overlap-hot branch, neither
\[
 {s_\alpha^2\Xi_iJ_{123}\over H_{ij}H_{ik}}
 \quad\hbox{nor}\quad
 {\Xi_i\over H_{ij}H_{ik}}
\]
is known to be \(O(s_\alpha)\) uniformly; these are respectively the
uncapped and capped junction quotients.  Thus no size assumption on
\(J_{123}\) alone closes the branch.  The only case not supplied with
the final \(s_\alpha\) debit is exactly \textup{(i)--(iii)}.  V50 pays
the complete joining coordinate and
V51 retains all predecessor projections, so no separate output
residual remains.
\end{proof}

\section{Exact upstream correction}
\label{sec:V52-upstream}

\begin{assessment}[The missing object is a junction allocation]
\label{ass:V52-upstream}
For overlap-hot centers the elementary inequalities give only
\[
 {\Xi_i\min\{s_\alpha^2J_{123},1\}
  \over H_{ij}H_{ik}}\le C,
\]
whereas cancellation of the one final payer requires an additional
\(s_\alpha\).  The local two-clock bound has been hidden inside the
unit cap.  It is not legitimate to assert that a dyadic
first-crossing level restores that clock without a coefficientwise
allocation theorem.  The missing statement is a deterministic
three-sequence inequality, or else a sharper retained Wilson
junction multiplier.  This is upstream of output fusion and
resolvent analysis.
\end{assessment}

\section{V52 result ledger and V53 programme}
\label{sec:V52-ledger}

\begin{theorem}[Fourth-series V52 result ledger]
\label{thm:V52-results}
Relative to V1--V51, V52 proves:
\begin{enumerate}[label=\textup{(V52.\arabic*)},leftmargin=6.3em]
\item exact factorization of every private three-row coordinate bank
      into its Wilson singleton multiplier product;
\item a support-uniform quadratic debit for that product;
\item an overlap-mass lower bound for one incident pair stock;
\item the fully capped complete ternary numerator;
\item uniform payment of every centered two-pair-bond tree in the
      thermal, overlap-hot, and private-hot regimes;
\item rigorous payment of thermal and private-hot junction terms;
\item reduction of the only unproved \(3+2\) source to an
      overlap-hot junction allocation problem at arbitrary junction
      stock size; and
\item rejection of an unstated dyadic first-crossing argument.
\end{enumerate}
\end{theorem}

\begin{proof}
These are
\Cref{lem:V52-private-factor,
rec:V52-private-debit,
lem:V52-overlap-stock,
lem:V52-capped-junction,
lem:V52-pair-payment,
lem:V52-junction-payment,
thm:V52-pair-source-close,
thm:V52-rigorous-reduction,
ass:V52-upstream}.
\end{proof}

\begin{programme}[V53: junction allocation inequality]
\label{prog:V52-V53}
\begin{enumerate}[leftmargin=3em]
\item Write the triple-active data as three nonnegative sequences
      \(x_e,y_e,z_e\ge1\) and compare
      \(\min\{s_\alpha^2\sum_ex_ey_ez_e,1\}\) with the three products
      \((\sum_ex_ey_e)(\sum_ex_ez_e)\).
\item Test the claimed extra \(s_\alpha\) on one-atom, diffuse,
      imbalanced, and alternating-height configurations.
\item If it is false as a deterministic inequality, retain the exact
      Wilson junction multiplier rather than its unit cap and identify
      the missing heat moment.
\item If it is true after a necessary mass hypothesis, prove the
      hypothesis from the overlap/private dichotomy.
\item Only after an exact inequality is available, insert it into raw
      pinching and decide the complete \(3+2\) source.
\end{enumerate}
\end{programme}

\begin{firewall}[Claim boundary after fourth-series V52]
\label{fire:V52-claim-boundary}
V52 closes every pair-bond realization in the complete
projection-free \(3+2\) prefix and closes the thermal and
private-hot junction branches.  It does not close the
overlap-hot junction branch, every \(3+2\) source, the full
quintic MTA, an all-order atlas, WUFC, CPM, cutoff removal, reflection
positivity, Osterwalder--Schrader reconstruction, construction of the
continuum Yang--Mills measure, or a uniform positive continuum
spectral gap.  The full Yang--Mills mass gap has not been proved.
\end{firewall}

\paragraph{Parent-version record.}
V52 was formed from
\texttt{Renewal\_Yang\_Mills\_Mass\_Gap\_Research\_Notes\_V51.tex}.
The V51 parent had \(24\,948\) logical source lines,
\(869\,944\) bytes, and SHA--256
\[
 \texttt{1d3bef0bbd29d2e7a8858a84fdfa549eb95bed49f741687f13bd6047567cc497}.
\]
The next compulsory companion audit is V55.

% =============================================================================
% END APPENDED FOURTH-SERIES V52 MATERIAL
% =============================================================================

% =============================================================================
% BEGIN APPENDED FOURTH-SERIES V53 MATERIAL
% =============================================================================

\chapter{V53: Correct clock accounting closes every \(3+2\) source}
\label{chap:V53-32-closure}

\section{Correction of the V52 frontier}
\label{sec:V53-correction}

V52 correctly proves all of its stated estimates, but its final
assessment asks the overlap-hot junction to supply one Wilson clock
too many.  In a \(3+2\) predecessor, the binary component already
contributes \(s_\alpha\theta_{ce}\), while the complete final
mass--resolvent--joining carrier costs \(C/s_\alpha\).  These two
factors cancel.  The ternary junction quotient need only be uniformly
bounded by a ternary tree sum; it need not be \(O(s_\alpha)\).

\begin{firewall}[Status of the V52 statements]
\label{fire:V53-V52-status}
\Cref{lem:V52-private-factor,
lem:V52-overlap-stock,
lem:V52-capped-junction,
lem:V52-pair-payment,
lem:V52-junction-payment,
thm:V52-pair-source-close,
thm:V52-junction-partial-close}
remain valid.  The alleged residual in
\Cref{thm:V52-rigorous-reduction,ass:V52-upstream} is superseded:
it is an overrestriction caused by the extra-clock demand, not a
counterexample to any proved estimate.
\end{firewall}

\section{A deterministic junction allocation}
\label{sec:V53-allocation}

\begin{lemma}[Three-stock allocation inequality]
\label{lem:V53-three-stock}
Let \(A,B,C,J,X_1,X_2,X_3\ge0\), with \(X_i>0\), and suppose
\begin{align}
 J&\le AB,\qquad J\le AC,\qquad J\le BC,
 \label{eq:V53-J-products}\\
 X_1&\le K(A+B),\qquad
 X_2\le K(A+C),\qquad
 X_3\le K(B+C)
 \label{eq:V53-X-incidence}
\end{align}
for some \(K\ge1\).  Then, for \(0<s\le1\),
\begin{equation}
 \boxed{
 \min\{s^2J,1\}
 \le 2Ks\left(
 {AB\over X_1}+{AC\over X_2}+{BC\over X_3}
 \right).}
 \label{eq:V53-three-stock}
\end{equation}
In particular the same bound without the factor \(s\) also holds.
\end{lemma}

\begin{proof}
Relabel \(A,B,C\) so that \(A\le B\le C\); the three terms and
hypotheses are symmetric under the corresponding relabelling of the
vertices.  From \eqref{eq:V53-X-incidence},
\begin{align*}
 {AB\over X_1}
 +{AC\over X_2}
 +{BC\over X_3}
 &\ge {1\over K}\left(
 {AB\over A+B}
 +{AC\over A+C}
 +{BC\over B+C}\right)\\
 &\ge {1\over K}\left({A\over2}+{A\over2}+{B\over2}\right)
 \ge {B\over2K}.
\end{align*}
The first inequality in \eqref{eq:V53-J-products} gives
\(J\le AB\le B^2\).  If \(sB\le1\), then
\[
 \min\{s^2J,1\}\le s^2B^2\le sB.
\]
If \(sB>1\), then
\[
 \min\{s^2J,1\}\le1<sB.
\]
In either case the left side is at most \(sB\), which is bounded by
the right side of \eqref{eq:V53-three-stock}.
\end{proof}

\begin{lemma}[Overlap masses meet the allocation hypotheses]
\label{lem:V53-overlap-hypotheses}
Set
\[
 A=H_{12},\qquad B=H_{13},\qquad C=H_{23},
 \qquad J=J_{123},\qquad X_i=\Xi_i.
\]
If \(Z_i\ge\Xi_i/2\) for all \(i\), then
\eqref{eq:V53-J-products}--\eqref{eq:V53-X-incidence} hold with
\(K=2\).
\end{lemma}

\begin{proof}
\Cref{cor:V51-junction-tree} applied at each of the three centers
gives
\[
 J_{123}\le H_{12}H_{13},\qquad
 J_{123}\le H_{12}H_{23},\qquad
 J_{123}\le H_{13}H_{23}.
\]
By \Cref{lem:V52-overlap-stock},
\[
 \Xi_1\le2(H_{12}+H_{13}),\quad
 \Xi_2\le2(H_{12}+H_{23}),\quad
 \Xi_3\le2(H_{13}+H_{23}).
\]
These are exactly the required hypotheses.
\end{proof}

\begin{corollary}[Normalized overlap-hot junction tree]
\label{cor:V53-normalized-junction}
In the overlap branch \(Z_i\ge\Xi_i/2\) for all three rows,
\begin{equation}
 {\min\{s_\alpha^2J_{123},1\}
  \over\Xi_1\Xi_2\Xi_3}
 \le
 Cs_\alpha
 \sum_{\{i,j,k\}=\{1,2,3\}}
 \theta_{ij}\theta_{ik}.
 \label{eq:V53-normalized-junction}
\end{equation}
\end{corollary}

\begin{proof}
Apply
\Cref{lem:V53-three-stock,lem:V53-overlap-hypotheses} and divide by
\(\Xi_1\Xi_2\Xi_3\).  For example,
\[
 {H_{12}H_{13}/\Xi_1\over\Xi_1\Xi_2\Xi_3}
 =\theta_{12}\theta_{13}.
\]
The other two terms are identical after relabelling.
\end{proof}

\section{Exact one-payer ledger}
\label{sec:V53-ledger-proof}

\begin{lemma}[The binary clock cancels the final payer once]
\label{lem:V53-clock-cancellation}
For the \(3+2\) predecessor
\(\{\{a,b,d\},\{c,e\}\}\), let
\(\mathcal N_3\) be a normalized ternary prefix numerator.  The
complete binary component, joining covariance, output BFA mass, and
one final reduced resolvent multiply it by at most
\begin{equation}
 C\theta_{ce}.
 \label{eq:V53-clock-cancellation}
\end{equation}
No residual power of \(s_\alpha\) or \(s_\alpha^{-1}\) remains.
\end{lemma}

\begin{proof}
\Cref{lem:V51-binary-prefix} contributes
\(Cs_\alpha\theta_{ce}\).  The complete relative joining covariance,
including final output mass and the unique reduced resolvent, costs
at most \(C/s_\alpha\) by the V50 weighted joining theorem.  Their
product is \(C\theta_{ce}\).  Both steps use raw pinching, so their
successive output sums remain subprobability sums.  Introducing any
additional inverse clock would count the same final resolvent twice.
\end{proof}

\begin{theorem}[Closure of the overlap-hot junction branch]
\label{thm:V53-overlap-junction-close}
The overlap-hot junction branch left open in
\Cref{thm:V52-rigorous-reduction} satisfies the complete quintic
marked-tree payment.
\end{theorem}

\begin{proof}
The projection-complete ternary numerator of
\Cref{lem:V52-capped-junction}, after input normalization and raw
pinching, contributes
\[
 {C\min\{s_\alpha^2J_{abd},1\}\over
   \Xi_a\Xi_b\Xi_d}.
\]
In the residual branch every one of \(a,b,d\) has at least half of
its mass on overlap coordinates.  Apply
\Cref{cor:V53-normalized-junction} on these three rows.  Then apply
\Cref{lem:V53-clock-cancellation}.  The result is
\[
 C
 \sum_{\{i,j,k\}=\{a,b,d\}}
 \theta_{ij}\theta_{ik}\theta_{ce}.
\]
Each displayed monomial is the internal predecessor forest.  The
complete first-connectivity covariance at \(r\) supplies a quotient
edge between the two components, producing an ordinary five-row
tree.  Its output projections, including scalar prefix outputs, were
retained in V51.  Hence the full branch is bounded by the global
quintic tree atlas.
\end{proof}

\begin{theorem}[Complete \(3+2\) first-connectivity theorem]
\label{thm:V53-complete-32}
All ten quintic first-connectivity source families whose predecessor
partition has type \(3+2\) satisfy the support- and
representation-uniform BFA marked-tree bound.  The conclusion holds
for arbitrary coordinate incidence, coincident local roles, growing
representations, scalar or nontrivial predecessor outputs, payer
coincidence, and arbitrary output multiplicities.
\end{theorem}

\begin{proof}
V50 assembles the complete joining partition family into one relative
covariance and pays its output sum without splitting the cofinal pair.
V51 proves the projection-complete ternary prefix numerator and raw
coefficient lifting.  V52 closes every two-pair-bond realization and
the thermal and private-hot junction branches.  The only remaining
overlap-hot junction branch is
\Cref{thm:V53-overlap-junction-close}.  These cases exhaust the exact
connected-coordinate expansion.  Relabelling the three- and two-row
blocks covers the ten set partitions of type \(3+2\).
\end{proof}

\begin{corollary}[The V50 predecessor residual is discharged]
\label{cor:V53-V50-discharged}
No complete \(3+2\) first-connectivity fibre remains in the quintic
residual.
\end{corollary}

\begin{proof}
This is exactly \Cref{thm:V53-complete-32}.
\end{proof}

\section{The next upstream source}
\label{sec:V53-next}

\begin{assessment}[Projection-complete quartic numerator is next]
\label{ass:V53-next}
The other two-block predecessor type is \(4+1\).  Archive f3 proves
the complete quartic MTA, but that theorem carries its own final
output payer and cannot be inserted unchanged below a quintic joining
coordinate.  The next nonduplicative task is therefore the quartic
analogue of V51: extract a projection-complete four-row numerator,
including its scalar output, and retain the exact internal
three-edge tree stocks before the one quintic payer.
\end{assessment}

\section{V53 result ledger and V54 programme}
\label{sec:V53-results}

\begin{theorem}[Fourth-series V53 result ledger]
\label{thm:V53-results}
Relative to V1--V52, V53 proves:
\begin{enumerate}[label=\textup{(V53.\arabic*)},leftmargin=6.3em]
\item correction of the extra-clock demand in the V52 assessment;
\item the deterministic three-stock junction allocation
      \eqref{eq:V53-three-stock};
\item verification of its hypotheses from the three overlap-mass
      inequalities and the junction-to-tree inequalities;
\item a normalized projection-complete junction bound with a surplus
      Wilson clock;
\item exact cancellation of the binary predecessor clock against the
      single final payer;
\item closure of the overlap-hot junction branch; and
\item closure of all ten \(3+2\) first-connectivity source families.
\end{enumerate}
\end{theorem}

\begin{proof}
These are
\Cref{fire:V53-V52-status,
lem:V53-three-stock,
lem:V53-overlap-hypotheses,
cor:V53-normalized-junction,
lem:V53-clock-cancellation,
thm:V53-overlap-junction-close,
thm:V53-complete-32}.
\end{proof}

\begin{programme}[V54: projection-complete quartic prefix]
\label{prog:V53-V54}
\begin{enumerate}[leftmargin=3em]
\item Express the full four-row product cumulant by connected local
      hyperedge assignments on arbitrary supports.
\item Retain the complete three-edge tree skeleton before taking
      operator norms; include coincident ternary and quartic local
      junctions.
\item Extract an unresolvented raw BFA coefficient inequality from
      the proved quartic archive, with the scalar output restored.
\item Separate private banks exactly and retain their Wilson
      multiplier products.
\item Insert the result into the \(4+1\) first-connectivity covariance
      with only the final quintic payer.
\item If the singleton predecessor supplies no cancelling clock,
      identify the exact additional clock required from the joining
      covariance rather than importing a second resolvent.
\end{enumerate}
\end{programme}

\begin{firewall}[Claim boundary after fourth-series V53]
\label{fire:V53-claim-boundary}
V53 proves every quintic \(3+2\) first-connectivity source.  It does
not prove the \(4+1\), \(3+1+1\), \(2+2+1\),
\(2+1+1+1\), or discrete source families, the full quintic MTA, an
all-order atlas, WUFC, CPM, cutoff removal, reflection positivity,
Osterwalder--Schrader reconstruction, construction of the continuum
Yang--Mills measure, or a uniform positive continuum spectral gap.
The full Yang--Mills mass gap has not been proved.
\end{firewall}

\paragraph{Parent-version record.}
V53 was formed from
\texttt{Renewal\_Yang\_Mills\_Mass\_Gap\_Research\_Notes\_V52.tex}.
The V52 parent had \(25\,372\) logical source lines,
\(884\,455\) bytes, and SHA--256
\[
 \texttt{332f19d5cc773fd97f4f79d67db219ae99b4891ab4c4824e9d8aa27e0eec7bed}.
\]
The next compulsory companion audit is V55.

% =============================================================================
% END APPENDED FOURTH-SERIES V53 MATERIAL
% =============================================================================

% =============================================================================
% BEGIN APPENDED FOURTH-SERIES V54 MATERIAL
% =============================================================================

\chapter{V54: Payer stripping and a scalar lift close \(4+1\)}
\label{chap:V54-quartic-prefix}

\section{The extraction problem}
\label{sec:V54-purpose}

Archive f3
\texttt{thm:V96-global-quartic-MTA} proves the complete connected
four-row BFA carrier with one allocated payer:
\[
 \kappa_\otimes^{\rm abs}(\mathcal C_4^{\rm pay})
 \le C\sum_{\tau\in\mathfrak T_4}
          \prod_{ij\in E(\tau)}\theta_{ij}.
\]
The theorem includes the repaired \(3+1\) and \(2+1+1\) rescue
atlases, the \(2+2\) source, and the selected discrete source.  It
must not be inserted directly into a quintic prefix, because the
quintic word has its own final payer.  This chapter removes the old
payer coefficientwise.  The only exceptional output is the scalar
channel, which is handled by an injective private-spectator lift.

\section{Nontrivial-output payer stripping}
\label{sec:V54-strip}

\begin{lemma}[Upper product-gap card]
\label{lem:V54-upper-gap}
For every nontrivial product representation \(\boldsymbol T\),
\begin{equation}
 0<1-q_{\alpha,\boldsymbol T}
 \le Cs_\alpha\Xi(\boldsymbol T).
 \label{eq:V54-upper-gap}
\end{equation}
The constant is uniform in support and representation.
\end{lemma}

\begin{proof}
For one \(SU(2)\) factor, the exact Bessel multiplier and the
Casimir upper-gap estimate in archive f1 give
\[
 1-q_{\alpha,T_e}\le Cs_\alpha a_{T_e}.
\]
For numbers \(0\le q_e\le1\),
\[
 1-\prod_eq_e\le\sum_e(1-q_e).
\]
Summing the one-coordinate estimates proves
\eqref{eq:V54-upper-gap}.  Nontriviality makes the left side
positive.
\end{proof}

Let \(N_{\boldsymbol T,\mu}\) denote a raw normalized quartic
numerator coefficient and let
\[
 P_{\boldsymbol T,\mu}
 :={\Xi(\boldsymbol T)\over1-q_{\alpha,\boldsymbol T}}\,
   N_{\boldsymbol T,\mu}
\]
denote the same coefficient with the archived payer, for
\(\boldsymbol T\ne\boldsymbol1\).  All dimension, height, trace-norm,
and multiplicity weights are included in these symbols.

\begin{theorem}[Nontrivial quartic numerator extraction]
\label{thm:V54-nontrivial-strip}
For four BFA-normalized input rows,
\begin{equation}
 \sum_{\substack{\boldsymbol T\ne\boldsymbol1\\\mu}}
 |N_{\boldsymbol T,\mu}|
 \le
 Cs_\alpha
 \sum_{\tau\in\mathfrak T_4}
 \prod_{ij\in E(\tau)}\theta_{ij}.
 \label{eq:V54-nontrivial-strip}
\end{equation}
\end{theorem}

\begin{proof}
Coefficientwise,
\[
 |N_{\boldsymbol T,\mu}|
 ={1-q_{\alpha,\boldsymbol T}\over\Xi(\boldsymbol T)}
  |P_{\boldsymbol T,\mu}|
 \le Cs_\alpha|P_{\boldsymbol T,\mu}|
\]
by \Cref{lem:V54-upper-gap}.  Sum over raw isotypic copies and apply
the restored quartic MTA in archive f3.  Because the multiplier is
inserted after physical output compression, no representation count
is introduced.
\end{proof}

\begin{firewall}[Stripping is not a second inverse]
\label{fire:V54-strip}
\Cref{thm:V54-nontrivial-strip} algebraically multiplies the proved
paid coefficient by
\((1-q_{\alpha,\boldsymbol T})/\Xi(\boldsymbol T)\).
Its conclusion contains no resolvent and no output mass.  Later use
of the quintic final payer therefore spends exactly one inverse.
\end{firewall}

\section{A private spectator lifts the scalar output}
\label{sec:V54-scalar}

Fix a fundamental representation \(R_*\) on one new coordinate
\(*\), disjoint from all four original supports.  Attach it only to
row \(1\).  Write \(a_*>0\) for its fixed balanced mass and
\(q_*:=q_{\alpha,R_*}\).

\begin{lemma}[Exact scalar-channel lift]
\label{lem:V54-scalar-lift}
The original quartic scalar-output coefficient is mapped injectively
to the augmented output which is trivial on every old coordinate and
equal to \(R_*\) on the new coordinate.  Its raw numerator is
\begin{equation}
 N_*^{\rm aug}=q_*N_{\boldsymbol1}^{\rm old}
 \label{eq:V54-scalar-lift}
\end{equation}
up to the fixed spectator normalization constant.
\end{lemma}

\begin{proof}
The new coordinate is private to the first row.  As in
\Cref{lem:V52-private-factor}, every term of the fourth
moment--cumulant formula contains row \(1\) exactly once, so
expectation over the new coordinate factors as \(q_*I\).
On that coordinate there is no fusion choice: the output is \(R_*\)
with multiplicity one.  On all old coordinates restrict the raw
pinching sum to the original scalar projection.  Orthogonality of
product-group isotypic outputs makes this map injective.  Choosing a
fixed trace-one spectator coefficient accounts for the harmless
normalization constant.
\end{proof}

\begin{lemma}[The lifted payer is one clock]
\label{lem:V54-lifted-payer}
For all sufficiently large \(\alpha\),
\begin{equation}
 q_*\ge {1\over2},
 \qquad
 {a_*\over1-q_*}\ge {c\over s_\alpha}.
 \label{eq:V54-lifted-payer}
\end{equation}
\end{lemma}

\begin{proof}
The fixed-representation Wilson expansion gives
\[
 q_*=1-c_*s_\alpha+O(s_\alpha^2)
\]
with \(c_*>0\).  The first assertion follows for large \(\alpha\);
the second follows from
\(1-q_*\le Cs_\alpha a_*\).
\end{proof}

\begin{lemma}[Augmented input normalization does not enlarge a tree]
\label{lem:V54-augmented-normalization}
Let \(\Xi_1'=\Xi_1+a_*\) and let
\(\theta_{1j}'=H_{1j}/(\Xi_1'\Xi_j)\); all other pair marks are
unchanged.  The augmented first-row BFA norm is at most
\[
 C_*{\Xi_1'\over\Xi_1}
\]
times its old normalized norm.  For every spanning tree
\(\tau\in\mathfrak T_4\),
\begin{equation}
 {\Xi_1'\over\Xi_1}
 \prod_{ij\in E(\tau)}\theta_{ij}'
 \le
 \prod_{ij\in E(\tau)}\theta_{ij}.
 \label{eq:V54-tree-normalization}
\end{equation}
\end{lemma}

\begin{proof}
The fixed spectator changes the height and dimension factors by one
constant \(C_*\) and replaces the row mass \(\Xi_1\) by
\(\Xi_1'\).  It creates no pair stock.  If
\(\deg_\tau(1)=d\ge1\), then
\[
 \prod_{ij\in E(\tau)}\theta_{ij}'
 =
 \left({\Xi_1\over\Xi_1'}\right)^d
 \prod_{ij\in E(\tau)}\theta_{ij}.
\]
Multiplication by \(\Xi_1'/\Xi_1\) leaves exponent \(d-1\ge0\),
proving \eqref{eq:V54-tree-normalization}.
\end{proof}

\begin{theorem}[Scalar quartic numerator extraction]
\label{thm:V54-scalar-strip}
For four BFA-normalized input rows, the complete scalar-output
quartic numerator obeys
\begin{equation}
 |N_{\boldsymbol1}|
 \le
 Cs_\alpha
 \sum_{\tau\in\mathfrak T_4}
 \prod_{ij\in E(\tau)}\theta_{ij}.
 \label{eq:V54-scalar-strip}
\end{equation}
The statement includes every scalar multiplicity block.
\end{theorem}

\begin{proof}
Apply the restored paid quartic MTA to the augmented four-row
interaction.  Restrict its nonnegative absolute-capacity sum to the
output described in \Cref{lem:V54-scalar-lift}.  On that output the
payer is \(a_* /(1-q_*)\ge c/s_\alpha\), while
\(q_*\ge1/2\).  Hence the paid estimate implies
\[
 |N_{\boldsymbol1}^{\rm old}|
 \le Cs_\alpha
 \left({\Xi_1'\over\Xi_1}\right)
 \sum_{\tau\in\mathfrak T_4}
 \prod_{ij\in E(\tau)}\theta_{ij}'.
\]
Use \Cref{lem:V54-augmented-normalization}.  Raw pinching on the old
scalar isotypic sector sums its physical multiplicity blocks before
the fixed spectator is attached, so no count is created.
\end{proof}

\begin{firewall}[The spectator proves an inequality only]
\label{fire:V54-spectator-scope}
The fresh coordinate is an auxiliary test used to bound the scalar
coefficient by an already-proved nontrivial-output theorem.  It is
not added to the physical lattice, the quintic word, or the continuum
theory.  After \Cref{thm:V54-scalar-strip} it is discarded.
\end{firewall}

\section{Projection-complete quartic numerator}
\label{sec:V54-complete}

\begin{theorem}[Unresolvented projection-complete quartic atlas]
\label{thm:V54-quartic-numerator}
The complete connected four-row numerator, with arbitrary finite
supports, representations, output channels, and physical
multiplicities, satisfies
\begin{equation}
 \boxed{
 \kappa_\otimes^{\rm abs}(\mathcal C_4^{\rm num})
 \le
 Cs_\alpha
 \sum_{\tau\in\mathfrak T_4}
 \prod_{ij\in E(\tau)}\theta_{ij}.}
 \label{eq:V54-quartic-numerator}
\end{equation}
The left side includes the scalar output and contains neither an
output mass nor a reduced resolvent.
\end{theorem}

\begin{proof}
Split the raw output pinching into
\(\boldsymbol T\ne\boldsymbol1\) and
\(\boldsymbol T=\boldsymbol1\).  Apply
\Cref{thm:V54-nontrivial-strip} to the first part and
\Cref{thm:V54-scalar-strip} to the second.  Their right sides are the
same sixteen-tree sum.  Adding them changes only the constant.
\end{proof}

\section{Closure of all \(4+1\) quintic sources}
\label{sec:V54-41}

\begin{theorem}[Complete \(4+1\) first-connectivity theorem]
\label{thm:V54-complete-41}
All five quintic first-connectivity source families whose predecessor
partition has type \(4+1\) satisfy the support- and
representation-uniform BFA marked-tree bound.
\end{theorem}

\begin{proof}
Fix a predecessor
\(\rho=\{\{1,2,3,4\},\{5\}\}\).  The complete connected prefix on
the four-row block satisfies
\eqref{eq:V54-quartic-numerator}, contributing
\[
 Cs_\alpha
 \sum_{\tau\in\mathfrak T_4}
 \prod_{ij\in E(\tau)}\theta_{ij}.
\]
The singleton prefix is a complete moment contraction and costs at
most one.  At the first-connectivity coordinate, the exact relative
partition identity of V50 assembles all joining letters into one
covariance of the four-row block-product with row \(5\).  Its final
output BFA mass and unique reduced resolvent cost at most
\(C/s_\alpha\), with raw output pinching.  This cancels the single
clock in \eqref{eq:V54-quartic-numerator}.

Every quartic tree \(\tau\) has three edges.  The joining covariance
retains an actual quotient edge from the four-row component to row
\(5\); adjoining it to \(\tau\) gives an ordinary five-row spanning
tree.  Thus the resulting sum is bounded by the global quintic tree
atlas.  Scalar quartic prefix outputs are included by
\Cref{thm:V54-scalar-strip}; payer coincidence and multiplicities are
included by the V50 weighted joining theorem.  Relabelling the
singleton block gives all five partitions of type \(4+1\).
\end{proof}

\begin{corollary}[Both two-block quintic predecessor types close]
\label{cor:V54-two-block}
Every quintic first-connectivity source whose predecessor partition
has two blocks is closed.
\end{corollary}

\begin{proof}
The two possible block-size types are \(4+1\) and \(3+2\).
Apply
\Cref{thm:V54-complete-41,thm:V53-complete-32}.
\end{proof}

\section{Next source type}
\label{sec:V54-next}

\begin{assessment}[Three predecessor blocks]
\label{ass:V54-next}
The next quotient partition has three blocks, beginning with
\(3+1+1\) and \(2+2+1\).  The prefix connected carriers are already
projection-complete by V51 and the binary covariance card.  What is
new is the three-block local joining sum: V50's two-block identity
reduces to one covariance, whereas three block-products assemble into
their complete third cumulant.  V55 must audit the companion notes,
then prove a reducible, scalar-inclusive, one-payer third-cumulant
joining card before inserting the known prefix numerators.
\end{assessment}

\section{V54 result ledger and V55 programme}
\label{sec:V54-results}

\begin{theorem}[Fourth-series V54 result ledger]
\label{thm:V54-results}
Relative to V1--V53, V54 proves:
\begin{enumerate}[label=\textup{(V54.\arabic*)},leftmargin=6.3em]
\item the support-uniform upper product-gap estimate;
\item coefficientwise removal of the archived quartic payer on every
      nontrivial output;
\item an injective fresh-coordinate lift of the scalar quartic output;
\item compensation of the augmented input normalization by the
      mandatory incidence of row \(1\) in every spanning tree;
\item the unresolvented, scalar-inclusive, projection-complete
      quartic numerator atlas with one Wilson clock;
\item closure of all five \(4+1\) first-connectivity sources; and
\item closure of every two-block quintic predecessor partition.
\end{enumerate}
\end{theorem}

\begin{proof}
These are
\Cref{lem:V54-upper-gap,
thm:V54-nontrivial-strip,
lem:V54-scalar-lift,
lem:V54-augmented-normalization,
thm:V54-scalar-strip,
thm:V54-quartic-numerator,
thm:V54-complete-41,
cor:V54-two-block}.
\end{proof}

\begin{programme}[V55: audit and three-block joining]
\label{prog:V54-V55}
\begin{enumerate}[leftmargin=3em]
\item Perform the compulsory read-only audit of the newest active
      SM--Einstein and Navier--Stokes notes.
\item Apply the relative product-cumulant identity of V50 to three
      predecessor block-products and retain their complete third
      cumulant.
\item Extend the projection-complete ternary numerator of V51 from
      irreducible row products to arbitrary reducible block-products
      by spectator amplification and raw pinching.
\item Keep the one final mass--resolvent payer and identify which
      joining clock cancels it.
\item Insert the \(3+1+1\) and \(2+2+1\) prefix carriers without
      spending an intermediate inverse.
\end{enumerate}
\end{programme}

\begin{firewall}[Claim boundary after fourth-series V54]
\label{fire:V54-claim-boundary}
V54 proves every two-block quintic first-connectivity source.  It
does not prove the three-, four-, or five-block predecessor sources,
the full quintic MTA, an all-order atlas, WUFC, CPM, cutoff removal,
reflection positivity, Osterwalder--Schrader reconstruction,
construction of the continuum Yang--Mills measure, or a uniform
positive continuum spectral gap.  The full Yang--Mills mass gap has
not been proved.
\end{firewall}

\paragraph{Parent-version record.}
V54 was formed from
\texttt{Renewal\_Yang\_Mills\_Mass\_Gap\_Research\_Notes\_V53.tex}.
The V53 parent had \(25\,695\) logical source lines,
\(895\,675\) bytes, and SHA--256
\[
 \texttt{3a5a3405e6b5cc680d64e2d86c30edaad0a413cea5a5a26a3206cb5a07948bf3}.
\]
The next compulsory companion audit is V55.

% =============================================================================
% END APPENDED FOURTH-SERIES V54 MATERIAL
% =============================================================================

% =============================================================================
% BEGIN APPENDED FOURTH-SERIES V55 MATERIAL
% =============================================================================

\chapter{V55: Companion audit and closure of three-block sources}
\label{chap:V55-three-block}

\section{Scheduled SM/NS companion audit}
\label{sec:V55-audit}

The compulsory read-only audit was performed before choosing the V55
proof.  The newest active companion files were:
\begin{center}
\begin{tabular}{p{0.47\textwidth}rr}
\toprule
file&logical lines&bytes\\
\midrule
\texttt{Renewal\_SM\_Einstein\_Continuum\_Research\_Notes\_V46.tex}
&\(17\,834\)&\(612\,623\)\\
\texttt{Renewal\_NS\_Stability\_Research\_Notes\_V31.tex}
&\(23\,052\)&\(887\,537\)\\
\bottomrule
\end{tabular}
\end{center}
Their SHA--256 digests are respectively
\[
\begin{split}
 &\texttt{05c5cf57eca46612504281eda54e4b37e3563fa48debff21d970308e00678bf0},\\
 &\texttt{f1121b62238ad5491bb7cf03635ea9934942edf573ed8faaa9ee79e1d38a6696}.
\end{split}
\]

SM V46 computes the exact connection between normalized Wick frames.
Constant frame changes form a group and telescope before a norm is
taken; norming every adjacent connection separately would spend a
radius gap at every shell.  It also proves that field degree and graph
provenance require different factorial grades.  The relevant YM
lesson is structural: assemble the three predecessor block-products
inside their exact relative cumulant before applying absolute values,
and do not confuse prefix-tree degree with the final-resolvent grade.

NS V31 is unchanged.  Its first-moment formation ladder is
multiplicity-free, but replacing the restricted signed formation
source by a positive global stock returns the continuation-class
quantity.  The relevant YM lesson remains to keep the complete signed
third cumulant and its joining marks until after the final payer.

\begin{firewall}[Companion-transfer boundary]
\label{fire:V55-audit}
Only the two assembly principles just stated are transferred.  The SM
Wick connection supplies no compact-group multiplier bound, and the
NS formation ladder supplies no Peter--Weyl pinching or Casimir
estimate.  Conversely, the YM theorem below proves no moving-frame
certificate and no Navier--Stokes formation correlation.  The next
compulsory audit is V60.
\end{firewall}

\section{A generic numerator extraction below the final payer}
\label{sec:V55-generic-strip}

V54 used the restored quartic paid atlas.  The same argument applies
at every already-proved lower row order.

\begin{theorem}[Scalar-inclusive payer-stripping principle]
\label{thm:V55-payer-stripping}
Let \(m\ge2\) be fixed.  Suppose the complete connected \(m\)-row
carrier satisfies a support- and representation-uniform once-paid BFA
marked-tree atlas
\begin{equation}
 \kappa_\otimes^{\rm abs}(\mathcal C_m^{\rm pay})
 \le C_m
 \sum_{\tau\in\mathfrak T_m}
 \prod_{ij\in E(\tau)}\theta_{ij}.
 \label{eq:V55-paid-hypothesis}
\end{equation}
Then its complete projection-inclusive numerator satisfies
\begin{equation}
 \boxed{
 \kappa_\otimes^{\rm abs}(\mathcal C_m^{\rm num})
 \le C_m's_\alpha
 \sum_{\tau\in\mathfrak T_m}
 \prod_{ij\in E(\tau)}\theta_{ij}.}
 \label{eq:V55-stripped-conclusion}
\end{equation}
The conclusion includes the scalar output and contains no reduced
resolvent.
\end{theorem}

\begin{proof}
On every nontrivial output, multiply the paid coefficient by
\[
 {1-q_{\alpha,\boldsymbol T}\over\Xi(\boldsymbol T)}
 \le Cs_\alpha
\]
from \Cref{lem:V54-upper-gap}, then sum the nonnegative raw pinching
masses.

For the scalar output, attach a fixed fundamental representation on
one fresh coordinate private to row \(1\).  The connected
\(m\)-cumulant factors by its fixed multiplier \(q_*\), the selected
augmented output has payer
\(a_* /(1-q_*)\ge c/s_\alpha\), and \(q_*\ge1/2\).
The first input norm grows by at most
\(C(\Xi_1+a_*)/\Xi_1\).  Every spanning tree on \(m\ge2\) is incident
to row \(1\), so the decrease of its incident normalized marks
absorbs this ratio exactly as in
\eqref{eq:V54-tree-normalization}.  Restrict the paid atlas to that
augmented output and remove the spectator after the estimate.  This
proves the scalar part.  Adding the two output sectors proves
\eqref{eq:V55-stripped-conclusion}.
\end{proof}

\begin{corollary}[Projection-complete binary, ternary, and quartic
numerators]
\label{cor:V55-lower-numerators}
The conclusion \eqref{eq:V55-stripped-conclusion} holds for
\(m=2,3,4\).
\end{corollary}

\begin{proof}
Archive f1 proves the complete binary fusion/covariance atlas and
\texttt{thm:V50-global-ternary-MTA}.  Archive f3 proves
\texttt{thm:V96-global-quartic-MTA}.  Apply
\Cref{thm:V55-payer-stripping}.  For \(m=4\) this recovers
\Cref{thm:V54-quartic-numerator}; for \(m=3\) it packages the direct
projection-complete construction of V51--V53 into a uniform
one-clock interface.
\end{proof}

\begin{firewall}[No circular use of a lower inverse]
\label{fire:V55-lower-inverse}
\Cref{thm:V55-payer-stripping} is a consequence of a completed
lower-order theorem, obtained by multiplying away its payer.  Its
output is a numerator inequality.  When used in a quintic prefix, no
lower-order inverse remains.
\end{firewall}

\section{The exact three-block joining carrier}
\label{sec:V55-joining}

Let
\(\rho=\{C_1,C_2,C_3\}\) be the predecessor partition immediately
below the first-connectivity coordinate \(r\), and set
\[
 Y_\nu:=\bigotimes_{i\in C_\nu}A_{i,r},
 \qquad \nu=1,2,3.
\]
Identity factors are inserted for rows inactive at \(r\).

\begin{theorem}[Exact three-block joining identity]
\label{thm:V55-exact-three-joining}
The complete local partition sum at \(r\) which joins the three
predecessor blocks is
\begin{equation}
 \boxed{
 \mathcal J_{\rho,r}
 =\kappa_3^\alpha(Y_1,Y_2,Y_3).}
 \label{eq:V55-exact-three-joining}
\end{equation}
It must not be replaced by a sequence of two pairwise covariances.
\end{theorem}

\begin{proof}
Apply the relative product-cumulant identity
\eqref{eq:V50-relative-cumulant} to the three blocks of \(\rho\).
The retained local partitions are exactly those satisfying
\(\pi\vee\rho=\widehat1\), and their coefficient-one sum is the
third cumulant of the three block-products.  A sequential covariance
expansion selects an intermediate bracketing and omits the remaining
M\"obius terms, so it is not equal to the complete joining sum.
\end{proof}

\begin{lemma}[Reducible block-products preserve the final ternary
atlas]
\label{lem:V55-reducible-joining}
The once-paid ternary marked-tree atlas applies to
\eqref{eq:V55-exact-three-joining} when each \(Y_\nu\) is a reducible
tensor-product representation and its input coefficient is an
arbitrary trace-class prefix output.  Its raw output sum has no
representation or multiplicity count.
\end{lemma}

\begin{proof}
Decompose each finite tensor-product carrier orthogonally into
irreducible isotypic blocks.  On every triple of input blocks apply
archive f1
\texttt{thm:V50-global-ternary-MTA}, whose constant is uniform in
the irreducible labels and physical multiplicities.  Perform the
input and output sums by complete raw pinching.  Orthogonal
projections and multiplicity partial traces are trace-norm
contractions; their raw masses sum to at most the input trace norm.
Thus no number of summands enters.  Equivalently, this is spectator
amplification of the V51 operator numerator followed by the one final
payer.
\end{proof}

\begin{lemma}[Row-level lift of a quotient tree]
\label{lem:V55-row-lift}
Choose an internal spanning tree in every nonsingleton predecessor
block and a two-edge tree on the three quotient vertices
\(\{C_1,C_2,C_3\}\).  Resolve each quotient edge to one actual row
pair crossing the corresponding blocks.  The union is a labelled
five-row spanning tree, and every normalized block-product mark
expands into a finite nonnegative sum of the corresponding row-pair
marks.
\end{lemma}

\begin{proof}
The graph assertion is archive f3
\texttt{lem:V97-tree-lift}: the internal trees use
\(\sum_\nu(|C_\nu|-1)=5-3=2\) edges, and the quotient tree uses two,
so their connected union has five vertices and four edges.

At the joining coordinate, the balanced mass of a block-product is
bounded, up to a fixed row-order constant, by the sum of the masses of
its constituent rows.  Hence a pair stock between two block-products
is bounded by
\[
 C\sum_{i\in C_\nu}\sum_{j\in C_\lambda}
 a_{i,r}a_{j,r}.
\]
After the original five input normalizations this is the corresponding
finite sum of row-pair marks.  There are at most five rows, so the
endpoint expansion changes only the arity constant.
\end{proof}

\begin{theorem}[Final three-block joining card]
\label{thm:V55-three-block-card}
With the exact prefix coefficients left inside their predecessor
blocks, the joining carrier
\eqref{eq:V55-exact-three-joining}, the complete moment suffix, the
final BFA output mass, and the unique reduced resolvent satisfy the
two-edge quotient-tree marked bound.  The estimate is stable under
scalar prefix outputs, payer coincidence, reducible block-products,
and all physical multiplicities.
\end{theorem}

\begin{proof}
Use
\Cref{thm:V55-exact-three-joining} before taking norms.  Apply the
once-paid ternary atlas through
\Cref{lem:V55-reducible-joining}.  Its two tree edges are the two
quotient edges.  Expand them into actual row endpoints by
\Cref{lem:V55-row-lift}.  The complete moment suffix is a contraction.
The ternary theorem already retains the output mass and the one
reduced resolvent, so no additional payer is inserted.  Scalar prefix
outputs are input coefficient blocks here, not final scalar outputs;
they are included by trace-class amplification and raw pinching.
\end{proof}

\section{Closure of \(3+1+1\) and \(2+2+1\)}
\label{sec:V55-closure}

\begin{theorem}[Complete \(3+1+1\) first-connectivity theorem]
\label{thm:V55-complete-311}
All ten quintic first-connectivity source families whose predecessor
partition has type \(3+1+1\) satisfy the uniform BFA marked-tree
bound.
\end{theorem}

\begin{proof}
The connected three-row predecessor block obeys the
projection-complete numerator atlas
\Cref{cor:V55-lower-numerators}, with one \(s_\alpha\) surplus and
an internal two-edge tree.  The two singleton prefixes are complete
moment contractions.  Apply the final three-block joining card
\Cref{thm:V55-three-block-card}; it supplies the two quotient-tree
edges and the sole final payer.  By
\Cref{lem:V55-row-lift}, the four edges form a five-row tree.
The surplus \(s_\alpha\le C\) may be discarded.  Relabelling the
three-row block gives the ten partitions.
\end{proof}

\begin{theorem}[Complete \(2+2+1\) first-connectivity theorem]
\label{thm:V55-complete-221}
All fifteen quintic first-connectivity source families whose
predecessor partition has type \(2+2+1\) satisfy the uniform BFA
marked-tree bound.
\end{theorem}

\begin{proof}
Apply \Cref{cor:V55-lower-numerators} to each of the two connected
binary predecessor blocks.  They contribute their two internal edges
and \(s_\alpha^2\).  The singleton prefix is contractive.  Apply
\Cref{thm:V55-three-block-card} for the two quotient edges and the one
final payer.  The tree lift gives an ordinary five-row tree, and the
two numerator clocks are surplus.  There are five choices of the
singleton and three pairings of the remaining labels, hence fifteen
families.
\end{proof}

\begin{corollary}[Every predecessor with at most three blocks closes]
\label{cor:V55-at-most-three}
Every quintic first-connectivity source whose predecessor partition
has two or three blocks satisfies the uniform marked-tree atlas.
\end{corollary}

\begin{proof}
Two-block sources are
\Cref{cor:V54-two-block}.  Three-block sources have types
\(3+1+1\) and \(2+2+1\), covered by
\Cref{thm:V55-complete-311,thm:V55-complete-221}.
\end{proof}

\section{Next source type}
\label{sec:V55-next}

\begin{assessment}[Four predecessor blocks]
\label{ass:V55-next}
The only nondiscrete predecessor type not yet addressed is
\(2+1+1+1\).  Its prefix has one binary connected block, for which
\Cref{cor:V55-lower-numerators} supplies a scalar-inclusive numerator
clock.  The joining coordinate is the complete fourth cumulant of
four block-products and carries the final payer.  V56 must prove that
the restored paid quartic atlas is stable under these reducible
block-products and that its three quotient edges expand to row-level
marks before raw pinching.
\end{assessment}

\section{V55 result ledger and V56 programme}
\label{sec:V55-results}

\begin{theorem}[Fourth-series V55 result ledger]
\label{thm:V55-results}
Relative to V1--V54, V55 proves:
\begin{enumerate}[label=\textup{(V55.\arabic*)},leftmargin=6.3em]
\item the compulsory read-only audit of SM V46 and NS V31;
\item a generic scalar-inclusive payer-stripping principle;
\item projection-complete one-clock numerator interfaces at row
      orders two, three, and four;
\item exact assembly of a three-block joining family into one third
      cumulant of block-products;
\item stability of the final paid ternary atlas under reducible
      block-products and raw multiplicity pinching;
\item exact lifting of internal and quotient trees to a five-row tree;
\item closure of all ten \(3+1+1\) and all fifteen \(2+2+1\)
      source families; and
\item closure of every quintic predecessor partition with at most
      three blocks.
\end{enumerate}
\end{theorem}

\begin{proof}
These are
\Cref{sec:V55-audit,
thm:V55-payer-stripping,
cor:V55-lower-numerators,
thm:V55-exact-three-joining,
lem:V55-reducible-joining,
lem:V55-row-lift,
thm:V55-three-block-card,
thm:V55-complete-311,
thm:V55-complete-221,
cor:V55-at-most-three}.
\end{proof}

\begin{programme}[V56: four-block joining]
\label{prog:V55-V56}
\begin{enumerate}[leftmargin=3em]
\item For
      \(\rho=\{\{a,b\},\{c\},\{d\},\{e\}\}\), assemble the complete
      joining partition family into
      \(\kappa_4(Y_{ab},Y_c,Y_d,Y_e)\).
\item Decompose the pair block-product into irreducibles only inside
      raw pinching; do not count its multiplicities.
\item Apply the restored paid quartic MTA as the final joining card,
      not as a prefix theorem.
\item Expand its three quotient-tree marks into actual row-pair
      endpoints and combine them with the binary prefix edge.
\item Check the scalar binary prefix output and every payer
      coincidence explicitly.
\end{enumerate}
\end{programme}

\begin{firewall}[Claim boundary after fourth-series V55]
\label{fire:V55-claim-boundary}
V55 proves every quintic source with a two- or three-block
predecessor partition.  It does not prove the \(2+1+1+1\) or
discrete predecessor sources, the full quintic MTA, an all-order
atlas, WUFC, CPM, cutoff removal, reflection positivity,
Osterwalder--Schrader reconstruction, construction of the continuum
Yang--Mills measure, or a uniform positive continuum spectral gap.
The full Yang--Mills mass gap has not been proved.
\end{firewall}

\paragraph{Parent-version record.}
V55 was formed from
\texttt{Renewal\_Yang\_Mills\_Mass\_Gap\_Research\_Notes\_V54.tex}.
The V54 parent had \(26\,096\) logical source lines,
\(909\,702\) bytes, and SHA--256
\[
 \texttt{c4bc9a826cee29c6af51bef168d2905f3b4e314ac23fb34a0b81f87ffac38932}.
\]
The next compulsory companion audit is V60.

% =============================================================================
% END APPENDED FOURTH-SERIES V55 MATERIAL
% =============================================================================

% =============================================================================
% BEGIN APPENDED FOURTH-SERIES V56 MATERIAL
% =============================================================================

\chapter{V56: The paid quartic joining card closes \(2+1+1+1\)}
\label{chap:V56-four-block}

\section{Exact four-block assembly}
\label{sec:V56-assembly}

Fix the predecessor partition
\[
 \rho=\{\{a,b\},\{c\},\{d\},\{e\}\}
\]
immediately below the first-connectivity coordinate \(r\).  Put
\[
 Y_{ab}:=A_{a,r}\otimes A_{b,r},
 \qquad
 Y_c:=A_{c,r},\quad Y_d:=A_{d,r},\quad Y_e:=A_{e,r},
\]
with identities on inactive row legs.

\begin{theorem}[Exact four-block joining identity]
\label{thm:V56-exact-four-joining}
The complete local partition family at \(r\) which joins the four
predecessor blocks is
\begin{equation}
 \boxed{
 \mathcal J_{\rho,r}
 =\kappa_4^\alpha(Y_{ab},Y_c,Y_d,Y_e).}
 \label{eq:V56-exact-four-joining}
\end{equation}
The equality includes every local block collision and every partition
which joins several predecessor blocks at once.
\end{theorem}

\begin{proof}
Apply the relative product-cumulant identity
\eqref{eq:V50-relative-cumulant} to the four blocks of \(\rho\).
Its left side is the coefficient-one sum over all local partitions
\(\pi\) satisfying
\(\pi\vee\rho=\widehat1\).  Its right side is the fourth cumulant of
the four block-products, which is
\eqref{eq:V56-exact-four-joining}.  No choice of a local spanning
tree is made before the signed cumulant has been assembled.
\end{proof}

\begin{firewall}[No iterated lower cumulants]
\label{fire:V56-no-iteration}
An iteration of three covariances depends on a bracketing and omits
fourth-cumulant M\"obius terms.  Likewise, a covariance of one
three-block cumulant with the fourth block spends an intermediate
inverse if the paid ternary theorem is used.  V56 uses only the
single complete fourth cumulant in
\eqref{eq:V56-exact-four-joining}.
\end{firewall}

\section{Reducible stability of the final quartic atlas}
\label{sec:V56-reducible}

\begin{lemma}[Paid quartic atlas under block-product amplification]
\label{lem:V56-reducible-quartic}
Archive f3
\texttt{thm:V96-global-quartic-MTA} remains valid when any of its four
input representations is a finite reducible tensor-product
representation and its coefficient is an arbitrary trace-class
prefix output.  The same constant is valid, up to a fixed row-order
factor, and the raw output sum has no multiplicity count.
\end{lemma}

\begin{proof}
Resolve each block-product carrier into orthogonal irreducible
isotypic summands.  Apply the restored quartic theorem on every tuple
of irreducible input blocks.  That theorem is uniform in
representations, physical multiplicities, payer placement, and
output-duality orientation.  Sum the input and output compressions by
raw pinching.  Orthogonal compression and multiplicity partial trace
are trace-norm contractions, so the total raw mass is at most the
trace norm of the original coefficient.  Only the finite number of
ways to assign at most five physical row legs to four block carriers
changes the constant.
\end{proof}

\begin{lemma}[Expansion of quartic quotient marks]
\label{lem:V56-quotient-marks}
Every three-edge quotient-tree monomial produced by
\Cref{lem:V56-reducible-quartic} is bounded by a fixed finite sum of
monomials obtained by selecting one actual row-pair endpoint across
each quotient edge.
\end{lemma}

\begin{proof}
For predecessor blocks \(C,D\), the coordinate mass of their
block-product pair is bounded by
\[
 C\left(\sum_{i\in C}a_{i,r}\right)
  \left(\sum_{j\in D}a_{j,r}\right)
 =
 C\sum_{i\in C}\sum_{j\in D}a_{i,r}a_{j,r}.
\]
Expand this nonnegative sum for each of the three quotient edges.
There are at most \(5^6\) endpoint choices, a fixed arity constant.
After the five original input mass normalizations, each factor is the
corresponding row-pair mark at \(r\).  No representation
multiplicity is part of this endpoint expansion.
\end{proof}

\begin{theorem}[Final four-block joining card]
\label{thm:V56-four-block-card}
The complete joining carrier
\eqref{eq:V56-exact-four-joining}, unrestricted moment suffix, final
BFA output mass, and unique reduced resolvent satisfy a uniform
three-edge quotient-tree marked bound.  This includes reducible
block-products, scalar prefix outputs, all payer coincidences, and
all physical multiplicities.
\end{theorem}

\begin{proof}
Apply
\Cref{thm:V56-exact-four-joining,lem:V56-reducible-quartic}.
The restored quartic theorem carries exactly one payer, placed at the
actual final output.  Expand its three quotient marks by
\Cref{lem:V56-quotient-marks}.  The suffix moment is a contraction,
and all remaining outputs sum by raw pinching.  A scalar prefix
coefficient is simply an allowed trace-class input block; the final
quartic output is centered by its own \(Q_0\), as required for the
unique reduced resolvent.
\end{proof}

\section{Closure of the ten four-block sources}
\label{sec:V56-closure}

\begin{lemma}[Binary prefix numerator, including its scalar output]
\label{lem:V56-binary-prefix}
The connected prefix on \(\{a,b\}\), before any final payer, obeys
\begin{equation}
 \kappa_\otimes^{\rm abs}(\mathcal C_{ab}^{\rm num})
 \le Cs_\alpha\theta_{ab},
 \label{eq:V56-binary-prefix}
\end{equation}
including both scalar and nontrivial prefix outputs.
\end{lemma}

\begin{proof}
This is the \(m=2\) case of
\Cref{cor:V55-lower-numerators}.  Its scalar sector is supplied by
the private-spectator lift in
\Cref{thm:V55-payer-stripping}, so no \(Q_0\) is inserted in the
prefix.
\end{proof}

\begin{theorem}[Complete \(2+1+1+1\) first-connectivity theorem]
\label{thm:V56-complete-2111}
All ten quintic first-connectivity source families whose predecessor
partition has type \(2+1+1+1\) satisfy the support- and
representation-uniform BFA marked-tree bound.
\end{theorem}

\begin{proof}
The binary predecessor block contributes its internal edge and the
one-clock numerator bound
\eqref{eq:V56-binary-prefix}.  The three singleton prefixes are
complete moment contractions.  Apply
\Cref{thm:V56-four-block-card}; it supplies the three quotient-tree
edges together with the sole final payer.  The internal edge and
three quotient edges form a five-row spanning tree by archive f3
\texttt{lem:V97-tree-lift}.  All coefficient sums are raw pinching
sums.  The remaining factor \(s_\alpha\) is surplus and may be
discarded for large \(\alpha\).  Choosing the two-row block gives
\(\binom52=10\) source families.
\end{proof}

\begin{corollary}[Every nondiscrete quintic predecessor closes]
\label{cor:V56-nondiscrete}
Every quintic first-connectivity source whose predecessor partition
is not the discrete partition satisfies the uniform BFA marked-tree
bound.
\end{corollary}

\begin{proof}
The two-block types \(4+1\) and \(3+2\) are
\Cref{cor:V54-two-block}.  The three-block types \(3+1+1\) and
\(2+2+1\) are \Cref{cor:V55-at-most-three}.  The sole four-block type
is \Cref{thm:V56-complete-2111}.  These are all five nondiscrete
types in the exact V97 census.
\end{proof}

\section{The unique remaining quintic source}
\label{sec:V56-discrete}

\begin{theorem}[Discrete-prefix identity]
\label{thm:V56-discrete-identity}
For the discrete predecessor
\(\rho=\{\{1\},\{2\},\{3\},\{4\},\{5\}\}\), the complete source is
\begin{equation}
 \mathcal S_{\widehat0}^{(5)}
 =
 \sum_r
 \left(\bigotimes_{i=1}^5M_{\{i\}}^{<r}\right)
 \otimes K_{5,r}\otimes M_{>r},
 \label{eq:V56-discrete-identity}
\end{equation}
where \(K_{5,r}\) is the complete signed one-coordinate fifth
cumulant.  Every word in the discrete source occurs exactly once.
\end{theorem}

\begin{proof}
Before the first-connectivity coordinate no pair of rows belongs to
one cumulative block, so the prefix factors into five complete
singleton moments.  At \(r\), the condition
\(\widehat0\vee\pi_r=\widehat1\) forces
\(\pi_r=\widehat1\), whose connected M\"obius coefficient is the
complete fifth cumulant.  The suffix is unrestricted.  Uniqueness of
the first-connectivity coordinate gives coefficient one and
disjointness.
\end{proof}

\begin{assessment}[Remaining analytic interface]
\label{ass:V56-discrete}
Archive f1
\texttt{thm:V56-uniform-one-link-quintic} gives a
representation-uniform one-coordinate fifth-cumulant tree-valence
decomposition.  Archive f3 V37--V41 tests its sparse transition
packets in the full BFA normalization.  What must still be checked is
that the local card already contains the final output mass and unique
resolvent for every payer placement, and that its coordinate-dependent
four-mark right side sums diagonally to the global \(125\)-tree atlas.
No lower-prefix theorem is missing in
\eqref{eq:V56-discrete-identity}.
\end{assessment}

\section{V56 result ledger and V57 programme}
\label{sec:V56-results}

\begin{theorem}[Fourth-series V56 result ledger]
\label{thm:V56-results}
Relative to V1--V55, V56 proves:
\begin{enumerate}[label=\textup{(V56.\arabic*)},leftmargin=6.3em]
\item exact assembly of the four-block joining family into one
      fourth cumulant of block-products;
\item stability of the restored paid quartic atlas under reducible
      block-product amplification and raw pinching;
\item expansion of its three quotient-tree marks into actual row
      endpoints;
\item a complete final four-block joining card with one payer;
\item closure of all ten \(2+1+1+1\) source families;
\item closure of every nondiscrete quintic predecessor partition; and
\item reduction of the full quintic MTA to the unique discrete source
      \eqref{eq:V56-discrete-identity}.
\end{enumerate}
\end{theorem}

\begin{proof}
These are
\Cref{thm:V56-exact-four-joining,
lem:V56-reducible-quartic,
lem:V56-quotient-marks,
thm:V56-four-block-card,
lem:V56-binary-prefix,
thm:V56-complete-2111,
cor:V56-nondiscrete,
thm:V56-discrete-identity}.
\end{proof}

\begin{programme}[V57: paid discrete fifth-cumulant source]
\label{prog:V56-V57}
\begin{enumerate}[leftmargin=3em]
\item Read the exact hypotheses and output weights of archive f1
      \texttt{thm:V56-uniform-one-link-quintic}.
\item Compare them term by term with the selected-junction interface
      used for the quartic discrete source in archive f3 V96.
\item Retain all \(125\) local labelled trees and prove the diagonal
      four-mark coordinate sum.
\item Verify final \(Q_0\), output mass, unique resolvent, dimension
      cancellation, and every payer coincidence.
\item If the archived local theorem is missing any one of those
      factors, isolate that factor rather than declaring the
      discrete source closed.
\end{enumerate}
\end{programme}

\begin{firewall}[Claim boundary after fourth-series V56]
\label{fire:V56-claim-boundary}
V56 proves every nondiscrete quintic first-connectivity source.  It
does not yet prove the discrete fifth-cumulant source, the full
quintic MTA, an all-order atlas, WUFC, CPM, cutoff removal, reflection
positivity, Osterwalder--Schrader reconstruction, construction of the
continuum Yang--Mills measure, or a uniform positive continuum
spectral gap.  The full Yang--Mills mass gap has not been proved.
\end{firewall}

\paragraph{Parent-version record.}
V56 was formed from
\texttt{Renewal\_Yang\_Mills\_Mass\_Gap\_Research\_Notes\_V55.tex}.
The V55 parent had \(26\,495\) logical source lines,
\(925\,184\) bytes, and SHA--256
\[
 \texttt{5553981b6623e33c6fcc669a247b402395647030aec48125f60554053d383f48}.
\]
The next compulsory companion audit is V60.

% =============================================================================
% END APPENDED FOURTH-SERIES V56 MATERIAL
% =============================================================================

% =============================================================================
% BEGIN APPENDED FOURTH-SERIES V57 MATERIAL
% =============================================================================

\chapter{V57: The exact paid gate for the discrete quintic source}
\label{chap:V57-discrete-gate}

\section{Archive correction and scope}
\label{sec:V57-archive}

The theorem cited at the end of V56 is located in frozen archive f2,
not f1:
\[
\texttt{Renewal\_Yang\_Mills\_Mass\_Gap\_Research\_Notes\_V100\_f2.tex},
\qquad
\texttt{thm:V56-uniform-one-link-quintic}.
\]
It decomposes the complete one-coordinate fifth cumulant as
\begin{equation}
 \mathbf K_5^\alpha
 =\sum_{\tau\in\mathfrak T_5}\mathbf K_{5,\tau}^\alpha,
 \qquad
 \|\mathbf K_{5,\tau}^\alpha\|_{\rm op}
 \le C_5s_\alpha^4
 \prod_{i=1}^5A_{i,r}^{\deg_\tau(i)/2},
 \label{eq:V57-local-card}
\end{equation}
uniformly in the five irreducible inputs and every output
compression.  It does not contain
\[
 {\Xi(\boldsymbol T)\over1-q_{\alpha,\boldsymbol T}}.
\]
Thus it is a local operator numerator card, not a paid selected-source
theorem.

\begin{firewall}[Correction of the V56 assessment]
\label{fire:V57-location}
The mathematical statement quoted in
\Cref{ass:V56-discrete} is accurate, but its archive attribution is
corrected from f1 to f2.  More importantly, output compression in
\eqref{eq:V57-local-card} is contractive and multiplicity-free, but
it does not by itself pay the output mass or the unique reduced
resolvent.
\end{firewall}

\section{Exact coefficient of a selected discrete fibre}
\label{sec:V57-fibre}

Fix \(r\), a labelled tree \(\tau\in\mathfrak T_5\), and write
\[
 d_i:=\deg_\tau(i),\qquad
 a_i:=A_{i,r},\qquad
 P_i:=\Xi_i-a_i,\qquad
 Y:=\sum_iP_i.
\]
In the discrete predecessor source, every coordinate strictly below
\(r\) is active on at most one row.  Its complete prefix multiplier is
\[
 q_P:=\prod_{i=1}^5
 \prod_{\substack{e<r\\e\ {\rm active\ on}\ i}}
 q_{\alpha,R_{i,e}}.
\]
There is no nontrivial connected prefix block.

\begin{lemma}[Exact private factorization of the discrete fibre]
\label{lem:V57-private-factor}
Before final output resolution, the tree fibre at \(r\) has the form
\begin{equation}
 q_P\,\mathbf K_{5,\tau,r}^\alpha\otimes\mathbf M_{>r},
 \qquad
 \|\mathbf M_{>r}\|_{\rm op}\le1.
 \label{eq:V57-private-factor}
\end{equation}
\end{lemma}

\begin{proof}
This is \Cref{thm:V56-discrete-identity} refined by the tree
decomposition \eqref{eq:V57-local-card}.  Below \(r\), discreteness
forces every active coordinate to be private, so its moment is the
central singleton multiplier.  Above \(r\), the unrestricted complete
moment is an average of tensor-product unitaries and is contractive.
Coordinate independence proves the tensor factorization.
\end{proof}

\begin{definition}[Paid selected quintic junction gate]
\label{def:V57-PGSEL5}
For a nontrivial final output \(\boldsymbol T\), define
\begin{align}
 \mathcal P_{\tau,r}(\boldsymbol T)
 &:=
 {\Xi(\boldsymbol T)\over1-q_{\alpha,\boldsymbol T}}\,
 q_P\,
 s_\alpha^4
 \prod_{i=1}^5a_i^{d_i/2},
 \label{eq:V57-paid-scalar}\\
 \mathcal T_{\tau,r}
 &:=
 \prod_{i=1}^5a_i^{d_i}\Xi_i^{1-d_i}
 =
 \left(\prod_i\Xi_i\right)
 \prod_{ij\in E(\tau)}
 {a_i\over\Xi_i}{a_j\over\Xi_j}.
 \label{eq:V57-target-scalar}
\end{align}
The gate \(\operatorname{PGSEL}_5\) is the coefficientwise inequality
\begin{equation}
 \boxed{
 \mathcal P_{\tau,r}(\boldsymbol T)
 \le C\mathcal T_{\tau,r}}
 \label{eq:V57-PGSEL5}
\end{equation}
after exact Fourier dimension cancellation and raw pinching, uniformly
in \(r,\tau,\boldsymbol T\), input representations, supports,
multiplicities, and payer placement.
\end{definition}

\begin{lemma}[PGSEL5 is exactly the missing analytic factor]
\label{lem:V57-exact-missing}
If \(\operatorname{PGSEL}_5\) holds, the complete discrete source
\eqref{eq:V56-discrete-identity} satisfies the quintic marked-tree
atlas.  Conversely, the local operator card
\eqref{eq:V57-local-card} and contractive suffix reduce the selected
fibre estimate precisely to \eqref{eq:V57-PGSEL5}; no further support
or coordinate count remains.
\end{lemma}

\begin{proof}
Normalize each of the five pure input blocks in BFA.  Exact Fourier
normalization cancels the output dimension against the factor
\(d_{\rm in}/d_{\boldsymbol T}\), and raw pinching sums all output
copies with total mass at most one.  By
\Cref{lem:V57-private-factor} and
\eqref{eq:V57-local-card}, the remaining scalar coefficient is bounded
by \(\mathcal P_{\tau,r}(\boldsymbol T)\), divided by
\(\prod_i\Xi_i\).  The right side
\(\mathcal T_{\tau,r}/\prod_i\Xi_i\) is exactly the local normalized
tree monomial.  Thus
\eqref{eq:V57-PGSEL5} is sufficient and is also the sole scalar
comparison left by the proved operator and pinching steps.
\end{proof}

\section{Coordinate aggregation is not the obstruction}
\label{sec:V57-summation}

Define the local row-pair mark
\[
 \vartheta_{ij,r}:=
 {a_{i,r}a_{j,r}\over\Xi_i\Xi_j}
\]
when all five rows are active at \(r\), and zero otherwise.

\begin{lemma}[Diagonal four-mark summation]
\label{lem:V57-diagonal-sum}
For every labelled five-row tree \(\tau\),
\begin{equation}
 \sum_r\prod_{ij\in E(\tau)}\vartheta_{ij,r}
 \le
 \prod_{ij\in E(\tau)}\theta_{ij}.
 \label{eq:V57-diagonal-sum}
\end{equation}
Consequently,
\[
 \sum_r\sum_{\tau\in\mathfrak T_5}
 \prod_{ij\in E(\tau)}\vartheta_{ij,r}
 \le
 \sum_{\tau\in\mathfrak T_5}
 \prod_{ij\in E(\tau)}\theta_{ij}.
\]
\end{lemma}

\begin{proof}
Fix the four edges \(f_1,\ldots,f_4\) of \(\tau\).  Nonnegativity
gives
\[
 \sum_r\prod_{\ell=1}^4\vartheta_{f_\ell,r}
 \le
 \prod_{\ell=1}^4\left(\sum_{r_\ell}
 \vartheta_{f_\ell,r_\ell}\right)
 \le
 \prod_{\ell=1}^4\theta_{f_\ell}.
\]
The last inequality holds because the global pair stock sums over all
coordinates, not only fivefold-active ones.  Sum over the fixed
\(5^{5-2}=125\) labelled trees.
\end{proof}

\begin{theorem}[Conditional discrete-source closure]
\label{thm:V57-conditional-discrete}
If \(\operatorname{PGSEL}_5\) holds for every selected output and
payer placement, then the complete discrete quintic source satisfies
the global \(125\)-tree MTA, uniformly in support, representations,
volume, and multiplicities.
\end{theorem}

\begin{proof}
Apply \Cref{lem:V57-exact-missing} to every selected fibre.  Sum first
over raw outputs, then over \(r\), and use
\Cref{lem:V57-diagonal-sum}.  The tree and payer-placement sets are
finite at order five, so only a fixed arity constant remains.
\end{proof}

\section{What the frozen private compiler already proves}
\label{sec:V57-private}

\begin{record}[Archived private-dominant branch]
\label{rec:V57-private}
Frozen f2
\texttt{thm:V54-private-scalar-compiler}, combined with
\texttt{thm:V56-uniform-one-link-quintic}, proves
\eqref{eq:V57-PGSEL5} whenever, for a fixed \(D\),
\begin{equation}
 \Xi(\boldsymbol T)\le DY,
 \qquad
 \prod_{i=1}^5\Xi_i^{d_i-1}\le DY^3.
 \label{eq:V57-private-dominance}
\end{equation}
It uses private Wilson product moments of orders three and four and
spends the final resolvent exactly once.
\end{record}

\begin{proof}
This is the \(m=5\) specialization of the cited compiler.  Its local
tree weight is exactly the right side of
\eqref{eq:V57-local-card}; its conclusion is exactly
\(\mathcal P_{\tau,r}\le C\mathcal T_{\tau,r}\).
\end{proof}

\begin{corollary}[Exact complement of the archived closure]
\label{cor:V57-complement}
Any discrete fibre not covered by
\Cref{rec:V57-private} lies in at least one of the two branches
\begin{align}
 \Xi(\boldsymbol T)&>DY,
 \label{eq:V57-output-dominant}\\
 \prod_{i=1}^5\Xi_i^{d_i-1}&>DY^3.
 \label{eq:V57-branching-dominant}
\end{align}
\end{corollary}

\begin{proof}
This is the logical negation of the conjunction in
\eqref{eq:V57-private-dominance}.
\end{proof}

\begin{firewall}[Why the local card cannot simply absorb the payer]
\label{fire:V57-no-premature}
The local estimate \eqref{eq:V57-local-card} is uniform under output
compression, but a contraction bound does not control the growing
factor
\(\Xi(\boldsymbol T)/(1-q_{\alpha,\boldsymbol T})\).
Nor does the diagonal summation
\Cref{lem:V57-diagonal-sum} repair a failed coefficientwise inequality.
Therefore neither theorem proves the discrete source outside
\eqref{eq:V57-private-dominance}.
\end{firewall}

\section{Regression against the active sparse packet}
\label{sec:V57-regression}

\begin{assessment}[Consistency with V37--V41]
\label{ass:V57-sparse}
The sparse transition packet of V37 has a hot payer and three fixed
prefix covariances.  V41 closes it only after exact Fourier dimension
cancellation, raw pinching, and the full mass--resolvent ledger,
obtaining an extra inverse root scale.  This is consistent with
\(\operatorname{PGSEL}_5\) but does not prove it: the discrete source
has no prefix covariance clocks, and its difficult outputs may lie in
\eqref{eq:V57-output-dominant} or
\eqref{eq:V57-branching-dominant}.
\end{assessment}

\section{V57 result ledger and V58 programme}
\label{sec:V57-results}

\begin{theorem}[Fourth-series V57 result ledger]
\label{thm:V57-results}
Relative to V1--V56, V57 proves:
\begin{enumerate}[label=\textup{(V57.\arabic*)},leftmargin=6.3em]
\item correction of the one-link quintic theorem's archive location
      and exact scope;
\item exact private factorization of every selected discrete fibre;
\item formulation of the paid selected quintic junction gate
      \(\operatorname{PGSEL}_5\);
\item proof that this gate is the sole remaining coefficientwise
      analytic factor after local operator control and raw pinching;
\item support-uniform diagonal summation of all four local tree marks;
\item conditional closure of the global discrete source;
\item identification of the already-proved private-dominant sector;
      and
\item reduction of its complement to the output-dominant and
      branching-dominant inequalities
      \eqref{eq:V57-output-dominant}--%
      \eqref{eq:V57-branching-dominant}.
\end{enumerate}
\end{theorem}

\begin{proof}
These are
\Cref{fire:V57-location,
lem:V57-private-factor,
def:V57-PGSEL5,
lem:V57-exact-missing,
lem:V57-diagonal-sum,
thm:V57-conditional-discrete,
rec:V57-private,
cor:V57-complement}.
\end{proof}

\begin{programme}[V58: output- and branching-dominant PGSEL5]
\label{prog:V57-V58}
\begin{enumerate}[leftmargin=3em]
\item On \eqref{eq:V57-output-dominant}, use the high-output product
      gap before bounding the output mass; test whether reverse fusion
      forces a hot input retained by the tree valence.
\item On \eqref{eq:V57-branching-dominant}, choose a maximal-degree
      tree center and compare its local mass \(a_i\) with its private
      mass \(P_i\).
\item Split local-dominant and private-dominant centers without
      comparing an input multiplier to a fused output multiplier.
\item Retain the hot-row displacement bound
      \texttt{lem:V56-hot-cumulant} when the local center is cofinal;
      do not replace it by the weaker polynomial star card.
\item Test the resulting inequalities on the equal-spin dual-singlet
      output and the V37 transition packet.
\end{enumerate}
\end{programme}

\begin{firewall}[Claim boundary after fourth-series V57]
\label{fire:V57-claim-boundary}
V57 reduces the last quintic source to the exact paid gate
\(\operatorname{PGSEL}_5\), proves its coordinate summation, and
records its private-dominant closure.  It does not prove PGSEL5 on the
output- or branching-dominant complement, the discrete source, the
full quintic MTA, an all-order atlas, WUFC, CPM, cutoff removal,
reflection positivity, Osterwalder--Schrader reconstruction,
construction of the continuum Yang--Mills measure, or a uniform
positive continuum spectral gap.  The full Yang--Mills mass gap has
not been proved.
\end{firewall}

\paragraph{Parent-version record.}
V57 was formed from
\texttt{Renewal\_Yang\_Mills\_Mass\_Gap\_Research\_Notes\_V56.tex}.
The V56 parent had \(26\,800\) logical source lines,
\(936\,843\) bytes, and SHA--256
\[
 \texttt{183a509ac967230360eb9d0ce000a73684fc1338e5cccb849f0714e55b6f0f21}.
\]
The next compulsory companion audit is V60.

% =============================================================================
% END APPENDED FOURTH-SERIES V57 MATERIAL
% =============================================================================

% =============================================================================
% BEGIN APPENDED FOURTH-SERIES V58 MATERIAL
% =============================================================================

\chapter{V58: The output-dominant discrete quintic branch}
\label{chap:V58-output}

\section{Root and private masses}
\label{sec:V58-masses}

Keep the selected discrete fibre of V57 and put
\[
 A:=\sum_{i=1}^5a_i,\qquad
 M:=\max_i a_i,\qquad
 Y:=\sum_iP_i.
\]
Choose \(i_*\) with \(a_{i_*}=M\), and let
\(\tau_*\) be the labelled star centered at \(i_*\).

\begin{lemma}[Final output mass is formed from root and private mass]
\label{lem:V58-output-formation}
Every resolved final output of the selected discrete fibre satisfies
\begin{equation}
 \Xi(\boldsymbol T)\le C_F(A+Y)
 \label{eq:V58-output-formation}
\end{equation}
for a fixed fusion constant \(C_F\).
\end{lemma}

\begin{proof}
At a prefix-private coordinate the output representation equals the
single active input representation, so its balanced mass contributes
exactly to \(Y\).  At \(r\), the output is contained in the tensor
product of the five root representations.  The fixed-order Casimir
fusion inequality bounds its balanced mass by
\(C_F\sum_i a_i=C_FA\).  The complete moment suffix does not alter
the Fourier representation.  Summing the coordinate masses proves
\eqref{eq:V58-output-formation}.
\end{proof}

\begin{lemma}[Output dominance forces root dominance]
\label{lem:V58-root-dominance}
Fix \(D\ge2C_F\).  If
\begin{equation}
 \Xi(\boldsymbol T)>DY,
 \label{eq:V58-output-dominant}
\end{equation}
then
\begin{equation}
 \Xi(\boldsymbol T)\le2C_FA,\qquad
 Y\le {2C_F\over D}A,\qquad
 \Xi_{i_*}\le C_D M.
 \label{eq:V58-root-dominance}
\end{equation}
\end{lemma}

\begin{proof}
By \eqref{eq:V58-output-formation} and
\eqref{eq:V58-output-dominant},
\[
 \Xi(\boldsymbol T)
 \le C_FA+{C_F\over D}\Xi(\boldsymbol T).
\]
Since \(C_F/D\le1/2\), absorb the last term to obtain the first
inequality.  The output-dominance hypothesis then gives the second.
Finally,
\[
 \Xi_{i_*}=M+P_{i_*}\le M+Y
 \le M+{2C_F\over D}A\le(1+10C_F/D)M,
\]
because \(A\le5M\).
\end{proof}

\section{A star target with no private denominator on the leaves}
\label{sec:V58-star}

\begin{lemma}[Largest-root star floor]
\label{lem:V58-star-floor}
Under \eqref{eq:V58-output-dominant}, the scalar tree target of
\(\tau_*\) satisfies
\begin{equation}
 \mathcal T_{\tau_*,r}
 ={M^4\over\Xi_{i_*}^3}
  \prod_{j\ne i_*}a_j
 \ge c_D M\prod_{j\ne i_*}a_j.
 \label{eq:V58-star-floor}
\end{equation}
\end{lemma}

\begin{proof}
The center of a five-row star has degree four and every leaf has
degree one.  Substitute these degrees into
\eqref{eq:V57-target-scalar}.  The leaf factors
\(\Xi_j^{1-d_j}\) equal one.  Use
\(\Xi_{i_*}\le C_DM\) from
\Cref{lem:V58-root-dominance}.
\end{proof}

\begin{lemma}[Centered displacement bound at the chosen star]
\label{lem:V58-star-displacement}
The complete one-coordinate fifth cumulant may be assigned to
\(\tau_*\) and obeys
\begin{equation}
 \|\mathbf K_{5,\tau_*}^\alpha\|_{\rm op}
 \le C\prod_{j\ne i_*}c_j,
 \qquad
 c_j:=\min\{1,\sqrt{s_\alpha a_j}\},
 \label{eq:V58-star-displacement}
\end{equation}
with all other tree fibres set to zero.
\end{lemma}

\begin{proof}
Frozen f2
\texttt{lem:V56-hot-cumulant} proves
\[
 \|\mathbf K_5^\alpha\|_{\rm op}
 \le C\prod_{j\ne i_*}c_j
\]
for every chosen center, not only for a thermally hot center.  Define
\(\mathbf K_{5,\tau_*}^\alpha:=\mathbf K_5^\alpha\) and the other
fibres to be zero.  This is an exact tree decomposition and retains
the complete signed cumulant before taking its norm.
\end{proof}

\section{Both output-gap regimes}
\label{sec:V58-gap}

\begin{theorem}[PGSEL5 on the output-dominant branch]
\label{thm:V58-output-PGSEL5}
For \(D\ge2C_F\), every selected discrete fibre satisfying
\eqref{eq:V58-output-dominant} obeys
\(\operatorname{PGSEL}_5\), uniformly in all supports,
representations, final outputs, multiplicities, and payer placements.
\end{theorem}

\begin{proof}
Assign the complete cumulant to \(\tau_*\) as in
\Cref{lem:V58-star-displacement}.  The private multiplier obeys
\(q_P\le1\).

On the low-output branch
\(s_\alpha\Xi(\boldsymbol T)\le c_0\), the product-gap split gives
\[
 {\Xi(\boldsymbol T)\over1-q_{\alpha,\boldsymbol T}}
 \le {C\over s_\alpha}.
\]
Using \(c_j\le\sqrt{s_\alpha a_j}\) on all four leaves,
\begin{align*}
 {\Xi(\boldsymbol T)\over1-q_{\alpha,\boldsymbol T}}\,
 q_P\|\mathbf K_5^\alpha\|_{\rm op}
 &\le {C\over s_\alpha}
       \prod_{j\ne i_*}\sqrt{s_\alpha a_j}\\
 &=Cs_\alpha\prod_{j\ne i_*}a_j^{1/2}\\
 &\le C M\prod_{j\ne i_*}a_j.
\end{align*}
The last inequality uses
\(s_\alpha\le1\), \(M\ge1\), and \(a_j\ge1\).
Apply \Cref{lem:V58-star-floor}.

On the high-output branch, the same gap split gives
\[
 {\Xi(\boldsymbol T)\over1-q_{\alpha,\boldsymbol T}}
 \le C\Xi(\boldsymbol T)\le C_DA\le C_DM
\]
by \Cref{lem:V58-root-dominance}.  Since \(c_j\le1\le a_j\),
\[
 {\Xi(\boldsymbol T)\over1-q_{\alpha,\boldsymbol T}}\,
 q_P\|\mathbf K_5^\alpha\|_{\rm op}
 \le C_DM\prod_{j\ne i_*}a_j.
\]
Again use \Cref{lem:V58-star-floor}.  Thus the paid scalar is bounded
by the selected star target in both branches.  Exact Fourier
dimension cancellation and raw pinching are the unchanged steps of
\Cref{lem:V57-exact-missing}, so the result is the full PGSEL5
statement.
\end{proof}

\begin{corollary}[Only branching dominance remains]
\label{cor:V58-only-branching}
Fix \(D\ge2C_F\).  Every unclosed selected discrete fibre must satisfy
\begin{equation}
 \Xi(\boldsymbol T)\le DY,
 \qquad
 \prod_{i=1}^5\Xi_i^{d_i-1}>DY^3
 \label{eq:V58-residual}
\end{equation}
for every tree fibre still used in its decomposition.
\end{corollary}

\begin{proof}
The output-dominant alternative is
\Cref{thm:V58-output-PGSEL5}.  On its complement, the first
inequality in \eqref{eq:V57-private-dominance} holds.  If the second
also holds, \Cref{rec:V57-private} closes the fibre.  Its failure is
the second inequality in \eqref{eq:V58-residual}.
\end{proof}

\section{Interpretation of the branching residual}
\label{sec:V58-residual}

\begin{lemma}[Branching dominance forces a root-dominant
high-degree row]
\label{lem:V58-root-hot-degree}
In the residual \eqref{eq:V58-residual}, at least one vertex \(i\)
with \(d_i\ge2\) satisfies
\begin{equation}
 \Xi_i>D^{1/3}Y,
 \qquad
 a_i>(D^{1/3}-1)Y.
 \label{eq:V58-root-hot-degree}
\end{equation}
\end{lemma}

\begin{proof}
Put \(e_i=d_i-1\ge0\).  A five-row tree satisfies
\[
 \sum_ie_i=3.
\]
If every vertex with \(e_i>0\) obeyed
\(\Xi_i\le D^{1/3}Y\), then
\[
 \prod_i\Xi_i^{e_i}
 \le(D^{1/3}Y)^{\sum_ie_i}=DY^3,
\]
contrary to \eqref{eq:V58-residual}.  Thus the first inequality holds
for some \(i\) with \(d_i\ge2\).  Since
\(P_i\le Y\) and \(\Xi_i=a_i+P_i\), the second follows.
\end{proof}

\begin{assessment}[Next exact object]
\label{ass:V58-next}
The output mass is now controlled by the private bank, while the
tree's three excess degree powers force at least one high-degree row
whose root mass dominates the entire private aggregate.  V59 should
retain the centered displacement bound at that row and split the
remaining \(125\) trees by degree profile.  A star has excess three
at one center; a fork has excess two plus one; a path has three
distinct excess-one vertices.  The exact private multiplier should
be charged only to excess degrees not already paid by the
root-dominant displacement.
\end{assessment}

\section{V58 result ledger and V59 programme}
\label{sec:V58-results}

\begin{theorem}[Fourth-series V58 result ledger]
\label{thm:V58-results}
Relative to V1--V57, V58 proves:
\begin{enumerate}[label=\textup{(V58.\arabic*)},leftmargin=6.3em]
\item formation of every final discrete-fibre output from root mass
      plus private mass;
\item root dominance under the output-dominant alternative;
\item a lower bound for the largest-root star target;
\item exact reassignment of the complete fifth cumulant to that star
      using the archived centered displacement bound;
\item PGSEL5 in both the low- and high-output gap regimes of the
      output-dominant branch;
\item reduction of every remaining discrete fibre to
      \eqref{eq:V58-residual}; and
\item identification of a root-dominant high-degree row in that
      residual.
\end{enumerate}
\end{theorem}

\begin{proof}
These are
\Cref{lem:V58-output-formation,
lem:V58-root-dominance,
lem:V58-star-floor,
lem:V58-star-displacement,
thm:V58-output-PGSEL5,
cor:V58-only-branching,
lem:V58-root-hot-degree}.
\end{proof}

\begin{programme}[V59: degree-resolved private moments]
\label{prog:V58-V59}
\begin{enumerate}[leftmargin=3em]
\item Split the \(125\) trees into stars and nonstars; the latter are
      paths or forks with maximum degree three.
\item Replace the aggregate private mass \(Y\) by rowwise private
      masses \(P_i\) in the exact multiplier product.
\item Use third- and fourth-order Wilson product moments only on the
      high-degree rows demanded by \(\sum(d_i-1)=3\).
\item Test whether weighted AM--GM allocates
      \(\prod\Xi_i^{d_i-1}\) to
      \((\sum P_i)^3\) after the local root masses have been removed.
\item Retain the centered displacement bound if a high-degree row is
      root-hot; otherwise charge its degree excess to its own private
      multiplier.
\end{enumerate}
\end{programme}

\begin{firewall}[Claim boundary after fourth-series V58]
\label{fire:V58-claim-boundary}
V58 proves PGSEL5 on the output-dominant branch and reduces the
unclosed discrete source to the branching-dominant condition
\eqref{eq:V58-residual}.  It does not prove that branch, the complete
discrete source, the full quintic MTA, an all-order atlas, WUFC, CPM,
cutoff removal, reflection positivity, Osterwalder--Schrader
reconstruction, construction of the continuum Yang--Mills measure,
or a uniform positive continuum spectral gap.  The full Yang--Mills
mass gap has not been proved.
\end{firewall}

\paragraph{Parent-version record.}
V58 was formed from
\texttt{Renewal\_Yang\_Mills\_Mass\_Gap\_Research\_Notes\_V57.tex}.
The V57 parent had \(27\,154\) logical source lines,
\(949\,167\) bytes, and SHA--256
\[
 \texttt{a975db03b686581ceeeff9c367fe6a234874524b6b49df4aebfb18de4a0ff51d}.
\]
The next compulsory companion audit is V60.

% =============================================================================
% END APPENDED FOURTH-SERIES V58 MATERIAL
% =============================================================================

% =============================================================================
% BEGIN APPENDED FOURTH-SERIES V59 MATERIAL
% =============================================================================

\chapter{V59: Closure of PGSEL5 and the global quintic MTA}
\label{chap:V59-quintic}

\section{Re-fibration by the root-dominant row}
\label{sec:V59-refibre}

In the residual \eqref{eq:V58-residual}, the tree which exposed
branching dominance need not remain the tree to which the complete
local cumulant is assigned.  The exact centered displacement bound
holds for every chosen center.  We therefore re-fibre the complete
cumulant by the root-dominant row supplied by
\Cref{lem:V58-root-hot-degree}.

\begin{lemma}[Root-dominant star floor]
\label{lem:V59-star-floor}
Suppose, for one row \(i\),
\begin{equation}
 a_i\ge\kappa Y
 \label{eq:V59-root-dominant}
\end{equation}
with \(\kappa>0\).  Let \(\tau_i\) be the five-row star centered at
\(i\).  Then
\begin{equation}
 \mathcal T_{\tau_i,r}
 ={a_i^4\over\Xi_i^3}\prod_{j\ne i}a_j
 \ge c_\kappa a_i\prod_{j\ne i}a_j.
 \label{eq:V59-star-floor}
\end{equation}
\end{lemma}

\begin{proof}
Since \(P_i\le Y\le a_i/\kappa\),
\[
 \Xi_i=a_i+P_i\le(1+\kappa^{-1})a_i.
\]
Insert this in the displayed star target.  The four leaves have
degree one and hence carry no private denominator.
\end{proof}

\begin{lemma}[Every branching residual has a root-dominant center]
\label{lem:V59-residual-center}
Choose \(D>8\).  Every fibre in
\eqref{eq:V58-residual} has a row \(i\) with \(d_i\ge2\) and
\begin{equation}
 a_i>\kappa_DY,\qquad
 \kappa_D:=D^{1/3}-1>1.
 \label{eq:V59-residual-center}
\end{equation}
\end{lemma}

\begin{proof}
This is \eqref{eq:V58-root-hot-degree}.  The numerical choice
\(D>8\) makes \(\kappa_D>1\).
\end{proof}

\begin{remark}[Why the original degree profile may be discarded]
\label{rem:V59-refibre}
Frozen f2 \texttt{lem:V56-hot-cumulant} bounds the complete signed
fifth cumulant for every chosen center.  Hence setting the star fibre
at the center in \Cref{lem:V59-residual-center} equal to the complete
cumulant and all other fibres to zero is an exact decomposition.
The output-dependent choice has only five possible values and is made
inside the raw output pinching; it creates no representation count.
\end{remark}

\section{Payment in the two output-gap regimes}
\label{sec:V59-payment}

\begin{theorem}[PGSEL5 on the branching-dominant residual]
\label{thm:V59-branching-PGSEL5}
Fix
\[
 D>\max\{8,2C_F\}.
\]
Every selected discrete fibre in the residual
\eqref{eq:V58-residual} obeys
\(\operatorname{PGSEL}_5\), uniformly in supports,
representations, final outputs, multiplicities, and payer placement.
\end{theorem}

\begin{proof}
Choose the row \(i\) from
\Cref{lem:V59-residual-center} and assign the complete cumulant to
the star \(\tau_i\).  The centered displacement theorem gives
\begin{equation}
 \|\mathbf K_5^\alpha\|_{\rm op}
 \le C\prod_{j\ne i}c_j,
 \qquad
 c_j\le\min\{1,\sqrt{s_\alpha a_j}\}.
 \label{eq:V59-centered-displacement}
\end{equation}
Also \(q_P\le1\).

On the low-output branch,
\[
 {\Xi(\boldsymbol T)\over1-q_{\alpha,\boldsymbol T}}
 \le {C\over s_\alpha}.
\]
Use the square-root alternative in
\eqref{eq:V59-centered-displacement} on all four leaves:
\begin{align*}
 {\Xi(\boldsymbol T)\over1-q_{\alpha,\boldsymbol T}}\,
 q_P\|\mathbf K_5^\alpha\|_{\rm op}
 &\le
 {C\over s_\alpha}
 \prod_{j\ne i}\sqrt{s_\alpha a_j}\\
 &=Cs_\alpha\prod_{j\ne i}a_j^{1/2}\\
 &\le C a_i\prod_{j\ne i}a_j.
\end{align*}
The last inequality uses \(s_\alpha\le1\) and
\(a_i,a_j\ge1\).  Apply \Cref{lem:V59-star-floor}.

On the high-output branch, the residual includes
\(\Xi(\boldsymbol T)\le DY\), so
\[
 {\Xi(\boldsymbol T)\over1-q_{\alpha,\boldsymbol T}}
 \le C\Xi(\boldsymbol T)
 \le CDY
 \le {CD\over\kappa_D}a_i.
\]
Use \(c_j\le1\le a_j\) in
\eqref{eq:V59-centered-displacement}:
\[
 {\Xi(\boldsymbol T)\over1-q_{\alpha,\boldsymbol T}}\,
 q_P\|\mathbf K_5^\alpha\|_{\rm op}
 \le C_Da_i\prod_{j\ne i}a_j.
\]
Again apply \Cref{lem:V59-star-floor}.  Exact dimension cancellation,
raw pinching, and the contractive suffix are already incorporated by
\Cref{lem:V57-exact-missing}.  This proves PGSEL5 in both regimes.
\end{proof}

\begin{firewall}[No multiplier monotonicity]
\label{fire:V59-no-monotonicity}
The proof uses only \(q_P\le1\), the final product-gap split, and the
complete centered displacement bound.  It never compares a fused
output multiplier with any input multiplier and never divides by an
intermediate gap.
\end{firewall}

\section{Completion of the discrete source}
\label{sec:V59-discrete}

\begin{theorem}[Complete paid selected fifth-junction gate]
\label{thm:V59-complete-PGSEL5}
The gate \(\operatorname{PGSEL}_5\) of
\Cref{def:V57-PGSEL5} holds on every selected discrete quintic fibre.
\end{theorem}

\begin{proof}
Fix \(D>\max\{8,2C_F\}\).  If
\(\Xi(\boldsymbol T)>DY\), apply
\Cref{thm:V58-output-PGSEL5}.  Otherwise
\(\Xi(\boldsymbol T)\le DY\).  If the branching inequality in
\eqref{eq:V57-private-dominance} holds, apply the archived
private-terminal compiler recorded in
\Cref{rec:V57-private}.  If it fails, the fibre lies in
\eqref{eq:V58-residual} and
\Cref{thm:V59-branching-PGSEL5} applies.  The three alternatives are
exhaustive.
\end{proof}

\begin{theorem}[Complete discrete quintic source]
\label{thm:V59-discrete-source}
The discrete predecessor source
\eqref{eq:V56-discrete-identity} satisfies
\begin{equation}
 \kappa_\otimes^{\rm abs}
 (\mathcal S_{\widehat0}^{(5),{\rm pay}})
 \le
 C\sum_{\tau\in\mathfrak T_5}
 \prod_{ij\in E(\tau)}\theta_{ij}.
 \label{eq:V59-discrete-source}
\end{equation}
The constant is uniform in support, volume, representations, output
multiplicities, and payer placement.
\end{theorem}

\begin{proof}
Apply \Cref{thm:V59-complete-PGSEL5} in the conditional source theorem
\Cref{thm:V57-conditional-discrete}.  The diagonal four-mark
summation \eqref{eq:V57-diagonal-sum} removes the coordinate \(r\)
without a support count.  Raw output pinching removes representation
and multiplicity counts.
\end{proof}

\section{Global quintic marked-tree atlas}
\label{sec:V59-global}

\begin{theorem}[Global quintic MTA]
\label{thm:V59-global-quintic-MTA}
For arbitrary finite five-row supports and trace-class
row-factorized Fourier coefficients, the complete once-paid connected
five-row physical carrier satisfies
\begin{equation}
 \boxed{
 \kappa_\otimes^{\rm abs}(\mathcal C_5^{\rm pay})
 \le
 C_5\sum_{\tau\in\mathfrak T_5}
 \prod_{ij\in E(\tau)}\theta_{ij}.}
 \label{eq:V59-global-quintic-MTA}
\end{equation}
The estimate is uniform in volume, supports, input representations,
physical output multiplicities, output-duality orientation, and payer
placement.
\end{theorem}

\begin{proof}
Use the exact coefficient-one first-connectivity census of archive f3
V97.  Every nondiscrete predecessor source is closed by
\Cref{cor:V56-nondiscrete}.  The unique discrete source is
\Cref{thm:V59-discrete-source}.  The \(51\) source families are
disjoint and exhaustive.  Each family is bounded by a finite sum of
five-row tree monomials through the exact internal/quotient tree lift;
enlarge each sum to the complete \(125\)-tree atlas and absorb the
fixed source and payer counts into \(C_5\).
\end{proof}

\begin{corollary}[The fixed-order quintic frontier is discharged]
\label{cor:V59-quintic-discharged}
No five-row first-connectivity source family remains open.
\end{corollary}

\begin{proof}
This is \Cref{thm:V59-global-quintic-MTA} together with the exhaustive
V97 source census.
\end{proof}

\section{What quintic closure does not imply}
\label{sec:V59-boundary}

\begin{assessment}[The next obstruction is arity growth]
\label{ass:V59-next}
The programme has now proved complete marked-tree atlases through row
order five.  The next implication in the route of
\eqref{eq:V1-route} is not automatic.  WUFC and the Perron/cluster
majorant require one constant \(K_0\) controlling every row order,
roughly \(K_0^{m-1}\), together with factorial or tree-enumeration
weights.  Fixed constants \(C_m<\infty\) for \(m\le5\) do not provide
such a bound.  V60 must perform the companion audit and then extract
the exact \(m\)-dependence of the displacement moments, private
product moments, first-connectivity census, and scalar-lift
normalization before proposing an all-order induction.
\end{assessment}

\section{V59 result ledger and V60 programme}
\label{sec:V59-results}

\begin{theorem}[Fourth-series V59 result ledger]
\label{thm:V59-results}
Relative to V1--V58, V59 proves:
\begin{enumerate}[label=\textup{(V59.\arabic*)},leftmargin=6.3em]
\item a root-dominant star floor for the branching residual;
\item exact re-fibration of the complete fifth cumulant by the
      root-dominant row;
\item PGSEL5 in both final-output gap regimes of that residual;
\item the complete paid selected fifth-junction gate;
\item closure of the unique discrete quintic source;
\item the support- and representation-uniform global quintic
      marked-tree atlas; and
\item discharge of all \(51\) quintic first-connectivity families.
\end{enumerate}
\end{theorem}

\begin{proof}
These are
\Cref{lem:V59-star-floor,
lem:V59-residual-center,
thm:V59-branching-PGSEL5,
thm:V59-complete-PGSEL5,
thm:V59-discrete-source,
thm:V59-global-quintic-MTA,
cor:V59-quintic-discharged}.
\end{proof}

\begin{programme}[V60: audit and all-order constant ledger]
\label{prog:V59-V60}
\begin{enumerate}[leftmargin=3em]
\item Perform the compulsory read-only audit of the newest active
      SM--Einstein and Navier--Stokes notes.
\item Record explicit arity dependence in the mixed displacement
      moments and the complete cumulant M\"obius sum.
\item Record the \(m\)-dependence of private Wilson product moments
      and the first-connectivity partition census.
\item Separate Cayley tree growth \(m^{m-2}\) from the per-edge
      analytic constant required by WUFC.
\item Test whether the scalar-spectator lift and payer stripping are
      uniform as \(m\to\infty\).
\item Formulate an all-order induction only if every constant is
      bounded by \(m!K_0^{m-1}\) in the precise WUFC normalization.
\end{enumerate}
\end{programme}

\begin{firewall}[Claim boundary after fourth-series V59]
\label{fire:V59-claim-boundary}
V59 proves the complete fixed-order global quintic marked-tree atlas.
It does not prove an all-order atlas, the required exponential/factorial
arity bound, WUFC, CPM, cutoff removal, reflection positivity,
Osterwalder--Schrader reconstruction, construction of the continuum
Yang--Mills measure, or a uniform positive continuum spectral gap.
The full Yang--Mills mass gap has not been proved.
\end{firewall}

\paragraph{Parent-version record.}
V59 was formed from
\texttt{Renewal\_Yang\_Mills\_Mass\_Gap\_Research\_Notes\_V58.tex}.
The V58 parent had \(27\,474\) logical source lines,
\(959\,645\) bytes, and SHA--256
\[
 \texttt{2032d94c57646e8f18d587a04ea3f74e6d112f8e88cfa8097cae99758da2bbcc}.
\]
The next compulsory companion audit is V60.

% =============================================================================
% END APPENDED FOURTH-SERIES V59 MATERIAL
% =============================================================================

% =============================================================================
% BEGIN APPENDED FOURTH-SERIES V60 MATERIAL
% =============================================================================

\chapter{V60: Companion audit and the all-order constant ledger}
\label{chap:V60-arity}

\section{Scheduled companion audit}
\label{sec:V60-audit}

The read-only audit was refreshed after the parallel SM programme
advanced during inspection.  The files actually audited were
\begin{center}
\begin{tabular}{p{0.47\textwidth}rr}
\toprule
file&logical lines&bytes\\
\midrule
\texttt{Renewal\_SM\_Einstein\_Continuum\_Research\_Notes\_V51.tex}
&\(20\,045\)&\(686\,539\)\\
\texttt{Renewal\_NS\_Stability\_Research\_Notes\_V31.tex}
&\(23\,052\)&\(887\,537\)\\
\bottomrule
\end{tabular}
\end{center}
Their SHA--256 digests are
\[
\begin{split}
 &\texttt{e64547c5ebcbf5488daf88ac0d22a2a5cfb2e9c5025aeaa3dc3ee1aab964998f},\\
 &\texttt{f1121b62238ad5491bb7cf03635ea9934942edf573ed8faaa9ee79e1d38a6696}.
\end{split}
\]

SM V51 proves in a massive scalar model that the one-site contraction
gap collapses like \(a^2\), while diffusive acceleration and
directional conjugation retain finite physical decay.  It explicitly
excludes massless gauge directions.  The YM lesson is that a
fixed-regulator row gap cannot be promoted to a continuum mass gap
without the correct physical scaling and reconstruction.  NS V31 is
unchanged: its signed formation source must be retained before
positive norms, and its present bound is still continuation-class.

\begin{firewall}[Audit type boundary]
\label{fire:V60-audit}
No scalar convexity theorem, NS stopping estimate, or physical mass is
imported into YM.  Only the warnings about assembly order and
continuum scaling are retained.  The next compulsory audit is V65.
\end{firewall}

\section{Factorial-exponential combinatorics}
\label{sec:V60-combinatorics}

Let
\[
 L_m:=\sum_{\pi\in\Pi_m}(|\pi|-1)!
     =\sum_{k=1}^m S(m,k)(k-1)!,
\]
the absolute M\"obius mass of the \(m\)-th cumulant.

\begin{theorem}[M\"obius mass has factorial-exponential growth]
\label{thm:V60-Mobius-growth}
There is \(K_{\rm M}<\infty\) such that
\begin{equation}
 L_m\le m!K_{\rm M}^{\,m}
 \qquad(m\ge1).
 \label{eq:V60-Mobius-growth}
\end{equation}
\end{theorem}

\begin{proof}
The exponential generating function is
\[
 \sum_{m\ge1}{L_m\over m!}z^m
 =\sum_{k\ge1}{(e^z-1)^k\over k}
 =-\log(2-e^z).
\]
For any fixed \(0<r<\log2\), one has
\(|e^z-1|\le e^r-1<1\) on \(|z|\le r\), so the last function is
bounded there.  Cauchy's estimate gives
\(L_m/m!\le M_rr^{-m}\).  Enlarge \(r^{-1}\) to absorb \(M_r\).
\end{proof}

\begin{lemma}[Partition and tree censuses fit the same scale]
\label{lem:V60-census-growth}
There is a universal \(K_{\rm C}\) such that
\begin{equation}
 B_m+m^{m-2}\le m!K_{\rm C}^{\,m}
 \qquad(m\ge2),
 \label{eq:V60-census-growth}
\end{equation}
where \(B_m\) is the Bell number and \(m^{m-2}\) is the number of
labelled spanning trees.
\end{lemma}

\begin{proof}
A set partition is encoded by a map from \([m]\) to \([m]\), so
\(B_m\le m^m\).  Cayley's number is also at most \(m^m\).
The elementary Stirling lower bound \(m!\ge(m/e)^m\) gives
\(m^m\le e^m m!\).
\end{proof}

\begin{corollary}[First-connectivity enumeration is not the
all-order obstruction]
\label{cor:V60-combinatorics}
The absolute cumulant expansion, the proper-prefix census, and the
labelled-tree atlas can all be absorbed into a bound of the WUFC form
\(m!K^m\).  None forces super-factorial growth.
\end{corollary}

\section{Analytic constants}
\label{sec:V60-analytic}

\begin{lemma}[Mixed displacement moments are factorial-exponential]
\label{lem:V60-displacement-growth}
The normal-coordinate proof of the archived displacement estimate may
be chosen so that its order-\(m\) constant obeys
\begin{equation}
 C_m^{\rm disp}\le m!K_{\rm D}^{\,m}.
 \label{eq:V60-displacement-growth}
\end{equation}
\end{lemma}

\begin{proof}
The pointwise proof bounds a product of \(m\) displacements by
\[
 K^m\prod_{\ell=1}^mc_{R_\ell}(1+\|Z\|)^m.
\]
Uniform Wilson radial domination gives
\(\mathbb E(1+\|Z\|)^m\le K_1^m m^{m/2}\).
Since \(m^{m/2}\le e^m m!\), the normal-neighbourhood contribution
has the stated bound.  The complement is exponentially small in
\(s_\alpha^{-1}\); optimizing the elementary exponential domination
at order \(m\) adds at most another \(m!K_2^m\).
\end{proof}

\begin{lemma}[The scalar lift is arity-neutral]
\label{lem:V60-scalar-lift}
The fixed private-spectator payer stripping of
\Cref{thm:V55-payer-stripping} introduces a constant independent of
\(m\), apart from the constant already present in the paid
\(m\)-row atlas.
\end{lemma}

\begin{proof}
The spectator representation, its dimension, height, multiplier, and
payer are fixed.  Every spanning tree on \(m\ge2\) has degree at least
one at the spectator row, so the augmented row-mass ratio is absorbed
without an \(m\)-dependent endpoint sum.
\end{proof}

\begin{assessment}[Exact all-order bottleneck]
\label{ass:V60-bottleneck}
\Cref{thm:V60-Mobius-growth,lem:V60-census-growth,
lem:V60-displacement-growth,lem:V60-scalar-lift}
remove the elementary combinatorial and spectator constants.  What is
not proved is an all-order operator decomposition
\[
 \mathbf K_m^\alpha
 =\sum_{\tau\in\mathfrak T_m}\mathbf K_{m,\tau}^\alpha,
 \qquad
 \|\mathbf K_{m,\tau}^\alpha\|_{\rm op}
 \le K_0^{m-1}s_\alpha^{m-1}
 \prod_iA_i^{\deg_\tau(i)/2},
\]
uniformly in \(m\), together with private product moments on the same
constant scale.  The fifth-order proof assigns the entire hot
cumulant to one star and uses a finite partition lattice; it does not
yet provide this arity-uniform decomposition.
\end{assessment}

\section{V60 ledger and V61 programme}
\label{sec:V60-results}

\begin{theorem}[Fourth-series V60 result ledger]
\label{thm:V60-results}
V60 proves the scheduled audit, factorial-exponential bounds for
cumulant M\"obius mass, partition and tree censuses, a
factorial-exponential displacement-moment bound, and arity-neutrality
of the scalar lift.  It localizes the all-order obstruction to the
uniform analytic tree decomposition and matching private moments.
\end{theorem}

\begin{proof}
Use
\Cref{sec:V60-audit,
thm:V60-Mobius-growth,
lem:V60-census-growth,
lem:V60-displacement-growth,
lem:V60-scalar-lift,
ass:V60-bottleneck}.
\end{proof}

\begin{programme}[V61: all-order hot-star numerator]
\label{prog:V60-V61}
\begin{enumerate}[leftmargin=3em]
\item Prove the centered \(m\)-cumulant displacement bound while
      tracking the absolute M\"obius mass by
      \eqref{eq:V60-Mobius-growth}.
\item Assign a hot tuple to the least hot-row star and determine the
      exact per-edge constant.
\item Derive private Wilson product moments with
      \(m!K^m\) constants from the same radial domination.
\item Separate the thermal all-order tree expansion from the hot-star
      branch; do not infer it from the fixed \(m=5\) theorem.
\end{enumerate}
\end{programme}

\begin{firewall}[Claim boundary after fourth-series V60]
\label{fire:V60-claim-boundary}
V60 proves the global quintic MTA and several required all-order
constant bounds.  It does not prove the all-order local tree
decomposition, all-order MTA, WUFC, CPM, cutoff removal, reflection
positivity, Osterwalder--Schrader reconstruction, construction of the
continuum Yang--Mills measure, or a uniform positive continuum
spectral gap.  The full Yang--Mills mass gap has not been proved.
\end{firewall}

\paragraph{Parent-version record.}
V60 was formed from
\texttt{Renewal\_Yang\_Mills\_Mass\_Gap\_Research\_Notes\_V59.tex}.
The V59 parent had \(27\,795\) logical source lines,
\(970\,860\) bytes, and SHA--256
\[
 \texttt{0c49a866abe5020fc8679a208182aebf6ecb520932ce610815a6b6acf4107369}.
\]
The next compulsory companion audit is V65.

% =============================================================================
% END APPENDED FOURTH-SERIES V60 MATERIAL
% =============================================================================

% =============================================================================
% BEGIN APPENDED FOURTH-SERIES V61 MATERIAL
% =============================================================================

\chapter{V61: An all-order hot-star cumulant theorem}
\label{chap:V61-hot-star}

\section{Weighted partition sums}
\label{sec:V61-weighted-partitions}

The separate bounds of V60 must not be multiplied in their worst
forms: doing so would create a spurious second factorial.  The block
moments and partition lattice are summed together.

\begin{lemma}[Weighted logarithmic partition bound]
\label{lem:V61-weighted-partition}
Let \(w_n\ge0\), \(n\ge1\), and suppose
\begin{equation}
 w_n\le n!K^n.
 \label{eq:V61-block-weight}
\end{equation}
Define
\begin{equation}
 S_m:=
 \sum_{\pi\in\Pi_m}(|\pi|-1)!
 \prod_{B\in\pi}w_{|B|}.
 \label{eq:V61-weighted-partition-sum}
\end{equation}
Then there is \(K'<\infty\), depending only on \(K\), such that
\begin{equation}
 S_m\le m!(K')^m.
 \label{eq:V61-weighted-partition-bound}
\end{equation}
\end{lemma}

\begin{proof}
Put
\[
 W(z):=\sum_{n\ge1}{w_n\over n!}z^n.
\]
The hypothesis makes \(W\) analytic for \(|z|<K^{-1}\), and
\(W(0)=0\).  The labelled exponential formula gives
\[
 \sum_{m\ge1}{S_m\over m!}z^m
 =
 \sum_{k\ge1}{W(z)^k\over k}
 =-\log(1-W(z)).
\]
Choose \(r>0\) so small that
\(\sup_{|z|\le r}|W(z)|<1\).  Cauchy's estimate on that circle gives
\(S_m/m!\le M_rr^{-m}\), which is
\eqref{eq:V61-weighted-partition-bound} after enlarging \(K'\).
\end{proof}

\begin{remark}[Why V60's two separate estimates were insufficient]
\label{rem:V61-no-double-factorial}
The crude product
\[
 \left(\sum_{\pi}(|\pi|-1)!\right)
 \max_\pi\prod_{B\in\pi}w_{|B|}
\]
can display two factorials.  The exact generating function in
\Cref{lem:V61-weighted-partition} proves that this is an artefact of
separating the partition and block-moment sums.
\end{remark}

\section{All-order centered displacement cumulants}
\label{sec:V61-cumulant}

At one Wilson coordinate let
\[
 X_i(U):=D^{R_i}(U),\qquad
 Y_i(U):=X_i(U)-I,\qquad
 b_i:=\sqrt{s_\alpha A_i},\qquad
 c_i:=\min\{1,b_i\}.
\]
All operators act on separate tensor legs.

\begin{lemma}[Uniform mixed displacement blocks with tracked order]
\label{lem:V61-mixed-block}
There is \(K_{\rm b}\) such that for every nonempty
\(B\subseteq[m]\),
\begin{equation}
 \left\|
 \mathbb E_{\nu_\alpha}\bigotimes_{j\in B}Y_j
 \right\|_{\rm op}
 \le
 w_{|B|}\prod_{j\in B}c_j,
 \qquad
 w_n\le n!K_{\rm b}^{\,n}.
 \label{eq:V61-mixed-block}
\end{equation}
The constants are independent of \(m\), the representations, and
their dimensions.
\end{lemma}

\begin{proof}
The normal-coordinate displacement estimate gives
\[
 \|Y_j\|_{\rm op}
 \le Kc_j(1+\|Z\|)
\]
on the Wilson well.  Hence the block moment is bounded by
\[
 K^n\mathbb E(1+\|Z\|)^n\prod_{j\in B}c_j
 \le K_1^n n^{n/2}\prod_{j\in B}c_j.
\]
The inequality \(n^{n/2}\le e^n n!\) gives the required \(w_n\).
Off the well, \(\|Y_j\|\le2\), while the Wilson mass is
exponentially small in \(s_\alpha^{-1}\).  The nontrivial Casimir
floor gives \(c_j\ge c\sqrt{s_\alpha}\); optimizing exponential
domination at order \(n\) adds at most \(n!K_2^n\).  Combine the two
regions.
\end{proof}

\begin{theorem}[All-order centered displacement cumulant]
\label{thm:V61-centered-cumulant}
For every \(m\ge2\) and every chosen center \(i\),
\begin{equation}
 \boxed{
 \left\|\kappa_m^{\nu_\alpha}
 (X_1,\ldots,X_m)\right\|_{\rm op}
 \le m!K_{\rm h}^{\,m}
 \prod_{j\ne i}c_j.}
 \label{eq:V61-centered-cumulant}
\end{equation}
The estimate is uniform in \(m,\alpha\), representations, dimensions,
and spectator amplification.
\end{theorem}

\begin{proof}
A cumulant of order at least two is invariant under adding a constant
to any argument, so replace every \(X_j\) by \(Y_j\).  Expand
\[
 \kappa_m(Y_1,\ldots,Y_m)
 =
 \sum_{\pi\in\Pi_m}
 (|\pi|-1)!(-1)^{|\pi|-1}
 \bigotimes_{B\in\pi}
 \mathbb E\bigotimes_{j\in B}Y_j.
\]
In each block containing \(i\), discard \(c_i\le1\); retain every
other \(c_j\) exactly once over the partition.  Taking norms and using
\Cref{lem:V61-mixed-block} leaves
\[
 \prod_{j\ne i}c_j
 \sum_{\pi\in\Pi_m}(|\pi|-1)!
 \prod_{B\in\pi}w_{|B|}.
\]
\Cref{lem:V61-weighted-partition} bounds the sum by
\(m!K_{\rm h}^m\).  Tensoring identities on spectator carriers
preserves every operator norm.
\end{proof}

\section{Hot-star assignment at arbitrary arity}
\label{sec:V61-star}

For a labelled tree \(\tau\) on \([m]\), define its local thermal
weight
\[
 \mathfrak w_\tau
 :=s_\alpha^{m-1}
 \prod_{j=1}^mA_j^{\deg_\tau(j)/2}
 =\prod_{j=1}^mb_j^{\deg_\tau(j)}.
\]

\begin{theorem}[All-order hot-star local tree card]
\label{thm:V61-hot-star}
Suppose \(b_i\ge1\) for at least one row, and choose the least such
label \(i_*\).  Let \(\tau_{i_*}\) be the star centered at \(i_*\).
Then the complete one-coordinate \(m\)-th Wilson cumulant admits the
exact decomposition
\[
 K_{m,\tau_{i_*}}:=K_m,\qquad K_{m,\tau}:=0
 \quad(\tau\ne\tau_{i_*}),
\]
and
\begin{equation}
 \boxed{
 \|K_{m,\tau_{i_*}}\|_{\rm op}
 \le m!K_{\rm h}^{\,m}\mathfrak w_{\tau_{i_*}}.}
 \label{eq:V61-hot-star}
\end{equation}
\end{theorem}

\begin{proof}
Apply \Cref{thm:V61-centered-cumulant} with center \(i_*\).
Since \(c_j\le b_j\) and \(b_{i_*}\ge1\),
\[
 \prod_{j\ne i_*}c_j
 \le b_{i_*}^{m-1}\prod_{j\ne i_*}b_j
 =\mathfrak w_{\tau_{i_*}}.
\]
Assigning the complete cumulant to one star is an exact equality, not
an estimate of separately normed partition terms.
\end{proof}

\begin{corollary}[The hot local sector has the WUFC arity scale]
\label{cor:V61-hot-WUFC-scale}
After absorbing one fixed factor into the base constant,
\eqref{eq:V61-hot-star} has the required form
\[
 m!K_0^{m-1}\mathfrak w_{\tau_{i_*}}.
\]
Thus no all-order obstruction remains in the thermally hot local
sector.
\end{corollary}

\begin{proof}
For \(m\ge2\),
\(K_{\rm h}^m\le
\max\{K_{\rm h},K_{\rm h}^2\}^{m-1}\).
\end{proof}

\section{Remaining thermal problem}
\label{sec:V61-thermal}

\begin{assessment}[The thermal sector needs connected excess]
\label{ass:V61-thermal}
If every \(b_i<1\), the bound
\eqref{eq:V61-centered-cumulant} retains only \(m-1\) displacement
factors, whereas a tree weight contains \(2(m-1)\) half-displacements.
At fixed order, Gaussian connectedness and Wilson-density corrections
supply the missing excess.  An all-order proof must sum those
connected contractions with \(m!K^m\) constants.  Absolute
moment--cumulant expansion alone cannot supply the extra small
factors.
\end{assessment}

\begin{firewall}[No all-order MTA yet]
\label{fire:V61-no-thermal}
\Cref{thm:V61-hot-star} proves only the local sector in which at least
one \(s_\alpha A_i\ge1\).  It does not prove the all-thermal local
tree decomposition, nor does it track the all-order constants in the
private product compiler.
\end{firewall}

\section{V61 ledger and V62 programme}
\label{sec:V61-results}

\begin{theorem}[Fourth-series V61 result ledger]
\label{thm:V61-results}
V61 proves a weighted partition summation which avoids a false second
factorial, a factorial-exponential mixed displacement-block bound, the
all-order centered cumulant estimate
\eqref{eq:V61-centered-cumulant}, and the all-order hot-star local
tree card at the exact WUFC arity scale.
\end{theorem}

\begin{proof}
Use
\Cref{lem:V61-weighted-partition,
lem:V61-mixed-block,
thm:V61-centered-cumulant,
thm:V61-hot-star,
cor:V61-hot-WUFC-scale}.
\end{proof}

\begin{programme}[V62: private moments with tracked arity]
\label{prog:V61-V62}
\begin{enumerate}[leftmargin=3em]
\item Reopen the Bernoulli product compiler used in frozen f2 V52
      and record its exact dependence on moment order.
\item Prove one-link failure debits with \(k!K^k\) constants.
\item Deduce
      \((s_\alpha Y)^kq_P\le k!K^k\) uniformly in the private bank.
\item Insert that estimate into the arbitrary-order private-terminal
      compiler and isolate the thermal local tree card as the sole
      remaining MTA input.
\end{enumerate}
\end{programme}

\begin{firewall}[Claim boundary after fourth-series V61]
\label{fire:V61-claim-boundary}
V61 proves the all-order thermally hot local cumulant tree card.  It
does not prove the all-thermal local card, tracked private-product
constants, all-order MTA, WUFC, CPM, cutoff removal, reflection
positivity, Osterwalder--Schrader reconstruction, construction of the
continuum Yang--Mills measure, or a uniform positive continuum
spectral gap.  The full Yang--Mills mass gap has not been proved.
\end{firewall}

\paragraph{Parent-version record.}
V61 was formed from
\texttt{Renewal\_Yang\_Mills\_Mass\_Gap\_Research\_Notes\_V60.tex}.
The V60 parent had \(28\,025\) logical source lines,
\(978\,872\) bytes, and SHA--256
\[
 \texttt{cee113cc7e36bfca363e515101af3669da7449f98605aee03edf20eccee7c024}.
\]
The next compulsory companion audit is V65.

% =============================================================================
% END APPENDED FOURTH-SERIES V61 MATERIAL
% =============================================================================

% =============================================================================
% BEGIN APPENDED FOURTH-SERIES V62 MATERIAL
% =============================================================================

\chapter{V62: The private-moment \(m!\) bound is false}
\label{chap:V62-private-growth}

\section{What the Bernoulli compiler would preserve}
\label{sec:V62-compiler}

\begin{lemma}[Tracked Bernoulli product compiler]
\label{lem:V62-tracked-Bernoulli}
Let \(x_f\ge0\), \(0\le r_f\le1\), and suppose, for every
\(k\ge1\),
\begin{equation}
 x_f^kr_f\le k!K^k(1-r_f).
 \label{eq:V62-one-site-hypothesis}
\end{equation}
Then, uniformly in the finite index set \(F\),
\begin{equation}
 \left(\sum_{f\in F}x_f\right)^m
 \prod_{f\in F}r_f
 \le m!(2K)^m.
 \label{eq:V62-product-conclusion}
\end{equation}
\end{lemma}

\begin{proof}
Expand the \(m\)-th power and group a monomial by its \(j\) distinct
indices and their positive ordered multiplicities
\(k_1+\cdots+k_j=m\).  Its multinomial coefficient is
\(m!/\prod_\ell k_\ell!\).  Applying
\eqref{eq:V62-one-site-hypothesis} at the distinct indices leaves the
factor
\[
 {m!\over\prod_\ell k_\ell!}
 \prod_{\ell=1}^j(k_\ell!K^{k_\ell})
 =m!K^m.
\]
For each multiplicity pattern, the sum over unordered distinct
indices is bounded by the probability of exactly \(j\) failures in
independent Bernoulli trials, hence by one.  There are
\(\binom{m-1}{j-1}\) positive ordered compositions.  Summing over
\(j\) gives \(2^{m-1}\), proving
\eqref{eq:V62-product-conclusion}.
\end{proof}

\begin{assessment}[The compiler is not the source of bad growth]
\label{ass:V62-compiler}
The frozen f2 Bernoulli argument is compatible with the desired
factorial-exponential scale.  The issue is whether its one-site
hypothesis has constants \(k!K^k\) uniformly in \(k\).
\end{assessment}

\section{Exact \(SU(2)\) regression}
\label{sec:V62-SU2}

Fix once and for all a finite \(\alpha>0\).  For spin
\(j=(n-1)/2\), where \(n\) is a positive odd integer, V40 gives
\begin{equation}
 \eta_j(\alpha)={I_n(\alpha)\over I_1(\alpha)}.
 \label{eq:V62-Bessel-multiplier}
\end{equation}
The modified Bessel series implies
\begin{equation}
 I_n(\alpha)
 =\sum_{r\ge0}{(\alpha/2)^{n+2r}\over r!(n+r)!}
 \ge {(\alpha/2)^n\over n!}.
 \label{eq:V62-Bessel-lower}
\end{equation}

\begin{theorem}[No factorial-exponential scaled-Casimir moments]
\label{thm:V62-no-factorial-moments}
For every fixed \(\alpha>0\) and every \(K<\infty\),
\begin{equation}
 \sup_{\substack{j\in\frac12\mathbb N_0\\j>0}}
 \left[s_\alpha(1+C_2(j))\right]^m
 \eta_j(\alpha)
 >m!K^m
 \label{eq:V62-failure}
\end{equation}
for all sufficiently large \(m\).  Consequently there is no
\(\alpha_0<\infty\) and \(K<\infty\) for which the one-link
scaled-Casimir moments are bounded by \(m!K^m\) simultaneously for
all \(m\) and all \(\alpha\ge\alpha_0\).
\end{theorem}

\begin{proof}
For odd \(n\ge3\),
\[
 1+C_2((n-1)/2)\ge c n^2.
\]
Using \eqref{eq:V62-Bessel-lower} and \(n!\le n^n\), the left side of
\eqref{eq:V62-failure} is at least
\begin{equation}
 C_\alpha
 (c_\alpha n^2)^m
 {(\alpha/2)^n\over n^n}.
 \label{eq:V62-lower-moment}
\end{equation}
Choose an odd integer \(n_m\) with
\[
 {m\over2\log m}\le n_m\le {2m\over\log m}
\]
for all large \(m\).  Taking logarithms in
\eqref{eq:V62-lower-moment} gives
\begin{align*}
 \log(\hbox{lower bound})
 &\ge
 2m\log n_m-n_m\log n_m-O_\alpha(m+n_m)\\
 &\ge
 2m\log m-2m\log\log m-O_\alpha(m).
 \label{eq:V62-log-lower}
\end{align*}
On the other hand,
\[
 \log(m!K^m)\le m\log m+O_K(m).
\]
The difference tends to \(+\infty\), proving the first assertion.
If a uniform \(\alpha_0,K\) existed, set \(\alpha=\alpha_0\) and
contradict the first assertion.
\end{proof}

\begin{corollary}[The desired one-site failure debit is also false]
\label{cor:V62-no-failure-debit}
There is no fixed \(K\) such that
\begin{equation}
 x_j^m\eta_j(\alpha)
 \le m!K^m(1-\eta_j(\alpha))
 \label{eq:V62-false-failure-debit}
\end{equation}
holds for all \(m\), all spins, and all
\(\alpha\ge\alpha_0\).
\end{corollary}

\begin{proof}
Since \(1-\eta_j(\alpha)\le1\), the claimed failure debit would imply
the moment bound contradicted by
\Cref{thm:V62-no-factorial-moments}.
\end{proof}

\begin{corollary}[The absolute private-moment route cannot prove
WUFC]
\label{cor:V62-private-route-fails}
The frozen f2 proof of
\((s_\alpha Y)^mq_P\le C_m^{\rm priv}\), combined with absolute
one-link Casimir moments, cannot yield
\(C_m^{\rm priv}\le m!K^m\) uniformly in \(m\).
\end{corollary}

\begin{proof}
Take a private bank consisting of one coordinate.  Then the product
moment is exactly the one-link quantity in
\Cref{thm:V62-no-factorial-moments}.  No product compiler can improve
that one-coordinate lower bound.
\end{proof}

\section{What remains valid}
\label{sec:V62-valid}

\begin{firewall}[The hot-star theorem survives]
\label{fire:V62-hot-star}
\Cref{thm:V61-hot-star} uses bounded representation displacements
\(c_R=\min\{1,\sqrt{s_\alpha A_R}\}\), not positive powers of
\(A_R\eta_R\).  Therefore
\Cref{thm:V62-no-factorial-moments} does not affect the all-order
hot-star local cumulant card.
\end{firewall}

\begin{assessment}[Necessary direction change]
\label{ass:V62-direction}
An all-order proof must avoid charging arbitrarily high powers of a
private Casimir to one Wilson multiplier at a fixed regulator.  Two
routes remain logically possible:
\begin{enumerate}[label=\textup{(\roman*)},leftmargin=3em]
\item re-fibre every high-degree private terminal into a hot local
      star, as V59 did at order five, so only bounded displacements
      rather than Casimir moments are used; or
\item change the analytic arity/representation grade to the sharp
      Gevrey scale dictated by the Bessel tail, and reprove WUFC and
      the Perron majorant in that grade.
\end{enumerate}
The first route preserves the present \(m!K^m\) target and must be
tested before changing the norm.
\end{assessment}

\section{V62 ledger and V63 programme}
\label{sec:V62-results}

\begin{theorem}[Fourth-series V62 result ledger]
\label{thm:V62-results}
V62 proves that the Bernoulli compiler preserves
factorial-exponential constants under a factorial-exponential
one-site debit, and then proves by the exact \(SU(2)\) Bessel lower
bound that this one-site debit is false.  It follows that the absolute
private-Casimir-moment route cannot establish the current WUFC arity
constant.
\end{theorem}

\begin{proof}
Use
\Cref{lem:V62-tracked-Bernoulli,
thm:V62-no-factorial-moments,
cor:V62-no-failure-debit,
cor:V62-private-route-fails}.
\end{proof}

\begin{programme}[V63: all-order terminal re-fibration]
\label{prog:V62-V63}
\begin{enumerate}[leftmargin=3em]
\item Generalize the V58--V59 output/private/branching trichotomy to
      an \(m\)-row selected junction and a tree with
      \(\sum_i(\deg_i-1)=m-2\).
\item In the failed branching branch, use weighted pigeonhole to
      find a root-dominant internal row.
\item Reassign the complete cumulant to the star centered at that row
      and apply \Cref{thm:V61-centered-cumulant}.
\item Track the dependence on the dominance threshold so it remains
      exponential per edge.
\item Determine whether this eliminates every use of private moments
      whose order grows with \(m\).
\end{enumerate}
\end{programme}

\begin{firewall}[Claim boundary after fourth-series V62]
\label{fire:V62-claim-boundary}
V62 disproves one proposed all-order private-moment estimate and
redirects the proof to terminal re-fibration.  It does not prove that
re-fibration, the all-thermal local card, all-order MTA, WUFC, CPM,
cutoff removal, reflection positivity, Osterwalder--Schrader
reconstruction, construction of the continuum Yang--Mills measure,
or a uniform positive continuum spectral gap.  The full Yang--Mills
mass gap has not been proved.
\end{firewall}

\paragraph{Parent-version record.}
V62 was formed from
\texttt{Renewal\_Yang\_Mills\_Mass\_Gap\_Research\_Notes\_V61.tex}.
The V61 parent had \(28\,321\) logical source lines,
\(987\,930\) bytes, and SHA--256
\[
 \texttt{7040cdfce1b68f750efaf5ab84f178cc77cd78e5119a58db0d0956f6bddfbc93}.
\]
The next compulsory companion audit is V65.

% =============================================================================
% END APPENDED FOURTH-SERIES V62 MATERIAL
% =============================================================================

% =============================================================================
% BEGIN APPENDED FOURTH-SERIES V63 MATERIAL
% =============================================================================

\chapter{V63: All-order output and branching re-fibration}
\label{chap:V63-refibration}

\section{Selected \(m\)-junction data}
\label{sec:V63-data}

Fix \(m\ge3\), one selected coordinate common to all \(m\) rows, and
a labelled tree \(\tau\) with degrees \(d_i\).  Let
\[
 a_i:=1+C_2(R_{i,r}),\quad
 \Xi_i=a_i+P_i,\quad
 A:=\sum_i a_i,\quad
 Y:=\sum_iP_i,\quad
 B_\tau:=\prod_i\Xi_i^{d_i-1}.
\]
Then
\begin{equation}
 \sum_i(d_i-1)=m-2.
 \label{eq:V63-excess}
\end{equation}
Let \(q_P\) be the exact product of all private prefix multipliers and
let \(\Xi_T\) be the final output mass.

\begin{lemma}[Arity-tracked output formation]
\label{lem:V63-output-formation}
There is \(K_{\rm F}\) such that
\begin{equation}
 \Xi_T\le K_{\rm F}^{\,m}(A+Y).
 \label{eq:V63-output-formation}
\end{equation}
\end{lemma}

\begin{proof}
Private-coordinate outputs contribute \(Y\).  At the selected
coordinate, repeated fixed-group Casimir fusion bounds the final mass
by at most an exponential-per-input constant times
\(\sum_i a_i\).  Iterating the binary fusion inequality \(m-1\)
times gives \(K_{\rm F}^mA\).  Absorb the private part into the same
constant.
\end{proof}

Choose fixed \(L>2\) and set
\begin{equation}
 D_m:=\max\{(4K_{\rm F})^m,L^{m-2}\}.
 \label{eq:V63-threshold}
\end{equation}
Thus \(D_m\le K_D^m\) for a fixed \(K_D\).

\section{Output-dominant re-fibration}
\label{sec:V63-output}

\begin{lemma}[Output dominance makes the largest root a star center]
\label{lem:V63-output-root}
If \(\Xi_T>D_mY\), choose \(i_*\) with
\(a_{i_*}=M:=\max_i a_i\).  Then
\begin{equation}
 Y\le K_1^{-m}A,\qquad
 \Xi_{i_*}\le C M,
 \label{eq:V63-output-root}
\end{equation}
after increasing the fixed base in \(D_m\).  Consequently the star
\(\tau_{i_*}\) satisfies
\begin{equation}
 \mathcal T_{\tau_{i_*}}
 =\left(\prod_ja_j\right)
  \left({a_{i_*}\over\Xi_{i_*}}\right)^{m-2}
 \ge K_2^{-m}\prod_ja_j.
 \label{eq:V63-output-star-floor}
\end{equation}
\end{lemma}

\begin{proof}
Combine
\eqref{eq:V63-output-formation} with
\(Y<\Xi_T/D_m\) and absorb the resulting
\(K_{\rm F}^m\Xi_T/D_m\) term.  This gives
\(\Xi_T\le2K_{\rm F}^mA\) and then
\(Y\le2K_{\rm F}^mA/D_m\).  Since \(A\le mM\) and
\(m\le2^m\), the chosen exponential threshold makes \(Y\le CM\).
Thus \(\Xi_{i_*}\le CM\).  Substitute in the star target.
\end{proof}

\begin{theorem}[All-order output-dominant selected gate]
\label{thm:V63-output-gate}
On \(\Xi_T>D_mY\), the complete paid selected \(m\)-junction obeys
its marked-tree bound with constant
\begin{equation}
 m!K_{\rm out}^{\,m}.
 \label{eq:V63-output-gate}
\end{equation}
\end{theorem}

\begin{proof}
Assign the complete cumulant to the star centered at \(i_*\).
\Cref{thm:V61-centered-cumulant} gives
\[
 \|K_m\|_{\rm op}
 \le m!K_{\rm h}^m\prod_{j\ne i_*}c_j.
\]
On low final output the payer is at most \(C/s_\alpha\).  Using
\(c_j\le\sqrt{s_\alpha a_j}\),
\[
 {C\over s_\alpha}\|K_m\|
 \le m!K^m
 s_\alpha^{(m-3)/2}
 \prod_{j\ne i_*}a_j^{1/2}
 \le m!K^m\prod_ja_j.
\]
On high final output the payer is at most \(C\Xi_T\), and
\Cref{lem:V63-output-root} gives
\(\Xi_T\le K^mA\le K_3^mM\).  Since \(c_j\le1\),
\[
 C\Xi_T\|K_m\|
 \le m!K_4^mM
 \le m!K_4^m\prod_ja_j.
\]
Use the star floor
\eqref{eq:V63-output-star-floor}, \(q_P\le1\), exact Fourier
normalization, and raw pinching.
\end{proof}

\section{Branching-dominant re-fibration}
\label{sec:V63-branching}

\begin{lemma}[All-order weighted pigeonhole]
\label{lem:V63-pigeonhole}
If
\begin{equation}
 B_\tau>D_mY^{m-2},
 \label{eq:V63-branching}
\end{equation}
then some internal vertex \(i\), \(d_i\ge2\), obeys
\begin{equation}
 \Xi_i>D_m^{1/(m-2)}Y\ge LY,
 \qquad
 a_i>(L-1)Y.
 \label{eq:V63-root-dominant}
\end{equation}
\end{lemma}

\begin{proof}
Put \(e_i=d_i-1\).  These exponents are nonnegative and sum to
\(m-2\) by \eqref{eq:V63-excess}.  If every \(e_i>0\) vertex had
\(\Xi_i\le D_m^{1/(m-2)}Y\), their weighted product would be at most
\(D_mY^{m-2}\), contradicting
\eqref{eq:V63-branching}.  Finally
\(a_i=\Xi_i-P_i\) and \(P_i\le Y\).
\end{proof}

\begin{lemma}[Root-dominant all-order star floor]
\label{lem:V63-root-star}
For the row in \Cref{lem:V63-pigeonhole},
\begin{equation}
 \mathcal T_{\tau_i}
 =\left(\prod_ja_j\right)
  \left({a_i\over\Xi_i}\right)^{m-2}
 \ge K_L^{-m}\prod_ja_j.
 \label{eq:V63-root-star}
\end{equation}
\end{lemma}

\begin{proof}
Since \(a_i>(L-1)Y\) and \(P_i\le Y\),
\[
 {\Xi_i\over a_i}\le1+{1\over L-1}.
\]
Raise the reciprocal to power \(m-2\).
\end{proof}

\begin{theorem}[All-order branching-dominant selected gate]
\label{thm:V63-branching-gate}
Suppose
\[
 \Xi_T\le D_mY,\qquad B_\tau>D_mY^{m-2}.
\]
Then the complete paid selected \(m\)-junction obeys its marked-tree
bound with constant \(m!K_{\rm br}^{\,m}\).
\end{theorem}

\begin{proof}
Reassign the complete cumulant to the star centered at the row from
\Cref{lem:V63-pigeonhole}.  The low-output estimate is identical to
the first half of \Cref{thm:V63-output-gate}.  On high output,
\[
 {\Xi_T\over1-q_T}\le C\Xi_T
 \le CD_mY
 \le {CD_m\over L-1}a_i.
\]
The centered displacement estimate and \(c_j\le1\le a_j\) give
\[
 {\Xi_T\over1-q_T}\,q_P\|K_m\|
 \le m!K^mD_m\prod_ja_j
 \le m!K_1^m\prod_ja_j.
\]
Apply \eqref{eq:V63-root-star}.  Dimension cancellation and raw
pinching are unchanged.
\end{proof}

\section{The reduced all-order residual}
\label{sec:V63-residual}

\begin{theorem}[All-order selected-junction trichotomy]
\label{thm:V63-trichotomy}
At every arity \(m\ge3\), the output-dominant and
branching-dominant selected-junction branches satisfy the desired
\(m!K^m\) marked-tree bound.  The only branch not closed by
re-fibration is
\begin{equation}
 \boxed{
 \Xi_T\le D_mY,\qquad
 B_\tau\le D_mY^{m-2}.}
 \label{eq:V63-private-residual}
\end{equation}
\end{theorem}

\begin{proof}
The output-dominant branch is
\Cref{thm:V63-output-gate}.  On its complement, branching dominance
is \Cref{thm:V63-branching-gate}.  The simultaneous complement is
\eqref{eq:V63-private-residual}.
\end{proof}

\begin{assessment}[What re-fibration did and did not remove]
\label{ass:V63-residual}
V63 eliminates every growing private moment from the output- and
branching-dominant branches.  The residual
\eqref{eq:V63-private-residual} is precisely the branch on which the
old compiler charges \(m-2\) or \(m-1\) powers to the private
multiplier.  V62 proves that this charge cannot have the desired
factorial-exponential constant by absolute Casimir moments.  Further
progress requires either a multirow first-moment allocation or the
all-thermal connected tree expansion; repeating the old private
compiler is circular.
\end{assessment}

\section{V63 ledger and V64 programme}
\label{sec:V63-results}

\begin{theorem}[Fourth-series V63 result ledger]
\label{thm:V63-results}
V63 proves an exponential-per-edge output threshold, an all-order
weighted-pigeonhole lemma, and \(m!K^m\) selected-junction bounds on
both the output-dominant and branching-dominant branches.  It reduces
the all-order terminal problem to
\eqref{eq:V63-private-residual}.
\end{theorem}

\begin{proof}
Use
\Cref{lem:V63-output-formation,
thm:V63-output-gate,
lem:V63-pigeonhole,
thm:V63-branching-gate,
thm:V63-trichotomy}.
\end{proof}

\begin{programme}[V64: multirow private allocation]
\label{prog:V63-V64}
\begin{enumerate}[leftmargin=3em]
\item Split \(q_P=\prod_iq_{P_i}\) and retain rowwise private masses
      \(P_i\), rather than the aggregate \(Y\).
\item Choose the star center minimizing
      \(\Xi_i/a_i\).
\item Test whether first private moments on distinct leaf rows pay the
      center's \(m-2\) excess without a high moment on one coordinate.
\item Treat the case of only one nonempty private row separately; it
      should be made a leaf by re-fibration.
\item If all choices force one row to pay growing degree, construct
      the exact \(SU(2)\) counterpacket to the present MTA
      normalization.
\end{enumerate}
\end{programme}

\begin{firewall}[Claim boundary after fourth-series V63]
\label{fire:V63-claim-boundary}
V63 proves all-order selected-junction control outside the
private-dominant residual.  It does not prove that residual, the
all-thermal local card, all-order MTA, WUFC, CPM, cutoff removal,
reflection positivity, Osterwalder--Schrader reconstruction,
construction of the continuum Yang--Mills measure, or a uniform
positive continuum spectral gap.  The full Yang--Mills mass gap has
not been proved.
\end{firewall}

\paragraph{Parent-version record.}
V63 was formed from
\texttt{Renewal\_Yang\_Mills\_Mass\_Gap\_Research\_Notes\_V62.tex}.
The V62 parent had \(28\,570\) logical source lines,
\(996\,290\) bytes, and SHA--256
\[
 \texttt{c9642b20e450f9fbb7383558803a97962793475bd425af5c0fc56b9fc5479516}.
\]
The next compulsory companion audit is V65.

% =============================================================================
% END APPENDED FOURTH-SERIES V63 MATERIAL
% =============================================================================

% =============================================================================
% BEGIN APPENDED FOURTH-SERIES V64 MATERIAL
% =============================================================================

\chapter{V64: Prüfer resummation and rowwise private debits}
\label{chap:V64-rowwise}

\section{The summed tree budget}
\label{sec:V64-Pruefer}

Retain the selected \(m\)-junction notation of V63 and split the
private bank row by row:
\[
 P_i:=\sum_{f\in\mathcal P_i}(1+C_2(R_f)),
 \qquad
 q_i:=\prod_{f\in\mathcal P_i}\eta_{R_f}(\alpha),
 \qquad
 q_P=\prod_{i=1}^mq_i.
\]
An empty row has \(P_i=0\) and \(q_i=1\).  Put
\[
 s:=s_\alpha,\qquad
 t_i:=sa_i,\qquad p_i:=sP_i,\qquad
 x_i:={a_i\over\Xi_i}={t_i\over t_i+p_i},
\]
and set
\[
 H:=\sum_i\sqrt{t_i},\qquad S:=\sum_ix_i.
\]
As before \(0<s\le1\) and \(a_i\ge1\), hence \(t_i\ge s\).

\begin{lemma}[Exact Prüfer collapse of both tree polynomials]
\label{lem:V64-Pruefer}
For labelled trees on \([m]\),
\begin{align}
 \sum_{\tau\in\mathfrak T_m}\mathcal T_\tau
 &=\left(\prod_i a_i\right)S^{m-2},
 \label{eq:V64-target-Pruefer}\\
 \sum_{\tau\in\mathfrak T_m}
   \prod_i t_i^{d_i(\tau)/2}
 &=\left(\prod_i\sqrt{t_i}\right)H^{m-2}.
 \label{eq:V64-local-Pruefer}
\end{align}
\end{lemma}

\begin{proof}
The Prüfer word of a labelled tree has length \(m-2\), and label
\(i\) occurs exactly \(d_i(\tau)-1\) times.  Therefore, for arbitrary
commuting \(z_i\),
\[
 \sum_{\tau\in\mathfrak T_m}\prod_i z_i^{d_i(\tau)-1}
 =\left(\sum_i z_i\right)^{m-2}.
\]
Since
\[
 \mathcal T_\tau
 =\left(\prod_i a_i\right)
  \prod_i x_i^{d_i(\tau)-1},
\]
the choice \(z_i=x_i\) proves
\eqref{eq:V64-target-Pruefer}.  Factoring one
\(\sqrt{t_i}\) from each vertex and taking \(z_i=\sqrt{t_i}\)
proves \eqref{eq:V64-local-Pruefer}.
\end{proof}

\begin{remark}[Why the sum, rather than a fixed tree, matters]
\label{rem:V64-summed-budget}
The marked-tree atlas asks for the sum over labelled trees.  It is
therefore legitimate to assign a complete cumulant to any collection
of tree fibres whose sum is exact.  A coefficientwise private debit
for a prescribed high-degree tree is stronger than necessary and is
exactly what created the false growing one-row moment exposed in V62.
\end{remark}

\section{Only first private moments}
\label{sec:V64-first-moments}

\begin{lemma}[Uniform rowwise first debit]
\label{lem:V64-row-first}
There is \(C_1<\infty\), independent of \(m,i\), the support, and the
representations, such that
\begin{equation}
 p_iq_i\le C_1,\qquad 0\le q_i\le1.
 \label{eq:V64-row-first}
\end{equation}
The aggregate bank also satisfies
\begin{equation}
 sYq_P\le C_1.
 \label{eq:V64-aggregate-first}
\end{equation}
\end{lemma}

\begin{proof}
Apply the fixed first-order case of
\texttt{cor:V52-private-products-all-orders} separately to
\(\mathcal P_i\), and once to their disjoint union.  V62 invalidates
the claimed arity growth of higher moments; it does not affect this
fixed first-order statement.
\end{proof}

\begin{lemma}[Fractionalization of a first debit]
\label{lem:V64-fractional}
If \(0<x_i\le1/2\), then, for every \(0\le\gamma\le1\),
\begin{equation}
 q_i\le C_2
 \left({x_i\over t_i}\right)^\gamma .
 \label{eq:V64-fractional}
\end{equation}
\end{lemma}

\begin{proof}
The identity
\[
 p_i=t_i{x_i^{-1}-1}
\]
and \(x_i\le1/2\) give \(p_i\ge t_i/(2x_i)\).  Hence
\eqref{eq:V64-row-first} gives
\[
 q_i\le\min\left\{1,{2C_1x_i\over t_i}\right\}.
\]
For \(u\ge0\) and \(0\le\gamma\le1\),
\(\min\{1,u\}\le u^\gamma\).  Absorb
\((2C_1)^\gamma\) into \(C_2\).
\end{proof}

\section{Two unconditional residual closures}
\label{sec:V64-easy}

\begin{theorem}[A private-free row closes the residual]
\label{thm:V64-free-row}
Suppose the V63 private-dominant residual
\eqref{eq:V63-private-residual} holds and \(P_h=0\) for some row
\(h\).  Then the complete paid selected junction obeys the desired
marked-tree bound with constant \(m!K_{\rm free}^m\).
\end{theorem}

\begin{proof}
Assign the complete cumulant to the star centered at \(h\).  Since
\(x_h=1\),
\[
 \mathcal T_{\tau_h}=\prod_i a_i.
\]
On low output, \Cref{thm:V61-centered-cumulant} gives
\[
 {C\over s}q_P\|K_m\|
 \le m!K^m s^{(m-3)/2}
       \prod_{j\ne h}a_j^{1/2}
 \le m!K^m\prod_i a_i.
\]
On high output, the residual and
\eqref{eq:V64-aggregate-first} give
\[
 \Xi_Tq_P\le D_mYq_P\le {C_1D_m\over s}.
\]
The same centered-cumulant estimate therefore gives
\[
 C\Xi_Tq_P\|K_m\|
 \le m!K^mD_m s^{(m-3)/2}
       \prod_{j\ne h}a_j^{1/2}
 \le m!K_1^m\prod_i a_i,
\]
because \(D_m\le K_D^m\).  Exact Fourier normalization and raw
pinching are unchanged.
\end{proof}

\begin{theorem}[A weak private ratio closes the residual]
\label{thm:V64-weak-row}
The conclusion of \Cref{thm:V64-free-row} remains true if some row
\(h\) merely satisfies
\begin{equation}
 P_h\le a_h,\qquad\text{equivalently}\qquad x_h\ge{1\over2}.
 \label{eq:V64-weak-row}
\end{equation}
\end{theorem}

\begin{proof}
Use the star centered at \(h\).  Its target obeys
\[
 \mathcal T_{\tau_h}
 =\left(\prod_i a_i\right)x_h^{m-2}
 \ge2^{-(m-2)}\prod_i a_i.
\]
The two estimates in the proof of \Cref{thm:V64-free-row} did not use
\(P_h=0\), except to identify the star target with
\(\prod_i a_i\).  The displayed exponential loss is absorbed into
the per-edge base.
\end{proof}

\begin{corollary}[Exact hard rowwise residual]
\label{cor:V64-hard-rowwise}
After V63 and \Cref{thm:V64-weak-row}, the only unclosed
private-dominant selected fibres satisfy
\begin{equation}
 P_i>a_i\quad\hbox{and}\quad 0<x_i<1/2
 \qquad(1\le i\le m).
 \label{eq:V64-hard-rowwise}
\end{equation}
Thus every remaining row has a nonempty genuinely dominant private
bank.
\end{corollary}

\section{Balanced thermal closure conditional on the local card}
\label{sec:V64-balanced}

\begin{definition}[All-thermal local tree card]
\label{def:V64-thermal-card}
The all-thermal local card at arity \(m\) is an exact decomposition
\[
 K_m=\sum_{\tau\in\mathfrak T_m}K_{m,\tau},
\qquad t_i<1,
\]
such that
\begin{equation}
 \|K_{m,\tau}\|_{\rm op}
 \le m!K_0^m\prod_i t_i^{d_i(\tau)/2}.
 \label{eq:V64-thermal-card}
\end{equation}
\end{definition}

\begin{theorem}[Balanced all-thermal private closure]
\label{thm:V64-balanced}
Fix \(B<\infty\).  Assume
\eqref{eq:V64-hard-rowwise}, the local card
\eqref{eq:V64-thermal-card}, and
\begin{equation}
 t_{\max}\le B\,t_{\min}.
 \label{eq:V64-balance}
\end{equation}
Then both final-resolvent branches satisfy the summed marked-tree
bound with constant
\begin{equation}
 m!K_B^m,
 \label{eq:V64-balanced-conclusion}
\end{equation}
using no private moment of order greater than one.
\end{theorem}

\begin{proof}
Summing \eqref{eq:V64-thermal-card} and using
\eqref{eq:V64-local-Pruefer} bounds the local numerator by
\[
 m!K_0^m\left(\prod_i\sqrt{t_i}\right)H^{m-2}.
\]
The total target is
\[
 \left(\prod_i a_i\right)S^{m-2}
 =s^{-m}\left(\prod_i t_i\right)S^{m-2}.
\]

On low output, put \(\gamma=(m-2)/m\).  By
\Cref{lem:V64-fractional} and the arithmetic--geometric mean
inequality,
\[
 q_P\le C_2^m\prod_i\left({x_i\over t_i}\right)^\gamma,
 \qquad
 \prod_i x_i^\gamma
 \le\left({S\over m}\right)^{m-2}.
\]
After dividing the paid local bound by the total target, the scalar
ratio is at most
\[
 C^m
 {s^{m-1}(H/m)^{m-2}
  \over\prod_i t_i^{1/2+\gamma}}.
\]
Now \(H/m\le\sqrt{t_{\max}}\), while
\[
 \sum_i(1/2+\gamma)={3m\over2}-2.
\]
Since \(s\le t_{\min}\) and
\(t_{\max}\le Bt_{\min}\), the last ratio is at most
\[
 C^m
 {t_{\min}^{m-1}(Bt_{\min})^{(m-2)/2}
  \over t_{\min}^{3m/2-2}}
 \le (C\sqrt B)^m.
\]

On high output,
\(\Xi_T\le D_mY=D_ms^{-1}\sum_kp_k\).  For each summand use
\(p_kq_k\le C_1\), and put
\(\delta=(m-2)/(m-1)\) on the other \(m-1\) rows.  Then
\[
 p_kq_P
 \le C^m\prod_{i\ne k}
       \left({x_i\over t_i}\right)^\delta,
\qquad
 \prod_{i\ne k}x_i^\delta
 \le\left({S\over m-1}\right)^{m-2}.
\]
After summing \(k\), dividing by the total target, and using
\(H/(m-1)\le(3/2)\sqrt{t_{\max}}\) for \(m\ge3\), the ratio is at
most
\[
 C^mD_m
 {s^{m-1}t_{\max}^{(m-2)/2}
  \over t_{\min}^{m/2+m-2}}.
\]
The powers of \(t_{\min}\) cancel exactly as in the low-output
calculation.  The remaining factors \(D_m\), \(B^{(m-2)/2}\), the
sum over \(k\), and \(3/2\) are exponential per input.  This proves
\eqref{eq:V64-balanced-conclusion}.  Every use of a private row was
either fractionalized from \eqref{eq:V64-row-first} or was the single
factor \(p_kq_k\).
\end{proof}

\begin{assessment}[What remains after rowwise resummation]
\label{ass:V64-residual}
The false absolute moment route has now been removed from three
substantial sectors: a private-free row, a weak private row, and the
balanced all-thermal hard-row sector once its local thermal card is
available.  The unresolved selected-junction problem is the union of
\begin{enumerate}[label=\textup{(\roman*)},leftmargin=3em]
\item the proof of the all-thermal local card itself;
\item hard-row fibres containing a thermally hot selected row; and
\item all-thermal hard-row fibres with unbounded anisotropy
      \(t_{\max}/t_{\min}\).
\end{enumerate}
The last two cases must be attacked jointly: the large selected scale
which ruins the balanced comparison is also the natural center for
the V61 displacement estimate.
\end{assessment}

\section{V64 ledger and V65 programme}
\label{sec:V64-results}

\begin{theorem}[Fourth-series V64 result ledger]
\label{thm:V64-results}
V64 proves the exact Prüfer resummation of the local and target tree
polynomials, rowwise fractional first debits, unconditional closure
when any private ratio is weak, and conditional balanced all-thermal
closure without a growing private moment.
\end{theorem}

\begin{proof}
Use
\Cref{lem:V64-Pruefer,
lem:V64-row-first,
lem:V64-fractional,
thm:V64-free-row,
thm:V64-weak-row,
thm:V64-balanced}.
\end{proof}

\begin{programme}[V65: companion audit and anisotropic interpolation]
\label{prog:V64-V65}
\begin{enumerate}[leftmargin=3em]
\item Perform the compulsory read-only audit of the newest active
      SM--Einstein and Navier--Stokes notes.
\item Do not import a scalar or fluid gap into the gauge problem;
      extract only a type-compatible multiscale or interpolation
      mechanism.
\item Interpolate between the V61 centered displacement bound and the
      still-missing all-thermal connected tree card at the largest
      \(t_i\).
\item Test whether each factor of \(t_{\max}/t_{\min}\) can be paired
      with a distinct rowwise first debit.
\item If not, formulate the sharp anisotropic local lemma which must
      replace \eqref{eq:V64-thermal-card}.
\end{enumerate}
\end{programme}

\begin{firewall}[Claim boundary after fourth-series V64]
\label{fire:V64-claim-boundary}
V64 does not prove the all-thermal local tree card, the complete
all-order selected gate, all-order MTA, WUFC, CPM, cutoff removal,
reflection positivity, Osterwalder--Schrader reconstruction,
construction of the continuum Yang--Mills measure, or a uniform
positive continuum spectral gap.  The full Yang--Mills mass gap has
not been proved.
\end{firewall}

\paragraph{Parent-version record.}
V64 was formed from
\texttt{Renewal\_Yang\_Mills\_Mass\_Gap\_Research\_Notes\_V63.tex}.
The V63 parent had \(28\,872\) logical source lines,
\(1\,005\,459\) bytes, and SHA--256
\[
 \texttt{81974809305222c9c362bb25bb8675e2874ce12044f1ceeaee3f0db491ecb3cf}.
\]
The next compulsory companion audit is V65.

% =============================================================================
% END APPENDED FOURTH-SERIES V64 MATERIAL
% =============================================================================

% =============================================================================
% BEGIN APPENDED FOURTH-SERIES V65 MATERIAL
% =============================================================================

\chapter{V65: Companion audit and the anisotropic debit boundary}
\label{chap:V65-audit}

\section{Compulsory read-only companion audit}
\label{sec:V65-audit}

The active companion files inspected at this checkpoint were
\begin{center}
\begin{tabular}{p{0.54\textwidth}rr}
\toprule
file&logical lines&bytes\\
\midrule
\texttt{Renewal\_SM\_Einstein\_Continuum\_Research\_Notes\_V57.tex}
&\(22\,794\)&\(775\,204\)\\
\texttt{Renewal\_NS\_Stability\_Research\_Notes\_V31.tex}
&\(23\,052\)&\(887\,537\)\\
\bottomrule
\end{tabular}
\end{center}
Their SHA--256 digests are respectively
\[
\begin{split}
 &\texttt{b940f9d619395540f22227fca90ccc03306dd92f4ff83f213f1606e239010e78},\\
 &\texttt{f1121b62238ad5491bb7cf03635ea9934942edf573ed8faaa9ee79e1d38a6696}.
\end{split}
\]
Neither companion file was modified.

SM V55 gives an exact change between retained and conditional radial
charts by reconstructing the same joint density in both charts.  It
also proves that a collapsed conditional is justified only when its
leakage divided by its contraction gap tends to zero.  SM V56 retains
the zero-Higgs orbit-type stratum at which the principal gauge chart
degenerates.  SM V57 proves the usual one-generation local anomaly
cancellations, but explicitly separates them from determinant
positivity and from a cutoff-uniform average-phase lower bound.

NS V31 proves an exact dyadic first-moment identity and a
scale-matched heat-formation ladder.  Summing the geometric parent
ancestry before estimating removes shell multiplicity, but its
universal bound returns the enstrophy-square continuation stock.  The
note therefore retains the stopping-time/source correlation instead
of declaring closure from a correct combinatorial summation alone.

\begin{assessment}[Type-compatible transfer to the YM problem]
\label{ass:V65-audit-transfer}
Three structural lessons transfer, but no numerical companion
estimate does:
\begin{enumerate}[label=\textup{(\roman*)},leftmargin=3em]
\item a re-fibration must be an exact identity for the same operator,
      not a comparison of unrelated absolute majorants;
\item an anisotropic or singular stratum must be retained until its
      own estimate is proved; and
\item exact summation can eliminate multiplicity without generating
      the missing analytic smallness.
\end{enumerate}
Accordingly V64's Prüfer identity removes the labelled-tree census,
but it cannot be used as a substitute for a Wilson-tail or connected
thermal estimate.
\end{assessment}

\begin{firewall}[No cross-program gap transfer]
\label{fire:V65-no-transfer}
The SM radial contraction, anomaly sums, determinant phase, and the NS
heat and Leray estimates do not imply a Wilson cumulant bound, WUFC,
or a Yang--Mills spectral gap.  They are used only to determine the
order in which exact re-fibration, summation, and analytic estimates
must be assembled.
\end{firewall}

\section{First moments alone cannot remove thermal anisotropy}
\label{sec:V65-countermodel}

The next proposition concerns the scalar hypotheses used in V64.  It
does not assign the artificial multipliers below to actual Wilson
representations.

\begin{proposition}[Sharp countermodel to a first-debit-only proof]
\label{prop:V65-first-only-fails}
Fix \(m\ge6\).  For every sufficiently small
\(\varepsilon>0\), there are data
\[
 s\le t_i\le1,\qquad
 p_i=t_i(x_i^{-1}-1),\qquad
 0<x_i<1/2,\qquad
 0<q_i\le1,
\]
such that \(p_iq_i\le1\) for every row, but
\begin{equation}
 \mathcal R_m:=
 {q_Ps^{m-1}H^{m-2}
  \over
  \left(\prod_i\sqrt{t_i}\right)S^{m-2}}
 \longrightarrow\infty
 \qquad(\varepsilon\downarrow0).
 \label{eq:V65-counter-ratio}
\end{equation}
Thus the rowwise first-debit axioms and the thermal tree polynomial
alone do not give an anisotropy-uniform scalar compiler.
\end{proposition}

\begin{proof}
Take
\[
 s=\varepsilon,\qquad
 t_1=1,\qquad t_i=\varepsilon\ (i\ge2),\qquad
 x_i=\varepsilon\ (1\le i\le m).
\]
Then
\[
 p_1={1-\varepsilon\over\varepsilon},
 \qquad
 p_i=1-\varepsilon\quad(i\ge2).
\]
Set
\[
 q_1={\varepsilon\over1-\varepsilon},
 \qquad q_i=1\quad(i\ge2).
\]
All the asserted scalar constraints hold, with
\(p_iq_i=1\) for the first row and
\(p_iq_i<1\) on the remaining rows.  Moreover
\[
 q_P={\varepsilon\over1-\varepsilon},\qquad
 H=1+(m-1)\sqrt\varepsilon,\qquad
 S=m\varepsilon,\qquad
 \prod_i\sqrt{t_i}=\varepsilon^{(m-1)/2}.
\]
Substitution into \eqref{eq:V65-counter-ratio} gives
\[
 \mathcal R_m
 =
 {H^{m-2}\over(1-\varepsilon)m^{m-2}}\,
 \varepsilon^{(5-m)/2}.
\]
This diverges for \(m>5\); for \(m=6\) it is already of order
\(\varepsilon^{-1/2}\).
\end{proof}

\begin{corollary}[The missing datum is a genuine Wilson tail]
\label{cor:V65-tail-needed}
Removing the balance hypothesis
\eqref{eq:V64-balance} requires information stronger than
\(p_iq_i\le C\).  In particular, on a row with
\(t_i\gg s\) and \(p_i/t_i\gg1\), the actual multiplier must retain
its representation tail rather than being replaced by its first
moment.
\end{corollary}

\begin{proof}
The countermodel satisfies every algebraic identity and every
first-debit hypothesis used in the V64 balanced proof.  Its only
feature not certified as Wilson-realizable is the value of \(q_1\).
Therefore any successful extension must use a property which
distinguishes an actual Wilson multiplier from that abstract value.
\end{proof}

\begin{firewall}[The countermodel is methodological, not physical]
\label{fire:V65-countermodel-scope}
\Cref{prop:V65-first-only-fails} does not disprove the Yang--Mills
selected gate.  Indeed, in the displayed anisotropic scaling an
actual \(SU(2)\) private representation lies far into the Bessel
tail, and its multiplier is expected to be much smaller than the
artificial \(q_1\).  The proposition proves only that this stronger
tail must be retained explicitly.
\end{firewall}

\section{The exact upstream target}
\label{sec:V65-target}

For \(r\ge0\), define
\begin{equation}
 h(r):=r\,\operatorname{arsinh}r-\sqrt{1+r^2}+1.
 \label{eq:V65-rate}
\end{equation}
This is nonnegative, increasing, quadratic at the origin, and of
order \(r\log r\) at infinity.  The exact \(SU(2)\) formula in V40
suggests the one-link estimate
\begin{equation}
 {I_n(\alpha)\over I_1(\alpha)}
 \le
 \exp\{-c\alpha h(n/\alpha)\}.
 \label{eq:V65-desired-tail}
\end{equation}
Unlike a power moment in the Casimir, this bound remains adapted to
both the parabolic range \(n\lesssim\alpha\) and the factorial Bessel
tail \(n\gg\alpha\).

\begin{programme}[V66: contour-shift Wilson tail]
\label{prog:V65-V66}
\begin{enumerate}[leftmargin=3em]
\item Starting from the exact Bessel integral, shift the angular
      contour by \(-iu\) and prove
      \[
      I_n(\alpha)
      \le I_0(\alpha)
      \exp\{-nu+\alpha(\cosh u-1)\}.
      \]
\item Optimize at \(u=\operatorname{arsinh}(n/\alpha)\).
\item Absorb \(I_0/I_1\) without a per-coordinate constant, using the
      one-link gap in the small-deviation regime and the contour rate
      in the large-deviation regime.
\item Express the result in the scaled private mass
      \(p=s_\alpha(1+C_2)\), retaining the \(s_\alpha\)-dependent
      crossover rather than replacing it by a false polynomial
      moment.
\item Test the resulting tail on the exact scalar packet of
      \Cref{prop:V65-first-only-fails}.
\end{enumerate}
\end{programme}

\section{V65 ledger}
\label{sec:V65-results}

\begin{theorem}[Fourth-series V65 result ledger]
\label{thm:V65-results}
V65 completes the scheduled companion audit, proves that rowwise
first moments alone cannot close the anisotropic scalar compiler, and
identifies the exact uniform Bessel large-deviation estimate which
must be acquired next.
\end{theorem}

\begin{proof}
Use
\Cref{ass:V65-audit-transfer,
prop:V65-first-only-fails,
cor:V65-tail-needed}.
\end{proof}

\begin{firewall}[Claim boundary after fourth-series V65]
\label{fire:V65-claim-boundary}
V65 proves neither the proposed Wilson tail nor the all-thermal local
tree card.  It does not prove the complete all-order selected gate,
all-order MTA, WUFC, CPM, cutoff removal, reflection positivity,
Osterwalder--Schrader reconstruction, construction of the continuum
Yang--Mills measure, or a uniform positive continuum spectral gap.
The full Yang--Mills mass gap has not been proved.
\end{firewall}

\paragraph{Parent-version record.}
V65 was formed from
\texttt{Renewal\_Yang\_Mills\_Mass\_Gap\_Research\_Notes\_V64.tex}.
The V64 parent had \(29\,258\) logical source lines,
\(1\,017\,021\) bytes, and SHA--256
\[
 \texttt{23726429cdb87a07efa11e01896597b90c9a99a793ff0f413051e114c85aedde}.
\]
The next compulsory companion audit is V70.

% =============================================================================
% END APPENDED FOURTH-SERIES V65 MATERIAL
% =============================================================================

% =============================================================================
% BEGIN APPENDED FOURTH-SERIES V66 MATERIAL
% =============================================================================

\chapter{V66: A uniform \(SU(2)\) Wilson large-deviation debit}
\label{chap:V66-tail}

\section{The contour-shift rate}
\label{sec:V66-contour}

Retain
\[
 h(r)=r\,\operatorname{arsinh}r-\sqrt{1+r^2}+1.
\]
Then
\begin{equation}
 h'(r)=\operatorname{arsinh}r,\qquad
 h''(r)={1\over\sqrt{1+r^2}},\qquad h(0)=h'(0)=0.
 \label{eq:V66-h-derivatives}
\end{equation}
In particular \(h\) is nonnegative and increasing.

\begin{lemma}[Exact Bessel contour bound]
\label{lem:V66-contour}
For every \(\alpha>0\), integer \(n\ge0\), and \(u\ge0\),
\begin{equation}
 {I_n(\alpha)\over I_0(\alpha)}
 \le
 \exp\{-nu+\alpha(\cosh u-1)\}.
 \label{eq:V66-contour}
\end{equation}
Consequently,
\begin{equation}
 {I_n(\alpha)\over I_0(\alpha)}
 \le
 \exp\{-\alpha h(n/\alpha)\}.
 \label{eq:V66-optimized-contour}
\end{equation}
\end{lemma}

\begin{proof}
Use the periodic entire integrand representation
\[
 I_n(\alpha)
 ={1\over2\pi}\int_{-\pi}^{\pi}
 e^{\alpha\cos z}e^{-inz}\,dz.
\]
Shift the contour to \(z=\theta-iu\).  The two vertical sides cancel
because \(n\) is an integer.  On the shifted segment,
\[
 |e^{\alpha\cos(\theta-iu)}e^{-in(\theta-iu)}|
 =
 e^{-nu}e^{\alpha\cos\theta\cosh u}.
\]
It follows that
\[
 I_n(\alpha)\le e^{-nu}I_0(\alpha\cosh u).
\]
Moreover
\[
 {d\over d\beta}\log I_0(\beta)
 ={I_1(\beta)\over I_0(\beta)}\le1.
\]
Integrating from \(\alpha\) to \(\alpha\cosh u\) proves
\eqref{eq:V66-contour}.  Its exponent is minimized when
\(\sinh u=n/\alpha\).  Substitution of
\(u=\operatorname{arsinh}(n/\alpha)\) gives
\eqref{eq:V66-optimized-contour}.
\end{proof}

\begin{lemma}[Removing the \(I_0/I_1\) normalization]
\label{lem:V66-normalization}
There are \(c>0\) and \(\alpha_0<\infty\) such that, for every
\(\alpha\ge\alpha_0\) and integer \(n\ge2\),
\begin{equation}
 {I_n(\alpha)\over I_1(\alpha)}
 \le
 \exp\{-c\alpha h(n/\alpha)\}.
 \label{eq:V66-normalized-tail}
\end{equation}
\end{lemma}

\begin{proof}
The Bessel recurrence and positivity give
\[
 {I_0(\alpha)\over I_1(\alpha)}
 ={I_2(\alpha)\over I_1(\alpha)}+{2\over\alpha}
 \le1+{2\over\alpha}\le e^{2/\alpha}.
\]
Put \(E=\alpha h(n/\alpha)\).  If \(E\ge4/\alpha\), then
\Cref{lem:V66-contour} yields
\[
 {I_n\over I_1}
 \le e^{-E+2/\alpha}\le e^{-E/2}.
\]

It remains to consider \(E<4/\alpha\).  From
\eqref{eq:V66-h-derivatives},
\[
 h(r)=\int_0^r\operatorname{arsinh}v\,dv
 \le {r^2\over2},
\]
so \(E\le n^2/(2\alpha)\).  The established one-link gap and
Wilson-to-heat calibration give, for the spin
\(j=(n-1)/2\),
\[
 1-\eta_j(\alpha)
 \ge c_0\min\left\{1,{n^2\over\alpha}\right\}.
\]
Because \(0<\eta_j\le1\),
\[
 \eta_j\le\exp\{-(1-\eta_j)\}.
\]
If \(n^2/\alpha\le1\), the gap exponent is at least a fixed
multiple of \(E\).  If \(n^2/\alpha>1\), then the gap exponent is
fixed while \(E<4/\alpha\le1\), after increasing \(\alpha_0\).
Thus \(\eta_j\le e^{-c_1E}\) in the remaining case.  Combine the two
cases and use the exact identity
\(\eta_j=I_n/I_1\).
\end{proof}

\begin{remark}[Compatibility with the V62 obstruction]
\label{rem:V66-V62}
The exponent in \eqref{eq:V66-normalized-tail} is quadratic,
\(n^2/\alpha\), when \(n\ll\alpha\), but is only of order
\(n\log(n/\alpha)\) when \(n\gg\alpha\).  Hence the theorem does not
imply the false factorial-exponential moments of V62.  It retains
precisely the Bessel crossover which those moments discarded.
\end{remark}

\section{Scaled-mass form and private-bank aggregation}
\label{sec:V66-scaled}

Let \(a_j=1+C_2(j)\), \(p_j=s_\alpha a_j\), and define, for a fixed
sufficiently small \(c_*>0\),
\begin{equation}
 \Psi_s(p):={1\over s}
 h\!\left(c_*\sqrt{sp}\right).
 \label{eq:V66-Psi}
\end{equation}

\begin{lemma}[Elementary properties of the scaled rate]
\label{lem:V66-Psi-properties}
For fixed \(s>0\), \(\Psi_s\) is nonnegative, increasing, concave on
\([0,\infty)\), and \(\Psi_s(0)=0\).  Consequently,
\begin{equation}
 \sum_{f\in F}\Psi_s(p_f)
 \ge\Psi_s\left(\sum_{f\in F}p_f\right)
 \label{eq:V66-subadditive-rate}
\end{equation}
for every finite family \(p_f\ge0\).
\end{lemma}

\begin{proof}
Only concavity needs proof.  With
\(z=c_*\sqrt{sp}\), differentiation gives a positive constant
multiple of
\[
 {\operatorname{arsinh}z\over z}.
\]
This is decreasing because
\[
 \operatorname{arsinh}z>{z\over\sqrt{1+z^2}}
 \qquad(z>0).
\]
Thus \(\Psi_s'\) is decreasing.  If \(f\) is nonnegative, concave,
and \(f(0)=0\), then
\[
 f(x)\ge{x\over x+y}f(x+y),\qquad
 f(y)\ge{y\over x+y}f(x+y).
\]
Adding and iterating proves
\eqref{eq:V66-subadditive-rate}.
\end{proof}

\begin{theorem}[Uniform scaled \(SU(2)\) Wilson tail]
\label{thm:V66-scaled-tail}
After choosing \(c_*>0\) sufficiently small, there are
\(c>0\) and \(\alpha_0<\infty\) such that every nontrivial spin
\(j\) and every \(\alpha\ge\alpha_0\) satisfy
\begin{equation}
 \boxed{\quad
 \eta_j(\alpha)
 \le\exp\{-c\Psi_{s_\alpha}(p_j)\}.
 \quad}
 \label{eq:V66-scaled-tail}
\end{equation}
\end{theorem}

\begin{proof}
Write \(n=2j+1\).  The calibration
\(s_\alpha=2/\alpha+O(\alpha^{-2})\) makes \(\alpha\) uniformly
comparable with \(s_\alpha^{-1}\).  Also
\[
 a_j=1+{n^2-1\over4}
 \asymp n^2.
\]
Thus \(n/\alpha\) is uniformly comparable with
\(\sqrt{s_\alpha p_j}=s_\alpha\sqrt{a_j}\).
For every fixed \(0<\lambda\le1\), concavity of
\(\operatorname{arsinh}\) gives
\[
 h(\lambda r)
 =\lambda\int_0^r\operatorname{arsinh}(\lambda v)\,dv
 \ge\lambda^2h(r).
\]
Fixed multiplicative changes in the argument and prefactor of \(h\)
can therefore be absorbed into \(c_*\) and the exterior constant.
Apply \Cref{lem:V66-normalization}.
\end{proof}

\begin{corollary}[Support-uniform private-bank rate]
\label{cor:V66-bank-tail}
Let a private row contain finitely many nontrivial \(SU(2)\)
representations, with total scaled mass
\[
 p_i=s_\alpha P_i=\sum_{f\in\mathcal P_i}p_f
\]
and multiplier \(q_i=\prod_f\eta_{R_f}(\alpha)\).  Then
\begin{equation}
 \boxed{\quad
 q_i\le\exp\{-c\Psi_{s_\alpha}(p_i)\}.
 \quad}
 \label{eq:V66-bank-tail}
\end{equation}
No support-count factor occurs.
\end{corollary}

\begin{proof}
Multiply \eqref{eq:V66-scaled-tail} over the private coordinates and
use \eqref{eq:V66-subadditive-rate}.
\end{proof}

\section{Regression on the anisotropic countermodel}
\label{sec:V66-regression}

\begin{proposition}[The actual tail kills the V65 bad scaling]
\label{prop:V66-kills-countermodel}
In the scalar scaling of
\Cref{prop:V65-first-only-fails}, the first row has
\[
 p_1={1-\varepsilon\over\varepsilon},
 \qquad s=\varepsilon.
\]
Any actual \(SU(2)\) private bank with this scaled mass obeys
\begin{equation}
 q_1\le e^{-c/\varepsilon}.
 \label{eq:V66-counter-tail}
\end{equation}
Consequently its contribution to
\(\mathcal R_m\) tends to zero for every fixed \(m\).
\end{proposition}

\begin{proof}
Here \(sp_1=1-\varepsilon\), so
\[
 \Psi_s(p_1)
 ={1\over\varepsilon}
 h(c_*\sqrt{1-\varepsilon})
 \ge {c_1\over\varepsilon}
\]
for small \(\varepsilon\).  Apply
\Cref{cor:V66-bank-tail}.  The remaining factor in the V65
calculation is only a fixed power of \(\varepsilon^{-1}\), which is
dominated by \(e^{-c/\varepsilon}\).
\end{proof}

\begin{assessment}[Scope of the acquired debit]
\label{ass:V66-scope}
V66 supplies the exact information absent from the first-moment
countermodel: a concentrated anisotropic private bank is
super-polynomially suppressed in the regulator scale.  The theorem is
proved for the exact \(SU(2)\) Wilson multiplier.  Extending it to the
chosen general compact gauge group requires the corresponding
highest-weight Fourier-coefficient bound; that extension is not
silently assumed.
\end{assessment}

\section{V66 ledger and V67 programme}
\label{sec:V66-results}

\begin{theorem}[Fourth-series V66 result ledger]
\label{thm:V66-results}
V66 proves a uniform contour large-deviation estimate for exact
\(SU(2)\) Wilson multipliers, rewrites it in the scaled private mass,
aggregates it over an arbitrary private bank without a support factor,
and shows that it eliminates the explicit V65 anisotropic
countermodel.
\end{theorem}

\begin{proof}
Use
\Cref{lem:V66-contour,
lem:V66-normalization,
thm:V66-scaled-tail,
cor:V66-bank-tail,
prop:V66-kills-countermodel}.
\end{proof}

\begin{programme}[V67: anisotropic scalar optimization]
\label{prog:V66-V67}
\begin{enumerate}[leftmargin=3em]
\item Insert \eqref{eq:V66-bank-tail} into the exact ratio
      \eqref{eq:V65-counter-ratio} for arbitrary row data.
\item Split rows at \(s p_i=1\), the intrinsic crossover of
      \(\Psi_s\), rather than at an artificial balance constant.
\item Prove or refute an exponential-per-input supremum for both
      final-resolvent branches.
\item Keep the all-thermal local tree card as an explicit hypothesis;
      the scalar optimization cannot prove that operator identity.
\item Determine which part of the argument extends from \(SU(2)\) to
      a general fixed compact semisimple gauge group.
\end{enumerate}
\end{programme}

\begin{firewall}[Claim boundary after fourth-series V66]
\label{fire:V66-claim-boundary}
V66 proves a private Wilson-tail debit, not the anisotropic scalar
optimization or the all-thermal local tree card.  It does not prove
the complete all-order selected gate, all-order MTA, WUFC, CPM,
cutoff removal, reflection positivity, Osterwalder--Schrader
reconstruction, construction of the continuum Yang--Mills measure,
or a uniform positive continuum spectral gap.  The full Yang--Mills
mass gap has not been proved.
\end{firewall}

\paragraph{Parent-version record.}
V66 was formed from
\texttt{Renewal\_Yang\_Mills\_Mass\_Gap\_Research\_Notes\_V65.tex}.
The V65 parent had \(29\,515\) logical source lines,
\(1\,026\,035\) bytes, and SHA--256
\[
 \texttt{58665c5f01fa47b69f2d9553245e0a93a5cf1ce9eb30721efa7436c752c7e430}.
\]
The next compulsory companion audit is V70.

% =============================================================================
% END APPENDED FOURTH-SERIES V66 MATERIAL
% =============================================================================

% =============================================================================
% BEGIN APPENDED FOURTH-SERIES V67 MATERIAL
% =============================================================================

\chapter{V67: The exact Gaussian tangent marked-tree card}
\label{chap:V67-Gaussian}

\section{Replicated Gaussian interpolation}
\label{sec:V67-replica}

Let \(E\) be a fixed \(q\)-dimensional Euclidean space, and let
\(Z\) be centered Gaussian with covariance \(sI_E\).  For
\(m\) bounded smooth operator-valued functions
\[
 F_i:E\longrightarrow\mathcal B(\mathcal H_i),
\]
write \(\kappa_\otimes(F_1(Z),\ldots,F_m(Z))\) for the tensor
cumulant, with the operators acting on separate tensor legs.

For a labelled tree \(\tau\) and edge parameters
\(\boldsymbol u=(u_e)_{e\in E(\tau)}\in[0,1]^{m-1}\), define
\[
 w_{ii}^\tau(\boldsymbol u)=1,\qquad
 w_{ij}^\tau(\boldsymbol u)
 =\min_{e\in{\rm path}_\tau(i,j)}u_e
 \quad(i\ne j).
\]

\begin{lemma}[Positivity of the tree covariance]
\label{lem:V67-tree-covariance}
The matrix \(w^\tau(\boldsymbol u)\) is positive semidefinite.
Therefore there is a replicated Gaussian vector
\((Z_1,\ldots,Z_m)\) with covariance
\[
 \mathbb E[Z_i\otimes Z_j]
 =s w_{ij}^\tau(\boldsymbol u)I_E.
\]
\end{lemma}

\begin{proof}
For \(r\in[0,1]\), delete from \(\tau\) every edge with
\(u_e<r\), and let \(P_r\) be the matrix whose \(ij\) entry is one
when \(i,j\) lie in the same remaining component and zero otherwise.
Each \(P_r\) is block diagonal with all-one positive semidefinite
blocks.  Moreover
\[
 w^\tau(\boldsymbol u)=\int_0^1P_r\,dr.
\]
Indeed the \(ij\) entry of the integral is the minimum edge parameter
along their unique path.  An integral of positive semidefinite
matrices is positive semidefinite.
\end{proof}

For \(i\ne j\), put
\begin{equation}
 \Delta_{ij}:=
 s\sum_{a=1}^q
 {\partial\over\partial z_i^a}
 {\partial\over\partial z_j^a}.
 \label{eq:V67-edge-operator}
\end{equation}

\begin{theorem}[Gaussian tensor-cumulant tree formula]
\label{thm:V67-Gaussian-tree-formula}
One has the exact identity
\begin{equation}
 \boxed{\quad
 \kappa_\otimes(F_1(Z),\ldots,F_m(Z))
 =
 \sum_{\tau\in\mathfrak T_m}
 \int_{[0,1]^{m-1}}
 \mathbb E_{w^\tau(\boldsymbol u)}
 \left[
 \left(\prod_{ij\in E(\tau)}\Delta_{ij}\right)
 \bigotimes_{k=1}^mF_k(Z_k)
 \right]d\boldsymbol u.
 \quad}
 \label{eq:V67-Gaussian-tree-formula}
\end{equation}
\end{theorem}

\begin{proof}
First take scalar-valued \(F_i\) with bounded derivatives.  If
\(\mathbb E_w\) denotes Gaussian expectation with covariance
\(s(w_{ij}I_E)\), Gaussian differentiation under the integral gives
\begin{equation}
 {\partial\over\partial w_{ij}}
 \mathbb E_w\prod_kF_k(Z_k)
 =
 \mathbb E_w\left[
 \Delta_{ij}\prod_kF_k(Z_k)\right].
 \label{eq:V67-Gaussian-differentiation}
\end{equation}
The factor \(1/2\) in covariance differentiation is cancelled by the
two symmetric off-diagonal entries \(w_{ij},w_{ji}\).

We recall the connected forest interpolation identity and include
its induction.  Start at the diagonal covariance, where all replicas
are independent.  Apply the fundamental theorem of calculus to one
off-diagonal interpolation variable.  The zero endpoint factors over
the two components separated by that edge; the derivative endpoint
joins those components and contributes the corresponding partial
derivative.  Repeat within every remaining component.  A term which
ends with \(k\) components is indexed by a forest with \(k\)
components, and its covariance between two vertices is the minimum
edge parameter on their forest path, or zero if no path exists.
Möbius inversion on the partition lattice cancels every term with
\(k>1\).  The surviving connected forests are precisely the labelled
trees, each edge has been differentiated once, and their interpolation
covariance is \(w^\tau(\boldsymbol u)\).  Thus
\[
 \sum_{\pi\in\Pi_m}(|\pi|-1)!(-1)^{|\pi|-1}
 \prod_{B\in\pi}\mathbb E\prod_{i\in B}F_i(Z)
\]
equals the right side of
\eqref{eq:V67-Gaussian-tree-formula}, by
\eqref{eq:V67-Gaussian-differentiation}.  The left side is the
cumulant.

For finite-dimensional operator values, apply the scalar identity to
every matrix coefficient.  Multilinearity reconstructs the tensor
identity.  Bounded derivatives justify differentiation and Fubini;
the general stated case follows by the usual smooth truncation and
dominated convergence.
\end{proof}

\section{A normed tree theorem}
\label{sec:V67-norm}

\begin{theorem}[Gaussian marked-tree bound]
\label{thm:V67-Gaussian-card}
Suppose that, for every \(k\ge0\), every \(z\in E\), and unit
directions \(v_1,\ldots,v_k\),
\begin{equation}
 \|D^kF_i(z)[v_1,\ldots,v_k]\|_{\rm op}
 \le M_iL^kA_i^{k/2}.
 \label{eq:V67-derivative-hypothesis}
\end{equation}
Then the tree fibre in
\eqref{eq:V67-Gaussian-tree-formula} obeys
\begin{equation}
 \boxed{\quad
 \|K_{m,\tau}^{\rm G}\|_{\rm op}
 \le
 \left(qL^2\right)^{m-1}s^{m-1}
 \left(\prod_iM_i\right)
 \prod_iA_i^{d_i(\tau)/2}.
 \quad}
 \label{eq:V67-Gaussian-card}
\end{equation}
There is no factorial, representation dimension, or interpolation
volume loss.
\end{theorem}

\begin{proof}
Expand the \(q\)-term sum in every edge operator
\(\Delta_{ij}\).  There are \(q^{m-1}\) choices of edge-coordinate
labels.  At vertex \(i\), exactly \(d_i(\tau)\) derivatives hit
\(F_i\).  Apply \eqref{eq:V67-derivative-hypothesis}, use
\(\sum_i d_i(\tau)=2m-2\), take the operator norm inside the Gaussian
expectation, and note that the interpolation cube has volume one.
This gives
\[
 s^{m-1}q^{m-1}L^{2m-2}
 \prod_iM_iA_i^{d_i(\tau)/2},
\]
which is \eqref{eq:V67-Gaussian-card}.
\end{proof}

\section{Application to representation matrices}
\label{sec:V67-representations}

Fix an orthonormal basis \(X_1,\ldots,X_q\) of the compact Lie
algebra and define
\[
 F_R(z):=D^R\!\left(\exp\left(\sum_{a=1}^qz^aX_a\right)\right)
 =\exp\left(\sum_{a=1}^qz^a\,dR(X_a)\right).
\]

\begin{lemma}[Uniform representation derivative bound]
\label{lem:V67-representation-derivatives}
For every irreducible unitary \(R\),
\begin{equation}
 \|D^kF_R(z)[v_1,\ldots,v_k]\|_{\rm op}
 \le
 C_G^k(1+C_2(R))^{k/2}
 \prod_{\ell=1}^k\|v_\ell\|.
 \label{eq:V67-representation-derivatives}
\end{equation}
The constant depends only on the fixed compact group and the chosen
Lie-algebra metric.
\end{lemma}

\begin{proof}
The Casimir identity gives
\[
 \|dR(V)\|_{\rm op}
 \le C_G\sqrt{C_2(R)}\,\|V\|.
\]
For \(A(z)=\sum_az^adR(X_a)\), repeated Duhamel differentiation of
\(e^{A(z)}\) writes the \(k\)-th derivative as the sum over the
\(k!\) orderings of the inserted operators
\(dR(v_1),\ldots,dR(v_k)\), integrated over the ordered \(k\)-simplex.
All intervening exponentials are unitary because \(A(z)\) is
skew-adjoint.  Each ordered integral has volume \(1/k!\), so the
ordering count cancels exactly.  Multiplying the generator bounds and
using \(C_2(R)\le1+C_2(R)\) proves the claim.
\end{proof}

\begin{corollary}[Exact all-order tangent-model local card]
\label{cor:V67-tangent-card}
For the Gaussian tangent law on the Lie algebra and
\(A_i=1+C_2(R_i)\), the complete one-coordinate tensor cumulant has
an exact labelled-tree decomposition satisfying
\begin{equation}
 \|K_{m,\tau}^{\rm G}\|_{\rm op}
 \le
 K_G^{m-1}s^{m-1}
 \prod_iA_i^{d_i(\tau)/2}.
 \label{eq:V67-tangent-card}
\end{equation}
This is the all-thermal local tree card at the stronger
exponential-per-edge scale.
\end{corollary}

\begin{proof}
Apply \Cref{thm:V67-Gaussian-card} with
\(M_i=1\), \(L=C_G\), and
\Cref{lem:V67-representation-derivatives}.
\end{proof}

\section{What separates the Wilson law from its tangent model}
\label{sec:V67-Wilson-remainder}

In normal coordinates, the Wilson one-link law is a Gaussian tangent
density multiplied by the Haar Jacobian, the nonquadratic part of the
plaquette action, a normalizing factor, and a complement outside the
normal chart.  V52 gives fixed-order \(L^1\) derivative bounds for
that density.  Those bounds do not yet provide an all-order
connected interpolation in which every density correction is attached
to the same labelled input tree with an exponential-per-edge
constant.

\begin{assessment}[Exact new bottleneck]
\label{ass:V67-bottleneck}
\Cref{cor:V67-tangent-card} proves that neither Gaussian
connectedness, noncommuting representation matrices, nor arbitrary
arity is the missing local mechanism.  The remaining local problem is
the Wilson-density transfer:
\begin{enumerate}[label=\textup{(\roman*)},leftmargin=3em]
\item insert the nonquadratic density without creating an unattached
      vacuum component;
\item sum all correction insertions before taking operator norms;
\item retain one factor \(sA_i^{1/2}A_j^{1/2}\) on every input-tree
      edge; and
\item keep the correction-series constant exponential in \(m\).
\end{enumerate}
The private anisotropic optimization of V66 remains a separate scalar
step, conditional on this transfer.
\end{assessment}

\begin{firewall}[A tangent theorem is not the Wilson theorem]
\label{fire:V67-not-Wilson}
The actual Wilson measure is not Gaussian in normal coordinates and
has a compact-chart complement.  Therefore
\Cref{cor:V67-tangent-card} cannot be substituted for
\eqref{eq:V64-thermal-card} until the density transfer and complement
are proved uniformly in arity.
\end{firewall}

\section{V67 ledger and V68 programme}
\label{sec:V67-results}

\begin{theorem}[Fourth-series V67 result ledger]
\label{thm:V67-results}
V67 proves the exact replicated-Gaussian tensor-cumulant tree formula,
its operator marked-tree norm bound, uniform derivative bounds for
compact-group representation matrices, and hence the complete
all-order Gaussian tangent local card with no factorial loss.
\end{theorem}

\begin{proof}
Use
\Cref{thm:V67-Gaussian-tree-formula,
thm:V67-Gaussian-card,
lem:V67-representation-derivatives,
cor:V67-tangent-card}.
\end{proof}

\begin{programme}[V68: Wilson-density transfer]
\label{prog:V67-V68}
\begin{enumerate}[leftmargin=3em]
\item Write the exact normal-coordinate Wilson density as the
      Gaussian tangent density times a correction
      \(\rho_s(z)\), with a separately retained chart complement.
\item Introduce correction insertions before applying the connected
      forest formula, so vacuum components cancel by normalization.
\item Prove an exponential generating-function bound for the
      correction blocks; do not multiply a partition factorial by a
      separate derivative factorial.
\item Show that correction vertices contract onto the input tree and
      do not consume an additional final resolvent.
\item Test the complement term against the all-order thermal weights,
      using the representation-displacement floor rather than a
      private Casimir moment.
\end{enumerate}
\end{programme}

\begin{firewall}[Claim boundary after fourth-series V67]
\label{fire:V67-claim-boundary}
V67 proves the tangent Gaussian local card, not the Wilson-density
transfer, the anisotropic scalar optimization, or the complete
all-order selected gate.  It does not prove all-order MTA, WUFC, CPM,
cutoff removal, reflection positivity, Osterwalder--Schrader
reconstruction, construction of the continuum Yang--Mills measure,
or a uniform positive continuum spectral gap.  The full Yang--Mills
mass gap has not been proved.
\end{firewall}

\paragraph{Parent-version record.}
V67 was formed from
\texttt{Renewal\_Yang\_Mills\_Mass\_Gap\_Research\_Notes\_V66.tex}.
The V66 parent had \(29\,855\) logical source lines,
\(1\,036\,201\) bytes, and SHA--256
\[
 \texttt{606b5f33676a3215d1cbf262713ad92e6124c373d63c7d60fff792b63adb9078}.
\]
The next compulsory companion audit is V70.

% =============================================================================
% END APPENDED FOURTH-SERIES V67 MATERIAL
% =============================================================================

% =============================================================================
% BEGIN APPENDED FOURTH-SERIES V68 MATERIAL
% =============================================================================

\chapter{V68: Linked Wilson-density insertions}
\label{chap:V68-linked}

\section{Tensor cumulants under an exponential tilt}
\label{sec:V68-tilt}

Let \(\gamma\) be a probability measure, let \(V\) be a real scalar
function, and define
\[
 d\nu_\lambda
 ={e^{\lambda V}\over Z(\lambda)}\,d\gamma,
 \qquad
 Z(\lambda):=\mathbb E_\gamma e^{\lambda V}.
\]
The functions \(F_i\) remain operator-valued on separate tensor legs.

\begin{theorem}[Exact linked insertion identity]
\label{thm:V68-linked-insertion}
Suppose the following series converges absolutely in operator norm:
\begin{equation}
 \sum_{r=0}^\infty{1\over r!}
 \left\|
 \kappa_{\otimes,\gamma}
 (F_1,\ldots,F_m,\underbrace{V,\ldots,V}_{r})
 \right\|_{\rm op}<\infty.
 \label{eq:V68-insertion-convergence}
\end{equation}
Then
\begin{equation}
 \boxed{\quad
 \kappa_{\otimes,\nu_1}(F_1,\ldots,F_m)
 =
 \sum_{r=0}^\infty{1\over r!}
 \kappa_{\otimes,\gamma}
 (F_1,\ldots,F_m,\underbrace{V,\ldots,V}_{r}).
 \quad}
 \label{eq:V68-linked-insertion}
\end{equation}
Here a scalar \(V\) occupies a new tensor-trivial labelled slot; the
mixed cumulant is defined by the same partition formula.
\end{theorem}

\begin{proof}
Introduce commuting nilpotent variables \(t_i\), \(t_i^2=0\), and
form the finite formal polynomial
\[
 \mathscr Z(\lambda,\boldsymbol t)
 :=
 \mathbb E_\gamma\left[
 e^{\lambda V}
 \bigotimes_{i=1}^m(I+t_iF_i)\right].
\]
Its constant term in the \(t_i\) is the positive scalar
\(Z(\lambda)\), so its formal logarithm in the nilpotent variables is
well defined after factoring \(Z(\lambda)I\).  Direct expansion of
the logarithm shows
\begin{equation}
 [t_1\cdots t_m]\log\mathscr Z(\lambda,\boldsymbol t)
 =
 \kappa_{\otimes,\nu_\lambda}(F_1,\ldots,F_m).
 \label{eq:V68-log-source}
\end{equation}
Indeed a square-free monomial selects a set partition of
\([m]\), and the logarithmic coefficient of a partition with \(k\)
blocks is \((-1)^{k-1}(k-1)!\).

Expand the same formal logarithm jointly at
\((\lambda,\boldsymbol t)=(0,\boldsymbol0)\).  Its coefficient of
\(\lambda^rt_1\cdots t_m/r!\) is exactly the mixed tensor cumulant
with \(r\) labelled copies of \(V\), by the identical partition
calculation.  Evaluate the Taylor series at \(\lambda=1\).
Hypothesis \eqref{eq:V68-insertion-convergence} permits coefficientwise
summation in operator norm and gives
\eqref{eq:V68-linked-insertion}.
\end{proof}

\begin{corollary}[Small bounded tilts satisfy the analytic hypothesis]
\label{cor:V68-small-bounded}
If \(V\) is bounded and
\(\|V\|_\infty<\log2\), then the source logarithm is analytic in a
complex disc of radius greater than one, and
\eqref{eq:V68-linked-insertion} holds.
\end{corollary}

\begin{proof}
Choose \(R>1\) with \(R\|V\|_\infty<\log2\).  On
\(|\lambda|\le R\),
\[
 |Z(\lambda)-1|
 \le\mathbb E|e^{\lambda V}-1|
 \le e^{R\|V\|_\infty}-1<1.
\]
Thus \(Z\) has no zero in the disc.  The nilpotent source logarithm
is analytic there, so its Taylor coefficients converge absolutely at
\(\lambda=1\) by Cauchy's estimate.
\end{proof}

\begin{lemma}[Constant corrections cancel]
\label{lem:V68-constant-cancels}
Replacing \(V\) by \(V+c\) leaves every term with \(r\ge1\) in
\eqref{eq:V68-linked-insertion} unchanged after the full sum is
assembled.
\end{lemma}

\begin{proof}
A cumulant of total order at least two vanishes if any one argument
is constant.  Multilinearity therefore removes every occurrence of
\(c\).  This is the cumulant form of the cancellation between
normalization and vacuum insertions.
\end{proof}

\section{Insertion into the Gaussian forest formula}
\label{sec:V68-forest}

Take \(\gamma\) to be the Gaussian tangent law of V67.  Label the
\(m\) physical inputs by \(1,\ldots,m\) and the \(r\) correction
slots by \(m+1,\ldots,m+r\).

\begin{theorem}[Exact input-plus-correction tree expansion]
\label{thm:V68-correction-trees}
Under \eqref{eq:V68-insertion-convergence},
\begin{align}
 \kappa_{\otimes,\nu_1}(F_1,\ldots,F_m)
 &=
 \sum_{r=0}^\infty{1\over r!}
 \sum_{\tau\in\mathfrak T_{m+r}}
 \int_{[0,1]^{m+r-1}}
 \mathbb E_{w^\tau(\boldsymbol u)}
 \notag\\
 &\qquad\left[
 \left(\prod_{ij\in E(\tau)}\Delta_{ij}\right)
 \left\{
 \bigotimes_{i=1}^mF_i(Z_i)
 \prod_{\ell=m+1}^{m+r}V(Z_\ell)
 \right\}\right]d\boldsymbol u.
 \label{eq:V68-correction-trees}
\end{align}
Every nonzero correction contribution is connected to the physical
input set.  There is no correction-only vacuum component.
\end{theorem}

\begin{proof}
Apply \Cref{thm:V67-Gaussian-tree-formula} to every mixed cumulant
on the right side of \eqref{eq:V68-linked-insertion}.  Absolute
convergence permits the displayed rearrangement.  Each term is a
spanning tree on all \(m+r\) labels, hence every correction label has
a path to every physical input.  A correction-only connected
component would be a separate component of that graph and is
therefore absent.
\end{proof}

\begin{firewall}[Connectivity precedes the norm]
\label{fire:V68-connect-first}
Expanding \(e^V\) in the numerator and its normalizer separately and
then taking absolute values produces vacuum correction diagrams which
are not present in
\eqref{eq:V68-correction-trees}.  The exact logarithmic/cumulant
assembly must occur before operator norms, just as required by the
SM V55 and NS V31 audit lessons.
\end{firewall}

\section{The correction-tree census}
\label{sec:V68-census}

\begin{lemma}[Arbitrary small-activity correction counts are
exponential per input]
\label{lem:V68-correction-census}
If \(m\ge2\) and \(0\le\varepsilon<e^{-2}\), then
\begin{equation}
 \sum_{r=0}^\infty{\varepsilon^r\over r!}
 (m+r)^{m+r-2}
 \le C_\varepsilon^m m^{m-2}
 \label{eq:V68-correction-census}
\end{equation}
for a finite constant \(C_\varepsilon\) independent of \(m\).
\end{lemma}

\begin{proof}
For \(r\ge1\), \(r!\ge(r/e)^r\), while
\[
 \left(1+{m\over r}\right)^r\le e^m,
 \qquad
 \left(1+{r\over m}\right)^{m-2}\le e^r.
\]
Therefore
\begin{align*}
 {(m+r)^{m+r-2}\over r!}
 &=
 m^{m-2}
 \left(1+{r\over m}\right)^{m-2}
 {(m+r)^r\over r!}\\
 &\le
 m^{m-2}e^r
 e^{r+m}
 =
 e^m m^{m-2}e^{2r}.
\end{align*}
The \(r=0\) term is \(m^{m-2}\).  Sum the geometric series in
\(e^2\varepsilon\), and enlarge its fixed prefactor to the \(m\)-th
power.
\end{proof}

\begin{assessment}[What the census does and does not prove]
\label{ass:V68-census}
If every correction vertex can be assigned a genuine analytic
activity \(\varepsilon<e^{-2}\), then the unrestricted number of
Wilson-density insertions does not introduce a second factorial or a
support count.  This is only the census.  One must still prove the
activity from the exact Wilson density and contract every
input-plus-correction tree onto an input marked tree without losing
its endpoint factors.
\end{assessment}

\section{The normal-coordinate transfer datum}
\label{sec:V68-transfer}

On a normal chart the desired exact factorization has the form
\begin{equation}
 d\nu_\alpha(z)
 =
 {e^{V_s(z)}\over\mathbb E_{\gamma_s}e^{V_s}}\,
 d\gamma_s(z),
 \label{eq:V68-normal-factorization}
\end{equation}
with the complement retained separately.  The correction \(V_s\)
contains the nonquadratic Wilson action and the logarithm of the Haar
Jacobian.  A usable transfer certificate must control its mixed
Gaussian derivatives in the edge operators of
\eqref{eq:V68-correction-trees}; a bound on \(\|V_s\|_\infty\) over
an expanding normal chart is neither expected nor required.

\begin{definition}[Wilson correction activity certificate]
\label{def:V68-WCAC}
A WCAC consists of:
\begin{enumerate}
\item the exact factorization
      \eqref{eq:V68-normal-factorization} plus a disjoint chart
      complement;
\item absolute convergence of
      \eqref{eq:V68-correction-trees};
\item a contraction map from every input-plus-correction tree to a
      labelled input tree;
\item a norm estimate whose total correction activity is
      \(\varepsilon<e^{-2}\), uniformly for sufficiently small
      \(s\); and
\item a complement estimate with the same input endpoint weights.
\end{enumerate}
\end{definition}

\begin{assessment}[Upstream status]
\label{ass:V68-upstream}
V68 proves the exact algebraic part of the WCAC: linked insertion,
vacuum cancellation, the connected Gaussian forest expansion, and
the correction-count census.  The two analytic parts not yet proved
are the derivative activity and the endpoint-preserving contraction
to an input tree.
\end{assessment}

\section{V68 ledger and V69 programme}
\label{sec:V68-results}

\begin{theorem}[Fourth-series V68 result ledger]
\label{thm:V68-results}
V68 proves the exact exponential-tilt insertion identity for tensor
cumulants, cancellation of constant and vacuum corrections, the
input-plus-correction Gaussian tree expansion, and an
exponential-per-input census for arbitrary small-activity correction
counts.
\end{theorem}

\begin{proof}
Use
\Cref{thm:V68-linked-insertion,
lem:V68-constant-cancels,
thm:V68-correction-trees,
lem:V68-correction-census}.
\end{proof}

\begin{programme}[V69: correction-tree contraction]
\label{prog:V68-V69}
\begin{enumerate}[leftmargin=3em]
\item Take the minimal subtree spanning the \(m\) physical inputs in
      an input-plus-correction tree and separate its correction
      branches from its Steiner correction vertices.
\item Sum rooted correction branches into vertex activities before
      suppressing degree-two Steiner vertices.
\item Resolve every higher-degree Steiner vertex into an input tree
      by a deterministic nearest-input rule and track the endpoint
      powers.
\item Prove an exponential fibre bound for the number of correction
      trees mapping to one input tree.
\item Only after that combinatorial contraction, insert the
      normal-coordinate derivative estimates for \(V_s\).
\end{enumerate}
\end{programme}

\begin{firewall}[Claim boundary after fourth-series V68]
\label{fire:V68-claim-boundary}
V68 does not prove a WCAC, the Wilson all-thermal local card, the
anisotropic scalar optimization, or the complete all-order selected
gate.  It does not prove all-order MTA, WUFC, CPM, cutoff removal,
reflection positivity, Osterwalder--Schrader reconstruction,
construction of the continuum Yang--Mills measure, or a uniform
positive continuum spectral gap.  The full Yang--Mills mass gap has
not been proved.
\end{firewall}

\paragraph{Parent-version record.}
V68 was formed from
\texttt{Renewal\_Yang\_Mills\_Mass\_Gap\_Research\_Notes\_V67.tex}.
The V67 parent had \(30\,188\) logical source lines,
\(1\,048\,016\) bytes, and SHA--256
\[
 \texttt{11915945f092f6944dcecdcfa07b506e0d5ea9979006f2e9f127352323a2d4af}.
\]
The next compulsory companion audit is V70.

% =============================================================================
% END APPENDED FOURTH-SERIES V68 MATERIAL
% =============================================================================

% =============================================================================
% BEGIN APPENDED FOURTH-SERIES V69 MATERIAL
% =============================================================================

\chapter{V69: Steiner correction vertices and endpoint transfer}
\label{chap:V69-Steiner}

\section{The physical core of a correction tree}
\label{sec:V69-core}

Let \(T\) be a tree whose vertex set is the disjoint union of
\(m\ge2\) labelled physical vertices \(P=[m]\) and finitely many
correction vertices.  Its \emph{physical core} \(C(T)\) is the union
of all simple paths joining pairs of physical vertices.  Equivalently,
one obtains \(C(T)\) by repeatedly pruning correction leaves.  Let
\[
 \delta_i:=\deg_{C(T)}(i)\quad(i\in P)
\]
and let \(S(T)\) be the correction vertices which remain in the core,
with core degrees \(k_v\).

\begin{lemma}[Core degree facts]
\label{lem:V69-core-degrees}
Every \(\delta_i\ge1\), every \(k_v\ge2\), and
\begin{equation}
 \boxed{\quad
 D(T):=\sum_{v\in S(T)}(k_v-2)
 =2m-2-\sum_{i=1}^m\delta_i.
 \quad}
 \label{eq:V69-deficit}
\end{equation}
\end{lemma}

\begin{proof}
Since \(m\ge2\), each physical vertex lies on a path to another
physical vertex and therefore has positive core degree.  A correction
leaf would have been pruned, so \(k_v\ge2\).  If
\(r_c=|S(T)|\), then the core is a tree on \(m+r_c\) vertices and has
\(m+r_c-1\) edges.  The handshaking identity gives
\[
 \sum_i\delta_i+\sum_vk_v=2(m+r_c-1).
\]
Subtract \(2r_c\) and rearrange to obtain
\eqref{eq:V69-deficit}.
\end{proof}

\begin{lemma}[Resolution to a physical labelled tree]
\label{lem:V69-resolution}
There is a labelled tree \(\sigma(T)\) on \([m]\) and nonnegative
integers \(n_i\) such that
\begin{equation}
 d_i(\sigma(T))=\delta_i+n_i,
 \qquad
 \sum_i n_i=D(T).
 \label{eq:V69-resolution}
\end{equation}
One may make the choice deterministic.
\end{lemma}

\begin{proof}
Fix the least physical label, say \(1\), and set
\[
 n_1=D(T),\qquad n_i=0\quad(i>1).
\]
By \eqref{eq:V69-deficit},
\[
 \sum_i(\delta_i+n_i)=2m-2.
\]
Every proposed degree is positive.  Also the first degree is at most
\(m-1\), because the other \(m-1\) degrees are at least one.
The Prüfer degree theorem says that positive integers \(d_i\) are the
degrees of a labelled tree exactly when their sum is \(2m-2\): use a
Prüfer word in which label \(i\) occurs \(d_i-1\) times.  Choosing the
lexicographically least such word makes \(\sigma(T)\) deterministic.
\end{proof}

\begin{remark}[Meaning of the deficit]
\label{rem:V69-deficit}
A degree-two correction vertex merely subdivides a physical
connection and contributes no deficit.  A \(k\)-valent Steiner
correction vertex saves \(k-2\) physical endpoint incidences relative
to an ordinary tree on the physical labels.  Formula
\eqref{eq:V69-deficit} says there are no other missing endpoint
powers.
\end{remark}

\section{Thermal endpoint transfer}
\label{sec:V69-transfer}

Put
\[
 b_i:=\sqrt{sA_i}.
\]
In the all-thermal sector,
\begin{equation}
 \sqrt s\le b_i<1,
 \label{eq:V69-thermal-range}
\end{equation}
because \(A_i\ge1\).

\begin{theorem}[Steiner deficit pays the missing endpoint powers]
\label{thm:V69-Steiner-transfer}
For every physical core \(C(T)\), let
\[
 W(T):=
 \left(\prod_i b_i^{\delta_i}\right)
 \prod_{v\in S(T)}s^{(k_v-2)/2}.
 \label{eq:V69-core-weight}
\]
Then the resolved physical tree of
\Cref{lem:V69-resolution} satisfies
\begin{equation}
 \boxed{\quad
 W(T)\le
 \prod_i b_i^{d_i(\sigma(T))}.
 \quad}
 \label{eq:V69-endpoint-transfer}
\end{equation}
\end{theorem}

\begin{proof}
By \eqref{eq:V69-deficit},
\[
 \prod_{v\in S(T)}s^{(k_v-2)/2}
 =s^{D(T)/2}.
\]
The deterministic resolution used
\(n_1=D(T)\).  Since \(b_1\ge\sqrt s\),
\[
 s^{D(T)/2}\le b_1^{D(T)}
 =\prod_i b_i^{n_i}.
\]
Multiply by \(\prod_i b_i^{\delta_i}\) and use
\eqref{eq:V69-resolution}.
\end{proof}

\begin{corollary}[Extra pruned incidences are harmless when thermal]
\label{cor:V69-pruned}
Any physical--correction edges removed while pruning to \(C(T)\)
contribute additional physical factors \(b_i\le1\).  They may be
discarded in an upper bound.  Thus
\eqref{eq:V69-endpoint-transfer} remains valid after arbitrary
correction decorations have first been summed into their attachment
activities.
\end{corollary}

\begin{proof}
Every derivative edge incident to a physical vertex contributes one
factor \(b_i\).  An incidence outside the core only multiplies the
left side by \(b_i\le1\).  Correction-only factors are, by hypothesis,
included in the rooted decoration activity.
\end{proof}

\section{The exact analytic scale requested from the Wilson density}
\label{sec:V69-analytic-scale}

Use standard Gaussian coordinates \(y=z/\sqrt s\).  Physical
representation derivatives then cost \(b_i\).  A core correction
vertex of degree \(k\) must supply the factor
\begin{equation}
 s^{(k-2)/2}.
 \label{eq:V69-required-correction-scale}
\end{equation}
This scale is forced by the graph identity, not guessed from power
counting.

\begin{definition}[Steiner-scaled correction activity]
\label{def:V69-SSCA}
An SSCA with activity \(\varepsilon\) is a bound on every connected
normal-coordinate correction vertex which, after summing Lie-algebra
edge labels and Gaussian polynomial moments, has the form
\begin{equation}
 \mathcal V_k
 \le
 \varepsilon K^k s^{(k-2)/2}
 \qquad(k\ge2).
 \label{eq:V69-SSCA}
\end{equation}
Degree-zero pieces are removed by
\Cref{lem:V68-constant-cancels}; degree-one pieces belong to rooted
decorations and must be summed before the core is formed.
\end{definition}

\begin{theorem}[Conditional endpoint-preserving Wilson transfer]
\label{thm:V69-conditional-transfer}
Assume:
\begin{enumerate}[label=\textup{(\roman*)},leftmargin=3em]
\item the exact linked expansion
      \eqref{eq:V68-correction-trees};
\item an SSCA \eqref{eq:V69-SSCA};
\item a support-uniform rooted-decoration sum of size at most
      \(\varepsilon\); and
\item an exponential-per-input bound on the fibres of the map
      \(T\mapsto\sigma(T)\).
\end{enumerate}
Then every normal-chart all-thermal Wilson cumulant satisfies the
input marked-tree endpoint powers
\[
 \prod_i(sA_i)^{d_i(\sigma)/2}
\]
with an exponential-per-input constant.
\end{theorem}

\begin{proof}
Prune correction decorations and use the third hypothesis to charge
their activities.  On the remaining core, multiply the SSCA factors.
Their powers of \(s\) are exactly
\[
 \prod_vs^{(k_v-2)/2}=s^{D(T)/2}.
\]
\Cref{thm:V69-Steiner-transfer} converts this and the original
physical incidences into the endpoint weight of
\(\sigma(T)\).  The factors \(K^{k_v}\) are exponential in the total
edge count and are included in the declared fibre bound.  Finally
sum all trees in each fibre and then all labelled input trees.  No
private multiplier or final resolvent is used in this local
operation.
\end{proof}

\begin{assessment}[Why the SSCA matches the Taylor scaling]
\label{ass:V69-Taylor}
After the quadratic Gaussian part is removed, an order-\(k\)
derivative of the scaled Wilson action formally carries
\(\alpha s^{k/2}\asymp s^{(k-2)/2}\).  The quartic leading correction
has size \(s\), exactly the \(k=4\) case.  The Haar-Jacobian
correction starts at order \(s\) and is better than required at
\(k=2\).  This verifies the dimensional match, but not the uniform
Gaussian polynomial summation, rooted-decoration bound, or chart
complement.  Those analytic estimates remain to be proved.
\end{assessment}

\begin{assessment}[Combinatorial obstruction removed]
\label{ass:V69-combinatorial}
V68 left open how a branching correction vertex could be suppressed
without losing an input endpoint factor.  V69 removes that
combinatorial obstruction exactly.  The missing number of endpoints
is the Steiner excess \(D(T)\), and the natural derivative scale
supplies precisely \(s^{D(T)/2}\).  What remains is to prove the SSCA
and its fibre sum for the actual Wilson correction.
\end{assessment}

\section{V69 ledger and V70 programme}
\label{sec:V69-results}

\begin{theorem}[Fourth-series V69 result ledger]
\label{thm:V69-results}
V69 proves the Steiner endpoint-deficit identity, deterministic
resolution to a physical labelled tree, and exact transfer of the
natural correction derivative powers into all missing physical
endpoint weights.  It reduces the normal-chart Wilson transfer to
the analytic SSCA, rooted decorations, and an exponential fibre
bound.
\end{theorem}

\begin{proof}
Use
\Cref{lem:V69-core-degrees,
lem:V69-resolution,
thm:V69-Steiner-transfer,
thm:V69-conditional-transfer}.
\end{proof}

\begin{programme}[V70: audit and SSCA acquisition]
\label{prog:V69-V70}
\begin{enumerate}[leftmargin=3em]
\item Perform the compulsory read-only audit of the newest active
      SM--Einstein and Navier--Stokes notes.
\item Write the exact scaled \(SU(2)\) Wilson action and Haar
      Jacobian in the normal coordinate used by V40.
\item Prove Gaussian \(L^1\) derivative bounds at the SSCA scale
      \eqref{eq:V69-SSCA}, summing derivative partitions by one
      exponential generating function.
\item Sum degree-one rooted correction decorations before taking
      norms.
\item Keep the complement outside the normal chart as a separate
      stratum; do not infer its endpoint weights from exponential
      smallness alone.
\end{enumerate}
\end{programme}

\begin{firewall}[Claim boundary after fourth-series V69]
\label{fire:V69-claim-boundary}
V69 proves the endpoint-transfer combinatorics, not an SSCA, the
Wilson all-thermal local card, the anisotropic scalar optimization,
or the complete all-order selected gate.  It does not prove
all-order MTA, WUFC, CPM, cutoff removal, reflection positivity,
Osterwalder--Schrader reconstruction, construction of the continuum
Yang--Mills measure, or a uniform positive continuum spectral gap.
The full Yang--Mills mass gap has not been proved.
\end{firewall}

\paragraph{Parent-version record.}
V69 was formed from
\texttt{Renewal\_Yang\_Mills\_Mass\_Gap\_Research\_Notes\_V68.tex}.
The V68 parent had \(30\,513\) logical source lines,
\(1\,059\,030\) bytes, and SHA--256
\[
 \texttt{4040ad2bdb7706bda1e9a9efd1b771e270a245dce50816500ecafb0811fc05d2}.
\]
The next compulsory companion audit is V70.

% =============================================================================
% END APPENDED FOURTH-SERIES V69 MATERIAL
% =============================================================================

% =============================================================================
% BEGIN APPENDED FOURTH-SERIES V70 MATERIAL
% =============================================================================

\chapter{V70: Companion audit and the Wilson chart-scale obstruction}
\label{chap:V70-chart}

\section{Compulsory read-only companion audit}
\label{sec:V70-audit}

The active companion files inspected at this checkpoint were
\begin{center}
\begin{tabular}{p{0.54\textwidth}rr}
\toprule
file&logical lines&bytes\\
\midrule
\texttt{Renewal\_SM\_Einstein\_Continuum\_Research\_Notes\_V59.tex}
&\(23\,676\)&\(801\,839\)\\
\texttt{Renewal\_NS\_Stability\_Research\_Notes\_V31.tex}
&\(23\,052\)&\(887\,537\)\\
\bottomrule
\end{tabular}
\end{center}
Their SHA--256 digests are
\[
\begin{split}
 &\texttt{a651ab0907125be3e078e33d5ab020bc48c131ef5e147e399adc38200140310c},\\
 &\texttt{f1121b62238ad5491bb7cf03635ea9934942edf573ed8faaa9ee79e1d38a6696}.
\end{split}
\]
Neither companion file was modified.

SM V58 constructs an exact finite-cutoff Grassmann Schur complement
and proves that small or zero modes must be moved into the retained
sector before the inverse is used.  SM V59 proves a nonnormal
Combes--Thomas estimate from a singular-value gap and shows that
pointwise exponential decay is summable only when its exponent beats
the rooted graph growth.  It makes the loss of that inequality an
explicit re-fibration threshold.  NS V31 is unchanged; its exact
geometric ancestry sum removes multiplicity but returns a
continuation-class analytic stock.

\begin{assessment}[V70 audit transfer]
\label{ass:V70-audit-transfer}
The type-compatible lesson is an acquisition criterion, not a bound.
For YM, the correction activity must beat the correction-tree growth
in \Cref{lem:V68-correction-census}, and an expansion may be inverted
or normalized only on the exact density block for which its
nonvanishing and convergence have been proved.  If a chart loses that
inequality, the complement must be retained as a separate stratum.
\end{assessment}

\begin{firewall}[No imported locality or regularity estimate]
\label{fire:V70-no-import}
The fermionic singular gap, Combes--Thomas exponent, Grassmann norm,
and NS heat-formation stock have different carriers from the Wilson
one-link density.  None is imported numerically into the SSCA, WUFC,
or a Yang--Mills gap.
\end{firewall}

\section{Exact \(SU(2)\) normal-coordinate density}
\label{sec:V70-density}

Write \(U=\exp(z^aX_a)\), \(r=|z|\), on the exponential ball
\(0\le r<\pi\), with the normalization used in V40.  Haar measure has
Jacobian
\[
 J(z)=\left({\sin r\over r}\right)^2,
\]
and the Wilson factor is \(e^{\alpha\cos r}\).  Put
\[
 \epsilon:=\alpha^{-1},\qquad z=\sqrt\epsilon\,y.
\]

\begin{proposition}[Exact Gaussian factorization on the exponential
ball]
\label{prop:V70-exact-density}
Up to its scalar normalizer, the rescaled \(SU(2)\) Wilson density is
\begin{equation}
 \mathbf1_{\{|y|<\pi/\sqrt\epsilon\}}
 e^{-|y|^2/2}e^{V_\epsilon(y)}\,dy,
 \label{eq:V70-exact-density}
\end{equation}
where
\begin{align}
 V_\epsilon(y)
 &:=
 {1\over\epsilon}
 \left(
 \cos(\sqrt\epsilon|y|)-1
 +{\epsilon|y|^2\over2}
 \right)
 \notag\\
 &\qquad
 +2\log\left(
 {\sin(\sqrt\epsilon|y|)
  \over\sqrt\epsilon|y|}
 \right).
 \label{eq:V70-correction}
\end{align}
The value at \(y=0\) is defined by continuity.
\end{proposition}

\begin{proof}
In exponential coordinates,
\[
 e^{\alpha\cos r}J(z)\,dz
 =
 e^\alpha
 e^{-r^2/(2\epsilon)}
 \exp\left\{
 {\cos r-1+r^2/2\over\epsilon}
 +2\log{\sin r\over r}
 \right\}dz.
\]
Substitute \(z=\sqrt\epsilon y\).  The Jacobian
\(\epsilon^{3/2}\), the factor \(e^\alpha\), and the Gaussian
normalization are scalar and enter the final normalizer.  The
remaining exponent is exactly \eqref{eq:V70-correction}.
\end{proof}

\begin{lemma}[Fixed physical charts are not small tilts]
\label{lem:V70-fixed-chart}
For every sufficiently small fixed \(0<\delta<1\), there are
\(c_\delta>0\) and \(C_\delta<\infty\) such that
\begin{equation}
 \sup_{|y|\le\delta/\sqrt\epsilon}
 V_\epsilon(y)
 \ge {c_\delta\over\epsilon}-C_\delta.
 \label{eq:V70-fixed-chart}
\end{equation}
Thus the small-bounded-tilt criterion of
\Cref{cor:V68-small-bounded} cannot hold uniformly on a fixed
normal-coordinate ball.
\end{lemma}

\begin{proof}
Choose \(r_0=\delta/2\).  Taylor's theorem gives
\[
 \cos r_0-1+{r_0^2\over2}\ge c r_0^4
\]
for small fixed \(\delta\).  Evaluate
\eqref{eq:V70-correction} at
\(|y|=r_0/\sqrt\epsilon\).  The Haar logarithm is a finite constant
depending only on \(r_0\), while the first term is at least
\(cr_0^4/\epsilon\).
\end{proof}

\section{Mesoscopic charts and their exact tradeoff}
\label{sec:V70-mesoscopic}

\begin{lemma}[Small correction on a mesoscopic Gaussian ball]
\label{lem:V70-mesoscopic-small}
There are fixed \(c,C>0\) such that, whenever
\(\sqrt\epsilon R\le c\),
\begin{equation}
 \sup_{|y|\le R}|V_\epsilon(y)|
 \le C\left(\epsilon R^4+\epsilon R^2\right).
 \label{eq:V70-mesoscopic-small}
\end{equation}
In particular, if \(R=o(\epsilon^{-1/4})\), the correction tends
uniformly to zero on the ball.
\end{lemma}

\begin{proof}
For \(|r|\le c\),
\[
 |\cos r-1+r^2/2|\le Cr^4,
 \qquad
 \left|\log{\sin r\over r}\right|\le Cr^2.
\]
Set \(r=\sqrt\epsilon|y|\) in
\eqref{eq:V70-correction} and take the supremum.
\end{proof}

\begin{lemma}[Wilson mass outside a mesoscopic ball]
\label{lem:V70-mesoscopic-tail}
After decreasing the fixed chart radius if necessary, there are
\(c,C>0\) such that
\begin{equation}
 \nu_\alpha\{|y|>R\}
 \le Ce^{-cR^2}+Ce^{-c/\epsilon}
 \label{eq:V70-mesoscopic-tail}
\end{equation}
for \(1\le R\le c/\sqrt\epsilon\).
\end{lemma}

\begin{proof}
On the fixed inner exponential ball,
\[
 \cos r-1\le-c_0r^2,
\]
and the Haar Jacobian is bounded above.  The Wilson normalizer has the
standard lower Laplace bound \(c\epsilon^{3/2}e^{1/\epsilon}\), obtained
by integrating over \(|z|\le\sqrt\epsilon\).  Rescaling therefore
bounds the inner tail by a fixed Gaussian tail
\(Ce^{-cR^2}\).  On the complement of the fixed inner ball,
\(\cos r\le1-c_1\); comparison with the same normalizer gives
\(Ce^{-c_1/\epsilon}\).
\end{proof}

\begin{proposition}[One mesoscopic cutoff cannot prove the all-order
endpoint estimate]
\label{prop:V70-cutoff-no-go}
Suppose one tries to prove all arities by:
\begin{enumerate}[label=\textup{(\roman*)},leftmargin=3em]
\item treating \(|y|\le R\) as a uniformly small bounded tilt via
      \eqref{eq:V70-mesoscopic-small}; and
\item using only \eqref{eq:V70-mesoscopic-tail} to supply \(m\)
      missing thermal half-powers on the complement.
\end{enumerate}
No choice \(R=R(\epsilon)\) independent of \(m\) can make both steps
uniform for all \(m\).
\end{proposition}

\begin{proof}
Because every nontrivial thermal endpoint has
\(b_i\ge c\sqrt\epsilon\), paying \(m\) such endpoint factors from a
bare complement estimate requires, up to exponential-per-input
constants,
\[
 e^{-cR^2}\lesssim \epsilon^{m/2},
\]
and hence
\begin{equation}
 R^2\gtrsim m\log(1/\epsilon).
 \label{eq:V70-tail-required-radius}
\end{equation}
Uniform smallness of the inner tilt requires
\[
 \epsilon R^4=o(1).
\]
Combining with \eqref{eq:V70-tail-required-radius} forces
\[
 \epsilon m^2\log^2(1/\epsilon)=o(1),
\]
which fails when \(m\) is unbounded at fixed regulator.  Thus a
single arity-independent mesoscopic cutoff cannot establish both
claims.
\end{proof}

\begin{firewall}[Scope of the cutoff no-go]
\label{fire:V70-cutoff-scope}
\Cref{prop:V70-cutoff-no-go} does not disprove the Wilson
all-thermal tree card.  It rules out one proof architecture: uniform
small bounded tilt on one mesoscopic ball plus a bare tail estimate.
An exact full-density interpolation, an arity-adaptive decomposition
whose constants are summed before taking the supremum, or a direct
Brascamp--Lieb/Stein identity may still supply the card.
\end{firewall}

\section{Revised upstream target}
\label{sec:V70-target}

\begin{assessment}[Why V69's dimensional match is not yet an activity]
\label{ass:V70-SSCA-status}
The action derivatives have the formal scale
\(\epsilon^{(k-2)/2}\), exactly matching the Steiner deficit of V69.
However, that power is consumed when it is converted into missing
physical endpoint factors.  It cannot automatically be counted a
second time as a small correction-vertex activity.  V68's
small-activity census therefore requires a sharper linked estimate,
not merely Taylor power counting.
\end{assessment}

\begin{programme}[V71: nonperturbative Wilson interpolation]
\label{prog:V70-V71}
\begin{enumerate}[leftmargin=3em]
\item Differentiate the exact normalized Wilson expectations in
      \(\alpha\) or along a convex interpolation of the full action,
      rather than expanding \(e^{V_\epsilon}\) into unrestricted
      correction vertices.
\item Seek an integration-by-parts identity in which normalization
      cancels the radial vacuum term and every derivative joins two
      existing physical components.
\item Preserve the V69 Steiner deficit powers, but do not also charge
      them as correction activities.
\item Test the identity first on scalar class functions and then on
      tensor representation matrices.
\item Keep the chart complement inside the exact interpolation until
      its endpoint factors are visible.
\end{enumerate}
\end{programme}

\section{V70 ledger}
\label{sec:V70-results}

\begin{theorem}[Fourth-series V70 result ledger]
\label{thm:V70-results}
V70 completes the scheduled companion audit, derives the exact
Gaussian-times-correction formula for the \(SU(2)\) Wilson density,
proves the fixed-chart large-correction obstruction, quantifies the
mesoscopic smallness/tail tradeoff, and rules out a single-cutoff
small-tilt proof of the all-order endpoint estimate.
\end{theorem}

\begin{proof}
Use
\Cref{ass:V70-audit-transfer,
prop:V70-exact-density,
lem:V70-fixed-chart,
lem:V70-mesoscopic-small,
lem:V70-mesoscopic-tail,
prop:V70-cutoff-no-go}.
\end{proof}

\begin{firewall}[Claim boundary after fourth-series V70]
\label{fire:V70-claim-boundary}
V70 does not prove a nonperturbative Wilson interpolation, the Wilson
all-thermal local card, the anisotropic scalar optimization, or the
complete all-order selected gate.  It does not prove all-order MTA,
WUFC, CPM, cutoff removal, reflection positivity,
Osterwalder--Schrader reconstruction, construction of the continuum
Yang--Mills measure, or a uniform positive continuum spectral gap.
The full Yang--Mills mass gap has not been proved.
\end{firewall}

\paragraph{Parent-version record.}
V70 was formed from
\texttt{Renewal\_Yang\_Mills\_Mass\_Gap\_Research\_Notes\_V69.tex}.
The V69 parent had \(30\,818\) logical source lines,
\(1\,069\,437\) bytes, and SHA--256
\[
 \texttt{e910729e52e1709b6c540527b14a632045d463fdd31e78dd65e5e0ed8b94051c}.
\]
The next compulsory companion audit is V75.

% =============================================================================
% END APPENDED FOURTH-SERIES V70 MATERIAL
% =============================================================================

% =============================================================================
% BEGIN APPENDED FOURTH-SERIES V71 MATERIAL
% =============================================================================

\chapter{V71: Exact radial transport of the full \(SU(2)\) Wilson law}
\label{chap:V71-transport}

\section{The canonical radial quantile map}
\label{sec:V71-quantile}

Let \(\gamma_3\) be standard Gaussian measure on
\(\mathbb R^3\).  Its radial density and distribution function are
\[
 g(r)=\sqrt{2\over\pi}\,r^2e^{-r^2/2},
 \qquad
 G(r)=\int_0^rg(u)\,du.
\]
The radial part of the central \(SU(2)\) Wilson law is
\[
 w_\alpha(\theta)
 ={e^{\alpha\cos\theta}\sin^2\theta\over Z_\alpha^{\rm rad}},
 \qquad
 W_\alpha(\theta)=\int_0^\theta w_\alpha(u)\,du,
 \qquad0\le\theta\le\pi.
\]
Both angular directions are uniform on \(S^2\).

\begin{definition}[Wilson radial transport]
\label{def:V71-transport}
Define
\begin{equation}
 T_\alpha(r):=W_\alpha^{-1}(G(r)),
 \qquad0\le r<\infty,
 \label{eq:V71-quantile}
\end{equation}
and
\begin{equation}
 \Phi_\alpha(y):=
 \begin{cases}
 \exp\!\left(T_\alpha(|y|){y\over|y|}\right),&y\ne0,\\
 I,&y=0.
 \end{cases}
 \label{eq:V71-group-map}
\end{equation}
The Lie-algebra metric and exponential normalization are those of
V70.
\end{definition}

\begin{theorem}[Exact full-law pushforward]
\label{thm:V71-pushforward}
The map \(\Phi_\alpha\) pushes \(\gamma_3\) exactly to the
\(SU(2)\) Wilson one-link probability:
\begin{equation}
 (\Phi_\alpha)_\#\gamma_3=\nu_\alpha.
 \label{eq:V71-pushforward}
\end{equation}
No normal-chart complement or density normalizer remains.
\end{theorem}

\begin{proof}
The functions \(G\) and \(W_\alpha\) are continuous and strictly
increasing on the interiors of their domains, with range \([0,1]\).
Thus \(T_\alpha\) is well defined and
\[
 W_\alpha(T_\alpha(r))=G(r).
\]
It follows that \(T_\alpha(|Y|)\) has density \(w_\alpha\) when
\(Y\sim\gamma_3\).  The direction \(Y/|Y|\) is uniform and
independent of \(|Y|\).  The exponential polar coordinates on
\(SU(2)\) have the same uniform angular variable and radial Haar
factor \(\sin^2\theta\).  Hence the image law has exactly the stated
Wilson radial and angular density.
\end{proof}

\begin{corollary}[Exact replacement of every Wilson cumulant]
\label{cor:V71-cumulant-replacement}
For arbitrary representation matrices \(F_i(U)=D^{R_i}(U)\),
\begin{equation}
 \kappa_{\otimes,\nu_\alpha}(F_1,\ldots,F_m)
 =
 \kappa_{\otimes,\gamma_3}
 (F_1\circ\Phi_\alpha,\ldots,F_m\circ\Phi_\alpha).
 \label{eq:V71-cumulant-replacement}
\end{equation}
\end{corollary}

\begin{proof}
Every tensor moment on the two sides agrees by
\eqref{eq:V71-pushforward}.  Apply the common moment--cumulant
polynomial.
\end{proof}

\section{Endpoint behavior of the transport}
\label{sec:V71-endpoints}

\begin{lemma}[Correct parabolic scale at the origin]
\label{lem:V71-origin}
The radial transport extends differentiably to zero and
\begin{equation}
 c\alpha^{-1/2}
 \le T_\alpha'(0)
 \le C\alpha^{-1/2}
 \label{eq:V71-origin-scale}
\end{equation}
for all sufficiently large \(\alpha\).
More precisely \(T_\alpha(r)=\lambda_\alpha r+O_\alpha(r^3)\),
where \(\lambda_\alpha\asymp\alpha^{-1/2}\).
\end{lemma}

\begin{proof}
At the two origins,
\[
 G(r)={1\over3}\sqrt{2\over\pi}\,r^3+O(r^5),
\]
and
\[
 W_\alpha(\theta)
 ={e^\alpha\over3Z_\alpha^{\rm rad}}\theta^3
 +O_\alpha(\theta^5).
\]
The Laplace estimate already used in V40 and V70 gives
\[
 Z_\alpha^{\rm rad}\asymp e^\alpha\alpha^{-3/2}.
\]
Equating the two distribution functions yields
\[
 T_\alpha(r)^3
 =
 C_\alpha\alpha^{-3/2}r^3+O_\alpha(r^5),
\]
where \(C_\alpha\) stays between two positive constants.  Taking the
cubic root proves the expansion and
\eqref{eq:V71-origin-scale}.  The radial vector map
\(T_\alpha(r)y/r\) is therefore differentiable at zero.
\end{proof}

\begin{lemma}[The radial derivative vanishes at the antipodal end]
\label{lem:V71-infinity}
For each fixed \(\alpha>0\),
\begin{equation}
 T_\alpha(r)\longrightarrow\pi,
 \qquad
 T_\alpha'(r)\longrightarrow0
 \quad(r\to\infty).
 \label{eq:V71-infinity}
\end{equation}
\end{lemma}

\begin{proof}
The first assertion follows from
\(G(r)\to1=W_\alpha(\pi)\).  Put
\(\delta(r)=\pi-T_\alpha(r)\).  Standard one-dimensional endpoint
expansions give
\[
 1-G(r)\asymp r e^{-r^2/2},
 \qquad
 1-W_\alpha(\pi-\delta)
 \asymp_\alpha \delta^3
 \quad(\delta\downarrow0).
\]
Consequently
\[
 \delta(r)\asymp_\alpha
 r^{1/3}e^{-r^2/6}.
\]
Differentiating
\(W_\alpha(T_\alpha(r))=G(r)\) gives
\[
 T_\alpha'(r)
 ={g(r)\over w_\alpha(T_\alpha(r))}.
\]
Using \(g(r)\asymp r^2e^{-r^2/2}\) and
\(w_\alpha(\pi-\delta)\asymp_\alpha\delta^2\), together with the
tail identity above, yields
\[
 T_\alpha'(r)\asymp_\alpha r\delta(r)\longrightarrow0.
\]
\end{proof}

\begin{assessment}[Why the complement is now internal]
\label{ass:V71-complement}
Large Gaussian radii are transported smoothly toward the antipode.
Their small probability is not estimated and discarded before the
cumulant is formed; it remains inside the exact Gaussian expectation
in \eqref{eq:V71-cumulant-replacement}.  Thus V70's incompatible
inner/outer cutoff requirements do not arise in this representation.
\end{assessment}

\section{The derivative certificate which would close the local card}
\label{sec:V71-certificate}

Put
\[
 \zeta_\alpha(y):=
 T_\alpha(|y|){y\over|y|}
\]
with its continuous value at zero.

\begin{definition}[Radial transport derivative certificate]
\label{def:V71-RTDC}
An RTDC is a constant \(K<\infty\) such that, for every
\(k\ge1\), sufficiently large \(\alpha\), and \(y\in\mathbb R^3\),
\begin{equation}
 \|D^k\zeta_\alpha(y)\|
 \le k!K^k s_\alpha^{k/2}.
 \label{eq:V71-RTDC}
\end{equation}
\end{definition}

\begin{lemma}[Composition bound under the RTDC]
\label{lem:V71-composition}
If \eqref{eq:V71-RTDC} holds, then
\[
 \widetilde F_R(y):=D^R(\Phi_\alpha(y))
\]
satisfies
\begin{equation}
 \|D^k\widetilde F_R(y)\|_{\rm op}
 \le k!K_1^k
 \bigl(s_\alpha(1+C_2(R))\bigr)^{k/2}.
 \label{eq:V71-composition}
\end{equation}
\end{lemma}

\begin{proof}
Apply the set-partition form of the Faà di Bruno formula to
\[
 \widetilde F_R(y)
 =
 \exp(dR(\zeta_\alpha(y))).
\]
If a partition has \(j\) blocks, the \(j\)-th derivative of the
representation exponential is bounded by
\[
 C_G^j(1+C_2(R))^{j/2}
\]
times its \(j\) Lie-algebra directions, by the Duhamel argument of
\Cref{lem:V67-representation-derivatives}.  The RTDC contributes
\[
 s_\alpha^{k/2}K^k
 \prod_{B\in\pi}|B|!.
\]
Since \(j\le k\) and \(1+C_2(R)\ge1\), replace the Casimir power by
\((1+C_2(R))^{k/2}\).  Finally,
\[
 \sum_{\pi\in\Pi_k}\prod_{B\in\pi}|B|!
 \le k!C_0^k.
\]
For completeness, its exponential generating function is
\[
 \exp\left(\sum_{\ell\ge1}z^\ell\right)
 =\exp\left({z\over1-z}\right),
\]
which is bounded on every fixed circle \(|z|<1\); Cauchy's estimate
gives the displayed bound.  Combining the factors proves
\eqref{eq:V71-composition}.
\end{proof}

\begin{lemma}[Tree degree factorial]
\label{lem:V71-degree-factorial}
For every labelled tree \(\tau\) on \(m\ge2\) vertices,
\begin{equation}
 \prod_{i=1}^m d_i(\tau)!\le(m-1)!.
 \label{eq:V71-degree-factorial}
\end{equation}
\end{lemma}

\begin{proof}
Put \(e_i=d_i-1\), so \(\sum_ie_i=m-2\).  For nonnegative integers,
\[
 (a+1)!(b+1)!\le(a+b+1)!
\]
whenever \(a,b\ge1\); zero entries contribute \(1!\).  Repeatedly
merge the positive \(e_i\).  The largest possible product is obtained
when all \(m-2\) Prüfer occurrences lie at one label, giving
\((m-1)!\).
\end{proof}

\begin{theorem}[RTDC implies the full \(SU(2)\) one-link local tree
card]
\label{thm:V71-RTDC-implies-card}
If the RTDC holds, then the exact Wilson cumulant admits a labelled
tree decomposition
\[
 K_m^\alpha=\sum_{\tau\in\mathfrak T_m}K_{m,\tau}^\alpha
\]
with
\begin{equation}
 \boxed{\quad
 \|K_{m,\tau}^\alpha\|_{\rm op}
 \le
 m!K_2^m
 s_\alpha^{m-1}
 \prod_i(1+C_2(R_i))^{d_i(\tau)/2}.
 \quad}
 \label{eq:V71-Wilson-card}
\end{equation}
The construction includes the entire Wilson law and is uniform in
representations and their dimensions.
\end{theorem}

\begin{proof}
Use the exact replacement
\eqref{eq:V71-cumulant-replacement} and apply the standard-covariance
version of \Cref{thm:V67-Gaussian-tree-formula}.  At vertex \(i\),
\Cref{lem:V71-composition} with \(k=d_i(\tau)\) contributes
\[
 d_i!\,K_1^{d_i}
 \bigl(s_\alpha(1+C_2(R_i))\bigr)^{d_i/2}.
\]
The edge-coordinate sum contributes only a fixed
group-dimensional constant per edge.  Since
\(\sum_i d_i=2m-2\), all fixed bases are exponential per input.
\Cref{lem:V71-degree-factorial} bounds the product of vertex
factorials by \((m-1)!\le m!\), proving
\eqref{eq:V71-Wilson-card}.
\end{proof}

\begin{firewall}[The RTDC is not yet proved]
\label{fire:V71-RTDC-open}
\Cref{lem:V71-origin,lem:V71-infinity} verify the correct first-scale
behavior at both endpoints, but they do not give uniform intermediate
or higher derivative bounds.  Therefore
\Cref{thm:V71-RTDC-implies-card} is a conditional implication, not a
proof of the Wilson local card.
\end{firewall}

\section{V71 ledger and V72 programme}
\label{sec:V71-results}

\begin{theorem}[Fourth-series V71 result ledger]
\label{thm:V71-results}
V71 constructs an exact radial transport from the full Gaussian law
to the full \(SU(2)\) Wilson law, proves exact cumulant replacement,
checks the sharp origin and far-tail endpoint behavior, and proves
that the RTDC would imply the all-order Wilson one-link marked-tree
card with the required factorial-exponential constant.
\end{theorem}

\begin{proof}
Use
\Cref{thm:V71-pushforward,
cor:V71-cumulant-replacement,
lem:V71-origin,
lem:V71-infinity,
thm:V71-RTDC-implies-card}.
\end{proof}

\begin{programme}[V72: derivative bounds for the radial quantile]
\label{prog:V71-V72}
\begin{enumerate}[leftmargin=3em]
\item Rescale \(u_\alpha(r)=\sqrt\alpha\,T_\alpha(r)\) and derive
      its exact first-order ODE from the two radial densities.
\item Prove a uniform bound for \(u_\alpha'\) by matched lower and
      upper radial-tail estimates, including the antipodal layer.
\item Differentiate the logarithmic ODE and test whether
      \(u_\alpha^{(k)}\) has an analytic \(k!K^k\) bound uniformly in
      \(\alpha\).
\item Convert one-dimensional radial bounds into Fréchet derivative
      bounds for \(T_\alpha(|y|)y/|y|\), checking cancellation at
      \(y=0\).
\item If the analytic bound fails, isolate the exact derivative order
      and radial regime and replace the RTDC by the weakest
      tree-summable certificate.
\end{enumerate}
\end{programme}

\begin{firewall}[Claim boundary after fourth-series V71]
\label{fire:V71-claim-boundary}
V71 does not prove the RTDC, the Wilson all-thermal local card, the
anisotropic scalar optimization, or the complete all-order selected
gate.  It does not prove all-order MTA, WUFC, CPM, cutoff removal,
reflection positivity, Osterwalder--Schrader reconstruction,
construction of the continuum Yang--Mills measure, or a uniform
positive continuum spectral gap.  The full Yang--Mills mass gap has
not been proved.
\end{firewall}

\paragraph{Parent-version record.}
V71 was formed from
\texttt{Renewal\_Yang\_Mills\_Mass\_Gap\_Research\_Notes\_V70.tex}.
The V70 parent had \(31\,149\) logical source lines,
\(1\,080\,510\) bytes, and SHA--256
\[
 \texttt{a52a8071bf12b5f62bf0986f16ff94389ca04e13ee0144e0ecb5d61d3596ae16}.
\]
The next compulsory companion audit is V75.

% =============================================================================
% END APPENDED FOURTH-SERIES V71 MATERIAL
% =============================================================================

% =============================================================================
% BEGIN APPENDED FOURTH-SERIES V72 MATERIAL
% =============================================================================

\chapter{V72: The antipodal caustic of radial Wilson transport}
\label{chap:V72-caustic}

\section{The two radial large-deviation rates}
\label{sec:V72-rates}

Let \(G\), \(W_\alpha\), and \(T_\alpha\) be as in V71.

\begin{lemma}[Gaussian radial tail rate]
\label{lem:V72-Gaussian-rate}
For every fixed \(c>0\),
\begin{equation}
 -{1\over\alpha}
 \log\bigl(1-G(c\sqrt\alpha)\bigr)
 \longrightarrow {c^2\over2}.
 \label{eq:V72-Gaussian-rate}
\end{equation}
\end{lemma}

\begin{proof}
The radial tail is a fixed constant times
\[
 \int_{c\sqrt\alpha}^\infty r^2e^{-r^2/2}\,dr.
\]
Integration by parts bounds it above and below by
\(e^{-c^2\alpha/2}\) times factors polynomial in \(\alpha\).
Their logarithms divided by \(\alpha\) tend to zero.
\end{proof}

\begin{lemma}[Wilson upper-tail rate]
\label{lem:V72-Wilson-rate}
For every fixed \(0<\theta<\pi\),
\begin{equation}
 -{1\over\alpha}
 \log\bigl(1-W_\alpha(\theta)\bigr)
 \longrightarrow 1-\cos\theta.
 \label{eq:V72-Wilson-rate}
\end{equation}
\end{lemma}

\begin{proof}
The denominator satisfies
\[
 {1\over\alpha}\log Z_\alpha^{\rm rad}\longrightarrow1
\]
by the Laplace estimate at \(\theta=0\).  In the numerator
\[
 \int_\theta^\pi e^{\alpha\cos u}\sin^2u\,du,
\]
the maximum of \(\cos u\) on \([\theta,\pi]\) is
\(\cos\theta\).  The integral is at most
\(\pi e^{\alpha\cos\theta}\).  Conversely, for every
\(\eta>0\), a sufficiently short fixed interval immediately to the
right of \(\theta\) has
\(\cos u\ge\cos\theta-\eta\) and a positive integral of
\(\sin^2u\).  Thus its logarithm divided by \(\alpha\) tends to
\(\cos\theta\).  Subtract the denominator rate.
\end{proof}

\section{The macroscopic transport profile}
\label{sec:V72-profile}

\begin{theorem}[Large-deviation profile of the exact transport]
\label{thm:V72-profile}
For every \(0<c<2\),
\begin{equation}
 T_\alpha(c\sqrt\alpha)
 \longrightarrow
 \Theta(c):=
 \arccos\left(1-{c^2\over2}\right)
 =2\arcsin{c\over2}.
 \label{eq:V72-profile}
\end{equation}
\end{theorem}

\begin{proof}
The defining quantile identity gives
\[
 1-W_\alpha(T_\alpha(c\sqrt\alpha))
 =
 1-G(c\sqrt\alpha).
\]
Let \(\theta_-<\Theta(c)<\theta_+\) be fixed interior angles.
By \Cref{lem:V72-Wilson-rate},
\[
 1-\cos\theta_-<{c^2\over2}
 <1-\cos\theta_+.
\]
Together with \Cref{lem:V72-Gaussian-rate}, these strict
inequalities imply, for all sufficiently large \(\alpha\),
\[
 1-W_\alpha(\theta_-)
 >
 1-G(c\sqrt\alpha)
 >
 1-W_\alpha(\theta_+).
\]
Since the Wilson upper tail is strictly decreasing, this is equivalent
to
\[
 \theta_-<T_\alpha(c\sqrt\alpha)<\theta_+.
\]
Let both brackets tend to \(\Theta(c)\).
\end{proof}

\begin{corollary}[Square-root antipodal caustic]
\label{cor:V72-square-root}
As \(c\uparrow2\),
\begin{equation}
 \pi-\Theta(c)\sim2\sqrt{2-c},
 \qquad
 \Theta'(c)={1\over\sqrt{1-c^2/4}}\longrightarrow\infty.
 \label{eq:V72-square-root}
\end{equation}
\end{corollary}

\begin{proof}
Write \(c=2-h\) and
\(\Theta(c)=\pi-\delta\).  Then
\[
 \cos(\pi-\delta)
 =-1+{\delta^2\over2}+O(\delta^4),
\]
whereas
\[
 1-{(2-h)^2\over2}
 =-1+2h+O(h^2).
\]
Thus \(\delta\sim2\sqrt h\).  Differentiating the arcsine expression
gives the second assertion.
\end{proof}

\section{Failure of the uniform radial transport certificate}
\label{sec:V72-failure}

\begin{theorem}[The V71 RTDC is false already at first order]
\label{thm:V72-RTDC-false}
There is no constant \(K\) such that
\begin{equation}
 \sup_{r\ge0}T_\alpha'(r)
 \le K\sqrt{s_\alpha}
 \label{eq:V72-false-Lipschitz}
\end{equation}
for all sufficiently large \(\alpha\).  Consequently the RTDC
\eqref{eq:V71-RTDC} fails for the radial quantile transport.
\end{theorem}

\begin{proof}
Suppose \eqref{eq:V72-false-Lipschitz} held.  The calibration
\(s_\alpha\asymp\alpha^{-1}\) would make
\[
 f_\alpha(c):=T_\alpha(c\sqrt\alpha)
\]
uniformly Lipschitz in \(c\), because
\[
 |f_\alpha'(c)|
 =\sqrt\alpha\,T_\alpha'(c\sqrt\alpha)
 \le C K.
\]
Passing to the pointwise limit in
\Cref{thm:V72-profile} would imply
\[
 |\Theta(c)-\Theta(d)|\le CK|c-d|
 \qquad(0<c,d<2).
\]
This is impossible by
\eqref{eq:V72-square-root}, or equivalently because
\(\Theta'\) is unbounded near \(2\).  Finally, the radial derivative
of
\(\zeta_\alpha(y)=T_\alpha(|y|)y/|y|\) has norm
\(T_\alpha'(|y|)\), so the \(k=1\) RTDC would imply
\eqref{eq:V72-false-Lipschitz}.
\end{proof}

\begin{firewall}[Conditional V71 theorem remains logically correct]
\label{fire:V72-V71}
\Cref{thm:V71-RTDC-implies-card} is a correct implication, but its
hypothesis is false for the proposed radial monotone transport.
Therefore it supplies no Wilson local card and must not be used in the
later MTA compiler.
\end{firewall}

\begin{assessment}[Geometric origin of the failure]
\label{ass:V72-origin}
At radius \(c\sqrt\alpha\), the Gaussian cost is \(c^2/2\).  The
Wilson cost to reach angle \(\theta\) is \(1-\cos\theta\), whose
derivative vanishes at the antipode.  Matching the two costs forces
the square-root fold in \eqref{eq:V72-square-root}.  The failure is
therefore not a rough density estimate; it is an intrinsic caustic of
the monotone radial parametrization.
\end{assessment}

\section{The viable weighted replacement}
\label{sec:V72-weighted}

Although the derivative becomes large near the caustic, that region
lies at Gaussian radius approximately \(2\sqrt\alpha\), with
probability of exponential order \(e^{-2\alpha}\).  A pointwise
transport bound discards this compensation before the Gaussian tree
expectation is taken.

\begin{definition}[Weighted radial transport tree certificate]
\label{def:V72-WRTTC}
A WRTTC is a direct bound, uniform in Gaussian tree covariances
\(w^\tau\), of the form
\begin{align}
 &\int
 \prod_{i=1}^m
 \left\|
 D^{d_i(\tau)}
 \bigl(D^{R_i}\circ\Phi_\alpha\bigr)(y_i)
 \right\|_{\rm op}
 d\gamma_{w^\tau}(\boldsymbol y)
 \notag\\
 &\qquad\le
 m!K^m
 \prod_i
 \bigl(s_\alpha(1+C_2(R_i))\bigr)^{d_i(\tau)/2}.
 \label{eq:V72-WRTTC}
\end{align}
Unlike the RTDC, this certificate keeps derivative growth and
Gaussian rarity in the same integral.
\end{definition}

\begin{proposition}[A WRTTC is sufficient for the local Wilson card]
\label{prop:V72-WRTTC-sufficient}
If \eqref{eq:V72-WRTTC} holds with at most one factorial over the
complete tree expansion, then the Gaussian forest identity applied to
\eqref{eq:V71-cumulant-replacement} gives the all-order one-link
Wilson local marked-tree card.
\end{proposition}

\begin{proof}
In the standard-Gaussian version of
\eqref{eq:V67-Gaussian-tree-formula}, take the operator norm inside
each interpolated Gaussian expectation and expand the fixed
three-dimensional edge-coordinate contractions.  The resulting
integral is exactly the left side of
\eqref{eq:V72-WRTTC}, up to a fixed constant per edge.  Its right
side is the required local tree monomial.  No chart split or private
moment is introduced.
\end{proof}

\begin{assessment}[The WRTTC is a target, not a theorem]
\label{ass:V72-WRTTC-open}
The WRTTC is strictly adapted to the failure just proved, but it is
not established here.  In particular, separate \(L^p\) bounds for
each derivative followed by Hölder may create an extra \(m^m\)
factor.  The caustic layer must be integrated jointly with the
tree-correlated Gaussian replicas.
\end{assessment}

\section{V72 ledger and V73 programme}
\label{sec:V72-results}

\begin{theorem}[Fourth-series V72 result ledger]
\label{thm:V72-results}
V72 proves the Gaussian and Wilson radial large-deviation rates, the
macroscopic profile of the exact radial quantile map, its
square-root antipodal caustic, and the failure of the pointwise RTDC
at first order.  It replaces that false route by the precisely
sufficient probability-weighted transport tree certificate.
\end{theorem}

\begin{proof}
Use
\Cref{lem:V72-Gaussian-rate,
lem:V72-Wilson-rate,
thm:V72-profile,
cor:V72-square-root,
thm:V72-RTDC-false,
prop:V72-WRTTC-sufficient}.
\end{proof}

\begin{programme}[V73: caustic-layer weighted derivatives]
\label{prog:V72-V73}
\begin{enumerate}[leftmargin=3em]
\item Derive the exact density-ratio formula
      \(T_\alpha'=g/(w_\alpha\circ T_\alpha)\) in the scaled radius
      \(c=r/\sqrt\alpha\).
\item Partition \(c<2\), the caustic layer \(c\approx2\), and
      \(c>2\) using rate differences, not an arbitrary physical
      chart.
\item Prove weighted first-derivative moments with the Gaussian
      probability retained.
\item Test whether the common tree covariance can concentrate several
      replicas simultaneously in the caustic layer and charge that
      event only once.
\item Abandon the transport route if the joint layer produces more
      than one factorial; then return to an intrinsic Wilson
      semigroup identity.
\end{enumerate}
\end{programme}

\begin{firewall}[Claim boundary after fourth-series V72]
\label{fire:V72-claim-boundary}
V72 disproves the RTDC but does not prove the WRTTC, the Wilson
all-thermal local card, the anisotropic scalar optimization, or the
complete all-order selected gate.  It does not prove all-order MTA,
WUFC, CPM, cutoff removal, reflection positivity,
Osterwalder--Schrader reconstruction, construction of the continuum
Yang--Mills measure, or a uniform positive continuum spectral gap.
The full Yang--Mills mass gap has not been proved.
\end{firewall}

\paragraph{Parent-version record.}
V72 was formed from
\texttt{Renewal\_Yang\_Mills\_Mass\_Gap\_Research\_Notes\_V71.tex}.
The V71 parent had \(31\,541\) logical source lines,
\(1\,092\,151\) bytes, and SHA--256
\[
 \texttt{dfdc0e77dbc11b1df0155cbe75ff450a602cd1bee56fb22a965c9d20e4ca3a20}.
\]
The next compulsory companion audit is V75.

% =============================================================================
% END APPENDED FOURTH-SERIES V72 MATERIAL
% =============================================================================

% =============================================================================
% BEGIN APPENDED FOURTH-SERIES V73 MATERIAL
% =============================================================================

\chapter{V73: Quantitative width and slope of the antipodal caustic}
\label{chap:V73-caustic-scale}

\section{Critical quantile notation}
\label{sec:V73-notation}

Set
\[
 r_\alpha:=2\sqrt\alpha,\qquad
 \delta_\alpha:=\pi-T_\alpha(r_\alpha),\qquad
 D_\alpha:=\sqrt\alpha\,\delta_\alpha.
\]
The rate theorem V72 implies \(\delta_\alpha\to0\), but does not
determine its boundary-layer scale.

\begin{lemma}[Critical Gaussian tail with prefactor]
\label{lem:V73-Gaussian-critical}
There are fixed \(0<c<C<\infty\) such that
\begin{equation}
 c\sqrt\alpha\,e^{-2\alpha}
 \le
 1-G(2\sqrt\alpha)
 \le
 C\sqrt\alpha\,e^{-2\alpha}
 \label{eq:V73-Gaussian-critical}
\end{equation}
for all sufficiently large \(\alpha\).
\end{lemma}

\begin{proof}
For \(R\ge1\), integration by parts gives
\[
 \int_R^\infty r^2e^{-r^2/2}\,dr
 =
 Re^{-R^2/2}+\int_R^\infty e^{-r^2/2}\,dr.
\]
The second term lies between zero and
\(R^{-1}e^{-R^2/2}\).  Put \(R=2\sqrt\alpha\) and restore the fixed
radial Gaussian normalizer.
\end{proof}

\section{Wilson antipodal boundary layer}
\label{sec:V73-Wilson-layer}

\begin{lemma}[Uniform antipodal tail reduction]
\label{lem:V73-antipodal-reduction}
If \(\delta_\alpha\to0\) and
\(D_\alpha=\sqrt\alpha\,\delta_\alpha\) satisfies
\(D_\alpha^4/\alpha\to0\), then
\begin{align}
 1-W_\alpha(\pi-\delta_\alpha)
 &\asymp
 e^{-2\alpha}
 \int_0^{D_\alpha}e^{v^2/2}v^2\,dv,
 \label{eq:V73-antipodal-reduction}\\
 w_\alpha(\pi-\delta_\alpha)
 &\asymp
 \alpha^{1/2}e^{-2\alpha}
 e^{D_\alpha^2/2}D_\alpha^2.
 \label{eq:V73-antipodal-density}
\end{align}
The comparison constants are fixed.
\end{lemma}

\begin{proof}
Put \(u=\pi-\theta\).  On a shrinking antipodal interval,
\[
 \cos(\pi-u)=-\cos u
 =-1+{u^2\over2}+O(u^4),
\qquad
 \sin^2u=u^2(1+O(u^2)).
\]
With \(v=\sqrt\alpha\,u\), the unnormalized upper-tail integral is
\[
 e^{-\alpha}\alpha^{-3/2}
 \int_0^{D_\alpha}
 e^{v^2/2+O(v^4/\alpha)}
 v^2\left(1+O(v^2/\alpha)\right)dv.
\]
The assumed error tends uniformly to zero.  Divide by
\[
 Z_\alpha^{\rm rad}\asymp e^\alpha\alpha^{-3/2}
\]
to obtain \eqref{eq:V73-antipodal-reduction}.  Evaluating the
normalized density at \(u=\delta_\alpha\) gives
\eqref{eq:V73-antipodal-density}.
\end{proof}

\begin{lemma}[Elementary growing-endpoint integral]
\label{lem:V73-endpoint-integral}
For \(D\ge2\),
\begin{equation}
 cD e^{D^2/2}
 \le
 \int_0^D e^{v^2/2}v^2\,dv
 \le
 CD e^{D^2/2}.
 \label{eq:V73-endpoint-integral}
\end{equation}
\end{lemma}

\begin{proof}
Since
\[
 {d\over dv}\left(ve^{v^2/2}\right)
 =(1+v^2)e^{v^2/2},
\]
the upper bound follows by replacing \(v^2\) with \(1+v^2\).
For the lower bound integrate only over
\([D-D^{-1},D]\).  On this interval \(v\asymp D\) and
\[
 e^{v^2/2}\ge e^{(D-D^{-1})^2/2}
 \ge c e^{D^2/2}.
\]
The interval has length \(D^{-1}\), giving the lower bound.
\end{proof}

\begin{theorem}[Exact logarithmic width of the critical caustic]
\label{thm:V73-caustic-width}
One has
\begin{equation}
 D_\alpha^2
 =
 \log\alpha+O(\log\log\alpha),
 \label{eq:V73-D-scale}
\end{equation}
and hence
\begin{equation}
 \boxed{\quad
 \delta_\alpha
 \asymp
 \sqrt{\log\alpha\over\alpha}.
 \quad}
 \label{eq:V73-delta-scale}
\end{equation}
\end{theorem}

\begin{proof}
The quantile identity and
\Cref{lem:V73-Gaussian-critical} give
\[
 1-W_\alpha(\pi-\delta_\alpha)
 \asymp\sqrt\alpha\,e^{-2\alpha}.
\]
The rate result already gives \(\delta_\alpha\to0\).  Applying the
Taylor estimates in
\Cref{lem:V73-antipodal-reduction} first with fixed large upper and
lower trial values for \(D\) shows \(D_\alpha\to\infty\).  Trial
values \(D=C\sqrt{\log\alpha}\) and
\(D=c\sqrt{\log\alpha}\), using monotonicity of the tail, then place
\(D_\alpha\) between two such multiples.  In particular
\(D_\alpha^4/\alpha\to0\), justifying
\Cref{lem:V73-antipodal-reduction} without circularity.

We therefore have
\[
 \int_0^{D_\alpha}e^{v^2/2}v^2\,dv
 \asymp\sqrt\alpha.
\]
\Cref{lem:V73-endpoint-integral} yields
\[
 D_\alpha e^{D_\alpha^2/2}\asymp\sqrt\alpha.
\]
Taking logarithms gives
\[
 {D_\alpha^2\over2}+\log D_\alpha
 ={1\over2}\log\alpha+O(1),
\]
which is \eqref{eq:V73-D-scale}.  Divide by \(\alpha\) and take
square roots for \eqref{eq:V73-delta-scale}.
\end{proof}

\section{Critical derivative size}
\label{sec:V73-slope}

\begin{theorem}[Quantitative blow-up of the normalized transport
slope]
\label{thm:V73-slope}
At \(r_\alpha=2\sqrt\alpha\),
\begin{align}
 T_\alpha'(r_\alpha)
 &\asymp {1\over\sqrt{\log\alpha}},
 \label{eq:V73-unscaled-slope}\\
 \sqrt\alpha\,T_\alpha'(r_\alpha)
 &\asymp
 \sqrt{\alpha\over\log\alpha}.
 \label{eq:V73-scaled-slope}
\end{align}
\end{theorem}

\begin{proof}
Differentiate the exact quantile relation:
\[
 T_\alpha'(r_\alpha)
 ={g(r_\alpha)\over
   w_\alpha(\pi-\delta_\alpha)}.
\]
The Gaussian density satisfies
\[
 g(2\sqrt\alpha)\asymp\alpha e^{-2\alpha}.
\]
By \eqref{eq:V73-antipodal-density},
\[
 w_\alpha(\pi-\delta_\alpha)
 \asymp
 \alpha^{1/2}e^{-2\alpha}
 e^{D_\alpha^2/2}D_\alpha^2.
\]
The proof of \Cref{thm:V73-caustic-width} gives
\[
 D_\alpha e^{D_\alpha^2/2}\asymp\sqrt\alpha.
\]
Thus the Wilson density is
\(\asymp\alpha D_\alpha e^{-2\alpha}\), and the density ratio is
\(\asymp D_\alpha^{-1}\).  Use
\(D_\alpha\asymp\sqrt{\log\alpha}\).
\end{proof}

\begin{corollary}[The failure is large but exponentially rare]
\label{cor:V73-rare}
The pointwise RTDC misses the critical derivative by the factor
\[
 \sqrt{\alpha\over\log\alpha},
\]
while the Gaussian radial density at that location is of order
\[
 \alpha e^{-2\alpha}.
\]
Thus every fixed polynomial power of the normalized derivative is
dominated by the caustic probability as \(\alpha\to\infty\).
\end{corollary}

\begin{proof}
The first statement is
\eqref{eq:V73-scaled-slope}; the second is the displayed Gaussian
density in its proof.  For fixed \(p\),
\[
 \left({\alpha\over\log\alpha}\right)^{p/2}
 \alpha e^{-2\alpha}\longrightarrow0.
\]
\end{proof}

\begin{firewall}[Fixed moments do not prove the all-order WRTTC]
\label{fire:V73-fixed-not-all}
\Cref{cor:V73-rare} holds for each fixed power.  The WRTTC requires
simultaneous arity control and correlated replicas.  Optimizing
\(\alpha\) against a growing derivative power can change the
constant growth, so fixed-\(p\) rarity cannot be promoted to
\(m!K^m\) without a tracked moment calculation.
\end{firewall}

\section{V73 ledger and V74 programme}
\label{sec:V73-results}

\begin{theorem}[Fourth-series V73 result ledger]
\label{thm:V73-results}
V73 proves that the critical radial quantile lies
\(\asymp\sqrt{\log\alpha/\alpha}\) from the antipode and that its
unscaled derivative is
\(\asymp(\log\alpha)^{-1/2}\).  Hence the RTDC-normalized slope grows
as \(\sqrt{\alpha/\log\alpha}\), but occurs at Gaussian density
scale \(e^{-2\alpha}\).
\end{theorem}

\begin{proof}
Use
\Cref{thm:V73-caustic-width,
thm:V73-slope,
cor:V73-rare}.
\end{proof}

\begin{programme}[V74: tracked caustic moments]
\label{prog:V73-V74}
\begin{enumerate}[leftmargin=3em]
\item Prove upper bounds for \(T_\alpha'\) throughout the
      \(c<2\), caustic, and \(c>2\) radial regimes.
\item Optimize
      \(\mathbb E_{\gamma_3}
      [(\sqrt\alpha T_\alpha')^p]\)
      simultaneously in \(p\) and \(\alpha\).
\item Determine whether its growth is \(p!K^p\), a smaller
      \((p!)^{1/2}K^p\), or a larger Gevrey class.
\item Repeat for tangential derivatives \(T_\alpha(r)/r\).
\item Only if the one-replica moments have the correct scale, test
      joint tree-correlated Gaussian replicas without separate
      Hölder losses.
\end{enumerate}
\end{programme}

\begin{firewall}[Claim boundary after fourth-series V73]
\label{fire:V73-claim-boundary}
V73 quantifies the caustic but does not prove tracked derivative
moments, the WRTTC, the Wilson all-thermal local card, the anisotropic
scalar optimization, or the complete all-order selected gate.  It
does not prove all-order MTA, WUFC, CPM, cutoff removal, reflection
positivity, Osterwalder--Schrader reconstruction, construction of the
continuum Yang--Mills measure, or a uniform positive continuum
spectral gap.  The full Yang--Mills mass gap has not been proved.
\end{firewall}

\paragraph{Parent-version record.}
V73 was formed from
\texttt{Renewal\_Yang\_Mills\_Mass\_Gap\_Research\_Notes\_V72.tex}.
The V72 parent had \(31\,864\) logical source lines,
\(1\,102\,152\) bytes, and SHA--256
\[
 \texttt{12edd1987ce8b325bf6daa08aad0d4b2ee246188640b37397b3cd12be5c4a7d2}.
\]
The next compulsory companion audit is V75.

% =============================================================================
% END APPENDED FOURTH-SERIES V73 MATERIAL
% =============================================================================

% =============================================================================
% BEGIN APPENDED FOURTH-SERIES V74 MATERIAL
% =============================================================================

\chapter{V74: An optimized weighted-transport tree compiler}
\label{chap:V74-compiler}

\section{Degree-adapted Hölder allocation}
\label{sec:V74-Holder}

\begin{lemma}[One common moment order for an entire tree]
\label{lem:V74-common-moment}
Let \(L_1,\ldots,L_m\ge0\) have an arbitrary joint law and identical
one-variable moment bound
\begin{equation}
 \|L_i\|_{L^p}\le K\sqrt p
 \qquad(p\ge2).
 \label{eq:V74-subGaussian}
\end{equation}
If \(d_i\ge1\) and
\[
 Q:=\sum_i d_i,
\]
then
\begin{equation}
 \mathbb E\prod_iL_i^{d_i}
 \le (K\sqrt Q)^Q.
 \label{eq:V74-common-moment}
\end{equation}
No independence is required.
\end{lemma}

\begin{proof}
Use generalized Hölder with exponent
\[
 p_i={Q\over d_i}
\]
on the factor \(L_i^{d_i}\).  Since
\(\sum_i p_i^{-1}=1\),
\[
 \mathbb E\prod_iL_i^{d_i}
 \le
 \prod_i
 \left(
 \mathbb E L_i^{d_ip_i}
 \right)^{1/p_i}
 =
 \prod_i\|L_i\|_{L^Q}^{d_i}.
\]
Apply \eqref{eq:V74-subGaussian} at the single common order \(Q\)
and sum the exponents.
\end{proof}

\begin{corollary}[Tree-degree moments have one-factorial growth]
\label{cor:V74-tree-moment}
For tree degrees \(d_i=d_i(\tau)\),
\begin{equation}
 \mathbb E\prod_iL_i^{d_i}
 \le m!K_1^m.
 \label{eq:V74-tree-moment}
\end{equation}
\end{corollary}

\begin{proof}
A tree has \(Q=2m-2\), so
\[
 (K\sqrt Q)^Q
 \le K_2^{2m}(2m)^{m-1}.
\]
The elementary lower Stirling bound
\(m!\ge(m/e)^m\) absorbs
\((2m)^{m-1}\) into \(m!K_1^m\).
\end{proof}

\begin{remark}[Why uniform \(m\)-fold Hölder is wasteful]
\label{rem:V74-Hoelder-loss}
Taking exponent \(m\) for every vertex would require moments of order
\(md_i\) and produce
\[
 \prod_i(md_i)^{d_i/2}.
\]
On a star this has an avoidable extra half-factorial.  The choice
\(Q/d_i\) makes every powered factor use the same total order
\(Q=2m-2\), exactly matching the tree incidence count.
\end{remark}

\section{A sufficient derivative envelope}
\label{sec:V74-envelope}

For a Gaussian replica covariance \(w^\tau\), let \(Y_i\) denote the
\(i\)-th marginal.  Every marginal is standard Gaussian, even though
the replicas may be fully correlated.

\begin{definition}[Sub-Gaussian exponential derivative envelope]
\label{def:V74-SEDE}
A SEDE consists of nonnegative functions \(L_\alpha\) and a constant
\(K\) such that
\begin{align}
 \sup_\alpha\|L_\alpha(Y)\|_{L^p(\gamma_3)}
 &\le K\sqrt p
 &&(p\ge2),
 \label{eq:V74-envelope-moment}\\
 \left\|
 D^k(D^R\circ\Phi_\alpha)(y)
 \right\|_{\rm op}
 &\le
 K^k
 \bigl(s_\alpha(1+C_2(R))\bigr)^{k/2}
 L_\alpha(y)^k
 &&(k\ge1).
 \label{eq:V74-envelope-derivative}
\end{align}
The essential feature is exponential dependence on \(k\), with no
additional \(k!\).
\end{definition}

\begin{theorem}[SEDE implies the WRTTC]
\label{thm:V74-SEDE-WRTTC}
If a SEDE holds, then every labelled tree and every interpolated
Gaussian covariance in V67 satisfy
\begin{align}
 &\mathbb E_{w^\tau}
 \prod_i
 \left\|
 D^{d_i(\tau)}
 (D^{R_i}\circ\Phi_\alpha)(Y_i)
 \right\|_{\rm op}
 \notag\\
 &\qquad\le
 m!K_1^m
 \prod_i
 \bigl(s_\alpha(1+C_2(R_i))\bigr)^{d_i(\tau)/2}.
 \label{eq:V74-SEDE-WRTTC}
\end{align}
Thus the WRTTC and the full \(SU(2)\) one-link Wilson local card
follow.
\end{theorem}

\begin{proof}
Apply \eqref{eq:V74-envelope-derivative} at every vertex.  All
deterministic constants use
\(\sum_i d_i=2m-2\) and are exponential per input.  The remaining
random factor is
\[
 \mathbb E_{w^\tau}
 \prod_iL_\alpha(Y_i)^{d_i}.
\]
The marginals obey
\eqref{eq:V74-envelope-moment}; their dependence is irrelevant by
\Cref{lem:V74-common-moment}.  Apply
\Cref{cor:V74-tree-moment}.  This is precisely
\eqref{eq:V72-WRTTC}, and
\Cref{prop:V72-WRTTC-sufficient} gives the local card.
\end{proof}

\section{The derivative-order boundary}
\label{sec:V74-order}

\begin{proposition}[A vertex factorial creates a second factorial in
this compiler]
\label{prop:V74-factorial-boundary}
Suppose \eqref{eq:V74-envelope-derivative} is replaced by
\begin{equation}
 \|D^k(D^R\circ\Phi_\alpha)(y)\|
 \le
 (k!)^\beta K^k
 \bigl(s_\alpha(1+C_2(R))\bigr)^{k/2}
 L_\alpha(y)^k
 \label{eq:V74-Gevrey-envelope}
\end{equation}
with \(\beta>0\).  Then the preceding absolute Hölder compiler yields
at best
\begin{equation}
 (m!)^{1+\beta}K_1^m
 \label{eq:V74-second-factorial}
\end{equation}
on star trees, rather than \(m!K_1^m\).
\end{proposition}

\begin{proof}
The random envelope produces one factorial by
\Cref{cor:V74-tree-moment}.  The deterministic derivative factors
contribute
\[
 \prod_i(d_i!)^\beta.
\]
For a star this is exactly
\(((m-1)!)^\beta\).  Multiplication gives
\eqref{eq:V74-second-factorial} up to exponential bases.
\end{proof}

\begin{firewall}[Method boundary, not an impossibility theorem]
\label{fire:V74-factorial-scope}
\Cref{prop:V74-factorial-boundary} proves that the stated absolute
envelope/Hölder method cannot tolerate a vertex factorial.  It does
not prove that the Wilson local card is false when pointwise
derivatives have such growth.  Exact cancellations inside the
Gaussian expectation or a different forest assignment could remove
the apparent second factorial.
\end{firewall}

\section{How the caustic fits the envelope}
\label{sec:V74-caustic}

The V73 calculation suggests the normalized first-derivative envelope
\begin{equation}
 L_\alpha^{(1)}(y):=
 1+s_\alpha^{-1/2}
 \|D\zeta_\alpha(y)\|.
 \label{eq:V74-first-envelope}
\end{equation}
At the critical caustic it is of order
\(\sqrt{\alpha/\log\alpha}\), while the radial Gaussian density is
of order \(e^{-2\alpha}\).  Optimizing the model expression
\[
 \left({\alpha\over\log\alpha}\right)^{p/2}e^{-2\alpha}
\]
in \(\alpha\) gives a scale no larger than
\((Cp/\log(2+p))^{p/2}\), consistent with the sub-Gaussian target
\eqref{eq:V74-envelope-moment}.  This calculation motivates but does
not prove the uniform moment bound, because V73 controls only the
critical point rather than the full caustic layer.

\begin{assessment}[Exact remaining transport task]
\label{ass:V74-remaining}
The transport route is now reduced to two sharp questions:
\begin{enumerate}[label=\textup{(\roman*)},leftmargin=3em]
\item does a common envelope \(L_\alpha\) control every derivative
      order with only \(K^k\), as in
      \eqref{eq:V74-envelope-derivative}; and
\item does that envelope have uniform sub-Gaussian moments?
\end{enumerate}
The first question is stricter.  A routine analytic inverse-function
bound gives \(k!\), which
\Cref{prop:V74-factorial-boundary} shows is insufficient for this
absolute compiler.
\end{assessment}

\section{V74 ledger and V75 programme}
\label{sec:V74-results}

\begin{theorem}[Fourth-series V74 result ledger]
\label{thm:V74-results}
V74 proves the degree-adapted Hölder lemma, a one-factorial moment
bound for arbitrary correlated tree replicas, and the SEDE-to-WRTTC
compiler.  It also proves that a vertex \(k!\) loss creates a second
factorial in this method.
\end{theorem}

\begin{proof}
Use
\Cref{lem:V74-common-moment,
cor:V74-tree-moment,
thm:V74-SEDE-WRTTC,
prop:V74-factorial-boundary}.
\end{proof}

\begin{programme}[V75: audit and caustic envelope]
\label{prog:V74-V75}
\begin{enumerate}[leftmargin=3em]
\item Perform the compulsory read-only audit of the newest active
      SM--Einstein and Navier--Stokes notes.
\item Prove global two-sided density-ratio bounds for \(T_\alpha'\)
      in the pre-caustic, caustic, and post-caustic regimes.
\item Establish or refute the sub-Gaussian moment bound for
      \(L_\alpha^{(1)}\).
\item Differentiate the quantile ODE once and test whether the second
      derivative admits the same envelope without a \(2!\)-type
      iterative loss.
\item If factorial derivative growth is intrinsic, abandon the
      absolute transport compiler and return to the intrinsic Wilson
      forest/semigroup route.
\end{enumerate}
\end{programme}

\begin{firewall}[Claim boundary after fourth-series V74]
\label{fire:V74-claim-boundary}
V74 proves a conditional weighted compiler, not a SEDE, WRTTC, the
Wilson all-thermal local card, the anisotropic scalar optimization,
or the complete all-order selected gate.  It does not prove
all-order MTA, WUFC, CPM, cutoff removal, reflection positivity,
Osterwalder--Schrader reconstruction, construction of the continuum
Yang--Mills measure, or a uniform positive continuum spectral gap.
The full Yang--Mills mass gap has not been proved.
\end{firewall}

\paragraph{Parent-version record.}
V74 was formed from
\texttt{Renewal\_Yang\_Mills\_Mass\_Gap\_Research\_Notes\_V73.tex}.
The V73 parent had \(32\,182\) logical source lines,
\(1\,111\,020\) bytes, and SHA--256
\[
 \texttt{b7d436db743ef4ea4ab1d378bdf363b30a4d93672ddf799eac319f709895fdf2}.
\]
The next compulsory companion audit is V75.

% =============================================================================
% END APPENDED FOURTH-SERIES V74 MATERIAL
% =============================================================================

% =============================================================================
% BEGIN APPENDED FOURTH-SERIES V75 MATERIAL
% =============================================================================

\chapter{V75: Companion audit and the sharp caustic moment floor}
\label{chap:V75-moment-floor}

\section{Compulsory read-only companion audit}
\label{sec:V75-audit}

The active companion files inspected at this checkpoint were
\begin{center}
\begin{tabular}{p{0.54\textwidth}rr}
\toprule
file&logical lines&bytes\\
\midrule
\texttt{Renewal\_SM\_Einstein\_Continuum\_Research\_Notes\_V61.tex}
&\(24\,754\)&\(843\,640\)\\
\texttt{Renewal\_NS\_Stability\_Research\_Notes\_V31.tex}
&\(23\,052\)&\(887\,537\)\\
\bottomrule
\end{tabular}
\end{center}
Their SHA--256 digests are
\[
\begin{split}
 &\texttt{0a7843663300f938d1b65fbb7524a0efdc5e647ee40b92505699b8f965a6c108},\\
 &\texttt{f1121b62238ad5491bb7cf03635ea9934942edf573ed8faaa9ee79e1d38a6696}.
\end{split}
\]
Neither companion file was modified.

SM V60 proves an exact projective common-refinement calculus for its
finite-cutoff bosonic and Grassmann shell cells.  Coarse reductions
need not be invertible; two cells are compared by retaining their
union and integrating exactly toward each cell.  SM V61 shows that
the inherited determinant line respects those refinements.  It also
keeps harmonic Maxwell modes explicit, proves a volume-uniform bound
only for fixed local neutral currents, and retains the full
electroweak band where the photon projector ceases to be canonical.
NS V31 is unchanged.

\begin{assessment}[V75 audit transfer]
\label{ass:V75-audit-transfer}
The YM transport and intrinsic-Wilson descriptions should likewise be
compared through the exact common object
\eqref{eq:V71-cumulant-replacement}, not through a fictitious inverse
of a coarse norm estimate.  The caustic layer is a genuine transport
stratum and must remain inside the Gaussian expectation.  No SM
determinant, photon, or NS regularity constant acts on the Wilson
one-link carrier, so none is imported.
\end{assessment}

\begin{firewall}[Companion scope]
\label{fire:V75-companion-scope}
The finite Hodge splitting and neutral-current infrared bound do not
prove confinement or a pure-Yang--Mills mass gap.  The exact
common-refinement principle is used only as an assembly rule.
\end{firewall}

\section{An exact derivative-moment formula}
\label{sec:V75-moment-identity}

Let
\[
 R_\alpha(\theta):=
 G^{-1}(W_\alpha(\theta))
\]
be the inverse radial quantile.

\begin{lemma}[Transport derivative moments in Wilson radius]
\label{lem:V75-moment-identity}
For every \(p>0\),
\begin{equation}
 \boxed{\quad
 \mathbb E_{\gamma_3}
 \left[
 \bigl(\sqrt\alpha\,T_\alpha'(|Y|)\bigr)^p
 \right]
 =
 \alpha^{p/2}
 \int_0^\pi
 g(R_\alpha(\theta))^p
 w_\alpha(\theta)^{1-p}\,d\theta.
 \quad}
 \label{eq:V75-moment-identity}
\end{equation}
\end{lemma}

\begin{proof}
Differentiating
\(W_\alpha(T_\alpha(r))=G(r)\) gives
\[
 T_\alpha'(r)
 ={g(r)\over w_\alpha(T_\alpha(r))}.
\]
Insert this into the Gaussian radial moment and change variables
\(\theta=T_\alpha(r)\).  Since
\[
 d\theta={g(r)\over w_\alpha(\theta)}\,dr,
\]
the resulting integrand is exactly the right side of
\eqref{eq:V75-moment-identity}.
\end{proof}

\section{A whole critical interval, not one point}
\label{sec:V75-critical-interval}

Define
\[
 I_\alpha:=
 \left[
 2\sqrt\alpha,\,
 2\sqrt\alpha+\alpha^{-1/2}
 \right].
\]

\begin{lemma}[Uniform critical-layer comparability]
\label{lem:V75-critical-interval}
For all sufficiently large \(\alpha\) and every \(r\in I_\alpha\),
\begin{align}
 g(r)&\asymp\alpha e^{-2\alpha},
 \label{eq:V75-critical-density}\\
 \pi-T_\alpha(r)&\asymp
 \sqrt{\log\alpha\over\alpha},
 \label{eq:V75-critical-angle}\\
 T_\alpha'(r)&\asymp{1\over\sqrt{\log\alpha}}.
 \label{eq:V75-critical-slope}
\end{align}
All comparison constants are fixed.
\end{lemma}

\begin{proof}
Across \(I_\alpha\), the logarithm of the Gaussian density changes by
\[
 \int_{2\sqrt\alpha}^r
 \left({2\over u}-u\right)du=O(1),
\]
which proves \eqref{eq:V75-critical-density}.  The same calculation
or the Gaussian Mills bounds show that its upper tail changes by only
a fixed multiplicative factor.

Write
\[
 D(r):=\sqrt\alpha\{\pi-T_\alpha(r)\}.
\]
The quantile identity turns the preceding tail comparison into a
fixed multiplicative comparison between
\[
 \int_0^{D(r)}e^{v^2/2}v^2\,dv
\quad\hbox{and}\quad
 \int_0^{D_\alpha}e^{v^2/2}v^2\,dv.
\]
The uniform antipodal reduction of V73 applies because
\(D_\alpha^2=\log\alpha+O(\log\log\alpha)\).
By \Cref{lem:V73-endpoint-integral}, a fixed multiplicative change of
the integral changes \(D\) by \(O(D_\alpha^{-1})\).  Hence
\[
 D(r)=D_\alpha+O(D_\alpha^{-1})
 \asymp\sqrt{\log\alpha},
\]
proving \eqref{eq:V75-critical-angle}.

The Wilson density formula
\eqref{eq:V73-antipodal-density} changes by only a fixed factor under
this \(D\)-variation.  Divide
\eqref{eq:V75-critical-density} by that Wilson density in the exact
formula for \(T_\alpha'\).  The calculation in
\Cref{thm:V73-slope} is therefore uniform on \(I_\alpha\) and gives
\eqref{eq:V75-critical-slope}.
\end{proof}

\begin{corollary}[Gaussian mass of the critical interval]
\label{cor:V75-critical-mass}
\begin{equation}
 \gamma_3\{|Y|\in I_\alpha\}
 \asymp\sqrt\alpha\,e^{-2\alpha}.
 \label{eq:V75-critical-mass}
\end{equation}
\end{corollary}

\begin{proof}
Multiply the density estimate
\eqref{eq:V75-critical-density} by
\(|I_\alpha|=\alpha^{-1/2}\).
\end{proof}

\section{The necessary moment growth}
\label{sec:V75-lower}

\begin{theorem}[Sharp caustic lower bound for derivative moments]
\label{thm:V75-moment-floor}
There is \(c>0\) such that, for all sufficiently large \(p\),
\begin{equation}
 \boxed{\quad
 \sup_{\alpha\ge\alpha_0}
 \left\|
 \sqrt\alpha\,T_\alpha'(|Y|)
 \right\|_{L^p(\gamma_3)}
 \ge
 c\sqrt{p\over\log p}.
 \quad}
 \label{eq:V75-moment-floor}
\end{equation}
\end{theorem}

\begin{proof}
By \Cref{lem:V75-critical-interval,
cor:V75-critical-mass},
\begin{align*}
 \mathbb E
 \left[
 \bigl(\sqrt\alpha T_\alpha'(|Y|)\bigr)^p
 \right]
 &\ge
 c\sqrt\alpha e^{-2\alpha}
 \left(
 c\sqrt{\alpha\over\log\alpha}
 \right)^p.
\end{align*}
Choose \(\alpha=p/4\), rounded upward if necessary.  Taking the
\(p\)-th root gives
\[
 c\left(\sqrt\alpha e^{-2\alpha}\right)^{1/p}
 \sqrt{\alpha\over\log\alpha}
 \ge c_1\sqrt{p\over\log p},
\]
because the first factor stays bounded below by a positive constant.
\end{proof}

\begin{corollary}[Any common envelope is nearly sub-Gaussian sharp]
\label{cor:V75-envelope-floor}
If a prospective envelope \(L_\alpha\) pointwise dominates
\[
 c\sqrt\alpha\,T_\alpha'(|y|),
\]
then
\[
 \sup_\alpha\|L_\alpha\|_{L^p(\gamma_3)}
 \ge c_1\sqrt{p/\log p}.
\]
Thus a uniform bounded or subexponential-in-\(p\) envelope is
impossible.  The sub-Gaussian upper target \(K\sqrt p\) in V74, if
true, is optimal up to a logarithm.
\end{corollary}

\begin{firewall}[A lower bound is not the SEDE]
\label{fire:V75-lower-not-upper}
\Cref{thm:V75-moment-floor} neither proves nor refutes the
sub-Gaussian upper bound.  It proves that the caustic cannot be erased
by changing constants and that any successful WRTTC proof must spend
nearly the full permitted moment growth.
\end{firewall}

\section{V75 ledger and V76 programme}
\label{sec:V75-results}

\begin{theorem}[Fourth-series V75 result ledger]
\label{thm:V75-results}
V75 completes the compulsory companion audit, derives an exact
one-dimensional identity for radial transport derivative moments,
extends the V73 critical estimates to an interval of quantified
Gaussian mass, and proves the lower moment growth
\(\sqrt{p/\log p}\).
\end{theorem}

\begin{proof}
Use
\Cref{ass:V75-audit-transfer,
lem:V75-moment-identity,
lem:V75-critical-interval,
thm:V75-moment-floor}.
\end{proof}

\begin{programme}[V76: global first-derivative upper bound]
\label{prog:V75-V76}
\begin{enumerate}[leftmargin=3em]
\item Prove matching density-ratio bounds for
      \(T_\alpha'\) in the pre-caustic region
      \(r\le2\sqrt\alpha-\alpha^{-1/2}\).
\item Treat the full logarithmic caustic window by the endpoint
      integral of V73.
\item Use the antipodal tail identity to control
      \(r>2\sqrt\alpha+\alpha^{-1/2}\).
\item Insert the three bounds into
      \eqref{eq:V75-moment-identity} and test
      \[
      \sup_\alpha
      \|\sqrt\alpha T_\alpha'\|_{L^p}
      \le K\sqrt p.
      \]
\item Only after proving that estimate, examine tangential and
      higher derivatives; do not infer the SEDE from first order.
\end{enumerate}
\end{programme}

\begin{firewall}[Claim boundary after fourth-series V75]
\label{fire:V75-claim-boundary}
V75 proves a moment lower bound, not the matching upper bound, SEDE,
WRTTC, the Wilson all-thermal local card, the anisotropic scalar
optimization, or the complete all-order selected gate.  It does not
prove all-order MTA, WUFC, CPM, cutoff removal, reflection positivity,
Osterwalder--Schrader reconstruction, construction of the continuum
Yang--Mills measure, or a uniform positive continuum spectral gap.
The full Yang--Mills mass gap has not been proved.
\end{firewall}

\paragraph{Parent-version record.}
V75 was formed from
\texttt{Renewal\_Yang\_Mills\_Mass\_Gap\_Research\_Notes\_V74.tex}.
The V74 parent had \(32\,482\) logical source lines,
\(1\,120\,109\) bytes, and SHA--256
\[
 \texttt{3a96ec3f2426fff7078b7f023fc46d3e22d75827d42bdc66346c479cea97901b}.
\]
The next compulsory companion audit is V80.

% =============================================================================
% END APPENDED FOURTH-SERIES V75 MATERIAL
% =============================================================================

% =============================================================================
% BEGIN APPENDED FOURTH-SERIES V76 MATERIAL
% =============================================================================

\chapter{V76: Global first-derivative control by radial hazards}
\label{chap:V76-hazard}

\section{Gaussian radial hazards}
\label{sec:V76-Gaussian}

Fix the Gaussian median \(r_*>0\), so \(G(r_*)=1/2\).

\begin{lemma}[Gaussian lower and upper hazard bounds]
\label{lem:V76-Gaussian-hazards}
There is \(C<\infty\) such that
\begin{align}
 {g(r)\over G(r)}
 &\le {C\over r},
 &&0<r\le r_*,
 \label{eq:V76-Gaussian-lower-hazard}\\
 {g(r)\over1-G(r)}
 &\le C(1+r),
 &&r\ge r_*.
 \label{eq:V76-Gaussian-upper-hazard}
\end{align}
\end{lemma}

\begin{proof}
On the compact interval \((0,r_*]\),
\[
 g(r)\asymp r^2,\qquad G(r)\asymp r^3
\]
near zero, and their ratio is continuous away from zero.  This proves
\eqref{eq:V76-Gaussian-lower-hazard}.  For \(r\ge1\), integration by
parts gives
\[
 1-G(r)\asymp {g(r)\over r}
\]
up to fixed polynomial corrections.  On the remaining compact
interval \([r_*,1]\), the upper hazard is bounded.  This proves
\eqref{eq:V76-Gaussian-upper-hazard}.
\end{proof}

\section{Wilson radial hazards}
\label{sec:V76-Wilson}

Let \(\theta_{\alpha,*}\) be the Wilson radial median:
\[
 W_\alpha(\theta_{\alpha,*})={1\over2}.
\]

\begin{lemma}[Wilson median is on the parabolic scale]
\label{lem:V76-Wilson-median}
There are fixed \(0<c_*<C_*<\infty\) such that
\begin{equation}
 {c_*\over\sqrt\alpha}
 \le\theta_{\alpha,*}
 \le {C_*\over\sqrt\alpha}
 \label{eq:V76-Wilson-median}
\end{equation}
for all sufficiently large \(\alpha\).
\end{lemma}

\begin{proof}
Put \(v=\sqrt\alpha\,\theta\) in the radial density.  On every fixed
\(v\)-interval,
\[
 e^{\alpha(\cos(v/\sqrt\alpha)-1)}
 {\sin^2(v/\sqrt\alpha)\over v^2/\alpha}
 \longrightarrow e^{-v^2/2}.
\]
The Gaussian domination used in V40 makes the convergence valid for
the radial distribution functions.  The limiting density is a
positive constant times \(v^2e^{-v^2/2}\), whose median lies strictly
between zero and infinity.  Choose fixed brackets around that median.
\end{proof}

\begin{lemma}[Small-radius Wilson lower hazard]
\label{lem:V76-Wilson-lower-hazard}
For \(0<\theta\le\theta_{\alpha,*}\),
\begin{equation}
 {w_\alpha(\theta)\over W_\alpha(\theta)}
 \ge {c\over\theta}.
 \label{eq:V76-Wilson-lower-hazard}
\end{equation}
Moreover, if
\(W_\alpha(\theta)=G(r)\) with \(0<r\le r_*\), then
\begin{equation}
 \theta\asymp {r\over\sqrt\alpha}.
 \label{eq:V76-small-quantile}
\end{equation}
\end{lemma}

\begin{proof}
By \Cref{lem:V76-Wilson-median},
\(\theta\le C_*/\sqrt\alpha\).  On that interval,
\[
 e^{\alpha\cos\theta}\sin^2\theta
 \asymp e^\alpha\theta^2
\]
with comparison factors depending only on \(C_*\).  Integrating from
zero to \(\theta\) gives
\[
 W_\alpha(\theta)\asymp
 {e^\alpha\theta^3\over Z_\alpha^{\rm rad}},
\qquad
 w_\alpha(\theta)\asymp
 {e^\alpha\theta^2\over Z_\alpha^{\rm rad}},
\]
which proves the hazard bound.  Also
\[
 W_\alpha(\theta)\asymp\alpha^{3/2}\theta^3,
\qquad
 G(r)\asymp r^3
\]
for the declared median intervals.  Equality of the distribution
functions proves \eqref{eq:V76-small-quantile}.
\end{proof}

\begin{lemma}[Post-median Wilson upper hazard]
\label{lem:V76-Wilson-upper-hazard}
For every
\(\theta_{\alpha,*}\le\theta<\pi\),
\begin{equation}
 {w_\alpha(\theta)\over1-W_\alpha(\theta)}
 \ge c\sqrt\alpha.
 \label{eq:V76-Wilson-upper-hazard}
\end{equation}
\end{lemma}

\begin{proof}
The normalizer cancels in
\[
 {1-W_\alpha(\theta)\over w_\alpha(\theta)}
 =
 \int_\theta^\pi
 e^{\alpha(\cos u-\cos\theta)}
 \left({\sin u\over\sin\theta}\right)^2du.
 \label{eq:V76-hazard-integral}
\]
We bound this integral by \(C/\sqrt\alpha\).

First take
\(\theta\in[\theta_{\alpha,*},\delta]\) for a fixed small
\(\delta\).  Rescaling \(x=\sqrt\alpha\,u\), and using that the
lower endpoint is beyond a fixed positive median of
\(x^2e^{-x^2/2}\), gives the standard radial Gaussian tail-to-density
bound \(C/\sqrt\alpha\).  Uniform Taylor remainders hold on this
range.

On the fixed middle region
\([\delta,\pi-\delta]\), the exponent decreases linearly at the
lower endpoint:
\[
 \cos u-\cos\theta
 \le-c_\delta(u-\theta)
\]
until a fixed distance has been traversed; the remainder is
exponentially smaller.  Thus
\eqref{eq:V76-hazard-integral} is \(O(\alpha^{-1})\).

Finally put \(d=\pi-\theta\le\delta\).  With
\(v=\pi-u\), elementary sine and cosine comparisons give
\[
 {1-W_\alpha(\theta)\over w_\alpha(\theta)}
 \le C\min\left\{d,{1\over\alpha d}\right\}.
\]
Indeed the first bound discards the exponential on an interval of
length \(d\), while the second is the endpoint Laplace bound when
\(\alpha d^2\ge1\).  The minimum is at most
\(C/\sqrt\alpha\).  The three regions prove
\eqref{eq:V76-Wilson-upper-hazard}.
\end{proof}

\section{The global transport slope}
\label{sec:V76-slope}

\begin{theorem}[Global first-derivative transport bound]
\label{thm:V76-global-slope}
For all sufficiently large \(\alpha\) and all \(r\ge0\),
\begin{equation}
 \boxed{\quad
 T_\alpha'(r)
 \le {C(1+r)\over\sqrt\alpha}.
 \quad}
 \label{eq:V76-global-slope}
\end{equation}
\end{theorem}

\begin{proof}
If \(r\le r_*\), then
\(\theta=T_\alpha(r)\le\theta_{\alpha,*}\).  Use the lower-tail
identity
\[
 T_\alpha'(r)
 =
 {g(r)/G(r)\over
  w_\alpha(\theta)/W_\alpha(\theta)}
\]
together with
\Cref{lem:V76-Gaussian-hazards,
lem:V76-Wilson-lower-hazard}:
\[
 T_\alpha'(r)
 \le C{\theta\over r}
 \le {C\over\sqrt\alpha}.
\]

If \(r\ge r_*\), use the upper-tail identity
\[
 T_\alpha'(r)
 =
 {g(r)/(1-G(r))\over
  w_\alpha(\theta)/(1-W_\alpha(\theta))}.
\]
\Cref{lem:V76-Gaussian-hazards,
lem:V76-Wilson-upper-hazard} gives
\eqref{eq:V76-global-slope}.
\end{proof}

\begin{corollary}[Tangential transport bound]
\label{cor:V76-tangential}
For every \(r>0\),
\begin{equation}
 {T_\alpha(r)\over r}
 \le {C(1+r)\over\sqrt\alpha}.
 \label{eq:V76-tangential}
\end{equation}
\end{corollary}

\begin{proof}
Since \(T_\alpha(0)=0\), integrate
\eqref{eq:V76-global-slope}:
\[
 T_\alpha(r)
 \le {C\over\sqrt\alpha}
 \int_0^r(1+u)\,du
 \le {Cr(1+r)\over\sqrt\alpha}.
\]
\end{proof}

\section{First-order sub-Gaussian envelope}
\label{sec:V76-envelope}

\begin{theorem}[The full radial vector map has the first-order SEDE]
\label{thm:V76-first-SEDE}
For
\[
 \zeta_\alpha(y)=T_\alpha(|y|){y\over|y|},
\]
one has
\begin{equation}
 \|D\zeta_\alpha(y)\|
 \le {C(1+|y|)\over\sqrt\alpha}.
 \label{eq:V76-vector-derivative}
\end{equation}
Consequently, with
\[
 L(y):=C(1+|y|),
\]
\begin{equation}
 \sup_{\alpha\ge\alpha_0}
 \|L(Y)\|_{L^p(\gamma_3)}
 \le C_1\sqrt p
 \qquad(p\ge2),
 \label{eq:V76-first-moments}
\end{equation}
and the derivative bound
\eqref{eq:V74-envelope-derivative} holds at \(k=1\).
\end{theorem}

\begin{proof}
The derivative of a radial vector map has radial eigenvalue
\(T_\alpha'(r)\) and two tangential eigenvalues
\(T_\alpha(r)/r\).  Apply
\Cref{thm:V76-global-slope,cor:V76-tangential}, and use
\(s_\alpha\asymp\alpha^{-1}\).
The radial part of a standard three-dimensional Gaussian has
\[
 \||Y|\|_{L^p}\le C\sqrt p,
\]
which follows directly from its Gamma-function moments.  This proves
\eqref{eq:V76-first-moments}.  Composing with a representation uses
one generator and gives the factor
\((s_\alpha(1+C_2(R)))^{1/2}\), as in V67.
\end{proof}

\begin{assessment}[First-order status]
\label{ass:V76-status}
The upper bound \(K\sqrt p\) and the V75 lower bound
\(\sqrt{p/\log p}\) now bracket the optimal first-derivative moment
growth within a logarithm.  Thus the antipodal caustic does not
obstruct the first-order SEDE.  No conclusion about \(k\ge2\)
follows: differentiating the density-ratio ODE can introduce
derivatives of both hazard functions and potentially a vertex
factorial.
\end{assessment}

\section{V76 ledger and V77 programme}
\label{sec:V76-results}

\begin{theorem}[Fourth-series V76 result ledger]
\label{thm:V76-results}
V76 proves matched Gaussian and Wilson radial hazard bounds, the
global slope estimate
\(T_\alpha'(r)\le C(1+r)/\sqrt\alpha\), the corresponding tangential
bound, and the complete first-order sub-Gaussian transport envelope.
\end{theorem}

\begin{proof}
Use
\Cref{lem:V76-Gaussian-hazards,
lem:V76-Wilson-lower-hazard,
lem:V76-Wilson-upper-hazard,
thm:V76-global-slope,
thm:V76-first-SEDE}.
\end{proof}

\begin{programme}[V77: second derivative and recursive hazard costs]
\label{prog:V76-V77}
\begin{enumerate}[leftmargin=3em]
\item Differentiate
      \(T_\alpha'=g/(w_\alpha\circ T_\alpha)\) exactly and express
      \(T_\alpha''\) through the two log-density scores.
\item Use the quantile relation to expose cancellation between the
      Gaussian and Wilson scores in the parabolic region.
\item Bound radial and tangential Hessian components by a common
      sub-Gaussian envelope at scale \(\alpha^{-1}\).
\item Track whether differentiating again creates Bell/factorial
      growth; do not infer all orders from the Hessian.
\item If the Hessian already violates exponential derivative-order
      dependence, stop the absolute transport route and return to the
      intrinsic Wilson forest formula.
\end{enumerate}
\end{programme}

\begin{firewall}[Claim boundary after fourth-series V76]
\label{fire:V76-claim-boundary}
V76 proves only the first-order SEDE.  It does not prove the
all-order SEDE, WRTTC, the Wilson all-thermal local card, the
anisotropic scalar optimization, or the complete all-order selected
gate.  It does not prove all-order MTA, WUFC, CPM, cutoff removal,
reflection positivity, Osterwalder--Schrader reconstruction,
construction of the continuum Yang--Mills measure, or a uniform
positive continuum spectral gap.  The full Yang--Mills mass gap has
not been proved.
\end{firewall}

\paragraph{Parent-version record.}
V76 was formed from
\texttt{Renewal\_Yang\_Mills\_Mass\_Gap\_Research\_Notes\_V75.tex}.
The V75 parent had \(32\,803\) logical source lines,
\(1\,129\,787\) bytes, and SHA--256
\[
 \texttt{0b4f3ecc67f191aaf0906e9d087426e9803e9bc5c68c85f14219dd0ff9256a20}.
\]
The next compulsory companion audit is V80.

% =============================================================================
% END APPENDED FOURTH-SERIES V76 MATERIAL
% =============================================================================

% =============================================================================
% BEGIN APPENDED FOURTH-SERIES V77 MATERIAL
% =============================================================================

\chapter{V77: The second-derivative caustic rules out the SEDE}
\label{chap:V77-second-caustic}

\section{Exact score identity}
\label{sec:V77-score}

Let
\[
 a(r):={d\over dr}\log g(r)={2\over r}-r,
\]
and
\[
 b_\alpha(\theta):=
 {d\over d\theta}\log w_\alpha(\theta)
 =-\alpha\sin\theta+2\cot\theta.
\]

\begin{lemma}[Second derivative of the quantile]
\label{lem:V77-second-identity}
For every \(r>0\),
\begin{equation}
 \boxed{\quad
 T_\alpha''(r)
 =
 T_\alpha'(r)
 \left[
 a(r)-
 b_\alpha(T_\alpha(r))T_\alpha'(r)
 \right].
 \quad}
 \label{eq:V77-second-identity}
\end{equation}
\end{lemma}

\begin{proof}
Take the logarithmic derivative of
\[
 T_\alpha'(r)
 ={g(r)\over w_\alpha(T_\alpha(r))}.
\]
The chain rule gives
\[
 {T_\alpha''\over T_\alpha'}
 =(\log g)'-
 (\log w_\alpha)'(T_\alpha)T_\alpha',
\]
which is the claim.
\end{proof}

\section{Two hazard expansions at the critical layer}
\label{sec:V77-expansions}

Write upper hazards
\[
 H_G(r):={g(r)\over1-G(r)},
 \qquad
 H_{W,\alpha}(\theta):=
 {w_\alpha(\theta)\over1-W_\alpha(\theta)}.
\]
At matched quantiles,
\begin{equation}
 T_\alpha'(r)
 ={H_G(r)\over H_{W,\alpha}(T_\alpha(r))}.
 \label{eq:V77-hazard-ratio}
\end{equation}

\begin{lemma}[Critical hazard expansions]
\label{lem:V77-hazard-expansions}
Uniformly for \(r\in I_\alpha\), put
\[
 D(r):=\sqrt\alpha\{\pi-T_\alpha(r)\}.
\]
Then
\begin{align}
 H_G(r)
 &=r-r^{-1}+O(r^{-3}),
 \label{eq:V77-Gaussian-hazard-expansion}\\
 H_{W,\alpha}(T_\alpha(r))
 &=
 \sqrt\alpha
 \left[
 D(r)+D(r)^{-1}
 +O\left(D(r)^{-3}+{D(r)^3\over\alpha}\right)
 \right],
 \label{eq:V77-Wilson-hazard-expansion}\\
 b_\alpha(T_\alpha(r))
 &=
 -\sqrt\alpha
 \left[
 D(r)+{2\over D(r)}
 +O\left({D(r)^3\over\alpha}\right)
 \right].
 \label{eq:V77-score-expansion}
\end{align}
\end{lemma}

\begin{proof}
The Gaussian Mills expansion gives
\[
 1-G(r)
 =c e^{-r^2/2}
 \left(r+r^{-1}+O(r^{-3})\right),
\]
while
\(g(r)=c r^2e^{-r^2/2}\).  Division proves
\eqref{eq:V77-Gaussian-hazard-expansion}.

For the Wilson hazard, V73 reduces the upper tail to
\[
 F(D):=\int_0^De^{v^2/2}v^2\,dv.
\]
Repeated endpoint integration by parts gives
\[
 F(D)
 =
 e^{D^2/2}
 \left(D-D^{-1}+O(D^{-3})\right).
\]
The endpoint density in the \(D\)-coordinate is
\[
 e^{D^2/2}D^2
 \left(1+O(D^4/\alpha)\right),
\]
and conversion from \(D\) to \(\theta\) contributes
\(\sqrt\alpha\).  Division proves
\eqref{eq:V77-Wilson-hazard-expansion}.  Finally, with
\(\theta=\pi-D/\sqrt\alpha\),
\[
 -\alpha\sin\theta+2\cot\theta
 =
 -\sqrt\alpha
 \left(D+{2\over D}\right)
 +O(D^3/\sqrt\alpha),
\]
which is \eqref{eq:V77-score-expansion}.  The V75 interval theorem
gives
\[
 D(r)^2=\log\alpha+O(\log\log\alpha)
\]
uniformly, so every displayed remainder is uniform.
\end{proof}

\begin{corollary}[Two-term critical slope]
\label{cor:V77-two-term-slope}
Uniformly on \(I_\alpha\),
\begin{equation}
 T_\alpha'(r)
 =
 {2\over D(r)}
 \left[
 1-{1\over D(r)^2}
 +O\left(D(r)^{-4}+{D(r)^2\over\alpha}\right)
 \right].
 \label{eq:V77-two-term-slope}
\end{equation}
\end{corollary}

\begin{proof}
Use \(r=2\sqrt\alpha+O(\alpha^{-1/2})\) in
\eqref{eq:V77-Gaussian-hazard-expansion}, divide by
\eqref{eq:V77-Wilson-hazard-expansion}, and expand
\((1+D^{-2})^{-1}\).
\end{proof}

\section{The second-derivative layer}
\label{sec:V77-second-layer}

\begin{theorem}[Critical second derivative]
\label{thm:V77-critical-second}
Uniformly for \(r\in I_\alpha\),
\begin{equation}
 \boxed{\quad
 T_\alpha''(r)
 \asymp
 {\sqrt\alpha\over(\log\alpha)^{3/2}}.
 \quad}
 \label{eq:V77-critical-second}
\end{equation}
In particular the second derivative has a fixed sign throughout this
interval for sufficiently large \(\alpha\).
\end{theorem}

\begin{proof}
Insert
\eqref{eq:V77-score-expansion} and
\eqref{eq:V77-two-term-slope} into the bracket in
\eqref{eq:V77-second-identity}.  Since
\[
 a(r)=-2\sqrt\alpha+O(\alpha^{-1/2}),
\]
one obtains
\begin{align*}
 b_\alpha(T_\alpha(r))T_\alpha'(r)
 &=
 -2\sqrt\alpha
 \left[
 1+{1\over D(r)^2}
 +O\left(D(r)^{-4}+{D(r)^2\over\alpha}\right)
 \right].
\end{align*}
Therefore
\[
 a(r)-b_\alpha(T_\alpha(r))T_\alpha'(r)
 =
 {2\sqrt\alpha\over D(r)^2}
 \left(1+o(1)\right).
\]
Multiplication by
\(T_\alpha'(r)=2D(r)^{-1}(1+o(1))\) gives
\[
 T_\alpha''(r)
 ={4\sqrt\alpha\over D(r)^3}(1+o(1)).
\]
Use \(D(r)\asymp\sqrt{\log\alpha}\).
\end{proof}

\section{Failure of the common sub-Gaussian derivative envelope}
\label{sec:V77-SEDE-failure}

\begin{lemma}[A fixed representation detects the radial Hessian]
\label{lem:V77-fundamental-detects}
Let \(R_*\) be a fixed nontrivial \(SU(2)\) representation and
restrict \(D^{R_*}\circ\Phi_\alpha\) to a fixed radial ray.  There is
\(c_*>0\) such that
\begin{equation}
 \left\|
 {d^2\over dr^2}
 D^{R_*}(\exp(T_\alpha(r)X))
 \right\|_{\rm op}
 \ge c_*|T_\alpha''(r)|
 \label{eq:V77-detection}
\end{equation}
for a normalized nonzero generator \(X\).
\end{lemma}

\begin{proof}
Along the ray,
\[
 {d^2\over dr^2}e^{T_\alpha(r)dR_*(X)}
 =
 e^{T_\alpha dR_*(X)}
 \left[
 dR_*(X)T_\alpha''
 +dR_*(X)^2(T_\alpha')^2
 \right].
\]
In an irreducible \(SU(2)\) weight basis, the first term is
skew-adjoint and the second self-adjoint.  On a nonzero weight vector
their scalar components are respectively imaginary and real, so the
modulus of their sum is at least a fixed nonzero weight times
\(|T_\alpha''|\).  Left multiplication by the unitary exponential
preserves the norm.
\end{proof}

\begin{theorem}[No SEDE for the radial quantile transport]
\label{thm:V77-no-SEDE}
There is no family \(L_\alpha\) satisfying both the sub-Gaussian
moment condition
\[
 \sup_\alpha\|L_\alpha\|_{L^p(\gamma_3)}
 \le K\sqrt p
\]
and the derivative estimate
\eqref{eq:V74-envelope-derivative} for every \(k\), even when that
estimate is required only at \(k=2\) and for one fixed nontrivial
representation.
\end{theorem}

\begin{proof}
For the fixed representation in
\Cref{lem:V77-fundamental-detects}, the deterministic second-order
factor in \eqref{eq:V74-envelope-derivative} is comparable with
\(\alpha^{-1}L_\alpha(y)^2\).  Thus on
\(|y|\in I_\alpha\),
\Cref{thm:V77-critical-second} forces
\[
 L_\alpha(y)
 \ge
 c\left(\alpha|T_\alpha''(|y|)|\right)^{1/2}
 \ge
 {c\alpha^{3/4}\over(\log\alpha)^{3/4}}.
\]
The interval has Gaussian mass
\(\asymp\sqrt\alpha e^{-2\alpha}\) by V75.  Hence
\[
 \|L_\alpha\|_{L^p}
 \ge
 c\left(\sqrt\alpha e^{-2\alpha}\right)^{1/p}
 {\alpha^{3/4}\over(\log\alpha)^{3/4}}.
\]
Choose \(\alpha=p/4\).  For large \(p\), this gives
\begin{equation}
 \sup_\alpha\|L_\alpha\|_{L^p}
 \ge
 c\left({p\over\log p}\right)^{3/4},
 \label{eq:V77-envelope-lower}
\end{equation}
which eventually exceeds every \(K\sqrt p\).
\end{proof}

\begin{corollary}[The V74 sufficient route is closed]
\label{cor:V77-route-closed}
The SEDE-to-WRTTC theorem
\Cref{thm:V74-SEDE-WRTTC} cannot prove the Wilson local card for the
radial quantile transport.
\end{corollary}

\begin{firewall}[The WRTTC itself is not disproved]
\label{fire:V77-WRTTC-open}
V77 rules out domination by one absolute sub-Gaussian envelope.
It does not rule out the WRTTC, which may retain cancellations among
higher transport derivatives inside the operator-valued Gaussian
forest integral.  Nor does it rule out a nonradial transport.
Nevertheless, continuing to differentiate the radial quantile and
take separate absolute derivative envelopes would repeat a disproved
route.
\end{firewall}

\section{V77 ledger and V78 programme}
\label{sec:V77-results}

\begin{theorem}[Fourth-series V77 result ledger]
\label{thm:V77-results}
V77 derives the exact quantile score identity, proves the two-term
hazard expansion in the critical layer, obtains
\[
 T_\alpha''\asymp
 \sqrt\alpha(\log\alpha)^{-3/2},
\]
and proves that the common sub-Gaussian derivative envelope required
by V74 is impossible already at second order.
\end{theorem}

\begin{proof}
Use
\Cref{lem:V77-second-identity,
lem:V77-hazard-expansions,
thm:V77-critical-second,
thm:V77-no-SEDE}.
\end{proof}

\begin{programme}[V78: intrinsic Wilson connected identity]
\label{prog:V77-V78}
\begin{enumerate}[leftmargin=3em]
\item Stop the absolute radial-transport derivative programme.
\item On the full compact group, use Haar integration by parts for
      the normalized Wilson density to derive an exact score identity
      with no chart complement.
\item Introduce the reversible Wilson diffusion and its mean-zero
      resolvent only after proving the corresponding one-link
      spectral/Poincaré bound.
\item Express covariance as a carré-du-champ resolvent and determine
      whether differentiation of the resolvent can be organized by
      input trees without a vertex factorial.
\item Test the construction on the exact \(SU(2)\) Bessel channels
      of V40 before making an all-order claim.
\end{enumerate}
\end{programme}

\begin{firewall}[Claim boundary after fourth-series V77]
\label{fire:V77-claim-boundary}
V77 disproves the SEDE route but does not prove an intrinsic Wilson
forest identity, the Wilson all-thermal local card, the anisotropic
scalar optimization, or the complete all-order selected gate.  It
does not prove all-order MTA, WUFC, CPM, cutoff removal, reflection
positivity, Osterwalder--Schrader reconstruction, construction of the
continuum Yang--Mills measure, or a uniform positive continuum
spectral gap.  The full Yang--Mills mass gap has not been proved.
\end{firewall}

\paragraph{Parent-version record.}
V77 was formed from
\texttt{Renewal\_Yang\_Mills\_Mass\_Gap\_Research\_Notes\_V76.tex}.
The V76 parent had \(33\,165\) logical source lines,
\(1\,140\,127\) bytes, and SHA--256
\[
 \texttt{27873f642b70cf6515dbc1995f8bea2ade1c9b13aec3ae687e305edbf64455ac}.
\]
The next compulsory companion audit is V80.

% =============================================================================
% END APPENDED FOURTH-SERIES V77 MATERIAL
% =============================================================================

% =============================================================================
% BEGIN APPENDED FOURTH-SERIES V78 MATERIAL
% =============================================================================

\chapter{V78: Intrinsic Wilson diffusion and the one-link gap}
\label{chap:V78-Wilson-diffusion}

\section{Purpose and nonduplication}
\label{sec:V78-purpose}

V77 proved that the radial quantile transport has a genuine
second-derivative caustic and therefore cannot satisfy the absolute
sub-Gaussian derivative envelope proposed in V74.  The present note
does not weaken that envelope or differentiate the transport again.
It returns to the compact group and uses the Wilson density itself.

Earlier reduced resolvents in V1--V64 are Peter--Weyl multipliers
associated with one-link convolution.  The operator introduced below
is different: it is the local weighted Laplacian whose invariant law
is the one-link Wilson measure.  Thus the Poincare estimate and
carre-du-champ formula proved here do not duplicate the previous
fusion-multiplier calculations.  Their purpose is to construct an
intrinsic connected identity with no chart complement and no
transport derivative.

\section{Haar score and reversible generator}
\label{sec:V78-score-generator}

Retain the Wilson law \eqref{eq:V40-Wilson-law} and put
\[
 \phi(U):={1\over2}\operatorname{ReTr}U.
\]
Choose left-invariant vector fields \(X_1,X_2,X_3\) orthonormal in
the Casimir normalization
\[
 -\sum_{a=1}^3X_a^2\big|_{R_j}=c_j=j(j+1).
\]
For this normalization \(SU(2)\) is the round three-sphere of radius
two.  Define
\[
 S_a:=-\alpha X_a\phi.
\]

\begin{lemma}[Exact Wilson score identity]
\label{lem:V78-score}
For every smooth scalar \(f\),
\begin{equation}
 \boxed{\qquad
 \mathbb E_{\nu_\alpha}[fS_a]
 =
 \mathbb E_{\nu_\alpha}[X_af].
 \qquad}
 \label{eq:V78-score}
\end{equation}
In particular \(\mathbb E_{\nu_\alpha}S_a=0\).
The identity also holds entrywise for finite-dimensional
operator-valued \(f\).
\end{lemma}

\begin{proof}
Left-invariant vector fields have zero Haar divergence.  Therefore
\[
 0
 =
 \int_{SU(2)}X_a\{f e^{\alpha\phi}\}\,dU
 =
 Z_\alpha
 \left(
 \mathbb E_{\nu_\alpha}X_af
 +
 \alpha\mathbb E_{\nu_\alpha}[fX_a\phi]
 \right).
\]
Move the second term to the other side and use the definition of
\(S_a\).  Taking \(f=1\) proves centering.  Applying the scalar
identity to matrix entries proves the last assertion.
\end{proof}

In \(L^2(\nu_\alpha)\),
\[
 X_a^{*,\alpha}=-X_a-\alpha X_a\phi.
\]
Set
\begin{equation}
 A_\alpha
 :=
 \sum_{a=1}^3X_a^{*,\alpha}X_a
 =
 -\sum_aX_a^2-\alpha\sum_a(X_a\phi)X_a.
 \label{eq:V78-generator}
\end{equation}

\begin{lemma}[Dirichlet form and kernel]
\label{lem:V78-Dirichlet}
The closure of \(A_\alpha\) is nonnegative and self-adjoint, and
\begin{equation}
 \langle f,A_\alpha g\rangle_{\nu_\alpha}
 =
 \sum_{a=1}^3
 \langle X_af,X_ag\rangle_{\nu_\alpha}.
 \label{eq:V78-Dirichlet}
\end{equation}
Its kernel consists exactly of the constants.
\end{lemma}

\begin{proof}
The displayed identity follows from the definition by adjoints.
It gives nonnegativity and the Friedrichs self-adjoint realization.
If its quadratic form vanishes, every weak derivative \(X_af\)
vanishes.  The fields span the tangent bundle and \(SU(2)\) is
connected, so \(f\) is constant.  Constants plainly lie in the
kernel.
\end{proof}

\section{A full-group Poincare inequality at the Wilson scale}
\label{sec:V78-Poincare}

Write
\[
 U=\cos\theta\,{\bf1}
   +i\sin\theta\,\omega\cdot\sigma,
 \qquad
 0\le\theta\le\pi,\quad \omega\in S^2.
\]
Then
\[
 d\nu_\alpha=d\mu_\alpha(\theta)\,d\sigma(\omega),
 \qquad
 d\mu_\alpha=w_\alpha(\theta)\,d\theta,
\]
where \(\sigma\) is normalized uniform measure on \(S^2\).

\begin{lemma}[One-dimensional Cheeger consequence]
\label{lem:V78-Cheeger}
Let \(p\,dx\) be a probability measure on an interval, with
distribution function \(F\), and let \(m\) be a median.  If
\[
 {p(x)\over F(x)}\ge h\quad(x\le m),
 \qquad
 {p(x)\over1-F(x)}\ge h\quad(x\ge m),
\]
then every locally absolutely continuous \(u\) satisfies
\begin{equation}
 \operatorname{Var}(u)
 \le {4\over h^2}\int|u'|^2p\,dx.
 \label{eq:V78-Cheeger-Poincare}
\end{equation}
\end{lemma}

\begin{proof}
The two hazard inequalities give the one-dimensional Cheeger
isoperimetric bound
\[
 \mu^+(E)\ge h\min\{\mu(E),1-\mu(E)\}.
\]
Indeed, decomposing an open set into intervals and summing its
weighted boundary points reduces the assertion to the two displayed
tail inequalities; approximation gives the general case.  The
coarea formula consequently gives
\[
 \int|v-\operatorname{med}v|\,d\mu
 \le h^{-1}\int|v'|\,d\mu.
\]
Let \(c\) be a median of \(u\) and apply this inequality to
\(v=(u-c)|u-c|\), whose median is zero.  Cauchy--Schwarz gives
\[
 \int|u-c|^2d\mu
 \le {2\over h}\int|u-c||u'|d\mu
 \le {2\over h}
 \left(\int|u-c|^2d\mu\right)^{1/2}
 \left(\int|u'|^2d\mu\right)^{1/2}.
\]
After division and squaring, use
\(\operatorname{Var}(u)\le\int|u-c|^2d\mu\).
\end{proof}

\begin{lemma}[Radial Wilson Poincare inequality]
\label{lem:V78-radial-Poincare}
For all sufficiently large \(\alpha\),
\begin{equation}
 \operatorname{Var}_{\mu_\alpha}(u)
 \le {C\over\alpha}
 \int_0^\pi|u'(\theta)|^2\,d\mu_\alpha(\theta).
 \label{eq:V78-radial-Poincare}
\end{equation}
\end{lemma}

\begin{proof}
For \(0<\theta\le\theta_{\alpha,*}\),
\Cref{lem:V76-Wilson-lower-hazard} and
\(\theta_{\alpha,*}\le C_*/\sqrt\alpha\) give
\[
 {w_\alpha(\theta)\over W_\alpha(\theta)}
 \ge {c\over\theta}
 \ge c_1\sqrt\alpha.
\]
For \(\theta\ge\theta_{\alpha,*}\),
\Cref{lem:V76-Wilson-upper-hazard} gives
\[
 {w_\alpha(\theta)\over1-W_\alpha(\theta)}
 \ge c_2\sqrt\alpha.
\]
Apply \Cref{lem:V78-Cheeger} with
\(h=\min(c_1,c_2)\sqrt\alpha\).
\end{proof}

\begin{lemma}[Exact transverse moment]
\label{lem:V78-transverse-moment}
For every \(\alpha>0\),
\begin{equation}
 \boxed{\qquad
 \alpha\,
 \mathbb E_{\mu_\alpha}\sin^2\theta
 =
 3\mathbb E_{\mu_\alpha}\cos\theta.
 \qquad}
 \label{eq:V78-transverse-moment}
\end{equation}
Consequently
\[
 \mathbb E_{\mu_\alpha}\sin^2\theta\le {3\over\alpha}.
\]
\end{lemma}

\begin{proof}
The boundary values of
\(e^{\alpha\cos\theta}\sin^3\theta\) vanish.  Hence
\[
 0
 =
 \int_0^\pi
 {d\over d\theta}
 \left(e^{\alpha\cos\theta}\sin^3\theta\right)d\theta
 =
 \int_0^\pi e^{\alpha\cos\theta}
 \left(3\sin^2\theta\cos\theta
       -\alpha\sin^4\theta\right)d\theta.
\]
Divide by the radial normalizer.  The upper bound follows from
\(\cos\theta\le1\).
\end{proof}

\begin{theorem}[One-link Wilson spectral gap]
\label{thm:V78-Wilson-gap}
There are \(c>0\) and \(\alpha_0<\infty\) such that, for
\(\alpha\ge\alpha_0\) and every smooth \(f:SU(2)\to\mathbb C\),
\begin{equation}
 \boxed{\qquad
 \operatorname{Var}_{\nu_\alpha}(f)
 \le {C\over\alpha}
 \sum_{a=1}^3
 \mathbb E_{\nu_\alpha}|X_af|^2.
 \qquad}
 \label{eq:V78-full-Poincare}
\end{equation}
Equivalently, the first nonzero eigenvalue of \(A_\alpha\) obeys
\begin{equation}
 \lambda_1(A_\alpha)\ge c\alpha.
 \label{eq:V78-gap}
\end{equation}
\end{theorem}

\begin{proof}
Expand in an orthonormal spherical-harmonic basis on \(S^2\):
\[
 f(\theta,\omega)
 =
 \sum_{\ell=0}^\infty\sum_{m=-\ell}^{\ell}
 u_{\ell m}(\theta)Y_{\ell m}(\omega),
 \qquad
 -\Delta_{S^2}Y_{\ell m}
 =\ell(\ell+1)Y_{\ell m}.
\]
The Casimir normalization gives
\begin{equation}
 \sum_a|X_af|^2
 =
 {1\over4}
 \left(
 |\partial_\theta f|^2
 +{|\nabla_{S^2}f|^2\over\sin^2\theta}
 \right).
 \label{eq:V78-polar-energy}
\end{equation}
For the constant spherical mode,
\Cref{lem:V78-radial-Poincare} bounds its variance by
\(C\alpha^{-1}\int|u_{00}'|^2d\mu_\alpha\).

Fix now \(\ell\ge1\) and abbreviate \(u=u_{\ell m}\).
The radial Poincare inequality gives
\[
 \int|u|^2d\mu_\alpha
 \le
 {C\over\alpha}\int|u'|^2d\mu_\alpha
 +\left|\int u\,d\mu_\alpha\right|^2.
\]
By Cauchy--Schwarz and
\Cref{lem:V78-transverse-moment},
\[
 \left|\int u\,d\mu_\alpha\right|^2
 =
 \left|
 \int {u\over\sin\theta}\sin\theta\,d\mu_\alpha
 \right|^2
 \le
 {3\over\alpha}
 \int {|u|^2\over\sin^2\theta}\,d\mu_\alpha.
\]
Since \(\ell(\ell+1)\ge2\), this is bounded by
\[
 {C\over\alpha}
 \int
 \left(
 |u'|^2+
 {\ell(\ell+1)|u|^2\over\sin^2\theta}
 \right)d\mu_\alpha.
\]
Sum over \((\ell,m)\), use Parseval and
\eqref{eq:V78-polar-energy}, and absorb the fixed factor four.
Smooth approximation handles the polar endpoints.  This proves
\eqref{eq:V78-full-Poincare}; the variational characterization of
the first nonzero eigenvalue gives \eqref{eq:V78-gap}.
\end{proof}

\begin{remark}[Why global Bakry--Emery curvature is not the proof]
\label{rem:V78-BE-fails}
In the radius-two metric,
\[
 \operatorname{Ric}={1\over2}g,
 \qquad
 \operatorname{Hess}\phi=-{1\over4}\phi g.
\]
Writing \(d\nu_\alpha=e^{-V}dU/Z_\alpha\) gives
\(V=-\alpha\phi\), and hence
\[
 \operatorname{Ric}+\operatorname{Hess}V
 =
 \left({1\over2}+{\alpha\phi\over4}\right)g.
\]
At the antipode this equals
\((2-\alpha)g/4\), which is negative for \(\alpha>2\).
Thus the gap in \Cref{thm:V78-Wilson-gap} is a localization and
mode-decomposition fact, not a consequence of a false global
positive-curvature assertion.
\end{remark}

\section{Mean-zero resolvent and connected identity}
\label{sec:V78-resolvent}

Let
\[
 Q_\alpha f=f-\mathbb E_{\nu_\alpha}f,
 \qquad
 R_\alpha=A_\alpha^{-1}Q_\alpha.
\]
The inverse is understood on the orthogonal complement of constants.

\begin{theorem}[Intrinsic Wilson covariance identity]
\label{thm:V78-covariance}
For smooth scalar \(f,g\),
\begin{equation}
 \boxed{\qquad
 \operatorname{Cov}_{\nu_\alpha}(f,g)
 =
 \sum_{a=1}^3
 \left\langle X_af,X_aR_\alpha g\right\rangle_{\nu_\alpha}.
 \qquad}
 \label{eq:V78-covariance}
\end{equation}
Moreover,
\begin{align}
 \|R_\alpha g\|_2
 &\le {C\over\alpha}\|Q_\alpha g\|_2,
 \label{eq:V78-resolvent-L2}\\
 \left(
 \sum_a\|X_aR_\alpha g\|_2^2
 \right)^{1/2}
 &\le {C\over\sqrt\alpha}\|Q_\alpha g\|_2.
 \label{eq:V78-resolvent-gradient}
\end{align}
All three identities extend by polarization to finite-dimensional
operator entries.
\end{theorem}

\begin{proof}
Since \(A_\alpha R_\alpha g=Q_\alpha g\),
\Cref{lem:V78-Dirichlet} gives
\[
 \langle Q_\alpha f,Q_\alpha g\rangle
 =
 \langle f,A_\alpha R_\alpha g\rangle
 =
 \sum_a\langle X_af,X_aR_\alpha g\rangle.
\]
This is \eqref{eq:V78-covariance}.  The spectral theorem and
\eqref{eq:V78-gap} give
\[
 \|R_\alpha Q_\alpha\|_{2\to2}\le {C\over\alpha},
\]
which is \eqref{eq:V78-resolvent-L2}.  Finally,
\[
 \sum_a\|X_aR_\alpha g\|_2^2
 =
 \langle R_\alpha g,Q_\alpha g\rangle
 \le {C\over\alpha}\|Q_\alpha g\|_2^2.
\]
Entrywise application and polarization prove the last assertion.
\end{proof}

\begin{corollary}[Semigroup form]
\label{cor:V78-semigroup}
If \(P_t^\alpha=e^{-tA_\alpha}\), then
\begin{equation}
 \operatorname{Cov}_{\nu_\alpha}(f,g)
 =
 \int_0^\infty
 \sum_a
 \left\langle X_af,X_aP_t^\alpha Q_\alpha g\right\rangle
 dt.
 \label{eq:V78-semigroup}
\end{equation}
The integral converges absolutely after Cauchy--Schwarz in the
Dirichlet form and the spectral gap.
\end{corollary}

\begin{proof}
The spectral theorem gives
\[
 R_\alpha=\int_0^\infty P_t^\alpha Q_\alpha\,dt
\]
strongly on \(L^2\).  Insert this identity into
\eqref{eq:V78-covariance}.  The gap supplies exponential decay on
the mean-zero subspace and justifies passage under the integral.
\end{proof}

\section{The exact first differentiated-resolvent equation}
\label{sec:V78-differentiated-resolvent}

\begin{lemma}[One derivative of the Wilson resolvent]
\label{lem:V78-resolvent-derivative}
For a smooth \(q\), put \(u=R_\alpha q\).  Then
\begin{equation}
 \boxed{\qquad
 X_bu
 =
 \mathbb E_{\nu_\alpha}[uS_b]
 +
 R_\alpha
 \left(
 X_bq+[A_\alpha,X_b]u
 \right).
 \qquad}
 \label{eq:V78-resolvent-derivative}
\end{equation}
The expression in parentheses is automatically centered.  In the
left-invariant frame,
\begin{equation}
 [A_\alpha,X_b]
 =
 \alpha\sum_a
 \left\{
 (X_bX_a\phi)X_a
 -(X_a\phi)[X_a,X_b]
 \right\}.
 \label{eq:V78-commutator}
\end{equation}
\end{lemma}

\begin{proof}
The Casimir \(\sum_aX_a^2\) commutes with every left-invariant
\(X_b\).  Expanding the drift part of
\eqref{eq:V78-generator} therefore gives
\eqref{eq:V78-commutator}.  Since
\(A_\alpha u=Q_\alpha q\),
\[
 A_\alpha X_bu
 =
 X_bA_\alpha u+[A_\alpha,X_b]u
 =
 X_bq+[A_\alpha,X_b]u.
\]
The right-hand side lies in the range of \(A_\alpha\), hence is
centered.  Applying \(R_\alpha\) recovers the mean-zero part of
\(X_bu\).  Its constant part is
\[
 \mathbb E_{\nu_\alpha}X_bu
 =
 \mathbb E_{\nu_\alpha}[uS_b]
\]
by \Cref{lem:V78-score}.  Adding the two parts proves
\eqref{eq:V78-resolvent-derivative}.
\end{proof}

\begin{assessment}[The new all-order bottleneck]
\label{ass:V78-bottleneck}
The factors in \eqref{eq:V78-resolvent-derivative} have the correct
first scaling: the commutator carries one factor \(\alpha\), while
each mean-zero resolvent carries \(\alpha^{-1}\) in \(L^2\), and the
gradient resolvent carries \(\alpha^{-1/2}\).  Derivatives of
\(\phi\) grow only exponentially with their order because \(\phi\)
is a fixed matrix coefficient.  What is not yet proved is that
iterating \eqref{eq:V78-resolvent-derivative} can be indexed by input
trees with only exponential local multiplicity.  A naive Leibniz
expansion may still create a Bell or factorial loss.  This is the
precise successor of the transport bottleneck, now expressed in
intrinsic compact-group terms.
\end{assessment}

\section{Exact Bessel-channel consistency check}
\label{sec:V78-Bessel-check}

\begin{proposition}[The transverse identity is the fixed-channel
Bessel recurrence]
\label{prop:V78-Bessel-check}
The exact moment in \eqref{eq:V78-transverse-moment} satisfies
\begin{equation}
 \mathbb E_{\nu_\alpha}\cos\theta
 =\eta_{1/2}(\alpha)
 ={I_2(\alpha)\over I_1(\alpha)},
 \qquad
 \mathbb E_{\nu_\alpha}\sin^2\theta
 ={3\over\alpha}{I_2(\alpha)\over I_1(\alpha)}.
 \label{eq:V78-Bessel-check}
\end{equation}
In particular the full-group gap proof uses the exact Wilson clock
\(\mathbb E\sin^2\theta\sim3/\alpha\), consistently with the
calibration \(s_\alpha\sim2/\alpha\) in V40.
\end{proposition}

\begin{proof}
For the fundamental character,
\(\chi_{1/2}(\theta)=2\cos\theta\) and \(d_{1/2}=2\).
Thus \eqref{eq:V40-eta-def} gives
\(\mathbb E\cos\theta=\eta_{1/2}\), and
\Cref{lem:V40-Bessel-formula} gives the Bessel ratio.
The second identity is
\eqref{eq:V78-transverse-moment}.  The fixed-index asymptotic in
\Cref{lem:V40-transition-Bessel} makes the ratio tend to one.
\end{proof}

\section{V78 ledger and V79 programme}
\label{sec:V78-results}

\begin{theorem}[Fourth-series V78 result ledger]
\label{thm:V78-results}
V78 proves:
\begin{enumerate}[label=\textup{(V78.\arabic*)},leftmargin=6.3em]
\item the exact full-group Wilson score identity;
\item a radial Cheeger inequality at scale \(\sqrt\alpha\);
\item the exact transverse moment
      \(\alpha\mathbb E\sin^2\theta=3\mathbb E\cos\theta\);
\item the full one-link Poincare inequality and diffusion gap
      \(\lambda_1(A_\alpha)\ge c\alpha\);
\item the exact intrinsic covariance and semigroup-resolvent
      identities with resolvent scales
      \(\alpha^{-1}\) and \(\alpha^{-1/2}\);
\item the first differentiated-resolvent equation, including its
      centered commutator insertion; and
\item agreement of the scale with the exact Bessel channels of V40.
\end{enumerate}
\end{theorem}

\begin{proof}
Use
\Cref{lem:V78-score,
lem:V78-radial-Poincare,
lem:V78-transverse-moment,
thm:V78-Wilson-gap,
thm:V78-covariance,
cor:V78-semigroup,
lem:V78-resolvent-derivative,
prop:V78-Bessel-check}.
\end{proof}

\begin{programme}[V79: tree expansion of differentiated Wilson
resolvents]
\label{prog:V78-V79}
\begin{enumerate}[leftmargin=3em]
\item Derive a multilinear recursion from
      \eqref{eq:V78-resolvent-derivative} in which each new input
      either differentiates an observable, creates one commutator
      vertex, or closes one centered score leaf.
\item Prove an energy estimate for
      \(R_\alpha[A_\alpha,X_b]R_\alpha\) that retains the
      \(\alpha^{-1/2}\) gradient debit instead of bounding the
      coefficient and resolvents separately.
\item Quotient recursion histories by their underlying rooted input
      tree and determine whether the remaining local ordering count
      is exponential or factorial.
\item Test the second and third connected derivatives explicitly on
      fundamental and transition Bessel channels.
\item Only after the tree census closes, formulate an intrinsic
      all-order Wilson local card and return it to the V64 thermal
      endpoint ledger.
\end{enumerate}
\end{programme}

\begin{firewall}[Claim boundary after fourth-series V78]
\label{fire:V78-claim-boundary}
V78 proves a one-link diffusion spectral gap of order \(\alpha\);
this is not the continuum Hamiltonian mass gap.  It does not yet
prove the all-order differentiated-resolvent tree bound, the Wilson
all-thermal local card, the anisotropic scalar optimization, or the
complete all-order selected gate.  It does not prove all-order MTA,
WUFC, CPM, cutoff removal, reflection positivity,
Osterwalder--Schrader reconstruction, construction of the continuum
Yang--Mills measure, or a uniform positive continuum spectral gap.
The full Yang--Mills mass gap has not been proved.
\end{firewall}

\paragraph{Parent-version record.}
V78 was formed from
\texttt{Renewal\_Yang\_Mills\_Mass\_Gap\_Research\_Notes\_V77.tex}.
The V77 parent had \(33\,553\) logical source lines,
\(1\,150\,236\) bytes, and SHA--256
\[
 \texttt{870ecb0bdd16cb287b2780aedfd2a880aa176525aec3fbc2babddb34a7ade962}.
\]
The next compulsory companion audit is V80.

% =============================================================================
% END APPENDED FOURTH-SERIES V78 MATERIAL
% =============================================================================

% =============================================================================
% BEGIN APPENDED FOURTH-SERIES V79 MATERIAL
% =============================================================================

\chapter{V79: Supersymmetric resolvent resummation and tree census}
\label{chap:V79-supersymmetric-resummation}

\section{Why the V78 commutator must be resummed}
\label{sec:V79-purpose}

The differentiated scalar resolvent identity
\eqref{eq:V78-resolvent-derivative} is exact, but iterating it in a
fixed left-invariant frame obscures two cancellations.  First, the
commutator with the scalar generator is already part of the natural
weighted Laplacian on one-forms.  Second, every one-form produced by
differentiating a scalar is exact, and on exact forms a Lie-derivative
insertion returns to a scalar resolvent.  V79 proves both statements
before undertaking another estimate.

This is not a return to the radial transport.  All operators below
are intrinsic on the compact group and use the Wilson measure
directly.

\section{Weighted de Rham complex}
\label{sec:V79-de-Rham}

Let \(d_\alpha^*\) denote the adjoint of the exterior derivative in
\(L^2(\nu_\alpha)\).  Since
\(d\nu_\alpha=Z_\alpha^{-1}e^{\alpha\phi}dU\),
\begin{equation}
 d_\alpha^*
 =
 e^{-\alpha\phi}d^*e^{\alpha\phi}
 =
 d^*-\alpha\iota_{\nabla\phi}.
 \label{eq:V79-weighted-adjoint}
\end{equation}
On \(k\)-forms define the weighted Hodge Laplacian
\begin{equation}
 A_\alpha^{(k)}
 :=
 dd_\alpha^*+d_\alpha^*d.
 \label{eq:V79-Hodge-Laplacian}
\end{equation}
At \(k=0\), the first term vanishes and
\(A_\alpha^{(0)}=A_\alpha\) from
\eqref{eq:V78-generator}.

\begin{lemma}[Exact de Rham intertwining]
\label{lem:V79-intertwining}
On smooth forms,
\begin{equation}
 dA_\alpha^{(k)}
 =
 A_\alpha^{(k+1)}d.
 \label{eq:V79-intertwining}
\end{equation}
\end{lemma}

\begin{proof}
Using \(d^2=0\),
\[
 dA_\alpha^{(k)}
 =
 d d_\alpha^*d
 =
 (dd_\alpha^*+d_\alpha^*d)d
 =
 A_\alpha^{(k+1)}d.
\]
\end{proof}

Let \(\mathcal E_\alpha^1\) be the \(L^2(\nu_\alpha)\)-closure of
the exact smooth one-forms.  The operator \(A_\alpha^{(1)}\)
preserves this space by
\eqref{eq:V79-intertwining}; let
\(R_\alpha^{(1),{\rm ex}}\) denote its inverse there.

\begin{theorem}[Exact-form gap and resolvent intertwining]
\label{thm:V79-exact-gap}
The spectrum of \(A_\alpha^{(1)}\) on
\(\mathcal E_\alpha^1\) is the nonzero spectrum of
\(A_\alpha^{(0)}\), with multiplicity.  Consequently
\begin{equation}
 A_\alpha^{(1)}\big|_{\mathcal E_\alpha^1}
 \ge c\alpha,
 \label{eq:V79-exact-gap}
\end{equation}
and, for every smooth scalar \(q\),
\begin{equation}
 \boxed{\qquad
 dR_\alpha q
 =
 R_\alpha^{(1),{\rm ex}}\,dq.
 \qquad}
 \label{eq:V79-resolvent-intertwining}
\end{equation}
\end{theorem}

\begin{proof}
Because the compact weighted Laplacians have compact resolvent, it is
enough to use an eigenbasis.  If
\(A_\alpha^{(0)}f=\lambda f\) with \(\lambda>0\), then
\eqref{eq:V79-intertwining} gives
\[
 A_\alpha^{(1)}df=\lambda\,df,
\]
and \(df\ne0\).  Conversely, an exact eigenform
\(\omega=df\) with positive eigenvalue can be represented with
\(f\) orthogonal to constants.  Then
\[
 d(A_\alpha^{(0)}f-\lambda f)=0,
\]
so the expression in parentheses is constant.  Its mean is zero
after adding a constant to \(f\), and hence it vanishes.  This proves
the spectral correspondence.  The gap follows from
\Cref{thm:V78-Wilson-gap}.  Finally apply
\eqref{eq:V79-intertwining} to
\(A_\alpha^{(0)}R_\alpha q=Q_\alpha q\) and invert on exact
one-forms.
\end{proof}

\begin{corollary}[One-form version of the covariance identity]
\label{cor:V79-one-form-covariance}
For smooth \(f,g\),
\begin{equation}
 \operatorname{Cov}_{\nu_\alpha}(f,g)
 =
 \left\langle
 df,R_\alpha^{(1),{\rm ex}}dg
 \right\rangle_{\nu_\alpha}.
 \label{eq:V79-one-form-covariance}
\end{equation}
\end{corollary}

\begin{proof}
Use
\eqref{eq:V79-resolvent-intertwining} in
\eqref{eq:V78-covariance}.
\end{proof}

\section{Exponential tilts and exact inverse variation}
\label{sec:V79-tilts}

For fixed smooth real functions \(h_1,\ldots,h_N\), set
\[
 d\nu_{\alpha,\mathbf t}
 =
 {e^{\sum_i t_i h_i}\over
  \mathbb E_{\nu_\alpha}e^{\sum_i t_i h_i}}
 d\nu_\alpha.
\]
Normalization does not enter an adjoint, and therefore
\begin{equation}
 d_{\alpha,\mathbf t}^*
 =
 d_\alpha^*
 -\sum_i t_i\iota_{\nabla h_i}.
 \label{eq:V79-tilted-adjoint}
\end{equation}
Let \(A_{\alpha,\mathbf t}^{(k)}\) be the corresponding weighted
Hodge Laplacian.

\begin{lemma}[Affine Hodge variation]
\label{lem:V79-Hodge-variation}
For each tilt direction \(h_i\),
\begin{equation}
 \partial_{t_i}A_{\alpha,\mathbf t}^{(k)}
 =
 -\left(
 d\iota_{\nabla h_i}+\iota_{\nabla h_i}d
 \right)
 =
 -\mathcal L_{\nabla h_i}.
 \label{eq:V79-Hodge-variation}
\end{equation}
Here \(\mathcal L\) is the Lie derivative on forms.
\end{lemma}

\begin{proof}
Differentiate \eqref{eq:V79-tilted-adjoint} inside
\eqref{eq:V79-Hodge-Laplacian}.  Cartan's formula
\(\mathcal L_X=d\iota_X+\iota_Xd\) gives the last equality.
\end{proof}

For \(\mathbf t\) near zero, let
\(R_{\alpha,\mathbf t}^{(1),{\rm ex}}\) be the inverse of
\(A_{\alpha,\mathbf t}^{(1)}\) on exact one-forms.  Smooth bounded
tilts preserve a positive gap locally in \(\mathbf t\).

\begin{lemma}[Exact resolvent variation]
\label{lem:V79-resolvent-variation}
At every parameter for which the exact-form inverse exists,
\begin{equation}
 \partial_{t_i}R_{\alpha,\mathbf t}^{(1),{\rm ex}}
 =
 R_{\alpha,\mathbf t}^{(1),{\rm ex}}
 \mathcal L_{\nabla h_i}
 R_{\alpha,\mathbf t}^{(1),{\rm ex}}.
 \label{eq:V79-resolvent-variation}
\end{equation}
For distinct directions \(h_1,\ldots,h_k\), at \(\mathbf t=0\),
\begin{equation}
 \partial_{t_1}\cdots\partial_{t_k}R_\alpha^{(1),{\rm ex}}
 =
 \sum_{\pi\in S_k}
 R_\alpha^{(1),{\rm ex}}
 \mathcal L_{\nabla h_{\pi(1)}}
 R_\alpha^{(1),{\rm ex}}\cdots
 \mathcal L_{\nabla h_{\pi(k)}}
 R_\alpha^{(1),{\rm ex}}.
 \label{eq:V79-higher-resolvent-variation}
\end{equation}
\end{lemma}

\begin{proof}
Differentiate
\(A_{\alpha,\mathbf t}^{(1)}
 R_{\alpha,\mathbf t}^{(1),{\rm ex}}=I\) on the invariant exact
subspace and use
\eqref{eq:V79-Hodge-variation}.  This proves the first formula.
The operator \(A_{\alpha,\mathbf t}^{(1)}\) is affine in
\(\mathbf t\), so every later derivative acts on one of the
resolvent factors already present.  Induction gives every ordering
of the \(k\) labelled Lie derivatives exactly once, which is
\eqref{eq:V79-higher-resolvent-variation}.
\end{proof}

\section{Scalarization of every exact-form insertion}
\label{sec:V79-scalarization}

\begin{theorem}[Lie-resolvent insertion returns to scalars]
\label{thm:V79-scalarization}
For smooth scalars \(g,h\),
\begin{equation}
 \boxed{\qquad
 R_\alpha^{(1),{\rm ex}}
 \mathcal L_{\nabla h}
 R_\alpha^{(1),{\rm ex}}\,dg
 =
 dR_\alpha
 \left\langle dh,dR_\alpha g\right\rangle.
 \qquad}
 \label{eq:V79-scalarization}
\end{equation}
The scalar resolvent on the right automatically removes the mean of
the pointwise inner product.
\end{theorem}

\begin{proof}
By \eqref{eq:V79-resolvent-intertwining},
\[
 R_\alpha^{(1),{\rm ex}}dg=dR_\alpha g.
\]
Cartan's formula and \(d^2=0\) imply, for every scalar \(u\),
\[
 \mathcal L_{\nabla h}du
 =
 d\iota_{\nabla h}du
 =
 d\langle dh,du\rangle.
\]
Apply this with \(u=R_\alpha g\), and use
\eqref{eq:V79-resolvent-intertwining} once more.
\end{proof}

\begin{corollary}[The ordered one-form word is a scalar recursion]
\label{cor:V79-word-scalarization}
Every summand in
\eqref{eq:V79-higher-resolvent-variation} applied to an exact input
\(dg\) remains exact after each insertion.  It can therefore be
written using only \(d\), scalar contractions, and the scalar
mean-zero resolvent \(R_\alpha\).
\end{corollary}

\begin{proof}
\Cref{thm:V79-scalarization} performs one insertion and returns an
exact form.  Iterate.
\end{proof}

\section{The exact third connected identity}
\label{sec:V79-third-cumulant}

For real functions, write
\(\kappa_\alpha(f,g,h)\) for their joint third cumulant.

\begin{theorem}[Two-term intrinsic third cumulant]
\label{thm:V79-third-cumulant}
For smooth real \(f,g,h\),
\begin{align}
 \kappa_\alpha(f,g,h)
 &=
 \operatorname{Cov}_{\nu_\alpha}
 \left(
 \langle df,dR_\alpha g\rangle,h
 \right)
 \nonumber\\
 &\quad+
 \operatorname{Cov}_{\nu_\alpha}
 \left(
 f,\langle dh,dR_\alpha g\rangle
 \right)
 \label{eq:V79-third-covariance}\\
 &=
 \left\langle
 d\langle df,dR_\alpha g\rangle,
 dR_\alpha h
 \right\rangle_{\nu_\alpha}
 \nonumber\\
 &\quad+
 \left\langle
 df,
 dR_\alpha\langle dh,dR_\alpha g\rangle
 \right\rangle_{\nu_\alpha}.
 \label{eq:V79-third-tree}
\end{align}
Although the right-hand side is rooted at \(g\), its value is
symmetric in \(f,g,h\).
\end{theorem}

\begin{proof}
Tilt \(\nu_\alpha\) by \(e^{th}\).  Differentiation of a normalized
expectation gives
\[
 {d\over dt}\mathbb E_t F\big|_{t=0}
 =
 \operatorname{Cov}_{\nu_\alpha}(F,h).
\]
Differentiate the tilted version of
\eqref{eq:V79-one-form-covariance}.  Differentiating its inner
product measure gives the first covariance in
\eqref{eq:V79-third-covariance}.  Differentiating the exact-form
resolvent and using
\Cref{lem:V79-resolvent-variation,
thm:V79-scalarization} gives
\[
 \left\langle
 df,
 dR_\alpha\langle dh,dR_\alpha g\rangle
 \right\rangle
 =
 \operatorname{Cov}
 \left(
 f,\langle dh,dR_\alpha g\rangle
 \right),
\]
which is the second covariance.  Applying
\eqref{eq:V78-covariance} to the first one proves
\eqref{eq:V79-third-tree}.  The left-hand side is a joint cumulant,
so the equality itself proves the asserted symmetry.
\end{proof}

\section{A one-factorial census for the formal expansion}
\label{sec:V79-tree-census}

\begin{definition}[Intrinsic differentiation history]
\label{def:V79-history}
Start with the rooted covariance edge
\(\langle df,R_\alpha^{(1),{\rm ex}}dg\rangle\).
Insert each new labelled observable by either:
\begin{enumerate}[leftmargin=3em]
\item differentiating a normalized inner product, which creates one
      covariance edge and is immediately resolved by
      \eqref{eq:V79-one-form-covariance}; or
\item differentiating an exact-form resolvent, which creates one
      ordered Lie-derivative insertion and is immediately scalarized
      by \eqref{eq:V79-scalarization}.
\end{enumerate}
Product-rule differentiation of an existing scalar contraction
records the chosen factor as the attachment place.  The resulting
decorated rooted object is called an intrinsic differentiation
history.
\end{definition}

\begin{lemma}[Plane rooted tree encoding]
\label{lem:V79-plane-tree}
Every intrinsic differentiation history with \(m\) distinctly
labelled observables has a canonical encoding by a plane rooted tree
on those labels with one of finitely many local decorations at each
edge.  Conversely, its encoding and decorations determine every
choice made in the recursion.  Consequently the number \(H_m\) of
formal histories satisfies
\begin{equation}
 H_m\le C^{m-1}(m-1)!
 \label{eq:V79-history-count}
\end{equation}
for a universal finite \(C\).
\end{lemma}

\begin{proof}
Root the initial expression at its distinguished observable.  A
normalized-measure derivative joins its new label by a covariance
edge.  A resolvent derivative inserts the new label at a specified
ordered position on an existing edge.  A product-rule choice joins
it at the selected contraction vertex.  Ordering the children by
the order in which insertions occur makes the result a plane rooted
tree.  The operation type and the selected contraction side require
only finitely many decorations.  Removing the largest insertion-time
label reverses the construction, proving injectivity.

There are \(C_{m-1}\) plane rooted tree shapes with \(m\) vertices,
where \(C_n\) is the \(n\)-th Catalan number.  With the root label
fixed, the other labels can be assigned in \((m-1)!\) ways.  Since
\(C_{m-1}\le4^{m-1}\), and finitely many decorations contribute only
another exponential factor, \eqref{eq:V79-history-count} follows.
\end{proof}

\begin{corollary}[No second factorial from formal resolvent
differentiation]
\label{cor:V79-no-second-factorial}
The permutation sum in
\eqref{eq:V79-higher-resolvent-variation}, together with all
product-rule and normalized-measure choices in the connected
recursion, costs at most one global factorial times an exponential.
It does not by itself create a vertex factorial in addition to a
tree-count factorial.
\end{corollary}

\begin{proof}
Local orderings are precisely the plane orders already counted in
\Cref{lem:V79-plane-tree}; multiplying by another independent
ordering count would count the same histories twice.
\end{proof}

\begin{firewall}[The census is formal, not yet an analytic tree
bound]
\label{fire:V79-census-firewall}
\Cref{lem:V79-plane-tree} controls the number of exact algebraic
terms.  It does not bound an individual contraction containing
\[
 dR_\alpha\langle dh,dR_\alpha g\rangle.
\]
In particular, the \(L^2\) resolvent bounds of V78 do not alone
control products of gradients or their differentiated contractions.
The analytic local card for those products remains to be proved.
\end{firewall}

\section{Fixed Bessel cumulants as a scale test}
\label{sec:V79-Bessel-cumulants}

\begin{proposition}[One global factorial in the radial fundamental
channel]
\label{prop:V79-Bessel-cumulants}
For each fixed \(m\ge2\),
\begin{equation}
 \kappa_m(\phi,\ldots,\phi)
 =
 (-1)^m{3\over2}(m-1)!\,\alpha^{-m}
 +O_m(\alpha^{-m-1}).
 \label{eq:V79-Bessel-cumulants}
\end{equation}
\end{proposition}

\begin{proof}
The radial normalizer is a constant multiple of
\(I_1(\alpha)/\alpha\).  The fixed-order Bessel expansion gives
\[
 \log Z_\alpha
 =
 \alpha-{3\over2}\log\alpha+C
 -{3\over8\alpha}+O(\alpha^{-2}).
\]
The \(m\)-th derivative of \(\log Z_\alpha\) is the \(m\)-th
cumulant of \(\phi\).  For fixed \(m\ge2\), differentiating the full
classical asymptotic expansion termwise gives
\[
 -{3\over2}{d^m\over d\alpha^m}\log\alpha
 =
 (-1)^m{3\over2}(m-1)!\alpha^{-m},
\]
and the remaining terms have the stated order.
\end{proof}

\begin{assessment}[Interpretation of the channel test]
\label{ass:V79-channel-test}
The fundamental radial channel exhibits exactly one factorial and
one factor \(\alpha^{-1}\) per connected input.  This is consistent
with the formal census and with the Wilson clock, but it is a
fixed-channel test only.  It supplies no representation-uniform
bound for the transition and Cartan channels isolated in V40.
\end{assessment}

\section{V79 ledger and V80 programme}
\label{sec:V79-results}

\begin{theorem}[Fourth-series V79 result ledger]
\label{thm:V79-results}
V79 proves:
\begin{enumerate}[label=\textup{(V79.\arabic*)},leftmargin=6.3em]
\item exact weighted de Rham and resolvent intertwining;
\item inheritance of the \(c\alpha\) scalar gap on exact one-forms;
\item the affine tilted Hodge variation and its ordered inverse
      derivatives;
\item exact scalarization of every Lie-resolvent insertion on an
      exact form;
\item a two-term intrinsic formula for the third Wilson cumulant;
\item a plane-rooted-tree encoding with at most
      \(C^m(m-1)!\) formal histories; and
\item the matching one-factorial asymptotic in the fundamental
      radial Bessel channel.
\end{enumerate}
\end{theorem}

\begin{proof}
Use
\Cref{thm:V79-exact-gap,
lem:V79-Hodge-variation,
lem:V79-resolvent-variation,
thm:V79-scalarization,
thm:V79-third-cumulant,
lem:V79-plane-tree,
prop:V79-Bessel-cumulants}.
\end{proof}

\begin{programme}[V80: audited analytic contraction card]
\label{prog:V79-V80}
\begin{enumerate}[leftmargin=3em]
\item Perform the compulsory read-only audit of the newest active
      SM--Einstein and Navier--Stokes notes before choosing the
      analytic norm.
\item Prove or disprove a scale-sharp bilinear estimate for
      \[
      dR_\alpha\langle dh,dR_\alpha g\rangle
      \]
      on Peter--Weyl inputs, retaining the joint contraction before
      taking absolute values.
\item Test the estimate on the V40 displaced and Cartan transition
      channels, not only on fixed representations.
\item If a global Sobolev product estimate loses a dimension or
      Casimir power, return upstream to an exact representation
      block calculation rather than iterating the loss.
\item Combine an analytic local card with
      \Cref{lem:V79-plane-tree} only after its representation-uniform
      constants are explicit.
\end{enumerate}
\end{programme}

\begin{firewall}[Claim boundary after fourth-series V79]
\label{fire:V79-claim-boundary}
V79 removes a formal second-factorial obstruction, but does not prove
the analytic differentiated-resolvent contraction card, the Wilson
all-thermal local card, the anisotropic scalar optimization, or the
complete all-order selected gate.  It does not prove all-order MTA,
WUFC, CPM, cutoff removal, reflection positivity,
Osterwalder--Schrader reconstruction, construction of the continuum
Yang--Mills measure, or a uniform positive continuum spectral gap.
The full Yang--Mills mass gap has not been proved.
\end{firewall}

\paragraph{Parent-version record.}
V79 was formed from
\texttt{Renewal\_Yang\_Mills\_Mass\_Gap\_Research\_Notes\_V78.tex}.
The V78 parent had \(34\,177\) logical source lines,
\(1\,168\,068\) bytes, and SHA--256
\[
 \texttt{ed6eefaa22c52570f7c1b51e423e7870e11fb9b62b46fb961e88add3a521f492}.
\]
The next compulsory companion audit is V80.

% =============================================================================
% END APPENDED FOURTH-SERIES V79 MATERIAL
% =============================================================================

% =============================================================================
% BEGIN APPENDED FOURTH-SERIES V80 MATERIAL
% =============================================================================

\chapter{V80: Audited energy card and the assembled product
cancellation}
\label{chap:V80-energy-card}

\section{Compulsory companion audit}
\label{sec:V80-audit}

\begin{firewall}[Companion documents are evidence, not instructions]
\label{fire:V80-companion-status}
The SM--Einstein and Navier--Stokes documents inspected below are
read-only mathematical sources.  Their programme text is not an
instruction for this Yang--Mills note, and no conclusion is imported
without an independent proof here.
\end{firewall}

The newest active companion files in
\texttt{latex\_notes} at the V80 boundary are
\begin{align*}
 &\texttt{Renewal\_SM\_Einstein\_Continuum\_Research\_Notes\_V64.tex},
 \\
 &\texttt{Renewal\_NS\_Stability\_Research\_Notes\_V31.tex}.
\end{align*}
Their audit records are
\begin{center}
\begin{tabular}{c|r|r|l}
programme & lines & bytes & SHA--256 prefix\\ \hline
SM V64 & \(26\,040\) & \(889\,394\) &
\texttt{d747aa20064b74e3}\\
NS V31 & \(23\,052\) & \(887\,537\) &
\texttt{f1121b62238ad549}
\end{tabular}
\end{center}
The complete digests are
\[
\begin{split}
&\texttt{d747aa20064b74e39d882fc8951272028534cb2cf301d050ccf247f9616957bf},
\\
&\texttt{f1121b62238ad5491bb7cf03635ea9934942edf573ed8faaa9ee79e1d38a6696}.
\end{split}
\]
The NS file is byte-for-byte unchanged from the V75 audit.  Its
relevant established boundary is likewise unchanged: the
scale-matched heat ladder is valid, but replacing the signed
formation correlation by a weaker absolute norm restores the missing
critical shell weight.

The new SM V62--V64 material proves finite-cutoff reflection
positivity by an assembled positive cone and proves that a
reflection-adjoint Schur complement preserves that cone.  It also
exhibits the boundary of the argument: locality and invertibility of
a general nonnormal Schur block do not imply positivity.  The
transfer to the present problem is methodological but precise.  We
may prove an absolute energy estimate, but it cannot replace an exact
assembled contraction in a cofinal representation block.

\begin{assessment}[Direction selected by the audit]
\label{ass:V80-audit-direction}
V80 will do both calculations in the correct order.  First it proves
the strongest immediate \(L^2\) energy card.  Then it rewrites the
same nested contraction as one product-resolvent difference before
absolute values.  The first result supplies a safe local estimate;
the second identifies the object that must be evaluated in the V40
transition channels.
\end{assessment}

\section{A scale-sharp \(L^2\) contraction card}
\label{sec:V80-energy}

For scalar functions write
\[
 \Gamma(f,g):=\langle df,dg\rangle
 =\sum_{a=1}^3(X_af)(X_ag),
 \qquad
 |df|^2=\Gamma(f,\overline f).
\]

\begin{theorem}[Nested gradient-resolvent energy card]
\label{thm:V80-energy-card}
For all sufficiently large \(\alpha\) and smooth scalar \(h,g\),
\begin{equation}
 \boxed{\qquad
 \left\|
 dR_\alpha\Gamma(h,R_\alpha g)
 \right\|_{L^2(\nu_\alpha)}
 \le
 {C\over\alpha}
 \|dh\|_{L^\infty}
 \|Q_\alpha g\|_{L^2(\nu_\alpha)}.
 \qquad}
 \label{eq:V80-energy-card}
\end{equation}
No derivative norm of \(g\) occurs on the right.
\end{theorem}

\begin{proof}
The gradient resolvent estimate
\eqref{eq:V78-resolvent-gradient} gives, first with input
\(\Gamma(h,R_\alpha g)\) and then with input \(g\),
\begin{align*}
 \|dR_\alpha\Gamma(h,R_\alpha g)\|_2
 &\le {C\over\sqrt\alpha}
 \|Q_\alpha\Gamma(h,R_\alpha g)\|_2\\
 &\le {C\over\sqrt\alpha}
 \|\Gamma(h,R_\alpha g)\|_2\\
 &\le {C\over\sqrt\alpha}
 \|dh\|_\infty\|dR_\alpha g\|_2\\
 &\le {C\over\alpha}
 \|dh\|_\infty\|Q_\alpha g\|_2.
\end{align*}
This proves the claim.
\end{proof}

\begin{corollary}[Unitary representation input]
\label{cor:V80-representation-card}
Let \(D^R\) be a unitary irreducible representation and let
\(\Omega=(\Omega_a)_{a=1}^3\) be a compatible
matrix-valued one-form.  Pointwise,
\begin{equation}
 \left\|
 \sum_a(X_aD^R)\Omega_a
 \right\|_{\rm HS}
 \le
 \sqrt{c_R}
 \left(\sum_a\|\Omega_a\|_{\rm HS}^2\right)^{1/2}.
 \label{eq:V80-Casimir-row}
\end{equation}
Consequently the entrywise/operator-valued version of
\eqref{eq:V80-energy-card} obeys
\begin{equation}
 \left\|
 dR_\alpha
 \left\langle dD^R,dR_\alpha g\right\rangle
 \right\|_{L^2({\rm HS})}
 \le
 {C\sqrt{c_R}\over\alpha}
 \|Q_\alpha g\|_{L^2({\rm HS})}.
 \label{eq:V80-representation-energy}
\end{equation}
\end{corollary}

\begin{proof}
Put \(T_a=dR(X_a)\).  Unitarity and the Casimir identity give
\[
 \sum_aT_aT_a^*=c_RI.
\]
Thus the row operator
\((v_1,v_2,v_3)\mapsto\sum_aT_av_a\) has norm
\(\sqrt{c_R}\).  Since
\(X_aD^R(U)=T_aD^R(U)\), applying the row estimate column by column
proves \eqref{eq:V80-Casimir-row}.  Repeat the proof of
\Cref{thm:V80-energy-card} with Hilbert--Schmidt norms.
\end{proof}

\begin{corollary}[Fixed and cofinal scales]
\label{cor:V80-fixed-cofinal}
The representation card contributes
\[
 O(\alpha^{-1})
 \quad\hbox{when }c_R=O(1),
 \qquad
 O(\alpha^{-1/2})
 \quad\hbox{when }c_R\asymp\alpha.
\]
In particular, an absolute application of
\eqref{eq:V80-representation-energy} does not manufacture an
inverse-\(\alpha\) debit at a cofinal
\(d_R\asymp\sqrt\alpha\) vertex.
\end{corollary}

\begin{proof}
Substitute the two Casimir scales into
\(\sqrt{c_R}/\alpha\).
\end{proof}

\section{Exact product-resolvent cancellation}
\label{sec:V80-product-identity}

\begin{lemma}[Weighted diffusion product rule]
\label{lem:V80-product-rule}
For smooth scalar \(h,u\),
\begin{equation}
 A_\alpha(hu)
 =
 hA_\alpha u+uA_\alpha h-2\Gamma(h,u).
 \label{eq:V80-product-rule}
\end{equation}
\end{lemma}

\begin{proof}
The drift in \eqref{eq:V78-generator} is a derivation.  The second
order part obeys
\[
 -\sum_aX_a^2(hu)
 =
 -h\sum_aX_a^2u
 -u\sum_aX_a^2h
 -2\sum_a(X_ah)(X_au).
\]
Adding the drift terms proves the identity.
\end{proof}

\begin{theorem}[Assembled nested-contraction identity]
\label{thm:V80-assembled-identity}
For smooth scalar \(h,g\),
\begin{equation}
 \boxed{\qquad
 dR_\alpha\Gamma(h,R_\alpha g)
 =
 {1\over2}
 dR_\alpha
 \left\{
 hQ_\alpha g+(R_\alpha g)A_\alpha h
 \right\}
 -{1\over2}d\{hR_\alpha g\}.
 \qquad}
 \label{eq:V80-assembled-identity}
\end{equation}
The equality is exact; in particular the three terms on the
right may not be estimated separately when their leading symbols
cancel.
\end{theorem}

\begin{proof}
Set \(u=R_\alpha g\), so
\(A_\alpha u=Q_\alpha g\).  Rearranging
\eqref{eq:V80-product-rule} gives
\[
 \Gamma(h,u)
 =
 {1\over2}
 \left\{
 hQ_\alpha g+uA_\alpha h-A_\alpha(hu)
 \right\}.
\]
Apply \(dR_\alpha\).  Since
\[
 R_\alpha A_\alpha(hu)=Q_\alpha(hu),
 \qquad
 dQ_\alpha(hu)=d(hu),
\]
the result is \eqref{eq:V80-assembled-identity}.
\end{proof}

\begin{corollary}[No Hessian is required]
\label{cor:V80-no-Hessian}
The scalarized insertion of V79 can be evaluated using only products,
one application of \(A_\alpha\) to the declared input \(h\), scalar
mean-zero resolvents, and first derivatives.  No pointwise Hessian of
\(R_\alpha g\) is intrinsic to the expression.
\end{corollary}

\begin{proof}
This is the content of
\eqref{eq:V80-assembled-identity}.
\end{proof}

\section{What the absolute card sees in the V40 transition regimes}
\label{sec:V80-transition-test}

\begin{proposition}[The half-power boundary agrees with the exact
Wilson channels]
\label{prop:V80-transition-boundary}
On the parabolic representation scale
\(d_R\asymp\sqrt\alpha\), one has
\(c_R\asymp\alpha\), and the absolute card
\eqref{eq:V80-representation-energy} supplies only
\(\alpha^{-1/2}\).  This is the same scale as the exact displaced
one-cofinal Wilson quotient in
\eqref{eq:V40-one-cofinal-sqrt}.  It is strictly weaker than the
\(\alpha^{-1}\) undisplaced scale in
\eqref{eq:V40-central-alpha}, and it cannot by itself suppress the
nondecaying two-cofinal Cartan quotient in
\eqref{eq:V40-Cartan-limit}.
\end{proposition}

\begin{proof}
For \(SU(2)\),
\[
 c_R={d_R^2-1\over4}.
\]
Thus \(d_R\asymp\sqrt\alpha\) implies
\(\sqrt{c_R}/\alpha\asymp\alpha^{-1/2}\).
The three comparisons are the exact asymptotics already proved in
\Cref{thm:V40-one-cofinal,thm:V40-Cartan-profile}.
\end{proof}

\begin{firewall}[An \(L^2\) local card is not the cofinal
cancellation]
\label{fire:V80-L2-not-cofinal}
The bound \eqref{eq:V80-energy-card} is rigorous and useful at fixed
or weak representation scale.  In the displaced regime it is only
scale-matched, not improving; in the Cartan regime it omits the
fusion geometry entirely.  Applying it independently at every
vertex would repeat the failure diagnosed by the V80 companion
audit: the missing weight would be reintroduced after the exact
signed or Schur-like cancellation had been discarded.
\end{firewall}

\begin{assessment}[The object for the next exact block calculation]
\label{ass:V80-next-object}
The representation-uniform cofinal question is now localized to the
complete right-hand side of
\eqref{eq:V80-assembled-identity}.  In a Peter--Weyl fusion block,
the terms
\[
 dR_\alpha\{hQ_\alpha g\},
 \qquad
 dR_\alpha\{(R_\alpha g)A_\alpha h\},
 \qquad
 d\{hR_\alpha g\}
\]
must be combined before a norm is taken.  The exact product rule
guarantees their algebraic cancellation; what remains unknown is the
blockwise size after Wilson weighting and projection to the output
representation.
\end{assessment}

\section{V80 ledger and V81 programme}
\label{sec:V80-results}

\begin{theorem}[Fourth-series V80 result ledger]
\label{thm:V80-results}
V80:
\begin{enumerate}[label=\textup{(V80.\arabic*)},leftmargin=6.3em]
\item completes the compulsory read-only SM/NS audit;
\item proves the nested gradient-resolvent \(L^2\) card
      \eqref{eq:V80-energy-card};
\item proves its dimension-free unitary representation form with the
      exact Casimir cost \(\sqrt{c_R}/\alpha\);
\item derives the exact assembled product-resolvent identity
      \eqref{eq:V80-assembled-identity};
\item proves that the absolute card sees exactly the
      \(\alpha^{-1/2}\) cofinal boundary already detected by V40; and
\item localizes the next calculation to one complete Peter--Weyl
      block of the assembled identity.
\end{enumerate}
\end{theorem}

\begin{proof}
Use
\Cref{thm:V80-energy-card,
cor:V80-representation-card,
thm:V80-assembled-identity,
prop:V80-transition-boundary}.
\end{proof}

\begin{programme}[V81: exact assembled Peter--Weyl block]
\label{prog:V80-V81}
\begin{enumerate}[leftmargin=3em]
\item Project
      \eqref{eq:V80-assembled-identity} to one admissible
      \(p\otimes r\to K\) block without estimating its three terms
      separately.
\item Express the scalar coefficient through the exact Wilson
      multipliers \(\eta_p,\eta_r,\eta_K\), the diffusion resolvent
      eigen/block matrix, and the Casimir insertion.
\item Compute first the fixed-partner undisplaced and displaced
      limits of V40, then the two-cofinal Cartan endpoint.
\item Test whether the assembled coefficient is a
      reflection-adjoint positive Schur difference or a signed
      difference with no definite sign.  Do not infer positivity
      from locality or invertibility alone.
\item If the weighted diffusion is not diagonal on Peter--Weyl
      irreducibles, use its conjugation-sector Jacobi matrix rather
      than replacing it by the convolution resolvent of V40.
\item Return to the V79 tree expansion only if the complete block
      retains enough endpoint power uniformly in every transition
      regime.
\end{enumerate}
\end{programme}

\begin{firewall}[Claim boundary after fourth-series V80]
\label{fire:V80-claim-boundary}
V80 proves an analytic \(L^2\) contraction card and an exact
product-resolvent cancellation.  It does not prove the required
cofinal block estimate, the Wilson all-thermal local card, the
anisotropic scalar optimization, or the complete all-order selected
gate.  It does not prove all-order MTA, WUFC, CPM, cutoff removal,
reflection positivity, Osterwalder--Schrader reconstruction,
construction of the continuum Yang--Mills measure, or a uniform
positive continuum spectral gap.  The full Yang--Mills mass gap has
not been proved.
\end{firewall}

\paragraph{Parent-version record.}
V80 was formed from
\texttt{Renewal\_Yang\_Mills\_Mass\_Gap\_Research\_Notes\_V79.tex}.
The V79 parent had \(34\,740\) logical source lines,
\(1\,185\,647\) bytes, and SHA--256
\[
 \texttt{6fa314ff74d06fcc352e066e3f99296feea17b615e88b5e91b19798f5404f758}.
\]
The next compulsory companion audit is V85.

% =============================================================================
% END APPENDED FOURTH-SERIES V80 MATERIAL
% =============================================================================

% =============================================================================
% BEGIN APPENDED FOURTH-SERIES V81 MATERIAL
% =============================================================================

\chapter{V81: Exact radial resolvent and sharpness of the cofinal
half-power}
\label{chap:V81-radial-resolvent}

\section{Purpose}
\label{sec:V81-purpose}

V80 proved the upper bound
\[
 \|dR_\alpha\Gamma(D^R,R_\alpha g)\|_2
 \le C{\sqrt{c_R}\over\alpha}\|Q_\alpha g\|_2.
\]
When \(d_R\asymp\sqrt\alpha\), this is only
\(\alpha^{-1/2}\).  The present note determines whether that loss is
an artifact of taking absolute values.  It uses the exact central
Sturm--Liouville resolvent and the exact character scaling from V40.
The result is a no-go theorem: even the fully assembled nested
contraction is generically nonzero at order \(\alpha^{-1/2}\).

\section{The central Wilson Jacobi operator}
\label{sec:V81-Jacobi}

Put
\[
 C:=-\sum_aX_a^2.
\]
Then \(C\) acts on spin \(j\) by \(c_j=j(j+1)\), and
\[
 A_\alpha=C-\alpha\Gamma(\phi,\cdot).
\]

\begin{lemma}[Casimir expression for the Wilson drift]
\label{lem:V81-drift-Casimir}
If \(Cf=c_jf\), then
\begin{equation}
 \Gamma(\phi,f)
 =
 {1\over2}
 \left\{
 (c_{1/2}+c_j)\phi f-C(\phi f)
 \right\}.
 \label{eq:V81-drift-Casimir}
\end{equation}
\end{lemma}

\begin{proof}
Apply the product rule
\[
 C(\phi f)=\phi Cf+fC\phi-2\Gamma(\phi,f)
\]
and use \(C\phi=c_{1/2}\phi\).
\end{proof}

\begin{theorem}[Exact central Jacobi action]
\label{thm:V81-Jacobi}
With the convention \(\chi_{-1/2}=0\),
\begin{equation}
 \boxed{\qquad
 A_\alpha\chi_j
 =
 c_j\chi_j
 +{\alpha j\over4}\chi_{j+1/2}
 -{\alpha(j+1)\over4}\chi_{j-1/2}.
 \qquad}
 \label{eq:V81-Jacobi}
\end{equation}
The Wilson Gram matrix in this character basis is
\begin{equation}
 G_{jk}^{(\alpha)}
 =
 \langle\chi_j,\chi_k\rangle_{\nu_\alpha}
 =
 \sum_{\ell=|j-k|}^{j+k}
 d_\ell\eta_\ell(\alpha),
 \label{eq:V81-Gram}
\end{equation}
where the sum advances in unit spin steps.  The tridiagonal matrix in
\eqref{eq:V81-Jacobi} is self-adjoint with respect to this Gram
matrix, but not with respect to the Haar character basis.
\end{theorem}

\begin{proof}
The fusion rule gives
\[
 \phi\chi_j
 ={1\over2}
 \left(\chi_{j+1/2}+\chi_{j-1/2}\right).
\]
In \eqref{eq:V81-drift-Casimir}, the two Casimir differences are
\[
 c_j+c_{1/2}-c_{j+1/2}=-j,
 \qquad
 c_j+c_{1/2}-c_{j-1/2}=j+1.
\]
Substitution in
\(A_\alpha=C-\alpha\Gamma(\phi,\cdot)\) proves
\eqref{eq:V81-Jacobi}.  Also
\[
 \chi_j\chi_k
 =
 \sum_{\ell=|j-k|}^{j+k}\chi_\ell.
\]
Taking Wilson expectation and using
\(\mathbb E\chi_\ell=d_\ell\eta_\ell\) proves
\eqref{eq:V81-Gram}.  Self-adjointness follows independently from
\Cref{lem:V78-Dirichlet}, and hence gives the stated weighted matrix
relation.
\end{proof}

\begin{assessment}[Why the V40 convolution resolvent cannot be
substituted]
\label{ass:V81-two-resolvents}
The convolution multiplier of V40 is diagonal on the characters.
The diffusion generator in \eqref{eq:V81-Jacobi} couples
\(j-1/2,j,j+1/2\).  Therefore the two resolvents are not
interchangeable.  Any exact cofinal diffusion calculation must invert
the weighted Jacobi operator or its radial Sturm--Liouville form.
\end{assessment}

\section{Exact radial Green derivative}
\label{sec:V81-Green}

For a central function \(f(\theta)\), the Casimir normalization gives
\begin{equation}
 A_\alpha f
 =
 -{1\over4w_\alpha(\theta)}
 {d\over d\theta}
 \left(
 w_\alpha(\theta)f'(\theta)
 \right).
 \label{eq:V81-radial-generator}
\end{equation}

\begin{theorem}[Exact derivative of the radial resolvent]
\label{thm:V81-Green-derivative}
For a smooth central \(q\),
\begin{equation}
 \boxed{\qquad
 (R_\alpha q)'(\theta)
 =
 -{4\over w_\alpha(\theta)}
 \int_0^\theta
 Q_\alpha q(s)w_\alpha(s)\,ds.
 \qquad}
 \label{eq:V81-Green-derivative}
\end{equation}
The numerator also vanishes at \(\theta=\pi\), so the formula is
regular at both endpoints in the weighted-energy sense.
\end{theorem}

\begin{proof}
Let \(u=R_\alpha q\).  The equation
\(A_\alpha u=Q_\alpha q\), multiplied by
\(-4w_\alpha\), is
\[
 (w_\alpha u')'=-4w_\alpha Q_\alpha q.
\]
Regularity at zero gives
\(\lim_{\theta\downarrow0}w_\alpha(\theta)u'(\theta)=0\).
Integration proves \eqref{eq:V81-Green-derivative}.  At the other
endpoint, the integral over the whole interval is zero by the
definition of \(Q_\alpha\).
\end{proof}

\section{Parabolic resolvent limit}
\label{sec:V81-parabolic}

Let \(\mu_0\) be the Maxwell radial probability measure
\[
 d\mu_0(y)
 =
 \sqrt{2\over\pi}\,y^2e^{-y^2/2}\,dy,
 \qquad y\ge0,
\]
and define its positive radial Ornstein--Uhlenbeck generator
\begin{equation}
 \mathcal A
 :=
 -{1\over4}
 \left(
 {d^2\over dy^2}
 +{2\over y}{d\over dy}
 -y{d\over dy}
 \right).
 \label{eq:V81-OU-generator}
\end{equation}
Let \(\mathcal Q\) denote centering in \(\mu_0\), and let
\(\mathcal R=\mathcal A^{-1}\mathcal Q\).

\begin{lemma}[Scaled Wilson measure and Green kernel]
\label{lem:V81-scaled-Green}
Put
\[
 p_\alpha(y)
 :={1\over\sqrt\alpha}
 w_\alpha(y/\sqrt\alpha),
 \qquad
 0\le y\le\pi\sqrt\alpha.
\]
Then \(p_\alpha(y)dy\) converges to \(d\mu_0(y)\), locally smoothly
away from the endpoint and with a Gaussian dominating tail.  If
\(q_\alpha(y/\sqrt\alpha)\to q_0(y)\) locally uniformly and the
family is uniformly bounded, then
\begin{equation}
 \alpha
 (R_\alpha q_\alpha)(y/\sqrt\alpha)
 \longrightarrow
 \mathcal Rq_0(y)
 \label{eq:V81-resolvent-limit}
\end{equation}
in \(C^1\) on every compact \(y\)-interval, after the mean-zero
constants are fixed.
\end{lemma}

\begin{proof}
Taylor expansion gives
\[
 e^{\alpha(\cos(y/\sqrt\alpha)-1)}
 {\sin^2(y/\sqrt\alpha)\over y^2/\alpha}
 \longrightarrow e^{-y^2/2}.
\]
The localization and Gaussian domination used in
\Cref{lem:V40-transition-Bessel} give convergence of the normalized
density and of centered integrals.

Set
\[
 U_\alpha(y)
 :=
 \alpha(R_\alpha q_\alpha)(y/\sqrt\alpha).
\]
After changing variables in
\eqref{eq:V81-Green-derivative},
\[
 U_\alpha'(y)
 =
 -{4\over p_\alpha(y)}
 \int_0^y
 \left(q_\alpha(z/\sqrt\alpha)
 -\mathbb E_{\mu_\alpha}q_\alpha\right)
 p_\alpha(z)\,dz.
\]
This converges locally to the identical Green-derivative formula for
\(\mathcal Rq_0\).  Integration on a compact interval gives
convergence up to a constant, and the respective mean-zero
conditions determine that constant.  Gaussian tail domination makes
the mean normalization convergent.
\end{proof}

For \(x>0\), define
\[
 H_x(y):={\sin(xy)\over xy}.
\]

\begin{lemma}[Cofinal normalized characters]
\label{lem:V81-character-limit}
If \(j_\alpha\) satisfies
\[
 {d_{j_\alpha}\over\sqrt\alpha}\longrightarrow x\in(0,\infty),
\]
then
\begin{equation}
 {\chi_{j_\alpha}(y/\sqrt\alpha)\over d_{j_\alpha}}
 \longrightarrow H_x(y)
 \label{eq:V81-character-limit}
\end{equation}
in \(C^1\) on compact \(y\)-intervals.  Moreover
\[
 \mathbb E_{\mu_0}H_x=e^{-x^2/2},
\]
consistently with \Cref{cor:V40-Wilson-profile}.
\end{lemma}

\begin{proof}
Use
\[
 {\chi_j(\theta)\over d_j}
 =
 {\sin(d_j\theta)\over d_j\sin\theta}
\]
and substitute \(\theta=y/\sqrt\alpha\).  Differentiation gives the
same local limit.  The expectation identity is the radial Fourier
transform of the three-dimensional standard Gaussian; it is also
the limit of the exact Wilson expectation in
\Cref{cor:V40-Wilson-profile}.
\end{proof}

\section{A nonzero cofinal limiting contraction}
\label{sec:V81-sharpness}

For \(x,z>0\), put
\begin{equation}
 B_{x,z}(y)
 :=
 {1\over4}
 H_x'(y)(\mathcal RH_z)'(y),
 \qquad
 K_{x,z}(y)
 :=
 {1\over2}(\mathcal RB_{x,z})'(y).
 \label{eq:V81-limit-contraction}
\end{equation}
The factors \(1/4\) and \(1/2\) come from the radius-two group
metric.

\begin{theorem}[Parabolic limit of the assembled contraction]
\label{thm:V81-contraction-limit}
Let \(j_\alpha,k_\alpha\) be cofinal sequences with
\[
 {d_{j_\alpha}\over\sqrt\alpha}\to x,
 \qquad
 {d_{k_\alpha}\over\sqrt\alpha}\to z,
\]
and set
\[
 h_\alpha={\chi_{j_\alpha}\over d_{j_\alpha}},
 \qquad
 g_\alpha={\chi_{k_\alpha}\over d_{k_\alpha}}.
\]
For every compact interval \(I\Subset(0,\infty)\), the radial
magnitude satisfies
\begin{equation}
 \sqrt\alpha\,
 \left|
 dR_\alpha\Gamma(h_\alpha,R_\alpha g_\alpha)
 \right|(y/\sqrt\alpha)
 \longrightarrow
 |K_{x,z}(y)|
 \label{eq:V81-contraction-limit}
\end{equation}
uniformly for \(y\in I\).
\end{theorem}

\begin{proof}
By
\Cref{lem:V81-scaled-Green,lem:V81-character-limit},
\[
 \alpha R_\alpha g_\alpha(y/\sqrt\alpha)
 \longrightarrow\mathcal RH_z(y)
\]
in local \(C^1\).  Since a radial differential has squared norm
\(|df|^2=|f'(\theta)|^2/4\),
\[
 \Gamma(h_\alpha,R_\alpha g_\alpha)(y/\sqrt\alpha)
 \longrightarrow
 {1\over4}H_x'(y)(\mathcal RH_z)'(y)
 =B_{x,z}(y).
\]
This family is in fact bounded on the whole Wilson interval.
Indeed, \(|Q_\alpha g_\alpha|\le2\), so
\eqref{eq:V81-Green-derivative} and the two hazard estimates of V76
give
\[
 |(R_\alpha g_\alpha)'(\theta)|
 \le
 \begin{cases}
 C\theta,&0<\theta\le\theta_{\alpha,*},\\
 C\alpha^{-1/2},&\theta_{\alpha,*}\le\theta<\pi.
 \end{cases}
\]
Also
\[
 \left|
 {d\over d\theta}{\chi_{j_\alpha}(\theta)\over d_{j_\alpha}}
 \right|
 \le d_{j_\alpha}\le C\sqrt\alpha.
\]
The median bound
\(\theta_{\alpha,*}\le C/\sqrt\alpha\) now gives
\[
 \sup_{\theta\in(0,\pi)}
 |\Gamma(h_\alpha,R_\alpha g_\alpha)(\theta)|
 \le C.
\]
Apply the scaled Green lemma once more to this globally bounded
family; the Gaussian localization controls its tails.  If
\(v_\alpha=R_\alpha
\Gamma(h_\alpha,R_\alpha g_\alpha)\), then
\[
 \sqrt\alpha\,|dv_\alpha|(y/\sqrt\alpha)
 =
 {1\over2}
 \left|
 {d\over dy}
 \left\{
 \alpha v_\alpha(y/\sqrt\alpha)
 \right\}
 \right|,
\]
which converges to \(|K_{x,z}(y)|\).
\end{proof}

\begin{lemma}[The limiting contraction is nonzero]
\label{lem:V81-limit-nonzero}
For all sufficiently small fixed \(x,z>0\),
\(K_{x,z}\) is not identically zero.  More precisely, on every fixed
compact \(y\)-interval,
\begin{equation}
 K_{x,z}(y)
 =
 {x^2z^2\over9}y
 +O(x^4z^2+x^2z^4),
\label{eq:V81-small-parameters}
\end{equation}
where the remainder is uniform after multiplication by a fixed
polynomial in the compact endpoint.
\end{lemma}

\begin{proof}
The Taylor expansion gives
\[
 H_x(y)=1-{x^2y^2\over6}+O(x^4y^4).
\]
Under \(\mu_0\), \(\mathbb E y^2=3\), and direct calculation from
\eqref{eq:V81-OU-generator} gives
\[
 \mathcal A(y^2-3)={1\over2}(y^2-3),
 \qquad
 \mathcal R(y^2-3)=2(y^2-3).
\]
Therefore
\[
 (\mathcal RH_z)'(y)
 =
 -{2z^2\over3}y+O(z^4),
\]
locally with polynomial weights.  It follows that
\[
 B_{x,z}(y)
 =
 {x^2z^2\over18}y^2
 +O(x^4z^2+x^2z^4).
\]
Centering replaces \(y^2\) by \(y^2-3\), and applying
\(\mathcal R\), differentiating, and multiplying by \(1/2\) yields
\eqref{eq:V81-small-parameters}.  The leading function is nonzero.
\end{proof}

\begin{theorem}[Sharpness of the cofinal half-power]
\label{thm:V81-half-power-sharp}
There exist fixed \(x,z>0\), cofinal normalized-character sequences
\((h_\alpha,g_\alpha)\), and \(c_{x,z}>0\) such that
\begin{equation}
 \liminf_{\alpha\to\infty}
 \sqrt\alpha\,
 \left\|
 dR_\alpha\Gamma(h_\alpha,R_\alpha g_\alpha)
 \right\|_{L^2(\nu_\alpha)}
 \ge c_{x,z}.
 \label{eq:V81-half-power-lower}
\end{equation}
Consequently no representation-uniform improvement
\(o(\alpha^{-1/2})\) of the V80 cofinal energy card is possible,
even for central normalized characters and even though the exact
assembled identity \eqref{eq:V80-assembled-identity} is retained.
\end{theorem}

\begin{proof}
Choose small \(x,z>0\) so that
\Cref{lem:V81-limit-nonzero} makes
\(|K_{x,z}|\) bounded below on some compact interval
\(I\Subset(0,\infty)\).  By
\Cref{thm:V81-contraction-limit} and convergence of the rescaled
Wilson measure to \(\mu_0\),
\[
 \liminf_{\alpha\to\infty}
 \alpha
 \left\|
 dR_\alpha\Gamma(h_\alpha,R_\alpha g_\alpha)
 \right\|_2^2
 \ge
 \int_I|K_{x,z}(y)|^2\,d\mu_0(y)>0.
\]
Take square roots.  The centered \(L^2\) norm of \(g_\alpha\)
converges to the positive variance norm of \(H_z\), so the lower
bound is not caused by a diverging input normalization.
\end{proof}

\begin{corollary}[The missing cofinal debit is not local]
\label{cor:V81-debit-not-local}
The Wilson all-thermal closure cannot obtain an additional
\(\alpha^{-1/2}\) at every cofinal vertex from a stronger local
bound on
\(dR_\alpha\Gamma(h,R_\alpha g)\).  Any such debit must arise from
fusion-sector cancellation across several vertices, from the V64
tree capacity, or from a different global organization of the
connected word.
\end{corollary}

\begin{proof}
An additional local half-power would make the left side of
\eqref{eq:V81-half-power-lower} tend to zero, contradicting the
theorem.
\end{proof}

\section{V81 ledger and V82 programme}
\label{sec:V81-results}

\begin{theorem}[Fourth-series V81 result ledger]
\label{thm:V81-results}
V81 proves:
\begin{enumerate}[label=\textup{(V81.\arabic*)},leftmargin=6.3em]
\item the exact Wilson diffusion Jacobi action on central
      characters and its Bessel Gram matrix;
\item the distinction between the diffusion resolvent and the V40
      convolution resolvent;
\item the exact derivative formula for the radial Wilson Green
      operator;
\item convergence of its parabolic rescaling to the radial
      three-dimensional Ornstein--Uhlenbeck resolvent;
\item an explicit nonzero limiting formula for the nested assembled
      contraction on cofinal normalized characters; and
\item sharpness of the \(\alpha^{-1/2}\) local cofinal card.
\end{enumerate}
\end{theorem}

\begin{proof}
Use
\Cref{thm:V81-Jacobi,
thm:V81-Green-derivative,
lem:V81-scaled-Green,
thm:V81-contraction-limit,
lem:V81-limit-nonzero,
thm:V81-half-power-sharp}.
\end{proof}

\begin{programme}[V82: multi-vertex cofinal compensation]
\label{prog:V81-V82}
\begin{enumerate}[leftmargin=3em]
\item Stop seeking an inverse-\(\alpha\) local card at each cofinal
      vertex; \Cref{thm:V81-half-power-sharp} disproves it.
\item Insert the sharp local half-power into the V64
      Steiner/private-debit ledger and identify exactly which
      all-cofinal tree degree patterns remain unpaid.
\item Retain the complete fusion-sector sum across adjacent
      vertices before taking absolute values, especially when one
      Cartan output becomes an input to the next vertex.
\item In the unresolved degree patterns, test whether the exact
      central Jacobi flow has an orthogonality or telescoping identity
      across two consecutive vertices.
\item If no such identity exists, construct a normalized cofinal
      tree packet to decide whether the proposed intrinsic Wilson
      local card itself is false and return upstream to the selected
      gate formulation.
\end{enumerate}
\end{programme}

\begin{firewall}[Claim boundary after fourth-series V81]
\label{fire:V81-claim-boundary}
V81 proves a sharp local cofinal no-go theorem.  It does not decide
the multi-vertex fusion-sector cancellation, prove the Wilson
all-thermal local card, the anisotropic scalar optimization, or the
complete all-order selected gate.  It does not prove all-order MTA,
WUFC, CPM, cutoff removal, reflection positivity,
Osterwalder--Schrader reconstruction, construction of the continuum
Yang--Mills measure, or a uniform positive continuum spectral gap.
The full Yang--Mills mass gap has not been proved.
\end{firewall}

\paragraph{Parent-version record.}
V81 was formed from
\texttt{Renewal\_Yang\_Mills\_Mass\_Gap\_Research\_Notes\_V80.tex}.
The V80 parent had \(35\,142\) logical source lines,
\(1\,198\,456\) bytes, and SHA--256
\[
 \texttt{0e120cb5e187a074ac245256a86d60a068f265c4c5f21a52a42aad1f2c8f8d75}.
\]
The next compulsory companion audit is V85.

% =============================================================================
% END APPENDED FOURTH-SERIES V81 MATERIAL
% =============================================================================

% =============================================================================
% BEGIN APPENDED FOURTH-SERIES V82 MATERIAL
% =============================================================================

\chapter{V82: Exact odd cancellation closes the ternary
all-thermal card}
\label{chap:V82-ternary-card}

\section{The precise low-arity gap}
\label{sec:V82-purpose}

The V45 two-mark theorem bounds a ternary cumulant by the product of
two selected character gaps.  The V64 all-thermal local card asks for
more: a tree on three inputs has four incidences, and hence its
weight is
\[
 b_i^2b_jb_k,
 \qquad
 b_i:=\sqrt{s_\alpha A_i},
\]
when \(i\) is the degree-two vertex.  V81 shows that one rooted term
in the intrinsic resolvent identity need not contain the missing
half-power.  V82 proves that the complete third cumulant does contain
it.  The source is exact odd cancellation in the global logarithmic
coordinate, not a stronger bound on either rooted term separately.

\section{Global Wilson logarithmic coordinates}
\label{sec:V82-logarithm}

Identify \(\mathfrak{su}(2)\) with \(\mathbb R^3\) in the
radius-two Casimir normalization.  Away from the single Haar-null
point \(-{\bf1}\), write uniquely
\[
 U=\exp Z,
 \qquad
 Z=\theta\omega,
 \qquad
 0\le |Z|=\theta<\pi,\quad\omega\in S^2.
\]
The pushforward of \(\nu_\alpha\) to the open ball \(B_\pi(0)\) has
density proportional to
\begin{equation}
 e^{\alpha\cos|Z|}
 \left({\sin|Z|\over|Z|}\right)^2dZ.
 \label{eq:V82-log-density}
\end{equation}
It is exactly invariant under \(Z\mapsto-Z\).

\begin{lemma}[Uniform logarithmic-coordinate moments]
\label{lem:V82-log-moments}
There is \(C<\infty\) such that, for every sufficiently large
\(\alpha\) and every \(p\ge1\),
\begin{equation}
 \left(
 \mathbb E_{\nu_\alpha}|Z|^p
 \right)^{1/p}
 \le C\sqrt{p\over\alpha}.
 \label{eq:V82-log-moments}
\end{equation}
\end{lemma}

\begin{proof}
On \(0\le\theta\le\pi\),
\[
 1-\cos\theta\ge {2\theta^2\over\pi^2},
 \qquad
 \sin\theta\le\theta.
\]
The radial normalizer is bounded below by
\[
 c e^\alpha\int_0^{\alpha^{-1/2}}\theta^2\,d\theta
 \ge c_1e^\alpha\alpha^{-3/2}.
\]
Consequently
\[
 \mathbb E|Z|^p
 \le
 C\alpha^{3/2}
 \int_0^\infty
 \theta^{p+2}e^{-c\alpha\theta^2}\,d\theta
 \le
 C^p\alpha^{-p/2}
 \Gamma\left({p+3\over2}\right).
\]
The standard Gamma bound
\(\Gamma((p+3)/2)^{1/p}\le C\sqrt p\) proves the claim.
\end{proof}

\section{Odd linear part and quadratic remainder}
\label{sec:V82-linear-remainder}

Let \(R\) be a unitary irreducible representation.  Put
\[
 L_R(Z):=dR(Z),
 \qquad
 N_R(Z):=D^R(e^Z)-I-L_R(Z).
\]

\begin{lemma}[Uniform linear and remainder bounds]
\label{lem:V82-remainder}
Let \(A_R:=1+C_2(R)\) and
\(b_R=\sqrt{s_\alpha A_R}<1\).  For every fixed \(p\ge1\),
\begin{align}
 \|L_R(Z)\|_{L^p(\nu_\alpha;{\rm op})}
 &\le C\sqrt p\,b_R,
 \label{eq:V82-linear-bound}\\
 \|N_R(Z)\|_{L^p(\nu_\alpha;{\rm op})}
 &\le Cp\,b_R^2.
 \label{eq:V82-remainder-bound}
\end{align}
The constants are independent of the representation dimension.
\end{lemma}

\begin{proof}
For a unit Lie-algebra direction \(n\), the Casimir identity implies
\[
 -dR(n)^2\le C_2(R)I,
\]
and hence
\[
 \|L_R(Z)\|_{\rm op}
 \le\sqrt{C_2(R)}\,|Z|.
\]
Combine this with \Cref{lem:V82-log-moments} and
\(s_\alpha\asymp\alpha^{-1}\).

For a skew-adjoint operator \(B\), the integral Taylor remainder is
\[
 e^B-I-B
 =
 \int_0^1(1-t)e^{tB}B^2\,dt,
\]
so
\[
 \|e^B-I-B\|\le {1\over2}\|B\|^2.
\]
Take \(B=L_R(Z)\), use
\Cref{lem:V82-log-moments} at exponent \(2p\), and again use
\(s_\alpha\asymp\alpha^{-1}\).  This proves
\eqref{eq:V82-remainder-bound}.
\end{proof}

\begin{lemma}[Exact centered decomposition]
\label{lem:V82-centered-decomposition}
Let
\[
 \widetilde X_R
 :=
 D^R(U)-\mathbb E_{\nu_\alpha}D^R(U).
\]
Then
\begin{equation}
 \widetilde X_R
 =
 L_R(Z)+\widetilde N_R(Z),
 \qquad
 \widetilde N_R:=N_R-\mathbb E_{\nu_\alpha}N_R,
 \label{eq:V82-centered-decomposition}
\end{equation}
with
\begin{equation}
 \|L_R\|_{L^p}\le C_p b_R,
 \qquad
 \|\widetilde N_R\|_{L^p}\le C_p b_R^2.
 \label{eq:V82-centered-bounds}
\end{equation}
Moreover \(L_R(-Z)=-L_R(Z)\) and
\(\mathbb E_{\nu_\alpha}L_R=0\).
\end{lemma}

\begin{proof}
The logarithmic-coordinate law
\eqref{eq:V82-log-density} is even and \(L_R\) is linear, so its
expectation vanishes.  Centering
\[
 D^R(e^Z)=I+L_R(Z)+N_R(Z)
\]
therefore gives \eqref{eq:V82-centered-decomposition}.
The first bound is
\eqref{eq:V82-linear-bound}; the second follows from
\[
 \|N_R-\mathbb EN_R\|_{L^p}\le2\|N_R\|_{L^p}
\]
and \eqref{eq:V82-remainder-bound}.
\end{proof}

\section{The complete ternary cancellation}
\label{sec:V82-ternary}

For three representations on separate tensor legs put
\[
 X_i(U):=D^{R_i}(U),
 \qquad
 b_i:=\sqrt{s_\alpha A_i}<1,
 \qquad
 A_i=1+C_2(R_i).
\]

\begin{theorem}[Four-incidence ternary estimate]
\label{thm:V82-four-incidence}
There is a dimension- and representation-independent constant \(C\)
such that
\begin{equation}
 \boxed{\qquad
 \left\|
 \kappa_3^{\nu_\alpha}(X_1,X_2,X_3)
 \right\|_{\rm op}
 \le
 C\left(
 b_1^2b_2b_3+
 b_1b_2^2b_3+
 b_1b_2b_3^2
 \right).
 \qquad}
 \label{eq:V82-four-incidence}
\end{equation}
\end{theorem}

\begin{proof}
By \Cref{lem:V45-centered-triple},
\[
 \kappa_3(X_1,X_2,X_3)
 =
 \mathbb E[
 \widetilde X_1\otimes
 \widetilde X_2\otimes
 \widetilde X_3].
\]
Use
\eqref{eq:V82-centered-decomposition} on each tensor leg.  The term
with three linear factors vanishes exactly:
\[
 \mathbb E[
 L_1(Z)\otimes L_2(Z)\otimes L_3(Z)]
 =0,
\]
because its integrand is odd under \(Z\mapsto-Z\).

Every remaining one of the seven terms contains a nonempty set
\(S\subseteq\{1,2,3\}\) of centered remainders.  Holder's inequality
with exponent three and
\eqref{eq:V82-centered-bounds} give
\[
 \left\|
 \mathbb E
 \bigotimes_{i\in S}\widetilde N_i
 \otimes
 \bigotimes_{j\notin S}L_j
 \right\|
 \le
 C
 \left(\prod_{r=1}^3b_r\right)
 \left(\prod_{i\in S}b_i\right).
\]
Choose any \(i\in S\).  Since every \(b_j<1\), discard the other
extra factors to bound this by
\[
 Cb_i^2\prod_{j\ne i}b_j.
\]
Sum the seven terms and then the three possible selected indices.
This proves \eqref{eq:V82-four-incidence}.
\end{proof}

\begin{corollary}[The all-thermal local tree card holds at arity
three]
\label{cor:V82-ternary-tree-card}
Let \(\tau_i\) be the labelled tree on \(\{1,2,3\}\) whose
degree-two vertex is \(i\).  The complete ternary cumulant admits an
exact decomposition
\[
 K_3=\sum_{i=1}^3K_{3,\tau_i}
\]
such that
\begin{equation}
 \|K_{3,\tau_i}\|_{\rm op}
 \le 3!K_3^3
 \prod_{r=1}^3
 (s_\alpha A_r)^{d_r(\tau_i)/2}.
 \label{eq:V82-ternary-tree-card}
\end{equation}
Thus \eqref{eq:V64-thermal-card} is proved for \(m=3\).
\end{corollary}

\begin{proof}
Set
\[
 W_i=b_i^2\prod_{j\ne i}b_j,
 \qquad
 W=W_1+W_2+W_3.
\]
Every \(b_i>0\), so \(W>0\).  Define
\[
 K_{3,\tau_i}:={W_i\over W}K_3.
\]
This is an exact decomposition.  By
\Cref{thm:V82-four-incidence},
\[
 \|K_{3,\tau_i}\|
 \le {W_i\over W}CW=CW_i.
\]
Enlarge the fixed base \(K_3\) to write the constant in the declared
tree-card form.
\end{proof}

\begin{remark}[No chart complement]
\label{rem:V82-no-complement}
The principal logarithm omits only the single group point
\(-{\bf1}\), which has zero Wilson measure.  All moment estimates
were integrated on the complete ball \(|Z|<\pi\).  Thus the proof has
no mesoscopic chart cutoff and no unassigned complement of the type
that obstructed V70.
\end{remark}

\section{Reconciliation with the V81 sharp half-power}
\label{sec:V82-reconciliation}

\begin{proposition}[Rooted half-powers cancel only after assembly]
\label{prop:V82-rooted-assembly}
The lower bound in \Cref{thm:V81-half-power-sharp} is compatible with
\eqref{eq:V82-four-incidence}.  For fixed-scale thermal inputs,
each of the two rooted terms in
\eqref{eq:V79-third-covariance} can have the naive three-incidence
size \(b_1b_2b_3\), while their sum has the four-incidence size in
\eqref{eq:V82-four-incidence}.  At the logarithmic-coordinate level,
the canceled leading contribution is precisely the odd tensor
\(\mathbb E[L_1\otimes L_2\otimes L_3]\).
\end{proposition}

\begin{proof}
The two rooted terms sum exactly to the third cumulant by
\Cref{thm:V79-third-cumulant}.  Independently, the proof of
\Cref{thm:V82-four-incidence} identifies the only possible
three-incidence term and proves that it is zero.  No assertion that
either rooted term vanishes is needed.  When all \(b_i\asymp1\), the
four-incidence target is itself of order one, so the nonzero cofinal
half-power found in V81 violates no tree weight.
\end{proof}

\begin{assessment}[Updated all-thermal boundary]
\label{ass:V82-boundary}
The all-thermal local card is now proved at arities two and three:
arity two follows from the Wilson covariance identity and the gap,
while arity three is
\Cref{cor:V82-ternary-tree-card}.  At arity four, parity alone is
insufficient because the four-linear tensor is even.  The connected
fourth cumulant must cancel its three Gaussian pairings before
norms are taken.  This is the next genuinely new analytic step.
\end{assessment}

\section{V82 ledger and V83 programme}
\label{sec:V82-results}

\begin{theorem}[Fourth-series V82 result ledger]
\label{thm:V82-results}
V82 proves:
\begin{enumerate}[label=\textup{(V82.\arabic*)},leftmargin=6.3em]
\item a global, complement-free logarithmic-coordinate moment bound
      for the exact Wilson law;
\item a dimension-free decomposition of every centered
      representation into an odd linear part of size \(b_R\) and a
      centered remainder of size \(b_R^2\);
\item exact cancellation of the triple-linear tensor;
\item the full four-incidence ternary estimate
      \eqref{eq:V82-four-incidence}; and
\item the V64 all-thermal local tree card at arity three.
\end{enumerate}
\end{theorem}

\begin{proof}
Use
\Cref{lem:V82-log-moments,
lem:V82-remainder,
lem:V82-centered-decomposition,
thm:V82-four-incidence,
cor:V82-ternary-tree-card}.
\end{proof}

\begin{programme}[V83: fourth connected Gaussian-pairing defect]
\label{prog:V82-V83}
\begin{enumerate}[leftmargin=3em]
\item Expand the complete fourth cumulant in the global logarithmic
      coordinate before taking operator norms.
\item Separate the four-linear tensor from all terms containing a
      quadratic representation remainder.
\item Subtract the three covariance pairings inside the four-linear
      tensor and estimate the resulting radial fourth-cumulant defect
      relative to the Gaussian covariance at scale \(s_\alpha\).
\item Prove a six-incidence bound indexed by the sixteen labelled
      trees on four inputs, or identify the exact representation
      pattern that prevents it.
\item Use the global ball throughout; do not reintroduce a
      mesoscopic chart complement.
\end{enumerate}
\end{programme}

\begin{firewall}[Claim boundary after fourth-series V82]
\label{fire:V82-claim-boundary}
V82 proves the all-thermal local card only through arity three.  It
does not prove the fourth- or all-order card, the anisotropic scalar
optimization, or the complete all-order selected gate.  It does not
prove all-order MTA, WUFC, CPM, cutoff removal, reflection
positivity, Osterwalder--Schrader reconstruction, construction of
the continuum Yang--Mills measure, or a uniform positive continuum
spectral gap.  The full Yang--Mills mass gap has not been proved.
\end{firewall}

\paragraph{Parent-version record.}
V82 was formed from
\texttt{Renewal\_Yang\_Mills\_Mass\_Gap\_Research\_Notes\_V81.tex}.
The V81 parent had \(35\,720\) logical source lines,
\(1\,214\,307\) bytes, and SHA--256
\[
 \texttt{8c7e5a363a79133310e42c493617753a5d93c1ed5299956f5f00e07aa3b6966a}.
\]
The next compulsory companion audit is V85.

% =============================================================================
% END APPENDED FOURTH-SERIES V82 MATERIAL
% =============================================================================

% =============================================================================
% BEGIN APPENDED FOURTH-SERIES V83 MATERIAL
% =============================================================================

\chapter{V83: The fourth radial Wick defect closes the quartic
all-thermal card}
\label{chap:V83-quartic-card}

\section{Quadratic and cubic representation pieces}
\label{sec:V83-representation-expansion}

For a unitary irreducible representation \(R\), retain
\(L_R(Z)=dR(Z)\) and define
\[
 Q_R(Z):={1\over2}L_R(Z)^2,
 \qquad
 M_R(Z):=
 D^R(e^Z)-I-L_R(Z)-Q_R(Z).
\]
The quadratic term is even under \(Z\mapsto-Z\).

\begin{lemma}[Three-level centered expansion]
\label{lem:V83-three-level}
In the all-thermal range
\(b_R=\sqrt{s_\alpha(1+C_2(R))}<1\),
\begin{equation}
 D^R(U)-\mathbb E_{\nu_\alpha}D^R(U)
 =
 L_R+\widetilde Q_R+\widetilde M_R,
 \label{eq:V83-three-level}
\end{equation}
where
\[
 \widetilde Q_R=Q_R-\mathbb EQ_R,
 \qquad
 \widetilde M_R=M_R-\mathbb EM_R.
\]
For every fixed \(p\ge1\),
\begin{align}
 \|L_R\|_{L^p({\rm op})}
 &\le C_p b_R,
 \label{eq:V83-L-bound}\\
 \|\widetilde Q_R\|_{L^p({\rm op})}
 &\le C_p b_R^2,
 \label{eq:V83-Q-bound}\\
 \|\widetilde M_R\|_{L^p({\rm op})}
 &\le C_p b_R^3.
 \label{eq:V83-M-bound}
\end{align}
All constants are independent of the representation dimension.
\end{lemma}

\begin{proof}
The linear estimate is
\eqref{eq:V82-linear-bound}.  Since
\(\|L_R(Z)\|\le\sqrt{C_2(R)}|Z|\),
\Cref{lem:V82-log-moments} gives
\[
 \|Q_R\|_{L^p}\le C_p{C_2(R)\over\alpha}
 \le C_pb_R^2.
\]
Centering costs at most a factor two.

For skew-adjoint \(B\), Taylor's integral remainder gives
\[
 e^B-I-B-{1\over2}B^2
 =
 {1\over2}\int_0^1(1-t)^2e^{tB}B^3\,dt.
\]
Thus
\[
 \|M_R(Z)\|\le {1\over6}\|L_R(Z)\|^3.
\]
Use \Cref{lem:V82-log-moments} at exponent \(3p\), then center.
Finally, \(\mathbb EL_R=0\) by exact inversion symmetry, so centering
the representation gives \eqref{eq:V83-three-level}.
\end{proof}

\section{The logarithmic-coordinate fourth Wick defect}
\label{sec:V83-Wick-defect}

Write \(Z=(Z_1,Z_2,Z_3)\) in an orthonormal Lie-algebra basis and put
\[
 v_\alpha:={1\over3}\mathbb E|Z|^2,
 \qquad
 u_\alpha:={1\over15}\mathbb E|Z|^4.
\]

\begin{lemma}[One-order gain in radial fourth moments]
\label{lem:V83-radial-defect}
For all sufficiently large \(\alpha\),
\begin{align}
 v_\alpha
 &={1\over\alpha}+O(\alpha^{-2}),
 \label{eq:V83-second-moment}\\
 u_\alpha
 &={1\over\alpha^2}+O(\alpha^{-3}),
 \label{eq:V83-fourth-moment}\\
 |u_\alpha-v_\alpha^2|
 &\le C\alpha^{-3}.
 \label{eq:V83-Wick-defect}
\end{align}
\end{lemma}

\begin{proof}
Put \(y=\sqrt\alpha Z\) in
\eqref{eq:V82-log-density}.  Relative to Lebesgue measure in
\(\mathbb R^3\), the unnormalized density after removal of its
constant factor is
\[
 e^{\alpha(\cos(|y|/\sqrt\alpha)-1)}
 \left(
 {\sin(|y|/\sqrt\alpha)\over |y|/\sqrt\alpha}
 \right)^2.
\]
On every fixed ball this equals
\[
 e^{-|y|^2/2}
 \left[
 1+{1\over\alpha}
 \left({|y|^4\over24}-{|y|^2\over3}\right)
 +O\left({1+|y|^8\over\alpha^2}\right)
 \right].
\]
On \(|y|\le\alpha^{1/8}\), Taylor's theorem gives the displayed
remainder with a Gaussian majorant.  On the complement,
\(1-\cos\theta\ge2\theta^2/\pi^2\) gives an exponentially small
Gaussian tail, uniformly after multiplication by any fixed
polynomial.  Normalizing therefore yields
\[
 \mathbb E|y|^2=3+O(\alpha^{-1}),
 \qquad
 \mathbb E|y|^4=15+O(\alpha^{-1}).
\]
Rescaling proves
\eqref{eq:V83-second-moment} and
\eqref{eq:V83-fourth-moment}; subtracting their squares proves
\eqref{eq:V83-Wick-defect}.
\end{proof}

\begin{lemma}[Operator-valued four-linear cumulant]
\label{lem:V83-linear-cumulant}
For four unitary representations on separate tensor legs,
\begin{equation}
 \left\|
 \kappa_4^{\nu_\alpha}(L_1,L_2,L_3,L_4)
 \right\|_{\rm op}
 \le
 Cs_\alpha\prod_{i=1}^4b_i.
 \label{eq:V83-linear-cumulant}
\end{equation}
\end{lemma}

\begin{proof}
Rotational invariance gives
\begin{align*}
 \mathbb E[Z_aZ_b]
 &=v_\alpha\delta_{ab},\\
 \mathbb E[Z_aZ_bZ_cZ_d]
 &=u_\alpha
 \left(
 \delta_{ab}\delta_{cd}
 +\delta_{ac}\delta_{bd}
 +\delta_{ad}\delta_{bc}
 \right).
\end{align*}
Subtracting the three covariance pairings shows that the fourth
cumulant tensor is the same delta sum multiplied by
\(u_\alpha-v_\alpha^2\).

Write \(L_i(Z)=\sum_aZ_aT_{i,a}\).  The fixed three-dimensional
index sums and
\(\|T_{i,a}\|\le\sqrt{C_2(R_i)}\) give
\[
 \|\kappa_4(L_1,L_2,L_3,L_4)\|
 \le
 C|u_\alpha-v_\alpha^2|
 \prod_i\sqrt{C_2(R_i)}.
\]
Use \eqref{eq:V83-Wick-defect},
\(s_\alpha\asymp\alpha^{-1}\), and
\[
 {\sqrt{C_2(R_i)}\over\sqrt\alpha}\le Cb_i
\]
to obtain \eqref{eq:V83-linear-cumulant}.
\end{proof}

\section{Complete quartic power ledger}
\label{sec:V83-quartic-ledger}

For centered variables on separate tensor legs,
\begin{equation}
 \kappa_4(Y_1,Y_2,Y_3,Y_4)
 =
 \mathbb E[Y_1Y_2Y_3Y_4]
 -\sum_{\rm three\ pairings}
 \mathbb E[Y_iY_j]\,\mathbb E[Y_kY_\ell],
 \label{eq:V83-centered-kappa4}
\end{equation}
with the canonical tensor rebracketing understood.

\begin{lemma}[Fixed-order cumulant product bound]
\label{lem:V83-cumulant-product}
For centered operator-valued \(Y_1,\ldots,Y_4\),
\begin{equation}
 \|\kappa_4(Y_1,Y_2,Y_3,Y_4)\|
 \le4\prod_{i=1}^4\|Y_i\|_{L^4({\rm op})}.
 \label{eq:V83-cumulant-product}
\end{equation}
\end{lemma}

\begin{proof}
Holder bounds the fourfold moment by the displayed product.
Operator Cauchy--Schwarz bounds each covariance product by the same
quantity.  There are three pairings.
\end{proof}

\begin{lemma}[The one-quadratic terms vanish]
\label{lem:V83-one-Q-vanishes}
For distinct indices \(i,j,k,\ell\),
\begin{equation}
 \kappa_4(\widetilde Q_i,L_j,L_k,L_\ell)=0.
 \label{eq:V83-one-Q-vanishes}
\end{equation}
\end{lemma}

\begin{proof}
The fourfold moment has an even factor and three odd factors, so its
integrand is odd under \(Z\mapsto-Z\).  Every covariance between
\(\widetilde Q_i\) and an \(L\)-factor also vanishes by the same
symmetry.  Hence both the moment and all three pairing products in
\eqref{eq:V83-centered-kappa4} vanish.
\end{proof}

\begin{theorem}[Six-incidence quartic estimate]
\label{thm:V83-six-incidence}
For all-thermal inputs \(b_i<1\),
\begin{align}
 \left\|
 \kappa_4^{\nu_\alpha}(X_1,X_2,X_3,X_4)
 \right\|
 &\le
 C\left[
 \sum_{i=1}^4
 b_i^3\prod_{r\ne i}b_r
 \right.
 \nonumber\\
 &\hspace{3em}\left.
 +
 \sum_{1\le i<j\le4}
 b_i^2b_j^2
 \prod_{r\notin\{i,j\}}b_r
 \right].
 \label{eq:V83-six-incidence}
\end{align}
\end{theorem}

\begin{proof}
Expand each centered \(X_i\) by
\eqref{eq:V83-three-level} and use multilinearity of the fourth
cumulant.

The four-linear term is bounded by
\Cref{lem:V83-linear-cumulant}.  Since
\(b_ib_j\ge s_\alpha\) for every pair, it is bounded by any one of
the pair terms
\[
 \left(\prod_rb_r\right)b_ib_j.
\]

Every term containing exactly one \(\widetilde Q\) and three
linear factors vanishes by
\Cref{lem:V83-one-Q-vanishes}.  Consider the remaining terms.
If a term contains a factor \(\widetilde M_i\), then
\Cref{lem:V83-three-level,lem:V83-cumulant-product} bound it by
\[
 C\left(\prod_rb_r\right)b_i^2
\]
after discarding all other extra factors, which are at most one.
This is the star weight centered at \(i\).

If there is no \(\widetilde M\), every non-linear term not already
discarded contains at least two distinct quadratic factors, say at
\(i,j\).  Its bound is
\[
 C\left(\prod_rb_r\right)b_ib_j
\]
after discarding further factors.  This is the path weight whose
degree-two vertices are \(i,j\).  Summing the finitely many expansion
terms proves \eqref{eq:V83-six-incidence}.
\end{proof}

\begin{corollary}[The all-thermal local tree card holds at arity
four]
\label{cor:V83-quartic-tree-card}
The complete fourth Wilson cumulant admits an exact decomposition
\[
 K_4=\sum_{\tau\in\mathfrak T_4}K_{4,\tau}
\]
such that
\begin{equation}
 \|K_{4,\tau}\|_{\rm op}
 \le4!K_4^4
 \prod_{i=1}^4
 (s_\alpha A_i)^{d_i(\tau)/2}.
 \label{eq:V83-quartic-tree-card}
\end{equation}
Thus \eqref{eq:V64-thermal-card} is proved for \(m=4\).
\end{corollary}

\begin{proof}
Every tree on four labelled vertices is either a star, with weight
\[
 b_i^3\prod_{r\ne i}b_r,
\]
or a path, with weight
\[
 b_i^2b_j^2\prod_{r\notin\{i,j\}}b_r
\]
when \(i,j\) are its two internal vertices.  There are four stars
and, for each unordered internal pair, two paths with the same
weight.  Hence the right side of
\eqref{eq:V83-six-incidence} is bounded by a constant times
\[
 \mathcal W_4
 :=
 \sum_{\tau\in\mathfrak T_4}
 \prod_i b_i^{d_i(\tau)}.
\]
Define
\[
 K_{4,\tau}
 :=
 {\prod_i b_i^{d_i(\tau)}\over\mathcal W_4}K_4.
\]
This is exact and the preceding bound gives the declared estimate
after enlarging a fixed base.
\end{proof}

\begin{assessment}[Meaning of the Wick-defect gain]
\label{ass:V83-Wick-meaning}
The two extra incidences needed beyond the four centered linear
factors arise in exactly three ways:
\begin{enumerate}[leftmargin=3em]
\item one cubic representation remainder supplies two incidences at
      a star center;
\item two quadratic remainders supply one incidence at each internal
      path vertex; or
\item the connected fourth tensor of the logarithmic-coordinate law
      supplies \(s_\alpha\), which can be transferred to any pair of
      vertices because \(b_ib_j\ge s_\alpha\).
\end{enumerate}
Estimating the fourth moment before subtracting its three pairings
would lose the third mechanism completely.
\end{assessment}

\section{V83 ledger and V84 programme}
\label{sec:V83-results}

\begin{theorem}[Fourth-series V83 result ledger]
\label{thm:V83-results}
V83 proves:
\begin{enumerate}[label=\textup{(V83.\arabic*)},leftmargin=6.3em]
\item a dimension-free linear--quadratic--cubic decomposition of
      centered Wilson representation matrices;
\item the radial fourth Wick defect
      \(|u_\alpha-v_\alpha^2|\le C\alpha^{-3}\);
\item the resulting one-order gain for four linear representation
      factors;
\item exact disappearance of every one-quadratic/three-linear
      fourth cumulant;
\item the complete six-incidence quartic estimate; and
\item the V64 all-thermal local tree card at arity four.
\end{enumerate}
\end{theorem}

\begin{proof}
Use
\Cref{lem:V83-three-level,
lem:V83-radial-defect,
lem:V83-linear-cumulant,
lem:V83-one-Q-vanishes,
thm:V83-six-incidence,
cor:V83-quartic-tree-card}.
\end{proof}

\begin{programme}[V84: arbitrary-order Gaussian-defect expansion]
\label{prog:V83-V84}
\begin{enumerate}[leftmargin=3em]
\item Replace the order-four Taylor split by the complete exponential
      series in \(L_R(Z)\), retaining parity before norms.
\item Prove a uniform cumulant comparison between the scaled radial
      Wilson logarithm and a standard three-dimensional Gaussian,
      with one factor \(s_\alpha^{(k-2)/2}\) at every connected
      correction vertex of degree \(k\).
\item Encode Gaussian pair contractions and Wilson correction
      cumulants by the Steiner trees of V69.
\item Sum representation Taylor degrees and correction vertices with
      one exponential generating function, avoiding an independent
      factorial at each vertex.
\item If the uniform-in-order radial cumulant comparison fails,
      identify its first failing order and retain the proven
      arity-three and arity-four cards.
\end{enumerate}
\end{programme}

\begin{firewall}[Claim boundary after fourth-series V83]
\label{fire:V83-claim-boundary}
V83 proves the all-thermal local card only through arity four.  It
does not prove the arbitrary-order Gaussian-defect summation, the
anisotropic scalar optimization, or the complete all-order selected
gate.  It does not prove all-order MTA, WUFC, CPM, cutoff removal,
reflection positivity, Osterwalder--Schrader reconstruction,
construction of the continuum Yang--Mills measure, or a uniform
positive continuum spectral gap.  The full Yang--Mills mass gap has
not been proved.
\end{firewall}

\paragraph{Parent-version record.}
V83 was formed from
\texttt{Renewal\_Yang\_Mills\_Mass\_Gap\_Research\_Notes\_V82.tex}.
The V82 parent had \(36\,140\) logical source lines,
\(1\,226\,211\) bytes, and SHA--256
\[
 \texttt{2fa57efaf5baf1beb73c4244601cf5e7da63cebab728542fc33469826ef787e6}.
\]
The next compulsory companion audit is V85.

% =============================================================================
% END APPENDED FOURTH-SERIES V83 MATERIAL
% =============================================================================

% =============================================================================
% BEGIN APPENDED FOURTH-SERIES V84 MATERIAL
% =============================================================================

\chapter{V84: First Wilson--Gaussian defect closes the quintic
all-thermal card}
\label{chap:V84-quintic-card}

\section{A finite-degree Wilson--Gaussian comparison}
\label{sec:V84-comparison}

Let \(Y_\alpha=\sqrt\alpha Z\), where \(Z\) is the global
logarithmic coordinate of V82, and let \(\gamma_3\) be standard
Gaussian measure on \(\mathbb R^3\).

\begin{lemma}[Polynomial comparison through degree six]
\label{lem:V84-polynomial-comparison}
There is \(C<\infty\) such that, for every tensor-valued polynomial
\[
 P(y)=\sum_{|\beta|\le6}P_\beta y^\beta,
\]
\begin{equation}
 \left\|
 \mathbb E_{\nu_\alpha}P(Y_\alpha)
 -\mathbb E_{\gamma_3}P(Y)
 \right\|
 \le
 {C\over\alpha}
 \sum_{|\beta|\le6}\|P_\beta\|.
 \label{eq:V84-polynomial-comparison}
\end{equation}
The same assertion holds for every moment occurring in a joint
cumulant whose total polynomial degree is at most six.
\end{lemma}

\begin{proof}
The density expansion in the proof of
\Cref{lem:V83-radial-defect} is
\[
 {d\mathcal L(Y_\alpha)\over d\gamma_3}(y)
 =
 1+{1\over\alpha}
 \left(
 {|y|^4\over24}-{|y|^2\over3}-c_0
 \right)
 +O\left(
 {1+|y|^8\over\alpha^2}
 \right)
\]
on \(|y|\le\alpha^{1/8}\), where \(c_0\) is the Gaussian mean of
the polynomial in parentheses.  Both the exact density and the
remainder have a Gaussian dominating tail outside this ball, as
proved there.  Multiplication by a monomial of degree at most six
and integration therefore gives an \(O(\alpha^{-1})\) error.
Summing the finitely many coefficients proves
\eqref{eq:V84-polynomial-comparison}.  A block moment in the declared
cumulant is another polynomial of total degree at most six, so the
same estimate applies.
\end{proof}

\begin{lemma}[Fixed-degree cumulant comparison]
\label{lem:V84-cumulant-comparison}
Let \(P_1,\ldots,P_r\) be tensor-valued polynomials with
\(\sum_i\deg P_i\le6\).  Then
\begin{equation}
 \left\|
 \kappa_r^{\nu_\alpha}
 (P_1(Y_\alpha),\ldots,P_r(Y_\alpha))
 -
 \kappa_r^{\gamma_3}
 (P_1(Y),\ldots,P_r(Y))
 \right\|
 \le {C_r\over\alpha}
 \prod_i\|P_i\|_{\rm coef},
 \label{eq:V84-cumulant-comparison}
\end{equation}
where \(\|P_i\|_{\rm coef}\) is the sum of its coefficient norms.
\end{lemma}

\begin{proof}
Use the moment--cumulant formula
\[
 \kappa_r(P_1,\ldots,P_r)
 =
 \sum_{\pi\in\Pi_r}
 (|\pi|-1)!(-1)^{|\pi|-1}
 \prod_{B\in\pi}
 \mathbb E\prod_{i\in B}P_i.
\]
For each partition, subtract the Gaussian product of block moments
by a telescoping product identity.  Every changed block is
\(O(\alpha^{-1})\) by
\Cref{lem:V84-polynomial-comparison}; every unchanged fixed-degree
Gaussian or Wilson block moment is bounded by a constant times its
coefficient norm.  Since \(r\le6\), the finite partition sum is
absorbed in \(C_r\).
\end{proof}

\section{The unique low-excess quintic survivor}
\label{sec:V84-Gaussian-zero}

\begin{lemma}[Gaussian quadratic--linear connected zero]
\label{lem:V84-Gaussian-zero}
Let \(G\) be a centered Gaussian vector, let \(q(G)\) be a
quadratic polynomial, and let
\(\ell_2(G),\ldots,\ell_5(G)\) be linear.  Then
\begin{equation}
 \kappa_5^{\gamma_3}
 (q,\ell_2,\ell_3,\ell_4,\ell_5)=0.
 \label{eq:V84-Gaussian-zero}
\end{equation}
\end{lemma}

\begin{proof}
Constants in \(q\) do not affect a cumulant of order at least two,
so center its quadratic part.  Wick's formula represents the joint
cumulant by pairings whose contraction graph is connected on the
five declared polynomial vertices.  The quadratic vertex has two
Gaussian half-edges and each linear vertex has one, for six
half-edges in total.  A connected graph on five vertices needs at
least four edges, hence at least eight half-edges.  No connected
pairing exists, and the cumulant is zero.
\end{proof}

\begin{proposition}[One quadratic and four linear Wilson factors]
\label{prop:V84-one-Q-defect}
Let \(L_i=dR_i(Z)\) and \(Q_i=L_i^2/2\).  For all-thermal
\(b_i=\sqrt{s_\alpha A_i}<1\),
\begin{equation}
 \left\|
 \kappa_5^{\nu_\alpha}
 (Q_i,L_j,L_k,L_\ell,L_m)
 \right\|
 \le
 Cs_\alpha b_i^2b_jb_kb_\ell b_m.
 \label{eq:V84-one-Q-defect}
\end{equation}
\end{proposition}

\begin{proof}
Write \(Z=Y_\alpha/\sqrt\alpha\).  The coefficient norm of
\(L_r(Y_\alpha/\sqrt\alpha)\) is at most \(Cb_r\), and that of
\(Q_i(Y_\alpha/\sqrt\alpha)\) is at most \(Cb_i^2\).
Cumulants are unchanged by centering \(Q_i\).  Apply
\Cref{lem:V84-cumulant-comparison} at total degree six.  Its Gaussian
term vanishes by \Cref{lem:V84-Gaussian-zero}; its error is
\[
 {C\over\alpha}b_i^2b_jb_kb_\ell b_m.
\]
Use \(s_\alpha\asymp\alpha^{-1}\).
\end{proof}

\section{Fourth-order representation split}
\label{sec:V84-four-level}

Define
\[
 P_{R,3}:={1\over6}L_R^3,
\]
and
\[
 N_{R,4}
 :=
 D^R(e^Z)-I-L_R-{1\over2}L_R^2-{1\over6}L_R^3.
\]

\begin{lemma}[Centered expansion through cubic order]
\label{lem:V84-four-level}
For every all-thermal representation,
\begin{equation}
 D^R(U)-\mathbb ED^R(U)
 =
 L_R+\widetilde Q_R+P_{R,3}+\widetilde N_{R,4},
 \label{eq:V84-four-level}
\end{equation}
where tildes denote centering, and for fixed \(p\),
\begin{align}
 \|L_R\|_{L^p}&\le C_pb_R,\nonumber\\
 \|\widetilde Q_R\|_{L^p}&\le C_pb_R^2,\nonumber\\
 \|P_{R,3}\|_{L^p}&\le C_pb_R^3,\nonumber\\
 \|\widetilde N_{R,4}\|_{L^p}&\le C_pb_R^4.
 \label{eq:V84-four-level-bounds}
\end{align}
The linear and cubic pieces are odd; in particular their Wilson
means vanish.
\end{lemma}

\begin{proof}
The first two estimates are
\Cref{lem:V83-three-level}.  The cubic estimate follows from
\Cref{lem:V82-log-moments}.  Taylor's integral remainder and
unitarity give
\[
 \|N_{R,4}(Z)\|
 \le {1\over24}\|L_R(Z)\|^4.
\]
The moment lemma gives the fourth estimate, and centering costs at
most two.  Parity is immediate from homogeneity and the even
logarithmic-coordinate law.
\end{proof}

\section{Complete quintic power ledger}
\label{sec:V84-quintic-ledger}

\begin{lemma}[Fixed fifth-cumulant product bound]
\label{lem:V84-fifth-product}
For centered operator-valued \(Y_1,\ldots,Y_5\),
\begin{equation}
 \|\kappa_5(Y_1,\ldots,Y_5)\|
 \le C_5\prod_{i=1}^5\|Y_i\|_{L^5({\rm op})}.
 \label{eq:V84-fifth-product}
\end{equation}
\end{lemma}

\begin{proof}
Apply Holder to every block moment in the finite
moment--cumulant formula.  For each partition, the block exponents
can all be raised to five on the probability space.  Sum its fixed
coefficients.
\end{proof}

\begin{theorem}[Eight-incidence quintic estimate]
\label{thm:V84-eight-incidence}
Let \(X_i=D^{R_i}(U)\) and
\(b_i=\sqrt{s_\alpha A_i}<1\).  Then
\begin{equation}
 \boxed{\qquad
 \|\kappa_5^{\nu_\alpha}(X_1,\ldots,X_5)\|
 \le
 C
 \sum_{\substack{e_i\ge0\\\sum_i e_i=3}}
 \prod_{i=1}^5b_i^{1+e_i}.
 \qquad}
 \label{eq:V84-eight-incidence}
\end{equation}
Every exponent vector on the right is the degree vector minus one of
a labelled tree on five vertices.
\end{theorem}

\begin{proof}
Expand every centered representation with
\eqref{eq:V84-four-level} and use cumulant multilinearity.  Call
\(r-1\) the excess of a homogeneous Taylor term of degree \(r\).

The term with five linear factors vanishes because the joint law of
the five odd variables is invariant under simultaneous sign reversal.
There are two kinds of terms with total excess two:
one cubic and four linear factors, or two quadratic and three linear
factors.  The first contains five odd variables; the second contains
three.  In either case simultaneous inversion multiplies the
cumulant by \(-1\), so it vanishes.

The terms with total excess one are precisely one quadratic and four
linear factors.  By
\Cref{prop:V84-one-Q-defect}, such a term is bounded by
\[
 C\left(\prod_{r=1}^5b_r\right)b_i s_\alpha.
\]
Choose two distinct labels \(j,k\ne i\).  Since
\(s_\alpha\le b_jb_k\), this is bounded by
\[
 C\left(\prod_rb_r\right)b_ib_jb_k,
\]
one of the terms on the right of
\eqref{eq:V84-eight-incidence}.

Every term not yet treated has total Taylor excess at least three.
\Cref{lem:V84-four-level,lem:V84-fifth-product} bounds it by
\[
 C\left(\prod_rb_r\right)\prod_rb_r^{a_r},
 \qquad
 \sum_ra_r\ge3.
\]
Because \(b_r<1\), discard excess factors until nonnegative integers
\(e_r\le a_r\) with sum three remain.  This gives the corresponding
term in \eqref{eq:V84-eight-incidence}.  Sum the finitely many
Taylor choices.

Finally, positive integers \(d_i=1+e_i\) have sum eight.  By the
Pruefer degree theorem used in \Cref{lem:V69-resolution}, they are
the degrees of a labelled tree on five vertices.
\end{proof}

\begin{corollary}[The all-thermal local tree card holds at arity
five]
\label{cor:V84-quintic-tree-card}
The complete fifth Wilson cumulant has an exact decomposition
\[
 K_5=\sum_{\tau\in\mathfrak T_5}K_{5,\tau}
\]
such that
\begin{equation}
 \|K_{5,\tau}\|_{\rm op}
 \le5!K_5^5
 \prod_{i=1}^5
 (s_\alpha A_i)^{d_i(\tau)/2}.
 \label{eq:V84-quintic-tree-card}
\end{equation}
\end{corollary}

\begin{proof}
Every monomial on the right of
\eqref{eq:V84-eight-incidence} is a tree weight.  Therefore
\[
 \|K_5\|
 \le C\mathcal W_5,
 \qquad
 \mathcal W_5:=
 \sum_{\tau\in\mathfrak T_5}
 \prod_i b_i^{d_i(\tau)}.
\]
Define
\[
 K_{5,\tau}
 :=
 {\prod_i b_i^{d_i(\tau)}\over\mathcal W_5}K_5.
\]
This is exact, and the norm bound follows after enlarging one fixed
base.
\end{proof}

\begin{assessment}[What changed relative to the old quintic atlas]
\label{ass:V84-nonduplication}
V38--V54 analyzed several paid quintic fusion and predecessor words.
The present result is different: it proves the previously
hypothetical local all-thermal card
\eqref{eq:V64-thermal-card} itself at \(m=5\), before a final payer,
private bank, or selected fusion geometry is used.  It may therefore
be inserted into every later quintic ledger without spending those
resources twice.
\end{assessment}

\section{V84 ledger and V85 programme}
\label{sec:V84-results}

\begin{theorem}[Fourth-series V84 result ledger]
\label{thm:V84-results}
V84 proves:
\begin{enumerate}[label=\textup{(V84.\arabic*)},leftmargin=6.3em]
\item an \(O(\alpha^{-1})\) Wilson--Gaussian comparison for every
      tensor polynomial through degree six;
\item vanishing of the Gaussian connected word with one quadratic
      and four linear vertices;
\item the resulting \(s_\alpha\) gain for its Wilson counterpart;
\item cancellation of every quintic term with fewer than three
      Taylor excesses, except for that paid defect;
\item the complete eight-incidence quintic estimate; and
\item the V64 all-thermal local tree card at arity five.
\end{enumerate}
\end{theorem}

\begin{proof}
Use
\Cref{lem:V84-polynomial-comparison,
lem:V84-Gaussian-zero,
prop:V84-one-Q-defect,
lem:V84-four-level,
thm:V84-eight-incidence,
cor:V84-quintic-tree-card}.
\end{proof}

\begin{programme}[V85: audit and sixth-order connected defects]
\label{prog:V84-V85}
\begin{enumerate}[leftmargin=3em]
\item Perform the compulsory read-only audit of the newest active
      SM--Einstein and Navier--Stokes notes.
\item List every sixth-cumulant Taylor pattern with fewer than four
      excess incidences.
\item Use connected Gaussian half-edge counting to eliminate the
      patterns whose Gaussian cumulants vanish.
\item Prove the required first- or second-order Wilson--Gaussian
      comparison for each surviving pattern.
\item Close the ten-incidence arity-six tree card, or record the
      first pattern for which finite-degree comparison supplies too
      few powers.
\item Use the resulting pattern census to formulate the
      arbitrary-order connected-defect lemma without guessing its
      exponent.
\end{enumerate}
\end{programme}

\begin{firewall}[Claim boundary after fourth-series V84]
\label{fire:V84-claim-boundary}
V84 proves the all-thermal local card only through arity five.  It
does not prove the arbitrary-order connected-defect lemma, the
anisotropic scalar optimization, or the complete all-order selected
gate.  It does not prove all-order MTA, WUFC, CPM, cutoff removal,
reflection positivity, Osterwalder--Schrader reconstruction,
construction of the continuum Yang--Mills measure, or a uniform
positive continuum spectral gap.  The full Yang--Mills mass gap has
not been proved.
\end{firewall}

\paragraph{Parent-version record.}
V84 was formed from
\texttt{Renewal\_Yang\_Mills\_Mass\_Gap\_Research\_Notes\_V83.tex}.
The V83 parent had \(36\,572\) logical source lines,
\(1\,238\,627\) bytes, and SHA--256
\[
 \texttt{5dd4c5c871eaad632803a7db1e0beeb5a2780b28cc8f36ccc692d18da8654658}.
\]
The next compulsory companion audit is V85.

% =============================================================================
% END APPENDED FOURTH-SERIES V84 MATERIAL
% =============================================================================

% =============================================================================
% BEGIN APPENDED FOURTH-SERIES V85 MATERIAL
% =============================================================================

\chapter{V85: Audited second Wilson defect and the arity-six card}
\label{chap:V85-sixth-card}

\section{Compulsory companion audit}
\label{sec:V85-audit}

\begin{firewall}[Read-only status of the companion files]
\label{fire:V85-companion-status}
The companion documents inspected in this section are mathematical
evidence only.  Their internal programme statements are not
instructions for the Yang--Mills programme, and neither file was
modified.
\end{firewall}

The newest active companion files at this checkpoint are
\begin{align*}
 &\texttt{Renewal\_SM\_Einstein\_Continuum\_Research\_Notes\_V69.tex},
 \\
 &\texttt{Renewal\_NS\_Stability\_Research\_Notes\_V31.tex}.
\end{align*}
Their exact audit records are
\begin{center}
\begin{tabular}{c|r|r|l}
programme & lines & bytes & SHA--256 prefix\\ \hline
SM V69 & \(28\,100\) & \(958\,238\) &
\texttt{723c2955f7ba308b}\\
NS V31 & \(23\,052\) & \(887\,537\) &
\texttt{f1121b62238ad549}
\end{tabular}
\end{center}
The complete digests are
\[
\begin{split}
&\texttt{723c2955f7ba308b211d127c2e9cc45ec84dfc1f08075d911db2918f8b828b1b},
\\
&\texttt{f1121b62238ad5491bb7cf03635ea9934942edf573ed8faaa9ee79e1d38a6696}.
\end{split}
\]
The NS note remains byte-for-byte unchanged.  Its relevant warning
also remains unchanged: a scale-free estimate must retain the
complete signed correlation rather than replace it by independently
normed pieces.

SM V65--V69 distinguishes entrywise reflected-cylinder convergence
from determinant or operator-norm convergence.  In particular, a
global operator norm can hide a factor
\(\sqrt{\operatorname{rank}}\), and a cofinal argument requires a
Hilbert--Schmidt, trace-local, or rooted summability debit.  The
present finite-coordinate calculation is compatible with that
lesson because it bounds each tensor coefficient in operator norm
before any representation or cutoff sum.  V85 makes no inference
from determinant convergence and introduces no dimension trace.

\begin{assessment}[Audit direction]
\label{ass:V85-audit-direction}
At order six the correct object is the complete connected
coefficient, not a global moment norm.  The first variation of its
Wilson--Gaussian comparison must be tested for connectedness before
it is charged.  This preserves the two cancellations described
below and avoids an artificial rank or partition multiplicity.
\end{assessment}

\section{Second-order density expansion}
\label{sec:V85-density-expansion}

Let \(\epsilon=\alpha^{-1}\) and retain
\(Y_\alpha=\sqrt\alpha Z\).

\begin{lemma}[Normalized density through expansion]
\label{lem:V85-density-expansion}
Relative to \(\gamma_3\), the scaled logarithmic-coordinate law has
the fixed-degree expansion
\begin{equation}
 {d\mathcal L(Y_\alpha)\over d\gamma_3}(y)
 =
 1+\epsilon V_1(y)+\epsilon^2V_2(y)
 +O\left(
 \epsilon^3(1+|y|^{12})
 \right),
 \label{eq:V85-density-expansion}
\end{equation}
in every fixed polynomial moment.  Here
\begin{equation}
 V_1(y)
 =
 {|y|^4\over24}-{|y|^2\over3}-c_0,
 \qquad
 \mathbb E_{\gamma_3}V_1=0,
 \label{eq:V85-first-defect}
\end{equation}
and \(V_2\) is a centered even polynomial of degree at most eight.
More precisely, after multiplication by any fixed polynomial, the
integrated remainder is \(O(\epsilon^3)\).
\end{lemma}

\begin{proof}
Expand
\[
 \alpha\{\cos(\sqrt\epsilon|y|)-1\}
 =
 -{|y|^2\over2}
 +\epsilon{|y|^4\over24}
 -\epsilon^2{|y|^6\over720}
 +O(\epsilon^3|y|^8)
\]
and
\[
 2\log
 \left(
 {\sin(\sqrt\epsilon|y|)\over\sqrt\epsilon|y|}
 \right)
 =
 -\epsilon{|y|^2\over3}
 -\epsilon^2{|y|^4\over90}
 +O(\epsilon^3|y|^6).
\]
Exponentiating gives an even second coefficient of degree at most
eight.  Dividing by the normalizer subtracts the Gaussian means from
the first two coefficients and produces
\eqref{eq:V85-first-defect}.

On \(|y|\le\alpha^{1/10}\), Taylor's theorem gives the stated
polynomial remainder with a fixed Gaussian majorant after decreasing
its exponent.  Outside this ball,
\(1-\cos\theta\ge2\theta^2/\pi^2\) supplies an exponentially small
tail against every fixed polynomial.  This proves the integrated
expansion.
\end{proof}

\begin{lemma}[First variation of a joint cumulant]
\label{lem:V85-cumulant-variation}
Suppose probability laws \(\mu_\epsilon\) satisfy
\[
 {d\mu_\epsilon\over d\mu_0}
 =
 1+\epsilon V+O(\epsilon^2),
 \qquad
 \mathbb E_{\mu_0}V=0
\]
in all moments used below.  Then
\begin{equation}
 {d\over d\epsilon}
 \kappa_m^{\mu_\epsilon}(P_1,\ldots,P_m)
 \bigg|_{\epsilon=0}
 =
 \kappa_{m+1}^{\mu_0}(P_1,\ldots,P_m,V).
 \label{eq:V85-cumulant-variation}
\end{equation}
\end{lemma}

\begin{proof}
Introduce formal sources \(t_i\).  The cumulant is the coefficient of
\(t_1\cdots t_m\) in
\[
 \log\mathbb E_{\mu_\epsilon}
 \exp\left(\sum_it_iP_i\right).
\]
Differentiation at zero gives
\[
 {\mathbb E_{\mu_0}
 [V e^{\sum_it_iP_i}]
 \over
 \mathbb E_{\mu_0}e^{\sum_it_iP_i}},
\]
whose coefficient of \(t_1\cdots t_m\) is the joint cumulant with
the additional centered entry \(V\).
\end{proof}

\section{Gaussian half-edge deficits at order six}
\label{sec:V85-half-edge}

\begin{lemma}[Connected Gaussian half-edge criterion]
\label{lem:V85-half-edge}
Let \(P_1,\ldots,P_r\) be centered homogeneous polynomials of degrees
\(d_1,\ldots,d_r\) in one centered Gaussian vector.  If
\begin{equation}
 \sum_{i=1}^rd_i<2(r-1),
 \label{eq:V85-half-edge-deficit}
\end{equation}
then
\[
 \kappa_r^{\gamma_3}(P_1,\ldots,P_r)=0.
\]
\end{lemma}

\begin{proof}
Wick expansion of a Gaussian cumulant keeps exactly those pairings
whose contraction graph on the \(r\) declared polynomial vertices is
connected.  Such a graph needs at least \(r-1\) edges and therefore
at least \(2(r-1)\) polynomial half-edges.  Under
\eqref{eq:V85-half-edge-deficit} no connected pairing exists.
\end{proof}

\begin{lemma}[The six-linear cumulant begins at the second Wilson
defect]
\label{lem:V85-six-linear}
For linear operator-valued forms \(L_1,\ldots,L_6\) associated with
all-thermal representations,
\begin{equation}
 \left\|
 \kappa_6^{\nu_\alpha}(L_1,\ldots,L_6)
 \right\|
 \le
 Cs_\alpha^2\prod_{i=1}^6b_i.
 \label{eq:V85-six-linear}
\end{equation}
\end{lemma}

\begin{proof}
After writing \(Z=Y_\alpha/\sqrt\alpha\), factor the coefficient
norms \(Cb_i\) from the six linear forms.  Their Gaussian sixth
cumulant vanishes by
\Cref{lem:V85-half-edge}: six half-edges cannot connect six
vertices, which require ten.

By
\Cref{lem:V85-density-expansion,
lem:V85-cumulant-variation}, the coefficient of \(\epsilon\) is
\[
 \kappa_7^{\gamma_3}
 (L_1,\ldots,L_6,V_1).
\]
The degree-two part of \(V_1\) gives eight half-edges and the
degree-four part gives ten.  A connected graph on seven vertices
requires twelve half-edges.  Both contributions vanish by
\Cref{lem:V85-half-edge}.  The density remainder begins at
\(O(\epsilon^2)\) in every fixed moment, so the moment--cumulant
formula gives
\[
 \|\kappa_6(L_1,\ldots,L_6)\|
 \le C\epsilon^2\prod_ib_i.
\]
Use \(\epsilon\asymp s_\alpha\).
\end{proof}

\begin{lemma}[Every excess-two pattern gains one Wilson defect]
\label{lem:V85-excess-two}
Let six centered homogeneous representation terms have Taylor
degrees \(r_i\ge1\) satisfying
\[
 \sum_i(r_i-1)=2.
\]
Then their Wilson joint cumulant obeys
\begin{equation}
 \left\|
 \kappa_6^{\nu_\alpha}(P_{1,r_1},\ldots,P_{6,r_6})
 \right\|
 \le
 Cs_\alpha\prod_{i=1}^6b_i^{r_i}.
 \label{eq:V85-excess-two}
\end{equation}
\end{lemma}

\begin{proof}
The total Gaussian polynomial degree is eight.  A connected graph on
six vertices needs ten half-edges, so the Gaussian cumulant vanishes
by \Cref{lem:V85-half-edge}.  The first-order comparison following
from \Cref{lem:V85-density-expansion}, proved exactly as
\Cref{lem:V84-cumulant-comparison}, is \(O(\alpha^{-1})\) for tensor
polynomials through degree eight.  Factoring the coefficient norm
\(Cb_i^{r_i}\) from each entry proves the claim.
\end{proof}

\section{Five-level representation expansion}
\label{sec:V85-five-level}

Put
\[
 P_{R,r}:={L_R^r\over r!}
 \quad(1\le r\le4),
\]
and
\[
 N_{R,5}
 :=
 D^R(e^Z)-\sum_{r=0}^4{L_R^r\over r!}.
\]

\begin{lemma}[Expansion through quartic degree]
\label{lem:V85-five-level}
For an all-thermal representation,
\begin{equation}
 D^R(U)-\mathbb ED^R(U)
 =
 P_{R,1}+\widetilde P_{R,2}
 +P_{R,3}+\widetilde P_{R,4}
 +\widetilde N_{R,5},
 \label{eq:V85-five-level}
\end{equation}
where even-degree terms and the remainder are centered.  For fixed
\(p\),
\begin{equation}
 \|\widetilde P_{R,r}\|_{L^p}
 \le C_{p,r}b_R^r
 \quad(1\le r\le4),
 \qquad
 \|\widetilde N_{R,5}\|_{L^p}
 \le C_pb_R^5.
 \label{eq:V85-five-level-bounds}
\end{equation}
Odd homogeneous terms already have zero mean.
\end{lemma}

\begin{proof}
Use \Cref{lem:V82-log-moments} for each homogeneous term.  The
fifth-order unitary Taylor remainder obeys
\[
 \|N_{R,5}(Z)\|
 \le {1\over120}\|L_R(Z)\|^5.
\]
Centering costs at most a factor two.  Parity gives zero mean for
odd terms.
\end{proof}

\section{Complete arity-six ledger}
\label{sec:V85-sixth-ledger}

\begin{lemma}[Fixed sixth-cumulant product bound]
\label{lem:V85-sixth-product}
For centered operator-valued \(Y_1,\ldots,Y_6\),
\begin{equation}
 \|\kappa_6(Y_1,\ldots,Y_6)\|
 \le C_6\prod_{i=1}^6\|Y_i\|_{L^6({\rm op})}.
 \label{eq:V85-sixth-product}
\end{equation}
\end{lemma}

\begin{proof}
Apply Holder to the block moments in the fixed sixth
moment--cumulant formula and sum its finite coefficients.
\end{proof}

\begin{theorem}[Ten-incidence sixth-cumulant estimate]
\label{thm:V85-ten-incidence}
For all-thermal representation inputs \(X_i=D^{R_i}(U)\),
\begin{equation}
 \boxed{\qquad
 \|\kappa_6^{\nu_\alpha}(X_1,\ldots,X_6)\|
 \le
 C
 \sum_{\substack{e_i\ge0\\\sum_i e_i=4}}
 \prod_{i=1}^6b_i^{1+e_i}.
 \qquad}
 \label{eq:V85-ten-incidence}
\end{equation}
\end{theorem}

\begin{proof}
Expand by \eqref{eq:V85-five-level}.  A term has Taylor excess
\[
 E=\sum_i(r_i-1).
\]
If \(E\) is odd, the total polynomial degree \(6+E\) is odd.
Simultaneous inversion \(Z\mapsto-Z\) therefore changes the sign of
the joint cumulant, so every \(E=1\) and \(E=3\) term vanishes.

At \(E=0\), \Cref{lem:V85-six-linear} gives
\[
 C\left(\prod_ib_i\right)s_\alpha^2.
\]
Choose four distinct labels \(j_1,\ldots,j_4\).  Since
\[
 s_\alpha^2
 \le b_{j_1}b_{j_2}b_{j_3}b_{j_4},
\]
this is one term on the right of
\eqref{eq:V85-ten-incidence}.

At \(E=2\), \Cref{lem:V85-excess-two} gives
\[
 C\left(\prod_ib_i\right)
 \left(\prod_ib_i^{r_i-1}\right)s_\alpha.
\]
Choose two distinct labels \(j,k\).  The inequality
\(s_\alpha\le b_jb_k\) supplies the two remaining excess incidences.
The resulting exponent vector has total excess four.

Every term with \(E\ge4\) is bounded using
\Cref{lem:V85-five-level,lem:V85-sixth-product} by
\[
 C\left(\prod_ib_i\right)\prod_ib_i^{a_i},
 \qquad
 \sum_ia_i\ge4.
\]
Discard excess factors, all at most one, until integers
\(e_i\le a_i\) with sum four remain.  Summing the finite expansion
proves \eqref{eq:V85-ten-incidence}.
\end{proof}

\begin{corollary}[The all-thermal local tree card holds at arity
six]
\label{cor:V85-sixth-tree-card}
There is an exact decomposition
\[
 K_6=\sum_{\tau\in\mathfrak T_6}K_{6,\tau}
\]
such that
\begin{equation}
 \|K_{6,\tau}\|_{\rm op}
 \le6!K_6^6
 \prod_{i=1}^6
 (s_\alpha A_i)^{d_i(\tau)/2}.
 \label{eq:V85-sixth-tree-card}
\end{equation}
Thus \eqref{eq:V64-thermal-card} is proved for \(m=6\).
\end{corollary}

\begin{proof}
For every exponent vector in
\eqref{eq:V85-ten-incidence}, the positive integers
\(d_i=1+e_i\) sum to ten and therefore form the degree sequence of a
labelled tree on six vertices.  It follows that
\[
 \|K_6\|\le C
 \sum_{\tau\in\mathfrak T_6}\prod_i b_i^{d_i(\tau)}.
\]
Split \(K_6\) proportionally among these positive tree weights, as
in \Cref{cor:V84-quintic-tree-card}.  This gives the exact
decomposition and its bound.
\end{proof}

\section{The emerging arbitrary-order rule}
\label{sec:V85-general-rule}

\begin{proposition}[Defect count forced by Gaussian connectivity]
\label{prop:V85-defect-count}
Consider an \(m\)-input Taylor pattern with total Taylor excess
\[
 E=\sum_i(r_i-1).
\]
If \(m+E\) is odd, its Wilson cumulant vanishes exactly.  If
\(m+E\) is even, a connected Gaussian contraction requires at least
\begin{equation}
 q_{\min}
 =
 \max\left\{
 0,\left\lceil{m-2-E\over2}\right\rceil
 \right\}
 \label{eq:V85-min-defects}
\end{equation}
quartic first-defect insertions before the half-edge count can reach
the connectivity threshold.  The corresponding formal Wilson
factor is \(s_\alpha^{q_{\min}}\).
\end{proposition}

\begin{proof}
The \(m\) representation vertices contain \(m+E\) polynomial
half-edges.  Adding \(q\) first-defect vertices supplies at most
\(4q\) more half-edges and also adds \(q\) vertices.  A connected
graph then requires
\[
 m+E+4q\ge2(m+q-1),
\]
equivalently
\[
 2q\ge m-2-E.
\]
This proves \eqref{eq:V85-min-defects}.  The parity assertion follows
from exact inversion symmetry.
\end{proof}

\begin{firewall}[Half-edge necessity is not yet a uniform
all-order estimate]
\label{fire:V85-rule-firewall}
\Cref{prop:V85-defect-count} identifies how many density-defect
coefficients must vanish before a connected Gaussian graph can
exist.  It does not yet bound the order-\(q\) normalized density
coefficient uniformly in \(m,q\), nor sum all Taylor degrees with a
single factorial.  Those are the two remaining analytic ingredients
for an all-order local card.
\end{firewall}

\section{V85 ledger and V86 programme}
\label{sec:V85-results}

\begin{theorem}[Fourth-series V85 result ledger]
\label{thm:V85-results}
V85 proves:
\begin{enumerate}[label=\textup{(V85.\arabic*)},leftmargin=6.3em]
\item the compulsory read-only SM V69 and NS V31 audit;
\item the normalized Wilson logarithm density through second defect
      order;
\item the exact cumulant first-variation identity;
\item the connected Gaussian half-edge vanishing criterion;
\item the \(s_\alpha^2\) bound for the six-linear cumulant;
\item the \(s_\alpha\) bound for every excess-two sixth-order
      pattern;
\item the complete ten-incidence sixth-cumulant estimate and
      all-thermal tree card; and
\item the general minimum-defect count
      \eqref{eq:V85-min-defects}.
\end{enumerate}
\end{theorem}

\begin{proof}
Use
\Cref{lem:V85-density-expansion,
lem:V85-cumulant-variation,
lem:V85-half-edge,
lem:V85-six-linear,
lem:V85-excess-two,
thm:V85-ten-incidence,
cor:V85-sixth-tree-card,
prop:V85-defect-count}.
\end{proof}

\begin{programme}[V86: uniform defect-coefficient majorant]
\label{prog:V85-V86}
\begin{enumerate}[leftmargin=3em]
\item Write the normalized scaled Wilson density as a formal
      exponential of even interaction vertices of degrees
      \(2r\), each carrying \(s_\alpha^{r-1}\).
\item Prove a coefficientwise majorant for its normalized defect
      polynomials whose connected Gaussian contractions cost at most
      one global factorial times an exponential.
\item Combine the minimum-defect count
      \eqref{eq:V85-min-defects} with representation Taylor excesses
      to obtain exactly \(m-2\) extra input incidences.
\item Use the V69 Steiner resolution to send every correction-vertex
      excess to a labelled input tree.
\item Keep tensor coefficients in operator norm before summing any
      representation bank, in accordance with the V85 audit.
\item If the global logarithmic ball prevents a uniform coefficient
      majorant, isolate the first order at which the antipodal tail
      exceeds it rather than reverting to a mesoscopic cutoff.
\end{enumerate}
\end{programme}

\begin{firewall}[Claim boundary after fourth-series V85]
\label{fire:V85-claim-boundary}
V85 proves the all-thermal local card only through arity six and a
formal general defect count.  It does not prove the uniform
all-order defect majorant, the anisotropic scalar optimization, or
the complete all-order selected gate.  It does not prove all-order
MTA, WUFC, CPM, cutoff removal, reflection positivity,
Osterwalder--Schrader reconstruction, construction of the continuum
Yang--Mills measure, or a uniform positive continuum spectral gap.
The full Yang--Mills mass gap has not been proved.
\end{firewall}

\paragraph{Parent-version record.}
V85 was formed from
\texttt{Renewal\_Yang\_Mills\_Mass\_Gap\_Research\_Notes\_V84.tex}.
The V84 parent had \(36\,979\) logical source lines,
\(1\,251\,445\) bytes, and SHA--256
\[
 \texttt{f368960cfbcad8841ae1b94dfac8ab4db6786d56461fcfe930946a7274d93ed0}.
\]
The next compulsory companion audit is V90.

% =============================================================================
% END APPENDED FOURTH-SERIES V85 MATERIAL
% =============================================================================

% =============================================================================
% BEGIN APPENDED FOURTH-SERIES V86 MATERIAL
% =============================================================================

\chapter{V86: Haar-Jacobian factorial no-go and intrinsic
redirection}
\label{chap:V86-Haar-no-go}

\section{Separating the action and Haar corrections}
\label{sec:V86-two-corrections}

Put \(\epsilon=\alpha^{-1}\) and use the scaled logarithmic
coordinate \(Z=\sqrt\epsilon\,y\).  After the Gaussian quadratic
term has been removed, the Wilson-action correction is
\begin{equation}
 \mathcal W_\epsilon(y)
 :=
 \epsilon^{-1}
 \left\{
 \cos(\sqrt\epsilon|y|)-1
 \right\}
 +{|y|^2\over2}.
 \label{eq:V86-action-correction}
\end{equation}
The logarithm of the Haar Jacobian is
\begin{equation}
 \mathcal H_\epsilon(y)
 :=
 2\log
 \left(
 {\sin(\sqrt\epsilon|y|)
  \over\sqrt\epsilon|y|}
 \right).
 \label{eq:V86-Haar-correction}
\end{equation}
The finite-order arguments in V82--V85 expand both expressions only
to the order required by a fixed cumulant.  V86 asks whether their
derivatives admit one uniform exponential-in-order majorant.

\begin{lemma}[The action has the exact Steiner derivative scale]
\label{lem:V86-action-scale}
Along every unit Lie-algebra direction \(e\),
\begin{equation}
 {d^{2n}\over dt^{2n}}
 \mathcal W_\epsilon(te)\bigg|_{t=0}
 =
 (-1)^n\epsilon^{n-1}
 \qquad(n\ge2),
 \label{eq:V86-action-derivative}
\end{equation}
and all odd derivatives vanish.
\end{lemma}

\begin{proof}
The exact power series is
\[
 \mathcal W_\epsilon(y)
 =
 \sum_{n=2}^\infty
 {(-1)^n\epsilon^{n-1}|y|^{2n}\over(2n)!}.
\]
Restrict to \(y=te\) and differentiate.  The factorial in the
Taylor coefficient cancels the derivative factorial.
\end{proof}

\begin{remark}[Match to V69]
\label{rem:V86-action-match}
A correction vertex of degree \(k=2n\) carries
\(\epsilon^{n-1}=\epsilon^{(k-2)/2}\), exactly the scale in
\eqref{eq:V69-required-correction-scale}.  Since
\(s_\alpha\asymp\epsilon\), changing from \(\epsilon\) to
\(s_\alpha\) costs only an exponential base in \(k\).
\end{remark}

\section{Exact Haar coefficients}
\label{sec:V86-Haar-series}

\begin{lemma}[Haar product and logarithmic series]
\label{lem:V86-Haar-series}
For \(|r|<\pi\),
\begin{align}
 {\sin r\over r}
 &=
 \prod_{q=1}^\infty
 \left(1-{r^2\over\pi^2q^2}\right),
 \label{eq:V86-sine-product}\\
 2\log{\sin r\over r}
 &=
 -2\sum_{n=1}^\infty
 {\zeta(2n)\over n\pi^{2n}}r^{2n}.
 \label{eq:V86-Haar-series}
\end{align}
Consequently,
\begin{equation}
 \left|
 {d^{2n}\over dt^{2n}}
 \mathcal H_\epsilon(te)\bigg|_{t=0}
 \right|
 =
 {2(2n)!\zeta(2n)\over n\pi^{2n}}\epsilon^n.
 \label{eq:V86-Haar-derivative}
\end{equation}
\end{lemma}

\begin{proof}
Equation \eqref{eq:V86-sine-product} is the canonical product for
the sine function.  Taking the logarithm and using
\[
 \log(1-x)=-\sum_{n\ge1}{x^n\over n}
\]
gives
\[
 \log{\sin r\over r}
 =
 -\sum_{q\ge1}\sum_{n\ge1}
 {r^{2n}\over n\pi^{2n}q^{2n}},
\]
which converges absolutely for \(|r|<\pi\).  Summing \(q\) proves
\eqref{eq:V86-Haar-series}.  Substitute
\(r=\sqrt\epsilon\,t\) and differentiate.
\end{proof}

\begin{theorem}[No exponential all-order absolute Haar derivative
majorant]
\label{thm:V86-Haar-no-go}
Fix any \(\epsilon>0\).  There do not exist finite constants
\(C,K\), independent of \(n\), such that
\begin{equation}
 \left|
 {d^{2n}\over dt^{2n}}
 \mathcal H_\epsilon(te)\bigg|_{t=0}
 \right|
 \le
 CK^{2n}\epsilon^{n-1}
 \qquad(n\ge1).
 \label{eq:V86-false-Haar-majorant}
\end{equation}
\end{theorem}

\begin{proof}
By \eqref{eq:V86-Haar-derivative}, the ratio of the left side to
the right side without \(C\) is
\[
 {2\epsilon(2n)!\zeta(2n)
  \over n(\pi K)^{2n}}.
\]
Since \(\zeta(2n)\ge1\) and
\((2n)!^{1/(2n)}\sim2n/e\), this ratio tends to infinity for every
fixed \(K\).  No finite \(C\) can make
\eqref{eq:V86-false-Haar-majorant} hold.
\end{proof}

\begin{corollary}[Failure of the absolute Taylor implementation of
the SSCA]
\label{cor:V86-SSCA-no-go}
The global logarithmic-coordinate factorization cannot prove the
V69 SSCA by taking absolute values of each Haar-Jacobian derivative
and demanding an exponential local constant before the connected
Gaussian sum.  The extra factor \(\epsilon\) in
\eqref{eq:V86-Haar-derivative} does not dominate its derivative
factorial uniformly in the vertex degree.
\end{corollary}

\begin{proof}
At degree \(k=2n\), the SSCA scale is comparable with
\(\epsilon^{n-1}\) times an exponential in \(2n\).  Such a bound is
exactly \eqref{eq:V86-false-Haar-majorant}, which
\Cref{thm:V86-Haar-no-go} disproves.
\end{proof}

\begin{firewall}[Scope of the no-go theorem]
\label{fire:V86-no-go-scope}
The theorem rules out a specific absolute local-derivative
implementation.  It does not rule out cancellation after the
complete connected Gaussian contraction, a Gevrey resummation, or an
intrinsic group identity in which Haar measure is never converted
to a logarithmic potential.  In particular it does not invalidate
the fixed arity-three through arity-six cards: those use only
finitely many derivatives, with constants fixed at the declared
arity.
\end{firewall}

\section{The cut-locus interpretation}
\label{sec:V86-cut-locus}

\begin{proposition}[The factorial records the antipodal singularity]
\label{prop:V86-cut-locus}
The radius of analyticity of
\(\mathcal H_\epsilon(te)\) at \(t=0\) is
\[
 {\pi\over\sqrt\epsilon}.
\]
Its nearest singularities are the zeros of
\(\sin(\sqrt\epsilon\,t)\) at
\(t=\pm\pi/\sqrt\epsilon\).  Thus the factorial in
\eqref{eq:V86-Haar-derivative} is the analytic signature of the
antipodal cut locus, not a defect of the coefficient estimate.
\end{proposition}

\begin{proof}
The product
\eqref{eq:V86-sine-product} shows both the singularity locations and
the absence of a nearer one.  Cauchy's coefficient formula then
forces derivative growth on the scale
\((2n)!(\sqrt\epsilon/\pi)^{2n}\), in agreement with
\eqref{eq:V86-Haar-derivative}.
\end{proof}

\section{Why the intrinsic calculus avoids this factorial}
\label{sec:V86-intrinsic}

Retain
\[
 \phi(U)={1\over2}\operatorname{ReTr}U.
\]

\begin{lemma}[Exponential-order intrinsic Wilson derivatives]
\label{lem:V86-intrinsic-derivatives}
There is \(C_0<\infty\) such that for every word
\(\mathbf a=(a_1,\ldots,a_k)\) in the three normalized
left-invariant vector fields,
\begin{equation}
 \|X_{a_1}\cdots X_{a_k}\phi\|_{L^\infty(SU(2))}
 \le C_0^k.
 \label{eq:V86-intrinsic-derivatives}
\end{equation}
\end{lemma}

\begin{proof}
Let \(T_a=dD^{1/2}(X_a)\) in the fundamental representation.  Up to
the harmless real trace,
\[
 X_{a_1}\cdots X_{a_k}\phi(U)
 =
 {1\over2}\operatorname{ReTr}
 \left(
 T_{a_1}\cdots T_{a_k}D^{1/2}(U)
 \right).
\]
Unitarity bounds this by a fixed dimension factor times
\(\prod_i\|T_{a_i}\|\).  The three generators have one fixed maximum
norm; absorb it into \(C_0\).
\end{proof}

\begin{corollary}[Intrinsic commutator coefficients have no local
factorial]
\label{cor:V86-commutator-coefficients}
Every coefficient obtained by applying a word of left-invariant
derivatives to the commutator
\([A_\alpha,X_b]\) in
\eqref{eq:V78-commutator} is bounded by
\[
 \alpha C_1^k
\]
at word length \(k\), apart from the finite structure constants of
\(\mathfrak{su}(2)\).  No derivative of the Haar density occurs.
\end{corollary}

\begin{proof}
Formula \eqref{eq:V78-commutator} contains only first-order vector
fields, structure constants, and derivatives of \(\phi\).
Apply \Cref{lem:V86-intrinsic-derivatives}.  The left-invariant Haar
divergence is zero, so integration by parts introduces the Wilson
score from \Cref{lem:V78-score}, not a logarithmic-coordinate
Jacobian.
\end{proof}

\begin{theorem}[Coefficient-level intrinsic replacement]
\label{thm:V86-intrinsic-replacement}
At the algebraic coefficient level, the intrinsic Wilson
score/resolvent recursion has:
\begin{enumerate}[leftmargin=3em]
\item no Haar-Jacobian derivative vertex;
\item only exponential-order derivatives of the fixed potential
      \(\phi\);
\item one factor \(\alpha\) at each drift commutator, paired with a
      mean-zero resolvent of scale \(\alpha^{-1}\); and
\item at most \(C^m(m-1)!\) formal differentiation histories.
\end{enumerate}
Thus the factorial obstruction of
\Cref{thm:V86-Haar-no-go} is absent from the intrinsic coefficient
census.
\end{theorem}

\begin{proof}
The first assertion is the Haar integration-by-parts identity
\eqref{eq:V78-score}.  The second is
\Cref{lem:V86-intrinsic-derivatives}; the third follows from
\eqref{eq:V78-commutator} and
\eqref{eq:V78-resolvent-L2}.  The fourth is the plane-rooted-tree
census \eqref{eq:V79-history-count}.  These statements concern the
same recursion and therefore combine without an independent
factorial multiplication.
\end{proof}

\begin{firewall}[The remaining intrinsic analytic gate]
\label{fire:V86-intrinsic-gate}
\Cref{thm:V86-intrinsic-replacement} controls coefficients and formal
histories, not the norms of iterated gradient products.  A complete
all-order proof still requires an intrinsic energy or carré-du-champ
norm stable under
\[
 dR_\alpha\langle dh,dR_\alpha g\rangle
\]
and its higher tree analogues, with the full connected sum retained
when individual rooted terms have the sharp half-power of V81.
\end{firewall}

\section{V86 ledger and V87 programme}
\label{sec:V86-results}

\begin{theorem}[Fourth-series V86 result ledger]
\label{thm:V86-results}
V86 proves:
\begin{enumerate}[label=\textup{(V86.\arabic*)},leftmargin=6.3em]
\item exact all-order scaled derivatives of the Wilson-action
      correction at the V69 Steiner scale;
\item the canonical-product formula and exact derivatives of the
      Haar logarithmic Jacobian;
\item impossibility of an exponential-in-order absolute Haar
      derivative majorant;
\item identification of that factorial with the antipodal cut locus;
\item exponential-order intrinsic derivatives of the Wilson
      potential; and
\item absence of the Haar factorial from the intrinsic
      coefficient/history census.
\end{enumerate}
\end{theorem}

\begin{proof}
Use
\Cref{lem:V86-action-scale,
lem:V86-Haar-series,
thm:V86-Haar-no-go,
prop:V86-cut-locus,
lem:V86-intrinsic-derivatives,
thm:V86-intrinsic-replacement}.
\end{proof}

\begin{programme}[V87: intrinsic connected energy norm]
\label{prog:V86-V87}
\begin{enumerate}[leftmargin=3em]
\item Stop the absolute global logarithmic-Jacobian expansion at
      all orders; retain V82--V85 only as fixed-order cards.
\item Define a rooted energy norm for exact one-forms and their
      scalar contractions, using the \(c\alpha\) gap of V78.
\item Derive the complete fourth intrinsic cumulant before taking
      norms and identify its cancellation of paired rooted terms.
\item Test whether the energy norm is closed under one drift
      commutator plus resolvent without an \(L^\infty\) Hessian.
\item Compare the resulting fourth-order intrinsic expression with
      the already proved V83 card as a normalization check.
\item Generalize only after the norm reproduces every incidence of
      the fixed-order cards.
\end{enumerate}
\end{programme}

\begin{firewall}[Claim boundary after fourth-series V86]
\label{fire:V86-claim-boundary}
V86 proves a no-go theorem for an absolute Haar-Taylor route and an
intrinsic coefficient-level replacement.  It does not prove the
intrinsic all-order analytic norm, the all-order local card, the
anisotropic scalar optimization, or the complete all-order selected
gate.  It does not prove all-order MTA, WUFC, CPM, cutoff removal,
reflection positivity, Osterwalder--Schrader reconstruction,
construction of the continuum Yang--Mills measure, or a uniform
positive continuum spectral gap.  The full Yang--Mills mass gap has
not been proved.
\end{firewall}

\paragraph{Parent-version record.}
V86 was formed from
\texttt{Renewal\_Yang\_Mills\_Mass\_Gap\_Research\_Notes\_V85.tex}.
The V85 parent had \(37\,537\) logical source lines,
\(1\,268\,307\) bytes, and SHA--256
\[
 \texttt{af7092c0d935ae31da31d05c0f37d722cc7a85f150f0d362bd807b6c1a2345e5}.
\]
The next compulsory companion audit is V90.

% =============================================================================
% END APPENDED FOURTH-SERIES V86 MATERIAL
% =============================================================================

% =============================================================================
% BEGIN APPENDED FOURTH-SERIES V87 MATERIAL
% =============================================================================

\chapter{V87: Exact intrinsic chain formula and its norm boundary}
\label{chap:V87-chain-formula}

\section{The tilted scalar resolvent}
\label{sec:V87-tilted-resolvent}

For a smooth real \(h\), let
\[
 d\nu_{\alpha,t}
 =
 {e^{th}\over\mathbb E_{\nu_\alpha}e^{th}}d\nu_\alpha.
\]
Its scalar weighted generator and mean-zero resolvent are
\[
 A_{\alpha,t}
 =
 A_\alpha-t\Gamma(h,\cdot),
 \qquad
 R_{\alpha,t}=A_{\alpha,t}^{-1}Q_{\alpha,t}.
\]

\begin{lemma}[Differentiating a tilted scalar resolvent]
\label{lem:V87-tilted-resolvent}
Let \(q_t\) be a smooth differentiable family and put
\(u_t=R_{\alpha,t}q_t\).  At \(t=0\),
\begin{equation}
 d\dot u_0
 =
 dR_\alpha
 \left\{
 \dot q_0+\Gamma(h,R_\alpha q_0)
 \right\}.
 \label{eq:V87-tilted-resolvent}
\end{equation}
Only the differential is asserted; the derivative of the
mean-zero normalization may add a constant to \(\dot u_0\).
\end{lemma}

\begin{proof}
Differentiate
\[
 A_{\alpha,t}u_t
 =
 q_t-\mathbb E_{\nu_{\alpha,t}}q_t.
\]
Since
\[
 \partial_tA_{\alpha,t}|_{t=0}
 =-\Gamma(h,\cdot)
\]
and
\[
 {d\over dt}\mathbb E_{\nu_{\alpha,t}}q_t\bigg|_{t=0}
 =
 \mathbb E\dot q_0+\operatorname{Cov}(q_0,h),
\]
one obtains
\[
 A_\alpha\dot u_0
 =
 Q_\alpha\dot q_0
 +\Gamma(h,R_\alpha q_0)
 -\operatorname{Cov}(q_0,h).
\]
By the intrinsic covariance identity,
\[
 \mathbb E\Gamma(h,R_\alpha q_0)
 =
 \operatorname{Cov}(h,q_0).
\]
Thus the right side is centered.  Applying \(R_\alpha\) determines
the mean-zero part of \(\dot u_0\); applying \(d\) removes its
constant part and proves \eqref{eq:V87-tilted-resolvent}.
\end{proof}

\section{The intrinsic binary operation}
\label{sec:V87-binary-operation}

Define the measure-dependent bilinear operation
\begin{equation}
 \mathcal B_\alpha(f,g)
 :=
 \Gamma(f,R_\alpha g).
 \label{eq:V87-B}
\end{equation}
For operator-valued functions on separate tensor legs, derivatives
are placed on their declared legs and the contraction over the
three Lie-algebra directions is taken before the canonical tensor
rebracketing.

\begin{lemma}[Tilt derivative of the binary operation]
\label{lem:V87-B-variation}
If \(f_t,g_t\) are differentiable families, then
\begin{align}
 {d\over dt}
 \mathcal B_{\alpha,t}(f_t,g_t)\bigg|_{t=0}
 &=
 \mathcal B_\alpha(\dot f_0,g_0)
 +\mathcal B_\alpha(f_0,\dot g_0)
 \nonumber\\
 &\quad+
 \mathcal B_\alpha
 \left(
 f_0,\mathcal B_\alpha(h,g_0)
 \right).
 \label{eq:V87-B-variation}
\end{align}
\end{lemma}

\begin{proof}
Differentiate
\[
 \Gamma(f_t,R_{\alpha,t}g_t).
\]
The metric and \(\Gamma\) do not depend on \(t\).  The derivative of
the first argument gives the first term.  Apply
\Cref{lem:V87-tilted-resolvent} to the second argument; its
\(\dot g_0\) and \(\Gamma(h,R_\alpha g_0)\) pieces give the other two
terms.
\end{proof}

\begin{lemma}[Expectation derivative in oriented form]
\label{lem:V87-expectation-variation}
For a differentiable scalar family \(F_t\),
\begin{equation}
 {d\over dt}\mathbb E_{\nu_{\alpha,t}}F_t\bigg|_{t=0}
 =
 \mathbb E_{\nu_\alpha}
 \left[
 \dot F_0+\mathcal B_\alpha(h,F_0)
 \right].
 \label{eq:V87-expectation-variation}
\end{equation}
\end{lemma}

\begin{proof}
The derivative is
\(\mathbb E\dot F_0+\operatorname{Cov}(F_0,h)\).
Symmetry of covariance and
\eqref{eq:V78-covariance} give
\[
 \operatorname{Cov}(F_0,h)
 =
 \mathbb E\mathcal B_\alpha(h,F_0).
\]
\end{proof}

\section{All cumulants as ordered resolvent chains}
\label{sec:V87-chain}

For an ordered list \((i_1,\ldots,i_m)\), define the right-nested
chain
\begin{equation}
 \mathcal C_\alpha
 (f_{i_1},\ldots,f_{i_m})
 :=
 \mathcal B_\alpha
 \left(
 f_{i_1},
 \mathcal B_\alpha
 \left(
 f_{i_2},\ldots,
 \mathcal B_\alpha(f_{i_{m-1}},f_{i_m})
 \right)
 \right).
 \label{eq:V87-chain}
\end{equation}

\begin{theorem}[Exact intrinsic ordered-chain formula]
\label{thm:V87-chain-formula}
For every \(m\ge2\) and smooth scalar
\(f_1,\ldots,f_m\),
\begin{equation}
 \boxed{\qquad
 \kappa_m^{\nu_\alpha}(f_1,\ldots,f_m)
 =
 \sum_{\pi\in S_{m-1}}
 \mathbb E_{\nu_\alpha}
 \mathcal C_\alpha
 (f_{\pi(1)},\ldots,f_{\pi(m-1)},f_m).
 \qquad}
 \label{eq:V87-chain-formula}
\end{equation}
The formula extends entrywise to tensor-valued inputs on separate
legs, with canonical permutations restoring the declared leg order.
\end{theorem}

\begin{proof}
For \(m=2\), the formula is the covariance identity
\[
 \kappa_2(f_1,f_2)
 =
 \mathbb E\mathcal B_\alpha(f_1,f_2).
\]
Assume the assertion at order \(m\), and tilt by
\(e^{tf_{m+1}}\).  The derivative of a joint cumulant under a
normalized exponential tilt is
\[
 {d\over dt}
 \kappa_m^{\nu_{\alpha,t}}(f_1,\ldots,f_m)
 \bigg|_{t=0}
 =
 \kappa_{m+1}^{\nu_\alpha}
 (f_1,\ldots,f_m,f_{m+1}).
\]

Fix one chain in the induction hypothesis.  By
\Cref{lem:V87-expectation-variation}, differentiating its
expectation inserts \(f_{m+1}\) at the outermost position.  By
\Cref{lem:V87-B-variation}, differentiating each successive
resolvent inserts \(f_{m+1}\) immediately before its current right
argument.  These are exactly the \(m\) possible positions of the new
label before the fixed final entry \(f_m\).

As \(\pi\) runs through \(S_{m-1}\), insertion in those \(m\)
positions produces every permutation in \(S_m\) exactly once.
This proves \eqref{eq:V87-chain-formula} at order \(m+1\).
Entrywise differentiation proves the tensor statement; cumulant
symmetry supplies the canonical leg permutations.
\end{proof}

\begin{corollary}[No Bell partition sum]
\label{cor:V87-no-Bell}
The exact intrinsic formula has precisely \((m-1)!\) rooted chain
terms.  Its algebraic census contains one global factorial and no
independent Bell or vertex factorial.
\end{corollary}

\begin{proof}
The terms are indexed by \(S_{m-1}\) and occur exactly once in
\eqref{eq:V87-chain-formula}.
\end{proof}

\section{Dimension-free column energy}
\label{sec:V87-column-energy}

For a finite-dimensional operator-valued function \(F\), define
\[
 \|F\|_{2,c}
 :=
 \left\|
 \mathbb E_{\nu_\alpha}F^*F
 \right\|_{\rm op}^{1/2}.
\]
For an operator-valued one-form
\(\omega=(\omega_a)_{a=1}^3\), put
\[
 \|\omega\|_{2,c}
 :=
 \left\|
 \mathbb E_{\nu_\alpha}
 \sum_a\omega_a^*\omega_a
 \right\|_{\rm op}^{1/2}.
\]

\begin{lemma}[Complete column resolvent estimate]
\label{lem:V87-column-resolvent}
For every finite-dimensional operator-valued \(G\),
\begin{equation}
 \|dR_\alpha G\|_{2,c}
 \le {C\over\sqrt\alpha}\|Q_\alpha G\|_{2,c}.
 \label{eq:V87-column-resolvent}
\end{equation}
The constant is independent of the target dimension.
\end{lemma}

\begin{proof}
Fix a unit vector \(v\) in the input space.  The function
\((Q_\alpha G)v\) takes values in a Hilbert space, and
\Cref{thm:V78-Wilson-gap} tensorized with that Hilbert space gives
\[
 \mathbb E\sum_a
 \|X_aR_\alpha G\,v\|^2
 \le {C\over\alpha}
 \mathbb E\|Q_\alpha G\,v\|^2.
\]
Take the supremum over \(v\).  The two sides are precisely the
squares of the declared column norms.
\end{proof}

\begin{lemma}[One atomic-left chain step]
\label{lem:V87-chain-step}
Let \(D^R\) be a unitary representation with
\(b_R=\sqrt{s_\alpha(1+C_2(R))}<1\).  Then
\begin{equation}
 \|\mathcal B_\alpha(D^R,G)\|_{2,c}
 \le Cb_R\|Q_\alpha G\|_{2,c},
 \label{eq:V87-chain-step}
\end{equation}
and
\begin{equation}
 \|Q_\alpha\mathcal B_\alpha(D^R,G)\|_{2,c}
 \le Cb_R\|Q_\alpha G\|_{2,c}.
 \label{eq:V87-centered-chain-step}
\end{equation}
\end{lemma}

\begin{proof}
The Casimir row estimate
\eqref{eq:V80-Casimir-row}, in column norm, gives
\[
 \|\Gamma(D^R,R_\alpha G)\|_{2,c}
 \le
 \sqrt{C_2(R)}\,
 \|dR_\alpha G\|_{2,c}.
\]
Apply \Cref{lem:V87-column-resolvent} and
\(\sqrt{C_2(R)/\alpha}\le Cb_R\).
Centering a random operator costs at most two in column norm, proving
the second inequality.
\end{proof}

\begin{theorem}[Termwise chain bound]
\label{thm:V87-termwise-chain}
Let \(D_i=D^{R_i}\) be all-thermal unitary representations and set
\(b_i=\sqrt{s_\alpha A_i}<1\).  Every chain in
\eqref{eq:V87-chain-formula} satisfies
\begin{equation}
 \left\|
 \mathbb E
 \mathcal C_\alpha
 (D_{\pi(1)},\ldots,D_{\pi(m-1)},D_m)
 \right\|_{\rm op}
 \le
 C^m\prod_{i=1}^mb_i.
 \label{eq:V87-termwise-chain}
\end{equation}
Consequently
\begin{equation}
 \|\kappa_m^{\nu_\alpha}(D_1,\ldots,D_m)\|
 \le
 (m-1)!C^m\prod_{i=1}^mb_i.
 \label{eq:V87-chain-sum-bound}
\end{equation}
\end{theorem}

\begin{proof}
The centered input bound from
\Cref{lem:V82-centered-decomposition} gives
\[
 \|Q_\alpha D_i\|_{2,c}\le Cb_i.
\]
Starting at the innermost \(D_m\), apply
\eqref{eq:V87-centered-chain-step} successively to every non-outer
chain operation.  This gives
\[
 \|Q_\alpha G_2\|_{2,c}
 \le C^{m-1}\prod_{i=2}^mb_{\pi(i)}
\]
with the obvious relabelling of the final entry.  For the outer
operation, use
\[
 \|\mathbb E\mathcal B_\alpha(D_{\pi(1)},G_2)\|
 \le
 \|\mathcal B_\alpha(D_{\pi(1)},G_2)\|_{2,c}
\]
and \eqref{eq:V87-chain-step}.  This proves
\eqref{eq:V87-termwise-chain}.  Sum the \((m-1)!\) terms.
\end{proof}

\section{Exact location of the missing incidences}
\label{sec:V87-boundary}

\begin{proposition}[Termwise energy misses exactly \(m-2\)
incidences]
\label{prop:V87-missing-incidences}
A labelled tree on \(m\) inputs has \(2m-2\) endpoint incidences.
The termwise chain estimate
\eqref{eq:V87-termwise-chain} supplies \(m\), one at every input.
Thus it is short by exactly
\[
 (2m-2)-m=m-2
\]
thermal factors.  In the balanced fixed-representation regime
\(b_i\asymp\sqrt{s_\alpha}\), this is the missing factor
\[
 s_\alpha^{(m-2)/2}.
\]
\end{proposition}

\begin{proof}
The incidence count is the handshaking identity for a tree.  The
chain estimate contains the monomial \(\prod_ib_i\).  Subtraction
gives the assertion.
\end{proof}

\begin{corollary}[The finite-order cards are cancellations among
chains]
\label{cor:V87-finite-cards-as-chain-cancellation}
At orders three through six, the additional incidences proved in
V82--V85 cannot be recovered by bounding the individual terms in
\eqref{eq:V87-chain-formula}.  They arise only after the
permutation-chain sum is assembled.
\end{corollary}

\begin{proof}
\Cref{prop:V87-missing-incidences} shows that the separate chain
bounds lack \(m-2\) factors.  The complete cumulants satisfy the full
tree cards at these orders by
\Cref{cor:V82-ternary-tree-card,
cor:V83-quartic-tree-card,
cor:V84-quintic-tree-card,
cor:V85-sixth-tree-card}.
\end{proof}

\begin{firewall}[A closed chain norm is not an all-order tree norm]
\label{fire:V87-chain-firewall}
The column norm is stable under every right-nested chain operation
and is dimension-free, but applying it before summing permutations
loses the connected Gaussian/Wilson defects.  Therefore
\eqref{eq:V87-chain-sum-bound} is not the V64 all-thermal card and
must not be inserted into the endpoint ledger as though it had
\(2m-2\) incidences.
\end{firewall}

\section{V87 ledger and V88 programme}
\label{sec:V87-results}

\begin{theorem}[Fourth-series V87 result ledger]
\label{thm:V87-results}
V87 proves:
\begin{enumerate}[label=\textup{(V87.\arabic*)},leftmargin=6.3em]
\item the exact tilted scalar-resolvent insertion identity;
\item the exact all-order ordered-chain formula
      \eqref{eq:V87-chain-formula} with precisely one global
      factorial;
\item a dimension-free complete column resolvent estimate;
\item closure of that column norm under every atomic-left chain
      operation; and
\item the exact \(m-2\)-incidence deficit created by taking norms
      before summing the permutation chains.
\end{enumerate}
\end{theorem}

\begin{proof}
Use
\Cref{lem:V87-tilted-resolvent,
thm:V87-chain-formula,
lem:V87-column-resolvent,
lem:V87-chain-step,
thm:V87-termwise-chain,
prop:V87-missing-incidences}.
\end{proof}

\begin{programme}[V88: symmetrized-chain defect extraction]
\label{prog:V87-V88}
\begin{enumerate}[leftmargin=3em]
\item Keep the full permutation sum in
      \eqref{eq:V87-chain-formula} and group chains by their first
      common Gaussian contraction graph.
\item Prove intrinsically that every group with fewer than
      \(m-2\) extra incidences cancels, reproducing the half-edge
      criterion of V85 without differentiating the Haar Jacobian.
\item Seek a recursive projection which subtracts the Gaussian
      pair component at each chain level before the column norm is
      applied.
\item Verify the projection at orders three and four against the
      exact odd and Wick-defect proofs of V82--V83.
\item Bound the projected score coefficients using
      \eqref{eq:V86-intrinsic-derivatives} and the one-link gap.
\item Do not claim an all-order card until the projected
      permutation sum, rather than each chain, is norm-controlled.
\end{enumerate}
\end{programme}

\begin{firewall}[Claim boundary after fourth-series V87]
\label{fire:V87-claim-boundary}
V87 proves an exact all-order algebraic chain formula and a
dimension-free termwise energy bound.  It does not prove the
symmetrized-chain cancellation norm, the all-order local card, the
anisotropic scalar optimization, or the complete all-order selected
gate.  It does not prove all-order MTA, WUFC, CPM, cutoff removal,
reflection positivity, Osterwalder--Schrader reconstruction,
construction of the continuum Yang--Mills measure, or a uniform
positive continuum spectral gap.  The full Yang--Mills mass gap has
not been proved.
\end{firewall}

\paragraph{Parent-version record.}
V87 was formed from
\texttt{Renewal\_Yang\_Mills\_Mass\_Gap\_Research\_Notes\_V86.tex}.
The V86 parent had \(37\,921\) logical source lines,
\(1\,280\,648\) bytes, and SHA--256
\[
 \texttt{c00046841bb9cb17ea45308c17522a391c12568b970c5e36e114c77334c5b9bd}.
\]
The next compulsory companion audit is V90.

% =============================================================================
% END APPENDED FOURTH-SERIES V87 MATERIAL
% =============================================================================

% =============================================================================
% BEGIN APPENDED FOURTH-SERIES V88 MATERIAL
% =============================================================================

\chapter{V88: Symmetry-first projected chains and a single
all-order certificate}
\label{chap:V88-projected-chains}

\section{Center every internal contraction}
\label{sec:V88-centering}

Define
\begin{equation}
 \mathcal D_\alpha(f,g)
 :=
 Q_\alpha\mathcal B_\alpha(f,g)
 =
 Q_\alpha\Gamma(f,R_\alpha g).
 \label{eq:V88-D}
\end{equation}
Since \(R_\alpha\) kills constants,
\begin{equation}
 \mathcal B_\alpha(f,g)
 =
 \mathcal B_\alpha(f,Q_\alpha g),
 \qquad
 \mathcal B_\alpha
 (f,\mathcal B_\alpha(h,g))
 =
 \mathcal B_\alpha
 (f,\mathcal D_\alpha(h,g)).
 \label{eq:V88-internal-centering}
\end{equation}

\begin{definition}[Ordered projected tail]
\label{def:V88-projected-tail}
For an ordered tuple of distinct labels define recursively
\begin{align}
 \mathcal P_\alpha(f_i)
 &:=
 Q_\alpha f_i,
 \label{eq:V88-P-one}\\
 \mathcal P_\alpha(f_{i_1},\ldots,f_{i_r})
 &:=
 \mathcal D_\alpha
 \left(
 f_{i_1},
 \mathcal P_\alpha(f_{i_2},\ldots,f_{i_r})
 \right)
 \qquad(r\ge2).
 \label{eq:V88-P-recursion}
\end{align}
Every projected tail is centered.
\end{definition}

\begin{lemma}[Projected form of an ordered chain]
\label{lem:V88-projected-chain}
For every ordered list of length \(m\ge2\),
\begin{equation}
 \mathcal C_\alpha(f_{i_1},\ldots,f_{i_m})
 =
 \mathcal B_\alpha
 \left(
 f_{i_1},
 \mathcal P_\alpha(f_{i_2},\ldots,f_{i_m})
 \right).
 \label{eq:V88-projected-chain}
\end{equation}
\end{lemma}

\begin{proof}
Start at the innermost occurrence of
\(\mathcal B_\alpha\) in
\eqref{eq:V87-chain}.  Whenever its output becomes the second
argument of the next operation, replace it by its centered part
using \eqref{eq:V88-internal-centering}.  Iteration gives exactly
\eqref{eq:V88-P-recursion}.
\end{proof}

\section{Permutation sums as subset recursion}
\label{sec:V88-subset}

For every nonempty finite label set \(I\), define the fully
symmetrized projected tail
\begin{equation}
 \mathfrak P_\alpha(I)
 :=
 \sum_{\pi\in S_I}
 \mathcal P_\alpha
 (f_{\pi(1)},\ldots,f_{\pi(|I|)}).
 \label{eq:V88-sym-tail}
\end{equation}

\begin{theorem}[Exact subset recursion]
\label{thm:V88-subset-recursion}
For a singleton,
\[
 \mathfrak P_\alpha(\{i\})=Q_\alpha f_i.
\]
For \(|I|\ge2\),
\begin{equation}
 \boxed{\qquad
 \mathfrak P_\alpha(I)
 =
 \sum_{i\in I}
 \mathcal D_\alpha
 \left(
 f_i,\mathfrak P_\alpha(I\setminus\{i\})
 \right).
 \qquad}
 \label{eq:V88-subset-recursion}
\end{equation}
Thus the complete permutation family is represented by one centered
object for each nonempty subset.
\end{theorem}

\begin{proof}
In \eqref{eq:V88-sym-tail}, group permutations by their first label
\(i\).  The remaining ordered labels run through every permutation
of \(I\setminus\{i\}\).  Linearity of
\(\mathcal D_\alpha(f_i,\cdot)\) gives
\eqref{eq:V88-subset-recursion}.
\end{proof}

\begin{theorem}[Symmetry-first cumulant compression]
\label{thm:V88-cumulant-compression}
For \(m\ge2\),
\begin{equation}
 \boxed{\qquad
 \kappa_m^{\nu_\alpha}(f_1,\ldots,f_m)
 =
 {1\over m}
 \sum_{i=1}^m
 \mathbb E_{\nu_\alpha}
 \mathcal B_\alpha
 \left(
 f_i,
 \mathfrak P_\alpha([m]\setminus\{i\})
 \right).
 \qquad}
 \label{eq:V88-cumulant-compression}
\end{equation}
\end{theorem}

\begin{proof}
For a fixed final label \(r\), the chain formula
\eqref{eq:V87-chain-formula} sums all \((m-1)!\) orderings of the
other labels and equals the symmetric cumulant.  Sum that identity
over the \(m\) choices of \(r\).  The resulting \(m!\) chains are
all ordered lists of \([m]\), each exactly once.  Now group them by
their first label and use
\Cref{lem:V88-projected-chain} and
\eqref{eq:V88-sym-tail}.  The left side is \(m\kappa_m\);
division by \(m\) proves the formula.
\end{proof}

\begin{remark}[Dynamic rather than factorial representation]
\label{rem:V88-dynamic}
Formula \eqref{eq:V88-cumulant-compression} is algebraically equal to
the \((m-1)!\)-chain formula, but the recursion
\eqref{eq:V88-subset-recursion} stores only \(2^m-1\) centered
objects.  More importantly, every permutation cancellation inside a
fixed subset occurs before a norm is taken.
\end{remark}

\section{The symmetrized projected-chain certificate}
\label{sec:V88-certificate}

For thermal representation inputs set
\[
 b_i=\sqrt{s_\alpha A_i}<1,
 \qquad
 B_I:=\sum_{i\in I}b_i.
\]

\begin{definition}[Symmetrized projected-chain certificate]
\label{def:V88-SPCC}
The SPCC holds if there is \(K<\infty\), independent of the nonempty
label set \(I\), \(\alpha\), representation dimensions, and
spectator legs, such that
\begin{equation}
 \boxed{\qquad
 \|\mathfrak P_\alpha(I)\|_{2,c}
 \le
 |I|!K^{|I|}
 \left(\prod_{i\in I}b_i\right)
 B_I^{|I|-1}.
 \qquad}
 \label{eq:V88-SPCC}
\end{equation}
\end{definition}

\begin{theorem}[SPCC implies the all-order local tree card]
\label{thm:V88-SPCC-implies-card}
If \eqref{eq:V88-SPCC} holds, then for every \(m\ge2\),
\begin{equation}
 \|\kappa_m^{\nu_\alpha}(D^{R_1},\ldots,D^{R_m})\|
 \le
 m!K_1^m
 \left(\prod_{i=1}^mb_i\right)
 \left(\sum_{i=1}^mb_i\right)^{m-2}.
 \label{eq:V88-Pruefer-bound}
\end{equation}
Consequently the all-thermal local tree card
\eqref{eq:V64-thermal-card} holds at every arity.
\end{theorem}

\begin{proof}
Use
\eqref{eq:V88-cumulant-compression}.  The atomic-left chain step
\eqref{eq:V87-chain-step} gives
\[
 \left\|
 \mathbb E\mathcal B_\alpha
 (D^{R_i},\mathfrak P_\alpha([m]\setminus\{i\}))
 \right\|
 \le
 Cb_i
 \|\mathfrak P_\alpha([m]\setminus\{i\})\|_{2,c},
\]
because every symmetrized tail is centered.  Apply the SPCC with
\(|I|=m-1\):
\begin{align*}
 \|\kappa_m\|
 &\le
 {(m-1)!K^{m-1}C\over m}
 \left(\prod_{j=1}^mb_j\right)
 \sum_{i=1}^m
 \left(\sum_{j\ne i}b_j\right)^{m-2}\\
 &\le
 (m-1)!K_1^m
 \left(\prod_jb_j\right)
 \left(\sum_jb_j\right)^{m-2}.
\end{align*}
This is stronger than the displayed \(m!\) normalization.

The Pruefer identity gives
\[
 \sum_{\tau\in\mathfrak T_m}
 \prod_i b_i^{d_i(\tau)}
 =
 \left(\prod_i b_i\right)
 \left(\sum_i b_i\right)^{m-2}.
\]
Split the complete cumulant proportionally among these positive tree
weights, as in V82--V85.  This proves
\eqref{eq:V64-thermal-card}.
\end{proof}

\begin{assessment}[A single analytic gate]
\label{ass:V88-single-gate}
The all-order local problem is now reduced to one explicit estimate,
\eqref{eq:V88-SPCC}, on the recursively defined centered objects
\(\mathfrak P_\alpha(I)\).  The certificate is stronger than the
local card and is not yet proved.  Its advantage is that it names
the exact sum in which permutation-chain cancellation must occur;
there is no remaining ambiguity about which terms may be grouped
before the norm.
\end{assessment}

\section{Why ordinary triangle iteration cannot prove the SPCC}
\label{sec:V88-triangle-no-go}

\begin{proposition}[Termwise subset recursion loses all Prüfer
powers]
\label{prop:V88-triangle-bound}
Using only the triangle inequality in
\eqref{eq:V88-subset-recursion} and the atomic-left estimate
\eqref{eq:V87-centered-chain-step} gives
\begin{equation}
 \|\mathfrak P_\alpha(I)\|_{2,c}
 \le
 |I|!C^{|I|}
 \prod_{i\in I}b_i.
 \label{eq:V88-triangle-bound}
\end{equation}
It supplies none of the additional
\(B_I^{|I|-1}\) in \eqref{eq:V88-SPCC}.
\end{proposition}

\begin{proof}
The singleton estimate is
\(\|Q_\alpha D^{R_i}\|_{2,c}\le Cb_i\).
Assume \eqref{eq:V88-triangle-bound} for sets of size \(r-1\).
Then
\begin{align*}
 \|\mathfrak P_\alpha(I)\|_{2,c}
 &\le
 \sum_{i\in I}
 Cb_i
 \|\mathfrak P_\alpha(I\setminus\{i\})\|_{2,c}\\
 &\le
 \sum_{i\in I}
 C^r(r-1)!
 b_i\prod_{j\ne i}b_j\\
 &=
 C^rr!\prod_{j\in I}b_j.
\end{align*}
This proves the claim.
\end{proof}

\begin{corollary}[Norms must be taken after subset symmetrization]
\label{cor:V88-after-symmetry}
No proof which applies
\eqref{eq:V87-centered-chain-step} separately to the summands in
\eqref{eq:V88-subset-recursion} can establish the SPCC in the
balanced small-\(b\) regime.
\end{corollary}

\begin{proof}
If all \(b_i=b\ll1\) and \(|I|=r\), the SPCC right side contains the
additional factor
\[
 (rb)^{r-1}.
\]
The termwise estimate \eqref{eq:V88-triangle-bound} has no positive
power beyond \(b^r\).  A constant independent of \(b\) cannot bridge
the difference.
\end{proof}

\section{Low-order relation}
\label{sec:V88-low-order}

\begin{proposition}[The first unproved projected estimate]
\label{prop:V88-first-gate}
The SPCC is automatic for \(|I|=1\).  Its first new case is
\begin{equation}
 \left\|
 \mathcal D_\alpha(D^{R_1},D^{R_2})
 +
 \mathcal D_\alpha(D^{R_2},D^{R_1})
 \right\|_{2,c}
 \le
 CK^2b_1b_2(b_1+b_2).
 \label{eq:V88-pair-SPCC}
\end{equation}
The ternary cumulant card of V82 bounds the pairing of this
symmetrized tail against a third representation through
\eqref{eq:V88-cumulant-compression}; it does not by itself prove the
full column norm \eqref{eq:V88-pair-SPCC}.
\end{proposition}

\begin{proof}
For a two-label set, the recursion is
\[
 \mathfrak P_\alpha(\{1,2\})
 =
 \mathcal D_\alpha(D^{R_1},D^{R_2})
 +
 \mathcal D_\alpha(D^{R_2},D^{R_1}).
\]
Substitution in \eqref{eq:V88-SPCC} gives
\eqref{eq:V88-pair-SPCC}.  The V82 card tests only expectations
obtained after an additional outer
\(\mathcal B_\alpha\); a column norm takes a supremum over a larger
class of Hilbert-valued tests, so the implication cannot be reversed
without a density or duality theorem.
\end{proof}

\begin{firewall}[The SPCC has not been inferred from fixed orders]
\label{fire:V88-SPCC-open}
V82--V85 prove the complete local cumulant card at arities three
through six.  They do not prove the stronger projected-tail norm,
even at two labels.  V88 therefore records
\eqref{eq:V88-pair-SPCC} as the next analytic acquisition rather than
claiming the all-order card by induction.
\end{firewall}

\section{V88 ledger and V89 programme}
\label{sec:V88-results}

\begin{theorem}[Fourth-series V88 result ledger]
\label{thm:V88-results}
V88 proves:
\begin{enumerate}[label=\textup{(V88.\arabic*)},leftmargin=6.3em]
\item exact centering of every non-outer intrinsic contraction;
\item the subset recursion for fully symmetrized projected tails;
\item the symmetry-first cumulant compression
      \eqref{eq:V88-cumulant-compression};
\item a single SPCC sufficient for the full all-order Pruefer tree
      bound; and
\item the exact failure of termwise triangle iteration to supply the
      \(m-2\) missing incidence powers.
\end{enumerate}
\end{theorem}

\begin{proof}
Use
\Cref{lem:V88-projected-chain,
thm:V88-subset-recursion,
thm:V88-cumulant-compression,
thm:V88-SPCC-implies-card,
prop:V88-triangle-bound}.
\end{proof}

\begin{programme}[V89: the two-label projected-tail estimate]
\label{prog:V88-V89}
\begin{enumerate}[leftmargin=3em]
\item Analyze the exact symmetrized pair
      \eqref{eq:V88-pair-SPCC}, not its two summands separately.
\item Use the product-resolvent identity
      \eqref{eq:V80-assembled-identity} to expose cancellation of
      the constant Gaussian contraction before the column norm.
\item Prove the bound first for central characters using the radial
      Green formula \eqref{eq:V81-Green-derivative}.
\item Test the cofinal Ornstein--Uhlenbeck limit; the right side is
      order one there, so only the balanced low-\(b\) regime can
      obstruct the estimate.
\item Extend from central characters to matrix coefficients using
      conjugation-sector blocks only after the radial estimate is
      secure.
\item If \eqref{eq:V88-pair-SPCC} fails for an \(L^2\) column test,
      weaken the SPCC to the exact outer representation test class
      in \eqref{eq:V88-cumulant-compression}.
\end{enumerate}
\end{programme}

\begin{firewall}[Claim boundary after fourth-series V88]
\label{fire:V88-claim-boundary}
V88 gives an exact symmetry-first reduction and a sufficient
all-order certificate, not a proof of that certificate.  It does not
prove the all-order local card, the anisotropic scalar optimization,
or the complete all-order selected gate.  It does not prove
all-order MTA, WUFC, CPM, cutoff removal, reflection positivity,
Osterwalder--Schrader reconstruction, construction of the continuum
Yang--Mills measure, or a uniform positive continuum spectral gap.
The full Yang--Mills mass gap has not been proved.
\end{firewall}

\paragraph{Parent-version record.}
V88 was formed from
\texttt{Renewal\_Yang\_Mills\_Mass\_Gap\_Research\_Notes\_V87.tex}.
The V87 parent had \(38\,417\) logical source lines,
\(1\,294\,661\) bytes, and SHA--256
\[
 \texttt{e9c4007eaef699700ceae42d82f7109a4c11b9983d794301f8fd9acdf8097e4d}.
\]
The next compulsory companion audit is V90.

% =============================================================================
% END APPENDED FOURTH-SERIES V88 MATERIAL
% =============================================================================

% =============================================================================
% BEGIN APPENDED FOURTH-SERIES V89 MATERIAL
% =============================================================================

\chapter{V89: The two-label SPCC in the central character sector}
\label{chap:V89-central-pair}

\section{Normalized characters start quadratically}
\label{sec:V89-character-quadratic}

For an irreducible representation \(R\), put
\[
 h_R(U):={\chi_R(U)\over d_R},
 \qquad
 b_R:=\sqrt{s_\alpha(1+C_2(R))}<1.
\]

\begin{lemma}[Vanishing of the character linear term]
\label{lem:V89-trace-zero}
For every \(Z\in\mathfrak{su}(2)\),
\begin{equation}
 \operatorname{tr}dR(Z)=0.
 \label{eq:V89-trace-zero}
\end{equation}
Consequently
\begin{equation}
 h_R(e^Z)-1
 =
 {1\over d_R}
 \operatorname{tr}
 \left(
 e^{dR(Z)}-I-dR(Z)
 \right).
 \label{eq:V89-character-remainder}
\end{equation}
\end{lemma}

\begin{proof}
The map
\[
 Z\longmapsto\operatorname{tr}dR(Z)
\]
is a one-dimensional representation of the Lie algebra because the
trace of a commutator is zero.  The Lie algebra
\(\mathfrak{su}(2)\) is perfect, so this functional vanishes.
Expanding \(D^R(e^Z)=e^{dR(Z)}\) and taking the normalized trace
proves \eqref{eq:V89-character-remainder}.
\end{proof}

\begin{theorem}[Quadratic centered character norm]
\label{thm:V89-character-centered}
Uniformly over all all-thermal irreducible representations,
\begin{equation}
 \boxed{\qquad
 \|Q_\alpha h_R\|_{L^2(\nu_\alpha)}
 \le Cb_R^2.
 \qquad}
 \label{eq:V89-character-centered}
\end{equation}
The estimate is dimension-free.
\end{theorem}

\begin{proof}
For the skew-adjoint operator \(B=dR(Z)\),
\[
 \|e^B-I-B\|\le {1\over2}\|B\|^2.
\]
The normalized trace has absolute value at most the operator norm,
and the Casimir bound gives
\[
 |h_R(e^Z)-1|
 \le {1\over2}C_2(R)|Z|^2.
\]
Therefore
\[
 \|h_R-1\|_{L^2}
 \le
 {1\over2}C_2(R)
 \left(\mathbb E|Z|^4\right)^{1/2}
 \le {CC_2(R)\over\alpha}
\]
by \Cref{lem:V82-log-moments}.  Centering around the best constant
cannot increase the scalar \(L^2\) norm, and
\[
 {C_2(R)\over\alpha}
 \le C s_\alpha(1+C_2(R))
 =Cb_R^2.
\]
For the trivial representation both sides vanish.
\end{proof}

\begin{lemma}[Atomic-left estimate for a normalized character]
\label{lem:V89-character-chain-step}
For every centered scalar or compatible Hilbert-valued \(G\),
\begin{equation}
 \|\mathcal D_\alpha(h_R,G)\|_{2,c}
 \le Cb_R\|Q_\alpha G\|_{2,c}.
 \label{eq:V89-character-chain-step}
\end{equation}
\end{lemma}

\begin{proof}
The normalized trace derivative satisfies
\[
 \|dh_R\|_\infty\le\sqrt{C_2(R)}.
\]
Repeat the proof of
\eqref{eq:V87-centered-chain-step}, using the column resolvent
estimate \eqref{eq:V87-column-resolvent}.  The outer centering costs
at most two.
\end{proof}

\section{The central pair certificate}
\label{sec:V89-pair}

\begin{theorem}[Two-label central SPCC]
\label{thm:V89-central-SPCC}
For any two all-thermal normalized irreducible characters,
\begin{equation}
 \boxed{\qquad
 \left\|
 \mathcal D_\alpha(h_{R_1},h_{R_2})
 +
 \mathcal D_\alpha(h_{R_2},h_{R_1})
 \right\|_{L^2(\nu_\alpha)}
 \le
 Cb_1b_2(b_1+b_2).
 \qquad}
 \label{eq:V89-central-SPCC}
\end{equation}
Thus the two-label case of the SPCC
\eqref{eq:V88-pair-SPCC} holds on the central character sector.
\end{theorem}

\begin{proof}
By
\Cref{lem:V89-character-chain-step,
thm:V89-character-centered},
\begin{align*}
 \|\mathcal D_\alpha(h_{R_1},h_{R_2})\|_2
 &\le Cb_1\|Q_\alpha h_{R_2}\|_2
 \le Cb_1b_2^2,\\
 \|\mathcal D_\alpha(h_{R_2},h_{R_1})\|_2
 &\le Cb_2\|Q_\alpha h_{R_1}\|_2
 \le Cb_2b_1^2.
\end{align*}
Add the inequalities and factor \(b_1b_2\).
\end{proof}

\begin{remark}[The proof is stronger than pair cancellation]
\label{rem:V89-individual}
In the central sector each ordered projected pair already has the
required bound.  Symmetrization is not needed because the normalized
trace removes the linear logarithmic-coordinate component at the
observable level.
\end{remark}

\section{Tangent Gaussian check}
\label{sec:V89-Gaussian-check}

Let
\[
 \mathcal A_0
 =
 -{1\over4}(\Delta-y\cdot\nabla),
 \qquad
 \Gamma_0(f,g)
 ={1\over4}\nabla f\cdot\nabla g
\]
on standard Gaussian \(\mathbb R^3\), and let
\(\mathcal R_0=\mathcal A_0^{-1}\mathcal Q_0\).

\begin{proposition}[Every linear tangent pair is exactly constant]
\label{prop:V89-Gaussian-linear}
For operator-valued linear forms
\[
 \ell_T(y)=\sum_aT_ay_a,
 \qquad
 \ell_S(y)=\sum_aS_ay_a,
\]
\begin{equation}
 \Gamma_0(\ell_T,\mathcal R_0\ell_S)
 =
 \sum_aT_aS_a.
 \label{eq:V89-Gaussian-linear}
\end{equation}
In particular its centered projection is zero.
\end{proposition}

\begin{proof}
Since
\(\mathcal A_0y_a=y_a/4\),
\[
 \mathcal R_0\ell_S=4\ell_S.
\]
Substitution in \(\Gamma_0\) proves
\eqref{eq:V89-Gaussian-linear}.  Its right side is independent of
\(y\).
\end{proof}

\begin{assessment}[Meaning for noncentral matrix coefficients]
\label{ass:V89-noncentral}
A noncentral matrix coefficient generally has a nonzero linear
logarithmic term and only
\(\|QD^R\|_{2,c}\lesssim b_R\).  Nevertheless
\Cref{prop:V89-Gaussian-linear} shows that the leading tangent
contraction is constant and is removed by
\(\mathcal D_\alpha=Q_\alpha\mathcal B_\alpha\).
Thus the missing central proof does not signal a counterexample; it
identifies the estimate needed for the Wilson-versus-tangent
resolvent defect on matrix coefficients.
\end{assessment}

\section{A precise noncentral defect operator}
\label{sec:V89-defect}

For representation linearizations
\[
 \ell_R(y):=dR(y),
\]
define the rescaled projected-pair defect
\begin{equation}
 \mathfrak E_{\alpha}(R,S)
 :=
 Q_\alpha
 \Gamma
 \left(
 D^R,
 R_\alpha D^S
 \right)
 -
 \mathcal U_\alpha
 \mathcal Q_0
 \Gamma_0
 \left(
 \ell_R,\mathcal R_0\ell_S
 \right).
 \label{eq:V89-defect}
\end{equation}
Here \(\mathcal U_\alpha\) denotes pullback from
\(y=\sqrt\alpha Z\) with the representation generators scaled by
\(\alpha^{-1/2}\).  By
\Cref{prop:V89-Gaussian-linear}, the second term is zero; it is
retained in the definition to display the canceled tangent symbol.

\begin{definition}[Matrix projected-pair defect estimate]
\label{def:V89-MPDE}
The MPDE is the bound
\begin{equation}
 \|\mathfrak E_\alpha(R,S)\|_{2,c}
 \le
 Cb_Rb_S(b_R+b_S)
 \label{eq:V89-MPDE}
\end{equation}
for all all-thermal \(R,S\), uniformly in their dimensions.
\end{definition}

\begin{corollary}[MPDE closes the full two-label SPCC]
\label{cor:V89-MPDE-SPCC}
If the MPDE holds, then
\eqref{eq:V88-pair-SPCC} holds for arbitrary unitary representation
matrices.
\end{corollary}

\begin{proof}
Apply \eqref{eq:V89-MPDE} to \((R,S)\) and \((S,R)\), use the
vanishing tangent term, and add the two bounds.
\end{proof}

\begin{firewall}[No unproved pullback comparison]
\label{fire:V89-MPDE-open}
The notation \(\mathcal U_\alpha\) in
\eqref{eq:V89-defect} records the tangent normalization; V89 does
not assert a global transport map or compare the two measures by the
disproved SEDE.  The MPDE remains open outside the central character
sector and must be proved intrinsically or by a fixed two-input
comparison which does not require all-order Haar derivatives.
\end{firewall}

\section{V89 ledger and V90 programme}
\label{sec:V89-results}

\begin{theorem}[Fourth-series V89 result ledger]
\label{thm:V89-results}
V89 proves:
\begin{enumerate}[label=\textup{(V89.\arabic*)},leftmargin=6.3em]
\item exact absence of a linear logarithmic term in every normalized
      \(SU(2)\) character;
\item the dimension-free centered bound
      \(\|Qh_R\|_2\le Cb_R^2\);
\item the two-label SPCC in the complete central character sector;
\item exact constancy of every linear tangent-Gaussian projected
      pair; and
\item reduction of the noncentral two-label SPCC to the explicit
      matrix projected-pair defect estimate.
\end{enumerate}
\end{theorem}

\begin{proof}
Use
\Cref{lem:V89-trace-zero,
thm:V89-character-centered,
thm:V89-central-SPCC,
prop:V89-Gaussian-linear,
cor:V89-MPDE-SPCC}.
\end{proof}

\begin{programme}[V90: audit and the matrix projected-pair defect]
\label{prog:V89-V90}
\begin{enumerate}[leftmargin=3em]
\item Perform the compulsory read-only audit of the newest active
      SM--Einstein and Navier--Stokes notes.
\item Rewrite \(\mathfrak E_\alpha(R,S)\) through the intrinsic
      product-resolvent identity, keeping the constant tangent symbol
      subtracted before the column norm.
\item Use the exact one-form intertwining to compare only a
      two-input resolvent block; do not invoke the all-order radial
      transport.
\item Prove the MPDE first when
      \(b_R+b_S\le b_0\), where perturbation of the tangent generator
      has a fixed small parameter.
\item Use the crude chain bound when
      \(b_R+b_S>b_0\), where the desired extra factor is bounded
      below.
\item Check all estimates in the column norm before any
      representation trace or bank sum, as required by the V85
      audit.
\end{enumerate}
\end{programme}

\begin{firewall}[Claim boundary after fourth-series V89]
\label{fire:V89-claim-boundary}
V89 proves the first SPCC case only for normalized central
characters.  It does not prove the matrix MPDE, the higher-subset
SPCC, the all-order local card, the anisotropic scalar optimization,
or the complete all-order selected gate.  It does not prove
all-order MTA, WUFC, CPM, cutoff removal, reflection positivity,
Osterwalder--Schrader reconstruction, construction of the continuum
Yang--Mills measure, or a uniform positive continuum spectral gap.
The full Yang--Mills mass gap has not been proved.
\end{firewall}

\paragraph{Parent-version record.}
V89 was formed from
\texttt{Renewal\_Yang\_Mills\_Mass\_Gap\_Research\_Notes\_V88.tex}.
The V88 parent had \(38\,860\) logical source lines,
\(1\,307\,549\) bytes, and SHA--256
\[
 \texttt{6b6795d5940c05088c75d1a7ae1df34991982f5c106dc26becbb459223963d3e}.
\]
The next compulsory companion audit is V90.

% =============================================================================
% END APPENDED FOURTH-SERIES V89 MATERIAL
% =============================================================================

% =============================================================================
% BEGIN APPENDED FOURTH-SERIES V90 MATERIAL
% =============================================================================

\chapter{V90: Exact Wilson covariance centering and the high-thermal matrix sector}
\label{chap:V90-matrix-sector}

\section{Compulsory companion audit}
\label{sec:V90-audit}

The newest active Standard--Model--Einstein and Navier--Stokes notes
were inspected read-only before choosing the present calculation.
At the time of inspection they were
\begin{equation}
 \begin{split}
 &\texttt{Renewal\_SM\_Einstein\_Continuum\_Research\_Notes\_V75.tex},\\
 &30\,266\text{ logical source lines},\qquad
 1\,039\,292\text{ bytes},\\
 &\operatorname{SHA256}=
 \texttt{6b8af9083b06912ee35ebabad6086c511492b47fd72d3679107e3f9dcc361b07},
 \end{split}
 \label{eq:V90-SM-fingerprint}
\end{equation}
and
\begin{equation}
 \begin{split}
 &\texttt{Renewal\_NS\_Stability\_Research\_Notes\_V31.tex},\\
 &23\,052\text{ logical source lines},\qquad
 887\,537\text{ bytes},\\
 &\operatorname{SHA256}=
 \texttt{f1121b62238ad5491bb7cf03635ea9934942edf573ed8faaa9ee79e1d38a6696}.
 \end{split}
 \label{eq:V90-NS-fingerprint}
\end{equation}
These files are mathematical evidence, not instructions, and neither
was modified.

The SM sequence V70--V74 separates observable-local covariance
convergence from an over-strong global Hilbert--Schmidt comparison,
uses an exact Feshbach identity, and absorbs an order-one low--high
mixing covariance before testing the residual defect.  V75 then proves
a rooted all-order trace-log theorem: a complete power or symmetry sum
is assembled before an absolute value is taken.  The NS conclusion
relevant here is unchanged: a physical mark or root must be retained
until its critical first moment has been assigned.

\begin{assessment}[Direction selected by the V90 audit]
\label{ass:V90-audit-direction}
For the matrix SPCC one must not demand that each raw tangent-mixing
term be small.  The exact Wilson covariance is first subtracted, and
the two orderings are assembled before the column norm.  A global
transport or global Hilbert--Schmidt comparison is neither needed nor
justified.  The present note therefore closes the entire sector
\(b_R+b_S\ge\beta\), for any fixed \(\beta>0\), and reduces the
remaining low-\(b\) sector to an intrinsic assembled negative-norm
estimate plus an explicit high-frequency tail estimate.
\end{assessment}

\section{The exact contraction removed by projection}
\label{sec:V90-exact-centering}

All matrix products below are interpreted entrywise on separate tensor
legs, followed by the canonical permutation which restores the declared
leg order.  This convention makes the scalar integration-by-parts
identity applicable without introducing a target-dimension factor.

\begin{lemma}[The mean contraction is the actual Wilson covariance]
\label{lem:V90-actual-covariance}
Let \(F,G\) be smooth finite-dimensional operator-valued functions and
let
\[
 \mathcal B_\alpha(F,G)
 :=
 \Gamma(F,R_\alpha G).
\]
Then, entrywise,
\begin{equation}
 \mathbb E_{\nu_\alpha}\mathcal B_\alpha(F,G)
 =
 \mathbb E_{\nu_\alpha}
 \bigl\{(Q_\alpha F)(Q_\alpha G)\bigr\}.
 \label{eq:V90-actual-covariance}
\end{equation}
Consequently
\begin{equation}
 \mathcal D_\alpha(F,G)
 =
 \mathcal B_\alpha(F,G)
 -
 \mathbb E_{\nu_\alpha}\mathcal B_\alpha(F,G)
 \label{eq:V90-exact-subtraction}
\end{equation}
subtracts the exact Wilson covariance, not a Gaussian approximation to
it.
\end{lemma}

\begin{proof}
Put \(u=R_\alpha G\).  The Dirichlet identity for the self-adjoint
weighted generator gives
\[
 \mathbb E_{\nu_\alpha}\Gamma(F,u)
 =
 \mathbb E_{\nu_\alpha}F\,A_\alpha u.
\]
Since \(A_\alpha R_\alpha G=Q_\alpha G\), the right side is
\(\mathbb E FQ_\alpha G\).  Constants in \(F\) annihilate
\(Q_\alpha G\), so it equals
\(\mathbb E (Q_\alpha F)(Q_\alpha G)\).  This proves
\eqref{eq:V90-actual-covariance}.  The definition
\(\mathcal D_\alpha=Q_\alpha\mathcal B_\alpha\) gives
\eqref{eq:V90-exact-subtraction}.
\end{proof}

\begin{corollary}[The displayed tangent subtraction in V89 is
identically zero]
\label{cor:V90-tangent-subtraction-zero}
For representation matrices \(D^R,D^S\), the V89 defect satisfies
\begin{equation}
 \mathfrak E_\alpha(R,S)
 =
 \mathcal D_\alpha(D^R,D^S).
 \label{eq:V90-defect-is-D}
\end{equation}
Thus no global pullback between Gaussian and Wilson \(L^2\) spaces is
part of the MPDE.
\end{corollary}

\begin{proof}
By Proposition~\ref{prop:V89-Gaussian-linear},
\(\Gamma_0(\ell_R,\mathcal R_0\ell_S)\) is constant.  Its centered
projection \(\mathcal Q_0\) is therefore zero.  Substitute this fact
in \eqref{eq:V89-defect}; the remaining term is precisely
\(Q_\alpha\mathcal B_\alpha(D^R,D^S)\).
\end{proof}

\begin{remark}[Exact centering is the covariance analogue of
whitening]
\label{rem:V90-whitening}
The leading constant need not be compared with a guessed tangent
constant.  Equation~\eqref{eq:V90-exact-subtraction} removes the
finite-\(\alpha\) Wilson contraction itself.  This is the precise
analogue of the companion audit's instruction to absorb the actual
order-one covariance before estimating a remainder.
\end{remark}

\section{Complete closure of the high-\(b\) pair sector}
\label{sec:V90-high-b}

Fix once and for all a threshold
\[
 0<\beta<1.
\]
It will eventually be chosen small enough for a local tangent
perturbation.  No smallness is used in the following theorem.

\begin{theorem}[High-thermal matrix MPDE]
\label{thm:V90-high-b-MPDE}
Let \(R,S\) be arbitrary all-thermal unitary representations, so
\[
 b_R,b_S<1.
\]
If \(b_R+b_S\ge\beta\), then
\begin{equation}
 \boxed{\qquad
 \|\mathfrak E_\alpha(R,S)\|_{2,c}
 \le
 {C\over\beta}
 b_Rb_S(b_R+b_S).
 \qquad}
 \label{eq:V90-high-b-MPDE}
\end{equation}
The constant is independent of \(\alpha\), of the representation
dimensions, and of the target matrix size.
\end{theorem}

\begin{proof}
Corollary~\ref{cor:V90-tangent-subtraction-zero} and the atomic-left
column estimate \eqref{eq:V87-centered-chain-step} give
\[
 \|\mathfrak E_\alpha(R,S)\|_{2,c}
 =
 \|\mathcal D_\alpha(D^R,D^S)\|_{2,c}
 \le
 Cb_R\|Q_\alpha D^S\|_{2,c}.
\]
The centered representation estimate inherited in the proof of
Theorem~\ref{thm:V87-termwise-chain} yields
\[
 \|Q_\alpha D^S\|_{2,c}\le Cb_S.
\]
Hence the left side is at most \(Cb_Rb_S\).  The hypothesis implies
\[
 1\le\beta^{-1}(b_R+b_S),
\]
which proves \eqref{eq:V90-high-b-MPDE}.  Every estimate was taken in
column norm before a representation trace or bank sum.
\end{proof}

\begin{corollary}[High-thermal two-label SPCC]
\label{cor:V90-high-b-SPCC}
Under the hypotheses of Theorem~\ref{thm:V90-high-b-MPDE},
\begin{equation}
 \left\|
 \mathcal D_\alpha(D^R,D^S)
 +
 \mathcal D_\alpha(D^S,D^R)
 \right\|_{2,c}
 \le
 {C\over\beta}b_Rb_S(b_R+b_S).
 \label{eq:V90-high-b-SPCC}
\end{equation}
\end{corollary}

\begin{proof}
Apply Theorem~\ref{thm:V90-high-b-MPDE} to the two ordered pairs and
use the triangle inequality only after both exact Wilson covariance
subtractions have been made.
\end{proof}

\begin{corollary}[Exact location of the unresolved pair sector]
\label{cor:V90-pair-location}
The two-label SPCC is now proved
\begin{enumerate}
\item for every pair of normalized central characters, at every
      all-thermal scale, by Theorem~\ref{thm:V89-central-SPCC}; and
\item for arbitrary unitary representation matrices whenever
      \(b_R+b_S\ge\beta\), by
      Corollary~\ref{cor:V90-high-b-SPCC}.
\end{enumerate}
The only remaining two-label case is therefore
\begin{equation}
 b_R+b_S<\beta
 \quad\hbox{with a genuinely noncentral matrix component.}
 \label{eq:V90-only-low-b}
\end{equation}
\end{corollary}

\begin{proof}
The two cited results cover the stated sectors.  Their complement in
the all-thermal matrix problem is exactly
\eqref{eq:V90-only-low-b}.
\end{proof}

\section{The product identity lives first in a negative norm}
\label{sec:V90-negative-norm}

The product-resolvent identity of V80 supplies a comparison target
which lives entirely in the Wilson probability space.  Define
\begin{equation}
 \Theta_\alpha(F,G)
 :=
 {1\over2}dR_\alpha
 \left\{
 FQ_\alpha G+(R_\alpha G)A_\alpha F
 \right\}
 -
 {1\over2}d\{F R_\alpha G\}.
 \label{eq:V90-Theta}
\end{equation}
The three terms in this definition are an assembled expression; they
are not three independent error terms.

For a centered finite-dimensional operator-valued function \(H\),
define the complete negative norm
\begin{equation}
 \|H\|_{-1,c}
 :=
 \|dR_\alpha H\|_{2,c}.
 \label{eq:V90-negative-norm}
\end{equation}

\begin{lemma}[Exact assembled negative-norm identity]
\label{lem:V90-negative-identity}
For arbitrary smooth operator-valued \(F,G\),
\begin{equation}
 \boxed{\qquad
 \left\|
 \mathcal D_\alpha(F,G)+\mathcal D_\alpha(G,F)
 \right\|_{-1,c}
 =
 \left\|
 \Theta_\alpha(F,G)+\Theta_\alpha(G,F)
 \right\|_{2,c}.
 \qquad}
 \label{eq:V90-negative-identity}
\end{equation}
All swapped tensor legs are canonically restored before the sums are
formed.
\end{lemma}

\begin{proof}
The assembled product-resolvent identity
\eqref{eq:V80-assembled-identity}, applied entrywise, is
\[
 dR_\alpha\mathcal B_\alpha(F,G)=\Theta_\alpha(F,G).
\]
Because \(R_\alpha Q_\alpha=R_\alpha\) and
\(\mathcal D_\alpha=Q_\alpha\mathcal B_\alpha\),
\[
 dR_\alpha\mathcal D_\alpha(F,G)
 =
 dR_\alpha\mathcal B_\alpha(F,G).
\]
Add the two orderings and use definition
\eqref{eq:V90-negative-norm}.
\end{proof}

\begin{lemma}[The spectral gap has only the weak-to-strong direction]
\label{lem:V90-gap-direction}
For every centered operator-valued \(H\),
\begin{equation}
 \|H\|_{-1,c}
 \le
 {C\over\sqrt\alpha}\|H\|_{2,c}.
 \label{eq:V90-gap-direction}
\end{equation}
There is no converse with a constant independent of the spectral
frequency.
\end{lemma}

\begin{proof}
Fix a unit input vector \(v\).  The spectral theorem and the Wilson
gap \(A_\alpha\ge c\alpha\) on mean-zero functions give
\[
 \|dR_\alpha Hv\|_2^2
 =
 \langle R_\alpha Hv,Hv\rangle
 \le
 {C\over\alpha}\|Hv\|_2^2.
\]
Take the supremum over \(v\).  Conversely, on an
\(A_\alpha\)-eigenfunction of eigenvalue \(\lambda\), the ratio
\(\|H\|_2/\|H\|_{-1}\) is \(\sqrt\lambda\), and the spectrum is
unbounded.
\end{proof}

\begin{assessment}[A necessary correction to the V89 programme]
\label{ass:V90-negative-not-L2}
The product-resolvent identity controls
\(dR_\alpha\mathcal D_\alpha\), not
\(d\mathcal D_\alpha\).  It therefore supplies a negative-norm
cancellation and cannot, by the spectral gap alone, prove the
\(L^2\)-column SPCC.  Treating
\eqref{eq:V90-gap-direction} in the reverse direction would be a
false elliptic estimate.  A separate high-frequency tail bound is
required.
\end{assessment}

\section{A rigorous two-part low-\(b\) certificate}
\label{sec:V90-two-part-certificate}

Let \(P_{\le L\alpha}\) and \(P_{>L\alpha}\) denote the spectral
projections of \(A_\alpha\) on the mean-zero scalar space, amplified
entrywise to every finite matrix target.

\begin{proposition}[Spectral upgrade with an explicit tail]
\label{prop:V90-spectral-upgrade}
For every centered operator-valued \(H\) and every \(L>0\),
\begin{equation}
 \|H\|_{2,c}
 \le
 \sqrt{L\alpha}\,\|H\|_{-1,c}
 +
 \|P_{>L\alpha}H\|_{2,c}.
 \label{eq:V90-spectral-upgrade}
\end{equation}
\end{proposition}

\begin{proof}
For a unit vector \(v\), the spectral theorem gives
\[
 \|P_{\le L\alpha}Hv\|_2^2
 \le
 L\alpha\,
 \langle R_\alpha Hv,Hv\rangle
 =
 L\alpha\,\|dR_\alpha Hv\|_2^2.
\]
Take the supremum over \(v\), and then use the triangle inequality
for
\(H=P_{\le L\alpha}H+P_{>L\alpha}H\).
\end{proof}

For a representation pair write
\begin{equation}
 H_{R,S}
 :=
 \mathcal D_\alpha(D^R,D^S)
 +
 \mathcal D_\alpha(D^S,D^R).
 \label{eq:V90-pair-tail}
\end{equation}

\begin{definition}[Low-\(b\) spectral pair certificate]
\label{def:V90-LSPC}
The LSPC at thresholds \((\beta,L)\) consists of the two estimates
\begin{align}
 \left\|
 \Theta_\alpha(D^R,D^S)
 +
 \Theta_\alpha(D^S,D^R)
 \right\|_{2,c}
 &\le
 {C\over\sqrt\alpha}\,
 b_Rb_S(b_R+b_S),
 \label{eq:V90-LSPC-negative}\\
 \|P_{>L\alpha}H_{R,S}\|_{2,c}
 &\le
 Cb_Rb_S(b_R+b_S)
 \label{eq:V90-LSPC-tail}
\end{align}
whenever \(b_R+b_S<\beta\).  The constants and the fixed \(L\)
must be independent of \(\alpha,R,S\) and the matrix dimensions.
\end{definition}

\begin{corollary}[The LSPC suffices for the complete pair SPCC]
\label{cor:V90-LSPC-sufficient}
If the LSPC holds for some fixed
\(\beta\in(0,1)\) and \(L>0\), then the two-label SPCC
\eqref{eq:V88-pair-SPCC} holds for every all-thermal pair of unitary
representation matrices.
\end{corollary}

\begin{proof}
In the low-\(b\) sector, Lemma
\ref{lem:V90-negative-identity} and
\eqref{eq:V90-LSPC-negative} give
\[
 \|H_{R,S}\|_{-1,c}
 \le
 {C\over\sqrt\alpha}b_Rb_S(b_R+b_S).
\]
Insert this and \eqref{eq:V90-LSPC-tail} in
\eqref{eq:V90-spectral-upgrade}.  The result is
\[
 \|H_{R,S}\|_{2,c}
 \le
 C(1+\sqrt L)b_Rb_S(b_R+b_S).
\]
For \(b_R+b_S\ge\beta\), use
Corollary~\ref{cor:V90-high-b-SPCC}.  Since \(\beta,L\) are fixed,
enlarging the constant joins the two regimes.
\end{proof}

\begin{assessment}[Why the remaining certificate is local and
two-input]
\label{ass:V90-narrower-gate}
The V89 notation suggested comparing a Wilson contraction with a
Gaussian contraction through a pullback.  Corollary
\ref{cor:V90-tangent-subtraction-zero} removes that unnecessary
global comparison.  The first LSPC line asks for the assembled
two-input tangent cancellation in its natural negative norm.  The
second line asks only that this particular pair have no uncontrolled
spectral mass above \(L\alpha\).  Neither line requires an all-order
density transport or the disproved exponential majorant for all Haar
logarithmic derivatives.
\end{assessment}

\begin{firewall}[The LSPC is not proved in V90]
\label{fire:V90-LSPC-open}
Lemma~\ref{lem:V90-negative-identity} and
Proposition~\ref{prop:V90-spectral-upgrade} are exact proved
reductions, not proofs of
\eqref{eq:V90-LSPC-negative} or
\eqref{eq:V90-LSPC-tail}.  Estimating the three terms of
\(\Theta_\alpha\) separately discards the tangent cancellation.
Ignoring the second line would reverse
\eqref{eq:V90-gap-direction}.  V90 therefore does not claim the
low-\(b\) matrix pair, the higher-subset SPCC, or the all-order local
card.
\end{firewall}

\section{V90 ledger and V91 programme}
\label{sec:V90-results}

\begin{theorem}[Fourth-series V90 result ledger]
\label{thm:V90-results}
V90:
\begin{enumerate}[label=\textup{(V90.\arabic*)},leftmargin=6.3em]
\item completes the compulsory read-only SM/NS audit;
\item proves that \(\mathcal D_\alpha\) subtracts the exact
      finite-\(\alpha\) Wilson covariance;
\item removes the unnecessary global Gaussian pullback from the
      matrix defect;
\item proves the MPDE and the two-label SPCC for all arbitrary matrix
      pairs with \(b_R+b_S\ge\beta\);
\item localizes the only unresolved pair sector to low-\(b\),
      genuinely noncentral components; and
\item identifies the exact negative-norm content of the assembled
      product identity and proves the spectral-tail upgrade
      \eqref{eq:V90-spectral-upgrade}.
\end{enumerate}
\end{theorem}

\begin{proof}
Use
\Cref{lem:V90-actual-covariance,
cor:V90-tangent-subtraction-zero,
thm:V90-high-b-MPDE,
cor:V90-high-b-SPCC,
cor:V90-pair-location,
lem:V90-negative-identity,
prop:V90-spectral-upgrade}.
\end{proof}

\begin{programme}[V91: the low-\(b\) negative norm and spectral tail]
\label{prog:V90-V91}
\begin{enumerate}[leftmargin=3em]
\item Split the principal logarithmic chart into an inner ball and
      its exponentially small cut-locus tail.
\item On the inner ball, rescale \(y=\sqrt\alpha Z\) and expand only
      the fixed two-input assembled form \(\Theta_\alpha\), never an
      all-order Haar density.
\item Subtract the exactly constant tangent contraction before any
      column norm, retain the two matrix orderings together, and prove
      \eqref{eq:V90-LSPC-negative}.
\item Establish \eqref{eq:V90-LSPC-tail} independently, using a
      spectral multiplier or a fixed-order elliptic estimate; do not
      reverse the spectral-gap inequality.
\item Bound the cut-locus contribution using unitary representation
      bounds,
      Wilson concentration, and the spectral gap, without
      differentiating the logarithmic Haar Jacobian to growing
      order.
\item Complete both lines of the LSPC for a sufficiently small fixed
      \(\beta\) and a fixed \(L\).  If the normal-chart remainder is not uniform at the
      parabolic edge, replace it by the exact conjugation-sector
      Jacobi block rather than weakening the claim silently.
\item Only after the pair certificate is proved, test whether the
      subset recursion of V88 propagates the same assembled
      cancellation to three labels.
\end{enumerate}
\end{programme}

\begin{firewall}[Claim boundary after fourth-series V90]
\label{fire:V90-claim-boundary}
V90 proves exact Wilson covariance centering, the full high-\(b\)
matrix pair sector, and an intrinsic reduction of the remaining
low-\(b\) pair to the two-part LSPC.  It does not prove the LSPC, the complete
two-label matrix SPCC, the higher-subset SPCC, the all-order local
card, the anisotropic scalar optimization, or the complete all-order
selected gate.  It does not prove all-order MTA, WUFC, CPM, cutoff
removal, reflection positivity, Osterwalder--Schrader reconstruction,
construction of the continuum Yang--Mills measure, or a uniform
positive continuum spectral gap.  The full Yang--Mills mass gap has
not been proved.
\end{firewall}

\paragraph{Parent-version record.}
V90 was formed from
\texttt{Renewal\_Yang\_Mills\_Mass\_Gap\_Research\_Notes\_V89.tex}.
The V89 parent had \(39\,216\) logical source lines,
\(1\,317\,655\) bytes, and SHA--256
\[
 \texttt{56ead753b4e425fcbb2975bb3522cb580a45448720c32a9c1937e4b88bb2c3bf}.
\]
The next compulsory companion audit is V95.

% =============================================================================
% END APPENDED FOURTH-SERIES V90 MATERIAL
% =============================================================================

% =============================================================================
% BEGIN APPENDED FOURTH-SERIES V91 MATERIAL
% =============================================================================

\chapter{V91: The terminal negative-norm certificate for the all-order tree card}
\label{chap:V91-terminal-negative}

\section{Why the terminal norm may be weakened}
\label{sec:V91-terminal-norm}

V88 imposed an \(L^2\)-column estimate on every symmetrized
projected tail.  V90 showed that the product-resolvent identity
naturally controls the weaker norm
\[
 \|H\|_{-1,c}=\|dR_\alpha H\|_{2,c}.
\]
The apparent mismatch disappears at the terminal edge of the exact
cumulant compression: the outer representation derivative pairs
directly with \(dR_\alpha H\).  Thus the all-order local card does not
require an \(L^2\) upgrade of the final tail.

Retain the notation
\[
 b_i=\sqrt{s_\alpha A_i}<1,
 \qquad
 B_I=\sum_{i\in I}b_i
\]
from V88.

\begin{lemma}[Outer representation pairing uses only the negative
norm]
\label{lem:V91-outer-negative}
Let \(D^R\) be an all-thermal unitary representation and let \(G\)
be a centered finite-dimensional operator-valued function on the
remaining tensor legs.  Then
\begin{equation}
 \left\|
 \mathbb E_{\nu_\alpha}
 \mathcal B_\alpha(D^R,G)
 \right\|_{\rm op}
 \le
 \sqrt{C_2(R)}\,\|G\|_{-1,c}
 \le
 C\sqrt\alpha\,b_R\|G\|_{-1,c}.
 \label{eq:V91-outer-negative}
\end{equation}
The constant is independent of the representation and target
dimensions.
\end{lemma}

\begin{proof}
By definition,
\[
 \mathbb E\mathcal B_\alpha(D^R,G)
 =
 \mathbb E\sum_a
 (X_aD^R)(X_aR_\alpha G).
\]
The unitary Casimir row identity is
\[
 \sum_a(X_aD^R)^*(X_aD^R)=C_2(R)I.
\]
Apply the row--column Cauchy--Schwarz inequality before taking any
trace or bank sum.  This gives the first inequality in
\eqref{eq:V91-outer-negative}.  Since
\(s_\alpha\simeq\alpha^{-1}\) in the inherited Wilson range and
\(b_R^2=s_\alpha(1+C_2(R))\),
\[
 \sqrt{C_2(R)}
 \le C\sqrt\alpha\,b_R,
\]
which proves the second inequality.
\end{proof}

\section{The resolvent-norm symmetrized projected-chain certificate}
\label{sec:V91-RSPCC}

\begin{definition}[RSPCC]
\label{def:V91-RSPCC}
The resolvent-norm symmetrized projected-chain certificate holds if
there is \(K<\infty\), independent of the nonempty label set \(I\),
\(\alpha\), all representation dimensions, and spectator legs, such
that
\begin{equation}
 \boxed{\qquad
 \|\mathfrak P_\alpha(I)\|_{-1,c}
 \le
 {|I|!K^{|I|}\over\sqrt\alpha}
 \left(\prod_{i\in I}b_i\right)
 B_I^{|I|-1}.
 \qquad}
 \label{eq:V91-RSPCC}
\end{equation}
\end{definition}

\begin{theorem}[RSPCC implies the all-order local tree card]
\label{thm:V91-RSPCC-implies-card}
If \eqref{eq:V91-RSPCC} holds, then for every \(m\ge2\),
\begin{equation}
 \boxed{\qquad
 \|\kappa_m^{\nu_\alpha}
 (D^{R_1},\ldots,D^{R_m})\|
 \le
 m!K_1^m
 \left(\prod_{i=1}^mb_i\right)
 \left(\sum_{i=1}^mb_i\right)^{m-2}.
 \qquad}
 \label{eq:V91-Pruefer-bound}
\end{equation}
Consequently the all-thermal local tree card
\eqref{eq:V64-thermal-card} holds at every arity.
\end{theorem}

\begin{proof}
Use the exact symmetry-first compression
\eqref{eq:V88-cumulant-compression}:
\[
 \kappa_m
 =
 {1\over m}\sum_{i=1}^m
 \mathbb E\mathcal B_\alpha
 \left(
 D^{R_i},
 \mathfrak P_\alpha([m]\setminus\{i\})
 \right).
\]
Put \(I_i=[m]\setminus\{i\}\).  Lemma
\ref{lem:V91-outer-negative} and the RSPCC with
\(|I_i|=m-1\) give
\begin{align*}
 \|\kappa_m\|
 &\le
 {C\sqrt\alpha\over m}
 \sum_{i=1}^m
 b_i\,
 {(m-1)!K^{m-1}\over\sqrt\alpha}
 \left(\prod_{j\ne i}b_j\right)
 B_{I_i}^{m-2}\\
 &=
 {C(m-1)!K^{m-1}\over m}
 \left(\prod_{j=1}^mb_j\right)
 \sum_{i=1}^mB_{I_i}^{m-2}\\
 &\le
 C(m-1)!K^{m-1}
 \left(\prod_jb_j\right)
 \left(\sum_jb_j\right)^{m-2}.
\end{align*}
Absorb the fixed leading constant into \(K_1^m\) and enlarge
\((m-1)!\) to \(m!\).

Finally, the Pruefer identity
\[
 \sum_{\tau\in\mathfrak T_m}
 \prod_i b_i^{d_i(\tau)}
 =
 \left(\prod_i b_i\right)
 \left(\sum_i b_i\right)^{m-2}
\]
identifies the last factor with the sum of the labelled-tree
weights.  Split the complete cumulant proportionally among these
positive weights exactly as in the proof of
Theorem~\ref{thm:V88-SPCC-implies-card}.
\end{proof}

\begin{proposition}[The RSPCC is a weaker proof obligation than the former
proof obligation]
\label{prop:V91-weaker}
The \(L^2\) SPCC \eqref{eq:V88-SPCC} implies the RSPCC
\eqref{eq:V91-RSPCC}.  No frequency-independent converse follows
from the Wilson spectral gap.
\end{proposition}

\begin{proof}
Every projected tail is centered.  Apply
\eqref{eq:V90-gap-direction} to the V88 bound:
\[
 \|\mathfrak P_\alpha(I)\|_{-1,c}
 \le
 {C\over\sqrt\alpha}
 \|\mathfrak P_\alpha(I)\|_{2,c}.
\]
This is \eqref{eq:V91-RSPCC} after enlarging \(K\).  The absence of
a general converse is the eigenfunction calculation in
Lemma~\ref{lem:V90-gap-direction}: the reverse constant grows like
\(\sqrt\lambda\).
\end{proof}

\begin{corollary}[The spectral tail is optional for the local tree
card]
\label{cor:V91-tail-optional}
The high-frequency estimate \eqref{eq:V90-LSPC-tail} is needed only
if one insists on the stronger \(L^2\) pair SPCC.  It is not needed
to deduce the all-order local tree card from the RSPCC.
\end{corollary}

\begin{proof}
The proof of Theorem~\ref{thm:V91-RSPCC-implies-card} uses only
\eqref{eq:V91-RSPCC} and the outer pairing
\eqref{eq:V91-outer-negative}; no \(L^2\) norm of a terminal tail
occurs.
\end{proof}

\section{Exact all-subset one-form formulation}
\label{sec:V91-one-form-recursion}

\begin{theorem}[Symmetry-first one-form recursion]
\label{thm:V91-one-form-recursion}
For every label set \(I\) with \(|I|\ge2\),
\begin{equation}
 \boxed{\qquad
 dR_\alpha\mathfrak P_\alpha(I)
 =
 \sum_{i\in I}
 \Theta_\alpha
 \left(
 D^{R_i},
 \mathfrak P_\alpha(I\setminus\{i\})
 \right).
 \qquad}
 \label{eq:V91-one-form-recursion}
\end{equation}
The sum on the right is formed with canonical tensor-leg order before
its column norm is taken.
\end{theorem}

\begin{proof}
The exact subset recursion \eqref{eq:V88-subset-recursion} is
\[
 \mathfrak P_\alpha(I)
 =
 \sum_{i\in I}
 \mathcal D_\alpha
 \left(
 D^{R_i},
 \mathfrak P_\alpha(I\setminus\{i\})
 \right).
\]
For each summand, the proof of
Lemma~\ref{lem:V90-negative-identity} gives
\[
 dR_\alpha\mathcal D_\alpha(F,G)
 =
 \Theta_\alpha(F,G).
\]
Apply this identity without separating the terms and then sum over
\(i\).
\end{proof}

\begin{corollary}[Equivalent assembled RSPCC gate]
\label{cor:V91-equivalent-gate}
For \(|I|\ge2\), the RSPCC estimate is exactly the bound
\begin{equation}
 \left\|
 \sum_{i\in I}
 \Theta_\alpha
 \left(
 D^{R_i},
 \mathfrak P_\alpha(I\setminus\{i\})
 \right)
 \right\|_{2,c}
 \le
 {|I|!K^{|I|}\over\sqrt\alpha}
 \left(\prod_{i\in I}b_i\right)
 B_I^{|I|-1}.
 \label{eq:V91-assembled-gate}
\end{equation}
\end{corollary}

\begin{proof}
Use Theorem~\ref{thm:V91-one-form-recursion} and the definition
\(\|H\|_{-1,c}=\|dR_\alpha H\|_{2,c}\).
\end{proof}

\begin{assessment}[The correct all-order analytic target]
\label{ass:V91-correct-target}
The former \(L^2\) SPCC remains a valid sufficient condition, but it
is stronger than the terminal contraction requires.  The
type-matched all-order target is now
\eqref{eq:V91-assembled-gate}: one symmetrized sum of intrinsic
one-forms, with one explicit \(\alpha^{-1/2}\) factor.  Taking norms
term by term would again erase the \(B_I^{|I|-1}\) incidence
powers.  No global Gaussian--Wilson transport appears.
\end{assessment}

\section{Low-order ownership}
\label{sec:V91-low-order}

\begin{lemma}[The singleton RSPCC is automatic]
\label{lem:V91-singleton}
For every all-thermal representation \(R\),
\begin{equation}
 \|\mathfrak P_\alpha(\{R\})\|_{-1,c}
 \le
 {C\over\sqrt\alpha}b_R.
 \label{eq:V91-singleton}
\end{equation}
\end{lemma}

\begin{proof}
By definition
\(\mathfrak P_\alpha(\{R\})=Q_\alpha D^R\).  Hence
\[
 \|\mathfrak P_\alpha(\{R\})\|_{-1,c}
 =
 \|dR_\alpha D^R\|_{2,c}
 \le
 {C\over\sqrt\alpha}
 \|Q_\alpha D^R\|_{2,c}
 \le
 {C\over\sqrt\alpha}b_R
\]
by \eqref{eq:V87-column-resolvent} and the inherited centered
representation estimate.
\end{proof}

\begin{proposition}[Owned pair sectors for the RSPCC]
\label{prop:V91-owned-pairs}
The two-label RSPCC holds:
\begin{enumerate}
\item for all pairs of normalized central characters; and
\item for arbitrary unitary matrix pairs satisfying
      \(b_R+b_S\ge\beta\), for any fixed \(\beta>0\).
\end{enumerate}
\end{proposition}

\begin{proof}
In the central sector apply
Lemma~\ref{lem:V90-gap-direction} to
\eqref{eq:V89-central-SPCC}.  In the second sector apply the same
lemma to \eqref{eq:V90-high-b-SPCC}.  Both results acquire the
required factor \(\alpha^{-1/2}\).
\end{proof}

\begin{corollary}[The first unowned RSPCC case]
\label{cor:V91-first-unowned}
The only unproved two-label case needed by the weaker all-order route
is the genuinely noncentral sector \(b_R+b_S<\beta\), where the exact
target is
\begin{equation}
 \left\|
 \Theta_\alpha(D^R,D^S)
 +
 \Theta_\alpha(D^S,D^R)
 \right\|_{2,c}
 \le
 {C\over\sqrt\alpha}
 b_Rb_S(b_R+b_S).
 \label{eq:V91-low-pair-gate}
\end{equation}
No high-frequency tail estimate is part of this target.
\end{corollary}

\begin{proof}
For a two-label set, Theorem
\ref{thm:V91-one-form-recursion} gives exactly the expression on the
left.  The right side is \eqref{eq:V91-RSPCC} with \(|I|=2\).
Proposition~\ref{prop:V91-owned-pairs} covers the complementary
sectors.
\end{proof}

\begin{firewall}[The RSPCC is a reduction, not an all-order proof]
\label{fire:V91-RSPCC-open}
V91 proves that the RSPCC is sufficient and is implied by the
\(L^2\) SPCC, while no converse follows from the spectral gap.  It does not prove
\eqref{eq:V91-low-pair-gate}, the RSPCC for arbitrary label sets, or
the all-order local card.  The exact one-form recursion identifies
where cancellation must be proved; it does not estimate that sum.
\end{firewall}

\section{V91 ledger and V92 programme}
\label{sec:V91-results}

\begin{theorem}[Fourth-series V91 result ledger]
\label{thm:V91-results}
V91:
\begin{enumerate}[label=\textup{(V91.\arabic*)},leftmargin=6.3em]
\item proves the dimension-free outer pairing estimate in the
      terminal negative norm;
\item defines the weaker RSPCC and proves that it implies the full
      all-order Pruefer/tree card;
\item proves that the former \(L^2\) SPCC implies the RSPCC, while
      the spectral gap supplies no converse;
\item derives the exact all-subset symmetry-first one-form recursion;
\item removes the high-frequency-tail estimate from the main
      all-order local-card route; and
\item closes the singleton, central-pair, and arbitrary high-\(b\)
      pair sectors of the RSPCC.
\end{enumerate}
\end{theorem}

\begin{proof}
Use
\Cref{lem:V91-outer-negative,
thm:V91-RSPCC-implies-card,
prop:V91-weaker,
cor:V91-tail-optional,
thm:V91-one-form-recursion,
lem:V91-singleton,
prop:V91-owned-pairs}.
\end{proof}

\begin{programme}[V92: the low-\(b\) matrix pair in its natural norm]
\label{prog:V91-V92}
\begin{enumerate}[leftmargin=3em]
\item Prove \eqref{eq:V91-low-pair-gate}, not the stronger
      \(L^2\) pair estimate.
\item Split the principal logarithmic chart into an inner ball and
      its exponentially small cut-locus tail.
\item On the inner ball rescale \(y=\sqrt\alpha Z\), subtract the
      constant tangent contraction in the complete two-ordering
      expression, and estimate the resulting first-order
      coefficients.
\item On the outer set use unitary representation bounds and Wilson
      concentration; do not differentiate the logarithmic Haar
      Jacobian to an order which grows with the representation.
\item Keep every estimate in column norm and verify uniformity up to
      \(b_R+b_S=\beta\).
\item If the chart perturbation loses uniformity, compute the exact
      conjugation-sector Jacobi block of the assembled one-form
      rather than returning to the over-strong \(L^2\) target.
\end{enumerate}
\end{programme}

\begin{firewall}[Claim boundary after fourth-series V91]
\label{fire:V91-claim-boundary}
V91 proves a weaker sufficient all-order certificate and its exact
one-form formulation.  It does not prove the low-\(b\) noncentral
pair gate, the all-subset RSPCC, the all-order local card, the
anisotropic scalar optimization, or the complete all-order selected
gate.  It does not prove all-order MTA, WUFC, CPM, cutoff removal,
reflection positivity, Osterwalder--Schrader reconstruction,
construction of the continuum Yang--Mills measure, or a uniform
positive continuum spectral gap.  The full Yang--Mills mass gap has
not been proved.
\end{firewall}

\paragraph{Parent-version record.}
V91 was formed from
\texttt{Renewal\_Yang\_Mills\_Mass\_Gap\_Research\_Notes\_V90.tex}.
The V90 parent had \(39\,778\) logical source lines,
\(1\,335\,805\) bytes, and SHA--256
\[
 \texttt{1efaec15bd87b8e293c48821e50b3095e9274f8ede8919de786ae04475cb8397}.
\]
The next compulsory companion audit is V95.

% =============================================================================
% END APPENDED FOURTH-SERIES V91 MATERIAL
% =============================================================================

% =============================================================================
% BEGIN APPENDED FOURTH-SERIES V92 MATERIAL
% =============================================================================

\chapter{V92: Global quaternion coordinates close the matrix-pair RSPCC}
\label{chap:V92-quaternion-pair}

\section{Smooth tangent coordinates}
\label{sec:V92-coordinates}

Write \(SU(2)\) as the unit quaternions
\begin{equation}
 U=\phi(U)I+\sum_{a=1}^3x_a(U)J_a,
 \qquad
 \phi^2+\sum_ax_a^2=1,
 \label{eq:V92-quaternion}
\end{equation}
where \(\phi(U)=\frac12\operatorname{ReTr}U=\cos\theta\),
\((x_1,x_2,x_3)=\sin\theta\,n\), and the fixed matrices \(J_a\)
only set the normalization.  Unlike the principal logarithm, the
three \(x_a\) are smooth on the whole group.

In the radius-two normalization used in V78, the ambient-coordinate
identities are
\begin{align}
 -\sum_cX_c^2x_a&={3\over4}x_a,
 \label{eq:V92-coordinate-Casimir}\\
 \Gamma(\phi,x_a)&=-{1\over4}\phi x_a,
 \label{eq:V92-coordinate-cross}\\
 \Gamma(x_a,x_b)&={1\over4}
 \left(\delta_{ab}-x_ax_b\right).
 \label{eq:V92-coordinate-metric}
\end{align}

\begin{lemma}[Smooth-coordinate Wilson moments]
\label{lem:V92-coordinate-moments}
For every fixed \(p<\infty\),
\begin{equation}
 \|x\|_{L^p(\nu_\alpha)}
 \le {C_p\over\sqrt\alpha},
 \qquad
 \|1-\phi\|_{L^p(\nu_\alpha)}
 \le {C_p\over\alpha}.
 \label{eq:V92-coordinate-moments}
\end{equation}
\end{lemma}

\begin{proof}
On the principal ball,
\[
 |x|=\sin\theta\le\theta=|Z|,
 \qquad
 0\le1-\phi=1-\cos\theta\le{\theta^2\over2}.
\]
Apply the logarithmic moment estimate
\eqref{eq:V82-log-moments}.  The antipode has Haar measure zero, so
the almost-everywhere inequalities suffice.
\end{proof}

\section{An explicit coordinate-resolvent approximation}
\label{sec:V92-coordinate-resolvent}

\begin{lemma}[Exact residual for a smooth coordinate]
\label{lem:V92-coordinate-residual}
For each \(b\), set
\begin{equation}
 r_b:=(1-\phi)x_b-{3\over\alpha}x_b.
 \label{eq:V92-rb}
\end{equation}
Then \(r_b\) is centered,
\begin{equation}
 \|r_b\|_2\le C\alpha^{-3/2},
 \label{eq:V92-rb-bound}
\end{equation}
and
\begin{equation}
 R_\alpha x_b={4\over\alpha}x_b+W_b,
 \qquad
 W_b=R_\alpha r_b,
 \qquad
 \|dW_b\|_2\le C\alpha^{-2}.
 \label{eq:V92-coordinate-resolvent}
\end{equation}
\end{lemma}

\begin{proof}
Equations
\eqref{eq:V78-generator},
\eqref{eq:V92-coordinate-Casimir}, and
\eqref{eq:V92-coordinate-cross} give the exact identity
\[
 A_\alpha x_b
 =
 {3\over4}x_b+{\alpha\over4}\phi x_b.
\]
Consequently
\[
 A_\alpha\left({4\over\alpha}x_b\right)
 =
 \phi x_b+{3\over\alpha}x_b
 =
 x_b-r_b.
\]
The central Wilson law is invariant under
\((\phi,x)\mapsto(\phi,-x)\), so \(x_b\) and \(r_b\) are centered.
By Holder and Lemma~\ref{lem:V92-coordinate-moments},
\[
 \|r_b\|_2
 \le
 \|1-\phi\|_4\|x_b\|_4
 +{3\over\alpha}\|x_b\|_2
 \le C\alpha^{-3/2}.
\]
Apply \(R_\alpha\) to the residual equation.  The gradient resolvent
bound \eqref{eq:V78-resolvent-gradient} then gives
\[
 \|dW_b\|_2
 \le {C\over\sqrt\alpha}\|r_b\|_2
 \le C\alpha^{-2}.
\]
\end{proof}

\begin{theorem}[Universal smooth-coordinate pair]
\label{thm:V92-coordinate-pair}
For all \(a,b\in\{1,2,3\}\),
\begin{align}
 \|\mathcal D_\alpha(x_a,x_b)\|_2
 &\le C\alpha^{-2},
 \label{eq:V92-coordinate-pair-L2}\\
 \|\mathcal D_\alpha(x_a,x_b)\|_{-1}
 &\le C\alpha^{-5/2}.
 \label{eq:V92-coordinate-pair-negative}
\end{align}
\end{theorem}

\begin{proof}
Use \eqref{eq:V92-coordinate-resolvent} and then
\eqref{eq:V92-coordinate-metric}:
\begin{align*}
 \mathcal B_\alpha(x_a,x_b)
 &=
 \Gamma(x_a,R_\alpha x_b)\\
 &=
 {4\over\alpha}\Gamma(x_a,x_b)
 +\Gamma(x_a,W_b)\\
 &=
 {1\over\alpha}(\delta_{ab}-x_ax_b)
 +\Gamma(x_a,W_b).
\end{align*}
The differential of every ambient coordinate is uniformly bounded.
Therefore
\[
 \|\Gamma(x_a,W_b)\|_2
 \le C\|dW_b\|_2
 \le C\alpha^{-2}.
\]
Moreover,
\[
 {1\over\alpha}\|Q_\alpha(x_ax_b)\|_2
 \le
 {2\over\alpha}\|x_a\|_4\|x_b\|_4
 \le C\alpha^{-2}.
\]
The constant \(\delta_{ab}/\alpha\) is annihilated by
\(Q_\alpha\).  This proves
\eqref{eq:V92-coordinate-pair-L2}.  Apply
\eqref{eq:V90-gap-direction} to obtain
\eqref{eq:V92-coordinate-pair-negative}.
\end{proof}

\section{A global smooth representation decomposition}
\label{sec:V92-representation-split}

For an irreducible unitary representation \(R\), let \(T_a^R\) be
the differential matrices matched to the coordinates \(x_a\) at the
identity and put
\begin{equation}
 \Lambda_R(U):=\sum_{a=1}^3T_a^R x_a(U),
 \qquad
 N_R(U):=D^R(U)-I-\Lambda_R(U),
 \qquad
 \widetilde N_R:=Q_\alpha N_R.
 \label{eq:V92-smooth-split}
\end{equation}
Thus
\begin{equation}
 Q_\alpha D^R=\Lambda_R+\widetilde N_R.
 \label{eq:V92-centered-smooth-split}
\end{equation}
The coordinate normalization changes the \(T_a^R\) by only fixed
constants, and hence
\begin{equation}
 \left\|\sum_a(T_a^R)^*T_a^R\right\|_{\rm op}^{1/2}
 \le C\sqrt{C_2(R)}.
 \label{eq:V92-T-row}
\end{equation}

\begin{lemma}[Value and derivative bounds for the smooth remainder]
\label{lem:V92-smooth-remainder}
Uniformly for all all-thermal \(R\),
\begin{align}
 \|\widetilde N_R\|_{2,c}
 &\le Cb_R^2,
 \label{eq:V92-N-value}\\
 \|d\widetilde N_R\|_{2,c}
 &\le C\sqrt\alpha\,b_R^2,
 \label{eq:V92-N-derivative-L2}\\
 \left\|
 \sum_a(X_a\widetilde N_R)^*
       (X_a\widetilde N_R)
 \right\|_{L^\infty({\rm op})}^{1/2}
 &\le C\sqrt\alpha\,b_R.
 \label{eq:V92-N-derivative-infty}
\end{align}
Also
\begin{equation}
 \|d\Lambda_R\|_{\infty,\mathrm{row}}
 \le C\sqrt\alpha\,b_R.
 \label{eq:V92-Lambda-derivative}
\end{equation}
\end{lemma}

\begin{proof}
Write \(U=e^Z\), \(Z=\theta n\), away from the null antipode.
The tangent-normalized vectors \(z=\theta n\) and
\(x=\sin\theta\,n\) satisfy
\[
 |z-x|\le C|Z|^3.
\]
The unitary Taylor remainder gives
\[
 \|D^R(e^Z)-I-dR(Z)\|
 \le {1\over2}C_2(R)|Z|^2.
\]
Together with \eqref{eq:V92-T-row} and
\eqref{eq:V82-log-moments}, this yields
\[
 \|N_R\|_{2,c}
 \le
 {CC_2(R)\over\alpha}
 +{C\sqrt{C_2(R)}\over\alpha^{3/2}}
 \le Cb_R^2.
\]
Centering costs at most two, proving
\eqref{eq:V92-N-value}.

The one-form \(dD^R-d\Lambda_R\) vanishes at the identity.
Its covariant derivative is bounded by
\(C(C_2(R)+\sqrt{C_2(R)})\): one derivative may hit a
representation generator, and all coordinate derivatives are fixed
smooth tensors.  Integration along a minimizing geodesic gives
\[
 |dN_R(U)|_{\mathrm{row}}
 \le
 C(C_2(R)+\sqrt{C_2(R)})\,|Z|.
\]
Using \(\||Z|\|_2\le C\alpha^{-1/2}\) and
\(s_\alpha\simeq\alpha^{-1}\),
\[
 \|dN_R\|_{2,c}
 \le
 {C(1+C_2(R))\over\sqrt\alpha}
 \le C\sqrt\alpha\,b_R^2.
\]
Centering does not change a derivative.

Finally, the unitary Casimir row bound controls \(dD^R\) by
\(C\sqrt{C_2(R)}\), while the fixed smooth coordinate derivatives
and \eqref{eq:V92-T-row} give the same bound for
\(d\Lambda_R\).  This proves
\eqref{eq:V92-N-derivative-infty} and
\eqref{eq:V92-Lambda-derivative}, using
\(\sqrt{C_2(R)}\le C\sqrt\alpha b_R\).
\end{proof}

\section{A general atomic estimate in the negative norm}
\label{sec:V92-general-atom}

For a smooth operator-valued \(F\), set
\[
 M(F):=
 \left\|
 \sum_a(X_aF)^*(X_aF)
 \right\|_{L^\infty({\rm op})}^{1/2}.
\]

\begin{lemma}[General atomic-left negative estimate]
\label{lem:V92-general-atom}
For every smooth \(F\) and every centered compatible \(G\),
\begin{equation}
 \|\mathcal D_\alpha(F,G)\|_{-1,c}
 \le
 {C M(F)\over\alpha}\|G\|_{2,c}.
 \label{eq:V92-general-atom}
\end{equation}
\end{lemma}

\begin{proof}
The row--column Cauchy--Schwarz inequality and
\eqref{eq:V87-column-resolvent} give
\[
 \|\mathcal B_\alpha(F,G)\|_{2,c}
 \le
 M(F)\|dR_\alpha G\|_{2,c}
 \le
 {CM(F)\over\sqrt\alpha}\|G\|_{2,c}.
\]
Centering costs at most two.  The result is centered, so
\eqref{eq:V90-gap-direction} supplies one further
\(\alpha^{-1/2}\).
\end{proof}

\section{The formerly dangerous first-slot remainder}
\label{sec:V92-first-slot}

For a representation \(S\), define
\[
 W_S:=\sum_bT_b^S W_b.
\]
Lemma~\ref{lem:V92-coordinate-residual} and
\eqref{eq:V92-T-row} give
\begin{equation}
 R_\alpha\Lambda_S
 =
 {4\over\alpha}\Lambda_S+W_S,
 \qquad
 \|dW_S\|_{2,c}
 \le C\sqrt{C_2(S)}\,\alpha^{-2}.
 \label{eq:V92-Lambda-resolvent}
\end{equation}

\begin{lemma}[First-slot nonlinear gain]
\label{lem:V92-first-slot}
For all all-thermal \(R,S\),
\begin{equation}
 \|\mathcal D_\alpha(\widetilde N_R,\Lambda_S)\|_{-1,c}
 \le
 {C\over\sqrt\alpha}
 b_Rb_S(b_R+b_S).
 \label{eq:V92-first-slot}
\end{equation}
\end{lemma}

\begin{proof}
Using \eqref{eq:V92-Lambda-resolvent},
\[
 \mathcal B_\alpha(\widetilde N_R,\Lambda_S)
 =
 {4\over\alpha}\Gamma(\widetilde N_R,\Lambda_S)
 +
 \Gamma(\widetilde N_R,W_S).
\]
The first term is bounded in column \(L^2\) by
\begin{align*}
 {C\over\alpha}
 \|d\widetilde N_R\|_{2,c}
 \|d\Lambda_S\|_{\infty,\mathrm{row}}
 &\le
 {C\over\alpha}
 (\sqrt\alpha b_R^2)(\sqrt\alpha b_S)\\
 &\le Cb_R^2b_S.
\end{align*}
For the second term use
\eqref{eq:V92-N-derivative-infty} and
\eqref{eq:V92-Lambda-resolvent}:
\[
 \|\Gamma(\widetilde N_R,W_S)\|_{2,c}
 \le
 C(\sqrt\alpha b_R)
 (\sqrt{C_2(S)}\,\alpha^{-2})
 \le {Cb_Rb_S\over\alpha}.
\]
After centering and applying \eqref{eq:V90-gap-direction},
\[
 \|\mathcal D_\alpha(\widetilde N_R,\Lambda_S)\|_{-1,c}
 \le
 {C\over\sqrt\alpha}b_R^2b_S
 +{Cb_Rb_S\over\alpha^{3/2}}.
\]
The first term has the required form.  Since
\(b_R,b_S\ge\sqrt{s_\alpha}\ge c\alpha^{-1/2}\) and
\(\alpha\ge1\),
\[
 {1\over\alpha}\le C(b_R+b_S).
\]
This absorbs the second term and proves
\eqref{eq:V92-first-slot}.
\end{proof}

\section{Complete matrix-pair RSPCC}
\label{sec:V92-complete-pair}

\begin{lemma}[Smooth linear pair]
\label{lem:V92-linear-pair}
For all all-thermal \(R,S\),
\begin{equation}
 \left\|
 \mathcal D_\alpha(\Lambda_R,\Lambda_S)
 +
 \mathcal D_\alpha(\Lambda_S,\Lambda_R)
 \right\|_{-1,c}
 \le
 {C\over\sqrt\alpha}
 b_Rb_S(b_R+b_S).
 \label{eq:V92-linear-pair}
\end{equation}
\end{lemma}

\begin{proof}
Expand in the nine fixed coordinate pairs and apply
\eqref{eq:V92-coordinate-pair-negative} together with the two
Casimir row bounds \eqref{eq:V92-T-row}.  This gives
\[
 \left\|
 \mathcal D_\alpha(\Lambda_R,\Lambda_S)
 \right\|_{-1,c}
 \le
 C\alpha^{-5/2}
 \sqrt{C_2(R)C_2(S)}.
\]
Since
\(\sqrt{C_2(R)}\le C\sqrt\alpha b_R\) and similarly for \(S\),
the right side is at most
\[
 {C\over\alpha^{3/2}}b_Rb_S.
\]
As in the proof of Lemma~\ref{lem:V92-first-slot},
\(\alpha^{-1}\le C(b_R+b_S)\).  Add the reversed ordering.
\end{proof}

\begin{theorem}[Complete two-label matrix RSPCC]
\label{thm:V92-complete-pair}
For arbitrary all-thermal irreducible unitary representations
\(R,S\),
\begin{equation}
 \boxed{\qquad
 \left\|
 \mathcal D_\alpha(D^R,D^S)
 +
 \mathcal D_\alpha(D^S,D^R)
 \right\|_{-1,c}
 \le
 {C\over\sqrt\alpha}
 b_Rb_S(b_R+b_S).
 \qquad}
 \label{eq:V92-complete-pair}
\end{equation}
Thus the RSPCC \eqref{eq:V91-RSPCC} holds for every label set of
cardinality at most two.
\end{theorem}

\begin{proof}
Constants are invisible to both arguments of
\(\mathcal D_\alpha\), so
\eqref{eq:V92-centered-smooth-split} gives an exact expansion into
the two choices \(\Lambda\) and \(\widetilde N\) at each input.

The linear--linear sum is bounded by
Lemma~\ref{lem:V92-linear-pair}.  The two terms with a remainder in
the first slot and a linear function in the second are bounded by
Lemma~\ref{lem:V92-first-slot} and its label reversal.

For every remaining ordered term, the second input is a centered
remainder.  Lemma~\ref{lem:V92-general-atom},
\eqref{eq:V92-N-value},
\eqref{eq:V92-N-derivative-infty}, and
\eqref{eq:V92-Lambda-derivative} give either
\[
 {C\over\sqrt\alpha}b_Rb_S^2
 \quad\hbox{or}\quad
 {C\over\sqrt\alpha}b_R^2b_S.
\]
Each is bounded by the right side of
\eqref{eq:V92-complete-pair}.  Summing the finite expansion proves
the estimate.  The singleton case was proved in
Lemma~\ref{lem:V91-singleton}.
\end{proof}

\begin{corollary}[The arity-three local tree card follows from the
RSPCC route]
\label{cor:V92-ternary-card}
Theorem~\ref{thm:V92-complete-pair}, inserted in
Theorem~\ref{thm:V91-RSPCC-implies-card} at \(m=3\), proves the
complete ternary local tree card without the direct Taylor-cumulant
enumeration of V82.
\end{corollary}

\begin{proof}
For \(m=3\), every terminal set in
\eqref{eq:V88-cumulant-compression} has two labels.  Apply
\eqref{eq:V92-complete-pair} and then the outer pairing
\eqref{eq:V91-outer-negative}.
\end{proof}

\begin{assessment}[The first new subset is now three labels]
\label{ass:V92-next-subset}
The pair obstruction isolated in V88 has been closed in the
type-matched norm actually required by the cumulant formula.  The
first unproved RSPCC instance is
\[
 \|\mathfrak P_\alpha(\{1,2,3\})\|_{-1,c}
 \le
 {3!K^3\over\sqrt\alpha}
 b_1b_2b_3(b_1+b_2+b_3)^2.
\]
By \eqref{eq:V91-one-form-recursion}, this is one assembled sum of
three one-forms.  Its proof must retain their cyclic cancellation;
three separate applications of
Lemma~\ref{lem:V92-general-atom} lose both additional incidence
powers.
\end{assessment}

\begin{firewall}[Pair closure is not all-subset closure]
\label{fire:V92-pair-not-all}
Theorem~\ref{thm:V92-complete-pair} proves the full matrix pair in
the RSPCC norm.  It does not prove the stronger \(L^2\) pair SPCC,
the three-label RSPCC, or the all-order local tree card.  The smooth
coordinate decomposition is used at fixed pair arity; no
uniform-in-order Taylor expansion has been asserted.
\end{firewall}

\section{V92 ledger and V93 programme}
\label{sec:V92-results}

\begin{theorem}[Fourth-series V92 result ledger]
\label{thm:V92-results}
V92:
\begin{enumerate}[label=\textup{(V92.\arabic*)},leftmargin=6.3em]
\item replaces the cut-locus logarithmic linearization by three
      globally smooth quaternion coordinates;
\item proves their explicit Wilson resolvent approximation and
      universal projected-pair bound;
\item proves dimension-free value, \(L^2\)-derivative, and uniform
      derivative bounds for the smooth representation remainder;
\item proves the formerly dangerous first-slot nonlinear estimate;
\item closes the complete two-label matrix RSPCC at every
      all-thermal representation scale; and
\item recovers the ternary local tree card through the new
      terminal-negative-norm route.
\end{enumerate}
\end{theorem}

\begin{proof}
Use
\Cref{lem:V92-coordinate-residual,
thm:V92-coordinate-pair,
lem:V92-smooth-remainder,
lem:V92-first-slot,
thm:V92-complete-pair,
cor:V92-ternary-card}.
\end{proof}

\begin{programme}[V93: the three-label assembled one-form]
\label{prog:V92-V93}
\begin{enumerate}[leftmargin=3em]
\item Expand
      \(\sum_{i=1}^3
      \Theta_\alpha(D^{R_i},
      \mathfrak P_\alpha(\{1,2,3\}\setminus\{i\}))\)
      only after the cyclic sum is formed.
\item Insert the smooth decomposition
      \eqref{eq:V92-centered-smooth-split} at the three atomic legs,
      but keep each already-symmetrized pair tail intact.
\item Prove the all-linear cyclic term using the universal
      coordinate resolvent and rotational tensor decomposition.
\item Pay every term with a smooth remainder by its extra \(b_i\);
      treat a remainder in the first slot with the explicit
      coordinate-resolvent split rather than an \(L^\infty\)
      derivative bound alone.
\item Verify the required square
      \((b_1+b_2+b_3)^2\) before taking the column norm.
\item Only after the three-label RSPCC is proved, test the general
      subset recursion for an exponential, rather than factorial,
      local ordering count.
\end{enumerate}
\end{programme}

\begin{firewall}[Claim boundary after fourth-series V92]
\label{fire:V92-claim-boundary}
V92 proves the complete two-label matrix RSPCC.  It does not prove
the three-label or all-subset RSPCC, the all-order local tree card,
the anisotropic scalar optimization, or the complete all-order
selected gate.  It does not prove all-order MTA, WUFC, CPM, cutoff
removal, reflection positivity, Osterwalder--Schrader reconstruction,
construction of the continuum Yang--Mills measure, or a uniform
positive continuum spectral gap.  The full Yang--Mills mass gap has
not been proved.
\end{firewall}

\paragraph{Parent-version record.}
V92 was formed from
\texttt{Renewal\_Yang\_Mills\_Mass\_Gap\_Research\_Notes\_V91.tex}.
The V91 parent had \(40\,228\) logical source lines,
\(1\,348\,991\) bytes, and SHA--256
\[
 \texttt{cb78438e92bcc6936db1aa0ebf6b9f2196f326f731604af7865f7a72180f1398}.
\]
The next compulsory companion audit is V95.

% =============================================================================
% END APPENDED FOURTH-SERIES V92 MATERIAL
% =============================================================================

% =============================================================================
% BEGIN APPENDED FOURTH-SERIES V93 MATERIAL
% =============================================================================

\chapter{V93: Quadratic-coordinate resolvents and the first three-label sectors}
\label{chap:V93-quadratic-three}

\section{The centered quadratic coordinate space}
\label{sec:V93-quadratic-space}

Let \(\mathcal Q_2\) be the fixed finite-dimensional space spanned by
the centered products
\[
 Q_\alpha(x_ax_b),
 \qquad 1\le a,b\le3.
\]
For \(q\in\mathcal Q_2\), write \(|q|_{\mathrm{coef}}\) for any fixed
Euclidean norm of its coefficients in this nine-element spanning
family.  All such norms are equivalent with constants independent of
\(\alpha\).

Put
\begin{equation}
 m_\alpha:=\mathbb E_{\nu_\alpha}|x|^2,
 \qquad
 \rho_\alpha:=|x|^2-m_\alpha,
 \qquad
 q_{ab}^{\circ}:=
 x_ax_b-{\delta_{ab}\over3}|x|^2.
 \label{eq:V93-quadratic-basis}
\end{equation}
The \(q_{ab}^{\circ}\) span the traceless quadratic sector, while
\(\rho_\alpha\) spans the radial sector.  The transverse identity
\eqref{eq:V78-transverse-moment} is
\begin{equation}
 \alpha m_\alpha=3\mathbb E_{\nu_\alpha}\phi.
 \label{eq:V93-moment-identity}
\end{equation}

\begin{lemma}[Exact action on quadratic coordinates]
\label{lem:V93-quadratic-action}
One has
\begin{align}
 A_\alpha q_{ab}^{\circ}
 &=
 \left(2+{\alpha\over2}\phi\right)q_{ab}^{\circ},
 \label{eq:V93-traceless-action}\\
 A_\alpha|x|^2
 &=
 \left(2+{\alpha\over2}\phi\right)|x|^2
 -{3\over2}.
 \label{eq:V93-radial-action}
\end{align}
\end{lemma}

\begin{proof}
The diffusion product rule and
\eqref{eq:V92-coordinate-Casimir}--\eqref{eq:V92-coordinate-metric}
give
\begin{align*}
 A_\alpha(x_ax_b)
 &=
 x_aA_\alpha x_b+x_bA_\alpha x_a
 -2\Gamma(x_a,x_b)\\
 &=
 \left(2+{\alpha\over2}\phi\right)x_ax_b
 -{1\over2}\delta_{ab}.
\end{align*}
Summing \(a=b\) proves
\eqref{eq:V93-radial-action}.  Subtracting one third of its trace
from the preceding identity proves
\eqref{eq:V93-traceless-action}.
\end{proof}

\begin{theorem}[Quadratic-coordinate resolvent approximation]
\label{thm:V93-quadratic-resolvent}
For every \(q\in\mathcal Q_2\),
\begin{equation}
 R_\alpha q={2\over\alpha}q+W_q,
 \qquad
 \|dW_q\|_2
 \le
 C|q|_{\mathrm{coef}}\alpha^{-5/2}.
 \label{eq:V93-quadratic-resolvent}
\end{equation}
\end{theorem}

\begin{proof}
First take \(q=q_{ab}^{\circ}\).  By
\eqref{eq:V93-traceless-action},
\[
 q-A_\alpha\left({2\over\alpha}q\right)
 =
 (1-\phi)q-{4\over\alpha}q.
\]
Lemma~\ref{lem:V92-coordinate-moments} gives
\[
 \|q\|_4\le C\alpha^{-1},
 \qquad
 \|(1-\phi)q\|_2\le C\alpha^{-2}.
\]
Thus the residual has \(L^2\) norm at most \(C\alpha^{-2}\).

For the radial basis vector, use
\eqref{eq:V93-radial-action}:
\begin{align*}
 \rho_\alpha
 -
 A_\alpha\left({2\over\alpha}\rho_\alpha\right)
 &=
 (1-\phi)|x|^2
 -{4\over\alpha}|x|^2
 -m_\alpha+{3\over\alpha}\\
 &=
 (1-\phi)|x|^2
 -{4\over\alpha}|x|^2
 +{3\over\alpha}\mathbb E(1-\phi),
\end{align*}
where the last equality is
\eqref{eq:V93-moment-identity}.  Every term has \(L^2\) norm
\(O(\alpha^{-2})\) by
Lemma~\ref{lem:V92-coordinate-moments}.  The residuals are centered
because they equal \(q-A_\alpha(2q/\alpha)\).

Apply \(R_\alpha\) and then
\eqref{eq:V78-resolvent-gradient}.  This yields an additional
\(\alpha^{-1/2}\) and proves the theorem on the displayed basis.
Finite-dimensional norm equivalence proves the general statement.
\end{proof}

\section{The linear--quadratic pair normal form}
\label{sec:V93-linear-quadratic}

For \(q\in\mathcal Q_2\), the polynomial
\(\Gamma(x_a,q)\) has degree one plus degree three in the ambient
coordinates.  Denote its homogeneous linear part by
\[
 \operatorname{Lin}\Gamma(x_a,q)
\]
and define
\begin{equation}
 \mathscr L_a(q)
 :=
 6\,\operatorname{Lin}\Gamma(x_a,q).
 \label{eq:V93-linear-output}
\end{equation}
This is a linear combination of \(x_1,x_2,x_3\), and
\begin{equation}
 |\,\mathscr L_a(q)\,|_{\mathrm{coef}}
 \le C|q|_{\mathrm{coef}}.
 \label{eq:V93-linear-output-bound}
\end{equation}

\begin{theorem}[Linear--quadratic projected-pair normal form]
\label{thm:V93-linear-quadratic-normal}
For \(q\in\mathcal Q_2\) and \(a\in\{1,2,3\}\),
\begin{equation}
 \mathcal D_\alpha(q,x_a)
 +
 \mathcal D_\alpha(x_a,q)
 =
 {1\over\alpha}\mathscr L_a(q)+E_{\alpha,a}(q),
 \label{eq:V93-linear-quadratic-normal}
\end{equation}
where \(E_{\alpha,a}(q)\) is centered and
\begin{equation}
 \|E_{\alpha,a}(q)\|_2
 \le
 C|q|_{\mathrm{coef}}\alpha^{-2}.
 \label{eq:V93-linear-quadratic-error}
\end{equation}
\end{theorem}

\begin{proof}
Use
\eqref{eq:V92-coordinate-resolvent} and
\eqref{eq:V93-quadratic-resolvent}:
\begin{align*}
 \mathcal B_\alpha(q,x_a)
 +\mathcal B_\alpha(x_a,q)
 &=
 {4\over\alpha}\Gamma(q,x_a)
 +{2\over\alpha}\Gamma(x_a,q)\\
 &\quad+
 \Gamma(q,W_a)+\Gamma(x_a,W_q)\\
 &=
 {6\over\alpha}\Gamma(x_a,q)
 +\Gamma(q,W_a)+\Gamma(x_a,W_q).
\end{align*}
The first term equals
\(\alpha^{-1}\mathscr L_a(q)\) plus
\(\alpha^{-1}\) times a cubic coordinate polynomial.  The latter
has \(L^2\) norm \(O(\alpha^{-5/2})\) by
Lemma~\ref{lem:V92-coordinate-moments}.

The differential of \(q\) is uniformly bounded by
\(C|q|_{\mathrm{coef}}\), and \(dx_a\) is uniformly bounded.
Therefore
\[
 \|\Gamma(q,W_a)\|_2
 \le C|q|_{\mathrm{coef}}\alpha^{-2},
 \qquad
 \|\Gamma(x_a,W_q)\|_2
 \le C|q|_{\mathrm{coef}}\alpha^{-5/2}.
\]
Apply \(Q_\alpha\).  The linear term is already centered, and
centering the remaining terms costs at most two.  Since
\(\alpha\ge1\), all remainders are bounded as in
\eqref{eq:V93-linear-quadratic-error}.
\end{proof}

\section{Multilinear three-label tails}
\label{sec:V93-multilinear-tail}

For centered compatible inputs \(F_1,F_2,F_3\), define
\begin{equation}
 \mathfrak P_\alpha(F_1,F_2,F_3)
 :=
 \sum_{i=1}^3
 \mathcal D_\alpha
 \left(
 F_i,
 \mathcal D_\alpha(F_j,F_k)
 +
 \mathcal D_\alpha(F_k,F_j)
 \right),
 \label{eq:V93-three-tail}
\end{equation}
where \(\{i,j,k\}=\{1,2,3\}\), with the canonical tensor-leg order.
This is the trilinear extension of
\(\mathfrak P_\alpha(\{1,2,3\})\).

For a linear coordinate combination
\[
 L=\sum_aT_ax_a,
\]
write
\[
 |L|_1:=
 \left\|\sum_aT_a^*T_a\right\|_{\rm op}^{1/2}.
\]
For an operator-valued centered quadratic coordinate polynomial
\[
 Q=\sum_\gamma C_\gamma q_\gamma,
\]
where \((q_\gamma)\) is the fixed basis
\eqref{eq:V93-quadratic-basis}, use the analogous finite
row-coefficient norm \(|Q|_2\).

\begin{lemma}[Linear-pair \(L^2\) defect]
\label{lem:V93-linear-pair-L2}
For linear coordinate combinations \(L_1,L_2\),
\begin{equation}
 \left\|
 \mathcal D_\alpha(L_1,L_2)
 +
 \mathcal D_\alpha(L_2,L_1)
 \right\|_{2,c}
 \le
 C|L_1|_1|L_2|_1\alpha^{-2}.
 \label{eq:V93-linear-pair-L2}
\end{equation}
\end{lemma}

\begin{proof}
Expand in the nine coordinate pairs and apply
\eqref{eq:V92-coordinate-pair-L2}.  Fixed-dimensional
row--column Cauchy--Schwarz and norm equivalence introduce only an
absolute constant.
\end{proof}

\begin{theorem}[All-linear three-label sector]
\label{thm:V93-all-linear-three}
For three linear coordinate combinations,
\begin{equation}
 \|\mathfrak P_\alpha(L_1,L_2,L_3)\|_{-1,c}
 \le
 C|L_1|_1|L_2|_1|L_3|_1\alpha^{-3}.
 \label{eq:V93-all-linear-three}
\end{equation}
\end{theorem}

\begin{proof}
For each outer label \(i\), apply
Lemma~\ref{lem:V92-general-atom} to \(L_i\) and the symmetrized
inner pair.  Since \(M(L_i)\le C|L_i|_1\), Lemma
\ref{lem:V93-linear-pair-L2} gives
\[
 \left\|
 \mathcal D_\alpha
 \left(
 L_i,
 \mathcal D_\alpha(L_j,L_k)
 +\mathcal D_\alpha(L_k,L_j)
 \right)
 \right\|_{-1,c}
 \le
 C|L_i|_1|L_j|_1|L_k|_1\alpha^{-3}.
\]
There are three outer labels.  Sum the bounds.
\end{proof}

\begin{theorem}[One-quadratic/two-linear three-label sector]
\label{thm:V93-one-quadratic-three}
Let \(Q_1\) be a centered quadratic coordinate polynomial and let
\(L_2,L_3\) be linear coordinate combinations.  Then
\begin{equation}
 \|\mathfrak P_\alpha(Q_1,L_2,L_3)\|_{-1,c}
 \le
 C|Q_1|_2|L_2|_1|L_3|_1\alpha^{-3}.
 \label{eq:V93-one-quadratic-three}
\end{equation}
\end{theorem}

\begin{proof}
There are three terms in \eqref{eq:V93-three-tail}.
When \(Q_1\) is the outer input, use
\(M(Q_1)\le C|Q_1|_2\),
Lemma~\ref{lem:V92-general-atom}, and
\eqref{eq:V93-linear-pair-L2}.  This gives
\[
 C|Q_1|_2|L_2|_1|L_3|_1\alpha^{-3}.
\]

Consider the term with outer input \(L_2\).  By
Theorem~\ref{thm:V93-linear-quadratic-normal}, its inner pair has the
form
\[
 {1\over\alpha}L_{13}+E_{13},
\]
where
\[
 |L_{13}|_1\le C|Q_1|_2|L_3|_1,
 \qquad
 \|E_{13}\|_{2,c}
 \le C|Q_1|_2|L_3|_1\alpha^{-2}.
\]
For the first part, expand in coordinate pairs and use
\eqref{eq:V92-coordinate-pair-negative}:
\[
 {1\over\alpha}
 \|\mathcal D_\alpha(L_2,L_{13})\|_{-1,c}
 \le
 C|Q_1|_2|L_2|_1|L_3|_1\alpha^{-7/2}.
\]
For the error, Lemma~\ref{lem:V92-general-atom} gives
\[
 \|\mathcal D_\alpha(L_2,E_{13})\|_{-1,c}
 \le
 C|Q_1|_2|L_2|_1|L_3|_1\alpha^{-3}.
\]
The term with outer input \(L_3\) is identical after exchanging
labels.  Since \(\alpha^{-7/2}\le\alpha^{-3}\), the three estimates
prove the theorem.
\end{proof}

\section{Insertion of representation coefficients}
\label{sec:V93-representation-insertion}

Refine the smooth split of V92 by defining
\begin{equation}
 \mathsf Q_R
 :=
 {1\over2}Q_\alpha(\Lambda_R^2),
 \qquad
 \mathsf M_R
 :=
 Q_\alpha D^R-\Lambda_R-\mathsf Q_R.
 \label{eq:V93-three-level-split}
\end{equation}

\begin{lemma}[Global smooth split through quadratic degree]
\label{lem:V93-three-level}
For every all-thermal representation \(R\),
\begin{equation}
 Q_\alpha D^R
 =
 \Lambda_R+\mathsf Q_R+\mathsf M_R,
 \label{eq:V93-three-level}
\end{equation}
with
\begin{align}
 |\Lambda_R|_1
 &\le C\sqrt{C_2(R)}
 \le C\sqrt\alpha\,b_R,
 \label{eq:V93-linear-coefficient}\\
 |\mathsf Q_R|_2
 &\le CC_2(R)
 \le C\alpha b_R^2,
 \label{eq:V93-quadratic-coefficient}\\
 \|\mathsf M_R\|_{2,c}
 &\le Cb_R^3.
 \label{eq:V93-cubic-remainder}
\end{align}
\end{lemma}

\begin{proof}
The first two coefficient estimates follow from
\eqref{eq:V92-T-row}.  For the remainder, write \(U=e^Z\).
The unitary third-order Taylor formula gives
\[
 \left\|
 D^R(e^Z)-I-dR(Z)-{1\over2}dR(Z)^2
 \right\|
 \le
 {1\over6}C_2(R)^{3/2}|Z|^3.
\]
Replacing \(z=\theta n\) by \(x=\sin\theta\,n\) changes the linear
term by \(O(\sqrt{C_2(R)}|Z|^3)\) and the quadratic term by
\(O(C_2(R)|Z|^4)\).  Apply
\eqref{eq:V82-log-moments}.  The resulting bounds are
\[
 C\left[
 \left({C_2(R)\over\alpha}\right)^{3/2}
 +{\sqrt{C_2(R)}\over\alpha^{3/2}}
 +{C_2(R)\over\alpha^2}
 \right]
 \le Cb_R^3,
\]
where the last term uses
\(\alpha^{-1}\le Cb_R\).  Centering costs at most two.
\end{proof}

\begin{corollary}[The first two three-label degree sectors satisfy
the RSPCC scale]
\label{cor:V93-first-sectors}
Let \(B=b_1+b_2+b_3\).  The all-linear contribution and every
one-quadratic/two-linear contribution to
\(\mathfrak P_\alpha(D^{R_1},D^{R_2},D^{R_3})\) satisfy
\begin{equation}
 {C\over\sqrt\alpha}
 b_1b_2b_3B^2.
 \label{eq:V93-first-sectors}
\end{equation}
\end{corollary}

\begin{proof}
For the all-linear contribution,
Theorem~\ref{thm:V93-all-linear-three} and
\eqref{eq:V93-linear-coefficient} give
\[
 C\alpha^{-3}
 \prod_{i=1}^3(\sqrt\alpha b_i)
 =
 C\alpha^{-3/2}b_1b_2b_3.
\]
Since \(B\ge3\sqrt{s_\alpha}\ge c\alpha^{-1/2}\), this is bounded
by \eqref{eq:V93-first-sectors}.

If the quadratic input is at label one,
Theorem~\ref{thm:V93-one-quadratic-three} and
\eqref{eq:V93-linear-coefficient}--%
\eqref{eq:V93-quadratic-coefficient} give
\[
 C\alpha^{-3}
 (\alpha b_1^2)(\sqrt\alpha b_2)(\sqrt\alpha b_3)
 =
 {C\over\alpha}b_1^2b_2b_3.
\]
Because \(B\ge c\alpha^{-1/2}\) and \(b_1\le B\),
\[
 {1\over\alpha}b_1^2b_2b_3
 \le
 {C\over\sqrt\alpha}b_1b_2b_3B^2.
\]
The other quadratic locations follow by relabelling.
\end{proof}

\begin{assessment}[Exact remaining degree sectors]
\label{ass:V93-remaining-sectors}
By trilinearity and \eqref{eq:V93-three-level}, the three-label tail
is the sum of \(3^3\) degree choices.  V93 closes the choice
\((1,1,1)\) and the three permutations of \((2,1,1)\).  Every
remaining choice has total Taylor excess
\[
 (d_1-1)+(d_2-1)+(d_3-1)\ge2,
\]
where \(\mathsf M\) is assigned degree at least three.  Those terms
have enough value powers but not yet enough controlled
first-slot derivative powers.  They must be estimated using the
quadratic normal form before a global \(L^\infty\) derivative is
taken.
\end{assessment}

\begin{firewall}[The first sectors are not the complete
three-label RSPCC]
\label{fire:V93-not-complete-three}
Corollary~\ref{cor:V93-first-sectors} proves the two lowest degree
sectors of the three-label tail.  It does not estimate the sectors
containing two quadratic pieces or one cubic remainder.  V93
therefore does not claim the complete three-label RSPCC or infer the
all-order local card.
\end{firewall}

\section{V93 ledger and V94 programme}
\label{sec:V93-results}

\begin{theorem}[Fourth-series V93 result ledger]
\label{thm:V93-results}
V93:
\begin{enumerate}[label=\textup{(V93.\arabic*)},leftmargin=6.3em]
\item computes the exact Wilson generator on radial and traceless
      centered quadratic quaternion coordinates;
\item proves the uniform quadratic resolvent approximation
      \(R_\alpha q=(2/\alpha)q+W_q\);
\item derives a linear--quadratic projected-pair normal form with an
      \(O(\alpha^{-2})\) \(L^2\) remainder;
\item proves the three-label RSPCC scale in the all-linear sector;
\item proves the same scale in every one-quadratic/two-linear
      sector; and
\item reduces the remaining three-label problem to degree patterns
      of total Taylor excess at least two.
\end{enumerate}
\end{theorem}

\begin{proof}
Use
\Cref{lem:V93-quadratic-action,
thm:V93-quadratic-resolvent,
thm:V93-linear-quadratic-normal,
thm:V93-all-linear-three,
thm:V93-one-quadratic-three,
cor:V93-first-sectors}.
\end{proof}

\begin{programme}[V94: completion of the three-label degree ledger]
\label{prog:V93-V94}
\begin{enumerate}[leftmargin=3em]
\item Prove \(L^2\)-derivative bounds for
      \(\mathsf M_R\) at the scale
      \(\sqrt\alpha\,b_R^3\), retaining the smooth-coordinate
      degree.
\item Treat the permutations of \((3,1,1)\) by resolving the
      innermost linear coordinate explicitly before pairing it with
      the cubic remainder.
\item Treat \((2,2,1)\) using
      Theorem~\ref{thm:V93-quadratic-resolvent} on whichever
      quadratic input is in the second slot.
\item Bound every pattern of excess at least three by discarding
      surplus \(b_i<1\) only after two distinct incidence powers
      have been assigned.
\item Assemble the three outer-label terms before the final column
      norm whenever exactly two excess powers are present.
\item If a first-slot quadratic term still loses one power, derive
      its quadratic--quadratic pair normal form rather than invoking
      an over-strong global derivative bound.
\end{enumerate}
\end{programme}

\begin{firewall}[Claim boundary after fourth-series V93]
\label{fire:V93-claim-boundary}
V93 proves the all-linear and one-quadratic sectors of the
three-label RSPCC.  It does not prove the remaining degree sectors,
the complete three-label or all-subset RSPCC, the all-order local
tree card, the anisotropic scalar optimization, or the complete
all-order selected gate.  It does not prove all-order MTA, WUFC,
CPM, cutoff removal, reflection positivity,
Osterwalder--Schrader reconstruction, construction of the continuum
Yang--Mills measure, or a uniform positive continuum spectral gap.
The full Yang--Mills mass gap has not been proved.
\end{firewall}

\paragraph{Parent-version record.}
V93 was formed from
\texttt{Renewal\_Yang\_Mills\_Mass\_Gap\_Research\_Notes\_V92.tex}.
The V92 parent had \(40\,831\) logical source lines,
\(1\,365\,654\) bytes, and SHA--256
\[
 \texttt{f83a4a67010765b1e8427aff21c22099f91695190632cca559ae3b1f51055e71}.
\]
The next compulsory companion audit is V95.

% =============================================================================
% END APPENDED FOURTH-SERIES V93 MATERIAL
% =============================================================================

% =============================================================================
% BEGIN APPENDED FOURTH-SERIES V94 MATERIAL
% =============================================================================

\chapter{V94: Closure of the sharp excess-two three-label sectors}
\label{chap:V94-excess-two}

\section{Derivative scale of the smooth cubic remainder}
\label{sec:V94-cubic-derivative}

Retain the global smooth decomposition
\[
 Q_\alpha D^R=\Lambda_R+\mathsf Q_R+\mathsf M_R
\]
from \eqref{eq:V93-three-level}.  The value bound
\(\|\mathsf M_R\|_{2,c}\le Cb_R^3\) was proved in V93.

\begin{lemma}[Cubic-remainder derivative bounds]
\label{lem:V94-cubic-derivative}
For every all-thermal representation \(R\),
\begin{align}
 \|d\mathsf M_R\|_{2,c}
 &\le C\sqrt\alpha\,b_R^3,
 \label{eq:V94-cubic-derivative-L2}\\
 M(\mathsf M_R)
 &\le C\alpha b_R^2.
 \label{eq:V94-cubic-derivative-infty}
\end{align}
\end{lemma}

\begin{proof}
Differentiating the third-order unitary Taylor remainder used in
Lemma~\ref{lem:V93-three-level} leaves two powers of \(dR(Z)\).
Its row norm is bounded by
\[
 C C_2(R)^{3/2}|Z|^2.
\]
The replacement of \(z=\theta n\) by
\(x=\sin\theta\,n\) contributes terms bounded by
\[
 C\sqrt{C_2(R)}|Z|^2
 +
 CC_2(R)|Z|^3.
\]
Using \eqref{eq:V82-log-moments} gives
\begin{align*}
 \|d\mathsf M_R\|_{2,c}
 &\le
 C\left[
 {C_2(R)^{3/2}\over\alpha}
 +{\sqrt{C_2(R)}\over\alpha}
 +{C_2(R)\over\alpha^{3/2}}
 \right]\\
 &\le C\sqrt\alpha\,b_R^3.
\end{align*}
Centering does not change a derivative.

For the uniform bound, use
\[
 d\mathsf M_R
 =
 dD^R-d\Lambda_R-d\mathsf Q_R.
\]
The first two row norms are \(O(\sqrt{C_2(R)})\), while the
quadratic polynomial has row norm \(O(C_2(R))\).  Hence
\[
 M(\mathsf M_R)
 \le C(1+C_2(R))
 \le C\alpha b_R^2.
\]
\end{proof}

\section{One cubic remainder and two linear inputs}
\label{sec:V94-cubic-linear}

\begin{lemma}[Cubic--linear symmetrized pair in \(L^2\)]
\label{lem:V94-cubic-linear-pair}
For all-thermal \(R,S\),
\begin{equation}
 \left\|
 \mathcal D_\alpha(\mathsf M_R,\Lambda_S)
 +
 \mathcal D_\alpha(\Lambda_S,\mathsf M_R)
 \right\|_{2,c}
 \le
 Cb_R^3b_S.
 \label{eq:V94-cubic-linear-pair}
\end{equation}
\end{lemma}

\begin{proof}
Use \eqref{eq:V92-Lambda-resolvent} in the first ordering:
\[
 \mathcal B_\alpha(\mathsf M_R,\Lambda_S)
 =
 {4\over\alpha}\Gamma(\mathsf M_R,\Lambda_S)
 +
 \Gamma(\mathsf M_R,W_S).
\]
Equations
\eqref{eq:V94-cubic-derivative-L2} and
\eqref{eq:V92-Lambda-derivative} bound the first term by
\[
 {C\over\alpha}
 (\sqrt\alpha b_R^3)(\sqrt\alpha b_S)
 =
 Cb_R^3b_S.
\]
For the second term, use
\eqref{eq:V94-cubic-derivative-infty} and
\eqref{eq:V92-Lambda-resolvent}:
\[
 \|\Gamma(\mathsf M_R,W_S)\|_{2,c}
 \le
 C(\alpha b_R^2)
 (\sqrt\alpha b_S\,\alpha^{-2})
 =
 {C\over\sqrt\alpha}b_R^2b_S
 \le Cb_R^3b_S,
\]
because \(b_R\ge c\alpha^{-1/2}\).

For the reversed ordering, the pre-gap estimate in the proof of
Lemma~\ref{lem:V92-general-atom} gives
\[
 \|\mathcal D_\alpha(\Lambda_S,\mathsf M_R)\|_{2,c}
 \le
 {CM(\Lambda_S)\over\sqrt\alpha}
 \|\mathsf M_R\|_{2,c}
 \le Cb_Sb_R^3.
\]
Centering costs only an absolute factor.  Add the two orderings.
\end{proof}

\begin{theorem}[The \((3,1,1)\) three-label sector]
\label{thm:V94-311}
For every permutation of the labels,
\begin{equation}
 \left\|
 \mathfrak P_\alpha
 (\mathsf M_{R_1},\Lambda_{R_2},\Lambda_{R_3})
 \right\|_{-1,c}
 \le
 {C\over\sqrt\alpha}
 b_1^3b_2b_3.
 \label{eq:V94-311}
\end{equation}
\end{theorem}

\begin{proof}
There are three outer-label terms in
\eqref{eq:V93-three-tail}.  If \(\mathsf M_{R_1}\) is outer,
Lemma~\ref{lem:V92-general-atom},
\eqref{eq:V94-cubic-derivative-infty}, and
\eqref{eq:V93-linear-pair-L2} give
\begin{align*}
 \left\|
 \mathcal D_\alpha
 \left(
 \mathsf M_{R_1},
 \mathcal D_\alpha(\Lambda_{R_2},\Lambda_{R_3})
 +\mathcal D_\alpha(\Lambda_{R_3},\Lambda_{R_2})
 \right)
 \right\|_{-1,c}
 &\le
 {C\alpha b_1^2\over\alpha}
 (\alpha^{-2}\sqrt{C_2(R_2)C_2(R_3)})\\
 &\le
 {C\over\alpha}b_1^2b_2b_3\\
 &\le
 {C\over\sqrt\alpha}b_1^3b_2b_3.
\end{align*}
The last step again uses
\(\alpha^{-1/2}\le Cb_1\).

If a linear input, say \(\Lambda_{R_2}\), is outer, apply
Lemma~\ref{lem:V92-general-atom} and
\eqref{eq:V94-cubic-linear-pair}:
\[
 \left\|
 \mathcal D_\alpha
 \left(
 \Lambda_{R_2},
 \mathcal D_\alpha(\mathsf M_{R_1},\Lambda_{R_3})
 +\mathcal D_\alpha(\Lambda_{R_3},\mathsf M_{R_1})
 \right)
 \right\|_{-1,c}
 \le
 {C\over\sqrt\alpha}b_2b_1^3b_3.
\]
The third outer-label term is identical.
\end{proof}

\section{Ordered quadratic normal forms}
\label{sec:V94-ordered-quadratic}

\begin{lemma}[Ordered quadratic--linear normal form]
\label{lem:V94-ordered-QL}
Let \(Q\) be a centered operator-valued quadratic coordinate
polynomial and \(L\) an operator-valued linear coordinate
combination.  There is a linear coordinate combination
\(\mathscr L(Q,L)\) with
\[
 |\mathscr L(Q,L)|_1\le C|Q|_2|L|_1
\]
such that
\begin{equation}
 \mathcal D_\alpha(Q,L)
 =
 {1\over\alpha}\mathscr L(Q,L)+E_\alpha(Q,L),
 \qquad
 \|E_\alpha(Q,L)\|_{2,c}
 \le
 C|Q|_2|L|_1\alpha^{-2}.
 \label{eq:V94-ordered-QL}
\end{equation}
Moreover,
\begin{equation}
 \|\mathcal D_\alpha(Q,L)\|_{-1,c}
 \le
 C|Q|_2|L|_1\alpha^{-2}.
 \label{eq:V94-ordered-QL-negative}
\end{equation}
\end{lemma}

\begin{proof}
From \eqref{eq:V92-Lambda-resolvent},
\[
 \mathcal B_\alpha(Q,L)
 =
 {4\over\alpha}\Gamma(Q,L)+\Gamma(Q,W_L).
\]
The first term is \(4/\alpha\) times a degree-one plus degree-three
coordinate polynomial.  Define \(\mathscr L(Q,L)\) to be four times
its linear part.  The cubic part has \(L^2\) norm
\(O(|Q|_2|L|_1\alpha^{-3/2})\), and hence contributes
\(O(|Q|_2|L|_1\alpha^{-5/2})\).
The uniform derivative bound on a fixed quadratic polynomial and
\eqref{eq:V92-Lambda-resolvent} give
\[
 \|\Gamma(Q,W_L)\|_{2,c}
 \le C|Q|_2|L|_1\alpha^{-2}.
\]
Centering proves \eqref{eq:V94-ordered-QL}.

For a linear coordinate combination \(K\),
\eqref{eq:V92-coordinate-resolvent} implies
\[
 \|K\|_{-1,c}=\|dR_\alpha K\|_{2,c}
 \le C|K|_1\alpha^{-1}.
\]
Apply this to the leading term of
\eqref{eq:V94-ordered-QL}, and apply
\eqref{eq:V90-gap-direction} to the centered error.  This proves
\eqref{eq:V94-ordered-QL-negative}.
\end{proof}

\begin{lemma}[Quadratic--quadratic pair in \(L^2\)]
\label{lem:V94-QQ-pair}
For centered operator-valued quadratic coordinate polynomials
\(Q_1,Q_2\),
\begin{equation}
 \left\|
 \mathcal D_\alpha(Q_1,Q_2)
 +
 \mathcal D_\alpha(Q_2,Q_1)
 \right\|_{2,c}
 \le
 C|Q_1|_2|Q_2|_2\alpha^{-2}.
 \label{eq:V94-QQ-pair}
\end{equation}
\end{lemma}

\begin{proof}
Use \eqref{eq:V93-quadratic-resolvent}:
\[
 \mathcal B_\alpha(Q_1,Q_2)
 =
 {2\over\alpha}\Gamma(Q_1,Q_2)
 +
 \Gamma(Q_1,W_{Q_2}).
\]
The derivatives of two quadratic coordinate polynomials are linear
in \(x\), up to fixed smooth coefficient tensors.  Therefore
\[
 \|\Gamma(Q_1,Q_2)\|_{2,c}
 \le
 C|Q_1|_2|Q_2|_2\|x\|_4^2
 \le
 C|Q_1|_2|Q_2|_2\alpha^{-1}.
\]
The error is bounded by the uniform derivative of \(Q_1\) and
\eqref{eq:V93-quadratic-resolvent}:
\[
 \|\Gamma(Q_1,W_{Q_2})\|_{2,c}
 \le
 C|Q_1|_2|Q_2|_2\alpha^{-5/2}.
\]
Center, use \(\alpha\ge1\), and add the reversed ordering.
\end{proof}

\begin{theorem}[The \((2,2,1)\) three-label sector]
\label{thm:V94-221}
For two centered quadratic coordinate polynomials and one linear
coordinate combination,
\begin{equation}
 \|\mathfrak P_\alpha(Q_1,Q_2,L_3)\|_{-1,c}
 \le
 C|Q_1|_2|Q_2|_2|L_3|_1\alpha^{-3}.
 \label{eq:V94-221}
\end{equation}
\end{theorem}

\begin{proof}
When \(L_3\) is outer, combine
Lemma~\ref{lem:V92-general-atom} with
\eqref{eq:V94-QQ-pair}:
\[
 \left\|
 \mathcal D_\alpha
 \left(
 L_3,
 \mathcal D_\alpha(Q_1,Q_2)
 +\mathcal D_\alpha(Q_2,Q_1)
 \right)
 \right\|_{-1,c}
 \le
 C|Q_1|_2|Q_2|_2|L_3|_1\alpha^{-3}.
\]

When \(Q_1\) is outer, Theorem
\ref{thm:V93-linear-quadratic-normal} writes the inner pair as
\[
 {1\over\alpha}L_{23}+E_{23},
\]
where
\[
 |L_{23}|_1\le C|Q_2|_2|L_3|_1,
 \qquad
 \|E_{23}\|_{2,c}
 \le C|Q_2|_2|L_3|_1\alpha^{-2}.
\]
Equation~\eqref{eq:V94-ordered-QL-negative} gives
\[
 {1\over\alpha}
 \|\mathcal D_\alpha(Q_1,L_{23})\|_{-1,c}
 \le
 C|Q_1|_2|Q_2|_2|L_3|_1\alpha^{-3}.
\]
For the error use Lemma~\ref{lem:V92-general-atom} and
\(M(Q_1)\le C|Q_1|_2\):
\[
 \|\mathcal D_\alpha(Q_1,E_{23})\|_{-1,c}
 \le
 C|Q_1|_2|Q_2|_2|L_3|_1\alpha^{-3}.
\]
The outer-\(Q_2\) term is identical.
\end{proof}

\section{All sharp excess-two patterns are owned}
\label{sec:V94-sharp-ledger}

\begin{corollary}[Complete excess-two three-label bound]
\label{cor:V94-excess-two}
Every contribution to the three-label tail
\(\mathfrak P_\alpha(D^{R_1},D^{R_2},D^{R_3})\) having total smooth
Taylor excess exactly two satisfies
\begin{equation}
 {C\over\sqrt\alpha}
 b_1b_2b_3(b_1+b_2+b_3)^2.
 \label{eq:V94-excess-two}
\end{equation}
\end{corollary}

\begin{proof}
The excess-two patterns are precisely the permutations of
\((3,1,1)\) and \((2,2,1)\).
The first are bounded by Theorem~\ref{thm:V94-311}.

For the second, use Theorem~\ref{thm:V94-221} and
\eqref{eq:V93-linear-coefficient}--%
\eqref{eq:V93-quadratic-coefficient}:
\begin{align*}
 |\,\mathsf Q_{R_1}\,|_2
 |\,\mathsf Q_{R_2}\,|_2
 |\Lambda_{R_3}|_1\alpha^{-3}
 &\le
 C(\alpha b_1^2)(\alpha b_2^2)
 (\sqrt\alpha b_3)\alpha^{-3}\\
 &=
 {C\over\sqrt\alpha}b_1^2b_2^2b_3\\
 &\le
 {C\over\sqrt\alpha}b_1b_2b_3
 (b_1+b_2+b_3)^2.
\end{align*}
Relabeling covers all three placements.
\end{proof}

\begin{assessment}[Only surplus-excess patterns remain]
\label{ass:V94-surplus}
V93 owned total excess zero and one.  V94 owns every pattern of
total excess exactly two.  The remaining patterns are
\[
 (3,2,1),\quad
 (2,2,2),\quad
 (3,3,1),\quad
 (3,2,2),\quad
 (3,3,2),\quad
 (3,3,3),
\]
up to permutation, where degree three denotes the remainder
\(\mathsf M\).  Each has at least one surplus \(b\)-power beyond the
two required incidences.  What remains is to preserve those powers
when a high-degree piece occurs in the first slot; a crude uniform
derivative estimate may spend the surplus prematurely.
\end{assessment}

\begin{firewall}[Sharp sectors do not yet complete the degree ledger]
\label{fire:V94-not-complete}
Corollary~\ref{cor:V94-excess-two} closes every minimally powered
three-label sector.  It does not by itself estimate the six
surplus-excess patterns in
Assessment~\ref{ass:V94-surplus}.  V94 therefore does not yet claim
the complete three-label RSPCC.
\end{firewall}

\section{V94 ledger and V95 programme}
\label{sec:V94-results}

\begin{theorem}[Fourth-series V94 result ledger]
\label{thm:V94-results}
V94:
\begin{enumerate}[label=\textup{(V94.\arabic*)},leftmargin=6.3em]
\item proves the \(L^2\) and uniform derivative scales of the smooth
      cubic remainder;
\item proves the cubic--linear \(L^2\) pair estimate and closes every
      \((3,1,1)\) three-label sector;
\item derives an ordered quadratic--linear normal form in both
      \(L^2\) and the terminal negative norm;
\item proves a quadratic--quadratic projected-pair \(L^2\) bound;
\item closes every \((2,2,1)\) three-label sector; and
\item leaves only Taylor patterns with at least one surplus
      incidence power.
\end{enumerate}
\end{theorem}

\begin{proof}
Use
\Cref{lem:V94-cubic-derivative,
lem:V94-cubic-linear-pair,
thm:V94-311,
lem:V94-ordered-QL,
lem:V94-QQ-pair,
thm:V94-221,
cor:V94-excess-two}.
\end{proof}

\begin{programme}[V95: audit and surplus-excess closure]
\label{prog:V94-V95}
\begin{enumerate}[leftmargin=3em]
\item Perform the compulsory read-only audit of the newest active
      SM--Einstein and Navier--Stokes notes.
\item Prove pair \(L^2\) bounds for the degree combinations
      \((3,2)\) and \((3,3)\), using the explicit quadratic
      resolvent before a derivative of \(\mathsf M\) is estimated.
\item Close \((3,2,1)\) and \((2,2,2)\) first; these have exactly
      one surplus power.
\item Use \(b_i<1\) to discard only powers already assigned beyond
      two incidences, never a base factor.
\item Complete all remaining permutations and state the full
      three-label RSPCC only if every first-slot placement has been
      checked.
\item If the surplus is still lost, derive the cubic--quadratic
      pair normal form rather than appealing to the finite arity-four
      cumulant card.
\end{enumerate}
\end{programme}

\begin{firewall}[Claim boundary after fourth-series V94]
\label{fire:V94-claim-boundary}
V94 proves all three-label sectors of total Taylor excess at most
two.  It does not prove the surplus-excess sectors, the complete
three-label or all-subset RSPCC, the all-order local tree card, the
anisotropic scalar optimization, or the complete all-order selected
gate.  It does not prove all-order MTA, WUFC, CPM, cutoff removal,
reflection positivity, Osterwalder--Schrader reconstruction,
construction of the continuum Yang--Mills measure, or a uniform
positive continuum spectral gap.  The full Yang--Mills mass gap has
not been proved.
\end{firewall}

\paragraph{Parent-version record.}
V94 was formed from
\texttt{Renewal\_Yang\_Mills\_Mass\_Gap\_Research\_Notes\_V93.tex}.
The V93 parent had \(41\,407\) logical source lines,
\(1\,381\,797\) bytes, and SHA--256
\[
 \texttt{d7a094b43d42cadb69f30c42f612ed0303b215460a8cd475bed96d0acdd8b31a}.
\]
The next compulsory companion audit is V95.

% =============================================================================
% END APPENDED FOURTH-SERIES V94 MATERIAL
% =============================================================================

% =============================================================================
% BEGIN APPENDED FOURTH-SERIES V95 MATERIAL
% =============================================================================

\chapter{V95: Companion audit and the fixed-\(p\) surplus-sector gate}
\label{chap:V95-fixed-p}

\section{Compulsory companion audit}
\label{sec:V95-audit}

The newest active companion files were inspected read-only before the
present argument.  At the time of inspection they were
\begin{equation}
 \begin{split}
 &\texttt{Renewal\_SM\_Einstein\_Continuum\_Research\_Notes\_V86.tex},\\
 &33\,590\text{ logical source lines},\qquad
 1\,162\,800\text{ bytes},\\
 &\operatorname{SHA256}=
 \texttt{f6773d345386635bc7536a8b74618449900e7fd2febd070a7ef3ba892b2718ff},
 \end{split}
 \label{eq:V95-SM-fingerprint}
\end{equation}
and
\begin{equation}
 \begin{split}
 &\texttt{Renewal\_NS\_Stability\_Research\_Notes\_V31.tex},\\
 &23\,052\text{ logical source lines},\qquad
 887\,537\text{ bytes},\\
 &\operatorname{SHA256}=
 \texttt{f1121b62238ad5491bb7cf03635ea9934942edf573ed8faaa9ee79e1d38a6696}.
 \end{split}
 \label{eq:V95-NS-fingerprint}
\end{equation}
Neither file was modified.  They are evidence, not instructions.

SM V80 proves an exact negative-to-strong Hilbert-scale
interpolation, while explicitly warning that positive regularity is
an independent payment.  V81 identifies the negative norm through a
Poisson--Stein current, and V82 derives one positive derivative from
a Fisher/Lyapunov estimate in positive Gibbs sectors.  That
interpolation is sufficient for the SM cutoff telescope because any
positive geometric gain is summable.  It cannot be imported into
the present sharp incidence problem: interpolation would reduce the
power of \(b_1+b_2+b_3\).

SM V85--V86 makes a second point relevant here.  Moment divisibility
and Ward cancellation hold only for a complete assembled source or
packet; a single extracted vertex need not have the gain.  The NS
lesson remains unchanged: retain the physical root until its unique
critical moment is paid.

\begin{assessment}[Direction selected by the audit]
\label{ass:V95-audit-direction}
For the surplus Taylor sectors, the missing operation is not another
negative-to-strong interpolation.  It is a fixed-\(p\) gradient
resolvent estimate which permits the typical \(L^p\) derivative size
of a smooth degree piece to replace its much larger global
\(L^\infty\) derivative.  Only \(p=4,8\) are needed at three labels.
Once that estimate is available, all terms with at least two excess
powers close without any further cancellation claim.
\end{assessment}

\section{Complete fixed-\(p\) column norms}
\label{sec:V95-column-Lp}

For \(2\le p<\infty\), define
\begin{align}
 \|F\|_{p,c}
 &:=
 \sup_{\|v\|=1}
 \left(\mathbb E_{\nu_\alpha}\|Fv\|^p\right)^{1/p},
 \label{eq:V95-column-Lp}\\
 \|dF\|_{p,c}
 &:=
 \sup_{\|v\|=1}
 \left(
 \mathbb E_{\nu_\alpha}
 \left(\sum_a\|X_aFv\|^2\right)^{p/2}
 \right)^{1/p}.
 \label{eq:V95-column-gradient-Lp}
\end{align}
At \(p=2\) these agree with the column norms of V87.

\begin{lemma}[Fixed-\(p\) degree bounds]
\label{lem:V95-degree-Lp}
For every fixed \(2\le p\le8\), the smooth degree pieces of
\eqref{eq:V93-three-level} satisfy
\begin{align}
 \|\Lambda_R\|_{p,c}
 &\le C_pb_R,
 &
 \|d\Lambda_R\|_{p,c}
 &\le C_p\sqrt\alpha\,b_R,
 \label{eq:V95-L-degree}\\
 \|\mathsf Q_R\|_{p,c}
 &\le C_pb_R^2,
 &
 \|d\mathsf Q_R\|_{p,c}
 &\le C_p\sqrt\alpha\,b_R^2,
 \label{eq:V95-Q-degree}\\
 \|\mathsf M_R\|_{p,c}
 &\le C_pb_R^3,
 &
 \|d\mathsf M_R\|_{p,c}
 &\le C_p\sqrt\alpha\,b_R^3.
 \label{eq:V95-M-degree}
\end{align}
All constants are independent of the representation dimension.
\end{lemma}

\begin{proof}
For \(\Lambda_R\), combine the Casimir row estimate
\eqref{eq:V92-T-row} with
\(\|x\|_p\le C_p\alpha^{-1/2}\); its derivative is bounded by the
same coefficient row norm because the three \(x_a\) have fixed
smooth derivatives.

The value of \(\mathsf Q_R\) is a quadratic coordinate polynomial
with coefficient norm \(O(C_2(R))\).  Its derivative is coefficient
times one coordinate.  Therefore
\[
 \|\mathsf Q_R\|_{p,c}
 \le {C_pC_2(R)\over\alpha}
 \le C_pb_R^2,
 \qquad
 \|d\mathsf Q_R\|_{p,c}
 \le {C_pC_2(R)\over\sqrt\alpha}
 \le C_p\sqrt\alpha b_R^2.
\]

Finally repeat the Taylor-remainder proof of
\eqref{eq:V93-cubic-remainder} and
\eqref{eq:V94-cubic-derivative-L2}, replacing every fixed \(L^2\)
moment by the corresponding fixed \(L^{kp}\) logarithmic moment
from \eqref{eq:V82-log-moments}.  The representation coefficients
are unchanged, so the bounds are exactly
\eqref{eq:V95-M-degree}.  Centering costs at most two in every
\(L^p\) column norm.
\end{proof}

\section{The fixed-\(p\) Wilson resolvent card}
\label{sec:V95-FWRC}

\begin{definition}[FWRC]
\label{def:V95-FWRC}
The fixed-\(p\) Wilson resolvent card is the pair of complete
estimates
\begin{equation}
 \boxed{\qquad
 \|dR_\alpha G\|_{p,c}
 \le
 {C_p\over\sqrt\alpha}\|Q_\alpha G\|_{p,c},
 \qquad
 p\in\{4,8\},
 \qquad}
 \label{eq:V95-FWRC}
\end{equation}
for all smooth finite-dimensional operator-valued \(G\), with
constants independent of \(\alpha\) and the target dimension.
\end{definition}

\begin{remark}[What is and is not already known]
\label{rem:V95-FWRC-status}
At \(p=2\), \eqref{eq:V95-FWRC} is exactly the proved complete
column resolvent estimate
\eqref{eq:V87-column-resolvent}.  Neither the spectral gap nor
Hilbert-scale interpolation proves the \(p=4,8\) estimates.  They
are complete Riesz-transform bounds for the strongly concentrated
Wilson diffusion and require a separate proof.
\end{remark}

\begin{lemma}[Degree-adapted ordered pair under the FWRC]
\label{lem:V95-degree-pair}
Assume the FWRC.  Let
\[
 P_{R,d}\in\{\Lambda_R,\mathsf Q_R,\mathsf M_R\},
 \qquad d\in\{1,2,3\},
\]
with \(\mathsf M_R\) assigned degree three.  Then
\begin{align}
 \|\mathcal D_\alpha(P_{R,d},P_{S,e})\|_{4,c}
 &\le Cb_R^db_S^e,
 \label{eq:V95-pair-L4}\\
 \|\mathcal D_\alpha(P_{R,d},G)\|_{-1,c}
 &\le
 {C\over\sqrt\alpha}b_R^d\|G\|_{4,c}
 \label{eq:V95-outer-degree}
\end{align}
for every centered \(G\in L^4\) in the second line.
\end{lemma}

\begin{proof}
For the first line, fixed-dimensional row--column Holder gives
\[
 \|\mathcal B_\alpha(P_{R,d},P_{S,e})\|_{4,c}
 \le
 C\|dP_{R,d}\|_{8,c}
 \|dR_\alpha P_{S,e}\|_{8,c}.
\]
Use Lemma~\ref{lem:V95-degree-Lp} and the \(p=8\) FWRC.  The factors
\(\sqrt\alpha\) and \(\alpha^{-1/2}\) cancel, leaving
\(Cb_R^db_S^e\).  Centering costs at most two.

For the second line, Holder at exponents four gives
\[
 \|\mathcal B_\alpha(P_{R,d},G)\|_{2,c}
 \le
 C\|dP_{R,d}\|_{4,c}
 \|dR_\alpha G\|_{4,c}
 \le
 Cb_R^d\|G\|_{4,c}.
\]
Center and apply \eqref{eq:V90-gap-direction} to obtain the terminal
\(\alpha^{-1/2}\) factor.
\end{proof}

\section{Conditional closure of every surplus sector}
\label{sec:V95-surplus}

\begin{theorem}[FWRC closes all three-label surplus-excess patterns]
\label{thm:V95-FWRC-surplus}
Assume the FWRC.  For any degree choice
\[
 d_i\in\{1,2,3\},
 \qquad
 \sum_{i=1}^3(d_i-1)\ge2,
\]
one has
\begin{equation}
 \left\|
 \mathfrak P_\alpha
 (P_{R_1,d_1},P_{R_2,d_2},P_{R_3,d_3})
 \right\|_{-1,c}
 \le
 {C\over\sqrt\alpha}
 b_1^{d_1}b_2^{d_2}b_3^{d_3}.
 \label{eq:V95-degree-three-tail}
\end{equation}
Consequently
\begin{equation}
 \left\|
 \mathfrak P_\alpha
 (P_{R_1,d_1},P_{R_2,d_2},P_{R_3,d_3})
 \right\|_{-1,c}
 \le
 {C\over\sqrt\alpha}
 b_1b_2b_3(b_1+b_2+b_3)^2.
 \label{eq:V95-surplus-target}
\end{equation}
\end{theorem}

\begin{proof}
Expand \eqref{eq:V93-three-tail}.  For one ordered term, apply
\eqref{eq:V95-pair-L4} to its inner pair and then
\eqref{eq:V95-outer-degree} to the outer input.  The result is
\[
 {C\over\sqrt\alpha}
 b_1^{d_1}b_2^{d_2}b_3^{d_3},
\]
independently of the ordering.  Sum the six ordered chains to obtain
\eqref{eq:V95-degree-three-tail}.

Write \(e_i=d_i-1\).  At least two factors occur in the multiset
containing \(e_i\) copies of \(b_i\).  Retain any two of them and
discard every remaining surplus factor using \(b_i<1\).  The retained
product is at most
\((b_1+b_2+b_3)^2\).  Factoring one base \(b_i\) at every label
proves \eqref{eq:V95-surplus-target}.
\end{proof}

\begin{corollary}[FWRC implies the complete three-label RSPCC]
\label{cor:V95-complete-three}
If the FWRC holds, then
\begin{equation}
 \boxed{\qquad
 \|\mathfrak P_\alpha(\{1,2,3\})\|_{-1,c}
 \le
 {3!K^3\over\sqrt\alpha}
 b_1b_2b_3(b_1+b_2+b_3)^2.
 \qquad}
 \label{eq:V95-complete-three}
\end{equation}
\end{corollary}

\begin{proof}
The smooth split \eqref{eq:V93-three-level} has \(3^3\) sectors.
The excess-zero and excess-one sectors were proved in
Corollary~\ref{cor:V93-first-sectors}.  The sharp excess-two sectors
were proved in Corollary~\ref{cor:V94-excess-two}.  Theorem
\ref{thm:V95-FWRC-surplus} covers every remaining sector.  Sum this
fixed finite ledger and enlarge \(K\).
\end{proof}

\begin{assessment}[One precise analytic acquisition remains at
three labels]
\label{ass:V95-one-gate}
The algebraic and representation-degree ledger is complete.  The
only unproved input in Corollary~\ref{cor:V95-complete-three} is the
FWRC \eqref{eq:V95-FWRC}.  It is a fixed-exponent, one-link analytic
estimate, not an all-order Taylor comparison.  Proving it would close
the complete three-label RSPCC and make four labels the first unknown
subset size.
\end{assessment}

\begin{firewall}[The FWRC has not been inferred from the gap]
\label{fire:V95-FWRC-open}
V95 proves the fixed-\(p\) sizes of the representation degree pieces
and proves that the FWRC is sufficient.  It does not prove the
\(p=4,8\) gradient-resolvent estimates.  The \(p=2\) spectral gap
cannot be upgraded to them by reversing an interpolation inequality,
and the global Bakry--Emery curvature is negative near the antipode.
Thus the complete three-label RSPCC remains conditional.
\end{firewall}

\section{V95 ledger and V96 programme}
\label{sec:V95-results}

\begin{theorem}[Fourth-series V95 result ledger]
\label{thm:V95-results}
V95:
\begin{enumerate}[label=\textup{(V95.\arabic*)},leftmargin=6.3em]
\item completes the compulsory read-only SM/NS audit;
\item proves fixed-\(p\), dimension-free value and derivative bounds
      for every smooth representation degree piece through degree
      three;
\item identifies the exact \(p=4,8\) complete Wilson
      gradient-resolvent card;
\item proves degree-adapted ordered pair and outer negative-norm
      bounds from that card;
\item proves that the FWRC closes every surplus-excess three-label
      sector; and
\item reduces the complete three-label RSPCC to the single analytic
      acquisition \eqref{eq:V95-FWRC}.
\end{enumerate}
\end{theorem}

\begin{proof}
Use
\Cref{lem:V95-degree-Lp,
lem:V95-degree-pair,
thm:V95-FWRC-surplus,
cor:V95-complete-three}.
\end{proof}

\begin{programme}[V96: proof or refutation of the FWRC]
\label{prog:V95-V96}
\begin{enumerate}[leftmargin=3em]
\item Use the conjugation-sector decomposition of the Wilson
      generator to reduce \eqref{eq:V95-FWRC} to weighted radial
      Sturm--Liouville estimates.
\item Prove the \(p=4,8\) estimates first on every fixed angular
      sector by a one-dimensional Hardy inequality.
\item Obtain constants uniform in the angular momentum by retaining
      the positive centrifugal term rather than dropping it.
\item Treat the antipodal interval directly in the weighted
      Sturm--Liouville measure; do not invoke false global positive
      Bakry--Emery curvature.
\item Amplify the scalar estimate to Hilbert-valued functions before
      taking the column supremum.
\item If a uniform \(L^p\) Riesz bound fails, identify an explicit
      angular/radial countersequence and replace the FWRC by the
      smallest degree-generated test class needed in
      Lemma~\ref{lem:V95-degree-pair}.
\end{enumerate}
\end{programme}

\begin{firewall}[Claim boundary after fourth-series V95]
\label{fire:V95-claim-boundary}
V95 proves a conditional closure of all remaining three-label degree
sectors.  It does not prove the FWRC, the unconditional complete
three-label or all-subset RSPCC, the all-order local tree card, the
anisotropic scalar optimization, or the complete all-order selected
gate.  It does not prove all-order MTA, WUFC, CPM, cutoff removal,
reflection positivity, Osterwalder--Schrader reconstruction,
construction of the continuum Yang--Mills measure, or a uniform
positive continuum spectral gap.  The full Yang--Mills mass gap has
not been proved.
\end{firewall}

\paragraph{Parent-version record.}
V95 was formed from
\texttt{Renewal\_Yang\_Mills\_Mass\_Gap\_Research\_Notes\_V94.tex}.
The V94 parent had \(41\,912\) logical source lines,
\(1\,395\,257\) bytes, and SHA--256
\[
 \texttt{395e5dc06edd113ed812ec0baafa82e28ab0322c640177efdfd727350278e033}.
\]
The next compulsory companion audit is V100.

% =============================================================================
% END APPENDED FOURTH-SERIES V95 MATERIAL
% =============================================================================

% =============================================================================
% BEGIN APPENDED FOURTH-SERIES V96 MATERIAL
% =============================================================================

\chapter{V96: Semigroup proof of the complete fixed-\(p\) Wilson resolvent card}
\label{chap:V96-semigroup-FWRC}

\section{Curvature and gap have the same Wilson scale}
\label{sec:V96-curvature-gap}

Let
\[
 P_t:=e^{-tA_\alpha}
\]
be the Wilson Markov semigroup.  The diffusion generator in the
usual Bakry--Emery sign convention is \(-A_\alpha\), and its carré
du champ is the \(\Gamma\) used throughout these notes.

\begin{lemma}[Global curvature lower bound]
\label{lem:V96-curvature-lower}
There is an absolute \(\kappa<\infty\) such that, for
\(\alpha\ge\alpha_0\),
\begin{equation}
 \Gamma_2(f)
 \ge
 -\kappa\alpha\,\Gamma(f)
 \label{eq:V96-CD-negative}
\end{equation}
pointwise on \(SU(2)\).
\end{lemma}

\begin{proof}
For the density
\(d\nu_\alpha=Z_\alpha^{-1}e^{\alpha\phi}dU
=Z_\alpha^{-1}e^{-V_\alpha}dU\), one has
\(V_\alpha=-\alpha\phi\).  The radius-two geometric identities
recorded in Remark~\ref{rem:V78-BE-fails} are
\[
 \operatorname{Ric}={1\over2}g,
 \qquad
 \operatorname{Hess}\phi=-{1\over4}\phi g.
\]
Hence
\[
 \operatorname{Ric}+\operatorname{Hess}V_\alpha
 =
 \left({1\over2}+{\alpha\phi\over4}\right)g
 \ge
 \left({1\over2}-{\alpha\over4}\right)g
 \ge
 -\kappa\alpha g.
\]
The Bochner identity converts this tensor inequality into
\eqref{eq:V96-CD-negative}.
\end{proof}

\begin{remark}[Negative curvature is retained, not discarded]
\label{rem:V96-negative-curvature}
The lower bound is of order \(-\alpha\), not positive.  The argument
below succeeds because the independently proved spectral gap is of
order \(+\alpha\).  No global log-concavity assertion is made.
\end{remark}

\section{Reverse semigroup Poincare estimate}
\label{sec:V96-reverse-Poincare}

\begin{lemma}[Reverse Poincare under a negative curvature bound]
\label{lem:V96-reverse-Poincare}
Let \(H\) be a scalar or Hilbert-valued function.  Under
\eqref{eq:V96-CD-negative},
\begin{equation}
 |dP_tH|^2
 \le
 {\kappa\alpha\over1-e^{-2\kappa\alpha t}}
 \left(
 P_t\|H\|^2-\|P_tH\|^2
 \right)
 \le
 {\kappa\alpha\over1-e^{-2\kappa\alpha t}}
 P_t\|H\|^2.
 \label{eq:V96-reverse-Poincare}
\end{equation}
The scalar \(\kappa=0\) limit is interpreted as \(1/(2t)\).
\end{lemma}

\begin{proof}
For scalar \(H\), the diffusion identity is
\[
 P_t(H^2)-(P_tH)^2
 =
 2\int_0^t
 P_s\Gamma(P_{t-s}H)\,ds.
\]
The curvature lower bound gives the gradient propagation estimate
\[
 \Gamma(P_sF)
 \le
 e^{2\kappa\alpha s}P_s\Gamma(F).
\]
Apply it to \(F=P_{t-s}H\) and rearrange:
\[
 P_s\Gamma(P_{t-s}H)
 \ge
 e^{-2\kappa\alpha s}\Gamma(P_tH).
\]
Integration gives
\[
 P_t(H^2)-(P_tH)^2
 \ge
 {1-e^{-2\kappa\alpha t}\over\kappa\alpha}
 \Gamma(P_tH),
\]
which is \eqref{eq:V96-reverse-Poincare}.
For Hilbert-valued \(H\), apply the scalar identity to a finite
orthonormal expansion and sum its nonnegative coordinate
inequalities.  Approximation proves the general case.
\end{proof}

\begin{corollary}[Fixed-\(p\) semigroup gradient bound]
\label{cor:V96-gradient-bound}
For every \(p\ge2\),
\begin{equation}
 \|dP_tH\|_{L^p}
 \le
 \left(
 {\kappa\alpha\over1-e^{-2\kappa\alpha t}}
 \right)^{1/2}
 \|H\|_{L^p}.
 \label{eq:V96-gradient-bound}
\end{equation}
The same estimate holds for Hilbert-valued functions.
\end{corollary}

\begin{proof}
Take the \(L^{p/2}\) norm of
\eqref{eq:V96-reverse-Poincare}.  Since \(P_t\) is a Markov
contraction on \(L^{p/2}\),
\[
 \|P_t\|H\|^2\|_{p/2}
 \le
 \|\|H\|^2\|_{p/2}.
\]
Take square roots.
\end{proof}

\section{Mean-zero \(L^p\) decay}
\label{sec:V96-Lp-decay}

\begin{lemma}[Gap interpolation]
\label{lem:V96-Lp-decay}
Let \(\Pi H=\mathbb E_{\nu_\alpha}H\).  For every fixed
\(2\le p<\infty\),
\begin{equation}
 \|P_t(I-\Pi)H\|_{L^p}
 \le
 C_p e^{-c_p\alpha t}\|(I-\Pi)H\|_{L^p}.
 \label{eq:V96-Lp-decay}
\end{equation}
The constants are unchanged for Hilbert-valued \(H\).
\end{lemma}

\begin{proof}
The Wilson gap \eqref{eq:V78-gap} and self-adjointness give
\[
 \|P_t(I-\Pi)\|_{2\to2}\le e^{-c\alpha t}.
\]
Markov contraction and \(\|\Pi\|_{\infty\to\infty}\le1\) give
\[
 \|P_t(I-\Pi)\|_{\infty\to\infty}\le2.
\]
Riesz--Thorin interpolation between \(L^2\) and \(L^\infty\)
yields
\[
 \|P_t(I-\Pi)\|_{p\to p}
 \le
 2^{1-2/p}e^{-2c\alpha t/p},
\]
which is \eqref{eq:V96-Lp-decay}.  The \(L^2\) gap tensorizes with a
Hilbert target, and the Markov contraction applies to the Hilbert
norm by Jensen; the same interpolation proof follows.
\end{proof}

\section{The complete \(L^p\) gradient resolvent}
\label{sec:V96-resolvent}

\begin{theorem}[Complete fixed-\(p\) Wilson gradient-resolvent
estimate]
\label{thm:V96-FWRC}
For every fixed \(2\le p<\infty\) and every finite-dimensional
operator-valued \(G\),
\begin{equation}
 \boxed{\qquad
 \|dR_\alpha G\|_{p,c}
 \le
 {C_p\over\sqrt\alpha}
 \|Q_\alpha G\|_{p,c}.
 \qquad}
 \label{eq:V96-FWRC}
\end{equation}
The constant is independent of \(\alpha\), of the matrix dimension,
and of spectator tensor legs.  In particular, the FWRC
\eqref{eq:V95-FWRC} holds at \(p=4,8\).
\end{theorem}

\begin{proof}
First let \(G\) be Hilbert-valued and centered.  The semigroup
resolvent formula is
\[
 dR_\alpha G
 =
 \int_0^\infty dP_tG\,dt.
\]
Split the semigroup at half time:
\[
 dP_tG
 =
 dP_{t/2}(P_{t/2}G).
\]
Corollary~\ref{cor:V96-gradient-bound} at time \(t/2\), followed by
Lemma~\ref{lem:V96-Lp-decay}, gives
\begin{equation}
 \|dP_tG\|_p
 \le
 C_p
 \left(
 {\alpha\over1-e^{-\kappa\alpha t}}
 \right)^{1/2}
 e^{-c_p\alpha t}\|G\|_p.
 \label{eq:V96-integrand}
\end{equation}
The right side is integrable: after \(u=\alpha t\),
\begin{align*}
 \int_0^\infty
 \left(
 {\alpha\over1-e^{-\kappa\alpha t}}
 \right)^{1/2}
 e^{-c_p\alpha t}\,dt
 &=
 {1\over\sqrt\alpha}
 \int_0^\infty
 {e^{-c_pu}\over\sqrt{1-e^{-\kappa u}}}\,du\\
 &\le
 {C_p\over\sqrt\alpha}.
\end{align*}
Near zero the integrand in \(u\) is \(O(u^{-1/2})\); at infinity
it decays exponentially.  Minkowski's inequality proves
\eqref{eq:V96-FWRC} for centered Hilbert-valued \(G\).  The bound
\eqref{eq:V96-integrand} also justifies differentiation and
integration of the semigroup formula by approximation.

For an operator-valued \(G\), fix a unit input vector \(v\) and apply
the Hilbert-valued result to \((Q_\alpha G)v\).  The obtained
right-hand side is bounded by
\(\|Q_\alpha G\|_{p,c}\).  Taking the supremum over \(v\) proves the
complete column estimate without a target-dimension factor.
\end{proof}

\begin{corollary}[Unconditional complete three-label RSPCC]
\label{cor:V96-complete-three}
For arbitrary all-thermal irreducible unitary representations,
\begin{equation}
 \boxed{\qquad
 \|\mathfrak P_\alpha(\{1,2,3\})\|_{-1,c}
 \le
 {3!K^3\over\sqrt\alpha}
 b_1b_2b_3(b_1+b_2+b_3)^2.
 \qquad}
 \label{eq:V96-complete-three}
\end{equation}
\end{corollary}

\begin{proof}
Theorem~\ref{thm:V96-FWRC} proves the only hypothesis left open in
Corollary~\ref{cor:V95-complete-three}.  Apply that corollary.
\end{proof}

\begin{corollary}[A second independent proof of the arity-four local
tree card]
\label{cor:V96-arity-four}
The complete fourth Wilson cumulant satisfies
\begin{equation}
 \|\kappa_4^{\nu_\alpha}
 (D^{R_1},D^{R_2},D^{R_3},D^{R_4})\|
 \le
 4!K_1^4
 \left(\prod_{i=1}^4b_i\right)
 \left(\sum_{i=1}^4b_i\right)^2.
 \label{eq:V96-arity-four}
\end{equation}
Thus the arity-four local tree card follows through the
terminal-negative-norm route, independently of the direct
fourth-cumulant expansion of V83.
\end{corollary}

\begin{proof}
In the compression formula
\eqref{eq:V88-cumulant-compression} at \(m=4\), every terminal tail
has three labels.  Apply
\eqref{eq:V96-complete-three} and the outer pairing
\eqref{eq:V91-outer-negative}; then use the Pruefer identity exactly
as in Theorem~\ref{thm:V91-RSPCC-implies-card}.
\end{proof}

\begin{assessment}[First unproved subset size]
\label{ass:V96-next-subset}
The RSPCC is now proved for every label set of size at most three.
The first unproved instance has four labels:
\[
 \|\mathfrak P_\alpha(I)\|_{-1,c}
 \le
 {4!K^4\over\sqrt\alpha}
 \left(\prod_{i\in I}b_i\right)
 \left(\sum_{i\in I}b_i\right)^3,
 \qquad |I|=4.
\]
The FWRC removes fixed-\(p\) multiplication as a bottleneck, but it
does not by itself create the three incidence powers in the
low-degree sectors.  Those powers must again come from the assembled
coordinate-polynomial normal forms before norms.
\end{assessment}

\begin{firewall}[A one-link \(L^p\) resolvent is not a continuum
mass gap]
\label{fire:V96-not-continuum-gap}
Theorem~\ref{thm:V96-FWRC} is an intrinsic one-link Wilson diffusion
estimate.  It closes the analytic condition left by V95 and the
three-label RSPCC, but it does not prove the four-label or all-subset
RSPCC, the all-order local card, cutoff removal, reflection
positivity, Osterwalder--Schrader reconstruction, or a positive
continuum Hamiltonian gap.
\end{firewall}

\section{V96 ledger and V97 programme}
\label{sec:V96-results}

\begin{theorem}[Fourth-series V96 result ledger]
\label{thm:V96-results}
V96:
\begin{enumerate}[label=\textup{(V96.\arabic*)},leftmargin=6.3em]
\item combines the exact global curvature lower bound
      \(-\kappa\alpha\) with the independent \(c\alpha\) Wilson gap;
\item proves a Hilbert-valued reverse semigroup Poincare inequality;
\item proves exponential mean-zero \(L^p\) semigroup decay by
      interpolation;
\item proves the complete, dimension-free gradient-resolvent card
      for every fixed \(p\ge2\);
\item closes the complete three-label RSPCC unconditionally; and
\item recovers the arity-four local tree card through the
      symmetry-first terminal-tail method.
\end{enumerate}
\end{theorem}

\begin{proof}
Use
\Cref{lem:V96-curvature-lower,
lem:V96-reverse-Poincare,
lem:V96-Lp-decay,
thm:V96-FWRC,
cor:V96-complete-three,
cor:V96-arity-four}.
\end{proof}

\begin{programme}[V97: the four-label low-degree atlas]
\label{prog:V96-V97}
\begin{enumerate}[leftmargin=3em]
\item Write the exact four-label one-form recursion as the sum of
      four outer contractions with already-symmetrized three-label
      tails.
\item Use the FWRC at the finite Holder exponents required by four
      nested inputs.
\item Compute the all-linear coordinate sector first; every
      Gaussian tangent contraction is constant and must be removed
      before the next contraction.
\item Classify the degree patterns by total Taylor excess relative
      to the required three incidence powers.
\item Derive cubic-coordinate resolvent normal forms only for the
      patterns with fewer than three value powers.
\item Do not infer an all-subset induction until the local ordering
      count and the required \(p\)-exponents are uniform in the set
      size.
\end{enumerate}
\end{programme}

\begin{firewall}[Claim boundary after fourth-series V96]
\label{fire:V96-claim-boundary}
V96 proves the complete three-label RSPCC.  It does not prove the
four-label or all-subset RSPCC, the all-order local tree card, the
anisotropic scalar optimization, or the complete all-order selected
gate.  It does not prove all-order MTA, WUFC, CPM, cutoff removal,
reflection positivity, Osterwalder--Schrader reconstruction,
construction of the continuum Yang--Mills measure, or a uniform
positive continuum spectral gap.  The full Yang--Mills mass gap has
not been proved.
\end{firewall}

\paragraph{Parent-version record.}
V96 was formed from
\texttt{Renewal\_Yang\_Mills\_Mass\_Gap\_Research\_Notes\_V95.tex}.
The V95 parent had \(42\,305\) logical source lines,
\(1\,408\,274\) bytes, and SHA--256
\[
 \texttt{b8061407618ef78256d49fe994d35c7ecb20df6b10c1cbe788fd3010d9461b10}.
\]
The next compulsory companion audit is V100.

% =============================================================================
% END APPENDED FOURTH-SERIES V96 MATERIAL
% =============================================================================

% =============================================================================
% BEGIN APPENDED FOURTH-SERIES V97 MATERIAL
% =============================================================================

\chapter{V97: The four-label Holder atlas and high-excess closure}
\label{chap:V97-four-label}

\section{Exact four-label projected tail}
\label{sec:V97-four-tail}

For four centered compatible inputs define recursively
\begin{equation}
 \mathfrak P_\alpha(F_1,F_2,F_3,F_4)
 :=
 \sum_{i=1}^4
 \mathcal D_\alpha
 \left(
 F_i,
 \mathfrak P_\alpha
 (F_1,\ldots,\widehat F_i,\ldots,F_4)
 \right),
 \label{eq:V97-four-tail}
\end{equation}
with canonical restoration of tensor-leg order.  Expanding the
three-label tail shows that \eqref{eq:V97-four-tail} contains each of
the \(4!\) ordered projected chains exactly once.

\begin{lemma}[Four-label one-form recursion]
\label{lem:V97-four-one-form}
One has the exact identity
\begin{equation}
 dR_\alpha
 \mathfrak P_\alpha(F_1,F_2,F_3,F_4)
 =
 \sum_{i=1}^4
 \Theta_\alpha
 \left(
 F_i,
 \mathfrak P_\alpha
 (F_1,\ldots,\widehat F_i,\ldots,F_4)
 \right).
 \label{eq:V97-four-one-form}
\end{equation}
\end{lemma}

\begin{proof}
This is Theorem~\ref{thm:V91-one-form-recursion} with
\(|I|=4\), written in multilinear notation.
\end{proof}

\section{The finite Holder ladder}
\label{sec:V97-Holder-ladder}

V95 stated the degree estimates only through \(p=8\), because three
labels required no larger exponent.  Four nested inputs require one
more doubling.

\begin{lemma}[Degree bounds through \(p=16\)]
\label{lem:V97-degree-16}
For every fixed \(2\le p\le16\) and
\[
 P_{R,d}\in\{\Lambda_R,\mathsf Q_R,\mathsf M_R\},
 \qquad d\in\{1,2,3\},
\]
one has
\begin{equation}
 \|P_{R,d}\|_{p,c}\le C_pb_R^d,
 \qquad
 \|dP_{R,d}\|_{p,c}
 \le C_p\sqrt\alpha\,b_R^d.
 \label{eq:V97-degree-16}
\end{equation}
\end{lemma}

\begin{proof}
Repeat the proof of Lemma~\ref{lem:V95-degree-Lp}.  Every estimate
uses only a fixed logarithmic moment, now of order at most a fixed
multiple of sixteen.  Lemma~\ref{lem:V82-log-moments} supplies those
moments, and no constant depends on a representation dimension.
\end{proof}

For an ordering \(\pi\in S_4\), put
\begin{equation}
 \mathcal C_\alpha^\pi
 :=
 \mathcal D_\alpha
 \left(
 P_{\pi(1)},
 \mathcal D_\alpha
 \left(
 P_{\pi(2)},
 \mathcal D_\alpha(P_{\pi(3)},P_{\pi(4)})
 \right)
 \right),
 \label{eq:V97-ordered-four-chain}
\end{equation}
where every non-outer output is centered by
\(\mathcal D_\alpha\).

\begin{theorem}[Degree-preserving ordered four-chain]
\label{thm:V97-ordered-four-chain}
Let \(P_i=P_{R_i,d_i}\) with \(d_i\in\{1,2,3\}\).  For every
\(\pi\in S_4\),
\begin{equation}
 \|\mathcal C_\alpha^\pi\|_{-1,c}
 \le
 {C\over\sqrt\alpha}
 \prod_{i=1}^4b_i^{d_i}.
 \label{eq:V97-ordered-four-chain-bound}
\end{equation}
The constant is uniform in all representation and target
dimensions.
\end{theorem}

\begin{proof}
Let
\[
 G_{34}
 :=
 \mathcal D_\alpha(P_{\pi(3)},P_{\pi(4)}).
\]
Holder at exponent sixteen, the FWRC
\eqref{eq:V96-FWRC}, and
Lemma~\ref{lem:V97-degree-16} give
\[
 \|G_{34}\|_{8,c}
 \le
 C\|dP_{\pi(3)}\|_{16,c}
 \|dR_\alpha P_{\pi(4)}\|_{16,c}
 \le
 Cb_{\pi(3)}^{d_{\pi(3)}}
  b_{\pi(4)}^{d_{\pi(4)}}.
\]
Next put
\[
 G_{234}:=\mathcal D_\alpha(P_{\pi(2)},G_{34}).
\]
Holder at exponent eight and the \(p=8\) FWRC give
\[
 \|G_{234}\|_{4,c}
 \le
 Cb_{\pi(2)}^{d_{\pi(2)}}
   b_{\pi(3)}^{d_{\pi(3)}}
   b_{\pi(4)}^{d_{\pi(4)}}.
\]
At the outer operation, Holder at exponent four and the \(p=4\)
FWRC give the \(L^2\) bound
\[
 \left\|
 \mathcal D_\alpha(P_{\pi(1)},G_{234})
 \right\|_{2,c}
 \le
 C\prod_{i=1}^4b_i^{d_i}.
\]
The output is centered.  Apply
\eqref{eq:V90-gap-direction} to gain the terminal
\(\alpha^{-1/2}\).  At every stage, centering costs at most two and
the finite row--column Holder inequality introduces no target
dimension.
\end{proof}

\section{All sectors with three excess powers}
\label{sec:V97-high-excess}

\begin{theorem}[Four-label high-excess closure]
\label{thm:V97-high-excess}
Let
\[
 E:=\sum_{i=1}^4(d_i-1).
\]
If \(E\ge3\), then
\begin{equation}
 \boxed{\qquad
 \left\|
 \mathfrak P_\alpha
 (P_{R_1,d_1},P_{R_2,d_2},P_{R_3,d_3},P_{R_4,d_4})
 \right\|_{-1,c}
 \le
 {4!C\over\sqrt\alpha}
 \left(\prod_{i=1}^4b_i\right)
 \left(\sum_{i=1}^4b_i\right)^3.
 \qquad}
 \label{eq:V97-high-excess}
\end{equation}
\end{theorem}

\begin{proof}
Expand the symmetrized tail into its \(4!\) ordered chains and apply
Theorem~\ref{thm:V97-ordered-four-chain} to each.  Factor one
\(b_i\) at every label.  The remaining multiset contains \(E\ge3\)
factors, with \(d_i-1\) copies of \(b_i\).  Retain any three and
discard the others using \(b_i<1\).  Their product is at most
\[
 \left(\sum_{i=1}^4b_i\right)^3.
\]
Sum the \(4!\) ordered bounds.
\end{proof}

\begin{corollary}[Complete high-excess representation sector]
\label{cor:V97-high-excess-representation}
In the expansion of
\[
 \mathfrak P_\alpha(D^{R_1},D^{R_2},D^{R_3},D^{R_4})
\]
using the smooth split
\eqref{eq:V93-three-level} at every label, every degree sector with
total Taylor excess at least three satisfies the four-label RSPCC
target.
\end{corollary}

\begin{proof}
There are only \(3^4\) degree choices, and every such choice is
covered by Theorem~\ref{thm:V97-high-excess}.  Sum the fixed finite
number of sectors and absorb it into the constant.
\end{proof}

\section{The exact low-degree atlas}
\label{sec:V97-low-atlas}

\begin{proposition}[Only four low-degree pattern families remain]
\label{prop:V97-low-patterns}
Up to permutation, the degree sectors not covered by
Corollary~\ref{cor:V97-high-excess-representation} are exactly
\begin{equation}
 (1,1,1,1),\qquad
 (2,1,1,1),\qquad
 (3,1,1,1),\qquad
 (2,2,1,1).
 \label{eq:V97-low-patterns}
\end{equation}
They have total excess \(0,1,2,2\), respectively.
\end{proposition}

\begin{proof}
Each degree lies in \(\{1,2,3\}\).  Enumerating partitions of
\(E=0,1,2\) among four nonnegative integers \(d_i-1\le2\) gives
exactly the four patterns in
\eqref{eq:V97-low-patterns}.  Every other choice has \(E\ge3\).
\end{proof}

\begin{assessment}[Why V96 does not close the four low families]
\label{ass:V97-low-need-normal}
The degree-preserving ordered-chain estimate supplies
\[
 {C\over\sqrt\alpha}\prod_i b_i^{d_i}.
\]
For the four-label RSPCC, three further incidence powers are needed.
The patterns in \eqref{eq:V97-low-patterns} contain fewer than three
value powers, so at least one factor must come from an exact
coordinate-resolvent cancellation.  Applying the FWRC termwise
again would merely reproduce the displayed deficit.
\end{assessment}

\begin{assessment}[The all-linear normal-form ladder]
\label{ass:V97-linear-ladder}
The coordinate calculations already determine the required ladder:
\[
 \text{linear--linear pair}
 \longrightarrow
 {\alpha^{-1}}\times\text{centered quadratic},
\]
\[
 \text{linear acting on that pair}
 \longrightarrow
 {\alpha^{-2}}\times\text{linear},
\]
and a final linear contraction removes the constant tangent symbol.
To close \((1,1,1,1)\), V98 must make the
\(\alpha^{-2}\)-linear three-tail normal form quantitative in
\(L^4\), including its remainder.  The other three families then
require the corresponding insertion of one quadratic or cubic
smooth piece.
\end{assessment}

\begin{firewall}[Fixed arity is not an all-order Holder induction]
\label{fire:V97-not-all-order}
The proof uses the finite exponent ladder \(16,8,4,2\).  At label
set size \(r\), the analogous direct argument would require
exponents growing like \(2^r\).  The constants \(C_p\) in
Theorem~\ref{thm:V96-FWRC} have not been controlled uniformly along
that growth.  V97 therefore makes no all-subset induction claim.
\end{firewall}

\section{V97 ledger and V98 programme}
\label{sec:V97-results}

\begin{theorem}[Fourth-series V97 result ledger]
\label{thm:V97-results}
V97:
\begin{enumerate}[label=\textup{(V97.\arabic*)},leftmargin=6.3em]
\item writes the exact four-label projected-tail and one-form
      recursions;
\item extends the smooth degree bounds to the required exponent
      \(p=16\);
\item proves a dimension-free degree-preserving estimate for every
      ordered four-chain;
\item closes every four-label sector with total Taylor excess at
      least three;
\item proves that only the four low-degree families
      \eqref{eq:V97-low-patterns} remain; and
\item identifies the quantitative all-linear normal form required
      next.
\end{enumerate}
\end{theorem}

\begin{proof}
Use
\Cref{lem:V97-four-one-form,
lem:V97-degree-16,
thm:V97-ordered-four-chain,
thm:V97-high-excess,
prop:V97-low-patterns}.
\end{proof}

\begin{programme}[V98: the all-linear four-label normal form]
\label{prog:V97-V98}
\begin{enumerate}[leftmargin=3em]
\item Refine the coordinate-pair identity to
      \(\alpha^{-1}\) times a centered quadratic plus an
      \(L^8\)-controlled order-\(\alpha^{-2}\) remainder.
\item Insert the ordered linear--quadratic normal forms of V93--V94
      to derive the complete three-linear tail as
      \(\alpha^{-2}\) times a linear coordinate plus a sufficiently
      small \(L^4\) remainder.
\item Apply the last linear contraction before its negative norm;
      use the exact coordinate-pair defect rather than the FWRC
      bound on the leading linear term.
\item Prove the \((1,1,1,1)\) four-label RSPCC sector.
\item Test the same normal-form ladder with one quadratic insertion
      for the \((2,1,1,1)\) family.
\item Preserve all tensor-leg permutations before taking the column
      norm.
\end{enumerate}
\end{programme}

\begin{firewall}[Claim boundary after fourth-series V97]
\label{fire:V97-claim-boundary}
V97 closes every four-label degree sector of total Taylor excess at
least three.  It does not prove the four low-degree families, the
complete four-label or all-subset RSPCC, the all-order local tree
card, the anisotropic scalar optimization, or the complete
all-order selected gate.  It does not prove all-order MTA, WUFC,
CPM, cutoff removal, reflection positivity,
Osterwalder--Schrader reconstruction, construction of the continuum
Yang--Mills measure, or a uniform positive continuum spectral gap.
The full Yang--Mills mass gap has not been proved.
\end{firewall}

\paragraph{Parent-version record.}
V97 was formed from
\texttt{Renewal\_Yang\_Mills\_Mass\_Gap\_Research\_Notes\_V96.tex}.
The V96 parent had \(42\,714\) logical source lines,
\(1\,420\,244\) bytes, and SHA--256
\[
 \texttt{d8f972d600529aac4bf6b8b7206ad2d8ee002a4afc650252a056be27722465c9}.
\]
The next compulsory companion audit is V100.

% =============================================================================
% END APPENDED FOURTH-SERIES V97 MATERIAL
% =============================================================================

% =============================================================================
% BEGIN APPENDED FOURTH-SERIES V98 MATERIAL
% =============================================================================

\chapter{V98: Cubic triangular resolvents and the all-linear four-label sector}
\label{chap:V98-all-linear-four}

\section{An exact cubic coordinate identity}
\label{sec:V98-cubic-coordinate}

For \(b\in\{1,2,3\}\), put
\begin{equation}
 C_b:=x_b|x|^2.
 \label{eq:V98-Cb}
\end{equation}
This is a globally smooth odd cubic coordinate polynomial.

\begin{lemma}[Exact Wilson action on \(C_b\)]
\label{lem:V98-cubic-action}
One has
\begin{equation}
 A_\alpha C_b
 =
 \left({3\alpha\over4}\phi+{15\over4}\right)C_b
 -{5\over2}x_b.
 \label{eq:V98-cubic-action}
\end{equation}
\end{lemma}

\begin{proof}
Use the diffusion product rule with \(C_b=x_b|x|^2\).
Equations
\eqref{eq:V92-coordinate-metric} and
\eqref{eq:V93-radial-action} give
\[
 \Gamma(x_b,|x|^2)
 =
 {1\over2}x_b(1-|x|^2)
 =
 {1\over2}\phi^2x_b.
\]
Together with
\[
 A_\alpha x_b
 =
 \left({3\over4}+{\alpha\over4}\phi\right)x_b,
\]
this yields
\begin{align*}
 A_\alpha C_b
 &=
 x_bA_\alpha|x|^2+|x|^2A_\alpha x_b
 -2\Gamma(x_b,|x|^2)\\
 &=
 \left({11\over4}+{3\alpha\over4}\phi\right)C_b
 -{3\over2}x_b-\phi^2x_b.
\end{align*}
Since \(\phi^2=1-|x|^2\), the last term is
\(-x_b+C_b\), proving \eqref{eq:V98-cubic-action}.
\end{proof}

\begin{theorem}[Triangular cubic resolvent]
\label{thm:V98-cubic-resolvent}
For every fixed \(2\le p<\infty\),
\begin{equation}
 R_\alpha C_b
 =
 {4\over3\alpha}C_b
 +{40\over3\alpha^2}x_b
 +W_b^{(3)},
 \label{eq:V98-cubic-resolvent}
\end{equation}
where
\begin{equation}
 \|dW_b^{(3)}\|_{L^p(\nu_\alpha)}
 \le C_p\alpha^{-3}.
 \label{eq:V98-cubic-resolvent-error}
\end{equation}
\end{theorem}

\begin{proof}
Let
\[
 u_b^{(3)}
 :=
 {4\over3\alpha}C_b
 +{40\over3\alpha^2}x_b.
\]
Using Lemma~\ref{lem:V98-cubic-action} and the exact action on
\(x_b\), direct subtraction gives
\begin{align}
 C_b-A_\alpha u_b^{(3)}
 &=
 (1-\phi)C_b
 -{5\over\alpha}C_b
 +{10\over3\alpha}(1-\phi)x_b
 -{10\over\alpha^2}x_b.
 \label{eq:V98-cubic-residual}
\end{align}
Every term on the right has \(L^p\) norm
\(O_p(\alpha^{-5/2})\):
\[
 \|C_b\|_p\le C_p\alpha^{-3/2},
 \qquad
 \|(1-\phi)C_b\|_p\le C_p\alpha^{-5/2},
\]
and the remaining terms follow from
Lemma~\ref{lem:V92-coordinate-moments}.  The residual is centered
because it equals \(C_b-A_\alpha u_b^{(3)}\).
Apply \(R_\alpha\) and the complete \(L^p\) gradient-resolvent
estimate \eqref{eq:V96-FWRC}.  This supplies the additional
\(\alpha^{-1/2}\) in
\eqref{eq:V98-cubic-resolvent-error}.
\end{proof}

\section{A third-order expansion of the coordinate resolvent}
\label{sec:V98-coordinate-refinement}

The exact algebraic identity
\begin{equation}
 (1-\phi)x_b
 =
 {1\over2}C_b
 +{1\over2}(1-\phi)^2x_b
 \label{eq:V98-score-cubic-split}
\end{equation}
follows from
\(|x|^2=(1-\phi)(1+\phi)\).

\begin{lemma}[Refined coordinate resolvent]
\label{lem:V98-refined-coordinate}
Define
\begin{equation}
 a_\alpha
 :=
 {4\over\alpha+3}
 +{20\over3\alpha(\alpha+3)},
 \qquad
 c_\alpha:={2\over3(\alpha+3)}.
 \label{eq:V98-ac}
\end{equation}
Then, for every fixed \(p\ge2\),
\begin{equation}
 R_\alpha x_b
 =
 a_\alpha x_b+c_\alpha C_b+\widetilde W_b,
 \qquad
 \|d\widetilde W_b\|_p\le C_p\alpha^{-3}.
 \label{eq:V98-refined-coordinate}
\end{equation}
Moreover
\begin{equation}
 |a_\alpha|+|c_\alpha|\le {C\over\alpha}.
 \label{eq:V98-ac-bound}
\end{equation}
\end{lemma}

\begin{proof}
The residual identity in
Lemma~\ref{lem:V92-coordinate-residual} gives
\[
 R_\alpha x_b
 =
 {4\over\alpha}x_b
 +R_\alpha\left((1-\phi)x_b-{3\over\alpha}x_b\right).
\]
Move the last \(R_\alpha x_b\) term to the left:
\begin{equation}
 R_\alpha x_b
 =
 {4\over\alpha+3}x_b
 +{\alpha\over\alpha+3}
 R_\alpha\bigl((1-\phi)x_b\bigr).
 \label{eq:V98-coordinate-exact-rearrangement}
\end{equation}
Use \eqref{eq:V98-score-cubic-split} and
Theorem~\ref{thm:V98-cubic-resolvent}.  The remaining source
\((1-\phi)^2x_b\) has \(L^p\) norm
\(O_p(\alpha^{-5/2})\), so
\eqref{eq:V96-FWRC} bounds the gradient of its resolvent by
\(O_p(\alpha^{-3})\).  Collecting the \(x_b\) and \(C_b\)
coefficients gives exactly \eqref{eq:V98-ac}.
\end{proof}

\section{Coordinate pair to quadratic normal form}
\label{sec:V98-pair-normal}

\begin{lemma}[A fixed polynomial identity]
\label{lem:V98-Gamma-xC}
For all \(a,b\),
\begin{equation}
 \Gamma(x_a,C_b)
 =
 {1\over4}\delta_{ab}|x|^2
 +{1\over2}x_ax_b
 -{3\over4}|x|^2x_ax_b.
 \label{eq:V98-Gamma-xC}
\end{equation}
\end{lemma}

\begin{proof}
Use \(C_b=x_b|x|^2\),
\eqref{eq:V92-coordinate-metric}, and
\(\Gamma(x_a,|x|^2)=\frac12x_a(1-|x|^2)\).
The Leibniz rule gives
\begin{align*}
 \Gamma(x_a,C_b)
 &=
 |x|^2\Gamma(x_a,x_b)
 +x_b\Gamma(x_a,|x|^2),
\end{align*}
which simplifies to \eqref{eq:V98-Gamma-xC}.
\end{proof}

Define the centered quadratic polynomial
\begin{equation}
 \mathscr Q_{\alpha,ab}
 :=
 -{a_\alpha\over4}Q_\alpha(x_ax_b)
 +c_\alpha Q_\alpha
 \left(
 {1\over4}\delta_{ab}|x|^2+{1\over2}x_ax_b
 \right).
 \label{eq:V98-Qab}
\end{equation}
Its coefficient norm satisfies
\begin{equation}
 |\mathscr Q_{\alpha,ab}|_{\mathrm{coef}}
 \le {C\over\alpha}.
 \label{eq:V98-Qab-coefficient}
\end{equation}

\begin{theorem}[Coordinate-pair quadratic normal form]
\label{thm:V98-pair-normal}
For every fixed \(2\le p<\infty\),
\begin{equation}
 \mathcal D_\alpha(x_a,x_b)
 =
 \mathscr Q_{\alpha,ab}
 +E_{\alpha,ab}^{(2)},
 \qquad
 \|E_{\alpha,ab}^{(2)}\|_p
 \le C_p\alpha^{-3}.
 \label{eq:V98-pair-normal}
\end{equation}
\end{theorem}

\begin{proof}
Insert \eqref{eq:V98-refined-coordinate} into
\(\mathcal B_\alpha(x_a,x_b)=\Gamma(x_a,R_\alpha x_b)\):
\[
 \mathcal B_\alpha(x_a,x_b)
 =
 a_\alpha\Gamma(x_a,x_b)
 +c_\alpha\Gamma(x_a,C_b)
 +\Gamma(x_a,\widetilde W_b).
\]
After centering, the quadratic parts of the first two terms are
exactly \(\mathscr Q_{\alpha,ab}\), by
\eqref{eq:V92-coordinate-metric} and
\eqref{eq:V98-Gamma-xC}.  The remaining degree-four polynomial has
coefficient \(O(\alpha^{-1})\) and \(L^p\) norm
\(O_p(\alpha^{-3})\).  The last term is bounded by
\(C\|d\widetilde W_b\|_p\le C_p\alpha^{-3}\), because \(dx_a\)
is uniformly bounded.  Centering costs at most two.
\end{proof}

\section{Three-linear tail to a linear normal form}
\label{sec:V98-three-linear-normal}

\begin{lemma}[Ordered linear-on-quadratic normal form]
\label{lem:V98-LQ-ordered}
Let \(q\in\mathcal Q_2\).  For every coordinate \(x_a\) there is a
linear coordinate combination \(\mathscr J_a(q)\) such that
\begin{equation}
 \mathcal D_\alpha(x_a,q)
 =
 {1\over\alpha}\mathscr J_a(q)
 +F_{\alpha,a}(q),
 \label{eq:V98-LQ-ordered}
\end{equation}
with
\begin{equation}
 |\mathscr J_a(q)|_{\mathrm{coef}}
 \le C|q|_{\mathrm{coef}},
 \qquad
 \|F_{\alpha,a}(q)\|_p
 \le
 C_p|q|_{\mathrm{coef}}\alpha^{-5/2}
 \label{eq:V98-LQ-ordered-error}
\end{equation}
for every fixed \(p\ge2\).
\end{lemma}

\begin{proof}
The proof of Theorem~\ref{thm:V93-quadratic-resolvent}, combined
with \eqref{eq:V96-FWRC}, gives for every fixed \(p\)
\[
 R_\alpha q={2\over\alpha}q+W_q,
 \qquad
 \|dW_q\|_p
 \le C_p|q|_{\mathrm{coef}}\alpha^{-5/2}.
\]
Therefore
\[
 \mathcal B_\alpha(x_a,q)
 =
 {2\over\alpha}\Gamma(x_a,q)
 +\Gamma(x_a,W_q).
\]
The first term is \(\alpha^{-1}\) times a linear coordinate
polynomial plus a cubic coordinate polynomial.  Define
\(\mathscr J_a(q)\) from its linear part.  The cubic contribution is
bounded by
\[
 {C_p\over\alpha}|q|_{\mathrm{coef}}\alpha^{-3/2},
\]
and the \(W_q\) term has the same or smaller stated order.  Apply
\(Q_\alpha\); linear and cubic odd terms are already centered.
\end{proof}

\begin{theorem}[All-linear three-tail normal form]
\label{thm:V98-three-linear-normal}
For every \(a,b,c\), there is a linear coordinate combination
\(\mathscr L_{\alpha,abc}\) such that
\begin{equation}
 \mathfrak P_\alpha(x_a,x_b,x_c)
 =
 \mathscr L_{\alpha,abc}
 +E_{\alpha,abc}^{(3)},
 \label{eq:V98-three-linear-normal}
\end{equation}
where, for every fixed \(p\ge2\),
\begin{equation}
 |\mathscr L_{\alpha,abc}|_{\mathrm{coef}}
 \le C\alpha^{-2},
 \qquad
 \|E_{\alpha,abc}^{(3)}\|_p
 \le C_p\alpha^{-7/2}.
 \label{eq:V98-three-linear-error}
\end{equation}
\end{theorem}

\begin{proof}
Use the three-label recursion
\eqref{eq:V93-three-tail}.  By
Theorem~\ref{thm:V98-pair-normal}, each inner symmetrized pair is a
centered quadratic coordinate polynomial with coefficient norm
\(O(\alpha^{-1})\), plus an \(L^p\) remainder of order
\(\alpha^{-3}\).

Apply Lemma~\ref{lem:V98-LQ-ordered} to the quadratic part.  Its
linear output has coefficient \(O(\alpha^{-2})\), while its
remainder has \(L^p\) norm
\[
 O_p(\alpha^{-1}\alpha^{-5/2})
 =
 O_p(\alpha^{-7/2}).
\]
For an inner remainder \(E^{(2)}\), the FWRC and the uniform bound on
\(dx_a\) give
\[
 \|\mathcal D_\alpha(x_a,E^{(2)})\|_p
 \le
 C_p\alpha^{-1/2}\|E^{(2)}\|_p
 \le C_p\alpha^{-7/2}.
\]
Sum the three outer labels and both inner orderings.  All outputs are
odd, hence already centered.
\end{proof}

\section{Closure of the all-linear four-label sector}
\label{sec:V98-four-linear}

\begin{theorem}[Universal all-linear four-coordinate tail]
\label{thm:V98-four-coordinate}
For any \(a_1,a_2,a_3,a_4\),
\begin{equation}
 \left\|
 \mathfrak P_\alpha
 (x_{a_1},x_{a_2},x_{a_3},x_{a_4})
 \right\|_{-1}
 \le
 C\alpha^{-9/2}.
 \label{eq:V98-four-coordinate}
\end{equation}
\end{theorem}

\begin{proof}
Use the four-label recursion \eqref{eq:V97-four-tail}.  For each
outer coordinate, insert
\eqref{eq:V98-three-linear-normal} for the inner three-tail.

The leading part is a linear coordinate combination with coefficient
norm \(O(\alpha^{-2})\).  Expanding in coordinate pairs and using
\eqref{eq:V92-coordinate-pair-negative} bounds its outer
\(\mathcal D_\alpha\) contraction by
\[
 C\alpha^{-2}\alpha^{-5/2}
 =
 C\alpha^{-9/2}.
\]
For the remainder, Lemma~\ref{lem:V92-general-atom} and
\eqref{eq:V98-three-linear-error} give
\[
 \|\mathcal D_\alpha(x_{a_i},E^{(3)})\|_{-1}
 \le
 {C\over\alpha}\|E^{(3)}\|_2
 \le
 C\alpha^{-9/2}.
\]
Sum the four outer-label terms.
\end{proof}

\begin{corollary}[The \((1,1,1,1)\) representation sector]
\label{cor:V98-1111}
Let \(B=b_1+b_2+b_3+b_4\).  The all-linear contribution to
\(\mathfrak P_\alpha(D^{R_1},D^{R_2},D^{R_3},D^{R_4})\) satisfies
\begin{equation}
 \left\|
 \mathfrak P_\alpha
 (\Lambda_{R_1},\Lambda_{R_2},
  \Lambda_{R_3},\Lambda_{R_4})
 \right\|_{-1,c}
 \le
 {C\over\sqrt\alpha}
 \left(\prod_{i=1}^4b_i\right)B^3.
 \label{eq:V98-1111}
\end{equation}
\end{corollary}

\begin{proof}
Expand the four \(\Lambda\)'s in the fixed coordinate basis.
Theorem~\ref{thm:V98-four-coordinate} and the four Casimir row
bounds give
\[
 C\alpha^{-9/2}
 \prod_{i=1}^4\sqrt{C_2(R_i)}
 \le
 C\alpha^{-5/2}\prod_{i=1}^4b_i.
\]
Since \(B\ge4\sqrt{s_\alpha}\ge c\alpha^{-1/2}\),
\[
 \alpha^{-5/2}
 \le
 {C\over\sqrt\alpha}B^3.
\]
This proves \eqref{eq:V98-1111}.
\end{proof}

\begin{assessment}[Three low-degree families remain]
\label{ass:V98-three-families}
The all-linear family has been removed from the V97 atlas.  Up to
permutation, the unresolved four-label degree patterns are now
\[
 (2,1,1,1),\qquad
 (3,1,1,1),\qquad
 (2,2,1,1).
\]
Their normal forms can reuse the same triangular coordinate
resolvents, but the distinguished quadratic or cubic input must be
allowed in every one of the four chain positions before a norm is
taken.
\end{assessment}

\begin{firewall}[The coordinate normal form is fixed arity]
\label{fire:V98-fixed-arity}
Theorem~\ref{thm:V98-three-linear-normal} retains enough terms for
one further contraction.  It is not an all-order asymptotic series
and does not control constants as the number of nested labels grows.
V98 proves only the all-linear four-label family.
\end{firewall}

\section{V98 ledger and V99 programme}
\label{sec:V98-results}

\begin{theorem}[Fourth-series V98 result ledger]
\label{thm:V98-results}
V98:
\begin{enumerate}[label=\textup{(V98.\arabic*)},leftmargin=6.3em]
\item computes the exact triangular Wilson action on the cubic
      coordinate \(x_b|x|^2\);
\item proves its cubic resolvent with an explicit linear correction;
\item refines \(R_\alpha x_b\) through cubic coordinate degree;
\item upgrades every coordinate pair to a quadratic normal form
      with \(O_p(\alpha^{-3})\) remainder;
\item upgrades every all-linear three-tail to a linear normal form
      with \(O_p(\alpha^{-7/2})\) remainder; and
\item closes the all-linear four-label RSPCC sector.
\end{enumerate}
\end{theorem}

\begin{proof}
Use
\Cref{lem:V98-cubic-action,
thm:V98-cubic-resolvent,
lem:V98-refined-coordinate,
thm:V98-pair-normal,
thm:V98-three-linear-normal,
thm:V98-four-coordinate,
cor:V98-1111}.
\end{proof}

\begin{programme}[V99: distinguished quadratic and cubic insertions]
\label{prog:V98-V99}
\begin{enumerate}[leftmargin=3em]
\item Insert one quadratic input into each level of the
      pair-to-three-tail-to-four-tail normal-form ladder.
\item Prove the \((2,1,1,1)\) family with two incidence powers from
      coordinate cancellation and one from the quadratic value.
\item Treat \((3,1,1,1)\) using the cubic derivative bound of V94
      and the all-linear three-tail normal form.
\item Treat \((2,2,1,1)\) using the ordered quadratic--linear and
      quadratic--quadratic normal forms of V94.
\item Check all four placements of every distinguished input before
      declaring the four-label RSPCC complete.
\item Prepare the compulsory V100 companion audit and, if V99
      closes the low-degree atlas, use V100 to test the first
      arity-independent recursion.
\end{enumerate}
\end{programme}

\begin{firewall}[Claim boundary after fourth-series V98]
\label{fire:V98-claim-boundary}
V98 proves the all-linear four-label sector.  It does not prove the
three distinguished-input families, the complete four-label or
all-subset RSPCC, the all-order local tree card, the anisotropic
scalar optimization, or the complete all-order selected gate.  It
does not prove all-order MTA, WUFC, CPM, cutoff removal, reflection
positivity, Osterwalder--Schrader reconstruction, construction of
the continuum Yang--Mills measure, or a uniform positive continuum
spectral gap.  The full Yang--Mills mass gap has not been proved.
\end{firewall}

\paragraph{Parent-version record.}
V98 was formed from
\texttt{Renewal\_Yang\_Mills\_Mass\_Gap\_Research\_Notes\_V97.tex}.
The V97 parent had \(43\,076\) logical source lines,
\(1\,431\,002\) bytes, and SHA--256
\[
 \texttt{85e212938a833e2cdf76201c9995da07c4ae35aa3966432de2d640a15f2e2c0e}.
\]
The next compulsory companion audit is V100.

% =============================================================================
% END APPENDED FOURTH-SERIES V98 MATERIAL
% =============================================================================

% =============================================================================
% BEGIN APPENDED FOURTH-SERIES V99 MATERIAL
% =============================================================================

\chapter{V99: The finite polynomial tangent defect and complete
four-label closure}
\label{chap:V99-four-label-closure}

\section{Why the three remaining sectors share one cancellation}
\label{sec:V99-shared-cancellation}

The unresolved degree patterns from
\Cref{ass:V98-three-families} have total coordinate degree five or
six:
\[
 (2,1,1,1),\qquad (3,1,1,1),\qquad (2,2,1,1).
\]
It is important not to estimate their four possible distinguished
placements separately.  In the tangent Gaussian model every
projected four-tail of total degree at most six vanishes.  The
Wilson tail is therefore a single generator-and-measure defect.
This section makes that statement quantitative in the terminal
negative norm.

Let
\[
 \gamma_3(dy)=(2\pi)^{-3/2}e^{-|y|^2/2}\,dy
\]
and define the positive Ornstein--Uhlenbeck operator and its
carré du champ by
\begin{equation}
 A_0={1\over4}(-\Delta+y\cdot\nabla),
 \qquad
 \Gamma_0(p,q)={1\over4}\nabla p\cdot\nabla q.
 \label{eq:V99-OU}
\end{equation}
Write \(Q_0p=p-\mathbb E_{\gamma_3}p\),
\(R_0=A_0^{-1}Q_0\), and
\[
 \mathcal D_0(p,q):=Q_0\Gamma_0(p,R_0q).
\]
The symmetrized Gaussian projected tails \(\mathfrak P_0\) are
defined from \(\mathcal D_0\) by the same subset recursion as
\eqref{eq:V88-subset-recursion}.

\begin{lemma}[Fixed-degree comparison at every prescribed order]
\label{lem:V99-fixed-degree-comparison}
For each fixed integer \(M\), there are \(C_M,\alpha_M<\infty\)
such that every tensor polynomial
\[
 P(y)=\sum_{|\beta|\le M}P_\beta y^\beta
\]
satisfies, for \(\alpha\ge\alpha_M\),
\begin{equation}
 \left\|
 \mathbb E_{\nu_\alpha}P(\sqrt\alpha Z)
 -
 \mathbb E_{\gamma_3}P(Y)
 \right\|
 \le {C_M\over\alpha}
 \sum_{|\beta|\le M}\|P_\beta\|.
 \label{eq:V99-fixed-degree-comparison}
\end{equation}
The assertion is valid in complete column norm and after adjoining
fixed spectator tensor legs.
\end{lemma}

\begin{proof}
The proof of \Cref{lem:V84-polynomial-comparison} did not use the
number six except to bound a finite Gaussian moment.  On
\(|y|\le\alpha^{1/8}\) it gives
\[
 {d\mathcal L(\sqrt\alpha Z)\over d\gamma_3}(y)
 =
 1+{\Psi_4(y)\over\alpha}
 +O\left({1+|y|^8\over\alpha^2}\right),
\]
where \(\Psi_4\) is fixed and centered.  On the complement, both
the exact radial density and the displayed remainder have the
Gaussian domination established in that proof.  Multiplication by
\(|y|^k\), \(k\le M\), is integrable against this domination, with
a constant depending only on \(M\).  Integrate coefficientwise and
sum.  For an operator coefficient apply the scalar estimate to
\(P(y)v\), take the Euclidean norm by duality, and then take the
supremum over unit \(v\).  Spectator legs do not enter the scalar
density.
\end{proof}

\begin{lemma}[Gaussian connected-degree zero]
\label{lem:V99-Gaussian-tail-zero}
Let \(p_i\) be centered polynomials on \(\mathbb R^3\), each of
parity \(d_i\) and degree at most \(d_i\), where \(d_i\ge1\).
If
\[
 D:=d_1+d_2+d_3+d_4\le6,
\]
then
\begin{equation}
 \mathfrak P_0(p_1,p_2,p_3,p_4)=0.
 \label{eq:V99-Gaussian-tail-zero}
\end{equation}
\end{lemma}

\begin{proof}
Expand every \(p_i\) into multivariate Hermite polynomials.
On a Hermite polynomial of positive degree \(k\),
\(R_0\) multiplies by \(4/k\).  A monomial in
\(\Gamma_0(p,R_0q)\) joins one Gaussian half-edge of \(p\) to one
half-edge of \(q\).  Thus an ordered threefold contraction is a sum
of Wick diagrams on the four declared polynomial vertices whose
three contraction edges contain a spanning tree.  Such a diagram
uses six half-edges.

If \(D<6\), no diagram exists.  If \(D=6\), all half-edges are
consumed and the output is constant.  The outer occurrence of
\(Q_0\) in \(\mathcal D_0\) removes that constant.  This proves the
vanishing of each ordered projected chain, hence also of their
symmetrized sum.  Lower Hermite components only reduce the
available number of half-edges and obey the same argument.
\end{proof}

\begin{theorem}[Four-input Wilson--tangent defect]
\label{thm:V99-Wilson-tangent-defect}
Let \(p_i\) be tensor-valued homogeneous polynomials in the global
logarithmic coordinate \(Z\), of degrees
\(d_i\in\{1,2,3\}\), and suppose
\[
 4\le D:=\sum_{i=1}^4d_i\le6.
\]
Let \(|p_i|_{\rm row}\) be the complete row norm of its coefficient
tensor.  Then
\begin{equation}
 \boxed{\qquad
 \left\|
 \mathfrak P_\alpha
 (Q_\alpha p_1(Z),\ldots,Q_\alpha p_4(Z))
 \right\|_{-1,c}
 \le
 C\alpha^{-(D+3)/2}
 \prod_{i=1}^4|p_i|_{\rm row}.
 \qquad}
 \label{eq:V99-Wilson-tangent-defect}
\end{equation}
The constant is uniform in all coefficient and spectator
dimensions.  The same estimate holds for a finite sum of such
polynomials, coefficientwise.
\end{theorem}

\begin{proof}
We give the comparison argument because the extra
\(\alpha^{-1}\) is the entire content of the theorem.
Put \(y=\sqrt\alpha Z\), conjugate \(A_\alpha/\alpha\) to the
\(y\)-coordinate, and call the result \(\widehat A_\alpha\).
On the ball \(|y|\le\alpha^{1/8}\), the exact radial formula for
the \(SU(2)\) Laplacian and
\(\phi=\cos|Z|\) gives, on every fixed polynomial filtration,
\begin{equation}
 \widehat A_\alpha
 =
 A_0+\alpha^{-1}E_{\alpha,M},
 \qquad
 \widehat\Gamma_\alpha
 =
 \Gamma_0+\alpha^{-1}G_{\alpha,M}.
 \label{eq:V99-generator-comparison}
\end{equation}
Here \(E_{\alpha,M}\) maps degree at most \(M\) to a fixed
degree-\((M+2)\) filtration and \(G_{\alpha,M}\) maps two fixed
filtrations to their fixed product filtration; all coefficient
norms are uniformly bounded.  There is no
\(\alpha^{-1/2}\) term: both the radial metric coefficients and
\(\cos|Z|\) are even analytic functions of \(Z\).

For completeness, choose a smooth radial cutoff which is one on a
fixed identity neighbourhood.  The Wilson density on the support
of its derivatives is \(O(e^{-c\alpha})\), in every fixed Sobolev
norm.  Thus the cutoff error is smaller than the right side of
\eqref{eq:V99-Wilson-tangent-defect}.  Taylor's theorem with
remainder proves \eqref{eq:V99-generator-comparison} for the cutoff
polynomial.  Lemma~\ref{lem:V99-fixed-degree-comparison}, used up
to degree \(2(M+2)\), compares its Gram matrix and centering
projection with their Gaussian counterparts with error
\(O(\alpha^{-1})\).

On the nonconstant Gaussian polynomial block, \(A_0\ge1/4\).
The resolvent identity on this finite block therefore yields
\begin{equation}
 \widehat R_\alpha
 =
 R_0+\alpha^{-1}K_{\alpha,M},
 \label{eq:V99-resolvent-comparison}
\end{equation}
where \(K_{\alpha,M}\) is uniformly bounded between the required
fixed coefficient and \(L^p\) norms.  Equivalently, one may insert
the Gaussian approximate inverse into \(A_\alpha\); the residual
is \(O(\alpha^{-1})\), and
\eqref{eq:V96-FWRC} bounds the gradient of its exact resolvent.
This latter formulation also controls the component outside the
finite polynomial block.

Now expand the three occurrences of the centered contraction in an
ordered four-tail using
\eqref{eq:V99-generator-comparison},
\eqref{eq:V99-resolvent-comparison}, and the centering comparison.
The term in which all three factors are Gaussian is the
corresponding ordered \(\mathcal D_0\)-tail.  The sum of these
terms vanishes by
\Cref{lem:V99-Gaussian-tail-zero}.  Every surviving term contains
at least one of the three \(O(\alpha^{-1})\) defects.  The other
fixed-degree factors are uniformly bounded in the rescaled
polynomial norms.  Telescoping the finite product therefore gives
\[
 \|\widehat{\mathfrak P}_\alpha
       -\mathfrak P_0\|_{L^2,c}
 \le
 {C\over\alpha}\prod_i|p_i|_{\rm row}.
\]

Returning from \(y\) to \(Z\), the four input homogeneities supply
\(\alpha^{-D/2}\).  The terminal negative norm supplies a further
\(\alpha^{-1/2}\) by
\eqref{eq:V90-gap-direction}.  Hence
\[
 \alpha^{-D/2}\alpha^{-1}\alpha^{-1/2}
 =
 \alpha^{-(D+3)/2}.
\]
There are only \(4!\) orderings.  The complete row--column
argument used in
\Cref{thm:V96-FWRC,thm:V97-ordered-four-chain} makes every estimate
independent of coefficient and spectator dimensions.  Summing the
orderings proves \eqref{eq:V99-Wilson-tangent-defect}.
\end{proof}

\begin{remark}[The centering at total degree six is essential]
\label{rem:V99-degree-six-centering}
At \(D=6\), the tangent chain need not vanish before its final
projection: a connected tree can consume all six half-edges and
leave a scalar.  It is precisely the \(Q_0\) in the projected
contraction which removes this scalar.  Replacing
\(\mathcal D_\alpha\) by the uncentered
\(\mathcal B_\alpha\) would make
\Cref{thm:V99-Wilson-tangent-defect} false at the required order.
\end{remark}

\section{Separate the cubic jet from the representation remainder}
\label{sec:V99-cubic-split}

For the logarithmic Taylor expansion, write
\[
 L_R=dR(Z),\qquad
 \mathsf Q_R={1\over2}Q_\alpha L_R^2,\qquad
 \mathsf C_R={1\over6}L_R^3,
\]
and
\[
 \mathsf N_{R,4}
 :=
 Q_\alpha D^R-\Lambda_R-\mathsf Q_R-\mathsf C_R.
\]
This is the exact four-level split of
\Cref{lem:V84-four-level}; in the present notation
\(\Lambda_R=L_R\).

\begin{lemma}[Fourth-order value and derivative card]
\label{lem:V99-fourth-order-card}
For every fixed \(2\le p\le16\),
\begin{align}
 \|\mathsf N_{R,4}\|_{p,c}
 &\le C_pb_R^4,
 &
 \|d\mathsf N_{R,4}\|_{p,c}
 &\le C_p\sqrt\alpha\,b_R^4.
 \label{eq:V99-fourth-order-card}
\end{align}
The first three pieces obey
\begin{equation}
 \|P_{R,d}\|_{p,c}\le C_pb_R^d,
 \qquad
 \|dP_{R,d}\|_{p,c}\le C_p\sqrt\alpha\,b_R^d,
 \quad d=1,2,3.
 \label{eq:V99-first-three-card}
\end{equation}
\end{lemma}

\begin{proof}
The first value estimate is
\eqref{eq:V84-four-level-bounds}; the bounds in
\eqref{eq:V99-first-three-card} follow exactly as in
\Cref{lem:V97-degree-16}.  Differentiate Taylor's integral
remainder
\[
 e^{L_R}-I-L_R-\tfrac12L_R^2-\tfrac16L_R^3.
\]
Every differentiated summand contains one \(dL_R\) and at least
three undifferentiated \(L_R\)'s.  Complete row--column Hölder,
the Casimir row bound, and the fixed logarithmic moments give
\[
 \|d\mathsf N_{R,4}\|_{p,c}
 \le
 C_p\sqrt{C_2(R)}
 \left({C_2(R)\over\alpha}\right)^{3/2}
 \le C_p\sqrt\alpha\,b_R^4.
\]
The unitary integral remainder controls all higher powers without a
dimension factor.  Centering commutes with differentiation and
costs at most two in the value norm.
\end{proof}

\begin{corollary}[Every sector containing a fourth-order remainder]
\label{cor:V99-N4-sectors}
Every four-label sector containing at least one
\(\mathsf N_{R,4}\) satisfies the four-label RSPCC bound
\begin{equation}
 \|\mathfrak P_\alpha(P_1,P_2,P_3,P_4)\|_{-1,c}
 \le
 {C\over\sqrt\alpha}
 \left(\prod_{i=1}^4b_i\right)
 \left(\sum_{i=1}^4b_i\right)^3.
 \label{eq:V99-N4-sectors}
\end{equation}
\end{corollary}

\begin{proof}
Assign degree four to \(\mathsf N_{R,4}\).  The proof of
\Cref{thm:V97-ordered-four-chain} uses only the value and derivative
cards at \(p=16,8,4\).  Lemma~\ref{lem:V99-fourth-order-card}
supplies those cards for degree four as well.  Hence each ordered
chain is bounded by
\[
 {C\over\sqrt\alpha}\prod_i b_i^{d_i}.
\]
A degree-four label alone has three Taylor excess powers.  Factor
one \(b_i\) at each label, retain any three excess factors, and
discard the others using \(b_i<1\), exactly as in
\Cref{thm:V97-high-excess}.  Sum the \(4!\) orderings.
\end{proof}

\section{Closure of all three distinguished-input families}
\label{sec:V99-low-closure}

\begin{lemma}[Representation coefficient rows]
\label{lem:V99-representation-rows}
For \(d=1,2,3\), the homogeneous logarithmic polynomial
\(P_{R,d}\in\{\Lambda_R,\mathsf Q_R,\mathsf C_R\}\) satisfies
\begin{equation}
 |P_{R,d}|_{\rm row}
 \le C_d(1+C_2(R))^{d/2}
 \le C_d\alpha^{d/2}b_R^d.
 \label{eq:V99-representation-rows}
\end{equation}
\end{lemma}

\begin{proof}
The coefficient tensor is a symmetrized product of \(d\) derived
representation matrices.  Iterated complete row--column
Cauchy--Schwarz and
\(\sum_aT_a^*T_a=C_2(R)I\) give the first inequality.  Since
\(b_R^2=s_\alpha A_R\), \(A_R\ge1+C_2(R)\), and
\(s_\alpha\asymp\alpha^{-1}\), the second follows.
\end{proof}

\begin{theorem}[The three low-degree families]
\label{thm:V99-three-low-families}
Let \(B=b_1+b_2+b_3+b_4\).  For each of the degree patterns
\[
 (2,1,1,1),\qquad(3,1,1,1),\qquad(2,2,1,1),
\]
and every permutation of its distinguished labels,
\begin{equation}
 \boxed{\qquad
 \|\mathfrak P_\alpha
 (P_{R_1,d_1},P_{R_2,d_2},P_{R_3,d_3},P_{R_4,d_4})
 \|_{-1,c}
 \le
 {C\over\sqrt\alpha}
 \left(\prod_{i=1}^4b_i\right)B^3.
 \qquad}
 \label{eq:V99-three-low-families}
\end{equation}
\end{theorem}

\begin{proof}
Apply \Cref{thm:V99-Wilson-tangent-defect} and
\Cref{lem:V99-representation-rows}.  In all three cases the powers
of \(\alpha^{d_i/2}\) cancel the input homogeneity, leaving
\begin{equation}
 \|\mathfrak P_\alpha(P_{R_1,d_1},\ldots,P_{R_4,d_4})\|_{-1,c}
 \le
 C\alpha^{-3/2}\prod_i b_i^{d_i}.
 \label{eq:V99-common-low-bound}
\end{equation}
This estimate is permutation invariant, so it already covers every
possible chain position of every distinguished input.

For \((2,1,1,1)\), the factor remaining after extracting one
\(b_i\) at every label is \(b_j\).  Since
\[
 \alpha^{-1}\le CB^2,\qquad b_j\le B,
\]
\eqref{eq:V99-common-low-bound} implies the target.
For either excess-two pattern, the remaining product \(e\) is
\(b_j^2\) or \(b_jb_k\), hence \(e\le B^2\).  Also \(B\le4\)
because every label is all-thermal.  Consequently
\[
 \alpha^{-1}e\le CB^4\le4CB^3.
\]
Substitution in \eqref{eq:V99-common-low-bound} proves
\eqref{eq:V99-three-low-families}.
\end{proof}

\begin{theorem}[Complete four-label RSPCC]
\label{thm:V99-complete-four-RSPCC}
For arbitrary all-thermal irreducible unitary representations,
\begin{equation}
 \boxed{\qquad
 \|\mathfrak P_\alpha
 (D^{R_1},D^{R_2},D^{R_3},D^{R_4})\|_{-1,c}
 \le
 {4!K^4\over\sqrt\alpha}
 \left(\prod_{i=1}^4b_i\right)
 \left(\sum_{i=1}^4b_i\right)^3.
 \qquad}
 \label{eq:V99-complete-four-RSPCC}
\end{equation}
\end{theorem}

\begin{proof}
Use the exact four-level expansion
\[
 Q_\alpha D^{R_i}
 =
 \Lambda_{R_i}+\mathsf Q_{R_i}
 +\mathsf C_{R_i}+\mathsf N_{R_i,4}.
\]
Multilinearity creates \(4^4\) sectors.  Sectors containing
\(\mathsf N_4\) are covered by
\Cref{cor:V99-N4-sectors}.  Among the remaining sectors,
all patterns of excess at least three are covered by
\Cref{thm:V97-high-excess}; its proof applies unchanged to the
separate cubic piece.  The all-linear sector is
\Cref{cor:V98-1111}.  The only other patterns are the three in
\Cref{thm:V99-three-low-families}.  Sum this fixed finite ledger and
enlarge \(K\).
\end{proof}

\begin{corollary}[The arity-five all-thermal local tree card]
\label{cor:V99-arity-five-card}
The fifth Wilson cumulant satisfies the complete local tree bound
\begin{equation}
 \|\kappa_5^{\nu_\alpha}
 (D^{R_1},\ldots,D^{R_5})\|
 \le
 5!K_1^5
 \sum_{\tau\in\mathfrak T_5}
 \prod_{i=1}^5 b_i^{d_i(\tau)}.
 \label{eq:V99-arity-five-card}
\end{equation}
\end{corollary}

\begin{proof}
In the symmetry-first compression
\eqref{eq:V88-cumulant-compression}, every terminal projected tail
has four labels.  Apply
\Cref{thm:V99-complete-four-RSPCC}, then the outer negative pairing
\eqref{eq:V91-outer-negative}.  The Prüfer-code coefficient
identity in
\Cref{thm:V91-RSPCC-implies-card} converts
\(\prod b_i(\sum b_i)^3\) into the sum of labelled five-vertex tree
weights.  Sum the five possible outer labels and absorb the fixed
factor into \(K_1\).
\end{proof}

\begin{assessment}[The first unknown subset has moved to five]
\label{ass:V99-next-subset}
The RSPCC is now proved for every label set of size at most four,
and the local all-thermal tree card is proved through cumulant arity
five.  The first unknown projected estimate is
\[
 \|\mathfrak P_\alpha(I)\|_{-1,c}
 \le
 {5!K^5\over\sqrt\alpha}
 \left(\prod_{i\in I}b_i\right)
 \left(\sum_{i\in I}b_i\right)^4,
 \qquad |I|=5.
\]
The finite polynomial argument above does not iterate at arbitrary
arity: the degree threshold for a connected \(r\)-vertex Gaussian
tree is \(2(r-1)\), while the fixed-\(p\) Hölder ladder and its
constants also grow with \(r\).  V100 must therefore audit the
companion programmes and isolate an arity-independent replacement
for both growth mechanisms.
\end{assessment}

\begin{firewall}[Fixed-degree tangent cancellation is not the mass gap]
\label{fire:V99-not-mass-gap}
\Cref{thm:V99-complete-four-RSPCC} is a fixed-subset local
one-link estimate.  It does not prove the five-label or all-subset
RSPCC, an all-order local tree card, the anisotropic scalar
optimization, the complete selected gate, MTA, WUFC, CPM, cutoff
removal, reflection positivity, Osterwalder--Schrader
reconstruction, construction of the continuum Yang--Mills measure,
or a uniform positive continuum Hamiltonian gap.  The full
Yang--Mills mass gap has not been proved.
\end{firewall}

\section{V99 ledger and V100 programme}
\label{sec:V99-results}

\begin{theorem}[Fourth-series V99 result ledger]
\label{thm:V99-results}
V99:
\begin{enumerate}[label=\textup{(V99.\arabic*)},leftmargin=6.3em]
\item extends the Wilson--Gaussian moment comparison to every fixed
      prescribed polynomial degree;
\item proves that every four-input Gaussian projected tail of total
      degree at most six vanishes;
\item proves the terminal-negative Wilson tangent-defect estimate
      with one full \(\alpha^{-1}\) gain;
\item separates the cubic representation jet from a controlled
      fourth-order remainder;
\item closes all placements of the three remaining low-degree
      four-label families;
\item proves the complete four-label RSPCC; and
\item derives the arity-five all-thermal local tree card.
\end{enumerate}
\end{theorem}

\begin{proof}
Use
\Cref{lem:V99-fixed-degree-comparison,
lem:V99-Gaussian-tail-zero,
thm:V99-Wilson-tangent-defect,
lem:V99-fourth-order-card,
cor:V99-N4-sectors,
thm:V99-three-low-families,
thm:V99-complete-four-RSPCC,
cor:V99-arity-five-card}.
\end{proof}

\begin{programme}[V100: companion audit and the five-label gate]
\label{prog:V99-V100}
\begin{enumerate}[leftmargin=3em]
\item Inspect the newest active SM and NS notes read-only and record
      exact fingerprints before choosing the next route.
\item Write the exact five-label projected-tail recursion and its
      tangent connected-degree atlas.
\item Replace the exponentially growing direct Hölder ladder by a
      multilinear square-function or chaos-weighted norm if the
      companion audits support one.
\item Separate the issue of arity-uniform analytic constants from
      the issue of representation Taylor summation.
\item Prove a precise sufficient five-label certificate, or state
      the first unproved operator inequality without hiding it in a
      formal Gaussian expansion.
\item At the V100 boundary, freeze the fourth-series volume and
      restart a compact contextual volume containing the proved
      ledger and the live bottleneck.
\end{enumerate}
\end{programme}

\paragraph{Parent-version record.}
V99 was formed from
\texttt{Renewal\_Yang\_Mills\_Mass\_Gap\_Research\_Notes\_V98.tex}.
The V98 parent had \(43\,624\) logical source lines,
\(1\,445\,811\) bytes, and SHA--256
\[
 \texttt{a81edde9c8d0945c4c37fc78649d130ff1b0df3155e2baf141a55c696ad2cf4e}.
\]
The next compulsory companion audit is V100.

% =============================================================================
% END APPENDED FOURTH-SERIES V99 MATERIAL
% =============================================================================

% =============================================================================
% BEGIN APPENDED FOURTH-SERIES V100 MATERIAL
% =============================================================================

\chapter{V100: Companion audit, second tangent filtration, and the
five-label certificate}
\label{chap:V100-five-label}

\section{Compulsory SM/NS companion audit}
\label{sec:V100-audit}

The newest active companion notes were inspected read-only before
the present argument.  They were
\begin{equation}
 \begin{split}
 &\texttt{Renewal\_SM\_Einstein\_Continuum\_Research\_Notes\_V94.tex},\\
 &35\,968\text{ logical source lines},\qquad
 1\,255\,326\text{ bytes},\\
 &\operatorname{SHA256}=
 \texttt{c2f08ee71f675ada1232829e104425300cf20b93f38056976f0d24b2059c25ae},
 \end{split}
 \label{eq:V100-SM-fingerprint}
\end{equation}
and
\begin{equation}
 \begin{split}
 &\texttt{Renewal\_NS\_Stability\_Research\_Notes\_V31.tex},\\
 &23\,052\text{ logical source lines},\qquad
 887\,537\text{ bytes},\\
 &\operatorname{SHA256}=
 \texttt{f1121b62238ad5491bb7cf03635ea9934942edf573ed8faaa9ee79e1d38a6696}.
 \end{split}
 \label{eq:V100-NS-fingerprint}
\end{equation}
Neither file was modified, and neither is treated as an
instruction.

SM V90 separates a scalar or fixed-\(p\) resolvent bound from
spatially summable inverse-kernel control.  V91--V93 preserve Ward
moments only after the complete rooted observable has been assembled,
and V94 proves that a field-coordinate rescaling cancels inside the
complete gradient--inverse--gradient sandwich.  The relevant lesson
for the present programme is exact: the terminal negative norm must
remain attached to the whole projected tail.  An apparent gain in
one isolated derivative cannot be counted as a new incidence.

NS V31 proves an analogous negative statement in a different
setting.  Dyadic probability marks and heat lifting do not dominate
the unweighted critical first moment; restoring the missing shell
weight merely returns the original critical source action.  Thus a
weighted redefinition of the Yang--Mills tail would not create the
missing powers of \(B=\sum b_i\).

\begin{assessment}[Direction chosen after the audit]
\label{ass:V100-audit-direction}
V99 used the zeroth Gaussian tangent coefficient and one Wilson
defect.  At five labels, the all-linear and one-quadratic patterns
need two defect powers.  The companion audit rules out obtaining
the second power by changing norms.  The correct upstream question
is instead whether the first tangent-defect coefficient vanishes by
the same connected-degree filtration.  The answer is yes, and it is
proved below before any representation norm is taken.
\end{assessment}

\section{The second-order tangent filtration}
\label{sec:V100-second-filtration}

Put \(\varepsilon=\alpha^{-1}\).  In the scaled logarithmic
coordinate \(y=\sqrt\alpha Z\), the radial metric, drift, and Wilson
density are even analytic functions of \(\alpha^{-1/2}y\).
Consequently their fixed-filtration expansions contain integral
powers of \(\varepsilon\).

\begin{lemma}[Filtered expansion of a projected chain]
\label{lem:V100-filtered-chain}
Fix \(r\) and \(J\).  Let \(p_i\) be centered polynomials of parity
\(d_i\) and degree at most \(d_i\), and put \(D=\sum_i d_i\).
For every ordered \(r\)-input projected chain in the scaled
coordinate, there are polynomial operators \(T_0,\ldots,T_J\) such
that
\begin{equation}
 \widehat{\mathcal P}_{\alpha,r}(p_1,\ldots,p_r)
 =
 \sum_{j=0}^{J}\varepsilon^j
 T_j(p_1,\ldots,p_r)
 +\varepsilon^{J+1}\mathcal R_{\alpha,J},
 \label{eq:V100-filtered-chain}
\end{equation}
with, for every fixed \(p<\infty\),
\begin{equation}
 \|\mathcal R_{\alpha,J}\|_{L^p(\nu_\alpha),c}
 \le C_{r,J,p}\prod_i|p_i|_{\rm row}.
 \label{eq:V100-filtered-remainder}
\end{equation}
Moreover
\begin{equation}
 \deg T_j(p_1,\ldots,p_r)
 \le D-2(r-1)+2j,
 \label{eq:V100-degree-filtration}
\end{equation}
and \(T_j\) has the same parity as \(D\).  A negative degree means
that \(T_j=0\), and degree zero is removed by the final centering.
All constants are complete and independent of coefficient and
spectator dimensions.
\end{lemma}

\begin{proof}
Use a smooth identity-neighbourhood cutoff as in
\Cref{thm:V99-Wilson-tangent-defect}; its complement contributes
\(O(e^{-c\alpha})\) in every fixed Sobolev norm.  In the cutoff
chart, Taylor expansion of the exact radial coefficients gives
\begin{align}
 \widehat A_\alpha
 &=
 A_0+\sum_{j=1}^{J+1}\varepsilon^jE_j
 +\varepsilon^{J+2}E_{\alpha,J+2},
 \label{eq:V100-A-expansion}\\
 \widehat\Gamma_\alpha
 &=
 \Gamma_0+\sum_{j=1}^{J+1}\varepsilon^jG_j
 +\varepsilon^{J+2}G_{\alpha,J+2}.
 \label{eq:V100-Gamma-expansion}
\end{align}
On polynomials, \(E_j\) raises degree by at most \(2j\) relative
to \(A_0\), while \(G_j\) raises the output degree by at most
\(2j\) relative to the two-derivative contraction
\(\Gamma_0\).  This follows term by term: the order-\(j\) metric
coefficient has degree at most \(2j\), and the order-\(j\) drift
coefficient has degree at most \(2j+1\) followed by one
derivative.

On every fixed nonconstant polynomial filtration \(A_0\ge1/4\).
Repeated use of the finite-dimensional resolvent identity expands
\(\widehat R_\alpha\) through order \(J+1\).  A contribution of
total order \(j\) raises degree by at most \(2j\) relative to
\(R_0\).  The exact remainder is controlled by inserting the
truncated inverse into \(A_\alpha\) and applying the fixed-\(p\)
gradient resolvent
\eqref{eq:V96-FWRC}.  Lemma
\ref{lem:V99-fixed-degree-comparison}, at a sufficiently large but
fixed degree, controls the Gram matrices and the Wilson centering
projection.  Centering subtracts a constant and therefore never
increases degree.

Each of the \(r-1\) contractions differentiates both inputs once,
so the Gaussian order drops total degree by two.  Distribute total
defect order \(j\) among the generator, resolvent, metric, and
centering expansions.  The preceding paragraph permits at most
\(2j\) degrees to be restored.  This proves
\eqref{eq:V100-degree-filtration}.  Even radial coefficients preserve
parity.  Fixed-filtration Hölder estimates and the complete
row--column argument prove
\eqref{eq:V100-filtered-remainder}.  Finally the outer
\(Q_\alpha\) removes every scalar coefficient.
\end{proof}

\begin{corollary}[Five-input tangent orders]
\label{cor:V100-five-tangent-orders}
For \(r=5\), let \(E=D-5\in\{0,1,2,3\}\).  Define
\begin{equation}
 k(E):=
 \begin{cases}
 2,&E=0,1,\\
 1,&E=2,3.
 \end{cases}
 \label{eq:V100-kE}
\end{equation}
Then every coefficient of order strictly below
\(\varepsilon^{k(E)}\) in a centered five-input projected chain
vanishes.
\end{corollary}

\begin{proof}
Here \(2(r-1)=8\).  By
\eqref{eq:V100-degree-filtration}, the order-\(j\) coefficient has
degree at most
\[
 D-8+2j=E-3+2j.
\]
For \(E=0,1\), this degree is negative at \(j=0\) and at most zero
at \(j=1\); in the latter boundary case final centering removes the
constant.  For \(E=2,3\), the order-zero degree is at most zero.
This is precisely \eqref{eq:V100-kE}.
\end{proof}

\begin{theorem}[Five-input Wilson tangent-defect estimate]
\label{thm:V100-five-tangent-defect}
Let \(p_i(Z)\) be homogeneous tensor polynomials of degrees
\(d_i\ge1\), put \(D=\sum_i d_i=5+E\), and assume
\(0\le E\le3\).  Then
\begin{equation}
 \boxed{\qquad
 \|\mathfrak P_\alpha
 (Q_\alpha p_1,\ldots,Q_\alpha p_5)\|_{-1,c}
 \le
 C_E\,
 \alpha^{-(D+1)/2-k(E)}
 \prod_{i=1}^5|p_i|_{\rm row}.
 \qquad}
 \label{eq:V100-five-tangent-defect}
\end{equation}
\end{theorem}

\begin{proof}
Apply \Cref{lem:V100-filtered-chain} with \(J=1\).
For \(E=0,1\), Corollary~\ref{cor:V100-five-tangent-orders}
removes the coefficients of orders zero and one, leaving a uniform
\(O(\alpha^{-2})\) scaled \(L^2\) remainder.  For \(E=2,3\), use
the expansion with \(J=0\); the Gaussian coefficient vanishes and
the remainder is \(O(\alpha^{-1})\).

Returning from \(y\) to \(Z\) supplies
\(\alpha^{-D/2}\).  The output is centered, so
\eqref{eq:V90-gap-direction} converts its \(L^2\) bound to the
negative norm with another factor \(\alpha^{-1/2}\).  Sum the
\(5!\) ordered chains.  This is
\eqref{eq:V100-five-tangent-defect}.
\end{proof}

\section{The five-level representation split}
\label{sec:V100-five-level}

For \(1\le d\le4\), define the centered homogeneous logarithmic
piece
\begin{equation}
 \mathsf P_{R,d}:=
 Q_\alpha{L_R^d\over d!},
 \qquad L_R=dR(Z),
 \label{eq:V100-Pd}
\end{equation}
and define the exact remainder
\begin{equation}
 \mathsf N_{R,5}
 :=
 Q_\alpha D^R-\sum_{d=1}^4\mathsf P_{R,d}.
 \label{eq:V100-N5}
\end{equation}
Odd pieces are already centered by radial parity.

\begin{lemma}[Degree cards through five and \(p=32\)]
\label{lem:V100-degree-32}
For each fixed \(2\le p\le32\),
\begin{align}
 \|\mathsf P_{R,d}\|_{p,c}
 &\le C_pb_R^d,
 &
 \|d\mathsf P_{R,d}\|_{p,c}
 &\le C_p\sqrt\alpha\,b_R^d,
 \qquad 1\le d\le4,
 \label{eq:V100-Pd-card}\\
 \|\mathsf N_{R,5}\|_{p,c}
 &\le C_pb_R^5,
 &
 \|d\mathsf N_{R,5}\|_{p,c}
 &\le C_p\sqrt\alpha\,b_R^5.
 \label{eq:V100-N5-card}
\end{align}
Moreover
\begin{equation}
 |\mathsf P_{R,d}|_{\rm row}
 \le C_d\alpha^{d/2}b_R^d.
 \label{eq:V100-Pd-row}
\end{equation}
\end{lemma}

\begin{proof}
The proof of
\Cref{lem:V99-fourth-order-card,lem:V99-representation-rows}
applies verbatim at every fixed exponent through \(32\).
For the remainder, Taylor's integral formula starts with five
undifferentiated factors.  After one derivative, every term contains
one \(dL_R\) and at least four \(L_R\)'s.  Fixed logarithmic
moments and the Casimir row identity give
\[
 \|d\mathsf N_{R,5}\|_{p,c}
 \le
 C_p\sqrt{C_2(R)}
 \left({C_2(R)\over\alpha}\right)^2
 \le C_p\sqrt\alpha\,b_R^5.
\]
The value estimate and
\eqref{eq:V100-Pd-row} follow from the same computation without,
respectively with, the rescaling of coefficient tensors.
\end{proof}

\begin{theorem}[Five-label degree-preserving chain]
\label{thm:V100-five-degree-chain}
Let each input be one of the pieces in
\Cref{lem:V100-degree-32}, with assigned degree
\(d_i\in\{1,2,3,4,5\}\).  Then every ordered five-chain obeys
\begin{equation}
 \|\mathcal C^\pi_{\alpha,5}\|_{-1,c}
 \le
 {C\over\sqrt\alpha}\prod_{i=1}^5b_i^{d_i}.
 \label{eq:V100-five-degree-chain}
\end{equation}
Consequently every degree sector with
\(\sum_i(d_i-1)\ge4\) satisfies the five-label RSPCC target.
\end{theorem}

\begin{proof}
Start with the innermost pair.  Hölder and the \(p=32\) FWRC give
an \(L^{16}\) centered output bounded by the product of its two
degree weights.  The next three contractions use successively
\[
 (p_{\rm input},p_{\rm output})
 =
 (16,16),\ (8,8),\ (4,4),
\]
and the FWRC at \(16,8,4\), producing outputs in
\(L^8,L^4,L^2\).  At every step the input derivative contributes
\(\sqrt\alpha b_i^{d_i}\) and the resolvent gradient contributes
\(\alpha^{-1/2}\), so no power is lost.  The final centered
\(L^2\) output gains \(\alpha^{-1/2}\) in the negative norm.
This proves \eqref{eq:V100-five-degree-chain}.

If the total excess is at least four, factor one \(b_i\) at every
label, retain any four excess factors, and bound their product by
\((\sum_i b_i)^4\); discard all further factors using \(b_i<1\).
Sum the \(5!\) orderings.
\end{proof}

\section{Complete five-label closure}
\label{sec:V100-five-closure}

\begin{theorem}[Every low-excess five-label polynomial sector]
\label{thm:V100-five-low-sectors}
Let \(d_i\in\{1,2,3,4\}\) and
\(E=\sum_i(d_i-1)\le3\).  Put \(B=\sum_{i=1}^5b_i\).
Then
\begin{equation}
 \|\mathfrak P_\alpha
 (\mathsf P_{R_1,d_1},\ldots,\mathsf P_{R_5,d_5})
 \|_{-1,c}
 \le
 {C\over\sqrt\alpha}
 \left(\prod_{i=1}^5b_i\right)B^4.
 \label{eq:V100-five-low-sectors}
\end{equation}
\end{theorem}

\begin{proof}
Use \eqref{eq:V100-five-tangent-defect} and
\eqref{eq:V100-Pd-row}.  With \(D=5+E\), all input homogeneity
powers cancel, giving
\begin{equation}
 \|\mathfrak P_\alpha(\mathsf P_{R_1,d_1},\ldots,
 \mathsf P_{R_5,d_5})\|_{-1,c}
 \le
 {C\over\sqrt\alpha}\,
 \alpha^{-k(E)}
 \prod_i b_i^{d_i}.
 \label{eq:V100-five-common}
\end{equation}
After extracting the five base factors, the remaining product
\(e\) has degree \(E\) in the \(b_i\)'s, so \(e\le B^E\).
Also
\(\alpha^{-1}\le CB^2\) and \(B\le5\).  Therefore
\[
 \alpha^{-k(E)}e
 \le
 C B^{2k(E)+E}
 \le C'B^4,
\]
because \(2k(E)+E\ge4\); if the exponent exceeds four, use
\(B\le5\).  This proves the result simultaneously for every degree
pattern and every distinguished placement.
\end{proof}

\begin{theorem}[Complete five-label RSPCC]
\label{thm:V100-complete-five-RSPCC}
For arbitrary all-thermal irreducible unitary representations,
\begin{equation}
 \boxed{\qquad
 \|\mathfrak P_\alpha(D^{R_1},\ldots,D^{R_5})\|_{-1,c}
 \le
 {5!K^5\over\sqrt\alpha}
 \left(\prod_{i=1}^5b_i\right)
 \left(\sum_{i=1}^5b_i\right)^4.
 \qquad}
 \label{eq:V100-complete-five-RSPCC}
\end{equation}
\end{theorem}

\begin{proof}
Expand every input by the exact five-level split
\[
 Q_\alpha D^{R_i}
 =
 \sum_{d=1}^4\mathsf P_{R_i,d}+\mathsf N_{R_i,5}.
\]
Every sector of total excess at least four, including every sector
containing a degree-five remainder, is covered by
\Cref{thm:V100-five-degree-chain}.  All remaining polynomial sectors
are covered by
\Cref{thm:V100-five-low-sectors}.  Sum the fixed \(5^5\) sectors
and enlarge \(K\).
\end{proof}

\begin{corollary}[The arity-six all-thermal local tree card]
\label{cor:V100-arity-six}
The sixth Wilson cumulant obeys
\begin{equation}
 \|\kappa_6^{\nu_\alpha}(D^{R_1},\ldots,D^{R_6})\|
 \le
 6!K_1^6
 \sum_{\tau\in\mathfrak T_6}
 \prod_{i=1}^6 b_i^{d_i(\tau)}.
 \label{eq:V100-arity-six}
\end{equation}
\end{corollary}

\begin{proof}
Apply \Cref{thm:V100-complete-five-RSPCC} to each terminal
five-label tail in the symmetry-first compression, use the outer
negative pairing, and apply the Prüfer identity exactly as in
\Cref{cor:V99-arity-five-card}.
\end{proof}

\begin{assessment}[What has and has not become arity independent]
\label{ass:V100-arity-growth}
The degree filtration
\[
 \deg T_j\le D-2(r-1)+2j
\]
is arity independent as an algebraic statement.  It shows that for
each fixed \(r\), only finitely many tangent-defect coefficients and
finitely many Taylor degrees are needed.  The present proof still
uses \(L^p\) exponents up to \(2^r\), a Taylor sector count of order
\(r^r\), and constants \(C_{r,J,p}\) with no proved exponential
generating-function bound.  Thus it does not establish the uniform
\(r!K^r\) RSPCC required for all-order summation.
\end{assessment}

\begin{firewall}[Six cumulant orders are not a continuum mass gap]
\label{fire:V100-not-mass-gap}
\Cref{thm:V100-complete-five-RSPCC} proves a fixed five-label
one-link certificate and the local tree card through cumulant arity
six.  It does not prove an all-subset certificate with uniform
factorial constants, the full all-order local card, anisotropic
scalar optimization, the selected gate, MTA, WUFC, CPM, cutoff
removal, reflection positivity, Osterwalder--Schrader
reconstruction, construction of the continuum Yang--Mills measure,
or a uniform positive continuum Hamiltonian gap.  The full
Yang--Mills mass gap has not been proved.
\end{firewall}

\section{V100 result ledger and rollover programme}
\label{sec:V100-results}

\begin{theorem}[Fourth-series V100 result ledger]
\label{thm:V100-results}
V100:
\begin{enumerate}[label=\textup{(V100.\arabic*)},leftmargin=6.8em]
\item audits SM V94 and NS V31 and keeps the full terminal negative
      sandwich rather than assigning a coordinate or mark gain;
\item proves the all-order-in-degree filtered tangent expansion at
      every fixed chain length and expansion order;
\item proves the precise second-order vanishings needed at five
      labels;
\item derives the five-input tangent-defect estimate;
\item proves the five-level representation Taylor and derivative
      cards through \(p=32\);
\item closes every low- and high-excess five-label sector;
\item proves the complete five-label RSPCC; and
\item derives the arity-six all-thermal local tree card.
\end{enumerate}
\end{theorem}

\begin{proof}
Use
\Cref{lem:V100-filtered-chain,
cor:V100-five-tangent-orders,
thm:V100-five-tangent-defect,
lem:V100-degree-32,
thm:V100-five-degree-chain,
thm:V100-five-low-sectors,
thm:V100-complete-five-RSPCC,
cor:V100-arity-six}.
\end{proof}

\begin{programme}[Next compact volume: uniform analytic norms]
\label{prog:V100-rollover}
The next active volume will restart at V1 with a contextual summary
of the proved chain.  Its first live task is to replace the
fixed-arity \(2^r\) Hölder ladder by one analytic or
chaos-weighted complete norm in which:
\begin{enumerate}[leftmargin=3em]
\item multiplication and the Wilson gradient resolvent have
      exponential, rather than iterated-\(p\), constants;
\item the filtered tangent coefficients admit an exponential
      generating-function majorant;
\item the complete representation Taylor series is summed before
      individual sectors are normed; and
\item failure of any one of these statements is isolated as an
      explicit operator inequality, not concealed in a formal
      all-order expansion.
\end{enumerate}
\end{programme}

\paragraph{Parent-version record.}
V100 was formed from
\texttt{Renewal\_Yang\_Mills\_Mass\_Gap\_Research\_Notes\_V99.tex}.
The V99 parent had \(44\,199\) logical source lines,
\(1\,465\,401\) bytes, and SHA--256
\[
 \texttt{52f5ef6046e0e63fdbbf94eb9c2851f26e49b2af9c60d1cdc1f82cbab6444b9b}.
\]

% =============================================================================
% END APPENDED FOURTH-SERIES V100 MATERIAL
% =============================================================================

\end{document}
